\documentclass[9pt,a4paper,openany]{book}
\usepackage[T1]{fontenc}
\usepackage[utf8]{inputenc}
\usepackage{lmodern}
\usepackage[margin=25mm]{geometry}
\usepackage{longtable,booktabs,array,tabularx,hyperref,xcolor,makeidx}
\usepackage{hyperref}
\hypersetup{colorlinks=true,linkcolor=blue,urlcolor=blue}
\makeindex
\setcounter{tocdepth}{3}
\setcounter{secnumdepth}{3}
\newcommand{\reviewstatus}[1]{\par\smallskip\noindent\fcolorbox{gray}{gray!8}{\parbox{0.92\linewidth}{\textbf{Review:} #1}}\par\smallskip}
\newcommand{\reviewcomment}[1]{\par\smallskip\noindent\fcolorbox{gray}{gray!8}{\parbox{0.92\linewidth}{\textbf{Review Comment:} #1}}
\par\smallskip}
\newcommand{\enabledby}[1]{\par\smallskip\noindent\fcolorbox{gray}{gray!8}{\parbox{0.92\linewidth}{\textbf{EnabledBy:} #1}}
\par\smallskip}
\newcommand{\supports}[1]{\par\smallskip\noindent\fcolorbox{gray}{gray!8}{\parbox{0.92\linewidth}{\textbf{Supports:} #1}}
\par\smallskip}
\newcommand{\derivedfrom}[1]{\par\smallskip\noindent\fcolorbox{gray}{gray!8}{\parbox{0.92\linewidth}{\textbf{DerivedFrom:} #1}}
\par\smallskip}

\newcommand{\featuretable}[1]{%
{%
\small
\begin{longtable}{@{}p{.14\linewidth}|p{.15\linewidth}|p{.30\linewidth}|p{.32\linewidth}@{}}
\toprule Type & ID & Name & Description / Notes\\\midrule\endhead
#1
\bottomrule\end{longtable}
}%
\title{Harmonia\\Capability Architecture Reference Book\\\large Full-Candidate Working Draft, Version 0.4}
\author{Architecture Review Working Copy}
\date{15 October 2026}
\begin{document}
\frontmatter\maketitle\tableofcontents
\mainmatter
\chapter*{Document Purpose and Reading Rules}\addcontentsline{toc}{chapter}{Document Purpose and Reading Rules}
This is a deliberately inclusive, reviewable candidate reference architecture, not an approved replacement for the governed Harmonia repository. It combines established catalogue labels, proposed ancestry, newly inferred Features, and existing Enterprise Capabilities retained alongside fresh derivations. Contradictions and overlaps are intentionally visible.

\begin{itemize}
\item The Business Capability part describes organisational abilities, not Harmonia ownership of clinical or operational work.
\item The Business Enabling part retains established Features and proposes additional Features under candidate CT3 frames. Existing BEC CT1--CT3 ancestry remains candidate where the repository has not established it.
\item The Enterprise part is derived from repeated semantic enabling behaviours, then augmented by the pre-existing EC-01 through EC-14 catalogue without conflict suppression.
\item Names do not identify architectural elements: reference IDs do. IDs are deterministic for the identity key, independent of the printed chapter/section number; they are provisional reference IDs, not allocated canonical RADS IDs.
\item Similarity prompts inspection, not symmetry. Repeated Feature names are allowed under distinct CT3s; recurring patterns do not automatically imply identical responsibility.
\item FHIR R5 is conceptual evidence about healthcare responsibilities, relationships and lifecycle; FHIR fields and resource models are not Feature definitions.
\item Existing Functions, Services and Processes appear beneath Features. Empty tables are retained to make missing decomposition visible. New Functions, Services and Processes have not been fabricated.
\item Source-established content, inferred descriptions, and generated candidates are distinguished. Review annotations identify uncertainty and possible overreach.
\item Legacy documentation is excluded. Current governed Markdown, the working capability map, the RADS behaviour catalogue represented there, current scratchpad reasoning and FHIR R5 conceptual references are the intended evidence classes.
\end{itemize}
\paragraph{Review convention.} Treat each proposed Feature as an invitation to keep, rename, move, merge or delete it. Do not promote a candidate into governed architecture merely because it appears in this book.
\paragraph{Source limitation.} This generation uses the current v0.4 working mind-map as the complete machine-readable hierarchy and Feature-to-Function/Service/Process association source. The repository catalogue and Design Note sources were inspected for scope, but not all free-form RADS and scratchpad paragraphs have been transcribed verbatim. Consequently, generated explanatory prose is labelled candidate interpretation; absence of a note here does not imply absence from the RADS. Reconcile with governed sources before approving the full candidate.

\part{Business Capabilities}
\chapter{CT1 - Individual Care Delivery}\label{cap:bc-e85e18f4}
\noindent\textbf{Reference ID:} \texttt{BC-E85E18F4}\quad \textbf{Status:} Established working hierarchy\par
\noindent\textit{Source: Canonical Business Capability model / working map}\par
\textbf{Intent:} Individual Care Delivery identifies the organisational ability to undertake the stated health-service responsibility. The capability describes what the enterprise must be able to do; it does not assert that Harmonia performs or owns that activity.

\textbf{Boundary:} Harmonia relevance may arise because this activity establishes the context in which information is created, needed, interpreted, related, exchanged or presented. Responsibility for the underlying business decision or care activity remains with the appropriate actor or organisation.

\reviewstatus{Check scope and information relevance; no business authority transfers to Harmonia.}
\section{CT2 - Individual Health Assessment}\label{cap:bc-aa0d9b50}
\vspace{-0.5em}
\begin{table}[htbp]
  \centering
  \footnotesize
  \begin{tabularx}{1\textwidth}{|l|l|X|}
    \hline
    \multicolumn{3}{c}{\textbf{Capability Metadata}} \\
    \hline
    \textbf{Name} & \textbf{Value} & \textbf{Comment and/or Description} \\
    \hline
    \textbf{Reference Id} & \texttt{BC-AA0D9B50} & \\
    \hline
    \textbf{Status} & Established working hierarchy & \\
    \hline
    \textbf{Ancestry} &  & \\
    \hline
    \textbf{Source} &  & \\
    \hline
    \textbf{Enabled By} &  & \\
    \hline
    \textbf{Supports} &  & \\
    \hline
    \textbf{Derived From} &  & \\
    \hline
    \textbf{Reviewed Status} &  & \\
    \hline
  \end{tabularx}
\end{table}

\begin{table}[htbp]
  \centering
  \footnotesize
  \begin{tabularx}{1\textwidth}{|l|l|X|}
    \hline
    \multicolumn{3}{c}{\textbf{Contained Capabilities}} \\
    \hline
    \textbf{Capability Name} & \textbf{Capability Id} & \textbf{Comment} \\
    \hline
    Health History Assessment & \texttt{cap:bc-6956b3fa} (at \ref{cap:bc-6956b3fa}) & \\
    \hline
    Clinical Examination & \texttt{cap:bc-69ce36b4} (at \ref{cap:bc-69ce36b4}) & \\
    \hline
    Health Status Measurement & \texttt{cap:bc-3f88f485} (at \ref{cap:bc-3f88f485}) & \\
    \hline
  \end{tabularx}
\end{table}

\textbf{Intent:} Individual Health Assessment identifies the organisational ability to undertake the stated health-service responsibility. The capability describes what the enterprise must be able to do; it does not assert that Harmonia performs or owns that activity.

\textbf{Boundary:} Harmonia relevance may arise because this activity establishes the context in which information is created, needed, interpreted, related, exchanged or presented. Responsibility for the underlying business decision or care activity remains with the appropriate actor or organisation.

\subsection{CT3 - Health History Assessment}\label{cap:bc-6956b3fa}
\vspace{-0.5em}
\begin{table}[htbp]
  \centering
  \footnotesize
  \begin{tabularx}{1\textwidth}{|l|l|X|}
    \hline
    \multicolumn{3}{c}{\textbf{Capability Metadata}} \\
    \hline
    \textbf{Name} & \textbf{Value} & \textbf{Comment and/or Description} \\
    \hline
    \textbf{Reference Id} & \texttt{BC-6956B3FA} & \\
    \hline
    \textbf{Status} & Established working hierarchy & \\
    \hline
    \textbf{Ancestry} &  & \\
    \hline
    \textbf{Source} &  & \\
    \hline
    \textbf{Enabled By} &  & \\
    \hline
    \textbf{Supports} &  & \\
    \hline
    \textbf{Derived From} &  & \\
    \hline
    \textbf{Reviewed Status} & Review meaning, business boundary, information context and any unresolved ancestry. & \\
    \hline
  \end{tabularx}
\end{table}

\begin{table}[htbp]
  \centering
  \footnotesize
  \begin{tabularx}{1\textwidth}{|l|l|X|}
    \hline
    \multicolumn{3}{c}{\textbf{Contained Features}} \\
    \hline
    \textbf{Feature Name} & \textbf{Feature Id} & \textbf{Comment} \\
    \hline
    \multicolumn{3}{|l|}{\emph{No contained Features are currently identified.}} \\
    \hline
  \end{tabularx}
\end{table}

\textbf{Intent:} Health History Assessment identifies the organisational ability to undertake the stated health-service responsibility. The capability describes what the enterprise must be able to do; it does not assert that Harmonia performs or owns that activity.

\textbf{Boundary:} Harmonia relevance may arise because this activity establishes the context in which information is created, needed, interpreted, related, exchanged or presented. Responsibility for the underlying business decision or care activity remains with the appropriate actor or organisation.

\subsection{CT3 - Clinical Examination}\label{cap:bc-69ce36b4}
\vspace{-0.5em}
\begin{table}[htbp]
  \centering
  \footnotesize
  \begin{tabularx}{1\textwidth}{|l|l|X|}
    \hline
    \multicolumn{3}{c}{\textbf{Capability Metadata}} \\
    \hline
    \textbf{Name} & \textbf{Value} & \textbf{Comment and/or Description} \\
    \hline
    \textbf{Reference Id} & \texttt{BC-69CE36B4} & \\
    \hline
    \textbf{Status} & Established working hierarchy & \\
    \hline
    \textbf{Ancestry} &  & \\
    \hline
    \textbf{Source} &  & \\
    \hline
    \textbf{Enabled By} &  & \\
    \hline
    \textbf{Supports} &  & \\
    \hline
    \textbf{Derived From} &  & \\
    \hline
    \textbf{Reviewed Status} & Review meaning, business boundary, information context and any unresolved ancestry. & \\
    \hline
  \end{tabularx}
\end{table}

\begin{table}[htbp]
  \centering
  \footnotesize
  \begin{tabularx}{1\textwidth}{|l|l|X|}
    \hline
    \multicolumn{3}{c}{\textbf{Contained Features}} \\
    \hline
    \textbf{Feature Name} & \textbf{Feature Id} & \textbf{Comment} \\
    \hline
    \multicolumn{3}{|l|}{\emph{No contained Features are currently identified.}} \\
    \hline
  \end{tabularx}
\end{table}

\textbf{Intent:} Clinical Examination identifies the organisational ability to undertake the stated health-service responsibility. The capability describes what the enterprise must be able to do; it does not assert that Harmonia performs or owns that activity.

\textbf{Boundary:} Harmonia relevance may arise because this activity establishes the context in which information is created, needed, interpreted, related, exchanged or presented. Responsibility for the underlying business decision or care activity remains with the appropriate actor or organisation.

\subsection{CT3 - Health Status Measurement}\label{cap:bc-3f88f485}
\vspace{-0.5em}
\begin{table}[htbp]
  \centering
  \footnotesize
  \begin{tabularx}{1\textwidth}{|l|l|X|}
    \hline
    \multicolumn{3}{c}{\textbf{Capability Metadata}} \\
    \hline
    \textbf{Name} & \textbf{Value} & \textbf{Comment and/or Description} \\
    \hline
    \textbf{Reference Id} & \texttt{BC-3F88F485} & \\
    \hline
    \textbf{Status} & Established working hierarchy & \\
    \hline
    \textbf{Ancestry} &  & \\
    \hline
    \textbf{Source} &  & \\
    \hline
    \textbf{Enabled By} &  & \\
    \hline
    \textbf{Supports} &  & \\
    \hline
    \textbf{Derived From} &  & \\
    \hline
    \textbf{Reviewed Status} & Review meaning, business boundary, information context and any unresolved ancestry. & \\
    \hline
  \end{tabularx}
\end{table}

\begin{table}[htbp]
  \centering
  \footnotesize
  \begin{tabularx}{1\textwidth}{|l|l|X|}
    \hline
    \multicolumn{3}{c}{\textbf{Contained Features}} \\
    \hline
    \textbf{Feature Name} & \textbf{Feature Id} & \textbf{Comment} \\
    \hline
    \multicolumn{3}{|l|}{\emph{No contained Features are currently identified.}} \\
    \hline
  \end{tabularx}
\end{table}

\textbf{Intent:} Health Status Measurement identifies the organisational ability to undertake the stated health-service responsibility. The capability describes what the enterprise must be able to do; it does not assert that Harmonia performs or owns that activity.

\textbf{Boundary:} Harmonia relevance may arise because this activity establishes the context in which information is created, needed, interpreted, related, exchanged or presented. Responsibility for the underlying business decision or care activity remains with the appropriate actor or organisation.

\section{CT2 - Clinical Determination}\label{cap:bc-201a989e}
\vspace{-0.5em}
\begin{table}[htbp]
  \centering
  \footnotesize
  \begin{tabularx}{1\textwidth}{|l|l|X|}
    \hline
    \multicolumn{3}{c}{\textbf{Capability Metadata}} \\
    \hline
    \textbf{Name} & \textbf{Value} & \textbf{Comment and/or Description} \\
    \hline
    \textbf{Reference Id} & \texttt{BC-201A989E} & \\
    \hline
    \textbf{Status} & Established working hierarchy & \\
    \hline
    \textbf{Ancestry} &  & \\
    \hline
    \textbf{Source} &  & \\
    \hline
    \textbf{Enabled By} &  & \\
    \hline
    \textbf{Supports} &  & \\
    \hline
    \textbf{Derived From} &  & \\
    \hline
    \textbf{Reviewed Status} &  & \\
    \hline
  \end{tabularx}
\end{table}

\begin{table}[htbp]
  \centering
  \footnotesize
  \begin{tabularx}{1\textwidth}{|l|l|X|}
    \hline
    \multicolumn{3}{c}{\textbf{Contained Capabilities}} \\
    \hline
    \textbf{Capability Name} & \textbf{Capability Id} & \textbf{Comment} \\
    \hline
    Clinical Finding Interpretation & \texttt{cap:bc-20ca91e6} (at \ref{cap:bc-20ca91e6}) & \\
    \hline
    Diagnosis \& Problem Determination & \texttt{cap:bc-a1ac4a2a} (at \ref{cap:bc-a1ac4a2a}) & \\
    \hline
    Health Risk \& Prognostic Determination & \texttt{cap:bc-75e01b0d} (at \ref{cap:bc-75e01b0d}) & \\
    \hline
  \end{tabularx}
\end{table}

\textbf{Intent:} Clinical Determination identifies the organisational ability to undertake the stated health-service responsibility. The capability describes what the enterprise must be able to do; it does not assert that Harmonia performs or owns that activity.

\textbf{Boundary:} Harmonia relevance may arise because this activity establishes the context in which information is created, needed, interpreted, related, exchanged or presented. Responsibility for the underlying business decision or care activity remains with the appropriate actor or organisation.

\subsection{CT3 - Clinical Finding Interpretation}\label{cap:bc-20ca91e6}
\vspace{-0.5em}
\begin{table}[htbp]
  \centering
  \footnotesize
  \begin{tabularx}{1\textwidth}{|l|l|X|}
    \hline
    \multicolumn{3}{c}{\textbf{Capability Metadata}} \\
    \hline
    \textbf{Name} & \textbf{Value} & \textbf{Comment and/or Description} \\
    \hline
    \textbf{Reference Id} & \texttt{BC-20CA91E6} & \\
    \hline
    \textbf{Status} & Established working hierarchy & \\
    \hline
    \textbf{Ancestry} &  & \\
    \hline
    \textbf{Source} &  & \\
    \hline
    \textbf{Enabled By} &  & \\
    \hline
    \textbf{Supports} &  & \\
    \hline
    \textbf{Derived From} &  & \\
    \hline
    \textbf{Reviewed Status} & Review meaning, business boundary, information context and any unresolved ancestry. & \\
    \hline
  \end{tabularx}
\end{table}

\begin{table}[htbp]
  \centering
  \footnotesize
  \begin{tabularx}{1\textwidth}{|l|l|X|}
    \hline
    \multicolumn{3}{c}{\textbf{Contained Features}} \\
    \hline
    \textbf{Feature Name} & \textbf{Feature Id} & \textbf{Comment} \\
    \hline
    \multicolumn{3}{|l|}{\emph{No contained Features are currently identified.}} \\
    \hline
  \end{tabularx}
\end{table}

\textbf{Intent:} Clinical Finding Interpretation identifies the organisational ability to undertake the stated health-service responsibility. The capability describes what the enterprise must be able to do; it does not assert that Harmonia performs or owns that activity.

\textbf{Boundary:} Harmonia relevance may arise because this activity establishes the context in which information is created, needed, interpreted, related, exchanged or presented. Responsibility for the underlying business decision or care activity remains with the appropriate actor or organisation.

\subsection{CT3 - Diagnosis \& Problem Determination}\label{cap:bc-a1ac4a2a}
\vspace{-0.5em}
\begin{table}[htbp]
  \centering
  \footnotesize
  \begin{tabularx}{1\textwidth}{|l|l|X|}
    \hline
    \multicolumn{3}{c}{\textbf{Capability Metadata}} \\
    \hline
    \textbf{Name} & \textbf{Value} & \textbf{Comment and/or Description} \\
    \hline
    \textbf{Reference Id} & \texttt{BC-A1AC4A2A} & \\
    \hline
    \textbf{Status} & Established working hierarchy & \\
    \hline
    \textbf{Ancestry} &  & \\
    \hline
    \textbf{Source} &  & \\
    \hline
    \textbf{Enabled By} &  & \\
    \hline
    \textbf{Supports} &  & \\
    \hline
    \textbf{Derived From} &  & \\
    \hline
    \textbf{Reviewed Status} & Review meaning, business boundary, information context and any unresolved ancestry. & \\
    \hline
  \end{tabularx}
\end{table}

\begin{table}[htbp]
  \centering
  \footnotesize
  \begin{tabularx}{1\textwidth}{|l|l|X|}
    \hline
    \multicolumn{3}{c}{\textbf{Contained Features}} \\
    \hline
    \textbf{Feature Name} & \textbf{Feature Id} & \textbf{Comment} \\
    \hline
    \multicolumn{3}{|l|}{\emph{No contained Features are currently identified.}} \\
    \hline
  \end{tabularx}
\end{table}

\textbf{Intent:} Diagnosis \& Problem Determination identifies the organisational ability to undertake the stated health-service responsibility. The capability describes what the enterprise must be able to do; it does not assert that Harmonia performs or owns that activity.

\textbf{Boundary:} Harmonia relevance may arise because this activity establishes the context in which information is created, needed, interpreted, related, exchanged or presented. Responsibility for the underlying business decision or care activity remains with the appropriate actor or organisation.

\subsection{CT3 - Health Risk \& Prognostic Determination}\label{cap:bc-75e01b0d}
\vspace{-0.5em}
\begin{table}[htbp]
  \centering
  \footnotesize
  \begin{tabularx}{1\textwidth}{|l|l|X|}
    \hline
    \multicolumn{3}{c}{\textbf{Capability Metadata}} \\
    \hline
    \textbf{Name} & \textbf{Value} & \textbf{Comment and/or Description} \\
    \hline
    \textbf{Reference Id} & \texttt{BC-75E01B0D} & \\
    \hline
    \textbf{Status} & Established working hierarchy & \\
    \hline
    \textbf{Ancestry} &  & \\
    \hline
    \textbf{Source} &  & \\
    \hline
    \textbf{Enabled By} &  & \\
    \hline
    \textbf{Supports} &  & \\
    \hline
    \textbf{Derived From} &  & \\
    \hline
    \textbf{Reviewed Status} & Review meaning, business boundary, information context and any unresolved ancestry. & \\
    \hline
  \end{tabularx}
\end{table}

\begin{table}[htbp]
  \centering
  \footnotesize
  \begin{tabularx}{1\textwidth}{|l|l|X|}
    \hline
    \multicolumn{3}{c}{\textbf{Contained Features}} \\
    \hline
    \textbf{Feature Name} & \textbf{Feature Id} & \textbf{Comment} \\
    \hline
    \multicolumn{3}{|l|}{\emph{No contained Features are currently identified.}} \\
    \hline
  \end{tabularx}
\end{table}

\textbf{Intent:} Health Risk \& Prognostic Determination identifies the organisational ability to undertake the stated health-service responsibility. The capability describes what the enterprise must be able to do; it does not assert that Harmonia performs or owns that activity.

\textbf{Boundary:} Harmonia relevance may arise because this activity establishes the context in which information is created, needed, interpreted, related, exchanged or presented. Responsibility for the underlying business decision or care activity remains with the appropriate actor or organisation.

\section{CT2 - Individual Care Planning}\label{cap:bc-f657b57a}
\vspace{-0.5em}
\begin{table}[htbp]
  \centering
  \footnotesize
  \begin{tabularx}{1\textwidth}{|l|l|X|}
    \hline
    \multicolumn{3}{c}{\textbf{Capability Metadata}} \\
    \hline
    \textbf{Name} & \textbf{Value} & \textbf{Comment and/or Description} \\
    \hline
    \textbf{Reference Id} & \texttt{BC-F657B57A} & \\
    \hline
    \textbf{Status} & Established working hierarchy & \\
    \hline
    \textbf{Ancestry} &  & \\
    \hline
    \textbf{Source} &  & \\
    \hline
    \textbf{Enabled By} &  & \\
    \hline
    \textbf{Supports} &  & \\
    \hline
    \textbf{Derived From} &  & \\
    \hline
    \textbf{Reviewed Status} &  & \\
    \hline
  \end{tabularx}
\end{table}

\begin{table}[htbp]
  \centering
  \footnotesize
  \begin{tabularx}{1\textwidth}{|l|l|X|}
    \hline
    \multicolumn{3}{c}{\textbf{Contained Capabilities}} \\
    \hline
    \textbf{Capability Name} & \textbf{Capability Id} & \textbf{Comment} \\
    \hline
    Care Goal Definition & \texttt{cap:bc-e8bba49d} (at \ref{cap:bc-e8bba49d}) & \\
    \hline
    Care Response Planning & \texttt{cap:bc-3626ad34} (at \ref{cap:bc-3626ad34}) & \\
    \hline
  \end{tabularx}
\end{table}

\textbf{Intent:} Individual Care Planning identifies the organisational ability to undertake the stated health-service responsibility. The capability describes what the enterprise must be able to do; it does not assert that Harmonia performs or owns that activity.

\textbf{Boundary:} Harmonia relevance may arise because this activity establishes the context in which information is created, needed, interpreted, related, exchanged or presented. Responsibility for the underlying business decision or care activity remains with the appropriate actor or organisation.

\subsection{CT3 - Care Goal Definition}\label{cap:bc-e8bba49d}
\vspace{-0.5em}
\begin{table}[htbp]
  \centering
  \footnotesize
  \begin{tabularx}{1\textwidth}{|l|l|X|}
    \hline
    \multicolumn{3}{c}{\textbf{Capability Metadata}} \\
    \hline
    \textbf{Name} & \textbf{Value} & \textbf{Comment and/or Description} \\
    \hline
    \textbf{Reference Id} & \texttt{BC-E8BBA49D} & \\
    \hline
    \textbf{Status} & Established working hierarchy & \\
    \hline
    \textbf{Ancestry} &  & \\
    \hline
    \textbf{Source} &  & \\
    \hline
    \textbf{Enabled By} &  & \\
    \hline
    \textbf{Supports} &  & \\
    \hline
    \textbf{Derived From} &  & \\
    \hline
    \textbf{Reviewed Status} & Review meaning, business boundary, information context and any unresolved ancestry. & \\
    \hline
  \end{tabularx}
\end{table}

\begin{table}[htbp]
  \centering
  \footnotesize
  \begin{tabularx}{1\textwidth}{|l|l|X|}
    \hline
    \multicolumn{3}{c}{\textbf{Contained Features}} \\
    \hline
    \textbf{Feature Name} & \textbf{Feature Id} & \textbf{Comment} \\
    \hline
    \multicolumn{3}{|l|}{\emph{No contained Features are currently identified.}} \\
    \hline
  \end{tabularx}
\end{table}

\textbf{Intent:} Care Goal Definition identifies the organisational ability to undertake the stated health-service responsibility. The capability describes what the enterprise must be able to do; it does not assert that Harmonia performs or owns that activity.

\textbf{Boundary:} Harmonia relevance may arise because this activity establishes the context in which information is created, needed, interpreted, related, exchanged or presented. Responsibility for the underlying business decision or care activity remains with the appropriate actor or organisation.

\subsection{CT3 - Care Response Planning}\label{cap:bc-3626ad34}
\vspace{-0.5em}
\begin{table}[htbp]
  \centering
  \footnotesize
  \begin{tabularx}{1\textwidth}{|l|l|X|}
    \hline
    \multicolumn{3}{c}{\textbf{Capability Metadata}} \\
    \hline
    \textbf{Name} & \textbf{Value} & \textbf{Comment and/or Description} \\
    \hline
    \textbf{Reference Id} & \texttt{BC-3626AD34} & \\
    \hline
    \textbf{Status} & Established working hierarchy & \\
    \hline
    \textbf{Ancestry} &  & \\
    \hline
    \textbf{Source} &  & \\
    \hline
    \textbf{Enabled By} &  & \\
    \hline
    \textbf{Supports} &  & \\
    \hline
    \textbf{Derived From} &  & \\
    \hline
    \textbf{Reviewed Status} & Review meaning, business boundary, information context and any unresolved ancestry. & \\
    \hline
  \end{tabularx}
\end{table}

\begin{table}[htbp]
  \centering
  \footnotesize
  \begin{tabularx}{1\textwidth}{|l|l|X|}
    \hline
    \multicolumn{3}{c}{\textbf{Contained Features}} \\
    \hline
    \textbf{Feature Name} & \textbf{Feature Id} & \textbf{Comment} \\
    \hline
    \multicolumn{3}{|l|}{\emph{No contained Features are currently identified.}} \\
    \hline
  \end{tabularx}
\end{table}

\textbf{Intent:} Care Response Planning identifies the organisational ability to undertake the stated health-service responsibility. The capability describes what the enterprise must be able to do; it does not assert that Harmonia performs or owns that activity.

\textbf{Boundary:} Harmonia relevance may arise because this activity establishes the context in which information is created, needed, interpreted, related, exchanged or presented. Responsibility for the underlying business decision or care activity remains with the appropriate actor or organisation.

\section{CT2 - Clinical Intervention}\label{cap:bc-3cbc77c1}
\vspace{-0.5em}
\begin{table}[htbp]
  \centering
  \footnotesize
  \begin{tabularx}{1\textwidth}{|l|l|X|}
    \hline
    \multicolumn{3}{c}{\textbf{Capability Metadata}} \\
    \hline
    \textbf{Name} & \textbf{Value} & \textbf{Comment and/or Description} \\
    \hline
    \textbf{Reference Id} & \texttt{BC-3CBC77C1} & \\
    \hline
    \textbf{Status} & Established working hierarchy & \\
    \hline
    \textbf{Ancestry} &  & \\
    \hline
    \textbf{Source} &  & \\
    \hline
    \textbf{Enabled By} &  & \\
    \hline
    \textbf{Supports} &  & \\
    \hline
    \textbf{Derived From} &  & \\
    \hline
    \textbf{Reviewed Status} &  & \\
    \hline
  \end{tabularx}
\end{table}

\begin{table}[htbp]
  \centering
  \footnotesize
  \begin{tabularx}{1\textwidth}{|l|l|X|}
    \hline
    \multicolumn{3}{c}{\textbf{Contained Capabilities}} \\
    \hline
    \textbf{Capability Name} & \textbf{Capability Id} & \textbf{Comment} \\
    \hline
    \multicolumn{3}{|l|}{\emph{No contained capabilities are currently identified.}} \\
    \hline
  \end{tabularx}
\end{table}

\textbf{Intent:} Clinical Intervention identifies the organisational ability to undertake the stated health-service responsibility. The capability describes what the enterprise must be able to do; it does not assert that Harmonia performs or owns that activity.

\textbf{Boundary:} Harmonia relevance may arise because this activity establishes the context in which information is created, needed, interpreted, related, exchanged or presented. Responsibility for the underlying business decision or care activity remains with the appropriate actor or organisation.

Unresolved / stopping-point note: [STOPPING POINT] CT3 taxonomy intentionally not required for R1/R2

\section{CT2 - Individual Health Monitoring}\label{cap:bc-dd466a23}
\vspace{-0.5em}
\begin{table}[htbp]
  \centering
  \footnotesize
  \begin{tabularx}{1\textwidth}{|l|l|X|}
    \hline
    \multicolumn{3}{c}{\textbf{Capability Metadata}} \\
    \hline
    \textbf{Name} & \textbf{Value} & \textbf{Comment and/or Description} \\
    \hline
    \textbf{Reference Id} & \texttt{BC-DD466A23} & \\
    \hline
    \textbf{Status} & Established working hierarchy & \\
    \hline
    \textbf{Ancestry} &  & \\
    \hline
    \textbf{Source} &  & \\
    \hline
    \textbf{Enabled By} &  & \\
    \hline
    \textbf{Supports} &  & \\
    \hline
    \textbf{Derived From} &  & \\
    \hline
    \textbf{Reviewed Status} &  & \\
    \hline
  \end{tabularx}
\end{table}

\begin{table}[htbp]
  \centering
  \footnotesize
  \begin{tabularx}{1\textwidth}{|l|l|X|}
    \hline
    \multicolumn{3}{c}{\textbf{Contained Capabilities}} \\
    \hline
    \textbf{Capability Name} & \textbf{Capability Id} & \textbf{Comment} \\
    \hline
    Health State Monitoring & \texttt{cap:bc-7e2744d3} (at \ref{cap:bc-7e2744d3}) & \\
    \hline
    Care Response Monitoring & \texttt{cap:bc-02c5dbc7} (at \ref{cap:bc-02c5dbc7}) & \\
    \hline
  \end{tabularx}
\end{table}

\textbf{Intent:} Individual Health Monitoring identifies the organisational ability to undertake the stated health-service responsibility. The capability describes what the enterprise must be able to do; it does not assert that Harmonia performs or owns that activity.

\textbf{Boundary:} Harmonia relevance may arise because this activity establishes the context in which information is created, needed, interpreted, related, exchanged or presented. Responsibility for the underlying business decision or care activity remains with the appropriate actor or organisation.

\subsection{CT3 - Health State Monitoring}\label{cap:bc-7e2744d3}
\vspace{-0.5em}
\begin{table}[htbp]
  \centering
  \footnotesize
  \begin{tabularx}{1\textwidth}{|l|l|X|}
    \hline
    \multicolumn{3}{c}{\textbf{Capability Metadata}} \\
    \hline
    \textbf{Name} & \textbf{Value} & \textbf{Comment and/or Description} \\
    \hline
    \textbf{Reference Id} & \texttt{BC-7E2744D3} & \\
    \hline
    \textbf{Status} & Established working hierarchy & \\
    \hline
    \textbf{Ancestry} &  & \\
    \hline
    \textbf{Source} &  & \\
    \hline
    \textbf{Enabled By} &  & \\
    \hline
    \textbf{Supports} &  & \\
    \hline
    \textbf{Derived From} &  & \\
    \hline
    \textbf{Reviewed Status} & Review meaning, business boundary, information context and any unresolved ancestry. & \\
    \hline
  \end{tabularx}
\end{table}

\begin{table}[htbp]
  \centering
  \footnotesize
  \begin{tabularx}{1\textwidth}{|l|l|X|}
    \hline
    \multicolumn{3}{c}{\textbf{Contained Features}} \\
    \hline
    \textbf{Feature Name} & \textbf{Feature Id} & \textbf{Comment} \\
    \hline
    \multicolumn{3}{|l|}{\emph{No contained Features are currently identified.}} \\
    \hline
  \end{tabularx}
\end{table}

\textbf{Intent:} Health State Monitoring identifies the organisational ability to undertake the stated health-service responsibility. The capability describes what the enterprise must be able to do; it does not assert that Harmonia performs or owns that activity.

\textbf{Boundary:} Harmonia relevance may arise because this activity establishes the context in which information is created, needed, interpreted, related, exchanged or presented. Responsibility for the underlying business decision or care activity remains with the appropriate actor or organisation.

\subsection{CT3 - Care Response Monitoring}\label{cap:bc-02c5dbc7}
\vspace{-0.5em}
\begin{table}[htbp]
  \centering
  \footnotesize
  \begin{tabularx}{1\textwidth}{|l|l|X|}
    \hline
    \multicolumn{3}{c}{\textbf{Capability Metadata}} \\
    \hline
    \textbf{Name} & \textbf{Value} & \textbf{Comment and/or Description} \\
    \hline
    \textbf{Reference Id} & \texttt{BC-02C5DBC7} & \\
    \hline
    \textbf{Status} & Established working hierarchy & \\
    \hline
    \textbf{Ancestry} &  & \\
    \hline
    \textbf{Source} &  & \\
    \hline
    \textbf{Enabled By} &  & \\
    \hline
    \textbf{Supports} &  & \\
    \hline
    \textbf{Derived From} &  & \\
    \hline
    \textbf{Reviewed Status} & Review meaning, business boundary, information context and any unresolved ancestry. & \\
    \hline
  \end{tabularx}
\end{table}

\begin{table}[htbp]
  \centering
  \footnotesize
  \begin{tabularx}{1\textwidth}{|l|l|X|}
    \hline
    \multicolumn{3}{c}{\textbf{Contained Features}} \\
    \hline
    \textbf{Feature Name} & \textbf{Feature Id} & \textbf{Comment} \\
    \hline
    \multicolumn{3}{|l|}{\emph{No contained Features are currently identified.}} \\
    \hline
  \end{tabularx}
\end{table}

\textbf{Intent:} Care Response Monitoring identifies the organisational ability to undertake the stated health-service responsibility. The capability describes what the enterprise must be able to do; it does not assert that Harmonia performs or owns that activity.

\textbf{Boundary:} Harmonia relevance may arise because this activity establishes the context in which information is created, needed, interpreted, related, exchanged or presented. Responsibility for the underlying business decision or care activity remains with the appropriate actor or organisation.

\section{CT2 - Individual Care Review}\label{cap:bc-391ec3b9}
\vspace{-0.5em}
\begin{table}[htbp]
  \centering
  \footnotesize
  \begin{tabularx}{1\textwidth}{|l|l|X|}
    \hline
    \multicolumn{3}{c}{\textbf{Capability Metadata}} \\
    \hline
    \textbf{Name} & \textbf{Value} & \textbf{Comment and/or Description} \\
    \hline
    \textbf{Reference Id} & \texttt{BC-391EC3B9} & \\
    \hline
    \textbf{Status} & Established working hierarchy & \\
    \hline
    \textbf{Ancestry} &  & \\
    \hline
    \textbf{Source} &  & \\
    \hline
    \textbf{Enabled By} &  & \\
    \hline
    \textbf{Supports} &  & \\
    \hline
    \textbf{Derived From} &  & \\
    \hline
    \textbf{Reviewed Status} &  & \\
    \hline
  \end{tabularx}
\end{table}

\begin{table}[htbp]
  \centering
  \footnotesize
  \begin{tabularx}{1\textwidth}{|l|l|X|}
    \hline
    \multicolumn{3}{c}{\textbf{Contained Capabilities}} \\
    \hline
    \textbf{Capability Name} & \textbf{Capability Id} & \textbf{Comment} \\
    \hline
    Care Appropriateness Review & \texttt{cap:bc-a52e0a33} (at \ref{cap:bc-a52e0a33}) & \\
    \hline
    Care Effectiveness Review & \texttt{cap:bc-31953e77} (at \ref{cap:bc-31953e77}) & \\
    \hline
  \end{tabularx}
\end{table}

\textbf{Intent:} Individual Care Review identifies the organisational ability to undertake the stated health-service responsibility. The capability describes what the enterprise must be able to do; it does not assert that Harmonia performs or owns that activity.

\textbf{Boundary:} Harmonia relevance may arise because this activity establishes the context in which information is created, needed, interpreted, related, exchanged or presented. Responsibility for the underlying business decision or care activity remains with the appropriate actor or organisation.

\subsection{CT3 - Care Appropriateness Review}\label{cap:bc-a52e0a33}
\vspace{-0.5em}
\begin{table}[htbp]
  \centering
  \footnotesize
  \begin{tabularx}{1\textwidth}{|l|l|X|}
    \hline
    \multicolumn{3}{c}{\textbf{Capability Metadata}} \\
    \hline
    \textbf{Name} & \textbf{Value} & \textbf{Comment and/or Description} \\
    \hline
    \textbf{Reference Id} & \texttt{BC-A52E0A33} & \\
    \hline
    \textbf{Status} & Established working hierarchy & \\
    \hline
    \textbf{Ancestry} &  & \\
    \hline
    \textbf{Source} &  & \\
    \hline
    \textbf{Enabled By} &  & \\
    \hline
    \textbf{Supports} &  & \\
    \hline
    \textbf{Derived From} &  & \\
    \hline
    \textbf{Reviewed Status} & Review meaning, business boundary, information context and any unresolved ancestry. & \\
    \hline
  \end{tabularx}
\end{table}

\begin{table}[htbp]
  \centering
  \footnotesize
  \begin{tabularx}{1\textwidth}{|l|l|X|}
    \hline
    \multicolumn{3}{c}{\textbf{Contained Features}} \\
    \hline
    \textbf{Feature Name} & \textbf{Feature Id} & \textbf{Comment} \\
    \hline
    \multicolumn{3}{|l|}{\emph{No contained Features are currently identified.}} \\
    \hline
  \end{tabularx}
\end{table}

\textbf{Intent:} Care Appropriateness Review identifies the organisational ability to undertake the stated health-service responsibility. The capability describes what the enterprise must be able to do; it does not assert that Harmonia performs or owns that activity.

\textbf{Boundary:} Harmonia relevance may arise because this activity establishes the context in which information is created, needed, interpreted, related, exchanged or presented. Responsibility for the underlying business decision or care activity remains with the appropriate actor or organisation.

\subsection{CT3 - Care Effectiveness Review}\label{cap:bc-31953e77}
\vspace{-0.5em}
\begin{table}[htbp]
  \centering
  \footnotesize
  \begin{tabularx}{1\textwidth}{|l|l|X|}
    \hline
    \multicolumn{3}{c}{\textbf{Capability Metadata}} \\
    \hline
    \textbf{Name} & \textbf{Value} & \textbf{Comment and/or Description} \\
    \hline
    \textbf{Reference Id} & \texttt{BC-31953E77} & \\
    \hline
    \textbf{Status} & Established working hierarchy & \\
    \hline
    \textbf{Ancestry} &  & \\
    \hline
    \textbf{Source} &  & \\
    \hline
    \textbf{Enabled By} &  & \\
    \hline
    \textbf{Supports} &  & \\
    \hline
    \textbf{Derived From} &  & \\
    \hline
    \textbf{Reviewed Status} & Review meaning, business boundary, information context and any unresolved ancestry. & \\
    \hline
  \end{tabularx}
\end{table}

\begin{table}[htbp]
  \centering
  \footnotesize
  \begin{tabularx}{1\textwidth}{|l|l|X|}
    \hline
    \multicolumn{3}{c}{\textbf{Contained Features}} \\
    \hline
    \textbf{Feature Name} & \textbf{Feature Id} & \textbf{Comment} \\
    \hline
    \multicolumn{3}{|l|}{\emph{No contained Features are currently identified.}} \\
    \hline
  \end{tabularx}
\end{table}

\textbf{Intent:} Care Effectiveness Review identifies the organisational ability to undertake the stated health-service responsibility. The capability describes what the enterprise must be able to do; it does not assert that Harmonia performs or owns that activity.

\textbf{Boundary:} Harmonia relevance may arise because this activity establishes the context in which information is created, needed, interpreted, related, exchanged or presented. Responsibility for the underlying business decision or care activity remains with the appropriate actor or organisation.

\chapter{CT1 - Care Access \& Coordination}\label{cap:bc-2c769a48}
\noindent\textbf{Reference ID:} \texttt{BC-2C769A48}\quad \textbf{Status:} Established working hierarchy\par
\noindent\textit{Source: Canonical Business Capability model / working map}\par
\textbf{Intent:} Care Access \& Coordination identifies the organisational ability to undertake the stated health-service responsibility. The capability describes what the enterprise must be able to do; it does not assert that Harmonia performs or owns that activity.

\textbf{Boundary:} Harmonia relevance may arise because this activity establishes the context in which information is created, needed, interpreted, related, exchanged or presented. Responsibility for the underlying business decision or care activity remains with the appropriate actor or organisation.

\reviewstatus{Check scope and information relevance; no business authority transfers to Harmonia.}
\section{CT2 - Care Access Management}\label{cap:bc-fa41a886}
\vspace{-0.5em}
\begin{table}[htbp]
  \centering
  \footnotesize
  \begin{tabularx}{1\textwidth}{|l|l|X|}
    \hline
    \multicolumn{3}{c}{\textbf{Capability Metadata}} \\
    \hline
    \textbf{Name} & \textbf{Value} & \textbf{Comment and/or Description} \\
    \hline
    \textbf{Reference Id} & \texttt{BC-FA41A886} & \\
    \hline
    \textbf{Status} & Established working hierarchy & \\
    \hline
    \textbf{Ancestry} &  & \\
    \hline
    \textbf{Source} &  & \\
    \hline
    \textbf{Enabled By} &  & \\
    \hline
    \textbf{Supports} &  & \\
    \hline
    \textbf{Derived From} &  & \\
    \hline
    \textbf{Reviewed Status} &  & \\
    \hline
  \end{tabularx}
\end{table}

\begin{table}[htbp]
  \centering
  \footnotesize
  \begin{tabularx}{1\textwidth}{|l|l|X|}
    \hline
    \multicolumn{3}{c}{\textbf{Contained Capabilities}} \\
    \hline
    \textbf{Capability Name} & \textbf{Capability Id} & \textbf{Comment} \\
    \hline
    Care Intake Management & \texttt{cap:bc-267f1d1b} (at \ref{cap:bc-267f1d1b}) & \\
    \hline
    Care Referral Management & \texttt{cap:bc-ea6b7eb9} (at \ref{cap:bc-ea6b7eb9}) & \\
    \hline
    Care Eligibility Determination & \texttt{cap:bc-4f461126} (at \ref{cap:bc-4f461126}) & \\
    \hline
  \end{tabularx}
\end{table}

\textbf{Intent:} Care Access Management identifies the organisational ability to undertake the stated health-service responsibility. The capability describes what the enterprise must be able to do; it does not assert that Harmonia performs or owns that activity.

\textbf{Boundary:} Harmonia relevance may arise because this activity establishes the context in which information is created, needed, interpreted, related, exchanged or presented. Responsibility for the underlying business decision or care activity remains with the appropriate actor or organisation.

\subsection{CT3 - Care Intake Management}\label{cap:bc-267f1d1b}
\vspace{-0.5em}
\begin{table}[htbp]
  \centering
  \footnotesize
  \begin{tabularx}{1\textwidth}{|l|l|X|}
    \hline
    \multicolumn{3}{c}{\textbf{Capability Metadata}} \\
    \hline
    \textbf{Name} & \textbf{Value} & \textbf{Comment and/or Description} \\
    \hline
    \textbf{Reference Id} & \texttt{BC-267F1D1B} & \\
    \hline
    \textbf{Status} & Established working hierarchy & \\
    \hline
    \textbf{Ancestry} &  & \\
    \hline
    \textbf{Source} &  & \\
    \hline
    \textbf{Enabled By} &  & \\
    \hline
    \textbf{Supports} &  & \\
    \hline
    \textbf{Derived From} &  & \\
    \hline
    \textbf{Reviewed Status} & Review meaning, business boundary, information context and any unresolved ancestry. & \\
    \hline
  \end{tabularx}
\end{table}

\begin{table}[htbp]
  \centering
  \footnotesize
  \begin{tabularx}{1\textwidth}{|l|l|X|}
    \hline
    \multicolumn{3}{c}{\textbf{Contained Features}} \\
    \hline
    \textbf{Feature Name} & \textbf{Feature Id} & \textbf{Comment} \\
    \hline
    \multicolumn{3}{|l|}{\emph{No contained Features are currently identified.}} \\
    \hline
  \end{tabularx}
\end{table}

\textbf{Intent:} Care Intake Management identifies the organisational ability to undertake the stated health-service responsibility. The capability describes what the enterprise must be able to do; it does not assert that Harmonia performs or owns that activity.

\textbf{Boundary:} Harmonia relevance may arise because this activity establishes the context in which information is created, needed, interpreted, related, exchanged or presented. Responsibility for the underlying business decision or care activity remains with the appropriate actor or organisation.

\subsection{CT3 - Care Referral Management}\label{cap:bc-ea6b7eb9}
\vspace{-0.5em}
\begin{table}[htbp]
  \centering
  \footnotesize
  \begin{tabularx}{1\textwidth}{|l|l|X|}
    \hline
    \multicolumn{3}{c}{\textbf{Capability Metadata}} \\
    \hline
    \textbf{Name} & \textbf{Value} & \textbf{Comment and/or Description} \\
    \hline
    \textbf{Reference Id} & \texttt{BC-EA6B7EB9} & \\
    \hline
    \textbf{Status} & Established working hierarchy & \\
    \hline
    \textbf{Ancestry} &  & \\
    \hline
    \textbf{Source} &  & \\
    \hline
    \textbf{Enabled By} &  & \\
    \hline
    \textbf{Supports} &  & \\
    \hline
    \textbf{Derived From} &  & \\
    \hline
    \textbf{Reviewed Status} & Review meaning, business boundary, information context and any unresolved ancestry. & \\
    \hline
  \end{tabularx}
\end{table}

\begin{table}[htbp]
  \centering
  \footnotesize
  \begin{tabularx}{1\textwidth}{|l|l|X|}
    \hline
    \multicolumn{3}{c}{\textbf{Contained Features}} \\
    \hline
    \textbf{Feature Name} & \textbf{Feature Id} & \textbf{Comment} \\
    \hline
    \multicolumn{3}{|l|}{\emph{No contained Features are currently identified.}} \\
    \hline
  \end{tabularx}
\end{table}

\textbf{Intent:} Care Referral Management identifies the organisational ability to undertake the stated health-service responsibility. The capability describes what the enterprise must be able to do; it does not assert that Harmonia performs or owns that activity.

\textbf{Boundary:} Harmonia relevance may arise because this activity establishes the context in which information is created, needed, interpreted, related, exchanged or presented. Responsibility for the underlying business decision or care activity remains with the appropriate actor or organisation.

\subsection{CT3 - Care Eligibility Determination}\label{cap:bc-4f461126}
\vspace{-0.5em}
\begin{table}[htbp]
  \centering
  \footnotesize
  \begin{tabularx}{1\textwidth}{|l|l|X|}
    \hline
    \multicolumn{3}{c}{\textbf{Capability Metadata}} \\
    \hline
    \textbf{Name} & \textbf{Value} & \textbf{Comment and/or Description} \\
    \hline
    \textbf{Reference Id} & \texttt{BC-4F461126} & \\
    \hline
    \textbf{Status} & Established working hierarchy & \\
    \hline
    \textbf{Ancestry} &  & \\
    \hline
    \textbf{Source} &  & \\
    \hline
    \textbf{Enabled By} &  & \\
    \hline
    \textbf{Supports} &  & \\
    \hline
    \textbf{Derived From} &  & \\
    \hline
    \textbf{Reviewed Status} & Review meaning, business boundary, information context and any unresolved ancestry. & \\
    \hline
  \end{tabularx}
\end{table}

\begin{table}[htbp]
  \centering
  \footnotesize
  \begin{tabularx}{1\textwidth}{|l|l|X|}
    \hline
    \multicolumn{3}{c}{\textbf{Contained Features}} \\
    \hline
    \textbf{Feature Name} & \textbf{Feature Id} & \textbf{Comment} \\
    \hline
    \multicolumn{3}{|l|}{\emph{No contained Features are currently identified.}} \\
    \hline
  \end{tabularx}
\end{table}

\textbf{Intent:} Care Eligibility Determination identifies the organisational ability to undertake the stated health-service responsibility. The capability describes what the enterprise must be able to do; it does not assert that Harmonia performs or owns that activity.

\textbf{Boundary:} Harmonia relevance may arise because this activity establishes the context in which information is created, needed, interpreted, related, exchanged or presented. Responsibility for the underlying business decision or care activity remains with the appropriate actor or organisation.

\section{CT2 - Care Scheduling \& Capacity Management}\label{cap:bc-fe20c7c9}
\vspace{-0.5em}
\begin{table}[htbp]
  \centering
  \footnotesize
  \begin{tabularx}{1\textwidth}{|l|l|X|}
    \hline
    \multicolumn{3}{c}{\textbf{Capability Metadata}} \\
    \hline
    \textbf{Name} & \textbf{Value} & \textbf{Comment and/or Description} \\
    \hline
    \textbf{Reference Id} & \texttt{BC-FE20C7C9} & \\
    \hline
    \textbf{Status} & Established working hierarchy & \\
    \hline
    \textbf{Ancestry} &  & \\
    \hline
    \textbf{Source} &  & \\
    \hline
    \textbf{Enabled By} &  & \\
    \hline
    \textbf{Supports} &  & \\
    \hline
    \textbf{Derived From} &  & \\
    \hline
    \textbf{Reviewed Status} &  & \\
    \hline
  \end{tabularx}
\end{table}

\begin{table}[htbp]
  \centering
  \footnotesize
  \begin{tabularx}{1\textwidth}{|l|l|X|}
    \hline
    \multicolumn{3}{c}{\textbf{Contained Capabilities}} \\
    \hline
    \textbf{Capability Name} & \textbf{Capability Id} & \textbf{Comment} \\
    \hline
    Care Appointment Management & \texttt{cap:bc-9000e197} (at \ref{cap:bc-9000e197}) & \\
    \hline
    Care-Place Allocation & \texttt{cap:bc-a1f26127} (at \ref{cap:bc-a1f26127}) & \\
    \hline
    Care-Place Readiness Management & \texttt{cap:bc-3284993b} (at \ref{cap:bc-3284993b}) & \\
    \hline
    Care Resource Deployment & \texttt{cap:bc-b53cd721} (at \ref{cap:bc-b53cd721}) & \\
    \hline
    Operational Care Flow Coordination & \texttt{cap:bc-e24e4abe} (at \ref{cap:bc-e24e4abe}) & \\
    \hline
    Operational Care Capacity Management & \texttt{cap:bc-c2d4c33e} (at \ref{cap:bc-c2d4c33e}) & \\
    \hline
  \end{tabularx}
\end{table}

\textbf{Intent:} Care Scheduling \& Capacity Management identifies the organisational ability to undertake the stated health-service responsibility. The capability describes what the enterprise must be able to do; it does not assert that Harmonia performs or owns that activity.

\textbf{Boundary:} Harmonia relevance may arise because this activity establishes the context in which information is created, needed, interpreted, related, exchanged or presented. Responsibility for the underlying business decision or care activity remains with the appropriate actor or organisation.

\subsection{CT3 - Care Appointment Management}\label{cap:bc-9000e197}
\vspace{-0.5em}
\begin{table}[htbp]
  \centering
  \footnotesize
  \begin{tabularx}{1\textwidth}{|l|l|X|}
    \hline
    \multicolumn{3}{c}{\textbf{Capability Metadata}} \\
    \hline
    \textbf{Name} & \textbf{Value} & \textbf{Comment and/or Description} \\
    \hline
    \textbf{Reference Id} & \texttt{BC-9000E197} & \\
    \hline
    \textbf{Status} & Established working hierarchy & \\
    \hline
    \textbf{Ancestry} &  & \\
    \hline
    \textbf{Source} &  & \\
    \hline
    \textbf{Enabled By} &  & \\
    \hline
    \textbf{Supports} &  & \\
    \hline
    \textbf{Derived From} &  & \\
    \hline
    \textbf{Reviewed Status} & Review meaning, business boundary, information context and any unresolved ancestry. & \\
    \hline
  \end{tabularx}
\end{table}

\begin{table}[htbp]
  \centering
  \footnotesize
  \begin{tabularx}{1\textwidth}{|l|l|X|}
    \hline
    \multicolumn{3}{c}{\textbf{Contained Features}} \\
    \hline
    \textbf{Feature Name} & \textbf{Feature Id} & \textbf{Comment} \\
    \hline
    \multicolumn{3}{|l|}{\emph{No contained Features are currently identified.}} \\
    \hline
  \end{tabularx}
\end{table}

\textbf{Intent:} Care Appointment Management identifies the organisational ability to undertake the stated health-service responsibility. The capability describes what the enterprise must be able to do; it does not assert that Harmonia performs or owns that activity.

\textbf{Boundary:} Harmonia relevance may arise because this activity establishes the context in which information is created, needed, interpreted, related, exchanged or presented. Responsibility for the underlying business decision or care activity remains with the appropriate actor or organisation.

\subsection{CT3 - Care-Place Allocation}\label{cap:bc-a1f26127}
\vspace{-0.5em}
\begin{table}[htbp]
  \centering
  \footnotesize
  \begin{tabularx}{1\textwidth}{|l|l|X|}
    \hline
    \multicolumn{3}{c}{\textbf{Capability Metadata}} \\
    \hline
    \textbf{Name} & \textbf{Value} & \textbf{Comment and/or Description} \\
    \hline
    \textbf{Reference Id} & \texttt{BC-A1F26127} & \\
    \hline
    \textbf{Status} & Established working hierarchy & \\
    \hline
    \textbf{Ancestry} &  & \\
    \hline
    \textbf{Source} &  & \\
    \hline
    \textbf{Enabled By} &  & \\
    \hline
    \textbf{Supports} &  & \\
    \hline
    \textbf{Derived From} &  & \\
    \hline
    \textbf{Reviewed Status} & Review meaning, business boundary, information context and any unresolved ancestry. & \\
    \hline
  \end{tabularx}
\end{table}

\begin{table}[htbp]
  \centering
  \footnotesize
  \begin{tabularx}{1\textwidth}{|l|l|X|}
    \hline
    \multicolumn{3}{c}{\textbf{Contained Features}} \\
    \hline
    \textbf{Feature Name} & \textbf{Feature Id} & \textbf{Comment} \\
    \hline
    \multicolumn{3}{|l|}{\emph{No contained Features are currently identified.}} \\
    \hline
  \end{tabularx}
\end{table}

\textbf{Intent:} Care-Place Allocation identifies the organisational ability to undertake the stated health-service responsibility. The capability describes what the enterprise must be able to do; it does not assert that Harmonia performs or owns that activity.

\textbf{Boundary:} Harmonia relevance may arise because this activity establishes the context in which information is created, needed, interpreted, related, exchanged or presented. Responsibility for the underlying business decision or care activity remains with the appropriate actor or organisation.

\subsection{CT3 - Care-Place Readiness Management}\label{cap:bc-3284993b}
\vspace{-0.5em}
\begin{table}[htbp]
  \centering
  \footnotesize
  \begin{tabularx}{1\textwidth}{|l|l|X|}
    \hline
    \multicolumn{3}{c}{\textbf{Capability Metadata}} \\
    \hline
    \textbf{Name} & \textbf{Value} & \textbf{Comment and/or Description} \\
    \hline
    \textbf{Reference Id} & \texttt{BC-3284993B} & \\
    \hline
    \textbf{Status} & Established working hierarchy & \\
    \hline
    \textbf{Ancestry} &  & \\
    \hline
    \textbf{Source} &  & \\
    \hline
    \textbf{Enabled By} &  & \\
    \hline
    \textbf{Supports} &  & \\
    \hline
    \textbf{Derived From} &  & \\
    \hline
    \textbf{Reviewed Status} & Review meaning, business boundary, information context and any unresolved ancestry. & \\
    \hline
  \end{tabularx}
\end{table}

\begin{table}[htbp]
  \centering
  \footnotesize
  \begin{tabularx}{1\textwidth}{|l|l|X|}
    \hline
    \multicolumn{3}{c}{\textbf{Contained Features}} \\
    \hline
    \textbf{Feature Name} & \textbf{Feature Id} & \textbf{Comment} \\
    \hline
    \multicolumn{3}{|l|}{\emph{No contained Features are currently identified.}} \\
    \hline
  \end{tabularx}
\end{table}

\textbf{Intent:} Care-Place Readiness Management identifies the organisational ability to undertake the stated health-service responsibility. The capability describes what the enterprise must be able to do; it does not assert that Harmonia performs or owns that activity.

\textbf{Boundary:} Harmonia relevance may arise because this activity establishes the context in which information is created, needed, interpreted, related, exchanged or presented. Responsibility for the underlying business decision or care activity remains with the appropriate actor or organisation.

\subsection{CT3 - Care Resource Deployment}\label{cap:bc-b53cd721}
\vspace{-0.5em}
\begin{table}[htbp]
  \centering
  \footnotesize
  \begin{tabularx}{1\textwidth}{|l|l|X|}
    \hline
    \multicolumn{3}{c}{\textbf{Capability Metadata}} \\
    \hline
    \textbf{Name} & \textbf{Value} & \textbf{Comment and/or Description} \\
    \hline
    \textbf{Reference Id} & \texttt{BC-B53CD721} & \\
    \hline
    \textbf{Status} & Established working hierarchy & \\
    \hline
    \textbf{Ancestry} &  & \\
    \hline
    \textbf{Source} &  & \\
    \hline
    \textbf{Enabled By} &  & \\
    \hline
    \textbf{Supports} &  & \\
    \hline
    \textbf{Derived From} &  & \\
    \hline
    \textbf{Reviewed Status} & Review meaning, business boundary, information context and any unresolved ancestry. & \\
    \hline
  \end{tabularx}
\end{table}

\begin{table}[htbp]
  \centering
  \footnotesize
  \begin{tabularx}{1\textwidth}{|l|l|X|}
    \hline
    \multicolumn{3}{c}{\textbf{Contained Features}} \\
    \hline
    \textbf{Feature Name} & \textbf{Feature Id} & \textbf{Comment} \\
    \hline
    \multicolumn{3}{|l|}{\emph{No contained Features are currently identified.}} \\
    \hline
  \end{tabularx}
\end{table}

\textbf{Intent:} Care Resource Deployment identifies the organisational ability to undertake the stated health-service responsibility. The capability describes what the enterprise must be able to do; it does not assert that Harmonia performs or owns that activity.

\textbf{Boundary:} Harmonia relevance may arise because this activity establishes the context in which information is created, needed, interpreted, related, exchanged or presented. Responsibility for the underlying business decision or care activity remains with the appropriate actor or organisation.

\subsection{CT3 - Operational Care Flow Coordination}\label{cap:bc-e24e4abe}
\vspace{-0.5em}
\begin{table}[htbp]
  \centering
  \footnotesize
  \begin{tabularx}{1\textwidth}{|l|l|X|}
    \hline
    \multicolumn{3}{c}{\textbf{Capability Metadata}} \\
    \hline
    \textbf{Name} & \textbf{Value} & \textbf{Comment and/or Description} \\
    \hline
    \textbf{Reference Id} & \texttt{BC-E24E4ABE} & \\
    \hline
    \textbf{Status} & Established working hierarchy & \\
    \hline
    \textbf{Ancestry} &  & \\
    \hline
    \textbf{Source} &  & \\
    \hline
    \textbf{Enabled By} &  & \\
    \hline
    \textbf{Supports} &  & \\
    \hline
    \textbf{Derived From} &  & \\
    \hline
    \textbf{Reviewed Status} & Review meaning, business boundary, information context and any unresolved ancestry. & \\
    \hline
  \end{tabularx}
\end{table}

\begin{table}[htbp]
  \centering
  \footnotesize
  \begin{tabularx}{1\textwidth}{|l|l|X|}
    \hline
    \multicolumn{3}{c}{\textbf{Contained Features}} \\
    \hline
    \textbf{Feature Name} & \textbf{Feature Id} & \textbf{Comment} \\
    \hline
    \multicolumn{3}{|l|}{\emph{No contained Features are currently identified.}} \\
    \hline
  \end{tabularx}
\end{table}

\textbf{Intent:} Operational Care Flow Coordination identifies the organisational ability to undertake the stated health-service responsibility. The capability describes what the enterprise must be able to do; it does not assert that Harmonia performs or owns that activity.

\textbf{Boundary:} Harmonia relevance may arise because this activity establishes the context in which information is created, needed, interpreted, related, exchanged or presented. Responsibility for the underlying business decision or care activity remains with the appropriate actor or organisation.

\subsection{CT3 - Operational Care Capacity Management}\label{cap:bc-c2d4c33e}
\vspace{-0.5em}
\begin{table}[htbp]
  \centering
  \footnotesize
  \begin{tabularx}{1\textwidth}{|l|l|X|}
    \hline
    \multicolumn{3}{c}{\textbf{Capability Metadata}} \\
    \hline
    \textbf{Name} & \textbf{Value} & \textbf{Comment and/or Description} \\
    \hline
    \textbf{Reference Id} & \texttt{BC-C2D4C33E} & \\
    \hline
    \textbf{Status} & Established working hierarchy & \\
    \hline
    \textbf{Ancestry} &  & \\
    \hline
    \textbf{Source} &  & \\
    \hline
    \textbf{Enabled By} &  & \\
    \hline
    \textbf{Supports} &  & \\
    \hline
    \textbf{Derived From} &  & \\
    \hline
    \textbf{Reviewed Status} & Review meaning, business boundary, information context and any unresolved ancestry. & \\
    \hline
  \end{tabularx}
\end{table}

\begin{table}[htbp]
  \centering
  \footnotesize
  \begin{tabularx}{1\textwidth}{|l|l|X|}
    \hline
    \multicolumn{3}{c}{\textbf{Contained Features}} \\
    \hline
    \textbf{Feature Name} & \textbf{Feature Id} & \textbf{Comment} \\
    \hline
    \multicolumn{3}{|l|}{\emph{No contained Features are currently identified.}} \\
    \hline
  \end{tabularx}
\end{table}

\textbf{Intent:} Operational Care Capacity Management identifies the organisational ability to undertake the stated health-service responsibility. The capability describes what the enterprise must be able to do; it does not assert that Harmonia performs or owns that activity.

\textbf{Boundary:} Harmonia relevance may arise because this activity establishes the context in which information is created, needed, interpreted, related, exchanged or presented. Responsibility for the underlying business decision or care activity remains with the appropriate actor or organisation.

\section{CT2 - Care Coordination \& Transition}\label{cap:bc-59cffb65}
\vspace{-0.5em}
\begin{table}[htbp]
  \centering
  \footnotesize
  \begin{tabularx}{1\textwidth}{|l|l|X|}
    \hline
    \multicolumn{3}{c}{\textbf{Capability Metadata}} \\
    \hline
    \textbf{Name} & \textbf{Value} & \textbf{Comment and/or Description} \\
    \hline
    \textbf{Reference Id} & \texttt{BC-59CFFB65} & \\
    \hline
    \textbf{Status} & Established working hierarchy & \\
    \hline
    \textbf{Ancestry} &  & \\
    \hline
    \textbf{Source} &  & \\
    \hline
    \textbf{Enabled By} &  & \\
    \hline
    \textbf{Supports} &  & \\
    \hline
    \textbf{Derived From} &  & \\
    \hline
    \textbf{Reviewed Status} &  & \\
    \hline
  \end{tabularx}
\end{table}

\begin{table}[htbp]
  \centering
  \footnotesize
  \begin{tabularx}{1\textwidth}{|l|l|X|}
    \hline
    \multicolumn{3}{c}{\textbf{Contained Capabilities}} \\
    \hline
    \textbf{Capability Name} & \textbf{Capability Id} & \textbf{Comment} \\
    \hline
    Care-Team Coordination & \texttt{cap:bc-9c9bd5f2} (at \ref{cap:bc-9c9bd5f2}) & \\
    \hline
    Care Responsibility Transition & \texttt{cap:bc-aa12612c} (at \ref{cap:bc-aa12612c}) & \\
    \hline
    Care Discharge Coordination & \texttt{cap:bc-d9b48397} (at \ref{cap:bc-d9b48397}) & \\
    \hline
  \end{tabularx}
\end{table}

\textbf{Intent:} Care Coordination \& Transition identifies the organisational ability to undertake the stated health-service responsibility. The capability describes what the enterprise must be able to do; it does not assert that Harmonia performs or owns that activity.

\textbf{Boundary:} Harmonia relevance may arise because this activity establishes the context in which information is created, needed, interpreted, related, exchanged or presented. Responsibility for the underlying business decision or care activity remains with the appropriate actor or organisation.

\subsection{CT3 - Care-Team Coordination}\label{cap:bc-9c9bd5f2}
\vspace{-0.5em}
\begin{table}[htbp]
  \centering
  \footnotesize
  \begin{tabularx}{1\textwidth}{|l|l|X|}
    \hline
    \multicolumn{3}{c}{\textbf{Capability Metadata}} \\
    \hline
    \textbf{Name} & \textbf{Value} & \textbf{Comment and/or Description} \\
    \hline
    \textbf{Reference Id} & \texttt{BC-9C9BD5F2} & \\
    \hline
    \textbf{Status} & Established working hierarchy & \\
    \hline
    \textbf{Ancestry} &  & \\
    \hline
    \textbf{Source} &  & \\
    \hline
    \textbf{Enabled By} &  & \\
    \hline
    \textbf{Supports} &  & \\
    \hline
    \textbf{Derived From} &  & \\
    \hline
    \textbf{Reviewed Status} & Review meaning, business boundary, information context and any unresolved ancestry. & \\
    \hline
  \end{tabularx}
\end{table}

\begin{table}[htbp]
  \centering
  \footnotesize
  \begin{tabularx}{1\textwidth}{|l|l|X|}
    \hline
    \multicolumn{3}{c}{\textbf{Contained Features}} \\
    \hline
    \textbf{Feature Name} & \textbf{Feature Id} & \textbf{Comment} \\
    \hline
    \multicolumn{3}{|l|}{\emph{No contained Features are currently identified.}} \\
    \hline
  \end{tabularx}
\end{table}

\textbf{Intent:} Care-Team Coordination identifies the organisational ability to undertake the stated health-service responsibility. The capability describes what the enterprise must be able to do; it does not assert that Harmonia performs or owns that activity.

\textbf{Boundary:} Harmonia relevance may arise because this activity establishes the context in which information is created, needed, interpreted, related, exchanged or presented. Responsibility for the underlying business decision or care activity remains with the appropriate actor or organisation.

\subsection{CT3 - Care Responsibility Transition}\label{cap:bc-aa12612c}
\vspace{-0.5em}
\begin{table}[htbp]
  \centering
  \footnotesize
  \begin{tabularx}{1\textwidth}{|l|l|X|}
    \hline
    \multicolumn{3}{c}{\textbf{Capability Metadata}} \\
    \hline
    \textbf{Name} & \textbf{Value} & \textbf{Comment and/or Description} \\
    \hline
    \textbf{Reference Id} & \texttt{BC-AA12612C} & \\
    \hline
    \textbf{Status} & Established working hierarchy & \\
    \hline
    \textbf{Ancestry} &  & \\
    \hline
    \textbf{Source} &  & \\
    \hline
    \textbf{Enabled By} &  & \\
    \hline
    \textbf{Supports} &  & \\
    \hline
    \textbf{Derived From} &  & \\
    \hline
    \textbf{Reviewed Status} & Review meaning, business boundary, information context and any unresolved ancestry. & \\
    \hline
  \end{tabularx}
\end{table}

\begin{table}[htbp]
  \centering
  \footnotesize
  \begin{tabularx}{1\textwidth}{|l|l|X|}
    \hline
    \multicolumn{3}{c}{\textbf{Contained Features}} \\
    \hline
    \textbf{Feature Name} & \textbf{Feature Id} & \textbf{Comment} \\
    \hline
    \multicolumn{3}{|l|}{\emph{No contained Features are currently identified.}} \\
    \hline
  \end{tabularx}
\end{table}

\textbf{Intent:} Care Responsibility Transition identifies the organisational ability to undertake the stated health-service responsibility. The capability describes what the enterprise must be able to do; it does not assert that Harmonia performs or owns that activity.

\textbf{Boundary:} Harmonia relevance may arise because this activity establishes the context in which information is created, needed, interpreted, related, exchanged or presented. Responsibility for the underlying business decision or care activity remains with the appropriate actor or organisation.

\subsection{CT3 - Care Discharge Coordination}\label{cap:bc-d9b48397}
\vspace{-0.5em}
\begin{table}[htbp]
  \centering
  \footnotesize
  \begin{tabularx}{1\textwidth}{|l|l|X|}
    \hline
    \multicolumn{3}{c}{\textbf{Capability Metadata}} \\
    \hline
    \textbf{Name} & \textbf{Value} & \textbf{Comment and/or Description} \\
    \hline
    \textbf{Reference Id} & \texttt{BC-D9B48397} & \\
    \hline
    \textbf{Status} & Established working hierarchy & \\
    \hline
    \textbf{Ancestry} &  & \\
    \hline
    \textbf{Source} &  & \\
    \hline
    \textbf{Enabled By} &  & \\
    \hline
    \textbf{Supports} &  & \\
    \hline
    \textbf{Derived From} &  & \\
    \hline
    \textbf{Reviewed Status} & Review meaning, business boundary, information context and any unresolved ancestry. & \\
    \hline
  \end{tabularx}
\end{table}

\begin{table}[htbp]
  \centering
  \footnotesize
  \begin{tabularx}{1\textwidth}{|l|l|X|}
    \hline
    \multicolumn{3}{c}{\textbf{Contained Features}} \\
    \hline
    \textbf{Feature Name} & \textbf{Feature Id} & \textbf{Comment} \\
    \hline
    \multicolumn{3}{|l|}{\emph{No contained Features are currently identified.}} \\
    \hline
  \end{tabularx}
\end{table}

\textbf{Intent:} Care Discharge Coordination identifies the organisational ability to undertake the stated health-service responsibility. The capability describes what the enterprise must be able to do; it does not assert that Harmonia performs or owns that activity.

\textbf{Boundary:} Harmonia relevance may arise because this activity establishes the context in which information is created, needed, interpreted, related, exchanged or presented. Responsibility for the underlying business decision or care activity remains with the appropriate actor or organisation.

\section{CT2 - Care Support Logistics Management}\label{cap:bc-c313e0f4}
\vspace{-0.5em}
\begin{table}[htbp]
  \centering
  \footnotesize
  \begin{tabularx}{1\textwidth}{|l|l|X|}
    \hline
    \multicolumn{3}{c}{\textbf{Capability Metadata}} \\
    \hline
    \textbf{Name} & \textbf{Value} & \textbf{Comment and/or Description} \\
    \hline
    \textbf{Reference Id} & \texttt{BC-C313E0F4} & \\
    \hline
    \textbf{Status} & Established working hierarchy & \\
    \hline
    \textbf{Ancestry} &  & \\
    \hline
    \textbf{Source} &  & \\
    \hline
    \textbf{Enabled By} &  & \\
    \hline
    \textbf{Supports} &  & \\
    \hline
    \textbf{Derived From} &  & \\
    \hline
    \textbf{Reviewed Status} &  & \\
    \hline
  \end{tabularx}
\end{table}

\begin{table}[htbp]
  \centering
  \footnotesize
  \begin{tabularx}{1\textwidth}{|l|l|X|}
    \hline
    \multicolumn{3}{c}{\textbf{Contained Capabilities}} \\
    \hline
    \textbf{Capability Name} & \textbf{Capability Id} & \textbf{Comment} \\
    \hline
    Patient Movement Coordination & \texttt{cap:bc-ce1271e6} (at \ref{cap:bc-ce1271e6}) & \\
    \hline
    Clinical Material Logistics Coordination & \texttt{cap:bc-8924ed80} (at \ref{cap:bc-8924ed80}) & \\
    \hline
    Care Support Work Coordination & \texttt{cap:bc-e0fd8d9d} (at \ref{cap:bc-e0fd8d9d}) & \\
    \hline
  \end{tabularx}
\end{table}

\textbf{Intent:} Care Support Logistics Management identifies the organisational ability to undertake the stated health-service responsibility. The capability describes what the enterprise must be able to do; it does not assert that Harmonia performs or owns that activity.

\textbf{Boundary:} Harmonia relevance may arise because this activity establishes the context in which information is created, needed, interpreted, related, exchanged or presented. Responsibility for the underlying business decision or care activity remains with the appropriate actor or organisation.

\subsection{CT3 - Patient Movement Coordination}\label{cap:bc-ce1271e6}
\vspace{-0.5em}
\begin{table}[htbp]
  \centering
  \footnotesize
  \begin{tabularx}{1\textwidth}{|l|l|X|}
    \hline
    \multicolumn{3}{c}{\textbf{Capability Metadata}} \\
    \hline
    \textbf{Name} & \textbf{Value} & \textbf{Comment and/or Description} \\
    \hline
    \textbf{Reference Id} & \texttt{BC-CE1271E6} & \\
    \hline
    \textbf{Status} & Established working hierarchy & \\
    \hline
    \textbf{Ancestry} &  & \\
    \hline
    \textbf{Source} &  & \\
    \hline
    \textbf{Enabled By} &  & \\
    \hline
    \textbf{Supports} &  & \\
    \hline
    \textbf{Derived From} &  & \\
    \hline
    \textbf{Reviewed Status} & Review meaning, business boundary, information context and any unresolved ancestry. & \\
    \hline
  \end{tabularx}
\end{table}

\begin{table}[htbp]
  \centering
  \footnotesize
  \begin{tabularx}{1\textwidth}{|l|l|X|}
    \hline
    \multicolumn{3}{c}{\textbf{Contained Features}} \\
    \hline
    \textbf{Feature Name} & \textbf{Feature Id} & \textbf{Comment} \\
    \hline
    \multicolumn{3}{|l|}{\emph{No contained Features are currently identified.}} \\
    \hline
  \end{tabularx}
\end{table}

\textbf{Intent:} Patient Movement Coordination identifies the organisational ability to undertake the stated health-service responsibility. The capability describes what the enterprise must be able to do; it does not assert that Harmonia performs or owns that activity.

\textbf{Boundary:} Harmonia relevance may arise because this activity establishes the context in which information is created, needed, interpreted, related, exchanged or presented. Responsibility for the underlying business decision or care activity remains with the appropriate actor or organisation.

\subsection{CT3 - Clinical Material Logistics Coordination}\label{cap:bc-8924ed80}
\vspace{-0.5em}
\begin{table}[htbp]
  \centering
  \footnotesize
  \begin{tabularx}{1\textwidth}{|l|l|X|}
    \hline
    \multicolumn{3}{c}{\textbf{Capability Metadata}} \\
    \hline
    \textbf{Name} & \textbf{Value} & \textbf{Comment and/or Description} \\
    \hline
    \textbf{Reference Id} & \texttt{BC-8924ED80} & \\
    \hline
    \textbf{Status} & Established working hierarchy & \\
    \hline
    \textbf{Ancestry} &  & \\
    \hline
    \textbf{Source} &  & \\
    \hline
    \textbf{Enabled By} &  & \\
    \hline
    \textbf{Supports} &  & \\
    \hline
    \textbf{Derived From} &  & \\
    \hline
    \textbf{Reviewed Status} & Review meaning, business boundary, information context and any unresolved ancestry. & \\
    \hline
  \end{tabularx}
\end{table}

\begin{table}[htbp]
  \centering
  \footnotesize
  \begin{tabularx}{1\textwidth}{|l|l|X|}
    \hline
    \multicolumn{3}{c}{\textbf{Contained Features}} \\
    \hline
    \textbf{Feature Name} & \textbf{Feature Id} & \textbf{Comment} \\
    \hline
    \multicolumn{3}{|l|}{\emph{No contained Features are currently identified.}} \\
    \hline
  \end{tabularx}
\end{table}

\textbf{Intent:} Clinical Material Logistics Coordination identifies the organisational ability to undertake the stated health-service responsibility. The capability describes what the enterprise must be able to do; it does not assert that Harmonia performs or owns that activity.

\textbf{Boundary:} Harmonia relevance may arise because this activity establishes the context in which information is created, needed, interpreted, related, exchanged or presented. Responsibility for the underlying business decision or care activity remains with the appropriate actor or organisation.

\subsection{CT3 - Care Support Work Coordination}\label{cap:bc-e0fd8d9d}
\vspace{-0.5em}
\begin{table}[htbp]
  \centering
  \footnotesize
  \begin{tabularx}{1\textwidth}{|l|l|X|}
    \hline
    \multicolumn{3}{c}{\textbf{Capability Metadata}} \\
    \hline
    \textbf{Name} & \textbf{Value} & \textbf{Comment and/or Description} \\
    \hline
    \textbf{Reference Id} & \texttt{BC-E0FD8D9D} & \\
    \hline
    \textbf{Status} & Established working hierarchy & \\
    \hline
    \textbf{Ancestry} &  & \\
    \hline
    \textbf{Source} &  & \\
    \hline
    \textbf{Enabled By} &  & \\
    \hline
    \textbf{Supports} &  & \\
    \hline
    \textbf{Derived From} &  & \\
    \hline
    \textbf{Reviewed Status} & Review meaning, business boundary, information context and any unresolved ancestry. & \\
    \hline
  \end{tabularx}
\end{table}

\begin{table}[htbp]
  \centering
  \footnotesize
  \begin{tabularx}{1\textwidth}{|l|l|X|}
    \hline
    \multicolumn{3}{c}{\textbf{Contained Features}} \\
    \hline
    \textbf{Feature Name} & \textbf{Feature Id} & \textbf{Comment} \\
    \hline
    \multicolumn{3}{|l|}{\emph{No contained Features are currently identified.}} \\
    \hline
  \end{tabularx}
\end{table}

\textbf{Intent:} Care Support Work Coordination identifies the organisational ability to undertake the stated health-service responsibility. The capability describes what the enterprise must be able to do; it does not assert that Harmonia performs or owns that activity.

\textbf{Boundary:} Harmonia relevance may arise because this activity establishes the context in which information is created, needed, interpreted, related, exchanged or presented. Responsibility for the underlying business decision or care activity remains with the appropriate actor or organisation.

\chapter{CT1 - Health Rights, Advocacy \& Participation}\label{cap:bc-3e1a6dca}
\noindent\textbf{Reference ID:} \texttt{BC-3E1A6DCA}\quad \textbf{Status:} Established working hierarchy\par
\noindent\textit{Source: Canonical Business Capability model / working map}\par
\textbf{Intent:} Health Rights, Advocacy \& Participation identifies the organisational ability to undertake the stated health-service responsibility. The capability describes what the enterprise must be able to do; it does not assert that Harmonia performs or owns that activity.

\textbf{Boundary:} Harmonia relevance may arise because this activity establishes the context in which information is created, needed, interpreted, related, exchanged or presented. Responsibility for the underlying business decision or care activity remains with the appropriate actor or organisation.

\reviewstatus{Check scope and information relevance; no business authority transfers to Harmonia.}
\chapter{CT1 - Diagnostic, Therapeutic \& Clinical Support Services}\label{cap:bc-c705c8ad}
\noindent\textbf{Reference ID:} \texttt{BC-C705C8AD}\quad \textbf{Status:} Established working hierarchy\par
\noindent\textit{Source: Canonical Business Capability model / working map}\par
\textbf{Intent:} Diagnostic, Therapeutic \& Clinical Support Services identifies the organisational ability to undertake the stated health-service responsibility. The capability describes what the enterprise must be able to do; it does not assert that Harmonia performs or owns that activity.

\textbf{Boundary:} Harmonia relevance may arise because this activity establishes the context in which information is created, needed, interpreted, related, exchanged or presented. Responsibility for the underlying business decision or care activity remains with the appropriate actor or organisation.

\reviewstatus{Check scope and information relevance; no business authority transfers to Harmonia.}
\section{CT2 - Medication \& Pharmacy Services}\label{cap:bc-e0be8a8a}
\vspace{-0.5em}
\begin{table}[htbp]
  \centering
  \footnotesize
  \begin{tabularx}{1\textwidth}{|l|l|X|}
    \hline
    \multicolumn{3}{c}{\textbf{Capability Metadata}} \\
    \hline
    \textbf{Name} & \textbf{Value} & \textbf{Comment and/or Description} \\
    \hline
    \textbf{Reference Id} & \texttt{BC-E0BE8A8A} & \\
    \hline
    \textbf{Status} & Established working hierarchy & \\
    \hline
    \textbf{Ancestry} &  & \\
    \hline
    \textbf{Source} &  & \\
    \hline
    \textbf{Enabled By} &  & \\
    \hline
    \textbf{Supports} &  & \\
    \hline
    \textbf{Derived From} &  & \\
    \hline
    \textbf{Reviewed Status} &  & \\
    \hline
  \end{tabularx}
\end{table}

\begin{table}[htbp]
  \centering
  \footnotesize
  \begin{tabularx}{1\textwidth}{|l|l|X|}
    \hline
    \multicolumn{3}{c}{\textbf{Contained Capabilities}} \\
    \hline
    \textbf{Capability Name} & \textbf{Capability Id} & \textbf{Comment} \\
    \hline
    Prescription Medication Dispensing & \texttt{cap:bc-4da49aee} (at \ref{cap:bc-4da49aee}) & \\
    \hline
    Non-Prescription Medication Provision & \texttt{cap:bc-c67d1e73} (at \ref{cap:bc-c67d1e73}) & \\
    \hline
    Medication Advice & \texttt{cap:bc-4e4274d4} (at \ref{cap:bc-4e4274d4}) & \\
    \hline
    Medication Dose Administration Aid Provision & \texttt{cap:bc-f8e65cd9} (at \ref{cap:bc-f8e65cd9}) & \\
    \hline
  \end{tabularx}
\end{table}

\textbf{Intent:} Medication \& Pharmacy Services identifies the organisational ability to undertake the stated health-service responsibility. The capability describes what the enterprise must be able to do; it does not assert that Harmonia performs or owns that activity.

\textbf{Boundary:} Harmonia relevance may arise because this activity establishes the context in which information is created, needed, interpreted, related, exchanged or presented. Responsibility for the underlying business decision or care activity remains with the appropriate actor or organisation.

\subsection{CT3 - Prescription Medication Dispensing}\label{cap:bc-4da49aee}
\vspace{-0.5em}
\begin{table}[htbp]
  \centering
  \footnotesize
  \begin{tabularx}{1\textwidth}{|l|l|X|}
    \hline
    \multicolumn{3}{c}{\textbf{Capability Metadata}} \\
    \hline
    \textbf{Name} & \textbf{Value} & \textbf{Comment and/or Description} \\
    \hline
    \textbf{Reference Id} & \texttt{BC-4DA49AEE} & \\
    \hline
    \textbf{Status} & Established working hierarchy & \\
    \hline
    \textbf{Ancestry} &  & \\
    \hline
    \textbf{Source} &  & \\
    \hline
    \textbf{Enabled By} &  & \\
    \hline
    \textbf{Supports} &  & \\
    \hline
    \textbf{Derived From} &  & \\
    \hline
    \textbf{Reviewed Status} & Review meaning, business boundary, information context and any unresolved ancestry. & \\
    \hline
  \end{tabularx}
\end{table}

\begin{table}[htbp]
  \centering
  \footnotesize
  \begin{tabularx}{1\textwidth}{|l|l|X|}
    \hline
    \multicolumn{3}{c}{\textbf{Contained Features}} \\
    \hline
    \textbf{Feature Name} & \textbf{Feature Id} & \textbf{Comment} \\
    \hline
    \multicolumn{3}{|l|}{\emph{No contained Features are currently identified.}} \\
    \hline
  \end{tabularx}
\end{table}

\textbf{Intent:} Prescription Medication Dispensing identifies the organisational ability to undertake the stated health-service responsibility. The capability describes what the enterprise must be able to do; it does not assert that Harmonia performs or owns that activity.

\textbf{Boundary:} Harmonia relevance may arise because this activity establishes the context in which information is created, needed, interpreted, related, exchanged or presented. Responsibility for the underlying business decision or care activity remains with the appropriate actor or organisation.

\subsection{CT3 - Non-Prescription Medication Provision}\label{cap:bc-c67d1e73}
\vspace{-0.5em}
\begin{table}[htbp]
  \centering
  \footnotesize
  \begin{tabularx}{1\textwidth}{|l|l|X|}
    \hline
    \multicolumn{3}{c}{\textbf{Capability Metadata}} \\
    \hline
    \textbf{Name} & \textbf{Value} & \textbf{Comment and/or Description} \\
    \hline
    \textbf{Reference Id} & \texttt{BC-C67D1E73} & \\
    \hline
    \textbf{Status} & Established working hierarchy & \\
    \hline
    \textbf{Ancestry} &  & \\
    \hline
    \textbf{Source} &  & \\
    \hline
    \textbf{Enabled By} &  & \\
    \hline
    \textbf{Supports} &  & \\
    \hline
    \textbf{Derived From} &  & \\
    \hline
    \textbf{Reviewed Status} & Review meaning, business boundary, information context and any unresolved ancestry. & \\
    \hline
  \end{tabularx}
\end{table}

\begin{table}[htbp]
  \centering
  \footnotesize
  \begin{tabularx}{1\textwidth}{|l|l|X|}
    \hline
    \multicolumn{3}{c}{\textbf{Contained Features}} \\
    \hline
    \textbf{Feature Name} & \textbf{Feature Id} & \textbf{Comment} \\
    \hline
    \multicolumn{3}{|l|}{\emph{No contained Features are currently identified.}} \\
    \hline
  \end{tabularx}
\end{table}

\textbf{Intent:} Non-Prescription Medication Provision identifies the organisational ability to undertake the stated health-service responsibility. The capability describes what the enterprise must be able to do; it does not assert that Harmonia performs or owns that activity.

\textbf{Boundary:} Harmonia relevance may arise because this activity establishes the context in which information is created, needed, interpreted, related, exchanged or presented. Responsibility for the underlying business decision or care activity remains with the appropriate actor or organisation.

\subsection{CT3 - Medication Advice}\label{cap:bc-4e4274d4}
\vspace{-0.5em}
\begin{table}[htbp]
  \centering
  \footnotesize
  \begin{tabularx}{1\textwidth}{|l|l|X|}
    \hline
    \multicolumn{3}{c}{\textbf{Capability Metadata}} \\
    \hline
    \textbf{Name} & \textbf{Value} & \textbf{Comment and/or Description} \\
    \hline
    \textbf{Reference Id} & \texttt{BC-4E4274D4} & \\
    \hline
    \textbf{Status} & Established working hierarchy & \\
    \hline
    \textbf{Ancestry} &  & \\
    \hline
    \textbf{Source} &  & \\
    \hline
    \textbf{Enabled By} &  & \\
    \hline
    \textbf{Supports} &  & \\
    \hline
    \textbf{Derived From} &  & \\
    \hline
    \textbf{Reviewed Status} & Review meaning, business boundary, information context and any unresolved ancestry. & \\
    \hline
  \end{tabularx}
\end{table}

\begin{table}[htbp]
  \centering
  \footnotesize
  \begin{tabularx}{1\textwidth}{|l|l|X|}
    \hline
    \multicolumn{3}{c}{\textbf{Contained Features}} \\
    \hline
    \textbf{Feature Name} & \textbf{Feature Id} & \textbf{Comment} \\
    \hline
    \multicolumn{3}{|l|}{\emph{No contained Features are currently identified.}} \\
    \hline
  \end{tabularx}
\end{table}

\textbf{Intent:} Medication Advice identifies the organisational ability to undertake the stated health-service responsibility. The capability describes what the enterprise must be able to do; it does not assert that Harmonia performs or owns that activity.

\textbf{Boundary:} Harmonia relevance may arise because this activity establishes the context in which information is created, needed, interpreted, related, exchanged or presented. Responsibility for the underlying business decision or care activity remains with the appropriate actor or organisation.

\subsection{CT3 - Medication Dose Administration Aid Provision}\label{cap:bc-f8e65cd9}
\vspace{-0.5em}
\begin{table}[htbp]
  \centering
  \footnotesize
  \begin{tabularx}{1\textwidth}{|l|l|X|}
    \hline
    \multicolumn{3}{c}{\textbf{Capability Metadata}} \\
    \hline
    \textbf{Name} & \textbf{Value} & \textbf{Comment and/or Description} \\
    \hline
    \textbf{Reference Id} & \texttt{BC-F8E65CD9} & \\
    \hline
    \textbf{Status} & Established working hierarchy & \\
    \hline
    \textbf{Ancestry} &  & \\
    \hline
    \textbf{Source} &  & \\
    \hline
    \textbf{Enabled By} &  & \\
    \hline
    \textbf{Supports} &  & \\
    \hline
    \textbf{Derived From} &  & \\
    \hline
    \textbf{Reviewed Status} & Review meaning, business boundary, information context and any unresolved ancestry. & \\
    \hline
  \end{tabularx}
\end{table}

\begin{table}[htbp]
  \centering
  \footnotesize
  \begin{tabularx}{1\textwidth}{|l|l|X|}
    \hline
    \multicolumn{3}{c}{\textbf{Contained Features}} \\
    \hline
    \textbf{Feature Name} & \textbf{Feature Id} & \textbf{Comment} \\
    \hline
    \multicolumn{3}{|l|}{\emph{No contained Features are currently identified.}} \\
    \hline
  \end{tabularx}
\end{table}

\textbf{Intent:} Medication Dose Administration Aid Provision identifies the organisational ability to undertake the stated health-service responsibility. The capability describes what the enterprise must be able to do; it does not assert that Harmonia performs or owns that activity.

\textbf{Boundary:} Harmonia relevance may arise because this activity establishes the context in which information is created, needed, interpreted, related, exchanged or presented. Responsibility for the underlying business decision or care activity remains with the appropriate actor or organisation.

Unresolved / stopping-point note: [OI] Medication Compounding and Preparation

\section{CT2 - Diagnostic Support Provision}\label{cap:bc-7b292540}
\vspace{-0.5em}
\begin{table}[htbp]
  \centering
  \footnotesize
  \begin{tabularx}{1\textwidth}{|l|l|X|}
    \hline
    \multicolumn{3}{c}{\textbf{Capability Metadata}} \\
    \hline
    \textbf{Name} & \textbf{Value} & \textbf{Comment and/or Description} \\
    \hline
    \textbf{Reference Id} & \texttt{BC-7B292540} & \\
    \hline
    \textbf{Status} & Established working hierarchy & \\
    \hline
    \textbf{Ancestry} &  & \\
    \hline
    \textbf{Source} &  & \\
    \hline
    \textbf{Enabled By} &  & \\
    \hline
    \textbf{Supports} &  & \\
    \hline
    \textbf{Derived From} &  & \\
    \hline
    \textbf{Reviewed Status} &  & \\
    \hline
  \end{tabularx}
\end{table}

\begin{table}[htbp]
  \centering
  \footnotesize
  \begin{tabularx}{1\textwidth}{|l|l|X|}
    \hline
    \multicolumn{3}{c}{\textbf{Contained Capabilities}} \\
    \hline
    \textbf{Capability Name} & \textbf{Capability Id} & \textbf{Comment} \\
    \hline
    Diagnostic Service Appointment Management & \texttt{cap:bc-d7778ef0} (at \ref{cap:bc-d7778ef0}) & \\
    \hline
    Diagnostic Service Delivery & \texttt{cap:bc-20131f3b} (at \ref{cap:bc-20131f3b}) & \\
    \hline
    Diagnostic Specimen Management & \texttt{cap:bc-d29ea624} (at \ref{cap:bc-d29ea624}) & \\
    \hline
    Diagnostic Data Management & \texttt{cap:bc-47be7761} (at \ref{cap:bc-47be7761}) & \\
    \hline
    Diagnostic Interpretation & \texttt{cap:bc-30ab08bb} (at \ref{cap:bc-30ab08bb}) & \\
    \hline
    Diagnostic Report Management & \texttt{cap:bc-e3c0477c} (at \ref{cap:bc-e3c0477c}) & \\
    \hline
    Diagnostic Report Distribution & \texttt{cap:bc-e001550a} (at \ref{cap:bc-e001550a}) & \\
    \hline
  \end{tabularx}
\end{table}

\textbf{Intent:} Diagnostic Support Provision identifies the organisational ability to undertake the stated health-service responsibility. The capability describes what the enterprise must be able to do; it does not assert that Harmonia performs or owns that activity.

\textbf{Boundary:} Harmonia relevance may arise because this activity establishes the context in which information is created, needed, interpreted, related, exchanged or presented. Responsibility for the underlying business decision or care activity remains with the appropriate actor or organisation.

\subsection{CT3 - Diagnostic Service Appointment Management}\label{cap:bc-d7778ef0}
\vspace{-0.5em}
\begin{table}[htbp]
  \centering
  \footnotesize
  \begin{tabularx}{1\textwidth}{|l|l|X|}
    \hline
    \multicolumn{3}{c}{\textbf{Capability Metadata}} \\
    \hline
    \textbf{Name} & \textbf{Value} & \textbf{Comment and/or Description} \\
    \hline
    \textbf{Reference Id} & \texttt{BC-D7778EF0} & \\
    \hline
    \textbf{Status} & Established working hierarchy & \\
    \hline
    \textbf{Ancestry} &  & \\
    \hline
    \textbf{Source} &  & \\
    \hline
    \textbf{Enabled By} &  & \\
    \hline
    \textbf{Supports} &  & \\
    \hline
    \textbf{Derived From} &  & \\
    \hline
    \textbf{Reviewed Status} & Review meaning, business boundary, information context and any unresolved ancestry. & \\
    \hline
  \end{tabularx}
\end{table}

\begin{table}[htbp]
  \centering
  \footnotesize
  \begin{tabularx}{1\textwidth}{|l|l|X|}
    \hline
    \multicolumn{3}{c}{\textbf{Contained Features}} \\
    \hline
    \textbf{Feature Name} & \textbf{Feature Id} & \textbf{Comment} \\
    \hline
    \multicolumn{3}{|l|}{\emph{No contained Features are currently identified.}} \\
    \hline
  \end{tabularx}
\end{table}

\textbf{Intent:} Diagnostic Service Appointment Management identifies the organisational ability to undertake the stated health-service responsibility. The capability describes what the enterprise must be able to do; it does not assert that Harmonia performs or owns that activity.

\textbf{Boundary:} Harmonia relevance may arise because this activity establishes the context in which information is created, needed, interpreted, related, exchanged or presented. Responsibility for the underlying business decision or care activity remains with the appropriate actor or organisation.

\subsection{CT3 - Diagnostic Service Delivery}\label{cap:bc-20131f3b}
\vspace{-0.5em}
\begin{table}[htbp]
  \centering
  \footnotesize
  \begin{tabularx}{1\textwidth}{|l|l|X|}
    \hline
    \multicolumn{3}{c}{\textbf{Capability Metadata}} \\
    \hline
    \textbf{Name} & \textbf{Value} & \textbf{Comment and/or Description} \\
    \hline
    \textbf{Reference Id} & \texttt{BC-20131F3B} & \\
    \hline
    \textbf{Status} & Established working hierarchy & \\
    \hline
    \textbf{Ancestry} &  & \\
    \hline
    \textbf{Source} &  & \\
    \hline
    \textbf{Enabled By} &  & \\
    \hline
    \textbf{Supports} &  & \\
    \hline
    \textbf{Derived From} &  & \\
    \hline
    \textbf{Reviewed Status} & Review meaning, business boundary, information context and any unresolved ancestry. & \\
    \hline
  \end{tabularx}
\end{table}

\begin{table}[htbp]
  \centering
  \footnotesize
  \begin{tabularx}{1\textwidth}{|l|l|X|}
    \hline
    \multicolumn{3}{c}{\textbf{Contained Features}} \\
    \hline
    \textbf{Feature Name} & \textbf{Feature Id} & \textbf{Comment} \\
    \hline
    \multicolumn{3}{|l|}{\emph{No contained Features are currently identified.}} \\
    \hline
  \end{tabularx}
\end{table}

\textbf{Intent:} Diagnostic Service Delivery identifies the organisational ability to undertake the stated health-service responsibility. The capability describes what the enterprise must be able to do; it does not assert that Harmonia performs or owns that activity.

\textbf{Boundary:} Harmonia relevance may arise because this activity establishes the context in which information is created, needed, interpreted, related, exchanged or presented. Responsibility for the underlying business decision or care activity remains with the appropriate actor or organisation.

\subsection{CT3 - Diagnostic Specimen Management}\label{cap:bc-d29ea624}
\vspace{-0.5em}
\begin{table}[htbp]
  \centering
  \footnotesize
  \begin{tabularx}{1\textwidth}{|l|l|X|}
    \hline
    \multicolumn{3}{c}{\textbf{Capability Metadata}} \\
    \hline
    \textbf{Name} & \textbf{Value} & \textbf{Comment and/or Description} \\
    \hline
    \textbf{Reference Id} & \texttt{BC-D29EA624} & \\
    \hline
    \textbf{Status} & Established working hierarchy & \\
    \hline
    \textbf{Ancestry} &  & \\
    \hline
    \textbf{Source} &  & \\
    \hline
    \textbf{Enabled By} &  & \\
    \hline
    \textbf{Supports} &  & \\
    \hline
    \textbf{Derived From} &  & \\
    \hline
    \textbf{Reviewed Status} & Review meaning, business boundary, information context and any unresolved ancestry. & \\
    \hline
  \end{tabularx}
\end{table}

\begin{table}[htbp]
  \centering
  \footnotesize
  \begin{tabularx}{1\textwidth}{|l|l|X|}
    \hline
    \multicolumn{3}{c}{\textbf{Contained Features}} \\
    \hline
    \textbf{Feature Name} & \textbf{Feature Id} & \textbf{Comment} \\
    \hline
    \multicolumn{3}{|l|}{\emph{No contained Features are currently identified.}} \\
    \hline
  \end{tabularx}
\end{table}

\textbf{Intent:} Diagnostic Specimen Management identifies the organisational ability to undertake the stated health-service responsibility. The capability describes what the enterprise must be able to do; it does not assert that Harmonia performs or owns that activity.

\textbf{Boundary:} Harmonia relevance may arise because this activity establishes the context in which information is created, needed, interpreted, related, exchanged or presented. Responsibility for the underlying business decision or care activity remains with the appropriate actor or organisation.

\subsection{CT3 - Diagnostic Data Management}\label{cap:bc-47be7761}
\vspace{-0.5em}
\begin{table}[htbp]
  \centering
  \footnotesize
  \begin{tabularx}{1\textwidth}{|l|l|X|}
    \hline
    \multicolumn{3}{c}{\textbf{Capability Metadata}} \\
    \hline
    \textbf{Name} & \textbf{Value} & \textbf{Comment and/or Description} \\
    \hline
    \textbf{Reference Id} & \texttt{BC-47BE7761} & \\
    \hline
    \textbf{Status} & Established working hierarchy & \\
    \hline
    \textbf{Ancestry} &  & \\
    \hline
    \textbf{Source} &  & \\
    \hline
    \textbf{Enabled By} &  & \\
    \hline
    \textbf{Supports} &  & \\
    \hline
    \textbf{Derived From} &  & \\
    \hline
    \textbf{Reviewed Status} & Review meaning, business boundary, information context and any unresolved ancestry. & \\
    \hline
  \end{tabularx}
\end{table}

\begin{table}[htbp]
  \centering
  \footnotesize
  \begin{tabularx}{1\textwidth}{|l|l|X|}
    \hline
    \multicolumn{3}{c}{\textbf{Contained Features}} \\
    \hline
    \textbf{Feature Name} & \textbf{Feature Id} & \textbf{Comment} \\
    \hline
    \multicolumn{3}{|l|}{\emph{No contained Features are currently identified.}} \\
    \hline
  \end{tabularx}
\end{table}

\textbf{Intent:} Diagnostic Data Management identifies the organisational ability to undertake the stated health-service responsibility. The capability describes what the enterprise must be able to do; it does not assert that Harmonia performs or owns that activity.

\textbf{Boundary:} Harmonia relevance may arise because this activity establishes the context in which information is created, needed, interpreted, related, exchanged or presented. Responsibility for the underlying business decision or care activity remains with the appropriate actor or organisation.

\subsection{CT3 - Diagnostic Interpretation}\label{cap:bc-30ab08bb}
\vspace{-0.5em}
\begin{table}[htbp]
  \centering
  \footnotesize
  \begin{tabularx}{1\textwidth}{|l|l|X|}
    \hline
    \multicolumn{3}{c}{\textbf{Capability Metadata}} \\
    \hline
    \textbf{Name} & \textbf{Value} & \textbf{Comment and/or Description} \\
    \hline
    \textbf{Reference Id} & \texttt{BC-30AB08BB} & \\
    \hline
    \textbf{Status} & Established working hierarchy & \\
    \hline
    \textbf{Ancestry} &  & \\
    \hline
    \textbf{Source} &  & \\
    \hline
    \textbf{Enabled By} &  & \\
    \hline
    \textbf{Supports} &  & \\
    \hline
    \textbf{Derived From} &  & \\
    \hline
    \textbf{Reviewed Status} & Review meaning, business boundary, information context and any unresolved ancestry. & \\
    \hline
  \end{tabularx}
\end{table}

\begin{table}[htbp]
  \centering
  \footnotesize
  \begin{tabularx}{1\textwidth}{|l|l|X|}
    \hline
    \multicolumn{3}{c}{\textbf{Contained Features}} \\
    \hline
    \textbf{Feature Name} & \textbf{Feature Id} & \textbf{Comment} \\
    \hline
    \multicolumn{3}{|l|}{\emph{No contained Features are currently identified.}} \\
    \hline
  \end{tabularx}
\end{table}

\textbf{Intent:} Diagnostic Interpretation identifies the organisational ability to undertake the stated health-service responsibility. The capability describes what the enterprise must be able to do; it does not assert that Harmonia performs or owns that activity.

\textbf{Boundary:} Harmonia relevance may arise because this activity establishes the context in which information is created, needed, interpreted, related, exchanged or presented. Responsibility for the underlying business decision or care activity remains with the appropriate actor or organisation.

\subsection{CT3 - Diagnostic Report Management}\label{cap:bc-e3c0477c}
\vspace{-0.5em}
\begin{table}[htbp]
  \centering
  \footnotesize
  \begin{tabularx}{1\textwidth}{|l|l|X|}
    \hline
    \multicolumn{3}{c}{\textbf{Capability Metadata}} \\
    \hline
    \textbf{Name} & \textbf{Value} & \textbf{Comment and/or Description} \\
    \hline
    \textbf{Reference Id} & \texttt{BC-E3C0477C} & \\
    \hline
    \textbf{Status} & Established working hierarchy & \\
    \hline
    \textbf{Ancestry} &  & \\
    \hline
    \textbf{Source} &  & \\
    \hline
    \textbf{Enabled By} &  & \\
    \hline
    \textbf{Supports} &  & \\
    \hline
    \textbf{Derived From} &  & \\
    \hline
    \textbf{Reviewed Status} & Review meaning, business boundary, information context and any unresolved ancestry. & \\
    \hline
  \end{tabularx}
\end{table}

\begin{table}[htbp]
  \centering
  \footnotesize
  \begin{tabularx}{1\textwidth}{|l|l|X|}
    \hline
    \multicolumn{3}{c}{\textbf{Contained Features}} \\
    \hline
    \textbf{Feature Name} & \textbf{Feature Id} & \textbf{Comment} \\
    \hline
    \multicolumn{3}{|l|}{\emph{No contained Features are currently identified.}} \\
    \hline
  \end{tabularx}
\end{table}

\textbf{Intent:} Diagnostic Report Management identifies the organisational ability to undertake the stated health-service responsibility. The capability describes what the enterprise must be able to do; it does not assert that Harmonia performs or owns that activity.

\textbf{Boundary:} Harmonia relevance may arise because this activity establishes the context in which information is created, needed, interpreted, related, exchanged or presented. Responsibility for the underlying business decision or care activity remains with the appropriate actor or organisation.

\subsection{CT3 - Diagnostic Report Distribution}\label{cap:bc-e001550a}
\vspace{-0.5em}
\begin{table}[htbp]
  \centering
  \footnotesize
  \begin{tabularx}{1\textwidth}{|l|l|X|}
    \hline
    \multicolumn{3}{c}{\textbf{Capability Metadata}} \\
    \hline
    \textbf{Name} & \textbf{Value} & \textbf{Comment and/or Description} \\
    \hline
    \textbf{Reference Id} & \texttt{BC-E001550A} & \\
    \hline
    \textbf{Status} & Established working hierarchy & \\
    \hline
    \textbf{Ancestry} &  & \\
    \hline
    \textbf{Source} &  & \\
    \hline
    \textbf{Enabled By} &  & \\
    \hline
    \textbf{Supports} &  & \\
    \hline
    \textbf{Derived From} &  & \\
    \hline
    \textbf{Reviewed Status} & Review meaning, business boundary, information context and any unresolved ancestry. & \\
    \hline
  \end{tabularx}
\end{table}

\begin{table}[htbp]
  \centering
  \footnotesize
  \begin{tabularx}{1\textwidth}{|l|l|X|}
    \hline
    \multicolumn{3}{c}{\textbf{Contained Features}} \\
    \hline
    \textbf{Feature Name} & \textbf{Feature Id} & \textbf{Comment} \\
    \hline
    \multicolumn{3}{|l|}{\emph{No contained Features are currently identified.}} \\
    \hline
  \end{tabularx}
\end{table}

\textbf{Intent:} Diagnostic Report Distribution identifies the organisational ability to undertake the stated health-service responsibility. The capability describes what the enterprise must be able to do; it does not assert that Harmonia performs or owns that activity.

\textbf{Boundary:} Harmonia relevance may arise because this activity establishes the context in which information is created, needed, interpreted, related, exchanged or presented. Responsibility for the underlying business decision or care activity remains with the appropriate actor or organisation.

\chapter{CT1 - Health Products \& Clinical Technology}\label{cap:bc-c82d3cc6}
\noindent\textbf{Reference ID:} \texttt{BC-C82D3CC6}\quad \textbf{Status:} Established working hierarchy\par
\noindent\textit{Source: Canonical Business Capability model / working map}\par
\textbf{Intent:} Health Products \& Clinical Technology identifies the organisational ability to undertake the stated health-service responsibility. The capability describes what the enterprise must be able to do; it does not assert that Harmonia performs or owns that activity.

\textbf{Boundary:} Harmonia relevance may arise because this activity establishes the context in which information is created, needed, interpreted, related, exchanged or presented. Responsibility for the underlying business decision or care activity remains with the appropriate actor or organisation.

\reviewstatus{Check scope and information relevance; no business authority transfers to Harmonia.}
\section{CT2 - Health Product Management}\label{cap:bc-d7ffdedc}
\vspace{-0.5em}
\begin{table}[htbp]
  \centering
  \footnotesize
  \begin{tabularx}{1\textwidth}{|l|l|X|}
    \hline
    \multicolumn{3}{c}{\textbf{Capability Metadata}} \\
    \hline
    \textbf{Name} & \textbf{Value} & \textbf{Comment and/or Description} \\
    \hline
    \textbf{Reference Id} & \texttt{BC-D7FFDEDC} & \\
    \hline
    \textbf{Status} & Established working hierarchy & \\
    \hline
    \textbf{Ancestry} &  & \\
    \hline
    \textbf{Source} &  & \\
    \hline
    \textbf{Enabled By} &  & \\
    \hline
    \textbf{Supports} &  & \\
    \hline
    \textbf{Derived From} &  & \\
    \hline
    \textbf{Reviewed Status} &  & \\
    \hline
  \end{tabularx}
\end{table}

\begin{table}[htbp]
  \centering
  \footnotesize
  \begin{tabularx}{1\textwidth}{|l|l|X|}
    \hline
    \multicolumn{3}{c}{\textbf{Contained Capabilities}} \\
    \hline
    \textbf{Capability Name} & \textbf{Capability Id} & \textbf{Comment} \\
    \hline
    \multicolumn{3}{|l|}{\emph{No contained capabilities are currently identified.}} \\
    \hline
  \end{tabularx}
\end{table}

\textbf{Intent:} Health Product Management identifies the organisational ability to undertake the stated health-service responsibility. The capability describes what the enterprise must be able to do; it does not assert that Harmonia performs or owns that activity.

\textbf{Boundary:} Harmonia relevance may arise because this activity establishes the context in which information is created, needed, interpreted, related, exchanged or presented. Responsibility for the underlying business decision or care activity remains with the appropriate actor or organisation.

\section{CT2 - Clinical Technology Management}\label{cap:bc-377feb37}
\vspace{-0.5em}
\begin{table}[htbp]
  \centering
  \footnotesize
  \begin{tabularx}{1\textwidth}{|l|l|X|}
    \hline
    \multicolumn{3}{c}{\textbf{Capability Metadata}} \\
    \hline
    \textbf{Name} & \textbf{Value} & \textbf{Comment and/or Description} \\
    \hline
    \textbf{Reference Id} & \texttt{BC-377FEB37} & \\
    \hline
    \textbf{Status} & Established working hierarchy & \\
    \hline
    \textbf{Ancestry} &  & \\
    \hline
    \textbf{Source} &  & \\
    \hline
    \textbf{Enabled By} &  & \\
    \hline
    \textbf{Supports} &  & \\
    \hline
    \textbf{Derived From} &  & \\
    \hline
    \textbf{Reviewed Status} &  & \\
    \hline
  \end{tabularx}
\end{table}

\begin{table}[htbp]
  \centering
  \footnotesize
  \begin{tabularx}{1\textwidth}{|l|l|X|}
    \hline
    \multicolumn{3}{c}{\textbf{Contained Capabilities}} \\
    \hline
    \textbf{Capability Name} & \textbf{Capability Id} & \textbf{Comment} \\
    \hline
    Clinical Device Inventory Management & \texttt{cap:bc-097846f2} (at \ref{cap:bc-097846f2}) & \\
    \hline
    Clinical Device Feature \& Release Management & \texttt{cap:bc-e511b2db} (at \ref{cap:bc-e511b2db}) & \\
    \hline
    Clinical Device Loan \& Deployment Management & \texttt{cap:bc-e2e1c257} (at \ref{cap:bc-e2e1c257}) & \\
    \hline
    Clinical Device Utilisation Support & \texttt{cap:bc-bbe81736} (at \ref{cap:bc-bbe81736}) & \\
    \hline
    Clinical Device Integration Provision & \texttt{cap:bc-33ca15f2} (at \ref{cap:bc-33ca15f2}) & \\
    \hline
    Clinical Device Data Integration & \texttt{cap:bc-4f2a8587} (at \ref{cap:bc-4f2a8587}) & \\
    \hline
    Clinical Device Maintenance Management & \texttt{cap:bc-6cb8ac63} (at \ref{cap:bc-6cb8ac63}) & \\
    \hline
    Clinical Device Lifecycle Management & \texttt{cap:bc-1bff3f5b} (at \ref{cap:bc-1bff3f5b}) & \\
    \hline
  \end{tabularx}
\end{table}

\textbf{Intent:} Clinical Technology Management identifies the organisational ability to undertake the stated health-service responsibility. The capability describes what the enterprise must be able to do; it does not assert that Harmonia performs or owns that activity.

\textbf{Boundary:} Harmonia relevance may arise because this activity establishes the context in which information is created, needed, interpreted, related, exchanged or presented. Responsibility for the underlying business decision or care activity remains with the appropriate actor or organisation.

\subsection{CT3 - Clinical Device Inventory Management}\label{cap:bc-097846f2}
\vspace{-0.5em}
\begin{table}[htbp]
  \centering
  \footnotesize
  \begin{tabularx}{1\textwidth}{|l|l|X|}
    \hline
    \multicolumn{3}{c}{\textbf{Capability Metadata}} \\
    \hline
    \textbf{Name} & \textbf{Value} & \textbf{Comment and/or Description} \\
    \hline
    \textbf{Reference Id} & \texttt{BC-097846F2} & \\
    \hline
    \textbf{Status} & Established working hierarchy & \\
    \hline
    \textbf{Ancestry} &  & \\
    \hline
    \textbf{Source} &  & \\
    \hline
    \textbf{Enabled By} &  & \\
    \hline
    \textbf{Supports} &  & \\
    \hline
    \textbf{Derived From} &  & \\
    \hline
    \textbf{Reviewed Status} & Review meaning, business boundary, information context and any unresolved ancestry. & \\
    \hline
  \end{tabularx}
\end{table}

\begin{table}[htbp]
  \centering
  \footnotesize
  \begin{tabularx}{1\textwidth}{|l|l|X|}
    \hline
    \multicolumn{3}{c}{\textbf{Contained Features}} \\
    \hline
    \textbf{Feature Name} & \textbf{Feature Id} & \textbf{Comment} \\
    \hline
    \multicolumn{3}{|l|}{\emph{No contained Features are currently identified.}} \\
    \hline
  \end{tabularx}
\end{table}

\textbf{Intent:} Clinical Device Inventory Management identifies the organisational ability to undertake the stated health-service responsibility. The capability describes what the enterprise must be able to do; it does not assert that Harmonia performs or owns that activity.

\textbf{Boundary:} Harmonia relevance may arise because this activity establishes the context in which information is created, needed, interpreted, related, exchanged or presented. Responsibility for the underlying business decision or care activity remains with the appropriate actor or organisation.

\subsection{CT3 - Clinical Device Feature \& Release Management}\label{cap:bc-e511b2db}
\vspace{-0.5em}
\begin{table}[htbp]
  \centering
  \footnotesize
  \begin{tabularx}{1\textwidth}{|l|l|X|}
    \hline
    \multicolumn{3}{c}{\textbf{Capability Metadata}} \\
    \hline
    \textbf{Name} & \textbf{Value} & \textbf{Comment and/or Description} \\
    \hline
    \textbf{Reference Id} & \texttt{BC-E511B2DB} & \\
    \hline
    \textbf{Status} & Established working hierarchy & \\
    \hline
    \textbf{Ancestry} &  & \\
    \hline
    \textbf{Source} &  & \\
    \hline
    \textbf{Enabled By} &  & \\
    \hline
    \textbf{Supports} &  & \\
    \hline
    \textbf{Derived From} &  & \\
    \hline
    \textbf{Reviewed Status} & Review meaning, business boundary, information context and any unresolved ancestry. & \\
    \hline
  \end{tabularx}
\end{table}

\begin{table}[htbp]
  \centering
  \footnotesize
  \begin{tabularx}{1\textwidth}{|l|l|X|}
    \hline
    \multicolumn{3}{c}{\textbf{Contained Features}} \\
    \hline
    \textbf{Feature Name} & \textbf{Feature Id} & \textbf{Comment} \\
    \hline
    \multicolumn{3}{|l|}{\emph{No contained Features are currently identified.}} \\
    \hline
  \end{tabularx}
\end{table}

\textbf{Intent:} Clinical Device Feature \& Release Management identifies the organisational ability to undertake the stated health-service responsibility. The capability describes what the enterprise must be able to do; it does not assert that Harmonia performs or owns that activity.

\textbf{Boundary:} Harmonia relevance may arise because this activity establishes the context in which information is created, needed, interpreted, related, exchanged or presented. Responsibility for the underlying business decision or care activity remains with the appropriate actor or organisation.

\subsection{CT3 - Clinical Device Loan \& Deployment Management}\label{cap:bc-e2e1c257}
\vspace{-0.5em}
\begin{table}[htbp]
  \centering
  \footnotesize
  \begin{tabularx}{1\textwidth}{|l|l|X|}
    \hline
    \multicolumn{3}{c}{\textbf{Capability Metadata}} \\
    \hline
    \textbf{Name} & \textbf{Value} & \textbf{Comment and/or Description} \\
    \hline
    \textbf{Reference Id} & \texttt{BC-E2E1C257} & \\
    \hline
    \textbf{Status} & Established working hierarchy & \\
    \hline
    \textbf{Ancestry} &  & \\
    \hline
    \textbf{Source} &  & \\
    \hline
    \textbf{Enabled By} &  & \\
    \hline
    \textbf{Supports} &  & \\
    \hline
    \textbf{Derived From} &  & \\
    \hline
    \textbf{Reviewed Status} & Review meaning, business boundary, information context and any unresolved ancestry. & \\
    \hline
  \end{tabularx}
\end{table}

\begin{table}[htbp]
  \centering
  \footnotesize
  \begin{tabularx}{1\textwidth}{|l|l|X|}
    \hline
    \multicolumn{3}{c}{\textbf{Contained Features}} \\
    \hline
    \textbf{Feature Name} & \textbf{Feature Id} & \textbf{Comment} \\
    \hline
    \multicolumn{3}{|l|}{\emph{No contained Features are currently identified.}} \\
    \hline
  \end{tabularx}
\end{table}

\textbf{Intent:} Clinical Device Loan \& Deployment Management identifies the organisational ability to undertake the stated health-service responsibility. The capability describes what the enterprise must be able to do; it does not assert that Harmonia performs or owns that activity.

\textbf{Boundary:} Harmonia relevance may arise because this activity establishes the context in which information is created, needed, interpreted, related, exchanged or presented. Responsibility for the underlying business decision or care activity remains with the appropriate actor or organisation.

\subsection{CT3 - Clinical Device Utilisation Support}\label{cap:bc-bbe81736}
\vspace{-0.5em}
\begin{table}[htbp]
  \centering
  \footnotesize
  \begin{tabularx}{1\textwidth}{|l|l|X|}
    \hline
    \multicolumn{3}{c}{\textbf{Capability Metadata}} \\
    \hline
    \textbf{Name} & \textbf{Value} & \textbf{Comment and/or Description} \\
    \hline
    \textbf{Reference Id} & \texttt{BC-BBE81736} & \\
    \hline
    \textbf{Status} & Established working hierarchy & \\
    \hline
    \textbf{Ancestry} &  & \\
    \hline
    \textbf{Source} &  & \\
    \hline
    \textbf{Enabled By} &  & \\
    \hline
    \textbf{Supports} &  & \\
    \hline
    \textbf{Derived From} &  & \\
    \hline
    \textbf{Reviewed Status} & Review meaning, business boundary, information context and any unresolved ancestry. & \\
    \hline
  \end{tabularx}
\end{table}

\begin{table}[htbp]
  \centering
  \footnotesize
  \begin{tabularx}{1\textwidth}{|l|l|X|}
    \hline
    \multicolumn{3}{c}{\textbf{Contained Features}} \\
    \hline
    \textbf{Feature Name} & \textbf{Feature Id} & \textbf{Comment} \\
    \hline
    \multicolumn{3}{|l|}{\emph{No contained Features are currently identified.}} \\
    \hline
  \end{tabularx}
\end{table}

\textbf{Intent:} Clinical Device Utilisation Support identifies the organisational ability to undertake the stated health-service responsibility. The capability describes what the enterprise must be able to do; it does not assert that Harmonia performs or owns that activity.

\textbf{Boundary:} Harmonia relevance may arise because this activity establishes the context in which information is created, needed, interpreted, related, exchanged or presented. Responsibility for the underlying business decision or care activity remains with the appropriate actor or organisation.

\subsection{CT3 - Clinical Device Integration Provision}\label{cap:bc-33ca15f2}
\vspace{-0.5em}
\begin{table}[htbp]
  \centering
  \footnotesize
  \begin{tabularx}{1\textwidth}{|l|l|X|}
    \hline
    \multicolumn{3}{c}{\textbf{Capability Metadata}} \\
    \hline
    \textbf{Name} & \textbf{Value} & \textbf{Comment and/or Description} \\
    \hline
    \textbf{Reference Id} & \texttt{BC-33CA15F2} & \\
    \hline
    \textbf{Status} & Established working hierarchy & \\
    \hline
    \textbf{Ancestry} &  & \\
    \hline
    \textbf{Source} &  & \\
    \hline
    \textbf{Enabled By} &  & \\
    \hline
    \textbf{Supports} &  & \\
    \hline
    \textbf{Derived From} &  & \\
    \hline
    \textbf{Reviewed Status} & Review meaning, business boundary, information context and any unresolved ancestry. & \\
    \hline
  \end{tabularx}
\end{table}

\begin{table}[htbp]
  \centering
  \footnotesize
  \begin{tabularx}{1\textwidth}{|l|l|X|}
    \hline
    \multicolumn{3}{c}{\textbf{Contained Features}} \\
    \hline
    \textbf{Feature Name} & \textbf{Feature Id} & \textbf{Comment} \\
    \hline
    \multicolumn{3}{|l|}{\emph{No contained Features are currently identified.}} \\
    \hline
  \end{tabularx}
\end{table}

\textbf{Intent:} Clinical Device Integration Provision identifies the organisational ability to undertake the stated health-service responsibility. The capability describes what the enterprise must be able to do; it does not assert that Harmonia performs or owns that activity.

\textbf{Boundary:} Harmonia relevance may arise because this activity establishes the context in which information is created, needed, interpreted, related, exchanged or presented. Responsibility for the underlying business decision or care activity remains with the appropriate actor or organisation.

\subsection{CT3 - Clinical Device Data Integration}\label{cap:bc-4f2a8587}
\vspace{-0.5em}
\begin{table}[htbp]
  \centering
  \footnotesize
  \begin{tabularx}{1\textwidth}{|l|l|X|}
    \hline
    \multicolumn{3}{c}{\textbf{Capability Metadata}} \\
    \hline
    \textbf{Name} & \textbf{Value} & \textbf{Comment and/or Description} \\
    \hline
    \textbf{Reference Id} & \texttt{BC-4F2A8587} & \\
    \hline
    \textbf{Status} & Established working hierarchy & \\
    \hline
    \textbf{Ancestry} &  & \\
    \hline
    \textbf{Source} &  & \\
    \hline
    \textbf{Enabled By} &  & \\
    \hline
    \textbf{Supports} &  & \\
    \hline
    \textbf{Derived From} &  & \\
    \hline
    \textbf{Reviewed Status} & Review meaning, business boundary, information context and any unresolved ancestry. & \\
    \hline
  \end{tabularx}
\end{table}

\begin{table}[htbp]
  \centering
  \footnotesize
  \begin{tabularx}{1\textwidth}{|l|l|X|}
    \hline
    \multicolumn{3}{c}{\textbf{Contained Features}} \\
    \hline
    \textbf{Feature Name} & \textbf{Feature Id} & \textbf{Comment} \\
    \hline
    \multicolumn{3}{|l|}{\emph{No contained Features are currently identified.}} \\
    \hline
  \end{tabularx}
\end{table}

\textbf{Intent:} Clinical Device Data Integration identifies the organisational ability to undertake the stated health-service responsibility. The capability describes what the enterprise must be able to do; it does not assert that Harmonia performs or owns that activity.

\textbf{Boundary:} Harmonia relevance may arise because this activity establishes the context in which information is created, needed, interpreted, related, exchanged or presented. Responsibility for the underlying business decision or care activity remains with the appropriate actor or organisation.

\subsection{CT3 - Clinical Device Maintenance Management}\label{cap:bc-6cb8ac63}
\vspace{-0.5em}
\begin{table}[htbp]
  \centering
  \footnotesize
  \begin{tabularx}{1\textwidth}{|l|l|X|}
    \hline
    \multicolumn{3}{c}{\textbf{Capability Metadata}} \\
    \hline
    \textbf{Name} & \textbf{Value} & \textbf{Comment and/or Description} \\
    \hline
    \textbf{Reference Id} & \texttt{BC-6CB8AC63} & \\
    \hline
    \textbf{Status} & Established working hierarchy & \\
    \hline
    \textbf{Ancestry} &  & \\
    \hline
    \textbf{Source} &  & \\
    \hline
    \textbf{Enabled By} &  & \\
    \hline
    \textbf{Supports} &  & \\
    \hline
    \textbf{Derived From} &  & \\
    \hline
    \textbf{Reviewed Status} & Review meaning, business boundary, information context and any unresolved ancestry. & \\
    \hline
  \end{tabularx}
\end{table}

\begin{table}[htbp]
  \centering
  \footnotesize
  \begin{tabularx}{1\textwidth}{|l|l|X|}
    \hline
    \multicolumn{3}{c}{\textbf{Contained Features}} \\
    \hline
    \textbf{Feature Name} & \textbf{Feature Id} & \textbf{Comment} \\
    \hline
    \multicolumn{3}{|l|}{\emph{No contained Features are currently identified.}} \\
    \hline
  \end{tabularx}
\end{table}

\textbf{Intent:} Clinical Device Maintenance Management identifies the organisational ability to undertake the stated health-service responsibility. The capability describes what the enterprise must be able to do; it does not assert that Harmonia performs or owns that activity.

\textbf{Boundary:} Harmonia relevance may arise because this activity establishes the context in which information is created, needed, interpreted, related, exchanged or presented. Responsibility for the underlying business decision or care activity remains with the appropriate actor or organisation.

\subsection{CT3 - Clinical Device Lifecycle Management}\label{cap:bc-1bff3f5b}
\vspace{-0.5em}
\begin{table}[htbp]
  \centering
  \footnotesize
  \begin{tabularx}{1\textwidth}{|l|l|X|}
    \hline
    \multicolumn{3}{c}{\textbf{Capability Metadata}} \\
    \hline
    \textbf{Name} & \textbf{Value} & \textbf{Comment and/or Description} \\
    \hline
    \textbf{Reference Id} & \texttt{BC-1BFF3F5B} & \\
    \hline
    \textbf{Status} & Established working hierarchy & \\
    \hline
    \textbf{Ancestry} &  & \\
    \hline
    \textbf{Source} &  & \\
    \hline
    \textbf{Enabled By} &  & \\
    \hline
    \textbf{Supports} &  & \\
    \hline
    \textbf{Derived From} &  & \\
    \hline
    \textbf{Reviewed Status} & Review meaning, business boundary, information context and any unresolved ancestry. & \\
    \hline
  \end{tabularx}
\end{table}

\begin{table}[htbp]
  \centering
  \footnotesize
  \begin{tabularx}{1\textwidth}{|l|l|X|}
    \hline
    \multicolumn{3}{c}{\textbf{Contained Features}} \\
    \hline
    \textbf{Feature Name} & \textbf{Feature Id} & \textbf{Comment} \\
    \hline
    \multicolumn{3}{|l|}{\emph{No contained Features are currently identified.}} \\
    \hline
  \end{tabularx}
\end{table}

\textbf{Intent:} Clinical Device Lifecycle Management identifies the organisational ability to undertake the stated health-service responsibility. The capability describes what the enterprise must be able to do; it does not assert that Harmonia performs or owns that activity.

\textbf{Boundary:} Harmonia relevance may arise because this activity establishes the context in which information is created, needed, interpreted, related, exchanged or presented. Responsibility for the underlying business decision or care activity remains with the appropriate actor or organisation.

Unresolved CT2 framing: [CT2?] Health Product Management

\chapter{CT1 - Clinical Quality, Safety \& Improvement}\label{cap:bc-e95c383f}
\noindent\textbf{Reference ID:} \texttt{BC-E95C383F}\quad \textbf{Status:} Established working hierarchy\par
\noindent\textit{Source: Canonical Business Capability model / working map}\par
\textbf{Intent:} Clinical Quality, Safety \& Improvement identifies the organisational ability to undertake the stated health-service responsibility. The capability describes what the enterprise must be able to do; it does not assert that Harmonia performs or owns that activity.

\textbf{Boundary:} Harmonia relevance may arise because this activity establishes the context in which information is created, needed, interpreted, related, exchanged or presented. Responsibility for the underlying business decision or care activity remains with the appropriate actor or organisation.

\reviewstatus{Check scope and information relevance; no business authority transfers to Harmonia.}
\section{CT2 - Clinical Safety Management}\label{cap:bc-5b2a0e55}
\vspace{-0.5em}
\begin{table}[htbp]
  \centering
  \footnotesize
  \begin{tabularx}{1\textwidth}{|l|l|X|}
    \hline
    \multicolumn{3}{c}{\textbf{Capability Metadata}} \\
    \hline
    \textbf{Name} & \textbf{Value} & \textbf{Comment and/or Description} \\
    \hline
    \textbf{Reference Id} & \texttt{BC-5B2A0E55} & \\
    \hline
    \textbf{Status} & Established working hierarchy & \\
    \hline
    \textbf{Ancestry} &  & \\
    \hline
    \textbf{Source} &  & \\
    \hline
    \textbf{Enabled By} &  & \\
    \hline
    \textbf{Supports} &  & \\
    \hline
    \textbf{Derived From} &  & \\
    \hline
    \textbf{Reviewed Status} &  & \\
    \hline
  \end{tabularx}
\end{table}

\begin{table}[htbp]
  \centering
  \footnotesize
  \begin{tabularx}{1\textwidth}{|l|l|X|}
    \hline
    \multicolumn{3}{c}{\textbf{Contained Capabilities}} \\
    \hline
    \textbf{Capability Name} & \textbf{Capability Id} & \textbf{Comment} \\
    \hline
    Clinical Risk Identification & \texttt{cap:bc-7e3fbc64} (at \ref{cap:bc-7e3fbc64}) & \\
    \hline
    Clinical Risk Mitigation Management & \texttt{cap:bc-5121c6f1} (at \ref{cap:bc-5121c6f1}) & \\
    \hline
    Clinical Risk Response Management & \texttt{cap:bc-72dcf354} (at \ref{cap:bc-72dcf354}) & \\
    \hline
  \end{tabularx}
\end{table}

\textbf{Intent:} Clinical Safety Management identifies the organisational ability to undertake the stated health-service responsibility. The capability describes what the enterprise must be able to do; it does not assert that Harmonia performs or owns that activity.

\textbf{Boundary:} Harmonia relevance may arise because this activity establishes the context in which information is created, needed, interpreted, related, exchanged or presented. Responsibility for the underlying business decision or care activity remains with the appropriate actor or organisation.

\subsection{CT3 - Clinical Risk Identification}\label{cap:bc-7e3fbc64}
\vspace{-0.5em}
\begin{table}[htbp]
  \centering
  \footnotesize
  \begin{tabularx}{1\textwidth}{|l|l|X|}
    \hline
    \multicolumn{3}{c}{\textbf{Capability Metadata}} \\
    \hline
    \textbf{Name} & \textbf{Value} & \textbf{Comment and/or Description} \\
    \hline
    \textbf{Reference Id} & \texttt{BC-7E3FBC64} & \\
    \hline
    \textbf{Status} & Established working hierarchy & \\
    \hline
    \textbf{Ancestry} &  & \\
    \hline
    \textbf{Source} &  & \\
    \hline
    \textbf{Enabled By} &  & \\
    \hline
    \textbf{Supports} &  & \\
    \hline
    \textbf{Derived From} &  & \\
    \hline
    \textbf{Reviewed Status} & Review meaning, business boundary, information context and any unresolved ancestry. & \\
    \hline
  \end{tabularx}
\end{table}

\begin{table}[htbp]
  \centering
  \footnotesize
  \begin{tabularx}{1\textwidth}{|l|l|X|}
    \hline
    \multicolumn{3}{c}{\textbf{Contained Features}} \\
    \hline
    \textbf{Feature Name} & \textbf{Feature Id} & \textbf{Comment} \\
    \hline
    \multicolumn{3}{|l|}{\emph{No contained Features are currently identified.}} \\
    \hline
  \end{tabularx}
\end{table}

\textbf{Intent:} Clinical Risk Identification identifies the organisational ability to undertake the stated health-service responsibility. The capability describes what the enterprise must be able to do; it does not assert that Harmonia performs or owns that activity.

\textbf{Boundary:} Harmonia relevance may arise because this activity establishes the context in which information is created, needed, interpreted, related, exchanged or presented. Responsibility for the underlying business decision or care activity remains with the appropriate actor or organisation.

\subsection{CT3 - Clinical Risk Mitigation Management}\label{cap:bc-5121c6f1}
\vspace{-0.5em}
\begin{table}[htbp]
  \centering
  \footnotesize
  \begin{tabularx}{1\textwidth}{|l|l|X|}
    \hline
    \multicolumn{3}{c}{\textbf{Capability Metadata}} \\
    \hline
    \textbf{Name} & \textbf{Value} & \textbf{Comment and/or Description} \\
    \hline
    \textbf{Reference Id} & \texttt{BC-5121C6F1} & \\
    \hline
    \textbf{Status} & Established working hierarchy & \\
    \hline
    \textbf{Ancestry} &  & \\
    \hline
    \textbf{Source} &  & \\
    \hline
    \textbf{Enabled By} &  & \\
    \hline
    \textbf{Supports} &  & \\
    \hline
    \textbf{Derived From} &  & \\
    \hline
    \textbf{Reviewed Status} & Review meaning, business boundary, information context and any unresolved ancestry. & \\
    \hline
  \end{tabularx}
\end{table}

\begin{table}[htbp]
  \centering
  \footnotesize
  \begin{tabularx}{1\textwidth}{|l|l|X|}
    \hline
    \multicolumn{3}{c}{\textbf{Contained Features}} \\
    \hline
    \textbf{Feature Name} & \textbf{Feature Id} & \textbf{Comment} \\
    \hline
    \multicolumn{3}{|l|}{\emph{No contained Features are currently identified.}} \\
    \hline
  \end{tabularx}
\end{table}

\textbf{Intent:} Clinical Risk Mitigation Management identifies the organisational ability to undertake the stated health-service responsibility. The capability describes what the enterprise must be able to do; it does not assert that Harmonia performs or owns that activity.

\textbf{Boundary:} Harmonia relevance may arise because this activity establishes the context in which information is created, needed, interpreted, related, exchanged or presented. Responsibility for the underlying business decision or care activity remains with the appropriate actor or organisation.

\subsection{CT3 - Clinical Risk Response Management}\label{cap:bc-72dcf354}
\vspace{-0.5em}
\begin{table}[htbp]
  \centering
  \footnotesize
  \begin{tabularx}{1\textwidth}{|l|l|X|}
    \hline
    \multicolumn{3}{c}{\textbf{Capability Metadata}} \\
    \hline
    \textbf{Name} & \textbf{Value} & \textbf{Comment and/or Description} \\
    \hline
    \textbf{Reference Id} & \texttt{BC-72DCF354} & \\
    \hline
    \textbf{Status} & Established working hierarchy & \\
    \hline
    \textbf{Ancestry} &  & \\
    \hline
    \textbf{Source} &  & \\
    \hline
    \textbf{Enabled By} &  & \\
    \hline
    \textbf{Supports} &  & \\
    \hline
    \textbf{Derived From} &  & \\
    \hline
    \textbf{Reviewed Status} & Review meaning, business boundary, information context and any unresolved ancestry. & \\
    \hline
  \end{tabularx}
\end{table}

\begin{table}[htbp]
  \centering
  \footnotesize
  \begin{tabularx}{1\textwidth}{|l|l|X|}
    \hline
    \multicolumn{3}{c}{\textbf{Contained Features}} \\
    \hline
    \textbf{Feature Name} & \textbf{Feature Id} & \textbf{Comment} \\
    \hline
    \multicolumn{3}{|l|}{\emph{No contained Features are currently identified.}} \\
    \hline
  \end{tabularx}
\end{table}

\textbf{Intent:} Clinical Risk Response Management identifies the organisational ability to undertake the stated health-service responsibility. The capability describes what the enterprise must be able to do; it does not assert that Harmonia performs or owns that activity.

\textbf{Boundary:} Harmonia relevance may arise because this activity establishes the context in which information is created, needed, interpreted, related, exchanged or presented. Responsibility for the underlying business decision or care activity remains with the appropriate actor or organisation.

\section{CT2 - Clinical Quality Management}\label{cap:bc-077229cb}
\vspace{-0.5em}
\begin{table}[htbp]
  \centering
  \footnotesize
  \begin{tabularx}{1\textwidth}{|l|l|X|}
    \hline
    \multicolumn{3}{c}{\textbf{Capability Metadata}} \\
    \hline
    \textbf{Name} & \textbf{Value} & \textbf{Comment and/or Description} \\
    \hline
    \textbf{Reference Id} & \texttt{BC-077229CB} & \\
    \hline
    \textbf{Status} & Established working hierarchy & \\
    \hline
    \textbf{Ancestry} &  & \\
    \hline
    \textbf{Source} &  & \\
    \hline
    \textbf{Enabled By} &  & \\
    \hline
    \textbf{Supports} &  & \\
    \hline
    \textbf{Derived From} &  & \\
    \hline
    \textbf{Reviewed Status} &  & \\
    \hline
  \end{tabularx}
\end{table}

\begin{table}[htbp]
  \centering
  \footnotesize
  \begin{tabularx}{1\textwidth}{|l|l|X|}
    \hline
    \multicolumn{3}{c}{\textbf{Contained Capabilities}} \\
    \hline
    \textbf{Capability Name} & \textbf{Capability Id} & \textbf{Comment} \\
    \hline
    Clinical Quality Criteria Management & \texttt{cap:bc-807007a1} (at \ref{cap:bc-807007a1}) & \\
    \hline
    Clinical Quality Evidence Management & \texttt{cap:bc-fc14d2e2} (at \ref{cap:bc-fc14d2e2}) & \\
    \hline
    Clinical Quality Review \& Assessment Management & \texttt{cap:bc-6884dabf} (at \ref{cap:bc-6884dabf}) & \\
    \hline
    Clinical Quality Reporting & \texttt{cap:bc-6eeb86b4} (at \ref{cap:bc-6eeb86b4}) & \\
    \hline
  \end{tabularx}
\end{table}

\textbf{Intent:} Clinical Quality Management identifies the organisational ability to undertake the stated health-service responsibility. The capability describes what the enterprise must be able to do; it does not assert that Harmonia performs or owns that activity.

\textbf{Boundary:} Harmonia relevance may arise because this activity establishes the context in which information is created, needed, interpreted, related, exchanged or presented. Responsibility for the underlying business decision or care activity remains with the appropriate actor or organisation.

\subsection{CT3 - Clinical Quality Criteria Management}\label{cap:bc-807007a1}
\vspace{-0.5em}
\begin{table}[htbp]
  \centering
  \footnotesize
  \begin{tabularx}{1\textwidth}{|l|l|X|}
    \hline
    \multicolumn{3}{c}{\textbf{Capability Metadata}} \\
    \hline
    \textbf{Name} & \textbf{Value} & \textbf{Comment and/or Description} \\
    \hline
    \textbf{Reference Id} & \texttt{BC-807007A1} & \\
    \hline
    \textbf{Status} & Established working hierarchy & \\
    \hline
    \textbf{Ancestry} &  & \\
    \hline
    \textbf{Source} &  & \\
    \hline
    \textbf{Enabled By} &  & \\
    \hline
    \textbf{Supports} &  & \\
    \hline
    \textbf{Derived From} &  & \\
    \hline
    \textbf{Reviewed Status} & Review meaning, business boundary, information context and any unresolved ancestry. & \\
    \hline
  \end{tabularx}
\end{table}

\begin{table}[htbp]
  \centering
  \footnotesize
  \begin{tabularx}{1\textwidth}{|l|l|X|}
    \hline
    \multicolumn{3}{c}{\textbf{Contained Features}} \\
    \hline
    \textbf{Feature Name} & \textbf{Feature Id} & \textbf{Comment} \\
    \hline
    \multicolumn{3}{|l|}{\emph{No contained Features are currently identified.}} \\
    \hline
  \end{tabularx}
\end{table}

\textbf{Intent:} Clinical Quality Criteria Management identifies the organisational ability to undertake the stated health-service responsibility. The capability describes what the enterprise must be able to do; it does not assert that Harmonia performs or owns that activity.

\textbf{Boundary:} Harmonia relevance may arise because this activity establishes the context in which information is created, needed, interpreted, related, exchanged or presented. Responsibility for the underlying business decision or care activity remains with the appropriate actor or organisation.

\subsection{CT3 - Clinical Quality Evidence Management}\label{cap:bc-fc14d2e2}
\vspace{-0.5em}
\begin{table}[htbp]
  \centering
  \footnotesize
  \begin{tabularx}{1\textwidth}{|l|l|X|}
    \hline
    \multicolumn{3}{c}{\textbf{Capability Metadata}} \\
    \hline
    \textbf{Name} & \textbf{Value} & \textbf{Comment and/or Description} \\
    \hline
    \textbf{Reference Id} & \texttt{BC-FC14D2E2} & \\
    \hline
    \textbf{Status} & Established working hierarchy & \\
    \hline
    \textbf{Ancestry} &  & \\
    \hline
    \textbf{Source} &  & \\
    \hline
    \textbf{Enabled By} &  & \\
    \hline
    \textbf{Supports} &  & \\
    \hline
    \textbf{Derived From} &  & \\
    \hline
    \textbf{Reviewed Status} & Review meaning, business boundary, information context and any unresolved ancestry. & \\
    \hline
  \end{tabularx}
\end{table}

\begin{table}[htbp]
  \centering
  \footnotesize
  \begin{tabularx}{1\textwidth}{|l|l|X|}
    \hline
    \multicolumn{3}{c}{\textbf{Contained Features}} \\
    \hline
    \textbf{Feature Name} & \textbf{Feature Id} & \textbf{Comment} \\
    \hline
    \multicolumn{3}{|l|}{\emph{No contained Features are currently identified.}} \\
    \hline
  \end{tabularx}
\end{table}

\textbf{Intent:} Clinical Quality Evidence Management identifies the organisational ability to undertake the stated health-service responsibility. The capability describes what the enterprise must be able to do; it does not assert that Harmonia performs or owns that activity.

\textbf{Boundary:} Harmonia relevance may arise because this activity establishes the context in which information is created, needed, interpreted, related, exchanged or presented. Responsibility for the underlying business decision or care activity remains with the appropriate actor or organisation.

\subsection{CT3 - Clinical Quality Review \& Assessment Management}\label{cap:bc-6884dabf}
\vspace{-0.5em}
\begin{table}[htbp]
  \centering
  \footnotesize
  \begin{tabularx}{1\textwidth}{|l|l|X|}
    \hline
    \multicolumn{3}{c}{\textbf{Capability Metadata}} \\
    \hline
    \textbf{Name} & \textbf{Value} & \textbf{Comment and/or Description} \\
    \hline
    \textbf{Reference Id} & \texttt{BC-6884DABF} & \\
    \hline
    \textbf{Status} & Established working hierarchy & \\
    \hline
    \textbf{Ancestry} &  & \\
    \hline
    \textbf{Source} &  & \\
    \hline
    \textbf{Enabled By} &  & \\
    \hline
    \textbf{Supports} &  & \\
    \hline
    \textbf{Derived From} &  & \\
    \hline
    \textbf{Reviewed Status} & Review meaning, business boundary, information context and any unresolved ancestry. & \\
    \hline
  \end{tabularx}
\end{table}

\begin{table}[htbp]
  \centering
  \footnotesize
  \begin{tabularx}{1\textwidth}{|l|l|X|}
    \hline
    \multicolumn{3}{c}{\textbf{Contained Features}} \\
    \hline
    \textbf{Feature Name} & \textbf{Feature Id} & \textbf{Comment} \\
    \hline
    \multicolumn{3}{|l|}{\emph{No contained Features are currently identified.}} \\
    \hline
  \end{tabularx}
\end{table}

\textbf{Intent:} Clinical Quality Review \& Assessment Management identifies the organisational ability to undertake the stated health-service responsibility. The capability describes what the enterprise must be able to do; it does not assert that Harmonia performs or owns that activity.

\textbf{Boundary:} Harmonia relevance may arise because this activity establishes the context in which information is created, needed, interpreted, related, exchanged or presented. Responsibility for the underlying business decision or care activity remains with the appropriate actor or organisation.

\subsection{CT3 - Clinical Quality Reporting}\label{cap:bc-6eeb86b4}
\vspace{-0.5em}
\begin{table}[htbp]
  \centering
  \footnotesize
  \begin{tabularx}{1\textwidth}{|l|l|X|}
    \hline
    \multicolumn{3}{c}{\textbf{Capability Metadata}} \\
    \hline
    \textbf{Name} & \textbf{Value} & \textbf{Comment and/or Description} \\
    \hline
    \textbf{Reference Id} & \texttt{BC-6EEB86B4} & \\
    \hline
    \textbf{Status} & Established working hierarchy & \\
    \hline
    \textbf{Ancestry} &  & \\
    \hline
    \textbf{Source} &  & \\
    \hline
    \textbf{Enabled By} &  & \\
    \hline
    \textbf{Supports} &  & \\
    \hline
    \textbf{Derived From} &  & \\
    \hline
    \textbf{Reviewed Status} & Review meaning, business boundary, information context and any unresolved ancestry. & \\
    \hline
  \end{tabularx}
\end{table}

\begin{table}[htbp]
  \centering
  \footnotesize
  \begin{tabularx}{1\textwidth}{|l|l|X|}
    \hline
    \multicolumn{3}{c}{\textbf{Contained Features}} \\
    \hline
    \textbf{Feature Name} & \textbf{Feature Id} & \textbf{Comment} \\
    \hline
    \multicolumn{3}{|l|}{\emph{No contained Features are currently identified.}} \\
    \hline
  \end{tabularx}
\end{table}

\textbf{Intent:} Clinical Quality Reporting identifies the organisational ability to undertake the stated health-service responsibility. The capability describes what the enterprise must be able to do; it does not assert that Harmonia performs or owns that activity.

\textbf{Boundary:} Harmonia relevance may arise because this activity establishes the context in which information is created, needed, interpreted, related, exchanged or presented. Responsibility for the underlying business decision or care activity remains with the appropriate actor or organisation.

\chapter{CT1 - Community Health \& Wellbeing}\label{cap:bc-31c8dc2c}
\noindent\textbf{Reference ID:} \texttt{BC-31C8DC2C}\quad \textbf{Status:} Established working hierarchy\par
\noindent\textit{Source: Canonical Business Capability model / working map}\par
\textbf{Intent:} Community Health \& Wellbeing identifies the organisational ability to undertake the stated health-service responsibility. The capability describes what the enterprise must be able to do; it does not assert that Harmonia performs or owns that activity.

\textbf{Boundary:} Harmonia relevance may arise because this activity establishes the context in which information is created, needed, interpreted, related, exchanged or presented. Responsibility for the underlying business decision or care activity remains with the appropriate actor or organisation.

\reviewstatus{Check scope and information relevance; no business authority transfers to Harmonia.}
\section{CT2 - Preventative Health}\label{cap:bc-80da636c}
\vspace{-0.5em}
\begin{table}[htbp]
  \centering
  \footnotesize
  \begin{tabularx}{1\textwidth}{|l|l|X|}
    \hline
    \multicolumn{3}{c}{\textbf{Capability Metadata}} \\
    \hline
    \textbf{Name} & \textbf{Value} & \textbf{Comment and/or Description} \\
    \hline
    \textbf{Reference Id} & \texttt{BC-80DA636C} & \\
    \hline
    \textbf{Status} & Established working hierarchy & \\
    \hline
    \textbf{Ancestry} &  & \\
    \hline
    \textbf{Source} &  & \\
    \hline
    \textbf{Enabled By} &  & \\
    \hline
    \textbf{Supports} &  & \\
    \hline
    \textbf{Derived From} &  & \\
    \hline
    \textbf{Reviewed Status} &  & \\
    \hline
  \end{tabularx}
\end{table}

\begin{table}[htbp]
  \centering
  \footnotesize
  \begin{tabularx}{1\textwidth}{|l|l|X|}
    \hline
    \multicolumn{3}{c}{\textbf{Contained Capabilities}} \\
    \hline
    \textbf{Capability Name} & \textbf{Capability Id} & \textbf{Comment} \\
    \hline
    Preventative Health Problem Identification & \texttt{cap:bc-1fad6263} (at \ref{cap:bc-1fad6263}) & \\
    \hline
    Preventative Health Programme Management & \texttt{cap:bc-b3bdd9a1} (at \ref{cap:bc-b3bdd9a1}) & \\
    \hline
    Preventative Health Education Services & \texttt{cap:bc-57b54ebe} (at \ref{cap:bc-57b54ebe}) & \\
    \hline
    Preventative Health Programme Monitoring & \texttt{cap:bc-5202aef0} (at \ref{cap:bc-5202aef0}) & \\
    \hline
    Preventative Health Reporting & \texttt{cap:bc-e55d2726} (at \ref{cap:bc-e55d2726}) & \\
    \hline
    Preventative Health Cohort Management & \texttt{cap:bc-1800bcd9} (at \ref{cap:bc-1800bcd9}) & \\
    \hline
  \end{tabularx}
\end{table}

\textbf{Intent:} Preventative Health identifies the organisational ability to undertake the stated health-service responsibility. The capability describes what the enterprise must be able to do; it does not assert that Harmonia performs or owns that activity.

\textbf{Boundary:} Harmonia relevance may arise because this activity establishes the context in which information is created, needed, interpreted, related, exchanged or presented. Responsibility for the underlying business decision or care activity remains with the appropriate actor or organisation.

\subsection{CT3 - Preventative Health Problem Identification}\label{cap:bc-1fad6263}
\vspace{-0.5em}
\begin{table}[htbp]
  \centering
  \footnotesize
  \begin{tabularx}{1\textwidth}{|l|l|X|}
    \hline
    \multicolumn{3}{c}{\textbf{Capability Metadata}} \\
    \hline
    \textbf{Name} & \textbf{Value} & \textbf{Comment and/or Description} \\
    \hline
    \textbf{Reference Id} & \texttt{BC-1FAD6263} & \\
    \hline
    \textbf{Status} & Established working hierarchy & \\
    \hline
    \textbf{Ancestry} &  & \\
    \hline
    \textbf{Source} &  & \\
    \hline
    \textbf{Enabled By} &  & \\
    \hline
    \textbf{Supports} &  & \\
    \hline
    \textbf{Derived From} &  & \\
    \hline
    \textbf{Reviewed Status} & Review meaning, business boundary, information context and any unresolved ancestry. & \\
    \hline
  \end{tabularx}
\end{table}

\begin{table}[htbp]
  \centering
  \footnotesize
  \begin{tabularx}{1\textwidth}{|l|l|X|}
    \hline
    \multicolumn{3}{c}{\textbf{Contained Features}} \\
    \hline
    \textbf{Feature Name} & \textbf{Feature Id} & \textbf{Comment} \\
    \hline
    \multicolumn{3}{|l|}{\emph{No contained Features are currently identified.}} \\
    \hline
  \end{tabularx}
\end{table}

\textbf{Intent:} Preventative Health Problem Identification identifies the organisational ability to undertake the stated health-service responsibility. The capability describes what the enterprise must be able to do; it does not assert that Harmonia performs or owns that activity.

\textbf{Boundary:} Harmonia relevance may arise because this activity establishes the context in which information is created, needed, interpreted, related, exchanged or presented. Responsibility for the underlying business decision or care activity remains with the appropriate actor or organisation.

\subsection{CT3 - Preventative Health Programme Management}\label{cap:bc-b3bdd9a1}
\vspace{-0.5em}
\begin{table}[htbp]
  \centering
  \footnotesize
  \begin{tabularx}{1\textwidth}{|l|l|X|}
    \hline
    \multicolumn{3}{c}{\textbf{Capability Metadata}} \\
    \hline
    \textbf{Name} & \textbf{Value} & \textbf{Comment and/or Description} \\
    \hline
    \textbf{Reference Id} & \texttt{BC-B3BDD9A1} & \\
    \hline
    \textbf{Status} & Established working hierarchy & \\
    \hline
    \textbf{Ancestry} &  & \\
    \hline
    \textbf{Source} &  & \\
    \hline
    \textbf{Enabled By} &  & \\
    \hline
    \textbf{Supports} &  & \\
    \hline
    \textbf{Derived From} &  & \\
    \hline
    \textbf{Reviewed Status} & Review meaning, business boundary, information context and any unresolved ancestry. & \\
    \hline
  \end{tabularx}
\end{table}

\begin{table}[htbp]
  \centering
  \footnotesize
  \begin{tabularx}{1\textwidth}{|l|l|X|}
    \hline
    \multicolumn{3}{c}{\textbf{Contained Features}} \\
    \hline
    \textbf{Feature Name} & \textbf{Feature Id} & \textbf{Comment} \\
    \hline
    \multicolumn{3}{|l|}{\emph{No contained Features are currently identified.}} \\
    \hline
  \end{tabularx}
\end{table}

\textbf{Intent:} Preventative Health Programme Management identifies the organisational ability to undertake the stated health-service responsibility. The capability describes what the enterprise must be able to do; it does not assert that Harmonia performs or owns that activity.

\textbf{Boundary:} Harmonia relevance may arise because this activity establishes the context in which information is created, needed, interpreted, related, exchanged or presented. Responsibility for the underlying business decision or care activity remains with the appropriate actor or organisation.

\subsection{CT3 - Preventative Health Education Services}\label{cap:bc-57b54ebe}
\vspace{-0.5em}
\begin{table}[htbp]
  \centering
  \footnotesize
  \begin{tabularx}{1\textwidth}{|l|l|X|}
    \hline
    \multicolumn{3}{c}{\textbf{Capability Metadata}} \\
    \hline
    \textbf{Name} & \textbf{Value} & \textbf{Comment and/or Description} \\
    \hline
    \textbf{Reference Id} & \texttt{BC-57B54EBE} & \\
    \hline
    \textbf{Status} & Established working hierarchy & \\
    \hline
    \textbf{Ancestry} &  & \\
    \hline
    \textbf{Source} &  & \\
    \hline
    \textbf{Enabled By} &  & \\
    \hline
    \textbf{Supports} &  & \\
    \hline
    \textbf{Derived From} &  & \\
    \hline
    \textbf{Reviewed Status} & Review meaning, business boundary, information context and any unresolved ancestry. & \\
    \hline
  \end{tabularx}
\end{table}

\begin{table}[htbp]
  \centering
  \footnotesize
  \begin{tabularx}{1\textwidth}{|l|l|X|}
    \hline
    \multicolumn{3}{c}{\textbf{Contained Features}} \\
    \hline
    \textbf{Feature Name} & \textbf{Feature Id} & \textbf{Comment} \\
    \hline
    \multicolumn{3}{|l|}{\emph{No contained Features are currently identified.}} \\
    \hline
  \end{tabularx}
\end{table}

\textbf{Intent:} Preventative Health Education Services identifies the organisational ability to undertake the stated health-service responsibility. The capability describes what the enterprise must be able to do; it does not assert that Harmonia performs or owns that activity.

\textbf{Boundary:} Harmonia relevance may arise because this activity establishes the context in which information is created, needed, interpreted, related, exchanged or presented. Responsibility for the underlying business decision or care activity remains with the appropriate actor or organisation.

\subsection{CT3 - Preventative Health Programme Monitoring}\label{cap:bc-5202aef0}
\vspace{-0.5em}
\begin{table}[htbp]
  \centering
  \footnotesize
  \begin{tabularx}{1\textwidth}{|l|l|X|}
    \hline
    \multicolumn{3}{c}{\textbf{Capability Metadata}} \\
    \hline
    \textbf{Name} & \textbf{Value} & \textbf{Comment and/or Description} \\
    \hline
    \textbf{Reference Id} & \texttt{BC-5202AEF0} & \\
    \hline
    \textbf{Status} & Established working hierarchy & \\
    \hline
    \textbf{Ancestry} &  & \\
    \hline
    \textbf{Source} &  & \\
    \hline
    \textbf{Enabled By} &  & \\
    \hline
    \textbf{Supports} &  & \\
    \hline
    \textbf{Derived From} &  & \\
    \hline
    \textbf{Reviewed Status} & Review meaning, business boundary, information context and any unresolved ancestry. & \\
    \hline
  \end{tabularx}
\end{table}

\begin{table}[htbp]
  \centering
  \footnotesize
  \begin{tabularx}{1\textwidth}{|l|l|X|}
    \hline
    \multicolumn{3}{c}{\textbf{Contained Features}} \\
    \hline
    \textbf{Feature Name} & \textbf{Feature Id} & \textbf{Comment} \\
    \hline
    \multicolumn{3}{|l|}{\emph{No contained Features are currently identified.}} \\
    \hline
  \end{tabularx}
\end{table}

\textbf{Intent:} Preventative Health Programme Monitoring identifies the organisational ability to undertake the stated health-service responsibility. The capability describes what the enterprise must be able to do; it does not assert that Harmonia performs or owns that activity.

\textbf{Boundary:} Harmonia relevance may arise because this activity establishes the context in which information is created, needed, interpreted, related, exchanged or presented. Responsibility for the underlying business decision or care activity remains with the appropriate actor or organisation.

\subsection{CT3 - Preventative Health Reporting}\label{cap:bc-e55d2726}
\vspace{-0.5em}
\begin{table}[htbp]
  \centering
  \footnotesize
  \begin{tabularx}{1\textwidth}{|l|l|X|}
    \hline
    \multicolumn{3}{c}{\textbf{Capability Metadata}} \\
    \hline
    \textbf{Name} & \textbf{Value} & \textbf{Comment and/or Description} \\
    \hline
    \textbf{Reference Id} & \texttt{BC-E55D2726} & \\
    \hline
    \textbf{Status} & Established working hierarchy & \\
    \hline
    \textbf{Ancestry} &  & \\
    \hline
    \textbf{Source} &  & \\
    \hline
    \textbf{Enabled By} &  & \\
    \hline
    \textbf{Supports} &  & \\
    \hline
    \textbf{Derived From} &  & \\
    \hline
    \textbf{Reviewed Status} & Review meaning, business boundary, information context and any unresolved ancestry. & \\
    \hline
  \end{tabularx}
\end{table}

\begin{table}[htbp]
  \centering
  \footnotesize
  \begin{tabularx}{1\textwidth}{|l|l|X|}
    \hline
    \multicolumn{3}{c}{\textbf{Contained Features}} \\
    \hline
    \textbf{Feature Name} & \textbf{Feature Id} & \textbf{Comment} \\
    \hline
    \multicolumn{3}{|l|}{\emph{No contained Features are currently identified.}} \\
    \hline
  \end{tabularx}
\end{table}

\textbf{Intent:} Preventative Health Reporting identifies the organisational ability to undertake the stated health-service responsibility. The capability describes what the enterprise must be able to do; it does not assert that Harmonia performs or owns that activity.

\textbf{Boundary:} Harmonia relevance may arise because this activity establishes the context in which information is created, needed, interpreted, related, exchanged or presented. Responsibility for the underlying business decision or care activity remains with the appropriate actor or organisation.

\subsection{CT3 - Preventative Health Cohort Management}\label{cap:bc-1800bcd9}
\vspace{-0.5em}
\begin{table}[htbp]
  \centering
  \footnotesize
  \begin{tabularx}{1\textwidth}{|l|l|X|}
    \hline
    \multicolumn{3}{c}{\textbf{Capability Metadata}} \\
    \hline
    \textbf{Name} & \textbf{Value} & \textbf{Comment and/or Description} \\
    \hline
    \textbf{Reference Id} & \texttt{BC-1800BCD9} & \\
    \hline
    \textbf{Status} & Established working hierarchy & \\
    \hline
    \textbf{Ancestry} &  & \\
    \hline
    \textbf{Source} &  & \\
    \hline
    \textbf{Enabled By} &  & \\
    \hline
    \textbf{Supports} &  & \\
    \hline
    \textbf{Derived From} &  & \\
    \hline
    \textbf{Reviewed Status} & Review meaning, business boundary, information context and any unresolved ancestry. & \\
    \hline
  \end{tabularx}
\end{table}

\begin{table}[htbp]
  \centering
  \footnotesize
  \begin{tabularx}{1\textwidth}{|l|l|X|}
    \hline
    \multicolumn{3}{c}{\textbf{Contained Features}} \\
    \hline
    \textbf{Feature Name} & \textbf{Feature Id} & \textbf{Comment} \\
    \hline
    \multicolumn{3}{|l|}{\emph{No contained Features are currently identified.}} \\
    \hline
  \end{tabularx}
\end{table}

\textbf{Intent:} Preventative Health Cohort Management identifies the organisational ability to undertake the stated health-service responsibility. The capability describes what the enterprise must be able to do; it does not assert that Harmonia performs or owns that activity.

\textbf{Boundary:} Harmonia relevance may arise because this activity establishes the context in which information is created, needed, interpreted, related, exchanged or presented. Responsibility for the underlying business decision or care activity remains with the appropriate actor or organisation.

\chapter{CT1 - Population Health \& Health-System Planning}\label{cap:bc-28ccee27}
\noindent\textbf{Reference ID:} \texttt{BC-28CCEE27}\quad \textbf{Status:} Established working hierarchy\par
\noindent\textit{Source: Canonical Business Capability model / working map}\par
\textbf{Intent:} Population Health \& Health-System Planning identifies the organisational ability to undertake the stated health-service responsibility. The capability describes what the enterprise must be able to do; it does not assert that Harmonia performs or owns that activity.

\textbf{Boundary:} Harmonia relevance may arise because this activity establishes the context in which information is created, needed, interpreted, related, exchanged or presented. Responsibility for the underlying business decision or care activity remains with the appropriate actor or organisation.

\reviewstatus{Check scope and information relevance; no business authority transfers to Harmonia.}
\section{CT2 - Population Health Surveillance}\label{cap:bc-a6f4ca49}
\vspace{-0.5em}
\begin{table}[htbp]
  \centering
  \footnotesize
  \begin{tabularx}{1\textwidth}{|l|l|X|}
    \hline
    \multicolumn{3}{c}{\textbf{Capability Metadata}} \\
    \hline
    \textbf{Name} & \textbf{Value} & \textbf{Comment and/or Description} \\
    \hline
    \textbf{Reference Id} & \texttt{BC-A6F4CA49} & \\
    \hline
    \textbf{Status} & Established working hierarchy & \\
    \hline
    \textbf{Ancestry} &  & \\
    \hline
    \textbf{Source} &  & \\
    \hline
    \textbf{Enabled By} &  & \\
    \hline
    \textbf{Supports} &  & \\
    \hline
    \textbf{Derived From} &  & \\
    \hline
    \textbf{Reviewed Status} &  & \\
    \hline
  \end{tabularx}
\end{table}

\begin{table}[htbp]
  \centering
  \footnotesize
  \begin{tabularx}{1\textwidth}{|l|l|X|}
    \hline
    \multicolumn{3}{c}{\textbf{Contained Capabilities}} \\
    \hline
    \textbf{Capability Name} & \textbf{Capability Id} & \textbf{Comment} \\
    \hline
    Population Health Longitudinal Assessment & \texttt{cap:bc-f1d5d26c} (at \ref{cap:bc-f1d5d26c}) & \\
    \hline
    Population Health Surveillance Management & \texttt{cap:bc-afc31f01} (at \ref{cap:bc-afc31f01}) & \\
    \hline
    Population Health Event Identification & \texttt{cap:bc-bca30d47} (at \ref{cap:bc-bca30d47}) & \\
    \hline
    Notifiable Health Event Management & \texttt{cap:bc-5eaa6955} (at \ref{cap:bc-5eaa6955}) & \\
    \hline
    Population Health Reporting & \texttt{cap:bc-d8c15299} (at \ref{cap:bc-d8c15299}) & \\
    \hline
  \end{tabularx}
\end{table}

\textbf{Intent:} Population Health Surveillance identifies the organisational ability to undertake the stated health-service responsibility. The capability describes what the enterprise must be able to do; it does not assert that Harmonia performs or owns that activity.

\textbf{Boundary:} Harmonia relevance may arise because this activity establishes the context in which information is created, needed, interpreted, related, exchanged or presented. Responsibility for the underlying business decision or care activity remains with the appropriate actor or organisation.

\subsection{CT3 - Population Health Longitudinal Assessment}\label{cap:bc-f1d5d26c}
\vspace{-0.5em}
\begin{table}[htbp]
  \centering
  \footnotesize
  \begin{tabularx}{1\textwidth}{|l|l|X|}
    \hline
    \multicolumn{3}{c}{\textbf{Capability Metadata}} \\
    \hline
    \textbf{Name} & \textbf{Value} & \textbf{Comment and/or Description} \\
    \hline
    \textbf{Reference Id} & \texttt{BC-F1D5D26C} & \\
    \hline
    \textbf{Status} & Established working hierarchy & \\
    \hline
    \textbf{Ancestry} &  & \\
    \hline
    \textbf{Source} &  & \\
    \hline
    \textbf{Enabled By} &  & \\
    \hline
    \textbf{Supports} &  & \\
    \hline
    \textbf{Derived From} &  & \\
    \hline
    \textbf{Reviewed Status} & Review meaning, business boundary, information context and any unresolved ancestry. & \\
    \hline
  \end{tabularx}
\end{table}

\begin{table}[htbp]
  \centering
  \footnotesize
  \begin{tabularx}{1\textwidth}{|l|l|X|}
    \hline
    \multicolumn{3}{c}{\textbf{Contained Features}} \\
    \hline
    \textbf{Feature Name} & \textbf{Feature Id} & \textbf{Comment} \\
    \hline
    \multicolumn{3}{|l|}{\emph{No contained Features are currently identified.}} \\
    \hline
  \end{tabularx}
\end{table}

\textbf{Intent:} Population Health Longitudinal Assessment identifies the organisational ability to undertake the stated health-service responsibility. The capability describes what the enterprise must be able to do; it does not assert that Harmonia performs or owns that activity.

\textbf{Boundary:} Harmonia relevance may arise because this activity establishes the context in which information is created, needed, interpreted, related, exchanged or presented. Responsibility for the underlying business decision or care activity remains with the appropriate actor or organisation.

\subsection{CT3 - Population Health Surveillance Management}\label{cap:bc-afc31f01}
\vspace{-0.5em}
\begin{table}[htbp]
  \centering
  \footnotesize
  \begin{tabularx}{1\textwidth}{|l|l|X|}
    \hline
    \multicolumn{3}{c}{\textbf{Capability Metadata}} \\
    \hline
    \textbf{Name} & \textbf{Value} & \textbf{Comment and/or Description} \\
    \hline
    \textbf{Reference Id} & \texttt{BC-AFC31F01} & \\
    \hline
    \textbf{Status} & Established working hierarchy & \\
    \hline
    \textbf{Ancestry} &  & \\
    \hline
    \textbf{Source} &  & \\
    \hline
    \textbf{Enabled By} &  & \\
    \hline
    \textbf{Supports} &  & \\
    \hline
    \textbf{Derived From} &  & \\
    \hline
    \textbf{Reviewed Status} & Review meaning, business boundary, information context and any unresolved ancestry. & \\
    \hline
  \end{tabularx}
\end{table}

\begin{table}[htbp]
  \centering
  \footnotesize
  \begin{tabularx}{1\textwidth}{|l|l|X|}
    \hline
    \multicolumn{3}{c}{\textbf{Contained Features}} \\
    \hline
    \textbf{Feature Name} & \textbf{Feature Id} & \textbf{Comment} \\
    \hline
    \multicolumn{3}{|l|}{\emph{No contained Features are currently identified.}} \\
    \hline
  \end{tabularx}
\end{table}

\textbf{Intent:} Population Health Surveillance Management identifies the organisational ability to undertake the stated health-service responsibility. The capability describes what the enterprise must be able to do; it does not assert that Harmonia performs or owns that activity.

\textbf{Boundary:} Harmonia relevance may arise because this activity establishes the context in which information is created, needed, interpreted, related, exchanged or presented. Responsibility for the underlying business decision or care activity remains with the appropriate actor or organisation.

\subsection{CT3 - Population Health Event Identification}\label{cap:bc-bca30d47}
\vspace{-0.5em}
\begin{table}[htbp]
  \centering
  \footnotesize
  \begin{tabularx}{1\textwidth}{|l|l|X|}
    \hline
    \multicolumn{3}{c}{\textbf{Capability Metadata}} \\
    \hline
    \textbf{Name} & \textbf{Value} & \textbf{Comment and/or Description} \\
    \hline
    \textbf{Reference Id} & \texttt{BC-BCA30D47} & \\
    \hline
    \textbf{Status} & Established working hierarchy & \\
    \hline
    \textbf{Ancestry} &  & \\
    \hline
    \textbf{Source} &  & \\
    \hline
    \textbf{Enabled By} &  & \\
    \hline
    \textbf{Supports} &  & \\
    \hline
    \textbf{Derived From} &  & \\
    \hline
    \textbf{Reviewed Status} & Review meaning, business boundary, information context and any unresolved ancestry. & \\
    \hline
  \end{tabularx}
\end{table}

\begin{table}[htbp]
  \centering
  \footnotesize
  \begin{tabularx}{1\textwidth}{|l|l|X|}
    \hline
    \multicolumn{3}{c}{\textbf{Contained Features}} \\
    \hline
    \textbf{Feature Name} & \textbf{Feature Id} & \textbf{Comment} \\
    \hline
    \multicolumn{3}{|l|}{\emph{No contained Features are currently identified.}} \\
    \hline
  \end{tabularx}
\end{table}

\textbf{Intent:} Population Health Event Identification identifies the organisational ability to undertake the stated health-service responsibility. The capability describes what the enterprise must be able to do; it does not assert that Harmonia performs or owns that activity.

\textbf{Boundary:} Harmonia relevance may arise because this activity establishes the context in which information is created, needed, interpreted, related, exchanged or presented. Responsibility for the underlying business decision or care activity remains with the appropriate actor or organisation.

\subsection{CT3 - Notifiable Health Event Management}\label{cap:bc-5eaa6955}
\vspace{-0.5em}
\begin{table}[htbp]
  \centering
  \footnotesize
  \begin{tabularx}{1\textwidth}{|l|l|X|}
    \hline
    \multicolumn{3}{c}{\textbf{Capability Metadata}} \\
    \hline
    \textbf{Name} & \textbf{Value} & \textbf{Comment and/or Description} \\
    \hline
    \textbf{Reference Id} & \texttt{BC-5EAA6955} & \\
    \hline
    \textbf{Status} & Established working hierarchy & \\
    \hline
    \textbf{Ancestry} &  & \\
    \hline
    \textbf{Source} &  & \\
    \hline
    \textbf{Enabled By} &  & \\
    \hline
    \textbf{Supports} &  & \\
    \hline
    \textbf{Derived From} &  & \\
    \hline
    \textbf{Reviewed Status} & Review meaning, business boundary, information context and any unresolved ancestry. & \\
    \hline
  \end{tabularx}
\end{table}

\begin{table}[htbp]
  \centering
  \footnotesize
  \begin{tabularx}{1\textwidth}{|l|l|X|}
    \hline
    \multicolumn{3}{c}{\textbf{Contained Features}} \\
    \hline
    \textbf{Feature Name} & \textbf{Feature Id} & \textbf{Comment} \\
    \hline
    \multicolumn{3}{|l|}{\emph{No contained Features are currently identified.}} \\
    \hline
  \end{tabularx}
\end{table}

\textbf{Intent:} Notifiable Health Event Management identifies the organisational ability to undertake the stated health-service responsibility. The capability describes what the enterprise must be able to do; it does not assert that Harmonia performs or owns that activity.

\textbf{Boundary:} Harmonia relevance may arise because this activity establishes the context in which information is created, needed, interpreted, related, exchanged or presented. Responsibility for the underlying business decision or care activity remains with the appropriate actor or organisation.

\subsection{CT3 - Population Health Reporting}\label{cap:bc-d8c15299}
\vspace{-0.5em}
\begin{table}[htbp]
  \centering
  \footnotesize
  \begin{tabularx}{1\textwidth}{|l|l|X|}
    \hline
    \multicolumn{3}{c}{\textbf{Capability Metadata}} \\
    \hline
    \textbf{Name} & \textbf{Value} & \textbf{Comment and/or Description} \\
    \hline
    \textbf{Reference Id} & \texttt{BC-D8C15299} & \\
    \hline
    \textbf{Status} & Established working hierarchy & \\
    \hline
    \textbf{Ancestry} &  & \\
    \hline
    \textbf{Source} &  & \\
    \hline
    \textbf{Enabled By} &  & \\
    \hline
    \textbf{Supports} &  & \\
    \hline
    \textbf{Derived From} &  & \\
    \hline
    \textbf{Reviewed Status} & Review meaning, business boundary, information context and any unresolved ancestry. & \\
    \hline
  \end{tabularx}
\end{table}

\begin{table}[htbp]
  \centering
  \footnotesize
  \begin{tabularx}{1\textwidth}{|l|l|X|}
    \hline
    \multicolumn{3}{c}{\textbf{Contained Features}} \\
    \hline
    \textbf{Feature Name} & \textbf{Feature Id} & \textbf{Comment} \\
    \hline
    \multicolumn{3}{|l|}{\emph{No contained Features are currently identified.}} \\
    \hline
  \end{tabularx}
\end{table}

\textbf{Intent:} Population Health Reporting identifies the organisational ability to undertake the stated health-service responsibility. The capability describes what the enterprise must be able to do; it does not assert that Harmonia performs or owns that activity.

\textbf{Boundary:} Harmonia relevance may arise because this activity establishes the context in which information is created, needed, interpreted, related, exchanged or presented. Responsibility for the underlying business decision or care activity remains with the appropriate actor or organisation.

\section{CT2 - Health-System Planning}\label{cap:bc-40b6a2c7}
\vspace{-0.5em}
\begin{table}[htbp]
  \centering
  \footnotesize
  \begin{tabularx}{1\textwidth}{|l|l|X|}
    \hline
    \multicolumn{3}{c}{\textbf{Capability Metadata}} \\
    \hline
    \textbf{Name} & \textbf{Value} & \textbf{Comment and/or Description} \\
    \hline
    \textbf{Reference Id} & \texttt{BC-40B6A2C7} & \\
    \hline
    \textbf{Status} & Established working hierarchy & \\
    \hline
    \textbf{Ancestry} &  & \\
    \hline
    \textbf{Source} &  & \\
    \hline
    \textbf{Enabled By} &  & \\
    \hline
    \textbf{Supports} &  & \\
    \hline
    \textbf{Derived From} &  & \\
    \hline
    \textbf{Reviewed Status} &  & \\
    \hline
  \end{tabularx}
\end{table}

\begin{table}[htbp]
  \centering
  \footnotesize
  \begin{tabularx}{1\textwidth}{|l|l|X|}
    \hline
    \multicolumn{3}{c}{\textbf{Contained Capabilities}} \\
    \hline
    \textbf{Capability Name} & \textbf{Capability Id} & \textbf{Comment} \\
    \hline
    Health Service Demand Assessment & \texttt{cap:bc-770f60dc} (at \ref{cap:bc-770f60dc}) & \\
    \hline
    Health Service Capacity Planning & \texttt{cap:bc-3742724c} (at \ref{cap:bc-3742724c}) & \\
    \hline
    Health Service Distribution Planning & \texttt{cap:bc-33a5de96} (at \ref{cap:bc-33a5de96}) & \\
    \hline
  \end{tabularx}
\end{table}

\textbf{Intent:} Health-System Planning identifies the organisational ability to undertake the stated health-service responsibility. The capability describes what the enterprise must be able to do; it does not assert that Harmonia performs or owns that activity.

\textbf{Boundary:} Harmonia relevance may arise because this activity establishes the context in which information is created, needed, interpreted, related, exchanged or presented. Responsibility for the underlying business decision or care activity remains with the appropriate actor or organisation.

\subsection{CT3 - Health Service Demand Assessment}\label{cap:bc-770f60dc}
\vspace{-0.5em}
\begin{table}[htbp]
  \centering
  \footnotesize
  \begin{tabularx}{1\textwidth}{|l|l|X|}
    \hline
    \multicolumn{3}{c}{\textbf{Capability Metadata}} \\
    \hline
    \textbf{Name} & \textbf{Value} & \textbf{Comment and/or Description} \\
    \hline
    \textbf{Reference Id} & \texttt{BC-770F60DC} & \\
    \hline
    \textbf{Status} & Established working hierarchy & \\
    \hline
    \textbf{Ancestry} &  & \\
    \hline
    \textbf{Source} &  & \\
    \hline
    \textbf{Enabled By} &  & \\
    \hline
    \textbf{Supports} &  & \\
    \hline
    \textbf{Derived From} &  & \\
    \hline
    \textbf{Reviewed Status} & Review meaning, business boundary, information context and any unresolved ancestry. & \\
    \hline
  \end{tabularx}
\end{table}

\begin{table}[htbp]
  \centering
  \footnotesize
  \begin{tabularx}{1\textwidth}{|l|l|X|}
    \hline
    \multicolumn{3}{c}{\textbf{Contained Features}} \\
    \hline
    \textbf{Feature Name} & \textbf{Feature Id} & \textbf{Comment} \\
    \hline
    \multicolumn{3}{|l|}{\emph{No contained Features are currently identified.}} \\
    \hline
  \end{tabularx}
\end{table}

\textbf{Intent:} Health Service Demand Assessment identifies the organisational ability to undertake the stated health-service responsibility. The capability describes what the enterprise must be able to do; it does not assert that Harmonia performs or owns that activity.

\textbf{Boundary:} Harmonia relevance may arise because this activity establishes the context in which information is created, needed, interpreted, related, exchanged or presented. Responsibility for the underlying business decision or care activity remains with the appropriate actor or organisation.

\subsection{CT3 - Health Service Capacity Planning}\label{cap:bc-3742724c}
\vspace{-0.5em}
\begin{table}[htbp]
  \centering
  \footnotesize
  \begin{tabularx}{1\textwidth}{|l|l|X|}
    \hline
    \multicolumn{3}{c}{\textbf{Capability Metadata}} \\
    \hline
    \textbf{Name} & \textbf{Value} & \textbf{Comment and/or Description} \\
    \hline
    \textbf{Reference Id} & \texttt{BC-3742724C} & \\
    \hline
    \textbf{Status} & Established working hierarchy & \\
    \hline
    \textbf{Ancestry} &  & \\
    \hline
    \textbf{Source} &  & \\
    \hline
    \textbf{Enabled By} &  & \\
    \hline
    \textbf{Supports} &  & \\
    \hline
    \textbf{Derived From} &  & \\
    \hline
    \textbf{Reviewed Status} & Review meaning, business boundary, information context and any unresolved ancestry. & \\
    \hline
  \end{tabularx}
\end{table}

\begin{table}[htbp]
  \centering
  \footnotesize
  \begin{tabularx}{1\textwidth}{|l|l|X|}
    \hline
    \multicolumn{3}{c}{\textbf{Contained Features}} \\
    \hline
    \textbf{Feature Name} & \textbf{Feature Id} & \textbf{Comment} \\
    \hline
    \multicolumn{3}{|l|}{\emph{No contained Features are currently identified.}} \\
    \hline
  \end{tabularx}
\end{table}

\textbf{Intent:} Health Service Capacity Planning identifies the organisational ability to undertake the stated health-service responsibility. The capability describes what the enterprise must be able to do; it does not assert that Harmonia performs or owns that activity.

\textbf{Boundary:} Harmonia relevance may arise because this activity establishes the context in which information is created, needed, interpreted, related, exchanged or presented. Responsibility for the underlying business decision or care activity remains with the appropriate actor or organisation.

\subsection{CT3 - Health Service Distribution Planning}\label{cap:bc-33a5de96}
\vspace{-0.5em}
\begin{table}[htbp]
  \centering
  \footnotesize
  \begin{tabularx}{1\textwidth}{|l|l|X|}
    \hline
    \multicolumn{3}{c}{\textbf{Capability Metadata}} \\
    \hline
    \textbf{Name} & \textbf{Value} & \textbf{Comment and/or Description} \\
    \hline
    \textbf{Reference Id} & \texttt{BC-33A5DE96} & \\
    \hline
    \textbf{Status} & Established working hierarchy & \\
    \hline
    \textbf{Ancestry} &  & \\
    \hline
    \textbf{Source} &  & \\
    \hline
    \textbf{Enabled By} &  & \\
    \hline
    \textbf{Supports} &  & \\
    \hline
    \textbf{Derived From} &  & \\
    \hline
    \textbf{Reviewed Status} & Review meaning, business boundary, information context and any unresolved ancestry. & \\
    \hline
  \end{tabularx}
\end{table}

\begin{table}[htbp]
  \centering
  \footnotesize
  \begin{tabularx}{1\textwidth}{|l|l|X|}
    \hline
    \multicolumn{3}{c}{\textbf{Contained Features}} \\
    \hline
    \textbf{Feature Name} & \textbf{Feature Id} & \textbf{Comment} \\
    \hline
    \multicolumn{3}{|l|}{\emph{No contained Features are currently identified.}} \\
    \hline
  \end{tabularx}
\end{table}

\textbf{Intent:} Health Service Distribution Planning identifies the organisational ability to undertake the stated health-service responsibility. The capability describes what the enterprise must be able to do; it does not assert that Harmonia performs or owns that activity.

\textbf{Boundary:} Harmonia relevance may arise because this activity establishes the context in which information is created, needed, interpreted, related, exchanged or presented. Responsibility for the underlying business decision or care activity remains with the appropriate actor or organisation.

\chapter{CT1 - Health Information \& Knowledge Management}\label{cap:bc-d37dbd07}
\noindent\textbf{Reference ID:} \texttt{BC-D37DBD07}\quad \textbf{Status:} Established working hierarchy\par
\noindent\textit{Source: Canonical Business Capability model / working map}\par
\textbf{Intent:} Health Information \& Knowledge Management identifies the organisational ability to undertake the stated health-service responsibility. The capability describes what the enterprise must be able to do; it does not assert that Harmonia performs or owns that activity.

\textbf{Boundary:} Harmonia relevance may arise because this activity establishes the context in which information is created, needed, interpreted, related, exchanged or presented. Responsibility for the underlying business decision or care activity remains with the appropriate actor or organisation.

\reviewstatus{Check scope and information relevance; no business authority transfers to Harmonia.}
\section{CT2 - Health Information Management}\label{cap:bc-e21b0f21}
\vspace{-0.5em}
\begin{table}[htbp]
  \centering
  \footnotesize
  \begin{tabularx}{1\textwidth}{|l|l|X|}
    \hline
    \multicolumn{3}{c}{\textbf{Capability Metadata}} \\
    \hline
    \textbf{Name} & \textbf{Value} & \textbf{Comment and/or Description} \\
    \hline
    \textbf{Reference Id} & \texttt{BC-E21B0F21} & \\
    \hline
    \textbf{Status} & Established working hierarchy & \\
    \hline
    \textbf{Ancestry} &  & \\
    \hline
    \textbf{Source} &  & \\
    \hline
    \textbf{Enabled By} &  & \\
    \hline
    \textbf{Supports} &  & \\
    \hline
    \textbf{Derived From} &  & \\
    \hline
    \textbf{Reviewed Status} &  & \\
    \hline
  \end{tabularx}
\end{table}

\begin{table}[htbp]
  \centering
  \footnotesize
  \begin{tabularx}{1\textwidth}{|l|l|X|}
    \hline
    \multicolumn{3}{c}{\textbf{Contained Capabilities}} \\
    \hline
    \textbf{Capability Name} & \textbf{Capability Id} & \textbf{Comment} \\
    \hline
    Health Information Data Management & \texttt{cap:bc-05a55e6b} (at \ref{cap:bc-05a55e6b}) & \\
    \hline
    Health Information Structure Management & \texttt{cap:bc-822903b0} (at \ref{cap:bc-822903b0}) & \\
    \hline
    Health Information Lifecycle Management & \texttt{cap:bc-65531f2e} (at \ref{cap:bc-65531f2e}) & \\
    \hline
    Health Information Catalogue and Index Management & \texttt{cap:bc-2c668f82} (at \ref{cap:bc-2c668f82}) & \\
    \hline
    Health Information Compartmentalisation & \texttt{cap:bc-ad69e2a8} (at \ref{cap:bc-ad69e2a8}) & \\
    \hline
    Health Information Data Correlation & \texttt{cap:bc-39f5cc28} (at \ref{cap:bc-39f5cc28}) & \\
    \hline
    Health Information Data Retrieval & \texttt{cap:bc-c6e88499} (at \ref{cap:bc-c6e88499}) & \\
    \hline
  \end{tabularx}
\end{table}

\textbf{Intent:} Health Information Management identifies the organisational ability to undertake the stated health-service responsibility. The capability describes what the enterprise must be able to do; it does not assert that Harmonia performs or owns that activity.

\textbf{Boundary:} Harmonia relevance may arise because this activity establishes the context in which information is created, needed, interpreted, related, exchanged or presented. Responsibility for the underlying business decision or care activity remains with the appropriate actor or organisation.

\subsection{CT3 - Health Information Data Management}\label{cap:bc-05a55e6b}
\vspace{-0.5em}
\begin{table}[htbp]
  \centering
  \footnotesize
  \begin{tabularx}{1\textwidth}{|l|l|X|}
    \hline
    \multicolumn{3}{c}{\textbf{Capability Metadata}} \\
    \hline
    \textbf{Name} & \textbf{Value} & \textbf{Comment and/or Description} \\
    \hline
    \textbf{Reference Id} & \texttt{BC-05A55E6B} & \\
    \hline
    \textbf{Status} & Established working hierarchy & \\
    \hline
    \textbf{Ancestry} &  & \\
    \hline
    \textbf{Source} &  & \\
    \hline
    \textbf{Enabled By} &  & \\
    \hline
    \textbf{Supports} &  & \\
    \hline
    \textbf{Derived From} &  & \\
    \hline
    \textbf{Reviewed Status} & Review meaning, business boundary, information context and any unresolved ancestry. & \\
    \hline
  \end{tabularx}
\end{table}

\begin{table}[htbp]
  \centering
  \footnotesize
  \begin{tabularx}{1\textwidth}{|l|l|X|}
    \hline
    \multicolumn{3}{c}{\textbf{Contained Features}} \\
    \hline
    \textbf{Feature Name} & \textbf{Feature Id} & \textbf{Comment} \\
    \hline
    \multicolumn{3}{|l|}{\emph{No contained Features are currently identified.}} \\
    \hline
  \end{tabularx}
\end{table}

\textbf{Intent:} Health Information Data Management identifies the organisational ability to undertake the stated health-service responsibility. The capability describes what the enterprise must be able to do; it does not assert that Harmonia performs or owns that activity.

\textbf{Boundary:} Harmonia relevance may arise because this activity establishes the context in which information is created, needed, interpreted, related, exchanged or presented. Responsibility for the underlying business decision or care activity remains with the appropriate actor or organisation.

\subsection{CT3 - Health Information Structure Management}\label{cap:bc-822903b0}
\vspace{-0.5em}
\begin{table}[htbp]
  \centering
  \footnotesize
  \begin{tabularx}{1\textwidth}{|l|l|X|}
    \hline
    \multicolumn{3}{c}{\textbf{Capability Metadata}} \\
    \hline
    \textbf{Name} & \textbf{Value} & \textbf{Comment and/or Description} \\
    \hline
    \textbf{Reference Id} & \texttt{BC-822903B0} & \\
    \hline
    \textbf{Status} & Established working hierarchy & \\
    \hline
    \textbf{Ancestry} &  & \\
    \hline
    \textbf{Source} &  & \\
    \hline
    \textbf{Enabled By} &  & \\
    \hline
    \textbf{Supports} &  & \\
    \hline
    \textbf{Derived From} &  & \\
    \hline
    \textbf{Reviewed Status} & Review meaning, business boundary, information context and any unresolved ancestry. & \\
    \hline
  \end{tabularx}
\end{table}

\begin{table}[htbp]
  \centering
  \footnotesize
  \begin{tabularx}{1\textwidth}{|l|l|X|}
    \hline
    \multicolumn{3}{c}{\textbf{Contained Features}} \\
    \hline
    \textbf{Feature Name} & \textbf{Feature Id} & \textbf{Comment} \\
    \hline
    \multicolumn{3}{|l|}{\emph{No contained Features are currently identified.}} \\
    \hline
  \end{tabularx}
\end{table}

\textbf{Intent:} Health Information Structure Management identifies the organisational ability to undertake the stated health-service responsibility. The capability describes what the enterprise must be able to do; it does not assert that Harmonia performs or owns that activity.

\textbf{Boundary:} Harmonia relevance may arise because this activity establishes the context in which information is created, needed, interpreted, related, exchanged or presented. Responsibility for the underlying business decision or care activity remains with the appropriate actor or organisation.

\subsection{CT3 - Health Information Lifecycle Management}\label{cap:bc-65531f2e}
\vspace{-0.5em}
\begin{table}[htbp]
  \centering
  \footnotesize
  \begin{tabularx}{1\textwidth}{|l|l|X|}
    \hline
    \multicolumn{3}{c}{\textbf{Capability Metadata}} \\
    \hline
    \textbf{Name} & \textbf{Value} & \textbf{Comment and/or Description} \\
    \hline
    \textbf{Reference Id} & \texttt{BC-65531F2E} & \\
    \hline
    \textbf{Status} & Established working hierarchy & \\
    \hline
    \textbf{Ancestry} &  & \\
    \hline
    \textbf{Source} &  & \\
    \hline
    \textbf{Enabled By} &  & \\
    \hline
    \textbf{Supports} &  & \\
    \hline
    \textbf{Derived From} &  & \\
    \hline
    \textbf{Reviewed Status} & Review meaning, business boundary, information context and any unresolved ancestry. & \\
    \hline
  \end{tabularx}
\end{table}

\begin{table}[htbp]
  \centering
  \footnotesize
  \begin{tabularx}{1\textwidth}{|l|l|X|}
    \hline
    \multicolumn{3}{c}{\textbf{Contained Features}} \\
    \hline
    \textbf{Feature Name} & \textbf{Feature Id} & \textbf{Comment} \\
    \hline
    \multicolumn{3}{|l|}{\emph{No contained Features are currently identified.}} \\
    \hline
  \end{tabularx}
\end{table}

\textbf{Intent:} Health Information Lifecycle Management identifies the organisational ability to undertake the stated health-service responsibility. The capability describes what the enterprise must be able to do; it does not assert that Harmonia performs or owns that activity.

\textbf{Boundary:} Harmonia relevance may arise because this activity establishes the context in which information is created, needed, interpreted, related, exchanged or presented. Responsibility for the underlying business decision or care activity remains with the appropriate actor or organisation.

\subsection{CT3 - Health Information Catalogue and Index Management}\label{cap:bc-2c668f82}
\vspace{-0.5em}
\begin{table}[htbp]
  \centering
  \footnotesize
  \begin{tabularx}{1\textwidth}{|l|l|X|}
    \hline
    \multicolumn{3}{c}{\textbf{Capability Metadata}} \\
    \hline
    \textbf{Name} & \textbf{Value} & \textbf{Comment and/or Description} \\
    \hline
    \textbf{Reference Id} & \texttt{BC-2C668F82} & \\
    \hline
    \textbf{Status} & Established working hierarchy & \\
    \hline
    \textbf{Ancestry} &  & \\
    \hline
    \textbf{Source} &  & \\
    \hline
    \textbf{Enabled By} &  & \\
    \hline
    \textbf{Supports} &  & \\
    \hline
    \textbf{Derived From} &  & \\
    \hline
    \textbf{Reviewed Status} & Review meaning, business boundary, information context and any unresolved ancestry. & \\
    \hline
  \end{tabularx}
\end{table}

\begin{table}[htbp]
  \centering
  \footnotesize
  \begin{tabularx}{1\textwidth}{|l|l|X|}
    \hline
    \multicolumn{3}{c}{\textbf{Contained Features}} \\
    \hline
    \textbf{Feature Name} & \textbf{Feature Id} & \textbf{Comment} \\
    \hline
    \multicolumn{3}{|l|}{\emph{No contained Features are currently identified.}} \\
    \hline
  \end{tabularx}
\end{table}

\textbf{Intent:} Health Information Catalogue and Index Management identifies the organisational ability to undertake the stated health-service responsibility. The capability describes what the enterprise must be able to do; it does not assert that Harmonia performs or owns that activity.

\textbf{Boundary:} Harmonia relevance may arise because this activity establishes the context in which information is created, needed, interpreted, related, exchanged or presented. Responsibility for the underlying business decision or care activity remains with the appropriate actor or organisation.

\subsection{CT3 - Health Information Compartmentalisation}\label{cap:bc-ad69e2a8}
\vspace{-0.5em}
\begin{table}[htbp]
  \centering
  \footnotesize
  \begin{tabularx}{1\textwidth}{|l|l|X|}
    \hline
    \multicolumn{3}{c}{\textbf{Capability Metadata}} \\
    \hline
    \textbf{Name} & \textbf{Value} & \textbf{Comment and/or Description} \\
    \hline
    \textbf{Reference Id} & \texttt{BC-AD69E2A8} & \\
    \hline
    \textbf{Status} & Established working hierarchy & \\
    \hline
    \textbf{Ancestry} &  & \\
    \hline
    \textbf{Source} &  & \\
    \hline
    \textbf{Enabled By} &  & \\
    \hline
    \textbf{Supports} &  & \\
    \hline
    \textbf{Derived From} &  & \\
    \hline
    \textbf{Reviewed Status} & Review meaning, business boundary, information context and any unresolved ancestry. & \\
    \hline
  \end{tabularx}
\end{table}

\begin{table}[htbp]
  \centering
  \footnotesize
  \begin{tabularx}{1\textwidth}{|l|l|X|}
    \hline
    \multicolumn{3}{c}{\textbf{Contained Features}} \\
    \hline
    \textbf{Feature Name} & \textbf{Feature Id} & \textbf{Comment} \\
    \hline
    \multicolumn{3}{|l|}{\emph{No contained Features are currently identified.}} \\
    \hline
  \end{tabularx}
\end{table}

\textbf{Intent:} Health Information Compartmentalisation identifies the organisational ability to undertake the stated health-service responsibility. The capability describes what the enterprise must be able to do; it does not assert that Harmonia performs or owns that activity.

\textbf{Boundary:} Harmonia relevance may arise because this activity establishes the context in which information is created, needed, interpreted, related, exchanged or presented. Responsibility for the underlying business decision or care activity remains with the appropriate actor or organisation.

\subsection{CT3 - Health Information Data Correlation}\label{cap:bc-39f5cc28}
\vspace{-0.5em}
\begin{table}[htbp]
  \centering
  \footnotesize
  \begin{tabularx}{1\textwidth}{|l|l|X|}
    \hline
    \multicolumn{3}{c}{\textbf{Capability Metadata}} \\
    \hline
    \textbf{Name} & \textbf{Value} & \textbf{Comment and/or Description} \\
    \hline
    \textbf{Reference Id} & \texttt{BC-39F5CC28} & \\
    \hline
    \textbf{Status} & Established working hierarchy & \\
    \hline
    \textbf{Ancestry} &  & \\
    \hline
    \textbf{Source} &  & \\
    \hline
    \textbf{Enabled By} &  & \\
    \hline
    \textbf{Supports} &  & \\
    \hline
    \textbf{Derived From} &  & \\
    \hline
    \textbf{Reviewed Status} & Review meaning, business boundary, information context and any unresolved ancestry. & \\
    \hline
  \end{tabularx}
\end{table}

\begin{table}[htbp]
  \centering
  \footnotesize
  \begin{tabularx}{1\textwidth}{|l|l|X|}
    \hline
    \multicolumn{3}{c}{\textbf{Contained Features}} \\
    \hline
    \textbf{Feature Name} & \textbf{Feature Id} & \textbf{Comment} \\
    \hline
    \multicolumn{3}{|l|}{\emph{No contained Features are currently identified.}} \\
    \hline
  \end{tabularx}
\end{table}

\textbf{Intent:} Health Information Data Correlation identifies the organisational ability to undertake the stated health-service responsibility. The capability describes what the enterprise must be able to do; it does not assert that Harmonia performs or owns that activity.

\textbf{Boundary:} Harmonia relevance may arise because this activity establishes the context in which information is created, needed, interpreted, related, exchanged or presented. Responsibility for the underlying business decision or care activity remains with the appropriate actor or organisation.

\subsection{CT3 - Health Information Data Retrieval}\label{cap:bc-c6e88499}
\vspace{-0.5em}
\begin{table}[htbp]
  \centering
  \footnotesize
  \begin{tabularx}{1\textwidth}{|l|l|X|}
    \hline
    \multicolumn{3}{c}{\textbf{Capability Metadata}} \\
    \hline
    \textbf{Name} & \textbf{Value} & \textbf{Comment and/or Description} \\
    \hline
    \textbf{Reference Id} & \texttt{BC-C6E88499} & \\
    \hline
    \textbf{Status} & Established working hierarchy & \\
    \hline
    \textbf{Ancestry} &  & \\
    \hline
    \textbf{Source} &  & \\
    \hline
    \textbf{Enabled By} &  & \\
    \hline
    \textbf{Supports} &  & \\
    \hline
    \textbf{Derived From} &  & \\
    \hline
    \textbf{Reviewed Status} & Review meaning, business boundary, information context and any unresolved ancestry. & \\
    \hline
  \end{tabularx}
\end{table}

\begin{table}[htbp]
  \centering
  \footnotesize
  \begin{tabularx}{1\textwidth}{|l|l|X|}
    \hline
    \multicolumn{3}{c}{\textbf{Contained Features}} \\
    \hline
    \textbf{Feature Name} & \textbf{Feature Id} & \textbf{Comment} \\
    \hline
    \multicolumn{3}{|l|}{\emph{No contained Features are currently identified.}} \\
    \hline
  \end{tabularx}
\end{table}

\textbf{Intent:} Health Information Data Retrieval identifies the organisational ability to undertake the stated health-service responsibility. The capability describes what the enterprise must be able to do; it does not assert that Harmonia performs or owns that activity.

\textbf{Boundary:} Harmonia relevance may arise because this activity establishes the context in which information is created, needed, interpreted, related, exchanged or presented. Responsibility for the underlying business decision or care activity remains with the appropriate actor or organisation.

\section{CT2 - Clinical Knowledge Management}\label{cap:bc-55c02bd4}
\vspace{-0.5em}
\begin{table}[htbp]
  \centering
  \footnotesize
  \begin{tabularx}{1\textwidth}{|l|l|X|}
    \hline
    \multicolumn{3}{c}{\textbf{Capability Metadata}} \\
    \hline
    \textbf{Name} & \textbf{Value} & \textbf{Comment and/or Description} \\
    \hline
    \textbf{Reference Id} & \texttt{BC-55C02BD4} & \\
    \hline
    \textbf{Status} & Established working hierarchy & \\
    \hline
    \textbf{Ancestry} &  & \\
    \hline
    \textbf{Source} &  & \\
    \hline
    \textbf{Enabled By} &  & \\
    \hline
    \textbf{Supports} &  & \\
    \hline
    \textbf{Derived From} &  & \\
    \hline
    \textbf{Reviewed Status} &  & \\
    \hline
  \end{tabularx}
\end{table}

\begin{table}[htbp]
  \centering
  \footnotesize
  \begin{tabularx}{1\textwidth}{|l|l|X|}
    \hline
    \multicolumn{3}{c}{\textbf{Contained Capabilities}} \\
    \hline
    \textbf{Capability Name} & \textbf{Capability Id} & \textbf{Comment} \\
    \hline
    Clinical Knowledge Context Identification & \texttt{cap:bc-16799a80} (at \ref{cap:bc-16799a80}) & \\
    \hline
    Clinical Knowledge Context Extraction & \texttt{cap:bc-d45cd0aa} (at \ref{cap:bc-d45cd0aa}) & \\
    \hline
    Clinical Knowledge Information Formulation & \texttt{cap:bc-385fcdae} (at \ref{cap:bc-385fcdae}) & \\
    \hline
    Clinical Knowledge Concept Association & \texttt{cap:bc-ad456750} (at \ref{cap:bc-ad456750}) & \\
    \hline
    Clinical Knowledge Review \& Approval Management & \texttt{cap:bc-88d82704} (at \ref{cap:bc-88d82704}) & \\
    \hline
    Clinical Knowledge Publication \& Sharing & \texttt{cap:bc-e57a9c37} (at \ref{cap:bc-e57a9c37}) & \\
    \hline
  \end{tabularx}
\end{table}

\textbf{Intent:} Clinical Knowledge Management identifies the organisational ability to undertake the stated health-service responsibility. The capability describes what the enterprise must be able to do; it does not assert that Harmonia performs or owns that activity.

\textbf{Boundary:} Harmonia relevance may arise because this activity establishes the context in which information is created, needed, interpreted, related, exchanged or presented. Responsibility for the underlying business decision or care activity remains with the appropriate actor or organisation.

\subsection{CT3 - Clinical Knowledge Context Identification}\label{cap:bc-16799a80}
\vspace{-0.5em}
\begin{table}[htbp]
  \centering
  \footnotesize
  \begin{tabularx}{1\textwidth}{|l|l|X|}
    \hline
    \multicolumn{3}{c}{\textbf{Capability Metadata}} \\
    \hline
    \textbf{Name} & \textbf{Value} & \textbf{Comment and/or Description} \\
    \hline
    \textbf{Reference Id} & \texttt{BC-16799A80} & \\
    \hline
    \textbf{Status} & Established working hierarchy & \\
    \hline
    \textbf{Ancestry} &  & \\
    \hline
    \textbf{Source} &  & \\
    \hline
    \textbf{Enabled By} &  & \\
    \hline
    \textbf{Supports} &  & \\
    \hline
    \textbf{Derived From} &  & \\
    \hline
    \textbf{Reviewed Status} & Review meaning, business boundary, information context and any unresolved ancestry. & \\
    \hline
  \end{tabularx}
\end{table}

\begin{table}[htbp]
  \centering
  \footnotesize
  \begin{tabularx}{1\textwidth}{|l|l|X|}
    \hline
    \multicolumn{3}{c}{\textbf{Contained Features}} \\
    \hline
    \textbf{Feature Name} & \textbf{Feature Id} & \textbf{Comment} \\
    \hline
    \multicolumn{3}{|l|}{\emph{No contained Features are currently identified.}} \\
    \hline
  \end{tabularx}
\end{table}

\textbf{Intent:} Clinical Knowledge Context Identification identifies the organisational ability to undertake the stated health-service responsibility. The capability describes what the enterprise must be able to do; it does not assert that Harmonia performs or owns that activity.

\textbf{Boundary:} Harmonia relevance may arise because this activity establishes the context in which information is created, needed, interpreted, related, exchanged or presented. Responsibility for the underlying business decision or care activity remains with the appropriate actor or organisation.

\subsection{CT3 - Clinical Knowledge Context Extraction}\label{cap:bc-d45cd0aa}
\vspace{-0.5em}
\begin{table}[htbp]
  \centering
  \footnotesize
  \begin{tabularx}{1\textwidth}{|l|l|X|}
    \hline
    \multicolumn{3}{c}{\textbf{Capability Metadata}} \\
    \hline
    \textbf{Name} & \textbf{Value} & \textbf{Comment and/or Description} \\
    \hline
    \textbf{Reference Id} & \texttt{BC-D45CD0AA} & \\
    \hline
    \textbf{Status} & Established working hierarchy & \\
    \hline
    \textbf{Ancestry} &  & \\
    \hline
    \textbf{Source} &  & \\
    \hline
    \textbf{Enabled By} &  & \\
    \hline
    \textbf{Supports} &  & \\
    \hline
    \textbf{Derived From} &  & \\
    \hline
    \textbf{Reviewed Status} & Review meaning, business boundary, information context and any unresolved ancestry. & \\
    \hline
  \end{tabularx}
\end{table}

\begin{table}[htbp]
  \centering
  \footnotesize
  \begin{tabularx}{1\textwidth}{|l|l|X|}
    \hline
    \multicolumn{3}{c}{\textbf{Contained Features}} \\
    \hline
    \textbf{Feature Name} & \textbf{Feature Id} & \textbf{Comment} \\
    \hline
    \multicolumn{3}{|l|}{\emph{No contained Features are currently identified.}} \\
    \hline
  \end{tabularx}
\end{table}

\textbf{Intent:} Clinical Knowledge Context Extraction identifies the organisational ability to undertake the stated health-service responsibility. The capability describes what the enterprise must be able to do; it does not assert that Harmonia performs or owns that activity.

\textbf{Boundary:} Harmonia relevance may arise because this activity establishes the context in which information is created, needed, interpreted, related, exchanged or presented. Responsibility for the underlying business decision or care activity remains with the appropriate actor or organisation.

\subsection{CT3 - Clinical Knowledge Information Formulation}\label{cap:bc-385fcdae}
\vspace{-0.5em}
\begin{table}[htbp]
  \centering
  \footnotesize
  \begin{tabularx}{1\textwidth}{|l|l|X|}
    \hline
    \multicolumn{3}{c}{\textbf{Capability Metadata}} \\
    \hline
    \textbf{Name} & \textbf{Value} & \textbf{Comment and/or Description} \\
    \hline
    \textbf{Reference Id} & \texttt{BC-385FCDAE} & \\
    \hline
    \textbf{Status} & Established working hierarchy & \\
    \hline
    \textbf{Ancestry} &  & \\
    \hline
    \textbf{Source} &  & \\
    \hline
    \textbf{Enabled By} &  & \\
    \hline
    \textbf{Supports} &  & \\
    \hline
    \textbf{Derived From} &  & \\
    \hline
    \textbf{Reviewed Status} & Review meaning, business boundary, information context and any unresolved ancestry. & \\
    \hline
  \end{tabularx}
\end{table}

\begin{table}[htbp]
  \centering
  \footnotesize
  \begin{tabularx}{1\textwidth}{|l|l|X|}
    \hline
    \multicolumn{3}{c}{\textbf{Contained Features}} \\
    \hline
    \textbf{Feature Name} & \textbf{Feature Id} & \textbf{Comment} \\
    \hline
    \multicolumn{3}{|l|}{\emph{No contained Features are currently identified.}} \\
    \hline
  \end{tabularx}
\end{table}

\textbf{Intent:} Clinical Knowledge Information Formulation identifies the organisational ability to undertake the stated health-service responsibility. The capability describes what the enterprise must be able to do; it does not assert that Harmonia performs or owns that activity.

\textbf{Boundary:} Harmonia relevance may arise because this activity establishes the context in which information is created, needed, interpreted, related, exchanged or presented. Responsibility for the underlying business decision or care activity remains with the appropriate actor or organisation.

\subsection{CT3 - Clinical Knowledge Concept Association}\label{cap:bc-ad456750}
\vspace{-0.5em}
\begin{table}[htbp]
  \centering
  \footnotesize
  \begin{tabularx}{1\textwidth}{|l|l|X|}
    \hline
    \multicolumn{3}{c}{\textbf{Capability Metadata}} \\
    \hline
    \textbf{Name} & \textbf{Value} & \textbf{Comment and/or Description} \\
    \hline
    \textbf{Reference Id} & \texttt{BC-AD456750} & \\
    \hline
    \textbf{Status} & Established working hierarchy & \\
    \hline
    \textbf{Ancestry} &  & \\
    \hline
    \textbf{Source} &  & \\
    \hline
    \textbf{Enabled By} &  & \\
    \hline
    \textbf{Supports} &  & \\
    \hline
    \textbf{Derived From} &  & \\
    \hline
    \textbf{Reviewed Status} & Review meaning, business boundary, information context and any unresolved ancestry. & \\
    \hline
  \end{tabularx}
\end{table}

\begin{table}[htbp]
  \centering
  \footnotesize
  \begin{tabularx}{1\textwidth}{|l|l|X|}
    \hline
    \multicolumn{3}{c}{\textbf{Contained Features}} \\
    \hline
    \textbf{Feature Name} & \textbf{Feature Id} & \textbf{Comment} \\
    \hline
    \multicolumn{3}{|l|}{\emph{No contained Features are currently identified.}} \\
    \hline
  \end{tabularx}
\end{table}

\textbf{Intent:} Clinical Knowledge Concept Association identifies the organisational ability to undertake the stated health-service responsibility. The capability describes what the enterprise must be able to do; it does not assert that Harmonia performs or owns that activity.

\textbf{Boundary:} Harmonia relevance may arise because this activity establishes the context in which information is created, needed, interpreted, related, exchanged or presented. Responsibility for the underlying business decision or care activity remains with the appropriate actor or organisation.

\subsection{CT3 - Clinical Knowledge Review \& Approval Management}\label{cap:bc-88d82704}
\vspace{-0.5em}
\begin{table}[htbp]
  \centering
  \footnotesize
  \begin{tabularx}{1\textwidth}{|l|l|X|}
    \hline
    \multicolumn{3}{c}{\textbf{Capability Metadata}} \\
    \hline
    \textbf{Name} & \textbf{Value} & \textbf{Comment and/or Description} \\
    \hline
    \textbf{Reference Id} & \texttt{BC-88D82704} & \\
    \hline
    \textbf{Status} & Established working hierarchy & \\
    \hline
    \textbf{Ancestry} &  & \\
    \hline
    \textbf{Source} &  & \\
    \hline
    \textbf{Enabled By} &  & \\
    \hline
    \textbf{Supports} &  & \\
    \hline
    \textbf{Derived From} &  & \\
    \hline
    \textbf{Reviewed Status} & Review meaning, business boundary, information context and any unresolved ancestry. & \\
    \hline
  \end{tabularx}
\end{table}

\begin{table}[htbp]
  \centering
  \footnotesize
  \begin{tabularx}{1\textwidth}{|l|l|X|}
    \hline
    \multicolumn{3}{c}{\textbf{Contained Features}} \\
    \hline
    \textbf{Feature Name} & \textbf{Feature Id} & \textbf{Comment} \\
    \hline
    \multicolumn{3}{|l|}{\emph{No contained Features are currently identified.}} \\
    \hline
  \end{tabularx}
\end{table}

\textbf{Intent:} Clinical Knowledge Review \& Approval Management identifies the organisational ability to undertake the stated health-service responsibility. The capability describes what the enterprise must be able to do; it does not assert that Harmonia performs or owns that activity.

\textbf{Boundary:} Harmonia relevance may arise because this activity establishes the context in which information is created, needed, interpreted, related, exchanged or presented. Responsibility for the underlying business decision or care activity remains with the appropriate actor or organisation.

\subsection{CT3 - Clinical Knowledge Publication \& Sharing}\label{cap:bc-e57a9c37}
\vspace{-0.5em}
\begin{table}[htbp]
  \centering
  \footnotesize
  \begin{tabularx}{1\textwidth}{|l|l|X|}
    \hline
    \multicolumn{3}{c}{\textbf{Capability Metadata}} \\
    \hline
    \textbf{Name} & \textbf{Value} & \textbf{Comment and/or Description} \\
    \hline
    \textbf{Reference Id} & \texttt{BC-E57A9C37} & \\
    \hline
    \textbf{Status} & Established working hierarchy & \\
    \hline
    \textbf{Ancestry} &  & \\
    \hline
    \textbf{Source} &  & \\
    \hline
    \textbf{Enabled By} &  & \\
    \hline
    \textbf{Supports} &  & \\
    \hline
    \textbf{Derived From} &  & \\
    \hline
    \textbf{Reviewed Status} & Review meaning, business boundary, information context and any unresolved ancestry. & \\
    \hline
  \end{tabularx}
\end{table}

\begin{table}[htbp]
  \centering
  \footnotesize
  \begin{tabularx}{1\textwidth}{|l|l|X|}
    \hline
    \multicolumn{3}{c}{\textbf{Contained Features}} \\
    \hline
    \textbf{Feature Name} & \textbf{Feature Id} & \textbf{Comment} \\
    \hline
    \multicolumn{3}{|l|}{\emph{No contained Features are currently identified.}} \\
    \hline
  \end{tabularx}
\end{table}

\textbf{Intent:} Clinical Knowledge Publication \& Sharing identifies the organisational ability to undertake the stated health-service responsibility. The capability describes what the enterprise must be able to do; it does not assert that Harmonia performs or owns that activity.

\textbf{Boundary:} Harmonia relevance may arise because this activity establishes the context in which information is created, needed, interpreted, related, exchanged or presented. Responsibility for the underlying business decision or care activity remains with the appropriate actor or organisation.

\chapter{CT1 - Standards, Semantics \& Reference Governance}\label{cap:bc-014363f4}
\noindent\textbf{Reference ID:} \texttt{BC-014363F4}\quad \textbf{Status:} Established working hierarchy\par
\noindent\textit{Source: Canonical Business Capability model / working map}\par
\textbf{Intent:} Standards, Semantics \& Reference Governance identifies the organisational ability to undertake the stated health-service responsibility. The capability describes what the enterprise must be able to do; it does not assert that Harmonia performs or owns that activity.

\textbf{Boundary:} Harmonia relevance may arise because this activity establishes the context in which information is created, needed, interpreted, related, exchanged or presented. Responsibility for the underlying business decision or care activity remains with the appropriate actor or organisation.

\reviewstatus{Check scope and information relevance; no business authority transfers to Harmonia.}
\section{CT2 - Health Information Standards Governance}\label{cap:bc-c01c01b9}
\vspace{-0.5em}
\begin{table}[htbp]
  \centering
  \footnotesize
  \begin{tabularx}{1\textwidth}{|l|l|X|}
    \hline
    \multicolumn{3}{c}{\textbf{Capability Metadata}} \\
    \hline
    \textbf{Name} & \textbf{Value} & \textbf{Comment and/or Description} \\
    \hline
    \textbf{Reference Id} & \texttt{BC-C01C01B9} & \\
    \hline
    \textbf{Status} & Established working hierarchy & \\
    \hline
    \textbf{Ancestry} &  & \\
    \hline
    \textbf{Source} &  & \\
    \hline
    \textbf{Enabled By} &  & \\
    \hline
    \textbf{Supports} &  & \\
    \hline
    \textbf{Derived From} &  & \\
    \hline
    \textbf{Reviewed Status} &  & \\
    \hline
  \end{tabularx}
\end{table}

\begin{table}[htbp]
  \centering
  \footnotesize
  \begin{tabularx}{1\textwidth}{|l|l|X|}
    \hline
    \multicolumn{3}{c}{\textbf{Contained Capabilities}} \\
    \hline
    \textbf{Capability Name} & \textbf{Capability Id} & \textbf{Comment} \\
    \hline
    Information Model Governance & \texttt{cap:bc-c9d7e7c2} (at \ref{cap:bc-c9d7e7c2}) & \\
    \hline
    Information  Model Discovery & \texttt{cap:bc-f55f16a6} (at \ref{cap:bc-f55f16a6}) & \\
    \hline
    Information Model Adoption & \texttt{cap:bc-b27b41ed} (at \ref{cap:bc-b27b41ed}) & \\
    \hline
    Information Model Education & \texttt{cap:bc-c5e3e979} (at \ref{cap:bc-c5e3e979}) & \\
    \hline
    Information Model Compliance & \texttt{cap:bc-ff2bf6be} (at \ref{cap:bc-ff2bf6be}) & \\
    \hline
    Standards Discovery & \texttt{cap:bc-4486449f} (at \ref{cap:bc-4486449f}) & \\
    \hline
    Standards Impact Identification & \texttt{cap:bc-98a00c71} (at \ref{cap:bc-98a00c71}) & \\
    \hline
    Standards Adoption & \texttt{cap:bc-cb858814} (at \ref{cap:bc-cb858814}) & \\
    \hline
    Standards Education & \texttt{cap:bc-cb56b8bc} (at \ref{cap:bc-cb56b8bc}) & \\
    \hline
    Standards Bodies Engagement & \texttt{cap:bc-bfc63cc1} (at \ref{cap:bc-bfc63cc1}) & \\
    \hline
    Standards Compliance & \texttt{cap:bc-e515be41} (at \ref{cap:bc-e515be41}) & \\
    \hline
  \end{tabularx}
\end{table}

\textbf{Intent:} Health Information Standards Governance identifies the organisational ability to undertake the stated health-service responsibility. The capability describes what the enterprise must be able to do; it does not assert that Harmonia performs or owns that activity.

\textbf{Boundary:} Harmonia relevance may arise because this activity establishes the context in which information is created, needed, interpreted, related, exchanged or presented. Responsibility for the underlying business decision or care activity remains with the appropriate actor or organisation.

\subsection{CT3 - Information Model Governance}\label{cap:bc-c9d7e7c2}
\vspace{-0.5em}
\begin{table}[htbp]
  \centering
  \footnotesize
  \begin{tabularx}{1\textwidth}{|l|l|X|}
    \hline
    \multicolumn{3}{c}{\textbf{Capability Metadata}} \\
    \hline
    \textbf{Name} & \textbf{Value} & \textbf{Comment and/or Description} \\
    \hline
    \textbf{Reference Id} & \texttt{BC-C9D7E7C2} & \\
    \hline
    \textbf{Status} & Established working hierarchy & \\
    \hline
    \textbf{Ancestry} &  & \\
    \hline
    \textbf{Source} &  & \\
    \hline
    \textbf{Enabled By} &  & \\
    \hline
    \textbf{Supports} &  & \\
    \hline
    \textbf{Derived From} &  & \\
    \hline
    \textbf{Reviewed Status} & Review meaning, business boundary, information context and any unresolved ancestry. & \\
    \hline
  \end{tabularx}
\end{table}

\begin{table}[htbp]
  \centering
  \footnotesize
  \begin{tabularx}{1\textwidth}{|l|l|X|}
    \hline
    \multicolumn{3}{c}{\textbf{Contained Features}} \\
    \hline
    \textbf{Feature Name} & \textbf{Feature Id} & \textbf{Comment} \\
    \hline
    \multicolumn{3}{|l|}{\emph{No contained Features are currently identified.}} \\
    \hline
  \end{tabularx}
\end{table}

\textbf{Intent:} Information Model Governance identifies the organisational ability to undertake the stated health-service responsibility. The capability describes what the enterprise must be able to do; it does not assert that Harmonia performs or owns that activity.

\textbf{Boundary:} Harmonia relevance may arise because this activity establishes the context in which information is created, needed, interpreted, related, exchanged or presented. Responsibility for the underlying business decision or care activity remains with the appropriate actor or organisation.

\subsection{CT3 - Information  Model Discovery}\label{cap:bc-f55f16a6}
\vspace{-0.5em}
\begin{table}[htbp]
  \centering
  \footnotesize
  \begin{tabularx}{1\textwidth}{|l|l|X|}
    \hline
    \multicolumn{3}{c}{\textbf{Capability Metadata}} \\
    \hline
    \textbf{Name} & \textbf{Value} & \textbf{Comment and/or Description} \\
    \hline
    \textbf{Reference Id} & \texttt{BC-F55F16A6} & \\
    \hline
    \textbf{Status} & Established working hierarchy & \\
    \hline
    \textbf{Ancestry} &  & \\
    \hline
    \textbf{Source} &  & \\
    \hline
    \textbf{Enabled By} &  & \\
    \hline
    \textbf{Supports} &  & \\
    \hline
    \textbf{Derived From} &  & \\
    \hline
    \textbf{Reviewed Status} & Review meaning, business boundary, information context and any unresolved ancestry. & \\
    \hline
  \end{tabularx}
\end{table}

\begin{table}[htbp]
  \centering
  \footnotesize
  \begin{tabularx}{1\textwidth}{|l|l|X|}
    \hline
    \multicolumn{3}{c}{\textbf{Contained Features}} \\
    \hline
    \textbf{Feature Name} & \textbf{Feature Id} & \textbf{Comment} \\
    \hline
    \multicolumn{3}{|l|}{\emph{No contained Features are currently identified.}} \\
    \hline
  \end{tabularx}
\end{table}

\textbf{Intent:} Information  Model Discovery identifies the organisational ability to undertake the stated health-service responsibility. The capability describes what the enterprise must be able to do; it does not assert that Harmonia performs or owns that activity.

\textbf{Boundary:} Harmonia relevance may arise because this activity establishes the context in which information is created, needed, interpreted, related, exchanged or presented. Responsibility for the underlying business decision or care activity remains with the appropriate actor or organisation.

\subsection{CT3 - Information Model Adoption}\label{cap:bc-b27b41ed}
\vspace{-0.5em}
\begin{table}[htbp]
  \centering
  \footnotesize
  \begin{tabularx}{1\textwidth}{|l|l|X|}
    \hline
    \multicolumn{3}{c}{\textbf{Capability Metadata}} \\
    \hline
    \textbf{Name} & \textbf{Value} & \textbf{Comment and/or Description} \\
    \hline
    \textbf{Reference Id} & \texttt{BC-B27B41ED} & \\
    \hline
    \textbf{Status} & Established working hierarchy & \\
    \hline
    \textbf{Ancestry} &  & \\
    \hline
    \textbf{Source} &  & \\
    \hline
    \textbf{Enabled By} &  & \\
    \hline
    \textbf{Supports} &  & \\
    \hline
    \textbf{Derived From} &  & \\
    \hline
    \textbf{Reviewed Status} & Review meaning, business boundary, information context and any unresolved ancestry. & \\
    \hline
  \end{tabularx}
\end{table}

\begin{table}[htbp]
  \centering
  \footnotesize
  \begin{tabularx}{1\textwidth}{|l|l|X|}
    \hline
    \multicolumn{3}{c}{\textbf{Contained Features}} \\
    \hline
    \textbf{Feature Name} & \textbf{Feature Id} & \textbf{Comment} \\
    \hline
    \multicolumn{3}{|l|}{\emph{No contained Features are currently identified.}} \\
    \hline
  \end{tabularx}
\end{table}

\textbf{Intent:} Information Model Adoption identifies the organisational ability to undertake the stated health-service responsibility. The capability describes what the enterprise must be able to do; it does not assert that Harmonia performs or owns that activity.

\textbf{Boundary:} Harmonia relevance may arise because this activity establishes the context in which information is created, needed, interpreted, related, exchanged or presented. Responsibility for the underlying business decision or care activity remains with the appropriate actor or organisation.

\subsection{CT3 - Information Model Education}\label{cap:bc-c5e3e979}
\vspace{-0.5em}
\begin{table}[htbp]
  \centering
  \footnotesize
  \begin{tabularx}{1\textwidth}{|l|l|X|}
    \hline
    \multicolumn{3}{c}{\textbf{Capability Metadata}} \\
    \hline
    \textbf{Name} & \textbf{Value} & \textbf{Comment and/or Description} \\
    \hline
    \textbf{Reference Id} & \texttt{BC-C5E3E979} & \\
    \hline
    \textbf{Status} & Established working hierarchy & \\
    \hline
    \textbf{Ancestry} &  & \\
    \hline
    \textbf{Source} &  & \\
    \hline
    \textbf{Enabled By} &  & \\
    \hline
    \textbf{Supports} &  & \\
    \hline
    \textbf{Derived From} &  & \\
    \hline
    \textbf{Reviewed Status} & Review meaning, business boundary, information context and any unresolved ancestry. & \\
    \hline
  \end{tabularx}
\end{table}

\begin{table}[htbp]
  \centering
  \footnotesize
  \begin{tabularx}{1\textwidth}{|l|l|X|}
    \hline
    \multicolumn{3}{c}{\textbf{Contained Features}} \\
    \hline
    \textbf{Feature Name} & \textbf{Feature Id} & \textbf{Comment} \\
    \hline
    \multicolumn{3}{|l|}{\emph{No contained Features are currently identified.}} \\
    \hline
  \end{tabularx}
\end{table}

\textbf{Intent:} Information Model Education identifies the organisational ability to undertake the stated health-service responsibility. The capability describes what the enterprise must be able to do; it does not assert that Harmonia performs or owns that activity.

\textbf{Boundary:} Harmonia relevance may arise because this activity establishes the context in which information is created, needed, interpreted, related, exchanged or presented. Responsibility for the underlying business decision or care activity remains with the appropriate actor or organisation.

\subsection{CT3 - Information Model Compliance}\label{cap:bc-ff2bf6be}
\vspace{-0.5em}
\begin{table}[htbp]
  \centering
  \footnotesize
  \begin{tabularx}{1\textwidth}{|l|l|X|}
    \hline
    \multicolumn{3}{c}{\textbf{Capability Metadata}} \\
    \hline
    \textbf{Name} & \textbf{Value} & \textbf{Comment and/or Description} \\
    \hline
    \textbf{Reference Id} & \texttt{BC-FF2BF6BE} & \\
    \hline
    \textbf{Status} & Established working hierarchy & \\
    \hline
    \textbf{Ancestry} &  & \\
    \hline
    \textbf{Source} &  & \\
    \hline
    \textbf{Enabled By} &  & \\
    \hline
    \textbf{Supports} &  & \\
    \hline
    \textbf{Derived From} &  & \\
    \hline
    \textbf{Reviewed Status} & Review meaning, business boundary, information context and any unresolved ancestry. & \\
    \hline
  \end{tabularx}
\end{table}

\begin{table}[htbp]
  \centering
  \footnotesize
  \begin{tabularx}{1\textwidth}{|l|l|X|}
    \hline
    \multicolumn{3}{c}{\textbf{Contained Features}} \\
    \hline
    \textbf{Feature Name} & \textbf{Feature Id} & \textbf{Comment} \\
    \hline
    \multicolumn{3}{|l|}{\emph{No contained Features are currently identified.}} \\
    \hline
  \end{tabularx}
\end{table}

\textbf{Intent:} Information Model Compliance identifies the organisational ability to undertake the stated health-service responsibility. The capability describes what the enterprise must be able to do; it does not assert that Harmonia performs or owns that activity.

\textbf{Boundary:} Harmonia relevance may arise because this activity establishes the context in which information is created, needed, interpreted, related, exchanged or presented. Responsibility for the underlying business decision or care activity remains with the appropriate actor or organisation.

\subsection{CT3 - Standards Discovery}\label{cap:bc-4486449f}
\vspace{-0.5em}
\begin{table}[htbp]
  \centering
  \footnotesize
  \begin{tabularx}{1\textwidth}{|l|l|X|}
    \hline
    \multicolumn{3}{c}{\textbf{Capability Metadata}} \\
    \hline
    \textbf{Name} & \textbf{Value} & \textbf{Comment and/or Description} \\
    \hline
    \textbf{Reference Id} & \texttt{BC-4486449F} & \\
    \hline
    \textbf{Status} & Established working hierarchy & \\
    \hline
    \textbf{Ancestry} &  & \\
    \hline
    \textbf{Source} &  & \\
    \hline
    \textbf{Enabled By} &  & \\
    \hline
    \textbf{Supports} &  & \\
    \hline
    \textbf{Derived From} &  & \\
    \hline
    \textbf{Reviewed Status} & Review meaning, business boundary, information context and any unresolved ancestry. & \\
    \hline
  \end{tabularx}
\end{table}

\begin{table}[htbp]
  \centering
  \footnotesize
  \begin{tabularx}{1\textwidth}{|l|l|X|}
    \hline
    \multicolumn{3}{c}{\textbf{Contained Features}} \\
    \hline
    \textbf{Feature Name} & \textbf{Feature Id} & \textbf{Comment} \\
    \hline
    \multicolumn{3}{|l|}{\emph{No contained Features are currently identified.}} \\
    \hline
  \end{tabularx}
\end{table}

\textbf{Intent:} Standards Discovery identifies the organisational ability to undertake the stated health-service responsibility. The capability describes what the enterprise must be able to do; it does not assert that Harmonia performs or owns that activity.

\textbf{Boundary:} Harmonia relevance may arise because this activity establishes the context in which information is created, needed, interpreted, related, exchanged or presented. Responsibility for the underlying business decision or care activity remains with the appropriate actor or organisation.

\subsection{CT3 - Standards Impact Identification}\label{cap:bc-98a00c71}
\vspace{-0.5em}
\begin{table}[htbp]
  \centering
  \footnotesize
  \begin{tabularx}{1\textwidth}{|l|l|X|}
    \hline
    \multicolumn{3}{c}{\textbf{Capability Metadata}} \\
    \hline
    \textbf{Name} & \textbf{Value} & \textbf{Comment and/or Description} \\
    \hline
    \textbf{Reference Id} & \texttt{BC-98A00C71} & \\
    \hline
    \textbf{Status} & Established working hierarchy & \\
    \hline
    \textbf{Ancestry} &  & \\
    \hline
    \textbf{Source} &  & \\
    \hline
    \textbf{Enabled By} &  & \\
    \hline
    \textbf{Supports} &  & \\
    \hline
    \textbf{Derived From} &  & \\
    \hline
    \textbf{Reviewed Status} & Review meaning, business boundary, information context and any unresolved ancestry. & \\
    \hline
  \end{tabularx}
\end{table}

\begin{table}[htbp]
  \centering
  \footnotesize
  \begin{tabularx}{1\textwidth}{|l|l|X|}
    \hline
    \multicolumn{3}{c}{\textbf{Contained Features}} \\
    \hline
    \textbf{Feature Name} & \textbf{Feature Id} & \textbf{Comment} \\
    \hline
    \multicolumn{3}{|l|}{\emph{No contained Features are currently identified.}} \\
    \hline
  \end{tabularx}
\end{table}

\textbf{Intent:} Standards Impact Identification identifies the organisational ability to undertake the stated health-service responsibility. The capability describes what the enterprise must be able to do; it does not assert that Harmonia performs or owns that activity.

\textbf{Boundary:} Harmonia relevance may arise because this activity establishes the context in which information is created, needed, interpreted, related, exchanged or presented. Responsibility for the underlying business decision or care activity remains with the appropriate actor or organisation.

\subsection{CT3 - Standards Adoption}\label{cap:bc-cb858814}
\vspace{-0.5em}
\begin{table}[htbp]
  \centering
  \footnotesize
  \begin{tabularx}{1\textwidth}{|l|l|X|}
    \hline
    \multicolumn{3}{c}{\textbf{Capability Metadata}} \\
    \hline
    \textbf{Name} & \textbf{Value} & \textbf{Comment and/or Description} \\
    \hline
    \textbf{Reference Id} & \texttt{BC-CB858814} & \\
    \hline
    \textbf{Status} & Established working hierarchy & \\
    \hline
    \textbf{Ancestry} &  & \\
    \hline
    \textbf{Source} &  & \\
    \hline
    \textbf{Enabled By} &  & \\
    \hline
    \textbf{Supports} &  & \\
    \hline
    \textbf{Derived From} &  & \\
    \hline
    \textbf{Reviewed Status} & Review meaning, business boundary, information context and any unresolved ancestry. & \\
    \hline
  \end{tabularx}
\end{table}

\begin{table}[htbp]
  \centering
  \footnotesize
  \begin{tabularx}{1\textwidth}{|l|l|X|}
    \hline
    \multicolumn{3}{c}{\textbf{Contained Features}} \\
    \hline
    \textbf{Feature Name} & \textbf{Feature Id} & \textbf{Comment} \\
    \hline
    \multicolumn{3}{|l|}{\emph{No contained Features are currently identified.}} \\
    \hline
  \end{tabularx}
\end{table}

\textbf{Intent:} Standards Adoption identifies the organisational ability to undertake the stated health-service responsibility. The capability describes what the enterprise must be able to do; it does not assert that Harmonia performs or owns that activity.

\textbf{Boundary:} Harmonia relevance may arise because this activity establishes the context in which information is created, needed, interpreted, related, exchanged or presented. Responsibility for the underlying business decision or care activity remains with the appropriate actor or organisation.

\subsection{CT3 - Standards Education}\label{cap:bc-cb56b8bc}
\vspace{-0.5em}
\begin{table}[htbp]
  \centering
  \footnotesize
  \begin{tabularx}{1\textwidth}{|l|l|X|}
    \hline
    \multicolumn{3}{c}{\textbf{Capability Metadata}} \\
    \hline
    \textbf{Name} & \textbf{Value} & \textbf{Comment and/or Description} \\
    \hline
    \textbf{Reference Id} & \texttt{BC-CB56B8BC} & \\
    \hline
    \textbf{Status} & Established working hierarchy & \\
    \hline
    \textbf{Ancestry} &  & \\
    \hline
    \textbf{Source} &  & \\
    \hline
    \textbf{Enabled By} &  & \\
    \hline
    \textbf{Supports} &  & \\
    \hline
    \textbf{Derived From} &  & \\
    \hline
    \textbf{Reviewed Status} & Review meaning, business boundary, information context and any unresolved ancestry. & \\
    \hline
  \end{tabularx}
\end{table}

\begin{table}[htbp]
  \centering
  \footnotesize
  \begin{tabularx}{1\textwidth}{|l|l|X|}
    \hline
    \multicolumn{3}{c}{\textbf{Contained Features}} \\
    \hline
    \textbf{Feature Name} & \textbf{Feature Id} & \textbf{Comment} \\
    \hline
    \multicolumn{3}{|l|}{\emph{No contained Features are currently identified.}} \\
    \hline
  \end{tabularx}
\end{table}

\textbf{Intent:} Standards Education identifies the organisational ability to undertake the stated health-service responsibility. The capability describes what the enterprise must be able to do; it does not assert that Harmonia performs or owns that activity.

\textbf{Boundary:} Harmonia relevance may arise because this activity establishes the context in which information is created, needed, interpreted, related, exchanged or presented. Responsibility for the underlying business decision or care activity remains with the appropriate actor or organisation.

\subsection{CT3 - Standards Bodies Engagement}\label{cap:bc-bfc63cc1}
\vspace{-0.5em}
\begin{table}[htbp]
  \centering
  \footnotesize
  \begin{tabularx}{1\textwidth}{|l|l|X|}
    \hline
    \multicolumn{3}{c}{\textbf{Capability Metadata}} \\
    \hline
    \textbf{Name} & \textbf{Value} & \textbf{Comment and/or Description} \\
    \hline
    \textbf{Reference Id} & \texttt{BC-BFC63CC1} & \\
    \hline
    \textbf{Status} & Established working hierarchy & \\
    \hline
    \textbf{Ancestry} &  & \\
    \hline
    \textbf{Source} &  & \\
    \hline
    \textbf{Enabled By} &  & \\
    \hline
    \textbf{Supports} &  & \\
    \hline
    \textbf{Derived From} &  & \\
    \hline
    \textbf{Reviewed Status} & Review meaning, business boundary, information context and any unresolved ancestry. & \\
    \hline
  \end{tabularx}
\end{table}

\begin{table}[htbp]
  \centering
  \footnotesize
  \begin{tabularx}{1\textwidth}{|l|l|X|}
    \hline
    \multicolumn{3}{c}{\textbf{Contained Features}} \\
    \hline
    \textbf{Feature Name} & \textbf{Feature Id} & \textbf{Comment} \\
    \hline
    \multicolumn{3}{|l|}{\emph{No contained Features are currently identified.}} \\
    \hline
  \end{tabularx}
\end{table}

\textbf{Intent:} Standards Bodies Engagement identifies the organisational ability to undertake the stated health-service responsibility. The capability describes what the enterprise must be able to do; it does not assert that Harmonia performs or owns that activity.

\textbf{Boundary:} Harmonia relevance may arise because this activity establishes the context in which information is created, needed, interpreted, related, exchanged or presented. Responsibility for the underlying business decision or care activity remains with the appropriate actor or organisation.

\subsection{CT3 - Standards Compliance}\label{cap:bc-e515be41}
\vspace{-0.5em}
\begin{table}[htbp]
  \centering
  \footnotesize
  \begin{tabularx}{1\textwidth}{|l|l|X|}
    \hline
    \multicolumn{3}{c}{\textbf{Capability Metadata}} \\
    \hline
    \textbf{Name} & \textbf{Value} & \textbf{Comment and/or Description} \\
    \hline
    \textbf{Reference Id} & \texttt{BC-E515BE41} & \\
    \hline
    \textbf{Status} & Established working hierarchy & \\
    \hline
    \textbf{Ancestry} &  & \\
    \hline
    \textbf{Source} &  & \\
    \hline
    \textbf{Enabled By} &  & \\
    \hline
    \textbf{Supports} &  & \\
    \hline
    \textbf{Derived From} &  & \\
    \hline
    \textbf{Reviewed Status} & Review meaning, business boundary, information context and any unresolved ancestry. & \\
    \hline
  \end{tabularx}
\end{table}

\begin{table}[htbp]
  \centering
  \footnotesize
  \begin{tabularx}{1\textwidth}{|l|l|X|}
    \hline
    \multicolumn{3}{c}{\textbf{Contained Features}} \\
    \hline
    \textbf{Feature Name} & \textbf{Feature Id} & \textbf{Comment} \\
    \hline
    \multicolumn{3}{|l|}{\emph{No contained Features are currently identified.}} \\
    \hline
  \end{tabularx}
\end{table}

\textbf{Intent:} Standards Compliance identifies the organisational ability to undertake the stated health-service responsibility. The capability describes what the enterprise must be able to do; it does not assert that Harmonia performs or owns that activity.

\textbf{Boundary:} Harmonia relevance may arise because this activity establishes the context in which information is created, needed, interpreted, related, exchanged or presented. Responsibility for the underlying business decision or care activity remains with the appropriate actor or organisation.

\section{CT2 - Clinical Semantics Governance}\label{cap:bc-f9f83f32}
\vspace{-0.5em}
\begin{table}[htbp]
  \centering
  \footnotesize
  \begin{tabularx}{1\textwidth}{|l|l|X|}
    \hline
    \multicolumn{3}{c}{\textbf{Capability Metadata}} \\
    \hline
    \textbf{Name} & \textbf{Value} & \textbf{Comment and/or Description} \\
    \hline
    \textbf{Reference Id} & \texttt{BC-F9F83F32} & \\
    \hline
    \textbf{Status} & Established working hierarchy & \\
    \hline
    \textbf{Ancestry} &  & \\
    \hline
    \textbf{Source} &  & \\
    \hline
    \textbf{Enabled By} &  & \\
    \hline
    \textbf{Supports} &  & \\
    \hline
    \textbf{Derived From} &  & \\
    \hline
    \textbf{Reviewed Status} &  & \\
    \hline
  \end{tabularx}
\end{table}

\begin{table}[htbp]
  \centering
  \footnotesize
  \begin{tabularx}{1\textwidth}{|l|l|X|}
    \hline
    \multicolumn{3}{c}{\textbf{Contained Capabilities}} \\
    \hline
    \textbf{Capability Name} & \textbf{Capability Id} & \textbf{Comment} \\
    \hline
    Terminology Governance & \texttt{cap:bc-9d2ac8a5} (at \ref{cap:bc-9d2ac8a5}) & \\
    \hline
    Classification Governance & \texttt{cap:bc-9dbb18d4} (at \ref{cap:bc-9dbb18d4}) & \\
    \hline
    Semantic Mapping Governance & \texttt{cap:bc-05135118} (at \ref{cap:bc-05135118}) & \\
    \hline
    Terminology Utilisation Oversight & \texttt{cap:bc-18835f13} (at \ref{cap:bc-18835f13}) & \\
    \hline
    Terminology Localisation & \texttt{cap:bc-d4198837} (at \ref{cap:bc-d4198837}) & \\
    \hline
    Classification Utilisation Oversight & \texttt{cap:bc-a6fde957} (at \ref{cap:bc-a6fde957}) & \\
    \hline
  \end{tabularx}
\end{table}

\textbf{Intent:} Clinical Semantics Governance identifies the organisational ability to undertake the stated health-service responsibility. The capability describes what the enterprise must be able to do; it does not assert that Harmonia performs or owns that activity.

\textbf{Boundary:} Harmonia relevance may arise because this activity establishes the context in which information is created, needed, interpreted, related, exchanged or presented. Responsibility for the underlying business decision or care activity remains with the appropriate actor or organisation.

\subsection{CT3 - Terminology Governance}\label{cap:bc-9d2ac8a5}
\vspace{-0.5em}
\begin{table}[htbp]
  \centering
  \footnotesize
  \begin{tabularx}{1\textwidth}{|l|l|X|}
    \hline
    \multicolumn{3}{c}{\textbf{Capability Metadata}} \\
    \hline
    \textbf{Name} & \textbf{Value} & \textbf{Comment and/or Description} \\
    \hline
    \textbf{Reference Id} & \texttt{BC-9D2AC8A5} & \\
    \hline
    \textbf{Status} & Established working hierarchy & \\
    \hline
    \textbf{Ancestry} &  & \\
    \hline
    \textbf{Source} &  & \\
    \hline
    \textbf{Enabled By} &  & \\
    \hline
    \textbf{Supports} &  & \\
    \hline
    \textbf{Derived From} &  & \\
    \hline
    \textbf{Reviewed Status} & Review meaning, business boundary, information context and any unresolved ancestry. & \\
    \hline
  \end{tabularx}
\end{table}

\begin{table}[htbp]
  \centering
  \footnotesize
  \begin{tabularx}{1\textwidth}{|l|l|X|}
    \hline
    \multicolumn{3}{c}{\textbf{Contained Features}} \\
    \hline
    \textbf{Feature Name} & \textbf{Feature Id} & \textbf{Comment} \\
    \hline
    \multicolumn{3}{|l|}{\emph{No contained Features are currently identified.}} \\
    \hline
  \end{tabularx}
\end{table}

\textbf{Intent:} Terminology Governance identifies the organisational ability to undertake the stated health-service responsibility. The capability describes what the enterprise must be able to do; it does not assert that Harmonia performs or owns that activity.

\textbf{Boundary:} Harmonia relevance may arise because this activity establishes the context in which information is created, needed, interpreted, related, exchanged or presented. Responsibility for the underlying business decision or care activity remains with the appropriate actor or organisation.

\subsection{CT3 - Classification Governance}\label{cap:bc-9dbb18d4}
\vspace{-0.5em}
\begin{table}[htbp]
  \centering
  \footnotesize
  \begin{tabularx}{1\textwidth}{|l|l|X|}
    \hline
    \multicolumn{3}{c}{\textbf{Capability Metadata}} \\
    \hline
    \textbf{Name} & \textbf{Value} & \textbf{Comment and/or Description} \\
    \hline
    \textbf{Reference Id} & \texttt{BC-9DBB18D4} & \\
    \hline
    \textbf{Status} & Established working hierarchy & \\
    \hline
    \textbf{Ancestry} &  & \\
    \hline
    \textbf{Source} &  & \\
    \hline
    \textbf{Enabled By} &  & \\
    \hline
    \textbf{Supports} &  & \\
    \hline
    \textbf{Derived From} &  & \\
    \hline
    \textbf{Reviewed Status} & Review meaning, business boundary, information context and any unresolved ancestry. & \\
    \hline
  \end{tabularx}
\end{table}

\begin{table}[htbp]
  \centering
  \footnotesize
  \begin{tabularx}{1\textwidth}{|l|l|X|}
    \hline
    \multicolumn{3}{c}{\textbf{Contained Features}} \\
    \hline
    \textbf{Feature Name} & \textbf{Feature Id} & \textbf{Comment} \\
    \hline
    \multicolumn{3}{|l|}{\emph{No contained Features are currently identified.}} \\
    \hline
  \end{tabularx}
\end{table}

\textbf{Intent:} Classification Governance identifies the organisational ability to undertake the stated health-service responsibility. The capability describes what the enterprise must be able to do; it does not assert that Harmonia performs or owns that activity.

\textbf{Boundary:} Harmonia relevance may arise because this activity establishes the context in which information is created, needed, interpreted, related, exchanged or presented. Responsibility for the underlying business decision or care activity remains with the appropriate actor or organisation.

\subsection{CT3 - Semantic Mapping Governance}\label{cap:bc-05135118}
\vspace{-0.5em}
\begin{table}[htbp]
  \centering
  \footnotesize
  \begin{tabularx}{1\textwidth}{|l|l|X|}
    \hline
    \multicolumn{3}{c}{\textbf{Capability Metadata}} \\
    \hline
    \textbf{Name} & \textbf{Value} & \textbf{Comment and/or Description} \\
    \hline
    \textbf{Reference Id} & \texttt{BC-05135118} & \\
    \hline
    \textbf{Status} & Established working hierarchy & \\
    \hline
    \textbf{Ancestry} &  & \\
    \hline
    \textbf{Source} &  & \\
    \hline
    \textbf{Enabled By} &  & \\
    \hline
    \textbf{Supports} &  & \\
    \hline
    \textbf{Derived From} &  & \\
    \hline
    \textbf{Reviewed Status} & Review meaning, business boundary, information context and any unresolved ancestry. & \\
    \hline
  \end{tabularx}
\end{table}

\begin{table}[htbp]
  \centering
  \footnotesize
  \begin{tabularx}{1\textwidth}{|l|l|X|}
    \hline
    \multicolumn{3}{c}{\textbf{Contained Features}} \\
    \hline
    \textbf{Feature Name} & \textbf{Feature Id} & \textbf{Comment} \\
    \hline
    \multicolumn{3}{|l|}{\emph{No contained Features are currently identified.}} \\
    \hline
  \end{tabularx}
\end{table}

\textbf{Intent:} Semantic Mapping Governance identifies the organisational ability to undertake the stated health-service responsibility. The capability describes what the enterprise must be able to do; it does not assert that Harmonia performs or owns that activity.

\textbf{Boundary:} Harmonia relevance may arise because this activity establishes the context in which information is created, needed, interpreted, related, exchanged or presented. Responsibility for the underlying business decision or care activity remains with the appropriate actor or organisation.

\subsection{CT3 - Terminology Utilisation Oversight}\label{cap:bc-18835f13}
\vspace{-0.5em}
\begin{table}[htbp]
  \centering
  \footnotesize
  \begin{tabularx}{1\textwidth}{|l|l|X|}
    \hline
    \multicolumn{3}{c}{\textbf{Capability Metadata}} \\
    \hline
    \textbf{Name} & \textbf{Value} & \textbf{Comment and/or Description} \\
    \hline
    \textbf{Reference Id} & \texttt{BC-18835F13} & \\
    \hline
    \textbf{Status} & Established working hierarchy & \\
    \hline
    \textbf{Ancestry} &  & \\
    \hline
    \textbf{Source} &  & \\
    \hline
    \textbf{Enabled By} &  & \\
    \hline
    \textbf{Supports} &  & \\
    \hline
    \textbf{Derived From} &  & \\
    \hline
    \textbf{Reviewed Status} & Review meaning, business boundary, information context and any unresolved ancestry. & \\
    \hline
  \end{tabularx}
\end{table}

\begin{table}[htbp]
  \centering
  \footnotesize
  \begin{tabularx}{1\textwidth}{|l|l|X|}
    \hline
    \multicolumn{3}{c}{\textbf{Contained Features}} \\
    \hline
    \textbf{Feature Name} & \textbf{Feature Id} & \textbf{Comment} \\
    \hline
    \multicolumn{3}{|l|}{\emph{No contained Features are currently identified.}} \\
    \hline
  \end{tabularx}
\end{table}

\textbf{Intent:} Terminology Utilisation Oversight identifies the organisational ability to undertake the stated health-service responsibility. The capability describes what the enterprise must be able to do; it does not assert that Harmonia performs or owns that activity.

\textbf{Boundary:} Harmonia relevance may arise because this activity establishes the context in which information is created, needed, interpreted, related, exchanged or presented. Responsibility for the underlying business decision or care activity remains with the appropriate actor or organisation.

\subsection{CT3 - Terminology Localisation}\label{cap:bc-d4198837}
\vspace{-0.5em}
\begin{table}[htbp]
  \centering
  \footnotesize
  \begin{tabularx}{1\textwidth}{|l|l|X|}
    \hline
    \multicolumn{3}{c}{\textbf{Capability Metadata}} \\
    \hline
    \textbf{Name} & \textbf{Value} & \textbf{Comment and/or Description} \\
    \hline
    \textbf{Reference Id} & \texttt{BC-D4198837} & \\
    \hline
    \textbf{Status} & Established working hierarchy & \\
    \hline
    \textbf{Ancestry} &  & \\
    \hline
    \textbf{Source} &  & \\
    \hline
    \textbf{Enabled By} &  & \\
    \hline
    \textbf{Supports} &  & \\
    \hline
    \textbf{Derived From} &  & \\
    \hline
    \textbf{Reviewed Status} & Review meaning, business boundary, information context and any unresolved ancestry. & \\
    \hline
  \end{tabularx}
\end{table}

\begin{table}[htbp]
  \centering
  \footnotesize
  \begin{tabularx}{1\textwidth}{|l|l|X|}
    \hline
    \multicolumn{3}{c}{\textbf{Contained Features}} \\
    \hline
    \textbf{Feature Name} & \textbf{Feature Id} & \textbf{Comment} \\
    \hline
    \multicolumn{3}{|l|}{\emph{No contained Features are currently identified.}} \\
    \hline
  \end{tabularx}
\end{table}

\textbf{Intent:} Terminology Localisation identifies the organisational ability to undertake the stated health-service responsibility. The capability describes what the enterprise must be able to do; it does not assert that Harmonia performs or owns that activity.

\textbf{Boundary:} Harmonia relevance may arise because this activity establishes the context in which information is created, needed, interpreted, related, exchanged or presented. Responsibility for the underlying business decision or care activity remains with the appropriate actor or organisation.

\subsection{CT3 - Classification Utilisation Oversight}\label{cap:bc-a6fde957}
\vspace{-0.5em}
\begin{table}[htbp]
  \centering
  \footnotesize
  \begin{tabularx}{1\textwidth}{|l|l|X|}
    \hline
    \multicolumn{3}{c}{\textbf{Capability Metadata}} \\
    \hline
    \textbf{Name} & \textbf{Value} & \textbf{Comment and/or Description} \\
    \hline
    \textbf{Reference Id} & \texttt{BC-A6FDE957} & \\
    \hline
    \textbf{Status} & Established working hierarchy & \\
    \hline
    \textbf{Ancestry} &  & \\
    \hline
    \textbf{Source} &  & \\
    \hline
    \textbf{Enabled By} &  & \\
    \hline
    \textbf{Supports} &  & \\
    \hline
    \textbf{Derived From} &  & \\
    \hline
    \textbf{Reviewed Status} & Review meaning, business boundary, information context and any unresolved ancestry. & \\
    \hline
  \end{tabularx}
\end{table}

\begin{table}[htbp]
  \centering
  \footnotesize
  \begin{tabularx}{1\textwidth}{|l|l|X|}
    \hline
    \multicolumn{3}{c}{\textbf{Contained Features}} \\
    \hline
    \textbf{Feature Name} & \textbf{Feature Id} & \textbf{Comment} \\
    \hline
    \multicolumn{3}{|l|}{\emph{No contained Features are currently identified.}} \\
    \hline
  \end{tabularx}
\end{table}

\textbf{Intent:} Classification Utilisation Oversight identifies the organisational ability to undertake the stated health-service responsibility. The capability describes what the enterprise must be able to do; it does not assert that Harmonia performs or owns that activity.

\textbf{Boundary:} Harmonia relevance may arise because this activity establishes the context in which information is created, needed, interpreted, related, exchanged or presented. Responsibility for the underlying business decision or care activity remains with the appropriate actor or organisation.

\chapter{CT1 - Connected Health Services}\label{cap:bc-d0174b04}
\noindent\textbf{Reference ID:} \texttt{BC-D0174B04}\quad \textbf{Status:} Established working hierarchy\par
\noindent\textit{Source: Canonical Business Capability model / working map}\par
\textbf{Intent:} Connected Health Services identifies the organisational ability to undertake the stated health-service responsibility. The capability describes what the enterprise must be able to do; it does not assert that Harmonia performs or owns that activity.

\textbf{Boundary:} Harmonia relevance may arise because this activity establishes the context in which information is created, needed, interpreted, related, exchanged or presented. Responsibility for the underlying business decision or care activity remains with the appropriate actor or organisation.

\reviewstatus{Check scope and information relevance; no business authority transfers to Harmonia.}
\section{CT2 - Health Information Exchange}\label{cap:bc-1a2c0473}
\vspace{-0.5em}
\begin{table}[htbp]
  \centering
  \footnotesize
  \begin{tabularx}{1\textwidth}{|l|l|X|}
    \hline
    \multicolumn{3}{c}{\textbf{Capability Metadata}} \\
    \hline
    \textbf{Name} & \textbf{Value} & \textbf{Comment and/or Description} \\
    \hline
    \textbf{Reference Id} & \texttt{BC-1A2C0473} & \\
    \hline
    \textbf{Status} & Established working hierarchy & \\
    \hline
    \textbf{Ancestry} &  & \\
    \hline
    \textbf{Source} &  & \\
    \hline
    \textbf{Enabled By} &  & \\
    \hline
    \textbf{Supports} &  & \\
    \hline
    \textbf{Derived From} &  & \\
    \hline
    \textbf{Reviewed Status} &  & \\
    \hline
  \end{tabularx}
\end{table}

\begin{table}[htbp]
  \centering
  \footnotesize
  \begin{tabularx}{1\textwidth}{|l|l|X|}
    \hline
    \multicolumn{3}{c}{\textbf{Contained Capabilities}} \\
    \hline
    \textbf{Capability Name} & \textbf{Capability Id} & \textbf{Comment} \\
    \hline
    Standards-Based Interoperability Exploration & \texttt{cap:bc-25f9b730} (at \ref{cap:bc-25f9b730}) & \\
    \hline
    Standards-Based Interoperability Support & \texttt{cap:bc-5ca9556e} (at \ref{cap:bc-5ca9556e}) & \\
    \hline
    Standards-Based Interoperability Publication & \texttt{cap:bc-9d0c5f85} (at \ref{cap:bc-9d0c5f85}) & \\
    \hline
    Standards-Based Interoperability Client-Support & \texttt{cap:bc-ea4804ca} (at \ref{cap:bc-ea4804ca}) & \\
    \hline
    Health Information Sharing Framework Oversight & \texttt{cap:bc-eae75ce5} (at \ref{cap:bc-eae75ce5}) & \\
    \hline
    Health Information Sharing Client-Support & \texttt{cap:bc-f6df2559} (at \ref{cap:bc-f6df2559}) & \\
    \hline
    Health Information Sharing Release Management & \texttt{cap:bc-93bae720} (at \ref{cap:bc-93bae720}) & \\
    \hline
  \end{tabularx}
\end{table}

\textbf{Intent:} Health Information Exchange identifies the organisational ability to undertake the stated health-service responsibility. The capability describes what the enterprise must be able to do; it does not assert that Harmonia performs or owns that activity.

\textbf{Boundary:} Harmonia relevance may arise because this activity establishes the context in which information is created, needed, interpreted, related, exchanged or presented. Responsibility for the underlying business decision or care activity remains with the appropriate actor or organisation.

\subsection{CT3 - Standards-Based Interoperability Exploration}\label{cap:bc-25f9b730}
\vspace{-0.5em}
\begin{table}[htbp]
  \centering
  \footnotesize
  \begin{tabularx}{1\textwidth}{|l|l|X|}
    \hline
    \multicolumn{3}{c}{\textbf{Capability Metadata}} \\
    \hline
    \textbf{Name} & \textbf{Value} & \textbf{Comment and/or Description} \\
    \hline
    \textbf{Reference Id} & \texttt{BC-25F9B730} & \\
    \hline
    \textbf{Status} & Established working hierarchy & \\
    \hline
    \textbf{Ancestry} &  & \\
    \hline
    \textbf{Source} &  & \\
    \hline
    \textbf{Enabled By} &  & \\
    \hline
    \textbf{Supports} &  & \\
    \hline
    \textbf{Derived From} &  & \\
    \hline
    \textbf{Reviewed Status} & Review meaning, business boundary, information context and any unresolved ancestry. & \\
    \hline
  \end{tabularx}
\end{table}

\begin{table}[htbp]
  \centering
  \footnotesize
  \begin{tabularx}{1\textwidth}{|l|l|X|}
    \hline
    \multicolumn{3}{c}{\textbf{Contained Features}} \\
    \hline
    \textbf{Feature Name} & \textbf{Feature Id} & \textbf{Comment} \\
    \hline
    \multicolumn{3}{|l|}{\emph{No contained Features are currently identified.}} \\
    \hline
  \end{tabularx}
\end{table}

\textbf{Intent:} Standards-Based Interoperability Exploration identifies the organisational ability to undertake the stated health-service responsibility. The capability describes what the enterprise must be able to do; it does not assert that Harmonia performs or owns that activity.

\textbf{Boundary:} Harmonia relevance may arise because this activity establishes the context in which information is created, needed, interpreted, related, exchanged or presented. Responsibility for the underlying business decision or care activity remains with the appropriate actor or organisation.

\subsection{CT3 - Standards-Based Interoperability Support}\label{cap:bc-5ca9556e}
\vspace{-0.5em}
\begin{table}[htbp]
  \centering
  \footnotesize
  \begin{tabularx}{1\textwidth}{|l|l|X|}
    \hline
    \multicolumn{3}{c}{\textbf{Capability Metadata}} \\
    \hline
    \textbf{Name} & \textbf{Value} & \textbf{Comment and/or Description} \\
    \hline
    \textbf{Reference Id} & \texttt{BC-5CA9556E} & \\
    \hline
    \textbf{Status} & Established working hierarchy & \\
    \hline
    \textbf{Ancestry} &  & \\
    \hline
    \textbf{Source} &  & \\
    \hline
    \textbf{Enabled By} &  & \\
    \hline
    \textbf{Supports} &  & \\
    \hline
    \textbf{Derived From} &  & \\
    \hline
    \textbf{Reviewed Status} & Review meaning, business boundary, information context and any unresolved ancestry. & \\
    \hline
  \end{tabularx}
\end{table}

\begin{table}[htbp]
  \centering
  \footnotesize
  \begin{tabularx}{1\textwidth}{|l|l|X|}
    \hline
    \multicolumn{3}{c}{\textbf{Contained Features}} \\
    \hline
    \textbf{Feature Name} & \textbf{Feature Id} & \textbf{Comment} \\
    \hline
    \multicolumn{3}{|l|}{\emph{No contained Features are currently identified.}} \\
    \hline
  \end{tabularx}
\end{table}

\textbf{Intent:} Standards-Based Interoperability Support identifies the organisational ability to undertake the stated health-service responsibility. The capability describes what the enterprise must be able to do; it does not assert that Harmonia performs or owns that activity.

\textbf{Boundary:} Harmonia relevance may arise because this activity establishes the context in which information is created, needed, interpreted, related, exchanged or presented. Responsibility for the underlying business decision or care activity remains with the appropriate actor or organisation.

\subsection{CT3 - Standards-Based Interoperability Publication}\label{cap:bc-9d0c5f85}
\vspace{-0.5em}
\begin{table}[htbp]
  \centering
  \footnotesize
  \begin{tabularx}{1\textwidth}{|l|l|X|}
    \hline
    \multicolumn{3}{c}{\textbf{Capability Metadata}} \\
    \hline
    \textbf{Name} & \textbf{Value} & \textbf{Comment and/or Description} \\
    \hline
    \textbf{Reference Id} & \texttt{BC-9D0C5F85} & \\
    \hline
    \textbf{Status} & Established working hierarchy & \\
    \hline
    \textbf{Ancestry} &  & \\
    \hline
    \textbf{Source} &  & \\
    \hline
    \textbf{Enabled By} &  & \\
    \hline
    \textbf{Supports} &  & \\
    \hline
    \textbf{Derived From} &  & \\
    \hline
    \textbf{Reviewed Status} & Review meaning, business boundary, information context and any unresolved ancestry. & \\
    \hline
  \end{tabularx}
\end{table}

\begin{table}[htbp]
  \centering
  \footnotesize
  \begin{tabularx}{1\textwidth}{|l|l|X|}
    \hline
    \multicolumn{3}{c}{\textbf{Contained Features}} \\
    \hline
    \textbf{Feature Name} & \textbf{Feature Id} & \textbf{Comment} \\
    \hline
    \multicolumn{3}{|l|}{\emph{No contained Features are currently identified.}} \\
    \hline
  \end{tabularx}
\end{table}

\textbf{Intent:} Standards-Based Interoperability Publication identifies the organisational ability to undertake the stated health-service responsibility. The capability describes what the enterprise must be able to do; it does not assert that Harmonia performs or owns that activity.

\textbf{Boundary:} Harmonia relevance may arise because this activity establishes the context in which information is created, needed, interpreted, related, exchanged or presented. Responsibility for the underlying business decision or care activity remains with the appropriate actor or organisation.

\subsection{CT3 - Standards-Based Interoperability Client-Support}\label{cap:bc-ea4804ca}
\vspace{-0.5em}
\begin{table}[htbp]
  \centering
  \footnotesize
  \begin{tabularx}{1\textwidth}{|l|l|X|}
    \hline
    \multicolumn{3}{c}{\textbf{Capability Metadata}} \\
    \hline
    \textbf{Name} & \textbf{Value} & \textbf{Comment and/or Description} \\
    \hline
    \textbf{Reference Id} & \texttt{BC-EA4804CA} & \\
    \hline
    \textbf{Status} & Established working hierarchy & \\
    \hline
    \textbf{Ancestry} &  & \\
    \hline
    \textbf{Source} &  & \\
    \hline
    \textbf{Enabled By} &  & \\
    \hline
    \textbf{Supports} &  & \\
    \hline
    \textbf{Derived From} &  & \\
    \hline
    \textbf{Reviewed Status} & Review meaning, business boundary, information context and any unresolved ancestry. & \\
    \hline
  \end{tabularx}
\end{table}

\begin{table}[htbp]
  \centering
  \footnotesize
  \begin{tabularx}{1\textwidth}{|l|l|X|}
    \hline
    \multicolumn{3}{c}{\textbf{Contained Features}} \\
    \hline
    \textbf{Feature Name} & \textbf{Feature Id} & \textbf{Comment} \\
    \hline
    \multicolumn{3}{|l|}{\emph{No contained Features are currently identified.}} \\
    \hline
  \end{tabularx}
\end{table}

\textbf{Intent:} Standards-Based Interoperability Client-Support identifies the organisational ability to undertake the stated health-service responsibility. The capability describes what the enterprise must be able to do; it does not assert that Harmonia performs or owns that activity.

\textbf{Boundary:} Harmonia relevance may arise because this activity establishes the context in which information is created, needed, interpreted, related, exchanged or presented. Responsibility for the underlying business decision or care activity remains with the appropriate actor or organisation.

\subsection{CT3 - Health Information Sharing Framework Oversight}\label{cap:bc-eae75ce5}
\vspace{-0.5em}
\begin{table}[htbp]
  \centering
  \footnotesize
  \begin{tabularx}{1\textwidth}{|l|l|X|}
    \hline
    \multicolumn{3}{c}{\textbf{Capability Metadata}} \\
    \hline
    \textbf{Name} & \textbf{Value} & \textbf{Comment and/or Description} \\
    \hline
    \textbf{Reference Id} & \texttt{BC-EAE75CE5} & \\
    \hline
    \textbf{Status} & Established working hierarchy & \\
    \hline
    \textbf{Ancestry} &  & \\
    \hline
    \textbf{Source} &  & \\
    \hline
    \textbf{Enabled By} &  & \\
    \hline
    \textbf{Supports} &  & \\
    \hline
    \textbf{Derived From} &  & \\
    \hline
    \textbf{Reviewed Status} & Review meaning, business boundary, information context and any unresolved ancestry. & \\
    \hline
  \end{tabularx}
\end{table}

\begin{table}[htbp]
  \centering
  \footnotesize
  \begin{tabularx}{1\textwidth}{|l|l|X|}
    \hline
    \multicolumn{3}{c}{\textbf{Contained Features}} \\
    \hline
    \textbf{Feature Name} & \textbf{Feature Id} & \textbf{Comment} \\
    \hline
    \multicolumn{3}{|l|}{\emph{No contained Features are currently identified.}} \\
    \hline
  \end{tabularx}
\end{table}

\textbf{Intent:} Health Information Sharing Framework Oversight identifies the organisational ability to undertake the stated health-service responsibility. The capability describes what the enterprise must be able to do; it does not assert that Harmonia performs or owns that activity.

\textbf{Boundary:} Harmonia relevance may arise because this activity establishes the context in which information is created, needed, interpreted, related, exchanged or presented. Responsibility for the underlying business decision or care activity remains with the appropriate actor or organisation.

\subsection{CT3 - Health Information Sharing Client-Support}\label{cap:bc-f6df2559}
\vspace{-0.5em}
\begin{table}[htbp]
  \centering
  \footnotesize
  \begin{tabularx}{1\textwidth}{|l|l|X|}
    \hline
    \multicolumn{3}{c}{\textbf{Capability Metadata}} \\
    \hline
    \textbf{Name} & \textbf{Value} & \textbf{Comment and/or Description} \\
    \hline
    \textbf{Reference Id} & \texttt{BC-F6DF2559} & \\
    \hline
    \textbf{Status} & Established working hierarchy & \\
    \hline
    \textbf{Ancestry} &  & \\
    \hline
    \textbf{Source} &  & \\
    \hline
    \textbf{Enabled By} &  & \\
    \hline
    \textbf{Supports} &  & \\
    \hline
    \textbf{Derived From} &  & \\
    \hline
    \textbf{Reviewed Status} & Review meaning, business boundary, information context and any unresolved ancestry. & \\
    \hline
  \end{tabularx}
\end{table}

\begin{table}[htbp]
  \centering
  \footnotesize
  \begin{tabularx}{1\textwidth}{|l|l|X|}
    \hline
    \multicolumn{3}{c}{\textbf{Contained Features}} \\
    \hline
    \textbf{Feature Name} & \textbf{Feature Id} & \textbf{Comment} \\
    \hline
    \multicolumn{3}{|l|}{\emph{No contained Features are currently identified.}} \\
    \hline
  \end{tabularx}
\end{table}

\textbf{Intent:} Health Information Sharing Client-Support identifies the organisational ability to undertake the stated health-service responsibility. The capability describes what the enterprise must be able to do; it does not assert that Harmonia performs or owns that activity.

\textbf{Boundary:} Harmonia relevance may arise because this activity establishes the context in which information is created, needed, interpreted, related, exchanged or presented. Responsibility for the underlying business decision or care activity remains with the appropriate actor or organisation.

\subsection{CT3 - Health Information Sharing Release Management}\label{cap:bc-93bae720}
\vspace{-0.5em}
\begin{table}[htbp]
  \centering
  \footnotesize
  \begin{tabularx}{1\textwidth}{|l|l|X|}
    \hline
    \multicolumn{3}{c}{\textbf{Capability Metadata}} \\
    \hline
    \textbf{Name} & \textbf{Value} & \textbf{Comment and/or Description} \\
    \hline
    \textbf{Reference Id} & \texttt{BC-93BAE720} & \\
    \hline
    \textbf{Status} & Established working hierarchy & \\
    \hline
    \textbf{Ancestry} &  & \\
    \hline
    \textbf{Source} &  & \\
    \hline
    \textbf{Enabled By} &  & \\
    \hline
    \textbf{Supports} &  & \\
    \hline
    \textbf{Derived From} &  & \\
    \hline
    \textbf{Reviewed Status} & Review meaning, business boundary, information context and any unresolved ancestry. & \\
    \hline
  \end{tabularx}
\end{table}

\begin{table}[htbp]
  \centering
  \footnotesize
  \begin{tabularx}{1\textwidth}{|l|l|X|}
    \hline
    \multicolumn{3}{c}{\textbf{Contained Features}} \\
    \hline
    \textbf{Feature Name} & \textbf{Feature Id} & \textbf{Comment} \\
    \hline
    \multicolumn{3}{|l|}{\emph{No contained Features are currently identified.}} \\
    \hline
  \end{tabularx}
\end{table}

\textbf{Intent:} Health Information Sharing Release Management identifies the organisational ability to undertake the stated health-service responsibility. The capability describes what the enterprise must be able to do; it does not assert that Harmonia performs or owns that activity.

\textbf{Boundary:} Harmonia relevance may arise because this activity establishes the context in which information is created, needed, interpreted, related, exchanged or presented. Responsibility for the underlying business decision or care activity remains with the appropriate actor or organisation.

\section{CT2 - Digital Communication Channel}\label{cap:bc-5fd7002b}
\vspace{-0.5em}
\begin{table}[htbp]
  \centering
  \footnotesize
  \begin{tabularx}{1\textwidth}{|l|l|X|}
    \hline
    \multicolumn{3}{c}{\textbf{Capability Metadata}} \\
    \hline
    \textbf{Name} & \textbf{Value} & \textbf{Comment and/or Description} \\
    \hline
    \textbf{Reference Id} & \texttt{BC-5FD7002B} & \\
    \hline
    \textbf{Status} & Established working hierarchy & \\
    \hline
    \textbf{Ancestry} &  & \\
    \hline
    \textbf{Source} &  & \\
    \hline
    \textbf{Enabled By} &  & \\
    \hline
    \textbf{Supports} &  & \\
    \hline
    \textbf{Derived From} &  & \\
    \hline
    \textbf{Reviewed Status} &  & \\
    \hline
  \end{tabularx}
\end{table}

\begin{table}[htbp]
  \centering
  \footnotesize
  \begin{tabularx}{1\textwidth}{|l|l|X|}
    \hline
    \multicolumn{3}{c}{\textbf{Contained Capabilities}} \\
    \hline
    \textbf{Capability Name} & \textbf{Capability Id} & \textbf{Comment} \\
    \hline
    Clinical Message Conveyance & \texttt{cap:bc-1fa8828a} (at \ref{cap:bc-1fa8828a}) & \\
    \hline
    Health Event Notification & \texttt{cap:bc-baeb9e26} (at \ref{cap:bc-baeb9e26}) & \\
    \hline
    Health Service Collaboration & \texttt{cap:bc-b4211fe9} (at \ref{cap:bc-b4211fe9}) & \\
    \hline
  \end{tabularx}
\end{table}

\textbf{Intent:} Digital Communication Channel identifies the organisational ability to undertake the stated health-service responsibility. The capability describes what the enterprise must be able to do; it does not assert that Harmonia performs or owns that activity.

\textbf{Boundary:} Harmonia relevance may arise because this activity establishes the context in which information is created, needed, interpreted, related, exchanged or presented. Responsibility for the underlying business decision or care activity remains with the appropriate actor or organisation.

\subsection{CT3 - Clinical Message Conveyance}\label{cap:bc-1fa8828a}
\vspace{-0.5em}
\begin{table}[htbp]
  \centering
  \footnotesize
  \begin{tabularx}{1\textwidth}{|l|l|X|}
    \hline
    \multicolumn{3}{c}{\textbf{Capability Metadata}} \\
    \hline
    \textbf{Name} & \textbf{Value} & \textbf{Comment and/or Description} \\
    \hline
    \textbf{Reference Id} & \texttt{BC-1FA8828A} & \\
    \hline
    \textbf{Status} & Established working hierarchy & \\
    \hline
    \textbf{Ancestry} &  & \\
    \hline
    \textbf{Source} &  & \\
    \hline
    \textbf{Enabled By} &  & \\
    \hline
    \textbf{Supports} &  & \\
    \hline
    \textbf{Derived From} &  & \\
    \hline
    \textbf{Reviewed Status} & Review meaning, business boundary, information context and any unresolved ancestry. & \\
    \hline
  \end{tabularx}
\end{table}

\begin{table}[htbp]
  \centering
  \footnotesize
  \begin{tabularx}{1\textwidth}{|l|l|X|}
    \hline
    \multicolumn{3}{c}{\textbf{Contained Features}} \\
    \hline
    \textbf{Feature Name} & \textbf{Feature Id} & \textbf{Comment} \\
    \hline
    \multicolumn{3}{|l|}{\emph{No contained Features are currently identified.}} \\
    \hline
  \end{tabularx}
\end{table}

\textbf{Intent:} Clinical Message Conveyance identifies the organisational ability to undertake the stated health-service responsibility. The capability describes what the enterprise must be able to do; it does not assert that Harmonia performs or owns that activity.

\textbf{Boundary:} Harmonia relevance may arise because this activity establishes the context in which information is created, needed, interpreted, related, exchanged or presented. Responsibility for the underlying business decision or care activity remains with the appropriate actor or organisation.

\subsection{CT3 - Health Event Notification}\label{cap:bc-baeb9e26}
\vspace{-0.5em}
\begin{table}[htbp]
  \centering
  \footnotesize
  \begin{tabularx}{1\textwidth}{|l|l|X|}
    \hline
    \multicolumn{3}{c}{\textbf{Capability Metadata}} \\
    \hline
    \textbf{Name} & \textbf{Value} & \textbf{Comment and/or Description} \\
    \hline
    \textbf{Reference Id} & \texttt{BC-BAEB9E26} & \\
    \hline
    \textbf{Status} & Established working hierarchy & \\
    \hline
    \textbf{Ancestry} &  & \\
    \hline
    \textbf{Source} &  & \\
    \hline
    \textbf{Enabled By} &  & \\
    \hline
    \textbf{Supports} &  & \\
    \hline
    \textbf{Derived From} &  & \\
    \hline
    \textbf{Reviewed Status} & Review meaning, business boundary, information context and any unresolved ancestry. & \\
    \hline
  \end{tabularx}
\end{table}

\begin{table}[htbp]
  \centering
  \footnotesize
  \begin{tabularx}{1\textwidth}{|l|l|X|}
    \hline
    \multicolumn{3}{c}{\textbf{Contained Features}} \\
    \hline
    \textbf{Feature Name} & \textbf{Feature Id} & \textbf{Comment} \\
    \hline
    \multicolumn{3}{|l|}{\emph{No contained Features are currently identified.}} \\
    \hline
  \end{tabularx}
\end{table}

\textbf{Intent:} Health Event Notification identifies the organisational ability to undertake the stated health-service responsibility. The capability describes what the enterprise must be able to do; it does not assert that Harmonia performs or owns that activity.

\textbf{Boundary:} Harmonia relevance may arise because this activity establishes the context in which information is created, needed, interpreted, related, exchanged or presented. Responsibility for the underlying business decision or care activity remains with the appropriate actor or organisation.

\subsection{CT3 - Health Service Collaboration}\label{cap:bc-b4211fe9}
\vspace{-0.5em}
\begin{table}[htbp]
  \centering
  \footnotesize
  \begin{tabularx}{1\textwidth}{|l|l|X|}
    \hline
    \multicolumn{3}{c}{\textbf{Capability Metadata}} \\
    \hline
    \textbf{Name} & \textbf{Value} & \textbf{Comment and/or Description} \\
    \hline
    \textbf{Reference Id} & \texttt{BC-B4211FE9} & \\
    \hline
    \textbf{Status} & Established working hierarchy & \\
    \hline
    \textbf{Ancestry} &  & \\
    \hline
    \textbf{Source} &  & \\
    \hline
    \textbf{Enabled By} &  & \\
    \hline
    \textbf{Supports} &  & \\
    \hline
    \textbf{Derived From} &  & \\
    \hline
    \textbf{Reviewed Status} & Review meaning, business boundary, information context and any unresolved ancestry. & \\
    \hline
  \end{tabularx}
\end{table}

\begin{table}[htbp]
  \centering
  \footnotesize
  \begin{tabularx}{1\textwidth}{|l|l|X|}
    \hline
    \multicolumn{3}{c}{\textbf{Contained Features}} \\
    \hline
    \textbf{Feature Name} & \textbf{Feature Id} & \textbf{Comment} \\
    \hline
    \multicolumn{3}{|l|}{\emph{No contained Features are currently identified.}} \\
    \hline
  \end{tabularx}
\end{table}

\textbf{Intent:} Health Service Collaboration identifies the organisational ability to undertake the stated health-service responsibility. The capability describes what the enterprise must be able to do; it does not assert that Harmonia performs or owns that activity.

\textbf{Boundary:} Harmonia relevance may arise because this activity establishes the context in which information is created, needed, interpreted, related, exchanged or presented. Responsibility for the underlying business decision or care activity remains with the appropriate actor or organisation.

\section{CT2 - Domain Health Information Services — ADJUDICATED}\label{cap:bc-1079b553}
\vspace{-0.5em}
\begin{table}[htbp]
  \centering
  \footnotesize
  \begin{tabularx}{1\textwidth}{|l|l|X|}
    \hline
    \multicolumn{3}{c}{\textbf{Capability Metadata}} \\
    \hline
    \textbf{Name} & \textbf{Value} & \textbf{Comment and/or Description} \\
    \hline
    \textbf{Reference Id} & \texttt{BC-1079B553} & \\
    \hline
    \textbf{Status} & Established working hierarchy & \\
    \hline
    \textbf{Ancestry} &  & \\
    \hline
    \textbf{Source} &  & \\
    \hline
    \textbf{Enabled By} &  & \\
    \hline
    \textbf{Supports} &  & \\
    \hline
    \textbf{Derived From} &  & \\
    \hline
    \textbf{Reviewed Status} &  & \\
    \hline
  \end{tabularx}
\end{table}

\begin{table}[htbp]
  \centering
  \footnotesize
  \begin{tabularx}{1\textwidth}{|l|l|X|}
    \hline
    \multicolumn{3}{c}{\textbf{Contained Capabilities}} \\
    \hline
    \textbf{Capability Name} & \textbf{Capability Id} & \textbf{Comment} \\
    \hline
    Provider Directory Services & \texttt{cap:bc-34251ab8} (at \ref{cap:bc-34251ab8}) & \\
    \hline
    Enterprise Master Patient Index Services & \texttt{cap:bc-667fbaa5} (at \ref{cap:bc-667fbaa5}) & \\
    \hline
    Immunisation Register Services & \texttt{cap:bc-e3f3d922} (at \ref{cap:bc-e3f3d922}) & \\
    \hline
    Longitudinal Health Record Services & \texttt{cap:bc-edb9a9d1} (at \ref{cap:bc-edb9a9d1}) & \\
    \hline
    Collaboration Services & \texttt{cap:bc-a60c3d92} (at \ref{cap:bc-a60c3d92}) & \\
    \hline
  \end{tabularx}
\end{table}

\textbf{Intent:} Domain Health Information Services — ADJUDICATED identifies the organisational ability to undertake the stated health-service responsibility. The capability describes what the enterprise must be able to do; it does not assert that Harmonia performs or owns that activity.

\textbf{Boundary:} Harmonia relevance may arise because this activity establishes the context in which information is created, needed, interpreted, related, exchanged or presented. Responsibility for the underlying business decision or care activity remains with the appropriate actor or organisation.

\subsection{CT3 - Provider Directory Services}\label{cap:bc-34251ab8}
\vspace{-0.5em}
\begin{table}[htbp]
  \centering
  \footnotesize
  \begin{tabularx}{1\textwidth}{|l|l|X|}
    \hline
    \multicolumn{3}{c}{\textbf{Capability Metadata}} \\
    \hline
    \textbf{Name} & \textbf{Value} & \textbf{Comment and/or Description} \\
    \hline
    \textbf{Reference Id} & \texttt{BC-34251AB8} & \\
    \hline
    \textbf{Status} & Established working hierarchy & \\
    \hline
    \textbf{Ancestry} &  & \\
    \hline
    \textbf{Source} &  & \\
    \hline
    \textbf{Enabled By} &  & \\
    \hline
    \textbf{Supports} &  & \\
    \hline
    \textbf{Derived From} &  & \\
    \hline
    \textbf{Reviewed Status} & Review meaning, business boundary, information context and any unresolved ancestry. & \\
    \hline
  \end{tabularx}
\end{table}

\begin{table}[htbp]
  \centering
  \footnotesize
  \begin{tabularx}{1\textwidth}{|l|l|X|}
    \hline
    \multicolumn{3}{c}{\textbf{Contained Features}} \\
    \hline
    \textbf{Feature Name} & \textbf{Feature Id} & \textbf{Comment} \\
    \hline
    \multicolumn{3}{|l|}{\emph{No contained Features are currently identified.}} \\
    \hline
  \end{tabularx}
\end{table}

\textbf{Intent:} Provider Directory Services identifies the organisational ability to undertake the stated health-service responsibility. The capability describes what the enterprise must be able to do; it does not assert that Harmonia performs or owns that activity.

\textbf{Boundary:} Harmonia relevance may arise because this activity establishes the context in which information is created, needed, interpreted, related, exchanged or presented. Responsibility for the underlying business decision or care activity remains with the appropriate actor or organisation.

\subsection{CT3 - Enterprise Master Patient Index Services}\label{cap:bc-667fbaa5}
\vspace{-0.5em}
\begin{table}[htbp]
  \centering
  \footnotesize
  \begin{tabularx}{1\textwidth}{|l|l|X|}
    \hline
    \multicolumn{3}{c}{\textbf{Capability Metadata}} \\
    \hline
    \textbf{Name} & \textbf{Value} & \textbf{Comment and/or Description} \\
    \hline
    \textbf{Reference Id} & \texttt{BC-667FBAA5} & \\
    \hline
    \textbf{Status} & Established working hierarchy & \\
    \hline
    \textbf{Ancestry} &  & \\
    \hline
    \textbf{Source} &  & \\
    \hline
    \textbf{Enabled By} &  & \\
    \hline
    \textbf{Supports} &  & \\
    \hline
    \textbf{Derived From} &  & \\
    \hline
    \textbf{Reviewed Status} & Review meaning, business boundary, information context and any unresolved ancestry. & \\
    \hline
  \end{tabularx}
\end{table}

\begin{table}[htbp]
  \centering
  \footnotesize
  \begin{tabularx}{1\textwidth}{|l|l|X|}
    \hline
    \multicolumn{3}{c}{\textbf{Contained Features}} \\
    \hline
    \textbf{Feature Name} & \textbf{Feature Id} & \textbf{Comment} \\
    \hline
    \multicolumn{3}{|l|}{\emph{No contained Features are currently identified.}} \\
    \hline
  \end{tabularx}
\end{table}

\textbf{Intent:} Enterprise Master Patient Index Services identifies the organisational ability to undertake the stated health-service responsibility. The capability describes what the enterprise must be able to do; it does not assert that Harmonia performs or owns that activity.

\textbf{Boundary:} Harmonia relevance may arise because this activity establishes the context in which information is created, needed, interpreted, related, exchanged or presented. Responsibility for the underlying business decision or care activity remains with the appropriate actor or organisation.

\subsection{CT3 - Immunisation Register Services}\label{cap:bc-e3f3d922}
\vspace{-0.5em}
\begin{table}[htbp]
  \centering
  \footnotesize
  \begin{tabularx}{1\textwidth}{|l|l|X|}
    \hline
    \multicolumn{3}{c}{\textbf{Capability Metadata}} \\
    \hline
    \textbf{Name} & \textbf{Value} & \textbf{Comment and/or Description} \\
    \hline
    \textbf{Reference Id} & \texttt{BC-E3F3D922} & \\
    \hline
    \textbf{Status} & Established working hierarchy & \\
    \hline
    \textbf{Ancestry} &  & \\
    \hline
    \textbf{Source} &  & \\
    \hline
    \textbf{Enabled By} &  & \\
    \hline
    \textbf{Supports} &  & \\
    \hline
    \textbf{Derived From} &  & \\
    \hline
    \textbf{Reviewed Status} & Review meaning, business boundary, information context and any unresolved ancestry. & \\
    \hline
  \end{tabularx}
\end{table}

\begin{table}[htbp]
  \centering
  \footnotesize
  \begin{tabularx}{1\textwidth}{|l|l|X|}
    \hline
    \multicolumn{3}{c}{\textbf{Contained Features}} \\
    \hline
    \textbf{Feature Name} & \textbf{Feature Id} & \textbf{Comment} \\
    \hline
    \multicolumn{3}{|l|}{\emph{No contained Features are currently identified.}} \\
    \hline
  \end{tabularx}
\end{table}

\textbf{Intent:} Immunisation Register Services identifies the organisational ability to undertake the stated health-service responsibility. The capability describes what the enterprise must be able to do; it does not assert that Harmonia performs or owns that activity.

\textbf{Boundary:} Harmonia relevance may arise because this activity establishes the context in which information is created, needed, interpreted, related, exchanged or presented. Responsibility for the underlying business decision or care activity remains with the appropriate actor or organisation.

\subsection{CT3 - Longitudinal Health Record Services}\label{cap:bc-edb9a9d1}
\vspace{-0.5em}
\begin{table}[htbp]
  \centering
  \footnotesize
  \begin{tabularx}{1\textwidth}{|l|l|X|}
    \hline
    \multicolumn{3}{c}{\textbf{Capability Metadata}} \\
    \hline
    \textbf{Name} & \textbf{Value} & \textbf{Comment and/or Description} \\
    \hline
    \textbf{Reference Id} & \texttt{BC-EDB9A9D1} & \\
    \hline
    \textbf{Status} & Established working hierarchy & \\
    \hline
    \textbf{Ancestry} &  & \\
    \hline
    \textbf{Source} &  & \\
    \hline
    \textbf{Enabled By} &  & \\
    \hline
    \textbf{Supports} &  & \\
    \hline
    \textbf{Derived From} &  & \\
    \hline
    \textbf{Reviewed Status} & Review meaning, business boundary, information context and any unresolved ancestry. & \\
    \hline
  \end{tabularx}
\end{table}

\begin{table}[htbp]
  \centering
  \footnotesize
  \begin{tabularx}{1\textwidth}{|l|l|X|}
    \hline
    \multicolumn{3}{c}{\textbf{Contained Features}} \\
    \hline
    \textbf{Feature Name} & \textbf{Feature Id} & \textbf{Comment} \\
    \hline
    \multicolumn{3}{|l|}{\emph{No contained Features are currently identified.}} \\
    \hline
  \end{tabularx}
\end{table}

\textbf{Intent:} Longitudinal Health Record Services identifies the organisational ability to undertake the stated health-service responsibility. The capability describes what the enterprise must be able to do; it does not assert that Harmonia performs or owns that activity.

\textbf{Boundary:} Harmonia relevance may arise because this activity establishes the context in which information is created, needed, interpreted, related, exchanged or presented. Responsibility for the underlying business decision or care activity remains with the appropriate actor or organisation.

\subsection{CT3 - Collaboration Services}\label{cap:bc-a60c3d92}
\vspace{-0.5em}
\begin{table}[htbp]
  \centering
  \footnotesize
  \begin{tabularx}{1\textwidth}{|l|l|X|}
    \hline
    \multicolumn{3}{c}{\textbf{Capability Metadata}} \\
    \hline
    \textbf{Name} & \textbf{Value} & \textbf{Comment and/or Description} \\
    \hline
    \textbf{Reference Id} & \texttt{BC-A60C3D92} & \\
    \hline
    \textbf{Status} & Established working hierarchy & \\
    \hline
    \textbf{Ancestry} &  & \\
    \hline
    \textbf{Source} &  & \\
    \hline
    \textbf{Enabled By} &  & \\
    \hline
    \textbf{Supports} &  & \\
    \hline
    \textbf{Derived From} &  & \\
    \hline
    \textbf{Reviewed Status} & Review meaning, business boundary, information context and any unresolved ancestry. & \\
    \hline
  \end{tabularx}
\end{table}

\begin{table}[htbp]
  \centering
  \footnotesize
  \begin{tabularx}{1\textwidth}{|l|l|X|}
    \hline
    \multicolumn{3}{c}{\textbf{Contained Features}} \\
    \hline
    \textbf{Feature Name} & \textbf{Feature Id} & \textbf{Comment} \\
    \hline
    \multicolumn{3}{|l|}{\emph{No contained Features are currently identified.}} \\
    \hline
  \end{tabularx}
\end{table}

\textbf{Intent:} Collaboration Services identifies the organisational ability to undertake the stated health-service responsibility. The capability describes what the enterprise must be able to do; it does not assert that Harmonia performs or owns that activity.

\textbf{Boundary:} Harmonia relevance may arise because this activity establishes the context in which information is created, needed, interpreted, related, exchanged or presented. Responsibility for the underlying business decision or care activity remains with the appropriate actor or organisation.

\chapter{CT1 - Security, Privacy \& Digital Trust}\label{cap:bc-3525b6fe}
\noindent\textbf{Reference ID:} \texttt{BC-3525B6FE}\quad \textbf{Status:} Established working hierarchy\par
\noindent\textit{Source: Canonical Business Capability model / working map}\par
\textbf{Intent:} Security, Privacy \& Digital Trust identifies the organisational ability to undertake the stated health-service responsibility. The capability describes what the enterprise must be able to do; it does not assert that Harmonia performs or owns that activity.

\textbf{Boundary:} Harmonia relevance may arise because this activity establishes the context in which information is created, needed, interpreted, related, exchanged or presented. Responsibility for the underlying business decision or care activity remains with the appropriate actor or organisation.

\reviewstatus{Check scope and information relevance; no business authority transfers to Harmonia.}
\section{CT2 - Digital Identity \& Access Governance}\label{cap:bc-0d570829}
\vspace{-0.5em}
\begin{table}[htbp]
  \centering
  \footnotesize
  \begin{tabularx}{1\textwidth}{|l|l|X|}
    \hline
    \multicolumn{3}{c}{\textbf{Capability Metadata}} \\
    \hline
    \textbf{Name} & \textbf{Value} & \textbf{Comment and/or Description} \\
    \hline
    \textbf{Reference Id} & \texttt{BC-0D570829} & \\
    \hline
    \textbf{Status} & Established working hierarchy & \\
    \hline
    \textbf{Ancestry} &  & \\
    \hline
    \textbf{Source} &  & \\
    \hline
    \textbf{Enabled By} &  & \\
    \hline
    \textbf{Supports} &  & \\
    \hline
    \textbf{Derived From} &  & \\
    \hline
    \textbf{Reviewed Status} &  & \\
    \hline
  \end{tabularx}
\end{table}

\begin{table}[htbp]
  \centering
  \footnotesize
  \begin{tabularx}{1\textwidth}{|l|l|X|}
    \hline
    \multicolumn{3}{c}{\textbf{Contained Capabilities}} \\
    \hline
    \textbf{Capability Name} & \textbf{Capability Id} & \textbf{Comment} \\
    \hline
    Identity Assurance & \texttt{cap:bc-51f3c14e} (at \ref{cap:bc-51f3c14e}) & \\
    \hline
    Identity Modelling \& Interoperability & \texttt{cap:bc-e2872cbe} (at \ref{cap:bc-e2872cbe}) & \\
    \hline
    Identity Qualification Control & \texttt{cap:bc-d6f6bd86} (at \ref{cap:bc-d6f6bd86}) & \\
    \hline
    Identity Identification \& Authentication & \texttt{cap:bc-2c37a4f9} (at \ref{cap:bc-2c37a4f9}) & \\
    \hline
    Access Control Context Management & \texttt{cap:bc-92e77aa5} (at \ref{cap:bc-92e77aa5}) & \\
    \hline
    Access Control Assurance \& Reporting & \texttt{cap:bc-0e134f87} (at \ref{cap:bc-0e134f87}) & \\
    \hline
    Access Control Collaboration & \texttt{cap:bc-3a9db312} (at \ref{cap:bc-3a9db312}) & \\
    \hline
    Access Control Enforcement & \texttt{cap:bc-a78294cd} (at \ref{cap:bc-a78294cd}) & \\
    \hline
  \end{tabularx}
\end{table}

\textbf{Intent:} Digital Identity \& Access Governance identifies the organisational ability to undertake the stated health-service responsibility. The capability describes what the enterprise must be able to do; it does not assert that Harmonia performs or owns that activity.

\textbf{Boundary:} Harmonia relevance may arise because this activity establishes the context in which information is created, needed, interpreted, related, exchanged or presented. Responsibility for the underlying business decision or care activity remains with the appropriate actor or organisation.

\subsection{CT3 - Identity Assurance}\label{cap:bc-51f3c14e}
\vspace{-0.5em}
\begin{table}[htbp]
  \centering
  \footnotesize
  \begin{tabularx}{1\textwidth}{|l|l|X|}
    \hline
    \multicolumn{3}{c}{\textbf{Capability Metadata}} \\
    \hline
    \textbf{Name} & \textbf{Value} & \textbf{Comment and/or Description} \\
    \hline
    \textbf{Reference Id} & \texttt{BC-51F3C14E} & \\
    \hline
    \textbf{Status} & Established working hierarchy & \\
    \hline
    \textbf{Ancestry} &  & \\
    \hline
    \textbf{Source} &  & \\
    \hline
    \textbf{Enabled By} &  & \\
    \hline
    \textbf{Supports} &  & \\
    \hline
    \textbf{Derived From} &  & \\
    \hline
    \textbf{Reviewed Status} & Review meaning, business boundary, information context and any unresolved ancestry. & \\
    \hline
  \end{tabularx}
\end{table}

\begin{table}[htbp]
  \centering
  \footnotesize
  \begin{tabularx}{1\textwidth}{|l|l|X|}
    \hline
    \multicolumn{3}{c}{\textbf{Contained Features}} \\
    \hline
    \textbf{Feature Name} & \textbf{Feature Id} & \textbf{Comment} \\
    \hline
    \multicolumn{3}{|l|}{\emph{No contained Features are currently identified.}} \\
    \hline
  \end{tabularx}
\end{table}

\textbf{Intent:} Identity Assurance identifies the organisational ability to undertake the stated health-service responsibility. The capability describes what the enterprise must be able to do; it does not assert that Harmonia performs or owns that activity.

\textbf{Boundary:} Harmonia relevance may arise because this activity establishes the context in which information is created, needed, interpreted, related, exchanged or presented. Responsibility for the underlying business decision or care activity remains with the appropriate actor or organisation.

\subsection{CT3 - Identity Modelling \& Interoperability}\label{cap:bc-e2872cbe}
\vspace{-0.5em}
\begin{table}[htbp]
  \centering
  \footnotesize
  \begin{tabularx}{1\textwidth}{|l|l|X|}
    \hline
    \multicolumn{3}{c}{\textbf{Capability Metadata}} \\
    \hline
    \textbf{Name} & \textbf{Value} & \textbf{Comment and/or Description} \\
    \hline
    \textbf{Reference Id} & \texttt{BC-E2872CBE} & \\
    \hline
    \textbf{Status} & Established working hierarchy & \\
    \hline
    \textbf{Ancestry} &  & \\
    \hline
    \textbf{Source} &  & \\
    \hline
    \textbf{Enabled By} &  & \\
    \hline
    \textbf{Supports} &  & \\
    \hline
    \textbf{Derived From} &  & \\
    \hline
    \textbf{Reviewed Status} & Review meaning, business boundary, information context and any unresolved ancestry. & \\
    \hline
  \end{tabularx}
\end{table}

\begin{table}[htbp]
  \centering
  \footnotesize
  \begin{tabularx}{1\textwidth}{|l|l|X|}
    \hline
    \multicolumn{3}{c}{\textbf{Contained Features}} \\
    \hline
    \textbf{Feature Name} & \textbf{Feature Id} & \textbf{Comment} \\
    \hline
    \multicolumn{3}{|l|}{\emph{No contained Features are currently identified.}} \\
    \hline
  \end{tabularx}
\end{table}

\textbf{Intent:} Identity Modelling \& Interoperability identifies the organisational ability to undertake the stated health-service responsibility. The capability describes what the enterprise must be able to do; it does not assert that Harmonia performs or owns that activity.

\textbf{Boundary:} Harmonia relevance may arise because this activity establishes the context in which information is created, needed, interpreted, related, exchanged or presented. Responsibility for the underlying business decision or care activity remains with the appropriate actor or organisation.

\subsection{CT3 - Identity Qualification Control}\label{cap:bc-d6f6bd86}
\vspace{-0.5em}
\begin{table}[htbp]
  \centering
  \footnotesize
  \begin{tabularx}{1\textwidth}{|l|l|X|}
    \hline
    \multicolumn{3}{c}{\textbf{Capability Metadata}} \\
    \hline
    \textbf{Name} & \textbf{Value} & \textbf{Comment and/or Description} \\
    \hline
    \textbf{Reference Id} & \texttt{BC-D6F6BD86} & \\
    \hline
    \textbf{Status} & Established working hierarchy & \\
    \hline
    \textbf{Ancestry} &  & \\
    \hline
    \textbf{Source} &  & \\
    \hline
    \textbf{Enabled By} &  & \\
    \hline
    \textbf{Supports} &  & \\
    \hline
    \textbf{Derived From} &  & \\
    \hline
    \textbf{Reviewed Status} & Review meaning, business boundary, information context and any unresolved ancestry. & \\
    \hline
  \end{tabularx}
\end{table}

\begin{table}[htbp]
  \centering
  \footnotesize
  \begin{tabularx}{1\textwidth}{|l|l|X|}
    \hline
    \multicolumn{3}{c}{\textbf{Contained Features}} \\
    \hline
    \textbf{Feature Name} & \textbf{Feature Id} & \textbf{Comment} \\
    \hline
    \multicolumn{3}{|l|}{\emph{No contained Features are currently identified.}} \\
    \hline
  \end{tabularx}
\end{table}

\textbf{Intent:} Identity Qualification Control identifies the organisational ability to undertake the stated health-service responsibility. The capability describes what the enterprise must be able to do; it does not assert that Harmonia performs or owns that activity.

\textbf{Boundary:} Harmonia relevance may arise because this activity establishes the context in which information is created, needed, interpreted, related, exchanged or presented. Responsibility for the underlying business decision or care activity remains with the appropriate actor or organisation.

\subsection{CT3 - Identity Identification \& Authentication}\label{cap:bc-2c37a4f9}
\vspace{-0.5em}
\begin{table}[htbp]
  \centering
  \footnotesize
  \begin{tabularx}{1\textwidth}{|l|l|X|}
    \hline
    \multicolumn{3}{c}{\textbf{Capability Metadata}} \\
    \hline
    \textbf{Name} & \textbf{Value} & \textbf{Comment and/or Description} \\
    \hline
    \textbf{Reference Id} & \texttt{BC-2C37A4F9} & \\
    \hline
    \textbf{Status} & Established working hierarchy & \\
    \hline
    \textbf{Ancestry} &  & \\
    \hline
    \textbf{Source} &  & \\
    \hline
    \textbf{Enabled By} &  & \\
    \hline
    \textbf{Supports} &  & \\
    \hline
    \textbf{Derived From} &  & \\
    \hline
    \textbf{Reviewed Status} & Review meaning, business boundary, information context and any unresolved ancestry. & \\
    \hline
  \end{tabularx}
\end{table}

\begin{table}[htbp]
  \centering
  \footnotesize
  \begin{tabularx}{1\textwidth}{|l|l|X|}
    \hline
    \multicolumn{3}{c}{\textbf{Contained Features}} \\
    \hline
    \textbf{Feature Name} & \textbf{Feature Id} & \textbf{Comment} \\
    \hline
    \multicolumn{3}{|l|}{\emph{No contained Features are currently identified.}} \\
    \hline
  \end{tabularx}
\end{table}

\textbf{Intent:} Identity Identification \& Authentication identifies the organisational ability to undertake the stated health-service responsibility. The capability describes what the enterprise must be able to do; it does not assert that Harmonia performs or owns that activity.

\textbf{Boundary:} Harmonia relevance may arise because this activity establishes the context in which information is created, needed, interpreted, related, exchanged or presented. Responsibility for the underlying business decision or care activity remains with the appropriate actor or organisation.

\subsection{CT3 - Access Control Context Management}\label{cap:bc-92e77aa5}
\vspace{-0.5em}
\begin{table}[htbp]
  \centering
  \footnotesize
  \begin{tabularx}{1\textwidth}{|l|l|X|}
    \hline
    \multicolumn{3}{c}{\textbf{Capability Metadata}} \\
    \hline
    \textbf{Name} & \textbf{Value} & \textbf{Comment and/or Description} \\
    \hline
    \textbf{Reference Id} & \texttt{BC-92E77AA5} & \\
    \hline
    \textbf{Status} & Established working hierarchy & \\
    \hline
    \textbf{Ancestry} &  & \\
    \hline
    \textbf{Source} &  & \\
    \hline
    \textbf{Enabled By} &  & \\
    \hline
    \textbf{Supports} &  & \\
    \hline
    \textbf{Derived From} &  & \\
    \hline
    \textbf{Reviewed Status} & Review meaning, business boundary, information context and any unresolved ancestry. & \\
    \hline
  \end{tabularx}
\end{table}

\begin{table}[htbp]
  \centering
  \footnotesize
  \begin{tabularx}{1\textwidth}{|l|l|X|}
    \hline
    \multicolumn{3}{c}{\textbf{Contained Features}} \\
    \hline
    \textbf{Feature Name} & \textbf{Feature Id} & \textbf{Comment} \\
    \hline
    \multicolumn{3}{|l|}{\emph{No contained Features are currently identified.}} \\
    \hline
  \end{tabularx}
\end{table}

\textbf{Intent:} Access Control Context Management identifies the organisational ability to undertake the stated health-service responsibility. The capability describes what the enterprise must be able to do; it does not assert that Harmonia performs or owns that activity.

\textbf{Boundary:} Harmonia relevance may arise because this activity establishes the context in which information is created, needed, interpreted, related, exchanged or presented. Responsibility for the underlying business decision or care activity remains with the appropriate actor or organisation.

\subsection{CT3 - Access Control Assurance \& Reporting}\label{cap:bc-0e134f87}
\vspace{-0.5em}
\begin{table}[htbp]
  \centering
  \footnotesize
  \begin{tabularx}{1\textwidth}{|l|l|X|}
    \hline
    \multicolumn{3}{c}{\textbf{Capability Metadata}} \\
    \hline
    \textbf{Name} & \textbf{Value} & \textbf{Comment and/or Description} \\
    \hline
    \textbf{Reference Id} & \texttt{BC-0E134F87} & \\
    \hline
    \textbf{Status} & Established working hierarchy & \\
    \hline
    \textbf{Ancestry} &  & \\
    \hline
    \textbf{Source} &  & \\
    \hline
    \textbf{Enabled By} &  & \\
    \hline
    \textbf{Supports} &  & \\
    \hline
    \textbf{Derived From} &  & \\
    \hline
    \textbf{Reviewed Status} & Review meaning, business boundary, information context and any unresolved ancestry. & \\
    \hline
  \end{tabularx}
\end{table}

\begin{table}[htbp]
  \centering
  \footnotesize
  \begin{tabularx}{1\textwidth}{|l|l|X|}
    \hline
    \multicolumn{3}{c}{\textbf{Contained Features}} \\
    \hline
    \textbf{Feature Name} & \textbf{Feature Id} & \textbf{Comment} \\
    \hline
    \multicolumn{3}{|l|}{\emph{No contained Features are currently identified.}} \\
    \hline
  \end{tabularx}
\end{table}

\textbf{Intent:} Access Control Assurance \& Reporting identifies the organisational ability to undertake the stated health-service responsibility. The capability describes what the enterprise must be able to do; it does not assert that Harmonia performs or owns that activity.

\textbf{Boundary:} Harmonia relevance may arise because this activity establishes the context in which information is created, needed, interpreted, related, exchanged or presented. Responsibility for the underlying business decision or care activity remains with the appropriate actor or organisation.

\subsection{CT3 - Access Control Collaboration}\label{cap:bc-3a9db312}
\vspace{-0.5em}
\begin{table}[htbp]
  \centering
  \footnotesize
  \begin{tabularx}{1\textwidth}{|l|l|X|}
    \hline
    \multicolumn{3}{c}{\textbf{Capability Metadata}} \\
    \hline
    \textbf{Name} & \textbf{Value} & \textbf{Comment and/or Description} \\
    \hline
    \textbf{Reference Id} & \texttt{BC-3A9DB312} & \\
    \hline
    \textbf{Status} & Established working hierarchy & \\
    \hline
    \textbf{Ancestry} &  & \\
    \hline
    \textbf{Source} &  & \\
    \hline
    \textbf{Enabled By} &  & \\
    \hline
    \textbf{Supports} &  & \\
    \hline
    \textbf{Derived From} &  & \\
    \hline
    \textbf{Reviewed Status} & Review meaning, business boundary, information context and any unresolved ancestry. & \\
    \hline
  \end{tabularx}
\end{table}

\begin{table}[htbp]
  \centering
  \footnotesize
  \begin{tabularx}{1\textwidth}{|l|l|X|}
    \hline
    \multicolumn{3}{c}{\textbf{Contained Features}} \\
    \hline
    \textbf{Feature Name} & \textbf{Feature Id} & \textbf{Comment} \\
    \hline
    \multicolumn{3}{|l|}{\emph{No contained Features are currently identified.}} \\
    \hline
  \end{tabularx}
\end{table}

\textbf{Intent:} Access Control Collaboration identifies the organisational ability to undertake the stated health-service responsibility. The capability describes what the enterprise must be able to do; it does not assert that Harmonia performs or owns that activity.

\textbf{Boundary:} Harmonia relevance may arise because this activity establishes the context in which information is created, needed, interpreted, related, exchanged or presented. Responsibility for the underlying business decision or care activity remains with the appropriate actor or organisation.

\subsection{CT3 - Access Control Enforcement}\label{cap:bc-a78294cd}
\vspace{-0.5em}
\begin{table}[htbp]
  \centering
  \footnotesize
  \begin{tabularx}{1\textwidth}{|l|l|X|}
    \hline
    \multicolumn{3}{c}{\textbf{Capability Metadata}} \\
    \hline
    \textbf{Name} & \textbf{Value} & \textbf{Comment and/or Description} \\
    \hline
    \textbf{Reference Id} & \texttt{BC-A78294CD} & \\
    \hline
    \textbf{Status} & Established working hierarchy & \\
    \hline
    \textbf{Ancestry} &  & \\
    \hline
    \textbf{Source} &  & \\
    \hline
    \textbf{Enabled By} &  & \\
    \hline
    \textbf{Supports} &  & \\
    \hline
    \textbf{Derived From} &  & \\
    \hline
    \textbf{Reviewed Status} & Review meaning, business boundary, information context and any unresolved ancestry. & \\
    \hline
  \end{tabularx}
\end{table}

\begin{table}[htbp]
  \centering
  \footnotesize
  \begin{tabularx}{1\textwidth}{|l|l|X|}
    \hline
    \multicolumn{3}{c}{\textbf{Contained Features}} \\
    \hline
    \textbf{Feature Name} & \textbf{Feature Id} & \textbf{Comment} \\
    \hline
    \multicolumn{3}{|l|}{\emph{No contained Features are currently identified.}} \\
    \hline
  \end{tabularx}
\end{table}

\textbf{Intent:} Access Control Enforcement identifies the organisational ability to undertake the stated health-service responsibility. The capability describes what the enterprise must be able to do; it does not assert that Harmonia performs or owns that activity.

\textbf{Boundary:} Harmonia relevance may arise because this activity establishes the context in which information is created, needed, interpreted, related, exchanged or presented. Responsibility for the underlying business decision or care activity remains with the appropriate actor or organisation.

\section{CT2 - Health Information Protection}\label{cap:bc-448ee064}
\vspace{-0.5em}
\begin{table}[htbp]
  \centering
  \footnotesize
  \begin{tabularx}{1\textwidth}{|l|l|X|}
    \hline
    \multicolumn{3}{c}{\textbf{Capability Metadata}} \\
    \hline
    \textbf{Name} & \textbf{Value} & \textbf{Comment and/or Description} \\
    \hline
    \textbf{Reference Id} & \texttt{BC-448EE064} & \\
    \hline
    \textbf{Status} & Established working hierarchy & \\
    \hline
    \textbf{Ancestry} &  & \\
    \hline
    \textbf{Source} &  & \\
    \hline
    \textbf{Enabled By} &  & \\
    \hline
    \textbf{Supports} &  & \\
    \hline
    \textbf{Derived From} &  & \\
    \hline
    \textbf{Reviewed Status} &  & \\
    \hline
  \end{tabularx}
\end{table}

\begin{table}[htbp]
  \centering
  \footnotesize
  \begin{tabularx}{1\textwidth}{|l|l|X|}
    \hline
    \multicolumn{3}{c}{\textbf{Contained Capabilities}} \\
    \hline
    \textbf{Capability Name} & \textbf{Capability Id} & \textbf{Comment} \\
    \hline
    Health Privacy Safeguard Management & \texttt{cap:bc-f8de8e42} (at \ref{cap:bc-f8de8e42}) & \\
    \hline
    Health Information Safeguarding & \texttt{cap:bc-cbd883da} (at \ref{cap:bc-cbd883da}) & \\
    \hline
    Health Information Access Monitoring & \texttt{cap:bc-88ad2c7f} (at \ref{cap:bc-88ad2c7f}) & \\
    \hline
    Health Information Compartment Assurance & \texttt{cap:bc-69a112eb} (at \ref{cap:bc-69a112eb}) & \\
    \hline
    Health Information Classification Assurance & \texttt{cap:bc-c163719d} (at \ref{cap:bc-c163719d}) & \\
    \hline
  \end{tabularx}
\end{table}

\textbf{Intent:} Health Information Protection identifies the organisational ability to undertake the stated health-service responsibility. The capability describes what the enterprise must be able to do; it does not assert that Harmonia performs or owns that activity.

\textbf{Boundary:} Harmonia relevance may arise because this activity establishes the context in which information is created, needed, interpreted, related, exchanged or presented. Responsibility for the underlying business decision or care activity remains with the appropriate actor or organisation.

\subsection{CT3 - Health Privacy Safeguard Management}\label{cap:bc-f8de8e42}
\vspace{-0.5em}
\begin{table}[htbp]
  \centering
  \footnotesize
  \begin{tabularx}{1\textwidth}{|l|l|X|}
    \hline
    \multicolumn{3}{c}{\textbf{Capability Metadata}} \\
    \hline
    \textbf{Name} & \textbf{Value} & \textbf{Comment and/or Description} \\
    \hline
    \textbf{Reference Id} & \texttt{BC-F8DE8E42} & \\
    \hline
    \textbf{Status} & Established working hierarchy & \\
    \hline
    \textbf{Ancestry} &  & \\
    \hline
    \textbf{Source} &  & \\
    \hline
    \textbf{Enabled By} &  & \\
    \hline
    \textbf{Supports} &  & \\
    \hline
    \textbf{Derived From} &  & \\
    \hline
    \textbf{Reviewed Status} & Review meaning, business boundary, information context and any unresolved ancestry. & \\
    \hline
  \end{tabularx}
\end{table}

\begin{table}[htbp]
  \centering
  \footnotesize
  \begin{tabularx}{1\textwidth}{|l|l|X|}
    \hline
    \multicolumn{3}{c}{\textbf{Contained Features}} \\
    \hline
    \textbf{Feature Name} & \textbf{Feature Id} & \textbf{Comment} \\
    \hline
    \multicolumn{3}{|l|}{\emph{No contained Features are currently identified.}} \\
    \hline
  \end{tabularx}
\end{table}

\textbf{Intent:} Health Privacy Safeguard Management identifies the organisational ability to undertake the stated health-service responsibility. The capability describes what the enterprise must be able to do; it does not assert that Harmonia performs or owns that activity.

\textbf{Boundary:} Harmonia relevance may arise because this activity establishes the context in which information is created, needed, interpreted, related, exchanged or presented. Responsibility for the underlying business decision or care activity remains with the appropriate actor or organisation.

\subsection{CT3 - Health Information Safeguarding}\label{cap:bc-cbd883da}
\vspace{-0.5em}
\begin{table}[htbp]
  \centering
  \footnotesize
  \begin{tabularx}{1\textwidth}{|l|l|X|}
    \hline
    \multicolumn{3}{c}{\textbf{Capability Metadata}} \\
    \hline
    \textbf{Name} & \textbf{Value} & \textbf{Comment and/or Description} \\
    \hline
    \textbf{Reference Id} & \texttt{BC-CBD883DA} & \\
    \hline
    \textbf{Status} & Established working hierarchy & \\
    \hline
    \textbf{Ancestry} &  & \\
    \hline
    \textbf{Source} &  & \\
    \hline
    \textbf{Enabled By} &  & \\
    \hline
    \textbf{Supports} &  & \\
    \hline
    \textbf{Derived From} &  & \\
    \hline
    \textbf{Reviewed Status} & Review meaning, business boundary, information context and any unresolved ancestry. & \\
    \hline
  \end{tabularx}
\end{table}

\begin{table}[htbp]
  \centering
  \footnotesize
  \begin{tabularx}{1\textwidth}{|l|l|X|}
    \hline
    \multicolumn{3}{c}{\textbf{Contained Features}} \\
    \hline
    \textbf{Feature Name} & \textbf{Feature Id} & \textbf{Comment} \\
    \hline
    \multicolumn{3}{|l|}{\emph{No contained Features are currently identified.}} \\
    \hline
  \end{tabularx}
\end{table}

\textbf{Intent:} Health Information Safeguarding identifies the organisational ability to undertake the stated health-service responsibility. The capability describes what the enterprise must be able to do; it does not assert that Harmonia performs or owns that activity.

\textbf{Boundary:} Harmonia relevance may arise because this activity establishes the context in which information is created, needed, interpreted, related, exchanged or presented. Responsibility for the underlying business decision or care activity remains with the appropriate actor or organisation.

\subsection{CT3 - Health Information Access Monitoring}\label{cap:bc-88ad2c7f}
\vspace{-0.5em}
\begin{table}[htbp]
  \centering
  \footnotesize
  \begin{tabularx}{1\textwidth}{|l|l|X|}
    \hline
    \multicolumn{3}{c}{\textbf{Capability Metadata}} \\
    \hline
    \textbf{Name} & \textbf{Value} & \textbf{Comment and/or Description} \\
    \hline
    \textbf{Reference Id} & \texttt{BC-88AD2C7F} & \\
    \hline
    \textbf{Status} & Established working hierarchy & \\
    \hline
    \textbf{Ancestry} &  & \\
    \hline
    \textbf{Source} &  & \\
    \hline
    \textbf{Enabled By} &  & \\
    \hline
    \textbf{Supports} &  & \\
    \hline
    \textbf{Derived From} &  & \\
    \hline
    \textbf{Reviewed Status} & Review meaning, business boundary, information context and any unresolved ancestry. & \\
    \hline
  \end{tabularx}
\end{table}

\begin{table}[htbp]
  \centering
  \footnotesize
  \begin{tabularx}{1\textwidth}{|l|l|X|}
    \hline
    \multicolumn{3}{c}{\textbf{Contained Features}} \\
    \hline
    \textbf{Feature Name} & \textbf{Feature Id} & \textbf{Comment} \\
    \hline
    \multicolumn{3}{|l|}{\emph{No contained Features are currently identified.}} \\
    \hline
  \end{tabularx}
\end{table}

\textbf{Intent:} Health Information Access Monitoring identifies the organisational ability to undertake the stated health-service responsibility. The capability describes what the enterprise must be able to do; it does not assert that Harmonia performs or owns that activity.

\textbf{Boundary:} Harmonia relevance may arise because this activity establishes the context in which information is created, needed, interpreted, related, exchanged or presented. Responsibility for the underlying business decision or care activity remains with the appropriate actor or organisation.

\subsection{CT3 - Health Information Compartment Assurance}\label{cap:bc-69a112eb}
\vspace{-0.5em}
\begin{table}[htbp]
  \centering
  \footnotesize
  \begin{tabularx}{1\textwidth}{|l|l|X|}
    \hline
    \multicolumn{3}{c}{\textbf{Capability Metadata}} \\
    \hline
    \textbf{Name} & \textbf{Value} & \textbf{Comment and/or Description} \\
    \hline
    \textbf{Reference Id} & \texttt{BC-69A112EB} & \\
    \hline
    \textbf{Status} & Established working hierarchy & \\
    \hline
    \textbf{Ancestry} &  & \\
    \hline
    \textbf{Source} &  & \\
    \hline
    \textbf{Enabled By} &  & \\
    \hline
    \textbf{Supports} &  & \\
    \hline
    \textbf{Derived From} &  & \\
    \hline
    \textbf{Reviewed Status} & Review meaning, business boundary, information context and any unresolved ancestry. & \\
    \hline
  \end{tabularx}
\end{table}

\begin{table}[htbp]
  \centering
  \footnotesize
  \begin{tabularx}{1\textwidth}{|l|l|X|}
    \hline
    \multicolumn{3}{c}{\textbf{Contained Features}} \\
    \hline
    \textbf{Feature Name} & \textbf{Feature Id} & \textbf{Comment} \\
    \hline
    \multicolumn{3}{|l|}{\emph{No contained Features are currently identified.}} \\
    \hline
  \end{tabularx}
\end{table}

\textbf{Intent:} Health Information Compartment Assurance identifies the organisational ability to undertake the stated health-service responsibility. The capability describes what the enterprise must be able to do; it does not assert that Harmonia performs or owns that activity.

\textbf{Boundary:} Harmonia relevance may arise because this activity establishes the context in which information is created, needed, interpreted, related, exchanged or presented. Responsibility for the underlying business decision or care activity remains with the appropriate actor or organisation.

\subsection{CT3 - Health Information Classification Assurance}\label{cap:bc-c163719d}
\vspace{-0.5em}
\begin{table}[htbp]
  \centering
  \footnotesize
  \begin{tabularx}{1\textwidth}{|l|l|X|}
    \hline
    \multicolumn{3}{c}{\textbf{Capability Metadata}} \\
    \hline
    \textbf{Name} & \textbf{Value} & \textbf{Comment and/or Description} \\
    \hline
    \textbf{Reference Id} & \texttt{BC-C163719D} & \\
    \hline
    \textbf{Status} & Established working hierarchy & \\
    \hline
    \textbf{Ancestry} &  & \\
    \hline
    \textbf{Source} &  & \\
    \hline
    \textbf{Enabled By} &  & \\
    \hline
    \textbf{Supports} &  & \\
    \hline
    \textbf{Derived From} &  & \\
    \hline
    \textbf{Reviewed Status} & Review meaning, business boundary, information context and any unresolved ancestry. & \\
    \hline
  \end{tabularx}
\end{table}

\begin{table}[htbp]
  \centering
  \footnotesize
  \begin{tabularx}{1\textwidth}{|l|l|X|}
    \hline
    \multicolumn{3}{c}{\textbf{Contained Features}} \\
    \hline
    \textbf{Feature Name} & \textbf{Feature Id} & \textbf{Comment} \\
    \hline
    \multicolumn{3}{|l|}{\emph{No contained Features are currently identified.}} \\
    \hline
  \end{tabularx}
\end{table}

\textbf{Intent:} Health Information Classification Assurance identifies the organisational ability to undertake the stated health-service responsibility. The capability describes what the enterprise must be able to do; it does not assert that Harmonia performs or owns that activity.

\textbf{Boundary:} Harmonia relevance may arise because this activity establishes the context in which information is created, needed, interpreted, related, exchanged or presented. Responsibility for the underlying business decision or care activity remains with the appropriate actor or organisation.

\section{CT2 - Digital Accountability \& Trust}\label{cap:bc-db00a738}
\vspace{-0.5em}
\begin{table}[htbp]
  \centering
  \footnotesize
  \begin{tabularx}{1\textwidth}{|l|l|X|}
    \hline
    \multicolumn{3}{c}{\textbf{Capability Metadata}} \\
    \hline
    \textbf{Name} & \textbf{Value} & \textbf{Comment and/or Description} \\
    \hline
    \textbf{Reference Id} & \texttt{BC-DB00A738} & \\
    \hline
    \textbf{Status} & Established working hierarchy & \\
    \hline
    \textbf{Ancestry} &  & \\
    \hline
    \textbf{Source} &  & \\
    \hline
    \textbf{Enabled By} &  & \\
    \hline
    \textbf{Supports} &  & \\
    \hline
    \textbf{Derived From} &  & \\
    \hline
    \textbf{Reviewed Status} &  & \\
    \hline
  \end{tabularx}
\end{table}

\begin{table}[htbp]
  \centering
  \footnotesize
  \begin{tabularx}{1\textwidth}{|l|l|X|}
    \hline
    \multicolumn{3}{c}{\textbf{Contained Capabilities}} \\
    \hline
    \textbf{Capability Name} & \textbf{Capability Id} & \textbf{Comment} \\
    \hline
    Governed Activity Attribution & \texttt{cap:bc-b2fba48b} (at \ref{cap:bc-b2fba48b}) & \\
    \hline
    Security \& Information-Use Accountability & \texttt{cap:bc-c26e8578} (at \ref{cap:bc-c26e8578}) & \\
    \hline
  \end{tabularx}
\end{table}

\textbf{Intent:} Digital Accountability \& Trust identifies the organisational ability to undertake the stated health-service responsibility. The capability describes what the enterprise must be able to do; it does not assert that Harmonia performs or owns that activity.

\textbf{Boundary:} Harmonia relevance may arise because this activity establishes the context in which information is created, needed, interpreted, related, exchanged or presented. Responsibility for the underlying business decision or care activity remains with the appropriate actor or organisation.

\subsection{CT3 - Governed Activity Attribution}\label{cap:bc-b2fba48b}
\vspace{-0.5em}
\begin{table}[htbp]
  \centering
  \footnotesize
  \begin{tabularx}{1\textwidth}{|l|l|X|}
    \hline
    \multicolumn{3}{c}{\textbf{Capability Metadata}} \\
    \hline
    \textbf{Name} & \textbf{Value} & \textbf{Comment and/or Description} \\
    \hline
    \textbf{Reference Id} & \texttt{BC-B2FBA48B} & \\
    \hline
    \textbf{Status} & Established working hierarchy & \\
    \hline
    \textbf{Ancestry} &  & \\
    \hline
    \textbf{Source} &  & \\
    \hline
    \textbf{Enabled By} &  & \\
    \hline
    \textbf{Supports} &  & \\
    \hline
    \textbf{Derived From} &  & \\
    \hline
    \textbf{Reviewed Status} & Review meaning, business boundary, information context and any unresolved ancestry. & \\
    \hline
  \end{tabularx}
\end{table}

\begin{table}[htbp]
  \centering
  \footnotesize
  \begin{tabularx}{1\textwidth}{|l|l|X|}
    \hline
    \multicolumn{3}{c}{\textbf{Contained Features}} \\
    \hline
    \textbf{Feature Name} & \textbf{Feature Id} & \textbf{Comment} \\
    \hline
    \multicolumn{3}{|l|}{\emph{No contained Features are currently identified.}} \\
    \hline
  \end{tabularx}
\end{table}

\textbf{Intent:} Governed Activity Attribution identifies the organisational ability to undertake the stated health-service responsibility. The capability describes what the enterprise must be able to do; it does not assert that Harmonia performs or owns that activity.

\textbf{Boundary:} Harmonia relevance may arise because this activity establishes the context in which information is created, needed, interpreted, related, exchanged or presented. Responsibility for the underlying business decision or care activity remains with the appropriate actor or organisation.

\subsection{CT3 - Security \& Information-Use Accountability}\label{cap:bc-c26e8578}
\vspace{-0.5em}
\begin{table}[htbp]
  \centering
  \footnotesize
  \begin{tabularx}{1\textwidth}{|l|l|X|}
    \hline
    \multicolumn{3}{c}{\textbf{Capability Metadata}} \\
    \hline
    \textbf{Name} & \textbf{Value} & \textbf{Comment and/or Description} \\
    \hline
    \textbf{Reference Id} & \texttt{BC-C26E8578} & \\
    \hline
    \textbf{Status} & Established working hierarchy & \\
    \hline
    \textbf{Ancestry} &  & \\
    \hline
    \textbf{Source} &  & \\
    \hline
    \textbf{Enabled By} &  & \\
    \hline
    \textbf{Supports} &  & \\
    \hline
    \textbf{Derived From} &  & \\
    \hline
    \textbf{Reviewed Status} & Review meaning, business boundary, information context and any unresolved ancestry. & \\
    \hline
  \end{tabularx}
\end{table}

\begin{table}[htbp]
  \centering
  \footnotesize
  \begin{tabularx}{1\textwidth}{|l|l|X|}
    \hline
    \multicolumn{3}{c}{\textbf{Contained Features}} \\
    \hline
    \textbf{Feature Name} & \textbf{Feature Id} & \textbf{Comment} \\
    \hline
    \multicolumn{3}{|l|}{\emph{No contained Features are currently identified.}} \\
    \hline
  \end{tabularx}
\end{table}

\textbf{Intent:} Security \& Information-Use Accountability identifies the organisational ability to undertake the stated health-service responsibility. The capability describes what the enterprise must be able to do; it does not assert that Harmonia performs or owns that activity.

\textbf{Boundary:} Harmonia relevance may arise because this activity establishes the context in which information is created, needed, interpreted, related, exchanged or presented. Responsibility for the underlying business decision or care activity remains with the appropriate actor or organisation.

\chapter{CT1 - Health Research \& Innovation}\label{cap:bc-c5ed32f7}
\noindent\textbf{Reference ID:} \texttt{BC-C5ED32F7}\quad \textbf{Status:} Established working hierarchy\par
\noindent\textit{Source: Canonical Business Capability model / working map}\par
\textbf{Intent:} Health Research \& Innovation identifies the organisational ability to undertake the stated health-service responsibility. The capability describes what the enterprise must be able to do; it does not assert that Harmonia performs or owns that activity.

\textbf{Boundary:} Harmonia relevance may arise because this activity establishes the context in which information is created, needed, interpreted, related, exchanged or presented. Responsibility for the underlying business decision or care activity remains with the appropriate actor or organisation.

\reviewstatus{Check scope and information relevance; no business authority transfers to Harmonia.}
\section{CT2 - Research Information Support}\label{cap:bc-5fd1c13a}
\vspace{-0.5em}
\begin{table}[htbp]
  \centering
  \footnotesize
  \begin{tabularx}{1\textwidth}{|l|l|X|}
    \hline
    \multicolumn{3}{c}{\textbf{Capability Metadata}} \\
    \hline
    \textbf{Name} & \textbf{Value} & \textbf{Comment and/or Description} \\
    \hline
    \textbf{Reference Id} & \texttt{BC-5FD1C13A} & \\
    \hline
    \textbf{Status} & Established working hierarchy & \\
    \hline
    \textbf{Ancestry} &  & \\
    \hline
    \textbf{Source} &  & \\
    \hline
    \textbf{Enabled By} &  & \\
    \hline
    \textbf{Supports} &  & \\
    \hline
    \textbf{Derived From} &  & \\
    \hline
    \textbf{Reviewed Status} &  & \\
    \hline
  \end{tabularx}
\end{table}

\begin{table}[htbp]
  \centering
  \footnotesize
  \begin{tabularx}{1\textwidth}{|l|l|X|}
    \hline
    \multicolumn{3}{c}{\textbf{Contained Capabilities}} \\
    \hline
    \textbf{Capability Name} & \textbf{Capability Id} & \textbf{Comment} \\
    \hline
    Governed Research Information Provision & \texttt{cap:bc-3c3f6346} (at \ref{cap:bc-3c3f6346}) & \\
    \hline
  \end{tabularx}
\end{table}

\textbf{Intent:} Research Information Support identifies the organisational ability to undertake the stated health-service responsibility. The capability describes what the enterprise must be able to do; it does not assert that Harmonia performs or owns that activity.

\textbf{Boundary:} Harmonia relevance may arise because this activity establishes the context in which information is created, needed, interpreted, related, exchanged or presented. Responsibility for the underlying business decision or care activity remains with the appropriate actor or organisation.

\subsection{CT3 - Governed Research Information Provision}\label{cap:bc-3c3f6346}
\vspace{-0.5em}
\begin{table}[htbp]
  \centering
  \footnotesize
  \begin{tabularx}{1\textwidth}{|l|l|X|}
    \hline
    \multicolumn{3}{c}{\textbf{Capability Metadata}} \\
    \hline
    \textbf{Name} & \textbf{Value} & \textbf{Comment and/or Description} \\
    \hline
    \textbf{Reference Id} & \texttt{BC-3C3F6346} & \\
    \hline
    \textbf{Status} & Established working hierarchy & \\
    \hline
    \textbf{Ancestry} &  & \\
    \hline
    \textbf{Source} &  & \\
    \hline
    \textbf{Enabled By} &  & \\
    \hline
    \textbf{Supports} &  & \\
    \hline
    \textbf{Derived From} &  & \\
    \hline
    \textbf{Reviewed Status} & Review meaning, business boundary, information context and any unresolved ancestry. & \\
    \hline
  \end{tabularx}
\end{table}

\begin{table}[htbp]
  \centering
  \footnotesize
  \begin{tabularx}{1\textwidth}{|l|l|X|}
    \hline
    \multicolumn{3}{c}{\textbf{Contained Features}} \\
    \hline
    \textbf{Feature Name} & \textbf{Feature Id} & \textbf{Comment} \\
    \hline
    \multicolumn{3}{|l|}{\emph{No contained Features are currently identified.}} \\
    \hline
  \end{tabularx}
\end{table}

\textbf{Intent:} Governed Research Information Provision identifies the organisational ability to undertake the stated health-service responsibility. The capability describes what the enterprise must be able to do; it does not assert that Harmonia performs or owns that activity.

\textbf{Boundary:} Harmonia relevance may arise because this activity establishes the context in which information is created, needed, interpreted, related, exchanged or presented. Responsibility for the underlying business decision or care activity remains with the appropriate actor or organisation.

\section{CT2 - Health Innovation}\label{cap:bc-5b3ee7bd}
\vspace{-0.5em}
\begin{table}[htbp]
  \centering
  \footnotesize
  \begin{tabularx}{1\textwidth}{|l|l|X|}
    \hline
    \multicolumn{3}{c}{\textbf{Capability Metadata}} \\
    \hline
    \textbf{Name} & \textbf{Value} & \textbf{Comment and/or Description} \\
    \hline
    \textbf{Reference Id} & \texttt{BC-5B3EE7BD} & \\
    \hline
    \textbf{Status} & Established working hierarchy & \\
    \hline
    \textbf{Ancestry} &  & \\
    \hline
    \textbf{Source} &  & \\
    \hline
    \textbf{Enabled By} &  & \\
    \hline
    \textbf{Supports} &  & \\
    \hline
    \textbf{Derived From} &  & \\
    \hline
    \textbf{Reviewed Status} &  & \\
    \hline
  \end{tabularx}
\end{table}

\begin{table}[htbp]
  \centering
  \footnotesize
  \begin{tabularx}{1\textwidth}{|l|l|X|}
    \hline
    \multicolumn{3}{c}{\textbf{Contained Capabilities}} \\
    \hline
    \textbf{Capability Name} & \textbf{Capability Id} & \textbf{Comment} \\
    \hline
    Health Innovation Evaluation & \texttt{cap:bc-e804a440} (at \ref{cap:bc-e804a440}) & \\
    \hline
  \end{tabularx}
\end{table}

\textbf{Intent:} Health Innovation identifies the organisational ability to undertake the stated health-service responsibility. The capability describes what the enterprise must be able to do; it does not assert that Harmonia performs or owns that activity.

\textbf{Boundary:} Harmonia relevance may arise because this activity establishes the context in which information is created, needed, interpreted, related, exchanged or presented. Responsibility for the underlying business decision or care activity remains with the appropriate actor or organisation.

\subsection{CT3 - Health Innovation Evaluation}\label{cap:bc-e804a440}
\vspace{-0.5em}
\begin{table}[htbp]
  \centering
  \footnotesize
  \begin{tabularx}{1\textwidth}{|l|l|X|}
    \hline
    \multicolumn{3}{c}{\textbf{Capability Metadata}} \\
    \hline
    \textbf{Name} & \textbf{Value} & \textbf{Comment and/or Description} \\
    \hline
    \textbf{Reference Id} & \texttt{BC-E804A440} & \\
    \hline
    \textbf{Status} & Established working hierarchy & \\
    \hline
    \textbf{Ancestry} &  & \\
    \hline
    \textbf{Source} &  & \\
    \hline
    \textbf{Enabled By} &  & \\
    \hline
    \textbf{Supports} &  & \\
    \hline
    \textbf{Derived From} &  & \\
    \hline
    \textbf{Reviewed Status} & Review meaning, business boundary, information context and any unresolved ancestry. & \\
    \hline
  \end{tabularx}
\end{table}

\begin{table}[htbp]
  \centering
  \footnotesize
  \begin{tabularx}{1\textwidth}{|l|l|X|}
    \hline
    \multicolumn{3}{c}{\textbf{Contained Features}} \\
    \hline
    \textbf{Feature Name} & \textbf{Feature Id} & \textbf{Comment} \\
    \hline
    \multicolumn{3}{|l|}{\emph{No contained Features are currently identified.}} \\
    \hline
  \end{tabularx}
\end{table}

\textbf{Intent:} Health Innovation Evaluation identifies the organisational ability to undertake the stated health-service responsibility. The capability describes what the enterprise must be able to do; it does not assert that Harmonia performs or owns that activity.

\textbf{Boundary:} Harmonia relevance may arise because this activity establishes the context in which information is created, needed, interpreted, related, exchanged or presented. Responsibility for the underlying business decision or care activity remains with the appropriate actor or organisation.

\chapter{CT1 - Health Service Direction \& Stewardship}\label{cap:bc-a7b13259}
\noindent\textbf{Reference ID:} \texttt{BC-A7B13259}\quad \textbf{Status:} Established working hierarchy\par
\noindent\textit{Source: Canonical Business Capability model / working map}\par
\textbf{Intent:} Health Service Direction \& Stewardship identifies the organisational ability to undertake the stated health-service responsibility. The capability describes what the enterprise must be able to do; it does not assert that Harmonia performs or owns that activity.

\textbf{Boundary:} Harmonia relevance may arise because this activity establishes the context in which information is created, needed, interpreted, related, exchanged or presented. Responsibility for the underlying business decision or care activity remains with the appropriate actor or organisation.

\reviewstatus{Check scope and information relevance; no business authority transfers to Harmonia.}
\section{CT2 - Health Service Modelling}\label{cap:bc-25bc222b}
\vspace{-0.5em}
\begin{table}[htbp]
  \centering
  \footnotesize
  \begin{tabularx}{1\textwidth}{|l|l|X|}
    \hline
    \multicolumn{3}{c}{\textbf{Capability Metadata}} \\
    \hline
    \textbf{Name} & \textbf{Value} & \textbf{Comment and/or Description} \\
    \hline
    \textbf{Reference Id} & \texttt{BC-25BC222B} & \\
    \hline
    \textbf{Status} & Established working hierarchy & \\
    \hline
    \textbf{Ancestry} &  & \\
    \hline
    \textbf{Source} &  & \\
    \hline
    \textbf{Enabled By} &  & \\
    \hline
    \textbf{Supports} &  & \\
    \hline
    \textbf{Derived From} &  & \\
    \hline
    \textbf{Reviewed Status} &  & \\
    \hline
  \end{tabularx}
\end{table}

\begin{table}[htbp]
  \centering
  \footnotesize
  \begin{tabularx}{1\textwidth}{|l|l|X|}
    \hline
    \multicolumn{3}{c}{\textbf{Contained Capabilities}} \\
    \hline
    \textbf{Capability Name} & \textbf{Capability Id} & \textbf{Comment} \\
    \hline
    Health Service Need Identification & \texttt{cap:bc-df112315} (at \ref{cap:bc-df112315}) & \\
    \hline
    Health Service Capability Identification & \texttt{cap:bc-a7a90ee0} (at \ref{cap:bc-a7a90ee0}) & \\
    \hline
    Health Service Implementation Support & \texttt{cap:bc-eb94288c} (at \ref{cap:bc-eb94288c}) & \\
    \hline
    Health Service Continuous Improvement & \texttt{cap:bc-b1242610} (at \ref{cap:bc-b1242610}) & \\
    \hline
  \end{tabularx}
\end{table}

\textbf{Intent:} Health Service Modelling identifies the organisational ability to undertake the stated health-service responsibility. The capability describes what the enterprise must be able to do; it does not assert that Harmonia performs or owns that activity.

\textbf{Boundary:} Harmonia relevance may arise because this activity establishes the context in which information is created, needed, interpreted, related, exchanged or presented. Responsibility for the underlying business decision or care activity remains with the appropriate actor or organisation.

\subsection{CT3 - Health Service Need Identification}\label{cap:bc-df112315}
\vspace{-0.5em}
\begin{table}[htbp]
  \centering
  \footnotesize
  \begin{tabularx}{1\textwidth}{|l|l|X|}
    \hline
    \multicolumn{3}{c}{\textbf{Capability Metadata}} \\
    \hline
    \textbf{Name} & \textbf{Value} & \textbf{Comment and/or Description} \\
    \hline
    \textbf{Reference Id} & \texttt{BC-DF112315} & \\
    \hline
    \textbf{Status} & Established working hierarchy & \\
    \hline
    \textbf{Ancestry} &  & \\
    \hline
    \textbf{Source} &  & \\
    \hline
    \textbf{Enabled By} &  & \\
    \hline
    \textbf{Supports} &  & \\
    \hline
    \textbf{Derived From} &  & \\
    \hline
    \textbf{Reviewed Status} & Review meaning, business boundary, information context and any unresolved ancestry. & \\
    \hline
  \end{tabularx}
\end{table}

\begin{table}[htbp]
  \centering
  \footnotesize
  \begin{tabularx}{1\textwidth}{|l|l|X|}
    \hline
    \multicolumn{3}{c}{\textbf{Contained Features}} \\
    \hline
    \textbf{Feature Name} & \textbf{Feature Id} & \textbf{Comment} \\
    \hline
    \multicolumn{3}{|l|}{\emph{No contained Features are currently identified.}} \\
    \hline
  \end{tabularx}
\end{table}

\textbf{Intent:} Health Service Need Identification identifies the organisational ability to undertake the stated health-service responsibility. The capability describes what the enterprise must be able to do; it does not assert that Harmonia performs or owns that activity.

\textbf{Boundary:} Harmonia relevance may arise because this activity establishes the context in which information is created, needed, interpreted, related, exchanged or presented. Responsibility for the underlying business decision or care activity remains with the appropriate actor or organisation.

\subsection{CT3 - Health Service Capability Identification}\label{cap:bc-a7a90ee0}
\vspace{-0.5em}
\begin{table}[htbp]
  \centering
  \footnotesize
  \begin{tabularx}{1\textwidth}{|l|l|X|}
    \hline
    \multicolumn{3}{c}{\textbf{Capability Metadata}} \\
    \hline
    \textbf{Name} & \textbf{Value} & \textbf{Comment and/or Description} \\
    \hline
    \textbf{Reference Id} & \texttt{BC-A7A90EE0} & \\
    \hline
    \textbf{Status} & Established working hierarchy & \\
    \hline
    \textbf{Ancestry} &  & \\
    \hline
    \textbf{Source} &  & \\
    \hline
    \textbf{Enabled By} &  & \\
    \hline
    \textbf{Supports} &  & \\
    \hline
    \textbf{Derived From} &  & \\
    \hline
    \textbf{Reviewed Status} & Review meaning, business boundary, information context and any unresolved ancestry. & \\
    \hline
  \end{tabularx}
\end{table}

\begin{table}[htbp]
  \centering
  \footnotesize
  \begin{tabularx}{1\textwidth}{|l|l|X|}
    \hline
    \multicolumn{3}{c}{\textbf{Contained Features}} \\
    \hline
    \textbf{Feature Name} & \textbf{Feature Id} & \textbf{Comment} \\
    \hline
    \multicolumn{3}{|l|}{\emph{No contained Features are currently identified.}} \\
    \hline
  \end{tabularx}
\end{table}

\textbf{Intent:} Health Service Capability Identification identifies the organisational ability to undertake the stated health-service responsibility. The capability describes what the enterprise must be able to do; it does not assert that Harmonia performs or owns that activity.

\textbf{Boundary:} Harmonia relevance may arise because this activity establishes the context in which information is created, needed, interpreted, related, exchanged or presented. Responsibility for the underlying business decision or care activity remains with the appropriate actor or organisation.

\subsection{CT3 - Health Service Implementation Support}\label{cap:bc-eb94288c}
\vspace{-0.5em}
\begin{table}[htbp]
  \centering
  \footnotesize
  \begin{tabularx}{1\textwidth}{|l|l|X|}
    \hline
    \multicolumn{3}{c}{\textbf{Capability Metadata}} \\
    \hline
    \textbf{Name} & \textbf{Value} & \textbf{Comment and/or Description} \\
    \hline
    \textbf{Reference Id} & \texttt{BC-EB94288C} & \\
    \hline
    \textbf{Status} & Established working hierarchy & \\
    \hline
    \textbf{Ancestry} &  & \\
    \hline
    \textbf{Source} &  & \\
    \hline
    \textbf{Enabled By} &  & \\
    \hline
    \textbf{Supports} &  & \\
    \hline
    \textbf{Derived From} &  & \\
    \hline
    \textbf{Reviewed Status} & Review meaning, business boundary, information context and any unresolved ancestry. & \\
    \hline
  \end{tabularx}
\end{table}

\begin{table}[htbp]
  \centering
  \footnotesize
  \begin{tabularx}{1\textwidth}{|l|l|X|}
    \hline
    \multicolumn{3}{c}{\textbf{Contained Features}} \\
    \hline
    \textbf{Feature Name} & \textbf{Feature Id} & \textbf{Comment} \\
    \hline
    \multicolumn{3}{|l|}{\emph{No contained Features are currently identified.}} \\
    \hline
  \end{tabularx}
\end{table}

\textbf{Intent:} Health Service Implementation Support identifies the organisational ability to undertake the stated health-service responsibility. The capability describes what the enterprise must be able to do; it does not assert that Harmonia performs or owns that activity.

\textbf{Boundary:} Harmonia relevance may arise because this activity establishes the context in which information is created, needed, interpreted, related, exchanged or presented. Responsibility for the underlying business decision or care activity remains with the appropriate actor or organisation.

\subsection{CT3 - Health Service Continuous Improvement}\label{cap:bc-b1242610}
\vspace{-0.5em}
\begin{table}[htbp]
  \centering
  \footnotesize
  \begin{tabularx}{1\textwidth}{|l|l|X|}
    \hline
    \multicolumn{3}{c}{\textbf{Capability Metadata}} \\
    \hline
    \textbf{Name} & \textbf{Value} & \textbf{Comment and/or Description} \\
    \hline
    \textbf{Reference Id} & \texttt{BC-B1242610} & \\
    \hline
    \textbf{Status} & Established working hierarchy & \\
    \hline
    \textbf{Ancestry} &  & \\
    \hline
    \textbf{Source} &  & \\
    \hline
    \textbf{Enabled By} &  & \\
    \hline
    \textbf{Supports} &  & \\
    \hline
    \textbf{Derived From} &  & \\
    \hline
    \textbf{Reviewed Status} & Review meaning, business boundary, information context and any unresolved ancestry. & \\
    \hline
  \end{tabularx}
\end{table}

\begin{table}[htbp]
  \centering
  \footnotesize
  \begin{tabularx}{1\textwidth}{|l|l|X|}
    \hline
    \multicolumn{3}{c}{\textbf{Contained Features}} \\
    \hline
    \textbf{Feature Name} & \textbf{Feature Id} & \textbf{Comment} \\
    \hline
    \multicolumn{3}{|l|}{\emph{No contained Features are currently identified.}} \\
    \hline
  \end{tabularx}
\end{table}

\textbf{Intent:} Health Service Continuous Improvement identifies the organisational ability to undertake the stated health-service responsibility. The capability describes what the enterprise must be able to do; it does not assert that Harmonia performs or owns that activity.

\textbf{Boundary:} Harmonia relevance may arise because this activity establishes the context in which information is created, needed, interpreted, related, exchanged or presented. Responsibility for the underlying business decision or care activity remains with the appropriate actor or organisation.

\section{CT2 - Health Service Alignment}\label{cap:bc-8229abfd}
\vspace{-0.5em}
\begin{table}[htbp]
  \centering
  \footnotesize
  \begin{tabularx}{1\textwidth}{|l|l|X|}
    \hline
    \multicolumn{3}{c}{\textbf{Capability Metadata}} \\
    \hline
    \textbf{Name} & \textbf{Value} & \textbf{Comment and/or Description} \\
    \hline
    \textbf{Reference Id} & \texttt{BC-8229ABFD} & \\
    \hline
    \textbf{Status} & Established working hierarchy & \\
    \hline
    \textbf{Ancestry} &  & \\
    \hline
    \textbf{Source} &  & \\
    \hline
    \textbf{Enabled By} &  & \\
    \hline
    \textbf{Supports} &  & \\
    \hline
    \textbf{Derived From} &  & \\
    \hline
    \textbf{Reviewed Status} &  & \\
    \hline
  \end{tabularx}
\end{table}

\begin{table}[htbp]
  \centering
  \footnotesize
  \begin{tabularx}{1\textwidth}{|l|l|X|}
    \hline
    \multicolumn{3}{c}{\textbf{Contained Capabilities}} \\
    \hline
    \textbf{Capability Name} & \textbf{Capability Id} & \textbf{Comment} \\
    \hline
    Health Service Co-Dependency Management & \texttt{cap:bc-60910383} (at \ref{cap:bc-60910383}) & \\
    \hline
    Health Service Resource Identification & \texttt{cap:bc-ac53a525} (at \ref{cap:bc-ac53a525}) & \\
    \hline
    Health Service Financial Management & \texttt{cap:bc-93fdb4f7} (at \ref{cap:bc-93fdb4f7}) & \\
    \hline
    Health Service Implementation Planning & \texttt{cap:bc-ef0176f3} (at \ref{cap:bc-ef0176f3}) & \\
    \hline
  \end{tabularx}
\end{table}

\textbf{Intent:} Health Service Alignment identifies the organisational ability to undertake the stated health-service responsibility. The capability describes what the enterprise must be able to do; it does not assert that Harmonia performs or owns that activity.

\textbf{Boundary:} Harmonia relevance may arise because this activity establishes the context in which information is created, needed, interpreted, related, exchanged or presented. Responsibility for the underlying business decision or care activity remains with the appropriate actor or organisation.

\subsection{CT3 - Health Service Co-Dependency Management}\label{cap:bc-60910383}
\vspace{-0.5em}
\begin{table}[htbp]
  \centering
  \footnotesize
  \begin{tabularx}{1\textwidth}{|l|l|X|}
    \hline
    \multicolumn{3}{c}{\textbf{Capability Metadata}} \\
    \hline
    \textbf{Name} & \textbf{Value} & \textbf{Comment and/or Description} \\
    \hline
    \textbf{Reference Id} & \texttt{BC-60910383} & \\
    \hline
    \textbf{Status} & Established working hierarchy & \\
    \hline
    \textbf{Ancestry} &  & \\
    \hline
    \textbf{Source} &  & \\
    \hline
    \textbf{Enabled By} &  & \\
    \hline
    \textbf{Supports} &  & \\
    \hline
    \textbf{Derived From} &  & \\
    \hline
    \textbf{Reviewed Status} & Review meaning, business boundary, information context and any unresolved ancestry. & \\
    \hline
  \end{tabularx}
\end{table}

\begin{table}[htbp]
  \centering
  \footnotesize
  \begin{tabularx}{1\textwidth}{|l|l|X|}
    \hline
    \multicolumn{3}{c}{\textbf{Contained Features}} \\
    \hline
    \textbf{Feature Name} & \textbf{Feature Id} & \textbf{Comment} \\
    \hline
    \multicolumn{3}{|l|}{\emph{No contained Features are currently identified.}} \\
    \hline
  \end{tabularx}
\end{table}

\textbf{Intent:} Health Service Co-Dependency Management identifies the organisational ability to undertake the stated health-service responsibility. The capability describes what the enterprise must be able to do; it does not assert that Harmonia performs or owns that activity.

\textbf{Boundary:} Harmonia relevance may arise because this activity establishes the context in which information is created, needed, interpreted, related, exchanged or presented. Responsibility for the underlying business decision or care activity remains with the appropriate actor or organisation.

\subsection{CT3 - Health Service Resource Identification}\label{cap:bc-ac53a525}
\vspace{-0.5em}
\begin{table}[htbp]
  \centering
  \footnotesize
  \begin{tabularx}{1\textwidth}{|l|l|X|}
    \hline
    \multicolumn{3}{c}{\textbf{Capability Metadata}} \\
    \hline
    \textbf{Name} & \textbf{Value} & \textbf{Comment and/or Description} \\
    \hline
    \textbf{Reference Id} & \texttt{BC-AC53A525} & \\
    \hline
    \textbf{Status} & Established working hierarchy & \\
    \hline
    \textbf{Ancestry} &  & \\
    \hline
    \textbf{Source} &  & \\
    \hline
    \textbf{Enabled By} &  & \\
    \hline
    \textbf{Supports} &  & \\
    \hline
    \textbf{Derived From} &  & \\
    \hline
    \textbf{Reviewed Status} & Review meaning, business boundary, information context and any unresolved ancestry. & \\
    \hline
  \end{tabularx}
\end{table}

\begin{table}[htbp]
  \centering
  \footnotesize
  \begin{tabularx}{1\textwidth}{|l|l|X|}
    \hline
    \multicolumn{3}{c}{\textbf{Contained Features}} \\
    \hline
    \textbf{Feature Name} & \textbf{Feature Id} & \textbf{Comment} \\
    \hline
    \multicolumn{3}{|l|}{\emph{No contained Features are currently identified.}} \\
    \hline
  \end{tabularx}
\end{table}

\textbf{Intent:} Health Service Resource Identification identifies the organisational ability to undertake the stated health-service responsibility. The capability describes what the enterprise must be able to do; it does not assert that Harmonia performs or owns that activity.

\textbf{Boundary:} Harmonia relevance may arise because this activity establishes the context in which information is created, needed, interpreted, related, exchanged or presented. Responsibility for the underlying business decision or care activity remains with the appropriate actor or organisation.

\subsection{CT3 - Health Service Financial Management}\label{cap:bc-93fdb4f7}
\vspace{-0.5em}
\begin{table}[htbp]
  \centering
  \footnotesize
  \begin{tabularx}{1\textwidth}{|l|l|X|}
    \hline
    \multicolumn{3}{c}{\textbf{Capability Metadata}} \\
    \hline
    \textbf{Name} & \textbf{Value} & \textbf{Comment and/or Description} \\
    \hline
    \textbf{Reference Id} & \texttt{BC-93FDB4F7} & \\
    \hline
    \textbf{Status} & Established working hierarchy & \\
    \hline
    \textbf{Ancestry} &  & \\
    \hline
    \textbf{Source} &  & \\
    \hline
    \textbf{Enabled By} &  & \\
    \hline
    \textbf{Supports} &  & \\
    \hline
    \textbf{Derived From} &  & \\
    \hline
    \textbf{Reviewed Status} & Review meaning, business boundary, information context and any unresolved ancestry. & \\
    \hline
  \end{tabularx}
\end{table}

\begin{table}[htbp]
  \centering
  \footnotesize
  \begin{tabularx}{1\textwidth}{|l|l|X|}
    \hline
    \multicolumn{3}{c}{\textbf{Contained Features}} \\
    \hline
    \textbf{Feature Name} & \textbf{Feature Id} & \textbf{Comment} \\
    \hline
    \multicolumn{3}{|l|}{\emph{No contained Features are currently identified.}} \\
    \hline
  \end{tabularx}
\end{table}

\textbf{Intent:} Health Service Financial Management identifies the organisational ability to undertake the stated health-service responsibility. The capability describes what the enterprise must be able to do; it does not assert that Harmonia performs or owns that activity.

\textbf{Boundary:} Harmonia relevance may arise because this activity establishes the context in which information is created, needed, interpreted, related, exchanged or presented. Responsibility for the underlying business decision or care activity remains with the appropriate actor or organisation.

\subsection{CT3 - Health Service Implementation Planning}\label{cap:bc-ef0176f3}
\vspace{-0.5em}
\begin{table}[htbp]
  \centering
  \footnotesize
  \begin{tabularx}{1\textwidth}{|l|l|X|}
    \hline
    \multicolumn{3}{c}{\textbf{Capability Metadata}} \\
    \hline
    \textbf{Name} & \textbf{Value} & \textbf{Comment and/or Description} \\
    \hline
    \textbf{Reference Id} & \texttt{BC-EF0176F3} & \\
    \hline
    \textbf{Status} & Established working hierarchy & \\
    \hline
    \textbf{Ancestry} &  & \\
    \hline
    \textbf{Source} &  & \\
    \hline
    \textbf{Enabled By} &  & \\
    \hline
    \textbf{Supports} &  & \\
    \hline
    \textbf{Derived From} &  & \\
    \hline
    \textbf{Reviewed Status} & Review meaning, business boundary, information context and any unresolved ancestry. & \\
    \hline
  \end{tabularx}
\end{table}

\begin{table}[htbp]
  \centering
  \footnotesize
  \begin{tabularx}{1\textwidth}{|l|l|X|}
    \hline
    \multicolumn{3}{c}{\textbf{Contained Features}} \\
    \hline
    \textbf{Feature Name} & \textbf{Feature Id} & \textbf{Comment} \\
    \hline
    \multicolumn{3}{|l|}{\emph{No contained Features are currently identified.}} \\
    \hline
  \end{tabularx}
\end{table}

\textbf{Intent:} Health Service Implementation Planning identifies the organisational ability to undertake the stated health-service responsibility. The capability describes what the enterprise must be able to do; it does not assert that Harmonia performs or owns that activity.

\textbf{Boundary:} Harmonia relevance may arise because this activity establishes the context in which information is created, needed, interpreted, related, exchanged or presented. Responsibility for the underlying business decision or care activity remains with the appropriate actor or organisation.

\chapter{CT1 - Workforce \& Organisational Capability}\label{cap:bc-f02f310b}
\noindent\textbf{Reference ID:} \texttt{BC-F02F310B}\quad \textbf{Status:} Established working hierarchy\par
\noindent\textit{Source: Canonical Business Capability model / working map}\par
\textbf{Intent:} Workforce \& Organisational Capability identifies the organisational ability to undertake the stated health-service responsibility. The capability describes what the enterprise must be able to do; it does not assert that Harmonia performs or owns that activity.

\textbf{Boundary:} Harmonia relevance may arise because this activity establishes the context in which information is created, needed, interpreted, related, exchanged or presented. Responsibility for the underlying business decision or care activity remains with the appropriate actor or organisation.

\reviewstatus{Check scope and information relevance; no business authority transfers to Harmonia.}
\section{CT2 - Health Workforce Management}\label{cap:bc-b7de9919}
\vspace{-0.5em}
\begin{table}[htbp]
  \centering
  \footnotesize
  \begin{tabularx}{1\textwidth}{|l|l|X|}
    \hline
    \multicolumn{3}{c}{\textbf{Capability Metadata}} \\
    \hline
    \textbf{Name} & \textbf{Value} & \textbf{Comment and/or Description} \\
    \hline
    \textbf{Reference Id} & \texttt{BC-B7DE9919} & \\
    \hline
    \textbf{Status} & Established working hierarchy & \\
    \hline
    \textbf{Ancestry} &  & \\
    \hline
    \textbf{Source} &  & \\
    \hline
    \textbf{Enabled By} &  & \\
    \hline
    \textbf{Supports} &  & \\
    \hline
    \textbf{Derived From} &  & \\
    \hline
    \textbf{Reviewed Status} &  & \\
    \hline
  \end{tabularx}
\end{table}

\begin{table}[htbp]
  \centering
  \footnotesize
  \begin{tabularx}{1\textwidth}{|l|l|X|}
    \hline
    \multicolumn{3}{c}{\textbf{Contained Capabilities}} \\
    \hline
    \textbf{Capability Name} & \textbf{Capability Id} & \textbf{Comment} \\
    \hline
    Practitioner Qualification Management & \texttt{cap:bc-875377cd} (at \ref{cap:bc-875377cd}) & \\
    \hline
    Practitioner Engagement \& Privilege Management & \texttt{cap:bc-93af8cf1} (at \ref{cap:bc-93af8cf1}) & \\
    \hline
    Workforce Service Coverage Management & \texttt{cap:bc-db75d205} (at \ref{cap:bc-db75d205}) & \\
    \hline
  \end{tabularx}
\end{table}

\textbf{Intent:} Health Workforce Management identifies the organisational ability to undertake the stated health-service responsibility. The capability describes what the enterprise must be able to do; it does not assert that Harmonia performs or owns that activity.

\textbf{Boundary:} Harmonia relevance may arise because this activity establishes the context in which information is created, needed, interpreted, related, exchanged or presented. Responsibility for the underlying business decision or care activity remains with the appropriate actor or organisation.

\subsection{CT3 - Practitioner Qualification Management}\label{cap:bc-875377cd}
\vspace{-0.5em}
\begin{table}[htbp]
  \centering
  \footnotesize
  \begin{tabularx}{1\textwidth}{|l|l|X|}
    \hline
    \multicolumn{3}{c}{\textbf{Capability Metadata}} \\
    \hline
    \textbf{Name} & \textbf{Value} & \textbf{Comment and/or Description} \\
    \hline
    \textbf{Reference Id} & \texttt{BC-875377CD} & \\
    \hline
    \textbf{Status} & Established working hierarchy & \\
    \hline
    \textbf{Ancestry} &  & \\
    \hline
    \textbf{Source} &  & \\
    \hline
    \textbf{Enabled By} &  & \\
    \hline
    \textbf{Supports} &  & \\
    \hline
    \textbf{Derived From} &  & \\
    \hline
    \textbf{Reviewed Status} & Review meaning, business boundary, information context and any unresolved ancestry. & \\
    \hline
  \end{tabularx}
\end{table}

\begin{table}[htbp]
  \centering
  \footnotesize
  \begin{tabularx}{1\textwidth}{|l|l|X|}
    \hline
    \multicolumn{3}{c}{\textbf{Contained Features}} \\
    \hline
    \textbf{Feature Name} & \textbf{Feature Id} & \textbf{Comment} \\
    \hline
    \multicolumn{3}{|l|}{\emph{No contained Features are currently identified.}} \\
    \hline
  \end{tabularx}
\end{table}

\textbf{Intent:} Practitioner Qualification Management identifies the organisational ability to undertake the stated health-service responsibility. The capability describes what the enterprise must be able to do; it does not assert that Harmonia performs or owns that activity.

\textbf{Boundary:} Harmonia relevance may arise because this activity establishes the context in which information is created, needed, interpreted, related, exchanged or presented. Responsibility for the underlying business decision or care activity remains with the appropriate actor or organisation.

\subsection{CT3 - Practitioner Engagement \& Privilege Management}\label{cap:bc-93af8cf1}
\vspace{-0.5em}
\begin{table}[htbp]
  \centering
  \footnotesize
  \begin{tabularx}{1\textwidth}{|l|l|X|}
    \hline
    \multicolumn{3}{c}{\textbf{Capability Metadata}} \\
    \hline
    \textbf{Name} & \textbf{Value} & \textbf{Comment and/or Description} \\
    \hline
    \textbf{Reference Id} & \texttt{BC-93AF8CF1} & \\
    \hline
    \textbf{Status} & Established working hierarchy & \\
    \hline
    \textbf{Ancestry} &  & \\
    \hline
    \textbf{Source} &  & \\
    \hline
    \textbf{Enabled By} &  & \\
    \hline
    \textbf{Supports} &  & \\
    \hline
    \textbf{Derived From} &  & \\
    \hline
    \textbf{Reviewed Status} & Review meaning, business boundary, information context and any unresolved ancestry. & \\
    \hline
  \end{tabularx}
\end{table}

\begin{table}[htbp]
  \centering
  \footnotesize
  \begin{tabularx}{1\textwidth}{|l|l|X|}
    \hline
    \multicolumn{3}{c}{\textbf{Contained Features}} \\
    \hline
    \textbf{Feature Name} & \textbf{Feature Id} & \textbf{Comment} \\
    \hline
    \multicolumn{3}{|l|}{\emph{No contained Features are currently identified.}} \\
    \hline
  \end{tabularx}
\end{table}

\textbf{Intent:} Practitioner Engagement \& Privilege Management identifies the organisational ability to undertake the stated health-service responsibility. The capability describes what the enterprise must be able to do; it does not assert that Harmonia performs or owns that activity.

\textbf{Boundary:} Harmonia relevance may arise because this activity establishes the context in which information is created, needed, interpreted, related, exchanged or presented. Responsibility for the underlying business decision or care activity remains with the appropriate actor or organisation.

\subsection{CT3 - Workforce Service Coverage Management}\label{cap:bc-db75d205}
\vspace{-0.5em}
\begin{table}[htbp]
  \centering
  \footnotesize
  \begin{tabularx}{1\textwidth}{|l|l|X|}
    \hline
    \multicolumn{3}{c}{\textbf{Capability Metadata}} \\
    \hline
    \textbf{Name} & \textbf{Value} & \textbf{Comment and/or Description} \\
    \hline
    \textbf{Reference Id} & \texttt{BC-DB75D205} & \\
    \hline
    \textbf{Status} & Established working hierarchy & \\
    \hline
    \textbf{Ancestry} &  & \\
    \hline
    \textbf{Source} &  & \\
    \hline
    \textbf{Enabled By} &  & \\
    \hline
    \textbf{Supports} &  & \\
    \hline
    \textbf{Derived From} &  & \\
    \hline
    \textbf{Reviewed Status} & Review meaning, business boundary, information context and any unresolved ancestry. & \\
    \hline
  \end{tabularx}
\end{table}

\begin{table}[htbp]
  \centering
  \footnotesize
  \begin{tabularx}{1\textwidth}{|l|l|X|}
    \hline
    \multicolumn{3}{c}{\textbf{Contained Features}} \\
    \hline
    \textbf{Feature Name} & \textbf{Feature Id} & \textbf{Comment} \\
    \hline
    \multicolumn{3}{|l|}{\emph{No contained Features are currently identified.}} \\
    \hline
  \end{tabularx}
\end{table}

\textbf{Intent:} Workforce Service Coverage Management identifies the organisational ability to undertake the stated health-service responsibility. The capability describes what the enterprise must be able to do; it does not assert that Harmonia performs or owns that activity.

\textbf{Boundary:} Harmonia relevance may arise because this activity establishes the context in which information is created, needed, interpreted, related, exchanged or presented. Responsibility for the underlying business decision or care activity remains with the appropriate actor or organisation.

\section{CT2 - Healthcare Organisation \& Service Structure Management}\label{cap:bc-dc260b97}
\vspace{-0.5em}
\begin{table}[htbp]
  \centering
  \footnotesize
  \begin{tabularx}{1\textwidth}{|l|l|X|}
    \hline
    \multicolumn{3}{c}{\textbf{Capability Metadata}} \\
    \hline
    \textbf{Name} & \textbf{Value} & \textbf{Comment and/or Description} \\
    \hline
    \textbf{Reference Id} & \texttt{BC-DC260B97} & \\
    \hline
    \textbf{Status} & Established working hierarchy & \\
    \hline
    \textbf{Ancestry} &  & \\
    \hline
    \textbf{Source} &  & \\
    \hline
    \textbf{Enabled By} &  & \\
    \hline
    \textbf{Supports} &  & \\
    \hline
    \textbf{Derived From} &  & \\
    \hline
    \textbf{Reviewed Status} &  & \\
    \hline
  \end{tabularx}
\end{table}

\begin{table}[htbp]
  \centering
  \footnotesize
  \begin{tabularx}{1\textwidth}{|l|l|X|}
    \hline
    \multicolumn{3}{c}{\textbf{Contained Capabilities}} \\
    \hline
    \textbf{Capability Name} & \textbf{Capability Id} & \textbf{Comment} \\
    \hline
    Healthcare Organisation Identity \& Structure Management & \texttt{cap:bc-a170cfa2} (at \ref{cap:bc-a170cfa2}) & \\
    \hline
    Healthcare Location Definition Management & \texttt{cap:bc-65079689} (at \ref{cap:bc-65079689}) & \\
    \hline
    Healthcare Service Offering Management & \texttt{cap:bc-40a6eee3} (at \ref{cap:bc-40a6eee3}) & \\
    \hline
  \end{tabularx}
\end{table}

\textbf{Intent:} Healthcare Organisation \& Service Structure Management identifies the organisational ability to undertake the stated health-service responsibility. The capability describes what the enterprise must be able to do; it does not assert that Harmonia performs or owns that activity.

\textbf{Boundary:} Harmonia relevance may arise because this activity establishes the context in which information is created, needed, interpreted, related, exchanged or presented. Responsibility for the underlying business decision or care activity remains with the appropriate actor or organisation.

\subsection{CT3 - Healthcare Organisation Identity \& Structure Management}\label{cap:bc-a170cfa2}
\vspace{-0.5em}
\begin{table}[htbp]
  \centering
  \footnotesize
  \begin{tabularx}{1\textwidth}{|l|l|X|}
    \hline
    \multicolumn{3}{c}{\textbf{Capability Metadata}} \\
    \hline
    \textbf{Name} & \textbf{Value} & \textbf{Comment and/or Description} \\
    \hline
    \textbf{Reference Id} & \texttt{BC-A170CFA2} & \\
    \hline
    \textbf{Status} & Established working hierarchy & \\
    \hline
    \textbf{Ancestry} &  & \\
    \hline
    \textbf{Source} &  & \\
    \hline
    \textbf{Enabled By} &  & \\
    \hline
    \textbf{Supports} &  & \\
    \hline
    \textbf{Derived From} &  & \\
    \hline
    \textbf{Reviewed Status} & Review meaning, business boundary, information context and any unresolved ancestry. & \\
    \hline
  \end{tabularx}
\end{table}

\begin{table}[htbp]
  \centering
  \footnotesize
  \begin{tabularx}{1\textwidth}{|l|l|X|}
    \hline
    \multicolumn{3}{c}{\textbf{Contained Features}} \\
    \hline
    \textbf{Feature Name} & \textbf{Feature Id} & \textbf{Comment} \\
    \hline
    \multicolumn{3}{|l|}{\emph{No contained Features are currently identified.}} \\
    \hline
  \end{tabularx}
\end{table}

\textbf{Intent:} Healthcare Organisation Identity \& Structure Management identifies the organisational ability to undertake the stated health-service responsibility. The capability describes what the enterprise must be able to do; it does not assert that Harmonia performs or owns that activity.

\textbf{Boundary:} Harmonia relevance may arise because this activity establishes the context in which information is created, needed, interpreted, related, exchanged or presented. Responsibility for the underlying business decision or care activity remains with the appropriate actor or organisation.

\subsection{CT3 - Healthcare Location Definition Management}\label{cap:bc-65079689}
\vspace{-0.5em}
\begin{table}[htbp]
  \centering
  \footnotesize
  \begin{tabularx}{1\textwidth}{|l|l|X|}
    \hline
    \multicolumn{3}{c}{\textbf{Capability Metadata}} \\
    \hline
    \textbf{Name} & \textbf{Value} & \textbf{Comment and/or Description} \\
    \hline
    \textbf{Reference Id} & \texttt{BC-65079689} & \\
    \hline
    \textbf{Status} & Established working hierarchy & \\
    \hline
    \textbf{Ancestry} &  & \\
    \hline
    \textbf{Source} &  & \\
    \hline
    \textbf{Enabled By} &  & \\
    \hline
    \textbf{Supports} &  & \\
    \hline
    \textbf{Derived From} &  & \\
    \hline
    \textbf{Reviewed Status} & Review meaning, business boundary, information context and any unresolved ancestry. & \\
    \hline
  \end{tabularx}
\end{table}

\begin{table}[htbp]
  \centering
  \footnotesize
  \begin{tabularx}{1\textwidth}{|l|l|X|}
    \hline
    \multicolumn{3}{c}{\textbf{Contained Features}} \\
    \hline
    \textbf{Feature Name} & \textbf{Feature Id} & \textbf{Comment} \\
    \hline
    \multicolumn{3}{|l|}{\emph{No contained Features are currently identified.}} \\
    \hline
  \end{tabularx}
\end{table}

\textbf{Intent:} Healthcare Location Definition Management identifies the organisational ability to undertake the stated health-service responsibility. The capability describes what the enterprise must be able to do; it does not assert that Harmonia performs or owns that activity.

\textbf{Boundary:} Harmonia relevance may arise because this activity establishes the context in which information is created, needed, interpreted, related, exchanged or presented. Responsibility for the underlying business decision or care activity remains with the appropriate actor or organisation.

\subsection{CT3 - Healthcare Service Offering Management}\label{cap:bc-40a6eee3}
\vspace{-0.5em}
\begin{table}[htbp]
  \centering
  \footnotesize
  \begin{tabularx}{1\textwidth}{|l|l|X|}
    \hline
    \multicolumn{3}{c}{\textbf{Capability Metadata}} \\
    \hline
    \textbf{Name} & \textbf{Value} & \textbf{Comment and/or Description} \\
    \hline
    \textbf{Reference Id} & \texttt{BC-40A6EEE3} & \\
    \hline
    \textbf{Status} & Established working hierarchy & \\
    \hline
    \textbf{Ancestry} &  & \\
    \hline
    \textbf{Source} &  & \\
    \hline
    \textbf{Enabled By} &  & \\
    \hline
    \textbf{Supports} &  & \\
    \hline
    \textbf{Derived From} &  & \\
    \hline
    \textbf{Reviewed Status} & Review meaning, business boundary, information context and any unresolved ancestry. & \\
    \hline
  \end{tabularx}
\end{table}

\begin{table}[htbp]
  \centering
  \footnotesize
  \begin{tabularx}{1\textwidth}{|l|l|X|}
    \hline
    \multicolumn{3}{c}{\textbf{Contained Features}} \\
    \hline
    \textbf{Feature Name} & \textbf{Feature Id} & \textbf{Comment} \\
    \hline
    \multicolumn{3}{|l|}{\emph{No contained Features are currently identified.}} \\
    \hline
  \end{tabularx}
\end{table}

\textbf{Intent:} Healthcare Service Offering Management identifies the organisational ability to undertake the stated health-service responsibility. The capability describes what the enterprise must be able to do; it does not assert that Harmonia performs or owns that activity.

\textbf{Boundary:} Harmonia relevance may arise because this activity establishes the context in which information is created, needed, interpreted, related, exchanged or presented. Responsibility for the underlying business decision or care activity remains with the appropriate actor or organisation.

\chapter{CT1 - Corporate Resources \& Enterprise Services}\label{cap:bc-2a4887a2}
\noindent\textbf{Reference ID:} \texttt{BC-2A4887A2}\quad \textbf{Status:} Established working hierarchy\par
\noindent\textit{Source: Canonical Business Capability model / working map}\par
\textbf{Intent:} Corporate Resources \& Enterprise Services identifies the organisational ability to undertake the stated health-service responsibility. The capability describes what the enterprise must be able to do; it does not assert that Harmonia performs or owns that activity.

\textbf{Boundary:} Harmonia relevance may arise because this activity establishes the context in which information is created, needed, interpreted, related, exchanged or presented. Responsibility for the underlying business decision or care activity remains with the appropriate actor or organisation.

\reviewstatus{Check scope and information relevance; no business authority transfers to Harmonia.}
\section{CT2 - Financial \& Commercial Services}\label{cap:bc-3618d6d1}
\vspace{-0.5em}
\begin{table}[htbp]
  \centering
  \footnotesize
  \begin{tabularx}{1\textwidth}{|l|l|X|}
    \hline
    \multicolumn{3}{c}{\textbf{Capability Metadata}} \\
    \hline
    \textbf{Name} & \textbf{Value} & \textbf{Comment and/or Description} \\
    \hline
    \textbf{Reference Id} & \texttt{BC-3618D6D1} & \\
    \hline
    \textbf{Status} & Established working hierarchy & \\
    \hline
    \textbf{Ancestry} &  & \\
    \hline
    \textbf{Source} &  & \\
    \hline
    \textbf{Enabled By} &  & \\
    \hline
    \textbf{Supports} &  & \\
    \hline
    \textbf{Derived From} &  & \\
    \hline
    \textbf{Reviewed Status} &  & \\
    \hline
  \end{tabularx}
\end{table}

\begin{table}[htbp]
  \centering
  \footnotesize
  \begin{tabularx}{1\textwidth}{|l|l|X|}
    \hline
    \multicolumn{3}{c}{\textbf{Contained Capabilities}} \\
    \hline
    \textbf{Capability Name} & \textbf{Capability Id} & \textbf{Comment} \\
    \hline
    Clinical Activity Charging Information Provision & \texttt{cap:bc-924a98c9} (at \ref{cap:bc-924a98c9}) & \\
    \hline
  \end{tabularx}
\end{table}

\textbf{Intent:} Financial \& Commercial Services identifies the organisational ability to undertake the stated health-service responsibility. The capability describes what the enterprise must be able to do; it does not assert that Harmonia performs or owns that activity.

\textbf{Boundary:} Harmonia relevance may arise because this activity establishes the context in which information is created, needed, interpreted, related, exchanged or presented. Responsibility for the underlying business decision or care activity remains with the appropriate actor or organisation.

\subsection{CT3 - Clinical Activity Charging Information Provision}\label{cap:bc-924a98c9}
\vspace{-0.5em}
\begin{table}[htbp]
  \centering
  \footnotesize
  \begin{tabularx}{1\textwidth}{|l|l|X|}
    \hline
    \multicolumn{3}{c}{\textbf{Capability Metadata}} \\
    \hline
    \textbf{Name} & \textbf{Value} & \textbf{Comment and/or Description} \\
    \hline
    \textbf{Reference Id} & \texttt{BC-924A98C9} & \\
    \hline
    \textbf{Status} & Established working hierarchy & \\
    \hline
    \textbf{Ancestry} &  & \\
    \hline
    \textbf{Source} &  & \\
    \hline
    \textbf{Enabled By} &  & \\
    \hline
    \textbf{Supports} &  & \\
    \hline
    \textbf{Derived From} &  & \\
    \hline
    \textbf{Reviewed Status} & Review meaning, business boundary, information context and any unresolved ancestry. & \\
    \hline
  \end{tabularx}
\end{table}

\begin{table}[htbp]
  \centering
  \footnotesize
  \begin{tabularx}{1\textwidth}{|l|l|X|}
    \hline
    \multicolumn{3}{c}{\textbf{Contained Features}} \\
    \hline
    \textbf{Feature Name} & \textbf{Feature Id} & \textbf{Comment} \\
    \hline
    \multicolumn{3}{|l|}{\emph{No contained Features are currently identified.}} \\
    \hline
  \end{tabularx}
\end{table}

\textbf{Intent:} Clinical Activity Charging Information Provision identifies the organisational ability to undertake the stated health-service responsibility. The capability describes what the enterprise must be able to do; it does not assert that Harmonia performs or owns that activity.

\textbf{Boundary:} Harmonia relevance may arise because this activity establishes the context in which information is created, needed, interpreted, related, exchanged or presented. Responsibility for the underlying business decision or care activity remains with the appropriate actor or organisation.

\section{CT2 - Corporate Support Services}\label{cap:bc-f411af90}
\vspace{-0.5em}
\begin{table}[htbp]
  \centering
  \footnotesize
  \begin{tabularx}{1\textwidth}{|l|l|X|}
    \hline
    \multicolumn{3}{c}{\textbf{Capability Metadata}} \\
    \hline
    \textbf{Name} & \textbf{Value} & \textbf{Comment and/or Description} \\
    \hline
    \textbf{Reference Id} & \texttt{BC-F411AF90} & \\
    \hline
    \textbf{Status} & Established working hierarchy & \\
    \hline
    \textbf{Ancestry} &  & \\
    \hline
    \textbf{Source} &  & \\
    \hline
    \textbf{Enabled By} &  & \\
    \hline
    \textbf{Supports} &  & \\
    \hline
    \textbf{Derived From} &  & \\
    \hline
    \textbf{Reviewed Status} &  & \\
    \hline
  \end{tabularx}
\end{table}

\begin{table}[htbp]
  \centering
  \footnotesize
  \begin{tabularx}{1\textwidth}{|l|l|X|}
    \hline
    \multicolumn{3}{c}{\textbf{Contained Capabilities}} \\
    \hline
    \textbf{Capability Name} & \textbf{Capability Id} & \textbf{Comment} \\
    \hline
    Back-Office Services & \texttt{cap:bc-fa7fbbad} (at \ref{cap:bc-fa7fbbad}) & \\
    \hline
    Facilities Management & \texttt{cap:bc-c8c4c54f} (at \ref{cap:bc-c8c4c54f}) & \\
    \hline
    IT Infrastructure Management & \texttt{cap:bc-edd1575c} (at \ref{cap:bc-edd1575c}) & \\
    \hline
    Communications \& Publicity & \texttt{cap:bc-08ae4743} (at \ref{cap:bc-08ae4743}) & \\
    \hline
    People Management & \texttt{cap:bc-4406e5b2} (at \ref{cap:bc-4406e5b2}) & \\
    \hline
    Operational Oversight & \texttt{cap:bc-cf3a8b92} (at \ref{cap:bc-cf3a8b92}) & \\
    \hline
  \end{tabularx}
\end{table}

\textbf{Intent:} Corporate Support Services identifies the organisational ability to undertake the stated health-service responsibility. The capability describes what the enterprise must be able to do; it does not assert that Harmonia performs or owns that activity.

\textbf{Boundary:} Harmonia relevance may arise because this activity establishes the context in which information is created, needed, interpreted, related, exchanged or presented. Responsibility for the underlying business decision or care activity remains with the appropriate actor or organisation.

\subsection{CT3 - Back-Office Services}\label{cap:bc-fa7fbbad}
\vspace{-0.5em}
\begin{table}[htbp]
  \centering
  \footnotesize
  \begin{tabularx}{1\textwidth}{|l|l|X|}
    \hline
    \multicolumn{3}{c}{\textbf{Capability Metadata}} \\
    \hline
    \textbf{Name} & \textbf{Value} & \textbf{Comment and/or Description} \\
    \hline
    \textbf{Reference Id} & \texttt{BC-FA7FBBAD} & \\
    \hline
    \textbf{Status} & Established working hierarchy & \\
    \hline
    \textbf{Ancestry} &  & \\
    \hline
    \textbf{Source} &  & \\
    \hline
    \textbf{Enabled By} &  & \\
    \hline
    \textbf{Supports} &  & \\
    \hline
    \textbf{Derived From} &  & \\
    \hline
    \textbf{Reviewed Status} & Review meaning, business boundary, information context and any unresolved ancestry. & \\
    \hline
  \end{tabularx}
\end{table}

\begin{table}[htbp]
  \centering
  \footnotesize
  \begin{tabularx}{1\textwidth}{|l|l|X|}
    \hline
    \multicolumn{3}{c}{\textbf{Contained Features}} \\
    \hline
    \textbf{Feature Name} & \textbf{Feature Id} & \textbf{Comment} \\
    \hline
    \multicolumn{3}{|l|}{\emph{No contained Features are currently identified.}} \\
    \hline
  \end{tabularx}
\end{table}

\textbf{Intent:} Back-Office Services identifies the organisational ability to undertake the stated health-service responsibility. The capability describes what the enterprise must be able to do; it does not assert that Harmonia performs or owns that activity.

\textbf{Boundary:} Harmonia relevance may arise because this activity establishes the context in which information is created, needed, interpreted, related, exchanged or presented. Responsibility for the underlying business decision or care activity remains with the appropriate actor or organisation.

\subsection{CT3 - Facilities Management}\label{cap:bc-c8c4c54f}
\vspace{-0.5em}
\begin{table}[htbp]
  \centering
  \footnotesize
  \begin{tabularx}{1\textwidth}{|l|l|X|}
    \hline
    \multicolumn{3}{c}{\textbf{Capability Metadata}} \\
    \hline
    \textbf{Name} & \textbf{Value} & \textbf{Comment and/or Description} \\
    \hline
    \textbf{Reference Id} & \texttt{BC-C8C4C54F} & \\
    \hline
    \textbf{Status} & Established working hierarchy & \\
    \hline
    \textbf{Ancestry} &  & \\
    \hline
    \textbf{Source} &  & \\
    \hline
    \textbf{Enabled By} &  & \\
    \hline
    \textbf{Supports} &  & \\
    \hline
    \textbf{Derived From} &  & \\
    \hline
    \textbf{Reviewed Status} & Review meaning, business boundary, information context and any unresolved ancestry. & \\
    \hline
  \end{tabularx}
\end{table}

\begin{table}[htbp]
  \centering
  \footnotesize
  \begin{tabularx}{1\textwidth}{|l|l|X|}
    \hline
    \multicolumn{3}{c}{\textbf{Contained Features}} \\
    \hline
    \textbf{Feature Name} & \textbf{Feature Id} & \textbf{Comment} \\
    \hline
    \multicolumn{3}{|l|}{\emph{No contained Features are currently identified.}} \\
    \hline
  \end{tabularx}
\end{table}

\textbf{Intent:} Facilities Management identifies the organisational ability to undertake the stated health-service responsibility. The capability describes what the enterprise must be able to do; it does not assert that Harmonia performs or owns that activity.

\textbf{Boundary:} Harmonia relevance may arise because this activity establishes the context in which information is created, needed, interpreted, related, exchanged or presented. Responsibility for the underlying business decision or care activity remains with the appropriate actor or organisation.

\subsection{CT3 - IT Infrastructure Management}\label{cap:bc-edd1575c}
\vspace{-0.5em}
\begin{table}[htbp]
  \centering
  \footnotesize
  \begin{tabularx}{1\textwidth}{|l|l|X|}
    \hline
    \multicolumn{3}{c}{\textbf{Capability Metadata}} \\
    \hline
    \textbf{Name} & \textbf{Value} & \textbf{Comment and/or Description} \\
    \hline
    \textbf{Reference Id} & \texttt{BC-EDD1575C} & \\
    \hline
    \textbf{Status} & Established working hierarchy & \\
    \hline
    \textbf{Ancestry} &  & \\
    \hline
    \textbf{Source} &  & \\
    \hline
    \textbf{Enabled By} &  & \\
    \hline
    \textbf{Supports} &  & \\
    \hline
    \textbf{Derived From} &  & \\
    \hline
    \textbf{Reviewed Status} & Review meaning, business boundary, information context and any unresolved ancestry. & \\
    \hline
  \end{tabularx}
\end{table}

\begin{table}[htbp]
  \centering
  \footnotesize
  \begin{tabularx}{1\textwidth}{|l|l|X|}
    \hline
    \multicolumn{3}{c}{\textbf{Contained Features}} \\
    \hline
    \textbf{Feature Name} & \textbf{Feature Id} & \textbf{Comment} \\
    \hline
    \multicolumn{3}{|l|}{\emph{No contained Features are currently identified.}} \\
    \hline
  \end{tabularx}
\end{table}

\textbf{Intent:} IT Infrastructure Management identifies the organisational ability to undertake the stated health-service responsibility. The capability describes what the enterprise must be able to do; it does not assert that Harmonia performs or owns that activity.

\textbf{Boundary:} Harmonia relevance may arise because this activity establishes the context in which information is created, needed, interpreted, related, exchanged or presented. Responsibility for the underlying business decision or care activity remains with the appropriate actor or organisation.

\subsection{CT3 - Communications \& Publicity}\label{cap:bc-08ae4743}
\vspace{-0.5em}
\begin{table}[htbp]
  \centering
  \footnotesize
  \begin{tabularx}{1\textwidth}{|l|l|X|}
    \hline
    \multicolumn{3}{c}{\textbf{Capability Metadata}} \\
    \hline
    \textbf{Name} & \textbf{Value} & \textbf{Comment and/or Description} \\
    \hline
    \textbf{Reference Id} & \texttt{BC-08AE4743} & \\
    \hline
    \textbf{Status} & Established working hierarchy & \\
    \hline
    \textbf{Ancestry} &  & \\
    \hline
    \textbf{Source} &  & \\
    \hline
    \textbf{Enabled By} &  & \\
    \hline
    \textbf{Supports} &  & \\
    \hline
    \textbf{Derived From} &  & \\
    \hline
    \textbf{Reviewed Status} & Review meaning, business boundary, information context and any unresolved ancestry. & \\
    \hline
  \end{tabularx}
\end{table}

\begin{table}[htbp]
  \centering
  \footnotesize
  \begin{tabularx}{1\textwidth}{|l|l|X|}
    \hline
    \multicolumn{3}{c}{\textbf{Contained Features}} \\
    \hline
    \textbf{Feature Name} & \textbf{Feature Id} & \textbf{Comment} \\
    \hline
    \multicolumn{3}{|l|}{\emph{No contained Features are currently identified.}} \\
    \hline
  \end{tabularx}
\end{table}

\textbf{Intent:} Communications \& Publicity identifies the organisational ability to undertake the stated health-service responsibility. The capability describes what the enterprise must be able to do; it does not assert that Harmonia performs or owns that activity.

\textbf{Boundary:} Harmonia relevance may arise because this activity establishes the context in which information is created, needed, interpreted, related, exchanged or presented. Responsibility for the underlying business decision or care activity remains with the appropriate actor or organisation.

\subsection{CT3 - People Management}\label{cap:bc-4406e5b2}
\vspace{-0.5em}
\begin{table}[htbp]
  \centering
  \footnotesize
  \begin{tabularx}{1\textwidth}{|l|l|X|}
    \hline
    \multicolumn{3}{c}{\textbf{Capability Metadata}} \\
    \hline
    \textbf{Name} & \textbf{Value} & \textbf{Comment and/or Description} \\
    \hline
    \textbf{Reference Id} & \texttt{BC-4406E5B2} & \\
    \hline
    \textbf{Status} & Established working hierarchy & \\
    \hline
    \textbf{Ancestry} &  & \\
    \hline
    \textbf{Source} &  & \\
    \hline
    \textbf{Enabled By} &  & \\
    \hline
    \textbf{Supports} &  & \\
    \hline
    \textbf{Derived From} &  & \\
    \hline
    \textbf{Reviewed Status} & Review meaning, business boundary, information context and any unresolved ancestry. & \\
    \hline
  \end{tabularx}
\end{table}

\begin{table}[htbp]
  \centering
  \footnotesize
  \begin{tabularx}{1\textwidth}{|l|l|X|}
    \hline
    \multicolumn{3}{c}{\textbf{Contained Features}} \\
    \hline
    \textbf{Feature Name} & \textbf{Feature Id} & \textbf{Comment} \\
    \hline
    \multicolumn{3}{|l|}{\emph{No contained Features are currently identified.}} \\
    \hline
  \end{tabularx}
\end{table}

\textbf{Intent:} People Management identifies the organisational ability to undertake the stated health-service responsibility. The capability describes what the enterprise must be able to do; it does not assert that Harmonia performs or owns that activity.

\textbf{Boundary:} Harmonia relevance may arise because this activity establishes the context in which information is created, needed, interpreted, related, exchanged or presented. Responsibility for the underlying business decision or care activity remains with the appropriate actor or organisation.

\subsection{CT3 - Operational Oversight}\label{cap:bc-cf3a8b92}
\vspace{-0.5em}
\begin{table}[htbp]
  \centering
  \footnotesize
  \begin{tabularx}{1\textwidth}{|l|l|X|}
    \hline
    \multicolumn{3}{c}{\textbf{Capability Metadata}} \\
    \hline
    \textbf{Name} & \textbf{Value} & \textbf{Comment and/or Description} \\
    \hline
    \textbf{Reference Id} & \texttt{BC-CF3A8B92} & \\
    \hline
    \textbf{Status} & Established working hierarchy & \\
    \hline
    \textbf{Ancestry} &  & \\
    \hline
    \textbf{Source} &  & \\
    \hline
    \textbf{Enabled By} &  & \\
    \hline
    \textbf{Supports} &  & \\
    \hline
    \textbf{Derived From} &  & \\
    \hline
    \textbf{Reviewed Status} & Review meaning, business boundary, information context and any unresolved ancestry. & \\
    \hline
  \end{tabularx}
\end{table}

\begin{table}[htbp]
  \centering
  \footnotesize
  \begin{tabularx}{1\textwidth}{|l|l|X|}
    \hline
    \multicolumn{3}{c}{\textbf{Contained Features}} \\
    \hline
    \textbf{Feature Name} & \textbf{Feature Id} & \textbf{Comment} \\
    \hline
    \multicolumn{3}{|l|}{\emph{No contained Features are currently identified.}} \\
    \hline
  \end{tabularx}
\end{table}

\textbf{Intent:} Operational Oversight identifies the organisational ability to undertake the stated health-service responsibility. The capability describes what the enterprise must be able to do; it does not assert that Harmonia performs or owns that activity.

\textbf{Boundary:} Harmonia relevance may arise because this activity establishes the context in which information is created, needed, interpreted, related, exchanged or presented. Responsibility for the underlying business decision or care activity remains with the appropriate actor or organisation.

\chapter{CT1 - Health Service Governance — Harmonia-Relevant}\label{cap:bc-1ef62f79}
\noindent\textbf{Reference ID:} \texttt{BC-1EF62F79}\quad \textbf{Status:} Established working hierarchy\par
\noindent\textit{Source: Canonical Business Capability model / working map}\par
\textbf{Intent:} Health Service Governance — Harmonia-Relevant identifies the organisational ability to undertake the stated health-service responsibility. The capability describes what the enterprise must be able to do; it does not assert that Harmonia performs or owns that activity.

\textbf{Boundary:} Harmonia relevance may arise because this activity establishes the context in which information is created, needed, interpreted, related, exchanged or presented. Responsibility for the underlying business decision or care activity remains with the appropriate actor or organisation.

\reviewstatus{Check scope and information relevance; no business authority transfers to Harmonia.}
\section{CT2 - Governance Framework Management}\label{cap:bc-5998319b}
\vspace{-0.5em}
\begin{table}[htbp]
  \centering
  \footnotesize
  \begin{tabularx}{1\textwidth}{|l|l|X|}
    \hline
    \multicolumn{3}{c}{\textbf{Capability Metadata}} \\
    \hline
    \textbf{Name} & \textbf{Value} & \textbf{Comment and/or Description} \\
    \hline
    \textbf{Reference Id} & \texttt{BC-5998319B} & \\
    \hline
    \textbf{Status} & Established working hierarchy & \\
    \hline
    \textbf{Ancestry} &  & \\
    \hline
    \textbf{Source} &  & \\
    \hline
    \textbf{Enabled By} &  & \\
    \hline
    \textbf{Supports} &  & \\
    \hline
    \textbf{Derived From} &  & \\
    \hline
    \textbf{Reviewed Status} &  & \\
    \hline
  \end{tabularx}
\end{table}

\begin{table}[htbp]
  \centering
  \footnotesize
  \begin{tabularx}{1\textwidth}{|l|l|X|}
    \hline
    \multicolumn{3}{c}{\textbf{Contained Capabilities}} \\
    \hline
    \textbf{Capability Name} & \textbf{Capability Id} & \textbf{Comment} \\
    \hline
    Policy / Obligation Governance & \texttt{cap:bc-c01ff756} (at \ref{cap:bc-c01ff756}) & \\
    \hline
    Authority \& Constraint Governance & \texttt{cap:bc-b3ae7e9a} (at \ref{cap:bc-b3ae7e9a}) & \\
    \hline
  \end{tabularx}
\end{table}

\textbf{Intent:} Governance Framework Management identifies the organisational ability to undertake the stated health-service responsibility. The capability describes what the enterprise must be able to do; it does not assert that Harmonia performs or owns that activity.

\textbf{Boundary:} Harmonia relevance may arise because this activity establishes the context in which information is created, needed, interpreted, related, exchanged or presented. Responsibility for the underlying business decision or care activity remains with the appropriate actor or organisation.

\subsection{CT3 - Policy / Obligation Governance}\label{cap:bc-c01ff756}
\vspace{-0.5em}
\begin{table}[htbp]
  \centering
  \footnotesize
  \begin{tabularx}{1\textwidth}{|l|l|X|}
    \hline
    \multicolumn{3}{c}{\textbf{Capability Metadata}} \\
    \hline
    \textbf{Name} & \textbf{Value} & \textbf{Comment and/or Description} \\
    \hline
    \textbf{Reference Id} & \texttt{BC-C01FF756} & \\
    \hline
    \textbf{Status} & Established working hierarchy & \\
    \hline
    \textbf{Ancestry} &  & \\
    \hline
    \textbf{Source} &  & \\
    \hline
    \textbf{Enabled By} &  & \\
    \hline
    \textbf{Supports} &  & \\
    \hline
    \textbf{Derived From} &  & \\
    \hline
    \textbf{Reviewed Status} & Review meaning, business boundary, information context and any unresolved ancestry. & \\
    \hline
  \end{tabularx}
\end{table}

\begin{table}[htbp]
  \centering
  \footnotesize
  \begin{tabularx}{1\textwidth}{|l|l|X|}
    \hline
    \multicolumn{3}{c}{\textbf{Contained Features}} \\
    \hline
    \textbf{Feature Name} & \textbf{Feature Id} & \textbf{Comment} \\
    \hline
    \multicolumn{3}{|l|}{\emph{No contained Features are currently identified.}} \\
    \hline
  \end{tabularx}
\end{table}

\textbf{Intent:} Policy / Obligation Governance identifies the organisational ability to undertake the stated health-service responsibility. The capability describes what the enterprise must be able to do; it does not assert that Harmonia performs or owns that activity.

\textbf{Boundary:} Harmonia relevance may arise because this activity establishes the context in which information is created, needed, interpreted, related, exchanged or presented. Responsibility for the underlying business decision or care activity remains with the appropriate actor or organisation.

\subsection{CT3 - Authority \& Constraint Governance}\label{cap:bc-b3ae7e9a}
\vspace{-0.5em}
\begin{table}[htbp]
  \centering
  \footnotesize
  \begin{tabularx}{1\textwidth}{|l|l|X|}
    \hline
    \multicolumn{3}{c}{\textbf{Capability Metadata}} \\
    \hline
    \textbf{Name} & \textbf{Value} & \textbf{Comment and/or Description} \\
    \hline
    \textbf{Reference Id} & \texttt{BC-B3AE7E9A} & \\
    \hline
    \textbf{Status} & Established working hierarchy & \\
    \hline
    \textbf{Ancestry} &  & \\
    \hline
    \textbf{Source} &  & \\
    \hline
    \textbf{Enabled By} &  & \\
    \hline
    \textbf{Supports} &  & \\
    \hline
    \textbf{Derived From} &  & \\
    \hline
    \textbf{Reviewed Status} & Review meaning, business boundary, information context and any unresolved ancestry. & \\
    \hline
  \end{tabularx}
\end{table}

\begin{table}[htbp]
  \centering
  \footnotesize
  \begin{tabularx}{1\textwidth}{|l|l|X|}
    \hline
    \multicolumn{3}{c}{\textbf{Contained Features}} \\
    \hline
    \textbf{Feature Name} & \textbf{Feature Id} & \textbf{Comment} \\
    \hline
    \multicolumn{3}{|l|}{\emph{No contained Features are currently identified.}} \\
    \hline
  \end{tabularx}
\end{table}

\textbf{Intent:} Authority \& Constraint Governance identifies the organisational ability to undertake the stated health-service responsibility. The capability describes what the enterprise must be able to do; it does not assert that Harmonia performs or owns that activity.

\textbf{Boundary:} Harmonia relevance may arise because this activity establishes the context in which information is created, needed, interpreted, related, exchanged or presented. Responsibility for the underlying business decision or care activity remains with the appropriate actor or organisation.

\section{CT2 - Compliance Governance}\label{cap:bc-15ebfff4}
\vspace{-0.5em}
\begin{table}[htbp]
  \centering
  \footnotesize
  \begin{tabularx}{1\textwidth}{|l|l|X|}
    \hline
    \multicolumn{3}{c}{\textbf{Capability Metadata}} \\
    \hline
    \textbf{Name} & \textbf{Value} & \textbf{Comment and/or Description} \\
    \hline
    \textbf{Reference Id} & \texttt{BC-15EBFFF4} & \\
    \hline
    \textbf{Status} & Established working hierarchy & \\
    \hline
    \textbf{Ancestry} &  & \\
    \hline
    \textbf{Source} &  & \\
    \hline
    \textbf{Enabled By} &  & \\
    \hline
    \textbf{Supports} &  & \\
    \hline
    \textbf{Derived From} &  & \\
    \hline
    \textbf{Reviewed Status} &  & \\
    \hline
  \end{tabularx}
\end{table}

\begin{table}[htbp]
  \centering
  \footnotesize
  \begin{tabularx}{1\textwidth}{|l|l|X|}
    \hline
    \multicolumn{3}{c}{\textbf{Contained Capabilities}} \\
    \hline
    \textbf{Capability Name} & \textbf{Capability Id} & \textbf{Comment} \\
    \hline
    Compliance Assessment Definition & \texttt{cap:bc-9bebd16e} (at \ref{cap:bc-9bebd16e}) & \\
    \hline
    Governance Reporting \& Checkpoint Definition & \texttt{cap:bc-351a6321} (at \ref{cap:bc-351a6321}) & \\
    \hline
  \end{tabularx}
\end{table}

\textbf{Intent:} Compliance Governance identifies the organisational ability to undertake the stated health-service responsibility. The capability describes what the enterprise must be able to do; it does not assert that Harmonia performs or owns that activity.

\textbf{Boundary:} Harmonia relevance may arise because this activity establishes the context in which information is created, needed, interpreted, related, exchanged or presented. Responsibility for the underlying business decision or care activity remains with the appropriate actor or organisation.

\subsection{CT3 - Compliance Assessment Definition}\label{cap:bc-9bebd16e}
\vspace{-0.5em}
\begin{table}[htbp]
  \centering
  \footnotesize
  \begin{tabularx}{1\textwidth}{|l|l|X|}
    \hline
    \multicolumn{3}{c}{\textbf{Capability Metadata}} \\
    \hline
    \textbf{Name} & \textbf{Value} & \textbf{Comment and/or Description} \\
    \hline
    \textbf{Reference Id} & \texttt{BC-9BEBD16E} & \\
    \hline
    \textbf{Status} & Established working hierarchy & \\
    \hline
    \textbf{Ancestry} &  & \\
    \hline
    \textbf{Source} &  & \\
    \hline
    \textbf{Enabled By} &  & \\
    \hline
    \textbf{Supports} &  & \\
    \hline
    \textbf{Derived From} &  & \\
    \hline
    \textbf{Reviewed Status} & Review meaning, business boundary, information context and any unresolved ancestry. & \\
    \hline
  \end{tabularx}
\end{table}

\begin{table}[htbp]
  \centering
  \footnotesize
  \begin{tabularx}{1\textwidth}{|l|l|X|}
    \hline
    \multicolumn{3}{c}{\textbf{Contained Features}} \\
    \hline
    \textbf{Feature Name} & \textbf{Feature Id} & \textbf{Comment} \\
    \hline
    \multicolumn{3}{|l|}{\emph{No contained Features are currently identified.}} \\
    \hline
  \end{tabularx}
\end{table}

\textbf{Intent:} Compliance Assessment Definition identifies the organisational ability to undertake the stated health-service responsibility. The capability describes what the enterprise must be able to do; it does not assert that Harmonia performs or owns that activity.

\textbf{Boundary:} Harmonia relevance may arise because this activity establishes the context in which information is created, needed, interpreted, related, exchanged or presented. Responsibility for the underlying business decision or care activity remains with the appropriate actor or organisation.

\subsection{CT3 - Governance Reporting \& Checkpoint Definition}\label{cap:bc-351a6321}
\vspace{-0.5em}
\begin{table}[htbp]
  \centering
  \footnotesize
  \begin{tabularx}{1\textwidth}{|l|l|X|}
    \hline
    \multicolumn{3}{c}{\textbf{Capability Metadata}} \\
    \hline
    \textbf{Name} & \textbf{Value} & \textbf{Comment and/or Description} \\
    \hline
    \textbf{Reference Id} & \texttt{BC-351A6321} & \\
    \hline
    \textbf{Status} & Established working hierarchy & \\
    \hline
    \textbf{Ancestry} &  & \\
    \hline
    \textbf{Source} &  & \\
    \hline
    \textbf{Enabled By} &  & \\
    \hline
    \textbf{Supports} &  & \\
    \hline
    \textbf{Derived From} &  & \\
    \hline
    \textbf{Reviewed Status} & Review meaning, business boundary, information context and any unresolved ancestry. & \\
    \hline
  \end{tabularx}
\end{table}

\begin{table}[htbp]
  \centering
  \footnotesize
  \begin{tabularx}{1\textwidth}{|l|l|X|}
    \hline
    \multicolumn{3}{c}{\textbf{Contained Features}} \\
    \hline
    \textbf{Feature Name} & \textbf{Feature Id} & \textbf{Comment} \\
    \hline
    \multicolumn{3}{|l|}{\emph{No contained Features are currently identified.}} \\
    \hline
  \end{tabularx}
\end{table}

\textbf{Intent:} Governance Reporting \& Checkpoint Definition identifies the organisational ability to undertake the stated health-service responsibility. The capability describes what the enterprise must be able to do; it does not assert that Harmonia performs or owns that activity.

\textbf{Boundary:} Harmonia relevance may arise because this activity establishes the context in which information is created, needed, interpreted, related, exchanged or presented. Responsibility for the underlying business decision or care activity remains with the appropriate actor or organisation.

\chapter{CT1 - Health Service Assurance}\label{cap:bc-62520466}
\noindent\textbf{Reference ID:} \texttt{BC-62520466}\quad \textbf{Status:} Established working hierarchy\par
\noindent\textit{Source: Canonical Business Capability model / working map}\par
\textbf{Intent:} Health Service Assurance identifies the organisational ability to undertake the stated health-service responsibility. The capability describes what the enterprise must be able to do; it does not assert that Harmonia performs or owns that activity.

\textbf{Boundary:} Harmonia relevance may arise because this activity establishes the context in which information is created, needed, interpreted, related, exchanged or presented. Responsibility for the underlying business decision or care activity remains with the appropriate actor or organisation.

\reviewstatus{Check scope and information relevance; no business authority transfers to Harmonia.}
\section{CT2 - Health Service Assurance Specification}\label{cap:bc-a0266334}
\vspace{-0.5em}
\begin{table}[htbp]
  \centering
  \footnotesize
  \begin{tabularx}{1\textwidth}{|l|l|X|}
    \hline
    \multicolumn{3}{c}{\textbf{Capability Metadata}} \\
    \hline
    \textbf{Name} & \textbf{Value} & \textbf{Comment and/or Description} \\
    \hline
    \textbf{Reference Id} & \texttt{BC-A0266334} & \\
    \hline
    \textbf{Status} & Established working hierarchy & \\
    \hline
    \textbf{Ancestry} &  & \\
    \hline
    \textbf{Source} &  & \\
    \hline
    \textbf{Enabled By} &  & \\
    \hline
    \textbf{Supports} &  & \\
    \hline
    \textbf{Derived From} &  & \\
    \hline
    \textbf{Reviewed Status} &  & \\
    \hline
  \end{tabularx}
\end{table}

\begin{table}[htbp]
  \centering
  \footnotesize
  \begin{tabularx}{1\textwidth}{|l|l|X|}
    \hline
    \multicolumn{3}{c}{\textbf{Contained Capabilities}} \\
    \hline
    \textbf{Capability Name} & \textbf{Capability Id} & \textbf{Comment} \\
    \hline
    Assurance Approach Definition & \texttt{cap:bc-ba00f585} (at \ref{cap:bc-ba00f585}) & \\
    \hline
    Assurance Evaluation Criteria Governance & \texttt{cap:bc-28f181c3} (at \ref{cap:bc-28f181c3}) & \\
    \hline
  \end{tabularx}
\end{table}

\textbf{Intent:} Health Service Assurance Specification identifies the organisational ability to undertake the stated health-service responsibility. The capability describes what the enterprise must be able to do; it does not assert that Harmonia performs or owns that activity.

\textbf{Boundary:} Harmonia relevance may arise because this activity establishes the context in which information is created, needed, interpreted, related, exchanged or presented. Responsibility for the underlying business decision or care activity remains with the appropriate actor or organisation.

\subsection{CT3 - Assurance Approach Definition}\label{cap:bc-ba00f585}
\vspace{-0.5em}
\begin{table}[htbp]
  \centering
  \footnotesize
  \begin{tabularx}{1\textwidth}{|l|l|X|}
    \hline
    \multicolumn{3}{c}{\textbf{Capability Metadata}} \\
    \hline
    \textbf{Name} & \textbf{Value} & \textbf{Comment and/or Description} \\
    \hline
    \textbf{Reference Id} & \texttt{BC-BA00F585} & \\
    \hline
    \textbf{Status} & Established working hierarchy & \\
    \hline
    \textbf{Ancestry} &  & \\
    \hline
    \textbf{Source} &  & \\
    \hline
    \textbf{Enabled By} &  & \\
    \hline
    \textbf{Supports} &  & \\
    \hline
    \textbf{Derived From} &  & \\
    \hline
    \textbf{Reviewed Status} & Review meaning, business boundary, information context and any unresolved ancestry. & \\
    \hline
  \end{tabularx}
\end{table}

\begin{table}[htbp]
  \centering
  \footnotesize
  \begin{tabularx}{1\textwidth}{|l|l|X|}
    \hline
    \multicolumn{3}{c}{\textbf{Contained Features}} \\
    \hline
    \textbf{Feature Name} & \textbf{Feature Id} & \textbf{Comment} \\
    \hline
    \multicolumn{3}{|l|}{\emph{No contained Features are currently identified.}} \\
    \hline
  \end{tabularx}
\end{table}

\textbf{Intent:} Assurance Approach Definition identifies the organisational ability to undertake the stated health-service responsibility. The capability describes what the enterprise must be able to do; it does not assert that Harmonia performs or owns that activity.

\textbf{Boundary:} Harmonia relevance may arise because this activity establishes the context in which information is created, needed, interpreted, related, exchanged or presented. Responsibility for the underlying business decision or care activity remains with the appropriate actor or organisation.

\subsection{CT3 - Assurance Evaluation Criteria Governance}\label{cap:bc-28f181c3}
\vspace{-0.5em}
\begin{table}[htbp]
  \centering
  \footnotesize
  \begin{tabularx}{1\textwidth}{|l|l|X|}
    \hline
    \multicolumn{3}{c}{\textbf{Capability Metadata}} \\
    \hline
    \textbf{Name} & \textbf{Value} & \textbf{Comment and/or Description} \\
    \hline
    \textbf{Reference Id} & \texttt{BC-28F181C3} & \\
    \hline
    \textbf{Status} & Established working hierarchy & \\
    \hline
    \textbf{Ancestry} &  & \\
    \hline
    \textbf{Source} &  & \\
    \hline
    \textbf{Enabled By} &  & \\
    \hline
    \textbf{Supports} &  & \\
    \hline
    \textbf{Derived From} &  & \\
    \hline
    \textbf{Reviewed Status} & Review meaning, business boundary, information context and any unresolved ancestry. & \\
    \hline
  \end{tabularx}
\end{table}

\begin{table}[htbp]
  \centering
  \footnotesize
  \begin{tabularx}{1\textwidth}{|l|l|X|}
    \hline
    \multicolumn{3}{c}{\textbf{Contained Features}} \\
    \hline
    \textbf{Feature Name} & \textbf{Feature Id} & \textbf{Comment} \\
    \hline
    \multicolumn{3}{|l|}{\emph{No contained Features are currently identified.}} \\
    \hline
  \end{tabularx}
\end{table}

\textbf{Intent:} Assurance Evaluation Criteria Governance identifies the organisational ability to undertake the stated health-service responsibility. The capability describes what the enterprise must be able to do; it does not assert that Harmonia performs or owns that activity.

\textbf{Boundary:} Harmonia relevance may arise because this activity establishes the context in which information is created, needed, interpreted, related, exchanged or presented. Responsibility for the underlying business decision or care activity remains with the appropriate actor or organisation.

\section{CT2 - Assurance Evaluation}\label{cap:bc-11ace749}
\vspace{-0.5em}
\begin{table}[htbp]
  \centering
  \footnotesize
  \begin{tabularx}{1\textwidth}{|l|l|X|}
    \hline
    \multicolumn{3}{c}{\textbf{Capability Metadata}} \\
    \hline
    \textbf{Name} & \textbf{Value} & \textbf{Comment and/or Description} \\
    \hline
    \textbf{Reference Id} & \texttt{BC-11ACE749} & \\
    \hline
    \textbf{Status} & Established working hierarchy & \\
    \hline
    \textbf{Ancestry} &  & \\
    \hline
    \textbf{Source} &  & \\
    \hline
    \textbf{Enabled By} &  & \\
    \hline
    \textbf{Supports} &  & \\
    \hline
    \textbf{Derived From} &  & \\
    \hline
    \textbf{Reviewed Status} &  & \\
    \hline
  \end{tabularx}
\end{table}

\begin{table}[htbp]
  \centering
  \footnotesize
  \begin{tabularx}{1\textwidth}{|l|l|X|}
    \hline
    \multicolumn{3}{c}{\textbf{Contained Capabilities}} \\
    \hline
    \textbf{Capability Name} & \textbf{Capability Id} & \textbf{Comment} \\
    \hline
    Assurance Evidence Assessment & \texttt{cap:bc-ef506e30} (at \ref{cap:bc-ef506e30}) & \\
    \hline
    Assurance Conclusion Determination & \texttt{cap:bc-bc9e1e00} (at \ref{cap:bc-bc9e1e00}) & \\
    \hline
  \end{tabularx}
\end{table}

\textbf{Intent:} Assurance Evaluation identifies the organisational ability to undertake the stated health-service responsibility. The capability describes what the enterprise must be able to do; it does not assert that Harmonia performs or owns that activity.

\textbf{Boundary:} Harmonia relevance may arise because this activity establishes the context in which information is created, needed, interpreted, related, exchanged or presented. Responsibility for the underlying business decision or care activity remains with the appropriate actor or organisation.

\subsection{CT3 - Assurance Evidence Assessment}\label{cap:bc-ef506e30}
\vspace{-0.5em}
\begin{table}[htbp]
  \centering
  \footnotesize
  \begin{tabularx}{1\textwidth}{|l|l|X|}
    \hline
    \multicolumn{3}{c}{\textbf{Capability Metadata}} \\
    \hline
    \textbf{Name} & \textbf{Value} & \textbf{Comment and/or Description} \\
    \hline
    \textbf{Reference Id} & \texttt{BC-EF506E30} & \\
    \hline
    \textbf{Status} & Established working hierarchy & \\
    \hline
    \textbf{Ancestry} &  & \\
    \hline
    \textbf{Source} &  & \\
    \hline
    \textbf{Enabled By} &  & \\
    \hline
    \textbf{Supports} &  & \\
    \hline
    \textbf{Derived From} &  & \\
    \hline
    \textbf{Reviewed Status} & Review meaning, business boundary, information context and any unresolved ancestry. & \\
    \hline
  \end{tabularx}
\end{table}

\begin{table}[htbp]
  \centering
  \footnotesize
  \begin{tabularx}{1\textwidth}{|l|l|X|}
    \hline
    \multicolumn{3}{c}{\textbf{Contained Features}} \\
    \hline
    \textbf{Feature Name} & \textbf{Feature Id} & \textbf{Comment} \\
    \hline
    \multicolumn{3}{|l|}{\emph{No contained Features are currently identified.}} \\
    \hline
  \end{tabularx}
\end{table}

\textbf{Intent:} Assurance Evidence Assessment identifies the organisational ability to undertake the stated health-service responsibility. The capability describes what the enterprise must be able to do; it does not assert that Harmonia performs or owns that activity.

\textbf{Boundary:} Harmonia relevance may arise because this activity establishes the context in which information is created, needed, interpreted, related, exchanged or presented. Responsibility for the underlying business decision or care activity remains with the appropriate actor or organisation.

\subsection{CT3 - Assurance Conclusion Determination}\label{cap:bc-bc9e1e00}
\vspace{-0.5em}
\begin{table}[htbp]
  \centering
  \footnotesize
  \begin{tabularx}{1\textwidth}{|l|l|X|}
    \hline
    \multicolumn{3}{c}{\textbf{Capability Metadata}} \\
    \hline
    \textbf{Name} & \textbf{Value} & \textbf{Comment and/or Description} \\
    \hline
    \textbf{Reference Id} & \texttt{BC-BC9E1E00} & \\
    \hline
    \textbf{Status} & Established working hierarchy & \\
    \hline
    \textbf{Ancestry} &  & \\
    \hline
    \textbf{Source} &  & \\
    \hline
    \textbf{Enabled By} &  & \\
    \hline
    \textbf{Supports} &  & \\
    \hline
    \textbf{Derived From} &  & \\
    \hline
    \textbf{Reviewed Status} & Review meaning, business boundary, information context and any unresolved ancestry. & \\
    \hline
  \end{tabularx}
\end{table}

\begin{table}[htbp]
  \centering
  \footnotesize
  \begin{tabularx}{1\textwidth}{|l|l|X|}
    \hline
    \multicolumn{3}{c}{\textbf{Contained Features}} \\
    \hline
    \textbf{Feature Name} & \textbf{Feature Id} & \textbf{Comment} \\
    \hline
    \multicolumn{3}{|l|}{\emph{No contained Features are currently identified.}} \\
    \hline
  \end{tabularx}
\end{table}

\textbf{Intent:} Assurance Conclusion Determination identifies the organisational ability to undertake the stated health-service responsibility. The capability describes what the enterprise must be able to do; it does not assert that Harmonia performs or owns that activity.

\textbf{Boundary:} Harmonia relevance may arise because this activity establishes the context in which information is created, needed, interpreted, related, exchanged or presented. Responsibility for the underlying business decision or care activity remains with the appropriate actor or organisation.

\part{Business Enabling Capabilities and Features}
\chapter{CT1 - Entity Information Enablement}\label{cap:bec-f49a0d09}
\noindent\textbf{Reference ID:} \texttt{BEC-F49A0D09}\quad \textbf{Status:} Candidate ancestry\par
Intent. Entity Information Enablement is a candidate enabling capability within Business Enabling Capability layer. It frames the reusable system-supported behaviours needed to support this healthcare context without prescribing a product, interface, technology or deployment boundary.

Boundary. This capability supports the business activity and does not inherit the authority or accountability of the Business Capability it enables. Its child Features are review candidates until ancestry and semantics are adjudicated.

\section{CT2 - Client Administration}\label{cap:bec-75cec238}
\vspace{-0.5em}
\begin{table}[htbp]
  \centering
  \footnotesize
  \begin{tabularx}{1\textwidth}{|l|l|X|}
    \hline
    \multicolumn{3}{c}{\textbf{Capability Metadata}} \\
    \hline
    \textbf{Name} & \textbf{Value} & \textbf{Comment and/or Description} \\
    \hline
    \textbf{Reference Id} & \texttt{BEC-75CEC238} & \\
    \hline
    \textbf{Status} & Candidate ancestry & \\
    \hline
    \textbf{Ancestry} &  & \\
    \hline
    \textbf{Source} &  & \\
    \hline
    \textbf{Enabled By} &  & \\
    \hline
    \textbf{Supports} &  & \\
    \hline
    \textbf{Derived From} &  & \\
    \hline
    \textbf{Reviewed Status} &  & \\
    \hline
  \end{tabularx}
\end{table}

\begin{table}[htbp]
  \centering
  \footnotesize
  \begin{tabularx}{1\textwidth}{|l|l|X|}
    \hline
    \multicolumn{3}{c}{\textbf{Contained Capabilities}} \\
    \hline
    \textbf{Capability Name} & \textbf{Capability Id} & \textbf{Comment} \\
    \hline
    Person Identity & \texttt{cap:bec-31e69256} (at \ref{cap:bec-31e69256}) & \\
    \hline
    Healthcare Subject & \texttt{cap:bec-edce7c88} (at \ref{cap:bec-edce7c88}) & \\
    \hline
    Client Relationship & \texttt{cap:bec-2494865f} (at \ref{cap:bec-2494865f}) & \\
    \hline
    Client Privacy & \texttt{cap:bec-adc0cf33} (at \ref{cap:bec-adc0cf33}) & \\
    \hline
  \end{tabularx}
\end{table}

Intent. Client Administration is a candidate enabling capability within Entity Information Enablement. It frames the reusable system-supported behaviours needed to support this healthcare context without prescribing a product, interface, technology or deployment boundary.

Boundary. This capability supports the business activity and does not inherit the authority or accountability of the Business Capability it enables. Its child Features are review candidates until ancestry and semantics are adjudicated.

\subsection{CT3 - Person Identity}\label{cap:bec-31e69256}
\vspace{-0.5em}
\begin{table}[htbp]
  \centering
  \footnotesize
  \begin{tabularx}{1\textwidth}{|l|l|X|}
    \hline
    \multicolumn{3}{c}{\textbf{Capability Metadata}} \\
    \hline
    \textbf{Name} & \textbf{Value} & \textbf{Comment and/or Description} \\
    \hline
    \textbf{Reference Id} & \texttt{BEC-31E69256} & \\
    \hline
    \textbf{Status} & Candidate ancestry & \\
    \hline
    \textbf{Ancestry} &  & \\
    \hline
    \textbf{Source} &  & \\
    \hline
    \textbf{Enabled By} &  & \\
    \hline
    \textbf{Supports} &  & \\
    \hline
    \textbf{Derived From} &  & \\
    \hline
    \textbf{Reviewed Status} &  & \\
    \hline
  \end{tabularx}
\end{table}

\begin{table}[htbp]
  \centering
  \footnotesize
  \begin{tabularx}{1\textwidth}{|l|l|X|}
    \hline
    \multicolumn{3}{c}{\textbf{Contained Features}} \\
    \hline
    \textbf{Feature Name} & \textbf{Feature Id} & \textbf{Comment} \\
    \hline
    Identifier Resolution & \texttt{cap:bec-1acd822d} (at \ref{cap:bec-1acd822d}) & \\
    \hline
    Cross-Authority Identifier Correlation & \texttt{cap:bec-3e1d6ba5} (at \ref{cap:bec-3e1d6ba5}) & \\
    \hline
    Identity Alias Association & \texttt{cap:bec-905c884b} (at \ref{cap:bec-905c884b}) & \\
    \hline
    Governed Identity Correction & \texttt{cap:bec-f44a13c3} (at \ref{cap:bec-f44a13c3}) & \\
    \hline
    Identity Management & \texttt{cap:bec-60f846cb} (at \ref{cap:bec-60f846cb}) & \\
    \hline
    Identity Correlation & \texttt{cap:bec-d5d59967} (at \ref{cap:bec-d5d59967}) & \\
    \hline
    Identity Correction & \texttt{cap:bec-44123e71} (at \ref{cap:bec-44123e71}) & \\
    \hline
    Relationship Management & \texttt{cap:bec-c7981054} (at \ref{cap:bec-c7981054}) & \\
    \hline
    Identity Discovery & \texttt{cap:bec-07eb6044} (at \ref{cap:bec-07eb6044}) & \\
    \hline
  \end{tabularx}
\end{table}

Intent. Person Identity is a candidate enabling capability within Client Administration. It frames the reusable system-supported behaviours needed to support this healthcare context without prescribing a product, interface, technology or deployment boundary.

Boundary. This capability supports the business activity and does not inherit the authority or accountability of the Business Capability it enables. Its child Features are review candidates until ancestry and semantics are adjudicated.

FHIR R5 conceptual analogues considered (not implementation mappings): Person, Patient, RelatedPerson.

\subsubsection{Feature - Identifier Resolution}\label{cap:bec-1acd822d}
\vspace{-1em}
\begin{table}[htbp]
  \centering
  \footnotesize
  \begin{tabularx}{1\textwidth}{|l|l|X|}
    \hline
    \textbf{Metadata Name} & \textbf{Metadata Value} & \textbf{Description} \\
    \hline
    \textbf{Reference Id} & \texttt{BEC-1ACD822D} & \\
    \hline
    \textbf{Status} &  & \\
    \hline
    \textbf{Ancestry} &  & \\
    \hline
    \textbf{Source} & FEAT-EM-01 & \\
    \hline
    \textbf{Enabled By} & \ref{cap:ec-461b8f05 & \\
    \hline
    \textbf{Supports} &  & \\
    \hline
    \textbf{Derived From} &  & \\
    \hline
    \textbf{Reviewed Status} & MH-Approved & \\
    \hline
  \end{tabularx}
\end{table}

Intent. Identifier Resolution provides a distinct enabling behaviour in the context of Person Identity. This is an existing named Feature; its placement under the candidate CT3 remains subject to review.

\noindent\textbf{Feature boundary.} This behaviour is scoped to Person Identity; an identically named Feature elsewhere may carry different business meaning.\par
\featuretable{
Function & \texttt{FN-D225B2B0} & Resolve Person Identifier & Catalogue \\
Service & \texttt{SV-63F14B3F} & Person Identifier Resolution & Linked to preceding Function in source \\
}
}
\reviewcomment{None.}
\subsubsection{Feature - Cross-Authority Identifier Correlation}\label{cap:bec-3e1d6ba5}
\vspace{-1em}
\begin{table}[htbp]
  \centering
  \footnotesize
  \begin{tabularx}{1\textwidth}{|l|l|X|}
    \hline
    \textbf{Metadata Name} & \textbf{Metadata Value} & \textbf{Description} \\
    \hline
    \textbf{Reference Id} & \texttt{BEC-3E1D6BA5} & \\
    \hline
    \textbf{Status} & Existing Feature; candidate CT3 placement & \\
    \hline
    \textbf{Ancestry} &  & \\
    \hline
    \textbf{Source} & FEAT-EM-02 & \\
    \hline
    \textbf{Enabled By} &  & \\
    \hline
    \textbf{Supports} &  & \\
    \hline
    \textbf{Derived From} &  & \\
    \hline
    \textbf{Reviewed Status} & Confirm existing Feature meaning and proposed CT3 ownership. & \\
    \hline
  \end{tabularx}
\end{table}

Intent. Cross-Authority Identifier Correlation provides a distinct enabling behaviour in the context of Person Identity. This is an existing named Feature; its placement under the candidate CT3 remains subject to review.

\noindent\textbf{Feature boundary.} This behaviour is scoped to Person Identity; an identically named Feature elsewhere may carry different business meaning.\par
\featuretable{
Function & \texttt{FN-2ECE71F2} & Correlate Person Identifiers & Catalogue \\
Service & \texttt{SV-81DA437F} & Person Identifier Correlation & Linked to preceding Function in source \\
}
\subsubsection{Feature - Identity Alias Association}\label{cap:bec-905c884b}
\vspace{-1em}
\begin{table}[htbp]
  \centering
  \footnotesize
  \begin{tabularx}{1\textwidth}{|l|l|X|}
    \hline
    \textbf{Metadata Name} & \textbf{Metadata Value} & \textbf{Description} \\
    \hline
    \textbf{Reference Id} & \texttt{BEC-905C884B} & \\
    \hline
    \textbf{Status} & Existing Feature; candidate CT3 placement & \\
    \hline
    \textbf{Ancestry} &  & \\
    \hline
    \textbf{Source} & FEAT-EM-03 & \\
    \hline
    \textbf{Enabled By} &  & \\
    \hline
    \textbf{Supports} &  & \\
    \hline
    \textbf{Derived From} &  & \\
    \hline
    \textbf{Reviewed Status} & Confirm existing Feature meaning and proposed CT3 ownership. & \\
    \hline
  \end{tabularx}
\end{table}

Intent. Identity Alias Association provides a distinct enabling behaviour in the context of Person Identity. This is an existing named Feature; its placement under the candidate CT3 remains subject to review.

\noindent\textbf{Feature boundary.} This behaviour is scoped to Person Identity; an identically named Feature elsewhere may carry different business meaning.\par
\featuretable{
Function & \texttt{FN-FD73BCD6} & Maintain Person Identity Aliases & Catalogue \\
}
\subsubsection{Feature - Governed Identity Correction}\label{cap:bec-f44a13c3}
\vspace{-1em}
\begin{table}[htbp]
  \centering
  \footnotesize
  \begin{tabularx}{1\textwidth}{|l|l|X|}
    \hline
    \textbf{Metadata Name} & \textbf{Metadata Value} & \textbf{Description} \\
    \hline
    \textbf{Reference Id} & \texttt{BEC-F44A13C3} & \\
    \hline
    \textbf{Status} & Existing Feature; candidate CT3 placement & \\
    \hline
    \textbf{Ancestry} &  & \\
    \hline
    \textbf{Source} & FEAT-EM-04 & \\
    \hline
    \textbf{Enabled By} &  & \\
    \hline
    \textbf{Supports} &  & \\
    \hline
    \textbf{Derived From} &  & \\
    \hline
    \textbf{Reviewed Status} & Confirm existing Feature meaning and proposed CT3 ownership. & \\
    \hline
  \end{tabularx}
\end{table}

Intent. Governed Identity Correction provides a distinct enabling behaviour in the context of Person Identity. This is an existing named Feature; its placement under the candidate CT3 remains subject to review.

\noindent\textbf{Feature boundary.} This behaviour is scoped to Person Identity; an identically named Feature elsewhere may carry different business meaning.\par
\featuretable{
Function & \texttt{FN-82D55078} & Apply Governed Person Identity Correction & Catalogue \\
Service & \texttt{SV-468C2686} & Person Identity Correction & Linked to preceding Function in source \\
Process & \texttt{PR-D828E7C2} & Governed Person Identity Correction Process & Linked to preceding Function in source \\
}
\subsubsection{Feature - Identity Management}\label{cap:bec-60f846cb}
\vspace{-1em}
\begin{table}[htbp]
  \centering
  \footnotesize
  \begin{tabularx}{1\textwidth}{|l|l|X|}
    \hline
    \textbf{Metadata Name} & \textbf{Metadata Value} & \textbf{Description} \\
    \hline
    \textbf{Reference Id} & \texttt{BEC-60F846CB} & \\
    \hline
    \textbf{Status} & Generated candidate & \\
    \hline
    \textbf{Ancestry} &  & \\
    \hline
    \textbf{Source} & FHIR-informed candidate & \\
    \hline
    \textbf{Enabled By} &  & \\
    \hline
    \textbf{Supports} &  & \\
    \hline
    \textbf{Derived From} &  & \\
    \hline
    \textbf{Reviewed Status} & Generated candidate: keep / modify / relocate / reject; verify whether independently testable. & \\
    \hline
  \end{tabularx}
\end{table}

Intent. Identity Management is proposed to support the underlying healthcare concept represented by Person Identity. The proposal concerns behaviour and semantic responsibility, not a FHIR field, operation, resource instance or prescribed data model.

Conceptual evidence: HL7 FHIR R5 Person, Patient, RelatedPerson.

\noindent\textbf{Feature boundary.} This behaviour is scoped to Person Identity; an identically named Feature elsewhere may carry different business meaning.\par
\featuretable{
--- & --- & \emph{Not yet captured} & \emph{Complete during review} \\
}
\subsubsection{Feature - Identity Correlation}\label{cap:bec-d5d59967}
\vspace{-1em}
\begin{table}[htbp]
  \centering
  \footnotesize
  \begin{tabularx}{1\textwidth}{|l|l|X|}
    \hline
    \textbf{Metadata Name} & \textbf{Metadata Value} & \textbf{Description} \\
    \hline
    \textbf{Reference Id} & \texttt{BEC-D5D59967} & \\
    \hline
    \textbf{Status} & Generated candidate & \\
    \hline
    \textbf{Ancestry} &  & \\
    \hline
    \textbf{Source} & FHIR-informed candidate & \\
    \hline
    \textbf{Enabled By} &  & \\
    \hline
    \textbf{Supports} &  & \\
    \hline
    \textbf{Derived From} &  & \\
    \hline
    \textbf{Reviewed Status} & Generated candidate: keep / modify / relocate / reject; verify whether independently testable. & \\
    \hline
  \end{tabularx}
\end{table}

Intent. Identity Correlation is proposed to support the underlying healthcare concept represented by Person Identity. The proposal concerns behaviour and semantic responsibility, not a FHIR field, operation, resource instance or prescribed data model.

Conceptual evidence: HL7 FHIR R5 Person, Patient, RelatedPerson.

\noindent\textbf{Feature boundary.} This behaviour is scoped to Person Identity; an identically named Feature elsewhere may carry different business meaning.\par
\featuretable{
--- & --- & \emph{Not yet captured} & \emph{Complete during review} \\
}
\subsubsection{Feature - Identity Correction}\label{cap:bec-44123e71}
\vspace{-1em}
\begin{table}[htbp]
  \centering
  \footnotesize
  \begin{tabularx}{1\textwidth}{|l|l|X|}
    \hline
    \textbf{Metadata Name} & \textbf{Metadata Value} & \textbf{Description} \\
    \hline
    \textbf{Reference Id} & \texttt{BEC-44123E71} & \\
    \hline
    \textbf{Status} & Generated candidate & \\
    \hline
    \textbf{Ancestry} &  & \\
    \hline
    \textbf{Source} & FHIR-informed candidate & \\
    \hline
    \textbf{Enabled By} &  & \\
    \hline
    \textbf{Supports} &  & \\
    \hline
    \textbf{Derived From} &  & \\
    \hline
    \textbf{Reviewed Status} & Generated candidate: keep / modify / relocate / reject; verify whether independently testable. & \\
    \hline
  \end{tabularx}
\end{table}

Intent. Identity Correction is proposed to support the underlying healthcare concept represented by Person Identity. The proposal concerns behaviour and semantic responsibility, not a FHIR field, operation, resource instance or prescribed data model.

Conceptual evidence: HL7 FHIR R5 Person, Patient, RelatedPerson.

\noindent\textbf{Feature boundary.} This behaviour is scoped to Person Identity; an identically named Feature elsewhere may carry different business meaning.\par
\featuretable{
--- & --- & \emph{Not yet captured} & \emph{Complete during review} \\
}
\subsubsection{Feature - Relationship Management}\label{cap:bec-c7981054}
\vspace{-1em}
\begin{table}[htbp]
  \centering
  \footnotesize
  \begin{tabularx}{1\textwidth}{|l|l|X|}
    \hline
    \textbf{Metadata Name} & \textbf{Metadata Value} & \textbf{Description} \\
    \hline
    \textbf{Reference Id} & \texttt{BEC-C7981054} & \\
    \hline
    \textbf{Status} & Generated candidate & \\
    \hline
    \textbf{Ancestry} &  & \\
    \hline
    \textbf{Source} & FHIR-informed candidate & \\
    \hline
    \textbf{Enabled By} &  & \\
    \hline
    \textbf{Supports} &  & \\
    \hline
    \textbf{Derived From} &  & \\
    \hline
    \textbf{Reviewed Status} & Generated candidate: keep / modify / relocate / reject; verify whether independently testable. & \\
    \hline
  \end{tabularx}
\end{table}

Intent. Relationship Management is proposed to support the underlying healthcare concept represented by Person Identity. The proposal concerns behaviour and semantic responsibility, not a FHIR field, operation, resource instance or prescribed data model.

Conceptual evidence: HL7 FHIR R5 Person, Patient, RelatedPerson.

\noindent\textbf{Feature boundary.} This behaviour is scoped to Person Identity; an identically named Feature elsewhere may carry different business meaning.\par
\featuretable{
--- & --- & \emph{Not yet captured} & \emph{Complete during review} \\
}
\subsubsection{Feature - Identity Discovery}\label{cap:bec-07eb6044}
\vspace{-1em}
\begin{table}[htbp]
  \centering
  \footnotesize
  \begin{tabularx}{1\textwidth}{|l|l|X|}
    \hline
    \textbf{Metadata Name} & \textbf{Metadata Value} & \textbf{Description} \\
    \hline
    \textbf{Reference Id} & \texttt{BEC-07EB6044} & \\
    \hline
    \textbf{Status} & Generated candidate & \\
    \hline
    \textbf{Ancestry} &  & \\
    \hline
    \textbf{Source} & FHIR-informed candidate & \\
    \hline
    \textbf{Enabled By} &  & \\
    \hline
    \textbf{Supports} &  & \\
    \hline
    \textbf{Derived From} &  & \\
    \hline
    \textbf{Reviewed Status} & Generated candidate: keep / modify / relocate / reject; verify whether independently testable. & \\
    \hline
  \end{tabularx}
\end{table}

Intent. Identity Discovery is proposed to support the underlying healthcare concept represented by Person Identity. The proposal concerns behaviour and semantic responsibility, not a FHIR field, operation, resource instance or prescribed data model.

Conceptual evidence: HL7 FHIR R5 Person, Patient, RelatedPerson.

\noindent\textbf{Feature boundary.} This behaviour is scoped to Person Identity; an identically named Feature elsewhere may carry different business meaning.\par
\featuretable{
--- & --- & \emph{Not yet captured} & \emph{Complete during review} \\
}

\subsection{CT3 - Healthcare Subject}\label{cap:bec-edce7c88}
\vspace{-0.5em}
\begin{table}[htbp]
  \centering
  \footnotesize
  \begin{tabularx}{1\textwidth}{|l|l|X|}
    \hline
    \multicolumn{3}{c}{\textbf{Capability Metadata}} \\
    \hline
    \textbf{Name} & \textbf{Value} & \textbf{Comment and/or Description} \\
    \hline
    \textbf{Reference Id} & \texttt{BEC-EDCE7C88} & \\
    \hline
    \textbf{Status} & Candidate ancestry & \\
    \hline
    \textbf{Ancestry} &  & \\
    \hline
    \textbf{Source} &  & \\
    \hline
    \textbf{Enabled By} &  & \\
    \hline
    \textbf{Supports} &  & \\
    \hline
    \textbf{Derived From} &  & \\
    \hline
    \textbf{Reviewed Status} &  & \\
    \hline
  \end{tabularx}
\end{table}

\begin{table}[htbp]
  \centering
  \footnotesize
  \begin{tabularx}{1\textwidth}{|l|l|X|}
    \hline
    \multicolumn{3}{c}{\textbf{Contained Features}} \\
    \hline
    \textbf{Feature Name} & \textbf{Feature Id} & \textbf{Comment} \\
    \hline
    Subject Context Binding & \texttt{cap:bec-19b283f0} (at \ref{cap:bec-19b283f0}) & \\
    \hline
    Demographic Context Governance & \texttt{cap:bec-6b8c0fdf} (at \ref{cap:bec-6b8c0fdf}) & \\
    \hline
    Context Establishment & \texttt{cap:bec-98ed0fc6} (at \ref{cap:bec-98ed0fc6}) & \\
    \hline
    State Monitoring & \texttt{cap:bec-b3c6efa0} (at \ref{cap:bec-b3c6efa0}) & \\
    \hline
    Exception Management & \texttt{cap:bec-7c07a1c7} (at \ref{cap:bec-7c07a1c7}) & \\
    \hline
    Information Retrieval & \texttt{cap:bec-54eae8ca} (at \ref{cap:bec-54eae8ca}) & \\
    \hline
  \end{tabularx}
\end{table}

Intent. Healthcare Subject is a candidate enabling capability within Client Administration. It frames the reusable system-supported behaviours needed to support this healthcare context without prescribing a product, interface, technology or deployment boundary.

Boundary. This capability supports the business activity and does not inherit the authority or accountability of the Business Capability it enables. Its child Features are review candidates until ancestry and semantics are adjudicated.

\subsubsection{Feature - Subject Context Binding}\label{cap:bec-19b283f0}
\vspace{-1em}
\begin{table}[htbp]
  \centering
  \footnotesize
  \begin{tabularx}{1\textwidth}{|l|l|X|}
    \hline
    \textbf{Metadata Name} & \textbf{Metadata Value} & \textbf{Description} \\
    \hline
    \textbf{Reference Id} & \texttt{BEC-19B283F0} & \\
    \hline
    \textbf{Status} & Existing Feature; candidate CT3 placement & \\
    \hline
    \textbf{Ancestry} &  & \\
    \hline
    \textbf{Source} & FEAT-EM-05 & \\
    \hline
    \textbf{Enabled By} &  & \\
    \hline
    \textbf{Supports} &  & \\
    \hline
    \textbf{Derived From} &  & \\
    \hline
    \textbf{Reviewed Status} & Confirm existing Feature meaning and proposed CT3 ownership. & \\
    \hline
  \end{tabularx}
\end{table}

Intent. Subject Context Binding provides a distinct enabling behaviour in the context of Healthcare Subject. This is an existing named Feature; its placement under the candidate CT3 remains subject to review.

\noindent\textbf{Feature boundary.} This behaviour is scoped to Healthcare Subject; an identically named Feature elsewhere may carry different business meaning.\par
\featuretable{
Function & \texttt{FN-0D216244} & Establish Healthcare Subject Context & Catalogue \\
Service & \texttt{SV-A558105A} & Healthcare Subject Context Resolution & Linked to preceding Function in source \\
}
\subsubsection{Feature - Demographic Context Governance}\label{cap:bec-6b8c0fdf}
\vspace{-1em}
\begin{table}[htbp]
  \centering
  \footnotesize
  \begin{tabularx}{1\textwidth}{|l|l|X|}
    \hline
    \textbf{Metadata Name} & \textbf{Metadata Value} & \textbf{Description} \\
    \hline
    \textbf{Reference Id} & \texttt{BEC-6B8C0FDF} & \\
    \hline
    \textbf{Status} & Existing Feature; candidate CT3 placement & \\
    \hline
    \textbf{Ancestry} &  & \\
    \hline
    \textbf{Source} & FEAT-EM-06 & \\
    \hline
    \textbf{Enabled By} &  & \\
    \hline
    \textbf{Supports} &  & \\
    \hline
    \textbf{Derived From} &  & \\
    \hline
    \textbf{Reviewed Status} & Confirm existing Feature meaning and proposed CT3 ownership. & \\
    \hline
  \end{tabularx}
\end{table}

Intent. Demographic Context Governance provides a distinct enabling behaviour in the context of Healthcare Subject. This is an existing named Feature; its placement under the candidate CT3 remains subject to review.

\noindent\textbf{Feature boundary.} This behaviour is scoped to Healthcare Subject; an identically named Feature elsewhere may carry different business meaning.\par
\featuretable{
Function & \texttt{FN-1D465543} & Govern Subject Demographic Context & Catalogue \\
Service & \texttt{SV-AA64D660} & Demographic Change Distribution & Linked to preceding Function in source \\
}
\subsubsection{Feature - Context Establishment}\label{cap:bec-98ed0fc6}
\vspace{-1em}
\begin{table}[htbp]
  \centering
  \footnotesize
  \begin{tabularx}{1\textwidth}{|l|l|X|}
    \hline
    \textbf{Metadata Name} & \textbf{Metadata Value} & \textbf{Description} \\
    \hline
    \textbf{Reference Id} & \texttt{BEC-98ED0FC6} & \\
    \hline
    \textbf{Status} & Generated candidate & \\
    \hline
    \textbf{Ancestry} &  & \\
    \hline
    \textbf{Source} & Semantic-pattern candidate & \\
    \hline
    \textbf{Enabled By} &  & \\
    \hline
    \textbf{Supports} &  & \\
    \hline
    \textbf{Derived From} &  & \\
    \hline
    \textbf{Reviewed Status} & Generated candidate: keep / modify / relocate / reject; verify whether independently testable. & \\
    \hline
  \end{tabularx}
\end{table}

Intent. Context Establishment is proposed as a discrete, testable enabling behaviour for Healthcare Subject. It is included for architectural review, not as an established implementation obligation.

\noindent\textbf{Feature boundary.} This behaviour is scoped to Healthcare Subject; an identically named Feature elsewhere may carry different business meaning.\par
\featuretable{
--- & --- & \emph{Not yet captured} & \emph{Complete during review} \\
}
\subsubsection{Feature - State Monitoring}\label{cap:bec-b3c6efa0}
\vspace{-1em}
\begin{table}[htbp]
  \centering
  \footnotesize
  \begin{tabularx}{1\textwidth}{|l|l|X|}
    \hline
    \textbf{Metadata Name} & \textbf{Metadata Value} & \textbf{Description} \\
    \hline
    \textbf{Reference Id} & \texttt{BEC-B3C6EFA0} & \\
    \hline
    \textbf{Status} & Generated candidate & \\
    \hline
    \textbf{Ancestry} &  & \\
    \hline
    \textbf{Source} & Semantic-pattern candidate & \\
    \hline
    \textbf{Enabled By} &  & \\
    \hline
    \textbf{Supports} &  & \\
    \hline
    \textbf{Derived From} &  & \\
    \hline
    \textbf{Reviewed Status} & Generated candidate: keep / modify / relocate / reject; verify whether independently testable. & \\
    \hline
  \end{tabularx}
\end{table}

Intent. State Monitoring is proposed as a discrete, testable enabling behaviour for Healthcare Subject. It is included for architectural review, not as an established implementation obligation.

\noindent\textbf{Feature boundary.} This behaviour is scoped to Healthcare Subject; an identically named Feature elsewhere may carry different business meaning.\par
\featuretable{
--- & --- & \emph{Not yet captured} & \emph{Complete during review} \\
}
\subsubsection{Feature - Exception Management}\label{cap:bec-7c07a1c7}
\vspace{-1em}
\begin{table}[htbp]
  \centering
  \footnotesize
  \begin{tabularx}{1\textwidth}{|l|l|X|}
    \hline
    \textbf{Metadata Name} & \textbf{Metadata Value} & \textbf{Description} \\
    \hline
    \textbf{Reference Id} & \texttt{BEC-7C07A1C7} & \\
    \hline
    \textbf{Status} & Generated candidate & \\
    \hline
    \textbf{Ancestry} &  & \\
    \hline
    \textbf{Source} & Semantic-pattern candidate & \\
    \hline
    \textbf{Enabled By} &  & \\
    \hline
    \textbf{Supports} &  & \\
    \hline
    \textbf{Derived From} &  & \\
    \hline
    \textbf{Reviewed Status} & Generated candidate: keep / modify / relocate / reject; verify whether independently testable. & \\
    \hline
  \end{tabularx}
\end{table}

Intent. Exception Management is proposed as a discrete, testable enabling behaviour for Healthcare Subject. It is included for architectural review, not as an established implementation obligation.

\noindent\textbf{Feature boundary.} This behaviour is scoped to Healthcare Subject; an identically named Feature elsewhere may carry different business meaning.\par
\featuretable{
--- & --- & \emph{Not yet captured} & \emph{Complete during review} \\
}
\subsubsection{Feature - Information Retrieval}\label{cap:bec-54eae8ca}
\vspace{-1em}
\begin{table}[htbp]
  \centering
  \footnotesize
  \begin{tabularx}{1\textwidth}{|l|l|X|}
    \hline
    \textbf{Metadata Name} & \textbf{Metadata Value} & \textbf{Description} \\
    \hline
    \textbf{Reference Id} & \texttt{BEC-54EAE8CA} & \\
    \hline
    \textbf{Status} & Generated candidate & \\
    \hline
    \textbf{Ancestry} &  & \\
    \hline
    \textbf{Source} & Semantic-pattern candidate & \\
    \hline
    \textbf{Enabled By} &  & \\
    \hline
    \textbf{Supports} &  & \\
    \hline
    \textbf{Derived From} &  & \\
    \hline
    \textbf{Reviewed Status} & Generated candidate: keep / modify / relocate / reject; verify whether independently testable. & \\
    \hline
  \end{tabularx}
\end{table}

Intent. Information Retrieval is proposed as a discrete, testable enabling behaviour for Healthcare Subject. It is included for architectural review, not as an established implementation obligation.

\noindent\textbf{Feature boundary.} This behaviour is scoped to Healthcare Subject; an identically named Feature elsewhere may carry different business meaning.\par
\featuretable{
--- & --- & \emph{Not yet captured} & \emph{Complete during review} \\
}

\subsection{CT3 - Client Relationship}\label{cap:bec-2494865f}
\vspace{-0.5em}
\begin{table}[htbp]
  \centering
  \footnotesize
  \begin{tabularx}{1\textwidth}{|l|l|X|}
    \hline
    \multicolumn{3}{c}{\textbf{Capability Metadata}} \\
    \hline
    \textbf{Name} & \textbf{Value} & \textbf{Comment and/or Description} \\
    \hline
    \textbf{Reference Id} & \texttt{BEC-2494865F} & \\
    \hline
    \textbf{Status} & Candidate ancestry & \\
    \hline
    \textbf{Ancestry} &  & \\
    \hline
    \textbf{Source} &  & \\
    \hline
    \textbf{Enabled By} &  & \\
    \hline
    \textbf{Supports} &  & \\
    \hline
    \textbf{Derived From} &  & \\
    \hline
    \textbf{Reviewed Status} &  & \\
    \hline
  \end{tabularx}
\end{table}

\begin{table}[htbp]
  \centering
  \footnotesize
  \begin{tabularx}{1\textwidth}{|l|l|X|}
    \hline
    \multicolumn{3}{c}{\textbf{Contained Features}} \\
    \hline
    \textbf{Feature Name} & \textbf{Feature Id} & \textbf{Comment} \\
    \hline
    Next-of-Kin \& Guardian Association & \texttt{cap:bec-bd118344} (at \ref{cap:bec-bd118344}) & \\
    \hline
    Relationship Validity Tracking & \texttt{cap:bec-6811fcdb} (at \ref{cap:bec-6811fcdb}) & \\
    \hline
    Context Establishment & \texttt{cap:bec-cdb17807} (at \ref{cap:bec-cdb17807}) & \\
    \hline
    State Monitoring & \texttt{cap:bec-d61b3c3b} (at \ref{cap:bec-d61b3c3b}) & \\
    \hline
    Exception Management & \texttt{cap:bec-4aacc1b3} (at \ref{cap:bec-4aacc1b3}) & \\
    \hline
    Information Retrieval & \texttt{cap:bec-c9abfa6f} (at \ref{cap:bec-c9abfa6f}) & \\
    \hline
  \end{tabularx}
\end{table}

Intent. Client Relationship is a candidate enabling capability within Client Administration. It frames the reusable system-supported behaviours needed to support this healthcare context without prescribing a product, interface, technology or deployment boundary.

Boundary. This capability supports the business activity and does not inherit the authority or accountability of the Business Capability it enables. Its child Features are review candidates until ancestry and semantics are adjudicated.

FHIR R5 conceptual analogues considered (not implementation mappings): Patient, Person, RelatedPerson.

\subsubsection{Feature - Next-of-Kin \& Guardian Association}\label{cap:bec-bd118344}
\vspace{-1em}
\begin{table}[htbp]
  \centering
  \footnotesize
  \begin{tabularx}{1\textwidth}{|l|l|X|}
    \hline
    \textbf{Metadata Name} & \textbf{Metadata Value} & \textbf{Description} \\
    \hline
    \textbf{Reference Id} & \texttt{BEC-BD118344} & \\
    \hline
    \textbf{Status} & Existing Feature; candidate CT3 placement & \\
    \hline
    \textbf{Ancestry} &  & \\
    \hline
    \textbf{Source} & FEAT-EM-07 & \\
    \hline
    \textbf{Enabled By} &  & \\
    \hline
    \textbf{Supports} &  & \\
    \hline
    \textbf{Derived From} &  & \\
    \hline
    \textbf{Reviewed Status} & Confirm existing Feature meaning and proposed CT3 ownership. & \\
    \hline
  \end{tabularx}
\end{table}

Intent. Next-of-Kin \& Guardian Association provides a distinct enabling behaviour in the context of Client Relationship. This is an existing named Feature; its placement under the candidate CT3 remains subject to review.

\noindent\textbf{Feature boundary.} This behaviour is scoped to Client Relationship; an identically named Feature elsewhere may carry different business meaning.\par
\featuretable{
Function & \texttt{FN-7FCDD3C7} & Maintain Client Relationships & Catalogue \\
Service & \texttt{SV-55132022} & Client Relationship Query & Linked to preceding Function in source \\
}
\subsubsection{Feature - Relationship Validity Tracking}\label{cap:bec-6811fcdb}
\vspace{-1em}
\begin{table}[htbp]
  \centering
  \footnotesize
  \begin{tabularx}{1\textwidth}{|l|l|X|}
    \hline
    \textbf{Metadata Name} & \textbf{Metadata Value} & \textbf{Description} \\
    \hline
    \textbf{Reference Id} & \texttt{BEC-6811FCDB} & \\
    \hline
    \textbf{Status} & Existing Feature; candidate CT3 placement & \\
    \hline
    \textbf{Ancestry} &  & \\
    \hline
    \textbf{Source} & FEAT-EM-08 & \\
    \hline
    \textbf{Enabled By} &  & \\
    \hline
    \textbf{Supports} &  & \\
    \hline
    \textbf{Derived From} &  & \\
    \hline
    \textbf{Reviewed Status} & Confirm existing Feature meaning and proposed CT3 ownership. & \\
    \hline
  \end{tabularx}
\end{table}

Intent. Relationship Validity Tracking provides a distinct enabling behaviour in the context of Client Relationship. This is an existing named Feature; its placement under the candidate CT3 remains subject to review.

\noindent\textbf{Feature boundary.} This behaviour is scoped to Client Relationship; an identically named Feature elsewhere may carry different business meaning.\par
\featuretable{
Function & \texttt{FN-B70CA519} & Govern Client Relationship Validity & Catalogue \\
Service & \texttt{SV-662E0211} & Representative Authority Verification & Linked to preceding Function in source \\
}
\subsubsection{Feature - Context Establishment}\label{cap:bec-cdb17807}
\vspace{-1em}
\begin{table}[htbp]
  \centering
  \footnotesize
  \begin{tabularx}{1\textwidth}{|l|l|X|}
    \hline
    \textbf{Metadata Name} & \textbf{Metadata Value} & \textbf{Description} \\
    \hline
    \textbf{Reference Id} & \texttt{BEC-CDB17807} & \\
    \hline
    \textbf{Status} & Generated candidate & \\
    \hline
    \textbf{Ancestry} &  & \\
    \hline
    \textbf{Source} & Semantic-pattern candidate & \\
    \hline
    \textbf{Enabled By} &  & \\
    \hline
    \textbf{Supports} &  & \\
    \hline
    \textbf{Derived From} &  & \\
    \hline
    \textbf{Reviewed Status} & Generated candidate: keep / modify / relocate / reject; verify whether independently testable. & \\
    \hline
  \end{tabularx}
\end{table}

Intent. Context Establishment is proposed as a discrete, testable enabling behaviour for Client Relationship. It is included for architectural review, not as an established implementation obligation.

\noindent\textbf{Feature boundary.} This behaviour is scoped to Client Relationship; an identically named Feature elsewhere may carry different business meaning.\par
\featuretable{
--- & --- & \emph{Not yet captured} & \emph{Complete during review} \\
}
\subsubsection{Feature - State Monitoring}\label{cap:bec-d61b3c3b}
\vspace{-1em}
\begin{table}[htbp]
  \centering
  \footnotesize
  \begin{tabularx}{1\textwidth}{|l|l|X|}
    \hline
    \textbf{Metadata Name} & \textbf{Metadata Value} & \textbf{Description} \\
    \hline
    \textbf{Reference Id} & \texttt{BEC-D61B3C3B} & \\
    \hline
    \textbf{Status} & Generated candidate & \\
    \hline
    \textbf{Ancestry} &  & \\
    \hline
    \textbf{Source} & Semantic-pattern candidate & \\
    \hline
    \textbf{Enabled By} &  & \\
    \hline
    \textbf{Supports} &  & \\
    \hline
    \textbf{Derived From} &  & \\
    \hline
    \textbf{Reviewed Status} & Generated candidate: keep / modify / relocate / reject; verify whether independently testable. & \\
    \hline
  \end{tabularx}
\end{table}

Intent. State Monitoring is proposed as a discrete, testable enabling behaviour for Client Relationship. It is included for architectural review, not as an established implementation obligation.

\noindent\textbf{Feature boundary.} This behaviour is scoped to Client Relationship; an identically named Feature elsewhere may carry different business meaning.\par
\featuretable{
--- & --- & \emph{Not yet captured} & \emph{Complete during review} \\
}
\subsubsection{Feature - Exception Management}\label{cap:bec-4aacc1b3}
\vspace{-1em}
\begin{table}[htbp]
  \centering
  \footnotesize
  \begin{tabularx}{1\textwidth}{|l|l|X|}
    \hline
    \textbf{Metadata Name} & \textbf{Metadata Value} & \textbf{Description} \\
    \hline
    \textbf{Reference Id} & \texttt{BEC-4AACC1B3} & \\
    \hline
    \textbf{Status} & Generated candidate & \\
    \hline
    \textbf{Ancestry} &  & \\
    \hline
    \textbf{Source} & Semantic-pattern candidate & \\
    \hline
    \textbf{Enabled By} &  & \\
    \hline
    \textbf{Supports} &  & \\
    \hline
    \textbf{Derived From} &  & \\
    \hline
    \textbf{Reviewed Status} & Generated candidate: keep / modify / relocate / reject; verify whether independently testable. & \\
    \hline
  \end{tabularx}
\end{table}

Intent. Exception Management is proposed as a discrete, testable enabling behaviour for Client Relationship. It is included for architectural review, not as an established implementation obligation.

\noindent\textbf{Feature boundary.} This behaviour is scoped to Client Relationship; an identically named Feature elsewhere may carry different business meaning.\par
\featuretable{
--- & --- & \emph{Not yet captured} & \emph{Complete during review} \\
}
\subsubsection{Feature - Information Retrieval}\label{cap:bec-c9abfa6f}
\vspace{-1em}
\begin{table}[htbp]
  \centering
  \footnotesize
  \begin{tabularx}{1\textwidth}{|l|l|X|}
    \hline
    \textbf{Metadata Name} & \textbf{Metadata Value} & \textbf{Description} \\
    \hline
    \textbf{Reference Id} & \texttt{BEC-C9ABFA6F} & \\
    \hline
    \textbf{Status} & Generated candidate & \\
    \hline
    \textbf{Ancestry} &  & \\
    \hline
    \textbf{Source} & Semantic-pattern candidate & \\
    \hline
    \textbf{Enabled By} &  & \\
    \hline
    \textbf{Supports} &  & \\
    \hline
    \textbf{Derived From} &  & \\
    \hline
    \textbf{Reviewed Status} & Generated candidate: keep / modify / relocate / reject; verify whether independently testable. & \\
    \hline
  \end{tabularx}
\end{table}

Intent. Information Retrieval is proposed as a discrete, testable enabling behaviour for Client Relationship. It is included for architectural review, not as an established implementation obligation.

\noindent\textbf{Feature boundary.} This behaviour is scoped to Client Relationship; an identically named Feature elsewhere may carry different business meaning.\par
\featuretable{
--- & --- & \emph{Not yet captured} & \emph{Complete during review} \\
}

\subsection{CT3 - Client Privacy}\label{cap:bec-adc0cf33}
\vspace{-0.5em}
\begin{table}[htbp]
  \centering
  \footnotesize
  \begin{tabularx}{1\textwidth}{|l|l|X|}
    \hline
    \multicolumn{3}{c}{\textbf{Capability Metadata}} \\
    \hline
    \textbf{Name} & \textbf{Value} & \textbf{Comment and/or Description} \\
    \hline
    \textbf{Reference Id} & \texttt{BEC-ADC0CF33} & \\
    \hline
    \textbf{Status} & Candidate ancestry & \\
    \hline
    \textbf{Ancestry} &  & \\
    \hline
    \textbf{Source} &  & \\
    \hline
    \textbf{Enabled By} &  & \\
    \hline
    \textbf{Supports} &  & \\
    \hline
    \textbf{Derived From} &  & \\
    \hline
    \textbf{Reviewed Status} &  & \\
    \hline
  \end{tabularx}
\end{table}

\begin{table}[htbp]
  \centering
  \footnotesize
  \begin{tabularx}{1\textwidth}{|l|l|X|}
    \hline
    \multicolumn{3}{c}{\textbf{Contained Features}} \\
    \hline
    \textbf{Feature Name} & \textbf{Feature Id} & \textbf{Comment} \\
    \hline
    Consent Directive Evaluation & \texttt{cap:bec-f41b7e86} (at \ref{cap:bec-f41b7e86}) & \\
    \hline
    Confidentiality Flag Enforcement & \texttt{cap:bec-4d74b3d2} (at \ref{cap:bec-4d74b3d2}) & \\
    \hline
    Criteria Management & \texttt{cap:bec-aa3166d7} (at \ref{cap:bec-aa3166d7}) & \\
    \hline
    Evidence Association & \texttt{cap:bec-5a06ba3c} (at \ref{cap:bec-5a06ba3c}) & \\
    \hline
    Policy Evaluation & \texttt{cap:bec-ef1af8f9} (at \ref{cap:bec-ef1af8f9}) & \\
    \hline
    Decision Traceability & \texttt{cap:bec-7d40e369} (at \ref{cap:bec-7d40e369}) & \\
    \hline
    Exception Reporting & \texttt{cap:bec-faf82b77} (at \ref{cap:bec-faf82b77}) & \\
    \hline
  \end{tabularx}
\end{table}

Intent. Client Privacy is a candidate enabling capability within Client Administration. It frames the reusable system-supported behaviours needed to support this healthcare context without prescribing a product, interface, technology or deployment boundary.

Boundary. This capability supports the business activity and does not inherit the authority or accountability of the Business Capability it enables. Its child Features are review candidates until ancestry and semantics are adjudicated.

FHIR R5 conceptual analogues considered (not implementation mappings): Patient, Person, RelatedPerson.

\subsubsection{Feature - Consent Directive Evaluation}\label{cap:bec-f41b7e86}
\vspace{-1em}
\begin{table}[htbp]
  \centering
  \footnotesize
  \begin{tabularx}{1\textwidth}{|l|l|X|}
    \hline
    \textbf{Metadata Name} & \textbf{Metadata Value} & \textbf{Description} \\
    \hline
    \textbf{Reference Id} & \texttt{BEC-F41B7E86} & \\
    \hline
    \textbf{Status} & Existing Feature; candidate CT3 placement & \\
    \hline
    \textbf{Ancestry} &  & \\
    \hline
    \textbf{Source} & FEAT-EM-09 & \\
    \hline
    \textbf{Enabled By} &  & \\
    \hline
    \textbf{Supports} &  & \\
    \hline
    \textbf{Derived From} &  & \\
    \hline
    \textbf{Reviewed Status} & Confirm existing Feature meaning and proposed CT3 ownership. & \\
    \hline
  \end{tabularx}
\end{table}

Intent. Consent Directive Evaluation provides a distinct enabling behaviour in the context of Client Privacy. This is an existing named Feature; its placement under the candidate CT3 remains subject to review.

\noindent\textbf{Feature boundary.} This behaviour is scoped to Client Privacy; an identically named Feature elsewhere may carry different business meaning.\par
\featuretable{
Function & \texttt{FN-25C6914F} & Evaluate Client Consent & Catalogue \\
Service & \texttt{SV-8545EF98} & Consent Decision Service & Linked to preceding Function in source \\
}
\subsubsection{Feature - Confidentiality Flag Enforcement}\label{cap:bec-4d74b3d2}
\vspace{-1em}
\begin{table}[htbp]
  \centering
  \footnotesize
  \begin{tabularx}{1\textwidth}{|l|l|X|}
    \hline
    \textbf{Metadata Name} & \textbf{Metadata Value} & \textbf{Description} \\
    \hline
    \textbf{Reference Id} & \texttt{BEC-4D74B3D2} & \\
    \hline
    \textbf{Status} & Existing Feature; candidate CT3 placement & \\
    \hline
    \textbf{Ancestry} &  & \\
    \hline
    \textbf{Source} & FEAT-EM-10 & \\
    \hline
    \textbf{Enabled By} &  & \\
    \hline
    \textbf{Supports} &  & \\
    \hline
    \textbf{Derived From} &  & \\
    \hline
    \textbf{Reviewed Status} & Confirm existing Feature meaning and proposed CT3 ownership. & \\
    \hline
  \end{tabularx}
\end{table}

Intent. Confidentiality Flag Enforcement provides a distinct enabling behaviour in the context of Client Privacy. This is an existing named Feature; its placement under the candidate CT3 remains subject to review.

\noindent\textbf{Feature boundary.} This behaviour is scoped to Client Privacy; an identically named Feature elsewhere may carry different business meaning.\par
\featuretable{
Function & \texttt{FN-9F5DA17D} & Evaluate Client Confidentiality Restrictions & Catalogue \\
Service & \texttt{SV-BC5007C0} & Confidentiality Restriction Evaluation & Linked to preceding Function in source \\
}
\subsubsection{Feature - Criteria Management}\label{cap:bec-aa3166d7}
\vspace{-1em}
\begin{table}[htbp]
  \centering
  \footnotesize
  \begin{tabularx}{1\textwidth}{|l|l|X|}
    \hline
    \textbf{Metadata Name} & \textbf{Metadata Value} & \textbf{Description} \\
    \hline
    \textbf{Reference Id} & \texttt{BEC-AA3166D7} & \\
    \hline
    \textbf{Status} & Generated candidate & \\
    \hline
    \textbf{Ancestry} &  & \\
    \hline
    \textbf{Source} & Semantic-pattern candidate & \\
    \hline
    \textbf{Enabled By} &  & \\
    \hline
    \textbf{Supports} &  & \\
    \hline
    \textbf{Derived From} &  & \\
    \hline
    \textbf{Reviewed Status} & Generated candidate: keep / modify / relocate / reject; verify whether independently testable. & \\
    \hline
  \end{tabularx}
\end{table}

Intent. Criteria Management is proposed as a discrete, testable enabling behaviour for Client Privacy. It is included for architectural review, not as an established implementation obligation.

\noindent\textbf{Feature boundary.} This behaviour is scoped to Client Privacy; an identically named Feature elsewhere may carry different business meaning.\par
\featuretable{
--- & --- & \emph{Not yet captured} & \emph{Complete during review} \\
}
\subsubsection{Feature - Evidence Association}\label{cap:bec-5a06ba3c}
\vspace{-1em}
\begin{table}[htbp]
  \centering
  \footnotesize
  \begin{tabularx}{1\textwidth}{|l|l|X|}
    \hline
    \textbf{Metadata Name} & \textbf{Metadata Value} & \textbf{Description} \\
    \hline
    \textbf{Reference Id} & \texttt{BEC-5A06BA3C} & \\
    \hline
    \textbf{Status} & Generated candidate & \\
    \hline
    \textbf{Ancestry} &  & \\
    \hline
    \textbf{Source} & Semantic-pattern candidate & \\
    \hline
    \textbf{Enabled By} &  & \\
    \hline
    \textbf{Supports} &  & \\
    \hline
    \textbf{Derived From} &  & \\
    \hline
    \textbf{Reviewed Status} & Generated candidate: keep / modify / relocate / reject; verify whether independently testable. & \\
    \hline
  \end{tabularx}
\end{table}

Intent. Evidence Association is proposed as a discrete, testable enabling behaviour for Client Privacy. It is included for architectural review, not as an established implementation obligation.

\noindent\textbf{Feature boundary.} This behaviour is scoped to Client Privacy; an identically named Feature elsewhere may carry different business meaning.\par
\featuretable{
--- & --- & \emph{Not yet captured} & \emph{Complete during review} \\
}
\subsubsection{Feature - Policy Evaluation}\label{cap:bec-ef1af8f9}
\vspace{-1em}
\begin{table}[htbp]
  \centering
  \footnotesize
  \begin{tabularx}{1\textwidth}{|l|l|X|}
    \hline
    \textbf{Metadata Name} & \textbf{Metadata Value} & \textbf{Description} \\
    \hline
    \textbf{Reference Id} & \texttt{BEC-EF1AF8F9} & \\
    \hline
    \textbf{Status} & Generated candidate & \\
    \hline
    \textbf{Ancestry} &  & \\
    \hline
    \textbf{Source} & Semantic-pattern candidate & \\
    \hline
    \textbf{Enabled By} &  & \\
    \hline
    \textbf{Supports} &  & \\
    \hline
    \textbf{Derived From} &  & \\
    \hline
    \textbf{Reviewed Status} & Generated candidate: keep / modify / relocate / reject; verify whether independently testable. & \\
    \hline
  \end{tabularx}
\end{table}

Intent. Policy Evaluation is proposed as a discrete, testable enabling behaviour for Client Privacy. It is included for architectural review, not as an established implementation obligation.

\noindent\textbf{Feature boundary.} This behaviour is scoped to Client Privacy; an identically named Feature elsewhere may carry different business meaning.\par
\featuretable{
--- & --- & \emph{Not yet captured} & \emph{Complete during review} \\
}
\subsubsection{Feature - Decision Traceability}\label{cap:bec-7d40e369}
\vspace{-1em}
\begin{table}[htbp]
  \centering
  \footnotesize
  \begin{tabularx}{1\textwidth}{|l|l|X|}
    \hline
    \textbf{Metadata Name} & \textbf{Metadata Value} & \textbf{Description} \\
    \hline
    \textbf{Reference Id} & \texttt{BEC-7D40E369} & \\
    \hline
    \textbf{Status} & Generated candidate & \\
    \hline
    \textbf{Ancestry} &  & \\
    \hline
    \textbf{Source} & Semantic-pattern candidate & \\
    \hline
    \textbf{Enabled By} &  & \\
    \hline
    \textbf{Supports} &  & \\
    \hline
    \textbf{Derived From} &  & \\
    \hline
    \textbf{Reviewed Status} & Generated candidate: keep / modify / relocate / reject; verify whether independently testable. & \\
    \hline
  \end{tabularx}
\end{table}

Intent. Decision Traceability is proposed as a discrete, testable enabling behaviour for Client Privacy. It is included for architectural review, not as an established implementation obligation.

\noindent\textbf{Feature boundary.} This behaviour is scoped to Client Privacy; an identically named Feature elsewhere may carry different business meaning.\par
\featuretable{
--- & --- & \emph{Not yet captured} & \emph{Complete during review} \\
}
\subsubsection{Feature - Exception Reporting}\label{cap:bec-faf82b77}
\vspace{-1em}
\begin{table}[htbp]
  \centering
  \footnotesize
  \begin{tabularx}{1\textwidth}{|l|l|X|}
    \hline
    \textbf{Metadata Name} & \textbf{Metadata Value} & \textbf{Description} \\
    \hline
    \textbf{Reference Id} & \texttt{BEC-FAF82B77} & \\
    \hline
    \textbf{Status} & Generated candidate & \\
    \hline
    \textbf{Ancestry} &  & \\
    \hline
    \textbf{Source} & Semantic-pattern candidate & \\
    \hline
    \textbf{Enabled By} &  & \\
    \hline
    \textbf{Supports} &  & \\
    \hline
    \textbf{Derived From} &  & \\
    \hline
    \textbf{Reviewed Status} & Generated candidate: keep / modify / relocate / reject; verify whether independently testable. & \\
    \hline
  \end{tabularx}
\end{table}

Intent. Exception Reporting is proposed as a discrete, testable enabling behaviour for Client Privacy. It is included for architectural review, not as an established implementation obligation.

\noindent\textbf{Feature boundary.} This behaviour is scoped to Client Privacy; an identically named Feature elsewhere may carry different business meaning.\par
\featuretable{
--- & --- & \emph{Not yet captured} & \emph{Complete during review} \\
}

\section{CT2 - Provider Administration}\label{cap:bec-a4e759e2}
\vspace{-0.5em}
\begin{table}[htbp]
  \centering
  \footnotesize
  \begin{tabularx}{1\textwidth}{|l|l|X|}
    \hline
    \multicolumn{3}{c}{\textbf{Capability Metadata}} \\
    \hline
    \textbf{Name} & \textbf{Value} & \textbf{Comment and/or Description} \\
    \hline
    \textbf{Reference Id} & \texttt{BEC-A4E759E2} & \\
    \hline
    \textbf{Status} & Candidate ancestry & \\
    \hline
    \textbf{Ancestry} &  & \\
    \hline
    \textbf{Source} &  & \\
    \hline
    \textbf{Enabled By} &  & \\
    \hline
    \textbf{Supports} &  & \\
    \hline
    \textbf{Derived From} &  & \\
    \hline
    \textbf{Reviewed Status} &  & \\
    \hline
  \end{tabularx}
\end{table}

\begin{table}[htbp]
  \centering
  \footnotesize
  \begin{tabularx}{1\textwidth}{|l|l|X|}
    \hline
    \multicolumn{3}{c}{\textbf{Contained Capabilities}} \\
    \hline
    \textbf{Capability Name} & \textbf{Capability Id} & \textbf{Comment} \\
    \hline
    Practitioner & \texttt{cap:bec-5bb55046} (at \ref{cap:bec-5bb55046}) & \\
    \hline
    Practitioner Role & \texttt{cap:bec-2dec073a} (at \ref{cap:bec-2dec073a}) & \\
    \hline
    Provider Relationship & \texttt{cap:bec-c410cc51} (at \ref{cap:bec-c410cc51}) & \\
    \hline
    Provider Engagement & \texttt{cap:bec-f58dba08} (at \ref{cap:bec-f58dba08}) & \\
    \hline
  \end{tabularx}
\end{table}

Intent. Provider Administration is a candidate enabling capability within Entity Information Enablement. It frames the reusable system-supported behaviours needed to support this healthcare context without prescribing a product, interface, technology or deployment boundary.

Boundary. This capability supports the business activity and does not inherit the authority or accountability of the Business Capability it enables. Its child Features are review candidates until ancestry and semantics are adjudicated.

\subsection{CT3 - Practitioner}\label{cap:bec-5bb55046}
\vspace{-0.5em}
\begin{table}[htbp]
  \centering
  \footnotesize
  \begin{tabularx}{1\textwidth}{|l|l|X|}
    \hline
    \multicolumn{3}{c}{\textbf{Capability Metadata}} \\
    \hline
    \textbf{Name} & \textbf{Value} & \textbf{Comment and/or Description} \\
    \hline
    \textbf{Reference Id} & \texttt{BEC-5BB55046} & \\
    \hline
    \textbf{Status} & Candidate ancestry & \\
    \hline
    \textbf{Ancestry} &  & \\
    \hline
    \textbf{Source} &  & \\
    \hline
    \textbf{Enabled By} &  & \\
    \hline
    \textbf{Supports} &  & \\
    \hline
    \textbf{Derived From} &  & \\
    \hline
    \textbf{Reviewed Status} &  & \\
    \hline
  \end{tabularx}
\end{table}

\begin{table}[htbp]
  \centering
  \footnotesize
  \begin{tabularx}{1\textwidth}{|l|l|X|}
    \hline
    \multicolumn{3}{c}{\textbf{Contained Features}} \\
    \hline
    \textbf{Feature Name} & \textbf{Feature Id} & \textbf{Comment} \\
    \hline
    National Practitioner Verification & \texttt{cap:bec-95b0a32b} (at \ref{cap:bec-95b0a32b}) & \\
    \hline
    Practitioner Profile Governance & \texttt{cap:bec-7672710f} (at \ref{cap:bec-7672710f}) & \\
    \hline
    Identity Management & \texttt{cap:bec-a4cae7d6} (at \ref{cap:bec-a4cae7d6}) & \\
    \hline
    Identity Correlation & \texttt{cap:bec-c7bbd105} (at \ref{cap:bec-c7bbd105}) & \\
    \hline
    Qualification Association & \texttt{cap:bec-f9fa105b} (at \ref{cap:bec-f9fa105b}) & \\
    \hline
    Professional Relationship Management & \texttt{cap:bec-028abd17} (at \ref{cap:bec-028abd17}) & \\
    \hline
    Practitioner Discovery & \texttt{cap:bec-e9c766cb} (at \ref{cap:bec-e9c766cb}) & \\
    \hline
  \end{tabularx}
\end{table}

Intent. Practitioner is a candidate enabling capability within Provider Administration. It frames the reusable system-supported behaviours needed to support this healthcare context without prescribing a product, interface, technology or deployment boundary.

Boundary. This capability supports the business activity and does not inherit the authority or accountability of the Business Capability it enables. Its child Features are review candidates until ancestry and semantics are adjudicated.

FHIR R5 conceptual analogues considered (not implementation mappings): Practitioner, PractitionerRole.

\subsubsection{Feature - National Practitioner Verification}\label{cap:bec-95b0a32b}
\vspace{-1em}
\begin{table}[htbp]
  \centering
  \footnotesize
  \begin{tabularx}{1\textwidth}{|l|l|X|}
    \hline
    \textbf{Metadata Name} & \textbf{Metadata Value} & \textbf{Description} \\
    \hline
    \textbf{Reference Id} & \texttt{BEC-95B0A32B} & \\
    \hline
    \textbf{Status} & Existing Feature; candidate CT3 placement & \\
    \hline
    \textbf{Ancestry} &  & \\
    \hline
    \textbf{Source} & FEAT-EM-11 & \\
    \hline
    \textbf{Enabled By} &  & \\
    \hline
    \textbf{Supports} &  & \\
    \hline
    \textbf{Derived From} &  & \\
    \hline
    \textbf{Reviewed Status} & Confirm existing Feature meaning and proposed CT3 ownership. & \\
    \hline
  \end{tabularx}
\end{table}

Intent. National Practitioner Verification provides a distinct enabling behaviour in the context of Practitioner. This is an existing named Feature; its placement under the candidate CT3 remains subject to review.

\noindent\textbf{Feature boundary.} This behaviour is scoped to Practitioner; an identically named Feature elsewhere may carry different business meaning.\par
\featuretable{
Function & \texttt{FN-1D3A0AFC} & Verify Practitioner Registration & Catalogue \\
Service & \texttt{SV-7B03EECC} & Practitioner Verification Service & Linked to preceding Function in source \\
Process & \texttt{PR-5BC4701D} & Practitioner Verification Process & Linked to preceding Function in source \\
}
\subsubsection{Feature - Practitioner Profile Governance}\label{cap:bec-7672710f}
\vspace{-1em}
\begin{table}[htbp]
  \centering
  \footnotesize
  \begin{tabularx}{1\textwidth}{|l|l|X|}
    \hline
    \textbf{Metadata Name} & \textbf{Metadata Value} & \textbf{Description} \\
    \hline
    \textbf{Reference Id} & \texttt{BEC-7672710F} & \\
    \hline
    \textbf{Status} & Existing Feature; candidate CT3 placement & \\
    \hline
    \textbf{Ancestry} &  & \\
    \hline
    \textbf{Source} & FEAT-EM-12 & \\
    \hline
    \textbf{Enabled By} &  & \\
    \hline
    \textbf{Supports} &  & \\
    \hline
    \textbf{Derived From} &  & \\
    \hline
    \textbf{Reviewed Status} & Confirm existing Feature meaning and proposed CT3 ownership. & \\
    \hline
  \end{tabularx}
\end{table}

Intent. Practitioner Profile Governance provides a distinct enabling behaviour in the context of Practitioner. This is an existing named Feature; its placement under the candidate CT3 remains subject to review.

\noindent\textbf{Feature boundary.} This behaviour is scoped to Practitioner; an identically named Feature elsewhere may carry different business meaning.\par
\featuretable{
Function & \texttt{FN-A0A16EBE} & Govern Practitioner Profile & Catalogue \\
Service & \texttt{SV-43927F89} & Practitioner Directory Query & Linked to preceding Function in source \\
}
\subsubsection{Feature - Identity Management}\label{cap:bec-a4cae7d6}
\vspace{-1em}
\begin{table}[htbp]
  \centering
  \footnotesize
  \begin{tabularx}{1\textwidth}{|l|l|X|}
    \hline
    \textbf{Metadata Name} & \textbf{Metadata Value} & \textbf{Description} \\
    \hline
    \textbf{Reference Id} & \texttt{BEC-A4CAE7D6} & \\
    \hline
    \textbf{Status} & Generated candidate & \\
    \hline
    \textbf{Ancestry} &  & \\
    \hline
    \textbf{Source} & FHIR-informed candidate & \\
    \hline
    \textbf{Enabled By} &  & \\
    \hline
    \textbf{Supports} &  & \\
    \hline
    \textbf{Derived From} &  & \\
    \hline
    \textbf{Reviewed Status} & Generated candidate: keep / modify / relocate / reject; verify whether independently testable. & \\
    \hline
  \end{tabularx}
\end{table}

Intent. Identity Management is proposed to support the underlying healthcare concept represented by Practitioner. The proposal concerns behaviour and semantic responsibility, not a FHIR field, operation, resource instance or prescribed data model.

Conceptual evidence: HL7 FHIR R5 Practitioner, PractitionerRole.

\noindent\textbf{Feature boundary.} This behaviour is scoped to Practitioner; an identically named Feature elsewhere may carry different business meaning.\par
\featuretable{
--- & --- & \emph{Not yet captured} & \emph{Complete during review} \\
}
\subsubsection{Feature - Identity Correlation}\label{cap:bec-c7bbd105}
\vspace{-1em}
\begin{table}[htbp]
  \centering
  \footnotesize
  \begin{tabularx}{1\textwidth}{|l|l|X|}
    \hline
    \textbf{Metadata Name} & \textbf{Metadata Value} & \textbf{Description} \\
    \hline
    \textbf{Reference Id} & \texttt{BEC-C7BBD105} & \\
    \hline
    \textbf{Status} & Generated candidate & \\
    \hline
    \textbf{Ancestry} &  & \\
    \hline
    \textbf{Source} & FHIR-informed candidate & \\
    \hline
    \textbf{Enabled By} &  & \\
    \hline
    \textbf{Supports} &  & \\
    \hline
    \textbf{Derived From} &  & \\
    \hline
    \textbf{Reviewed Status} & Generated candidate: keep / modify / relocate / reject; verify whether independently testable. & \\
    \hline
  \end{tabularx}
\end{table}

Intent. Identity Correlation is proposed to support the underlying healthcare concept represented by Practitioner. The proposal concerns behaviour and semantic responsibility, not a FHIR field, operation, resource instance or prescribed data model.

Conceptual evidence: HL7 FHIR R5 Practitioner, PractitionerRole.

\noindent\textbf{Feature boundary.} This behaviour is scoped to Practitioner; an identically named Feature elsewhere may carry different business meaning.\par
\featuretable{
--- & --- & \emph{Not yet captured} & \emph{Complete during review} \\
}
\subsubsection{Feature - Qualification Association}\label{cap:bec-f9fa105b}
\vspace{-1em}
\begin{table}[htbp]
  \centering
  \footnotesize
  \begin{tabularx}{1\textwidth}{|l|l|X|}
    \hline
    \textbf{Metadata Name} & \textbf{Metadata Value} & \textbf{Description} \\
    \hline
    \textbf{Reference Id} & \texttt{BEC-F9FA105B} & \\
    \hline
    \textbf{Status} & Generated candidate & \\
    \hline
    \textbf{Ancestry} &  & \\
    \hline
    \textbf{Source} & FHIR-informed candidate & \\
    \hline
    \textbf{Enabled By} &  & \\
    \hline
    \textbf{Supports} &  & \\
    \hline
    \textbf{Derived From} &  & \\
    \hline
    \textbf{Reviewed Status} & Generated candidate: keep / modify / relocate / reject; verify whether independently testable. & \\
    \hline
  \end{tabularx}
\end{table}

Intent. Qualification Association is proposed to support the underlying healthcare concept represented by Practitioner. The proposal concerns behaviour and semantic responsibility, not a FHIR field, operation, resource instance or prescribed data model.

Conceptual evidence: HL7 FHIR R5 Practitioner, PractitionerRole.

\noindent\textbf{Feature boundary.} This behaviour is scoped to Practitioner; an identically named Feature elsewhere may carry different business meaning.\par
\featuretable{
--- & --- & \emph{Not yet captured} & \emph{Complete during review} \\
}
\subsubsection{Feature - Professional Relationship Management}\label{cap:bec-028abd17}
\vspace{-1em}
\begin{table}[htbp]
  \centering
  \footnotesize
  \begin{tabularx}{1\textwidth}{|l|l|X|}
    \hline
    \textbf{Metadata Name} & \textbf{Metadata Value} & \textbf{Description} \\
    \hline
    \textbf{Reference Id} & \texttt{BEC-028ABD17} & \\
    \hline
    \textbf{Status} & Generated candidate & \\
    \hline
    \textbf{Ancestry} &  & \\
    \hline
    \textbf{Source} & FHIR-informed candidate & \\
    \hline
    \textbf{Enabled By} &  & \\
    \hline
    \textbf{Supports} &  & \\
    \hline
    \textbf{Derived From} &  & \\
    \hline
    \textbf{Reviewed Status} & Generated candidate: keep / modify / relocate / reject; verify whether independently testable. & \\
    \hline
  \end{tabularx}
\end{table}

Intent. Professional Relationship Management is proposed to support the underlying healthcare concept represented by Practitioner. The proposal concerns behaviour and semantic responsibility, not a FHIR field, operation, resource instance or prescribed data model.

Conceptual evidence: HL7 FHIR R5 Practitioner, PractitionerRole.

\noindent\textbf{Feature boundary.} This behaviour is scoped to Practitioner; an identically named Feature elsewhere may carry different business meaning.\par
\featuretable{
--- & --- & \emph{Not yet captured} & \emph{Complete during review} \\
}
\subsubsection{Feature - Practitioner Discovery}\label{cap:bec-e9c766cb}
\vspace{-1em}
\begin{table}[htbp]
  \centering
  \footnotesize
  \begin{tabularx}{1\textwidth}{|l|l|X|}
    \hline
    \textbf{Metadata Name} & \textbf{Metadata Value} & \textbf{Description} \\
    \hline
    \textbf{Reference Id} & \texttt{BEC-E9C766CB} & \\
    \hline
    \textbf{Status} & Generated candidate & \\
    \hline
    \textbf{Ancestry} &  & \\
    \hline
    \textbf{Source} & FHIR-informed candidate & \\
    \hline
    \textbf{Enabled By} &  & \\
    \hline
    \textbf{Supports} &  & \\
    \hline
    \textbf{Derived From} &  & \\
    \hline
    \textbf{Reviewed Status} & Generated candidate: keep / modify / relocate / reject; verify whether independently testable. & \\
    \hline
  \end{tabularx}
\end{table}

Intent. Practitioner Discovery is proposed to support the underlying healthcare concept represented by Practitioner. The proposal concerns behaviour and semantic responsibility, not a FHIR field, operation, resource instance or prescribed data model.

Conceptual evidence: HL7 FHIR R5 Practitioner, PractitionerRole.

\noindent\textbf{Feature boundary.} This behaviour is scoped to Practitioner; an identically named Feature elsewhere may carry different business meaning.\par
\featuretable{
--- & --- & \emph{Not yet captured} & \emph{Complete during review} \\
}

\subsection{CT3 - Practitioner Role}\label{cap:bec-2dec073a}
\vspace{-0.5em}
\begin{table}[htbp]
  \centering
  \footnotesize
  \begin{tabularx}{1\textwidth}{|l|l|X|}
    \hline
    \multicolumn{3}{c}{\textbf{Capability Metadata}} \\
    \hline
    \textbf{Name} & \textbf{Value} & \textbf{Comment and/or Description} \\
    \hline
    \textbf{Reference Id} & \texttt{BEC-2DEC073A} & \\
    \hline
    \textbf{Status} & Candidate ancestry & \\
    \hline
    \textbf{Ancestry} &  & \\
    \hline
    \textbf{Source} &  & \\
    \hline
    \textbf{Enabled By} &  & \\
    \hline
    \textbf{Supports} &  & \\
    \hline
    \textbf{Derived From} &  & \\
    \hline
    \textbf{Reviewed Status} &  & \\
    \hline
  \end{tabularx}
\end{table}

\begin{table}[htbp]
  \centering
  \footnotesize
  \begin{tabularx}{1\textwidth}{|l|l|X|}
    \hline
    \multicolumn{3}{c}{\textbf{Contained Features}} \\
    \hline
    \textbf{Feature Name} & \textbf{Feature Id} & \textbf{Comment} \\
    \hline
    Role \& Specialty Modeling & \texttt{cap:bec-116f8108} (at \ref{cap:bec-116f8108}) & \\
    \hline
    Electronic Endpoint Binding & \texttt{cap:bec-cbde3b8c} (at \ref{cap:bec-cbde3b8c}) & \\
    \hline
    Identity Management & \texttt{cap:bec-7f9bcb4f} (at \ref{cap:bec-7f9bcb4f}) & \\
    \hline
    Identity Correlation & \texttt{cap:bec-314be2d1} (at \ref{cap:bec-314be2d1}) & \\
    \hline
    Qualification Association & \texttt{cap:bec-a4596b27} (at \ref{cap:bec-a4596b27}) & \\
    \hline
    Professional Relationship Management & \texttt{cap:bec-549d792b} (at \ref{cap:bec-549d792b}) & \\
    \hline
    Practitioner Discovery & \texttt{cap:bec-e24857d9} (at \ref{cap:bec-e24857d9}) & \\
    \hline
  \end{tabularx}
\end{table}

Intent. Practitioner Role is a candidate enabling capability within Provider Administration. It frames the reusable system-supported behaviours needed to support this healthcare context without prescribing a product, interface, technology or deployment boundary.

Boundary. This capability supports the business activity and does not inherit the authority or accountability of the Business Capability it enables. Its child Features are review candidates until ancestry and semantics are adjudicated.

FHIR R5 conceptual analogues considered (not implementation mappings): Practitioner, PractitionerRole.

\subsubsection{Feature - Role \& Specialty Modeling}\label{cap:bec-116f8108}
\vspace{-1em}
\begin{table}[htbp]
  \centering
  \footnotesize
  \begin{tabularx}{1\textwidth}{|l|l|X|}
    \hline
    \textbf{Metadata Name} & \textbf{Metadata Value} & \textbf{Description} \\
    \hline
    \textbf{Reference Id} & \texttt{BEC-116F8108} & \\
    \hline
    \textbf{Status} & Existing Feature; candidate CT3 placement & \\
    \hline
    \textbf{Ancestry} &  & \\
    \hline
    \textbf{Source} & FEAT-EM-13 & \\
    \hline
    \textbf{Enabled By} &  & \\
    \hline
    \textbf{Supports} &  & \\
    \hline
    \textbf{Derived From} &  & \\
    \hline
    \textbf{Reviewed Status} & Confirm existing Feature meaning and proposed CT3 ownership. & \\
    \hline
  \end{tabularx}
\end{table}

Intent. Role \& Specialty Modeling provides a distinct enabling behaviour in the context of Practitioner Role. This is an existing named Feature; its placement under the candidate CT3 remains subject to review.

\noindent\textbf{Feature boundary.} This behaviour is scoped to Practitioner Role; an identically named Feature elsewhere may carry different business meaning.\par
\featuretable{
Function & \texttt{FN-5984D30A} & Maintain Practitioner Roles & Catalogue \\
Service & \texttt{SV-4DA91572} & Practitioner Role Resolution & Linked to preceding Function in source \\
}
\subsubsection{Feature - Electronic Endpoint Binding}\label{cap:bec-cbde3b8c}
\vspace{-1em}
\begin{table}[htbp]
  \centering
  \footnotesize
  \begin{tabularx}{1\textwidth}{|l|l|X|}
    \hline
    \textbf{Metadata Name} & \textbf{Metadata Value} & \textbf{Description} \\
    \hline
    \textbf{Reference Id} & \texttt{BEC-CBDE3B8C} & \\
    \hline
    \textbf{Status} & Existing Feature; candidate CT3 placement & \\
    \hline
    \textbf{Ancestry} &  & \\
    \hline
    \textbf{Source} & FEAT-EM-14 & \\
    \hline
    \textbf{Enabled By} &  & \\
    \hline
    \textbf{Supports} &  & \\
    \hline
    \textbf{Derived From} &  & \\
    \hline
    \textbf{Reviewed Status} & Confirm existing Feature meaning and proposed CT3 ownership. & \\
    \hline
  \end{tabularx}
\end{table}

Intent. Electronic Endpoint Binding provides a distinct enabling behaviour in the context of Practitioner Role. This is an existing named Feature; its placement under the candidate CT3 remains subject to review.

\noindent\textbf{Feature boundary.} This behaviour is scoped to Practitioner Role; an identically named Feature elsewhere may carry different business meaning.\par
\featuretable{
Function & \texttt{FN-C8DE7F63} & Bind Practitioner Electronic Endpoints & Catalogue \\
Service & \texttt{SV-D37B1512} & Practitioner Endpoint Lookup & Linked to preceding Function in source \\
}
\subsubsection{Feature - Identity Management}\label{cap:bec-7f9bcb4f}
\vspace{-1em}
\begin{table}[htbp]
  \centering
  \footnotesize
  \begin{tabularx}{1\textwidth}{|l|l|X|}
    \hline
    \textbf{Metadata Name} & \textbf{Metadata Value} & \textbf{Description} \\
    \hline
    \textbf{Reference Id} & \texttt{BEC-7F9BCB4F} & \\
    \hline
    \textbf{Status} & Generated candidate & \\
    \hline
    \textbf{Ancestry} &  & \\
    \hline
    \textbf{Source} & FHIR-informed candidate & \\
    \hline
    \textbf{Enabled By} &  & \\
    \hline
    \textbf{Supports} &  & \\
    \hline
    \textbf{Derived From} &  & \\
    \hline
    \textbf{Reviewed Status} & Generated candidate: keep / modify / relocate / reject; verify whether independently testable. & \\
    \hline
  \end{tabularx}
\end{table}

Intent. Identity Management is proposed to support the underlying healthcare concept represented by Practitioner Role. The proposal concerns behaviour and semantic responsibility, not a FHIR field, operation, resource instance or prescribed data model.

Conceptual evidence: HL7 FHIR R5 Practitioner, PractitionerRole.

\noindent\textbf{Feature boundary.} This behaviour is scoped to Practitioner Role; an identically named Feature elsewhere may carry different business meaning.\par
\featuretable{
--- & --- & \emph{Not yet captured} & \emph{Complete during review} \\
}
\subsubsection{Feature - Identity Correlation}\label{cap:bec-314be2d1}
\vspace{-1em}
\begin{table}[htbp]
  \centering
  \footnotesize
  \begin{tabularx}{1\textwidth}{|l|l|X|}
    \hline
    \textbf{Metadata Name} & \textbf{Metadata Value} & \textbf{Description} \\
    \hline
    \textbf{Reference Id} & \texttt{BEC-314BE2D1} & \\
    \hline
    \textbf{Status} & Generated candidate & \\
    \hline
    \textbf{Ancestry} &  & \\
    \hline
    \textbf{Source} & FHIR-informed candidate & \\
    \hline
    \textbf{Enabled By} &  & \\
    \hline
    \textbf{Supports} &  & \\
    \hline
    \textbf{Derived From} &  & \\
    \hline
    \textbf{Reviewed Status} & Generated candidate: keep / modify / relocate / reject; verify whether independently testable. & \\
    \hline
  \end{tabularx}
\end{table}

Intent. Identity Correlation is proposed to support the underlying healthcare concept represented by Practitioner Role. The proposal concerns behaviour and semantic responsibility, not a FHIR field, operation, resource instance or prescribed data model.

Conceptual evidence: HL7 FHIR R5 Practitioner, PractitionerRole.

\noindent\textbf{Feature boundary.} This behaviour is scoped to Practitioner Role; an identically named Feature elsewhere may carry different business meaning.\par
\featuretable{
--- & --- & \emph{Not yet captured} & \emph{Complete during review} \\
}
\subsubsection{Feature - Qualification Association}\label{cap:bec-a4596b27}
\vspace{-1em}
\begin{table}[htbp]
  \centering
  \footnotesize
  \begin{tabularx}{1\textwidth}{|l|l|X|}
    \hline
    \textbf{Metadata Name} & \textbf{Metadata Value} & \textbf{Description} \\
    \hline
    \textbf{Reference Id} & \texttt{BEC-A4596B27} & \\
    \hline
    \textbf{Status} & Generated candidate & \\
    \hline
    \textbf{Ancestry} &  & \\
    \hline
    \textbf{Source} & FHIR-informed candidate & \\
    \hline
    \textbf{Enabled By} &  & \\
    \hline
    \textbf{Supports} &  & \\
    \hline
    \textbf{Derived From} &  & \\
    \hline
    \textbf{Reviewed Status} & Generated candidate: keep / modify / relocate / reject; verify whether independently testable. & \\
    \hline
  \end{tabularx}
\end{table}

Intent. Qualification Association is proposed to support the underlying healthcare concept represented by Practitioner Role. The proposal concerns behaviour and semantic responsibility, not a FHIR field, operation, resource instance or prescribed data model.

Conceptual evidence: HL7 FHIR R5 Practitioner, PractitionerRole.

\noindent\textbf{Feature boundary.} This behaviour is scoped to Practitioner Role; an identically named Feature elsewhere may carry different business meaning.\par
\featuretable{
--- & --- & \emph{Not yet captured} & \emph{Complete during review} \\
}
\subsubsection{Feature - Professional Relationship Management}\label{cap:bec-549d792b}
\vspace{-1em}
\begin{table}[htbp]
  \centering
  \footnotesize
  \begin{tabularx}{1\textwidth}{|l|l|X|}
    \hline
    \textbf{Metadata Name} & \textbf{Metadata Value} & \textbf{Description} \\
    \hline
    \textbf{Reference Id} & \texttt{BEC-549D792B} & \\
    \hline
    \textbf{Status} & Generated candidate & \\
    \hline
    \textbf{Ancestry} &  & \\
    \hline
    \textbf{Source} & FHIR-informed candidate & \\
    \hline
    \textbf{Enabled By} &  & \\
    \hline
    \textbf{Supports} &  & \\
    \hline
    \textbf{Derived From} &  & \\
    \hline
    \textbf{Reviewed Status} & Generated candidate: keep / modify / relocate / reject; verify whether independently testable. & \\
    \hline
  \end{tabularx}
\end{table}

Intent. Professional Relationship Management is proposed to support the underlying healthcare concept represented by Practitioner Role. The proposal concerns behaviour and semantic responsibility, not a FHIR field, operation, resource instance or prescribed data model.

Conceptual evidence: HL7 FHIR R5 Practitioner, PractitionerRole.

\noindent\textbf{Feature boundary.} This behaviour is scoped to Practitioner Role; an identically named Feature elsewhere may carry different business meaning.\par
\featuretable{
--- & --- & \emph{Not yet captured} & \emph{Complete during review} \\
}
\subsubsection{Feature - Practitioner Discovery}\label{cap:bec-e24857d9}
\vspace{-1em}
\begin{table}[htbp]
  \centering
  \footnotesize
  \begin{tabularx}{1\textwidth}{|l|l|X|}
    \hline
    \textbf{Metadata Name} & \textbf{Metadata Value} & \textbf{Description} \\
    \hline
    \textbf{Reference Id} & \texttt{BEC-E24857D9} & \\
    \hline
    \textbf{Status} & Generated candidate & \\
    \hline
    \textbf{Ancestry} &  & \\
    \hline
    \textbf{Source} & FHIR-informed candidate & \\
    \hline
    \textbf{Enabled By} &  & \\
    \hline
    \textbf{Supports} &  & \\
    \hline
    \textbf{Derived From} &  & \\
    \hline
    \textbf{Reviewed Status} & Generated candidate: keep / modify / relocate / reject; verify whether independently testable. & \\
    \hline
  \end{tabularx}
\end{table}

Intent. Practitioner Discovery is proposed to support the underlying healthcare concept represented by Practitioner Role. The proposal concerns behaviour and semantic responsibility, not a FHIR field, operation, resource instance or prescribed data model.

Conceptual evidence: HL7 FHIR R5 Practitioner, PractitionerRole.

\noindent\textbf{Feature boundary.} This behaviour is scoped to Practitioner Role; an identically named Feature elsewhere may carry different business meaning.\par
\featuretable{
--- & --- & \emph{Not yet captured} & \emph{Complete during review} \\
}

\subsection{CT3 - Provider Relationship}\label{cap:bec-c410cc51}
\vspace{-0.5em}
\begin{table}[htbp]
  \centering
  \footnotesize
  \begin{tabularx}{1\textwidth}{|l|l|X|}
    \hline
    \multicolumn{3}{c}{\textbf{Capability Metadata}} \\
    \hline
    \textbf{Name} & \textbf{Value} & \textbf{Comment and/or Description} \\
    \hline
    \textbf{Reference Id} & \texttt{BEC-C410CC51} & \\
    \hline
    \textbf{Status} & Candidate ancestry & \\
    \hline
    \textbf{Ancestry} &  & \\
    \hline
    \textbf{Source} &  & \\
    \hline
    \textbf{Enabled By} &  & \\
    \hline
    \textbf{Supports} &  & \\
    \hline
    \textbf{Derived From} &  & \\
    \hline
    \textbf{Reviewed Status} &  & \\
    \hline
  \end{tabularx}
\end{table}

\begin{table}[htbp]
  \centering
  \footnotesize
  \begin{tabularx}{1\textwidth}{|l|l|X|}
    \hline
    \multicolumn{3}{c}{\textbf{Contained Features}} \\
    \hline
    \textbf{Feature Name} & \textbf{Feature Id} & \textbf{Comment} \\
    \hline
    Practitioner-Organisation Affiliation & \texttt{cap:bec-f5495dd1} (at \ref{cap:bec-f5495dd1}) & \\
    \hline
    Context Establishment & \texttt{cap:bec-0027885b} (at \ref{cap:bec-0027885b}) & \\
    \hline
    State Monitoring & \texttt{cap:bec-2cf899d2} (at \ref{cap:bec-2cf899d2}) & \\
    \hline
    Exception Management & \texttt{cap:bec-9806e183} (at \ref{cap:bec-9806e183}) & \\
    \hline
    Information Retrieval & \texttt{cap:bec-737e5a3f} (at \ref{cap:bec-737e5a3f}) & \\
    \hline
  \end{tabularx}
\end{table}

Intent. Provider Relationship is a candidate enabling capability within Provider Administration. It frames the reusable system-supported behaviours needed to support this healthcare context without prescribing a product, interface, technology or deployment boundary.

Boundary. This capability supports the business activity and does not inherit the authority or accountability of the Business Capability it enables. Its child Features are review candidates until ancestry and semantics are adjudicated.

FHIR R5 conceptual analogues considered (not implementation mappings): Practitioner, PractitionerRole, Organization.

\subsubsection{Feature - Practitioner-Organisation Affiliation}\label{cap:bec-f5495dd1}
\vspace{-1em}
\begin{table}[htbp]
  \centering
  \footnotesize
  \begin{tabularx}{1\textwidth}{|l|l|X|}
    \hline
    \textbf{Metadata Name} & \textbf{Metadata Value} & \textbf{Description} \\
    \hline
    \textbf{Reference Id} & \texttt{BEC-F5495DD1} & \\
    \hline
    \textbf{Status} & Existing Feature; candidate CT3 placement & \\
    \hline
    \textbf{Ancestry} &  & \\
    \hline
    \textbf{Source} & FEAT-EM-15 & \\
    \hline
    \textbf{Enabled By} &  & \\
    \hline
    \textbf{Supports} &  & \\
    \hline
    \textbf{Derived From} &  & \\
    \hline
    \textbf{Reviewed Status} & Confirm existing Feature meaning and proposed CT3 ownership. & \\
    \hline
  \end{tabularx}
\end{table}

Intent. Practitioner-Organisation Affiliation provides a distinct enabling behaviour in the context of Provider Relationship. This is an existing named Feature; its placement under the candidate CT3 remains subject to review.

\noindent\textbf{Feature boundary.} This behaviour is scoped to Provider Relationship; an identically named Feature elsewhere may carry different business meaning.\par
\featuretable{
Function & \texttt{FN-241752DB} & Maintain Practitioner Affiliations & Catalogue \\
Service & \texttt{SV-719DFDA5} & Practitioner Affiliation Resolution & Linked to preceding Function in source \\
}
\subsubsection{Feature - Context Establishment}\label{cap:bec-0027885b}
\vspace{-1em}
\begin{table}[htbp]
  \centering
  \footnotesize
  \begin{tabularx}{1\textwidth}{|l|l|X|}
    \hline
    \textbf{Metadata Name} & \textbf{Metadata Value} & \textbf{Description} \\
    \hline
    \textbf{Reference Id} & \texttt{BEC-0027885B} & \\
    \hline
    \textbf{Status} & Generated candidate & \\
    \hline
    \textbf{Ancestry} &  & \\
    \hline
    \textbf{Source} & Semantic-pattern candidate & \\
    \hline
    \textbf{Enabled By} &  & \\
    \hline
    \textbf{Supports} &  & \\
    \hline
    \textbf{Derived From} &  & \\
    \hline
    \textbf{Reviewed Status} & Generated candidate: keep / modify / relocate / reject; verify whether independently testable. & \\
    \hline
  \end{tabularx}
\end{table}

Intent. Context Establishment is proposed as a discrete, testable enabling behaviour for Provider Relationship. It is included for architectural review, not as an established implementation obligation.

\noindent\textbf{Feature boundary.} This behaviour is scoped to Provider Relationship; an identically named Feature elsewhere may carry different business meaning.\par
\featuretable{
--- & --- & \emph{Not yet captured} & \emph{Complete during review} \\
}
\subsubsection{Feature - State Monitoring}\label{cap:bec-2cf899d2}
\vspace{-1em}
\begin{table}[htbp]
  \centering
  \footnotesize
  \begin{tabularx}{1\textwidth}{|l|l|X|}
    \hline
    \textbf{Metadata Name} & \textbf{Metadata Value} & \textbf{Description} \\
    \hline
    \textbf{Reference Id} & \texttt{BEC-2CF899D2} & \\
    \hline
    \textbf{Status} & Generated candidate & \\
    \hline
    \textbf{Ancestry} &  & \\
    \hline
    \textbf{Source} & Semantic-pattern candidate & \\
    \hline
    \textbf{Enabled By} &  & \\
    \hline
    \textbf{Supports} &  & \\
    \hline
    \textbf{Derived From} &  & \\
    \hline
    \textbf{Reviewed Status} & Generated candidate: keep / modify / relocate / reject; verify whether independently testable. & \\
    \hline
  \end{tabularx}
\end{table}

Intent. State Monitoring is proposed as a discrete, testable enabling behaviour for Provider Relationship. It is included for architectural review, not as an established implementation obligation.

\noindent\textbf{Feature boundary.} This behaviour is scoped to Provider Relationship; an identically named Feature elsewhere may carry different business meaning.\par
\featuretable{
--- & --- & \emph{Not yet captured} & \emph{Complete during review} \\
}
\subsubsection{Feature - Exception Management}\label{cap:bec-9806e183}
\vspace{-1em}
\begin{table}[htbp]
  \centering
  \footnotesize
  \begin{tabularx}{1\textwidth}{|l|l|X|}
    \hline
    \textbf{Metadata Name} & \textbf{Metadata Value} & \textbf{Description} \\
    \hline
    \textbf{Reference Id} & \texttt{BEC-9806E183} & \\
    \hline
    \textbf{Status} & Generated candidate & \\
    \hline
    \textbf{Ancestry} &  & \\
    \hline
    \textbf{Source} & Semantic-pattern candidate & \\
    \hline
    \textbf{Enabled By} &  & \\
    \hline
    \textbf{Supports} &  & \\
    \hline
    \textbf{Derived From} &  & \\
    \hline
    \textbf{Reviewed Status} & Generated candidate: keep / modify / relocate / reject; verify whether independently testable. & \\
    \hline
  \end{tabularx}
\end{table}

Intent. Exception Management is proposed as a discrete, testable enabling behaviour for Provider Relationship. It is included for architectural review, not as an established implementation obligation.

\noindent\textbf{Feature boundary.} This behaviour is scoped to Provider Relationship; an identically named Feature elsewhere may carry different business meaning.\par
\featuretable{
--- & --- & \emph{Not yet captured} & \emph{Complete during review} \\
}
\subsubsection{Feature - Information Retrieval}\label{cap:bec-737e5a3f}
\vspace{-1em}
\begin{table}[htbp]
  \centering
  \footnotesize
  \begin{tabularx}{1\textwidth}{|l|l|X|}
    \hline
    \textbf{Metadata Name} & \textbf{Metadata Value} & \textbf{Description} \\
    \hline
    \textbf{Reference Id} & \texttt{BEC-737E5A3F} & \\
    \hline
    \textbf{Status} & Generated candidate & \\
    \hline
    \textbf{Ancestry} &  & \\
    \hline
    \textbf{Source} & Semantic-pattern candidate & \\
    \hline
    \textbf{Enabled By} &  & \\
    \hline
    \textbf{Supports} &  & \\
    \hline
    \textbf{Derived From} &  & \\
    \hline
    \textbf{Reviewed Status} & Generated candidate: keep / modify / relocate / reject; verify whether independently testable. & \\
    \hline
  \end{tabularx}
\end{table}

Intent. Information Retrieval is proposed as a discrete, testable enabling behaviour for Provider Relationship. It is included for architectural review, not as an established implementation obligation.

\noindent\textbf{Feature boundary.} This behaviour is scoped to Provider Relationship; an identically named Feature elsewhere may carry different business meaning.\par
\featuretable{
--- & --- & \emph{Not yet captured} & \emph{Complete during review} \\
}

\subsection{CT3 - Provider Engagement}\label{cap:bec-f58dba08}
\vspace{-0.5em}
\begin{table}[htbp]
  \centering
  \footnotesize
  \begin{tabularx}{1\textwidth}{|l|l|X|}
    \hline
    \multicolumn{3}{c}{\textbf{Capability Metadata}} \\
    \hline
    \textbf{Name} & \textbf{Value} & \textbf{Comment and/or Description} \\
    \hline
    \textbf{Reference Id} & \texttt{BEC-F58DBA08} & \\
    \hline
    \textbf{Status} & Candidate ancestry & \\
    \hline
    \textbf{Ancestry} &  & \\
    \hline
    \textbf{Source} &  & \\
    \hline
    \textbf{Enabled By} &  & \\
    \hline
    \textbf{Supports} &  & \\
    \hline
    \textbf{Derived From} &  & \\
    \hline
    \textbf{Reviewed Status} &  & \\
    \hline
  \end{tabularx}
\end{table}

\begin{table}[htbp]
  \centering
  \footnotesize
  \begin{tabularx}{1\textwidth}{|l|l|X|}
    \hline
    \multicolumn{3}{c}{\textbf{Contained Features}} \\
    \hline
    \textbf{Feature Name} & \textbf{Feature Id} & \textbf{Comment} \\
    \hline
    Clinical Privileging Record & \texttt{cap:bec-e0307698} (at \ref{cap:bec-e0307698}) & \\
    \hline
    Context Establishment & \texttt{cap:bec-0c362f3d} (at \ref{cap:bec-0c362f3d}) & \\
    \hline
    State Monitoring & \texttt{cap:bec-141e9549} (at \ref{cap:bec-141e9549}) & \\
    \hline
    Exception Management & \texttt{cap:bec-788a21bf} (at \ref{cap:bec-788a21bf}) & \\
    \hline
    Information Retrieval & \texttt{cap:bec-167af187} (at \ref{cap:bec-167af187}) & \\
    \hline
  \end{tabularx}
\end{table}

Intent. Provider Engagement is a candidate enabling capability within Provider Administration. It frames the reusable system-supported behaviours needed to support this healthcare context without prescribing a product, interface, technology or deployment boundary.

Boundary. This capability supports the business activity and does not inherit the authority or accountability of the Business Capability it enables. Its child Features are review candidates until ancestry and semantics are adjudicated.

FHIR R5 conceptual analogues considered (not implementation mappings): Practitioner, PractitionerRole, Organization.

\subsubsection{Feature - Clinical Privileging Record}\label{cap:bec-e0307698}
\vspace{-1em}
\begin{table}[htbp]
  \centering
  \footnotesize
  \begin{tabularx}{1\textwidth}{|l|l|X|}
    \hline
    \textbf{Metadata Name} & \textbf{Metadata Value} & \textbf{Description} \\
    \hline
    \textbf{Reference Id} & \texttt{BEC-E0307698} & \\
    \hline
    \textbf{Status} & Existing Feature; candidate CT3 placement & \\
    \hline
    \textbf{Ancestry} &  & \\
    \hline
    \textbf{Source} & FEAT-EM-16 & \\
    \hline
    \textbf{Enabled By} &  & \\
    \hline
    \textbf{Supports} &  & \\
    \hline
    \textbf{Derived From} &  & \\
    \hline
    \textbf{Reviewed Status} & Confirm existing Feature meaning and proposed CT3 ownership. & \\
    \hline
  \end{tabularx}
\end{table}

Intent. Clinical Privileging Record provides a distinct enabling behaviour in the context of Provider Engagement. This is an existing named Feature; its placement under the candidate CT3 remains subject to review.

\noindent\textbf{Feature boundary.} This behaviour is scoped to Provider Engagement; an identically named Feature elsewhere may carry different business meaning.\par
\featuretable{
Function & \texttt{FN-794014E6} & Maintain Clinical Privileges & Catalogue \\
Service & \texttt{SV-D4EF110C} & Clinical Privilege Verification & Linked to preceding Function in source \\
}
\subsubsection{Feature - Context Establishment}\label{cap:bec-0c362f3d}
\vspace{-1em}
\begin{table}[htbp]
  \centering
  \footnotesize
  \begin{tabularx}{1\textwidth}{|l|l|X|}
    \hline
    \textbf{Metadata Name} & \textbf{Metadata Value} & \textbf{Description} \\
    \hline
    \textbf{Reference Id} & \texttt{BEC-0C362F3D} & \\
    \hline
    \textbf{Status} & Generated candidate & \\
    \hline
    \textbf{Ancestry} &  & \\
    \hline
    \textbf{Source} & Semantic-pattern candidate & \\
    \hline
    \textbf{Enabled By} &  & \\
    \hline
    \textbf{Supports} &  & \\
    \hline
    \textbf{Derived From} &  & \\
    \hline
    \textbf{Reviewed Status} & Generated candidate: keep / modify / relocate / reject; verify whether independently testable. & \\
    \hline
  \end{tabularx}
\end{table}

Intent. Context Establishment is proposed as a discrete, testable enabling behaviour for Provider Engagement. It is included for architectural review, not as an established implementation obligation.

\noindent\textbf{Feature boundary.} This behaviour is scoped to Provider Engagement; an identically named Feature elsewhere may carry different business meaning.\par
\featuretable{
--- & --- & \emph{Not yet captured} & \emph{Complete during review} \\
}
\subsubsection{Feature - State Monitoring}\label{cap:bec-141e9549}
\vspace{-1em}
\begin{table}[htbp]
  \centering
  \footnotesize
  \begin{tabularx}{1\textwidth}{|l|l|X|}
    \hline
    \textbf{Metadata Name} & \textbf{Metadata Value} & \textbf{Description} \\
    \hline
    \textbf{Reference Id} & \texttt{BEC-141E9549} & \\
    \hline
    \textbf{Status} & Generated candidate & \\
    \hline
    \textbf{Ancestry} &  & \\
    \hline
    \textbf{Source} & Semantic-pattern candidate & \\
    \hline
    \textbf{Enabled By} &  & \\
    \hline
    \textbf{Supports} &  & \\
    \hline
    \textbf{Derived From} &  & \\
    \hline
    \textbf{Reviewed Status} & Generated candidate: keep / modify / relocate / reject; verify whether independently testable. & \\
    \hline
  \end{tabularx}
\end{table}

Intent. State Monitoring is proposed as a discrete, testable enabling behaviour for Provider Engagement. It is included for architectural review, not as an established implementation obligation.

\noindent\textbf{Feature boundary.} This behaviour is scoped to Provider Engagement; an identically named Feature elsewhere may carry different business meaning.\par
\featuretable{
--- & --- & \emph{Not yet captured} & \emph{Complete during review} \\
}
\subsubsection{Feature - Exception Management}\label{cap:bec-788a21bf}
\vspace{-1em}
\begin{table}[htbp]
  \centering
  \footnotesize
  \begin{tabularx}{1\textwidth}{|l|l|X|}
    \hline
    \textbf{Metadata Name} & \textbf{Metadata Value} & \textbf{Description} \\
    \hline
    \textbf{Reference Id} & \texttt{BEC-788A21BF} & \\
    \hline
    \textbf{Status} & Generated candidate & \\
    \hline
    \textbf{Ancestry} &  & \\
    \hline
    \textbf{Source} & Semantic-pattern candidate & \\
    \hline
    \textbf{Enabled By} &  & \\
    \hline
    \textbf{Supports} &  & \\
    \hline
    \textbf{Derived From} &  & \\
    \hline
    \textbf{Reviewed Status} & Generated candidate: keep / modify / relocate / reject; verify whether independently testable. & \\
    \hline
  \end{tabularx}
\end{table}

Intent. Exception Management is proposed as a discrete, testable enabling behaviour for Provider Engagement. It is included for architectural review, not as an established implementation obligation.

\noindent\textbf{Feature boundary.} This behaviour is scoped to Provider Engagement; an identically named Feature elsewhere may carry different business meaning.\par
\featuretable{
--- & --- & \emph{Not yet captured} & \emph{Complete during review} \\
}
\subsubsection{Feature - Information Retrieval}\label{cap:bec-167af187}
\vspace{-1em}
\begin{table}[htbp]
  \centering
  \footnotesize
  \begin{tabularx}{1\textwidth}{|l|l|X|}
    \hline
    \textbf{Metadata Name} & \textbf{Metadata Value} & \textbf{Description} \\
    \hline
    \textbf{Reference Id} & \texttt{BEC-167AF187} & \\
    \hline
    \textbf{Status} & Generated candidate & \\
    \hline
    \textbf{Ancestry} &  & \\
    \hline
    \textbf{Source} & Semantic-pattern candidate & \\
    \hline
    \textbf{Enabled By} &  & \\
    \hline
    \textbf{Supports} &  & \\
    \hline
    \textbf{Derived From} &  & \\
    \hline
    \textbf{Reviewed Status} & Generated candidate: keep / modify / relocate / reject; verify whether independently testable. & \\
    \hline
  \end{tabularx}
\end{table}

Intent. Information Retrieval is proposed as a discrete, testable enabling behaviour for Provider Engagement. It is included for architectural review, not as an established implementation obligation.

\noindent\textbf{Feature boundary.} This behaviour is scoped to Provider Engagement; an identically named Feature elsewhere may carry different business meaning.\par
\featuretable{
--- & --- & \emph{Not yet captured} & \emph{Complete during review} \\
}

\section{CT2 - Organisation Administration}\label{cap:bec-1ff88a71}
\vspace{-0.5em}
\begin{table}[htbp]
  \centering
  \footnotesize
  \begin{tabularx}{1\textwidth}{|l|l|X|}
    \hline
    \multicolumn{3}{c}{\textbf{Capability Metadata}} \\
    \hline
    \textbf{Name} & \textbf{Value} & \textbf{Comment and/or Description} \\
    \hline
    \textbf{Reference Id} & \texttt{BEC-1FF88A71} & \\
    \hline
    \textbf{Status} & Candidate ancestry & \\
    \hline
    \textbf{Ancestry} &  & \\
    \hline
    \textbf{Source} &  & \\
    \hline
    \textbf{Enabled By} &  & \\
    \hline
    \textbf{Supports} &  & \\
    \hline
    \textbf{Derived From} &  & \\
    \hline
    \textbf{Reviewed Status} &  & \\
    \hline
  \end{tabularx}
\end{table}

\begin{table}[htbp]
  \centering
  \footnotesize
  \begin{tabularx}{1\textwidth}{|l|l|X|}
    \hline
    \multicolumn{3}{c}{\textbf{Contained Capabilities}} \\
    \hline
    \textbf{Capability Name} & \textbf{Capability Id} & \textbf{Comment} \\
    \hline
    Healthcare Organisation & \texttt{cap:bec-ace86b9b} (at \ref{cap:bec-ace86b9b}) & \\
    \hline
    Organisational Structure & \texttt{cap:bec-fa161a55} (at \ref{cap:bec-fa161a55}) & \\
    \hline
    Organisation Contact & \texttt{cap:bec-5c12bf05} (at \ref{cap:bec-5c12bf05}) & \\
    \hline
  \end{tabularx}
\end{table}

Intent. Organisation Administration is a candidate enabling capability within Entity Information Enablement. It frames the reusable system-supported behaviours needed to support this healthcare context without prescribing a product, interface, technology or deployment boundary.

Boundary. This capability supports the business activity and does not inherit the authority or accountability of the Business Capability it enables. Its child Features are review candidates until ancestry and semantics are adjudicated.

\subsection{CT3 - Healthcare Organisation}\label{cap:bec-ace86b9b}
\vspace{-0.5em}
\begin{table}[htbp]
  \centering
  \footnotesize
  \begin{tabularx}{1\textwidth}{|l|l|X|}
    \hline
    \multicolumn{3}{c}{\textbf{Capability Metadata}} \\
    \hline
    \textbf{Name} & \textbf{Value} & \textbf{Comment and/or Description} \\
    \hline
    \textbf{Reference Id} & \texttt{BEC-ACE86B9B} & \\
    \hline
    \textbf{Status} & Candidate ancestry & \\
    \hline
    \textbf{Ancestry} &  & \\
    \hline
    \textbf{Source} &  & \\
    \hline
    \textbf{Enabled By} &  & \\
    \hline
    \textbf{Supports} &  & \\
    \hline
    \textbf{Derived From} &  & \\
    \hline
    \textbf{Reviewed Status} &  & \\
    \hline
  \end{tabularx}
\end{table}

\begin{table}[htbp]
  \centering
  \footnotesize
  \begin{tabularx}{1\textwidth}{|l|l|X|}
    \hline
    \multicolumn{3}{c}{\textbf{Contained Features}} \\
    \hline
    \textbf{Feature Name} & \textbf{Feature Id} & \textbf{Comment} \\
    \hline
    National Organisation Verification & \texttt{cap:bec-28102e35} (at \ref{cap:bec-28102e35}) & \\
    \hline
    Organisation Profile Management & \texttt{cap:bec-df977eab} (at \ref{cap:bec-df977eab}) & \\
    \hline
    Identity Management & \texttt{cap:bec-157386a8} (at \ref{cap:bec-157386a8}) & \\
    \hline
    Identity Correlation & \texttt{cap:bec-9262fb93} (at \ref{cap:bec-9262fb93}) & \\
    \hline
    Organisational Relationship Management & \texttt{cap:bec-ad09d42a} (at \ref{cap:bec-ad09d42a}) & \\
    \hline
    Organisational Structure Management & \texttt{cap:bec-6b51a123} (at \ref{cap:bec-6b51a123}) & \\
    \hline
    Organisation Discovery & \texttt{cap:bec-d6fb2614} (at \ref{cap:bec-d6fb2614}) & \\
    \hline
  \end{tabularx}
\end{table}

Intent. Healthcare Organisation is a candidate enabling capability within Organisation Administration. It frames the reusable system-supported behaviours needed to support this healthcare context without prescribing a product, interface, technology or deployment boundary.

Boundary. This capability supports the business activity and does not inherit the authority or accountability of the Business Capability it enables. Its child Features are review candidates until ancestry and semantics are adjudicated.

FHIR R5 conceptual analogues considered (not implementation mappings): Organization, OrganizationAffiliation.

\subsubsection{Feature - National Organisation Verification}\label{cap:bec-28102e35}
\vspace{-1em}
\begin{table}[htbp]
  \centering
  \footnotesize
  \begin{tabularx}{1\textwidth}{|l|l|X|}
    \hline
    \textbf{Metadata Name} & \textbf{Metadata Value} & \textbf{Description} \\
    \hline
    \textbf{Reference Id} & \texttt{BEC-28102E35} & \\
    \hline
    \textbf{Status} & Existing Feature; candidate CT3 placement & \\
    \hline
    \textbf{Ancestry} &  & \\
    \hline
    \textbf{Source} & FEAT-EM-17 & \\
    \hline
    \textbf{Enabled By} &  & \\
    \hline
    \textbf{Supports} &  & \\
    \hline
    \textbf{Derived From} &  & \\
    \hline
    \textbf{Reviewed Status} & Confirm existing Feature meaning and proposed CT3 ownership. & \\
    \hline
  \end{tabularx}
\end{table}

Intent. National Organisation Verification provides a distinct enabling behaviour in the context of Healthcare Organisation. This is an existing named Feature; its placement under the candidate CT3 remains subject to review.

\noindent\textbf{Feature boundary.} This behaviour is scoped to Healthcare Organisation; an identically named Feature elsewhere may carry different business meaning.\par
\featuretable{
Function & \texttt{FN-48725D18} & Verify Healthcare Organisation & Catalogue \\
Service & \texttt{SV-66391EE5} & Organisation Verification Service & Linked to preceding Function in source \\
}
\subsubsection{Feature - Organisation Profile Management}\label{cap:bec-df977eab}
\vspace{-1em}
\begin{table}[htbp]
  \centering
  \footnotesize
  \begin{tabularx}{1\textwidth}{|l|l|X|}
    \hline
    \textbf{Metadata Name} & \textbf{Metadata Value} & \textbf{Description} \\
    \hline
    \textbf{Reference Id} & \texttt{BEC-DF977EAB} & \\
    \hline
    \textbf{Status} & Existing Feature; candidate CT3 placement & \\
    \hline
    \textbf{Ancestry} &  & \\
    \hline
    \textbf{Source} & FEAT-EM-18 & \\
    \hline
    \textbf{Enabled By} &  & \\
    \hline
    \textbf{Supports} &  & \\
    \hline
    \textbf{Derived From} &  & \\
    \hline
    \textbf{Reviewed Status} & Confirm existing Feature meaning and proposed CT3 ownership. & \\
    \hline
  \end{tabularx}
\end{table}

Intent. Organisation Profile Management provides a distinct enabling behaviour in the context of Healthcare Organisation. This is an existing named Feature; its placement under the candidate CT3 remains subject to review.

\noindent\textbf{Feature boundary.} This behaviour is scoped to Healthcare Organisation; an identically named Feature elsewhere may carry different business meaning.\par
\featuretable{
Function & \texttt{FN-BE033910} & Govern Organisation Profile & Catalogue \\
Service & \texttt{SV-9EE2C43D} & Organisation Directory Query & Linked to preceding Function in source \\
}
\subsubsection{Feature - Identity Management}\label{cap:bec-157386a8}
\vspace{-1em}
\begin{table}[htbp]
  \centering
  \footnotesize
  \begin{tabularx}{1\textwidth}{|l|l|X|}
    \hline
    \textbf{Metadata Name} & \textbf{Metadata Value} & \textbf{Description} \\
    \hline
    \textbf{Reference Id} & \texttt{BEC-157386A8} & \\
    \hline
    \textbf{Status} & Generated candidate & \\
    \hline
    \textbf{Ancestry} &  & \\
    \hline
    \textbf{Source} & FHIR-informed candidate & \\
    \hline
    \textbf{Enabled By} &  & \\
    \hline
    \textbf{Supports} &  & \\
    \hline
    \textbf{Derived From} &  & \\
    \hline
    \textbf{Reviewed Status} & Generated candidate: keep / modify / relocate / reject; verify whether independently testable. & \\
    \hline
  \end{tabularx}
\end{table}

Intent. Identity Management is proposed to support the underlying healthcare concept represented by Healthcare Organisation. The proposal concerns behaviour and semantic responsibility, not a FHIR field, operation, resource instance or prescribed data model.

Conceptual evidence: HL7 FHIR R5 Organization, OrganizationAffiliation.

\noindent\textbf{Feature boundary.} This behaviour is scoped to Healthcare Organisation; an identically named Feature elsewhere may carry different business meaning.\par
\featuretable{
--- & --- & \emph{Not yet captured} & \emph{Complete during review} \\
}
\subsubsection{Feature - Identity Correlation}\label{cap:bec-9262fb93}
\vspace{-1em}
\begin{table}[htbp]
  \centering
  \footnotesize
  \begin{tabularx}{1\textwidth}{|l|l|X|}
    \hline
    \textbf{Metadata Name} & \textbf{Metadata Value} & \textbf{Description} \\
    \hline
    \textbf{Reference Id} & \texttt{BEC-9262FB93} & \\
    \hline
    \textbf{Status} & Generated candidate & \\
    \hline
    \textbf{Ancestry} &  & \\
    \hline
    \textbf{Source} & FHIR-informed candidate & \\
    \hline
    \textbf{Enabled By} &  & \\
    \hline
    \textbf{Supports} &  & \\
    \hline
    \textbf{Derived From} &  & \\
    \hline
    \textbf{Reviewed Status} & Generated candidate: keep / modify / relocate / reject; verify whether independently testable. & \\
    \hline
  \end{tabularx}
\end{table}

Intent. Identity Correlation is proposed to support the underlying healthcare concept represented by Healthcare Organisation. The proposal concerns behaviour and semantic responsibility, not a FHIR field, operation, resource instance or prescribed data model.

Conceptual evidence: HL7 FHIR R5 Organization, OrganizationAffiliation.

\noindent\textbf{Feature boundary.} This behaviour is scoped to Healthcare Organisation; an identically named Feature elsewhere may carry different business meaning.\par
\featuretable{
--- & --- & \emph{Not yet captured} & \emph{Complete during review} \\
}
\subsubsection{Feature - Organisational Relationship Management}\label{cap:bec-ad09d42a}
\vspace{-1em}
\begin{table}[htbp]
  \centering
  \footnotesize
  \begin{tabularx}{1\textwidth}{|l|l|X|}
    \hline
    \textbf{Metadata Name} & \textbf{Metadata Value} & \textbf{Description} \\
    \hline
    \textbf{Reference Id} & \texttt{BEC-AD09D42A} & \\
    \hline
    \textbf{Status} & Generated candidate & \\
    \hline
    \textbf{Ancestry} &  & \\
    \hline
    \textbf{Source} & FHIR-informed candidate & \\
    \hline
    \textbf{Enabled By} &  & \\
    \hline
    \textbf{Supports} &  & \\
    \hline
    \textbf{Derived From} &  & \\
    \hline
    \textbf{Reviewed Status} & Generated candidate: keep / modify / relocate / reject; verify whether independently testable. & \\
    \hline
  \end{tabularx}
\end{table}

Intent. Organisational Relationship Management is proposed to support the underlying healthcare concept represented by Healthcare Organisation. The proposal concerns behaviour and semantic responsibility, not a FHIR field, operation, resource instance or prescribed data model.

Conceptual evidence: HL7 FHIR R5 Organization, OrganizationAffiliation.

\noindent\textbf{Feature boundary.} This behaviour is scoped to Healthcare Organisation; an identically named Feature elsewhere may carry different business meaning.\par
\featuretable{
--- & --- & \emph{Not yet captured} & \emph{Complete during review} \\
}
\subsubsection{Feature - Organisational Structure Management}\label{cap:bec-6b51a123}
\vspace{-1em}
\begin{table}[htbp]
  \centering
  \footnotesize
  \begin{tabularx}{1\textwidth}{|l|l|X|}
    \hline
    \textbf{Metadata Name} & \textbf{Metadata Value} & \textbf{Description} \\
    \hline
    \textbf{Reference Id} & \texttt{BEC-6B51A123} & \\
    \hline
    \textbf{Status} & Generated candidate & \\
    \hline
    \textbf{Ancestry} &  & \\
    \hline
    \textbf{Source} & FHIR-informed candidate & \\
    \hline
    \textbf{Enabled By} &  & \\
    \hline
    \textbf{Supports} &  & \\
    \hline
    \textbf{Derived From} &  & \\
    \hline
    \textbf{Reviewed Status} & Generated candidate: keep / modify / relocate / reject; verify whether independently testable. & \\
    \hline
  \end{tabularx}
\end{table}

Intent. Organisational Structure Management is proposed to support the underlying healthcare concept represented by Healthcare Organisation. The proposal concerns behaviour and semantic responsibility, not a FHIR field, operation, resource instance or prescribed data model.

Conceptual evidence: HL7 FHIR R5 Organization, OrganizationAffiliation.

\noindent\textbf{Feature boundary.} This behaviour is scoped to Healthcare Organisation; an identically named Feature elsewhere may carry different business meaning.\par
\featuretable{
--- & --- & \emph{Not yet captured} & \emph{Complete during review} \\
}
\subsubsection{Feature - Organisation Discovery}\label{cap:bec-d6fb2614}
\vspace{-1em}
\begin{table}[htbp]
  \centering
  \footnotesize
  \begin{tabularx}{1\textwidth}{|l|l|X|}
    \hline
    \textbf{Metadata Name} & \textbf{Metadata Value} & \textbf{Description} \\
    \hline
    \textbf{Reference Id} & \texttt{BEC-D6FB2614} & \\
    \hline
    \textbf{Status} & Generated candidate & \\
    \hline
    \textbf{Ancestry} &  & \\
    \hline
    \textbf{Source} & FHIR-informed candidate & \\
    \hline
    \textbf{Enabled By} &  & \\
    \hline
    \textbf{Supports} &  & \\
    \hline
    \textbf{Derived From} &  & \\
    \hline
    \textbf{Reviewed Status} & Generated candidate: keep / modify / relocate / reject; verify whether independently testable. & \\
    \hline
  \end{tabularx}
\end{table}

Intent. Organisation Discovery is proposed to support the underlying healthcare concept represented by Healthcare Organisation. The proposal concerns behaviour and semantic responsibility, not a FHIR field, operation, resource instance or prescribed data model.

Conceptual evidence: HL7 FHIR R5 Organization, OrganizationAffiliation.

\noindent\textbf{Feature boundary.} This behaviour is scoped to Healthcare Organisation; an identically named Feature elsewhere may carry different business meaning.\par
\featuretable{
--- & --- & \emph{Not yet captured} & \emph{Complete during review} \\
}

\subsection{CT3 - Organisational Structure}\label{cap:bec-fa161a55}
\vspace{-0.5em}
\begin{table}[htbp]
  \centering
  \footnotesize
  \begin{tabularx}{1\textwidth}{|l|l|X|}
    \hline
    \multicolumn{3}{c}{\textbf{Capability Metadata}} \\
    \hline
    \textbf{Name} & \textbf{Value} & \textbf{Comment and/or Description} \\
    \hline
    \textbf{Reference Id} & \texttt{BEC-FA161A55} & \\
    \hline
    \textbf{Status} & Candidate ancestry & \\
    \hline
    \textbf{Ancestry} &  & \\
    \hline
    \textbf{Source} &  & \\
    \hline
    \textbf{Enabled By} &  & \\
    \hline
    \textbf{Supports} &  & \\
    \hline
    \textbf{Derived From} &  & \\
    \hline
    \textbf{Reviewed Status} &  & \\
    \hline
  \end{tabularx}
\end{table}

\begin{table}[htbp]
  \centering
  \footnotesize
  \begin{tabularx}{1\textwidth}{|l|l|X|}
    \hline
    \multicolumn{3}{c}{\textbf{Contained Features}} \\
    \hline
    \textbf{Feature Name} & \textbf{Feature Id} & \textbf{Comment} \\
    \hline
    Hierarchical Structure Navigation & \texttt{cap:bec-f5068c04} (at \ref{cap:bec-f5068c04}) & \\
    \hline
    Identity Management & \texttt{cap:bec-48e4c563} (at \ref{cap:bec-48e4c563}) & \\
    \hline
    Identity Correlation & \texttt{cap:bec-a50f5059} (at \ref{cap:bec-a50f5059}) & \\
    \hline
    Organisational Relationship Management & \texttt{cap:bec-676858f6} (at \ref{cap:bec-676858f6}) & \\
    \hline
    Organisational Structure Management & \texttt{cap:bec-acab1168} (at \ref{cap:bec-acab1168}) & \\
    \hline
    Organisation Discovery & \texttt{cap:bec-721b38d1} (at \ref{cap:bec-721b38d1}) & \\
    \hline
  \end{tabularx}
\end{table}

Intent. Organisational Structure is a candidate enabling capability within Organisation Administration. It frames the reusable system-supported behaviours needed to support this healthcare context without prescribing a product, interface, technology or deployment boundary.

Boundary. This capability supports the business activity and does not inherit the authority or accountability of the Business Capability it enables. Its child Features are review candidates until ancestry and semantics are adjudicated.

FHIR R5 conceptual analogues considered (not implementation mappings): Organization, OrganizationAffiliation.

\subsubsection{Feature - Hierarchical Structure Navigation}\label{cap:bec-f5068c04}
\vspace{-1em}
\begin{table}[htbp]
  \centering
  \footnotesize
  \begin{tabularx}{1\textwidth}{|l|l|X|}
    \hline
    \textbf{Metadata Name} & \textbf{Metadata Value} & \textbf{Description} \\
    \hline
    \textbf{Reference Id} & \texttt{BEC-F5068C04} & \\
    \hline
    \textbf{Status} & Existing Feature; candidate CT3 placement & \\
    \hline
    \textbf{Ancestry} &  & \\
    \hline
    \textbf{Source} & FEAT-EM-19 & \\
    \hline
    \textbf{Enabled By} &  & \\
    \hline
    \textbf{Supports} &  & \\
    \hline
    \textbf{Derived From} &  & \\
    \hline
    \textbf{Reviewed Status} & Confirm existing Feature meaning and proposed CT3 ownership. & \\
    \hline
  \end{tabularx}
\end{table}

Intent. Hierarchical Structure Navigation provides a distinct enabling behaviour in the context of Organisational Structure. This is an existing named Feature; its placement under the candidate CT3 remains subject to review.

\noindent\textbf{Feature boundary.} This behaviour is scoped to Organisational Structure; an identically named Feature elsewhere may carry different business meaning.\par
\featuretable{
Function & \texttt{FN-9D01DB36} & Maintain Organisational Structure & Catalogue \\
Service & \texttt{SV-4765132F} & Organisational Hierarchy Resolution & Linked to preceding Function in source \\
}
\subsubsection{Feature - Identity Management}\label{cap:bec-48e4c563}
\vspace{-1em}
\begin{table}[htbp]
  \centering
  \footnotesize
  \begin{tabularx}{1\textwidth}{|l|l|X|}
    \hline
    \textbf{Metadata Name} & \textbf{Metadata Value} & \textbf{Description} \\
    \hline
    \textbf{Reference Id} & \texttt{BEC-48E4C563} & \\
    \hline
    \textbf{Status} & Generated candidate & \\
    \hline
    \textbf{Ancestry} &  & \\
    \hline
    \textbf{Source} & FHIR-informed candidate & \\
    \hline
    \textbf{Enabled By} &  & \\
    \hline
    \textbf{Supports} &  & \\
    \hline
    \textbf{Derived From} &  & \\
    \hline
    \textbf{Reviewed Status} & Generated candidate: keep / modify / relocate / reject; verify whether independently testable. & \\
    \hline
  \end{tabularx}
\end{table}

Intent. Identity Management is proposed to support the underlying healthcare concept represented by Organisational Structure. The proposal concerns behaviour and semantic responsibility, not a FHIR field, operation, resource instance or prescribed data model.

Conceptual evidence: HL7 FHIR R5 Organization, OrganizationAffiliation.

\noindent\textbf{Feature boundary.} This behaviour is scoped to Organisational Structure; an identically named Feature elsewhere may carry different business meaning.\par
\featuretable{
--- & --- & \emph{Not yet captured} & \emph{Complete during review} \\
}
\subsubsection{Feature - Identity Correlation}\label{cap:bec-a50f5059}
\vspace{-1em}
\begin{table}[htbp]
  \centering
  \footnotesize
  \begin{tabularx}{1\textwidth}{|l|l|X|}
    \hline
    \textbf{Metadata Name} & \textbf{Metadata Value} & \textbf{Description} \\
    \hline
    \textbf{Reference Id} & \texttt{BEC-A50F5059} & \\
    \hline
    \textbf{Status} & Generated candidate & \\
    \hline
    \textbf{Ancestry} &  & \\
    \hline
    \textbf{Source} & FHIR-informed candidate & \\
    \hline
    \textbf{Enabled By} &  & \\
    \hline
    \textbf{Supports} &  & \\
    \hline
    \textbf{Derived From} &  & \\
    \hline
    \textbf{Reviewed Status} & Generated candidate: keep / modify / relocate / reject; verify whether independently testable. & \\
    \hline
  \end{tabularx}
\end{table}

Intent. Identity Correlation is proposed to support the underlying healthcare concept represented by Organisational Structure. The proposal concerns behaviour and semantic responsibility, not a FHIR field, operation, resource instance or prescribed data model.

Conceptual evidence: HL7 FHIR R5 Organization, OrganizationAffiliation.

\noindent\textbf{Feature boundary.} This behaviour is scoped to Organisational Structure; an identically named Feature elsewhere may carry different business meaning.\par
\featuretable{
--- & --- & \emph{Not yet captured} & \emph{Complete during review} \\
}
\subsubsection{Feature - Organisational Relationship Management}\label{cap:bec-676858f6}
\vspace{-1em}
\begin{table}[htbp]
  \centering
  \footnotesize
  \begin{tabularx}{1\textwidth}{|l|l|X|}
    \hline
    \textbf{Metadata Name} & \textbf{Metadata Value} & \textbf{Description} \\
    \hline
    \textbf{Reference Id} & \texttt{BEC-676858F6} & \\
    \hline
    \textbf{Status} & Generated candidate & \\
    \hline
    \textbf{Ancestry} &  & \\
    \hline
    \textbf{Source} & FHIR-informed candidate & \\
    \hline
    \textbf{Enabled By} &  & \\
    \hline
    \textbf{Supports} &  & \\
    \hline
    \textbf{Derived From} &  & \\
    \hline
    \textbf{Reviewed Status} & Generated candidate: keep / modify / relocate / reject; verify whether independently testable. & \\
    \hline
  \end{tabularx}
\end{table}

Intent. Organisational Relationship Management is proposed to support the underlying healthcare concept represented by Organisational Structure. The proposal concerns behaviour and semantic responsibility, not a FHIR field, operation, resource instance or prescribed data model.

Conceptual evidence: HL7 FHIR R5 Organization, OrganizationAffiliation.

\noindent\textbf{Feature boundary.} This behaviour is scoped to Organisational Structure; an identically named Feature elsewhere may carry different business meaning.\par
\featuretable{
--- & --- & \emph{Not yet captured} & \emph{Complete during review} \\
}
\subsubsection{Feature - Organisational Structure Management}\label{cap:bec-acab1168}
\vspace{-1em}
\begin{table}[htbp]
  \centering
  \footnotesize
  \begin{tabularx}{1\textwidth}{|l|l|X|}
    \hline
    \textbf{Metadata Name} & \textbf{Metadata Value} & \textbf{Description} \\
    \hline
    \textbf{Reference Id} & \texttt{BEC-ACAB1168} & \\
    \hline
    \textbf{Status} & Generated candidate & \\
    \hline
    \textbf{Ancestry} &  & \\
    \hline
    \textbf{Source} & FHIR-informed candidate & \\
    \hline
    \textbf{Enabled By} &  & \\
    \hline
    \textbf{Supports} &  & \\
    \hline
    \textbf{Derived From} &  & \\
    \hline
    \textbf{Reviewed Status} & Generated candidate: keep / modify / relocate / reject; verify whether independently testable. & \\
    \hline
  \end{tabularx}
\end{table}

Intent. Organisational Structure Management is proposed to support the underlying healthcare concept represented by Organisational Structure. The proposal concerns behaviour and semantic responsibility, not a FHIR field, operation, resource instance or prescribed data model.

Conceptual evidence: HL7 FHIR R5 Organization, OrganizationAffiliation.

\noindent\textbf{Feature boundary.} This behaviour is scoped to Organisational Structure; an identically named Feature elsewhere may carry different business meaning.\par
\featuretable{
--- & --- & \emph{Not yet captured} & \emph{Complete during review} \\
}
\subsubsection{Feature - Organisation Discovery}\label{cap:bec-721b38d1}
\vspace{-1em}
\begin{table}[htbp]
  \centering
  \footnotesize
  \begin{tabularx}{1\textwidth}{|l|l|X|}
    \hline
    \textbf{Metadata Name} & \textbf{Metadata Value} & \textbf{Description} \\
    \hline
    \textbf{Reference Id} & \texttt{BEC-721B38D1} & \\
    \hline
    \textbf{Status} & Generated candidate & \\
    \hline
    \textbf{Ancestry} &  & \\
    \hline
    \textbf{Source} & FHIR-informed candidate & \\
    \hline
    \textbf{Enabled By} &  & \\
    \hline
    \textbf{Supports} &  & \\
    \hline
    \textbf{Derived From} &  & \\
    \hline
    \textbf{Reviewed Status} & Generated candidate: keep / modify / relocate / reject; verify whether independently testable. & \\
    \hline
  \end{tabularx}
\end{table}

Intent. Organisation Discovery is proposed to support the underlying healthcare concept represented by Organisational Structure. The proposal concerns behaviour and semantic responsibility, not a FHIR field, operation, resource instance or prescribed data model.

Conceptual evidence: HL7 FHIR R5 Organization, OrganizationAffiliation.

\noindent\textbf{Feature boundary.} This behaviour is scoped to Organisational Structure; an identically named Feature elsewhere may carry different business meaning.\par
\featuretable{
--- & --- & \emph{Not yet captured} & \emph{Complete during review} \\
}

\subsection{CT3 - Organisation Contact}\label{cap:bec-5c12bf05}
\vspace{-0.5em}
\begin{table}[htbp]
  \centering
  \footnotesize
  \begin{tabularx}{1\textwidth}{|l|l|X|}
    \hline
    \multicolumn{3}{c}{\textbf{Capability Metadata}} \\
    \hline
    \textbf{Name} & \textbf{Value} & \textbf{Comment and/or Description} \\
    \hline
    \textbf{Reference Id} & \texttt{BEC-5C12BF05} & \\
    \hline
    \textbf{Status} & Candidate ancestry & \\
    \hline
    \textbf{Ancestry} &  & \\
    \hline
    \textbf{Source} &  & \\
    \hline
    \textbf{Enabled By} &  & \\
    \hline
    \textbf{Supports} &  & \\
    \hline
    \textbf{Derived From} &  & \\
    \hline
    \textbf{Reviewed Status} &  & \\
    \hline
  \end{tabularx}
\end{table}

\begin{table}[htbp]
  \centering
  \footnotesize
  \begin{tabularx}{1\textwidth}{|l|l|X|}
    \hline
    \multicolumn{3}{c}{\textbf{Contained Features}} \\
    \hline
    \textbf{Feature Name} & \textbf{Feature Id} & \textbf{Comment} \\
    \hline
    Directory Contact Management & \texttt{cap:bec-d6fa9fdf} (at \ref{cap:bec-d6fa9fdf}) & \\
    \hline
    Identity Management & \texttt{cap:bec-3fd95735} (at \ref{cap:bec-3fd95735}) & \\
    \hline
    Identity Correlation & \texttt{cap:bec-737c2f3e} (at \ref{cap:bec-737c2f3e}) & \\
    \hline
    Organisational Relationship Management & \texttt{cap:bec-efb7f137} (at \ref{cap:bec-efb7f137}) & \\
    \hline
    Organisational Structure Management & \texttt{cap:bec-272efc75} (at \ref{cap:bec-272efc75}) & \\
    \hline
    Organisation Discovery & \texttt{cap:bec-73ef81cd} (at \ref{cap:bec-73ef81cd}) & \\
    \hline
  \end{tabularx}
\end{table}

Intent. Organisation Contact is a candidate enabling capability within Organisation Administration. It frames the reusable system-supported behaviours needed to support this healthcare context without prescribing a product, interface, technology or deployment boundary.

Boundary. This capability supports the business activity and does not inherit the authority or accountability of the Business Capability it enables. Its child Features are review candidates until ancestry and semantics are adjudicated.

FHIR R5 conceptual analogues considered (not implementation mappings): Organization, OrganizationAffiliation.

\subsubsection{Feature - Directory Contact Management}\label{cap:bec-d6fa9fdf}
\vspace{-1em}
\begin{table}[htbp]
  \centering
  \footnotesize
  \begin{tabularx}{1\textwidth}{|l|l|X|}
    \hline
    \textbf{Metadata Name} & \textbf{Metadata Value} & \textbf{Description} \\
    \hline
    \textbf{Reference Id} & \texttt{BEC-D6FA9FDF} & \\
    \hline
    \textbf{Status} & Existing Feature; candidate CT3 placement & \\
    \hline
    \textbf{Ancestry} &  & \\
    \hline
    \textbf{Source} & FEAT-EM-20 & \\
    \hline
    \textbf{Enabled By} &  & \\
    \hline
    \textbf{Supports} &  & \\
    \hline
    \textbf{Derived From} &  & \\
    \hline
    \textbf{Reviewed Status} & Confirm existing Feature meaning and proposed CT3 ownership. & \\
    \hline
  \end{tabularx}
\end{table}

Intent. Directory Contact Management provides a distinct enabling behaviour in the context of Organisation Contact. This is an existing named Feature; its placement under the candidate CT3 remains subject to review.

\noindent\textbf{Feature boundary.} This behaviour is scoped to Organisation Contact; an identically named Feature elsewhere may carry different business meaning.\par
\featuretable{
Function & \texttt{FN-8B2BFE12} & Maintain Organisation Contacts & Catalogue \\
Service & \texttt{SV-207AC507} & Organisation Contact Lookup & Linked to preceding Function in source \\
}
\subsubsection{Feature - Identity Management}\label{cap:bec-3fd95735}
\vspace{-1em}
\begin{table}[htbp]
  \centering
  \footnotesize
  \begin{tabularx}{1\textwidth}{|l|l|X|}
    \hline
    \textbf{Metadata Name} & \textbf{Metadata Value} & \textbf{Description} \\
    \hline
    \textbf{Reference Id} & \texttt{BEC-3FD95735} & \\
    \hline
    \textbf{Status} & Generated candidate & \\
    \hline
    \textbf{Ancestry} &  & \\
    \hline
    \textbf{Source} & FHIR-informed candidate & \\
    \hline
    \textbf{Enabled By} &  & \\
    \hline
    \textbf{Supports} &  & \\
    \hline
    \textbf{Derived From} &  & \\
    \hline
    \textbf{Reviewed Status} & Generated candidate: keep / modify / relocate / reject; verify whether independently testable. & \\
    \hline
  \end{tabularx}
\end{table}

Intent. Identity Management is proposed to support the underlying healthcare concept represented by Organisation Contact. The proposal concerns behaviour and semantic responsibility, not a FHIR field, operation, resource instance or prescribed data model.

Conceptual evidence: HL7 FHIR R5 Organization, OrganizationAffiliation.

\noindent\textbf{Feature boundary.} This behaviour is scoped to Organisation Contact; an identically named Feature elsewhere may carry different business meaning.\par
\featuretable{
--- & --- & \emph{Not yet captured} & \emph{Complete during review} \\
}
\subsubsection{Feature - Identity Correlation}\label{cap:bec-737c2f3e}
\vspace{-1em}
\begin{table}[htbp]
  \centering
  \footnotesize
  \begin{tabularx}{1\textwidth}{|l|l|X|}
    \hline
    \textbf{Metadata Name} & \textbf{Metadata Value} & \textbf{Description} \\
    \hline
    \textbf{Reference Id} & \texttt{BEC-737C2F3E} & \\
    \hline
    \textbf{Status} & Generated candidate & \\
    \hline
    \textbf{Ancestry} &  & \\
    \hline
    \textbf{Source} & FHIR-informed candidate & \\
    \hline
    \textbf{Enabled By} &  & \\
    \hline
    \textbf{Supports} &  & \\
    \hline
    \textbf{Derived From} &  & \\
    \hline
    \textbf{Reviewed Status} & Generated candidate: keep / modify / relocate / reject; verify whether independently testable. & \\
    \hline
  \end{tabularx}
\end{table}

Intent. Identity Correlation is proposed to support the underlying healthcare concept represented by Organisation Contact. The proposal concerns behaviour and semantic responsibility, not a FHIR field, operation, resource instance or prescribed data model.

Conceptual evidence: HL7 FHIR R5 Organization, OrganizationAffiliation.

\noindent\textbf{Feature boundary.} This behaviour is scoped to Organisation Contact; an identically named Feature elsewhere may carry different business meaning.\par
\featuretable{
--- & --- & \emph{Not yet captured} & \emph{Complete during review} \\
}
\subsubsection{Feature - Organisational Relationship Management}\label{cap:bec-efb7f137}
\vspace{-1em}
\begin{table}[htbp]
  \centering
  \footnotesize
  \begin{tabularx}{1\textwidth}{|l|l|X|}
    \hline
    \textbf{Metadata Name} & \textbf{Metadata Value} & \textbf{Description} \\
    \hline
    \textbf{Reference Id} & \texttt{BEC-EFB7F137} & \\
    \hline
    \textbf{Status} & Generated candidate & \\
    \hline
    \textbf{Ancestry} &  & \\
    \hline
    \textbf{Source} & FHIR-informed candidate & \\
    \hline
    \textbf{Enabled By} &  & \\
    \hline
    \textbf{Supports} &  & \\
    \hline
    \textbf{Derived From} &  & \\
    \hline
    \textbf{Reviewed Status} & Generated candidate: keep / modify / relocate / reject; verify whether independently testable. & \\
    \hline
  \end{tabularx}
\end{table}

Intent. Organisational Relationship Management is proposed to support the underlying healthcare concept represented by Organisation Contact. The proposal concerns behaviour and semantic responsibility, not a FHIR field, operation, resource instance or prescribed data model.

Conceptual evidence: HL7 FHIR R5 Organization, OrganizationAffiliation.

\noindent\textbf{Feature boundary.} This behaviour is scoped to Organisation Contact; an identically named Feature elsewhere may carry different business meaning.\par
\featuretable{
--- & --- & \emph{Not yet captured} & \emph{Complete during review} \\
}
\subsubsection{Feature - Organisational Structure Management}\label{cap:bec-272efc75}
\vspace{-1em}
\begin{table}[htbp]
  \centering
  \footnotesize
  \begin{tabularx}{1\textwidth}{|l|l|X|}
    \hline
    \textbf{Metadata Name} & \textbf{Metadata Value} & \textbf{Description} \\
    \hline
    \textbf{Reference Id} & \texttt{BEC-272EFC75} & \\
    \hline
    \textbf{Status} & Generated candidate & \\
    \hline
    \textbf{Ancestry} &  & \\
    \hline
    \textbf{Source} & FHIR-informed candidate & \\
    \hline
    \textbf{Enabled By} &  & \\
    \hline
    \textbf{Supports} &  & \\
    \hline
    \textbf{Derived From} &  & \\
    \hline
    \textbf{Reviewed Status} & Generated candidate: keep / modify / relocate / reject; verify whether independently testable. & \\
    \hline
  \end{tabularx}
\end{table}

Intent. Organisational Structure Management is proposed to support the underlying healthcare concept represented by Organisation Contact. The proposal concerns behaviour and semantic responsibility, not a FHIR field, operation, resource instance or prescribed data model.

Conceptual evidence: HL7 FHIR R5 Organization, OrganizationAffiliation.

\noindent\textbf{Feature boundary.} This behaviour is scoped to Organisation Contact; an identically named Feature elsewhere may carry different business meaning.\par
\featuretable{
--- & --- & \emph{Not yet captured} & \emph{Complete during review} \\
}
\subsubsection{Feature - Organisation Discovery}\label{cap:bec-73ef81cd}
\vspace{-1em}
\begin{table}[htbp]
  \centering
  \footnotesize
  \begin{tabularx}{1\textwidth}{|l|l|X|}
    \hline
    \textbf{Metadata Name} & \textbf{Metadata Value} & \textbf{Description} \\
    \hline
    \textbf{Reference Id} & \texttt{BEC-73EF81CD} & \\
    \hline
    \textbf{Status} & Generated candidate & \\
    \hline
    \textbf{Ancestry} &  & \\
    \hline
    \textbf{Source} & FHIR-informed candidate & \\
    \hline
    \textbf{Enabled By} &  & \\
    \hline
    \textbf{Supports} &  & \\
    \hline
    \textbf{Derived From} &  & \\
    \hline
    \textbf{Reviewed Status} & Generated candidate: keep / modify / relocate / reject; verify whether independently testable. & \\
    \hline
  \end{tabularx}
\end{table}

Intent. Organisation Discovery is proposed to support the underlying healthcare concept represented by Organisation Contact. The proposal concerns behaviour and semantic responsibility, not a FHIR field, operation, resource instance or prescribed data model.

Conceptual evidence: HL7 FHIR R5 Organization, OrganizationAffiliation.

\noindent\textbf{Feature boundary.} This behaviour is scoped to Organisation Contact; an identically named Feature elsewhere may carry different business meaning.\par
\featuretable{
--- & --- & \emph{Not yet captured} & \emph{Complete during review} \\
}

\section{CT2 - Location Administration}\label{cap:bec-d4b0d057}
\vspace{-0.5em}
\begin{table}[htbp]
  \centering
  \footnotesize
  \begin{tabularx}{1\textwidth}{|l|l|X|}
    \hline
    \multicolumn{3}{c}{\textbf{Capability Metadata}} \\
    \hline
    \textbf{Name} & \textbf{Value} & \textbf{Comment and/or Description} \\
    \hline
    \textbf{Reference Id} & \texttt{BEC-D4B0D057} & \\
    \hline
    \textbf{Status} & Candidate ancestry & \\
    \hline
    \textbf{Ancestry} &  & \\
    \hline
    \textbf{Source} &  & \\
    \hline
    \textbf{Enabled By} &  & \\
    \hline
    \textbf{Supports} &  & \\
    \hline
    \textbf{Derived From} &  & \\
    \hline
    \textbf{Reviewed Status} &  & \\
    \hline
  \end{tabularx}
\end{table}

\begin{table}[htbp]
  \centering
  \footnotesize
  \begin{tabularx}{1\textwidth}{|l|l|X|}
    \hline
    \multicolumn{3}{c}{\textbf{Contained Capabilities}} \\
    \hline
    \textbf{Capability Name} & \textbf{Capability Id} & \textbf{Comment} \\
    \hline
    Healthcare Location & \texttt{cap:bec-7e83f611} (at \ref{cap:bec-7e83f611}) & \\
    \hline
    Location Structure & \texttt{cap:bec-a66483f0} (at \ref{cap:bec-a66483f0}) & \\
    \hline
    Care-Place Definition & \texttt{cap:bec-0c8d2717} (at \ref{cap:bec-0c8d2717}) & \\
    \hline
  \end{tabularx}
\end{table}

Intent. Location Administration is a candidate enabling capability within Entity Information Enablement. It frames the reusable system-supported behaviours needed to support this healthcare context without prescribing a product, interface, technology or deployment boundary.

Boundary. This capability supports the business activity and does not inherit the authority or accountability of the Business Capability it enables. Its child Features are review candidates until ancestry and semantics are adjudicated.

\subsection{CT3 - Healthcare Location}\label{cap:bec-7e83f611}
\vspace{-0.5em}
\begin{table}[htbp]
  \centering
  \footnotesize
  \begin{tabularx}{1\textwidth}{|l|l|X|}
    \hline
    \multicolumn{3}{c}{\textbf{Capability Metadata}} \\
    \hline
    \textbf{Name} & \textbf{Value} & \textbf{Comment and/or Description} \\
    \hline
    \textbf{Reference Id} & \texttt{BEC-7E83F611} & \\
    \hline
    \textbf{Status} & Candidate ancestry & \\
    \hline
    \textbf{Ancestry} &  & \\
    \hline
    \textbf{Source} &  & \\
    \hline
    \textbf{Enabled By} &  & \\
    \hline
    \textbf{Supports} &  & \\
    \hline
    \textbf{Derived From} &  & \\
    \hline
    \textbf{Reviewed Status} &  & \\
    \hline
  \end{tabularx}
\end{table}

\begin{table}[htbp]
  \centering
  \footnotesize
  \begin{tabularx}{1\textwidth}{|l|l|X|}
    \hline
    \multicolumn{3}{c}{\textbf{Contained Features}} \\
    \hline
    \textbf{Feature Name} & \textbf{Feature Id} & \textbf{Comment} \\
    \hline
    Location Geospatial \& Address Resolution & \texttt{cap:bec-885ce86b} (at \ref{cap:bec-885ce86b}) & \\
    \hline
    Location Identity Management & \texttt{cap:bec-15887e90} (at \ref{cap:bec-15887e90}) & \\
    \hline
    Location Hierarchy Management & \texttt{cap:bec-c52417ee} (at \ref{cap:bec-c52417ee}) & \\
    \hline
    Location Availability Management & \texttt{cap:bec-0538c23a} (at \ref{cap:bec-0538c23a}) & \\
    \hline
    Location Discovery & \texttt{cap:bec-be1b64cd} (at \ref{cap:bec-be1b64cd}) & \\
    \hline
  \end{tabularx}
\end{table}

Intent. Healthcare Location is a candidate enabling capability within Location Administration. It frames the reusable system-supported behaviours needed to support this healthcare context without prescribing a product, interface, technology or deployment boundary.

Boundary. This capability supports the business activity and does not inherit the authority or accountability of the Business Capability it enables. Its child Features are review candidates until ancestry and semantics are adjudicated.

FHIR R5 conceptual analogues considered (not implementation mappings): Location.

\subsubsection{Feature - Location Geospatial \& Address Resolution}\label{cap:bec-885ce86b}
\vspace{-1em}
\begin{table}[htbp]
  \centering
  \footnotesize
  \begin{tabularx}{1\textwidth}{|l|l|X|}
    \hline
    \textbf{Metadata Name} & \textbf{Metadata Value} & \textbf{Description} \\
    \hline
    \textbf{Reference Id} & \texttt{BEC-885CE86B} & \\
    \hline
    \textbf{Status} & Existing Feature; candidate CT3 placement & \\
    \hline
    \textbf{Ancestry} &  & \\
    \hline
    \textbf{Source} & FEAT-EM-21 & \\
    \hline
    \textbf{Enabled By} &  & \\
    \hline
    \textbf{Supports} &  & \\
    \hline
    \textbf{Derived From} &  & \\
    \hline
    \textbf{Reviewed Status} & Confirm existing Feature meaning and proposed CT3 ownership. & \\
    \hline
  \end{tabularx}
\end{table}

Intent. Location Geospatial \& Address Resolution provides a distinct enabling behaviour in the context of Healthcare Location. This is an existing named Feature; its placement under the candidate CT3 remains subject to review.

\noindent\textbf{Feature boundary.} This behaviour is scoped to Healthcare Location; an identically named Feature elsewhere may carry different business meaning.\par
\featuretable{
Function & \texttt{FN-E85A06B2} & Maintain Healthcare Location & Catalogue \\
Service & \texttt{SV-74104935} & Location Directory Query & Linked to preceding Function in source \\
}
\subsubsection{Feature - Location Identity Management}\label{cap:bec-15887e90}
\vspace{-1em}
\begin{table}[htbp]
  \centering
  \footnotesize
  \begin{tabularx}{1\textwidth}{|l|l|X|}
    \hline
    \textbf{Metadata Name} & \textbf{Metadata Value} & \textbf{Description} \\
    \hline
    \textbf{Reference Id} & \texttt{BEC-15887E90} & \\
    \hline
    \textbf{Status} & Generated candidate & \\
    \hline
    \textbf{Ancestry} &  & \\
    \hline
    \textbf{Source} & FHIR-informed candidate & \\
    \hline
    \textbf{Enabled By} &  & \\
    \hline
    \textbf{Supports} &  & \\
    \hline
    \textbf{Derived From} &  & \\
    \hline
    \textbf{Reviewed Status} & Generated candidate: keep / modify / relocate / reject; verify whether independently testable. & \\
    \hline
  \end{tabularx}
\end{table}

Intent. Location Identity Management is proposed to support the underlying healthcare concept represented by Healthcare Location. The proposal concerns behaviour and semantic responsibility, not a FHIR field, operation, resource instance or prescribed data model.

Conceptual evidence: HL7 FHIR R5 Location.

\noindent\textbf{Feature boundary.} This behaviour is scoped to Healthcare Location; an identically named Feature elsewhere may carry different business meaning.\par
\featuretable{
--- & --- & \emph{Not yet captured} & \emph{Complete during review} \\
}
\subsubsection{Feature - Location Hierarchy Management}\label{cap:bec-c52417ee}
\vspace{-1em}
\begin{table}[htbp]
  \centering
  \footnotesize
  \begin{tabularx}{1\textwidth}{|l|l|X|}
    \hline
    \textbf{Metadata Name} & \textbf{Metadata Value} & \textbf{Description} \\
    \hline
    \textbf{Reference Id} & \texttt{BEC-C52417EE} & \\
    \hline
    \textbf{Status} & Generated candidate & \\
    \hline
    \textbf{Ancestry} &  & \\
    \hline
    \textbf{Source} & FHIR-informed candidate & \\
    \hline
    \textbf{Enabled By} &  & \\
    \hline
    \textbf{Supports} &  & \\
    \hline
    \textbf{Derived From} &  & \\
    \hline
    \textbf{Reviewed Status} & Generated candidate: keep / modify / relocate / reject; verify whether independently testable. & \\
    \hline
  \end{tabularx}
\end{table}

Intent. Location Hierarchy Management is proposed to support the underlying healthcare concept represented by Healthcare Location. The proposal concerns behaviour and semantic responsibility, not a FHIR field, operation, resource instance or prescribed data model.

Conceptual evidence: HL7 FHIR R5 Location.

\noindent\textbf{Feature boundary.} This behaviour is scoped to Healthcare Location; an identically named Feature elsewhere may carry different business meaning.\par
\featuretable{
--- & --- & \emph{Not yet captured} & \emph{Complete during review} \\
}
\subsubsection{Feature - Location Availability Management}\label{cap:bec-0538c23a}
\vspace{-1em}
\begin{table}[htbp]
  \centering
  \footnotesize
  \begin{tabularx}{1\textwidth}{|l|l|X|}
    \hline
    \textbf{Metadata Name} & \textbf{Metadata Value} & \textbf{Description} \\
    \hline
    \textbf{Reference Id} & \texttt{BEC-0538C23A} & \\
    \hline
    \textbf{Status} & Generated candidate & \\
    \hline
    \textbf{Ancestry} &  & \\
    \hline
    \textbf{Source} & FHIR-informed candidate & \\
    \hline
    \textbf{Enabled By} &  & \\
    \hline
    \textbf{Supports} &  & \\
    \hline
    \textbf{Derived From} &  & \\
    \hline
    \textbf{Reviewed Status} & Generated candidate: keep / modify / relocate / reject; verify whether independently testable. & \\
    \hline
  \end{tabularx}
\end{table}

Intent. Location Availability Management is proposed to support the underlying healthcare concept represented by Healthcare Location. The proposal concerns behaviour and semantic responsibility, not a FHIR field, operation, resource instance or prescribed data model.

Conceptual evidence: HL7 FHIR R5 Location.

\noindent\textbf{Feature boundary.} This behaviour is scoped to Healthcare Location; an identically named Feature elsewhere may carry different business meaning.\par
\featuretable{
--- & --- & \emph{Not yet captured} & \emph{Complete during review} \\
}
\subsubsection{Feature - Location Discovery}\label{cap:bec-be1b64cd}
\vspace{-1em}
\begin{table}[htbp]
  \centering
  \footnotesize
  \begin{tabularx}{1\textwidth}{|l|l|X|}
    \hline
    \textbf{Metadata Name} & \textbf{Metadata Value} & \textbf{Description} \\
    \hline
    \textbf{Reference Id} & \texttt{BEC-BE1B64CD} & \\
    \hline
    \textbf{Status} & Generated candidate & \\
    \hline
    \textbf{Ancestry} &  & \\
    \hline
    \textbf{Source} & FHIR-informed candidate & \\
    \hline
    \textbf{Enabled By} &  & \\
    \hline
    \textbf{Supports} &  & \\
    \hline
    \textbf{Derived From} &  & \\
    \hline
    \textbf{Reviewed Status} & Generated candidate: keep / modify / relocate / reject; verify whether independently testable. & \\
    \hline
  \end{tabularx}
\end{table}

Intent. Location Discovery is proposed to support the underlying healthcare concept represented by Healthcare Location. The proposal concerns behaviour and semantic responsibility, not a FHIR field, operation, resource instance or prescribed data model.

Conceptual evidence: HL7 FHIR R5 Location.

\noindent\textbf{Feature boundary.} This behaviour is scoped to Healthcare Location; an identically named Feature elsewhere may carry different business meaning.\par
\featuretable{
--- & --- & \emph{Not yet captured} & \emph{Complete during review} \\
}

\subsection{CT3 - Location Structure}\label{cap:bec-a66483f0}
\vspace{-0.5em}
\begin{table}[htbp]
  \centering
  \footnotesize
  \begin{tabularx}{1\textwidth}{|l|l|X|}
    \hline
    \multicolumn{3}{c}{\textbf{Capability Metadata}} \\
    \hline
    \textbf{Name} & \textbf{Value} & \textbf{Comment and/or Description} \\
    \hline
    \textbf{Reference Id} & \texttt{BEC-A66483F0} & \\
    \hline
    \textbf{Status} & Candidate ancestry & \\
    \hline
    \textbf{Ancestry} &  & \\
    \hline
    \textbf{Source} &  & \\
    \hline
    \textbf{Enabled By} &  & \\
    \hline
    \textbf{Supports} &  & \\
    \hline
    \textbf{Derived From} &  & \\
    \hline
    \textbf{Reviewed Status} &  & \\
    \hline
  \end{tabularx}
\end{table}

\begin{table}[htbp]
  \centering
  \footnotesize
  \begin{tabularx}{1\textwidth}{|l|l|X|}
    \hline
    \multicolumn{3}{c}{\textbf{Contained Features}} \\
    \hline
    \textbf{Feature Name} & \textbf{Feature Id} & \textbf{Comment} \\
    \hline
    Physical Hierarchy Modeling & \texttt{cap:bec-c82c8d0a} (at \ref{cap:bec-c82c8d0a}) & \\
    \hline
    Location Identity Management & \texttt{cap:bec-fc2dd94f} (at \ref{cap:bec-fc2dd94f}) & \\
    \hline
    Location Hierarchy Management & \texttt{cap:bec-6dcc28f3} (at \ref{cap:bec-6dcc28f3}) & \\
    \hline
    Location Availability Management & \texttt{cap:bec-310ff956} (at \ref{cap:bec-310ff956}) & \\
    \hline
    Location Discovery & \texttt{cap:bec-4bcd6b9c} (at \ref{cap:bec-4bcd6b9c}) & \\
    \hline
  \end{tabularx}
\end{table}

Intent. Location Structure is a candidate enabling capability within Location Administration. It frames the reusable system-supported behaviours needed to support this healthcare context without prescribing a product, interface, technology or deployment boundary.

Boundary. This capability supports the business activity and does not inherit the authority or accountability of the Business Capability it enables. Its child Features are review candidates until ancestry and semantics are adjudicated.

FHIR R5 conceptual analogues considered (not implementation mappings): Location.

\subsubsection{Feature - Physical Hierarchy Modeling}\label{cap:bec-c82c8d0a}
\vspace{-1em}
\begin{table}[htbp]
  \centering
  \footnotesize
  \begin{tabularx}{1\textwidth}{|l|l|X|}
    \hline
    \textbf{Metadata Name} & \textbf{Metadata Value} & \textbf{Description} \\
    \hline
    \textbf{Reference Id} & \texttt{BEC-C82C8D0A} & \\
    \hline
    \textbf{Status} & Existing Feature; candidate CT3 placement & \\
    \hline
    \textbf{Ancestry} &  & \\
    \hline
    \textbf{Source} & FEAT-EM-22 & \\
    \hline
    \textbf{Enabled By} &  & \\
    \hline
    \textbf{Supports} &  & \\
    \hline
    \textbf{Derived From} &  & \\
    \hline
    \textbf{Reviewed Status} & Confirm existing Feature meaning and proposed CT3 ownership. & \\
    \hline
  \end{tabularx}
\end{table}

Intent. Physical Hierarchy Modeling provides a distinct enabling behaviour in the context of Location Structure. This is an existing named Feature; its placement under the candidate CT3 remains subject to review.

\noindent\textbf{Feature boundary.} This behaviour is scoped to Location Structure; an identically named Feature elsewhere may carry different business meaning.\par
\featuretable{
Function & \texttt{FN-A464CD8F} & Maintain Physical Location Hierarchy & Catalogue \\
Service & \texttt{SV-2EF94929} & Location Hierarchy Resolution & Linked to preceding Function in source \\
}
\subsubsection{Feature - Location Identity Management}\label{cap:bec-fc2dd94f}
\vspace{-1em}
\begin{table}[htbp]
  \centering
  \footnotesize
  \begin{tabularx}{1\textwidth}{|l|l|X|}
    \hline
    \textbf{Metadata Name} & \textbf{Metadata Value} & \textbf{Description} \\
    \hline
    \textbf{Reference Id} & \texttt{BEC-FC2DD94F} & \\
    \hline
    \textbf{Status} & Generated candidate & \\
    \hline
    \textbf{Ancestry} &  & \\
    \hline
    \textbf{Source} & FHIR-informed candidate & \\
    \hline
    \textbf{Enabled By} &  & \\
    \hline
    \textbf{Supports} &  & \\
    \hline
    \textbf{Derived From} &  & \\
    \hline
    \textbf{Reviewed Status} & Generated candidate: keep / modify / relocate / reject; verify whether independently testable. & \\
    \hline
  \end{tabularx}
\end{table}

Intent. Location Identity Management is proposed to support the underlying healthcare concept represented by Location Structure. The proposal concerns behaviour and semantic responsibility, not a FHIR field, operation, resource instance or prescribed data model.

Conceptual evidence: HL7 FHIR R5 Location.

\noindent\textbf{Feature boundary.} This behaviour is scoped to Location Structure; an identically named Feature elsewhere may carry different business meaning.\par
\featuretable{
--- & --- & \emph{Not yet captured} & \emph{Complete during review} \\
}
\subsubsection{Feature - Location Hierarchy Management}\label{cap:bec-6dcc28f3}
\vspace{-1em}
\begin{table}[htbp]
  \centering
  \footnotesize
  \begin{tabularx}{1\textwidth}{|l|l|X|}
    \hline
    \textbf{Metadata Name} & \textbf{Metadata Value} & \textbf{Description} \\
    \hline
    \textbf{Reference Id} & \texttt{BEC-6DCC28F3} & \\
    \hline
    \textbf{Status} & Generated candidate & \\
    \hline
    \textbf{Ancestry} &  & \\
    \hline
    \textbf{Source} & FHIR-informed candidate & \\
    \hline
    \textbf{Enabled By} &  & \\
    \hline
    \textbf{Supports} &  & \\
    \hline
    \textbf{Derived From} &  & \\
    \hline
    \textbf{Reviewed Status} & Generated candidate: keep / modify / relocate / reject; verify whether independently testable. & \\
    \hline
  \end{tabularx}
\end{table}

Intent. Location Hierarchy Management is proposed to support the underlying healthcare concept represented by Location Structure. The proposal concerns behaviour and semantic responsibility, not a FHIR field, operation, resource instance or prescribed data model.

Conceptual evidence: HL7 FHIR R5 Location.

\noindent\textbf{Feature boundary.} This behaviour is scoped to Location Structure; an identically named Feature elsewhere may carry different business meaning.\par
\featuretable{
--- & --- & \emph{Not yet captured} & \emph{Complete during review} \\
}
\subsubsection{Feature - Location Availability Management}\label{cap:bec-310ff956}
\vspace{-1em}
\begin{table}[htbp]
  \centering
  \footnotesize
  \begin{tabularx}{1\textwidth}{|l|l|X|}
    \hline
    \textbf{Metadata Name} & \textbf{Metadata Value} & \textbf{Description} \\
    \hline
    \textbf{Reference Id} & \texttt{BEC-310FF956} & \\
    \hline
    \textbf{Status} & Generated candidate & \\
    \hline
    \textbf{Ancestry} &  & \\
    \hline
    \textbf{Source} & FHIR-informed candidate & \\
    \hline
    \textbf{Enabled By} &  & \\
    \hline
    \textbf{Supports} &  & \\
    \hline
    \textbf{Derived From} &  & \\
    \hline
    \textbf{Reviewed Status} & Generated candidate: keep / modify / relocate / reject; verify whether independently testable. & \\
    \hline
  \end{tabularx}
\end{table}

Intent. Location Availability Management is proposed to support the underlying healthcare concept represented by Location Structure. The proposal concerns behaviour and semantic responsibility, not a FHIR field, operation, resource instance or prescribed data model.

Conceptual evidence: HL7 FHIR R5 Location.

\noindent\textbf{Feature boundary.} This behaviour is scoped to Location Structure; an identically named Feature elsewhere may carry different business meaning.\par
\featuretable{
--- & --- & \emph{Not yet captured} & \emph{Complete during review} \\
}
\subsubsection{Feature - Location Discovery}\label{cap:bec-4bcd6b9c}
\vspace{-1em}
\begin{table}[htbp]
  \centering
  \footnotesize
  \begin{tabularx}{1\textwidth}{|l|l|X|}
    \hline
    \textbf{Metadata Name} & \textbf{Metadata Value} & \textbf{Description} \\
    \hline
    \textbf{Reference Id} & \texttt{BEC-4BCD6B9C} & \\
    \hline
    \textbf{Status} & Generated candidate & \\
    \hline
    \textbf{Ancestry} &  & \\
    \hline
    \textbf{Source} & FHIR-informed candidate & \\
    \hline
    \textbf{Enabled By} &  & \\
    \hline
    \textbf{Supports} &  & \\
    \hline
    \textbf{Derived From} &  & \\
    \hline
    \textbf{Reviewed Status} & Generated candidate: keep / modify / relocate / reject; verify whether independently testable. & \\
    \hline
  \end{tabularx}
\end{table}

Intent. Location Discovery is proposed to support the underlying healthcare concept represented by Location Structure. The proposal concerns behaviour and semantic responsibility, not a FHIR field, operation, resource instance or prescribed data model.

Conceptual evidence: HL7 FHIR R5 Location.

\noindent\textbf{Feature boundary.} This behaviour is scoped to Location Structure; an identically named Feature elsewhere may carry different business meaning.\par
\featuretable{
--- & --- & \emph{Not yet captured} & \emph{Complete during review} \\
}

\subsection{CT3 - Care-Place Definition}\label{cap:bec-0c8d2717}
\vspace{-0.5em}
\begin{table}[htbp]
  \centering
  \footnotesize
  \begin{tabularx}{1\textwidth}{|l|l|X|}
    \hline
    \multicolumn{3}{c}{\textbf{Capability Metadata}} \\
    \hline
    \textbf{Name} & \textbf{Value} & \textbf{Comment and/or Description} \\
    \hline
    \textbf{Reference Id} & \texttt{BEC-0C8D2717} & \\
    \hline
    \textbf{Status} & Candidate ancestry & \\
    \hline
    \textbf{Ancestry} &  & \\
    \hline
    \textbf{Source} &  & \\
    \hline
    \textbf{Enabled By} &  & \\
    \hline
    \textbf{Supports} &  & \\
    \hline
    \textbf{Derived From} &  & \\
    \hline
    \textbf{Reviewed Status} &  & \\
    \hline
  \end{tabularx}
\end{table}

\begin{table}[htbp]
  \centering
  \footnotesize
  \begin{tabularx}{1\textwidth}{|l|l|X|}
    \hline
    \multicolumn{3}{c}{\textbf{Contained Features}} \\
    \hline
    \textbf{Feature Name} & \textbf{Feature Id} & \textbf{Comment} \\
    \hline
    Atomic Care-Place Modeling & \texttt{cap:bec-6322e089} (at \ref{cap:bec-6322e089}) & \\
    \hline
    Context Establishment & \texttt{cap:bec-705cd00c} (at \ref{cap:bec-705cd00c}) & \\
    \hline
    State Monitoring & \texttt{cap:bec-67dff33a} (at \ref{cap:bec-67dff33a}) & \\
    \hline
    Exception Management & \texttt{cap:bec-e79afe62} (at \ref{cap:bec-e79afe62}) & \\
    \hline
    Information Retrieval & \texttt{cap:bec-4b7196d5} (at \ref{cap:bec-4b7196d5}) & \\
    \hline
  \end{tabularx}
\end{table}

Intent. Care-Place Definition is a candidate enabling capability within Location Administration. It frames the reusable system-supported behaviours needed to support this healthcare context without prescribing a product, interface, technology or deployment boundary.

Boundary. This capability supports the business activity and does not inherit the authority or accountability of the Business Capability it enables. Its child Features are review candidates until ancestry and semantics are adjudicated.

FHIR R5 conceptual analogues considered (not implementation mappings): Location, Encounter.

\subsubsection{Feature - Atomic Care-Place Modeling}\label{cap:bec-6322e089}
\vspace{-1em}
\begin{table}[htbp]
  \centering
  \footnotesize
  \begin{tabularx}{1\textwidth}{|l|l|X|}
    \hline
    \textbf{Metadata Name} & \textbf{Metadata Value} & \textbf{Description} \\
    \hline
    \textbf{Reference Id} & \texttt{BEC-6322E089} & \\
    \hline
    \textbf{Status} & Existing Feature; candidate CT3 placement & \\
    \hline
    \textbf{Ancestry} &  & \\
    \hline
    \textbf{Source} & FEAT-EM-23 & \\
    \hline
    \textbf{Enabled By} &  & \\
    \hline
    \textbf{Supports} &  & \\
    \hline
    \textbf{Derived From} &  & \\
    \hline
    \textbf{Reviewed Status} & Confirm existing Feature meaning and proposed CT3 ownership. & \\
    \hline
  \end{tabularx}
\end{table}

Intent. Atomic Care-Place Modeling provides a distinct enabling behaviour in the context of Care-Place Definition. This is an existing named Feature; its placement under the candidate CT3 remains subject to review.

\noindent\textbf{Feature boundary.} This behaviour is scoped to Care-Place Definition; an identically named Feature elsewhere may carry different business meaning.\par
\featuretable{
Function & \texttt{FN-B14AEA53} & Maintain Care-Place Definition & Catalogue \\
Service & \texttt{SV-8CFC2D44} & Care-Place Specification Lookup & Linked to preceding Function in source \\
}
\subsubsection{Feature - Context Establishment}\label{cap:bec-705cd00c}
\vspace{-1em}
\begin{table}[htbp]
  \centering
  \footnotesize
  \begin{tabularx}{1\textwidth}{|l|l|X|}
    \hline
    \textbf{Metadata Name} & \textbf{Metadata Value} & \textbf{Description} \\
    \hline
    \textbf{Reference Id} & \texttt{BEC-705CD00C} & \\
    \hline
    \textbf{Status} & Generated candidate & \\
    \hline
    \textbf{Ancestry} &  & \\
    \hline
    \textbf{Source} & Semantic-pattern candidate & \\
    \hline
    \textbf{Enabled By} &  & \\
    \hline
    \textbf{Supports} &  & \\
    \hline
    \textbf{Derived From} &  & \\
    \hline
    \textbf{Reviewed Status} & Generated candidate: keep / modify / relocate / reject; verify whether independently testable. & \\
    \hline
  \end{tabularx}
\end{table}

Intent. Context Establishment is proposed as a discrete, testable enabling behaviour for Care-Place Definition. It is included for architectural review, not as an established implementation obligation.

\noindent\textbf{Feature boundary.} This behaviour is scoped to Care-Place Definition; an identically named Feature elsewhere may carry different business meaning.\par
\featuretable{
--- & --- & \emph{Not yet captured} & \emph{Complete during review} \\
}
\subsubsection{Feature - State Monitoring}\label{cap:bec-67dff33a}
\vspace{-1em}
\begin{table}[htbp]
  \centering
  \footnotesize
  \begin{tabularx}{1\textwidth}{|l|l|X|}
    \hline
    \textbf{Metadata Name} & \textbf{Metadata Value} & \textbf{Description} \\
    \hline
    \textbf{Reference Id} & \texttt{BEC-67DFF33A} & \\
    \hline
    \textbf{Status} & Generated candidate & \\
    \hline
    \textbf{Ancestry} &  & \\
    \hline
    \textbf{Source} & Semantic-pattern candidate & \\
    \hline
    \textbf{Enabled By} &  & \\
    \hline
    \textbf{Supports} &  & \\
    \hline
    \textbf{Derived From} &  & \\
    \hline
    \textbf{Reviewed Status} & Generated candidate: keep / modify / relocate / reject; verify whether independently testable. & \\
    \hline
  \end{tabularx}
\end{table}

Intent. State Monitoring is proposed as a discrete, testable enabling behaviour for Care-Place Definition. It is included for architectural review, not as an established implementation obligation.

\noindent\textbf{Feature boundary.} This behaviour is scoped to Care-Place Definition; an identically named Feature elsewhere may carry different business meaning.\par
\featuretable{
--- & --- & \emph{Not yet captured} & \emph{Complete during review} \\
}
\subsubsection{Feature - Exception Management}\label{cap:bec-e79afe62}
\vspace{-1em}
\begin{table}[htbp]
  \centering
  \footnotesize
  \begin{tabularx}{1\textwidth}{|l|l|X|}
    \hline
    \textbf{Metadata Name} & \textbf{Metadata Value} & \textbf{Description} \\
    \hline
    \textbf{Reference Id} & \texttt{BEC-E79AFE62} & \\
    \hline
    \textbf{Status} & Generated candidate & \\
    \hline
    \textbf{Ancestry} &  & \\
    \hline
    \textbf{Source} & Semantic-pattern candidate & \\
    \hline
    \textbf{Enabled By} &  & \\
    \hline
    \textbf{Supports} &  & \\
    \hline
    \textbf{Derived From} &  & \\
    \hline
    \textbf{Reviewed Status} & Generated candidate: keep / modify / relocate / reject; verify whether independently testable. & \\
    \hline
  \end{tabularx}
\end{table}

Intent. Exception Management is proposed as a discrete, testable enabling behaviour for Care-Place Definition. It is included for architectural review, not as an established implementation obligation.

\noindent\textbf{Feature boundary.} This behaviour is scoped to Care-Place Definition; an identically named Feature elsewhere may carry different business meaning.\par
\featuretable{
--- & --- & \emph{Not yet captured} & \emph{Complete during review} \\
}
\subsubsection{Feature - Information Retrieval}\label{cap:bec-4b7196d5}
\vspace{-1em}
\begin{table}[htbp]
  \centering
  \footnotesize
  \begin{tabularx}{1\textwidth}{|l|l|X|}
    \hline
    \textbf{Metadata Name} & \textbf{Metadata Value} & \textbf{Description} \\
    \hline
    \textbf{Reference Id} & \texttt{BEC-4B7196D5} & \\
    \hline
    \textbf{Status} & Generated candidate & \\
    \hline
    \textbf{Ancestry} &  & \\
    \hline
    \textbf{Source} & Semantic-pattern candidate & \\
    \hline
    \textbf{Enabled By} &  & \\
    \hline
    \textbf{Supports} &  & \\
    \hline
    \textbf{Derived From} &  & \\
    \hline
    \textbf{Reviewed Status} & Generated candidate: keep / modify / relocate / reject; verify whether independently testable. & \\
    \hline
  \end{tabularx}
\end{table}

Intent. Information Retrieval is proposed as a discrete, testable enabling behaviour for Care-Place Definition. It is included for architectural review, not as an established implementation obligation.

\noindent\textbf{Feature boundary.} This behaviour is scoped to Care-Place Definition; an identically named Feature elsewhere may carry different business meaning.\par
\featuretable{
--- & --- & \emph{Not yet captured} & \emph{Complete during review} \\
}

\section{CT2 - Health Service Administration}\label{cap:bec-b3668209}
\vspace{-0.5em}
\begin{table}[htbp]
  \centering
  \footnotesize
  \begin{tabularx}{1\textwidth}{|l|l|X|}
    \hline
    \multicolumn{3}{c}{\textbf{Capability Metadata}} \\
    \hline
    \textbf{Name} & \textbf{Value} & \textbf{Comment and/or Description} \\
    \hline
    \textbf{Reference Id} & \texttt{BEC-B3668209} & \\
    \hline
    \textbf{Status} & Candidate ancestry & \\
    \hline
    \textbf{Ancestry} &  & \\
    \hline
    \textbf{Source} &  & \\
    \hline
    \textbf{Enabled By} &  & \\
    \hline
    \textbf{Supports} &  & \\
    \hline
    \textbf{Derived From} &  & \\
    \hline
    \textbf{Reviewed Status} &  & \\
    \hline
  \end{tabularx}
\end{table}

\begin{table}[htbp]
  \centering
  \footnotesize
  \begin{tabularx}{1\textwidth}{|l|l|X|}
    \hline
    \multicolumn{3}{c}{\textbf{Contained Capabilities}} \\
    \hline
    \textbf{Capability Name} & \textbf{Capability Id} & \textbf{Comment} \\
    \hline
    Healthcare Service Definition & \texttt{cap:bec-1864d570} (at \ref{cap:bec-1864d570}) & \\
    \hline
    Service Provision Relationship & \texttt{cap:bec-cc9fd951} (at \ref{cap:bec-cc9fd951}) & \\
    \hline
    Service Availability Definition & \texttt{cap:bec-cf667cc1} (at \ref{cap:bec-cf667cc1}) & \\
    \hline
    Service Eligibility Definition & \texttt{cap:bec-8858c964} (at \ref{cap:bec-8858c964}) & \\
    \hline
  \end{tabularx}
\end{table}

Intent. Health Service Administration is a candidate enabling capability within Entity Information Enablement. It frames the reusable system-supported behaviours needed to support this healthcare context without prescribing a product, interface, technology or deployment boundary.

Boundary. This capability supports the business activity and does not inherit the authority or accountability of the Business Capability it enables. Its child Features are review candidates until ancestry and semantics are adjudicated.

\subsection{CT3 - Healthcare Service Definition}\label{cap:bec-1864d570}
\vspace{-0.5em}
\begin{table}[htbp]
  \centering
  \footnotesize
  \begin{tabularx}{1\textwidth}{|l|l|X|}
    \hline
    \multicolumn{3}{c}{\textbf{Capability Metadata}} \\
    \hline
    \textbf{Name} & \textbf{Value} & \textbf{Comment and/or Description} \\
    \hline
    \textbf{Reference Id} & \texttt{BEC-1864D570} & \\
    \hline
    \textbf{Status} & Candidate ancestry & \\
    \hline
    \textbf{Ancestry} &  & \\
    \hline
    \textbf{Source} &  & \\
    \hline
    \textbf{Enabled By} &  & \\
    \hline
    \textbf{Supports} &  & \\
    \hline
    \textbf{Derived From} &  & \\
    \hline
    \textbf{Reviewed Status} &  & \\
    \hline
  \end{tabularx}
\end{table}

\begin{table}[htbp]
  \centering
  \footnotesize
  \begin{tabularx}{1\textwidth}{|l|l|X|}
    \hline
    \multicolumn{3}{c}{\textbf{Contained Features}} \\
    \hline
    \textbf{Feature Name} & \textbf{Feature Id} & \textbf{Comment} \\
    \hline
    Service Catalogue Publishing & \texttt{cap:bec-5af74389} (at \ref{cap:bec-5af74389}) & \\
    \hline
    Service Definition Management & \texttt{cap:bec-8cdb717f} (at \ref{cap:bec-8cdb717f}) & \\
    \hline
    Service Eligibility Management & \texttt{cap:bec-766061be} (at \ref{cap:bec-766061be}) & \\
    \hline
    Service Availability Management & \texttt{cap:bec-22cf71df} (at \ref{cap:bec-22cf71df}) & \\
    \hline
    Service Discovery & \texttt{cap:bec-dee34cb8} (at \ref{cap:bec-dee34cb8}) & \\
    \hline
  \end{tabularx}
\end{table}

Intent. Healthcare Service Definition is a candidate enabling capability within Health Service Administration. It frames the reusable system-supported behaviours needed to support this healthcare context without prescribing a product, interface, technology or deployment boundary.

Boundary. This capability supports the business activity and does not inherit the authority or accountability of the Business Capability it enables. Its child Features are review candidates until ancestry and semantics are adjudicated.

FHIR R5 conceptual analogues considered (not implementation mappings): HealthcareService, ServiceRequest.

\subsubsection{Feature - Service Catalogue Publishing}\label{cap:bec-5af74389}
\vspace{-1em}
\begin{table}[htbp]
  \centering
  \footnotesize
  \begin{tabularx}{1\textwidth}{|l|l|X|}
    \hline
    \textbf{Metadata Name} & \textbf{Metadata Value} & \textbf{Description} \\
    \hline
    \textbf{Reference Id} & \texttt{BEC-5AF74389} & \\
    \hline
    \textbf{Status} & Existing Feature; candidate CT3 placement & \\
    \hline
    \textbf{Ancestry} &  & \\
    \hline
    \textbf{Source} & FEAT-EM-24 & \\
    \hline
    \textbf{Enabled By} &  & \\
    \hline
    \textbf{Supports} &  & \\
    \hline
    \textbf{Derived From} &  & \\
    \hline
    \textbf{Reviewed Status} & Confirm existing Feature meaning and proposed CT3 ownership. & \\
    \hline
  \end{tabularx}
\end{table}

Intent. Service Catalogue Publishing provides a distinct enabling behaviour in the context of Healthcare Service Definition. This is an existing named Feature; its placement under the candidate CT3 remains subject to review.

\noindent\textbf{Feature boundary.} This behaviour is scoped to Healthcare Service Definition; an identically named Feature elsewhere may carry different business meaning.\par
\featuretable{
Function & \texttt{FN-1D8F7F27} & Maintain Healthcare Service Definition & Catalogue \\
Service & \texttt{SV-E52D8F6C} & Healthcare Service Directory Query & Linked to preceding Function in source \\
}
\subsubsection{Feature - Service Definition Management}\label{cap:bec-8cdb717f}
\vspace{-1em}
\begin{table}[htbp]
  \centering
  \footnotesize
  \begin{tabularx}{1\textwidth}{|l|l|X|}
    \hline
    \textbf{Metadata Name} & \textbf{Metadata Value} & \textbf{Description} \\
    \hline
    \textbf{Reference Id} & \texttt{BEC-8CDB717F} & \\
    \hline
    \textbf{Status} & Generated candidate & \\
    \hline
    \textbf{Ancestry} &  & \\
    \hline
    \textbf{Source} & FHIR-informed candidate & \\
    \hline
    \textbf{Enabled By} &  & \\
    \hline
    \textbf{Supports} &  & \\
    \hline
    \textbf{Derived From} &  & \\
    \hline
    \textbf{Reviewed Status} & Generated candidate: keep / modify / relocate / reject; verify whether independently testable. & \\
    \hline
  \end{tabularx}
\end{table}

Intent. Service Definition Management is proposed to support the underlying healthcare concept represented by Healthcare Service Definition. The proposal concerns behaviour and semantic responsibility, not a FHIR field, operation, resource instance or prescribed data model.

Conceptual evidence: HL7 FHIR R5 HealthcareService, ServiceRequest.

\noindent\textbf{Feature boundary.} This behaviour is scoped to Healthcare Service Definition; an identically named Feature elsewhere may carry different business meaning.\par
\featuretable{
--- & --- & \emph{Not yet captured} & \emph{Complete during review} \\
}
\subsubsection{Feature - Service Eligibility Management}\label{cap:bec-766061be}
\vspace{-1em}
\begin{table}[htbp]
  \centering
  \footnotesize
  \begin{tabularx}{1\textwidth}{|l|l|X|}
    \hline
    \textbf{Metadata Name} & \textbf{Metadata Value} & \textbf{Description} \\
    \hline
    \textbf{Reference Id} & \texttt{BEC-766061BE} & \\
    \hline
    \textbf{Status} & Generated candidate & \\
    \hline
    \textbf{Ancestry} &  & \\
    \hline
    \textbf{Source} & FHIR-informed candidate & \\
    \hline
    \textbf{Enabled By} &  & \\
    \hline
    \textbf{Supports} &  & \\
    \hline
    \textbf{Derived From} &  & \\
    \hline
    \textbf{Reviewed Status} & Generated candidate: keep / modify / relocate / reject; verify whether independently testable. & \\
    \hline
  \end{tabularx}
\end{table}

Intent. Service Eligibility Management is proposed to support the underlying healthcare concept represented by Healthcare Service Definition. The proposal concerns behaviour and semantic responsibility, not a FHIR field, operation, resource instance or prescribed data model.

Conceptual evidence: HL7 FHIR R5 HealthcareService, ServiceRequest.

\noindent\textbf{Feature boundary.} This behaviour is scoped to Healthcare Service Definition; an identically named Feature elsewhere may carry different business meaning.\par
\featuretable{
--- & --- & \emph{Not yet captured} & \emph{Complete during review} \\
}
\subsubsection{Feature - Service Availability Management}\label{cap:bec-22cf71df}
\vspace{-1em}
\begin{table}[htbp]
  \centering
  \footnotesize
  \begin{tabularx}{1\textwidth}{|l|l|X|}
    \hline
    \textbf{Metadata Name} & \textbf{Metadata Value} & \textbf{Description} \\
    \hline
    \textbf{Reference Id} & \texttt{BEC-22CF71DF} & \\
    \hline
    \textbf{Status} & Generated candidate & \\
    \hline
    \textbf{Ancestry} &  & \\
    \hline
    \textbf{Source} & FHIR-informed candidate & \\
    \hline
    \textbf{Enabled By} &  & \\
    \hline
    \textbf{Supports} &  & \\
    \hline
    \textbf{Derived From} &  & \\
    \hline
    \textbf{Reviewed Status} & Generated candidate: keep / modify / relocate / reject; verify whether independently testable. & \\
    \hline
  \end{tabularx}
\end{table}

Intent. Service Availability Management is proposed to support the underlying healthcare concept represented by Healthcare Service Definition. The proposal concerns behaviour and semantic responsibility, not a FHIR field, operation, resource instance or prescribed data model.

Conceptual evidence: HL7 FHIR R5 HealthcareService, ServiceRequest.

\noindent\textbf{Feature boundary.} This behaviour is scoped to Healthcare Service Definition; an identically named Feature elsewhere may carry different business meaning.\par
\featuretable{
--- & --- & \emph{Not yet captured} & \emph{Complete during review} \\
}
\subsubsection{Feature - Service Discovery}\label{cap:bec-dee34cb8}
\vspace{-1em}
\begin{table}[htbp]
  \centering
  \footnotesize
  \begin{tabularx}{1\textwidth}{|l|l|X|}
    \hline
    \textbf{Metadata Name} & \textbf{Metadata Value} & \textbf{Description} \\
    \hline
    \textbf{Reference Id} & \texttt{BEC-DEE34CB8} & \\
    \hline
    \textbf{Status} & Generated candidate & \\
    \hline
    \textbf{Ancestry} &  & \\
    \hline
    \textbf{Source} & FHIR-informed candidate & \\
    \hline
    \textbf{Enabled By} &  & \\
    \hline
    \textbf{Supports} &  & \\
    \hline
    \textbf{Derived From} &  & \\
    \hline
    \textbf{Reviewed Status} & Generated candidate: keep / modify / relocate / reject; verify whether independently testable. & \\
    \hline
  \end{tabularx}
\end{table}

Intent. Service Discovery is proposed to support the underlying healthcare concept represented by Healthcare Service Definition. The proposal concerns behaviour and semantic responsibility, not a FHIR field, operation, resource instance or prescribed data model.

Conceptual evidence: HL7 FHIR R5 HealthcareService, ServiceRequest.

\noindent\textbf{Feature boundary.} This behaviour is scoped to Healthcare Service Definition; an identically named Feature elsewhere may carry different business meaning.\par
\featuretable{
--- & --- & \emph{Not yet captured} & \emph{Complete during review} \\
}

\subsection{CT3 - Service Provision Relationship}\label{cap:bec-cc9fd951}
\vspace{-0.5em}
\begin{table}[htbp]
  \centering
  \footnotesize
  \begin{tabularx}{1\textwidth}{|l|l|X|}
    \hline
    \multicolumn{3}{c}{\textbf{Capability Metadata}} \\
    \hline
    \textbf{Name} & \textbf{Value} & \textbf{Comment and/or Description} \\
    \hline
    \textbf{Reference Id} & \texttt{BEC-CC9FD951} & \\
    \hline
    \textbf{Status} & Candidate ancestry & \\
    \hline
    \textbf{Ancestry} &  & \\
    \hline
    \textbf{Source} &  & \\
    \hline
    \textbf{Enabled By} &  & \\
    \hline
    \textbf{Supports} &  & \\
    \hline
    \textbf{Derived From} &  & \\
    \hline
    \textbf{Reviewed Status} &  & \\
    \hline
  \end{tabularx}
\end{table}

\begin{table}[htbp]
  \centering
  \footnotesize
  \begin{tabularx}{1\textwidth}{|l|l|X|}
    \hline
    \multicolumn{3}{c}{\textbf{Contained Features}} \\
    \hline
    \textbf{Feature Name} & \textbf{Feature Id} & \textbf{Comment} \\
    \hline
    Service-Location-Provider Binding & \texttt{cap:bec-c68165b1} (at \ref{cap:bec-c68165b1}) & \\
    \hline
    Service Definition Management & \texttt{cap:bec-46b095fb} (at \ref{cap:bec-46b095fb}) & \\
    \hline
    Service Eligibility Management & \texttt{cap:bec-99951d40} (at \ref{cap:bec-99951d40}) & \\
    \hline
    Service Availability Management & \texttt{cap:bec-f964694f} (at \ref{cap:bec-f964694f}) & \\
    \hline
    Service Discovery & \texttt{cap:bec-7a743ddd} (at \ref{cap:bec-7a743ddd}) & \\
    \hline
  \end{tabularx}
\end{table}

Intent. Service Provision Relationship is a candidate enabling capability within Health Service Administration. It frames the reusable system-supported behaviours needed to support this healthcare context without prescribing a product, interface, technology or deployment boundary.

Boundary. This capability supports the business activity and does not inherit the authority or accountability of the Business Capability it enables. Its child Features are review candidates until ancestry and semantics are adjudicated.

FHIR R5 conceptual analogues considered (not implementation mappings): HealthcareService, ServiceRequest.

\subsubsection{Feature - Service-Location-Provider Binding}\label{cap:bec-c68165b1}
\vspace{-1em}
\begin{table}[htbp]
  \centering
  \footnotesize
  \begin{tabularx}{1\textwidth}{|l|l|X|}
    \hline
    \textbf{Metadata Name} & \textbf{Metadata Value} & \textbf{Description} \\
    \hline
    \textbf{Reference Id} & \texttt{BEC-C68165B1} & \\
    \hline
    \textbf{Status} & Existing Feature; candidate CT3 placement & \\
    \hline
    \textbf{Ancestry} &  & \\
    \hline
    \textbf{Source} & FEAT-EM-25 & \\
    \hline
    \textbf{Enabled By} &  & \\
    \hline
    \textbf{Supports} &  & \\
    \hline
    \textbf{Derived From} &  & \\
    \hline
    \textbf{Reviewed Status} & Confirm existing Feature meaning and proposed CT3 ownership. & \\
    \hline
  \end{tabularx}
\end{table}

Intent. Service-Location-Provider Binding provides a distinct enabling behaviour in the context of Service Provision Relationship. This is an existing named Feature; its placement under the candidate CT3 remains subject to review.

\noindent\textbf{Feature boundary.} This behaviour is scoped to Service Provision Relationship; an identically named Feature elsewhere may carry different business meaning.\par
\featuretable{
Function & \texttt{FN-77492C30} & Maintain Service / Location / Provider Map & Catalogue \\
Service & \texttt{SV-67205926} & Service Provision Resolution & Linked to preceding Function in source \\
}
\subsubsection{Feature - Service Definition Management}\label{cap:bec-46b095fb}
\vspace{-1em}
\begin{table}[htbp]
  \centering
  \footnotesize
  \begin{tabularx}{1\textwidth}{|l|l|X|}
    \hline
    \textbf{Metadata Name} & \textbf{Metadata Value} & \textbf{Description} \\
    \hline
    \textbf{Reference Id} & \texttt{BEC-46B095FB} & \\
    \hline
    \textbf{Status} & Generated candidate & \\
    \hline
    \textbf{Ancestry} &  & \\
    \hline
    \textbf{Source} & FHIR-informed candidate & \\
    \hline
    \textbf{Enabled By} &  & \\
    \hline
    \textbf{Supports} &  & \\
    \hline
    \textbf{Derived From} &  & \\
    \hline
    \textbf{Reviewed Status} & Generated candidate: keep / modify / relocate / reject; verify whether independently testable. & \\
    \hline
  \end{tabularx}
\end{table}

Intent. Service Definition Management is proposed to support the underlying healthcare concept represented by Service Provision Relationship. The proposal concerns behaviour and semantic responsibility, not a FHIR field, operation, resource instance or prescribed data model.

Conceptual evidence: HL7 FHIR R5 HealthcareService, ServiceRequest.

\noindent\textbf{Feature boundary.} This behaviour is scoped to Service Provision Relationship; an identically named Feature elsewhere may carry different business meaning.\par
\featuretable{
--- & --- & \emph{Not yet captured} & \emph{Complete during review} \\
}
\subsubsection{Feature - Service Eligibility Management}\label{cap:bec-99951d40}
\vspace{-1em}
\begin{table}[htbp]
  \centering
  \footnotesize
  \begin{tabularx}{1\textwidth}{|l|l|X|}
    \hline
    \textbf{Metadata Name} & \textbf{Metadata Value} & \textbf{Description} \\
    \hline
    \textbf{Reference Id} & \texttt{BEC-99951D40} & \\
    \hline
    \textbf{Status} & Generated candidate & \\
    \hline
    \textbf{Ancestry} &  & \\
    \hline
    \textbf{Source} & FHIR-informed candidate & \\
    \hline
    \textbf{Enabled By} &  & \\
    \hline
    \textbf{Supports} &  & \\
    \hline
    \textbf{Derived From} &  & \\
    \hline
    \textbf{Reviewed Status} & Generated candidate: keep / modify / relocate / reject; verify whether independently testable. & \\
    \hline
  \end{tabularx}
\end{table}

Intent. Service Eligibility Management is proposed to support the underlying healthcare concept represented by Service Provision Relationship. The proposal concerns behaviour and semantic responsibility, not a FHIR field, operation, resource instance or prescribed data model.

Conceptual evidence: HL7 FHIR R5 HealthcareService, ServiceRequest.

\noindent\textbf{Feature boundary.} This behaviour is scoped to Service Provision Relationship; an identically named Feature elsewhere may carry different business meaning.\par
\featuretable{
--- & --- & \emph{Not yet captured} & \emph{Complete during review} \\
}
\subsubsection{Feature - Service Availability Management}\label{cap:bec-f964694f}
\vspace{-1em}
\begin{table}[htbp]
  \centering
  \footnotesize
  \begin{tabularx}{1\textwidth}{|l|l|X|}
    \hline
    \textbf{Metadata Name} & \textbf{Metadata Value} & \textbf{Description} \\
    \hline
    \textbf{Reference Id} & \texttt{BEC-F964694F} & \\
    \hline
    \textbf{Status} & Generated candidate & \\
    \hline
    \textbf{Ancestry} &  & \\
    \hline
    \textbf{Source} & FHIR-informed candidate & \\
    \hline
    \textbf{Enabled By} &  & \\
    \hline
    \textbf{Supports} &  & \\
    \hline
    \textbf{Derived From} &  & \\
    \hline
    \textbf{Reviewed Status} & Generated candidate: keep / modify / relocate / reject; verify whether independently testable. & \\
    \hline
  \end{tabularx}
\end{table}

Intent. Service Availability Management is proposed to support the underlying healthcare concept represented by Service Provision Relationship. The proposal concerns behaviour and semantic responsibility, not a FHIR field, operation, resource instance or prescribed data model.

Conceptual evidence: HL7 FHIR R5 HealthcareService, ServiceRequest.

\noindent\textbf{Feature boundary.} This behaviour is scoped to Service Provision Relationship; an identically named Feature elsewhere may carry different business meaning.\par
\featuretable{
--- & --- & \emph{Not yet captured} & \emph{Complete during review} \\
}
\subsubsection{Feature - Service Discovery}\label{cap:bec-7a743ddd}
\vspace{-1em}
\begin{table}[htbp]
  \centering
  \footnotesize
  \begin{tabularx}{1\textwidth}{|l|l|X|}
    \hline
    \textbf{Metadata Name} & \textbf{Metadata Value} & \textbf{Description} \\
    \hline
    \textbf{Reference Id} & \texttt{BEC-7A743DDD} & \\
    \hline
    \textbf{Status} & Generated candidate & \\
    \hline
    \textbf{Ancestry} &  & \\
    \hline
    \textbf{Source} & FHIR-informed candidate & \\
    \hline
    \textbf{Enabled By} &  & \\
    \hline
    \textbf{Supports} &  & \\
    \hline
    \textbf{Derived From} &  & \\
    \hline
    \textbf{Reviewed Status} & Generated candidate: keep / modify / relocate / reject; verify whether independently testable. & \\
    \hline
  \end{tabularx}
\end{table}

Intent. Service Discovery is proposed to support the underlying healthcare concept represented by Service Provision Relationship. The proposal concerns behaviour and semantic responsibility, not a FHIR field, operation, resource instance or prescribed data model.

Conceptual evidence: HL7 FHIR R5 HealthcareService, ServiceRequest.

\noindent\textbf{Feature boundary.} This behaviour is scoped to Service Provision Relationship; an identically named Feature elsewhere may carry different business meaning.\par
\featuretable{
--- & --- & \emph{Not yet captured} & \emph{Complete during review} \\
}

\subsection{CT3 - Service Availability Definition}\label{cap:bec-cf667cc1}
\vspace{-0.5em}
\begin{table}[htbp]
  \centering
  \footnotesize
  \begin{tabularx}{1\textwidth}{|l|l|X|}
    \hline
    \multicolumn{3}{c}{\textbf{Capability Metadata}} \\
    \hline
    \textbf{Name} & \textbf{Value} & \textbf{Comment and/or Description} \\
    \hline
    \textbf{Reference Id} & \texttt{BEC-CF667CC1} & \\
    \hline
    \textbf{Status} & Candidate ancestry & \\
    \hline
    \textbf{Ancestry} &  & \\
    \hline
    \textbf{Source} &  & \\
    \hline
    \textbf{Enabled By} &  & \\
    \hline
    \textbf{Supports} &  & \\
    \hline
    \textbf{Derived From} &  & \\
    \hline
    \textbf{Reviewed Status} &  & \\
    \hline
  \end{tabularx}
\end{table}

\begin{table}[htbp]
  \centering
  \footnotesize
  \begin{tabularx}{1\textwidth}{|l|l|X|}
    \hline
    \multicolumn{3}{c}{\textbf{Contained Features}} \\
    \hline
    \textbf{Feature Name} & \textbf{Feature Id} & \textbf{Comment} \\
    \hline
    Operating Schedule Specification & \texttt{cap:bec-a4abc149} (at \ref{cap:bec-a4abc149}) & \\
    \hline
    Service Definition Management & \texttt{cap:bec-61f28669} (at \ref{cap:bec-61f28669}) & \\
    \hline
    Service Eligibility Management & \texttt{cap:bec-0aa34d18} (at \ref{cap:bec-0aa34d18}) & \\
    \hline
    Service Availability Management & \texttt{cap:bec-235c211d} (at \ref{cap:bec-235c211d}) & \\
    \hline
    Service Discovery & \texttt{cap:bec-b8a2142a} (at \ref{cap:bec-b8a2142a}) & \\
    \hline
  \end{tabularx}
\end{table}

Intent. Service Availability Definition is a candidate enabling capability within Health Service Administration. It frames the reusable system-supported behaviours needed to support this healthcare context without prescribing a product, interface, technology or deployment boundary.

Boundary. This capability supports the business activity and does not inherit the authority or accountability of the Business Capability it enables. Its child Features are review candidates until ancestry and semantics are adjudicated.

FHIR R5 conceptual analogues considered (not implementation mappings): HealthcareService, ServiceRequest.

\subsubsection{Feature - Operating Schedule Specification}\label{cap:bec-a4abc149}
\vspace{-1em}
\begin{table}[htbp]
  \centering
  \footnotesize
  \begin{tabularx}{1\textwidth}{|l|l|X|}
    \hline
    \textbf{Metadata Name} & \textbf{Metadata Value} & \textbf{Description} \\
    \hline
    \textbf{Reference Id} & \texttt{BEC-A4ABC149} & \\
    \hline
    \textbf{Status} & Existing Feature; candidate CT3 placement & \\
    \hline
    \textbf{Ancestry} &  & \\
    \hline
    \textbf{Source} & FEAT-EM-26 & \\
    \hline
    \textbf{Enabled By} &  & \\
    \hline
    \textbf{Supports} &  & \\
    \hline
    \textbf{Derived From} &  & \\
    \hline
    \textbf{Reviewed Status} & Confirm existing Feature meaning and proposed CT3 ownership. & \\
    \hline
  \end{tabularx}
\end{table}

Intent. Operating Schedule Specification provides a distinct enabling behaviour in the context of Service Availability Definition. This is an existing named Feature; its placement under the candidate CT3 remains subject to review.

\noindent\textbf{Feature boundary.} This behaviour is scoped to Service Availability Definition; an identically named Feature elsewhere may carry different business meaning.\par
\featuretable{
Function & \texttt{FN-C2FB32EA} & Maintain Service Availability & Catalogue \\
Service & \texttt{SV-BC04899F} & Service Availability Query & Linked to preceding Function in source \\
}
\subsubsection{Feature - Service Definition Management}\label{cap:bec-61f28669}
\vspace{-1em}
\begin{table}[htbp]
  \centering
  \footnotesize
  \begin{tabularx}{1\textwidth}{|l|l|X|}
    \hline
    \textbf{Metadata Name} & \textbf{Metadata Value} & \textbf{Description} \\
    \hline
    \textbf{Reference Id} & \texttt{BEC-61F28669} & \\
    \hline
    \textbf{Status} & Generated candidate & \\
    \hline
    \textbf{Ancestry} &  & \\
    \hline
    \textbf{Source} & FHIR-informed candidate & \\
    \hline
    \textbf{Enabled By} &  & \\
    \hline
    \textbf{Supports} &  & \\
    \hline
    \textbf{Derived From} &  & \\
    \hline
    \textbf{Reviewed Status} & Generated candidate: keep / modify / relocate / reject; verify whether independently testable. & \\
    \hline
  \end{tabularx}
\end{table}

Intent. Service Definition Management is proposed to support the underlying healthcare concept represented by Service Availability Definition. The proposal concerns behaviour and semantic responsibility, not a FHIR field, operation, resource instance or prescribed data model.

Conceptual evidence: HL7 FHIR R5 HealthcareService, ServiceRequest.

\noindent\textbf{Feature boundary.} This behaviour is scoped to Service Availability Definition; an identically named Feature elsewhere may carry different business meaning.\par
\featuretable{
--- & --- & \emph{Not yet captured} & \emph{Complete during review} \\
}
\subsubsection{Feature - Service Eligibility Management}\label{cap:bec-0aa34d18}
\vspace{-1em}
\begin{table}[htbp]
  \centering
  \footnotesize
  \begin{tabularx}{1\textwidth}{|l|l|X|}
    \hline
    \textbf{Metadata Name} & \textbf{Metadata Value} & \textbf{Description} \\
    \hline
    \textbf{Reference Id} & \texttt{BEC-0AA34D18} & \\
    \hline
    \textbf{Status} & Generated candidate & \\
    \hline
    \textbf{Ancestry} &  & \\
    \hline
    \textbf{Source} & FHIR-informed candidate & \\
    \hline
    \textbf{Enabled By} &  & \\
    \hline
    \textbf{Supports} &  & \\
    \hline
    \textbf{Derived From} &  & \\
    \hline
    \textbf{Reviewed Status} & Generated candidate: keep / modify / relocate / reject; verify whether independently testable. & \\
    \hline
  \end{tabularx}
\end{table}

Intent. Service Eligibility Management is proposed to support the underlying healthcare concept represented by Service Availability Definition. The proposal concerns behaviour and semantic responsibility, not a FHIR field, operation, resource instance or prescribed data model.

Conceptual evidence: HL7 FHIR R5 HealthcareService, ServiceRequest.

\noindent\textbf{Feature boundary.} This behaviour is scoped to Service Availability Definition; an identically named Feature elsewhere may carry different business meaning.\par
\featuretable{
--- & --- & \emph{Not yet captured} & \emph{Complete during review} \\
}
\subsubsection{Feature - Service Availability Management}\label{cap:bec-235c211d}
\vspace{-1em}
\begin{table}[htbp]
  \centering
  \footnotesize
  \begin{tabularx}{1\textwidth}{|l|l|X|}
    \hline
    \textbf{Metadata Name} & \textbf{Metadata Value} & \textbf{Description} \\
    \hline
    \textbf{Reference Id} & \texttt{BEC-235C211D} & \\
    \hline
    \textbf{Status} & Generated candidate & \\
    \hline
    \textbf{Ancestry} &  & \\
    \hline
    \textbf{Source} & FHIR-informed candidate & \\
    \hline
    \textbf{Enabled By} &  & \\
    \hline
    \textbf{Supports} &  & \\
    \hline
    \textbf{Derived From} &  & \\
    \hline
    \textbf{Reviewed Status} & Generated candidate: keep / modify / relocate / reject; verify whether independently testable. & \\
    \hline
  \end{tabularx}
\end{table}

Intent. Service Availability Management is proposed to support the underlying healthcare concept represented by Service Availability Definition. The proposal concerns behaviour and semantic responsibility, not a FHIR field, operation, resource instance or prescribed data model.

Conceptual evidence: HL7 FHIR R5 HealthcareService, ServiceRequest.

\noindent\textbf{Feature boundary.} This behaviour is scoped to Service Availability Definition; an identically named Feature elsewhere may carry different business meaning.\par
\featuretable{
--- & --- & \emph{Not yet captured} & \emph{Complete during review} \\
}
\subsubsection{Feature - Service Discovery}\label{cap:bec-b8a2142a}
\vspace{-1em}
\begin{table}[htbp]
  \centering
  \footnotesize
  \begin{tabularx}{1\textwidth}{|l|l|X|}
    \hline
    \textbf{Metadata Name} & \textbf{Metadata Value} & \textbf{Description} \\
    \hline
    \textbf{Reference Id} & \texttt{BEC-B8A2142A} & \\
    \hline
    \textbf{Status} & Generated candidate & \\
    \hline
    \textbf{Ancestry} &  & \\
    \hline
    \textbf{Source} & FHIR-informed candidate & \\
    \hline
    \textbf{Enabled By} &  & \\
    \hline
    \textbf{Supports} &  & \\
    \hline
    \textbf{Derived From} &  & \\
    \hline
    \textbf{Reviewed Status} & Generated candidate: keep / modify / relocate / reject; verify whether independently testable. & \\
    \hline
  \end{tabularx}
\end{table}

Intent. Service Discovery is proposed to support the underlying healthcare concept represented by Service Availability Definition. The proposal concerns behaviour and semantic responsibility, not a FHIR field, operation, resource instance or prescribed data model.

Conceptual evidence: HL7 FHIR R5 HealthcareService, ServiceRequest.

\noindent\textbf{Feature boundary.} This behaviour is scoped to Service Availability Definition; an identically named Feature elsewhere may carry different business meaning.\par
\featuretable{
--- & --- & \emph{Not yet captured} & \emph{Complete during review} \\
}

\subsection{CT3 - Service Eligibility Definition}\label{cap:bec-8858c964}
\vspace{-0.5em}
\begin{table}[htbp]
  \centering
  \footnotesize
  \begin{tabularx}{1\textwidth}{|l|l|X|}
    \hline
    \multicolumn{3}{c}{\textbf{Capability Metadata}} \\
    \hline
    \textbf{Name} & \textbf{Value} & \textbf{Comment and/or Description} \\
    \hline
    \textbf{Reference Id} & \texttt{BEC-8858C964} & \\
    \hline
    \textbf{Status} & Candidate ancestry & \\
    \hline
    \textbf{Ancestry} &  & \\
    \hline
    \textbf{Source} &  & \\
    \hline
    \textbf{Enabled By} &  & \\
    \hline
    \textbf{Supports} &  & \\
    \hline
    \textbf{Derived From} &  & \\
    \hline
    \textbf{Reviewed Status} &  & \\
    \hline
  \end{tabularx}
\end{table}

\begin{table}[htbp]
  \centering
  \footnotesize
  \begin{tabularx}{1\textwidth}{|l|l|X|}
    \hline
    \multicolumn{3}{c}{\textbf{Contained Features}} \\
    \hline
    \textbf{Feature Name} & \textbf{Feature Id} & \textbf{Comment} \\
    \hline
    Eligibility \& Catchment Specification & \texttt{cap:bec-748b065d} (at \ref{cap:bec-748b065d}) & \\
    \hline
    Service Definition Management & \texttt{cap:bec-271154e8} (at \ref{cap:bec-271154e8}) & \\
    \hline
    Service Eligibility Management & \texttt{cap:bec-1f711951} (at \ref{cap:bec-1f711951}) & \\
    \hline
    Service Availability Management & \texttt{cap:bec-fa68f289} (at \ref{cap:bec-fa68f289}) & \\
    \hline
    Service Discovery & \texttt{cap:bec-9b1e4ac5} (at \ref{cap:bec-9b1e4ac5}) & \\
    \hline
  \end{tabularx}
\end{table}

Intent. Service Eligibility Definition is a candidate enabling capability within Health Service Administration. It frames the reusable system-supported behaviours needed to support this healthcare context without prescribing a product, interface, technology or deployment boundary.

Boundary. This capability supports the business activity and does not inherit the authority or accountability of the Business Capability it enables. Its child Features are review candidates until ancestry and semantics are adjudicated.

FHIR R5 conceptual analogues considered (not implementation mappings): HealthcareService, ServiceRequest.

\subsubsection{Feature - Eligibility \& Catchment Specification}\label{cap:bec-748b065d}
\vspace{-1em}
\begin{table}[htbp]
  \centering
  \footnotesize
  \begin{tabularx}{1\textwidth}{|l|l|X|}
    \hline
    \textbf{Metadata Name} & \textbf{Metadata Value} & \textbf{Description} \\
    \hline
    \textbf{Reference Id} & \texttt{BEC-748B065D} & \\
    \hline
    \textbf{Status} & Existing Feature; candidate CT3 placement & \\
    \hline
    \textbf{Ancestry} &  & \\
    \hline
    \textbf{Source} & FEAT-EM-27 & \\
    \hline
    \textbf{Enabled By} &  & \\
    \hline
    \textbf{Supports} &  & \\
    \hline
    \textbf{Derived From} &  & \\
    \hline
    \textbf{Reviewed Status} & Confirm existing Feature meaning and proposed CT3 ownership. & \\
    \hline
  \end{tabularx}
\end{table}

Intent. Eligibility \& Catchment Specification provides a distinct enabling behaviour in the context of Service Eligibility Definition. This is an existing named Feature; its placement under the candidate CT3 remains subject to review.

\noindent\textbf{Feature boundary.} This behaviour is scoped to Service Eligibility Definition; an identically named Feature elsewhere may carry different business meaning.\par
\featuretable{
Function & \texttt{FN-6C76DAE7} & Maintain Service Eligibility Rules & Catalogue \\
Service & \texttt{SV-0C1B1639} & Service Eligibility Evaluation & Linked to preceding Function in source \\
}
\subsubsection{Feature - Service Definition Management}\label{cap:bec-271154e8}
\vspace{-1em}
\begin{table}[htbp]
  \centering
  \footnotesize
  \begin{tabularx}{1\textwidth}{|l|l|X|}
    \hline
    \textbf{Metadata Name} & \textbf{Metadata Value} & \textbf{Description} \\
    \hline
    \textbf{Reference Id} & \texttt{BEC-271154E8} & \\
    \hline
    \textbf{Status} & Generated candidate & \\
    \hline
    \textbf{Ancestry} &  & \\
    \hline
    \textbf{Source} & FHIR-informed candidate & \\
    \hline
    \textbf{Enabled By} &  & \\
    \hline
    \textbf{Supports} &  & \\
    \hline
    \textbf{Derived From} &  & \\
    \hline
    \textbf{Reviewed Status} & Generated candidate: keep / modify / relocate / reject; verify whether independently testable. & \\
    \hline
  \end{tabularx}
\end{table}

Intent. Service Definition Management is proposed to support the underlying healthcare concept represented by Service Eligibility Definition. The proposal concerns behaviour and semantic responsibility, not a FHIR field, operation, resource instance or prescribed data model.

Conceptual evidence: HL7 FHIR R5 HealthcareService, ServiceRequest.

\noindent\textbf{Feature boundary.} This behaviour is scoped to Service Eligibility Definition; an identically named Feature elsewhere may carry different business meaning.\par
\featuretable{
--- & --- & \emph{Not yet captured} & \emph{Complete during review} \\
}
\subsubsection{Feature - Service Eligibility Management}\label{cap:bec-1f711951}
\vspace{-1em}
\begin{table}[htbp]
  \centering
  \footnotesize
  \begin{tabularx}{1\textwidth}{|l|l|X|}
    \hline
    \textbf{Metadata Name} & \textbf{Metadata Value} & \textbf{Description} \\
    \hline
    \textbf{Reference Id} & \texttt{BEC-1F711951} & \\
    \hline
    \textbf{Status} & Generated candidate & \\
    \hline
    \textbf{Ancestry} &  & \\
    \hline
    \textbf{Source} & FHIR-informed candidate & \\
    \hline
    \textbf{Enabled By} &  & \\
    \hline
    \textbf{Supports} &  & \\
    \hline
    \textbf{Derived From} &  & \\
    \hline
    \textbf{Reviewed Status} & Generated candidate: keep / modify / relocate / reject; verify whether independently testable. & \\
    \hline
  \end{tabularx}
\end{table}

Intent. Service Eligibility Management is proposed to support the underlying healthcare concept represented by Service Eligibility Definition. The proposal concerns behaviour and semantic responsibility, not a FHIR field, operation, resource instance or prescribed data model.

Conceptual evidence: HL7 FHIR R5 HealthcareService, ServiceRequest.

\noindent\textbf{Feature boundary.} This behaviour is scoped to Service Eligibility Definition; an identically named Feature elsewhere may carry different business meaning.\par
\featuretable{
--- & --- & \emph{Not yet captured} & \emph{Complete during review} \\
}
\subsubsection{Feature - Service Availability Management}\label{cap:bec-fa68f289}
\vspace{-1em}
\begin{table}[htbp]
  \centering
  \footnotesize
  \begin{tabularx}{1\textwidth}{|l|l|X|}
    \hline
    \textbf{Metadata Name} & \textbf{Metadata Value} & \textbf{Description} \\
    \hline
    \textbf{Reference Id} & \texttt{BEC-FA68F289} & \\
    \hline
    \textbf{Status} & Generated candidate & \\
    \hline
    \textbf{Ancestry} &  & \\
    \hline
    \textbf{Source} & FHIR-informed candidate & \\
    \hline
    \textbf{Enabled By} &  & \\
    \hline
    \textbf{Supports} &  & \\
    \hline
    \textbf{Derived From} &  & \\
    \hline
    \textbf{Reviewed Status} & Generated candidate: keep / modify / relocate / reject; verify whether independently testable. & \\
    \hline
  \end{tabularx}
\end{table}

Intent. Service Availability Management is proposed to support the underlying healthcare concept represented by Service Eligibility Definition. The proposal concerns behaviour and semantic responsibility, not a FHIR field, operation, resource instance or prescribed data model.

Conceptual evidence: HL7 FHIR R5 HealthcareService, ServiceRequest.

\noindent\textbf{Feature boundary.} This behaviour is scoped to Service Eligibility Definition; an identically named Feature elsewhere may carry different business meaning.\par
\featuretable{
--- & --- & \emph{Not yet captured} & \emph{Complete during review} \\
}
\subsubsection{Feature - Service Discovery}\label{cap:bec-9b1e4ac5}
\vspace{-1em}
\begin{table}[htbp]
  \centering
  \footnotesize
  \begin{tabularx}{1\textwidth}{|l|l|X|}
    \hline
    \textbf{Metadata Name} & \textbf{Metadata Value} & \textbf{Description} \\
    \hline
    \textbf{Reference Id} & \texttt{BEC-9B1E4AC5} & \\
    \hline
    \textbf{Status} & Generated candidate & \\
    \hline
    \textbf{Ancestry} &  & \\
    \hline
    \textbf{Source} & FHIR-informed candidate & \\
    \hline
    \textbf{Enabled By} &  & \\
    \hline
    \textbf{Supports} &  & \\
    \hline
    \textbf{Derived From} &  & \\
    \hline
    \textbf{Reviewed Status} & Generated candidate: keep / modify / relocate / reject; verify whether independently testable. & \\
    \hline
  \end{tabularx}
\end{table}

Intent. Service Discovery is proposed to support the underlying healthcare concept represented by Service Eligibility Definition. The proposal concerns behaviour and semantic responsibility, not a FHIR field, operation, resource instance or prescribed data model.

Conceptual evidence: HL7 FHIR R5 HealthcareService, ServiceRequest.

\noindent\textbf{Feature boundary.} This behaviour is scoped to Service Eligibility Definition; an identically named Feature elsewhere may carry different business meaning.\par
\featuretable{
--- & --- & \emph{Not yet captured} & \emph{Complete during review} \\
}

\section{CT2 - Health Product Administration}\label{cap:bec-60f2add5}
\vspace{-0.5em}
\begin{table}[htbp]
  \centering
  \footnotesize
  \begin{tabularx}{1\textwidth}{|l|l|X|}
    \hline
    \multicolumn{3}{c}{\textbf{Capability Metadata}} \\
    \hline
    \textbf{Name} & \textbf{Value} & \textbf{Comment and/or Description} \\
    \hline
    \textbf{Reference Id} & \texttt{BEC-60F2ADD5} & \\
    \hline
    \textbf{Status} & Candidate ancestry & \\
    \hline
    \textbf{Ancestry} &  & \\
    \hline
    \textbf{Source} &  & \\
    \hline
    \textbf{Enabled By} &  & \\
    \hline
    \textbf{Supports} &  & \\
    \hline
    \textbf{Derived From} &  & \\
    \hline
    \textbf{Reviewed Status} &  & \\
    \hline
  \end{tabularx}
\end{table}

\begin{table}[htbp]
  \centering
  \footnotesize
  \begin{tabularx}{1\textwidth}{|l|l|X|}
    \hline
    \multicolumn{3}{c}{\textbf{Contained Capabilities}} \\
    \hline
    \textbf{Capability Name} & \textbf{Capability Id} & \textbf{Comment} \\
    \hline
    Health Product Definition & \texttt{cap:bec-1be2add6} (at \ref{cap:bec-1be2add6}) & \\
    \hline
    Product Formulary \& Catalogue & \texttt{cap:bec-f3f076f6} (at \ref{cap:bec-f3f076f6}) & \\
    \hline
    Product Batch \& Serial Tracking & \texttt{cap:bec-0ace1412} (at \ref{cap:bec-0ace1412}) & \\
    \hline
  \end{tabularx}
\end{table}

Intent. Health Product Administration is a candidate enabling capability within Entity Information Enablement. It frames the reusable system-supported behaviours needed to support this healthcare context without prescribing a product, interface, technology or deployment boundary.

Boundary. This capability supports the business activity and does not inherit the authority or accountability of the Business Capability it enables. Its child Features are review candidates until ancestry and semantics are adjudicated.

\subsection{CT3 - Health Product Definition}\label{cap:bec-1be2add6}
\vspace{-0.5em}
\begin{table}[htbp]
  \centering
  \footnotesize
  \begin{tabularx}{1\textwidth}{|l|l|X|}
    \hline
    \multicolumn{3}{c}{\textbf{Capability Metadata}} \\
    \hline
    \textbf{Name} & \textbf{Value} & \textbf{Comment and/or Description} \\
    \hline
    \textbf{Reference Id} & \texttt{BEC-1BE2ADD6} & \\
    \hline
    \textbf{Status} & Candidate ancestry & \\
    \hline
    \textbf{Ancestry} &  & \\
    \hline
    \textbf{Source} &  & \\
    \hline
    \textbf{Enabled By} &  & \\
    \hline
    \textbf{Supports} &  & \\
    \hline
    \textbf{Derived From} &  & \\
    \hline
    \textbf{Reviewed Status} &  & \\
    \hline
  \end{tabularx}
\end{table}

\begin{table}[htbp]
  \centering
  \footnotesize
  \begin{tabularx}{1\textwidth}{|l|l|X|}
    \hline
    \multicolumn{3}{c}{\textbf{Contained Features}} \\
    \hline
    \textbf{Feature Name} & \textbf{Feature Id} & \textbf{Comment} \\
    \hline
    No established Feature decomposition (Adjacent / Reference context) & \texttt{cap:bec-3700fbfb} (at \ref{cap:bec-3700fbfb}) & \\
    \hline
    Context Establishment & \texttt{cap:bec-0eecbf8c} (at \ref{cap:bec-0eecbf8c}) & \\
    \hline
    State Monitoring & \texttt{cap:bec-8f3d543e} (at \ref{cap:bec-8f3d543e}) & \\
    \hline
    Exception Management & \texttt{cap:bec-6f386b2a} (at \ref{cap:bec-6f386b2a}) & \\
    \hline
    Information Retrieval & \texttt{cap:bec-f792b00b} (at \ref{cap:bec-f792b00b}) & \\
    \hline
  \end{tabularx}
\end{table}

Intent. Health Product Definition is a candidate enabling capability within Health Product Administration. It frames the reusable system-supported behaviours needed to support this healthcare context without prescribing a product, interface, technology or deployment boundary.

Boundary. This capability supports the business activity and does not inherit the authority or accountability of the Business Capability it enables. Its child Features are review candidates until ancestry and semantics are adjudicated.

FHIR R5 conceptual analogues considered (not implementation mappings): Medication, Device, NutritionProduct.

\subsubsection{Feature - No established Feature decomposition (Adjacent / Reference context)}\label{cap:bec-3700fbfb}
\vspace{-1em}
\begin{table}[htbp]
  \centering
  \footnotesize
  \begin{tabularx}{1\textwidth}{|l|l|X|}
    \hline
    \textbf{Metadata Name} & \textbf{Metadata Value} & \textbf{Description} \\
    \hline
    \textbf{Reference Id} & \texttt{BEC-3700FBFB} & \\
    \hline
    \textbf{Status} & Generated candidate & \\
    \hline
    \textbf{Ancestry} &  & \\
    \hline
    \textbf{Source} & Existing catalogue & \\
    \hline
    \textbf{Enabled By} &  & \\
    \hline
    \textbf{Supports} &  & \\
    \hline
    \textbf{Derived From} &  & \\
    \hline
    \textbf{Reviewed Status} & Generated candidate: keep / modify / relocate / reject; verify whether independently testable. & \\
    \hline
  \end{tabularx}
\end{table}

Intent. No established Feature decomposition (Adjacent / Reference context) provides a distinct enabling behaviour in the context of Health Product Definition. This is an existing named Feature; its placement under the candidate CT3 remains subject to review.

\noindent\textbf{Feature boundary.} This behaviour is scoped to Health Product Definition; an identically named Feature elsewhere may carry different business meaning.\par
\featuretable{
--- & --- & \emph{Not yet captured} & \emph{Complete during review} \\
}
\subsubsection{Feature - Context Establishment}\label{cap:bec-0eecbf8c}
\vspace{-1em}
\begin{table}[htbp]
  \centering
  \footnotesize
  \begin{tabularx}{1\textwidth}{|l|l|X|}
    \hline
    \textbf{Metadata Name} & \textbf{Metadata Value} & \textbf{Description} \\
    \hline
    \textbf{Reference Id} & \texttt{BEC-0EECBF8C} & \\
    \hline
    \textbf{Status} & Generated candidate & \\
    \hline
    \textbf{Ancestry} &  & \\
    \hline
    \textbf{Source} & Semantic-pattern candidate & \\
    \hline
    \textbf{Enabled By} &  & \\
    \hline
    \textbf{Supports} &  & \\
    \hline
    \textbf{Derived From} &  & \\
    \hline
    \textbf{Reviewed Status} & Generated candidate: keep / modify / relocate / reject; verify whether independently testable. & \\
    \hline
  \end{tabularx}
\end{table}

Intent. Context Establishment is proposed as a discrete, testable enabling behaviour for Health Product Definition. It is included for architectural review, not as an established implementation obligation.

\noindent\textbf{Feature boundary.} This behaviour is scoped to Health Product Definition; an identically named Feature elsewhere may carry different business meaning.\par
\featuretable{
--- & --- & \emph{Not yet captured} & \emph{Complete during review} \\
}
\subsubsection{Feature - State Monitoring}\label{cap:bec-8f3d543e}
\vspace{-1em}
\begin{table}[htbp]
  \centering
  \footnotesize
  \begin{tabularx}{1\textwidth}{|l|l|X|}
    \hline
    \textbf{Metadata Name} & \textbf{Metadata Value} & \textbf{Description} \\
    \hline
    \textbf{Reference Id} & \texttt{BEC-8F3D543E} & \\
    \hline
    \textbf{Status} & Generated candidate & \\
    \hline
    \textbf{Ancestry} &  & \\
    \hline
    \textbf{Source} & Semantic-pattern candidate & \\
    \hline
    \textbf{Enabled By} &  & \\
    \hline
    \textbf{Supports} &  & \\
    \hline
    \textbf{Derived From} &  & \\
    \hline
    \textbf{Reviewed Status} & Generated candidate: keep / modify / relocate / reject; verify whether independently testable. & \\
    \hline
  \end{tabularx}
\end{table}

Intent. State Monitoring is proposed as a discrete, testable enabling behaviour for Health Product Definition. It is included for architectural review, not as an established implementation obligation.

\noindent\textbf{Feature boundary.} This behaviour is scoped to Health Product Definition; an identically named Feature elsewhere may carry different business meaning.\par
\featuretable{
--- & --- & \emph{Not yet captured} & \emph{Complete during review} \\
}
\subsubsection{Feature - Exception Management}\label{cap:bec-6f386b2a}
\vspace{-1em}
\begin{table}[htbp]
  \centering
  \footnotesize
  \begin{tabularx}{1\textwidth}{|l|l|X|}
    \hline
    \textbf{Metadata Name} & \textbf{Metadata Value} & \textbf{Description} \\
    \hline
    \textbf{Reference Id} & \texttt{BEC-6F386B2A} & \\
    \hline
    \textbf{Status} & Generated candidate & \\
    \hline
    \textbf{Ancestry} &  & \\
    \hline
    \textbf{Source} & Semantic-pattern candidate & \\
    \hline
    \textbf{Enabled By} &  & \\
    \hline
    \textbf{Supports} &  & \\
    \hline
    \textbf{Derived From} &  & \\
    \hline
    \textbf{Reviewed Status} & Generated candidate: keep / modify / relocate / reject; verify whether independently testable. & \\
    \hline
  \end{tabularx}
\end{table}

Intent. Exception Management is proposed as a discrete, testable enabling behaviour for Health Product Definition. It is included for architectural review, not as an established implementation obligation.

\noindent\textbf{Feature boundary.} This behaviour is scoped to Health Product Definition; an identically named Feature elsewhere may carry different business meaning.\par
\featuretable{
--- & --- & \emph{Not yet captured} & \emph{Complete during review} \\
}
\subsubsection{Feature - Information Retrieval}\label{cap:bec-f792b00b}
\vspace{-1em}
\begin{table}[htbp]
  \centering
  \footnotesize
  \begin{tabularx}{1\textwidth}{|l|l|X|}
    \hline
    \textbf{Metadata Name} & \textbf{Metadata Value} & \textbf{Description} \\
    \hline
    \textbf{Reference Id} & \texttt{BEC-F792B00B} & \\
    \hline
    \textbf{Status} & Generated candidate & \\
    \hline
    \textbf{Ancestry} &  & \\
    \hline
    \textbf{Source} & Semantic-pattern candidate & \\
    \hline
    \textbf{Enabled By} &  & \\
    \hline
    \textbf{Supports} &  & \\
    \hline
    \textbf{Derived From} &  & \\
    \hline
    \textbf{Reviewed Status} & Generated candidate: keep / modify / relocate / reject; verify whether independently testable. & \\
    \hline
  \end{tabularx}
\end{table}

Intent. Information Retrieval is proposed as a discrete, testable enabling behaviour for Health Product Definition. It is included for architectural review, not as an established implementation obligation.

\noindent\textbf{Feature boundary.} This behaviour is scoped to Health Product Definition; an identically named Feature elsewhere may carry different business meaning.\par
\featuretable{
--- & --- & \emph{Not yet captured} & \emph{Complete during review} \\
}

\subsection{CT3 - Product Formulary \& Catalogue}\label{cap:bec-f3f076f6}
\vspace{-0.5em}
\begin{table}[htbp]
  \centering
  \footnotesize
  \begin{tabularx}{1\textwidth}{|l|l|X|}
    \hline
    \multicolumn{3}{c}{\textbf{Capability Metadata}} \\
    \hline
    \textbf{Name} & \textbf{Value} & \textbf{Comment and/or Description} \\
    \hline
    \textbf{Reference Id} & \texttt{BEC-F3F076F6} & \\
    \hline
    \textbf{Status} & Candidate ancestry & \\
    \hline
    \textbf{Ancestry} &  & \\
    \hline
    \textbf{Source} &  & \\
    \hline
    \textbf{Enabled By} &  & \\
    \hline
    \textbf{Supports} &  & \\
    \hline
    \textbf{Derived From} &  & \\
    \hline
    \textbf{Reviewed Status} &  & \\
    \hline
  \end{tabularx}
\end{table}

\begin{table}[htbp]
  \centering
  \footnotesize
  \begin{tabularx}{1\textwidth}{|l|l|X|}
    \hline
    \multicolumn{3}{c}{\textbf{Contained Features}} \\
    \hline
    \textbf{Feature Name} & \textbf{Feature Id} & \textbf{Comment} \\
    \hline
    No established Feature decomposition (Adjacent / Reference context) & \texttt{cap:bec-4937ca4f} (at \ref{cap:bec-4937ca4f}) & \\
    \hline
    Identity Resolution & \texttt{cap:bec-d5cd867c} (at \ref{cap:bec-d5cd867c}) & \\
    \hline
    Identity Correlation & \texttt{cap:bec-a315e2f3} (at \ref{cap:bec-a315e2f3}) & \\
    \hline
    Relationship Management & \texttt{cap:bec-3dbacc5a} (at \ref{cap:bec-3dbacc5a}) & \\
    \hline
    Lifecycle Management & \texttt{cap:bec-5c469882} (at \ref{cap:bec-5c469882}) & \\
    \hline
    Discovery & \texttt{cap:bec-9b063e58} (at \ref{cap:bec-9b063e58}) & \\
    \hline
  \end{tabularx}
\end{table}

Intent. Product Formulary \& Catalogue is a candidate enabling capability within Health Product Administration. It frames the reusable system-supported behaviours needed to support this healthcare context without prescribing a product, interface, technology or deployment boundary.

Boundary. This capability supports the business activity and does not inherit the authority or accountability of the Business Capability it enables. Its child Features are review candidates until ancestry and semantics are adjudicated.

\subsubsection{Feature - No established Feature decomposition (Adjacent / Reference context)}\label{cap:bec-4937ca4f}
\vspace{-1em}
\begin{table}[htbp]
  \centering
  \footnotesize
  \begin{tabularx}{1\textwidth}{|l|l|X|}
    \hline
    \textbf{Metadata Name} & \textbf{Metadata Value} & \textbf{Description} \\
    \hline
    \textbf{Reference Id} & \texttt{BEC-4937CA4F} & \\
    \hline
    \textbf{Status} & Generated candidate & \\
    \hline
    \textbf{Ancestry} &  & \\
    \hline
    \textbf{Source} & Existing catalogue & \\
    \hline
    \textbf{Enabled By} &  & \\
    \hline
    \textbf{Supports} &  & \\
    \hline
    \textbf{Derived From} &  & \\
    \hline
    \textbf{Reviewed Status} & Generated candidate: keep / modify / relocate / reject; verify whether independently testable. & \\
    \hline
  \end{tabularx}
\end{table}

Intent. No established Feature decomposition (Adjacent / Reference context) provides a distinct enabling behaviour in the context of Product Formulary \& Catalogue. This is an existing named Feature; its placement under the candidate CT3 remains subject to review.

\noindent\textbf{Feature boundary.} This behaviour is scoped to Product Formulary \& Catalogue; an identically named Feature elsewhere may carry different business meaning.\par
\featuretable{
--- & --- & \emph{Not yet captured} & \emph{Complete during review} \\
}
\subsubsection{Feature - Identity Resolution}\label{cap:bec-d5cd867c}
\vspace{-1em}
\begin{table}[htbp]
  \centering
  \footnotesize
  \begin{tabularx}{1\textwidth}{|l|l|X|}
    \hline
    \textbf{Metadata Name} & \textbf{Metadata Value} & \textbf{Description} \\
    \hline
    \textbf{Reference Id} & \texttt{BEC-D5CD867C} & \\
    \hline
    \textbf{Status} & Generated candidate & \\
    \hline
    \textbf{Ancestry} &  & \\
    \hline
    \textbf{Source} & Semantic-pattern candidate & \\
    \hline
    \textbf{Enabled By} &  & \\
    \hline
    \textbf{Supports} &  & \\
    \hline
    \textbf{Derived From} &  & \\
    \hline
    \textbf{Reviewed Status} & Generated candidate: keep / modify / relocate / reject; verify whether independently testable. & \\
    \hline
  \end{tabularx}
\end{table}

Intent. Identity Resolution is proposed as a discrete, testable enabling behaviour for Product Formulary \& Catalogue. It is included for architectural review, not as an established implementation obligation.

\noindent\textbf{Feature boundary.} This behaviour is scoped to Product Formulary \& Catalogue; an identically named Feature elsewhere may carry different business meaning.\par
\featuretable{
--- & --- & \emph{Not yet captured} & \emph{Complete during review} \\
}
\subsubsection{Feature - Identity Correlation}\label{cap:bec-a315e2f3}
\vspace{-1em}
\begin{table}[htbp]
  \centering
  \footnotesize
  \begin{tabularx}{1\textwidth}{|l|l|X|}
    \hline
    \textbf{Metadata Name} & \textbf{Metadata Value} & \textbf{Description} \\
    \hline
    \textbf{Reference Id} & \texttt{BEC-A315E2F3} & \\
    \hline
    \textbf{Status} & Generated candidate & \\
    \hline
    \textbf{Ancestry} &  & \\
    \hline
    \textbf{Source} & Semantic-pattern candidate & \\
    \hline
    \textbf{Enabled By} &  & \\
    \hline
    \textbf{Supports} &  & \\
    \hline
    \textbf{Derived From} &  & \\
    \hline
    \textbf{Reviewed Status} & Generated candidate: keep / modify / relocate / reject; verify whether independently testable. & \\
    \hline
  \end{tabularx}
\end{table}

Intent. Identity Correlation is proposed as a discrete, testable enabling behaviour for Product Formulary \& Catalogue. It is included for architectural review, not as an established implementation obligation.

\noindent\textbf{Feature boundary.} This behaviour is scoped to Product Formulary \& Catalogue; an identically named Feature elsewhere may carry different business meaning.\par
\featuretable{
--- & --- & \emph{Not yet captured} & \emph{Complete during review} \\
}
\subsubsection{Feature - Relationship Management}\label{cap:bec-3dbacc5a}
\vspace{-1em}
\begin{table}[htbp]
  \centering
  \footnotesize
  \begin{tabularx}{1\textwidth}{|l|l|X|}
    \hline
    \textbf{Metadata Name} & \textbf{Metadata Value} & \textbf{Description} \\
    \hline
    \textbf{Reference Id} & \texttt{BEC-3DBACC5A} & \\
    \hline
    \textbf{Status} & Generated candidate & \\
    \hline
    \textbf{Ancestry} &  & \\
    \hline
    \textbf{Source} & Semantic-pattern candidate & \\
    \hline
    \textbf{Enabled By} &  & \\
    \hline
    \textbf{Supports} &  & \\
    \hline
    \textbf{Derived From} &  & \\
    \hline
    \textbf{Reviewed Status} & Generated candidate: keep / modify / relocate / reject; verify whether independently testable. & \\
    \hline
  \end{tabularx}
\end{table}

Intent. Relationship Management is proposed as a discrete, testable enabling behaviour for Product Formulary \& Catalogue. It is included for architectural review, not as an established implementation obligation.

\noindent\textbf{Feature boundary.} This behaviour is scoped to Product Formulary \& Catalogue; an identically named Feature elsewhere may carry different business meaning.\par
\featuretable{
--- & --- & \emph{Not yet captured} & \emph{Complete during review} \\
}
\subsubsection{Feature - Lifecycle Management}\label{cap:bec-5c469882}
\vspace{-1em}
\begin{table}[htbp]
  \centering
  \footnotesize
  \begin{tabularx}{1\textwidth}{|l|l|X|}
    \hline
    \textbf{Metadata Name} & \textbf{Metadata Value} & \textbf{Description} \\
    \hline
    \textbf{Reference Id} & \texttt{BEC-5C469882} & \\
    \hline
    \textbf{Status} & Generated candidate & \\
    \hline
    \textbf{Ancestry} &  & \\
    \hline
    \textbf{Source} & Semantic-pattern candidate & \\
    \hline
    \textbf{Enabled By} &  & \\
    \hline
    \textbf{Supports} &  & \\
    \hline
    \textbf{Derived From} &  & \\
    \hline
    \textbf{Reviewed Status} & Generated candidate: keep / modify / relocate / reject; verify whether independently testable. & \\
    \hline
  \end{tabularx}
\end{table}

Intent. Lifecycle Management is proposed as a discrete, testable enabling behaviour for Product Formulary \& Catalogue. It is included for architectural review, not as an established implementation obligation.

\noindent\textbf{Feature boundary.} This behaviour is scoped to Product Formulary \& Catalogue; an identically named Feature elsewhere may carry different business meaning.\par
\featuretable{
--- & --- & \emph{Not yet captured} & \emph{Complete during review} \\
}
\subsubsection{Feature - Discovery}\label{cap:bec-9b063e58}
\vspace{-1em}
\begin{table}[htbp]
  \centering
  \footnotesize
  \begin{tabularx}{1\textwidth}{|l|l|X|}
    \hline
    \textbf{Metadata Name} & \textbf{Metadata Value} & \textbf{Description} \\
    \hline
    \textbf{Reference Id} & \texttt{BEC-9B063E58} & \\
    \hline
    \textbf{Status} & Generated candidate & \\
    \hline
    \textbf{Ancestry} &  & \\
    \hline
    \textbf{Source} & Semantic-pattern candidate & \\
    \hline
    \textbf{Enabled By} &  & \\
    \hline
    \textbf{Supports} &  & \\
    \hline
    \textbf{Derived From} &  & \\
    \hline
    \textbf{Reviewed Status} & Generated candidate: keep / modify / relocate / reject; verify whether independently testable. & \\
    \hline
  \end{tabularx}
\end{table}

Intent. Discovery is proposed as a discrete, testable enabling behaviour for Product Formulary \& Catalogue. It is included for architectural review, not as an established implementation obligation.

\noindent\textbf{Feature boundary.} This behaviour is scoped to Product Formulary \& Catalogue; an identically named Feature elsewhere may carry different business meaning.\par
\featuretable{
--- & --- & \emph{Not yet captured} & \emph{Complete during review} \\
}

\subsection{CT3 - Product Batch \& Serial Tracking}\label{cap:bec-0ace1412}
\vspace{-0.5em}
\begin{table}[htbp]
  \centering
  \footnotesize
  \begin{tabularx}{1\textwidth}{|l|l|X|}
    \hline
    \multicolumn{3}{c}{\textbf{Capability Metadata}} \\
    \hline
    \textbf{Name} & \textbf{Value} & \textbf{Comment and/or Description} \\
    \hline
    \textbf{Reference Id} & \texttt{BEC-0ACE1412} & \\
    \hline
    \textbf{Status} & Candidate ancestry & \\
    \hline
    \textbf{Ancestry} &  & \\
    \hline
    \textbf{Source} &  & \\
    \hline
    \textbf{Enabled By} &  & \\
    \hline
    \textbf{Supports} &  & \\
    \hline
    \textbf{Derived From} &  & \\
    \hline
    \textbf{Reviewed Status} &  & \\
    \hline
  \end{tabularx}
\end{table}

\begin{table}[htbp]
  \centering
  \footnotesize
  \begin{tabularx}{1\textwidth}{|l|l|X|}
    \hline
    \multicolumn{3}{c}{\textbf{Contained Features}} \\
    \hline
    \textbf{Feature Name} & \textbf{Feature Id} & \textbf{Comment} \\
    \hline
    No established Feature decomposition (Adjacent / Reference context) & \texttt{cap:bec-4614d02a} (at \ref{cap:bec-4614d02a}) & \\
    \hline
    Context Establishment & \texttt{cap:bec-a2361427} (at \ref{cap:bec-a2361427}) & \\
    \hline
    State Monitoring & \texttt{cap:bec-c2e441ed} (at \ref{cap:bec-c2e441ed}) & \\
    \hline
    Exception Management & \texttt{cap:bec-db97a930} (at \ref{cap:bec-db97a930}) & \\
    \hline
    Information Retrieval & \texttt{cap:bec-4c113522} (at \ref{cap:bec-4c113522}) & \\
    \hline
  \end{tabularx}
\end{table}

Intent. Product Batch \& Serial Tracking is a candidate enabling capability within Health Product Administration. It frames the reusable system-supported behaviours needed to support this healthcare context without prescribing a product, interface, technology or deployment boundary.

Boundary. This capability supports the business activity and does not inherit the authority or accountability of the Business Capability it enables. Its child Features are review candidates until ancestry and semantics are adjudicated.

\subsubsection{Feature - No established Feature decomposition (Adjacent / Reference context)}\label{cap:bec-4614d02a}
\vspace{-1em}
\begin{table}[htbp]
  \centering
  \footnotesize
  \begin{tabularx}{1\textwidth}{|l|l|X|}
    \hline
    \textbf{Metadata Name} & \textbf{Metadata Value} & \textbf{Description} \\
    \hline
    \textbf{Reference Id} & \texttt{BEC-4614D02A} & \\
    \hline
    \textbf{Status} & Generated candidate & \\
    \hline
    \textbf{Ancestry} &  & \\
    \hline
    \textbf{Source} & Existing catalogue & \\
    \hline
    \textbf{Enabled By} &  & \\
    \hline
    \textbf{Supports} &  & \\
    \hline
    \textbf{Derived From} &  & \\
    \hline
    \textbf{Reviewed Status} & Generated candidate: keep / modify / relocate / reject; verify whether independently testable. & \\
    \hline
  \end{tabularx}
\end{table}

Intent. No established Feature decomposition (Adjacent / Reference context) provides a distinct enabling behaviour in the context of Product Batch \& Serial Tracking. This is an existing named Feature; its placement under the candidate CT3 remains subject to review.

\noindent\textbf{Feature boundary.} This behaviour is scoped to Product Batch \& Serial Tracking; an identically named Feature elsewhere may carry different business meaning.\par
\featuretable{
--- & --- & \emph{Not yet captured} & \emph{Complete during review} \\
}
\subsubsection{Feature - Context Establishment}\label{cap:bec-a2361427}
\vspace{-1em}
\begin{table}[htbp]
  \centering
  \footnotesize
  \begin{tabularx}{1\textwidth}{|l|l|X|}
    \hline
    \textbf{Metadata Name} & \textbf{Metadata Value} & \textbf{Description} \\
    \hline
    \textbf{Reference Id} & \texttt{BEC-A2361427} & \\
    \hline
    \textbf{Status} & Generated candidate & \\
    \hline
    \textbf{Ancestry} &  & \\
    \hline
    \textbf{Source} & Semantic-pattern candidate & \\
    \hline
    \textbf{Enabled By} &  & \\
    \hline
    \textbf{Supports} &  & \\
    \hline
    \textbf{Derived From} &  & \\
    \hline
    \textbf{Reviewed Status} & Generated candidate: keep / modify / relocate / reject; verify whether independently testable. & \\
    \hline
  \end{tabularx}
\end{table}

Intent. Context Establishment is proposed as a discrete, testable enabling behaviour for Product Batch \& Serial Tracking. It is included for architectural review, not as an established implementation obligation.

\noindent\textbf{Feature boundary.} This behaviour is scoped to Product Batch \& Serial Tracking; an identically named Feature elsewhere may carry different business meaning.\par
\featuretable{
--- & --- & \emph{Not yet captured} & \emph{Complete during review} \\
}
\subsubsection{Feature - State Monitoring}\label{cap:bec-c2e441ed}
\vspace{-1em}
\begin{table}[htbp]
  \centering
  \footnotesize
  \begin{tabularx}{1\textwidth}{|l|l|X|}
    \hline
    \textbf{Metadata Name} & \textbf{Metadata Value} & \textbf{Description} \\
    \hline
    \textbf{Reference Id} & \texttt{BEC-C2E441ED} & \\
    \hline
    \textbf{Status} & Generated candidate & \\
    \hline
    \textbf{Ancestry} &  & \\
    \hline
    \textbf{Source} & Semantic-pattern candidate & \\
    \hline
    \textbf{Enabled By} &  & \\
    \hline
    \textbf{Supports} &  & \\
    \hline
    \textbf{Derived From} &  & \\
    \hline
    \textbf{Reviewed Status} & Generated candidate: keep / modify / relocate / reject; verify whether independently testable. & \\
    \hline
  \end{tabularx}
\end{table}

Intent. State Monitoring is proposed as a discrete, testable enabling behaviour for Product Batch \& Serial Tracking. It is included for architectural review, not as an established implementation obligation.

\noindent\textbf{Feature boundary.} This behaviour is scoped to Product Batch \& Serial Tracking; an identically named Feature elsewhere may carry different business meaning.\par
\featuretable{
--- & --- & \emph{Not yet captured} & \emph{Complete during review} \\
}
\subsubsection{Feature - Exception Management}\label{cap:bec-db97a930}
\vspace{-1em}
\begin{table}[htbp]
  \centering
  \footnotesize
  \begin{tabularx}{1\textwidth}{|l|l|X|}
    \hline
    \textbf{Metadata Name} & \textbf{Metadata Value} & \textbf{Description} \\
    \hline
    \textbf{Reference Id} & \texttt{BEC-DB97A930} & \\
    \hline
    \textbf{Status} & Generated candidate & \\
    \hline
    \textbf{Ancestry} &  & \\
    \hline
    \textbf{Source} & Semantic-pattern candidate & \\
    \hline
    \textbf{Enabled By} &  & \\
    \hline
    \textbf{Supports} &  & \\
    \hline
    \textbf{Derived From} &  & \\
    \hline
    \textbf{Reviewed Status} & Generated candidate: keep / modify / relocate / reject; verify whether independently testable. & \\
    \hline
  \end{tabularx}
\end{table}

Intent. Exception Management is proposed as a discrete, testable enabling behaviour for Product Batch \& Serial Tracking. It is included for architectural review, not as an established implementation obligation.

\noindent\textbf{Feature boundary.} This behaviour is scoped to Product Batch \& Serial Tracking; an identically named Feature elsewhere may carry different business meaning.\par
\featuretable{
--- & --- & \emph{Not yet captured} & \emph{Complete during review} \\
}
\subsubsection{Feature - Information Retrieval}\label{cap:bec-4c113522}
\vspace{-1em}
\begin{table}[htbp]
  \centering
  \footnotesize
  \begin{tabularx}{1\textwidth}{|l|l|X|}
    \hline
    \textbf{Metadata Name} & \textbf{Metadata Value} & \textbf{Description} \\
    \hline
    \textbf{Reference Id} & \texttt{BEC-4C113522} & \\
    \hline
    \textbf{Status} & Generated candidate & \\
    \hline
    \textbf{Ancestry} &  & \\
    \hline
    \textbf{Source} & Semantic-pattern candidate & \\
    \hline
    \textbf{Enabled By} &  & \\
    \hline
    \textbf{Supports} &  & \\
    \hline
    \textbf{Derived From} &  & \\
    \hline
    \textbf{Reviewed Status} & Generated candidate: keep / modify / relocate / reject; verify whether independently testable. & \\
    \hline
  \end{tabularx}
\end{table}

Intent. Information Retrieval is proposed as a discrete, testable enabling behaviour for Product Batch \& Serial Tracking. It is included for architectural review, not as an established implementation obligation.

\noindent\textbf{Feature boundary.} This behaviour is scoped to Product Batch \& Serial Tracking; an identically named Feature elsewhere may carry different business meaning.\par
\featuretable{
--- & --- & \emph{Not yet captured} & \emph{Complete during review} \\
}

\section{CT2 - Clinical Device Administration}\label{cap:bec-39c0d38d}
\vspace{-0.5em}
\begin{table}[htbp]
  \centering
  \footnotesize
  \begin{tabularx}{1\textwidth}{|l|l|X|}
    \hline
    \multicolumn{3}{c}{\textbf{Capability Metadata}} \\
    \hline
    \textbf{Name} & \textbf{Value} & \textbf{Comment and/or Description} \\
    \hline
    \textbf{Reference Id} & \texttt{BEC-39C0D38D} & \\
    \hline
    \textbf{Status} & Candidate ancestry & \\
    \hline
    \textbf{Ancestry} &  & \\
    \hline
    \textbf{Source} &  & \\
    \hline
    \textbf{Enabled By} &  & \\
    \hline
    \textbf{Supports} &  & \\
    \hline
    \textbf{Derived From} &  & \\
    \hline
    \textbf{Reviewed Status} &  & \\
    \hline
  \end{tabularx}
\end{table}

\begin{table}[htbp]
  \centering
  \footnotesize
  \begin{tabularx}{1\textwidth}{|l|l|X|}
    \hline
    \multicolumn{3}{c}{\textbf{Contained Capabilities}} \\
    \hline
    \textbf{Capability Name} & \textbf{Capability Id} & \textbf{Comment} \\
    \hline
    Device Definition & \texttt{cap:bec-ef278d25} (at \ref{cap:bec-ef278d25}) & \\
    \hline
    Device Instance & \texttt{cap:bec-b3dedcd2} (at \ref{cap:bec-b3dedcd2}) & \\
    \hline
    Device Association & \texttt{cap:bec-868954a0} (at \ref{cap:bec-868954a0}) & \\
    \hline
    Device Connectivity & \texttt{cap:bec-6548f58a} (at \ref{cap:bec-6548f58a}) & \\
    \hline
  \end{tabularx}
\end{table}

Intent. Clinical Device Administration is a candidate enabling capability within Entity Information Enablement. It frames the reusable system-supported behaviours needed to support this healthcare context without prescribing a product, interface, technology or deployment boundary.

Boundary. This capability supports the business activity and does not inherit the authority or accountability of the Business Capability it enables. Its child Features are review candidates until ancestry and semantics are adjudicated.

\subsection{CT3 - Device Definition}\label{cap:bec-ef278d25}
\vspace{-0.5em}
\begin{table}[htbp]
  \centering
  \footnotesize
  \begin{tabularx}{1\textwidth}{|l|l|X|}
    \hline
    \multicolumn{3}{c}{\textbf{Capability Metadata}} \\
    \hline
    \textbf{Name} & \textbf{Value} & \textbf{Comment and/or Description} \\
    \hline
    \textbf{Reference Id} & \texttt{BEC-EF278D25} & \\
    \hline
    \textbf{Status} & Candidate ancestry & \\
    \hline
    \textbf{Ancestry} &  & \\
    \hline
    \textbf{Source} &  & \\
    \hline
    \textbf{Enabled By} &  & \\
    \hline
    \textbf{Supports} &  & \\
    \hline
    \textbf{Derived From} &  & \\
    \hline
    \textbf{Reviewed Status} &  & \\
    \hline
  \end{tabularx}
\end{table}

\begin{table}[htbp]
  \centering
  \footnotesize
  \begin{tabularx}{1\textwidth}{|l|l|X|}
    \hline
    \multicolumn{3}{c}{\textbf{Contained Features}} \\
    \hline
    \textbf{Feature Name} & \textbf{Feature Id} & \textbf{Comment} \\
    \hline
    Device Type Specification & \texttt{cap:bec-784e5062} (at \ref{cap:bec-784e5062}) & \\
    \hline
    Device Identity Management & \texttt{cap:bec-3117c7a3} (at \ref{cap:bec-3117c7a3}) & \\
    \hline
    Device State Management & \texttt{cap:bec-9f1fbf5c} (at \ref{cap:bec-9f1fbf5c}) & \\
    \hline
    Device Relationship Management & \texttt{cap:bec-018c4c3f} (at \ref{cap:bec-018c4c3f}) & \\
    \hline
    Device Provenance Tracking & \texttt{cap:bec-624e75ee} (at \ref{cap:bec-624e75ee}) & \\
    \hline
  \end{tabularx}
\end{table}

Intent. Device Definition is a candidate enabling capability within Clinical Device Administration. It frames the reusable system-supported behaviours needed to support this healthcare context without prescribing a product, interface, technology or deployment boundary.

Boundary. This capability supports the business activity and does not inherit the authority or accountability of the Business Capability it enables. Its child Features are review candidates until ancestry and semantics are adjudicated.

FHIR R5 conceptual analogues considered (not implementation mappings): Device, DeviceDefinition.

\subsubsection{Feature - Device Type Specification}\label{cap:bec-784e5062}
\vspace{-1em}
\begin{table}[htbp]
  \centering
  \footnotesize
  \begin{tabularx}{1\textwidth}{|l|l|X|}
    \hline
    \textbf{Metadata Name} & \textbf{Metadata Value} & \textbf{Description} \\
    \hline
    \textbf{Reference Id} & \texttt{BEC-784E5062} & \\
    \hline
    \textbf{Status} & Existing Feature; candidate CT3 placement & \\
    \hline
    \textbf{Ancestry} &  & \\
    \hline
    \textbf{Source} & FEAT-EM-28 & \\
    \hline
    \textbf{Enabled By} &  & \\
    \hline
    \textbf{Supports} &  & \\
    \hline
    \textbf{Derived From} &  & \\
    \hline
    \textbf{Reviewed Status} & Confirm existing Feature meaning and proposed CT3 ownership. & \\
    \hline
  \end{tabularx}
\end{table}

Intent. Device Type Specification provides a distinct enabling behaviour in the context of Device Definition. This is an existing named Feature; its placement under the candidate CT3 remains subject to review.

\noindent\textbf{Feature boundary.} This behaviour is scoped to Device Definition; an identically named Feature elsewhere may carry different business meaning.\par
\featuretable{
Function & \texttt{FN-979C3626} & Maintain Device Definition & Catalogue \\
Service & \texttt{SV-60A611FC} & Device Definition Query & Linked to preceding Function in source \\
}
\subsubsection{Feature - Device Identity Management}\label{cap:bec-3117c7a3}
\vspace{-1em}
\begin{table}[htbp]
  \centering
  \footnotesize
  \begin{tabularx}{1\textwidth}{|l|l|X|}
    \hline
    \textbf{Metadata Name} & \textbf{Metadata Value} & \textbf{Description} \\
    \hline
    \textbf{Reference Id} & \texttt{BEC-3117C7A3} & \\
    \hline
    \textbf{Status} & Generated candidate & \\
    \hline
    \textbf{Ancestry} &  & \\
    \hline
    \textbf{Source} & FHIR-informed candidate & \\
    \hline
    \textbf{Enabled By} &  & \\
    \hline
    \textbf{Supports} &  & \\
    \hline
    \textbf{Derived From} &  & \\
    \hline
    \textbf{Reviewed Status} & Generated candidate: keep / modify / relocate / reject; verify whether independently testable. & \\
    \hline
  \end{tabularx}
\end{table}

Intent. Device Identity Management is proposed to support the underlying healthcare concept represented by Device Definition. The proposal concerns behaviour and semantic responsibility, not a FHIR field, operation, resource instance or prescribed data model.

Conceptual evidence: HL7 FHIR R5 Device, DeviceDefinition.

\noindent\textbf{Feature boundary.} This behaviour is scoped to Device Definition; an identically named Feature elsewhere may carry different business meaning.\par
\featuretable{
--- & --- & \emph{Not yet captured} & \emph{Complete during review} \\
}
\subsubsection{Feature - Device State Management}\label{cap:bec-9f1fbf5c}
\vspace{-1em}
\begin{table}[htbp]
  \centering
  \footnotesize
  \begin{tabularx}{1\textwidth}{|l|l|X|}
    \hline
    \textbf{Metadata Name} & \textbf{Metadata Value} & \textbf{Description} \\
    \hline
    \textbf{Reference Id} & \texttt{BEC-9F1FBF5C} & \\
    \hline
    \textbf{Status} & Generated candidate & \\
    \hline
    \textbf{Ancestry} &  & \\
    \hline
    \textbf{Source} & FHIR-informed candidate & \\
    \hline
    \textbf{Enabled By} &  & \\
    \hline
    \textbf{Supports} &  & \\
    \hline
    \textbf{Derived From} &  & \\
    \hline
    \textbf{Reviewed Status} & Generated candidate: keep / modify / relocate / reject; verify whether independently testable. & \\
    \hline
  \end{tabularx}
\end{table}

Intent. Device State Management is proposed to support the underlying healthcare concept represented by Device Definition. The proposal concerns behaviour and semantic responsibility, not a FHIR field, operation, resource instance or prescribed data model.

Conceptual evidence: HL7 FHIR R5 Device, DeviceDefinition.

\noindent\textbf{Feature boundary.} This behaviour is scoped to Device Definition; an identically named Feature elsewhere may carry different business meaning.\par
\featuretable{
--- & --- & \emph{Not yet captured} & \emph{Complete during review} \\
}
\subsubsection{Feature - Device Relationship Management}\label{cap:bec-018c4c3f}
\vspace{-1em}
\begin{table}[htbp]
  \centering
  \footnotesize
  \begin{tabularx}{1\textwidth}{|l|l|X|}
    \hline
    \textbf{Metadata Name} & \textbf{Metadata Value} & \textbf{Description} \\
    \hline
    \textbf{Reference Id} & \texttt{BEC-018C4C3F} & \\
    \hline
    \textbf{Status} & Generated candidate & \\
    \hline
    \textbf{Ancestry} &  & \\
    \hline
    \textbf{Source} & FHIR-informed candidate & \\
    \hline
    \textbf{Enabled By} &  & \\
    \hline
    \textbf{Supports} &  & \\
    \hline
    \textbf{Derived From} &  & \\
    \hline
    \textbf{Reviewed Status} & Generated candidate: keep / modify / relocate / reject; verify whether independently testable. & \\
    \hline
  \end{tabularx}
\end{table}

Intent. Device Relationship Management is proposed to support the underlying healthcare concept represented by Device Definition. The proposal concerns behaviour and semantic responsibility, not a FHIR field, operation, resource instance or prescribed data model.

Conceptual evidence: HL7 FHIR R5 Device, DeviceDefinition.

\noindent\textbf{Feature boundary.} This behaviour is scoped to Device Definition; an identically named Feature elsewhere may carry different business meaning.\par
\featuretable{
--- & --- & \emph{Not yet captured} & \emph{Complete during review} \\
}
\subsubsection{Feature - Device Provenance Tracking}\label{cap:bec-624e75ee}
\vspace{-1em}
\begin{table}[htbp]
  \centering
  \footnotesize
  \begin{tabularx}{1\textwidth}{|l|l|X|}
    \hline
    \textbf{Metadata Name} & \textbf{Metadata Value} & \textbf{Description} \\
    \hline
    \textbf{Reference Id} & \texttt{BEC-624E75EE} & \\
    \hline
    \textbf{Status} & Generated candidate & \\
    \hline
    \textbf{Ancestry} &  & \\
    \hline
    \textbf{Source} & FHIR-informed candidate & \\
    \hline
    \textbf{Enabled By} &  & \\
    \hline
    \textbf{Supports} &  & \\
    \hline
    \textbf{Derived From} &  & \\
    \hline
    \textbf{Reviewed Status} & Generated candidate: keep / modify / relocate / reject; verify whether independently testable. & \\
    \hline
  \end{tabularx}
\end{table}

Intent. Device Provenance Tracking is proposed to support the underlying healthcare concept represented by Device Definition. The proposal concerns behaviour and semantic responsibility, not a FHIR field, operation, resource instance or prescribed data model.

Conceptual evidence: HL7 FHIR R5 Device, DeviceDefinition.

\noindent\textbf{Feature boundary.} This behaviour is scoped to Device Definition; an identically named Feature elsewhere may carry different business meaning.\par
\featuretable{
--- & --- & \emph{Not yet captured} & \emph{Complete during review} \\
}

\subsection{CT3 - Device Instance}\label{cap:bec-b3dedcd2}
\vspace{-0.5em}
\begin{table}[htbp]
  \centering
  \footnotesize
  \begin{tabularx}{1\textwidth}{|l|l|X|}
    \hline
    \multicolumn{3}{c}{\textbf{Capability Metadata}} \\
    \hline
    \textbf{Name} & \textbf{Value} & \textbf{Comment and/or Description} \\
    \hline
    \textbf{Reference Id} & \texttt{BEC-B3DEDCD2} & \\
    \hline
    \textbf{Status} & Candidate ancestry & \\
    \hline
    \textbf{Ancestry} &  & \\
    \hline
    \textbf{Source} &  & \\
    \hline
    \textbf{Enabled By} &  & \\
    \hline
    \textbf{Supports} &  & \\
    \hline
    \textbf{Derived From} &  & \\
    \hline
    \textbf{Reviewed Status} &  & \\
    \hline
  \end{tabularx}
\end{table}

\begin{table}[htbp]
  \centering
  \footnotesize
  \begin{tabularx}{1\textwidth}{|l|l|X|}
    \hline
    \multicolumn{3}{c}{\textbf{Contained Features}} \\
    \hline
    \textbf{Feature Name} & \textbf{Feature Id} & \textbf{Comment} \\
    \hline
    Device Serial Tracking & \texttt{cap:bec-8fb4909a} (at \ref{cap:bec-8fb4909a}) & \\
    \hline
    Device Identity Management & \texttt{cap:bec-a111f20f} (at \ref{cap:bec-a111f20f}) & \\
    \hline
    Device State Management & \texttt{cap:bec-8c9d3c83} (at \ref{cap:bec-8c9d3c83}) & \\
    \hline
    Device Relationship Management & \texttt{cap:bec-a8e9a77a} (at \ref{cap:bec-a8e9a77a}) & \\
    \hline
    Device Provenance Tracking & \texttt{cap:bec-eb9ae886} (at \ref{cap:bec-eb9ae886}) & \\
    \hline
  \end{tabularx}
\end{table}

Intent. Device Instance is a candidate enabling capability within Clinical Device Administration. It frames the reusable system-supported behaviours needed to support this healthcare context without prescribing a product, interface, technology or deployment boundary.

Boundary. This capability supports the business activity and does not inherit the authority or accountability of the Business Capability it enables. Its child Features are review candidates until ancestry and semantics are adjudicated.

FHIR R5 conceptual analogues considered (not implementation mappings): Device, DeviceDefinition.

\subsubsection{Feature - Device Serial Tracking}\label{cap:bec-8fb4909a}
\vspace{-1em}
\begin{table}[htbp]
  \centering
  \footnotesize
  \begin{tabularx}{1\textwidth}{|l|l|X|}
    \hline
    \textbf{Metadata Name} & \textbf{Metadata Value} & \textbf{Description} \\
    \hline
    \textbf{Reference Id} & \texttt{BEC-8FB4909A} & \\
    \hline
    \textbf{Status} & Existing Feature; candidate CT3 placement & \\
    \hline
    \textbf{Ancestry} &  & \\
    \hline
    \textbf{Source} & FEAT-EM-29 & \\
    \hline
    \textbf{Enabled By} &  & \\
    \hline
    \textbf{Supports} &  & \\
    \hline
    \textbf{Derived From} &  & \\
    \hline
    \textbf{Reviewed Status} & Confirm existing Feature meaning and proposed CT3 ownership. & \\
    \hline
  \end{tabularx}
\end{table}

Intent. Device Serial Tracking provides a distinct enabling behaviour in the context of Device Instance. This is an existing named Feature; its placement under the candidate CT3 remains subject to review.

\noindent\textbf{Feature boundary.} This behaviour is scoped to Device Instance; an identically named Feature elsewhere may carry different business meaning.\par
\featuretable{
Function & \texttt{FN-5A9A1E7C} & Maintain Device Instance Identity & Catalogue \\
Service & \texttt{SV-7FDB49E9} & Device Instance Resolution & Linked to preceding Function in source \\
}
\subsubsection{Feature - Device Identity Management}\label{cap:bec-a111f20f}
\vspace{-1em}
\begin{table}[htbp]
  \centering
  \footnotesize
  \begin{tabularx}{1\textwidth}{|l|l|X|}
    \hline
    \textbf{Metadata Name} & \textbf{Metadata Value} & \textbf{Description} \\
    \hline
    \textbf{Reference Id} & \texttt{BEC-A111F20F} & \\
    \hline
    \textbf{Status} & Generated candidate & \\
    \hline
    \textbf{Ancestry} &  & \\
    \hline
    \textbf{Source} & FHIR-informed candidate & \\
    \hline
    \textbf{Enabled By} &  & \\
    \hline
    \textbf{Supports} &  & \\
    \hline
    \textbf{Derived From} &  & \\
    \hline
    \textbf{Reviewed Status} & Generated candidate: keep / modify / relocate / reject; verify whether independently testable. & \\
    \hline
  \end{tabularx}
\end{table}

Intent. Device Identity Management is proposed to support the underlying healthcare concept represented by Device Instance. The proposal concerns behaviour and semantic responsibility, not a FHIR field, operation, resource instance or prescribed data model.

Conceptual evidence: HL7 FHIR R5 Device, DeviceDefinition.

\noindent\textbf{Feature boundary.} This behaviour is scoped to Device Instance; an identically named Feature elsewhere may carry different business meaning.\par
\featuretable{
--- & --- & \emph{Not yet captured} & \emph{Complete during review} \\
}
\subsubsection{Feature - Device State Management}\label{cap:bec-8c9d3c83}
\vspace{-1em}
\begin{table}[htbp]
  \centering
  \footnotesize
  \begin{tabularx}{1\textwidth}{|l|l|X|}
    \hline
    \textbf{Metadata Name} & \textbf{Metadata Value} & \textbf{Description} \\
    \hline
    \textbf{Reference Id} & \texttt{BEC-8C9D3C83} & \\
    \hline
    \textbf{Status} & Generated candidate & \\
    \hline
    \textbf{Ancestry} &  & \\
    \hline
    \textbf{Source} & FHIR-informed candidate & \\
    \hline
    \textbf{Enabled By} &  & \\
    \hline
    \textbf{Supports} &  & \\
    \hline
    \textbf{Derived From} &  & \\
    \hline
    \textbf{Reviewed Status} & Generated candidate: keep / modify / relocate / reject; verify whether independently testable. & \\
    \hline
  \end{tabularx}
\end{table}

Intent. Device State Management is proposed to support the underlying healthcare concept represented by Device Instance. The proposal concerns behaviour and semantic responsibility, not a FHIR field, operation, resource instance or prescribed data model.

Conceptual evidence: HL7 FHIR R5 Device, DeviceDefinition.

\noindent\textbf{Feature boundary.} This behaviour is scoped to Device Instance; an identically named Feature elsewhere may carry different business meaning.\par
\featuretable{
--- & --- & \emph{Not yet captured} & \emph{Complete during review} \\
}
\subsubsection{Feature - Device Relationship Management}\label{cap:bec-a8e9a77a}
\vspace{-1em}
\begin{table}[htbp]
  \centering
  \footnotesize
  \begin{tabularx}{1\textwidth}{|l|l|X|}
    \hline
    \textbf{Metadata Name} & \textbf{Metadata Value} & \textbf{Description} \\
    \hline
    \textbf{Reference Id} & \texttt{BEC-A8E9A77A} & \\
    \hline
    \textbf{Status} & Generated candidate & \\
    \hline
    \textbf{Ancestry} &  & \\
    \hline
    \textbf{Source} & FHIR-informed candidate & \\
    \hline
    \textbf{Enabled By} &  & \\
    \hline
    \textbf{Supports} &  & \\
    \hline
    \textbf{Derived From} &  & \\
    \hline
    \textbf{Reviewed Status} & Generated candidate: keep / modify / relocate / reject; verify whether independently testable. & \\
    \hline
  \end{tabularx}
\end{table}

Intent. Device Relationship Management is proposed to support the underlying healthcare concept represented by Device Instance. The proposal concerns behaviour and semantic responsibility, not a FHIR field, operation, resource instance or prescribed data model.

Conceptual evidence: HL7 FHIR R5 Device, DeviceDefinition.

\noindent\textbf{Feature boundary.} This behaviour is scoped to Device Instance; an identically named Feature elsewhere may carry different business meaning.\par
\featuretable{
--- & --- & \emph{Not yet captured} & \emph{Complete during review} \\
}
\subsubsection{Feature - Device Provenance Tracking}\label{cap:bec-eb9ae886}
\vspace{-1em}
\begin{table}[htbp]
  \centering
  \footnotesize
  \begin{tabularx}{1\textwidth}{|l|l|X|}
    \hline
    \textbf{Metadata Name} & \textbf{Metadata Value} & \textbf{Description} \\
    \hline
    \textbf{Reference Id} & \texttt{BEC-EB9AE886} & \\
    \hline
    \textbf{Status} & Generated candidate & \\
    \hline
    \textbf{Ancestry} &  & \\
    \hline
    \textbf{Source} & FHIR-informed candidate & \\
    \hline
    \textbf{Enabled By} &  & \\
    \hline
    \textbf{Supports} &  & \\
    \hline
    \textbf{Derived From} &  & \\
    \hline
    \textbf{Reviewed Status} & Generated candidate: keep / modify / relocate / reject; verify whether independently testable. & \\
    \hline
  \end{tabularx}
\end{table}

Intent. Device Provenance Tracking is proposed to support the underlying healthcare concept represented by Device Instance. The proposal concerns behaviour and semantic responsibility, not a FHIR field, operation, resource instance or prescribed data model.

Conceptual evidence: HL7 FHIR R5 Device, DeviceDefinition.

\noindent\textbf{Feature boundary.} This behaviour is scoped to Device Instance; an identically named Feature elsewhere may carry different business meaning.\par
\featuretable{
--- & --- & \emph{Not yet captured} & \emph{Complete during review} \\
}

\subsection{CT3 - Device Association}\label{cap:bec-868954a0}
\vspace{-0.5em}
\begin{table}[htbp]
  \centering
  \footnotesize
  \begin{tabularx}{1\textwidth}{|l|l|X|}
    \hline
    \multicolumn{3}{c}{\textbf{Capability Metadata}} \\
    \hline
    \textbf{Name} & \textbf{Value} & \textbf{Comment and/or Description} \\
    \hline
    \textbf{Reference Id} & \texttt{BEC-868954A0} & \\
    \hline
    \textbf{Status} & Candidate ancestry & \\
    \hline
    \textbf{Ancestry} &  & \\
    \hline
    \textbf{Source} &  & \\
    \hline
    \textbf{Enabled By} &  & \\
    \hline
    \textbf{Supports} &  & \\
    \hline
    \textbf{Derived From} &  & \\
    \hline
    \textbf{Reviewed Status} &  & \\
    \hline
  \end{tabularx}
\end{table}

\begin{table}[htbp]
  \centering
  \footnotesize
  \begin{tabularx}{1\textwidth}{|l|l|X|}
    \hline
    \multicolumn{3}{c}{\textbf{Contained Features}} \\
    \hline
    \textbf{Feature Name} & \textbf{Feature Id} & \textbf{Comment} \\
    \hline
    Device-Patient/Location Assignment & \texttt{cap:bec-22b624c5} (at \ref{cap:bec-22b624c5}) & \\
    \hline
    Device Identity Management & \texttt{cap:bec-dbe8dd8a} (at \ref{cap:bec-dbe8dd8a}) & \\
    \hline
    Device State Management & \texttt{cap:bec-d81970aa} (at \ref{cap:bec-d81970aa}) & \\
    \hline
    Device Relationship Management & \texttt{cap:bec-49a18e8c} (at \ref{cap:bec-49a18e8c}) & \\
    \hline
    Device Provenance Tracking & \texttt{cap:bec-fc971e86} (at \ref{cap:bec-fc971e86}) & \\
    \hline
  \end{tabularx}
\end{table}

Intent. Device Association is a candidate enabling capability within Clinical Device Administration. It frames the reusable system-supported behaviours needed to support this healthcare context without prescribing a product, interface, technology or deployment boundary.

Boundary. This capability supports the business activity and does not inherit the authority or accountability of the Business Capability it enables. Its child Features are review candidates until ancestry and semantics are adjudicated.

FHIR R5 conceptual analogues considered (not implementation mappings): Device, DeviceDefinition.

\subsubsection{Feature - Device-Patient/Location Assignment}\label{cap:bec-22b624c5}
\vspace{-1em}
\begin{table}[htbp]
  \centering
  \footnotesize
  \begin{tabularx}{1\textwidth}{|l|l|X|}
    \hline
    \textbf{Metadata Name} & \textbf{Metadata Value} & \textbf{Description} \\
    \hline
    \textbf{Reference Id} & \texttt{BEC-22B624C5} & \\
    \hline
    \textbf{Status} & Existing Feature; candidate CT3 placement & \\
    \hline
    \textbf{Ancestry} &  & \\
    \hline
    \textbf{Source} & FEAT-EM-30 & \\
    \hline
    \textbf{Enabled By} &  & \\
    \hline
    \textbf{Supports} &  & \\
    \hline
    \textbf{Derived From} &  & \\
    \hline
    \textbf{Reviewed Status} & Confirm existing Feature meaning and proposed CT3 ownership. & \\
    \hline
  \end{tabularx}
\end{table}

Intent. Device-Patient/Location Assignment provides a distinct enabling behaviour in the context of Device Association. This is an existing named Feature; its placement under the candidate CT3 remains subject to review.

\noindent\textbf{Feature boundary.} This behaviour is scoped to Device Association; an identically named Feature elsewhere may carry different business meaning.\par
\featuretable{
Function & \texttt{FN-4F6DB240} & Maintain Device Association & Catalogue \\
Service & \texttt{SV-C85EA8CC} & Device Association Query & Linked to preceding Function in source \\
}
\subsubsection{Feature - Device Identity Management}\label{cap:bec-dbe8dd8a}
\vspace{-1em}
\begin{table}[htbp]
  \centering
  \footnotesize
  \begin{tabularx}{1\textwidth}{|l|l|X|}
    \hline
    \textbf{Metadata Name} & \textbf{Metadata Value} & \textbf{Description} \\
    \hline
    \textbf{Reference Id} & \texttt{BEC-DBE8DD8A} & \\
    \hline
    \textbf{Status} & Generated candidate & \\
    \hline
    \textbf{Ancestry} &  & \\
    \hline
    \textbf{Source} & FHIR-informed candidate & \\
    \hline
    \textbf{Enabled By} &  & \\
    \hline
    \textbf{Supports} &  & \\
    \hline
    \textbf{Derived From} &  & \\
    \hline
    \textbf{Reviewed Status} & Generated candidate: keep / modify / relocate / reject; verify whether independently testable. & \\
    \hline
  \end{tabularx}
\end{table}

Intent. Device Identity Management is proposed to support the underlying healthcare concept represented by Device Association. The proposal concerns behaviour and semantic responsibility, not a FHIR field, operation, resource instance or prescribed data model.

Conceptual evidence: HL7 FHIR R5 Device, DeviceDefinition.

\noindent\textbf{Feature boundary.} This behaviour is scoped to Device Association; an identically named Feature elsewhere may carry different business meaning.\par
\featuretable{
--- & --- & \emph{Not yet captured} & \emph{Complete during review} \\
}
\subsubsection{Feature - Device State Management}\label{cap:bec-d81970aa}
\vspace{-1em}
\begin{table}[htbp]
  \centering
  \footnotesize
  \begin{tabularx}{1\textwidth}{|l|l|X|}
    \hline
    \textbf{Metadata Name} & \textbf{Metadata Value} & \textbf{Description} \\
    \hline
    \textbf{Reference Id} & \texttt{BEC-D81970AA} & \\
    \hline
    \textbf{Status} & Generated candidate & \\
    \hline
    \textbf{Ancestry} &  & \\
    \hline
    \textbf{Source} & FHIR-informed candidate & \\
    \hline
    \textbf{Enabled By} &  & \\
    \hline
    \textbf{Supports} &  & \\
    \hline
    \textbf{Derived From} &  & \\
    \hline
    \textbf{Reviewed Status} & Generated candidate: keep / modify / relocate / reject; verify whether independently testable. & \\
    \hline
  \end{tabularx}
\end{table}

Intent. Device State Management is proposed to support the underlying healthcare concept represented by Device Association. The proposal concerns behaviour and semantic responsibility, not a FHIR field, operation, resource instance or prescribed data model.

Conceptual evidence: HL7 FHIR R5 Device, DeviceDefinition.

\noindent\textbf{Feature boundary.} This behaviour is scoped to Device Association; an identically named Feature elsewhere may carry different business meaning.\par
\featuretable{
--- & --- & \emph{Not yet captured} & \emph{Complete during review} \\
}
\subsubsection{Feature - Device Relationship Management}\label{cap:bec-49a18e8c}
\vspace{-1em}
\begin{table}[htbp]
  \centering
  \footnotesize
  \begin{tabularx}{1\textwidth}{|l|l|X|}
    \hline
    \textbf{Metadata Name} & \textbf{Metadata Value} & \textbf{Description} \\
    \hline
    \textbf{Reference Id} & \texttt{BEC-49A18E8C} & \\
    \hline
    \textbf{Status} & Generated candidate & \\
    \hline
    \textbf{Ancestry} &  & \\
    \hline
    \textbf{Source} & FHIR-informed candidate & \\
    \hline
    \textbf{Enabled By} &  & \\
    \hline
    \textbf{Supports} &  & \\
    \hline
    \textbf{Derived From} &  & \\
    \hline
    \textbf{Reviewed Status} & Generated candidate: keep / modify / relocate / reject; verify whether independently testable. & \\
    \hline
  \end{tabularx}
\end{table}

Intent. Device Relationship Management is proposed to support the underlying healthcare concept represented by Device Association. The proposal concerns behaviour and semantic responsibility, not a FHIR field, operation, resource instance or prescribed data model.

Conceptual evidence: HL7 FHIR R5 Device, DeviceDefinition.

\noindent\textbf{Feature boundary.} This behaviour is scoped to Device Association; an identically named Feature elsewhere may carry different business meaning.\par
\featuretable{
--- & --- & \emph{Not yet captured} & \emph{Complete during review} \\
}
\subsubsection{Feature - Device Provenance Tracking}\label{cap:bec-fc971e86}
\vspace{-1em}
\begin{table}[htbp]
  \centering
  \footnotesize
  \begin{tabularx}{1\textwidth}{|l|l|X|}
    \hline
    \textbf{Metadata Name} & \textbf{Metadata Value} & \textbf{Description} \\
    \hline
    \textbf{Reference Id} & \texttt{BEC-FC971E86} & \\
    \hline
    \textbf{Status} & Generated candidate & \\
    \hline
    \textbf{Ancestry} &  & \\
    \hline
    \textbf{Source} & FHIR-informed candidate & \\
    \hline
    \textbf{Enabled By} &  & \\
    \hline
    \textbf{Supports} &  & \\
    \hline
    \textbf{Derived From} &  & \\
    \hline
    \textbf{Reviewed Status} & Generated candidate: keep / modify / relocate / reject; verify whether independently testable. & \\
    \hline
  \end{tabularx}
\end{table}

Intent. Device Provenance Tracking is proposed to support the underlying healthcare concept represented by Device Association. The proposal concerns behaviour and semantic responsibility, not a FHIR field, operation, resource instance or prescribed data model.

Conceptual evidence: HL7 FHIR R5 Device, DeviceDefinition.

\noindent\textbf{Feature boundary.} This behaviour is scoped to Device Association; an identically named Feature elsewhere may carry different business meaning.\par
\featuretable{
--- & --- & \emph{Not yet captured} & \emph{Complete during review} \\
}

\subsection{CT3 - Device Connectivity}\label{cap:bec-6548f58a}
\vspace{-0.5em}
\begin{table}[htbp]
  \centering
  \footnotesize
  \begin{tabularx}{1\textwidth}{|l|l|X|}
    \hline
    \multicolumn{3}{c}{\textbf{Capability Metadata}} \\
    \hline
    \textbf{Name} & \textbf{Value} & \textbf{Comment and/or Description} \\
    \hline
    \textbf{Reference Id} & \texttt{BEC-6548F58A} & \\
    \hline
    \textbf{Status} & Candidate ancestry & \\
    \hline
    \textbf{Ancestry} &  & \\
    \hline
    \textbf{Source} &  & \\
    \hline
    \textbf{Enabled By} &  & \\
    \hline
    \textbf{Supports} &  & \\
    \hline
    \textbf{Derived From} &  & \\
    \hline
    \textbf{Reviewed Status} &  & \\
    \hline
  \end{tabularx}
\end{table}

\begin{table}[htbp]
  \centering
  \footnotesize
  \begin{tabularx}{1\textwidth}{|l|l|X|}
    \hline
    \multicolumn{3}{c}{\textbf{Contained Features}} \\
    \hline
    \textbf{Feature Name} & \textbf{Feature Id} & \textbf{Comment} \\
    \hline
    Device Endpoint Registration & \texttt{cap:bec-45355492} (at \ref{cap:bec-45355492}) & \\
    \hline
    Device Identity Management & \texttt{cap:bec-d0a68c93} (at \ref{cap:bec-d0a68c93}) & \\
    \hline
    Device State Management & \texttt{cap:bec-b24562ce} (at \ref{cap:bec-b24562ce}) & \\
    \hline
    Device Relationship Management & \texttt{cap:bec-14a1b631} (at \ref{cap:bec-14a1b631}) & \\
    \hline
    Device Provenance Tracking & \texttt{cap:bec-41ede6c0} (at \ref{cap:bec-41ede6c0}) & \\
    \hline
  \end{tabularx}
\end{table}

Intent. Device Connectivity is a candidate enabling capability within Clinical Device Administration. It frames the reusable system-supported behaviours needed to support this healthcare context without prescribing a product, interface, technology or deployment boundary.

Boundary. This capability supports the business activity and does not inherit the authority or accountability of the Business Capability it enables. Its child Features are review candidates until ancestry and semantics are adjudicated.

FHIR R5 conceptual analogues considered (not implementation mappings): Device, DeviceDefinition.

\subsubsection{Feature - Device Endpoint Registration}\label{cap:bec-45355492}
\vspace{-1em}
\begin{table}[htbp]
  \centering
  \footnotesize
  \begin{tabularx}{1\textwidth}{|l|l|X|}
    \hline
    \textbf{Metadata Name} & \textbf{Metadata Value} & \textbf{Description} \\
    \hline
    \textbf{Reference Id} & \texttt{BEC-45355492} & \\
    \hline
    \textbf{Status} & Existing Feature; candidate CT3 placement & \\
    \hline
    \textbf{Ancestry} &  & \\
    \hline
    \textbf{Source} & FEAT-EM-31 & \\
    \hline
    \textbf{Enabled By} &  & \\
    \hline
    \textbf{Supports} &  & \\
    \hline
    \textbf{Derived From} &  & \\
    \hline
    \textbf{Reviewed Status} & Confirm existing Feature meaning and proposed CT3 ownership. & \\
    \hline
  \end{tabularx}
\end{table}

Intent. Device Endpoint Registration provides a distinct enabling behaviour in the context of Device Connectivity. This is an existing named Feature; its placement under the candidate CT3 remains subject to review.

\noindent\textbf{Feature boundary.} This behaviour is scoped to Device Connectivity; an identically named Feature elsewhere may carry different business meaning.\par
\featuretable{
Function & \texttt{FN-6B109517} & Maintain Device Communication Endpoints & Catalogue \\
Service & \texttt{SV-4728D31C} & Device Endpoint Resolution & Linked to preceding Function in source \\
}
\subsubsection{Feature - Device Identity Management}\label{cap:bec-d0a68c93}
\vspace{-1em}
\begin{table}[htbp]
  \centering
  \footnotesize
  \begin{tabularx}{1\textwidth}{|l|l|X|}
    \hline
    \textbf{Metadata Name} & \textbf{Metadata Value} & \textbf{Description} \\
    \hline
    \textbf{Reference Id} & \texttt{BEC-D0A68C93} & \\
    \hline
    \textbf{Status} & Generated candidate & \\
    \hline
    \textbf{Ancestry} &  & \\
    \hline
    \textbf{Source} & FHIR-informed candidate & \\
    \hline
    \textbf{Enabled By} &  & \\
    \hline
    \textbf{Supports} &  & \\
    \hline
    \textbf{Derived From} &  & \\
    \hline
    \textbf{Reviewed Status} & Generated candidate: keep / modify / relocate / reject; verify whether independently testable. & \\
    \hline
  \end{tabularx}
\end{table}

Intent. Device Identity Management is proposed to support the underlying healthcare concept represented by Device Connectivity. The proposal concerns behaviour and semantic responsibility, not a FHIR field, operation, resource instance or prescribed data model.

Conceptual evidence: HL7 FHIR R5 Device, DeviceDefinition.

\noindent\textbf{Feature boundary.} This behaviour is scoped to Device Connectivity; an identically named Feature elsewhere may carry different business meaning.\par
\featuretable{
--- & --- & \emph{Not yet captured} & \emph{Complete during review} \\
}
\subsubsection{Feature - Device State Management}\label{cap:bec-b24562ce}
\vspace{-1em}
\begin{table}[htbp]
  \centering
  \footnotesize
  \begin{tabularx}{1\textwidth}{|l|l|X|}
    \hline
    \textbf{Metadata Name} & \textbf{Metadata Value} & \textbf{Description} \\
    \hline
    \textbf{Reference Id} & \texttt{BEC-B24562CE} & \\
    \hline
    \textbf{Status} & Generated candidate & \\
    \hline
    \textbf{Ancestry} &  & \\
    \hline
    \textbf{Source} & FHIR-informed candidate & \\
    \hline
    \textbf{Enabled By} &  & \\
    \hline
    \textbf{Supports} &  & \\
    \hline
    \textbf{Derived From} &  & \\
    \hline
    \textbf{Reviewed Status} & Generated candidate: keep / modify / relocate / reject; verify whether independently testable. & \\
    \hline
  \end{tabularx}
\end{table}

Intent. Device State Management is proposed to support the underlying healthcare concept represented by Device Connectivity. The proposal concerns behaviour and semantic responsibility, not a FHIR field, operation, resource instance or prescribed data model.

Conceptual evidence: HL7 FHIR R5 Device, DeviceDefinition.

\noindent\textbf{Feature boundary.} This behaviour is scoped to Device Connectivity; an identically named Feature elsewhere may carry different business meaning.\par
\featuretable{
--- & --- & \emph{Not yet captured} & \emph{Complete during review} \\
}
\subsubsection{Feature - Device Relationship Management}\label{cap:bec-14a1b631}
\vspace{-1em}
\begin{table}[htbp]
  \centering
  \footnotesize
  \begin{tabularx}{1\textwidth}{|l|l|X|}
    \hline
    \textbf{Metadata Name} & \textbf{Metadata Value} & \textbf{Description} \\
    \hline
    \textbf{Reference Id} & \texttt{BEC-14A1B631} & \\
    \hline
    \textbf{Status} & Generated candidate & \\
    \hline
    \textbf{Ancestry} &  & \\
    \hline
    \textbf{Source} & FHIR-informed candidate & \\
    \hline
    \textbf{Enabled By} &  & \\
    \hline
    \textbf{Supports} &  & \\
    \hline
    \textbf{Derived From} &  & \\
    \hline
    \textbf{Reviewed Status} & Generated candidate: keep / modify / relocate / reject; verify whether independently testable. & \\
    \hline
  \end{tabularx}
\end{table}

Intent. Device Relationship Management is proposed to support the underlying healthcare concept represented by Device Connectivity. The proposal concerns behaviour and semantic responsibility, not a FHIR field, operation, resource instance or prescribed data model.

Conceptual evidence: HL7 FHIR R5 Device, DeviceDefinition.

\noindent\textbf{Feature boundary.} This behaviour is scoped to Device Connectivity; an identically named Feature elsewhere may carry different business meaning.\par
\featuretable{
--- & --- & \emph{Not yet captured} & \emph{Complete during review} \\
}
\subsubsection{Feature - Device Provenance Tracking}\label{cap:bec-41ede6c0}
\vspace{-1em}
\begin{table}[htbp]
  \centering
  \footnotesize
  \begin{tabularx}{1\textwidth}{|l|l|X|}
    \hline
    \textbf{Metadata Name} & \textbf{Metadata Value} & \textbf{Description} \\
    \hline
    \textbf{Reference Id} & \texttt{BEC-41EDE6C0} & \\
    \hline
    \textbf{Status} & Generated candidate & \\
    \hline
    \textbf{Ancestry} &  & \\
    \hline
    \textbf{Source} & FHIR-informed candidate & \\
    \hline
    \textbf{Enabled By} &  & \\
    \hline
    \textbf{Supports} &  & \\
    \hline
    \textbf{Derived From} &  & \\
    \hline
    \textbf{Reviewed Status} & Generated candidate: keep / modify / relocate / reject; verify whether independently testable. & \\
    \hline
  \end{tabularx}
\end{table}

Intent. Device Provenance Tracking is proposed to support the underlying healthcare concept represented by Device Connectivity. The proposal concerns behaviour and semantic responsibility, not a FHIR field, operation, resource instance or prescribed data model.

Conceptual evidence: HL7 FHIR R5 Device, DeviceDefinition.

\noindent\textbf{Feature boundary.} This behaviour is scoped to Device Connectivity; an identically named Feature elsewhere may carry different business meaning.\par
\featuretable{
--- & --- & \emph{Not yet captured} & \emph{Complete during review} \\
}

\chapter{CT1 - Care Administration Enablement}\label{cap:bec-ff45e08f}
\noindent\textbf{Reference ID:} \texttt{BEC-FF45E08F}\quad \textbf{Status:} Candidate ancestry\par
Intent. Care Administration Enablement is a candidate enabling capability within Business Enabling Capability layer. It frames the reusable system-supported behaviours needed to support this healthcare context without prescribing a product, interface, technology or deployment boundary.

Boundary. This capability supports the business activity and does not inherit the authority or accountability of the Business Capability it enables. Its child Features are review candidates until ancestry and semantics are adjudicated.

\section{CT2 - Client-Centred Administration}\label{cap:bec-dd7184df}
\vspace{-0.5em}
\begin{table}[htbp]
  \centering
  \footnotesize
  \begin{tabularx}{1\textwidth}{|l|l|X|}
    \hline
    \multicolumn{3}{c}{\textbf{Capability Metadata}} \\
    \hline
    \textbf{Name} & \textbf{Value} & \textbf{Comment and/or Description} \\
    \hline
    \textbf{Reference Id} & \texttt{BEC-DD7184DF} & \\
    \hline
    \textbf{Status} & Candidate ancestry & \\
    \hline
    \textbf{Ancestry} &  & \\
    \hline
    \textbf{Source} &  & \\
    \hline
    \textbf{Enabled By} &  & \\
    \hline
    \textbf{Supports} &  & \\
    \hline
    \textbf{Derived From} &  & \\
    \hline
    \textbf{Reviewed Status} &  & \\
    \hline
  \end{tabularx}
\end{table}

\begin{table}[htbp]
  \centering
  \footnotesize
  \begin{tabularx}{1\textwidth}{|l|l|X|}
    \hline
    \multicolumn{3}{c}{\textbf{Contained Capabilities}} \\
    \hline
    \textbf{Capability Name} & \textbf{Capability Id} & \textbf{Comment} \\
    \hline
    Client Intake \& Demographic Update & \texttt{cap:bec-25db15ba} (at \ref{cap:bec-25db15ba}) & \\
    \hline
  \end{tabularx}
\end{table}

Intent. Client-Centred Administration is a candidate enabling capability within Care Administration Enablement. It frames the reusable system-supported behaviours needed to support this healthcare context without prescribing a product, interface, technology or deployment boundary.

Boundary. This capability supports the business activity and does not inherit the authority or accountability of the Business Capability it enables. Its child Features are review candidates until ancestry and semantics are adjudicated.

\subsection{CT3 - Client Intake \& Demographic Update}\label{cap:bec-25db15ba}
\vspace{-0.5em}
\begin{table}[htbp]
  \centering
  \footnotesize
  \begin{tabularx}{1\textwidth}{|l|l|X|}
    \hline
    \multicolumn{3}{c}{\textbf{Capability Metadata}} \\
    \hline
    \textbf{Name} & \textbf{Value} & \textbf{Comment and/or Description} \\
    \hline
    \textbf{Reference Id} & \texttt{BEC-25DB15BA} & \\
    \hline
    \textbf{Status} & Candidate ancestry & \\
    \hline
    \textbf{Ancestry} &  & \\
    \hline
    \textbf{Source} &  & \\
    \hline
    \textbf{Enabled By} &  & \\
    \hline
    \textbf{Supports} &  & \\
    \hline
    \textbf{Derived From} &  & \\
    \hline
    \textbf{Reviewed Status} &  & \\
    \hline
  \end{tabularx}
\end{table}

\begin{table}[htbp]
  \centering
  \footnotesize
  \begin{tabularx}{1\textwidth}{|l|l|X|}
    \hline
    \multicolumn{3}{c}{\textbf{Contained Features}} \\
    \hline
    \textbf{Feature Name} & \textbf{Feature Id} & \textbf{Comment} \\
    \hline
    Demographic Update Ingestion & \texttt{cap:bec-4ab59d5a} (at \ref{cap:bec-4ab59d5a}) & \\
    \hline
    Demographic Change Distribution & \texttt{cap:bec-bdccd5ed} (at \ref{cap:bec-bdccd5ed}) & \\
    \hline
    Context Establishment & \texttt{cap:bec-163e98ad} (at \ref{cap:bec-163e98ad}) & \\
    \hline
    State Monitoring & \texttt{cap:bec-80dff5d9} (at \ref{cap:bec-80dff5d9}) & \\
    \hline
    Exception Management & \texttt{cap:bec-f29d899b} (at \ref{cap:bec-f29d899b}) & \\
    \hline
    Information Retrieval & \texttt{cap:bec-d55a3518} (at \ref{cap:bec-d55a3518}) & \\
    \hline
  \end{tabularx}
\end{table}

Intent. Client Intake \& Demographic Update is a candidate enabling capability within Client-Centred Administration. It frames the reusable system-supported behaviours needed to support this healthcare context without prescribing a product, interface, technology or deployment boundary.

Boundary. This capability supports the business activity and does not inherit the authority or accountability of the Business Capability it enables. Its child Features are review candidates until ancestry and semantics are adjudicated.

FHIR R5 conceptual analogues considered (not implementation mappings): Patient, Person, RelatedPerson.

\subsubsection{Feature - Demographic Update Ingestion}\label{cap:bec-4ab59d5a}
\vspace{-1em}
\begin{table}[htbp]
  \centering
  \footnotesize
  \begin{tabularx}{1\textwidth}{|l|l|X|}
    \hline
    \textbf{Metadata Name} & \textbf{Metadata Value} & \textbf{Description} \\
    \hline
    \textbf{Reference Id} & \texttt{BEC-4AB59D5A} & \\
    \hline
    \textbf{Status} & Existing Feature; candidate CT3 placement & \\
    \hline
    \textbf{Ancestry} &  & \\
    \hline
    \textbf{Source} & FEAT-SA-01 & \\
    \hline
    \textbf{Enabled By} &  & \\
    \hline
    \textbf{Supports} &  & \\
    \hline
    \textbf{Derived From} &  & \\
    \hline
    \textbf{Reviewed Status} & Confirm existing Feature meaning and proposed CT3 ownership. & \\
    \hline
  \end{tabularx}
\end{table}

Intent. Demographic Update Ingestion provides a distinct enabling behaviour in the context of Client Intake \& Demographic Update. This is an existing named Feature; its placement under the candidate CT3 remains subject to review.

\noindent\textbf{Feature boundary.} This behaviour is scoped to Client Intake \& Demographic Update; an identically named Feature elsewhere may carry different business meaning.\par
\featuretable{
Function & \texttt{FN-1118B1A2} & Receive Client Demographic Change & Catalogue \\
Service & \texttt{SV-6685441E} & Demographic Intake Service & Linked to preceding Function in source \\
}
\subsubsection{Feature - Demographic Change Distribution}\label{cap:bec-bdccd5ed}
\vspace{-1em}
\begin{table}[htbp]
  \centering
  \footnotesize
  \begin{tabularx}{1\textwidth}{|l|l|X|}
    \hline
    \textbf{Metadata Name} & \textbf{Metadata Value} & \textbf{Description} \\
    \hline
    \textbf{Reference Id} & \texttt{BEC-BDCCD5ED} & \\
    \hline
    \textbf{Status} & Existing Feature; candidate CT3 placement & \\
    \hline
    \textbf{Ancestry} &  & \\
    \hline
    \textbf{Source} & FEAT-SA-02 & \\
    \hline
    \textbf{Enabled By} &  & \\
    \hline
    \textbf{Supports} &  & \\
    \hline
    \textbf{Derived From} &  & \\
    \hline
    \textbf{Reviewed Status} & Confirm existing Feature meaning and proposed CT3 ownership. & \\
    \hline
  \end{tabularx}
\end{table}

Intent. Demographic Change Distribution provides a distinct enabling behaviour in the context of Client Intake \& Demographic Update. This is an existing named Feature; its placement under the candidate CT3 remains subject to review.

\noindent\textbf{Feature boundary.} This behaviour is scoped to Client Intake \& Demographic Update; an identically named Feature elsewhere may carry different business meaning.\par
\featuretable{
Function & \texttt{FN-976703A0} & Distribute Governed Demographic Change & Catalogue \\
Service & \texttt{SV-D38D4C72} & Demographic Update Notification & Linked to preceding Function in source \\
}
\subsubsection{Feature - Context Establishment}\label{cap:bec-163e98ad}
\vspace{-1em}
\begin{table}[htbp]
  \centering
  \footnotesize
  \begin{tabularx}{1\textwidth}{|l|l|X|}
    \hline
    \textbf{Metadata Name} & \textbf{Metadata Value} & \textbf{Description} \\
    \hline
    \textbf{Reference Id} & \texttt{BEC-163E98AD} & \\
    \hline
    \textbf{Status} & Generated candidate & \\
    \hline
    \textbf{Ancestry} &  & \\
    \hline
    \textbf{Source} & Semantic-pattern candidate & \\
    \hline
    \textbf{Enabled By} &  & \\
    \hline
    \textbf{Supports} &  & \\
    \hline
    \textbf{Derived From} &  & \\
    \hline
    \textbf{Reviewed Status} & Generated candidate: keep / modify / relocate / reject; verify whether independently testable. & \\
    \hline
  \end{tabularx}
\end{table}

Intent. Context Establishment is proposed as a discrete, testable enabling behaviour for Client Intake \& Demographic Update. It is included for architectural review, not as an established implementation obligation.

\noindent\textbf{Feature boundary.} This behaviour is scoped to Client Intake \& Demographic Update; an identically named Feature elsewhere may carry different business meaning.\par
\featuretable{
--- & --- & \emph{Not yet captured} & \emph{Complete during review} \\
}
\subsubsection{Feature - State Monitoring}\label{cap:bec-80dff5d9}
\vspace{-1em}
\begin{table}[htbp]
  \centering
  \footnotesize
  \begin{tabularx}{1\textwidth}{|l|l|X|}
    \hline
    \textbf{Metadata Name} & \textbf{Metadata Value} & \textbf{Description} \\
    \hline
    \textbf{Reference Id} & \texttt{BEC-80DFF5D9} & \\
    \hline
    \textbf{Status} & Generated candidate & \\
    \hline
    \textbf{Ancestry} &  & \\
    \hline
    \textbf{Source} & Semantic-pattern candidate & \\
    \hline
    \textbf{Enabled By} &  & \\
    \hline
    \textbf{Supports} &  & \\
    \hline
    \textbf{Derived From} &  & \\
    \hline
    \textbf{Reviewed Status} & Generated candidate: keep / modify / relocate / reject; verify whether independently testable. & \\
    \hline
  \end{tabularx}
\end{table}

Intent. State Monitoring is proposed as a discrete, testable enabling behaviour for Client Intake \& Demographic Update. It is included for architectural review, not as an established implementation obligation.

\noindent\textbf{Feature boundary.} This behaviour is scoped to Client Intake \& Demographic Update; an identically named Feature elsewhere may carry different business meaning.\par
\featuretable{
--- & --- & \emph{Not yet captured} & \emph{Complete during review} \\
}
\subsubsection{Feature - Exception Management}\label{cap:bec-f29d899b}
\vspace{-1em}
\begin{table}[htbp]
  \centering
  \footnotesize
  \begin{tabularx}{1\textwidth}{|l|l|X|}
    \hline
    \textbf{Metadata Name} & \textbf{Metadata Value} & \textbf{Description} \\
    \hline
    \textbf{Reference Id} & \texttt{BEC-F29D899B} & \\
    \hline
    \textbf{Status} & Generated candidate & \\
    \hline
    \textbf{Ancestry} &  & \\
    \hline
    \textbf{Source} & Semantic-pattern candidate & \\
    \hline
    \textbf{Enabled By} &  & \\
    \hline
    \textbf{Supports} &  & \\
    \hline
    \textbf{Derived From} &  & \\
    \hline
    \textbf{Reviewed Status} & Generated candidate: keep / modify / relocate / reject; verify whether independently testable. & \\
    \hline
  \end{tabularx}
\end{table}

Intent. Exception Management is proposed as a discrete, testable enabling behaviour for Client Intake \& Demographic Update. It is included for architectural review, not as an established implementation obligation.

\noindent\textbf{Feature boundary.} This behaviour is scoped to Client Intake \& Demographic Update; an identically named Feature elsewhere may carry different business meaning.\par
\featuretable{
--- & --- & \emph{Not yet captured} & \emph{Complete during review} \\
}
\subsubsection{Feature - Information Retrieval}\label{cap:bec-d55a3518}
\vspace{-1em}
\begin{table}[htbp]
  \centering
  \footnotesize
  \begin{tabularx}{1\textwidth}{|l|l|X|}
    \hline
    \textbf{Metadata Name} & \textbf{Metadata Value} & \textbf{Description} \\
    \hline
    \textbf{Reference Id} & \texttt{BEC-D55A3518} & \\
    \hline
    \textbf{Status} & Generated candidate & \\
    \hline
    \textbf{Ancestry} &  & \\
    \hline
    \textbf{Source} & Semantic-pattern candidate & \\
    \hline
    \textbf{Enabled By} &  & \\
    \hline
    \textbf{Supports} &  & \\
    \hline
    \textbf{Derived From} &  & \\
    \hline
    \textbf{Reviewed Status} & Generated candidate: keep / modify / relocate / reject; verify whether independently testable. & \\
    \hline
  \end{tabularx}
\end{table}

Intent. Information Retrieval is proposed as a discrete, testable enabling behaviour for Client Intake \& Demographic Update. It is included for architectural review, not as an established implementation obligation.

\noindent\textbf{Feature boundary.} This behaviour is scoped to Client Intake \& Demographic Update; an identically named Feature elsewhere may carry different business meaning.\par
\featuretable{
--- & --- & \emph{Not yet captured} & \emph{Complete during review} \\
}

\section{CT2 - Referral Administration}\label{cap:bec-c16b2202}
\vspace{-0.5em}
\begin{table}[htbp]
  \centering
  \footnotesize
  \begin{tabularx}{1\textwidth}{|l|l|X|}
    \hline
    \multicolumn{3}{c}{\textbf{Capability Metadata}} \\
    \hline
    \textbf{Name} & \textbf{Value} & \textbf{Comment and/or Description} \\
    \hline
    \textbf{Reference Id} & \texttt{BEC-C16B2202} & \\
    \hline
    \textbf{Status} & Candidate ancestry & \\
    \hline
    \textbf{Ancestry} &  & \\
    \hline
    \textbf{Source} &  & \\
    \hline
    \textbf{Enabled By} &  & \\
    \hline
    \textbf{Supports} &  & \\
    \hline
    \textbf{Derived From} &  & \\
    \hline
    \textbf{Reviewed Status} &  & \\
    \hline
  \end{tabularx}
\end{table}

\begin{table}[htbp]
  \centering
  \footnotesize
  \begin{tabularx}{1\textwidth}{|l|l|X|}
    \hline
    \multicolumn{3}{c}{\textbf{Contained Capabilities}} \\
    \hline
    \textbf{Capability Name} & \textbf{Capability Id} & \textbf{Comment} \\
    \hline
    Referral Intake \& Progression & \texttt{cap:bec-cd3f820d} (at \ref{cap:bec-cd3f820d}) & \\
    \hline
  \end{tabularx}
\end{table}

Intent. Referral Administration is a candidate enabling capability within Care Administration Enablement. It frames the reusable system-supported behaviours needed to support this healthcare context without prescribing a product, interface, technology or deployment boundary.

Boundary. This capability supports the business activity and does not inherit the authority or accountability of the Business Capability it enables. Its child Features are review candidates until ancestry and semantics are adjudicated.

\subsection{CT3 - Referral Intake \& Progression}\label{cap:bec-cd3f820d}
\vspace{-0.5em}
\begin{table}[htbp]
  \centering
  \footnotesize
  \begin{tabularx}{1\textwidth}{|l|l|X|}
    \hline
    \multicolumn{3}{c}{\textbf{Capability Metadata}} \\
    \hline
    \textbf{Name} & \textbf{Value} & \textbf{Comment and/or Description} \\
    \hline
    \textbf{Reference Id} & \texttt{BEC-CD3F820D} & \\
    \hline
    \textbf{Status} & Candidate ancestry & \\
    \hline
    \textbf{Ancestry} &  & \\
    \hline
    \textbf{Source} &  & \\
    \hline
    \textbf{Enabled By} &  & \\
    \hline
    \textbf{Supports} &  & \\
    \hline
    \textbf{Derived From} &  & \\
    \hline
    \textbf{Reviewed Status} &  & \\
    \hline
  \end{tabularx}
\end{table}

\begin{table}[htbp]
  \centering
  \footnotesize
  \begin{tabularx}{1\textwidth}{|l|l|X|}
    \hline
    \multicolumn{3}{c}{\textbf{Contained Features}} \\
    \hline
    \textbf{Feature Name} & \textbf{Feature Id} & \textbf{Comment} \\
    \hline
    Referral Document Ingestion & \texttt{cap:bec-8e9ad826} (at \ref{cap:bec-8e9ad826}) & \\
    \hline
    Referral Supporting Information & \texttt{cap:bec-2c8f0bd3} (at \ref{cap:bec-2c8f0bd3}) & \\
    \hline
    Referral Status \& Outcome Tracking & \texttt{cap:bec-6e903b7f} (at \ref{cap:bec-6e903b7f}) & \\
    \hline
    Referral Lifecycle Management & \texttt{cap:bec-7c40f1c9} (at \ref{cap:bec-7c40f1c9}) & \\
    \hline
    Referral Routing & \texttt{cap:bec-a551e3f4} (at \ref{cap:bec-a551e3f4}) & \\
    \hline
    Referral Status Tracking & \texttt{cap:bec-c4e89e34} (at \ref{cap:bec-c4e89e34}) & \\
    \hline
    Referral Outcome Association & \texttt{cap:bec-57720bc2} (at \ref{cap:bec-57720bc2}) & \\
    \hline
  \end{tabularx}
\end{table}

Intent. Referral Intake \& Progression is a candidate enabling capability within Referral Administration. It frames the reusable system-supported behaviours needed to support this healthcare context without prescribing a product, interface, technology or deployment boundary.

Boundary. This capability supports the business activity and does not inherit the authority or accountability of the Business Capability it enables. Its child Features are review candidates until ancestry and semantics are adjudicated.

FHIR R5 conceptual analogues considered (not implementation mappings): ServiceRequest, Task.

\subsubsection{Feature - Referral Document Ingestion}\label{cap:bec-8e9ad826}
\vspace{-1em}
\begin{table}[htbp]
  \centering
  \footnotesize
  \begin{tabularx}{1\textwidth}{|l|l|X|}
    \hline
    \textbf{Metadata Name} & \textbf{Metadata Value} & \textbf{Description} \\
    \hline
    \textbf{Reference Id} & \texttt{BEC-8E9AD826} & \\
    \hline
    \textbf{Status} & Existing Feature; candidate CT3 placement & \\
    \hline
    \textbf{Ancestry} &  & \\
    \hline
    \textbf{Source} & FEAT-SA-03 & \\
    \hline
    \textbf{Enabled By} &  & \\
    \hline
    \textbf{Supports} &  & \\
    \hline
    \textbf{Derived From} &  & \\
    \hline
    \textbf{Reviewed Status} & Confirm existing Feature meaning and proposed CT3 ownership. & \\
    \hline
  \end{tabularx}
\end{table}

Intent. Referral Document Ingestion provides a distinct enabling behaviour in the context of Referral Intake \& Progression. This is an existing named Feature; its placement under the candidate CT3 remains subject to review.

\noindent\textbf{Feature boundary.} This behaviour is scoped to Referral Intake \& Progression; an identically named Feature elsewhere may carry different business meaning.\par
\featuretable{
Function & \texttt{FN-38B78956} & Receive Referral & Catalogue \\
Service & \texttt{SV-BBAA13B9} & Referral Submission Service & Linked to preceding Function in source \\
}
\subsubsection{Feature - Referral Supporting Information}\label{cap:bec-2c8f0bd3}
\vspace{-1em}
\begin{table}[htbp]
  \centering
  \footnotesize
  \begin{tabularx}{1\textwidth}{|l|l|X|}
    \hline
    \textbf{Metadata Name} & \textbf{Metadata Value} & \textbf{Description} \\
    \hline
    \textbf{Reference Id} & \texttt{BEC-2C8F0BD3} & \\
    \hline
    \textbf{Status} & Existing Feature; candidate CT3 placement & \\
    \hline
    \textbf{Ancestry} &  & \\
    \hline
    \textbf{Source} & FEAT-SA-04 & \\
    \hline
    \textbf{Enabled By} &  & \\
    \hline
    \textbf{Supports} &  & \\
    \hline
    \textbf{Derived From} &  & \\
    \hline
    \textbf{Reviewed Status} & Confirm existing Feature meaning and proposed CT3 ownership. & \\
    \hline
  \end{tabularx}
\end{table}

Intent. Referral Supporting Information provides a distinct enabling behaviour in the context of Referral Intake \& Progression. This is an existing named Feature; its placement under the candidate CT3 remains subject to review.

\noindent\textbf{Feature boundary.} This behaviour is scoped to Referral Intake \& Progression; an identically named Feature elsewhere may carry different business meaning.\par
\featuretable{
Function & \texttt{FN-49C4407E} & Assemble Referral Context & Catalogue \\
Service & \texttt{SV-2AC745CC} & Referral Context Query & Linked to preceding Function in source \\
}
\subsubsection{Feature - Referral Status \& Outcome Tracking}\label{cap:bec-6e903b7f}
\vspace{-1em}
\begin{table}[htbp]
  \centering
  \footnotesize
  \begin{tabularx}{1\textwidth}{|l|l|X|}
    \hline
    \textbf{Metadata Name} & \textbf{Metadata Value} & \textbf{Description} \\
    \hline
    \textbf{Reference Id} & \texttt{BEC-6E903B7F} & \\
    \hline
    \textbf{Status} & Existing Feature; candidate CT3 placement & \\
    \hline
    \textbf{Ancestry} &  & \\
    \hline
    \textbf{Source} & FEAT-SA-05 & \\
    \hline
    \textbf{Enabled By} &  & \\
    \hline
    \textbf{Supports} &  & \\
    \hline
    \textbf{Derived From} &  & \\
    \hline
    \textbf{Reviewed Status} & Confirm existing Feature meaning and proposed CT3 ownership. & \\
    \hline
  \end{tabularx}
\end{table}

Intent. Referral Status \& Outcome Tracking provides a distinct enabling behaviour in the context of Referral Intake \& Progression. This is an existing named Feature; its placement under the candidate CT3 remains subject to review.

\noindent\textbf{Feature boundary.} This behaviour is scoped to Referral Intake \& Progression; an identically named Feature elsewhere may carry different business meaning.\par
\featuretable{
Function & \texttt{FN-764AF30E} & Manage Referral Progression & Catalogue \\
Service & \texttt{SV-3AE258E9} & Referral Status \& Outcome Service & Linked to preceding Function in source \\
Process & \texttt{PR-1E2327B8} & Referral Progression Process & Linked to preceding Function in source \\
}
\subsubsection{Feature - Referral Lifecycle Management}\label{cap:bec-7c40f1c9}
\vspace{-1em}
\begin{table}[htbp]
  \centering
  \footnotesize
  \begin{tabularx}{1\textwidth}{|l|l|X|}
    \hline
    \textbf{Metadata Name} & \textbf{Metadata Value} & \textbf{Description} \\
    \hline
    \textbf{Reference Id} & \texttt{BEC-7C40F1C9} & \\
    \hline
    \textbf{Status} & Generated candidate & \\
    \hline
    \textbf{Ancestry} &  & \\
    \hline
    \textbf{Source} & FHIR-informed candidate & \\
    \hline
    \textbf{Enabled By} &  & \\
    \hline
    \textbf{Supports} &  & \\
    \hline
    \textbf{Derived From} &  & \\
    \hline
    \textbf{Reviewed Status} & Generated candidate: keep / modify / relocate / reject; verify whether independently testable. & \\
    \hline
  \end{tabularx}
\end{table}

Intent. Referral Lifecycle Management is proposed to support the underlying healthcare concept represented by Referral Intake \& Progression. The proposal concerns behaviour and semantic responsibility, not a FHIR field, operation, resource instance or prescribed data model.

Conceptual evidence: HL7 FHIR R5 ServiceRequest, Task.

\noindent\textbf{Feature boundary.} This behaviour is scoped to Referral Intake \& Progression; an identically named Feature elsewhere may carry different business meaning.\par
\featuretable{
--- & --- & \emph{Not yet captured} & \emph{Complete during review} \\
}
\subsubsection{Feature - Referral Routing}\label{cap:bec-a551e3f4}
\vspace{-1em}
\begin{table}[htbp]
  \centering
  \footnotesize
  \begin{tabularx}{1\textwidth}{|l|l|X|}
    \hline
    \textbf{Metadata Name} & \textbf{Metadata Value} & \textbf{Description} \\
    \hline
    \textbf{Reference Id} & \texttt{BEC-A551E3F4} & \\
    \hline
    \textbf{Status} & Generated candidate & \\
    \hline
    \textbf{Ancestry} &  & \\
    \hline
    \textbf{Source} & FHIR-informed candidate & \\
    \hline
    \textbf{Enabled By} &  & \\
    \hline
    \textbf{Supports} &  & \\
    \hline
    \textbf{Derived From} &  & \\
    \hline
    \textbf{Reviewed Status} & Generated candidate: keep / modify / relocate / reject; verify whether independently testable. & \\
    \hline
  \end{tabularx}
\end{table}

Intent. Referral Routing is proposed to support the underlying healthcare concept represented by Referral Intake \& Progression. The proposal concerns behaviour and semantic responsibility, not a FHIR field, operation, resource instance or prescribed data model.

Conceptual evidence: HL7 FHIR R5 ServiceRequest, Task.

\noindent\textbf{Feature boundary.} This behaviour is scoped to Referral Intake \& Progression; an identically named Feature elsewhere may carry different business meaning.\par
\featuretable{
--- & --- & \emph{Not yet captured} & \emph{Complete during review} \\
}
\subsubsection{Feature - Referral Status Tracking}\label{cap:bec-c4e89e34}
\vspace{-1em}
\begin{table}[htbp]
  \centering
  \footnotesize
  \begin{tabularx}{1\textwidth}{|l|l|X|}
    \hline
    \textbf{Metadata Name} & \textbf{Metadata Value} & \textbf{Description} \\
    \hline
    \textbf{Reference Id} & \texttt{BEC-C4E89E34} & \\
    \hline
    \textbf{Status} & Generated candidate & \\
    \hline
    \textbf{Ancestry} &  & \\
    \hline
    \textbf{Source} & FHIR-informed candidate & \\
    \hline
    \textbf{Enabled By} &  & \\
    \hline
    \textbf{Supports} &  & \\
    \hline
    \textbf{Derived From} &  & \\
    \hline
    \textbf{Reviewed Status} & Generated candidate: keep / modify / relocate / reject; verify whether independently testable. & \\
    \hline
  \end{tabularx}
\end{table}

Intent. Referral Status Tracking is proposed to support the underlying healthcare concept represented by Referral Intake \& Progression. The proposal concerns behaviour and semantic responsibility, not a FHIR field, operation, resource instance or prescribed data model.

Conceptual evidence: HL7 FHIR R5 ServiceRequest, Task.

\noindent\textbf{Feature boundary.} This behaviour is scoped to Referral Intake \& Progression; an identically named Feature elsewhere may carry different business meaning.\par
\featuretable{
--- & --- & \emph{Not yet captured} & \emph{Complete during review} \\
}
\subsubsection{Feature - Referral Outcome Association}\label{cap:bec-57720bc2}
\vspace{-1em}
\begin{table}[htbp]
  \centering
  \footnotesize
  \begin{tabularx}{1\textwidth}{|l|l|X|}
    \hline
    \textbf{Metadata Name} & \textbf{Metadata Value} & \textbf{Description} \\
    \hline
    \textbf{Reference Id} & \texttt{BEC-57720BC2} & \\
    \hline
    \textbf{Status} & Generated candidate & \\
    \hline
    \textbf{Ancestry} &  & \\
    \hline
    \textbf{Source} & FHIR-informed candidate & \\
    \hline
    \textbf{Enabled By} &  & \\
    \hline
    \textbf{Supports} &  & \\
    \hline
    \textbf{Derived From} &  & \\
    \hline
    \textbf{Reviewed Status} & Generated candidate: keep / modify / relocate / reject; verify whether independently testable. & \\
    \hline
  \end{tabularx}
\end{table}

Intent. Referral Outcome Association is proposed to support the underlying healthcare concept represented by Referral Intake \& Progression. The proposal concerns behaviour and semantic responsibility, not a FHIR field, operation, resource instance or prescribed data model.

Conceptual evidence: HL7 FHIR R5 ServiceRequest, Task.

\noindent\textbf{Feature boundary.} This behaviour is scoped to Referral Intake \& Progression; an identically named Feature elsewhere may carry different business meaning.\par
\featuretable{
--- & --- & \emph{Not yet captured} & \emph{Complete during review} \\
}

\section{CT2 - Scheduling Administration}\label{cap:bec-2abe98e1}
\vspace{-0.5em}
\begin{table}[htbp]
  \centering
  \footnotesize
  \begin{tabularx}{1\textwidth}{|l|l|X|}
    \hline
    \multicolumn{3}{c}{\textbf{Capability Metadata}} \\
    \hline
    \textbf{Name} & \textbf{Value} & \textbf{Comment and/or Description} \\
    \hline
    \textbf{Reference Id} & \texttt{BEC-2ABE98E1} & \\
    \hline
    \textbf{Status} & Candidate ancestry & \\
    \hline
    \textbf{Ancestry} &  & \\
    \hline
    \textbf{Source} &  & \\
    \hline
    \textbf{Enabled By} &  & \\
    \hline
    \textbf{Supports} &  & \\
    \hline
    \textbf{Derived From} &  & \\
    \hline
    \textbf{Reviewed Status} &  & \\
    \hline
  \end{tabularx}
\end{table}

\begin{table}[htbp]
  \centering
  \footnotesize
  \begin{tabularx}{1\textwidth}{|l|l|X|}
    \hline
    \multicolumn{3}{c}{\textbf{Contained Capabilities}} \\
    \hline
    \textbf{Capability Name} & \textbf{Capability Id} & \textbf{Comment} \\
    \hline
    Appointment \& Booking Administration & \texttt{cap:bec-8fef95d4} (at \ref{cap:bec-8fef95d4}) & \\
    \hline
  \end{tabularx}
\end{table}

Intent. Scheduling Administration is a candidate enabling capability within Care Administration Enablement. It frames the reusable system-supported behaviours needed to support this healthcare context without prescribing a product, interface, technology or deployment boundary.

Boundary. This capability supports the business activity and does not inherit the authority or accountability of the Business Capability it enables. Its child Features are review candidates until ancestry and semantics are adjudicated.

\subsection{CT3 - Appointment \& Booking Administration}\label{cap:bec-8fef95d4}
\vspace{-0.5em}
\begin{table}[htbp]
  \centering
  \footnotesize
  \begin{tabularx}{1\textwidth}{|l|l|X|}
    \hline
    \multicolumn{3}{c}{\textbf{Capability Metadata}} \\
    \hline
    \textbf{Name} & \textbf{Value} & \textbf{Comment and/or Description} \\
    \hline
    \textbf{Reference Id} & \texttt{BEC-8FEF95D4} & \\
    \hline
    \textbf{Status} & Candidate ancestry & \\
    \hline
    \textbf{Ancestry} &  & \\
    \hline
    \textbf{Source} &  & \\
    \hline
    \textbf{Enabled By} &  & \\
    \hline
    \textbf{Supports} &  & \\
    \hline
    \textbf{Derived From} &  & \\
    \hline
    \textbf{Reviewed Status} &  & \\
    \hline
  \end{tabularx}
\end{table}

\begin{table}[htbp]
  \centering
  \footnotesize
  \begin{tabularx}{1\textwidth}{|l|l|X|}
    \hline
    \multicolumn{3}{c}{\textbf{Contained Features}} \\
    \hline
    \textbf{Feature Name} & \textbf{Feature Id} & \textbf{Comment} \\
    \hline
    Appointment Notification Ingestion & \texttt{cap:bec-8c2f56e0} (at \ref{cap:bec-8c2f56e0}) & \\
    \hline
    Appointment Status Synchronization & \texttt{cap:bec-4143272f} (at \ref{cap:bec-4143272f}) & \\
    \hline
    Appointment Lifecycle Management & \texttt{cap:bec-262653a0} (at \ref{cap:bec-262653a0}) & \\
    \hline
    Participant Coordination & \texttt{cap:bec-d9d79fa0} (at \ref{cap:bec-d9d79fa0}) & \\
    \hline
    Availability Validation & \texttt{cap:bec-425f9192} (at \ref{cap:bec-425f9192}) & \\
    \hline
    Appointment Exception Management & \texttt{cap:bec-4d29beff} (at \ref{cap:bec-4d29beff}) & \\
    \hline
  \end{tabularx}
\end{table}

Intent. Appointment \& Booking Administration is a candidate enabling capability within Scheduling Administration. It frames the reusable system-supported behaviours needed to support this healthcare context without prescribing a product, interface, technology or deployment boundary.

Boundary. This capability supports the business activity and does not inherit the authority or accountability of the Business Capability it enables. Its child Features are review candidates until ancestry and semantics are adjudicated.

FHIR R5 conceptual analogues considered (not implementation mappings): Appointment, Schedule, Slot.

\subsubsection{Feature - Appointment Notification Ingestion}\label{cap:bec-8c2f56e0}
\vspace{-1em}
\begin{table}[htbp]
  \centering
  \footnotesize
  \begin{tabularx}{1\textwidth}{|l|l|X|}
    \hline
    \textbf{Metadata Name} & \textbf{Metadata Value} & \textbf{Description} \\
    \hline
    \textbf{Reference Id} & \texttt{BEC-8C2F56E0} & \\
    \hline
    \textbf{Status} & Existing Feature; candidate CT3 placement & \\
    \hline
    \textbf{Ancestry} &  & \\
    \hline
    \textbf{Source} & FEAT-SA-06 & \\
    \hline
    \textbf{Enabled By} &  & \\
    \hline
    \textbf{Supports} &  & \\
    \hline
    \textbf{Derived From} &  & \\
    \hline
    \textbf{Reviewed Status} & Confirm existing Feature meaning and proposed CT3 ownership. & \\
    \hline
  \end{tabularx}
\end{table}

Intent. Appointment Notification Ingestion provides a distinct enabling behaviour in the context of Appointment \& Booking Administration. This is an existing named Feature; its placement under the candidate CT3 remains subject to review.

\noindent\textbf{Feature boundary.} This behaviour is scoped to Appointment \& Booking Administration; an identically named Feature elsewhere may carry different business meaning.\par
\featuretable{
Function & \texttt{FN-16301829} & Receive Appointment Notification & Catalogue \\
Service & \texttt{SV-3775B465} & Appointment Notification Ingress & Linked to preceding Function in source \\
}
\subsubsection{Feature - Appointment Status Synchronization}\label{cap:bec-4143272f}
\vspace{-1em}
\begin{table}[htbp]
  \centering
  \footnotesize
  \begin{tabularx}{1\textwidth}{|l|l|X|}
    \hline
    \textbf{Metadata Name} & \textbf{Metadata Value} & \textbf{Description} \\
    \hline
    \textbf{Reference Id} & \texttt{BEC-4143272F} & \\
    \hline
    \textbf{Status} & Existing Feature; candidate CT3 placement & \\
    \hline
    \textbf{Ancestry} &  & \\
    \hline
    \textbf{Source} & FEAT-SA-07 & \\
    \hline
    \textbf{Enabled By} &  & \\
    \hline
    \textbf{Supports} &  & \\
    \hline
    \textbf{Derived From} &  & \\
    \hline
    \textbf{Reviewed Status} & Confirm existing Feature meaning and proposed CT3 ownership. & \\
    \hline
  \end{tabularx}
\end{table}

Intent. Appointment Status Synchronization provides a distinct enabling behaviour in the context of Appointment \& Booking Administration. This is an existing named Feature; its placement under the candidate CT3 remains subject to review.

\noindent\textbf{Feature boundary.} This behaviour is scoped to Appointment \& Booking Administration; an identically named Feature elsewhere may carry different business meaning.\par
\featuretable{
Function & \texttt{FN-687F15C3} & Synchronise Appointment Status & Catalogue \\
Service & \texttt{SV-39B8CA51} & Appointment Schedule Query & Linked to preceding Function in source \\
}
\subsubsection{Feature - Appointment Lifecycle Management}\label{cap:bec-262653a0}
\vspace{-1em}
\begin{table}[htbp]
  \centering
  \footnotesize
  \begin{tabularx}{1\textwidth}{|l|l|X|}
    \hline
    \textbf{Metadata Name} & \textbf{Metadata Value} & \textbf{Description} \\
    \hline
    \textbf{Reference Id} & \texttt{BEC-262653A0} & \\
    \hline
    \textbf{Status} & Generated candidate & \\
    \hline
    \textbf{Ancestry} &  & \\
    \hline
    \textbf{Source} & FHIR-informed candidate & \\
    \hline
    \textbf{Enabled By} &  & \\
    \hline
    \textbf{Supports} &  & \\
    \hline
    \textbf{Derived From} &  & \\
    \hline
    \textbf{Reviewed Status} & Generated candidate: keep / modify / relocate / reject; verify whether independently testable. & \\
    \hline
  \end{tabularx}
\end{table}

Intent. Appointment Lifecycle Management is proposed to support the underlying healthcare concept represented by Appointment \& Booking Administration. The proposal concerns behaviour and semantic responsibility, not a FHIR field, operation, resource instance or prescribed data model.

Conceptual evidence: HL7 FHIR R5 Appointment, Schedule, Slot.

\noindent\textbf{Feature boundary.} This behaviour is scoped to Appointment \& Booking Administration; an identically named Feature elsewhere may carry different business meaning.\par
\featuretable{
--- & --- & \emph{Not yet captured} & \emph{Complete during review} \\
}
\subsubsection{Feature - Participant Coordination}\label{cap:bec-d9d79fa0}
\vspace{-1em}
\begin{table}[htbp]
  \centering
  \footnotesize
  \begin{tabularx}{1\textwidth}{|l|l|X|}
    \hline
    \textbf{Metadata Name} & \textbf{Metadata Value} & \textbf{Description} \\
    \hline
    \textbf{Reference Id} & \texttt{BEC-D9D79FA0} & \\
    \hline
    \textbf{Status} & Generated candidate & \\
    \hline
    \textbf{Ancestry} &  & \\
    \hline
    \textbf{Source} & FHIR-informed candidate & \\
    \hline
    \textbf{Enabled By} &  & \\
    \hline
    \textbf{Supports} &  & \\
    \hline
    \textbf{Derived From} &  & \\
    \hline
    \textbf{Reviewed Status} & Generated candidate: keep / modify / relocate / reject; verify whether independently testable. & \\
    \hline
  \end{tabularx}
\end{table}

Intent. Participant Coordination is proposed to support the underlying healthcare concept represented by Appointment \& Booking Administration. The proposal concerns behaviour and semantic responsibility, not a FHIR field, operation, resource instance or prescribed data model.

Conceptual evidence: HL7 FHIR R5 Appointment, Schedule, Slot.

\noindent\textbf{Feature boundary.} This behaviour is scoped to Appointment \& Booking Administration; an identically named Feature elsewhere may carry different business meaning.\par
\featuretable{
--- & --- & \emph{Not yet captured} & \emph{Complete during review} \\
}
\subsubsection{Feature - Availability Validation}\label{cap:bec-425f9192}
\vspace{-1em}
\begin{table}[htbp]
  \centering
  \footnotesize
  \begin{tabularx}{1\textwidth}{|l|l|X|}
    \hline
    \textbf{Metadata Name} & \textbf{Metadata Value} & \textbf{Description} \\
    \hline
    \textbf{Reference Id} & \texttt{BEC-425F9192} & \\
    \hline
    \textbf{Status} & Generated candidate & \\
    \hline
    \textbf{Ancestry} &  & \\
    \hline
    \textbf{Source} & FHIR-informed candidate & \\
    \hline
    \textbf{Enabled By} &  & \\
    \hline
    \textbf{Supports} &  & \\
    \hline
    \textbf{Derived From} &  & \\
    \hline
    \textbf{Reviewed Status} & Generated candidate: keep / modify / relocate / reject; verify whether independently testable. & \\
    \hline
  \end{tabularx}
\end{table}

Intent. Availability Validation is proposed to support the underlying healthcare concept represented by Appointment \& Booking Administration. The proposal concerns behaviour and semantic responsibility, not a FHIR field, operation, resource instance or prescribed data model.

Conceptual evidence: HL7 FHIR R5 Appointment, Schedule, Slot.

\noindent\textbf{Feature boundary.} This behaviour is scoped to Appointment \& Booking Administration; an identically named Feature elsewhere may carry different business meaning.\par
\featuretable{
--- & --- & \emph{Not yet captured} & \emph{Complete during review} \\
}
\subsubsection{Feature - Appointment Exception Management}\label{cap:bec-4d29beff}
\vspace{-1em}
\begin{table}[htbp]
  \centering
  \footnotesize
  \begin{tabularx}{1\textwidth}{|l|l|X|}
    \hline
    \textbf{Metadata Name} & \textbf{Metadata Value} & \textbf{Description} \\
    \hline
    \textbf{Reference Id} & \texttt{BEC-4D29BEFF} & \\
    \hline
    \textbf{Status} & Generated candidate & \\
    \hline
    \textbf{Ancestry} &  & \\
    \hline
    \textbf{Source} & FHIR-informed candidate & \\
    \hline
    \textbf{Enabled By} &  & \\
    \hline
    \textbf{Supports} &  & \\
    \hline
    \textbf{Derived From} &  & \\
    \hline
    \textbf{Reviewed Status} & Generated candidate: keep / modify / relocate / reject; verify whether independently testable. & \\
    \hline
  \end{tabularx}
\end{table}

Intent. Appointment Exception Management is proposed to support the underlying healthcare concept represented by Appointment \& Booking Administration. The proposal concerns behaviour and semantic responsibility, not a FHIR field, operation, resource instance or prescribed data model.

Conceptual evidence: HL7 FHIR R5 Appointment, Schedule, Slot.

\noindent\textbf{Feature boundary.} This behaviour is scoped to Appointment \& Booking Administration; an identically named Feature elsewhere may carry different business meaning.\par
\featuretable{
--- & --- & \emph{Not yet captured} & \emph{Complete during review} \\
}

\section{CT2 - Episode \& Encounter Administration}\label{cap:bec-63e23d60}
\vspace{-0.5em}
\begin{table}[htbp]
  \centering
  \footnotesize
  \begin{tabularx}{1\textwidth}{|l|l|X|}
    \hline
    \multicolumn{3}{c}{\textbf{Capability Metadata}} \\
    \hline
    \textbf{Name} & \textbf{Value} & \textbf{Comment and/or Description} \\
    \hline
    \textbf{Reference Id} & \texttt{BEC-63E23D60} & \\
    \hline
    \textbf{Status} & Candidate ancestry & \\
    \hline
    \textbf{Ancestry} &  & \\
    \hline
    \textbf{Source} &  & \\
    \hline
    \textbf{Enabled By} &  & \\
    \hline
    \textbf{Supports} &  & \\
    \hline
    \textbf{Derived From} &  & \\
    \hline
    \textbf{Reviewed Status} &  & \\
    \hline
  \end{tabularx}
\end{table}

\begin{table}[htbp]
  \centering
  \footnotesize
  \begin{tabularx}{1\textwidth}{|l|l|X|}
    \hline
    \multicolumn{3}{c}{\textbf{Contained Capabilities}} \\
    \hline
    \textbf{Capability Name} & \textbf{Capability Id} & \textbf{Comment} \\
    \hline
    Encounter Lifecycle Management & \texttt{cap:bec-1c528fc1} (at \ref{cap:bec-1c528fc1}) & \\
    \hline
  \end{tabularx}
\end{table}

Intent. Episode \& Encounter Administration is a candidate enabling capability within Care Administration Enablement. It frames the reusable system-supported behaviours needed to support this healthcare context without prescribing a product, interface, technology or deployment boundary.

Boundary. This capability supports the business activity and does not inherit the authority or accountability of the Business Capability it enables. Its child Features are review candidates until ancestry and semantics are adjudicated.

\subsection{CT3 - Encounter Lifecycle Management}\label{cap:bec-1c528fc1}
\vspace{-0.5em}
\begin{table}[htbp]
  \centering
  \footnotesize
  \begin{tabularx}{1\textwidth}{|l|l|X|}
    \hline
    \multicolumn{3}{c}{\textbf{Capability Metadata}} \\
    \hline
    \textbf{Name} & \textbf{Value} & \textbf{Comment and/or Description} \\
    \hline
    \textbf{Reference Id} & \texttt{BEC-1C528FC1} & \\
    \hline
    \textbf{Status} & Candidate ancestry & \\
    \hline
    \textbf{Ancestry} &  & \\
    \hline
    \textbf{Source} &  & \\
    \hline
    \textbf{Enabled By} &  & \\
    \hline
    \textbf{Supports} &  & \\
    \hline
    \textbf{Derived From} &  & \\
    \hline
    \textbf{Reviewed Status} &  & \\
    \hline
  \end{tabularx}
\end{table}

\begin{table}[htbp]
  \centering
  \footnotesize
  \begin{tabularx}{1\textwidth}{|l|l|X|}
    \hline
    \multicolumn{3}{c}{\textbf{Contained Features}} \\
    \hline
    \textbf{Feature Name} & \textbf{Feature Id} & \textbf{Comment} \\
    \hline
    Encounter Context Creation & \texttt{cap:bec-0dc29356} (at \ref{cap:bec-0dc29356}) & \\
    \hline
    Encounter State Progression & \texttt{cap:bec-5d8ca48c} (at \ref{cap:bec-5d8ca48c}) & \\
    \hline
    Encounter Bed/Location Association & \texttt{cap:bec-847fcfdb} (at \ref{cap:bec-847fcfdb}) & \\
    \hline
    Encounter Lifecycle Management & \texttt{cap:bec-3ec99787} (at \ref{cap:bec-3ec99787}) & \\
    \hline
    Care Context Association & \texttt{cap:bec-46431ec5} (at \ref{cap:bec-46431ec5}) & \\
    \hline
    Participant Association & \texttt{cap:bec-12eb85c6} (at \ref{cap:bec-12eb85c6}) & \\
    \hline
    Encounter History Retrieval & \texttt{cap:bec-d42c14e5} (at \ref{cap:bec-d42c14e5}) & \\
    \hline
  \end{tabularx}
\end{table}

Intent. Encounter Lifecycle Management is a candidate enabling capability within Episode \& Encounter Administration. It frames the reusable system-supported behaviours needed to support this healthcare context without prescribing a product, interface, technology or deployment boundary.

Boundary. This capability supports the business activity and does not inherit the authority or accountability of the Business Capability it enables. Its child Features are review candidates until ancestry and semantics are adjudicated.

FHIR R5 conceptual analogues considered (not implementation mappings): Encounter, EpisodeOfCare.

\subsubsection{Feature - Encounter Context Creation}\label{cap:bec-0dc29356}
\vspace{-1em}
\begin{table}[htbp]
  \centering
  \footnotesize
  \begin{tabularx}{1\textwidth}{|l|l|X|}
    \hline
    \textbf{Metadata Name} & \textbf{Metadata Value} & \textbf{Description} \\
    \hline
    \textbf{Reference Id} & \texttt{BEC-0DC29356} & \\
    \hline
    \textbf{Status} & Existing Feature; candidate CT3 placement & \\
    \hline
    \textbf{Ancestry} &  & \\
    \hline
    \textbf{Source} & FEAT-SA-08 & \\
    \hline
    \textbf{Enabled By} &  & \\
    \hline
    \textbf{Supports} &  & \\
    \hline
    \textbf{Derived From} &  & \\
    \hline
    \textbf{Reviewed Status} & Confirm existing Feature meaning and proposed CT3 ownership. & \\
    \hline
  \end{tabularx}
\end{table}

Intent. Encounter Context Creation provides a distinct enabling behaviour in the context of Encounter Lifecycle Management. This is an existing named Feature; its placement under the candidate CT3 remains subject to review.

\noindent\textbf{Feature boundary.} This behaviour is scoped to Encounter Lifecycle Management; an identically named Feature elsewhere may carry different business meaning.\par
\featuretable{
Function & \texttt{FN-CDC10B7B} & Establish Encounter Context & Catalogue \\
Service & \texttt{SV-55A86305} & Encounter Context Resolution & Linked to preceding Function in source \\
}
\subsubsection{Feature - Encounter State Progression}\label{cap:bec-5d8ca48c}
\vspace{-1em}
\begin{table}[htbp]
  \centering
  \footnotesize
  \begin{tabularx}{1\textwidth}{|l|l|X|}
    \hline
    \textbf{Metadata Name} & \textbf{Metadata Value} & \textbf{Description} \\
    \hline
    \textbf{Reference Id} & \texttt{BEC-5D8CA48C} & \\
    \hline
    \textbf{Status} & Existing Feature; candidate CT3 placement & \\
    \hline
    \textbf{Ancestry} &  & \\
    \hline
    \textbf{Source} & FEAT-SA-09 & \\
    \hline
    \textbf{Enabled By} &  & \\
    \hline
    \textbf{Supports} &  & \\
    \hline
    \textbf{Derived From} &  & \\
    \hline
    \textbf{Reviewed Status} & Confirm existing Feature meaning and proposed CT3 ownership. & \\
    \hline
  \end{tabularx}
\end{table}

Intent. Encounter State Progression provides a distinct enabling behaviour in the context of Encounter Lifecycle Management. This is an existing named Feature; its placement under the candidate CT3 remains subject to review.

\noindent\textbf{Feature boundary.} This behaviour is scoped to Encounter Lifecycle Management; an identically named Feature elsewhere may carry different business meaning.\par
\featuretable{
Function & \texttt{FN-51FA5B68} & Progress Encounter State & Catalogue \\
Service & \texttt{SV-4103B339} & Encounter Lifecycle Event Notification & Linked to preceding Function in source \\
Process & \texttt{PR-E9189DD7} & Encounter Lifecycle Process & Linked to preceding Function in source \\
}
\subsubsection{Feature - Encounter Bed/Location Association}\label{cap:bec-847fcfdb}
\vspace{-1em}
\begin{table}[htbp]
  \centering
  \footnotesize
  \begin{tabularx}{1\textwidth}{|l|l|X|}
    \hline
    \textbf{Metadata Name} & \textbf{Metadata Value} & \textbf{Description} \\
    \hline
    \textbf{Reference Id} & \texttt{BEC-847FCFDB} & \\
    \hline
    \textbf{Status} & Existing Feature; candidate CT3 placement & \\
    \hline
    \textbf{Ancestry} &  & \\
    \hline
    \textbf{Source} & FEAT-SA-10 & \\
    \hline
    \textbf{Enabled By} &  & \\
    \hline
    \textbf{Supports} &  & \\
    \hline
    \textbf{Derived From} &  & \\
    \hline
    \textbf{Reviewed Status} & Confirm existing Feature meaning and proposed CT3 ownership. & \\
    \hline
  \end{tabularx}
\end{table}

Intent. Encounter Bed/Location Association provides a distinct enabling behaviour in the context of Encounter Lifecycle Management. This is an existing named Feature; its placement under the candidate CT3 remains subject to review.

\noindent\textbf{Feature boundary.} This behaviour is scoped to Encounter Lifecycle Management; an identically named Feature elsewhere may carry different business meaning.\par
\featuretable{
Function & \texttt{FN-125F7505} & Maintain Encounter Care-Place Association & Catalogue \\
Service & \texttt{SV-11F6B242} & Encounter Location History Query & Linked to preceding Function in source \\
}
\subsubsection{Feature - Encounter Lifecycle Management}\label{cap:bec-3ec99787}
\vspace{-1em}
\begin{table}[htbp]
  \centering
  \footnotesize
  \begin{tabularx}{1\textwidth}{|l|l|X|}
    \hline
    \textbf{Metadata Name} & \textbf{Metadata Value} & \textbf{Description} \\
    \hline
    \textbf{Reference Id} & \texttt{BEC-3EC99787} & \\
    \hline
    \textbf{Status} & Generated candidate & \\
    \hline
    \textbf{Ancestry} &  & \\
    \hline
    \textbf{Source} & FHIR-informed candidate & \\
    \hline
    \textbf{Enabled By} &  & \\
    \hline
    \textbf{Supports} &  & \\
    \hline
    \textbf{Derived From} &  & \\
    \hline
    \textbf{Reviewed Status} & Generated candidate: keep / modify / relocate / reject; verify whether independently testable. & \\
    \hline
  \end{tabularx}
\end{table}

Intent. Encounter Lifecycle Management is proposed to support the underlying healthcare concept represented by Encounter Lifecycle Management. The proposal concerns behaviour and semantic responsibility, not a FHIR field, operation, resource instance or prescribed data model.

Conceptual evidence: HL7 FHIR R5 Encounter, EpisodeOfCare.

\noindent\textbf{Feature boundary.} This behaviour is scoped to Encounter Lifecycle Management; an identically named Feature elsewhere may carry different business meaning.\par
\featuretable{
--- & --- & \emph{Not yet captured} & \emph{Complete during review} \\
}
\subsubsection{Feature - Care Context Association}\label{cap:bec-46431ec5}
\vspace{-1em}
\begin{table}[htbp]
  \centering
  \footnotesize
  \begin{tabularx}{1\textwidth}{|l|l|X|}
    \hline
    \textbf{Metadata Name} & \textbf{Metadata Value} & \textbf{Description} \\
    \hline
    \textbf{Reference Id} & \texttt{BEC-46431EC5} & \\
    \hline
    \textbf{Status} & Generated candidate & \\
    \hline
    \textbf{Ancestry} &  & \\
    \hline
    \textbf{Source} & FHIR-informed candidate & \\
    \hline
    \textbf{Enabled By} &  & \\
    \hline
    \textbf{Supports} &  & \\
    \hline
    \textbf{Derived From} &  & \\
    \hline
    \textbf{Reviewed Status} & Generated candidate: keep / modify / relocate / reject; verify whether independently testable. & \\
    \hline
  \end{tabularx}
\end{table}

Intent. Care Context Association is proposed to support the underlying healthcare concept represented by Encounter Lifecycle Management. The proposal concerns behaviour and semantic responsibility, not a FHIR field, operation, resource instance or prescribed data model.

Conceptual evidence: HL7 FHIR R5 Encounter, EpisodeOfCare.

\noindent\textbf{Feature boundary.} This behaviour is scoped to Encounter Lifecycle Management; an identically named Feature elsewhere may carry different business meaning.\par
\featuretable{
--- & --- & \emph{Not yet captured} & \emph{Complete during review} \\
}
\subsubsection{Feature - Participant Association}\label{cap:bec-12eb85c6}
\vspace{-1em}
\begin{table}[htbp]
  \centering
  \footnotesize
  \begin{tabularx}{1\textwidth}{|l|l|X|}
    \hline
    \textbf{Metadata Name} & \textbf{Metadata Value} & \textbf{Description} \\
    \hline
    \textbf{Reference Id} & \texttt{BEC-12EB85C6} & \\
    \hline
    \textbf{Status} & Generated candidate & \\
    \hline
    \textbf{Ancestry} &  & \\
    \hline
    \textbf{Source} & FHIR-informed candidate & \\
    \hline
    \textbf{Enabled By} &  & \\
    \hline
    \textbf{Supports} &  & \\
    \hline
    \textbf{Derived From} &  & \\
    \hline
    \textbf{Reviewed Status} & Generated candidate: keep / modify / relocate / reject; verify whether independently testable. & \\
    \hline
  \end{tabularx}
\end{table}

Intent. Participant Association is proposed to support the underlying healthcare concept represented by Encounter Lifecycle Management. The proposal concerns behaviour and semantic responsibility, not a FHIR field, operation, resource instance or prescribed data model.

Conceptual evidence: HL7 FHIR R5 Encounter, EpisodeOfCare.

\noindent\textbf{Feature boundary.} This behaviour is scoped to Encounter Lifecycle Management; an identically named Feature elsewhere may carry different business meaning.\par
\featuretable{
--- & --- & \emph{Not yet captured} & \emph{Complete during review} \\
}
\subsubsection{Feature - Encounter History Retrieval}\label{cap:bec-d42c14e5}
\vspace{-1em}
\begin{table}[htbp]
  \centering
  \footnotesize
  \begin{tabularx}{1\textwidth}{|l|l|X|}
    \hline
    \textbf{Metadata Name} & \textbf{Metadata Value} & \textbf{Description} \\
    \hline
    \textbf{Reference Id} & \texttt{BEC-D42C14E5} & \\
    \hline
    \textbf{Status} & Generated candidate & \\
    \hline
    \textbf{Ancestry} &  & \\
    \hline
    \textbf{Source} & FHIR-informed candidate & \\
    \hline
    \textbf{Enabled By} &  & \\
    \hline
    \textbf{Supports} &  & \\
    \hline
    \textbf{Derived From} &  & \\
    \hline
    \textbf{Reviewed Status} & Generated candidate: keep / modify / relocate / reject; verify whether independently testable. & \\
    \hline
  \end{tabularx}
\end{table}

Intent. Encounter History Retrieval is proposed to support the underlying healthcare concept represented by Encounter Lifecycle Management. The proposal concerns behaviour and semantic responsibility, not a FHIR field, operation, resource instance or prescribed data model.

Conceptual evidence: HL7 FHIR R5 Encounter, EpisodeOfCare.

\noindent\textbf{Feature boundary.} This behaviour is scoped to Encounter Lifecycle Management; an identically named Feature elsewhere may carry different business meaning.\par
\featuretable{
--- & --- & \emph{Not yet captured} & \emph{Complete during review} \\
}

\section{CT2 - Order Administration}\label{cap:bec-b6be6e1d}
\vspace{-0.5em}
\begin{table}[htbp]
  \centering
  \footnotesize
  \begin{tabularx}{1\textwidth}{|l|l|X|}
    \hline
    \multicolumn{3}{c}{\textbf{Capability Metadata}} \\
    \hline
    \textbf{Name} & \textbf{Value} & \textbf{Comment and/or Description} \\
    \hline
    \textbf{Reference Id} & \texttt{BEC-B6BE6E1D} & \\
    \hline
    \textbf{Status} & Candidate ancestry & \\
    \hline
    \textbf{Ancestry} &  & \\
    \hline
    \textbf{Source} &  & \\
    \hline
    \textbf{Enabled By} &  & \\
    \hline
    \textbf{Supports} &  & \\
    \hline
    \textbf{Derived From} &  & \\
    \hline
    \textbf{Reviewed Status} &  & \\
    \hline
  \end{tabularx}
\end{table}

\begin{table}[htbp]
  \centering
  \footnotesize
  \begin{tabularx}{1\textwidth}{|l|l|X|}
    \hline
    \multicolumn{3}{c}{\textbf{Contained Capabilities}} \\
    \hline
    \textbf{Capability Name} & \textbf{Capability Id} & \textbf{Comment} \\
    \hline
    Closed-Loop Order Progression & \texttt{cap:bec-5e9ecea1} (at \ref{cap:bec-5e9ecea1}) & \\
    \hline
  \end{tabularx}
\end{table}

Intent. Order Administration is a candidate enabling capability within Care Administration Enablement. It frames the reusable system-supported behaviours needed to support this healthcare context without prescribing a product, interface, technology or deployment boundary.

Boundary. This capability supports the business activity and does not inherit the authority or accountability of the Business Capability it enables. Its child Features are review candidates until ancestry and semantics are adjudicated.

\subsection{CT3 - Closed-Loop Order Progression}\label{cap:bec-5e9ecea1}
\vspace{-0.5em}
\begin{table}[htbp]
  \centering
  \footnotesize
  \begin{tabularx}{1\textwidth}{|l|l|X|}
    \hline
    \multicolumn{3}{c}{\textbf{Capability Metadata}} \\
    \hline
    \textbf{Name} & \textbf{Value} & \textbf{Comment and/or Description} \\
    \hline
    \textbf{Reference Id} & \texttt{BEC-5E9ECEA1} & \\
    \hline
    \textbf{Status} & Candidate ancestry & \\
    \hline
    \textbf{Ancestry} &  & \\
    \hline
    \textbf{Source} &  & \\
    \hline
    \textbf{Enabled By} &  & \\
    \hline
    \textbf{Supports} &  & \\
    \hline
    \textbf{Derived From} &  & \\
    \hline
    \textbf{Reviewed Status} &  & \\
    \hline
  \end{tabularx}
\end{table}

\begin{table}[htbp]
  \centering
  \footnotesize
  \begin{tabularx}{1\textwidth}{|l|l|X|}
    \hline
    \multicolumn{3}{c}{\textbf{Contained Features}} \\
    \hline
    \textbf{Feature Name} & \textbf{Feature Id} & \textbf{Comment} \\
    \hline
    Order Request Ingestion & \texttt{cap:bec-d0729bfb} (at \ref{cap:bec-d0729bfb}) & \\
    \hline
    Order Destination Resolution \& Routing & \texttt{cap:bec-20beda79} (at \ref{cap:bec-20beda79}) & \\
    \hline
    Order Closed-Loop Progression Tracking & \texttt{cap:bec-e57851c6} (at \ref{cap:bec-e57851c6}) & \\
    \hline
    Order Cancellation \& Modification Coordination & \texttt{cap:bec-84eff3a3} (at \ref{cap:bec-84eff3a3}) & \\
    \hline
    Request Lifecycle Management & \texttt{cap:bec-142e39c3} (at \ref{cap:bec-142e39c3}) & \\
    \hline
    Request Fulfilment Tracking & \texttt{cap:bec-977d6a7d} (at \ref{cap:bec-977d6a7d}) & \\
    \hline
    Request Status Correlation & \texttt{cap:bec-81d316c2} (at \ref{cap:bec-81d316c2}) & \\
    \hline
    Request Exception Management & \texttt{cap:bec-aa804bef} (at \ref{cap:bec-aa804bef}) & \\
    \hline
  \end{tabularx}
\end{table}

Intent. Closed-Loop Order Progression is a candidate enabling capability within Order Administration. It frames the reusable system-supported behaviours needed to support this healthcare context without prescribing a product, interface, technology or deployment boundary.

Boundary. This capability supports the business activity and does not inherit the authority or accountability of the Business Capability it enables. Its child Features are review candidates until ancestry and semantics are adjudicated.

FHIR R5 conceptual analogues considered (not implementation mappings): ServiceRequest, Task.

\subsubsection{Feature - Order Request Ingestion}\label{cap:bec-d0729bfb}
\vspace{-1em}
\begin{table}[htbp]
  \centering
  \footnotesize
  \begin{tabularx}{1\textwidth}{|l|l|X|}
    \hline
    \textbf{Metadata Name} & \textbf{Metadata Value} & \textbf{Description} \\
    \hline
    \textbf{Reference Id} & \texttt{BEC-D0729BFB} & \\
    \hline
    \textbf{Status} & Existing Feature; candidate CT3 placement & \\
    \hline
    \textbf{Ancestry} &  & \\
    \hline
    \textbf{Source} & FEAT-SA-11 & \\
    \hline
    \textbf{Enabled By} &  & \\
    \hline
    \textbf{Supports} &  & \\
    \hline
    \textbf{Derived From} &  & \\
    \hline
    \textbf{Reviewed Status} & Confirm existing Feature meaning and proposed CT3 ownership. & \\
    \hline
  \end{tabularx}
\end{table}

Intent. Order Request Ingestion provides a distinct enabling behaviour in the context of Closed-Loop Order Progression. This is an existing named Feature; its placement under the candidate CT3 remains subject to review.

\noindent\textbf{Feature boundary.} This behaviour is scoped to Closed-Loop Order Progression; an identically named Feature elsewhere may carry different business meaning.\par
\featuretable{
Function & \texttt{FN-6013BCAF} & Receive Order Request & Catalogue \\
Service & \texttt{SV-F8B4F8D5} & Order Requisition Ingress & Linked to preceding Function in source \\
}
\subsubsection{Feature - Order Destination Resolution \& Routing}\label{cap:bec-20beda79}
\vspace{-1em}
\begin{table}[htbp]
  \centering
  \footnotesize
  \begin{tabularx}{1\textwidth}{|l|l|X|}
    \hline
    \textbf{Metadata Name} & \textbf{Metadata Value} & \textbf{Description} \\
    \hline
    \textbf{Reference Id} & \texttt{BEC-20BEDA79} & \\
    \hline
    \textbf{Status} & Existing Feature; candidate CT3 placement & \\
    \hline
    \textbf{Ancestry} &  & \\
    \hline
    \textbf{Source} & FEAT-SA-12 & \\
    \hline
    \textbf{Enabled By} &  & \\
    \hline
    \textbf{Supports} &  & \\
    \hline
    \textbf{Derived From} &  & \\
    \hline
    \textbf{Reviewed Status} & Confirm existing Feature meaning and proposed CT3 ownership. & \\
    \hline
  \end{tabularx}
\end{table}

Intent. Order Destination Resolution \& Routing provides a distinct enabling behaviour in the context of Closed-Loop Order Progression. This is an existing named Feature; its placement under the candidate CT3 remains subject to review.

\noindent\textbf{Feature boundary.} This behaviour is scoped to Closed-Loop Order Progression; an identically named Feature elsewhere may carry different business meaning.\par
\featuretable{
Function & \texttt{FN-726AA36A} & Resolve Order Destination & Catalogue \\
Service & \texttt{SV-1508B9FE} & Order Dispatch Service & Linked to preceding Function in source \\
}
\subsubsection{Feature - Order Closed-Loop Progression Tracking}\label{cap:bec-e57851c6}
\vspace{-1em}
\begin{table}[htbp]
  \centering
  \footnotesize
  \begin{tabularx}{1\textwidth}{|l|l|X|}
    \hline
    \textbf{Metadata Name} & \textbf{Metadata Value} & \textbf{Description} \\
    \hline
    \textbf{Reference Id} & \texttt{BEC-E57851C6} & \\
    \hline
    \textbf{Status} & Existing Feature; candidate CT3 placement & \\
    \hline
    \textbf{Ancestry} &  & \\
    \hline
    \textbf{Source} & FEAT-SA-13 & \\
    \hline
    \textbf{Enabled By} &  & \\
    \hline
    \textbf{Supports} &  & \\
    \hline
    \textbf{Derived From} &  & \\
    \hline
    \textbf{Reviewed Status} & Confirm existing Feature meaning and proposed CT3 ownership. & \\
    \hline
  \end{tabularx}
\end{table}

Intent. Order Closed-Loop Progression Tracking provides a distinct enabling behaviour in the context of Closed-Loop Order Progression. This is an existing named Feature; its placement under the candidate CT3 remains subject to review.

\noindent\textbf{Feature boundary.} This behaviour is scoped to Closed-Loop Order Progression; an identically named Feature elsewhere may carry different business meaning.\par
\featuretable{
Function & \texttt{FN-3A94F4DF} & Manage Order Progression & Catalogue \\
Service & \texttt{SV-492AD5A0} & Order Status \& Tracking Query & Linked to preceding Function in source \\
Process & \texttt{PR-D1B28670} & Closed-Loop Order Progression Process & Linked to preceding Function in source \\
}
\subsubsection{Feature - Order Cancellation \& Modification Coordination}\label{cap:bec-84eff3a3}
\vspace{-1em}
\begin{table}[htbp]
  \centering
  \footnotesize
  \begin{tabularx}{1\textwidth}{|l|l|X|}
    \hline
    \textbf{Metadata Name} & \textbf{Metadata Value} & \textbf{Description} \\
    \hline
    \textbf{Reference Id} & \texttt{BEC-84EFF3A3} & \\
    \hline
    \textbf{Status} & Existing Feature; candidate CT3 placement & \\
    \hline
    \textbf{Ancestry} &  & \\
    \hline
    \textbf{Source} & FEAT-SA-14 & \\
    \hline
    \textbf{Enabled By} &  & \\
    \hline
    \textbf{Supports} &  & \\
    \hline
    \textbf{Derived From} &  & \\
    \hline
    \textbf{Reviewed Status} & Confirm existing Feature meaning and proposed CT3 ownership. & \\
    \hline
  \end{tabularx}
\end{table}

Intent. Order Cancellation \& Modification Coordination provides a distinct enabling behaviour in the context of Closed-Loop Order Progression. This is an existing named Feature; its placement under the candidate CT3 remains subject to review.

\noindent\textbf{Feature boundary.} This behaviour is scoped to Closed-Loop Order Progression; an identically named Feature elsewhere may carry different business meaning.\par
\featuretable{
Function & \texttt{FN-05225DCB} & Coordinate Order Modification / Cancellation & Catalogue \\
Service & \texttt{SV-E490151A} & Order Cancellation Service & Linked to preceding Function in source \\
}
\subsubsection{Feature - Request Lifecycle Management}\label{cap:bec-142e39c3}
\vspace{-1em}
\begin{table}[htbp]
  \centering
  \footnotesize
  \begin{tabularx}{1\textwidth}{|l|l|X|}
    \hline
    \textbf{Metadata Name} & \textbf{Metadata Value} & \textbf{Description} \\
    \hline
    \textbf{Reference Id} & \texttt{BEC-142E39C3} & \\
    \hline
    \textbf{Status} & Generated candidate & \\
    \hline
    \textbf{Ancestry} &  & \\
    \hline
    \textbf{Source} & FHIR-informed candidate & \\
    \hline
    \textbf{Enabled By} &  & \\
    \hline
    \textbf{Supports} &  & \\
    \hline
    \textbf{Derived From} &  & \\
    \hline
    \textbf{Reviewed Status} & Generated candidate: keep / modify / relocate / reject; verify whether independently testable. & \\
    \hline
  \end{tabularx}
\end{table}

Intent. Request Lifecycle Management is proposed to support the underlying healthcare concept represented by Closed-Loop Order Progression. The proposal concerns behaviour and semantic responsibility, not a FHIR field, operation, resource instance or prescribed data model.

Conceptual evidence: HL7 FHIR R5 ServiceRequest, Task.

\noindent\textbf{Feature boundary.} This behaviour is scoped to Closed-Loop Order Progression; an identically named Feature elsewhere may carry different business meaning.\par
\featuretable{
--- & --- & \emph{Not yet captured} & \emph{Complete during review} \\
}
\subsubsection{Feature - Request Fulfilment Tracking}\label{cap:bec-977d6a7d}
\vspace{-1em}
\begin{table}[htbp]
  \centering
  \footnotesize
  \begin{tabularx}{1\textwidth}{|l|l|X|}
    \hline
    \textbf{Metadata Name} & \textbf{Metadata Value} & \textbf{Description} \\
    \hline
    \textbf{Reference Id} & \texttt{BEC-977D6A7D} & \\
    \hline
    \textbf{Status} & Generated candidate & \\
    \hline
    \textbf{Ancestry} &  & \\
    \hline
    \textbf{Source} & FHIR-informed candidate & \\
    \hline
    \textbf{Enabled By} &  & \\
    \hline
    \textbf{Supports} &  & \\
    \hline
    \textbf{Derived From} &  & \\
    \hline
    \textbf{Reviewed Status} & Generated candidate: keep / modify / relocate / reject; verify whether independently testable. & \\
    \hline
  \end{tabularx}
\end{table}

Intent. Request Fulfilment Tracking is proposed to support the underlying healthcare concept represented by Closed-Loop Order Progression. The proposal concerns behaviour and semantic responsibility, not a FHIR field, operation, resource instance or prescribed data model.

Conceptual evidence: HL7 FHIR R5 ServiceRequest, Task.

\noindent\textbf{Feature boundary.} This behaviour is scoped to Closed-Loop Order Progression; an identically named Feature elsewhere may carry different business meaning.\par
\featuretable{
--- & --- & \emph{Not yet captured} & \emph{Complete during review} \\
}
\subsubsection{Feature - Request Status Correlation}\label{cap:bec-81d316c2}
\vspace{-1em}
\begin{table}[htbp]
  \centering
  \footnotesize
  \begin{tabularx}{1\textwidth}{|l|l|X|}
    \hline
    \textbf{Metadata Name} & \textbf{Metadata Value} & \textbf{Description} \\
    \hline
    \textbf{Reference Id} & \texttt{BEC-81D316C2} & \\
    \hline
    \textbf{Status} & Generated candidate & \\
    \hline
    \textbf{Ancestry} &  & \\
    \hline
    \textbf{Source} & FHIR-informed candidate & \\
    \hline
    \textbf{Enabled By} &  & \\
    \hline
    \textbf{Supports} &  & \\
    \hline
    \textbf{Derived From} &  & \\
    \hline
    \textbf{Reviewed Status} & Generated candidate: keep / modify / relocate / reject; verify whether independently testable. & \\
    \hline
  \end{tabularx}
\end{table}

Intent. Request Status Correlation is proposed to support the underlying healthcare concept represented by Closed-Loop Order Progression. The proposal concerns behaviour and semantic responsibility, not a FHIR field, operation, resource instance or prescribed data model.

Conceptual evidence: HL7 FHIR R5 ServiceRequest, Task.

\noindent\textbf{Feature boundary.} This behaviour is scoped to Closed-Loop Order Progression; an identically named Feature elsewhere may carry different business meaning.\par
\featuretable{
--- & --- & \emph{Not yet captured} & \emph{Complete during review} \\
}
\subsubsection{Feature - Request Exception Management}\label{cap:bec-aa804bef}
\vspace{-1em}
\begin{table}[htbp]
  \centering
  \footnotesize
  \begin{tabularx}{1\textwidth}{|l|l|X|}
    \hline
    \textbf{Metadata Name} & \textbf{Metadata Value} & \textbf{Description} \\
    \hline
    \textbf{Reference Id} & \texttt{BEC-AA804BEF} & \\
    \hline
    \textbf{Status} & Generated candidate & \\
    \hline
    \textbf{Ancestry} &  & \\
    \hline
    \textbf{Source} & FHIR-informed candidate & \\
    \hline
    \textbf{Enabled By} &  & \\
    \hline
    \textbf{Supports} &  & \\
    \hline
    \textbf{Derived From} &  & \\
    \hline
    \textbf{Reviewed Status} & Generated candidate: keep / modify / relocate / reject; verify whether independently testable. & \\
    \hline
  \end{tabularx}
\end{table}

Intent. Request Exception Management is proposed to support the underlying healthcare concept represented by Closed-Loop Order Progression. The proposal concerns behaviour and semantic responsibility, not a FHIR field, operation, resource instance or prescribed data model.

Conceptual evidence: HL7 FHIR R5 ServiceRequest, Task.

\noindent\textbf{Feature boundary.} This behaviour is scoped to Closed-Loop Order Progression; an identically named Feature elsewhere may carry different business meaning.\par
\featuretable{
--- & --- & \emph{Not yet captured} & \emph{Complete during review} \\
}

\section{CT2 - Diagnostic Administration}\label{cap:bec-b67a0e82}
\vspace{-0.5em}
\begin{table}[htbp]
  \centering
  \footnotesize
  \begin{tabularx}{1\textwidth}{|l|l|X|}
    \hline
    \multicolumn{3}{c}{\textbf{Capability Metadata}} \\
    \hline
    \textbf{Name} & \textbf{Value} & \textbf{Comment and/or Description} \\
    \hline
    \textbf{Reference Id} & \texttt{BEC-B67A0E82} & \\
    \hline
    \textbf{Status} & Candidate ancestry & \\
    \hline
    \textbf{Ancestry} &  & \\
    \hline
    \textbf{Source} &  & \\
    \hline
    \textbf{Enabled By} &  & \\
    \hline
    \textbf{Supports} &  & \\
    \hline
    \textbf{Derived From} &  & \\
    \hline
    \textbf{Reviewed Status} &  & \\
    \hline
  \end{tabularx}
\end{table}

\begin{table}[htbp]
  \centering
  \footnotesize
  \begin{tabularx}{1\textwidth}{|l|l|X|}
    \hline
    \multicolumn{3}{c}{\textbf{Contained Capabilities}} \\
    \hline
    \textbf{Capability Name} & \textbf{Capability Id} & \textbf{Comment} \\
    \hline
    Diagnostic Workflow Management & \texttt{cap:bec-c0dd4810} (at \ref{cap:bec-c0dd4810}) & \\
    \hline
  \end{tabularx}
\end{table}

Intent. Diagnostic Administration is a candidate enabling capability within Care Administration Enablement. It frames the reusable system-supported behaviours needed to support this healthcare context without prescribing a product, interface, technology or deployment boundary.

Boundary. This capability supports the business activity and does not inherit the authority or accountability of the Business Capability it enables. Its child Features are review candidates until ancestry and semantics are adjudicated.

\subsection{CT3 - Diagnostic Workflow Management}\label{cap:bec-c0dd4810}
\vspace{-0.5em}
\begin{table}[htbp]
  \centering
  \footnotesize
  \begin{tabularx}{1\textwidth}{|l|l|X|}
    \hline
    \multicolumn{3}{c}{\textbf{Capability Metadata}} \\
    \hline
    \textbf{Name} & \textbf{Value} & \textbf{Comment and/or Description} \\
    \hline
    \textbf{Reference Id} & \texttt{BEC-C0DD4810} & \\
    \hline
    \textbf{Status} & Candidate ancestry & \\
    \hline
    \textbf{Ancestry} &  & \\
    \hline
    \textbf{Source} &  & \\
    \hline
    \textbf{Enabled By} &  & \\
    \hline
    \textbf{Supports} &  & \\
    \hline
    \textbf{Derived From} &  & \\
    \hline
    \textbf{Reviewed Status} &  & \\
    \hline
  \end{tabularx}
\end{table}

\begin{table}[htbp]
  \centering
  \footnotesize
  \begin{tabularx}{1\textwidth}{|l|l|X|}
    \hline
    \multicolumn{3}{c}{\textbf{Contained Features}} \\
    \hline
    \textbf{Feature Name} & \textbf{Feature Id} & \textbf{Comment} \\
    \hline
    Diagnostic Request Correlation & \texttt{cap:bec-53e5c256} (at \ref{cap:bec-53e5c256}) & \\
    \hline
    Diagnostic Report Ingestion \& Binding & \texttt{cap:bec-4b72f961} (at \ref{cap:bec-4b72f961}) & \\
    \hline
    Diagnostic Report Distribution & \texttt{cap:bec-322a3126} (at \ref{cap:bec-322a3126}) & \\
    \hline
    Diagnostic Result Association & \texttt{cap:bec-caa7ba59} (at \ref{cap:bec-caa7ba59}) & \\
    \hline
    Diagnostic Status Tracking & \texttt{cap:bec-67ff7ec6} (at \ref{cap:bec-67ff7ec6}) & \\
    \hline
    Diagnostic Provenance Capture & \texttt{cap:bec-47552fd9} (at \ref{cap:bec-47552fd9}) & \\
    \hline
  \end{tabularx}
\end{table}

Intent. Diagnostic Workflow Management is a candidate enabling capability within Diagnostic Administration. It frames the reusable system-supported behaviours needed to support this healthcare context without prescribing a product, interface, technology or deployment boundary.

Boundary. This capability supports the business activity and does not inherit the authority or accountability of the Business Capability it enables. Its child Features are review candidates until ancestry and semantics are adjudicated.

FHIR R5 conceptual analogues considered (not implementation mappings): ServiceRequest, DiagnosticReport, Observation, Specimen, Task, PlanDefinition, ActivityDefinition.

\subsubsection{Feature - Diagnostic Request Correlation}\label{cap:bec-53e5c256}
\vspace{-1em}
\begin{table}[htbp]
  \centering
  \footnotesize
  \begin{tabularx}{1\textwidth}{|l|l|X|}
    \hline
    \textbf{Metadata Name} & \textbf{Metadata Value} & \textbf{Description} \\
    \hline
    \textbf{Reference Id} & \texttt{BEC-53E5C256} & \\
    \hline
    \textbf{Status} & Existing Feature; candidate CT3 placement & \\
    \hline
    \textbf{Ancestry} &  & \\
    \hline
    \textbf{Source} & FEAT-SA-15 & \\
    \hline
    \textbf{Enabled By} &  & \\
    \hline
    \textbf{Supports} &  & \\
    \hline
    \textbf{Derived From} &  & \\
    \hline
    \textbf{Reviewed Status} & Confirm existing Feature meaning and proposed CT3 ownership. & \\
    \hline
  \end{tabularx}
\end{table}

Intent. Diagnostic Request Correlation provides a distinct enabling behaviour in the context of Diagnostic Workflow Management. This is an existing named Feature; its placement under the candidate CT3 remains subject to review.

\noindent\textbf{Feature boundary.} This behaviour is scoped to Diagnostic Workflow Management; an identically named Feature elsewhere may carry different business meaning.\par
\featuretable{
Function & \texttt{FN-EB6590EF} & Correlate Diagnostic Request & Catalogue \\
Service & \texttt{SV-67719EEB} & Diagnostic Correlation Resolution & Linked to preceding Function in source \\
}
\subsubsection{Feature - Diagnostic Report Ingestion \& Binding}\label{cap:bec-4b72f961}
\vspace{-1em}
\begin{table}[htbp]
  \centering
  \footnotesize
  \begin{tabularx}{1\textwidth}{|l|l|X|}
    \hline
    \textbf{Metadata Name} & \textbf{Metadata Value} & \textbf{Description} \\
    \hline
    \textbf{Reference Id} & \texttt{BEC-4B72F961} & \\
    \hline
    \textbf{Status} & Existing Feature; candidate CT3 placement & \\
    \hline
    \textbf{Ancestry} &  & \\
    \hline
    \textbf{Source} & FEAT-SA-16 & \\
    \hline
    \textbf{Enabled By} &  & \\
    \hline
    \textbf{Supports} &  & \\
    \hline
    \textbf{Derived From} &  & \\
    \hline
    \textbf{Reviewed Status} & Confirm existing Feature meaning and proposed CT3 ownership. & \\
    \hline
  \end{tabularx}
\end{table}

Intent. Diagnostic Report Ingestion \& Binding provides a distinct enabling behaviour in the context of Diagnostic Workflow Management. This is an existing named Feature; its placement under the candidate CT3 remains subject to review.

\noindent\textbf{Feature boundary.} This behaviour is scoped to Diagnostic Workflow Management; an identically named Feature elsewhere may carry different business meaning.\par
\featuretable{
Function & \texttt{FN-90273665} & Bind Diagnostic Report & Catalogue \\
Service & \texttt{SV-1B6935C7} & Diagnostic Report Ingress & Linked to preceding Function in source \\
}
\subsubsection{Feature - Diagnostic Report Distribution}\label{cap:bec-322a3126}
\vspace{-1em}
\begin{table}[htbp]
  \centering
  \footnotesize
  \begin{tabularx}{1\textwidth}{|l|l|X|}
    \hline
    \textbf{Metadata Name} & \textbf{Metadata Value} & \textbf{Description} \\
    \hline
    \textbf{Reference Id} & \texttt{BEC-322A3126} & \\
    \hline
    \textbf{Status} & Existing Feature; candidate CT3 placement & \\
    \hline
    \textbf{Ancestry} &  & \\
    \hline
    \textbf{Source} & FEAT-SA-17 & \\
    \hline
    \textbf{Enabled By} &  & \\
    \hline
    \textbf{Supports} &  & \\
    \hline
    \textbf{Derived From} &  & \\
    \hline
    \textbf{Reviewed Status} & Confirm existing Feature meaning and proposed CT3 ownership. & \\
    \hline
  \end{tabularx}
\end{table}

Intent. Diagnostic Report Distribution provides a distinct enabling behaviour in the context of Diagnostic Workflow Management. This is an existing named Feature; its placement under the candidate CT3 remains subject to review.

\noindent\textbf{Feature boundary.} This behaviour is scoped to Diagnostic Workflow Management; an identically named Feature elsewhere may carry different business meaning.\par
\featuretable{
Function & \texttt{FN-3F9CD0DD} & Distribute Diagnostic Report & Catalogue \\
Service & \texttt{SV-E7528349} & Diagnostic Result Delivery Service & Linked to preceding Function in source \\
}
\subsubsection{Feature - Diagnostic Result Association}\label{cap:bec-caa7ba59}
\vspace{-1em}
\begin{table}[htbp]
  \centering
  \footnotesize
  \begin{tabularx}{1\textwidth}{|l|l|X|}
    \hline
    \textbf{Metadata Name} & \textbf{Metadata Value} & \textbf{Description} \\
    \hline
    \textbf{Reference Id} & \texttt{BEC-CAA7BA59} & \\
    \hline
    \textbf{Status} & Generated candidate & \\
    \hline
    \textbf{Ancestry} &  & \\
    \hline
    \textbf{Source} & FHIR-informed candidate & \\
    \hline
    \textbf{Enabled By} &  & \\
    \hline
    \textbf{Supports} &  & \\
    \hline
    \textbf{Derived From} &  & \\
    \hline
    \textbf{Reviewed Status} & Generated candidate: keep / modify / relocate / reject; verify whether independently testable. & \\
    \hline
  \end{tabularx}
\end{table}

Intent. Diagnostic Result Association is proposed to support the underlying healthcare concept represented by Diagnostic Workflow Management. The proposal concerns behaviour and semantic responsibility, not a FHIR field, operation, resource instance or prescribed data model.

Conceptual evidence: HL7 FHIR R5 ServiceRequest, DiagnosticReport, Observation, Specimen, Task, PlanDefinition, ActivityDefinition.

\noindent\textbf{Feature boundary.} This behaviour is scoped to Diagnostic Workflow Management; an identically named Feature elsewhere may carry different business meaning.\par
\featuretable{
--- & --- & \emph{Not yet captured} & \emph{Complete during review} \\
}
\subsubsection{Feature - Diagnostic Status Tracking}\label{cap:bec-67ff7ec6}
\vspace{-1em}
\begin{table}[htbp]
  \centering
  \footnotesize
  \begin{tabularx}{1\textwidth}{|l|l|X|}
    \hline
    \textbf{Metadata Name} & \textbf{Metadata Value} & \textbf{Description} \\
    \hline
    \textbf{Reference Id} & \texttt{BEC-67FF7EC6} & \\
    \hline
    \textbf{Status} & Generated candidate & \\
    \hline
    \textbf{Ancestry} &  & \\
    \hline
    \textbf{Source} & FHIR-informed candidate & \\
    \hline
    \textbf{Enabled By} &  & \\
    \hline
    \textbf{Supports} &  & \\
    \hline
    \textbf{Derived From} &  & \\
    \hline
    \textbf{Reviewed Status} & Generated candidate: keep / modify / relocate / reject; verify whether independently testable. & \\
    \hline
  \end{tabularx}
\end{table}

Intent. Diagnostic Status Tracking is proposed to support the underlying healthcare concept represented by Diagnostic Workflow Management. The proposal concerns behaviour and semantic responsibility, not a FHIR field, operation, resource instance or prescribed data model.

Conceptual evidence: HL7 FHIR R5 ServiceRequest, DiagnosticReport, Observation, Specimen, Task, PlanDefinition, ActivityDefinition.

\noindent\textbf{Feature boundary.} This behaviour is scoped to Diagnostic Workflow Management; an identically named Feature elsewhere may carry different business meaning.\par
\featuretable{
--- & --- & \emph{Not yet captured} & \emph{Complete during review} \\
}
\subsubsection{Feature - Diagnostic Provenance Capture}\label{cap:bec-47552fd9}
\vspace{-1em}
\begin{table}[htbp]
  \centering
  \footnotesize
  \begin{tabularx}{1\textwidth}{|l|l|X|}
    \hline
    \textbf{Metadata Name} & \textbf{Metadata Value} & \textbf{Description} \\
    \hline
    \textbf{Reference Id} & \texttt{BEC-47552FD9} & \\
    \hline
    \textbf{Status} & Generated candidate & \\
    \hline
    \textbf{Ancestry} &  & \\
    \hline
    \textbf{Source} & FHIR-informed candidate & \\
    \hline
    \textbf{Enabled By} &  & \\
    \hline
    \textbf{Supports} &  & \\
    \hline
    \textbf{Derived From} &  & \\
    \hline
    \textbf{Reviewed Status} & Generated candidate: keep / modify / relocate / reject; verify whether independently testable. & \\
    \hline
  \end{tabularx}
\end{table}

Intent. Diagnostic Provenance Capture is proposed to support the underlying healthcare concept represented by Diagnostic Workflow Management. The proposal concerns behaviour and semantic responsibility, not a FHIR field, operation, resource instance or prescribed data model.

Conceptual evidence: HL7 FHIR R5 ServiceRequest, DiagnosticReport, Observation, Specimen, Task, PlanDefinition, ActivityDefinition.

\noindent\textbf{Feature boundary.} This behaviour is scoped to Diagnostic Workflow Management; an identically named Feature elsewhere may carry different business meaning.\par
\featuretable{
--- & --- & \emph{Not yet captured} & \emph{Complete during review} \\
}

\section{CT2 - Medication Administration}\label{cap:bec-b4278279}
\vspace{-0.5em}
\begin{table}[htbp]
  \centering
  \footnotesize
  \begin{tabularx}{1\textwidth}{|l|l|X|}
    \hline
    \multicolumn{3}{c}{\textbf{Capability Metadata}} \\
    \hline
    \textbf{Name} & \textbf{Value} & \textbf{Comment and/or Description} \\
    \hline
    \textbf{Reference Id} & \texttt{BEC-B4278279} & \\
    \hline
    \textbf{Status} & Candidate ancestry & \\
    \hline
    \textbf{Ancestry} &  & \\
    \hline
    \textbf{Source} &  & \\
    \hline
    \textbf{Enabled By} &  & \\
    \hline
    \textbf{Supports} &  & \\
    \hline
    \textbf{Derived From} &  & \\
    \hline
    \textbf{Reviewed Status} &  & \\
    \hline
  \end{tabularx}
\end{table}

\begin{table}[htbp]
  \centering
  \footnotesize
  \begin{tabularx}{1\textwidth}{|l|l|X|}
    \hline
    \multicolumn{3}{c}{\textbf{Contained Capabilities}} \\
    \hline
    \textbf{Capability Name} & \textbf{Capability Id} & \textbf{Comment} \\
    \hline
    Medication Order \& Dispense Management & \texttt{cap:bec-1efaa849} (at \ref{cap:bec-1efaa849}) & \\
    \hline
  \end{tabularx}
\end{table}

Intent. Medication Administration is a candidate enabling capability within Care Administration Enablement. It frames the reusable system-supported behaviours needed to support this healthcare context without prescribing a product, interface, technology or deployment boundary.

Boundary. This capability supports the business activity and does not inherit the authority or accountability of the Business Capability it enables. Its child Features are review candidates until ancestry and semantics are adjudicated.

\subsection{CT3 - Medication Order \& Dispense Management}\label{cap:bec-1efaa849}
\vspace{-0.5em}
\begin{table}[htbp]
  \centering
  \footnotesize
  \begin{tabularx}{1\textwidth}{|l|l|X|}
    \hline
    \multicolumn{3}{c}{\textbf{Capability Metadata}} \\
    \hline
    \textbf{Name} & \textbf{Value} & \textbf{Comment and/or Description} \\
    \hline
    \textbf{Reference Id} & \texttt{BEC-1EFAA849} & \\
    \hline
    \textbf{Status} & Candidate ancestry & \\
    \hline
    \textbf{Ancestry} &  & \\
    \hline
    \textbf{Source} &  & \\
    \hline
    \textbf{Enabled By} &  & \\
    \hline
    \textbf{Supports} &  & \\
    \hline
    \textbf{Derived From} &  & \\
    \hline
    \textbf{Reviewed Status} &  & \\
    \hline
  \end{tabularx}
\end{table}

\begin{table}[htbp]
  \centering
  \footnotesize
  \begin{tabularx}{1\textwidth}{|l|l|X|}
    \hline
    \multicolumn{3}{c}{\textbf{Contained Features}} \\
    \hline
    \textbf{Feature Name} & \textbf{Feature Id} & \textbf{Comment} \\
    \hline
    Medication Order Ingestion & \texttt{cap:bec-b26d4f17} (at \ref{cap:bec-b26d4f17}) & \\
    \hline
    Dispense Event Tracking & \texttt{cap:bec-0de47133} (at \ref{cap:bec-0de47133}) & \\
    \hline
    Medication Administration Record (MAR) Ingestion & \texttt{cap:bec-d93d1c27} (at \ref{cap:bec-d93d1c27}) & \\
    \hline
    Request Lifecycle Management & \texttt{cap:bec-ab5055a8} (at \ref{cap:bec-ab5055a8}) & \\
    \hline
    Request Fulfilment Tracking & \texttt{cap:bec-e0aba2ef} (at \ref{cap:bec-e0aba2ef}) & \\
    \hline
    Request Status Correlation & \texttt{cap:bec-d9110ce8} (at \ref{cap:bec-d9110ce8}) & \\
    \hline
    Request Exception Management & \texttt{cap:bec-b0e3d03a} (at \ref{cap:bec-b0e3d03a}) & \\
    \hline
    Medication Identity Resolution & \texttt{cap:bec-074e40cb} (at \ref{cap:bec-074e40cb}) & \\
    \hline
    Medication Safety Context Validation & \texttt{cap:bec-afa2a006} (at \ref{cap:bec-afa2a006}) & \\
    \hline
    Medication Activity Tracking & \texttt{cap:bec-9e3ba088} (at \ref{cap:bec-9e3ba088}) & \\
    \hline
    Medication Information Correlation & \texttt{cap:bec-7d9a630b} (at \ref{cap:bec-7d9a630b}) & \\
    \hline
  \end{tabularx}
\end{table}

Intent. Medication Order \& Dispense Management is a candidate enabling capability within Medication Administration. It frames the reusable system-supported behaviours needed to support this healthcare context without prescribing a product, interface, technology or deployment boundary.

Boundary. This capability supports the business activity and does not inherit the authority or accountability of the Business Capability it enables. Its child Features are review candidates until ancestry and semantics are adjudicated.

FHIR R5 conceptual analogues considered (not implementation mappings): ServiceRequest, Task, Medication, MedicationRequest, MedicationDispense, MedicationAdministration.

\subsubsection{Feature - Medication Order Ingestion}\label{cap:bec-b26d4f17}
\vspace{-1em}
\begin{table}[htbp]
  \centering
  \footnotesize
  \begin{tabularx}{1\textwidth}{|l|l|X|}
    \hline
    \textbf{Metadata Name} & \textbf{Metadata Value} & \textbf{Description} \\
    \hline
    \textbf{Reference Id} & \texttt{BEC-B26D4F17} & \\
    \hline
    \textbf{Status} & Existing Feature; candidate CT3 placement & \\
    \hline
    \textbf{Ancestry} &  & \\
    \hline
    \textbf{Source} & FEAT-SA-18 & \\
    \hline
    \textbf{Enabled By} &  & \\
    \hline
    \textbf{Supports} &  & \\
    \hline
    \textbf{Derived From} &  & \\
    \hline
    \textbf{Reviewed Status} & Confirm existing Feature meaning and proposed CT3 ownership. & \\
    \hline
  \end{tabularx}
\end{table}

Intent. Medication Order Ingestion provides a distinct enabling behaviour in the context of Medication Order \& Dispense Management. This is an existing named Feature; its placement under the candidate CT3 remains subject to review.

\noindent\textbf{Feature boundary.} This behaviour is scoped to Medication Order \& Dispense Management; an identically named Feature elsewhere may carry different business meaning.\par
\featuretable{
Function & \texttt{FN-C54DE274} & Receive Medication Order & Catalogue \\
Service & \texttt{SV-AF45BD12} & Medication Order Ingress & Linked to preceding Function in source \\
}
\subsubsection{Feature - Dispense Event Tracking}\label{cap:bec-0de47133}
\vspace{-1em}
\begin{table}[htbp]
  \centering
  \footnotesize
  \begin{tabularx}{1\textwidth}{|l|l|X|}
    \hline
    \textbf{Metadata Name} & \textbf{Metadata Value} & \textbf{Description} \\
    \hline
    \textbf{Reference Id} & \texttt{BEC-0DE47133} & \\
    \hline
    \textbf{Status} & Existing Feature; candidate CT3 placement & \\
    \hline
    \textbf{Ancestry} &  & \\
    \hline
    \textbf{Source} & FEAT-SA-19 & \\
    \hline
    \textbf{Enabled By} &  & \\
    \hline
    \textbf{Supports} &  & \\
    \hline
    \textbf{Derived From} &  & \\
    \hline
    \textbf{Reviewed Status} & Confirm existing Feature meaning and proposed CT3 ownership. & \\
    \hline
  \end{tabularx}
\end{table}

Intent. Dispense Event Tracking provides a distinct enabling behaviour in the context of Medication Order \& Dispense Management. This is an existing named Feature; its placement under the candidate CT3 remains subject to review.

\noindent\textbf{Feature boundary.} This behaviour is scoped to Medication Order \& Dispense Management; an identically named Feature elsewhere may carry different business meaning.\par
\featuretable{
Function & \texttt{FN-E2DC5DD3} & Track Dispense Event & Catalogue \\
Service & \texttt{SV-881532A4} & Dispense Event Ingress & Linked to preceding Function in source \\
}
\subsubsection{Feature - Medication Administration Record (MAR) Ingestion}\label{cap:bec-d93d1c27}
\vspace{-1em}
\begin{table}[htbp]
  \centering
  \footnotesize
  \begin{tabularx}{1\textwidth}{|l|l|X|}
    \hline
    \textbf{Metadata Name} & \textbf{Metadata Value} & \textbf{Description} \\
    \hline
    \textbf{Reference Id} & \texttt{BEC-D93D1C27} & \\
    \hline
    \textbf{Status} & Existing Feature; candidate CT3 placement & \\
    \hline
    \textbf{Ancestry} &  & \\
    \hline
    \textbf{Source} & FEAT-SA-20 & \\
    \hline
    \textbf{Enabled By} &  & \\
    \hline
    \textbf{Supports} &  & \\
    \hline
    \textbf{Derived From} &  & \\
    \hline
    \textbf{Reviewed Status} & Confirm existing Feature meaning and proposed CT3 ownership. & \\
    \hline
  \end{tabularx}
\end{table}

Intent. Medication Administration Record (MAR) Ingestion provides a distinct enabling behaviour in the context of Medication Order \& Dispense Management. This is an existing named Feature; its placement under the candidate CT3 remains subject to review.

\noindent\textbf{Feature boundary.} This behaviour is scoped to Medication Order \& Dispense Management; an identically named Feature elsewhere may carry different business meaning.\par
\featuretable{
Function & \texttt{FN-EA73EE3F} & Receive Medication Administration Event & Catalogue \\
Service & \texttt{SV-9FA97A53} & Medication Administration Ingress & Linked to preceding Function in source \\
}
\subsubsection{Feature - Request Lifecycle Management}\label{cap:bec-ab5055a8}
\vspace{-1em}
\begin{table}[htbp]
  \centering
  \footnotesize
  \begin{tabularx}{1\textwidth}{|l|l|X|}
    \hline
    \textbf{Metadata Name} & \textbf{Metadata Value} & \textbf{Description} \\
    \hline
    \textbf{Reference Id} & \texttt{BEC-AB5055A8} & \\
    \hline
    \textbf{Status} & Generated candidate & \\
    \hline
    \textbf{Ancestry} &  & \\
    \hline
    \textbf{Source} & FHIR-informed candidate & \\
    \hline
    \textbf{Enabled By} &  & \\
    \hline
    \textbf{Supports} &  & \\
    \hline
    \textbf{Derived From} &  & \\
    \hline
    \textbf{Reviewed Status} & Generated candidate: keep / modify / relocate / reject; verify whether independently testable. & \\
    \hline
  \end{tabularx}
\end{table}

Intent. Request Lifecycle Management is proposed to support the underlying healthcare concept represented by Medication Order \& Dispense Management. The proposal concerns behaviour and semantic responsibility, not a FHIR field, operation, resource instance or prescribed data model.

Conceptual evidence: HL7 FHIR R5 ServiceRequest, Task, Medication, MedicationRequest, MedicationDispense, MedicationAdministration.

\noindent\textbf{Feature boundary.} This behaviour is scoped to Medication Order \& Dispense Management; an identically named Feature elsewhere may carry different business meaning.\par
\featuretable{
--- & --- & \emph{Not yet captured} & \emph{Complete during review} \\
}
\subsubsection{Feature - Request Fulfilment Tracking}\label{cap:bec-e0aba2ef}
\vspace{-1em}
\begin{table}[htbp]
  \centering
  \footnotesize
  \begin{tabularx}{1\textwidth}{|l|l|X|}
    \hline
    \textbf{Metadata Name} & \textbf{Metadata Value} & \textbf{Description} \\
    \hline
    \textbf{Reference Id} & \texttt{BEC-E0ABA2EF} & \\
    \hline
    \textbf{Status} & Generated candidate & \\
    \hline
    \textbf{Ancestry} &  & \\
    \hline
    \textbf{Source} & FHIR-informed candidate & \\
    \hline
    \textbf{Enabled By} &  & \\
    \hline
    \textbf{Supports} &  & \\
    \hline
    \textbf{Derived From} &  & \\
    \hline
    \textbf{Reviewed Status} & Generated candidate: keep / modify / relocate / reject; verify whether independently testable. & \\
    \hline
  \end{tabularx}
\end{table}

Intent. Request Fulfilment Tracking is proposed to support the underlying healthcare concept represented by Medication Order \& Dispense Management. The proposal concerns behaviour and semantic responsibility, not a FHIR field, operation, resource instance or prescribed data model.

Conceptual evidence: HL7 FHIR R5 ServiceRequest, Task, Medication, MedicationRequest, MedicationDispense, MedicationAdministration.

\noindent\textbf{Feature boundary.} This behaviour is scoped to Medication Order \& Dispense Management; an identically named Feature elsewhere may carry different business meaning.\par
\featuretable{
--- & --- & \emph{Not yet captured} & \emph{Complete during review} \\
}
\subsubsection{Feature - Request Status Correlation}\label{cap:bec-d9110ce8}
\vspace{-1em}
\begin{table}[htbp]
  \centering
  \footnotesize
  \begin{tabularx}{1\textwidth}{|l|l|X|}
    \hline
    \textbf{Metadata Name} & \textbf{Metadata Value} & \textbf{Description} \\
    \hline
    \textbf{Reference Id} & \texttt{BEC-D9110CE8} & \\
    \hline
    \textbf{Status} & Generated candidate & \\
    \hline
    \textbf{Ancestry} &  & \\
    \hline
    \textbf{Source} & FHIR-informed candidate & \\
    \hline
    \textbf{Enabled By} &  & \\
    \hline
    \textbf{Supports} &  & \\
    \hline
    \textbf{Derived From} &  & \\
    \hline
    \textbf{Reviewed Status} & Generated candidate: keep / modify / relocate / reject; verify whether independently testable. & \\
    \hline
  \end{tabularx}
\end{table}

Intent. Request Status Correlation is proposed to support the underlying healthcare concept represented by Medication Order \& Dispense Management. The proposal concerns behaviour and semantic responsibility, not a FHIR field, operation, resource instance or prescribed data model.

Conceptual evidence: HL7 FHIR R5 ServiceRequest, Task, Medication, MedicationRequest, MedicationDispense, MedicationAdministration.

\noindent\textbf{Feature boundary.} This behaviour is scoped to Medication Order \& Dispense Management; an identically named Feature elsewhere may carry different business meaning.\par
\featuretable{
--- & --- & \emph{Not yet captured} & \emph{Complete during review} \\
}
\subsubsection{Feature - Request Exception Management}\label{cap:bec-b0e3d03a}
\vspace{-1em}
\begin{table}[htbp]
  \centering
  \footnotesize
  \begin{tabularx}{1\textwidth}{|l|l|X|}
    \hline
    \textbf{Metadata Name} & \textbf{Metadata Value} & \textbf{Description} \\
    \hline
    \textbf{Reference Id} & \texttt{BEC-B0E3D03A} & \\
    \hline
    \textbf{Status} & Generated candidate & \\
    \hline
    \textbf{Ancestry} &  & \\
    \hline
    \textbf{Source} & FHIR-informed candidate & \\
    \hline
    \textbf{Enabled By} &  & \\
    \hline
    \textbf{Supports} &  & \\
    \hline
    \textbf{Derived From} &  & \\
    \hline
    \textbf{Reviewed Status} & Generated candidate: keep / modify / relocate / reject; verify whether independently testable. & \\
    \hline
  \end{tabularx}
\end{table}

Intent. Request Exception Management is proposed to support the underlying healthcare concept represented by Medication Order \& Dispense Management. The proposal concerns behaviour and semantic responsibility, not a FHIR field, operation, resource instance or prescribed data model.

Conceptual evidence: HL7 FHIR R5 ServiceRequest, Task, Medication, MedicationRequest, MedicationDispense, MedicationAdministration.

\noindent\textbf{Feature boundary.} This behaviour is scoped to Medication Order \& Dispense Management; an identically named Feature elsewhere may carry different business meaning.\par
\featuretable{
--- & --- & \emph{Not yet captured} & \emph{Complete during review} \\
}
\subsubsection{Feature - Medication Identity Resolution}\label{cap:bec-074e40cb}
\vspace{-1em}
\begin{table}[htbp]
  \centering
  \footnotesize
  \begin{tabularx}{1\textwidth}{|l|l|X|}
    \hline
    \textbf{Metadata Name} & \textbf{Metadata Value} & \textbf{Description} \\
    \hline
    \textbf{Reference Id} & \texttt{BEC-074E40CB} & \\
    \hline
    \textbf{Status} & Generated candidate & \\
    \hline
    \textbf{Ancestry} &  & \\
    \hline
    \textbf{Source} & FHIR-informed candidate & \\
    \hline
    \textbf{Enabled By} &  & \\
    \hline
    \textbf{Supports} &  & \\
    \hline
    \textbf{Derived From} &  & \\
    \hline
    \textbf{Reviewed Status} & Generated candidate: keep / modify / relocate / reject; verify whether independently testable. & \\
    \hline
  \end{tabularx}
\end{table}

Intent. Medication Identity Resolution is proposed to support the underlying healthcare concept represented by Medication Order \& Dispense Management. The proposal concerns behaviour and semantic responsibility, not a FHIR field, operation, resource instance or prescribed data model.

Conceptual evidence: HL7 FHIR R5 ServiceRequest, Task, Medication, MedicationRequest, MedicationDispense, MedicationAdministration.

\noindent\textbf{Feature boundary.} This behaviour is scoped to Medication Order \& Dispense Management; an identically named Feature elsewhere may carry different business meaning.\par
\featuretable{
--- & --- & \emph{Not yet captured} & \emph{Complete during review} \\
}
\subsubsection{Feature - Medication Safety Context Validation}\label{cap:bec-afa2a006}
\vspace{-1em}
\begin{table}[htbp]
  \centering
  \footnotesize
  \begin{tabularx}{1\textwidth}{|l|l|X|}
    \hline
    \textbf{Metadata Name} & \textbf{Metadata Value} & \textbf{Description} \\
    \hline
    \textbf{Reference Id} & \texttt{BEC-AFA2A006} & \\
    \hline
    \textbf{Status} & Generated candidate & \\
    \hline
    \textbf{Ancestry} &  & \\
    \hline
    \textbf{Source} & FHIR-informed candidate & \\
    \hline
    \textbf{Enabled By} &  & \\
    \hline
    \textbf{Supports} &  & \\
    \hline
    \textbf{Derived From} &  & \\
    \hline
    \textbf{Reviewed Status} & Generated candidate: keep / modify / relocate / reject; verify whether independently testable. & \\
    \hline
  \end{tabularx}
\end{table}

Intent. Medication Safety Context Validation is proposed to support the underlying healthcare concept represented by Medication Order \& Dispense Management. The proposal concerns behaviour and semantic responsibility, not a FHIR field, operation, resource instance or prescribed data model.

Conceptual evidence: HL7 FHIR R5 ServiceRequest, Task, Medication, MedicationRequest, MedicationDispense, MedicationAdministration.

\noindent\textbf{Feature boundary.} This behaviour is scoped to Medication Order \& Dispense Management; an identically named Feature elsewhere may carry different business meaning.\par
\featuretable{
--- & --- & \emph{Not yet captured} & \emph{Complete during review} \\
}
\subsubsection{Feature - Medication Activity Tracking}\label{cap:bec-9e3ba088}
\vspace{-1em}
\begin{table}[htbp]
  \centering
  \footnotesize
  \begin{tabularx}{1\textwidth}{|l|l|X|}
    \hline
    \textbf{Metadata Name} & \textbf{Metadata Value} & \textbf{Description} \\
    \hline
    \textbf{Reference Id} & \texttt{BEC-9E3BA088} & \\
    \hline
    \textbf{Status} & Generated candidate & \\
    \hline
    \textbf{Ancestry} &  & \\
    \hline
    \textbf{Source} & FHIR-informed candidate & \\
    \hline
    \textbf{Enabled By} &  & \\
    \hline
    \textbf{Supports} &  & \\
    \hline
    \textbf{Derived From} &  & \\
    \hline
    \textbf{Reviewed Status} & Generated candidate: keep / modify / relocate / reject; verify whether independently testable. & \\
    \hline
  \end{tabularx}
\end{table}

Intent. Medication Activity Tracking is proposed to support the underlying healthcare concept represented by Medication Order \& Dispense Management. The proposal concerns behaviour and semantic responsibility, not a FHIR field, operation, resource instance or prescribed data model.

Conceptual evidence: HL7 FHIR R5 ServiceRequest, Task, Medication, MedicationRequest, MedicationDispense, MedicationAdministration.

\noindent\textbf{Feature boundary.} This behaviour is scoped to Medication Order \& Dispense Management; an identically named Feature elsewhere may carry different business meaning.\par
\featuretable{
--- & --- & \emph{Not yet captured} & \emph{Complete during review} \\
}
\subsubsection{Feature - Medication Information Correlation}\label{cap:bec-7d9a630b}
\vspace{-1em}
\begin{table}[htbp]
  \centering
  \footnotesize
  \begin{tabularx}{1\textwidth}{|l|l|X|}
    \hline
    \textbf{Metadata Name} & \textbf{Metadata Value} & \textbf{Description} \\
    \hline
    \textbf{Reference Id} & \texttt{BEC-7D9A630B} & \\
    \hline
    \textbf{Status} & Generated candidate & \\
    \hline
    \textbf{Ancestry} &  & \\
    \hline
    \textbf{Source} & FHIR-informed candidate & \\
    \hline
    \textbf{Enabled By} &  & \\
    \hline
    \textbf{Supports} &  & \\
    \hline
    \textbf{Derived From} &  & \\
    \hline
    \textbf{Reviewed Status} & Generated candidate: keep / modify / relocate / reject; verify whether independently testable. & \\
    \hline
  \end{tabularx}
\end{table}

Intent. Medication Information Correlation is proposed to support the underlying healthcare concept represented by Medication Order \& Dispense Management. The proposal concerns behaviour and semantic responsibility, not a FHIR field, operation, resource instance or prescribed data model.

Conceptual evidence: HL7 FHIR R5 ServiceRequest, Task, Medication, MedicationRequest, MedicationDispense, MedicationAdministration.

\noindent\textbf{Feature boundary.} This behaviour is scoped to Medication Order \& Dispense Management; an identically named Feature elsewhere may carry different business meaning.\par
\featuretable{
--- & --- & \emph{Not yet captured} & \emph{Complete during review} \\
}

\section{CT2 - Procedure Administration}\label{cap:bec-3233022c}
\vspace{-0.5em}
\begin{table}[htbp]
  \centering
  \footnotesize
  \begin{tabularx}{1\textwidth}{|l|l|X|}
    \hline
    \multicolumn{3}{c}{\textbf{Capability Metadata}} \\
    \hline
    \textbf{Name} & \textbf{Value} & \textbf{Comment and/or Description} \\
    \hline
    \textbf{Reference Id} & \texttt{BEC-3233022C} & \\
    \hline
    \textbf{Status} & Candidate ancestry & \\
    \hline
    \textbf{Ancestry} &  & \\
    \hline
    \textbf{Source} &  & \\
    \hline
    \textbf{Enabled By} &  & \\
    \hline
    \textbf{Supports} &  & \\
    \hline
    \textbf{Derived From} &  & \\
    \hline
    \textbf{Reviewed Status} &  & \\
    \hline
  \end{tabularx}
\end{table}

\begin{table}[htbp]
  \centering
  \footnotesize
  \begin{tabularx}{1\textwidth}{|l|l|X|}
    \hline
    \multicolumn{3}{c}{\textbf{Contained Capabilities}} \\
    \hline
    \textbf{Capability Name} & \textbf{Capability Id} & \textbf{Comment} \\
    \hline
    Procedure \& Surgical Administration & \texttt{cap:bec-489d3945} (at \ref{cap:bec-489d3945}) & \\
    \hline
  \end{tabularx}
\end{table}

Intent. Procedure Administration is a candidate enabling capability within Care Administration Enablement. It frames the reusable system-supported behaviours needed to support this healthcare context without prescribing a product, interface, technology or deployment boundary.

Boundary. This capability supports the business activity and does not inherit the authority or accountability of the Business Capability it enables. Its child Features are review candidates until ancestry and semantics are adjudicated.

\subsection{CT3 - Procedure \& Surgical Administration}\label{cap:bec-489d3945}
\vspace{-0.5em}
\begin{table}[htbp]
  \centering
  \footnotesize
  \begin{tabularx}{1\textwidth}{|l|l|X|}
    \hline
    \multicolumn{3}{c}{\textbf{Capability Metadata}} \\
    \hline
    \textbf{Name} & \textbf{Value} & \textbf{Comment and/or Description} \\
    \hline
    \textbf{Reference Id} & \texttt{BEC-489D3945} & \\
    \hline
    \textbf{Status} & Candidate ancestry & \\
    \hline
    \textbf{Ancestry} &  & \\
    \hline
    \textbf{Source} &  & \\
    \hline
    \textbf{Enabled By} &  & \\
    \hline
    \textbf{Supports} &  & \\
    \hline
    \textbf{Derived From} &  & \\
    \hline
    \textbf{Reviewed Status} &  & \\
    \hline
  \end{tabularx}
\end{table}

\begin{table}[htbp]
  \centering
  \footnotesize
  \begin{tabularx}{1\textwidth}{|l|l|X|}
    \hline
    \multicolumn{3}{c}{\textbf{Contained Features}} \\
    \hline
    \textbf{Feature Name} & \textbf{Feature Id} & \textbf{Comment} \\
    \hline
    Procedure Booking Ingestion & \texttt{cap:bec-f696528f} (at \ref{cap:bec-f696528f}) & \\
    \hline
    Procedural Documentation Ingestion & \texttt{cap:bec-75a8d392} (at \ref{cap:bec-75a8d392}) & \\
    \hline
    Context Establishment & \texttt{cap:bec-5dbeff67} (at \ref{cap:bec-5dbeff67}) & \\
    \hline
    State Monitoring & \texttt{cap:bec-6100d197} (at \ref{cap:bec-6100d197}) & \\
    \hline
    Exception Management & \texttt{cap:bec-82e3935c} (at \ref{cap:bec-82e3935c}) & \\
    \hline
    Information Retrieval & \texttt{cap:bec-c925ba63} (at \ref{cap:bec-c925ba63}) & \\
    \hline
  \end{tabularx}
\end{table}

Intent. Procedure \& Surgical Administration is a candidate enabling capability within Procedure Administration. It frames the reusable system-supported behaviours needed to support this healthcare context without prescribing a product, interface, technology or deployment boundary.

Boundary. This capability supports the business activity and does not inherit the authority or accountability of the Business Capability it enables. Its child Features are review candidates until ancestry and semantics are adjudicated.

\subsubsection{Feature - Procedure Booking Ingestion}\label{cap:bec-f696528f}
\vspace{-1em}
\begin{table}[htbp]
  \centering
  \footnotesize
  \begin{tabularx}{1\textwidth}{|l|l|X|}
    \hline
    \textbf{Metadata Name} & \textbf{Metadata Value} & \textbf{Description} \\
    \hline
    \textbf{Reference Id} & \texttt{BEC-F696528F} & \\
    \hline
    \textbf{Status} & Existing Feature; candidate CT3 placement & \\
    \hline
    \textbf{Ancestry} &  & \\
    \hline
    \textbf{Source} & FEAT-SA-21 & \\
    \hline
    \textbf{Enabled By} &  & \\
    \hline
    \textbf{Supports} &  & \\
    \hline
    \textbf{Derived From} &  & \\
    \hline
    \textbf{Reviewed Status} & Confirm existing Feature meaning and proposed CT3 ownership. & \\
    \hline
  \end{tabularx}
\end{table}

Intent. Procedure Booking Ingestion provides a distinct enabling behaviour in the context of Procedure \& Surgical Administration. This is an existing named Feature; its placement under the candidate CT3 remains subject to review.

\noindent\textbf{Feature boundary.} This behaviour is scoped to Procedure \& Surgical Administration; an identically named Feature elsewhere may carry different business meaning.\par
\featuretable{
Function & \texttt{FN-729F406C} & Receive Procedure Booking / Request & Catalogue \\
Service & \texttt{SV-29B7741A} & Procedure Booking Ingress & Linked to preceding Function in source \\
}
\subsubsection{Feature - Procedural Documentation Ingestion}\label{cap:bec-75a8d392}
\vspace{-1em}
\begin{table}[htbp]
  \centering
  \footnotesize
  \begin{tabularx}{1\textwidth}{|l|l|X|}
    \hline
    \textbf{Metadata Name} & \textbf{Metadata Value} & \textbf{Description} \\
    \hline
    \textbf{Reference Id} & \texttt{BEC-75A8D392} & \\
    \hline
    \textbf{Status} & Existing Feature; candidate CT3 placement & \\
    \hline
    \textbf{Ancestry} &  & \\
    \hline
    \textbf{Source} & FEAT-SA-22 & \\
    \hline
    \textbf{Enabled By} &  & \\
    \hline
    \textbf{Supports} &  & \\
    \hline
    \textbf{Derived From} &  & \\
    \hline
    \textbf{Reviewed Status} & Confirm existing Feature meaning and proposed CT3 ownership. & \\
    \hline
  \end{tabularx}
\end{table}

Intent. Procedural Documentation Ingestion provides a distinct enabling behaviour in the context of Procedure \& Surgical Administration. This is an existing named Feature; its placement under the candidate CT3 remains subject to review.

\noindent\textbf{Feature boundary.} This behaviour is scoped to Procedure \& Surgical Administration; an identically named Feature elsewhere may carry different business meaning.\par
\featuretable{
Function & \texttt{FN-68F8C215} & Receive Procedural Documentation & Catalogue \\
Service & \texttt{SV-C60ECC83} & Procedural Documentation Ingress & Linked to preceding Function in source \\
}
\subsubsection{Feature - Context Establishment}\label{cap:bec-5dbeff67}
\vspace{-1em}
\begin{table}[htbp]
  \centering
  \footnotesize
  \begin{tabularx}{1\textwidth}{|l|l|X|}
    \hline
    \textbf{Metadata Name} & \textbf{Metadata Value} & \textbf{Description} \\
    \hline
    \textbf{Reference Id} & \texttt{BEC-5DBEFF67} & \\
    \hline
    \textbf{Status} & Generated candidate & \\
    \hline
    \textbf{Ancestry} &  & \\
    \hline
    \textbf{Source} & Semantic-pattern candidate & \\
    \hline
    \textbf{Enabled By} &  & \\
    \hline
    \textbf{Supports} &  & \\
    \hline
    \textbf{Derived From} &  & \\
    \hline
    \textbf{Reviewed Status} & Generated candidate: keep / modify / relocate / reject; verify whether independently testable. & \\
    \hline
  \end{tabularx}
\end{table}

Intent. Context Establishment is proposed as a discrete, testable enabling behaviour for Procedure \& Surgical Administration. It is included for architectural review, not as an established implementation obligation.

\noindent\textbf{Feature boundary.} This behaviour is scoped to Procedure \& Surgical Administration; an identically named Feature elsewhere may carry different business meaning.\par
\featuretable{
--- & --- & \emph{Not yet captured} & \emph{Complete during review} \\
}
\subsubsection{Feature - State Monitoring}\label{cap:bec-6100d197}
\vspace{-1em}
\begin{table}[htbp]
  \centering
  \footnotesize
  \begin{tabularx}{1\textwidth}{|l|l|X|}
    \hline
    \textbf{Metadata Name} & \textbf{Metadata Value} & \textbf{Description} \\
    \hline
    \textbf{Reference Id} & \texttt{BEC-6100D197} & \\
    \hline
    \textbf{Status} & Generated candidate & \\
    \hline
    \textbf{Ancestry} &  & \\
    \hline
    \textbf{Source} & Semantic-pattern candidate & \\
    \hline
    \textbf{Enabled By} &  & \\
    \hline
    \textbf{Supports} &  & \\
    \hline
    \textbf{Derived From} &  & \\
    \hline
    \textbf{Reviewed Status} & Generated candidate: keep / modify / relocate / reject; verify whether independently testable. & \\
    \hline
  \end{tabularx}
\end{table}

Intent. State Monitoring is proposed as a discrete, testable enabling behaviour for Procedure \& Surgical Administration. It is included for architectural review, not as an established implementation obligation.

\noindent\textbf{Feature boundary.} This behaviour is scoped to Procedure \& Surgical Administration; an identically named Feature elsewhere may carry different business meaning.\par
\featuretable{
--- & --- & \emph{Not yet captured} & \emph{Complete during review} \\
}
\subsubsection{Feature - Exception Management}\label{cap:bec-82e3935c}
\vspace{-1em}
\begin{table}[htbp]
  \centering
  \footnotesize
  \begin{tabularx}{1\textwidth}{|l|l|X|}
    \hline
    \textbf{Metadata Name} & \textbf{Metadata Value} & \textbf{Description} \\
    \hline
    \textbf{Reference Id} & \texttt{BEC-82E3935C} & \\
    \hline
    \textbf{Status} & Generated candidate & \\
    \hline
    \textbf{Ancestry} &  & \\
    \hline
    \textbf{Source} & Semantic-pattern candidate & \\
    \hline
    \textbf{Enabled By} &  & \\
    \hline
    \textbf{Supports} &  & \\
    \hline
    \textbf{Derived From} &  & \\
    \hline
    \textbf{Reviewed Status} & Generated candidate: keep / modify / relocate / reject; verify whether independently testable. & \\
    \hline
  \end{tabularx}
\end{table}

Intent. Exception Management is proposed as a discrete, testable enabling behaviour for Procedure \& Surgical Administration. It is included for architectural review, not as an established implementation obligation.

\noindent\textbf{Feature boundary.} This behaviour is scoped to Procedure \& Surgical Administration; an identically named Feature elsewhere may carry different business meaning.\par
\featuretable{
--- & --- & \emph{Not yet captured} & \emph{Complete during review} \\
}
\subsubsection{Feature - Information Retrieval}\label{cap:bec-c925ba63}
\vspace{-1em}
\begin{table}[htbp]
  \centering
  \footnotesize
  \begin{tabularx}{1\textwidth}{|l|l|X|}
    \hline
    \textbf{Metadata Name} & \textbf{Metadata Value} & \textbf{Description} \\
    \hline
    \textbf{Reference Id} & \texttt{BEC-C925BA63} & \\
    \hline
    \textbf{Status} & Generated candidate & \\
    \hline
    \textbf{Ancestry} &  & \\
    \hline
    \textbf{Source} & Semantic-pattern candidate & \\
    \hline
    \textbf{Enabled By} &  & \\
    \hline
    \textbf{Supports} &  & \\
    \hline
    \textbf{Derived From} &  & \\
    \hline
    \textbf{Reviewed Status} & Generated candidate: keep / modify / relocate / reject; verify whether independently testable. & \\
    \hline
  \end{tabularx}
\end{table}

Intent. Information Retrieval is proposed as a discrete, testable enabling behaviour for Procedure \& Surgical Administration. It is included for architectural review, not as an established implementation obligation.

\noindent\textbf{Feature boundary.} This behaviour is scoped to Procedure \& Surgical Administration; an identically named Feature elsewhere may carry different business meaning.\par
\featuretable{
--- & --- & \emph{Not yet captured} & \emph{Complete during review} \\
}

\section{CT2 - Clinical Record Administration}\label{cap:bec-bf6253dc}
\vspace{-0.5em}
\begin{table}[htbp]
  \centering
  \footnotesize
  \begin{tabularx}{1\textwidth}{|l|l|X|}
    \hline
    \multicolumn{3}{c}{\textbf{Capability Metadata}} \\
    \hline
    \textbf{Name} & \textbf{Value} & \textbf{Comment and/or Description} \\
    \hline
    \textbf{Reference Id} & \texttt{BEC-BF6253DC} & \\
    \hline
    \textbf{Status} & Candidate ancestry & \\
    \hline
    \textbf{Ancestry} &  & \\
    \hline
    \textbf{Source} &  & \\
    \hline
    \textbf{Enabled By} &  & \\
    \hline
    \textbf{Supports} &  & \\
    \hline
    \textbf{Derived From} &  & \\
    \hline
    \textbf{Reviewed Status} &  & \\
    \hline
  \end{tabularx}
\end{table}

\begin{table}[htbp]
  \centering
  \footnotesize
  \begin{tabularx}{1\textwidth}{|l|l|X|}
    \hline
    \multicolumn{3}{c}{\textbf{Contained Capabilities}} \\
    \hline
    \textbf{Capability Name} & \textbf{Capability Id} & \textbf{Comment} \\
    \hline
    Clinical Documentation Lifecycle & \texttt{cap:bec-0ce70156} (at \ref{cap:bec-0ce70156}) & \\
    \hline
  \end{tabularx}
\end{table}

Intent. Clinical Record Administration is a candidate enabling capability within Care Administration Enablement. It frames the reusable system-supported behaviours needed to support this healthcare context without prescribing a product, interface, technology or deployment boundary.

Boundary. This capability supports the business activity and does not inherit the authority or accountability of the Business Capability it enables. Its child Features are review candidates until ancestry and semantics are adjudicated.

\subsection{CT3 - Clinical Documentation Lifecycle}\label{cap:bec-0ce70156}
\vspace{-0.5em}
\begin{table}[htbp]
  \centering
  \footnotesize
  \begin{tabularx}{1\textwidth}{|l|l|X|}
    \hline
    \multicolumn{3}{c}{\textbf{Capability Metadata}} \\
    \hline
    \textbf{Name} & \textbf{Value} & \textbf{Comment and/or Description} \\
    \hline
    \textbf{Reference Id} & \texttt{BEC-0CE70156} & \\
    \hline
    \textbf{Status} & Candidate ancestry & \\
    \hline
    \textbf{Ancestry} &  & \\
    \hline
    \textbf{Source} &  & \\
    \hline
    \textbf{Enabled By} &  & \\
    \hline
    \textbf{Supports} &  & \\
    \hline
    \textbf{Derived From} &  & \\
    \hline
    \textbf{Reviewed Status} &  & \\
    \hline
  \end{tabularx}
\end{table}

\begin{table}[htbp]
  \centering
  \footnotesize
  \begin{tabularx}{1\textwidth}{|l|l|X|}
    \hline
    \multicolumn{3}{c}{\textbf{Contained Features}} \\
    \hline
    \textbf{Feature Name} & \textbf{Feature Id} & \textbf{Comment} \\
    \hline
    Clinical Document Ingestion & \texttt{cap:bec-096b53fa} (at \ref{cap:bec-096b53fa}) & \\
    \hline
    Document Versioning \& Supersession & \texttt{cap:bec-a9a0c02b} (at \ref{cap:bec-a9a0c02b}) & \\
    \hline
    Document Metadata Indexing & \texttt{cap:bec-2963c113} (at \ref{cap:bec-2963c113}) & \\
    \hline
    Context Establishment & \texttt{cap:bec-836351ba} (at \ref{cap:bec-836351ba}) & \\
    \hline
    State Monitoring & \texttt{cap:bec-cf085c02} (at \ref{cap:bec-cf085c02}) & \\
    \hline
    Exception Management & \texttt{cap:bec-a27f8453} (at \ref{cap:bec-a27f8453}) & \\
    \hline
    Information Retrieval & \texttt{cap:bec-fdc92a6d} (at \ref{cap:bec-fdc92a6d}) & \\
    \hline
  \end{tabularx}
\end{table}

Intent. Clinical Documentation Lifecycle is a candidate enabling capability within Clinical Record Administration. It frames the reusable system-supported behaviours needed to support this healthcare context without prescribing a product, interface, technology or deployment boundary.

Boundary. This capability supports the business activity and does not inherit the authority or accountability of the Business Capability it enables. Its child Features are review candidates until ancestry and semantics are adjudicated.

\subsubsection{Feature - Clinical Document Ingestion}\label{cap:bec-096b53fa}
\vspace{-1em}
\begin{table}[htbp]
  \centering
  \footnotesize
  \begin{tabularx}{1\textwidth}{|l|l|X|}
    \hline
    \textbf{Metadata Name} & \textbf{Metadata Value} & \textbf{Description} \\
    \hline
    \textbf{Reference Id} & \texttt{BEC-096B53FA} & \\
    \hline
    \textbf{Status} & Existing Feature; candidate CT3 placement & \\
    \hline
    \textbf{Ancestry} &  & \\
    \hline
    \textbf{Source} & FEAT-SA-23 & \\
    \hline
    \textbf{Enabled By} &  & \\
    \hline
    \textbf{Supports} &  & \\
    \hline
    \textbf{Derived From} &  & \\
    \hline
    \textbf{Reviewed Status} & Confirm existing Feature meaning and proposed CT3 ownership. & \\
    \hline
  \end{tabularx}
\end{table}

Intent. Clinical Document Ingestion provides a distinct enabling behaviour in the context of Clinical Documentation Lifecycle. This is an existing named Feature; its placement under the candidate CT3 remains subject to review.

\noindent\textbf{Feature boundary.} This behaviour is scoped to Clinical Documentation Lifecycle; an identically named Feature elsewhere may carry different business meaning.\par
\featuretable{
Function & \texttt{FN-186C1C96} & Receive Clinical Document & Catalogue \\
Service & \texttt{SV-141449ED} & Clinical Document Ingress & Linked to preceding Function in source \\
}
\subsubsection{Feature - Document Versioning \& Supersession}\label{cap:bec-a9a0c02b}
\vspace{-1em}
\begin{table}[htbp]
  \centering
  \footnotesize
  \begin{tabularx}{1\textwidth}{|l|l|X|}
    \hline
    \textbf{Metadata Name} & \textbf{Metadata Value} & \textbf{Description} \\
    \hline
    \textbf{Reference Id} & \texttt{BEC-A9A0C02B} & \\
    \hline
    \textbf{Status} & Existing Feature; candidate CT3 placement & \\
    \hline
    \textbf{Ancestry} &  & \\
    \hline
    \textbf{Source} & FEAT-SA-24 & \\
    \hline
    \textbf{Enabled By} &  & \\
    \hline
    \textbf{Supports} &  & \\
    \hline
    \textbf{Derived From} &  & \\
    \hline
    \textbf{Reviewed Status} & Confirm existing Feature meaning and proposed CT3 ownership. & \\
    \hline
  \end{tabularx}
\end{table}

Intent. Document Versioning \& Supersession provides a distinct enabling behaviour in the context of Clinical Documentation Lifecycle. This is an existing named Feature; its placement under the candidate CT3 remains subject to review.

\noindent\textbf{Feature boundary.} This behaviour is scoped to Clinical Documentation Lifecycle; an identically named Feature elsewhere may carry different business meaning.\par
\featuretable{
Function & \texttt{FN-7494D50B} & Govern Clinical Document Lifecycle & Catalogue \\
Service & \texttt{SV-DB2EAD83} & Clinical Document Lifecycle Service & Linked to preceding Function in source \\
Process & \texttt{PR-0D877B3D} & Clinical Document Lifecycle Process & Linked to preceding Function in source \\
}
\subsubsection{Feature - Document Metadata Indexing}\label{cap:bec-2963c113}
\vspace{-1em}
\begin{table}[htbp]
  \centering
  \footnotesize
  \begin{tabularx}{1\textwidth}{|l|l|X|}
    \hline
    \textbf{Metadata Name} & \textbf{Metadata Value} & \textbf{Description} \\
    \hline
    \textbf{Reference Id} & \texttt{BEC-2963C113} & \\
    \hline
    \textbf{Status} & Existing Feature; candidate CT3 placement & \\
    \hline
    \textbf{Ancestry} &  & \\
    \hline
    \textbf{Source} & FEAT-SA-25 & \\
    \hline
    \textbf{Enabled By} &  & \\
    \hline
    \textbf{Supports} &  & \\
    \hline
    \textbf{Derived From} &  & \\
    \hline
    \textbf{Reviewed Status} & Confirm existing Feature meaning and proposed CT3 ownership. & \\
    \hline
  \end{tabularx}
\end{table}

Intent. Document Metadata Indexing provides a distinct enabling behaviour in the context of Clinical Documentation Lifecycle. This is an existing named Feature; its placement under the candidate CT3 remains subject to review.

\noindent\textbf{Feature boundary.} This behaviour is scoped to Clinical Documentation Lifecycle; an identically named Feature elsewhere may carry different business meaning.\par
\featuretable{
Function & \texttt{FN-B00AB3C2} & Index Clinical Document Metadata & Catalogue \\
Service & \texttt{SV-4EF821E0} & Document Metadata Registry Query & Linked to preceding Function in source \\
}
\subsubsection{Feature - Context Establishment}\label{cap:bec-836351ba}
\vspace{-1em}
\begin{table}[htbp]
  \centering
  \footnotesize
  \begin{tabularx}{1\textwidth}{|l|l|X|}
    \hline
    \textbf{Metadata Name} & \textbf{Metadata Value} & \textbf{Description} \\
    \hline
    \textbf{Reference Id} & \texttt{BEC-836351BA} & \\
    \hline
    \textbf{Status} & Generated candidate & \\
    \hline
    \textbf{Ancestry} &  & \\
    \hline
    \textbf{Source} & Semantic-pattern candidate & \\
    \hline
    \textbf{Enabled By} &  & \\
    \hline
    \textbf{Supports} &  & \\
    \hline
    \textbf{Derived From} &  & \\
    \hline
    \textbf{Reviewed Status} & Generated candidate: keep / modify / relocate / reject; verify whether independently testable. & \\
    \hline
  \end{tabularx}
\end{table}

Intent. Context Establishment is proposed as a discrete, testable enabling behaviour for Clinical Documentation Lifecycle. It is included for architectural review, not as an established implementation obligation.

\noindent\textbf{Feature boundary.} This behaviour is scoped to Clinical Documentation Lifecycle; an identically named Feature elsewhere may carry different business meaning.\par
\featuretable{
--- & --- & \emph{Not yet captured} & \emph{Complete during review} \\
}
\subsubsection{Feature - State Monitoring}\label{cap:bec-cf085c02}
\vspace{-1em}
\begin{table}[htbp]
  \centering
  \footnotesize
  \begin{tabularx}{1\textwidth}{|l|l|X|}
    \hline
    \textbf{Metadata Name} & \textbf{Metadata Value} & \textbf{Description} \\
    \hline
    \textbf{Reference Id} & \texttt{BEC-CF085C02} & \\
    \hline
    \textbf{Status} & Generated candidate & \\
    \hline
    \textbf{Ancestry} &  & \\
    \hline
    \textbf{Source} & Semantic-pattern candidate & \\
    \hline
    \textbf{Enabled By} &  & \\
    \hline
    \textbf{Supports} &  & \\
    \hline
    \textbf{Derived From} &  & \\
    \hline
    \textbf{Reviewed Status} & Generated candidate: keep / modify / relocate / reject; verify whether independently testable. & \\
    \hline
  \end{tabularx}
\end{table}

Intent. State Monitoring is proposed as a discrete, testable enabling behaviour for Clinical Documentation Lifecycle. It is included for architectural review, not as an established implementation obligation.

\noindent\textbf{Feature boundary.} This behaviour is scoped to Clinical Documentation Lifecycle; an identically named Feature elsewhere may carry different business meaning.\par
\featuretable{
--- & --- & \emph{Not yet captured} & \emph{Complete during review} \\
}
\subsubsection{Feature - Exception Management}\label{cap:bec-a27f8453}
\vspace{-1em}
\begin{table}[htbp]
  \centering
  \footnotesize
  \begin{tabularx}{1\textwidth}{|l|l|X|}
    \hline
    \textbf{Metadata Name} & \textbf{Metadata Value} & \textbf{Description} \\
    \hline
    \textbf{Reference Id} & \texttt{BEC-A27F8453} & \\
    \hline
    \textbf{Status} & Generated candidate & \\
    \hline
    \textbf{Ancestry} &  & \\
    \hline
    \textbf{Source} & Semantic-pattern candidate & \\
    \hline
    \textbf{Enabled By} &  & \\
    \hline
    \textbf{Supports} &  & \\
    \hline
    \textbf{Derived From} &  & \\
    \hline
    \textbf{Reviewed Status} & Generated candidate: keep / modify / relocate / reject; verify whether independently testable. & \\
    \hline
  \end{tabularx}
\end{table}

Intent. Exception Management is proposed as a discrete, testable enabling behaviour for Clinical Documentation Lifecycle. It is included for architectural review, not as an established implementation obligation.

\noindent\textbf{Feature boundary.} This behaviour is scoped to Clinical Documentation Lifecycle; an identically named Feature elsewhere may carry different business meaning.\par
\featuretable{
--- & --- & \emph{Not yet captured} & \emph{Complete during review} \\
}
\subsubsection{Feature - Information Retrieval}\label{cap:bec-fdc92a6d}
\vspace{-1em}
\begin{table}[htbp]
  \centering
  \footnotesize
  \begin{tabularx}{1\textwidth}{|l|l|X|}
    \hline
    \textbf{Metadata Name} & \textbf{Metadata Value} & \textbf{Description} \\
    \hline
    \textbf{Reference Id} & \texttt{BEC-FDC92A6D} & \\
    \hline
    \textbf{Status} & Generated candidate & \\
    \hline
    \textbf{Ancestry} &  & \\
    \hline
    \textbf{Source} & Semantic-pattern candidate & \\
    \hline
    \textbf{Enabled By} &  & \\
    \hline
    \textbf{Supports} &  & \\
    \hline
    \textbf{Derived From} &  & \\
    \hline
    \textbf{Reviewed Status} & Generated candidate: keep / modify / relocate / reject; verify whether independently testable. & \\
    \hline
  \end{tabularx}
\end{table}

Intent. Information Retrieval is proposed as a discrete, testable enabling behaviour for Clinical Documentation Lifecycle. It is included for architectural review, not as an established implementation obligation.

\noindent\textbf{Feature boundary.} This behaviour is scoped to Clinical Documentation Lifecycle; an identically named Feature elsewhere may carry different business meaning.\par
\featuretable{
--- & --- & \emph{Not yet captured} & \emph{Complete during review} \\
}

\section{CT2 - Care Plan Administration}\label{cap:bec-389062f6}
\vspace{-0.5em}
\begin{table}[htbp]
  \centering
  \footnotesize
  \begin{tabularx}{1\textwidth}{|l|l|X|}
    \hline
    \multicolumn{3}{c}{\textbf{Capability Metadata}} \\
    \hline
    \textbf{Name} & \textbf{Value} & \textbf{Comment and/or Description} \\
    \hline
    \textbf{Reference Id} & \texttt{BEC-389062F6} & \\
    \hline
    \textbf{Status} & Candidate ancestry & \\
    \hline
    \textbf{Ancestry} &  & \\
    \hline
    \textbf{Source} &  & \\
    \hline
    \textbf{Enabled By} &  & \\
    \hline
    \textbf{Supports} &  & \\
    \hline
    \textbf{Derived From} &  & \\
    \hline
    \textbf{Reviewed Status} &  & \\
    \hline
  \end{tabularx}
\end{table}

\begin{table}[htbp]
  \centering
  \footnotesize
  \begin{tabularx}{1\textwidth}{|l|l|X|}
    \hline
    \multicolumn{3}{c}{\textbf{Contained Capabilities}} \\
    \hline
    \textbf{Capability Name} & \textbf{Capability Id} & \textbf{Comment} \\
    \hline
    Care Plan Coordination & \texttt{cap:bec-e263e2d1} (at \ref{cap:bec-e263e2d1}) & \\
    \hline
  \end{tabularx}
\end{table}

Intent. Care Plan Administration is a candidate enabling capability within Care Administration Enablement. It frames the reusable system-supported behaviours needed to support this healthcare context without prescribing a product, interface, technology or deployment boundary.

Boundary. This capability supports the business activity and does not inherit the authority or accountability of the Business Capability it enables. Its child Features are review candidates until ancestry and semantics are adjudicated.

\subsection{CT3 - Care Plan Coordination}\label{cap:bec-e263e2d1}
\vspace{-0.5em}
\begin{table}[htbp]
  \centering
  \footnotesize
  \begin{tabularx}{1\textwidth}{|l|l|X|}
    \hline
    \multicolumn{3}{c}{\textbf{Capability Metadata}} \\
    \hline
    \textbf{Name} & \textbf{Value} & \textbf{Comment and/or Description} \\
    \hline
    \textbf{Reference Id} & \texttt{BEC-E263E2D1} & \\
    \hline
    \textbf{Status} & Candidate ancestry & \\
    \hline
    \textbf{Ancestry} &  & \\
    \hline
    \textbf{Source} &  & \\
    \hline
    \textbf{Enabled By} &  & \\
    \hline
    \textbf{Supports} &  & \\
    \hline
    \textbf{Derived From} &  & \\
    \hline
    \textbf{Reviewed Status} &  & \\
    \hline
  \end{tabularx}
\end{table}

\begin{table}[htbp]
  \centering
  \footnotesize
  \begin{tabularx}{1\textwidth}{|l|l|X|}
    \hline
    \multicolumn{3}{c}{\textbf{Contained Features}} \\
    \hline
    \textbf{Feature Name} & \textbf{Feature Id} & \textbf{Comment} \\
    \hline
    Care Plan Ingestion \& Distribution & \texttt{cap:bec-e3837ab5} (at \ref{cap:bec-e3837ab5}) & \\
    \hline
    Care Goal Association & \texttt{cap:bec-d4e51095} (at \ref{cap:bec-d4e51095}) & \\
    \hline
    Care Plan Lifecycle Tracking & \texttt{cap:bec-6ccf8e04} (at \ref{cap:bec-6ccf8e04}) & \\
    \hline
    Care Plan Context Retrieval & \texttt{cap:bec-065e070b} (at \ref{cap:bec-065e070b}) & \\
    \hline
  \end{tabularx}
\end{table}

Intent. Care Plan Coordination is a candidate enabling capability within Care Plan Administration. It frames the reusable system-supported behaviours needed to support this healthcare context without prescribing a product, interface, technology or deployment boundary.

Boundary. This capability supports the business activity and does not inherit the authority or accountability of the Business Capability it enables. Its child Features are review candidates until ancestry and semantics are adjudicated.

FHIR R5 conceptual analogues considered (not implementation mappings): CarePlan, Goal.

\subsubsection{Feature - Care Plan Ingestion \& Distribution}\label{cap:bec-e3837ab5}
\vspace{-1em}
\begin{table}[htbp]
  \centering
  \footnotesize
  \begin{tabularx}{1\textwidth}{|l|l|X|}
    \hline
    \textbf{Metadata Name} & \textbf{Metadata Value} & \textbf{Description} \\
    \hline
    \textbf{Reference Id} & \texttt{BEC-E3837AB5} & \\
    \hline
    \textbf{Status} & Existing Feature; candidate CT3 placement & \\
    \hline
    \textbf{Ancestry} &  & \\
    \hline
    \textbf{Source} & FEAT-SA-26 & \\
    \hline
    \textbf{Enabled By} &  & \\
    \hline
    \textbf{Supports} &  & \\
    \hline
    \textbf{Derived From} &  & \\
    \hline
    \textbf{Reviewed Status} & Confirm existing Feature meaning and proposed CT3 ownership. & \\
    \hline
  \end{tabularx}
\end{table}

Intent. Care Plan Ingestion \& Distribution provides a distinct enabling behaviour in the context of Care Plan Coordination. This is an existing named Feature; its placement under the candidate CT3 remains subject to review.

\noindent\textbf{Feature boundary.} This behaviour is scoped to Care Plan Coordination; an identically named Feature elsewhere may carry different business meaning.\par
\featuretable{
Function & \texttt{FN-B3DDFE19} & Receive Care Plan & Catalogue \\
Service & \texttt{SV-9FC888C0} & Care Plan Ingress & Linked to preceding Function in source \\
}
\subsubsection{Feature - Care Goal Association}\label{cap:bec-d4e51095}
\vspace{-1em}
\begin{table}[htbp]
  \centering
  \footnotesize
  \begin{tabularx}{1\textwidth}{|l|l|X|}
    \hline
    \textbf{Metadata Name} & \textbf{Metadata Value} & \textbf{Description} \\
    \hline
    \textbf{Reference Id} & \texttt{BEC-D4E51095} & \\
    \hline
    \textbf{Status} & Generated candidate & \\
    \hline
    \textbf{Ancestry} &  & \\
    \hline
    \textbf{Source} & FHIR-informed candidate & \\
    \hline
    \textbf{Enabled By} &  & \\
    \hline
    \textbf{Supports} &  & \\
    \hline
    \textbf{Derived From} &  & \\
    \hline
    \textbf{Reviewed Status} & Generated candidate: keep / modify / relocate / reject; verify whether independently testable. & \\
    \hline
  \end{tabularx}
\end{table}

Intent. Care Goal Association is proposed to support the underlying healthcare concept represented by Care Plan Coordination. The proposal concerns behaviour and semantic responsibility, not a FHIR field, operation, resource instance or prescribed data model.

Conceptual evidence: HL7 FHIR R5 CarePlan, Goal.

\noindent\textbf{Feature boundary.} This behaviour is scoped to Care Plan Coordination; an identically named Feature elsewhere may carry different business meaning.\par
\featuretable{
--- & --- & \emph{Not yet captured} & \emph{Complete during review} \\
}
\subsubsection{Feature - Care Plan Lifecycle Tracking}\label{cap:bec-6ccf8e04}
\vspace{-1em}
\begin{table}[htbp]
  \centering
  \footnotesize
  \begin{tabularx}{1\textwidth}{|l|l|X|}
    \hline
    \textbf{Metadata Name} & \textbf{Metadata Value} & \textbf{Description} \\
    \hline
    \textbf{Reference Id} & \texttt{BEC-6CCF8E04} & \\
    \hline
    \textbf{Status} & Generated candidate & \\
    \hline
    \textbf{Ancestry} &  & \\
    \hline
    \textbf{Source} & FHIR-informed candidate & \\
    \hline
    \textbf{Enabled By} &  & \\
    \hline
    \textbf{Supports} &  & \\
    \hline
    \textbf{Derived From} &  & \\
    \hline
    \textbf{Reviewed Status} & Generated candidate: keep / modify / relocate / reject; verify whether independently testable. & \\
    \hline
  \end{tabularx}
\end{table}

Intent. Care Plan Lifecycle Tracking is proposed to support the underlying healthcare concept represented by Care Plan Coordination. The proposal concerns behaviour and semantic responsibility, not a FHIR field, operation, resource instance or prescribed data model.

Conceptual evidence: HL7 FHIR R5 CarePlan, Goal.

\noindent\textbf{Feature boundary.} This behaviour is scoped to Care Plan Coordination; an identically named Feature elsewhere may carry different business meaning.\par
\featuretable{
--- & --- & \emph{Not yet captured} & \emph{Complete during review} \\
}
\subsubsection{Feature - Care Plan Context Retrieval}\label{cap:bec-065e070b}
\vspace{-1em}
\begin{table}[htbp]
  \centering
  \footnotesize
  \begin{tabularx}{1\textwidth}{|l|l|X|}
    \hline
    \textbf{Metadata Name} & \textbf{Metadata Value} & \textbf{Description} \\
    \hline
    \textbf{Reference Id} & \texttt{BEC-065E070B} & \\
    \hline
    \textbf{Status} & Generated candidate & \\
    \hline
    \textbf{Ancestry} &  & \\
    \hline
    \textbf{Source} & FHIR-informed candidate & \\
    \hline
    \textbf{Enabled By} &  & \\
    \hline
    \textbf{Supports} &  & \\
    \hline
    \textbf{Derived From} &  & \\
    \hline
    \textbf{Reviewed Status} & Generated candidate: keep / modify / relocate / reject; verify whether independently testable. & \\
    \hline
  \end{tabularx}
\end{table}

Intent. Care Plan Context Retrieval is proposed to support the underlying healthcare concept represented by Care Plan Coordination. The proposal concerns behaviour and semantic responsibility, not a FHIR field, operation, resource instance or prescribed data model.

Conceptual evidence: HL7 FHIR R5 CarePlan, Goal.

\noindent\textbf{Feature boundary.} This behaviour is scoped to Care Plan Coordination; an identically named Feature elsewhere may carry different business meaning.\par
\featuretable{
--- & --- & \emph{Not yet captured} & \emph{Complete during review} \\
}

\section{CT2 - Clinical Communication Administration}\label{cap:bec-0b71e5cb}
\vspace{-0.5em}
\begin{table}[htbp]
  \centering
  \footnotesize
  \begin{tabularx}{1\textwidth}{|l|l|X|}
    \hline
    \multicolumn{3}{c}{\textbf{Capability Metadata}} \\
    \hline
    \textbf{Name} & \textbf{Value} & \textbf{Comment and/or Description} \\
    \hline
    \textbf{Reference Id} & \texttt{BEC-0B71E5CB} & \\
    \hline
    \textbf{Status} & Candidate ancestry & \\
    \hline
    \textbf{Ancestry} &  & \\
    \hline
    \textbf{Source} &  & \\
    \hline
    \textbf{Enabled By} &  & \\
    \hline
    \textbf{Supports} &  & \\
    \hline
    \textbf{Derived From} &  & \\
    \hline
    \textbf{Reviewed Status} &  & \\
    \hline
  \end{tabularx}
\end{table}

\begin{table}[htbp]
  \centering
  \footnotesize
  \begin{tabularx}{1\textwidth}{|l|l|X|}
    \hline
    \multicolumn{3}{c}{\textbf{Contained Capabilities}} \\
    \hline
    \textbf{Capability Name} & \textbf{Capability Id} & \textbf{Comment} \\
    \hline
    Secure Clinical Communication & \texttt{cap:bec-6d84c2c1} (at \ref{cap:bec-6d84c2c1}) & \\
    \hline
  \end{tabularx}
\end{table}

Intent. Clinical Communication Administration is a candidate enabling capability within Care Administration Enablement. It frames the reusable system-supported behaviours needed to support this healthcare context without prescribing a product, interface, technology or deployment boundary.

Boundary. This capability supports the business activity and does not inherit the authority or accountability of the Business Capability it enables. Its child Features are review candidates until ancestry and semantics are adjudicated.

\subsection{CT3 - Secure Clinical Communication}\label{cap:bec-6d84c2c1}
\vspace{-0.5em}
\begin{table}[htbp]
  \centering
  \footnotesize
  \begin{tabularx}{1\textwidth}{|l|l|X|}
    \hline
    \multicolumn{3}{c}{\textbf{Capability Metadata}} \\
    \hline
    \textbf{Name} & \textbf{Value} & \textbf{Comment and/or Description} \\
    \hline
    \textbf{Reference Id} & \texttt{BEC-6D84C2C1} & \\
    \hline
    \textbf{Status} & Candidate ancestry & \\
    \hline
    \textbf{Ancestry} &  & \\
    \hline
    \textbf{Source} &  & \\
    \hline
    \textbf{Enabled By} &  & \\
    \hline
    \textbf{Supports} &  & \\
    \hline
    \textbf{Derived From} &  & \\
    \hline
    \textbf{Reviewed Status} &  & \\
    \hline
  \end{tabularx}
\end{table}

\begin{table}[htbp]
  \centering
  \footnotesize
  \begin{tabularx}{1\textwidth}{|l|l|X|}
    \hline
    \multicolumn{3}{c}{\textbf{Contained Features}} \\
    \hline
    \textbf{Feature Name} & \textbf{Feature Id} & \textbf{Comment} \\
    \hline
    Clinical Communication Distribution & \texttt{cap:bec-c55cedbd} (at \ref{cap:bec-c55cedbd}) & \\
    \hline
    Delivery Acknowledgement Tracking & \texttt{cap:bec-357afa04} (at \ref{cap:bec-357afa04}) & \\
    \hline
    Communication Audit Logging & \texttt{cap:bec-4966229d} (at \ref{cap:bec-4966229d}) & \\
    \hline
    Communication Context Management & \texttt{cap:bec-bee7ff8b} (at \ref{cap:bec-bee7ff8b}) & \\
    \hline
    Recipient Resolution & \texttt{cap:bec-b1e25071} (at \ref{cap:bec-b1e25071}) & \\
    \hline
    Communication Delivery Tracking & \texttt{cap:bec-206d85c4} (at \ref{cap:bec-206d85c4}) & \\
    \hline
  \end{tabularx}
\end{table}

Intent. Secure Clinical Communication is a candidate enabling capability within Clinical Communication Administration. It frames the reusable system-supported behaviours needed to support this healthcare context without prescribing a product, interface, technology or deployment boundary.

Boundary. This capability supports the business activity and does not inherit the authority or accountability of the Business Capability it enables. Its child Features are review candidates until ancestry and semantics are adjudicated.

FHIR R5 conceptual analogues considered (not implementation mappings): Communication, CommunicationRequest.

\subsubsection{Feature - Clinical Communication Distribution}\label{cap:bec-c55cedbd}
\vspace{-1em}
\begin{table}[htbp]
  \centering
  \footnotesize
  \begin{tabularx}{1\textwidth}{|l|l|X|}
    \hline
    \textbf{Metadata Name} & \textbf{Metadata Value} & \textbf{Description} \\
    \hline
    \textbf{Reference Id} & \texttt{BEC-C55CEDBD} & \\
    \hline
    \textbf{Status} & Existing Feature; candidate CT3 placement & \\
    \hline
    \textbf{Ancestry} &  & \\
    \hline
    \textbf{Source} & FEAT-SA-27 & \\
    \hline
    \textbf{Enabled By} &  & \\
    \hline
    \textbf{Supports} &  & \\
    \hline
    \textbf{Derived From} &  & \\
    \hline
    \textbf{Reviewed Status} & Confirm existing Feature meaning and proposed CT3 ownership. & \\
    \hline
  \end{tabularx}
\end{table}

Intent. Clinical Communication Distribution provides a distinct enabling behaviour in the context of Secure Clinical Communication. This is an existing named Feature; its placement under the candidate CT3 remains subject to review.

\noindent\textbf{Feature boundary.} This behaviour is scoped to Secure Clinical Communication; an identically named Feature elsewhere may carry different business meaning.\par
\featuretable{
Function & \texttt{FN-1752C25D} & Distribute Clinical Communication & Catalogue \\
Service & \texttt{SV-52F3D83D} & Secure Clinical Message Dispatch & Linked to preceding Function in source \\
}
\subsubsection{Feature - Delivery Acknowledgement Tracking}\label{cap:bec-357afa04}
\vspace{-1em}
\begin{table}[htbp]
  \centering
  \footnotesize
  \begin{tabularx}{1\textwidth}{|l|l|X|}
    \hline
    \textbf{Metadata Name} & \textbf{Metadata Value} & \textbf{Description} \\
    \hline
    \textbf{Reference Id} & \texttt{BEC-357AFA04} & \\
    \hline
    \textbf{Status} & Existing Feature; candidate CT3 placement & \\
    \hline
    \textbf{Ancestry} &  & \\
    \hline
    \textbf{Source} & FEAT-SA-28 & \\
    \hline
    \textbf{Enabled By} &  & \\
    \hline
    \textbf{Supports} &  & \\
    \hline
    \textbf{Derived From} &  & \\
    \hline
    \textbf{Reviewed Status} & Confirm existing Feature meaning and proposed CT3 ownership. & \\
    \hline
  \end{tabularx}
\end{table}

Intent. Delivery Acknowledgement Tracking provides a distinct enabling behaviour in the context of Secure Clinical Communication. This is an existing named Feature; its placement under the candidate CT3 remains subject to review.

\noindent\textbf{Feature boundary.} This behaviour is scoped to Secure Clinical Communication; an identically named Feature elsewhere may carry different business meaning.\par
\featuretable{
Function & \texttt{FN-EF7B4794} & Track Business Delivery Acknowledgement & Catalogue \\
Service & \texttt{SV-B8181DBE} & Communication Delivery Status Query & Linked to preceding Function in source \\
}
\subsubsection{Feature - Communication Audit Logging}\label{cap:bec-4966229d}
\vspace{-1em}
\begin{table}[htbp]
  \centering
  \footnotesize
  \begin{tabularx}{1\textwidth}{|l|l|X|}
    \hline
    \textbf{Metadata Name} & \textbf{Metadata Value} & \textbf{Description} \\
    \hline
    \textbf{Reference Id} & \texttt{BEC-4966229D} & \\
    \hline
    \textbf{Status} & Existing Feature; candidate CT3 placement & \\
    \hline
    \textbf{Ancestry} &  & \\
    \hline
    \textbf{Source} & FEAT-SA-29 & \\
    \hline
    \textbf{Enabled By} &  & \\
    \hline
    \textbf{Supports} &  & \\
    \hline
    \textbf{Derived From} &  & \\
    \hline
    \textbf{Reviewed Status} & Confirm existing Feature meaning and proposed CT3 ownership. & \\
    \hline
  \end{tabularx}
\end{table}

Intent. Communication Audit Logging provides a distinct enabling behaviour in the context of Secure Clinical Communication. This is an existing named Feature; its placement under the candidate CT3 remains subject to review.

\noindent\textbf{Feature boundary.} This behaviour is scoped to Secure Clinical Communication; an identically named Feature elsewhere may carry different business meaning.\par
\featuretable{
Function & \texttt{FN-9DEBD5A8} & Record Communication Transaction & Catalogue \\
Service & \texttt{SV-C31DD24C} & Communication Audit Query & Linked to preceding Function in source \\
}
\subsubsection{Feature - Communication Context Management}\label{cap:bec-bee7ff8b}
\vspace{-1em}
\begin{table}[htbp]
  \centering
  \footnotesize
  \begin{tabularx}{1\textwidth}{|l|l|X|}
    \hline
    \textbf{Metadata Name} & \textbf{Metadata Value} & \textbf{Description} \\
    \hline
    \textbf{Reference Id} & \texttt{BEC-BEE7FF8B} & \\
    \hline
    \textbf{Status} & Generated candidate & \\
    \hline
    \textbf{Ancestry} &  & \\
    \hline
    \textbf{Source} & FHIR-informed candidate & \\
    \hline
    \textbf{Enabled By} &  & \\
    \hline
    \textbf{Supports} &  & \\
    \hline
    \textbf{Derived From} &  & \\
    \hline
    \textbf{Reviewed Status} & Generated candidate: keep / modify / relocate / reject; verify whether independently testable. & \\
    \hline
  \end{tabularx}
\end{table}

Intent. Communication Context Management is proposed to support the underlying healthcare concept represented by Secure Clinical Communication. The proposal concerns behaviour and semantic responsibility, not a FHIR field, operation, resource instance or prescribed data model.

Conceptual evidence: HL7 FHIR R5 Communication, CommunicationRequest.

\noindent\textbf{Feature boundary.} This behaviour is scoped to Secure Clinical Communication; an identically named Feature elsewhere may carry different business meaning.\par
\featuretable{
--- & --- & \emph{Not yet captured} & \emph{Complete during review} \\
}
\subsubsection{Feature - Recipient Resolution}\label{cap:bec-b1e25071}
\vspace{-1em}
\begin{table}[htbp]
  \centering
  \footnotesize
  \begin{tabularx}{1\textwidth}{|l|l|X|}
    \hline
    \textbf{Metadata Name} & \textbf{Metadata Value} & \textbf{Description} \\
    \hline
    \textbf{Reference Id} & \texttt{BEC-B1E25071} & \\
    \hline
    \textbf{Status} & Generated candidate & \\
    \hline
    \textbf{Ancestry} &  & \\
    \hline
    \textbf{Source} & FHIR-informed candidate & \\
    \hline
    \textbf{Enabled By} &  & \\
    \hline
    \textbf{Supports} &  & \\
    \hline
    \textbf{Derived From} &  & \\
    \hline
    \textbf{Reviewed Status} & Generated candidate: keep / modify / relocate / reject; verify whether independently testable. & \\
    \hline
  \end{tabularx}
\end{table}

Intent. Recipient Resolution is proposed to support the underlying healthcare concept represented by Secure Clinical Communication. The proposal concerns behaviour and semantic responsibility, not a FHIR field, operation, resource instance or prescribed data model.

Conceptual evidence: HL7 FHIR R5 Communication, CommunicationRequest.

\noindent\textbf{Feature boundary.} This behaviour is scoped to Secure Clinical Communication; an identically named Feature elsewhere may carry different business meaning.\par
\featuretable{
--- & --- & \emph{Not yet captured} & \emph{Complete during review} \\
}
\subsubsection{Feature - Communication Delivery Tracking}\label{cap:bec-206d85c4}
\vspace{-1em}
\begin{table}[htbp]
  \centering
  \footnotesize
  \begin{tabularx}{1\textwidth}{|l|l|X|}
    \hline
    \textbf{Metadata Name} & \textbf{Metadata Value} & \textbf{Description} \\
    \hline
    \textbf{Reference Id} & \texttt{BEC-206D85C4} & \\
    \hline
    \textbf{Status} & Generated candidate & \\
    \hline
    \textbf{Ancestry} &  & \\
    \hline
    \textbf{Source} & FHIR-informed candidate & \\
    \hline
    \textbf{Enabled By} &  & \\
    \hline
    \textbf{Supports} &  & \\
    \hline
    \textbf{Derived From} &  & \\
    \hline
    \textbf{Reviewed Status} & Generated candidate: keep / modify / relocate / reject; verify whether independently testable. & \\
    \hline
  \end{tabularx}
\end{table}

Intent. Communication Delivery Tracking is proposed to support the underlying healthcare concept represented by Secure Clinical Communication. The proposal concerns behaviour and semantic responsibility, not a FHIR field, operation, resource instance or prescribed data model.

Conceptual evidence: HL7 FHIR R5 Communication, CommunicationRequest.

\noindent\textbf{Feature boundary.} This behaviour is scoped to Secure Clinical Communication; an identically named Feature elsewhere may carry different business meaning.\par
\featuretable{
--- & --- & \emph{Not yet captured} & \emph{Complete during review} \\
}

\section{CT2 - Funding \& Coverage Administration}\label{cap:bec-09d4a645}
\vspace{-0.5em}
\begin{table}[htbp]
  \centering
  \footnotesize
  \begin{tabularx}{1\textwidth}{|l|l|X|}
    \hline
    \multicolumn{3}{c}{\textbf{Capability Metadata}} \\
    \hline
    \textbf{Name} & \textbf{Value} & \textbf{Comment and/or Description} \\
    \hline
    \textbf{Reference Id} & \texttt{BEC-09D4A645} & \\
    \hline
    \textbf{Status} & Candidate ancestry & \\
    \hline
    \textbf{Ancestry} &  & \\
    \hline
    \textbf{Source} &  & \\
    \hline
    \textbf{Enabled By} &  & \\
    \hline
    \textbf{Supports} &  & \\
    \hline
    \textbf{Derived From} &  & \\
    \hline
    \textbf{Reviewed Status} &  & \\
    \hline
  \end{tabularx}
\end{table}

\begin{table}[htbp]
  \centering
  \footnotesize
  \begin{tabularx}{1\textwidth}{|l|l|X|}
    \hline
    \multicolumn{3}{c}{\textbf{Contained Capabilities}} \\
    \hline
    \textbf{Capability Name} & \textbf{Capability Id} & \textbf{Comment} \\
    \hline
    \multicolumn{3}{|l|}{\emph{No contained capabilities are currently identified.}} \\
    \hline
  \end{tabularx}
\end{table}

Intent. Funding \& Coverage Administration is a candidate enabling capability within Care Administration Enablement. It frames the reusable system-supported behaviours needed to support this healthcare context without prescribing a product, interface, technology or deployment boundary.

Boundary. This capability supports the business activity and does not inherit the authority or accountability of the Business Capability it enables. Its child Features are review candidates until ancestry and semantics are adjudicated.

\section{CT2 - Clinical Charging \& Claims Administration}\label{cap:bec-47d10eb4}
\vspace{-0.5em}
\begin{table}[htbp]
  \centering
  \footnotesize
  \begin{tabularx}{1\textwidth}{|l|l|X|}
    \hline
    \multicolumn{3}{c}{\textbf{Capability Metadata}} \\
    \hline
    \textbf{Name} & \textbf{Value} & \textbf{Comment and/or Description} \\
    \hline
    \textbf{Reference Id} & \texttt{BEC-47D10EB4} & \\
    \hline
    \textbf{Status} & Candidate ancestry & \\
    \hline
    \textbf{Ancestry} &  & \\
    \hline
    \textbf{Source} &  & \\
    \hline
    \textbf{Enabled By} &  & \\
    \hline
    \textbf{Supports} &  & \\
    \hline
    \textbf{Derived From} &  & \\
    \hline
    \textbf{Reviewed Status} &  & \\
    \hline
  \end{tabularx}
\end{table}

\begin{table}[htbp]
  \centering
  \footnotesize
  \begin{tabularx}{1\textwidth}{|l|l|X|}
    \hline
    \multicolumn{3}{c}{\textbf{Contained Capabilities}} \\
    \hline
    \textbf{Capability Name} & \textbf{Capability Id} & \textbf{Comment} \\
    \hline
    \multicolumn{3}{|l|}{\emph{No contained capabilities are currently identified.}} \\
    \hline
  \end{tabularx}
\end{table}

Intent. Clinical Charging \& Claims Administration is a candidate enabling capability within Care Administration Enablement. It frames the reusable system-supported behaviours needed to support this healthcare context without prescribing a product, interface, technology or deployment boundary.

Boundary. This capability supports the business activity and does not inherit the authority or accountability of the Business Capability it enables. Its child Features are review candidates until ancestry and semantics are adjudicated.

\chapter{CT1 - Clinical Service Context Enablement}\label{cap:bec-3a4524b6}
\noindent\textbf{Reference ID:} \texttt{BEC-3A4524B6}\quad \textbf{Status:} Candidate ancestry\par
Intent. Clinical Service Context Enablement is a candidate enabling capability within Business Enabling Capability layer. It frames the reusable system-supported behaviours needed to support this healthcare context without prescribing a product, interface, technology or deployment boundary.

Boundary. This capability supports the business activity and does not inherit the authority or accountability of the Business Capability it enables. Its child Features are review candidates until ancestry and semantics are adjudicated.

\section{CT2 - Clinical Assessment}\label{cap:bec-aeab9e52}
\vspace{-0.5em}
\begin{table}[htbp]
  \centering
  \footnotesize
  \begin{tabularx}{1\textwidth}{|l|l|X|}
    \hline
    \multicolumn{3}{c}{\textbf{Capability Metadata}} \\
    \hline
    \textbf{Name} & \textbf{Value} & \textbf{Comment and/or Description} \\
    \hline
    \textbf{Reference Id} & \texttt{BEC-AEAB9E52} & \\
    \hline
    \textbf{Status} & Candidate ancestry & \\
    \hline
    \textbf{Ancestry} &  & \\
    \hline
    \textbf{Source} &  & \\
    \hline
    \textbf{Enabled By} &  & \\
    \hline
    \textbf{Supports} &  & \\
    \hline
    \textbf{Derived From} &  & \\
    \hline
    \textbf{Reviewed Status} &  & \\
    \hline
  \end{tabularx}
\end{table}

\begin{table}[htbp]
  \centering
  \footnotesize
  \begin{tabularx}{1\textwidth}{|l|l|X|}
    \hline
    \multicolumn{3}{c}{\textbf{Contained Capabilities}} \\
    \hline
    \textbf{Capability Name} & \textbf{Capability Id} & \textbf{Comment} \\
    \hline
    \multicolumn{3}{|l|}{\emph{No contained capabilities are currently identified.}} \\
    \hline
  \end{tabularx}
\end{table}

Intent. Clinical Assessment is a candidate enabling capability within Clinical Service Context Enablement. It frames the reusable system-supported behaviours needed to support this healthcare context without prescribing a product, interface, technology or deployment boundary.

Boundary. This capability supports the business activity and does not inherit the authority or accountability of the Business Capability it enables. Its child Features are review candidates until ancestry and semantics are adjudicated.

\section{CT2 - Primary Care}\label{cap:bec-14fd8edd}
\vspace{-0.5em}
\begin{table}[htbp]
  \centering
  \footnotesize
  \begin{tabularx}{1\textwidth}{|l|l|X|}
    \hline
    \multicolumn{3}{c}{\textbf{Capability Metadata}} \\
    \hline
    \textbf{Name} & \textbf{Value} & \textbf{Comment and/or Description} \\
    \hline
    \textbf{Reference Id} & \texttt{BEC-14FD8EDD} & \\
    \hline
    \textbf{Status} & Candidate ancestry & \\
    \hline
    \textbf{Ancestry} &  & \\
    \hline
    \textbf{Source} &  & \\
    \hline
    \textbf{Enabled By} &  & \\
    \hline
    \textbf{Supports} &  & \\
    \hline
    \textbf{Derived From} &  & \\
    \hline
    \textbf{Reviewed Status} &  & \\
    \hline
  \end{tabularx}
\end{table}

\begin{table}[htbp]
  \centering
  \footnotesize
  \begin{tabularx}{1\textwidth}{|l|l|X|}
    \hline
    \multicolumn{3}{c}{\textbf{Contained Capabilities}} \\
    \hline
    \textbf{Capability Name} & \textbf{Capability Id} & \textbf{Comment} \\
    \hline
    Primary Care Coordination & \texttt{cap:bec-2d7a6dde} (at \ref{cap:bec-2d7a6dde}) & \\
    \hline
  \end{tabularx}
\end{table}

Intent. Primary Care is a candidate enabling capability within Clinical Service Context Enablement. It frames the reusable system-supported behaviours needed to support this healthcare context without prescribing a product, interface, technology or deployment boundary.

Boundary. This capability supports the business activity and does not inherit the authority or accountability of the Business Capability it enables. Its child Features are review candidates until ancestry and semantics are adjudicated.

\subsection{CT3 - Primary Care Coordination}\label{cap:bec-2d7a6dde}
\vspace{-0.5em}
\begin{table}[htbp]
  \centering
  \footnotesize
  \begin{tabularx}{1\textwidth}{|l|l|X|}
    \hline
    \multicolumn{3}{c}{\textbf{Capability Metadata}} \\
    \hline
    \textbf{Name} & \textbf{Value} & \textbf{Comment and/or Description} \\
    \hline
    \textbf{Reference Id} & \texttt{BEC-2D7A6DDE} & \\
    \hline
    \textbf{Status} & Candidate ancestry & \\
    \hline
    \textbf{Ancestry} &  & \\
    \hline
    \textbf{Source} &  & \\
    \hline
    \textbf{Enabled By} &  & \\
    \hline
    \textbf{Supports} &  & \\
    \hline
    \textbf{Derived From} &  & \\
    \hline
    \textbf{Reviewed Status} &  & \\
    \hline
  \end{tabularx}
\end{table}

\begin{table}[htbp]
  \centering
  \footnotesize
  \begin{tabularx}{1\textwidth}{|l|l|X|}
    \hline
    \multicolumn{3}{c}{\textbf{Contained Features}} \\
    \hline
    \textbf{Feature Name} & \textbf{Feature Id} & \textbf{Comment} \\
    \hline
    General Practice Shared Record Ingestion & \texttt{cap:bec-ab31e6bc} (at \ref{cap:bec-ab31e6bc}) & \\
    \hline
    Primary Care Notification Dispatch & \texttt{cap:bec-e7f50e85} (at \ref{cap:bec-e7f50e85}) & \\
    \hline
    Context Establishment & \texttt{cap:bec-db56036e} (at \ref{cap:bec-db56036e}) & \\
    \hline
    Eligibility Validation & \texttt{cap:bec-8a07fdb4} (at \ref{cap:bec-8a07fdb4}) & \\
    \hline
    Assignment \& Routing & \texttt{cap:bec-ca466b56} (at \ref{cap:bec-ca466b56}) & \\
    \hline
    State Tracking & \texttt{cap:bec-a945b4f5} (at \ref{cap:bec-a945b4f5}) & \\
    \hline
    Exception Management & \texttt{cap:bec-2737b30c} (at \ref{cap:bec-2737b30c}) & \\
    \hline
    Outcome Reporting & \texttt{cap:bec-f6fc078e} (at \ref{cap:bec-f6fc078e}) & \\
    \hline
  \end{tabularx}
\end{table}

Intent. Primary Care Coordination is a candidate enabling capability within Primary Care. It frames the reusable system-supported behaviours needed to support this healthcare context without prescribing a product, interface, technology or deployment boundary.

Boundary. This capability supports the business activity and does not inherit the authority or accountability of the Business Capability it enables. Its child Features are review candidates until ancestry and semantics are adjudicated.

\subsubsection{Feature - General Practice Shared Record Ingestion}\label{cap:bec-ab31e6bc}
\vspace{-1em}
\begin{table}[htbp]
  \centering
  \footnotesize
  \begin{tabularx}{1\textwidth}{|l|l|X|}
    \hline
    \textbf{Metadata Name} & \textbf{Metadata Value} & \textbf{Description} \\
    \hline
    \textbf{Reference Id} & \texttt{BEC-AB31E6BC} & \\
    \hline
    \textbf{Status} & Existing Feature; candidate CT3 placement & \\
    \hline
    \textbf{Ancestry} &  & \\
    \hline
    \textbf{Source} & FEAT-SD-01 & \\
    \hline
    \textbf{Enabled By} &  & \\
    \hline
    \textbf{Supports} &  & \\
    \hline
    \textbf{Derived From} &  & \\
    \hline
    \textbf{Reviewed Status} & Confirm existing Feature meaning and proposed CT3 ownership. & \\
    \hline
  \end{tabularx}
\end{table}

Intent. General Practice Shared Record Ingestion provides a distinct enabling behaviour in the context of Primary Care Coordination. This is an existing named Feature; its placement under the candidate CT3 remains subject to review.

\noindent\textbf{Feature boundary.} This behaviour is scoped to Primary Care Coordination; an identically named Feature elsewhere may carry different business meaning.\par
\featuretable{
Function & \texttt{FN-E574D667} & Correlate Primary Care Consultation Context & Catalogue \\
Service & \texttt{SV-CAE2DB8E} & Primary Care Summary Ingress & Linked to preceding Function in source \\
}
\subsubsection{Feature - Primary Care Notification Dispatch}\label{cap:bec-e7f50e85}
\vspace{-1em}
\begin{table}[htbp]
  \centering
  \footnotesize
  \begin{tabularx}{1\textwidth}{|l|l|X|}
    \hline
    \textbf{Metadata Name} & \textbf{Metadata Value} & \textbf{Description} \\
    \hline
    \textbf{Reference Id} & \texttt{BEC-E7F50E85} & \\
    \hline
    \textbf{Status} & Existing Feature; candidate CT3 placement & \\
    \hline
    \textbf{Ancestry} &  & \\
    \hline
    \textbf{Source} & FEAT-SD-02 & \\
    \hline
    \textbf{Enabled By} &  & \\
    \hline
    \textbf{Supports} &  & \\
    \hline
    \textbf{Derived From} &  & \\
    \hline
    \textbf{Reviewed Status} & Confirm existing Feature meaning and proposed CT3 ownership. & \\
    \hline
  \end{tabularx}
\end{table}

Intent. Primary Care Notification Dispatch provides a distinct enabling behaviour in the context of Primary Care Coordination. This is an existing named Feature; its placement under the candidate CT3 remains subject to review.

\noindent\textbf{Feature boundary.} This behaviour is scoped to Primary Care Coordination; an identically named Feature elsewhere may carry different business meaning.\par
\featuretable{
Function & \texttt{FN-C26BD265} & Coordinate Multi-Agency Chronic Care Plan & Catalogue \\
Service & \texttt{SV-D1870C50} & Chronic Care Timeline Query & Linked to preceding Function in source \\
}
\subsubsection{Feature - Context Establishment}\label{cap:bec-db56036e}
\vspace{-1em}
\begin{table}[htbp]
  \centering
  \footnotesize
  \begin{tabularx}{1\textwidth}{|l|l|X|}
    \hline
    \textbf{Metadata Name} & \textbf{Metadata Value} & \textbf{Description} \\
    \hline
    \textbf{Reference Id} & \texttt{BEC-DB56036E} & \\
    \hline
    \textbf{Status} & Generated candidate & \\
    \hline
    \textbf{Ancestry} &  & \\
    \hline
    \textbf{Source} & Semantic-pattern candidate & \\
    \hline
    \textbf{Enabled By} &  & \\
    \hline
    \textbf{Supports} &  & \\
    \hline
    \textbf{Derived From} &  & \\
    \hline
    \textbf{Reviewed Status} & Generated candidate: keep / modify / relocate / reject; verify whether independently testable. & \\
    \hline
  \end{tabularx}
\end{table}

Intent. Context Establishment is proposed as a discrete, testable enabling behaviour for Primary Care Coordination. It is included for architectural review, not as an established implementation obligation.

\noindent\textbf{Feature boundary.} This behaviour is scoped to Primary Care Coordination; an identically named Feature elsewhere may carry different business meaning.\par
\featuretable{
--- & --- & \emph{Not yet captured} & \emph{Complete during review} \\
}
\subsubsection{Feature - Eligibility Validation}\label{cap:bec-8a07fdb4}
\vspace{-1em}
\begin{table}[htbp]
  \centering
  \footnotesize
  \begin{tabularx}{1\textwidth}{|l|l|X|}
    \hline
    \textbf{Metadata Name} & \textbf{Metadata Value} & \textbf{Description} \\
    \hline
    \textbf{Reference Id} & \texttt{BEC-8A07FDB4} & \\
    \hline
    \textbf{Status} & Generated candidate & \\
    \hline
    \textbf{Ancestry} &  & \\
    \hline
    \textbf{Source} & Semantic-pattern candidate & \\
    \hline
    \textbf{Enabled By} &  & \\
    \hline
    \textbf{Supports} &  & \\
    \hline
    \textbf{Derived From} &  & \\
    \hline
    \textbf{Reviewed Status} & Generated candidate: keep / modify / relocate / reject; verify whether independently testable. & \\
    \hline
  \end{tabularx}
\end{table}

Intent. Eligibility Validation is proposed as a discrete, testable enabling behaviour for Primary Care Coordination. It is included for architectural review, not as an established implementation obligation.

\noindent\textbf{Feature boundary.} This behaviour is scoped to Primary Care Coordination; an identically named Feature elsewhere may carry different business meaning.\par
\featuretable{
--- & --- & \emph{Not yet captured} & \emph{Complete during review} \\
}
\subsubsection{Feature - Assignment \& Routing}\label{cap:bec-ca466b56}
\vspace{-1em}
\begin{table}[htbp]
  \centering
  \footnotesize
  \begin{tabularx}{1\textwidth}{|l|l|X|}
    \hline
    \textbf{Metadata Name} & \textbf{Metadata Value} & \textbf{Description} \\
    \hline
    \textbf{Reference Id} & \texttt{BEC-CA466B56} & \\
    \hline
    \textbf{Status} & Generated candidate & \\
    \hline
    \textbf{Ancestry} &  & \\
    \hline
    \textbf{Source} & Semantic-pattern candidate & \\
    \hline
    \textbf{Enabled By} &  & \\
    \hline
    \textbf{Supports} &  & \\
    \hline
    \textbf{Derived From} &  & \\
    \hline
    \textbf{Reviewed Status} & Generated candidate: keep / modify / relocate / reject; verify whether independently testable. & \\
    \hline
  \end{tabularx}
\end{table}

Intent. Assignment \& Routing is proposed as a discrete, testable enabling behaviour for Primary Care Coordination. It is included for architectural review, not as an established implementation obligation.

\noindent\textbf{Feature boundary.} This behaviour is scoped to Primary Care Coordination; an identically named Feature elsewhere may carry different business meaning.\par
\featuretable{
--- & --- & \emph{Not yet captured} & \emph{Complete during review} \\
}
\subsubsection{Feature - State Tracking}\label{cap:bec-a945b4f5}
\vspace{-1em}
\begin{table}[htbp]
  \centering
  \footnotesize
  \begin{tabularx}{1\textwidth}{|l|l|X|}
    \hline
    \textbf{Metadata Name} & \textbf{Metadata Value} & \textbf{Description} \\
    \hline
    \textbf{Reference Id} & \texttt{BEC-A945B4F5} & \\
    \hline
    \textbf{Status} & Generated candidate & \\
    \hline
    \textbf{Ancestry} &  & \\
    \hline
    \textbf{Source} & Semantic-pattern candidate & \\
    \hline
    \textbf{Enabled By} &  & \\
    \hline
    \textbf{Supports} &  & \\
    \hline
    \textbf{Derived From} &  & \\
    \hline
    \textbf{Reviewed Status} & Generated candidate: keep / modify / relocate / reject; verify whether independently testable. & \\
    \hline
  \end{tabularx}
\end{table}

Intent. State Tracking is proposed as a discrete, testable enabling behaviour for Primary Care Coordination. It is included for architectural review, not as an established implementation obligation.

\noindent\textbf{Feature boundary.} This behaviour is scoped to Primary Care Coordination; an identically named Feature elsewhere may carry different business meaning.\par
\featuretable{
--- & --- & \emph{Not yet captured} & \emph{Complete during review} \\
}
\subsubsection{Feature - Exception Management}\label{cap:bec-2737b30c}
\vspace{-1em}
\begin{table}[htbp]
  \centering
  \footnotesize
  \begin{tabularx}{1\textwidth}{|l|l|X|}
    \hline
    \textbf{Metadata Name} & \textbf{Metadata Value} & \textbf{Description} \\
    \hline
    \textbf{Reference Id} & \texttt{BEC-2737B30C} & \\
    \hline
    \textbf{Status} & Generated candidate & \\
    \hline
    \textbf{Ancestry} &  & \\
    \hline
    \textbf{Source} & Semantic-pattern candidate & \\
    \hline
    \textbf{Enabled By} &  & \\
    \hline
    \textbf{Supports} &  & \\
    \hline
    \textbf{Derived From} &  & \\
    \hline
    \textbf{Reviewed Status} & Generated candidate: keep / modify / relocate / reject; verify whether independently testable. & \\
    \hline
  \end{tabularx}
\end{table}

Intent. Exception Management is proposed as a discrete, testable enabling behaviour for Primary Care Coordination. It is included for architectural review, not as an established implementation obligation.

\noindent\textbf{Feature boundary.} This behaviour is scoped to Primary Care Coordination; an identically named Feature elsewhere may carry different business meaning.\par
\featuretable{
--- & --- & \emph{Not yet captured} & \emph{Complete during review} \\
}
\subsubsection{Feature - Outcome Reporting}\label{cap:bec-f6fc078e}
\vspace{-1em}
\begin{table}[htbp]
  \centering
  \footnotesize
  \begin{tabularx}{1\textwidth}{|l|l|X|}
    \hline
    \textbf{Metadata Name} & \textbf{Metadata Value} & \textbf{Description} \\
    \hline
    \textbf{Reference Id} & \texttt{BEC-F6FC078E} & \\
    \hline
    \textbf{Status} & Generated candidate & \\
    \hline
    \textbf{Ancestry} &  & \\
    \hline
    \textbf{Source} & Semantic-pattern candidate & \\
    \hline
    \textbf{Enabled By} &  & \\
    \hline
    \textbf{Supports} &  & \\
    \hline
    \textbf{Derived From} &  & \\
    \hline
    \textbf{Reviewed Status} & Generated candidate: keep / modify / relocate / reject; verify whether independently testable. & \\
    \hline
  \end{tabularx}
\end{table}

Intent. Outcome Reporting is proposed as a discrete, testable enabling behaviour for Primary Care Coordination. It is included for architectural review, not as an established implementation obligation.

\noindent\textbf{Feature boundary.} This behaviour is scoped to Primary Care Coordination; an identically named Feature elsewhere may carry different business meaning.\par
\featuretable{
--- & --- & \emph{Not yet captured} & \emph{Complete during review} \\
}

\section{CT2 - Acute Care}\label{cap:bec-4e6db5f9}
\vspace{-0.5em}
\begin{table}[htbp]
  \centering
  \footnotesize
  \begin{tabularx}{1\textwidth}{|l|l|X|}
    \hline
    \multicolumn{3}{c}{\textbf{Capability Metadata}} \\
    \hline
    \textbf{Name} & \textbf{Value} & \textbf{Comment and/or Description} \\
    \hline
    \textbf{Reference Id} & \texttt{BEC-4E6DB5F9} & \\
    \hline
    \textbf{Status} & Candidate ancestry & \\
    \hline
    \textbf{Ancestry} &  & \\
    \hline
    \textbf{Source} &  & \\
    \hline
    \textbf{Enabled By} &  & \\
    \hline
    \textbf{Supports} &  & \\
    \hline
    \textbf{Derived From} &  & \\
    \hline
    \textbf{Reviewed Status} &  & \\
    \hline
  \end{tabularx}
\end{table}

\begin{table}[htbp]
  \centering
  \footnotesize
  \begin{tabularx}{1\textwidth}{|l|l|X|}
    \hline
    \multicolumn{3}{c}{\textbf{Contained Capabilities}} \\
    \hline
    \textbf{Capability Name} & \textbf{Capability Id} & \textbf{Comment} \\
    \hline
    Acute Care Monitoring \& Event Coordination & \texttt{cap:bec-0fac6b5a} (at \ref{cap:bec-0fac6b5a}) & \\
    \hline
  \end{tabularx}
\end{table}

Intent. Acute Care is a candidate enabling capability within Clinical Service Context Enablement. It frames the reusable system-supported behaviours needed to support this healthcare context without prescribing a product, interface, technology or deployment boundary.

Boundary. This capability supports the business activity and does not inherit the authority or accountability of the Business Capability it enables. Its child Features are review candidates until ancestry and semantics are adjudicated.

\subsection{CT3 - Acute Care Monitoring \& Event Coordination}\label{cap:bec-0fac6b5a}
\vspace{-0.5em}
\begin{table}[htbp]
  \centering
  \footnotesize
  \begin{tabularx}{1\textwidth}{|l|l|X|}
    \hline
    \multicolumn{3}{c}{\textbf{Capability Metadata}} \\
    \hline
    \textbf{Name} & \textbf{Value} & \textbf{Comment and/or Description} \\
    \hline
    \textbf{Reference Id} & \texttt{BEC-0FAC6B5A} & \\
    \hline
    \textbf{Status} & Candidate ancestry & \\
    \hline
    \textbf{Ancestry} &  & \\
    \hline
    \textbf{Source} &  & \\
    \hline
    \textbf{Enabled By} &  & \\
    \hline
    \textbf{Supports} &  & \\
    \hline
    \textbf{Derived From} &  & \\
    \hline
    \textbf{Reviewed Status} &  & \\
    \hline
  \end{tabularx}
\end{table}

\begin{table}[htbp]
  \centering
  \footnotesize
  \begin{tabularx}{1\textwidth}{|l|l|X|}
    \hline
    \multicolumn{3}{c}{\textbf{Contained Features}} \\
    \hline
    \textbf{Feature Name} & \textbf{Feature Id} & \textbf{Comment} \\
    \hline
    Acute Clinical State Event Capture & \texttt{cap:bec-5dd077b5} (at \ref{cap:bec-5dd077b5}) & \\
    \hline
    Acute Care Clinical Context Provision & \texttt{cap:bec-870d4a16} (at \ref{cap:bec-870d4a16}) & \\
    \hline
    Context Establishment & \texttt{cap:bec-3fb23504} (at \ref{cap:bec-3fb23504}) & \\
    \hline
    Eligibility Validation & \texttt{cap:bec-ca660cb0} (at \ref{cap:bec-ca660cb0}) & \\
    \hline
    Assignment \& Routing & \texttt{cap:bec-caab941d} (at \ref{cap:bec-caab941d}) & \\
    \hline
    State Tracking & \texttt{cap:bec-dfc60092} (at \ref{cap:bec-dfc60092}) & \\
    \hline
    Exception Management & \texttt{cap:bec-1298d75b} (at \ref{cap:bec-1298d75b}) & \\
    \hline
    Outcome Reporting & \texttt{cap:bec-8ed2f7e4} (at \ref{cap:bec-8ed2f7e4}) & \\
    \hline
  \end{tabularx}
\end{table}

Intent. Acute Care Monitoring \& Event Coordination is a candidate enabling capability within Acute Care. It frames the reusable system-supported behaviours needed to support this healthcare context without prescribing a product, interface, technology or deployment boundary.

Boundary. This capability supports the business activity and does not inherit the authority or accountability of the Business Capability it enables. Its child Features are review candidates until ancestry and semantics are adjudicated.

\subsubsection{Feature - Acute Clinical State Event Capture}\label{cap:bec-5dd077b5}
\vspace{-1em}
\begin{table}[htbp]
  \centering
  \footnotesize
  \begin{tabularx}{1\textwidth}{|l|l|X|}
    \hline
    \textbf{Metadata Name} & \textbf{Metadata Value} & \textbf{Description} \\
    \hline
    \textbf{Reference Id} & \texttt{BEC-5DD077B5} & \\
    \hline
    \textbf{Status} & Existing Feature; candidate CT3 placement & \\
    \hline
    \textbf{Ancestry} &  & \\
    \hline
    \textbf{Source} & FEAT-SD-03 & \\
    \hline
    \textbf{Enabled By} &  & \\
    \hline
    \textbf{Supports} &  & \\
    \hline
    \textbf{Derived From} &  & \\
    \hline
    \textbf{Reviewed Status} & Confirm existing Feature meaning and proposed CT3 ownership. & \\
    \hline
  \end{tabularx}
\end{table}

Intent. Acute Clinical State Event Capture provides a distinct enabling behaviour in the context of Acute Care Monitoring \& Event Coordination. This is an existing named Feature; its placement under the candidate CT3 remains subject to review.

\noindent\textbf{Feature boundary.} This behaviour is scoped to Acute Care Monitoring \& Event Coordination; an identically named Feature elsewhere may carry different business meaning.\par
\featuretable{
Function & \texttt{FN-7A30C97C} & Route Acute Critical Result Alert & Catalogue \\
Service & \texttt{SV-3E7806CE} & Acute Clinical Alert Dispatch & Linked to preceding Function in source \\
}
\subsubsection{Feature - Acute Care Clinical Context Provision}\label{cap:bec-870d4a16}
\vspace{-1em}
\begin{table}[htbp]
  \centering
  \footnotesize
  \begin{tabularx}{1\textwidth}{|l|l|X|}
    \hline
    \textbf{Metadata Name} & \textbf{Metadata Value} & \textbf{Description} \\
    \hline
    \textbf{Reference Id} & \texttt{BEC-870D4A16} & \\
    \hline
    \textbf{Status} & Existing Feature; candidate CT3 placement & \\
    \hline
    \textbf{Ancestry} &  & \\
    \hline
    \textbf{Source} & FEAT-SD-04 & \\
    \hline
    \textbf{Enabled By} &  & \\
    \hline
    \textbf{Supports} &  & \\
    \hline
    \textbf{Derived From} &  & \\
    \hline
    \textbf{Reviewed Status} & Confirm existing Feature meaning and proposed CT3 ownership. & \\
    \hline
  \end{tabularx}
\end{table}

Intent. Acute Care Clinical Context Provision provides a distinct enabling behaviour in the context of Acute Care Monitoring \& Event Coordination. This is an existing named Feature; its placement under the candidate CT3 remains subject to review.

\noindent\textbf{Feature boundary.} This behaviour is scoped to Acute Care Monitoring \& Event Coordination; an identically named Feature elsewhere may carry different business meaning.\par
\featuretable{
Function & \texttt{FN-1150D721} & Assemble Acute Clinical Baseline & Catalogue \\
Service & \texttt{SV-F0D85F5F} & Acute Admission Dossier Query & Linked to preceding Function in source \\
}
\subsubsection{Feature - Context Establishment}\label{cap:bec-3fb23504}
\vspace{-1em}
\begin{table}[htbp]
  \centering
  \footnotesize
  \begin{tabularx}{1\textwidth}{|l|l|X|}
    \hline
    \textbf{Metadata Name} & \textbf{Metadata Value} & \textbf{Description} \\
    \hline
    \textbf{Reference Id} & \texttt{BEC-3FB23504} & \\
    \hline
    \textbf{Status} & Generated candidate & \\
    \hline
    \textbf{Ancestry} &  & \\
    \hline
    \textbf{Source} & Semantic-pattern candidate & \\
    \hline
    \textbf{Enabled By} &  & \\
    \hline
    \textbf{Supports} &  & \\
    \hline
    \textbf{Derived From} &  & \\
    \hline
    \textbf{Reviewed Status} & Generated candidate: keep / modify / relocate / reject; verify whether independently testable. & \\
    \hline
  \end{tabularx}
\end{table}

Intent. Context Establishment is proposed as a discrete, testable enabling behaviour for Acute Care Monitoring \& Event Coordination. It is included for architectural review, not as an established implementation obligation.

\noindent\textbf{Feature boundary.} This behaviour is scoped to Acute Care Monitoring \& Event Coordination; an identically named Feature elsewhere may carry different business meaning.\par
\featuretable{
--- & --- & \emph{Not yet captured} & \emph{Complete during review} \\
}
\subsubsection{Feature - Eligibility Validation}\label{cap:bec-ca660cb0}
\vspace{-1em}
\begin{table}[htbp]
  \centering
  \footnotesize
  \begin{tabularx}{1\textwidth}{|l|l|X|}
    \hline
    \textbf{Metadata Name} & \textbf{Metadata Value} & \textbf{Description} \\
    \hline
    \textbf{Reference Id} & \texttt{BEC-CA660CB0} & \\
    \hline
    \textbf{Status} & Generated candidate & \\
    \hline
    \textbf{Ancestry} &  & \\
    \hline
    \textbf{Source} & Semantic-pattern candidate & \\
    \hline
    \textbf{Enabled By} &  & \\
    \hline
    \textbf{Supports} &  & \\
    \hline
    \textbf{Derived From} &  & \\
    \hline
    \textbf{Reviewed Status} & Generated candidate: keep / modify / relocate / reject; verify whether independently testable. & \\
    \hline
  \end{tabularx}
\end{table}

Intent. Eligibility Validation is proposed as a discrete, testable enabling behaviour for Acute Care Monitoring \& Event Coordination. It is included for architectural review, not as an established implementation obligation.

\noindent\textbf{Feature boundary.} This behaviour is scoped to Acute Care Monitoring \& Event Coordination; an identically named Feature elsewhere may carry different business meaning.\par
\featuretable{
--- & --- & \emph{Not yet captured} & \emph{Complete during review} \\
}
\subsubsection{Feature - Assignment \& Routing}\label{cap:bec-caab941d}
\vspace{-1em}
\begin{table}[htbp]
  \centering
  \footnotesize
  \begin{tabularx}{1\textwidth}{|l|l|X|}
    \hline
    \textbf{Metadata Name} & \textbf{Metadata Value} & \textbf{Description} \\
    \hline
    \textbf{Reference Id} & \texttt{BEC-CAAB941D} & \\
    \hline
    \textbf{Status} & Generated candidate & \\
    \hline
    \textbf{Ancestry} &  & \\
    \hline
    \textbf{Source} & Semantic-pattern candidate & \\
    \hline
    \textbf{Enabled By} &  & \\
    \hline
    \textbf{Supports} &  & \\
    \hline
    \textbf{Derived From} &  & \\
    \hline
    \textbf{Reviewed Status} & Generated candidate: keep / modify / relocate / reject; verify whether independently testable. & \\
    \hline
  \end{tabularx}
\end{table}

Intent. Assignment \& Routing is proposed as a discrete, testable enabling behaviour for Acute Care Monitoring \& Event Coordination. It is included for architectural review, not as an established implementation obligation.

\noindent\textbf{Feature boundary.} This behaviour is scoped to Acute Care Monitoring \& Event Coordination; an identically named Feature elsewhere may carry different business meaning.\par
\featuretable{
--- & --- & \emph{Not yet captured} & \emph{Complete during review} \\
}
\subsubsection{Feature - State Tracking}\label{cap:bec-dfc60092}
\vspace{-1em}
\begin{table}[htbp]
  \centering
  \footnotesize
  \begin{tabularx}{1\textwidth}{|l|l|X|}
    \hline
    \textbf{Metadata Name} & \textbf{Metadata Value} & \textbf{Description} \\
    \hline
    \textbf{Reference Id} & \texttt{BEC-DFC60092} & \\
    \hline
    \textbf{Status} & Generated candidate & \\
    \hline
    \textbf{Ancestry} &  & \\
    \hline
    \textbf{Source} & Semantic-pattern candidate & \\
    \hline
    \textbf{Enabled By} &  & \\
    \hline
    \textbf{Supports} &  & \\
    \hline
    \textbf{Derived From} &  & \\
    \hline
    \textbf{Reviewed Status} & Generated candidate: keep / modify / relocate / reject; verify whether independently testable. & \\
    \hline
  \end{tabularx}
\end{table}

Intent. State Tracking is proposed as a discrete, testable enabling behaviour for Acute Care Monitoring \& Event Coordination. It is included for architectural review, not as an established implementation obligation.

\noindent\textbf{Feature boundary.} This behaviour is scoped to Acute Care Monitoring \& Event Coordination; an identically named Feature elsewhere may carry different business meaning.\par
\featuretable{
--- & --- & \emph{Not yet captured} & \emph{Complete during review} \\
}
\subsubsection{Feature - Exception Management}\label{cap:bec-1298d75b}
\vspace{-1em}
\begin{table}[htbp]
  \centering
  \footnotesize
  \begin{tabularx}{1\textwidth}{|l|l|X|}
    \hline
    \textbf{Metadata Name} & \textbf{Metadata Value} & \textbf{Description} \\
    \hline
    \textbf{Reference Id} & \texttt{BEC-1298D75B} & \\
    \hline
    \textbf{Status} & Generated candidate & \\
    \hline
    \textbf{Ancestry} &  & \\
    \hline
    \textbf{Source} & Semantic-pattern candidate & \\
    \hline
    \textbf{Enabled By} &  & \\
    \hline
    \textbf{Supports} &  & \\
    \hline
    \textbf{Derived From} &  & \\
    \hline
    \textbf{Reviewed Status} & Generated candidate: keep / modify / relocate / reject; verify whether independently testable. & \\
    \hline
  \end{tabularx}
\end{table}

Intent. Exception Management is proposed as a discrete, testable enabling behaviour for Acute Care Monitoring \& Event Coordination. It is included for architectural review, not as an established implementation obligation.

\noindent\textbf{Feature boundary.} This behaviour is scoped to Acute Care Monitoring \& Event Coordination; an identically named Feature elsewhere may carry different business meaning.\par
\featuretable{
--- & --- & \emph{Not yet captured} & \emph{Complete during review} \\
}
\subsubsection{Feature - Outcome Reporting}\label{cap:bec-8ed2f7e4}
\vspace{-1em}
\begin{table}[htbp]
  \centering
  \footnotesize
  \begin{tabularx}{1\textwidth}{|l|l|X|}
    \hline
    \textbf{Metadata Name} & \textbf{Metadata Value} & \textbf{Description} \\
    \hline
    \textbf{Reference Id} & \texttt{BEC-8ED2F7E4} & \\
    \hline
    \textbf{Status} & Generated candidate & \\
    \hline
    \textbf{Ancestry} &  & \\
    \hline
    \textbf{Source} & Semantic-pattern candidate & \\
    \hline
    \textbf{Enabled By} &  & \\
    \hline
    \textbf{Supports} &  & \\
    \hline
    \textbf{Derived From} &  & \\
    \hline
    \textbf{Reviewed Status} & Generated candidate: keep / modify / relocate / reject; verify whether independently testable. & \\
    \hline
  \end{tabularx}
\end{table}

Intent. Outcome Reporting is proposed as a discrete, testable enabling behaviour for Acute Care Monitoring \& Event Coordination. It is included for architectural review, not as an established implementation obligation.

\noindent\textbf{Feature boundary.} This behaviour is scoped to Acute Care Monitoring \& Event Coordination; an identically named Feature elsewhere may carry different business meaning.\par
\featuretable{
--- & --- & \emph{Not yet captured} & \emph{Complete during review} \\
}

\section{CT2 - Emergency Care}\label{cap:bec-729dfc66}
\vspace{-0.5em}
\begin{table}[htbp]
  \centering
  \footnotesize
  \begin{tabularx}{1\textwidth}{|l|l|X|}
    \hline
    \multicolumn{3}{c}{\textbf{Capability Metadata}} \\
    \hline
    \textbf{Name} & \textbf{Value} & \textbf{Comment and/or Description} \\
    \hline
    \textbf{Reference Id} & \texttt{BEC-729DFC66} & \\
    \hline
    \textbf{Status} & Candidate ancestry & \\
    \hline
    \textbf{Ancestry} &  & \\
    \hline
    \textbf{Source} &  & \\
    \hline
    \textbf{Enabled By} &  & \\
    \hline
    \textbf{Supports} &  & \\
    \hline
    \textbf{Derived From} &  & \\
    \hline
    \textbf{Reviewed Status} &  & \\
    \hline
  \end{tabularx}
\end{table}

\begin{table}[htbp]
  \centering
  \footnotesize
  \begin{tabularx}{1\textwidth}{|l|l|X|}
    \hline
    \multicolumn{3}{c}{\textbf{Contained Capabilities}} \\
    \hline
    \textbf{Capability Name} & \textbf{Capability Id} & \textbf{Comment} \\
    \hline
    Emergency Care Progression & \texttt{cap:bec-a9ebf8fe} (at \ref{cap:bec-a9ebf8fe}) & \\
    \hline
  \end{tabularx}
\end{table}

Intent. Emergency Care is a candidate enabling capability within Clinical Service Context Enablement. It frames the reusable system-supported behaviours needed to support this healthcare context without prescribing a product, interface, technology or deployment boundary.

Boundary. This capability supports the business activity and does not inherit the authority or accountability of the Business Capability it enables. Its child Features are review candidates until ancestry and semantics are adjudicated.

\subsection{CT3 - Emergency Care Progression}\label{cap:bec-a9ebf8fe}
\vspace{-0.5em}
\begin{table}[htbp]
  \centering
  \footnotesize
  \begin{tabularx}{1\textwidth}{|l|l|X|}
    \hline
    \multicolumn{3}{c}{\textbf{Capability Metadata}} \\
    \hline
    \textbf{Name} & \textbf{Value} & \textbf{Comment and/or Description} \\
    \hline
    \textbf{Reference Id} & \texttt{BEC-A9EBF8FE} & \\
    \hline
    \textbf{Status} & Candidate ancestry & \\
    \hline
    \textbf{Ancestry} &  & \\
    \hline
    \textbf{Source} &  & \\
    \hline
    \textbf{Enabled By} &  & \\
    \hline
    \textbf{Supports} &  & \\
    \hline
    \textbf{Derived From} &  & \\
    \hline
    \textbf{Reviewed Status} &  & \\
    \hline
  \end{tabularx}
\end{table}

\begin{table}[htbp]
  \centering
  \footnotesize
  \begin{tabularx}{1\textwidth}{|l|l|X|}
    \hline
    \multicolumn{3}{c}{\textbf{Contained Features}} \\
    \hline
    \textbf{Feature Name} & \textbf{Feature Id} & \textbf{Comment} \\
    \hline
    Emergency Department Arrival \& Triage Ingestion & \texttt{cap:bec-3e1ee6b4} (at \ref{cap:bec-3e1ee6b4}) & \\
    \hline
    Emergency Care Cross-Facility Correlation & \texttt{cap:bec-78f5eadb} (at \ref{cap:bec-78f5eadb}) & \\
    \hline
    Context Establishment & \texttt{cap:bec-ed681415} (at \ref{cap:bec-ed681415}) & \\
    \hline
    State Monitoring & \texttt{cap:bec-adc2ac77} (at \ref{cap:bec-adc2ac77}) & \\
    \hline
    Exception Management & \texttt{cap:bec-0ec08e00} (at \ref{cap:bec-0ec08e00}) & \\
    \hline
    Information Retrieval & \texttt{cap:bec-1723857c} (at \ref{cap:bec-1723857c}) & \\
    \hline
  \end{tabularx}
\end{table}

Intent. Emergency Care Progression is a candidate enabling capability within Emergency Care. It frames the reusable system-supported behaviours needed to support this healthcare context without prescribing a product, interface, technology or deployment boundary.

Boundary. This capability supports the business activity and does not inherit the authority or accountability of the Business Capability it enables. Its child Features are review candidates until ancestry and semantics are adjudicated.

\subsubsection{Feature - Emergency Department Arrival \& Triage Ingestion}\label{cap:bec-3e1ee6b4}
\vspace{-1em}
\begin{table}[htbp]
  \centering
  \footnotesize
  \begin{tabularx}{1\textwidth}{|l|l|X|}
    \hline
    \textbf{Metadata Name} & \textbf{Metadata Value} & \textbf{Description} \\
    \hline
    \textbf{Reference Id} & \texttt{BEC-3E1EE6B4} & \\
    \hline
    \textbf{Status} & Existing Feature; candidate CT3 placement & \\
    \hline
    \textbf{Ancestry} &  & \\
    \hline
    \textbf{Source} & FEAT-SD-05 & \\
    \hline
    \textbf{Enabled By} &  & \\
    \hline
    \textbf{Supports} &  & \\
    \hline
    \textbf{Derived From} &  & \\
    \hline
    \textbf{Reviewed Status} & Confirm existing Feature meaning and proposed CT3 ownership. & \\
    \hline
  \end{tabularx}
\end{table}

Intent. Emergency Department Arrival \& Triage Ingestion provides a distinct enabling behaviour in the context of Emergency Care Progression. This is an existing named Feature; its placement under the candidate CT3 remains subject to review.

\noindent\textbf{Feature boundary.} This behaviour is scoped to Emergency Care Progression; an identically named Feature elsewhere may carry different business meaning.\par
\featuretable{
Function & \texttt{FN-1DFA8DF2} & Bind ED Triage Assessment & Catalogue \\
Service & \texttt{SV-791CE8EE} & ED Triage Stream Ingress & Linked to preceding Function in source \\
}
\subsubsection{Feature - Emergency Care Cross-Facility Correlation}\label{cap:bec-78f5eadb}
\vspace{-1em}
\begin{table}[htbp]
  \centering
  \footnotesize
  \begin{tabularx}{1\textwidth}{|l|l|X|}
    \hline
    \textbf{Metadata Name} & \textbf{Metadata Value} & \textbf{Description} \\
    \hline
    \textbf{Reference Id} & \texttt{BEC-78F5EADB} & \\
    \hline
    \textbf{Status} & Existing Feature; candidate CT3 placement & \\
    \hline
    \textbf{Ancestry} &  & \\
    \hline
    \textbf{Source} & FEAT-SD-06 & \\
    \hline
    \textbf{Enabled By} &  & \\
    \hline
    \textbf{Supports} &  & \\
    \hline
    \textbf{Derived From} &  & \\
    \hline
    \textbf{Reviewed Status} & Confirm existing Feature meaning and proposed CT3 ownership. & \\
    \hline
  \end{tabularx}
\end{table}

Intent. Emergency Care Cross-Facility Correlation provides a distinct enabling behaviour in the context of Emergency Care Progression. This is an existing named Feature; its placement under the candidate CT3 remains subject to review.

\noindent\textbf{Feature boundary.} This behaviour is scoped to Emergency Care Progression; an identically named Feature elsewhere may carry different business meaning.\par
\featuretable{
Function & \texttt{FN-B6F63BC8} & Provide Emergency LHR Summary & Catalogue \\
Service & \texttt{SV-713D59A2} & Emergency Fast-Path LHR Query & Linked to preceding Function in source \\
}
\subsubsection{Feature - Context Establishment}\label{cap:bec-ed681415}
\vspace{-1em}
\begin{table}[htbp]
  \centering
  \footnotesize
  \begin{tabularx}{1\textwidth}{|l|l|X|}
    \hline
    \textbf{Metadata Name} & \textbf{Metadata Value} & \textbf{Description} \\
    \hline
    \textbf{Reference Id} & \texttt{BEC-ED681415} & \\
    \hline
    \textbf{Status} & Generated candidate & \\
    \hline
    \textbf{Ancestry} &  & \\
    \hline
    \textbf{Source} & Semantic-pattern candidate & \\
    \hline
    \textbf{Enabled By} &  & \\
    \hline
    \textbf{Supports} &  & \\
    \hline
    \textbf{Derived From} &  & \\
    \hline
    \textbf{Reviewed Status} & Generated candidate: keep / modify / relocate / reject; verify whether independently testable. & \\
    \hline
  \end{tabularx}
\end{table}

Intent. Context Establishment is proposed as a discrete, testable enabling behaviour for Emergency Care Progression. It is included for architectural review, not as an established implementation obligation.

\noindent\textbf{Feature boundary.} This behaviour is scoped to Emergency Care Progression; an identically named Feature elsewhere may carry different business meaning.\par
\featuretable{
--- & --- & \emph{Not yet captured} & \emph{Complete during review} \\
}
\subsubsection{Feature - State Monitoring}\label{cap:bec-adc2ac77}
\vspace{-1em}
\begin{table}[htbp]
  \centering
  \footnotesize
  \begin{tabularx}{1\textwidth}{|l|l|X|}
    \hline
    \textbf{Metadata Name} & \textbf{Metadata Value} & \textbf{Description} \\
    \hline
    \textbf{Reference Id} & \texttt{BEC-ADC2AC77} & \\
    \hline
    \textbf{Status} & Generated candidate & \\
    \hline
    \textbf{Ancestry} &  & \\
    \hline
    \textbf{Source} & Semantic-pattern candidate & \\
    \hline
    \textbf{Enabled By} &  & \\
    \hline
    \textbf{Supports} &  & \\
    \hline
    \textbf{Derived From} &  & \\
    \hline
    \textbf{Reviewed Status} & Generated candidate: keep / modify / relocate / reject; verify whether independently testable. & \\
    \hline
  \end{tabularx}
\end{table}

Intent. State Monitoring is proposed as a discrete, testable enabling behaviour for Emergency Care Progression. It is included for architectural review, not as an established implementation obligation.

\noindent\textbf{Feature boundary.} This behaviour is scoped to Emergency Care Progression; an identically named Feature elsewhere may carry different business meaning.\par
\featuretable{
--- & --- & \emph{Not yet captured} & \emph{Complete during review} \\
}
\subsubsection{Feature - Exception Management}\label{cap:bec-0ec08e00}
\vspace{-1em}
\begin{table}[htbp]
  \centering
  \footnotesize
  \begin{tabularx}{1\textwidth}{|l|l|X|}
    \hline
    \textbf{Metadata Name} & \textbf{Metadata Value} & \textbf{Description} \\
    \hline
    \textbf{Reference Id} & \texttt{BEC-0EC08E00} & \\
    \hline
    \textbf{Status} & Generated candidate & \\
    \hline
    \textbf{Ancestry} &  & \\
    \hline
    \textbf{Source} & Semantic-pattern candidate & \\
    \hline
    \textbf{Enabled By} &  & \\
    \hline
    \textbf{Supports} &  & \\
    \hline
    \textbf{Derived From} &  & \\
    \hline
    \textbf{Reviewed Status} & Generated candidate: keep / modify / relocate / reject; verify whether independently testable. & \\
    \hline
  \end{tabularx}
\end{table}

Intent. Exception Management is proposed as a discrete, testable enabling behaviour for Emergency Care Progression. It is included for architectural review, not as an established implementation obligation.

\noindent\textbf{Feature boundary.} This behaviour is scoped to Emergency Care Progression; an identically named Feature elsewhere may carry different business meaning.\par
\featuretable{
--- & --- & \emph{Not yet captured} & \emph{Complete during review} \\
}
\subsubsection{Feature - Information Retrieval}\label{cap:bec-1723857c}
\vspace{-1em}
\begin{table}[htbp]
  \centering
  \footnotesize
  \begin{tabularx}{1\textwidth}{|l|l|X|}
    \hline
    \textbf{Metadata Name} & \textbf{Metadata Value} & \textbf{Description} \\
    \hline
    \textbf{Reference Id} & \texttt{BEC-1723857C} & \\
    \hline
    \textbf{Status} & Generated candidate & \\
    \hline
    \textbf{Ancestry} &  & \\
    \hline
    \textbf{Source} & Semantic-pattern candidate & \\
    \hline
    \textbf{Enabled By} &  & \\
    \hline
    \textbf{Supports} &  & \\
    \hline
    \textbf{Derived From} &  & \\
    \hline
    \textbf{Reviewed Status} & Generated candidate: keep / modify / relocate / reject; verify whether independently testable. & \\
    \hline
  \end{tabularx}
\end{table}

Intent. Information Retrieval is proposed as a discrete, testable enabling behaviour for Emergency Care Progression. It is included for architectural review, not as an established implementation obligation.

\noindent\textbf{Feature boundary.} This behaviour is scoped to Emergency Care Progression; an identically named Feature elsewhere may carry different business meaning.\par
\featuretable{
--- & --- & \emph{Not yet captured} & \emph{Complete during review} \\
}

\section{CT2 - Inpatient Care}\label{cap:bec-ebfe3fb1}
\vspace{-0.5em}
\begin{table}[htbp]
  \centering
  \footnotesize
  \begin{tabularx}{1\textwidth}{|l|l|X|}
    \hline
    \multicolumn{3}{c}{\textbf{Capability Metadata}} \\
    \hline
    \textbf{Name} & \textbf{Value} & \textbf{Comment and/or Description} \\
    \hline
    \textbf{Reference Id} & \texttt{BEC-EBFE3FB1} & \\
    \hline
    \textbf{Status} & Candidate ancestry & \\
    \hline
    \textbf{Ancestry} &  & \\
    \hline
    \textbf{Source} &  & \\
    \hline
    \textbf{Enabled By} &  & \\
    \hline
    \textbf{Supports} &  & \\
    \hline
    \textbf{Derived From} &  & \\
    \hline
    \textbf{Reviewed Status} &  & \\
    \hline
  \end{tabularx}
\end{table}

\begin{table}[htbp]
  \centering
  \footnotesize
  \begin{tabularx}{1\textwidth}{|l|l|X|}
    \hline
    \multicolumn{3}{c}{\textbf{Contained Capabilities}} \\
    \hline
    \textbf{Capability Name} & \textbf{Capability Id} & \textbf{Comment} \\
    \hline
    Multidisciplinary Inpatient Care Coordination & \texttt{cap:bec-17776772} (at \ref{cap:bec-17776772}) & \\
    \hline
  \end{tabularx}
\end{table}

Intent. Inpatient Care is a candidate enabling capability within Clinical Service Context Enablement. It frames the reusable system-supported behaviours needed to support this healthcare context without prescribing a product, interface, technology or deployment boundary.

Boundary. This capability supports the business activity and does not inherit the authority or accountability of the Business Capability it enables. Its child Features are review candidates until ancestry and semantics are adjudicated.

\subsection{CT3 - Multidisciplinary Inpatient Care Coordination}\label{cap:bec-17776772}
\vspace{-0.5em}
\begin{table}[htbp]
  \centering
  \footnotesize
  \begin{tabularx}{1\textwidth}{|l|l|X|}
    \hline
    \multicolumn{3}{c}{\textbf{Capability Metadata}} \\
    \hline
    \textbf{Name} & \textbf{Value} & \textbf{Comment and/or Description} \\
    \hline
    \textbf{Reference Id} & \texttt{BEC-17776772} & \\
    \hline
    \textbf{Status} & Candidate ancestry & \\
    \hline
    \textbf{Ancestry} &  & \\
    \hline
    \textbf{Source} &  & \\
    \hline
    \textbf{Enabled By} &  & \\
    \hline
    \textbf{Supports} &  & \\
    \hline
    \textbf{Derived From} &  & \\
    \hline
    \textbf{Reviewed Status} &  & \\
    \hline
  \end{tabularx}
\end{table}

\begin{table}[htbp]
  \centering
  \footnotesize
  \begin{tabularx}{1\textwidth}{|l|l|X|}
    \hline
    \multicolumn{3}{c}{\textbf{Contained Features}} \\
    \hline
    \textbf{Feature Name} & \textbf{Feature Id} & \textbf{Comment} \\
    \hline
    Inpatient Rounding Context Provision & \texttt{cap:bec-6f99412c} (at \ref{cap:bec-6f99412c}) & \\
    \hline
    Inpatient Clinical Progression Tracking & \texttt{cap:bec-93c32399} (at \ref{cap:bec-93c32399}) & \\
    \hline
    Subject Identity Correlation & \texttt{cap:bec-826fe486} (at \ref{cap:bec-826fe486}) & \\
    \hline
    Subject Relationship Management & \texttt{cap:bec-941acc0e} (at \ref{cap:bec-941acc0e}) & \\
    \hline
    Subject Context Management & \texttt{cap:bec-152bd757} (at \ref{cap:bec-152bd757}) & \\
    \hline
    Subject Information Discovery & \texttt{cap:bec-c7242ab0} (at \ref{cap:bec-c7242ab0}) & \\
    \hline
  \end{tabularx}
\end{table}

Intent. Multidisciplinary Inpatient Care Coordination is a candidate enabling capability within Inpatient Care. It frames the reusable system-supported behaviours needed to support this healthcare context without prescribing a product, interface, technology or deployment boundary.

Boundary. This capability supports the business activity and does not inherit the authority or accountability of the Business Capability it enables. Its child Features are review candidates until ancestry and semantics are adjudicated.

FHIR R5 conceptual analogues considered (not implementation mappings): Patient, Person.

\subsubsection{Feature - Inpatient Rounding Context Provision}\label{cap:bec-6f99412c}
\vspace{-1em}
\begin{table}[htbp]
  \centering
  \footnotesize
  \begin{tabularx}{1\textwidth}{|l|l|X|}
    \hline
    \textbf{Metadata Name} & \textbf{Metadata Value} & \textbf{Description} \\
    \hline
    \textbf{Reference Id} & \texttt{BEC-6F99412C} & \\
    \hline
    \textbf{Status} & Existing Feature; candidate CT3 placement & \\
    \hline
    \textbf{Ancestry} &  & \\
    \hline
    \textbf{Source} & FEAT-SD-07 & \\
    \hline
    \textbf{Enabled By} &  & \\
    \hline
    \textbf{Supports} &  & \\
    \hline
    \textbf{Derived From} &  & \\
    \hline
    \textbf{Reviewed Status} & Confirm existing Feature meaning and proposed CT3 ownership. & \\
    \hline
  \end{tabularx}
\end{table}

Intent. Inpatient Rounding Context Provision provides a distinct enabling behaviour in the context of Multidisciplinary Inpatient Care Coordination. This is an existing named Feature; its placement under the candidate CT3 remains subject to review.

\noindent\textbf{Feature boundary.} This behaviour is scoped to Multidisciplinary Inpatient Care Coordination; an identically named Feature elsewhere may carry different business meaning.\par
\featuretable{
Function & \texttt{FN-B2ADF475} & Consolidate Inpatient Clinical Round Feed & Catalogue \\
Service & \texttt{SV-CE125781} & Ward Round Clinical Briefing Query & Linked to preceding Function in source \\
}
\subsubsection{Feature - Inpatient Clinical Progression Tracking}\label{cap:bec-93c32399}
\vspace{-1em}
\begin{table}[htbp]
  \centering
  \footnotesize
  \begin{tabularx}{1\textwidth}{|l|l|X|}
    \hline
    \textbf{Metadata Name} & \textbf{Metadata Value} & \textbf{Description} \\
    \hline
    \textbf{Reference Id} & \texttt{BEC-93C32399} & \\
    \hline
    \textbf{Status} & Existing Feature; candidate CT3 placement & \\
    \hline
    \textbf{Ancestry} &  & \\
    \hline
    \textbf{Source} & FEAT-SD-08 & \\
    \hline
    \textbf{Enabled By} &  & \\
    \hline
    \textbf{Supports} &  & \\
    \hline
    \textbf{Derived From} &  & \\
    \hline
    \textbf{Reviewed Status} & Confirm existing Feature meaning and proposed CT3 ownership. & \\
    \hline
  \end{tabularx}
\end{table}

Intent. Inpatient Clinical Progression Tracking provides a distinct enabling behaviour in the context of Multidisciplinary Inpatient Care Coordination. This is an existing named Feature; its placement under the candidate CT3 remains subject to review.

\noindent\textbf{Feature boundary.} This behaviour is scoped to Multidisciplinary Inpatient Care Coordination; an identically named Feature elsewhere may carry different business meaning.\par
\featuretable{
Function & \texttt{FN-E2708F5D} & Assemble Inpatient Discharge Dossier & Catalogue \\
Service & \texttt{SV-6432AAA3} & Discharge Dossier Packaging Service & Linked to preceding Function in source \\
}
\subsubsection{Feature - Subject Identity Correlation}\label{cap:bec-826fe486}
\vspace{-1em}
\begin{table}[htbp]
  \centering
  \footnotesize
  \begin{tabularx}{1\textwidth}{|l|l|X|}
    \hline
    \textbf{Metadata Name} & \textbf{Metadata Value} & \textbf{Description} \\
    \hline
    \textbf{Reference Id} & \texttt{BEC-826FE486} & \\
    \hline
    \textbf{Status} & Generated candidate & \\
    \hline
    \textbf{Ancestry} &  & \\
    \hline
    \textbf{Source} & FHIR-informed candidate & \\
    \hline
    \textbf{Enabled By} &  & \\
    \hline
    \textbf{Supports} &  & \\
    \hline
    \textbf{Derived From} &  & \\
    \hline
    \textbf{Reviewed Status} & Generated candidate: keep / modify / relocate / reject; verify whether independently testable. & \\
    \hline
  \end{tabularx}
\end{table}

Intent. Subject Identity Correlation is proposed to support the underlying healthcare concept represented by Multidisciplinary Inpatient Care Coordination. The proposal concerns behaviour and semantic responsibility, not a FHIR field, operation, resource instance or prescribed data model.

Conceptual evidence: HL7 FHIR R5 Patient, Person.

\noindent\textbf{Feature boundary.} This behaviour is scoped to Multidisciplinary Inpatient Care Coordination; an identically named Feature elsewhere may carry different business meaning.\par
\featuretable{
--- & --- & \emph{Not yet captured} & \emph{Complete during review} \\
}
\subsubsection{Feature - Subject Relationship Management}\label{cap:bec-941acc0e}
\vspace{-1em}
\begin{table}[htbp]
  \centering
  \footnotesize
  \begin{tabularx}{1\textwidth}{|l|l|X|}
    \hline
    \textbf{Metadata Name} & \textbf{Metadata Value} & \textbf{Description} \\
    \hline
    \textbf{Reference Id} & \texttt{BEC-941ACC0E} & \\
    \hline
    \textbf{Status} & Generated candidate & \\
    \hline
    \textbf{Ancestry} &  & \\
    \hline
    \textbf{Source} & FHIR-informed candidate & \\
    \hline
    \textbf{Enabled By} &  & \\
    \hline
    \textbf{Supports} &  & \\
    \hline
    \textbf{Derived From} &  & \\
    \hline
    \textbf{Reviewed Status} & Generated candidate: keep / modify / relocate / reject; verify whether independently testable. & \\
    \hline
  \end{tabularx}
\end{table}

Intent. Subject Relationship Management is proposed to support the underlying healthcare concept represented by Multidisciplinary Inpatient Care Coordination. The proposal concerns behaviour and semantic responsibility, not a FHIR field, operation, resource instance or prescribed data model.

Conceptual evidence: HL7 FHIR R5 Patient, Person.

\noindent\textbf{Feature boundary.} This behaviour is scoped to Multidisciplinary Inpatient Care Coordination; an identically named Feature elsewhere may carry different business meaning.\par
\featuretable{
--- & --- & \emph{Not yet captured} & \emph{Complete during review} \\
}
\subsubsection{Feature - Subject Context Management}\label{cap:bec-152bd757}
\vspace{-1em}
\begin{table}[htbp]
  \centering
  \footnotesize
  \begin{tabularx}{1\textwidth}{|l|l|X|}
    \hline
    \textbf{Metadata Name} & \textbf{Metadata Value} & \textbf{Description} \\
    \hline
    \textbf{Reference Id} & \texttt{BEC-152BD757} & \\
    \hline
    \textbf{Status} & Generated candidate & \\
    \hline
    \textbf{Ancestry} &  & \\
    \hline
    \textbf{Source} & FHIR-informed candidate & \\
    \hline
    \textbf{Enabled By} &  & \\
    \hline
    \textbf{Supports} &  & \\
    \hline
    \textbf{Derived From} &  & \\
    \hline
    \textbf{Reviewed Status} & Generated candidate: keep / modify / relocate / reject; verify whether independently testable. & \\
    \hline
  \end{tabularx}
\end{table}

Intent. Subject Context Management is proposed to support the underlying healthcare concept represented by Multidisciplinary Inpatient Care Coordination. The proposal concerns behaviour and semantic responsibility, not a FHIR field, operation, resource instance or prescribed data model.

Conceptual evidence: HL7 FHIR R5 Patient, Person.

\noindent\textbf{Feature boundary.} This behaviour is scoped to Multidisciplinary Inpatient Care Coordination; an identically named Feature elsewhere may carry different business meaning.\par
\featuretable{
--- & --- & \emph{Not yet captured} & \emph{Complete during review} \\
}
\subsubsection{Feature - Subject Information Discovery}\label{cap:bec-c7242ab0}
\vspace{-1em}
\begin{table}[htbp]
  \centering
  \footnotesize
  \begin{tabularx}{1\textwidth}{|l|l|X|}
    \hline
    \textbf{Metadata Name} & \textbf{Metadata Value} & \textbf{Description} \\
    \hline
    \textbf{Reference Id} & \texttt{BEC-C7242AB0} & \\
    \hline
    \textbf{Status} & Generated candidate & \\
    \hline
    \textbf{Ancestry} &  & \\
    \hline
    \textbf{Source} & FHIR-informed candidate & \\
    \hline
    \textbf{Enabled By} &  & \\
    \hline
    \textbf{Supports} &  & \\
    \hline
    \textbf{Derived From} &  & \\
    \hline
    \textbf{Reviewed Status} & Generated candidate: keep / modify / relocate / reject; verify whether independently testable. & \\
    \hline
  \end{tabularx}
\end{table}

Intent. Subject Information Discovery is proposed to support the underlying healthcare concept represented by Multidisciplinary Inpatient Care Coordination. The proposal concerns behaviour and semantic responsibility, not a FHIR field, operation, resource instance or prescribed data model.

Conceptual evidence: HL7 FHIR R5 Patient, Person.

\noindent\textbf{Feature boundary.} This behaviour is scoped to Multidisciplinary Inpatient Care Coordination; an identically named Feature elsewhere may carry different business meaning.\par
\featuretable{
--- & --- & \emph{Not yet captured} & \emph{Complete during review} \\
}

\section{CT2 - Ambulatory Care}\label{cap:bec-69dd19fd}
\vspace{-0.5em}
\begin{table}[htbp]
  \centering
  \footnotesize
  \begin{tabularx}{1\textwidth}{|l|l|X|}
    \hline
    \multicolumn{3}{c}{\textbf{Capability Metadata}} \\
    \hline
    \textbf{Name} & \textbf{Value} & \textbf{Comment and/or Description} \\
    \hline
    \textbf{Reference Id} & \texttt{BEC-69DD19FD} & \\
    \hline
    \textbf{Status} & Candidate ancestry & \\
    \hline
    \textbf{Ancestry} &  & \\
    \hline
    \textbf{Source} &  & \\
    \hline
    \textbf{Enabled By} &  & \\
    \hline
    \textbf{Supports} &  & \\
    \hline
    \textbf{Derived From} &  & \\
    \hline
    \textbf{Reviewed Status} &  & \\
    \hline
  \end{tabularx}
\end{table}

\begin{table}[htbp]
  \centering
  \footnotesize
  \begin{tabularx}{1\textwidth}{|l|l|X|}
    \hline
    \multicolumn{3}{c}{\textbf{Contained Capabilities}} \\
    \hline
    \textbf{Capability Name} & \textbf{Capability Id} & \textbf{Comment} \\
    \hline
    \multicolumn{3}{|l|}{\emph{No contained capabilities are currently identified.}} \\
    \hline
  \end{tabularx}
\end{table}

Intent. Ambulatory Care is a candidate enabling capability within Clinical Service Context Enablement. It frames the reusable system-supported behaviours needed to support this healthcare context without prescribing a product, interface, technology or deployment boundary.

Boundary. This capability supports the business activity and does not inherit the authority or accountability of the Business Capability it enables. Its child Features are review candidates until ancestry and semantics are adjudicated.

\section{CT2 - Surgical \& Procedural Care}\label{cap:bec-6c600ac1}
\vspace{-0.5em}
\begin{table}[htbp]
  \centering
  \footnotesize
  \begin{tabularx}{1\textwidth}{|l|l|X|}
    \hline
    \multicolumn{3}{c}{\textbf{Capability Metadata}} \\
    \hline
    \textbf{Name} & \textbf{Value} & \textbf{Comment and/or Description} \\
    \hline
    \textbf{Reference Id} & \texttt{BEC-6C600AC1} & \\
    \hline
    \textbf{Status} & Candidate ancestry & \\
    \hline
    \textbf{Ancestry} &  & \\
    \hline
    \textbf{Source} &  & \\
    \hline
    \textbf{Enabled By} &  & \\
    \hline
    \textbf{Supports} &  & \\
    \hline
    \textbf{Derived From} &  & \\
    \hline
    \textbf{Reviewed Status} &  & \\
    \hline
  \end{tabularx}
\end{table}

\begin{table}[htbp]
  \centering
  \footnotesize
  \begin{tabularx}{1\textwidth}{|l|l|X|}
    \hline
    \multicolumn{3}{c}{\textbf{Contained Capabilities}} \\
    \hline
    \textbf{Capability Name} & \textbf{Capability Id} & \textbf{Comment} \\
    \hline
    \multicolumn{3}{|l|}{\emph{No contained capabilities are currently identified.}} \\
    \hline
  \end{tabularx}
\end{table}

Intent. Surgical \& Procedural Care is a candidate enabling capability within Clinical Service Context Enablement. It frames the reusable system-supported behaviours needed to support this healthcare context without prescribing a product, interface, technology or deployment boundary.

Boundary. This capability supports the business activity and does not inherit the authority or accountability of the Business Capability it enables. Its child Features are review candidates until ancestry and semantics are adjudicated.

\section{CT2 - Diagnostic Services}\label{cap:bec-27e670bb}
\vspace{-0.5em}
\begin{table}[htbp]
  \centering
  \footnotesize
  \begin{tabularx}{1\textwidth}{|l|l|X|}
    \hline
    \multicolumn{3}{c}{\textbf{Capability Metadata}} \\
    \hline
    \textbf{Name} & \textbf{Value} & \textbf{Comment and/or Description} \\
    \hline
    \textbf{Reference Id} & \texttt{BEC-27E670BB} & \\
    \hline
    \textbf{Status} & Candidate ancestry & \\
    \hline
    \textbf{Ancestry} &  & \\
    \hline
    \textbf{Source} &  & \\
    \hline
    \textbf{Enabled By} &  & \\
    \hline
    \textbf{Supports} &  & \\
    \hline
    \textbf{Derived From} &  & \\
    \hline
    \textbf{Reviewed Status} &  & \\
    \hline
  \end{tabularx}
\end{table}

\begin{table}[htbp]
  \centering
  \footnotesize
  \begin{tabularx}{1\textwidth}{|l|l|X|}
    \hline
    \multicolumn{3}{c}{\textbf{Contained Capabilities}} \\
    \hline
    \textbf{Capability Name} & \textbf{Capability Id} & \textbf{Comment} \\
    \hline
    Diagnostic Service Integration & \texttt{cap:bec-bb867c0a} (at \ref{cap:bec-bb867c0a}) & \\
    \hline
  \end{tabularx}
\end{table}

Intent. Diagnostic Services is a candidate enabling capability within Clinical Service Context Enablement. It frames the reusable system-supported behaviours needed to support this healthcare context without prescribing a product, interface, technology or deployment boundary.

Boundary. This capability supports the business activity and does not inherit the authority or accountability of the Business Capability it enables. Its child Features are review candidates until ancestry and semantics are adjudicated.

\subsection{CT3 - Diagnostic Service Integration}\label{cap:bec-bb867c0a}
\vspace{-0.5em}
\begin{table}[htbp]
  \centering
  \footnotesize
  \begin{tabularx}{1\textwidth}{|l|l|X|}
    \hline
    \multicolumn{3}{c}{\textbf{Capability Metadata}} \\
    \hline
    \textbf{Name} & \textbf{Value} & \textbf{Comment and/or Description} \\
    \hline
    \textbf{Reference Id} & \texttt{BEC-BB867C0A} & \\
    \hline
    \textbf{Status} & Candidate ancestry & \\
    \hline
    \textbf{Ancestry} &  & \\
    \hline
    \textbf{Source} &  & \\
    \hline
    \textbf{Enabled By} &  & \\
    \hline
    \textbf{Supports} &  & \\
    \hline
    \textbf{Derived From} &  & \\
    \hline
    \textbf{Reviewed Status} &  & \\
    \hline
  \end{tabularx}
\end{table}

\begin{table}[htbp]
  \centering
  \footnotesize
  \begin{tabularx}{1\textwidth}{|l|l|X|}
    \hline
    \multicolumn{3}{c}{\textbf{Contained Features}} \\
    \hline
    \textbf{Feature Name} & \textbf{Feature Id} & \textbf{Comment} \\
    \hline
    Laboratory \& Imaging Result Distribution & \texttt{cap:bec-e7a7a26f} (at \ref{cap:bec-e7a7a26f}) & \\
    \hline
    Diagnostic History Consolidation & \texttt{cap:bec-821f8a39} (at \ref{cap:bec-821f8a39}) & \\
    \hline
    Service Definition Management & \texttt{cap:bec-48d54b88} (at \ref{cap:bec-48d54b88}) & \\
    \hline
    Service Eligibility Management & \texttt{cap:bec-e8e1f920} (at \ref{cap:bec-e8e1f920}) & \\
    \hline
    Service Availability Management & \texttt{cap:bec-870caea8} (at \ref{cap:bec-870caea8}) & \\
    \hline
    Service Discovery & \texttt{cap:bec-14c91d01} (at \ref{cap:bec-14c91d01}) & \\
    \hline
    Diagnostic Request Correlation & \texttt{cap:bec-bda9d14e} (at \ref{cap:bec-bda9d14e}) & \\
    \hline
    Diagnostic Result Association & \texttt{cap:bec-a299b6d0} (at \ref{cap:bec-a299b6d0}) & \\
    \hline
    Diagnostic Status Tracking & \texttt{cap:bec-d4c02a64} (at \ref{cap:bec-d4c02a64}) & \\
    \hline
    Diagnostic Provenance Capture & \texttt{cap:bec-fcd0d5e4} (at \ref{cap:bec-fcd0d5e4}) & \\
    \hline
  \end{tabularx}
\end{table}

Intent. Diagnostic Service Integration is a candidate enabling capability within Diagnostic Services. It frames the reusable system-supported behaviours needed to support this healthcare context without prescribing a product, interface, technology or deployment boundary.

Boundary. This capability supports the business activity and does not inherit the authority or accountability of the Business Capability it enables. Its child Features are review candidates until ancestry and semantics are adjudicated.

FHIR R5 conceptual analogues considered (not implementation mappings): HealthcareService, ServiceRequest, DiagnosticReport, Observation, Specimen.

\subsubsection{Feature - Laboratory \& Imaging Result Distribution}\label{cap:bec-e7a7a26f}
\vspace{-1em}
\begin{table}[htbp]
  \centering
  \footnotesize
  \begin{tabularx}{1\textwidth}{|l|l|X|}
    \hline
    \textbf{Metadata Name} & \textbf{Metadata Value} & \textbf{Description} \\
    \hline
    \textbf{Reference Id} & \texttt{BEC-E7A7A26F} & \\
    \hline
    \textbf{Status} & Existing Feature; candidate CT3 placement & \\
    \hline
    \textbf{Ancestry} &  & \\
    \hline
    \textbf{Source} & FEAT-SD-09 & \\
    \hline
    \textbf{Enabled By} &  & \\
    \hline
    \textbf{Supports} &  & \\
    \hline
    \textbf{Derived From} &  & \\
    \hline
    \textbf{Reviewed Status} & Confirm existing Feature meaning and proposed CT3 ownership. & \\
    \hline
  \end{tabularx}
\end{table}

Intent. Laboratory \& Imaging Result Distribution provides a distinct enabling behaviour in the context of Diagnostic Service Integration. This is an existing named Feature; its placement under the candidate CT3 remains subject to review.

\noindent\textbf{Feature boundary.} This behaviour is scoped to Diagnostic Service Integration; an identically named Feature elsewhere may carry different business meaning.\par
\featuretable{
Function & \texttt{FN-F282DF83} & Receive Diagnostic Requisition Feed & Catalogue \\
Service & \texttt{SV-69A46038} & Modality Worklist Distribution & Linked to preceding Function in source \\
}
\subsubsection{Feature - Diagnostic History Consolidation}\label{cap:bec-821f8a39}
\vspace{-1em}
\begin{table}[htbp]
  \centering
  \footnotesize
  \begin{tabularx}{1\textwidth}{|l|l|X|}
    \hline
    \textbf{Metadata Name} & \textbf{Metadata Value} & \textbf{Description} \\
    \hline
    \textbf{Reference Id} & \texttt{BEC-821F8A39} & \\
    \hline
    \textbf{Status} & Existing Feature; candidate CT3 placement & \\
    \hline
    \textbf{Ancestry} &  & \\
    \hline
    \textbf{Source} & FEAT-SD-10 & \\
    \hline
    \textbf{Enabled By} &  & \\
    \hline
    \textbf{Supports} &  & \\
    \hline
    \textbf{Derived From} &  & \\
    \hline
    \textbf{Reviewed Status} & Confirm existing Feature meaning and proposed CT3 ownership. & \\
    \hline
  \end{tabularx}
\end{table}

Intent. Diagnostic History Consolidation provides a distinct enabling behaviour in the context of Diagnostic Service Integration. This is an existing named Feature; its placement under the candidate CT3 remains subject to review.

\noindent\textbf{Feature boundary.} This behaviour is scoped to Diagnostic Service Integration; an identically named Feature elsewhere may carry different business meaning.\par
\featuretable{
Function & \texttt{FN-24FA9106} & Publish Governed Diagnostic Report & Catalogue \\
Service & \texttt{SV-1816C2F6} & Diagnostic Publication Service & Linked to preceding Function in source \\
}
\subsubsection{Feature - Service Definition Management}\label{cap:bec-48d54b88}
\vspace{-1em}
\begin{table}[htbp]
  \centering
  \footnotesize
  \begin{tabularx}{1\textwidth}{|l|l|X|}
    \hline
    \textbf{Metadata Name} & \textbf{Metadata Value} & \textbf{Description} \\
    \hline
    \textbf{Reference Id} & \texttt{BEC-48D54B88} & \\
    \hline
    \textbf{Status} & Generated candidate & \\
    \hline
    \textbf{Ancestry} &  & \\
    \hline
    \textbf{Source} & FHIR-informed candidate & \\
    \hline
    \textbf{Enabled By} &  & \\
    \hline
    \textbf{Supports} &  & \\
    \hline
    \textbf{Derived From} &  & \\
    \hline
    \textbf{Reviewed Status} & Generated candidate: keep / modify / relocate / reject; verify whether independently testable. & \\
    \hline
  \end{tabularx}
\end{table}

Intent. Service Definition Management is proposed to support the underlying healthcare concept represented by Diagnostic Service Integration. The proposal concerns behaviour and semantic responsibility, not a FHIR field, operation, resource instance or prescribed data model.

Conceptual evidence: HL7 FHIR R5 HealthcareService, ServiceRequest, DiagnosticReport, Observation, Specimen.

\noindent\textbf{Feature boundary.} This behaviour is scoped to Diagnostic Service Integration; an identically named Feature elsewhere may carry different business meaning.\par
\featuretable{
--- & --- & \emph{Not yet captured} & \emph{Complete during review} \\
}
\subsubsection{Feature - Service Eligibility Management}\label{cap:bec-e8e1f920}
\vspace{-1em}
\begin{table}[htbp]
  \centering
  \footnotesize
  \begin{tabularx}{1\textwidth}{|l|l|X|}
    \hline
    \textbf{Metadata Name} & \textbf{Metadata Value} & \textbf{Description} \\
    \hline
    \textbf{Reference Id} & \texttt{BEC-E8E1F920} & \\
    \hline
    \textbf{Status} & Generated candidate & \\
    \hline
    \textbf{Ancestry} &  & \\
    \hline
    \textbf{Source} & FHIR-informed candidate & \\
    \hline
    \textbf{Enabled By} &  & \\
    \hline
    \textbf{Supports} &  & \\
    \hline
    \textbf{Derived From} &  & \\
    \hline
    \textbf{Reviewed Status} & Generated candidate: keep / modify / relocate / reject; verify whether independently testable. & \\
    \hline
  \end{tabularx}
\end{table}

Intent. Service Eligibility Management is proposed to support the underlying healthcare concept represented by Diagnostic Service Integration. The proposal concerns behaviour and semantic responsibility, not a FHIR field, operation, resource instance or prescribed data model.

Conceptual evidence: HL7 FHIR R5 HealthcareService, ServiceRequest, DiagnosticReport, Observation, Specimen.

\noindent\textbf{Feature boundary.} This behaviour is scoped to Diagnostic Service Integration; an identically named Feature elsewhere may carry different business meaning.\par
\featuretable{
--- & --- & \emph{Not yet captured} & \emph{Complete during review} \\
}
\subsubsection{Feature - Service Availability Management}\label{cap:bec-870caea8}
\vspace{-1em}
\begin{table}[htbp]
  \centering
  \footnotesize
  \begin{tabularx}{1\textwidth}{|l|l|X|}
    \hline
    \textbf{Metadata Name} & \textbf{Metadata Value} & \textbf{Description} \\
    \hline
    \textbf{Reference Id} & \texttt{BEC-870CAEA8} & \\
    \hline
    \textbf{Status} & Generated candidate & \\
    \hline
    \textbf{Ancestry} &  & \\
    \hline
    \textbf{Source} & FHIR-informed candidate & \\
    \hline
    \textbf{Enabled By} &  & \\
    \hline
    \textbf{Supports} &  & \\
    \hline
    \textbf{Derived From} &  & \\
    \hline
    \textbf{Reviewed Status} & Generated candidate: keep / modify / relocate / reject; verify whether independently testable. & \\
    \hline
  \end{tabularx}
\end{table}

Intent. Service Availability Management is proposed to support the underlying healthcare concept represented by Diagnostic Service Integration. The proposal concerns behaviour and semantic responsibility, not a FHIR field, operation, resource instance or prescribed data model.

Conceptual evidence: HL7 FHIR R5 HealthcareService, ServiceRequest, DiagnosticReport, Observation, Specimen.

\noindent\textbf{Feature boundary.} This behaviour is scoped to Diagnostic Service Integration; an identically named Feature elsewhere may carry different business meaning.\par
\featuretable{
--- & --- & \emph{Not yet captured} & \emph{Complete during review} \\
}
\subsubsection{Feature - Service Discovery}\label{cap:bec-14c91d01}
\vspace{-1em}
\begin{table}[htbp]
  \centering
  \footnotesize
  \begin{tabularx}{1\textwidth}{|l|l|X|}
    \hline
    \textbf{Metadata Name} & \textbf{Metadata Value} & \textbf{Description} \\
    \hline
    \textbf{Reference Id} & \texttt{BEC-14C91D01} & \\
    \hline
    \textbf{Status} & Generated candidate & \\
    \hline
    \textbf{Ancestry} &  & \\
    \hline
    \textbf{Source} & FHIR-informed candidate & \\
    \hline
    \textbf{Enabled By} &  & \\
    \hline
    \textbf{Supports} &  & \\
    \hline
    \textbf{Derived From} &  & \\
    \hline
    \textbf{Reviewed Status} & Generated candidate: keep / modify / relocate / reject; verify whether independently testable. & \\
    \hline
  \end{tabularx}
\end{table}

Intent. Service Discovery is proposed to support the underlying healthcare concept represented by Diagnostic Service Integration. The proposal concerns behaviour and semantic responsibility, not a FHIR field, operation, resource instance or prescribed data model.

Conceptual evidence: HL7 FHIR R5 HealthcareService, ServiceRequest, DiagnosticReport, Observation, Specimen.

\noindent\textbf{Feature boundary.} This behaviour is scoped to Diagnostic Service Integration; an identically named Feature elsewhere may carry different business meaning.\par
\featuretable{
--- & --- & \emph{Not yet captured} & \emph{Complete during review} \\
}
\subsubsection{Feature - Diagnostic Request Correlation}\label{cap:bec-bda9d14e}
\vspace{-1em}
\begin{table}[htbp]
  \centering
  \footnotesize
  \begin{tabularx}{1\textwidth}{|l|l|X|}
    \hline
    \textbf{Metadata Name} & \textbf{Metadata Value} & \textbf{Description} \\
    \hline
    \textbf{Reference Id} & \texttt{BEC-BDA9D14E} & \\
    \hline
    \textbf{Status} & Generated candidate & \\
    \hline
    \textbf{Ancestry} &  & \\
    \hline
    \textbf{Source} & FHIR-informed candidate & \\
    \hline
    \textbf{Enabled By} &  & \\
    \hline
    \textbf{Supports} &  & \\
    \hline
    \textbf{Derived From} &  & \\
    \hline
    \textbf{Reviewed Status} & Generated candidate: keep / modify / relocate / reject; verify whether independently testable. & \\
    \hline
  \end{tabularx}
\end{table}

Intent. Diagnostic Request Correlation is proposed to support the underlying healthcare concept represented by Diagnostic Service Integration. The proposal concerns behaviour and semantic responsibility, not a FHIR field, operation, resource instance or prescribed data model.

Conceptual evidence: HL7 FHIR R5 HealthcareService, ServiceRequest, DiagnosticReport, Observation, Specimen.

\noindent\textbf{Feature boundary.} This behaviour is scoped to Diagnostic Service Integration; an identically named Feature elsewhere may carry different business meaning.\par
\featuretable{
--- & --- & \emph{Not yet captured} & \emph{Complete during review} \\
}
\subsubsection{Feature - Diagnostic Result Association}\label{cap:bec-a299b6d0}
\vspace{-1em}
\begin{table}[htbp]
  \centering
  \footnotesize
  \begin{tabularx}{1\textwidth}{|l|l|X|}
    \hline
    \textbf{Metadata Name} & \textbf{Metadata Value} & \textbf{Description} \\
    \hline
    \textbf{Reference Id} & \texttt{BEC-A299B6D0} & \\
    \hline
    \textbf{Status} & Generated candidate & \\
    \hline
    \textbf{Ancestry} &  & \\
    \hline
    \textbf{Source} & FHIR-informed candidate & \\
    \hline
    \textbf{Enabled By} &  & \\
    \hline
    \textbf{Supports} &  & \\
    \hline
    \textbf{Derived From} &  & \\
    \hline
    \textbf{Reviewed Status} & Generated candidate: keep / modify / relocate / reject; verify whether independently testable. & \\
    \hline
  \end{tabularx}
\end{table}

Intent. Diagnostic Result Association is proposed to support the underlying healthcare concept represented by Diagnostic Service Integration. The proposal concerns behaviour and semantic responsibility, not a FHIR field, operation, resource instance or prescribed data model.

Conceptual evidence: HL7 FHIR R5 HealthcareService, ServiceRequest, DiagnosticReport, Observation, Specimen.

\noindent\textbf{Feature boundary.} This behaviour is scoped to Diagnostic Service Integration; an identically named Feature elsewhere may carry different business meaning.\par
\featuretable{
--- & --- & \emph{Not yet captured} & \emph{Complete during review} \\
}
\subsubsection{Feature - Diagnostic Status Tracking}\label{cap:bec-d4c02a64}
\vspace{-1em}
\begin{table}[htbp]
  \centering
  \footnotesize
  \begin{tabularx}{1\textwidth}{|l|l|X|}
    \hline
    \textbf{Metadata Name} & \textbf{Metadata Value} & \textbf{Description} \\
    \hline
    \textbf{Reference Id} & \texttt{BEC-D4C02A64} & \\
    \hline
    \textbf{Status} & Generated candidate & \\
    \hline
    \textbf{Ancestry} &  & \\
    \hline
    \textbf{Source} & FHIR-informed candidate & \\
    \hline
    \textbf{Enabled By} &  & \\
    \hline
    \textbf{Supports} &  & \\
    \hline
    \textbf{Derived From} &  & \\
    \hline
    \textbf{Reviewed Status} & Generated candidate: keep / modify / relocate / reject; verify whether independently testable. & \\
    \hline
  \end{tabularx}
\end{table}

Intent. Diagnostic Status Tracking is proposed to support the underlying healthcare concept represented by Diagnostic Service Integration. The proposal concerns behaviour and semantic responsibility, not a FHIR field, operation, resource instance or prescribed data model.

Conceptual evidence: HL7 FHIR R5 HealthcareService, ServiceRequest, DiagnosticReport, Observation, Specimen.

\noindent\textbf{Feature boundary.} This behaviour is scoped to Diagnostic Service Integration; an identically named Feature elsewhere may carry different business meaning.\par
\featuretable{
--- & --- & \emph{Not yet captured} & \emph{Complete during review} \\
}
\subsubsection{Feature - Diagnostic Provenance Capture}\label{cap:bec-fcd0d5e4}
\vspace{-1em}
\begin{table}[htbp]
  \centering
  \footnotesize
  \begin{tabularx}{1\textwidth}{|l|l|X|}
    \hline
    \textbf{Metadata Name} & \textbf{Metadata Value} & \textbf{Description} \\
    \hline
    \textbf{Reference Id} & \texttt{BEC-FCD0D5E4} & \\
    \hline
    \textbf{Status} & Generated candidate & \\
    \hline
    \textbf{Ancestry} &  & \\
    \hline
    \textbf{Source} & FHIR-informed candidate & \\
    \hline
    \textbf{Enabled By} &  & \\
    \hline
    \textbf{Supports} &  & \\
    \hline
    \textbf{Derived From} &  & \\
    \hline
    \textbf{Reviewed Status} & Generated candidate: keep / modify / relocate / reject; verify whether independently testable. & \\
    \hline
  \end{tabularx}
\end{table}

Intent. Diagnostic Provenance Capture is proposed to support the underlying healthcare concept represented by Diagnostic Service Integration. The proposal concerns behaviour and semantic responsibility, not a FHIR field, operation, resource instance or prescribed data model.

Conceptual evidence: HL7 FHIR R5 HealthcareService, ServiceRequest, DiagnosticReport, Observation, Specimen.

\noindent\textbf{Feature boundary.} This behaviour is scoped to Diagnostic Service Integration; an identically named Feature elsewhere may carry different business meaning.\par
\featuretable{
--- & --- & \emph{Not yet captured} & \emph{Complete during review} \\
}

\section{CT2 - Therapeutic Services}\label{cap:bec-bfdb47ca}
\vspace{-0.5em}
\begin{table}[htbp]
  \centering
  \footnotesize
  \begin{tabularx}{1\textwidth}{|l|l|X|}
    \hline
    \multicolumn{3}{c}{\textbf{Capability Metadata}} \\
    \hline
    \textbf{Name} & \textbf{Value} & \textbf{Comment and/or Description} \\
    \hline
    \textbf{Reference Id} & \texttt{BEC-BFDB47CA} & \\
    \hline
    \textbf{Status} & Candidate ancestry & \\
    \hline
    \textbf{Ancestry} &  & \\
    \hline
    \textbf{Source} &  & \\
    \hline
    \textbf{Enabled By} &  & \\
    \hline
    \textbf{Supports} &  & \\
    \hline
    \textbf{Derived From} &  & \\
    \hline
    \textbf{Reviewed Status} &  & \\
    \hline
  \end{tabularx}
\end{table}

\begin{table}[htbp]
  \centering
  \footnotesize
  \begin{tabularx}{1\textwidth}{|l|l|X|}
    \hline
    \multicolumn{3}{c}{\textbf{Contained Capabilities}} \\
    \hline
    \textbf{Capability Name} & \textbf{Capability Id} & \textbf{Comment} \\
    \hline
    \multicolumn{3}{|l|}{\emph{No contained capabilities are currently identified.}} \\
    \hline
  \end{tabularx}
\end{table}

Intent. Therapeutic Services is a candidate enabling capability within Clinical Service Context Enablement. It frames the reusable system-supported behaviours needed to support this healthcare context without prescribing a product, interface, technology or deployment boundary.

Boundary. This capability supports the business activity and does not inherit the authority or accountability of the Business Capability it enables. Its child Features are review candidates until ancestry and semantics are adjudicated.

\section{CT2 - Medication Therapy}\label{cap:bec-b48bc336}
\vspace{-0.5em}
\begin{table}[htbp]
  \centering
  \footnotesize
  \begin{tabularx}{1\textwidth}{|l|l|X|}
    \hline
    \multicolumn{3}{c}{\textbf{Capability Metadata}} \\
    \hline
    \textbf{Name} & \textbf{Value} & \textbf{Comment and/or Description} \\
    \hline
    \textbf{Reference Id} & \texttt{BEC-B48BC336} & \\
    \hline
    \textbf{Status} & Candidate ancestry & \\
    \hline
    \textbf{Ancestry} &  & \\
    \hline
    \textbf{Source} &  & \\
    \hline
    \textbf{Enabled By} &  & \\
    \hline
    \textbf{Supports} &  & \\
    \hline
    \textbf{Derived From} &  & \\
    \hline
    \textbf{Reviewed Status} &  & \\
    \hline
  \end{tabularx}
\end{table}

\begin{table}[htbp]
  \centering
  \footnotesize
  \begin{tabularx}{1\textwidth}{|l|l|X|}
    \hline
    \multicolumn{3}{c}{\textbf{Contained Capabilities}} \\
    \hline
    \textbf{Capability Name} & \textbf{Capability Id} & \textbf{Comment} \\
    \hline
    Medication Therapy Provision and Monitoring & \texttt{cap:bec-0789035c} (at \ref{cap:bec-0789035c}) & \\
    \hline
  \end{tabularx}
\end{table}

Intent. Medication Therapy is a candidate enabling capability within Clinical Service Context Enablement. It frames the reusable system-supported behaviours needed to support this healthcare context without prescribing a product, interface, technology or deployment boundary.

Boundary. This capability supports the business activity and does not inherit the authority or accountability of the Business Capability it enables. Its child Features are review candidates until ancestry and semantics are adjudicated.

\subsection{CT3 - Medication Therapy Provision and Monitoring}\label{cap:bec-0789035c}
\vspace{-0.5em}
\begin{table}[htbp]
  \centering
  \footnotesize
  \begin{tabularx}{1\textwidth}{|l|l|X|}
    \hline
    \multicolumn{3}{c}{\textbf{Capability Metadata}} \\
    \hline
    \textbf{Name} & \textbf{Value} & \textbf{Comment and/or Description} \\
    \hline
    \textbf{Reference Id} & \texttt{BEC-0789035C} & \\
    \hline
    \textbf{Status} & Candidate ancestry & \\
    \hline
    \textbf{Ancestry} &  & \\
    \hline
    \textbf{Source} &  & \\
    \hline
    \textbf{Enabled By} &  & \\
    \hline
    \textbf{Supports} &  & \\
    \hline
    \textbf{Derived From} &  & \\
    \hline
    \textbf{Reviewed Status} &  & \\
    \hline
  \end{tabularx}
\end{table}

\begin{table}[htbp]
  \centering
  \footnotesize
  \begin{tabularx}{1\textwidth}{|l|l|X|}
    \hline
    \multicolumn{3}{c}{\textbf{Contained Features}} \\
    \hline
    \textbf{Feature Name} & \textbf{Feature Id} & \textbf{Comment} \\
    \hline
    Medication History Reconciliation Ingestion & \texttt{cap:bec-d1549684} (at \ref{cap:bec-d1549684}) & \\
    \hline
    Medication Administration Tracking & \texttt{cap:bec-95e0f2e9} (at \ref{cap:bec-95e0f2e9}) & \\
    \hline
    Medication Identity Resolution & \texttt{cap:bec-910194d7} (at \ref{cap:bec-910194d7}) & \\
    \hline
    Medication Safety Context Validation & \texttt{cap:bec-324d8bad} (at \ref{cap:bec-324d8bad}) & \\
    \hline
    Medication Activity Tracking & \texttt{cap:bec-0f34b38a} (at \ref{cap:bec-0f34b38a}) & \\
    \hline
    Medication Information Correlation & \texttt{cap:bec-b3d9c0f4} (at \ref{cap:bec-b3d9c0f4}) & \\
    \hline
  \end{tabularx}
\end{table}

Intent. Medication Therapy Provision and Monitoring is a candidate enabling capability within Medication Therapy. It frames the reusable system-supported behaviours needed to support this healthcare context without prescribing a product, interface, technology or deployment boundary.

Boundary. This capability supports the business activity and does not inherit the authority or accountability of the Business Capability it enables. Its child Features are review candidates until ancestry and semantics are adjudicated.

FHIR R5 conceptual analogues considered (not implementation mappings): Medication, MedicationRequest, MedicationDispense, MedicationAdministration.

\subsubsection{Feature - Medication History Reconciliation Ingestion}\label{cap:bec-d1549684}
\vspace{-1em}
\begin{table}[htbp]
  \centering
  \footnotesize
  \begin{tabularx}{1\textwidth}{|l|l|X|}
    \hline
    \textbf{Metadata Name} & \textbf{Metadata Value} & \textbf{Description} \\
    \hline
    \textbf{Reference Id} & \texttt{BEC-D1549684} & \\
    \hline
    \textbf{Status} & Existing Feature; candidate CT3 placement & \\
    \hline
    \textbf{Ancestry} &  & \\
    \hline
    \textbf{Source} & FEAT-SD-11 & \\
    \hline
    \textbf{Enabled By} &  & \\
    \hline
    \textbf{Supports} &  & \\
    \hline
    \textbf{Derived From} &  & \\
    \hline
    \textbf{Reviewed Status} & Confirm existing Feature meaning and proposed CT3 ownership. & \\
    \hline
  \end{tabularx}
\end{table}

Intent. Medication History Reconciliation Ingestion provides a distinct enabling behaviour in the context of Medication Therapy Provision and Monitoring. This is an existing named Feature; its placement under the candidate CT3 remains subject to review.

\noindent\textbf{Feature boundary.} This behaviour is scoped to Medication Therapy Provision and Monitoring; an identically named Feature elsewhere may carry different business meaning.\par
\featuretable{
Function & \texttt{FN-8E849508} & Project Longitudinal Medication History & Catalogue \\
Service & \texttt{SV-13563BEE} & Longitudinal Medication History Query & Linked to preceding Function in source \\
}
\subsubsection{Feature - Medication Administration Tracking}\label{cap:bec-95e0f2e9}
\vspace{-1em}
\begin{table}[htbp]
  \centering
  \footnotesize
  \begin{tabularx}{1\textwidth}{|l|l|X|}
    \hline
    \textbf{Metadata Name} & \textbf{Metadata Value} & \textbf{Description} \\
    \hline
    \textbf{Reference Id} & \texttt{BEC-95E0F2E9} & \\
    \hline
    \textbf{Status} & Existing Feature; candidate CT3 placement & \\
    \hline
    \textbf{Ancestry} &  & \\
    \hline
    \textbf{Source} & FEAT-SD-12 & \\
    \hline
    \textbf{Enabled By} &  & \\
    \hline
    \textbf{Supports} &  & \\
    \hline
    \textbf{Derived From} &  & \\
    \hline
    \textbf{Reviewed Status} & Confirm existing Feature meaning and proposed CT3 ownership. & \\
    \hline
  \end{tabularx}
\end{table}

Intent. Medication Administration Tracking provides a distinct enabling behaviour in the context of Medication Therapy Provision and Monitoring. This is an existing named Feature; its placement under the candidate CT3 remains subject to review.

\noindent\textbf{Feature boundary.} This behaviour is scoped to Medication Therapy Provision and Monitoring; an identically named Feature elsewhere may carry different business meaning.\par
\featuretable{
Function & \texttt{FN-10F116AB} & Correlate Adverse Reaction \& Allergy Profile & Catalogue \\
Service & \texttt{SV-95D71E1E} & Allergy \& Adverse Reaction Query & Linked to preceding Function in source \\
}
\subsubsection{Feature - Medication Identity Resolution}\label{cap:bec-910194d7}
\vspace{-1em}
\begin{table}[htbp]
  \centering
  \footnotesize
  \begin{tabularx}{1\textwidth}{|l|l|X|}
    \hline
    \textbf{Metadata Name} & \textbf{Metadata Value} & \textbf{Description} \\
    \hline
    \textbf{Reference Id} & \texttt{BEC-910194D7} & \\
    \hline
    \textbf{Status} & Generated candidate & \\
    \hline
    \textbf{Ancestry} &  & \\
    \hline
    \textbf{Source} & FHIR-informed candidate & \\
    \hline
    \textbf{Enabled By} &  & \\
    \hline
    \textbf{Supports} &  & \\
    \hline
    \textbf{Derived From} &  & \\
    \hline
    \textbf{Reviewed Status} & Generated candidate: keep / modify / relocate / reject; verify whether independently testable. & \\
    \hline
  \end{tabularx}
\end{table}

Intent. Medication Identity Resolution is proposed to support the underlying healthcare concept represented by Medication Therapy Provision and Monitoring. The proposal concerns behaviour and semantic responsibility, not a FHIR field, operation, resource instance or prescribed data model.

Conceptual evidence: HL7 FHIR R5 Medication, MedicationRequest, MedicationDispense, MedicationAdministration.

\noindent\textbf{Feature boundary.} This behaviour is scoped to Medication Therapy Provision and Monitoring; an identically named Feature elsewhere may carry different business meaning.\par
\featuretable{
--- & --- & \emph{Not yet captured} & \emph{Complete during review} \\
}
\subsubsection{Feature - Medication Safety Context Validation}\label{cap:bec-324d8bad}
\vspace{-1em}
\begin{table}[htbp]
  \centering
  \footnotesize
  \begin{tabularx}{1\textwidth}{|l|l|X|}
    \hline
    \textbf{Metadata Name} & \textbf{Metadata Value} & \textbf{Description} \\
    \hline
    \textbf{Reference Id} & \texttt{BEC-324D8BAD} & \\
    \hline
    \textbf{Status} & Generated candidate & \\
    \hline
    \textbf{Ancestry} &  & \\
    \hline
    \textbf{Source} & FHIR-informed candidate & \\
    \hline
    \textbf{Enabled By} &  & \\
    \hline
    \textbf{Supports} &  & \\
    \hline
    \textbf{Derived From} &  & \\
    \hline
    \textbf{Reviewed Status} & Generated candidate: keep / modify / relocate / reject; verify whether independently testable. & \\
    \hline
  \end{tabularx}
\end{table}

Intent. Medication Safety Context Validation is proposed to support the underlying healthcare concept represented by Medication Therapy Provision and Monitoring. The proposal concerns behaviour and semantic responsibility, not a FHIR field, operation, resource instance or prescribed data model.

Conceptual evidence: HL7 FHIR R5 Medication, MedicationRequest, MedicationDispense, MedicationAdministration.

\noindent\textbf{Feature boundary.} This behaviour is scoped to Medication Therapy Provision and Monitoring; an identically named Feature elsewhere may carry different business meaning.\par
\featuretable{
--- & --- & \emph{Not yet captured} & \emph{Complete during review} \\
}
\subsubsection{Feature - Medication Activity Tracking}\label{cap:bec-0f34b38a}
\vspace{-1em}
\begin{table}[htbp]
  \centering
  \footnotesize
  \begin{tabularx}{1\textwidth}{|l|l|X|}
    \hline
    \textbf{Metadata Name} & \textbf{Metadata Value} & \textbf{Description} \\
    \hline
    \textbf{Reference Id} & \texttt{BEC-0F34B38A} & \\
    \hline
    \textbf{Status} & Generated candidate & \\
    \hline
    \textbf{Ancestry} &  & \\
    \hline
    \textbf{Source} & FHIR-informed candidate & \\
    \hline
    \textbf{Enabled By} &  & \\
    \hline
    \textbf{Supports} &  & \\
    \hline
    \textbf{Derived From} &  & \\
    \hline
    \textbf{Reviewed Status} & Generated candidate: keep / modify / relocate / reject; verify whether independently testable. & \\
    \hline
  \end{tabularx}
\end{table}

Intent. Medication Activity Tracking is proposed to support the underlying healthcare concept represented by Medication Therapy Provision and Monitoring. The proposal concerns behaviour and semantic responsibility, not a FHIR field, operation, resource instance or prescribed data model.

Conceptual evidence: HL7 FHIR R5 Medication, MedicationRequest, MedicationDispense, MedicationAdministration.

\noindent\textbf{Feature boundary.} This behaviour is scoped to Medication Therapy Provision and Monitoring; an identically named Feature elsewhere may carry different business meaning.\par
\featuretable{
--- & --- & \emph{Not yet captured} & \emph{Complete during review} \\
}
\subsubsection{Feature - Medication Information Correlation}\label{cap:bec-b3d9c0f4}
\vspace{-1em}
\begin{table}[htbp]
  \centering
  \footnotesize
  \begin{tabularx}{1\textwidth}{|l|l|X|}
    \hline
    \textbf{Metadata Name} & \textbf{Metadata Value} & \textbf{Description} \\
    \hline
    \textbf{Reference Id} & \texttt{BEC-B3D9C0F4} & \\
    \hline
    \textbf{Status} & Generated candidate & \\
    \hline
    \textbf{Ancestry} &  & \\
    \hline
    \textbf{Source} & FHIR-informed candidate & \\
    \hline
    \textbf{Enabled By} &  & \\
    \hline
    \textbf{Supports} &  & \\
    \hline
    \textbf{Derived From} &  & \\
    \hline
    \textbf{Reviewed Status} & Generated candidate: keep / modify / relocate / reject; verify whether independently testable. & \\
    \hline
  \end{tabularx}
\end{table}

Intent. Medication Information Correlation is proposed to support the underlying healthcare concept represented by Medication Therapy Provision and Monitoring. The proposal concerns behaviour and semantic responsibility, not a FHIR field, operation, resource instance or prescribed data model.

Conceptual evidence: HL7 FHIR R5 Medication, MedicationRequest, MedicationDispense, MedicationAdministration.

\noindent\textbf{Feature boundary.} This behaviour is scoped to Medication Therapy Provision and Monitoring; an identically named Feature elsewhere may carry different business meaning.\par
\featuretable{
--- & --- & \emph{Not yet captured} & \emph{Complete during review} \\
}

\section{CT2 - Allied Health}\label{cap:bec-13b74601}
\vspace{-0.5em}
\begin{table}[htbp]
  \centering
  \footnotesize
  \begin{tabularx}{1\textwidth}{|l|l|X|}
    \hline
    \multicolumn{3}{c}{\textbf{Capability Metadata}} \\
    \hline
    \textbf{Name} & \textbf{Value} & \textbf{Comment and/or Description} \\
    \hline
    \textbf{Reference Id} & \texttt{BEC-13B74601} & \\
    \hline
    \textbf{Status} & Candidate ancestry & \\
    \hline
    \textbf{Ancestry} &  & \\
    \hline
    \textbf{Source} &  & \\
    \hline
    \textbf{Enabled By} &  & \\
    \hline
    \textbf{Supports} &  & \\
    \hline
    \textbf{Derived From} &  & \\
    \hline
    \textbf{Reviewed Status} &  & \\
    \hline
  \end{tabularx}
\end{table}

\begin{table}[htbp]
  \centering
  \footnotesize
  \begin{tabularx}{1\textwidth}{|l|l|X|}
    \hline
    \multicolumn{3}{c}{\textbf{Contained Capabilities}} \\
    \hline
    \textbf{Capability Name} & \textbf{Capability Id} & \textbf{Comment} \\
    \hline
    \multicolumn{3}{|l|}{\emph{No contained capabilities are currently identified.}} \\
    \hline
  \end{tabularx}
\end{table}

Intent. Allied Health is a candidate enabling capability within Clinical Service Context Enablement. It frames the reusable system-supported behaviours needed to support this healthcare context without prescribing a product, interface, technology or deployment boundary.

Boundary. This capability supports the business activity and does not inherit the authority or accountability of the Business Capability it enables. Its child Features are review candidates until ancestry and semantics are adjudicated.

\section{CT2 - Preventive Care}\label{cap:bec-8c53e055}
\vspace{-0.5em}
\begin{table}[htbp]
  \centering
  \footnotesize
  \begin{tabularx}{1\textwidth}{|l|l|X|}
    \hline
    \multicolumn{3}{c}{\textbf{Capability Metadata}} \\
    \hline
    \textbf{Name} & \textbf{Value} & \textbf{Comment and/or Description} \\
    \hline
    \textbf{Reference Id} & \texttt{BEC-8C53E055} & \\
    \hline
    \textbf{Status} & Candidate ancestry & \\
    \hline
    \textbf{Ancestry} &  & \\
    \hline
    \textbf{Source} &  & \\
    \hline
    \textbf{Enabled By} &  & \\
    \hline
    \textbf{Supports} &  & \\
    \hline
    \textbf{Derived From} &  & \\
    \hline
    \textbf{Reviewed Status} &  & \\
    \hline
  \end{tabularx}
\end{table}

\begin{table}[htbp]
  \centering
  \footnotesize
  \begin{tabularx}{1\textwidth}{|l|l|X|}
    \hline
    \multicolumn{3}{c}{\textbf{Contained Capabilities}} \\
    \hline
    \textbf{Capability Name} & \textbf{Capability Id} & \textbf{Comment} \\
    \hline
    Screening and Immunisation Enablement & \texttt{cap:bec-6ecb9608} (at \ref{cap:bec-6ecb9608}) & \\
    \hline
  \end{tabularx}
\end{table}

Intent. Preventive Care is a candidate enabling capability within Clinical Service Context Enablement. It frames the reusable system-supported behaviours needed to support this healthcare context without prescribing a product, interface, technology or deployment boundary.

Boundary. This capability supports the business activity and does not inherit the authority or accountability of the Business Capability it enables. Its child Features are review candidates until ancestry and semantics are adjudicated.

\subsection{CT3 - Screening and Immunisation Enablement}\label{cap:bec-6ecb9608}
\vspace{-0.5em}
\begin{table}[htbp]
  \centering
  \footnotesize
  \begin{tabularx}{1\textwidth}{|l|l|X|}
    \hline
    \multicolumn{3}{c}{\textbf{Capability Metadata}} \\
    \hline
    \textbf{Name} & \textbf{Value} & \textbf{Comment and/or Description} \\
    \hline
    \textbf{Reference Id} & \texttt{BEC-6ECB9608} & \\
    \hline
    \textbf{Status} & Candidate ancestry & \\
    \hline
    \textbf{Ancestry} &  & \\
    \hline
    \textbf{Source} &  & \\
    \hline
    \textbf{Enabled By} &  & \\
    \hline
    \textbf{Supports} &  & \\
    \hline
    \textbf{Derived From} &  & \\
    \hline
    \textbf{Reviewed Status} &  & \\
    \hline
  \end{tabularx}
\end{table}

\begin{table}[htbp]
  \centering
  \footnotesize
  \begin{tabularx}{1\textwidth}{|l|l|X|}
    \hline
    \multicolumn{3}{c}{\textbf{Contained Features}} \\
    \hline
    \textbf{Feature Name} & \textbf{Feature Id} & \textbf{Comment} \\
    \hline
    Immunisation Event Submission & \texttt{cap:bec-21483404} (at \ref{cap:bec-21483404}) & \\
    \hline
    Screening Recall Notification Distribution & \texttt{cap:bec-b609c8f5} (at \ref{cap:bec-b609c8f5}) & \\
    \hline
    Context Establishment & \texttt{cap:bec-261578c2} (at \ref{cap:bec-261578c2}) & \\
    \hline
    State Monitoring & \texttt{cap:bec-c577b7c8} (at \ref{cap:bec-c577b7c8}) & \\
    \hline
    Exception Management & \texttt{cap:bec-f9613d20} (at \ref{cap:bec-f9613d20}) & \\
    \hline
    Information Retrieval & \texttt{cap:bec-336aaf9d} (at \ref{cap:bec-336aaf9d}) & \\
    \hline
  \end{tabularx}
\end{table}

Intent. Screening and Immunisation Enablement is a candidate enabling capability within Preventive Care. It frames the reusable system-supported behaviours needed to support this healthcare context without prescribing a product, interface, technology or deployment boundary.

Boundary. This capability supports the business activity and does not inherit the authority or accountability of the Business Capability it enables. Its child Features are review candidates until ancestry and semantics are adjudicated.

FHIR R5 conceptual analogues considered (not implementation mappings): Immunization, ImmunizationRecommendation.

\subsubsection{Feature - Immunisation Event Submission}\label{cap:bec-21483404}
\vspace{-1em}
\begin{table}[htbp]
  \centering
  \footnotesize
  \begin{tabularx}{1\textwidth}{|l|l|X|}
    \hline
    \textbf{Metadata Name} & \textbf{Metadata Value} & \textbf{Description} \\
    \hline
    \textbf{Reference Id} & \texttt{BEC-21483404} & \\
    \hline
    \textbf{Status} & Existing Feature; candidate CT3 placement & \\
    \hline
    \textbf{Ancestry} &  & \\
    \hline
    \textbf{Source} & FEAT-SD-13 & \\
    \hline
    \textbf{Enabled By} &  & \\
    \hline
    \textbf{Supports} &  & \\
    \hline
    \textbf{Derived From} &  & \\
    \hline
    \textbf{Reviewed Status} & Confirm existing Feature meaning and proposed CT3 ownership. & \\
    \hline
  \end{tabularx}
\end{table}

Intent. Immunisation Event Submission provides a distinct enabling behaviour in the context of Screening and Immunisation Enablement. This is an existing named Feature; its placement under the candidate CT3 remains subject to review.

\noindent\textbf{Feature boundary.} This behaviour is scoped to Screening and Immunisation Enablement; an identically named Feature elsewhere may carry different business meaning.\par
\featuretable{
Function & \texttt{FN-EB1FF8ED} & Syndicate Immunisation Record & Catalogue \\
Service & \texttt{SV-61714186} & Immunisation History Query & Linked to preceding Function in source \\
}
\subsubsection{Feature - Screening Recall Notification Distribution}\label{cap:bec-b609c8f5}
\vspace{-1em}
\begin{table}[htbp]
  \centering
  \footnotesize
  \begin{tabularx}{1\textwidth}{|l|l|X|}
    \hline
    \textbf{Metadata Name} & \textbf{Metadata Value} & \textbf{Description} \\
    \hline
    \textbf{Reference Id} & \texttt{BEC-B609C8F5} & \\
    \hline
    \textbf{Status} & Existing Feature; candidate CT3 placement & \\
    \hline
    \textbf{Ancestry} &  & \\
    \hline
    \textbf{Source} & FEAT-SD-14 & \\
    \hline
    \textbf{Enabled By} &  & \\
    \hline
    \textbf{Supports} &  & \\
    \hline
    \textbf{Derived From} &  & \\
    \hline
    \textbf{Reviewed Status} & Confirm existing Feature meaning and proposed CT3 ownership. & \\
    \hline
  \end{tabularx}
\end{table}

Intent. Screening Recall Notification Distribution provides a distinct enabling behaviour in the context of Screening and Immunisation Enablement. This is an existing named Feature; its placement under the candidate CT3 remains subject to review.

\noindent\textbf{Feature boundary.} This behaviour is scoped to Screening and Immunisation Enablement; an identically named Feature elsewhere may carry different business meaning.\par
\featuretable{
Function & \texttt{FN-8E259DF2} & Project Screening Eligibility Timeline & Catalogue \\
Service & \texttt{SV-226A3300} & Preventive Screening Status Query & Linked to preceding Function in source \\
}
\subsubsection{Feature - Context Establishment}\label{cap:bec-261578c2}
\vspace{-1em}
\begin{table}[htbp]
  \centering
  \footnotesize
  \begin{tabularx}{1\textwidth}{|l|l|X|}
    \hline
    \textbf{Metadata Name} & \textbf{Metadata Value} & \textbf{Description} \\
    \hline
    \textbf{Reference Id} & \texttt{BEC-261578C2} & \\
    \hline
    \textbf{Status} & Generated candidate & \\
    \hline
    \textbf{Ancestry} &  & \\
    \hline
    \textbf{Source} & Semantic-pattern candidate & \\
    \hline
    \textbf{Enabled By} &  & \\
    \hline
    \textbf{Supports} &  & \\
    \hline
    \textbf{Derived From} &  & \\
    \hline
    \textbf{Reviewed Status} & Generated candidate: keep / modify / relocate / reject; verify whether independently testable. & \\
    \hline
  \end{tabularx}
\end{table}

Intent. Context Establishment is proposed as a discrete, testable enabling behaviour for Screening and Immunisation Enablement. It is included for architectural review, not as an established implementation obligation.

\noindent\textbf{Feature boundary.} This behaviour is scoped to Screening and Immunisation Enablement; an identically named Feature elsewhere may carry different business meaning.\par
\featuretable{
--- & --- & \emph{Not yet captured} & \emph{Complete during review} \\
}
\subsubsection{Feature - State Monitoring}\label{cap:bec-c577b7c8}
\vspace{-1em}
\begin{table}[htbp]
  \centering
  \footnotesize
  \begin{tabularx}{1\textwidth}{|l|l|X|}
    \hline
    \textbf{Metadata Name} & \textbf{Metadata Value} & \textbf{Description} \\
    \hline
    \textbf{Reference Id} & \texttt{BEC-C577B7C8} & \\
    \hline
    \textbf{Status} & Generated candidate & \\
    \hline
    \textbf{Ancestry} &  & \\
    \hline
    \textbf{Source} & Semantic-pattern candidate & \\
    \hline
    \textbf{Enabled By} &  & \\
    \hline
    \textbf{Supports} &  & \\
    \hline
    \textbf{Derived From} &  & \\
    \hline
    \textbf{Reviewed Status} & Generated candidate: keep / modify / relocate / reject; verify whether independently testable. & \\
    \hline
  \end{tabularx}
\end{table}

Intent. State Monitoring is proposed as a discrete, testable enabling behaviour for Screening and Immunisation Enablement. It is included for architectural review, not as an established implementation obligation.

\noindent\textbf{Feature boundary.} This behaviour is scoped to Screening and Immunisation Enablement; an identically named Feature elsewhere may carry different business meaning.\par
\featuretable{
--- & --- & \emph{Not yet captured} & \emph{Complete during review} \\
}
\subsubsection{Feature - Exception Management}\label{cap:bec-f9613d20}
\vspace{-1em}
\begin{table}[htbp]
  \centering
  \footnotesize
  \begin{tabularx}{1\textwidth}{|l|l|X|}
    \hline
    \textbf{Metadata Name} & \textbf{Metadata Value} & \textbf{Description} \\
    \hline
    \textbf{Reference Id} & \texttt{BEC-F9613D20} & \\
    \hline
    \textbf{Status} & Generated candidate & \\
    \hline
    \textbf{Ancestry} &  & \\
    \hline
    \textbf{Source} & Semantic-pattern candidate & \\
    \hline
    \textbf{Enabled By} &  & \\
    \hline
    \textbf{Supports} &  & \\
    \hline
    \textbf{Derived From} &  & \\
    \hline
    \textbf{Reviewed Status} & Generated candidate: keep / modify / relocate / reject; verify whether independently testable. & \\
    \hline
  \end{tabularx}
\end{table}

Intent. Exception Management is proposed as a discrete, testable enabling behaviour for Screening and Immunisation Enablement. It is included for architectural review, not as an established implementation obligation.

\noindent\textbf{Feature boundary.} This behaviour is scoped to Screening and Immunisation Enablement; an identically named Feature elsewhere may carry different business meaning.\par
\featuretable{
--- & --- & \emph{Not yet captured} & \emph{Complete during review} \\
}
\subsubsection{Feature - Information Retrieval}\label{cap:bec-336aaf9d}
\vspace{-1em}
\begin{table}[htbp]
  \centering
  \footnotesize
  \begin{tabularx}{1\textwidth}{|l|l|X|}
    \hline
    \textbf{Metadata Name} & \textbf{Metadata Value} & \textbf{Description} \\
    \hline
    \textbf{Reference Id} & \texttt{BEC-336AAF9D} & \\
    \hline
    \textbf{Status} & Generated candidate & \\
    \hline
    \textbf{Ancestry} &  & \\
    \hline
    \textbf{Source} & Semantic-pattern candidate & \\
    \hline
    \textbf{Enabled By} &  & \\
    \hline
    \textbf{Supports} &  & \\
    \hline
    \textbf{Derived From} &  & \\
    \hline
    \textbf{Reviewed Status} & Generated candidate: keep / modify / relocate / reject; verify whether independently testable. & \\
    \hline
  \end{tabularx}
\end{table}

Intent. Information Retrieval is proposed as a discrete, testable enabling behaviour for Screening and Immunisation Enablement. It is included for architectural review, not as an established implementation obligation.

\noindent\textbf{Feature boundary.} This behaviour is scoped to Screening and Immunisation Enablement; an identically named Feature elsewhere may carry different business meaning.\par
\featuretable{
--- & --- & \emph{Not yet captured} & \emph{Complete during review} \\
}

\section{CT2 - Rehabilitation}\label{cap:bec-a8289aed}
\vspace{-0.5em}
\begin{table}[htbp]
  \centering
  \footnotesize
  \begin{tabularx}{1\textwidth}{|l|l|X|}
    \hline
    \multicolumn{3}{c}{\textbf{Capability Metadata}} \\
    \hline
    \textbf{Name} & \textbf{Value} & \textbf{Comment and/or Description} \\
    \hline
    \textbf{Reference Id} & \texttt{BEC-A8289AED} & \\
    \hline
    \textbf{Status} & Candidate ancestry & \\
    \hline
    \textbf{Ancestry} &  & \\
    \hline
    \textbf{Source} &  & \\
    \hline
    \textbf{Enabled By} &  & \\
    \hline
    \textbf{Supports} &  & \\
    \hline
    \textbf{Derived From} &  & \\
    \hline
    \textbf{Reviewed Status} &  & \\
    \hline
  \end{tabularx}
\end{table}

\begin{table}[htbp]
  \centering
  \footnotesize
  \begin{tabularx}{1\textwidth}{|l|l|X|}
    \hline
    \multicolumn{3}{c}{\textbf{Contained Capabilities}} \\
    \hline
    \textbf{Capability Name} & \textbf{Capability Id} & \textbf{Comment} \\
    \hline
    \multicolumn{3}{|l|}{\emph{No contained capabilities are currently identified.}} \\
    \hline
  \end{tabularx}
\end{table}

Intent. Rehabilitation is a candidate enabling capability within Clinical Service Context Enablement. It frames the reusable system-supported behaviours needed to support this healthcare context without prescribing a product, interface, technology or deployment boundary.

Boundary. This capability supports the business activity and does not inherit the authority or accountability of the Business Capability it enables. Its child Features are review candidates until ancestry and semantics are adjudicated.

\section{CT2 - Community Care}\label{cap:bec-cdb89b75}
\vspace{-0.5em}
\begin{table}[htbp]
  \centering
  \footnotesize
  \begin{tabularx}{1\textwidth}{|l|l|X|}
    \hline
    \multicolumn{3}{c}{\textbf{Capability Metadata}} \\
    \hline
    \textbf{Name} & \textbf{Value} & \textbf{Comment and/or Description} \\
    \hline
    \textbf{Reference Id} & \texttt{BEC-CDB89B75} & \\
    \hline
    \textbf{Status} & Candidate ancestry & \\
    \hline
    \textbf{Ancestry} &  & \\
    \hline
    \textbf{Source} &  & \\
    \hline
    \textbf{Enabled By} &  & \\
    \hline
    \textbf{Supports} &  & \\
    \hline
    \textbf{Derived From} &  & \\
    \hline
    \textbf{Reviewed Status} &  & \\
    \hline
  \end{tabularx}
\end{table}

\begin{table}[htbp]
  \centering
  \footnotesize
  \begin{tabularx}{1\textwidth}{|l|l|X|}
    \hline
    \multicolumn{3}{c}{\textbf{Contained Capabilities}} \\
    \hline
    \textbf{Capability Name} & \textbf{Capability Id} & \textbf{Comment} \\
    \hline
    Community Care Coordination & \texttt{cap:bec-50748227} (at \ref{cap:bec-50748227}) & \\
    \hline
  \end{tabularx}
\end{table}

Intent. Community Care is a candidate enabling capability within Clinical Service Context Enablement. It frames the reusable system-supported behaviours needed to support this healthcare context without prescribing a product, interface, technology or deployment boundary.

Boundary. This capability supports the business activity and does not inherit the authority or accountability of the Business Capability it enables. Its child Features are review candidates until ancestry and semantics are adjudicated.

\subsection{CT3 - Community Care Coordination}\label{cap:bec-50748227}
\vspace{-0.5em}
\begin{table}[htbp]
  \centering
  \footnotesize
  \begin{tabularx}{1\textwidth}{|l|l|X|}
    \hline
    \multicolumn{3}{c}{\textbf{Capability Metadata}} \\
    \hline
    \textbf{Name} & \textbf{Value} & \textbf{Comment and/or Description} \\
    \hline
    \textbf{Reference Id} & \texttt{BEC-50748227} & \\
    \hline
    \textbf{Status} & Candidate ancestry & \\
    \hline
    \textbf{Ancestry} &  & \\
    \hline
    \textbf{Source} &  & \\
    \hline
    \textbf{Enabled By} &  & \\
    \hline
    \textbf{Supports} &  & \\
    \hline
    \textbf{Derived From} &  & \\
    \hline
    \textbf{Reviewed Status} &  & \\
    \hline
  \end{tabularx}
\end{table}

\begin{table}[htbp]
  \centering
  \footnotesize
  \begin{tabularx}{1\textwidth}{|l|l|X|}
    \hline
    \multicolumn{3}{c}{\textbf{Contained Features}} \\
    \hline
    \textbf{Feature Name} & \textbf{Feature Id} & \textbf{Comment} \\
    \hline
    Community Encounter Ingestion & \texttt{cap:bec-70eef999} (at \ref{cap:bec-70eef999}) & \\
    \hline
    Community Care Plan Synchronization & \texttt{cap:bec-5cb15f94} (at \ref{cap:bec-5cb15f94}) & \\
    \hline
    Context Establishment & \texttt{cap:bec-d96c0908} (at \ref{cap:bec-d96c0908}) & \\
    \hline
    Eligibility Validation & \texttt{cap:bec-a1012bd3} (at \ref{cap:bec-a1012bd3}) & \\
    \hline
    Assignment \& Routing & \texttt{cap:bec-5d53cb17} (at \ref{cap:bec-5d53cb17}) & \\
    \hline
    State Tracking & \texttt{cap:bec-52c8768a} (at \ref{cap:bec-52c8768a}) & \\
    \hline
    Exception Management & \texttt{cap:bec-a79a3ed2} (at \ref{cap:bec-a79a3ed2}) & \\
    \hline
    Outcome Reporting & \texttt{cap:bec-1725aa08} (at \ref{cap:bec-1725aa08}) & \\
    \hline
  \end{tabularx}
\end{table}

Intent. Community Care Coordination is a candidate enabling capability within Community Care. It frames the reusable system-supported behaviours needed to support this healthcare context without prescribing a product, interface, technology or deployment boundary.

Boundary. This capability supports the business activity and does not inherit the authority or accountability of the Business Capability it enables. Its child Features are review candidates until ancestry and semantics are adjudicated.

\subsubsection{Feature - Community Encounter Ingestion}\label{cap:bec-70eef999}
\vspace{-1em}
\begin{table}[htbp]
  \centering
  \footnotesize
  \begin{tabularx}{1\textwidth}{|l|l|X|}
    \hline
    \textbf{Metadata Name} & \textbf{Metadata Value} & \textbf{Description} \\
    \hline
    \textbf{Reference Id} & \texttt{BEC-70EEF999} & \\
    \hline
    \textbf{Status} & Existing Feature; candidate CT3 placement & \\
    \hline
    \textbf{Ancestry} &  & \\
    \hline
    \textbf{Source} & FEAT-SD-15 & \\
    \hline
    \textbf{Enabled By} &  & \\
    \hline
    \textbf{Supports} &  & \\
    \hline
    \textbf{Derived From} &  & \\
    \hline
    \textbf{Reviewed Status} & Confirm existing Feature meaning and proposed CT3 ownership. & \\
    \hline
  \end{tabularx}
\end{table}

Intent. Community Encounter Ingestion provides a distinct enabling behaviour in the context of Community Care Coordination. This is an existing named Feature; its placement under the candidate CT3 remains subject to review.

\noindent\textbf{Feature boundary.} This behaviour is scoped to Community Care Coordination; an identically named Feature elsewhere may carry different business meaning.\par
\featuretable{
Function & \texttt{FN-62234F7C} & Synchronize Community Visit Record & Catalogue \\
Service & \texttt{SV-36D96317} & Community Care Encounter Ingress & Linked to preceding Function in source \\
}
\subsubsection{Feature - Community Care Plan Synchronization}\label{cap:bec-5cb15f94}
\vspace{-1em}
\begin{table}[htbp]
  \centering
  \footnotesize
  \begin{tabularx}{1\textwidth}{|l|l|X|}
    \hline
    \textbf{Metadata Name} & \textbf{Metadata Value} & \textbf{Description} \\
    \hline
    \textbf{Reference Id} & \texttt{BEC-5CB15F94} & \\
    \hline
    \textbf{Status} & Existing Feature; candidate CT3 placement & \\
    \hline
    \textbf{Ancestry} &  & \\
    \hline
    \textbf{Source} & FEAT-SD-16 & \\
    \hline
    \textbf{Enabled By} &  & \\
    \hline
    \textbf{Supports} &  & \\
    \hline
    \textbf{Derived From} &  & \\
    \hline
    \textbf{Reviewed Status} & Confirm existing Feature meaning and proposed CT3 ownership. & \\
    \hline
  \end{tabularx}
\end{table}

Intent. Community Care Plan Synchronization provides a distinct enabling behaviour in the context of Community Care Coordination. This is an existing named Feature; its placement under the candidate CT3 remains subject to review.

\noindent\textbf{Feature boundary.} This behaviour is scoped to Community Care Coordination; an identically named Feature elsewhere may carry different business meaning.\par
\featuretable{
Function & \texttt{FN-26D511BC} & Bind Community Support Network Context & Catalogue \\
Service & \texttt{SV-9692370B} & Community Care Team Query & Linked to preceding Function in source \\
}
\subsubsection{Feature - Context Establishment}\label{cap:bec-d96c0908}
\vspace{-1em}
\begin{table}[htbp]
  \centering
  \footnotesize
  \begin{tabularx}{1\textwidth}{|l|l|X|}
    \hline
    \textbf{Metadata Name} & \textbf{Metadata Value} & \textbf{Description} \\
    \hline
    \textbf{Reference Id} & \texttt{BEC-D96C0908} & \\
    \hline
    \textbf{Status} & Generated candidate & \\
    \hline
    \textbf{Ancestry} &  & \\
    \hline
    \textbf{Source} & Semantic-pattern candidate & \\
    \hline
    \textbf{Enabled By} &  & \\
    \hline
    \textbf{Supports} &  & \\
    \hline
    \textbf{Derived From} &  & \\
    \hline
    \textbf{Reviewed Status} & Generated candidate: keep / modify / relocate / reject; verify whether independently testable. & \\
    \hline
  \end{tabularx}
\end{table}

Intent. Context Establishment is proposed as a discrete, testable enabling behaviour for Community Care Coordination. It is included for architectural review, not as an established implementation obligation.

\noindent\textbf{Feature boundary.} This behaviour is scoped to Community Care Coordination; an identically named Feature elsewhere may carry different business meaning.\par
\featuretable{
--- & --- & \emph{Not yet captured} & \emph{Complete during review} \\
}
\subsubsection{Feature - Eligibility Validation}\label{cap:bec-a1012bd3}
\vspace{-1em}
\begin{table}[htbp]
  \centering
  \footnotesize
  \begin{tabularx}{1\textwidth}{|l|l|X|}
    \hline
    \textbf{Metadata Name} & \textbf{Metadata Value} & \textbf{Description} \\
    \hline
    \textbf{Reference Id} & \texttt{BEC-A1012BD3} & \\
    \hline
    \textbf{Status} & Generated candidate & \\
    \hline
    \textbf{Ancestry} &  & \\
    \hline
    \textbf{Source} & Semantic-pattern candidate & \\
    \hline
    \textbf{Enabled By} &  & \\
    \hline
    \textbf{Supports} &  & \\
    \hline
    \textbf{Derived From} &  & \\
    \hline
    \textbf{Reviewed Status} & Generated candidate: keep / modify / relocate / reject; verify whether independently testable. & \\
    \hline
  \end{tabularx}
\end{table}

Intent. Eligibility Validation is proposed as a discrete, testable enabling behaviour for Community Care Coordination. It is included for architectural review, not as an established implementation obligation.

\noindent\textbf{Feature boundary.} This behaviour is scoped to Community Care Coordination; an identically named Feature elsewhere may carry different business meaning.\par
\featuretable{
--- & --- & \emph{Not yet captured} & \emph{Complete during review} \\
}
\subsubsection{Feature - Assignment \& Routing}\label{cap:bec-5d53cb17}
\vspace{-1em}
\begin{table}[htbp]
  \centering
  \footnotesize
  \begin{tabularx}{1\textwidth}{|l|l|X|}
    \hline
    \textbf{Metadata Name} & \textbf{Metadata Value} & \textbf{Description} \\
    \hline
    \textbf{Reference Id} & \texttt{BEC-5D53CB17} & \\
    \hline
    \textbf{Status} & Generated candidate & \\
    \hline
    \textbf{Ancestry} &  & \\
    \hline
    \textbf{Source} & Semantic-pattern candidate & \\
    \hline
    \textbf{Enabled By} &  & \\
    \hline
    \textbf{Supports} &  & \\
    \hline
    \textbf{Derived From} &  & \\
    \hline
    \textbf{Reviewed Status} & Generated candidate: keep / modify / relocate / reject; verify whether independently testable. & \\
    \hline
  \end{tabularx}
\end{table}

Intent. Assignment \& Routing is proposed as a discrete, testable enabling behaviour for Community Care Coordination. It is included for architectural review, not as an established implementation obligation.

\noindent\textbf{Feature boundary.} This behaviour is scoped to Community Care Coordination; an identically named Feature elsewhere may carry different business meaning.\par
\featuretable{
--- & --- & \emph{Not yet captured} & \emph{Complete during review} \\
}
\subsubsection{Feature - State Tracking}\label{cap:bec-52c8768a}
\vspace{-1em}
\begin{table}[htbp]
  \centering
  \footnotesize
  \begin{tabularx}{1\textwidth}{|l|l|X|}
    \hline
    \textbf{Metadata Name} & \textbf{Metadata Value} & \textbf{Description} \\
    \hline
    \textbf{Reference Id} & \texttt{BEC-52C8768A} & \\
    \hline
    \textbf{Status} & Generated candidate & \\
    \hline
    \textbf{Ancestry} &  & \\
    \hline
    \textbf{Source} & Semantic-pattern candidate & \\
    \hline
    \textbf{Enabled By} &  & \\
    \hline
    \textbf{Supports} &  & \\
    \hline
    \textbf{Derived From} &  & \\
    \hline
    \textbf{Reviewed Status} & Generated candidate: keep / modify / relocate / reject; verify whether independently testable. & \\
    \hline
  \end{tabularx}
\end{table}

Intent. State Tracking is proposed as a discrete, testable enabling behaviour for Community Care Coordination. It is included for architectural review, not as an established implementation obligation.

\noindent\textbf{Feature boundary.} This behaviour is scoped to Community Care Coordination; an identically named Feature elsewhere may carry different business meaning.\par
\featuretable{
--- & --- & \emph{Not yet captured} & \emph{Complete during review} \\
}
\subsubsection{Feature - Exception Management}\label{cap:bec-a79a3ed2}
\vspace{-1em}
\begin{table}[htbp]
  \centering
  \footnotesize
  \begin{tabularx}{1\textwidth}{|l|l|X|}
    \hline
    \textbf{Metadata Name} & \textbf{Metadata Value} & \textbf{Description} \\
    \hline
    \textbf{Reference Id} & \texttt{BEC-A79A3ED2} & \\
    \hline
    \textbf{Status} & Generated candidate & \\
    \hline
    \textbf{Ancestry} &  & \\
    \hline
    \textbf{Source} & Semantic-pattern candidate & \\
    \hline
    \textbf{Enabled By} &  & \\
    \hline
    \textbf{Supports} &  & \\
    \hline
    \textbf{Derived From} &  & \\
    \hline
    \textbf{Reviewed Status} & Generated candidate: keep / modify / relocate / reject; verify whether independently testable. & \\
    \hline
  \end{tabularx}
\end{table}

Intent. Exception Management is proposed as a discrete, testable enabling behaviour for Community Care Coordination. It is included for architectural review, not as an established implementation obligation.

\noindent\textbf{Feature boundary.} This behaviour is scoped to Community Care Coordination; an identically named Feature elsewhere may carry different business meaning.\par
\featuretable{
--- & --- & \emph{Not yet captured} & \emph{Complete during review} \\
}
\subsubsection{Feature - Outcome Reporting}\label{cap:bec-1725aa08}
\vspace{-1em}
\begin{table}[htbp]
  \centering
  \footnotesize
  \begin{tabularx}{1\textwidth}{|l|l|X|}
    \hline
    \textbf{Metadata Name} & \textbf{Metadata Value} & \textbf{Description} \\
    \hline
    \textbf{Reference Id} & \texttt{BEC-1725AA08} & \\
    \hline
    \textbf{Status} & Generated candidate & \\
    \hline
    \textbf{Ancestry} &  & \\
    \hline
    \textbf{Source} & Semantic-pattern candidate & \\
    \hline
    \textbf{Enabled By} &  & \\
    \hline
    \textbf{Supports} &  & \\
    \hline
    \textbf{Derived From} &  & \\
    \hline
    \textbf{Reviewed Status} & Generated candidate: keep / modify / relocate / reject; verify whether independently testable. & \\
    \hline
  \end{tabularx}
\end{table}

Intent. Outcome Reporting is proposed as a discrete, testable enabling behaviour for Community Care Coordination. It is included for architectural review, not as an established implementation obligation.

\noindent\textbf{Feature boundary.} This behaviour is scoped to Community Care Coordination; an identically named Feature elsewhere may carry different business meaning.\par
\featuretable{
--- & --- & \emph{Not yet captured} & \emph{Complete during review} \\
}

\section{CT2 - Outreach Care}\label{cap:bec-9ddd1bd3}
\vspace{-0.5em}
\begin{table}[htbp]
  \centering
  \footnotesize
  \begin{tabularx}{1\textwidth}{|l|l|X|}
    \hline
    \multicolumn{3}{c}{\textbf{Capability Metadata}} \\
    \hline
    \textbf{Name} & \textbf{Value} & \textbf{Comment and/or Description} \\
    \hline
    \textbf{Reference Id} & \texttt{BEC-9DDD1BD3} & \\
    \hline
    \textbf{Status} & Candidate ancestry & \\
    \hline
    \textbf{Ancestry} &  & \\
    \hline
    \textbf{Source} &  & \\
    \hline
    \textbf{Enabled By} &  & \\
    \hline
    \textbf{Supports} &  & \\
    \hline
    \textbf{Derived From} &  & \\
    \hline
    \textbf{Reviewed Status} &  & \\
    \hline
  \end{tabularx}
\end{table}

\begin{table}[htbp]
  \centering
  \footnotesize
  \begin{tabularx}{1\textwidth}{|l|l|X|}
    \hline
    \multicolumn{3}{c}{\textbf{Contained Capabilities}} \\
    \hline
    \textbf{Capability Name} & \textbf{Capability Id} & \textbf{Comment} \\
    \hline
    Outreach Care Continuity & \texttt{cap:bec-b8eb379a} (at \ref{cap:bec-b8eb379a}) & \\
    \hline
  \end{tabularx}
\end{table}

Intent. Outreach Care is a candidate enabling capability within Clinical Service Context Enablement. It frames the reusable system-supported behaviours needed to support this healthcare context without prescribing a product, interface, technology or deployment boundary.

Boundary. This capability supports the business activity and does not inherit the authority or accountability of the Business Capability it enables. Its child Features are review candidates until ancestry and semantics are adjudicated.

\subsection{CT3 - Outreach Care Continuity}\label{cap:bec-b8eb379a}
\vspace{-0.5em}
\begin{table}[htbp]
  \centering
  \footnotesize
  \begin{tabularx}{1\textwidth}{|l|l|X|}
    \hline
    \multicolumn{3}{c}{\textbf{Capability Metadata}} \\
    \hline
    \textbf{Name} & \textbf{Value} & \textbf{Comment and/or Description} \\
    \hline
    \textbf{Reference Id} & \texttt{BEC-B8EB379A} & \\
    \hline
    \textbf{Status} & Candidate ancestry & \\
    \hline
    \textbf{Ancestry} &  & \\
    \hline
    \textbf{Source} &  & \\
    \hline
    \textbf{Enabled By} &  & \\
    \hline
    \textbf{Supports} &  & \\
    \hline
    \textbf{Derived From} &  & \\
    \hline
    \textbf{Reviewed Status} &  & \\
    \hline
  \end{tabularx}
\end{table}

\begin{table}[htbp]
  \centering
  \footnotesize
  \begin{tabularx}{1\textwidth}{|l|l|X|}
    \hline
    \multicolumn{3}{c}{\textbf{Contained Features}} \\
    \hline
    \textbf{Feature Name} & \textbf{Feature Id} & \textbf{Comment} \\
    \hline
    Outreach Visit Context Provision & \texttt{cap:bec-6a90478e} (at \ref{cap:bec-6a90478e}) & \\
    \hline
    Outreach Encounter Outcome Ingestion & \texttt{cap:bec-8bcc3746} (at \ref{cap:bec-8bcc3746}) & \\
    \hline
    Context Establishment & \texttt{cap:bec-89f15aef} (at \ref{cap:bec-89f15aef}) & \\
    \hline
    State Monitoring & \texttt{cap:bec-59be4e9f} (at \ref{cap:bec-59be4e9f}) & \\
    \hline
    Exception Management & \texttt{cap:bec-8f38ae8d} (at \ref{cap:bec-8f38ae8d}) & \\
    \hline
    Information Retrieval & \texttt{cap:bec-607a7a67} (at \ref{cap:bec-607a7a67}) & \\
    \hline
  \end{tabularx}
\end{table}

Intent. Outreach Care Continuity is a candidate enabling capability within Outreach Care. It frames the reusable system-supported behaviours needed to support this healthcare context without prescribing a product, interface, technology or deployment boundary.

Boundary. This capability supports the business activity and does not inherit the authority or accountability of the Business Capability it enables. Its child Features are review candidates until ancestry and semantics are adjudicated.

\subsubsection{Feature - Outreach Visit Context Provision}\label{cap:bec-6a90478e}
\vspace{-1em}
\begin{table}[htbp]
  \centering
  \footnotesize
  \begin{tabularx}{1\textwidth}{|l|l|X|}
    \hline
    \textbf{Metadata Name} & \textbf{Metadata Value} & \textbf{Description} \\
    \hline
    \textbf{Reference Id} & \texttt{BEC-6A90478E} & \\
    \hline
    \textbf{Status} & Existing Feature; candidate CT3 placement & \\
    \hline
    \textbf{Ancestry} &  & \\
    \hline
    \textbf{Source} & FEAT-SD-17 & \\
    \hline
    \textbf{Enabled By} &  & \\
    \hline
    \textbf{Supports} &  & \\
    \hline
    \textbf{Derived From} &  & \\
    \hline
    \textbf{Reviewed Status} & Confirm existing Feature meaning and proposed CT3 ownership. & \\
    \hline
  \end{tabularx}
\end{table}

Intent. Outreach Visit Context Provision provides a distinct enabling behaviour in the context of Outreach Care Continuity. This is an existing named Feature; its placement under the candidate CT3 remains subject to review.

\noindent\textbf{Feature boundary.} This behaviour is scoped to Outreach Care Continuity; an identically named Feature elsewhere may carry different business meaning.\par
\featuretable{
Function & \texttt{FN-2393967C} & Provide Outreach Visit Context & Catalogue \\
}
\subsubsection{Feature - Outreach Encounter Outcome Ingestion}\label{cap:bec-8bcc3746}
\vspace{-1em}
\begin{table}[htbp]
  \centering
  \footnotesize
  \begin{tabularx}{1\textwidth}{|l|l|X|}
    \hline
    \textbf{Metadata Name} & \textbf{Metadata Value} & \textbf{Description} \\
    \hline
    \textbf{Reference Id} & \texttt{BEC-8BCC3746} & \\
    \hline
    \textbf{Status} & Existing Feature; candidate CT3 placement & \\
    \hline
    \textbf{Ancestry} &  & \\
    \hline
    \textbf{Source} & FEAT-SD-18 & \\
    \hline
    \textbf{Enabled By} &  & \\
    \hline
    \textbf{Supports} &  & \\
    \hline
    \textbf{Derived From} &  & \\
    \hline
    \textbf{Reviewed Status} & Confirm existing Feature meaning and proposed CT3 ownership. & \\
    \hline
  \end{tabularx}
\end{table}

Intent. Outreach Encounter Outcome Ingestion provides a distinct enabling behaviour in the context of Outreach Care Continuity. This is an existing named Feature; its placement under the candidate CT3 remains subject to review.

\noindent\textbf{Feature boundary.} This behaviour is scoped to Outreach Care Continuity; an identically named Feature elsewhere may carry different business meaning.\par
\featuretable{
Function & \texttt{FN-72154678} & Reconcile Outreach Clinical Batch & Catalogue \\
Service & \texttt{SV-C61F82DF} & Outreach Batch Reconciliation Service & Linked to preceding Function in source \\
}
\subsubsection{Feature - Context Establishment}\label{cap:bec-89f15aef}
\vspace{-1em}
\begin{table}[htbp]
  \centering
  \footnotesize
  \begin{tabularx}{1\textwidth}{|l|l|X|}
    \hline
    \textbf{Metadata Name} & \textbf{Metadata Value} & \textbf{Description} \\
    \hline
    \textbf{Reference Id} & \texttt{BEC-89F15AEF} & \\
    \hline
    \textbf{Status} & Generated candidate & \\
    \hline
    \textbf{Ancestry} &  & \\
    \hline
    \textbf{Source} & Semantic-pattern candidate & \\
    \hline
    \textbf{Enabled By} &  & \\
    \hline
    \textbf{Supports} &  & \\
    \hline
    \textbf{Derived From} &  & \\
    \hline
    \textbf{Reviewed Status} & Generated candidate: keep / modify / relocate / reject; verify whether independently testable. & \\
    \hline
  \end{tabularx}
\end{table}

Intent. Context Establishment is proposed as a discrete, testable enabling behaviour for Outreach Care Continuity. It is included for architectural review, not as an established implementation obligation.

\noindent\textbf{Feature boundary.} This behaviour is scoped to Outreach Care Continuity; an identically named Feature elsewhere may carry different business meaning.\par
\featuretable{
--- & --- & \emph{Not yet captured} & \emph{Complete during review} \\
}
\subsubsection{Feature - State Monitoring}\label{cap:bec-59be4e9f}
\vspace{-1em}
\begin{table}[htbp]
  \centering
  \footnotesize
  \begin{tabularx}{1\textwidth}{|l|l|X|}
    \hline
    \textbf{Metadata Name} & \textbf{Metadata Value} & \textbf{Description} \\
    \hline
    \textbf{Reference Id} & \texttt{BEC-59BE4E9F} & \\
    \hline
    \textbf{Status} & Generated candidate & \\
    \hline
    \textbf{Ancestry} &  & \\
    \hline
    \textbf{Source} & Semantic-pattern candidate & \\
    \hline
    \textbf{Enabled By} &  & \\
    \hline
    \textbf{Supports} &  & \\
    \hline
    \textbf{Derived From} &  & \\
    \hline
    \textbf{Reviewed Status} & Generated candidate: keep / modify / relocate / reject; verify whether independently testable. & \\
    \hline
  \end{tabularx}
\end{table}

Intent. State Monitoring is proposed as a discrete, testable enabling behaviour for Outreach Care Continuity. It is included for architectural review, not as an established implementation obligation.

\noindent\textbf{Feature boundary.} This behaviour is scoped to Outreach Care Continuity; an identically named Feature elsewhere may carry different business meaning.\par
\featuretable{
--- & --- & \emph{Not yet captured} & \emph{Complete during review} \\
}
\subsubsection{Feature - Exception Management}\label{cap:bec-8f38ae8d}
\vspace{-1em}
\begin{table}[htbp]
  \centering
  \footnotesize
  \begin{tabularx}{1\textwidth}{|l|l|X|}
    \hline
    \textbf{Metadata Name} & \textbf{Metadata Value} & \textbf{Description} \\
    \hline
    \textbf{Reference Id} & \texttt{BEC-8F38AE8D} & \\
    \hline
    \textbf{Status} & Generated candidate & \\
    \hline
    \textbf{Ancestry} &  & \\
    \hline
    \textbf{Source} & Semantic-pattern candidate & \\
    \hline
    \textbf{Enabled By} &  & \\
    \hline
    \textbf{Supports} &  & \\
    \hline
    \textbf{Derived From} &  & \\
    \hline
    \textbf{Reviewed Status} & Generated candidate: keep / modify / relocate / reject; verify whether independently testable. & \\
    \hline
  \end{tabularx}
\end{table}

Intent. Exception Management is proposed as a discrete, testable enabling behaviour for Outreach Care Continuity. It is included for architectural review, not as an established implementation obligation.

\noindent\textbf{Feature boundary.} This behaviour is scoped to Outreach Care Continuity; an identically named Feature elsewhere may carry different business meaning.\par
\featuretable{
--- & --- & \emph{Not yet captured} & \emph{Complete during review} \\
}
\subsubsection{Feature - Information Retrieval}\label{cap:bec-607a7a67}
\vspace{-1em}
\begin{table}[htbp]
  \centering
  \footnotesize
  \begin{tabularx}{1\textwidth}{|l|l|X|}
    \hline
    \textbf{Metadata Name} & \textbf{Metadata Value} & \textbf{Description} \\
    \hline
    \textbf{Reference Id} & \texttt{BEC-607A7A67} & \\
    \hline
    \textbf{Status} & Generated candidate & \\
    \hline
    \textbf{Ancestry} &  & \\
    \hline
    \textbf{Source} & Semantic-pattern candidate & \\
    \hline
    \textbf{Enabled By} &  & \\
    \hline
    \textbf{Supports} &  & \\
    \hline
    \textbf{Derived From} &  & \\
    \hline
    \textbf{Reviewed Status} & Generated candidate: keep / modify / relocate / reject; verify whether independently testable. & \\
    \hline
  \end{tabularx}
\end{table}

Intent. Information Retrieval is proposed as a discrete, testable enabling behaviour for Outreach Care Continuity. It is included for architectural review, not as an established implementation obligation.

\noindent\textbf{Feature boundary.} This behaviour is scoped to Outreach Care Continuity; an identically named Feature elsewhere may carry different business meaning.\par
\featuretable{
--- & --- & \emph{Not yet captured} & \emph{Complete during review} \\
}

\section{CT2 - Remote \& Virtual Care}\label{cap:bec-f9e9dd8e}
\vspace{-0.5em}
\begin{table}[htbp]
  \centering
  \footnotesize
  \begin{tabularx}{1\textwidth}{|l|l|X|}
    \hline
    \multicolumn{3}{c}{\textbf{Capability Metadata}} \\
    \hline
    \textbf{Name} & \textbf{Value} & \textbf{Comment and/or Description} \\
    \hline
    \textbf{Reference Id} & \texttt{BEC-F9E9DD8E} & \\
    \hline
    \textbf{Status} & Candidate ancestry & \\
    \hline
    \textbf{Ancestry} &  & \\
    \hline
    \textbf{Source} &  & \\
    \hline
    \textbf{Enabled By} &  & \\
    \hline
    \textbf{Supports} &  & \\
    \hline
    \textbf{Derived From} &  & \\
    \hline
    \textbf{Reviewed Status} &  & \\
    \hline
  \end{tabularx}
\end{table}

\begin{table}[htbp]
  \centering
  \footnotesize
  \begin{tabularx}{1\textwidth}{|l|l|X|}
    \hline
    \multicolumn{3}{c}{\textbf{Contained Capabilities}} \\
    \hline
    \textbf{Capability Name} & \textbf{Capability Id} & \textbf{Comment} \\
    \hline
    Remote Patient Monitoring & \texttt{cap:bec-68c27168} (at \ref{cap:bec-68c27168}) & \\
    \hline
  \end{tabularx}
\end{table}

Intent. Remote \& Virtual Care is a candidate enabling capability within Clinical Service Context Enablement. It frames the reusable system-supported behaviours needed to support this healthcare context without prescribing a product, interface, technology or deployment boundary.

Boundary. This capability supports the business activity and does not inherit the authority or accountability of the Business Capability it enables. Its child Features are review candidates until ancestry and semantics are adjudicated.

\subsection{CT3 - Remote Patient Monitoring}\label{cap:bec-68c27168}
\vspace{-0.5em}
\begin{table}[htbp]
  \centering
  \footnotesize
  \begin{tabularx}{1\textwidth}{|l|l|X|}
    \hline
    \multicolumn{3}{c}{\textbf{Capability Metadata}} \\
    \hline
    \textbf{Name} & \textbf{Value} & \textbf{Comment and/or Description} \\
    \hline
    \textbf{Reference Id} & \texttt{BEC-68C27168} & \\
    \hline
    \textbf{Status} & Candidate ancestry & \\
    \hline
    \textbf{Ancestry} &  & \\
    \hline
    \textbf{Source} &  & \\
    \hline
    \textbf{Enabled By} &  & \\
    \hline
    \textbf{Supports} &  & \\
    \hline
    \textbf{Derived From} &  & \\
    \hline
    \textbf{Reviewed Status} &  & \\
    \hline
  \end{tabularx}
\end{table}

\begin{table}[htbp]
  \centering
  \footnotesize
  \begin{tabularx}{1\textwidth}{|l|l|X|}
    \hline
    \multicolumn{3}{c}{\textbf{Contained Features}} \\
    \hline
    \textbf{Feature Name} & \textbf{Feature Id} & \textbf{Comment} \\
    \hline
    Remote Telemetry Ingestion & \texttt{cap:bec-fac2e484} (at \ref{cap:bec-fac2e484}) & \\
    \hline
    Virtual Care Session Context Binding & \texttt{cap:bec-d4a0665c} (at \ref{cap:bec-d4a0665c}) & \\
    \hline
    Subject Identity Correlation & \texttt{cap:bec-08d640dd} (at \ref{cap:bec-08d640dd}) & \\
    \hline
    Subject Relationship Management & \texttt{cap:bec-2a678be0} (at \ref{cap:bec-2a678be0}) & \\
    \hline
    Subject Context Management & \texttt{cap:bec-0b235f68} (at \ref{cap:bec-0b235f68}) & \\
    \hline
    Subject Information Discovery & \texttt{cap:bec-6b015be0} (at \ref{cap:bec-6b015be0}) & \\
    \hline
  \end{tabularx}
\end{table}

Intent. Remote Patient Monitoring is a candidate enabling capability within Remote \& Virtual Care. It frames the reusable system-supported behaviours needed to support this healthcare context without prescribing a product, interface, technology or deployment boundary.

Boundary. This capability supports the business activity and does not inherit the authority or accountability of the Business Capability it enables. Its child Features are review candidates until ancestry and semantics are adjudicated.

FHIR R5 conceptual analogues considered (not implementation mappings): Patient, Person.

\subsubsection{Feature - Remote Telemetry Ingestion}\label{cap:bec-fac2e484}
\vspace{-1em}
\begin{table}[htbp]
  \centering
  \footnotesize
  \begin{tabularx}{1\textwidth}{|l|l|X|}
    \hline
    \textbf{Metadata Name} & \textbf{Metadata Value} & \textbf{Description} \\
    \hline
    \textbf{Reference Id} & \texttt{BEC-FAC2E484} & \\
    \hline
    \textbf{Status} & Existing Feature; candidate CT3 placement & \\
    \hline
    \textbf{Ancestry} &  & \\
    \hline
    \textbf{Source} & FEAT-SD-19 & \\
    \hline
    \textbf{Enabled By} &  & \\
    \hline
    \textbf{Supports} &  & \\
    \hline
    \textbf{Derived From} &  & \\
    \hline
    \textbf{Reviewed Status} & Confirm existing Feature meaning and proposed CT3 ownership. & \\
    \hline
  \end{tabularx}
\end{table}

Intent. Remote Telemetry Ingestion provides a distinct enabling behaviour in the context of Remote Patient Monitoring. This is an existing named Feature; its placement under the candidate CT3 remains subject to review.

\noindent\textbf{Feature boundary.} This behaviour is scoped to Remote Patient Monitoring; an identically named Feature elsewhere may carry different business meaning.\par
\featuretable{
Function & \texttt{FN-D69C80D1} & Ingest Remote Patient Telemetry & Catalogue \\
Service & \texttt{SV-EA659DEE} & Remote Telemetry Ingress & Linked to preceding Function in source \\
}
\subsubsection{Feature - Virtual Care Session Context Binding}\label{cap:bec-d4a0665c}
\vspace{-1em}
\begin{table}[htbp]
  \centering
  \footnotesize
  \begin{tabularx}{1\textwidth}{|l|l|X|}
    \hline
    \textbf{Metadata Name} & \textbf{Metadata Value} & \textbf{Description} \\
    \hline
    \textbf{Reference Id} & \texttt{BEC-D4A0665C} & \\
    \hline
    \textbf{Status} & Existing Feature; candidate CT3 placement & \\
    \hline
    \textbf{Ancestry} &  & \\
    \hline
    \textbf{Source} & FEAT-SD-20 & \\
    \hline
    \textbf{Enabled By} &  & \\
    \hline
    \textbf{Supports} &  & \\
    \hline
    \textbf{Derived From} &  & \\
    \hline
    \textbf{Reviewed Status} & Confirm existing Feature meaning and proposed CT3 ownership. & \\
    \hline
  \end{tabularx}
\end{table}

Intent. Virtual Care Session Context Binding provides a distinct enabling behaviour in the context of Remote Patient Monitoring. This is an existing named Feature; its placement under the candidate CT3 remains subject to review.

\noindent\textbf{Feature boundary.} This behaviour is scoped to Remote Patient Monitoring; an identically named Feature elsewhere may carry different business meaning.\par
\featuretable{
Function & \texttt{FN-4FAD30E5} & Correlate Virtual Care Consultation Context & Catalogue \\
Service & \texttt{SV-FD3F0132} & Virtual Encounter Context Resolution & Linked to preceding Function in source \\
}
\subsubsection{Feature - Subject Identity Correlation}\label{cap:bec-08d640dd}
\vspace{-1em}
\begin{table}[htbp]
  \centering
  \footnotesize
  \begin{tabularx}{1\textwidth}{|l|l|X|}
    \hline
    \textbf{Metadata Name} & \textbf{Metadata Value} & \textbf{Description} \\
    \hline
    \textbf{Reference Id} & \texttt{BEC-08D640DD} & \\
    \hline
    \textbf{Status} & Generated candidate & \\
    \hline
    \textbf{Ancestry} &  & \\
    \hline
    \textbf{Source} & FHIR-informed candidate & \\
    \hline
    \textbf{Enabled By} &  & \\
    \hline
    \textbf{Supports} &  & \\
    \hline
    \textbf{Derived From} &  & \\
    \hline
    \textbf{Reviewed Status} & Generated candidate: keep / modify / relocate / reject; verify whether independently testable. & \\
    \hline
  \end{tabularx}
\end{table}

Intent. Subject Identity Correlation is proposed to support the underlying healthcare concept represented by Remote Patient Monitoring. The proposal concerns behaviour and semantic responsibility, not a FHIR field, operation, resource instance or prescribed data model.

Conceptual evidence: HL7 FHIR R5 Patient, Person.

\noindent\textbf{Feature boundary.} This behaviour is scoped to Remote Patient Monitoring; an identically named Feature elsewhere may carry different business meaning.\par
\featuretable{
--- & --- & \emph{Not yet captured} & \emph{Complete during review} \\
}
\subsubsection{Feature - Subject Relationship Management}\label{cap:bec-2a678be0}
\vspace{-1em}
\begin{table}[htbp]
  \centering
  \footnotesize
  \begin{tabularx}{1\textwidth}{|l|l|X|}
    \hline
    \textbf{Metadata Name} & \textbf{Metadata Value} & \textbf{Description} \\
    \hline
    \textbf{Reference Id} & \texttt{BEC-2A678BE0} & \\
    \hline
    \textbf{Status} & Generated candidate & \\
    \hline
    \textbf{Ancestry} &  & \\
    \hline
    \textbf{Source} & FHIR-informed candidate & \\
    \hline
    \textbf{Enabled By} &  & \\
    \hline
    \textbf{Supports} &  & \\
    \hline
    \textbf{Derived From} &  & \\
    \hline
    \textbf{Reviewed Status} & Generated candidate: keep / modify / relocate / reject; verify whether independently testable. & \\
    \hline
  \end{tabularx}
\end{table}

Intent. Subject Relationship Management is proposed to support the underlying healthcare concept represented by Remote Patient Monitoring. The proposal concerns behaviour and semantic responsibility, not a FHIR field, operation, resource instance or prescribed data model.

Conceptual evidence: HL7 FHIR R5 Patient, Person.

\noindent\textbf{Feature boundary.} This behaviour is scoped to Remote Patient Monitoring; an identically named Feature elsewhere may carry different business meaning.\par
\featuretable{
--- & --- & \emph{Not yet captured} & \emph{Complete during review} \\
}
\subsubsection{Feature - Subject Context Management}\label{cap:bec-0b235f68}
\vspace{-1em}
\begin{table}[htbp]
  \centering
  \footnotesize
  \begin{tabularx}{1\textwidth}{|l|l|X|}
    \hline
    \textbf{Metadata Name} & \textbf{Metadata Value} & \textbf{Description} \\
    \hline
    \textbf{Reference Id} & \texttt{BEC-0B235F68} & \\
    \hline
    \textbf{Status} & Generated candidate & \\
    \hline
    \textbf{Ancestry} &  & \\
    \hline
    \textbf{Source} & FHIR-informed candidate & \\
    \hline
    \textbf{Enabled By} &  & \\
    \hline
    \textbf{Supports} &  & \\
    \hline
    \textbf{Derived From} &  & \\
    \hline
    \textbf{Reviewed Status} & Generated candidate: keep / modify / relocate / reject; verify whether independently testable. & \\
    \hline
  \end{tabularx}
\end{table}

Intent. Subject Context Management is proposed to support the underlying healthcare concept represented by Remote Patient Monitoring. The proposal concerns behaviour and semantic responsibility, not a FHIR field, operation, resource instance or prescribed data model.

Conceptual evidence: HL7 FHIR R5 Patient, Person.

\noindent\textbf{Feature boundary.} This behaviour is scoped to Remote Patient Monitoring; an identically named Feature elsewhere may carry different business meaning.\par
\featuretable{
--- & --- & \emph{Not yet captured} & \emph{Complete during review} \\
}
\subsubsection{Feature - Subject Information Discovery}\label{cap:bec-6b015be0}
\vspace{-1em}
\begin{table}[htbp]
  \centering
  \footnotesize
  \begin{tabularx}{1\textwidth}{|l|l|X|}
    \hline
    \textbf{Metadata Name} & \textbf{Metadata Value} & \textbf{Description} \\
    \hline
    \textbf{Reference Id} & \texttt{BEC-6B015BE0} & \\
    \hline
    \textbf{Status} & Generated candidate & \\
    \hline
    \textbf{Ancestry} &  & \\
    \hline
    \textbf{Source} & FHIR-informed candidate & \\
    \hline
    \textbf{Enabled By} &  & \\
    \hline
    \textbf{Supports} &  & \\
    \hline
    \textbf{Derived From} &  & \\
    \hline
    \textbf{Reviewed Status} & Generated candidate: keep / modify / relocate / reject; verify whether independently testable. & \\
    \hline
  \end{tabularx}
\end{table}

Intent. Subject Information Discovery is proposed to support the underlying healthcare concept represented by Remote Patient Monitoring. The proposal concerns behaviour and semantic responsibility, not a FHIR field, operation, resource instance or prescribed data model.

Conceptual evidence: HL7 FHIR R5 Patient, Person.

\noindent\textbf{Feature boundary.} This behaviour is scoped to Remote Patient Monitoring; an identically named Feature elsewhere may carry different business meaning.\par
\featuretable{
--- & --- & \emph{Not yet captured} & \emph{Complete during review} \\
}

\chapter{CT1 - Health Service Operations Enablement}\label{cap:bec-4fa35177}
\noindent\textbf{Reference ID:} \texttt{BEC-4FA35177}\quad \textbf{Status:} Candidate ancestry\par
Intent. Health Service Operations Enablement is a candidate enabling capability within Business Enabling Capability layer. It frames the reusable system-supported behaviours needed to support this healthcare context without prescribing a product, interface, technology or deployment boundary.

Boundary. This capability supports the business activity and does not inherit the authority or accountability of the Business Capability it enables. Its child Features are review candidates until ancestry and semantics are adjudicated.

\section{CT2 - Clinic \& Practice Operations}\label{cap:bec-ed6d98e8}
\vspace{-0.5em}
\begin{table}[htbp]
  \centering
  \footnotesize
  \begin{tabularx}{1\textwidth}{|l|l|X|}
    \hline
    \multicolumn{3}{c}{\textbf{Capability Metadata}} \\
    \hline
    \textbf{Name} & \textbf{Value} & \textbf{Comment and/or Description} \\
    \hline
    \textbf{Reference Id} & \texttt{BEC-ED6D98E8} & \\
    \hline
    \textbf{Status} & Candidate ancestry & \\
    \hline
    \textbf{Ancestry} &  & \\
    \hline
    \textbf{Source} &  & \\
    \hline
    \textbf{Enabled By} &  & \\
    \hline
    \textbf{Supports} &  & \\
    \hline
    \textbf{Derived From} &  & \\
    \hline
    \textbf{Reviewed Status} &  & \\
    \hline
  \end{tabularx}
\end{table}

\begin{table}[htbp]
  \centering
  \footnotesize
  \begin{tabularx}{1\textwidth}{|l|l|X|}
    \hline
    \multicolumn{3}{c}{\textbf{Contained Capabilities}} \\
    \hline
    \textbf{Capability Name} & \textbf{Capability Id} & \textbf{Comment} \\
    \hline
    Session \& Queue Coordination & \texttt{cap:bec-fd9982fa} (at \ref{cap:bec-fd9982fa}) & \\
    \hline
  \end{tabularx}
\end{table}

Intent. Clinic \& Practice Operations is a candidate enabling capability within Health Service Operations Enablement. It frames the reusable system-supported behaviours needed to support this healthcare context without prescribing a product, interface, technology or deployment boundary.

Boundary. This capability supports the business activity and does not inherit the authority or accountability of the Business Capability it enables. Its child Features are review candidates until ancestry and semantics are adjudicated.

\subsection{CT3 - Session \& Queue Coordination}\label{cap:bec-fd9982fa}
\vspace{-0.5em}
\begin{table}[htbp]
  \centering
  \footnotesize
  \begin{tabularx}{1\textwidth}{|l|l|X|}
    \hline
    \multicolumn{3}{c}{\textbf{Capability Metadata}} \\
    \hline
    \textbf{Name} & \textbf{Value} & \textbf{Comment and/or Description} \\
    \hline
    \textbf{Reference Id} & \texttt{BEC-FD9982FA} & \\
    \hline
    \textbf{Status} & Candidate ancestry & \\
    \hline
    \textbf{Ancestry} &  & \\
    \hline
    \textbf{Source} &  & \\
    \hline
    \textbf{Enabled By} &  & \\
    \hline
    \textbf{Supports} &  & \\
    \hline
    \textbf{Derived From} &  & \\
    \hline
    \textbf{Reviewed Status} &  & \\
    \hline
  \end{tabularx}
\end{table}

\begin{table}[htbp]
  \centering
  \footnotesize
  \begin{tabularx}{1\textwidth}{|l|l|X|}
    \hline
    \multicolumn{3}{c}{\textbf{Contained Features}} \\
    \hline
    \textbf{Feature Name} & \textbf{Feature Id} & \textbf{Comment} \\
    \hline
    Clinic Session Tracking & \texttt{cap:bec-c271d00d} (at \ref{cap:bec-c271d00d}) & \\
    \hline
    Operational Queue Progression & \texttt{cap:bec-1d3dc25a} (at \ref{cap:bec-1d3dc25a}) & \\
    \hline
    Context Establishment & \texttt{cap:bec-f2dda45c} (at \ref{cap:bec-f2dda45c}) & \\
    \hline
    Eligibility Validation & \texttt{cap:bec-2e5f2477} (at \ref{cap:bec-2e5f2477}) & \\
    \hline
    Assignment \& Routing & \texttt{cap:bec-91e57a4c} (at \ref{cap:bec-91e57a4c}) & \\
    \hline
    State Tracking & \texttt{cap:bec-28a55b14} (at \ref{cap:bec-28a55b14}) & \\
    \hline
    Exception Management & \texttt{cap:bec-334d43b0} (at \ref{cap:bec-334d43b0}) & \\
    \hline
    Outcome Reporting & \texttt{cap:bec-4734d5e1} (at \ref{cap:bec-4734d5e1}) & \\
    \hline
  \end{tabularx}
\end{table}

Intent. Session \& Queue Coordination is a candidate enabling capability within Clinic \& Practice Operations. It frames the reusable system-supported behaviours needed to support this healthcare context without prescribing a product, interface, technology or deployment boundary.

Boundary. This capability supports the business activity and does not inherit the authority or accountability of the Business Capability it enables. Its child Features are review candidates until ancestry and semantics are adjudicated.

\subsubsection{Feature - Clinic Session Tracking}\label{cap:bec-c271d00d}
\vspace{-1em}
\begin{table}[htbp]
  \centering
  \footnotesize
  \begin{tabularx}{1\textwidth}{|l|l|X|}
    \hline
    \textbf{Metadata Name} & \textbf{Metadata Value} & \textbf{Description} \\
    \hline
    \textbf{Reference Id} & \texttt{BEC-C271D00D} & \\
    \hline
    \textbf{Status} & Existing Feature; candidate CT3 placement & \\
    \hline
    \textbf{Ancestry} &  & \\
    \hline
    \textbf{Source} & FEAT-HSO-01 & \\
    \hline
    \textbf{Enabled By} &  & \\
    \hline
    \textbf{Supports} &  & \\
    \hline
    \textbf{Derived From} &  & \\
    \hline
    \textbf{Reviewed Status} & Confirm existing Feature meaning and proposed CT3 ownership. & \\
    \hline
  \end{tabularx}
\end{table}

Intent. Clinic Session Tracking provides a distinct enabling behaviour in the context of Session \& Queue Coordination. This is an existing named Feature; its placement under the candidate CT3 remains subject to review.

\noindent\textbf{Feature boundary.} This behaviour is scoped to Session \& Queue Coordination; an identically named Feature elsewhere may carry different business meaning.\par
\featuretable{
Function & \texttt{FN-2004D2F5} & Track Clinic Operating Session & Catalogue \\
}
\subsubsection{Feature - Operational Queue Progression}\label{cap:bec-1d3dc25a}
\vspace{-1em}
\begin{table}[htbp]
  \centering
  \footnotesize
  \begin{tabularx}{1\textwidth}{|l|l|X|}
    \hline
    \textbf{Metadata Name} & \textbf{Metadata Value} & \textbf{Description} \\
    \hline
    \textbf{Reference Id} & \texttt{BEC-1D3DC25A} & \\
    \hline
    \textbf{Status} & Existing Feature; candidate CT3 placement & \\
    \hline
    \textbf{Ancestry} &  & \\
    \hline
    \textbf{Source} & FEAT-HSO-02 & \\
    \hline
    \textbf{Enabled By} &  & \\
    \hline
    \textbf{Supports} &  & \\
    \hline
    \textbf{Derived From} &  & \\
    \hline
    \textbf{Reviewed Status} & Confirm existing Feature meaning and proposed CT3 ownership. & \\
    \hline
  \end{tabularx}
\end{table}

Intent. Operational Queue Progression provides a distinct enabling behaviour in the context of Session \& Queue Coordination. This is an existing named Feature; its placement under the candidate CT3 remains subject to review.

\noindent\textbf{Feature boundary.} This behaviour is scoped to Session \& Queue Coordination; an identically named Feature elsewhere may carry different business meaning.\par
\featuretable{
Function & \texttt{FN-5228BB72} & Manage Clinic Queue Progression & Catalogue \\
Service & \texttt{SV-B60FBE79} & Clinic Queue Telemetry Service & Linked to preceding Function in source \\
Process & \texttt{PR-CCF51851} & Clinic / Operational Queue Progression Process & Linked to preceding Function in source \\
}
\subsubsection{Feature - Context Establishment}\label{cap:bec-f2dda45c}
\vspace{-1em}
\begin{table}[htbp]
  \centering
  \footnotesize
  \begin{tabularx}{1\textwidth}{|l|l|X|}
    \hline
    \textbf{Metadata Name} & \textbf{Metadata Value} & \textbf{Description} \\
    \hline
    \textbf{Reference Id} & \texttt{BEC-F2DDA45C} & \\
    \hline
    \textbf{Status} & Generated candidate & \\
    \hline
    \textbf{Ancestry} &  & \\
    \hline
    \textbf{Source} & Semantic-pattern candidate & \\
    \hline
    \textbf{Enabled By} &  & \\
    \hline
    \textbf{Supports} &  & \\
    \hline
    \textbf{Derived From} &  & \\
    \hline
    \textbf{Reviewed Status} & Generated candidate: keep / modify / relocate / reject; verify whether independently testable. & \\
    \hline
  \end{tabularx}
\end{table}

Intent. Context Establishment is proposed as a discrete, testable enabling behaviour for Session \& Queue Coordination. It is included for architectural review, not as an established implementation obligation.

\noindent\textbf{Feature boundary.} This behaviour is scoped to Session \& Queue Coordination; an identically named Feature elsewhere may carry different business meaning.\par
\featuretable{
--- & --- & \emph{Not yet captured} & \emph{Complete during review} \\
}
\subsubsection{Feature - Eligibility Validation}\label{cap:bec-2e5f2477}
\vspace{-1em}
\begin{table}[htbp]
  \centering
  \footnotesize
  \begin{tabularx}{1\textwidth}{|l|l|X|}
    \hline
    \textbf{Metadata Name} & \textbf{Metadata Value} & \textbf{Description} \\
    \hline
    \textbf{Reference Id} & \texttt{BEC-2E5F2477} & \\
    \hline
    \textbf{Status} & Generated candidate & \\
    \hline
    \textbf{Ancestry} &  & \\
    \hline
    \textbf{Source} & Semantic-pattern candidate & \\
    \hline
    \textbf{Enabled By} &  & \\
    \hline
    \textbf{Supports} &  & \\
    \hline
    \textbf{Derived From} &  & \\
    \hline
    \textbf{Reviewed Status} & Generated candidate: keep / modify / relocate / reject; verify whether independently testable. & \\
    \hline
  \end{tabularx}
\end{table}

Intent. Eligibility Validation is proposed as a discrete, testable enabling behaviour for Session \& Queue Coordination. It is included for architectural review, not as an established implementation obligation.

\noindent\textbf{Feature boundary.} This behaviour is scoped to Session \& Queue Coordination; an identically named Feature elsewhere may carry different business meaning.\par
\featuretable{
--- & --- & \emph{Not yet captured} & \emph{Complete during review} \\
}
\subsubsection{Feature - Assignment \& Routing}\label{cap:bec-91e57a4c}
\vspace{-1em}
\begin{table}[htbp]
  \centering
  \footnotesize
  \begin{tabularx}{1\textwidth}{|l|l|X|}
    \hline
    \textbf{Metadata Name} & \textbf{Metadata Value} & \textbf{Description} \\
    \hline
    \textbf{Reference Id} & \texttt{BEC-91E57A4C} & \\
    \hline
    \textbf{Status} & Generated candidate & \\
    \hline
    \textbf{Ancestry} &  & \\
    \hline
    \textbf{Source} & Semantic-pattern candidate & \\
    \hline
    \textbf{Enabled By} &  & \\
    \hline
    \textbf{Supports} &  & \\
    \hline
    \textbf{Derived From} &  & \\
    \hline
    \textbf{Reviewed Status} & Generated candidate: keep / modify / relocate / reject; verify whether independently testable. & \\
    \hline
  \end{tabularx}
\end{table}

Intent. Assignment \& Routing is proposed as a discrete, testable enabling behaviour for Session \& Queue Coordination. It is included for architectural review, not as an established implementation obligation.

\noindent\textbf{Feature boundary.} This behaviour is scoped to Session \& Queue Coordination; an identically named Feature elsewhere may carry different business meaning.\par
\featuretable{
--- & --- & \emph{Not yet captured} & \emph{Complete during review} \\
}
\subsubsection{Feature - State Tracking}\label{cap:bec-28a55b14}
\vspace{-1em}
\begin{table}[htbp]
  \centering
  \footnotesize
  \begin{tabularx}{1\textwidth}{|l|l|X|}
    \hline
    \textbf{Metadata Name} & \textbf{Metadata Value} & \textbf{Description} \\
    \hline
    \textbf{Reference Id} & \texttt{BEC-28A55B14} & \\
    \hline
    \textbf{Status} & Generated candidate & \\
    \hline
    \textbf{Ancestry} &  & \\
    \hline
    \textbf{Source} & Semantic-pattern candidate & \\
    \hline
    \textbf{Enabled By} &  & \\
    \hline
    \textbf{Supports} &  & \\
    \hline
    \textbf{Derived From} &  & \\
    \hline
    \textbf{Reviewed Status} & Generated candidate: keep / modify / relocate / reject; verify whether independently testable. & \\
    \hline
  \end{tabularx}
\end{table}

Intent. State Tracking is proposed as a discrete, testable enabling behaviour for Session \& Queue Coordination. It is included for architectural review, not as an established implementation obligation.

\noindent\textbf{Feature boundary.} This behaviour is scoped to Session \& Queue Coordination; an identically named Feature elsewhere may carry different business meaning.\par
\featuretable{
--- & --- & \emph{Not yet captured} & \emph{Complete during review} \\
}
\subsubsection{Feature - Exception Management}\label{cap:bec-334d43b0}
\vspace{-1em}
\begin{table}[htbp]
  \centering
  \footnotesize
  \begin{tabularx}{1\textwidth}{|l|l|X|}
    \hline
    \textbf{Metadata Name} & \textbf{Metadata Value} & \textbf{Description} \\
    \hline
    \textbf{Reference Id} & \texttt{BEC-334D43B0} & \\
    \hline
    \textbf{Status} & Generated candidate & \\
    \hline
    \textbf{Ancestry} &  & \\
    \hline
    \textbf{Source} & Semantic-pattern candidate & \\
    \hline
    \textbf{Enabled By} &  & \\
    \hline
    \textbf{Supports} &  & \\
    \hline
    \textbf{Derived From} &  & \\
    \hline
    \textbf{Reviewed Status} & Generated candidate: keep / modify / relocate / reject; verify whether independently testable. & \\
    \hline
  \end{tabularx}
\end{table}

Intent. Exception Management is proposed as a discrete, testable enabling behaviour for Session \& Queue Coordination. It is included for architectural review, not as an established implementation obligation.

\noindent\textbf{Feature boundary.} This behaviour is scoped to Session \& Queue Coordination; an identically named Feature elsewhere may carry different business meaning.\par
\featuretable{
--- & --- & \emph{Not yet captured} & \emph{Complete during review} \\
}
\subsubsection{Feature - Outcome Reporting}\label{cap:bec-4734d5e1}
\vspace{-1em}
\begin{table}[htbp]
  \centering
  \footnotesize
  \begin{tabularx}{1\textwidth}{|l|l|X|}
    \hline
    \textbf{Metadata Name} & \textbf{Metadata Value} & \textbf{Description} \\
    \hline
    \textbf{Reference Id} & \texttt{BEC-4734D5E1} & \\
    \hline
    \textbf{Status} & Generated candidate & \\
    \hline
    \textbf{Ancestry} &  & \\
    \hline
    \textbf{Source} & Semantic-pattern candidate & \\
    \hline
    \textbf{Enabled By} &  & \\
    \hline
    \textbf{Supports} &  & \\
    \hline
    \textbf{Derived From} &  & \\
    \hline
    \textbf{Reviewed Status} & Generated candidate: keep / modify / relocate / reject; verify whether independently testable. & \\
    \hline
  \end{tabularx}
\end{table}

Intent. Outcome Reporting is proposed as a discrete, testable enabling behaviour for Session \& Queue Coordination. It is included for architectural review, not as an established implementation obligation.

\noindent\textbf{Feature boundary.} This behaviour is scoped to Session \& Queue Coordination; an identically named Feature elsewhere may carry different business meaning.\par
\featuretable{
--- & --- & \emph{Not yet captured} & \emph{Complete during review} \\
}

\section{CT2 - Ward Operations}\label{cap:bec-25e1a896}
\vspace{-0.5em}
\begin{table}[htbp]
  \centering
  \footnotesize
  \begin{tabularx}{1\textwidth}{|l|l|X|}
    \hline
    \multicolumn{3}{c}{\textbf{Capability Metadata}} \\
    \hline
    \textbf{Name} & \textbf{Value} & \textbf{Comment and/or Description} \\
    \hline
    \textbf{Reference Id} & \texttt{BEC-25E1A896} & \\
    \hline
    \textbf{Status} & Candidate ancestry & \\
    \hline
    \textbf{Ancestry} &  & \\
    \hline
    \textbf{Source} &  & \\
    \hline
    \textbf{Enabled By} &  & \\
    \hline
    \textbf{Supports} &  & \\
    \hline
    \textbf{Derived From} &  & \\
    \hline
    \textbf{Reviewed Status} &  & \\
    \hline
  \end{tabularx}
\end{table}

\begin{table}[htbp]
  \centering
  \footnotesize
  \begin{tabularx}{1\textwidth}{|l|l|X|}
    \hline
    \multicolumn{3}{c}{\textbf{Contained Capabilities}} \\
    \hline
    \textbf{Capability Name} & \textbf{Capability Id} & \textbf{Comment} \\
    \hline
    Ward Logistics Coordination & \texttt{cap:bec-ce9bb804} (at \ref{cap:bec-ce9bb804}) & \\
    \hline
  \end{tabularx}
\end{table}

Intent. Ward Operations is a candidate enabling capability within Health Service Operations Enablement. It frames the reusable system-supported behaviours needed to support this healthcare context without prescribing a product, interface, technology or deployment boundary.

Boundary. This capability supports the business activity and does not inherit the authority or accountability of the Business Capability it enables. Its child Features are review candidates until ancestry and semantics are adjudicated.

\subsection{CT3 - Ward Logistics Coordination}\label{cap:bec-ce9bb804}
\vspace{-0.5em}
\begin{table}[htbp]
  \centering
  \footnotesize
  \begin{tabularx}{1\textwidth}{|l|l|X|}
    \hline
    \multicolumn{3}{c}{\textbf{Capability Metadata}} \\
    \hline
    \textbf{Name} & \textbf{Value} & \textbf{Comment and/or Description} \\
    \hline
    \textbf{Reference Id} & \texttt{BEC-CE9BB804} & \\
    \hline
    \textbf{Status} & Candidate ancestry & \\
    \hline
    \textbf{Ancestry} &  & \\
    \hline
    \textbf{Source} &  & \\
    \hline
    \textbf{Enabled By} &  & \\
    \hline
    \textbf{Supports} &  & \\
    \hline
    \textbf{Derived From} &  & \\
    \hline
    \textbf{Reviewed Status} &  & \\
    \hline
  \end{tabularx}
\end{table}

\begin{table}[htbp]
  \centering
  \footnotesize
  \begin{tabularx}{1\textwidth}{|l|l|X|}
    \hline
    \multicolumn{3}{c}{\textbf{Contained Features}} \\
    \hline
    \textbf{Feature Name} & \textbf{Feature Id} & \textbf{Comment} \\
    \hline
    Ward Patient Flow Coordination & \texttt{cap:bec-9408fffa} (at \ref{cap:bec-9408fffa}) & \\
    \hline
    Clinical Handover Context Assembly & \texttt{cap:bec-06c10c04} (at \ref{cap:bec-06c10c04}) & \\
    \hline
    Context Establishment & \texttt{cap:bec-0029bf2b} (at \ref{cap:bec-0029bf2b}) & \\
    \hline
    Eligibility Validation & \texttt{cap:bec-616ffdc2} (at \ref{cap:bec-616ffdc2}) & \\
    \hline
    Assignment \& Routing & \texttt{cap:bec-0da17c6a} (at \ref{cap:bec-0da17c6a}) & \\
    \hline
    State Tracking & \texttt{cap:bec-e6e33d5a} (at \ref{cap:bec-e6e33d5a}) & \\
    \hline
    Exception Management & \texttt{cap:bec-e8384db4} (at \ref{cap:bec-e8384db4}) & \\
    \hline
    Outcome Reporting & \texttt{cap:bec-4864c81b} (at \ref{cap:bec-4864c81b}) & \\
    \hline
  \end{tabularx}
\end{table}

Intent. Ward Logistics Coordination is a candidate enabling capability within Ward Operations. It frames the reusable system-supported behaviours needed to support this healthcare context without prescribing a product, interface, technology or deployment boundary.

Boundary. This capability supports the business activity and does not inherit the authority or accountability of the Business Capability it enables. Its child Features are review candidates until ancestry and semantics are adjudicated.

\subsubsection{Feature - Ward Patient Flow Coordination}\label{cap:bec-9408fffa}
\vspace{-1em}
\begin{table}[htbp]
  \centering
  \footnotesize
  \begin{tabularx}{1\textwidth}{|l|l|X|}
    \hline
    \textbf{Metadata Name} & \textbf{Metadata Value} & \textbf{Description} \\
    \hline
    \textbf{Reference Id} & \texttt{BEC-9408FFFA} & \\
    \hline
    \textbf{Status} & Existing Feature; candidate CT3 placement & \\
    \hline
    \textbf{Ancestry} &  & \\
    \hline
    \textbf{Source} & FEAT-HSO-03 & \\
    \hline
    \textbf{Enabled By} &  & \\
    \hline
    \textbf{Supports} &  & \\
    \hline
    \textbf{Derived From} &  & \\
    \hline
    \textbf{Reviewed Status} & Confirm existing Feature meaning and proposed CT3 ownership. & \\
    \hline
  \end{tabularx}
\end{table}

Intent. Ward Patient Flow Coordination provides a distinct enabling behaviour in the context of Ward Logistics Coordination. This is an existing named Feature; its placement under the candidate CT3 remains subject to review.

\noindent\textbf{Feature boundary.} This behaviour is scoped to Ward Logistics Coordination; an identically named Feature elsewhere may carry different business meaning.\par
\featuretable{
Function & \texttt{FN-D03C9CBA} & Association unresolved / not established & Catalogue \\
}
\subsubsection{Feature - Clinical Handover Context Assembly}\label{cap:bec-06c10c04}
\vspace{-1em}
\begin{table}[htbp]
  \centering
  \footnotesize
  \begin{tabularx}{1\textwidth}{|l|l|X|}
    \hline
    \textbf{Metadata Name} & \textbf{Metadata Value} & \textbf{Description} \\
    \hline
    \textbf{Reference Id} & \texttt{BEC-06C10C04} & \\
    \hline
    \textbf{Status} & Existing Feature; candidate CT3 placement & \\
    \hline
    \textbf{Ancestry} &  & \\
    \hline
    \textbf{Source} & FEAT-HSO-04 & \\
    \hline
    \textbf{Enabled By} &  & \\
    \hline
    \textbf{Supports} &  & \\
    \hline
    \textbf{Derived From} &  & \\
    \hline
    \textbf{Reviewed Status} & Confirm existing Feature meaning and proposed CT3 ownership. & \\
    \hline
  \end{tabularx}
\end{table}

Intent. Clinical Handover Context Assembly provides a distinct enabling behaviour in the context of Ward Logistics Coordination. This is an existing named Feature; its placement under the candidate CT3 remains subject to review.

\noindent\textbf{Feature boundary.} This behaviour is scoped to Ward Logistics Coordination; an identically named Feature elsewhere may carry different business meaning.\par
\featuretable{
Function & \texttt{FN-4DFCEC0F} & Assemble Ward Handover Context & Catalogue \\
}
\subsubsection{Feature - Context Establishment}\label{cap:bec-0029bf2b}
\vspace{-1em}
\begin{table}[htbp]
  \centering
  \footnotesize
  \begin{tabularx}{1\textwidth}{|l|l|X|}
    \hline
    \textbf{Metadata Name} & \textbf{Metadata Value} & \textbf{Description} \\
    \hline
    \textbf{Reference Id} & \texttt{BEC-0029BF2B} & \\
    \hline
    \textbf{Status} & Generated candidate & \\
    \hline
    \textbf{Ancestry} &  & \\
    \hline
    \textbf{Source} & Semantic-pattern candidate & \\
    \hline
    \textbf{Enabled By} &  & \\
    \hline
    \textbf{Supports} &  & \\
    \hline
    \textbf{Derived From} &  & \\
    \hline
    \textbf{Reviewed Status} & Generated candidate: keep / modify / relocate / reject; verify whether independently testable. & \\
    \hline
  \end{tabularx}
\end{table}

Intent. Context Establishment is proposed as a discrete, testable enabling behaviour for Ward Logistics Coordination. It is included for architectural review, not as an established implementation obligation.

\noindent\textbf{Feature boundary.} This behaviour is scoped to Ward Logistics Coordination; an identically named Feature elsewhere may carry different business meaning.\par
\featuretable{
--- & --- & \emph{Not yet captured} & \emph{Complete during review} \\
}
\subsubsection{Feature - Eligibility Validation}\label{cap:bec-616ffdc2}
\vspace{-1em}
\begin{table}[htbp]
  \centering
  \footnotesize
  \begin{tabularx}{1\textwidth}{|l|l|X|}
    \hline
    \textbf{Metadata Name} & \textbf{Metadata Value} & \textbf{Description} \\
    \hline
    \textbf{Reference Id} & \texttt{BEC-616FFDC2} & \\
    \hline
    \textbf{Status} & Generated candidate & \\
    \hline
    \textbf{Ancestry} &  & \\
    \hline
    \textbf{Source} & Semantic-pattern candidate & \\
    \hline
    \textbf{Enabled By} &  & \\
    \hline
    \textbf{Supports} &  & \\
    \hline
    \textbf{Derived From} &  & \\
    \hline
    \textbf{Reviewed Status} & Generated candidate: keep / modify / relocate / reject; verify whether independently testable. & \\
    \hline
  \end{tabularx}
\end{table}

Intent. Eligibility Validation is proposed as a discrete, testable enabling behaviour for Ward Logistics Coordination. It is included for architectural review, not as an established implementation obligation.

\noindent\textbf{Feature boundary.} This behaviour is scoped to Ward Logistics Coordination; an identically named Feature elsewhere may carry different business meaning.\par
\featuretable{
--- & --- & \emph{Not yet captured} & \emph{Complete during review} \\
}
\subsubsection{Feature - Assignment \& Routing}\label{cap:bec-0da17c6a}
\vspace{-1em}
\begin{table}[htbp]
  \centering
  \footnotesize
  \begin{tabularx}{1\textwidth}{|l|l|X|}
    \hline
    \textbf{Metadata Name} & \textbf{Metadata Value} & \textbf{Description} \\
    \hline
    \textbf{Reference Id} & \texttt{BEC-0DA17C6A} & \\
    \hline
    \textbf{Status} & Generated candidate & \\
    \hline
    \textbf{Ancestry} &  & \\
    \hline
    \textbf{Source} & Semantic-pattern candidate & \\
    \hline
    \textbf{Enabled By} &  & \\
    \hline
    \textbf{Supports} &  & \\
    \hline
    \textbf{Derived From} &  & \\
    \hline
    \textbf{Reviewed Status} & Generated candidate: keep / modify / relocate / reject; verify whether independently testable. & \\
    \hline
  \end{tabularx}
\end{table}

Intent. Assignment \& Routing is proposed as a discrete, testable enabling behaviour for Ward Logistics Coordination. It is included for architectural review, not as an established implementation obligation.

\noindent\textbf{Feature boundary.} This behaviour is scoped to Ward Logistics Coordination; an identically named Feature elsewhere may carry different business meaning.\par
\featuretable{
--- & --- & \emph{Not yet captured} & \emph{Complete during review} \\
}
\subsubsection{Feature - State Tracking}\label{cap:bec-e6e33d5a}
\vspace{-1em}
\begin{table}[htbp]
  \centering
  \footnotesize
  \begin{tabularx}{1\textwidth}{|l|l|X|}
    \hline
    \textbf{Metadata Name} & \textbf{Metadata Value} & \textbf{Description} \\
    \hline
    \textbf{Reference Id} & \texttt{BEC-E6E33D5A} & \\
    \hline
    \textbf{Status} & Generated candidate & \\
    \hline
    \textbf{Ancestry} &  & \\
    \hline
    \textbf{Source} & Semantic-pattern candidate & \\
    \hline
    \textbf{Enabled By} &  & \\
    \hline
    \textbf{Supports} &  & \\
    \hline
    \textbf{Derived From} &  & \\
    \hline
    \textbf{Reviewed Status} & Generated candidate: keep / modify / relocate / reject; verify whether independently testable. & \\
    \hline
  \end{tabularx}
\end{table}

Intent. State Tracking is proposed as a discrete, testable enabling behaviour for Ward Logistics Coordination. It is included for architectural review, not as an established implementation obligation.

\noindent\textbf{Feature boundary.} This behaviour is scoped to Ward Logistics Coordination; an identically named Feature elsewhere may carry different business meaning.\par
\featuretable{
--- & --- & \emph{Not yet captured} & \emph{Complete during review} \\
}
\subsubsection{Feature - Exception Management}\label{cap:bec-e8384db4}
\vspace{-1em}
\begin{table}[htbp]
  \centering
  \footnotesize
  \begin{tabularx}{1\textwidth}{|l|l|X|}
    \hline
    \textbf{Metadata Name} & \textbf{Metadata Value} & \textbf{Description} \\
    \hline
    \textbf{Reference Id} & \texttt{BEC-E8384DB4} & \\
    \hline
    \textbf{Status} & Generated candidate & \\
    \hline
    \textbf{Ancestry} &  & \\
    \hline
    \textbf{Source} & Semantic-pattern candidate & \\
    \hline
    \textbf{Enabled By} &  & \\
    \hline
    \textbf{Supports} &  & \\
    \hline
    \textbf{Derived From} &  & \\
    \hline
    \textbf{Reviewed Status} & Generated candidate: keep / modify / relocate / reject; verify whether independently testable. & \\
    \hline
  \end{tabularx}
\end{table}

Intent. Exception Management is proposed as a discrete, testable enabling behaviour for Ward Logistics Coordination. It is included for architectural review, not as an established implementation obligation.

\noindent\textbf{Feature boundary.} This behaviour is scoped to Ward Logistics Coordination; an identically named Feature elsewhere may carry different business meaning.\par
\featuretable{
--- & --- & \emph{Not yet captured} & \emph{Complete during review} \\
}
\subsubsection{Feature - Outcome Reporting}\label{cap:bec-4864c81b}
\vspace{-1em}
\begin{table}[htbp]
  \centering
  \footnotesize
  \begin{tabularx}{1\textwidth}{|l|l|X|}
    \hline
    \textbf{Metadata Name} & \textbf{Metadata Value} & \textbf{Description} \\
    \hline
    \textbf{Reference Id} & \texttt{BEC-4864C81B} & \\
    \hline
    \textbf{Status} & Generated candidate & \\
    \hline
    \textbf{Ancestry} &  & \\
    \hline
    \textbf{Source} & Semantic-pattern candidate & \\
    \hline
    \textbf{Enabled By} &  & \\
    \hline
    \textbf{Supports} &  & \\
    \hline
    \textbf{Derived From} &  & \\
    \hline
    \textbf{Reviewed Status} & Generated candidate: keep / modify / relocate / reject; verify whether independently testable. & \\
    \hline
  \end{tabularx}
\end{table}

Intent. Outcome Reporting is proposed as a discrete, testable enabling behaviour for Ward Logistics Coordination. It is included for architectural review, not as an established implementation obligation.

\noindent\textbf{Feature boundary.} This behaviour is scoped to Ward Logistics Coordination; an identically named Feature elsewhere may carry different business meaning.\par
\featuretable{
--- & --- & \emph{Not yet captured} & \emph{Complete during review} \\
}

\section{CT2 - Theatre Operations}\label{cap:bec-389897e6}
\vspace{-0.5em}
\begin{table}[htbp]
  \centering
  \footnotesize
  \begin{tabularx}{1\textwidth}{|l|l|X|}
    \hline
    \multicolumn{3}{c}{\textbf{Capability Metadata}} \\
    \hline
    \textbf{Name} & \textbf{Value} & \textbf{Comment and/or Description} \\
    \hline
    \textbf{Reference Id} & \texttt{BEC-389897E6} & \\
    \hline
    \textbf{Status} & Candidate ancestry & \\
    \hline
    \textbf{Ancestry} &  & \\
    \hline
    \textbf{Source} &  & \\
    \hline
    \textbf{Enabled By} &  & \\
    \hline
    \textbf{Supports} &  & \\
    \hline
    \textbf{Derived From} &  & \\
    \hline
    \textbf{Reviewed Status} &  & \\
    \hline
  \end{tabularx}
\end{table}

\begin{table}[htbp]
  \centering
  \footnotesize
  \begin{tabularx}{1\textwidth}{|l|l|X|}
    \hline
    \multicolumn{3}{c}{\textbf{Contained Capabilities}} \\
    \hline
    \textbf{Capability Name} & \textbf{Capability Id} & \textbf{Comment} \\
    \hline
    Perioperative Logistics & \texttt{cap:bec-d248d868} (at \ref{cap:bec-d248d868}) & \\
    \hline
  \end{tabularx}
\end{table}

Intent. Theatre Operations is a candidate enabling capability within Health Service Operations Enablement. It frames the reusable system-supported behaviours needed to support this healthcare context without prescribing a product, interface, technology or deployment boundary.

Boundary. This capability supports the business activity and does not inherit the authority or accountability of the Business Capability it enables. Its child Features are review candidates until ancestry and semantics are adjudicated.

\subsection{CT3 - Perioperative Logistics}\label{cap:bec-d248d868}
\vspace{-0.5em}
\begin{table}[htbp]
  \centering
  \footnotesize
  \begin{tabularx}{1\textwidth}{|l|l|X|}
    \hline
    \multicolumn{3}{c}{\textbf{Capability Metadata}} \\
    \hline
    \textbf{Name} & \textbf{Value} & \textbf{Comment and/or Description} \\
    \hline
    \textbf{Reference Id} & \texttt{BEC-D248D868} & \\
    \hline
    \textbf{Status} & Candidate ancestry & \\
    \hline
    \textbf{Ancestry} &  & \\
    \hline
    \textbf{Source} &  & \\
    \hline
    \textbf{Enabled By} &  & \\
    \hline
    \textbf{Supports} &  & \\
    \hline
    \textbf{Derived From} &  & \\
    \hline
    \textbf{Reviewed Status} &  & \\
    \hline
  \end{tabularx}
\end{table}

\begin{table}[htbp]
  \centering
  \footnotesize
  \begin{tabularx}{1\textwidth}{|l|l|X|}
    \hline
    \multicolumn{3}{c}{\textbf{Contained Features}} \\
    \hline
    \textbf{Feature Name} & \textbf{Feature Id} & \textbf{Comment} \\
    \hline
    Operating Theatre Case Progression & \texttt{cap:bec-98713abb} (at \ref{cap:bec-98713abb}) & \\
    \hline
    Perioperative Resource Notification & \texttt{cap:bec-b1e9630d} (at \ref{cap:bec-b1e9630d}) & \\
    \hline
    Context Establishment & \texttt{cap:bec-63da723b} (at \ref{cap:bec-63da723b}) & \\
    \hline
    State Monitoring & \texttt{cap:bec-9f573a4c} (at \ref{cap:bec-9f573a4c}) & \\
    \hline
    Exception Management & \texttt{cap:bec-1233b732} (at \ref{cap:bec-1233b732}) & \\
    \hline
    Information Retrieval & \texttt{cap:bec-792b9999} (at \ref{cap:bec-792b9999}) & \\
    \hline
  \end{tabularx}
\end{table}

Intent. Perioperative Logistics is a candidate enabling capability within Theatre Operations. It frames the reusable system-supported behaviours needed to support this healthcare context without prescribing a product, interface, technology or deployment boundary.

Boundary. This capability supports the business activity and does not inherit the authority or accountability of the Business Capability it enables. Its child Features are review candidates until ancestry and semantics are adjudicated.

\subsubsection{Feature - Operating Theatre Case Progression}\label{cap:bec-98713abb}
\vspace{-1em}
\begin{table}[htbp]
  \centering
  \footnotesize
  \begin{tabularx}{1\textwidth}{|l|l|X|}
    \hline
    \textbf{Metadata Name} & \textbf{Metadata Value} & \textbf{Description} \\
    \hline
    \textbf{Reference Id} & \texttt{BEC-98713ABB} & \\
    \hline
    \textbf{Status} & Existing Feature; candidate CT3 placement & \\
    \hline
    \textbf{Ancestry} &  & \\
    \hline
    \textbf{Source} & FEAT-HSO-05 & \\
    \hline
    \textbf{Enabled By} &  & \\
    \hline
    \textbf{Supports} &  & \\
    \hline
    \textbf{Derived From} &  & \\
    \hline
    \textbf{Reviewed Status} & Confirm existing Feature meaning and proposed CT3 ownership. & \\
    \hline
  \end{tabularx}
\end{table}

Intent. Operating Theatre Case Progression provides a distinct enabling behaviour in the context of Perioperative Logistics. This is an existing named Feature; its placement under the candidate CT3 remains subject to review.

\noindent\textbf{Feature boundary.} This behaviour is scoped to Perioperative Logistics; an identically named Feature elsewhere may carry different business meaning.\par
\featuretable{
Function & \texttt{FN-745DEF16} & Manage Theatre Case Progression & Catalogue \\
Service & \texttt{SV-A68B48AA} & Theatre Case Progression Telemetry & Linked to preceding Function in source \\
Process & \texttt{PR-C114B9F9} & Theatre Case Progression Process & Linked to preceding Function in source \\
}
\subsubsection{Feature - Perioperative Resource Notification}\label{cap:bec-b1e9630d}
\vspace{-1em}
\begin{table}[htbp]
  \centering
  \footnotesize
  \begin{tabularx}{1\textwidth}{|l|l|X|}
    \hline
    \textbf{Metadata Name} & \textbf{Metadata Value} & \textbf{Description} \\
    \hline
    \textbf{Reference Id} & \texttt{BEC-B1E9630D} & \\
    \hline
    \textbf{Status} & Existing Feature; candidate CT3 placement & \\
    \hline
    \textbf{Ancestry} &  & \\
    \hline
    \textbf{Source} & FEAT-HSO-06 & \\
    \hline
    \textbf{Enabled By} &  & \\
    \hline
    \textbf{Supports} &  & \\
    \hline
    \textbf{Derived From} &  & \\
    \hline
    \textbf{Reviewed Status} & Confirm existing Feature meaning and proposed CT3 ownership. & \\
    \hline
  \end{tabularx}
\end{table}

Intent. Perioperative Resource Notification provides a distinct enabling behaviour in the context of Perioperative Logistics. This is an existing named Feature; its placement under the candidate CT3 remains subject to review.

\noindent\textbf{Feature boundary.} This behaviour is scoped to Perioperative Logistics; an identically named Feature elsewhere may carry different business meaning.\par
\featuretable{
Function & \texttt{FN-4873088C} & Manage Theatre Case Progression & Catalogue \\
Process & \texttt{PR-BC593DC8} & Theatre Case Progression Process & Linked to preceding Function in source \\
}
\subsubsection{Feature - Context Establishment}\label{cap:bec-63da723b}
\vspace{-1em}
\begin{table}[htbp]
  \centering
  \footnotesize
  \begin{tabularx}{1\textwidth}{|l|l|X|}
    \hline
    \textbf{Metadata Name} & \textbf{Metadata Value} & \textbf{Description} \\
    \hline
    \textbf{Reference Id} & \texttt{BEC-63DA723B} & \\
    \hline
    \textbf{Status} & Generated candidate & \\
    \hline
    \textbf{Ancestry} &  & \\
    \hline
    \textbf{Source} & Semantic-pattern candidate & \\
    \hline
    \textbf{Enabled By} &  & \\
    \hline
    \textbf{Supports} &  & \\
    \hline
    \textbf{Derived From} &  & \\
    \hline
    \textbf{Reviewed Status} & Generated candidate: keep / modify / relocate / reject; verify whether independently testable. & \\
    \hline
  \end{tabularx}
\end{table}

Intent. Context Establishment is proposed as a discrete, testable enabling behaviour for Perioperative Logistics. It is included for architectural review, not as an established implementation obligation.

\noindent\textbf{Feature boundary.} This behaviour is scoped to Perioperative Logistics; an identically named Feature elsewhere may carry different business meaning.\par
\featuretable{
--- & --- & \emph{Not yet captured} & \emph{Complete during review} \\
}
\subsubsection{Feature - State Monitoring}\label{cap:bec-9f573a4c}
\vspace{-1em}
\begin{table}[htbp]
  \centering
  \footnotesize
  \begin{tabularx}{1\textwidth}{|l|l|X|}
    \hline
    \textbf{Metadata Name} & \textbf{Metadata Value} & \textbf{Description} \\
    \hline
    \textbf{Reference Id} & \texttt{BEC-9F573A4C} & \\
    \hline
    \textbf{Status} & Generated candidate & \\
    \hline
    \textbf{Ancestry} &  & \\
    \hline
    \textbf{Source} & Semantic-pattern candidate & \\
    \hline
    \textbf{Enabled By} &  & \\
    \hline
    \textbf{Supports} &  & \\
    \hline
    \textbf{Derived From} &  & \\
    \hline
    \textbf{Reviewed Status} & Generated candidate: keep / modify / relocate / reject; verify whether independently testable. & \\
    \hline
  \end{tabularx}
\end{table}

Intent. State Monitoring is proposed as a discrete, testable enabling behaviour for Perioperative Logistics. It is included for architectural review, not as an established implementation obligation.

\noindent\textbf{Feature boundary.} This behaviour is scoped to Perioperative Logistics; an identically named Feature elsewhere may carry different business meaning.\par
\featuretable{
--- & --- & \emph{Not yet captured} & \emph{Complete during review} \\
}
\subsubsection{Feature - Exception Management}\label{cap:bec-1233b732}
\vspace{-1em}
\begin{table}[htbp]
  \centering
  \footnotesize
  \begin{tabularx}{1\textwidth}{|l|l|X|}
    \hline
    \textbf{Metadata Name} & \textbf{Metadata Value} & \textbf{Description} \\
    \hline
    \textbf{Reference Id} & \texttt{BEC-1233B732} & \\
    \hline
    \textbf{Status} & Generated candidate & \\
    \hline
    \textbf{Ancestry} &  & \\
    \hline
    \textbf{Source} & Semantic-pattern candidate & \\
    \hline
    \textbf{Enabled By} &  & \\
    \hline
    \textbf{Supports} &  & \\
    \hline
    \textbf{Derived From} &  & \\
    \hline
    \textbf{Reviewed Status} & Generated candidate: keep / modify / relocate / reject; verify whether independently testable. & \\
    \hline
  \end{tabularx}
\end{table}

Intent. Exception Management is proposed as a discrete, testable enabling behaviour for Perioperative Logistics. It is included for architectural review, not as an established implementation obligation.

\noindent\textbf{Feature boundary.} This behaviour is scoped to Perioperative Logistics; an identically named Feature elsewhere may carry different business meaning.\par
\featuretable{
--- & --- & \emph{Not yet captured} & \emph{Complete during review} \\
}
\subsubsection{Feature - Information Retrieval}\label{cap:bec-792b9999}
\vspace{-1em}
\begin{table}[htbp]
  \centering
  \footnotesize
  \begin{tabularx}{1\textwidth}{|l|l|X|}
    \hline
    \textbf{Metadata Name} & \textbf{Metadata Value} & \textbf{Description} \\
    \hline
    \textbf{Reference Id} & \texttt{BEC-792B9999} & \\
    \hline
    \textbf{Status} & Generated candidate & \\
    \hline
    \textbf{Ancestry} &  & \\
    \hline
    \textbf{Source} & Semantic-pattern candidate & \\
    \hline
    \textbf{Enabled By} &  & \\
    \hline
    \textbf{Supports} &  & \\
    \hline
    \textbf{Derived From} &  & \\
    \hline
    \textbf{Reviewed Status} & Generated candidate: keep / modify / relocate / reject; verify whether independently testable. & \\
    \hline
  \end{tabularx}
\end{table}

Intent. Information Retrieval is proposed as a discrete, testable enabling behaviour for Perioperative Logistics. It is included for architectural review, not as an established implementation obligation.

\noindent\textbf{Feature boundary.} This behaviour is scoped to Perioperative Logistics; an identically named Feature elsewhere may carry different business meaning.\par
\featuretable{
--- & --- & \emph{Not yet captured} & \emph{Complete during review} \\
}

\section{CT2 - Emergency Department Operations}\label{cap:bec-78f79edb}
\vspace{-0.5em}
\begin{table}[htbp]
  \centering
  \footnotesize
  \begin{tabularx}{1\textwidth}{|l|l|X|}
    \hline
    \multicolumn{3}{c}{\textbf{Capability Metadata}} \\
    \hline
    \textbf{Name} & \textbf{Value} & \textbf{Comment and/or Description} \\
    \hline
    \textbf{Reference Id} & \texttt{BEC-78F79EDB} & \\
    \hline
    \textbf{Status} & Candidate ancestry & \\
    \hline
    \textbf{Ancestry} &  & \\
    \hline
    \textbf{Source} &  & \\
    \hline
    \textbf{Enabled By} &  & \\
    \hline
    \textbf{Supports} &  & \\
    \hline
    \textbf{Derived From} &  & \\
    \hline
    \textbf{Reviewed Status} &  & \\
    \hline
  \end{tabularx}
\end{table}

\begin{table}[htbp]
  \centering
  \footnotesize
  \begin{tabularx}{1\textwidth}{|l|l|X|}
    \hline
    \multicolumn{3}{c}{\textbf{Contained Capabilities}} \\
    \hline
    \textbf{Capability Name} & \textbf{Capability Id} & \textbf{Comment} \\
    \hline
    ED Operational Coordination & \texttt{cap:bec-469419ab} (at \ref{cap:bec-469419ab}) & \\
    \hline
  \end{tabularx}
\end{table}

Intent. Emergency Department Operations is a candidate enabling capability within Health Service Operations Enablement. It frames the reusable system-supported behaviours needed to support this healthcare context without prescribing a product, interface, technology or deployment boundary.

Boundary. This capability supports the business activity and does not inherit the authority or accountability of the Business Capability it enables. Its child Features are review candidates until ancestry and semantics are adjudicated.

\subsection{CT3 - ED Operational Coordination}\label{cap:bec-469419ab}
\vspace{-0.5em}
\begin{table}[htbp]
  \centering
  \footnotesize
  \begin{tabularx}{1\textwidth}{|l|l|X|}
    \hline
    \multicolumn{3}{c}{\textbf{Capability Metadata}} \\
    \hline
    \textbf{Name} & \textbf{Value} & \textbf{Comment and/or Description} \\
    \hline
    \textbf{Reference Id} & \texttt{BEC-469419AB} & \\
    \hline
    \textbf{Status} & Candidate ancestry & \\
    \hline
    \textbf{Ancestry} &  & \\
    \hline
    \textbf{Source} &  & \\
    \hline
    \textbf{Enabled By} &  & \\
    \hline
    \textbf{Supports} &  & \\
    \hline
    \textbf{Derived From} &  & \\
    \hline
    \textbf{Reviewed Status} &  & \\
    \hline
  \end{tabularx}
\end{table}

\begin{table}[htbp]
  \centering
  \footnotesize
  \begin{tabularx}{1\textwidth}{|l|l|X|}
    \hline
    \multicolumn{3}{c}{\textbf{Contained Features}} \\
    \hline
    \textbf{Feature Name} & \textbf{Feature Id} & \textbf{Comment} \\
    \hline
    ED Bed \& Bay Tracking & \texttt{cap:bec-2c9f0b1a} (at \ref{cap:bec-2c9f0b1a}) & \\
    \hline
    ED Length of Stay Tracking & \texttt{cap:bec-a7e6edd2} (at \ref{cap:bec-a7e6edd2}) & \\
    \hline
    Context Establishment & \texttt{cap:bec-4cfe53a4} (at \ref{cap:bec-4cfe53a4}) & \\
    \hline
    Eligibility Validation & \texttt{cap:bec-67497215} (at \ref{cap:bec-67497215}) & \\
    \hline
    Assignment \& Routing & \texttt{cap:bec-a2b49b7c} (at \ref{cap:bec-a2b49b7c}) & \\
    \hline
    State Tracking & \texttt{cap:bec-704e04ea} (at \ref{cap:bec-704e04ea}) & \\
    \hline
    Exception Management & \texttt{cap:bec-bce98f58} (at \ref{cap:bec-bce98f58}) & \\
    \hline
    Outcome Reporting & \texttt{cap:bec-5e9192cc} (at \ref{cap:bec-5e9192cc}) & \\
    \hline
  \end{tabularx}
\end{table}

Intent. ED Operational Coordination is a candidate enabling capability within Emergency Department Operations. It frames the reusable system-supported behaviours needed to support this healthcare context without prescribing a product, interface, technology or deployment boundary.

Boundary. This capability supports the business activity and does not inherit the authority or accountability of the Business Capability it enables. Its child Features are review candidates until ancestry and semantics are adjudicated.

\subsubsection{Feature - ED Bed \& Bay Tracking}\label{cap:bec-2c9f0b1a}
\vspace{-1em}
\begin{table}[htbp]
  \centering
  \footnotesize
  \begin{tabularx}{1\textwidth}{|l|l|X|}
    \hline
    \textbf{Metadata Name} & \textbf{Metadata Value} & \textbf{Description} \\
    \hline
    \textbf{Reference Id} & \texttt{BEC-2C9F0B1A} & \\
    \hline
    \textbf{Status} & Existing Feature; candidate CT3 placement & \\
    \hline
    \textbf{Ancestry} &  & \\
    \hline
    \textbf{Source} & FEAT-HSO-07 & \\
    \hline
    \textbf{Enabled By} &  & \\
    \hline
    \textbf{Supports} &  & \\
    \hline
    \textbf{Derived From} &  & \\
    \hline
    \textbf{Reviewed Status} & Confirm existing Feature meaning and proposed CT3 ownership. & \\
    \hline
  \end{tabularx}
\end{table}

Intent. ED Bed \& Bay Tracking provides a distinct enabling behaviour in the context of ED Operational Coordination. This is an existing named Feature; its placement under the candidate CT3 remains subject to review.

\noindent\textbf{Feature boundary.} This behaviour is scoped to ED Operational Coordination; an identically named Feature elsewhere may carry different business meaning.\par
\featuretable{
Function & \texttt{FN-667B5CDD} & Coordinate ED Departmental Capacity & Catalogue \\
}
\subsubsection{Feature - ED Length of Stay Tracking}\label{cap:bec-a7e6edd2}
\vspace{-1em}
\begin{table}[htbp]
  \centering
  \footnotesize
  \begin{tabularx}{1\textwidth}{|l|l|X|}
    \hline
    \textbf{Metadata Name} & \textbf{Metadata Value} & \textbf{Description} \\
    \hline
    \textbf{Reference Id} & \texttt{BEC-A7E6EDD2} & \\
    \hline
    \textbf{Status} & Existing Feature; candidate CT3 placement & \\
    \hline
    \textbf{Ancestry} &  & \\
    \hline
    \textbf{Source} & FEAT-HSO-08 & \\
    \hline
    \textbf{Enabled By} &  & \\
    \hline
    \textbf{Supports} &  & \\
    \hline
    \textbf{Derived From} &  & \\
    \hline
    \textbf{Reviewed Status} & Confirm existing Feature meaning and proposed CT3 ownership. & \\
    \hline
  \end{tabularx}
\end{table}

Intent. ED Length of Stay Tracking provides a distinct enabling behaviour in the context of ED Operational Coordination. This is an existing named Feature; its placement under the candidate CT3 remains subject to review.

\noindent\textbf{Feature boundary.} This behaviour is scoped to ED Operational Coordination; an identically named Feature elsewhere may carry different business meaning.\par
\featuretable{
Function & \texttt{FN-2C25236D} & Association unresolved / not established & Catalogue \\
}
\subsubsection{Feature - Context Establishment}\label{cap:bec-4cfe53a4}
\vspace{-1em}
\begin{table}[htbp]
  \centering
  \footnotesize
  \begin{tabularx}{1\textwidth}{|l|l|X|}
    \hline
    \textbf{Metadata Name} & \textbf{Metadata Value} & \textbf{Description} \\
    \hline
    \textbf{Reference Id} & \texttt{BEC-4CFE53A4} & \\
    \hline
    \textbf{Status} & Generated candidate & \\
    \hline
    \textbf{Ancestry} &  & \\
    \hline
    \textbf{Source} & Semantic-pattern candidate & \\
    \hline
    \textbf{Enabled By} &  & \\
    \hline
    \textbf{Supports} &  & \\
    \hline
    \textbf{Derived From} &  & \\
    \hline
    \textbf{Reviewed Status} & Generated candidate: keep / modify / relocate / reject; verify whether independently testable. & \\
    \hline
  \end{tabularx}
\end{table}

Intent. Context Establishment is proposed as a discrete, testable enabling behaviour for ED Operational Coordination. It is included for architectural review, not as an established implementation obligation.

\noindent\textbf{Feature boundary.} This behaviour is scoped to ED Operational Coordination; an identically named Feature elsewhere may carry different business meaning.\par
\featuretable{
--- & --- & \emph{Not yet captured} & \emph{Complete during review} \\
}
\subsubsection{Feature - Eligibility Validation}\label{cap:bec-67497215}
\vspace{-1em}
\begin{table}[htbp]
  \centering
  \footnotesize
  \begin{tabularx}{1\textwidth}{|l|l|X|}
    \hline
    \textbf{Metadata Name} & \textbf{Metadata Value} & \textbf{Description} \\
    \hline
    \textbf{Reference Id} & \texttt{BEC-67497215} & \\
    \hline
    \textbf{Status} & Generated candidate & \\
    \hline
    \textbf{Ancestry} &  & \\
    \hline
    \textbf{Source} & Semantic-pattern candidate & \\
    \hline
    \textbf{Enabled By} &  & \\
    \hline
    \textbf{Supports} &  & \\
    \hline
    \textbf{Derived From} &  & \\
    \hline
    \textbf{Reviewed Status} & Generated candidate: keep / modify / relocate / reject; verify whether independently testable. & \\
    \hline
  \end{tabularx}
\end{table}

Intent. Eligibility Validation is proposed as a discrete, testable enabling behaviour for ED Operational Coordination. It is included for architectural review, not as an established implementation obligation.

\noindent\textbf{Feature boundary.} This behaviour is scoped to ED Operational Coordination; an identically named Feature elsewhere may carry different business meaning.\par
\featuretable{
--- & --- & \emph{Not yet captured} & \emph{Complete during review} \\
}
\subsubsection{Feature - Assignment \& Routing}\label{cap:bec-a2b49b7c}
\vspace{-1em}
\begin{table}[htbp]
  \centering
  \footnotesize
  \begin{tabularx}{1\textwidth}{|l|l|X|}
    \hline
    \textbf{Metadata Name} & \textbf{Metadata Value} & \textbf{Description} \\
    \hline
    \textbf{Reference Id} & \texttt{BEC-A2B49B7C} & \\
    \hline
    \textbf{Status} & Generated candidate & \\
    \hline
    \textbf{Ancestry} &  & \\
    \hline
    \textbf{Source} & Semantic-pattern candidate & \\
    \hline
    \textbf{Enabled By} &  & \\
    \hline
    \textbf{Supports} &  & \\
    \hline
    \textbf{Derived From} &  & \\
    \hline
    \textbf{Reviewed Status} & Generated candidate: keep / modify / relocate / reject; verify whether independently testable. & \\
    \hline
  \end{tabularx}
\end{table}

Intent. Assignment \& Routing is proposed as a discrete, testable enabling behaviour for ED Operational Coordination. It is included for architectural review, not as an established implementation obligation.

\noindent\textbf{Feature boundary.} This behaviour is scoped to ED Operational Coordination; an identically named Feature elsewhere may carry different business meaning.\par
\featuretable{
--- & --- & \emph{Not yet captured} & \emph{Complete during review} \\
}
\subsubsection{Feature - State Tracking}\label{cap:bec-704e04ea}
\vspace{-1em}
\begin{table}[htbp]
  \centering
  \footnotesize
  \begin{tabularx}{1\textwidth}{|l|l|X|}
    \hline
    \textbf{Metadata Name} & \textbf{Metadata Value} & \textbf{Description} \\
    \hline
    \textbf{Reference Id} & \texttt{BEC-704E04EA} & \\
    \hline
    \textbf{Status} & Generated candidate & \\
    \hline
    \textbf{Ancestry} &  & \\
    \hline
    \textbf{Source} & Semantic-pattern candidate & \\
    \hline
    \textbf{Enabled By} &  & \\
    \hline
    \textbf{Supports} &  & \\
    \hline
    \textbf{Derived From} &  & \\
    \hline
    \textbf{Reviewed Status} & Generated candidate: keep / modify / relocate / reject; verify whether independently testable. & \\
    \hline
  \end{tabularx}
\end{table}

Intent. State Tracking is proposed as a discrete, testable enabling behaviour for ED Operational Coordination. It is included for architectural review, not as an established implementation obligation.

\noindent\textbf{Feature boundary.} This behaviour is scoped to ED Operational Coordination; an identically named Feature elsewhere may carry different business meaning.\par
\featuretable{
--- & --- & \emph{Not yet captured} & \emph{Complete during review} \\
}
\subsubsection{Feature - Exception Management}\label{cap:bec-bce98f58}
\vspace{-1em}
\begin{table}[htbp]
  \centering
  \footnotesize
  \begin{tabularx}{1\textwidth}{|l|l|X|}
    \hline
    \textbf{Metadata Name} & \textbf{Metadata Value} & \textbf{Description} \\
    \hline
    \textbf{Reference Id} & \texttt{BEC-BCE98F58} & \\
    \hline
    \textbf{Status} & Generated candidate & \\
    \hline
    \textbf{Ancestry} &  & \\
    \hline
    \textbf{Source} & Semantic-pattern candidate & \\
    \hline
    \textbf{Enabled By} &  & \\
    \hline
    \textbf{Supports} &  & \\
    \hline
    \textbf{Derived From} &  & \\
    \hline
    \textbf{Reviewed Status} & Generated candidate: keep / modify / relocate / reject; verify whether independently testable. & \\
    \hline
  \end{tabularx}
\end{table}

Intent. Exception Management is proposed as a discrete, testable enabling behaviour for ED Operational Coordination. It is included for architectural review, not as an established implementation obligation.

\noindent\textbf{Feature boundary.} This behaviour is scoped to ED Operational Coordination; an identically named Feature elsewhere may carry different business meaning.\par
\featuretable{
--- & --- & \emph{Not yet captured} & \emph{Complete during review} \\
}
\subsubsection{Feature - Outcome Reporting}\label{cap:bec-5e9192cc}
\vspace{-1em}
\begin{table}[htbp]
  \centering
  \footnotesize
  \begin{tabularx}{1\textwidth}{|l|l|X|}
    \hline
    \textbf{Metadata Name} & \textbf{Metadata Value} & \textbf{Description} \\
    \hline
    \textbf{Reference Id} & \texttt{BEC-5E9192CC} & \\
    \hline
    \textbf{Status} & Generated candidate & \\
    \hline
    \textbf{Ancestry} &  & \\
    \hline
    \textbf{Source} & Semantic-pattern candidate & \\
    \hline
    \textbf{Enabled By} &  & \\
    \hline
    \textbf{Supports} &  & \\
    \hline
    \textbf{Derived From} &  & \\
    \hline
    \textbf{Reviewed Status} & Generated candidate: keep / modify / relocate / reject; verify whether independently testable. & \\
    \hline
  \end{tabularx}
\end{table}

Intent. Outcome Reporting is proposed as a discrete, testable enabling behaviour for ED Operational Coordination. It is included for architectural review, not as an established implementation obligation.

\noindent\textbf{Feature boundary.} This behaviour is scoped to ED Operational Coordination; an identically named Feature elsewhere may carry different business meaning.\par
\featuretable{
--- & --- & \emph{Not yet captured} & \emph{Complete during review} \\
}

\section{CT2 - Outpatient Operations}\label{cap:bec-3ea18c92}
\vspace{-0.5em}
\begin{table}[htbp]
  \centering
  \footnotesize
  \begin{tabularx}{1\textwidth}{|l|l|X|}
    \hline
    \multicolumn{3}{c}{\textbf{Capability Metadata}} \\
    \hline
    \textbf{Name} & \textbf{Value} & \textbf{Comment and/or Description} \\
    \hline
    \textbf{Reference Id} & \texttt{BEC-3EA18C92} & \\
    \hline
    \textbf{Status} & Candidate ancestry & \\
    \hline
    \textbf{Ancestry} &  & \\
    \hline
    \textbf{Source} &  & \\
    \hline
    \textbf{Enabled By} &  & \\
    \hline
    \textbf{Supports} &  & \\
    \hline
    \textbf{Derived From} &  & \\
    \hline
    \textbf{Reviewed Status} &  & \\
    \hline
  \end{tabularx}
\end{table}

\begin{table}[htbp]
  \centering
  \footnotesize
  \begin{tabularx}{1\textwidth}{|l|l|X|}
    \hline
    \multicolumn{3}{c}{\textbf{Contained Capabilities}} \\
    \hline
    \textbf{Capability Name} & \textbf{Capability Id} & \textbf{Comment} \\
    \hline
    Outpatient Flow Enablement & \texttt{cap:bec-07398864} (at \ref{cap:bec-07398864}) & \\
    \hline
  \end{tabularx}
\end{table}

Intent. Outpatient Operations is a candidate enabling capability within Health Service Operations Enablement. It frames the reusable system-supported behaviours needed to support this healthcare context without prescribing a product, interface, technology or deployment boundary.

Boundary. This capability supports the business activity and does not inherit the authority or accountability of the Business Capability it enables. Its child Features are review candidates until ancestry and semantics are adjudicated.

\subsection{CT3 - Outpatient Flow Enablement}\label{cap:bec-07398864}
\vspace{-0.5em}
\begin{table}[htbp]
  \centering
  \footnotesize
  \begin{tabularx}{1\textwidth}{|l|l|X|}
    \hline
    \multicolumn{3}{c}{\textbf{Capability Metadata}} \\
    \hline
    \textbf{Name} & \textbf{Value} & \textbf{Comment and/or Description} \\
    \hline
    \textbf{Reference Id} & \texttt{BEC-07398864} & \\
    \hline
    \textbf{Status} & Candidate ancestry & \\
    \hline
    \textbf{Ancestry} &  & \\
    \hline
    \textbf{Source} &  & \\
    \hline
    \textbf{Enabled By} &  & \\
    \hline
    \textbf{Supports} &  & \\
    \hline
    \textbf{Derived From} &  & \\
    \hline
    \textbf{Reviewed Status} &  & \\
    \hline
  \end{tabularx}
\end{table}

\begin{table}[htbp]
  \centering
  \footnotesize
  \begin{tabularx}{1\textwidth}{|l|l|X|}
    \hline
    \multicolumn{3}{c}{\textbf{Contained Features}} \\
    \hline
    \textbf{Feature Name} & \textbf{Feature Id} & \textbf{Comment} \\
    \hline
    Outpatient Arrival Registration & \texttt{cap:bec-4382a684} (at \ref{cap:bec-4382a684}) & \\
    \hline
    Post-Clinic Order Coordination & \texttt{cap:bec-53eb5f3c} (at \ref{cap:bec-53eb5f3c}) & \\
    \hline
    Subject Identity Correlation & \texttt{cap:bec-0d49726e} (at \ref{cap:bec-0d49726e}) & \\
    \hline
    Subject Relationship Management & \texttt{cap:bec-fcb21dfc} (at \ref{cap:bec-fcb21dfc}) & \\
    \hline
    Subject Context Management & \texttt{cap:bec-cd0d5a3a} (at \ref{cap:bec-cd0d5a3a}) & \\
    \hline
    Subject Information Discovery & \texttt{cap:bec-17daea1c} (at \ref{cap:bec-17daea1c}) & \\
    \hline
  \end{tabularx}
\end{table}

Intent. Outpatient Flow Enablement is a candidate enabling capability within Outpatient Operations. It frames the reusable system-supported behaviours needed to support this healthcare context without prescribing a product, interface, technology or deployment boundary.

Boundary. This capability supports the business activity and does not inherit the authority or accountability of the Business Capability it enables. Its child Features are review candidates until ancestry and semantics are adjudicated.

FHIR R5 conceptual analogues considered (not implementation mappings): Patient, Person.

\subsubsection{Feature - Outpatient Arrival Registration}\label{cap:bec-4382a684}
\vspace{-1em}
\begin{table}[htbp]
  \centering
  \footnotesize
  \begin{tabularx}{1\textwidth}{|l|l|X|}
    \hline
    \textbf{Metadata Name} & \textbf{Metadata Value} & \textbf{Description} \\
    \hline
    \textbf{Reference Id} & \texttt{BEC-4382A684} & \\
    \hline
    \textbf{Status} & Existing Feature; candidate CT3 placement & \\
    \hline
    \textbf{Ancestry} &  & \\
    \hline
    \textbf{Source} & FEAT-HSO-09 & \\
    \hline
    \textbf{Enabled By} &  & \\
    \hline
    \textbf{Supports} &  & \\
    \hline
    \textbf{Derived From} &  & \\
    \hline
    \textbf{Reviewed Status} & Confirm existing Feature meaning and proposed CT3 ownership. & \\
    \hline
  \end{tabularx}
\end{table}

Intent. Outpatient Arrival Registration provides a distinct enabling behaviour in the context of Outpatient Flow Enablement. This is an existing named Feature; its placement under the candidate CT3 remains subject to review.

\noindent\textbf{Feature boundary.} This behaviour is scoped to Outpatient Flow Enablement; an identically named Feature elsewhere may carry different business meaning.\par
\featuretable{
Function & \texttt{FN-28A12573} & Manage Outpatient Capacity & Catalogue \\
}
\subsubsection{Feature - Post-Clinic Order Coordination}\label{cap:bec-53eb5f3c}
\vspace{-1em}
\begin{table}[htbp]
  \centering
  \footnotesize
  \begin{tabularx}{1\textwidth}{|l|l|X|}
    \hline
    \textbf{Metadata Name} & \textbf{Metadata Value} & \textbf{Description} \\
    \hline
    \textbf{Reference Id} & \texttt{BEC-53EB5F3C} & \\
    \hline
    \textbf{Status} & Existing Feature; candidate CT3 placement & \\
    \hline
    \textbf{Ancestry} &  & \\
    \hline
    \textbf{Source} & FEAT-HSO-10 & \\
    \hline
    \textbf{Enabled By} &  & \\
    \hline
    \textbf{Supports} &  & \\
    \hline
    \textbf{Derived From} &  & \\
    \hline
    \textbf{Reviewed Status} & Confirm existing Feature meaning and proposed CT3 ownership. & \\
    \hline
  \end{tabularx}
\end{table}

Intent. Post-Clinic Order Coordination provides a distinct enabling behaviour in the context of Outpatient Flow Enablement. This is an existing named Feature; its placement under the candidate CT3 remains subject to review.

\noindent\textbf{Feature boundary.} This behaviour is scoped to Outpatient Flow Enablement; an identically named Feature elsewhere may carry different business meaning.\par
\featuretable{
Function & \texttt{FN-B37F2F7D} & Capture Post-Clinic Follow-up Requirements & Catalogue \\
}
\subsubsection{Feature - Subject Identity Correlation}\label{cap:bec-0d49726e}
\vspace{-1em}
\begin{table}[htbp]
  \centering
  \footnotesize
  \begin{tabularx}{1\textwidth}{|l|l|X|}
    \hline
    \textbf{Metadata Name} & \textbf{Metadata Value} & \textbf{Description} \\
    \hline
    \textbf{Reference Id} & \texttt{BEC-0D49726E} & \\
    \hline
    \textbf{Status} & Generated candidate & \\
    \hline
    \textbf{Ancestry} &  & \\
    \hline
    \textbf{Source} & FHIR-informed candidate & \\
    \hline
    \textbf{Enabled By} &  & \\
    \hline
    \textbf{Supports} &  & \\
    \hline
    \textbf{Derived From} &  & \\
    \hline
    \textbf{Reviewed Status} & Generated candidate: keep / modify / relocate / reject; verify whether independently testable. & \\
    \hline
  \end{tabularx}
\end{table}

Intent. Subject Identity Correlation is proposed to support the underlying healthcare concept represented by Outpatient Flow Enablement. The proposal concerns behaviour and semantic responsibility, not a FHIR field, operation, resource instance or prescribed data model.

Conceptual evidence: HL7 FHIR R5 Patient, Person.

\noindent\textbf{Feature boundary.} This behaviour is scoped to Outpatient Flow Enablement; an identically named Feature elsewhere may carry different business meaning.\par
\featuretable{
--- & --- & \emph{Not yet captured} & \emph{Complete during review} \\
}
\subsubsection{Feature - Subject Relationship Management}\label{cap:bec-fcb21dfc}
\vspace{-1em}
\begin{table}[htbp]
  \centering
  \footnotesize
  \begin{tabularx}{1\textwidth}{|l|l|X|}
    \hline
    \textbf{Metadata Name} & \textbf{Metadata Value} & \textbf{Description} \\
    \hline
    \textbf{Reference Id} & \texttt{BEC-FCB21DFC} & \\
    \hline
    \textbf{Status} & Generated candidate & \\
    \hline
    \textbf{Ancestry} &  & \\
    \hline
    \textbf{Source} & FHIR-informed candidate & \\
    \hline
    \textbf{Enabled By} &  & \\
    \hline
    \textbf{Supports} &  & \\
    \hline
    \textbf{Derived From} &  & \\
    \hline
    \textbf{Reviewed Status} & Generated candidate: keep / modify / relocate / reject; verify whether independently testable. & \\
    \hline
  \end{tabularx}
\end{table}

Intent. Subject Relationship Management is proposed to support the underlying healthcare concept represented by Outpatient Flow Enablement. The proposal concerns behaviour and semantic responsibility, not a FHIR field, operation, resource instance or prescribed data model.

Conceptual evidence: HL7 FHIR R5 Patient, Person.

\noindent\textbf{Feature boundary.} This behaviour is scoped to Outpatient Flow Enablement; an identically named Feature elsewhere may carry different business meaning.\par
\featuretable{
--- & --- & \emph{Not yet captured} & \emph{Complete during review} \\
}
\subsubsection{Feature - Subject Context Management}\label{cap:bec-cd0d5a3a}
\vspace{-1em}
\begin{table}[htbp]
  \centering
  \footnotesize
  \begin{tabularx}{1\textwidth}{|l|l|X|}
    \hline
    \textbf{Metadata Name} & \textbf{Metadata Value} & \textbf{Description} \\
    \hline
    \textbf{Reference Id} & \texttt{BEC-CD0D5A3A} & \\
    \hline
    \textbf{Status} & Generated candidate & \\
    \hline
    \textbf{Ancestry} &  & \\
    \hline
    \textbf{Source} & FHIR-informed candidate & \\
    \hline
    \textbf{Enabled By} &  & \\
    \hline
    \textbf{Supports} &  & \\
    \hline
    \textbf{Derived From} &  & \\
    \hline
    \textbf{Reviewed Status} & Generated candidate: keep / modify / relocate / reject; verify whether independently testable. & \\
    \hline
  \end{tabularx}
\end{table}

Intent. Subject Context Management is proposed to support the underlying healthcare concept represented by Outpatient Flow Enablement. The proposal concerns behaviour and semantic responsibility, not a FHIR field, operation, resource instance or prescribed data model.

Conceptual evidence: HL7 FHIR R5 Patient, Person.

\noindent\textbf{Feature boundary.} This behaviour is scoped to Outpatient Flow Enablement; an identically named Feature elsewhere may carry different business meaning.\par
\featuretable{
--- & --- & \emph{Not yet captured} & \emph{Complete during review} \\
}
\subsubsection{Feature - Subject Information Discovery}\label{cap:bec-17daea1c}
\vspace{-1em}
\begin{table}[htbp]
  \centering
  \footnotesize
  \begin{tabularx}{1\textwidth}{|l|l|X|}
    \hline
    \textbf{Metadata Name} & \textbf{Metadata Value} & \textbf{Description} \\
    \hline
    \textbf{Reference Id} & \texttt{BEC-17DAEA1C} & \\
    \hline
    \textbf{Status} & Generated candidate & \\
    \hline
    \textbf{Ancestry} &  & \\
    \hline
    \textbf{Source} & FHIR-informed candidate & \\
    \hline
    \textbf{Enabled By} &  & \\
    \hline
    \textbf{Supports} &  & \\
    \hline
    \textbf{Derived From} &  & \\
    \hline
    \textbf{Reviewed Status} & Generated candidate: keep / modify / relocate / reject; verify whether independently testable. & \\
    \hline
  \end{tabularx}
\end{table}

Intent. Subject Information Discovery is proposed to support the underlying healthcare concept represented by Outpatient Flow Enablement. The proposal concerns behaviour and semantic responsibility, not a FHIR field, operation, resource instance or prescribed data model.

Conceptual evidence: HL7 FHIR R5 Patient, Person.

\noindent\textbf{Feature boundary.} This behaviour is scoped to Outpatient Flow Enablement; an identically named Feature elsewhere may carry different business meaning.\par
\featuretable{
--- & --- & \emph{Not yet captured} & \emph{Complete during review} \\
}

\section{CT2 - Bed \& Care-Place Management}\label{cap:bec-dba54d09}
\vspace{-0.5em}
\begin{table}[htbp]
  \centering
  \footnotesize
  \begin{tabularx}{1\textwidth}{|l|l|X|}
    \hline
    \multicolumn{3}{c}{\textbf{Capability Metadata}} \\
    \hline
    \textbf{Name} & \textbf{Value} & \textbf{Comment and/or Description} \\
    \hline
    \textbf{Reference Id} & \texttt{BEC-DBA54D09} & \\
    \hline
    \textbf{Status} & Candidate ancestry & \\
    \hline
    \textbf{Ancestry} &  & \\
    \hline
    \textbf{Source} &  & \\
    \hline
    \textbf{Enabled By} &  & \\
    \hline
    \textbf{Supports} &  & \\
    \hline
    \textbf{Derived From} &  & \\
    \hline
    \textbf{Reviewed Status} &  & \\
    \hline
  \end{tabularx}
\end{table}

\begin{table}[htbp]
  \centering
  \footnotesize
  \begin{tabularx}{1\textwidth}{|l|l|X|}
    \hline
    \multicolumn{3}{c}{\textbf{Contained Capabilities}} \\
    \hline
    \textbf{Capability Name} & \textbf{Capability Id} & \textbf{Comment} \\
    \hline
    Bed State \& Turnover Management & \texttt{cap:bec-9f88f277} (at \ref{cap:bec-9f88f277}) & \\
    \hline
  \end{tabularx}
\end{table}

Intent. Bed \& Care-Place Management is a candidate enabling capability within Health Service Operations Enablement. It frames the reusable system-supported behaviours needed to support this healthcare context without prescribing a product, interface, technology or deployment boundary.

Boundary. This capability supports the business activity and does not inherit the authority or accountability of the Business Capability it enables. Its child Features are review candidates until ancestry and semantics are adjudicated.

\subsection{CT3 - Bed State \& Turnover Management}\label{cap:bec-9f88f277}
\vspace{-0.5em}
\begin{table}[htbp]
  \centering
  \footnotesize
  \begin{tabularx}{1\textwidth}{|l|l|X|}
    \hline
    \multicolumn{3}{c}{\textbf{Capability Metadata}} \\
    \hline
    \textbf{Name} & \textbf{Value} & \textbf{Comment and/or Description} \\
    \hline
    \textbf{Reference Id} & \texttt{BEC-9F88F277} & \\
    \hline
    \textbf{Status} & Candidate ancestry & \\
    \hline
    \textbf{Ancestry} &  & \\
    \hline
    \textbf{Source} &  & \\
    \hline
    \textbf{Enabled By} &  & \\
    \hline
    \textbf{Supports} &  & \\
    \hline
    \textbf{Derived From} &  & \\
    \hline
    \textbf{Reviewed Status} &  & \\
    \hline
  \end{tabularx}
\end{table}

\begin{table}[htbp]
  \centering
  \footnotesize
  \begin{tabularx}{1\textwidth}{|l|l|X|}
    \hline
    \multicolumn{3}{c}{\textbf{Contained Features}} \\
    \hline
    \textbf{Feature Name} & \textbf{Feature Id} & \textbf{Comment} \\
    \hline
    Bed Availability Tracking & \texttt{cap:bec-1a021f75} (at \ref{cap:bec-1a021f75}) & \\
    \hline
    Bed State Progression & \texttt{cap:bec-abddba4a} (at \ref{cap:bec-abddba4a}) & \\
    \hline
    Bed Readiness Notification & \texttt{cap:bec-cb31c6e4} (at \ref{cap:bec-cb31c6e4}) & \\
    \hline
    Context Establishment & \texttt{cap:bec-21f00612} (at \ref{cap:bec-21f00612}) & \\
    \hline
    State Monitoring & \texttt{cap:bec-25aae44b} (at \ref{cap:bec-25aae44b}) & \\
    \hline
    Exception Management & \texttt{cap:bec-49bb68fa} (at \ref{cap:bec-49bb68fa}) & \\
    \hline
    Information Retrieval & \texttt{cap:bec-be0defdf} (at \ref{cap:bec-be0defdf}) & \\
    \hline
  \end{tabularx}
\end{table}

Intent. Bed State \& Turnover Management is a candidate enabling capability within Bed \& Care-Place Management. It frames the reusable system-supported behaviours needed to support this healthcare context without prescribing a product, interface, technology or deployment boundary.

Boundary. This capability supports the business activity and does not inherit the authority or accountability of the Business Capability it enables. Its child Features are review candidates until ancestry and semantics are adjudicated.

\subsubsection{Feature - Bed Availability Tracking}\label{cap:bec-1a021f75}
\vspace{-1em}
\begin{table}[htbp]
  \centering
  \footnotesize
  \begin{tabularx}{1\textwidth}{|l|l|X|}
    \hline
    \textbf{Metadata Name} & \textbf{Metadata Value} & \textbf{Description} \\
    \hline
    \textbf{Reference Id} & \texttt{BEC-1A021F75} & \\
    \hline
    \textbf{Status} & Existing Feature; candidate CT3 placement & \\
    \hline
    \textbf{Ancestry} &  & \\
    \hline
    \textbf{Source} & FEAT-HSO-11 & \\
    \hline
    \textbf{Enabled By} &  & \\
    \hline
    \textbf{Supports} &  & \\
    \hline
    \textbf{Derived From} &  & \\
    \hline
    \textbf{Reviewed Status} & Confirm existing Feature meaning and proposed CT3 ownership. & \\
    \hline
  \end{tabularx}
\end{table}

Intent. Bed Availability Tracking provides a distinct enabling behaviour in the context of Bed State \& Turnover Management. This is an existing named Feature; its placement under the candidate CT3 remains subject to review.

\noindent\textbf{Feature boundary.} This behaviour is scoped to Bed State \& Turnover Management; an identically named Feature elsewhere may carry different business meaning.\par
\featuretable{
Function & \texttt{FN-D9B0D559} & Maintain Operational Bed State & Catalogue \\
Service & \texttt{SV-8BFF8FD4} & Bed Availability Query & Linked to preceding Function in source \\
}
\subsubsection{Feature - Bed State Progression}\label{cap:bec-abddba4a}
\vspace{-1em}
\begin{table}[htbp]
  \centering
  \footnotesize
  \begin{tabularx}{1\textwidth}{|l|l|X|}
    \hline
    \textbf{Metadata Name} & \textbf{Metadata Value} & \textbf{Description} \\
    \hline
    \textbf{Reference Id} & \texttt{BEC-ABDDBA4A} & \\
    \hline
    \textbf{Status} & Existing Feature; candidate CT3 placement & \\
    \hline
    \textbf{Ancestry} &  & \\
    \hline
    \textbf{Source} & FEAT-HSO-12 & \\
    \hline
    \textbf{Enabled By} &  & \\
    \hline
    \textbf{Supports} &  & \\
    \hline
    \textbf{Derived From} &  & \\
    \hline
    \textbf{Reviewed Status} & Confirm existing Feature meaning and proposed CT3 ownership. & \\
    \hline
  \end{tabularx}
\end{table}

Intent. Bed State Progression provides a distinct enabling behaviour in the context of Bed State \& Turnover Management. This is an existing named Feature; its placement under the candidate CT3 remains subject to review.

\noindent\textbf{Feature boundary.} This behaviour is scoped to Bed State \& Turnover Management; an identically named Feature elsewhere may carry different business meaning.\par
\featuretable{
Function & \texttt{FN-65C150FA} & Manage Bed Turnover Progression & Catalogue \\
Service & \texttt{SV-4090386A} & Bed Turnover Lifecycle Service & Linked to preceding Function in source \\
Process & \texttt{PR-76C02538} & Bed Turnover Process & Linked to preceding Function in source \\
}
\subsubsection{Feature - Bed Readiness Notification}\label{cap:bec-cb31c6e4}
\vspace{-1em}
\begin{table}[htbp]
  \centering
  \footnotesize
  \begin{tabularx}{1\textwidth}{|l|l|X|}
    \hline
    \textbf{Metadata Name} & \textbf{Metadata Value} & \textbf{Description} \\
    \hline
    \textbf{Reference Id} & \texttt{BEC-CB31C6E4} & \\
    \hline
    \textbf{Status} & Existing Feature; candidate CT3 placement & \\
    \hline
    \textbf{Ancestry} &  & \\
    \hline
    \textbf{Source} & FEAT-HSO-13 & \\
    \hline
    \textbf{Enabled By} &  & \\
    \hline
    \textbf{Supports} &  & \\
    \hline
    \textbf{Derived From} &  & \\
    \hline
    \textbf{Reviewed Status} & Confirm existing Feature meaning and proposed CT3 ownership. & \\
    \hline
  \end{tabularx}
\end{table}

Intent. Bed Readiness Notification provides a distinct enabling behaviour in the context of Bed State \& Turnover Management. This is an existing named Feature; its placement under the candidate CT3 remains subject to review.

\noindent\textbf{Feature boundary.} This behaviour is scoped to Bed State \& Turnover Management; an identically named Feature elsewhere may carry different business meaning.\par
\featuretable{
Function & \texttt{FN-6E1D2F9C} & Manage Bed Turnover Progression & Catalogue \\
Process & \texttt{PR-F02CCB49} & Bed Turnover Process & Linked to preceding Function in source \\
}
\subsubsection{Feature - Context Establishment}\label{cap:bec-21f00612}
\vspace{-1em}
\begin{table}[htbp]
  \centering
  \footnotesize
  \begin{tabularx}{1\textwidth}{|l|l|X|}
    \hline
    \textbf{Metadata Name} & \textbf{Metadata Value} & \textbf{Description} \\
    \hline
    \textbf{Reference Id} & \texttt{BEC-21F00612} & \\
    \hline
    \textbf{Status} & Generated candidate & \\
    \hline
    \textbf{Ancestry} &  & \\
    \hline
    \textbf{Source} & Semantic-pattern candidate & \\
    \hline
    \textbf{Enabled By} &  & \\
    \hline
    \textbf{Supports} &  & \\
    \hline
    \textbf{Derived From} &  & \\
    \hline
    \textbf{Reviewed Status} & Generated candidate: keep / modify / relocate / reject; verify whether independently testable. & \\
    \hline
  \end{tabularx}
\end{table}

Intent. Context Establishment is proposed as a discrete, testable enabling behaviour for Bed State \& Turnover Management. It is included for architectural review, not as an established implementation obligation.

\noindent\textbf{Feature boundary.} This behaviour is scoped to Bed State \& Turnover Management; an identically named Feature elsewhere may carry different business meaning.\par
\featuretable{
--- & --- & \emph{Not yet captured} & \emph{Complete during review} \\
}
\subsubsection{Feature - State Monitoring}\label{cap:bec-25aae44b}
\vspace{-1em}
\begin{table}[htbp]
  \centering
  \footnotesize
  \begin{tabularx}{1\textwidth}{|l|l|X|}
    \hline
    \textbf{Metadata Name} & \textbf{Metadata Value} & \textbf{Description} \\
    \hline
    \textbf{Reference Id} & \texttt{BEC-25AAE44B} & \\
    \hline
    \textbf{Status} & Generated candidate & \\
    \hline
    \textbf{Ancestry} &  & \\
    \hline
    \textbf{Source} & Semantic-pattern candidate & \\
    \hline
    \textbf{Enabled By} &  & \\
    \hline
    \textbf{Supports} &  & \\
    \hline
    \textbf{Derived From} &  & \\
    \hline
    \textbf{Reviewed Status} & Generated candidate: keep / modify / relocate / reject; verify whether independently testable. & \\
    \hline
  \end{tabularx}
\end{table}

Intent. State Monitoring is proposed as a discrete, testable enabling behaviour for Bed State \& Turnover Management. It is included for architectural review, not as an established implementation obligation.

\noindent\textbf{Feature boundary.} This behaviour is scoped to Bed State \& Turnover Management; an identically named Feature elsewhere may carry different business meaning.\par
\featuretable{
--- & --- & \emph{Not yet captured} & \emph{Complete during review} \\
}
\subsubsection{Feature - Exception Management}\label{cap:bec-49bb68fa}
\vspace{-1em}
\begin{table}[htbp]
  \centering
  \footnotesize
  \begin{tabularx}{1\textwidth}{|l|l|X|}
    \hline
    \textbf{Metadata Name} & \textbf{Metadata Value} & \textbf{Description} \\
    \hline
    \textbf{Reference Id} & \texttt{BEC-49BB68FA} & \\
    \hline
    \textbf{Status} & Generated candidate & \\
    \hline
    \textbf{Ancestry} &  & \\
    \hline
    \textbf{Source} & Semantic-pattern candidate & \\
    \hline
    \textbf{Enabled By} &  & \\
    \hline
    \textbf{Supports} &  & \\
    \hline
    \textbf{Derived From} &  & \\
    \hline
    \textbf{Reviewed Status} & Generated candidate: keep / modify / relocate / reject; verify whether independently testable. & \\
    \hline
  \end{tabularx}
\end{table}

Intent. Exception Management is proposed as a discrete, testable enabling behaviour for Bed State \& Turnover Management. It is included for architectural review, not as an established implementation obligation.

\noindent\textbf{Feature boundary.} This behaviour is scoped to Bed State \& Turnover Management; an identically named Feature elsewhere may carry different business meaning.\par
\featuretable{
--- & --- & \emph{Not yet captured} & \emph{Complete during review} \\
}
\subsubsection{Feature - Information Retrieval}\label{cap:bec-be0defdf}
\vspace{-1em}
\begin{table}[htbp]
  \centering
  \footnotesize
  \begin{tabularx}{1\textwidth}{|l|l|X|}
    \hline
    \textbf{Metadata Name} & \textbf{Metadata Value} & \textbf{Description} \\
    \hline
    \textbf{Reference Id} & \texttt{BEC-BE0DEFDF} & \\
    \hline
    \textbf{Status} & Generated candidate & \\
    \hline
    \textbf{Ancestry} &  & \\
    \hline
    \textbf{Source} & Semantic-pattern candidate & \\
    \hline
    \textbf{Enabled By} &  & \\
    \hline
    \textbf{Supports} &  & \\
    \hline
    \textbf{Derived From} &  & \\
    \hline
    \textbf{Reviewed Status} & Generated candidate: keep / modify / relocate / reject; verify whether independently testable. & \\
    \hline
  \end{tabularx}
\end{table}

Intent. Information Retrieval is proposed as a discrete, testable enabling behaviour for Bed State \& Turnover Management. It is included for architectural review, not as an established implementation obligation.

\noindent\textbf{Feature boundary.} This behaviour is scoped to Bed State \& Turnover Management; an identically named Feature elsewhere may carry different business meaning.\par
\featuretable{
--- & --- & \emph{Not yet captured} & \emph{Complete during review} \\
}

\section{CT2 - Clinical Resource Management}\label{cap:bec-07a2fd29}
\vspace{-0.5em}
\begin{table}[htbp]
  \centering
  \footnotesize
  \begin{tabularx}{1\textwidth}{|l|l|X|}
    \hline
    \multicolumn{3}{c}{\textbf{Capability Metadata}} \\
    \hline
    \textbf{Name} & \textbf{Value} & \textbf{Comment and/or Description} \\
    \hline
    \textbf{Reference Id} & \texttt{BEC-07A2FD29} & \\
    \hline
    \textbf{Status} & Candidate ancestry & \\
    \hline
    \textbf{Ancestry} &  & \\
    \hline
    \textbf{Source} &  & \\
    \hline
    \textbf{Enabled By} &  & \\
    \hline
    \textbf{Supports} &  & \\
    \hline
    \textbf{Derived From} &  & \\
    \hline
    \textbf{Reviewed Status} &  & \\
    \hline
  \end{tabularx}
\end{table}

\begin{table}[htbp]
  \centering
  \footnotesize
  \begin{tabularx}{1\textwidth}{|l|l|X|}
    \hline
    \multicolumn{3}{c}{\textbf{Contained Capabilities}} \\
    \hline
    \textbf{Capability Name} & \textbf{Capability Id} & \textbf{Comment} \\
    \hline
    Specialized Equipment Coordination & \texttt{cap:bec-556b4a3f} (at \ref{cap:bec-556b4a3f}) & \\
    \hline
  \end{tabularx}
\end{table}

Intent. Clinical Resource Management is a candidate enabling capability within Health Service Operations Enablement. It frames the reusable system-supported behaviours needed to support this healthcare context without prescribing a product, interface, technology or deployment boundary.

Boundary. This capability supports the business activity and does not inherit the authority or accountability of the Business Capability it enables. Its child Features are review candidates until ancestry and semantics are adjudicated.

\subsection{CT3 - Specialized Equipment Coordination}\label{cap:bec-556b4a3f}
\vspace{-0.5em}
\begin{table}[htbp]
  \centering
  \footnotesize
  \begin{tabularx}{1\textwidth}{|l|l|X|}
    \hline
    \multicolumn{3}{c}{\textbf{Capability Metadata}} \\
    \hline
    \textbf{Name} & \textbf{Value} & \textbf{Comment and/or Description} \\
    \hline
    \textbf{Reference Id} & \texttt{BEC-556B4A3F} & \\
    \hline
    \textbf{Status} & Candidate ancestry & \\
    \hline
    \textbf{Ancestry} &  & \\
    \hline
    \textbf{Source} &  & \\
    \hline
    \textbf{Enabled By} &  & \\
    \hline
    \textbf{Supports} &  & \\
    \hline
    \textbf{Derived From} &  & \\
    \hline
    \textbf{Reviewed Status} &  & \\
    \hline
  \end{tabularx}
\end{table}

\begin{table}[htbp]
  \centering
  \footnotesize
  \begin{tabularx}{1\textwidth}{|l|l|X|}
    \hline
    \multicolumn{3}{c}{\textbf{Contained Features}} \\
    \hline
    \textbf{Feature Name} & \textbf{Feature Id} & \textbf{Comment} \\
    \hline
    Mobile Clinical Equipment Tracking & \texttt{cap:bec-5b685445} (at \ref{cap:bec-5b685445}) & \\
    \hline
    Context Establishment & \texttt{cap:bec-e7e5a452} (at \ref{cap:bec-e7e5a452}) & \\
    \hline
    Eligibility Validation & \texttt{cap:bec-aa245ff0} (at \ref{cap:bec-aa245ff0}) & \\
    \hline
    Assignment \& Routing & \texttt{cap:bec-a870eb50} (at \ref{cap:bec-a870eb50}) & \\
    \hline
    State Tracking & \texttt{cap:bec-f1b486e3} (at \ref{cap:bec-f1b486e3}) & \\
    \hline
    Exception Management & \texttt{cap:bec-99e2fa68} (at \ref{cap:bec-99e2fa68}) & \\
    \hline
    Outcome Reporting & \texttt{cap:bec-4d2a47ac} (at \ref{cap:bec-4d2a47ac}) & \\
    \hline
  \end{tabularx}
\end{table}

Intent. Specialized Equipment Coordination is a candidate enabling capability within Clinical Resource Management. It frames the reusable system-supported behaviours needed to support this healthcare context without prescribing a product, interface, technology or deployment boundary.

Boundary. This capability supports the business activity and does not inherit the authority or accountability of the Business Capability it enables. Its child Features are review candidates until ancestry and semantics are adjudicated.

\subsubsection{Feature - Mobile Clinical Equipment Tracking}\label{cap:bec-5b685445}
\vspace{-1em}
\begin{table}[htbp]
  \centering
  \footnotesize
  \begin{tabularx}{1\textwidth}{|l|l|X|}
    \hline
    \textbf{Metadata Name} & \textbf{Metadata Value} & \textbf{Description} \\
    \hline
    \textbf{Reference Id} & \texttt{BEC-5B685445} & \\
    \hline
    \textbf{Status} & Existing Feature; candidate CT3 placement & \\
    \hline
    \textbf{Ancestry} &  & \\
    \hline
    \textbf{Source} & FEAT-HSO-14 & \\
    \hline
    \textbf{Enabled By} &  & \\
    \hline
    \textbf{Supports} &  & \\
    \hline
    \textbf{Derived From} &  & \\
    \hline
    \textbf{Reviewed Status} & Confirm existing Feature meaning and proposed CT3 ownership. & \\
    \hline
  \end{tabularx}
\end{table}

Intent. Mobile Clinical Equipment Tracking provides a distinct enabling behaviour in the context of Specialized Equipment Coordination. This is an existing named Feature; its placement under the candidate CT3 remains subject to review.

\noindent\textbf{Feature boundary.} This behaviour is scoped to Specialized Equipment Coordination; an identically named Feature elsewhere may carry different business meaning.\par
\featuretable{
Function & \texttt{FN-A0F1E0D7} & Track Mobile Clinical Asset Operational State & Catalogue \\
Service & \texttt{SV-9608218D} & Clinical Resource Availability Query & Linked to preceding Function in source \\
}
\subsubsection{Feature - Context Establishment}\label{cap:bec-e7e5a452}
\vspace{-1em}
\begin{table}[htbp]
  \centering
  \footnotesize
  \begin{tabularx}{1\textwidth}{|l|l|X|}
    \hline
    \textbf{Metadata Name} & \textbf{Metadata Value} & \textbf{Description} \\
    \hline
    \textbf{Reference Id} & \texttt{BEC-E7E5A452} & \\
    \hline
    \textbf{Status} & Generated candidate & \\
    \hline
    \textbf{Ancestry} &  & \\
    \hline
    \textbf{Source} & Semantic-pattern candidate & \\
    \hline
    \textbf{Enabled By} &  & \\
    \hline
    \textbf{Supports} &  & \\
    \hline
    \textbf{Derived From} &  & \\
    \hline
    \textbf{Reviewed Status} & Generated candidate: keep / modify / relocate / reject; verify whether independently testable. & \\
    \hline
  \end{tabularx}
\end{table}

Intent. Context Establishment is proposed as a discrete, testable enabling behaviour for Specialized Equipment Coordination. It is included for architectural review, not as an established implementation obligation.

\noindent\textbf{Feature boundary.} This behaviour is scoped to Specialized Equipment Coordination; an identically named Feature elsewhere may carry different business meaning.\par
\featuretable{
--- & --- & \emph{Not yet captured} & \emph{Complete during review} \\
}
\subsubsection{Feature - Eligibility Validation}\label{cap:bec-aa245ff0}
\vspace{-1em}
\begin{table}[htbp]
  \centering
  \footnotesize
  \begin{tabularx}{1\textwidth}{|l|l|X|}
    \hline
    \textbf{Metadata Name} & \textbf{Metadata Value} & \textbf{Description} \\
    \hline
    \textbf{Reference Id} & \texttt{BEC-AA245FF0} & \\
    \hline
    \textbf{Status} & Generated candidate & \\
    \hline
    \textbf{Ancestry} &  & \\
    \hline
    \textbf{Source} & Semantic-pattern candidate & \\
    \hline
    \textbf{Enabled By} &  & \\
    \hline
    \textbf{Supports} &  & \\
    \hline
    \textbf{Derived From} &  & \\
    \hline
    \textbf{Reviewed Status} & Generated candidate: keep / modify / relocate / reject; verify whether independently testable. & \\
    \hline
  \end{tabularx}
\end{table}

Intent. Eligibility Validation is proposed as a discrete, testable enabling behaviour for Specialized Equipment Coordination. It is included for architectural review, not as an established implementation obligation.

\noindent\textbf{Feature boundary.} This behaviour is scoped to Specialized Equipment Coordination; an identically named Feature elsewhere may carry different business meaning.\par
\featuretable{
--- & --- & \emph{Not yet captured} & \emph{Complete during review} \\
}
\subsubsection{Feature - Assignment \& Routing}\label{cap:bec-a870eb50}
\vspace{-1em}
\begin{table}[htbp]
  \centering
  \footnotesize
  \begin{tabularx}{1\textwidth}{|l|l|X|}
    \hline
    \textbf{Metadata Name} & \textbf{Metadata Value} & \textbf{Description} \\
    \hline
    \textbf{Reference Id} & \texttt{BEC-A870EB50} & \\
    \hline
    \textbf{Status} & Generated candidate & \\
    \hline
    \textbf{Ancestry} &  & \\
    \hline
    \textbf{Source} & Semantic-pattern candidate & \\
    \hline
    \textbf{Enabled By} &  & \\
    \hline
    \textbf{Supports} &  & \\
    \hline
    \textbf{Derived From} &  & \\
    \hline
    \textbf{Reviewed Status} & Generated candidate: keep / modify / relocate / reject; verify whether independently testable. & \\
    \hline
  \end{tabularx}
\end{table}

Intent. Assignment \& Routing is proposed as a discrete, testable enabling behaviour for Specialized Equipment Coordination. It is included for architectural review, not as an established implementation obligation.

\noindent\textbf{Feature boundary.} This behaviour is scoped to Specialized Equipment Coordination; an identically named Feature elsewhere may carry different business meaning.\par
\featuretable{
--- & --- & \emph{Not yet captured} & \emph{Complete during review} \\
}
\subsubsection{Feature - State Tracking}\label{cap:bec-f1b486e3}
\vspace{-1em}
\begin{table}[htbp]
  \centering
  \footnotesize
  \begin{tabularx}{1\textwidth}{|l|l|X|}
    \hline
    \textbf{Metadata Name} & \textbf{Metadata Value} & \textbf{Description} \\
    \hline
    \textbf{Reference Id} & \texttt{BEC-F1B486E3} & \\
    \hline
    \textbf{Status} & Generated candidate & \\
    \hline
    \textbf{Ancestry} &  & \\
    \hline
    \textbf{Source} & Semantic-pattern candidate & \\
    \hline
    \textbf{Enabled By} &  & \\
    \hline
    \textbf{Supports} &  & \\
    \hline
    \textbf{Derived From} &  & \\
    \hline
    \textbf{Reviewed Status} & Generated candidate: keep / modify / relocate / reject; verify whether independently testable. & \\
    \hline
  \end{tabularx}
\end{table}

Intent. State Tracking is proposed as a discrete, testable enabling behaviour for Specialized Equipment Coordination. It is included for architectural review, not as an established implementation obligation.

\noindent\textbf{Feature boundary.} This behaviour is scoped to Specialized Equipment Coordination; an identically named Feature elsewhere may carry different business meaning.\par
\featuretable{
--- & --- & \emph{Not yet captured} & \emph{Complete during review} \\
}
\subsubsection{Feature - Exception Management}\label{cap:bec-99e2fa68}
\vspace{-1em}
\begin{table}[htbp]
  \centering
  \footnotesize
  \begin{tabularx}{1\textwidth}{|l|l|X|}
    \hline
    \textbf{Metadata Name} & \textbf{Metadata Value} & \textbf{Description} \\
    \hline
    \textbf{Reference Id} & \texttt{BEC-99E2FA68} & \\
    \hline
    \textbf{Status} & Generated candidate & \\
    \hline
    \textbf{Ancestry} &  & \\
    \hline
    \textbf{Source} & Semantic-pattern candidate & \\
    \hline
    \textbf{Enabled By} &  & \\
    \hline
    \textbf{Supports} &  & \\
    \hline
    \textbf{Derived From} &  & \\
    \hline
    \textbf{Reviewed Status} & Generated candidate: keep / modify / relocate / reject; verify whether independently testable. & \\
    \hline
  \end{tabularx}
\end{table}

Intent. Exception Management is proposed as a discrete, testable enabling behaviour for Specialized Equipment Coordination. It is included for architectural review, not as an established implementation obligation.

\noindent\textbf{Feature boundary.} This behaviour is scoped to Specialized Equipment Coordination; an identically named Feature elsewhere may carry different business meaning.\par
\featuretable{
--- & --- & \emph{Not yet captured} & \emph{Complete during review} \\
}
\subsubsection{Feature - Outcome Reporting}\label{cap:bec-4d2a47ac}
\vspace{-1em}
\begin{table}[htbp]
  \centering
  \footnotesize
  \begin{tabularx}{1\textwidth}{|l|l|X|}
    \hline
    \textbf{Metadata Name} & \textbf{Metadata Value} & \textbf{Description} \\
    \hline
    \textbf{Reference Id} & \texttt{BEC-4D2A47AC} & \\
    \hline
    \textbf{Status} & Generated candidate & \\
    \hline
    \textbf{Ancestry} &  & \\
    \hline
    \textbf{Source} & Semantic-pattern candidate & \\
    \hline
    \textbf{Enabled By} &  & \\
    \hline
    \textbf{Supports} &  & \\
    \hline
    \textbf{Derived From} &  & \\
    \hline
    \textbf{Reviewed Status} & Generated candidate: keep / modify / relocate / reject; verify whether independently testable. & \\
    \hline
  \end{tabularx}
\end{table}

Intent. Outcome Reporting is proposed as a discrete, testable enabling behaviour for Specialized Equipment Coordination. It is included for architectural review, not as an established implementation obligation.

\noindent\textbf{Feature boundary.} This behaviour is scoped to Specialized Equipment Coordination; an identically named Feature elsewhere may carry different business meaning.\par
\featuretable{
--- & --- & \emph{Not yet captured} & \emph{Complete during review} \\
}

\section{CT2 - Service Capacity Management}\label{cap:bec-db7874af}
\vspace{-0.5em}
\begin{table}[htbp]
  \centering
  \footnotesize
  \begin{tabularx}{1\textwidth}{|l|l|X|}
    \hline
    \multicolumn{3}{c}{\textbf{Capability Metadata}} \\
    \hline
    \textbf{Name} & \textbf{Value} & \textbf{Comment and/or Description} \\
    \hline
    \textbf{Reference Id} & \texttt{BEC-DB7874AF} & \\
    \hline
    \textbf{Status} & Candidate ancestry & \\
    \hline
    \textbf{Ancestry} &  & \\
    \hline
    \textbf{Source} &  & \\
    \hline
    \textbf{Enabled By} &  & \\
    \hline
    \textbf{Supports} &  & \\
    \hline
    \textbf{Derived From} &  & \\
    \hline
    \textbf{Reviewed Status} &  & \\
    \hline
  \end{tabularx}
\end{table}

\begin{table}[htbp]
  \centering
  \footnotesize
  \begin{tabularx}{1\textwidth}{|l|l|X|}
    \hline
    \multicolumn{3}{c}{\textbf{Contained Capabilities}} \\
    \hline
    \textbf{Capability Name} & \textbf{Capability Id} & \textbf{Comment} \\
    \hline
    Facility Demand \& Capacity Telemetry & \texttt{cap:bec-d90fab4c} (at \ref{cap:bec-d90fab4c}) & \\
    \hline
  \end{tabularx}
\end{table}

Intent. Service Capacity Management is a candidate enabling capability within Health Service Operations Enablement. It frames the reusable system-supported behaviours needed to support this healthcare context without prescribing a product, interface, technology or deployment boundary.

Boundary. This capability supports the business activity and does not inherit the authority or accountability of the Business Capability it enables. Its child Features are review candidates until ancestry and semantics are adjudicated.

\subsection{CT3 - Facility Demand \& Capacity Telemetry}\label{cap:bec-d90fab4c}
\vspace{-0.5em}
\begin{table}[htbp]
  \centering
  \footnotesize
  \begin{tabularx}{1\textwidth}{|l|l|X|}
    \hline
    \multicolumn{3}{c}{\textbf{Capability Metadata}} \\
    \hline
    \textbf{Name} & \textbf{Value} & \textbf{Comment and/or Description} \\
    \hline
    \textbf{Reference Id} & \texttt{BEC-D90FAB4C} & \\
    \hline
    \textbf{Status} & Candidate ancestry & \\
    \hline
    \textbf{Ancestry} &  & \\
    \hline
    \textbf{Source} &  & \\
    \hline
    \textbf{Enabled By} &  & \\
    \hline
    \textbf{Supports} &  & \\
    \hline
    \textbf{Derived From} &  & \\
    \hline
    \textbf{Reviewed Status} &  & \\
    \hline
  \end{tabularx}
\end{table}

\begin{table}[htbp]
  \centering
  \footnotesize
  \begin{tabularx}{1\textwidth}{|l|l|X|}
    \hline
    \multicolumn{3}{c}{\textbf{Contained Features}} \\
    \hline
    \textbf{Feature Name} & \textbf{Feature Id} & \textbf{Comment} \\
    \hline
    Operational Capacity Metric Aggregation & \texttt{cap:bec-a5cde7aa} (at \ref{cap:bec-a5cde7aa}) & \\
    \hline
    Context Establishment & \texttt{cap:bec-3559f81b} (at \ref{cap:bec-3559f81b}) & \\
    \hline
    State Monitoring & \texttt{cap:bec-f78fc8d3} (at \ref{cap:bec-f78fc8d3}) & \\
    \hline
    Exception Management & \texttt{cap:bec-0203e495} (at \ref{cap:bec-0203e495}) & \\
    \hline
    Information Retrieval & \texttt{cap:bec-a2be39bf} (at \ref{cap:bec-a2be39bf}) & \\
    \hline
  \end{tabularx}
\end{table}

Intent. Facility Demand \& Capacity Telemetry is a candidate enabling capability within Service Capacity Management. It frames the reusable system-supported behaviours needed to support this healthcare context without prescribing a product, interface, technology or deployment boundary.

Boundary. This capability supports the business activity and does not inherit the authority or accountability of the Business Capability it enables. Its child Features are review candidates until ancestry and semantics are adjudicated.

\subsubsection{Feature - Operational Capacity Metric Aggregation}\label{cap:bec-a5cde7aa}
\vspace{-1em}
\begin{table}[htbp]
  \centering
  \footnotesize
  \begin{tabularx}{1\textwidth}{|l|l|X|}
    \hline
    \textbf{Metadata Name} & \textbf{Metadata Value} & \textbf{Description} \\
    \hline
    \textbf{Reference Id} & \texttt{BEC-A5CDE7AA} & \\
    \hline
    \textbf{Status} & Existing Feature; candidate CT3 placement & \\
    \hline
    \textbf{Ancestry} &  & \\
    \hline
    \textbf{Source} & FEAT-HSO-15 & \\
    \hline
    \textbf{Enabled By} &  & \\
    \hline
    \textbf{Supports} &  & \\
    \hline
    \textbf{Derived From} &  & \\
    \hline
    \textbf{Reviewed Status} & Confirm existing Feature meaning and proposed CT3 ownership. & \\
    \hline
  \end{tabularx}
\end{table}

Intent. Operational Capacity Metric Aggregation provides a distinct enabling behaviour in the context of Facility Demand \& Capacity Telemetry. This is an existing named Feature; its placement under the candidate CT3 remains subject to review.

\noindent\textbf{Feature boundary.} This behaviour is scoped to Facility Demand \& Capacity Telemetry; an identically named Feature elsewhere may carry different business meaning.\par
\featuretable{
Function & \texttt{FN-5BC42C26} & Aggregate Service Capacity Telemetry & Catalogue \\
Service & \texttt{SV-FDA7E01C} & Regional Health System Capacity Query & Linked to preceding Function in source \\
}
\subsubsection{Feature - Context Establishment}\label{cap:bec-3559f81b}
\vspace{-1em}
\begin{table}[htbp]
  \centering
  \footnotesize
  \begin{tabularx}{1\textwidth}{|l|l|X|}
    \hline
    \textbf{Metadata Name} & \textbf{Metadata Value} & \textbf{Description} \\
    \hline
    \textbf{Reference Id} & \texttt{BEC-3559F81B} & \\
    \hline
    \textbf{Status} & Generated candidate & \\
    \hline
    \textbf{Ancestry} &  & \\
    \hline
    \textbf{Source} & Semantic-pattern candidate & \\
    \hline
    \textbf{Enabled By} &  & \\
    \hline
    \textbf{Supports} &  & \\
    \hline
    \textbf{Derived From} &  & \\
    \hline
    \textbf{Reviewed Status} & Generated candidate: keep / modify / relocate / reject; verify whether independently testable. & \\
    \hline
  \end{tabularx}
\end{table}

Intent. Context Establishment is proposed as a discrete, testable enabling behaviour for Facility Demand \& Capacity Telemetry. It is included for architectural review, not as an established implementation obligation.

\noindent\textbf{Feature boundary.} This behaviour is scoped to Facility Demand \& Capacity Telemetry; an identically named Feature elsewhere may carry different business meaning.\par
\featuretable{
--- & --- & \emph{Not yet captured} & \emph{Complete during review} \\
}
\subsubsection{Feature - State Monitoring}\label{cap:bec-f78fc8d3}
\vspace{-1em}
\begin{table}[htbp]
  \centering
  \footnotesize
  \begin{tabularx}{1\textwidth}{|l|l|X|}
    \hline
    \textbf{Metadata Name} & \textbf{Metadata Value} & \textbf{Description} \\
    \hline
    \textbf{Reference Id} & \texttt{BEC-F78FC8D3} & \\
    \hline
    \textbf{Status} & Generated candidate & \\
    \hline
    \textbf{Ancestry} &  & \\
    \hline
    \textbf{Source} & Semantic-pattern candidate & \\
    \hline
    \textbf{Enabled By} &  & \\
    \hline
    \textbf{Supports} &  & \\
    \hline
    \textbf{Derived From} &  & \\
    \hline
    \textbf{Reviewed Status} & Generated candidate: keep / modify / relocate / reject; verify whether independently testable. & \\
    \hline
  \end{tabularx}
\end{table}

Intent. State Monitoring is proposed as a discrete, testable enabling behaviour for Facility Demand \& Capacity Telemetry. It is included for architectural review, not as an established implementation obligation.

\noindent\textbf{Feature boundary.} This behaviour is scoped to Facility Demand \& Capacity Telemetry; an identically named Feature elsewhere may carry different business meaning.\par
\featuretable{
--- & --- & \emph{Not yet captured} & \emph{Complete during review} \\
}
\subsubsection{Feature - Exception Management}\label{cap:bec-0203e495}
\vspace{-1em}
\begin{table}[htbp]
  \centering
  \footnotesize
  \begin{tabularx}{1\textwidth}{|l|l|X|}
    \hline
    \textbf{Metadata Name} & \textbf{Metadata Value} & \textbf{Description} \\
    \hline
    \textbf{Reference Id} & \texttt{BEC-0203E495} & \\
    \hline
    \textbf{Status} & Generated candidate & \\
    \hline
    \textbf{Ancestry} &  & \\
    \hline
    \textbf{Source} & Semantic-pattern candidate & \\
    \hline
    \textbf{Enabled By} &  & \\
    \hline
    \textbf{Supports} &  & \\
    \hline
    \textbf{Derived From} &  & \\
    \hline
    \textbf{Reviewed Status} & Generated candidate: keep / modify / relocate / reject; verify whether independently testable. & \\
    \hline
  \end{tabularx}
\end{table}

Intent. Exception Management is proposed as a discrete, testable enabling behaviour for Facility Demand \& Capacity Telemetry. It is included for architectural review, not as an established implementation obligation.

\noindent\textbf{Feature boundary.} This behaviour is scoped to Facility Demand \& Capacity Telemetry; an identically named Feature elsewhere may carry different business meaning.\par
\featuretable{
--- & --- & \emph{Not yet captured} & \emph{Complete during review} \\
}
\subsubsection{Feature - Information Retrieval}\label{cap:bec-a2be39bf}
\vspace{-1em}
\begin{table}[htbp]
  \centering
  \footnotesize
  \begin{tabularx}{1\textwidth}{|l|l|X|}
    \hline
    \textbf{Metadata Name} & \textbf{Metadata Value} & \textbf{Description} \\
    \hline
    \textbf{Reference Id} & \texttt{BEC-A2BE39BF} & \\
    \hline
    \textbf{Status} & Generated candidate & \\
    \hline
    \textbf{Ancestry} &  & \\
    \hline
    \textbf{Source} & Semantic-pattern candidate & \\
    \hline
    \textbf{Enabled By} &  & \\
    \hline
    \textbf{Supports} &  & \\
    \hline
    \textbf{Derived From} &  & \\
    \hline
    \textbf{Reviewed Status} & Generated candidate: keep / modify / relocate / reject; verify whether independently testable. & \\
    \hline
  \end{tabularx}
\end{table}

Intent. Information Retrieval is proposed as a discrete, testable enabling behaviour for Facility Demand \& Capacity Telemetry. It is included for architectural review, not as an established implementation obligation.

\noindent\textbf{Feature boundary.} This behaviour is scoped to Facility Demand \& Capacity Telemetry; an identically named Feature elsewhere may carry different business meaning.\par
\featuretable{
--- & --- & \emph{Not yet captured} & \emph{Complete during review} \\
}

\section{CT2 - Clinical Workforce Management}\label{cap:bec-20c1e1b7}
\vspace{-0.5em}
\begin{table}[htbp]
  \centering
  \footnotesize
  \begin{tabularx}{1\textwidth}{|l|l|X|}
    \hline
    \multicolumn{3}{c}{\textbf{Capability Metadata}} \\
    \hline
    \textbf{Name} & \textbf{Value} & \textbf{Comment and/or Description} \\
    \hline
    \textbf{Reference Id} & \texttt{BEC-20C1E1B7} & \\
    \hline
    \textbf{Status} & Candidate ancestry & \\
    \hline
    \textbf{Ancestry} &  & \\
    \hline
    \textbf{Source} &  & \\
    \hline
    \textbf{Enabled By} &  & \\
    \hline
    \textbf{Supports} &  & \\
    \hline
    \textbf{Derived From} &  & \\
    \hline
    \textbf{Reviewed Status} &  & \\
    \hline
  \end{tabularx}
\end{table}

\begin{table}[htbp]
  \centering
  \footnotesize
  \begin{tabularx}{1\textwidth}{|l|l|X|}
    \hline
    \multicolumn{3}{c}{\textbf{Contained Capabilities}} \\
    \hline
    \textbf{Capability Name} & \textbf{Capability Id} & \textbf{Comment} \\
    \hline
    \multicolumn{3}{|l|}{\emph{No contained capabilities are currently identified.}} \\
    \hline
  \end{tabularx}
\end{table}

Intent. Clinical Workforce Management is a candidate enabling capability within Health Service Operations Enablement. It frames the reusable system-supported behaviours needed to support this healthcare context without prescribing a product, interface, technology or deployment boundary.

Boundary. This capability supports the business activity and does not inherit the authority or accountability of the Business Capability it enables. Its child Features are review candidates until ancestry and semantics are adjudicated.

\section{CT2 - Clinical Rostering}\label{cap:bec-f8c1b2df}
\vspace{-0.5em}
\begin{table}[htbp]
  \centering
  \footnotesize
  \begin{tabularx}{1\textwidth}{|l|l|X|}
    \hline
    \multicolumn{3}{c}{\textbf{Capability Metadata}} \\
    \hline
    \textbf{Name} & \textbf{Value} & \textbf{Comment and/or Description} \\
    \hline
    \textbf{Reference Id} & \texttt{BEC-F8C1B2DF} & \\
    \hline
    \textbf{Status} & Candidate ancestry & \\
    \hline
    \textbf{Ancestry} &  & \\
    \hline
    \textbf{Source} &  & \\
    \hline
    \textbf{Enabled By} &  & \\
    \hline
    \textbf{Supports} &  & \\
    \hline
    \textbf{Derived From} &  & \\
    \hline
    \textbf{Reviewed Status} &  & \\
    \hline
  \end{tabularx}
\end{table}

\begin{table}[htbp]
  \centering
  \footnotesize
  \begin{tabularx}{1\textwidth}{|l|l|X|}
    \hline
    \multicolumn{3}{c}{\textbf{Contained Capabilities}} \\
    \hline
    \textbf{Capability Name} & \textbf{Capability Id} & \textbf{Comment} \\
    \hline
    \multicolumn{3}{|l|}{\emph{No contained capabilities are currently identified.}} \\
    \hline
  \end{tabularx}
\end{table}

Intent. Clinical Rostering is a candidate enabling capability within Health Service Operations Enablement. It frames the reusable system-supported behaviours needed to support this healthcare context without prescribing a product, interface, technology or deployment boundary.

Boundary. This capability supports the business activity and does not inherit the authority or accountability of the Business Capability it enables. Its child Features are review candidates until ancestry and semantics are adjudicated.

\section{CT2 - Mobile Staff Management}\label{cap:bec-71740710}
\vspace{-0.5em}
\begin{table}[htbp]
  \centering
  \footnotesize
  \begin{tabularx}{1\textwidth}{|l|l|X|}
    \hline
    \multicolumn{3}{c}{\textbf{Capability Metadata}} \\
    \hline
    \textbf{Name} & \textbf{Value} & \textbf{Comment and/or Description} \\
    \hline
    \textbf{Reference Id} & \texttt{BEC-71740710} & \\
    \hline
    \textbf{Status} & Candidate ancestry & \\
    \hline
    \textbf{Ancestry} &  & \\
    \hline
    \textbf{Source} &  & \\
    \hline
    \textbf{Enabled By} &  & \\
    \hline
    \textbf{Supports} &  & \\
    \hline
    \textbf{Derived From} &  & \\
    \hline
    \textbf{Reviewed Status} &  & \\
    \hline
  \end{tabularx}
\end{table}

\begin{table}[htbp]
  \centering
  \footnotesize
  \begin{tabularx}{1\textwidth}{|l|l|X|}
    \hline
    \multicolumn{3}{c}{\textbf{Contained Capabilities}} \\
    \hline
    \textbf{Capability Name} & \textbf{Capability Id} & \textbf{Comment} \\
    \hline
    Mobile Worker Dispatch \& Safety & \texttt{cap:bec-ec43c4cb} (at \ref{cap:bec-ec43c4cb}) & \\
    \hline
  \end{tabularx}
\end{table}

Intent. Mobile Staff Management is a candidate enabling capability within Health Service Operations Enablement. It frames the reusable system-supported behaviours needed to support this healthcare context without prescribing a product, interface, technology or deployment boundary.

Boundary. This capability supports the business activity and does not inherit the authority or accountability of the Business Capability it enables. Its child Features are review candidates until ancestry and semantics are adjudicated.

\subsection{CT3 - Mobile Worker Dispatch \& Safety}\label{cap:bec-ec43c4cb}
\vspace{-0.5em}
\begin{table}[htbp]
  \centering
  \footnotesize
  \begin{tabularx}{1\textwidth}{|l|l|X|}
    \hline
    \multicolumn{3}{c}{\textbf{Capability Metadata}} \\
    \hline
    \textbf{Name} & \textbf{Value} & \textbf{Comment and/or Description} \\
    \hline
    \textbf{Reference Id} & \texttt{BEC-EC43C4CB} & \\
    \hline
    \textbf{Status} & Candidate ancestry & \\
    \hline
    \textbf{Ancestry} &  & \\
    \hline
    \textbf{Source} &  & \\
    \hline
    \textbf{Enabled By} &  & \\
    \hline
    \textbf{Supports} &  & \\
    \hline
    \textbf{Derived From} &  & \\
    \hline
    \textbf{Reviewed Status} &  & \\
    \hline
  \end{tabularx}
\end{table}

\begin{table}[htbp]
  \centering
  \footnotesize
  \begin{tabularx}{1\textwidth}{|l|l|X|}
    \hline
    \multicolumn{3}{c}{\textbf{Contained Features}} \\
    \hline
    \textbf{Feature Name} & \textbf{Feature Id} & \textbf{Comment} \\
    \hline
    Mobile Worker Task Dispatch & \texttt{cap:bec-d243c132} (at \ref{cap:bec-d243c132}) & \\
    \hline
    Mobile Visit Status Tracking & \texttt{cap:bec-caff9e49} (at \ref{cap:bec-caff9e49}) & \\
    \hline
    Context Establishment & \texttt{cap:bec-02254531} (at \ref{cap:bec-02254531}) & \\
    \hline
    Eligibility Validation & \texttt{cap:bec-04274fdd} (at \ref{cap:bec-04274fdd}) & \\
    \hline
    Assignment \& Routing & \texttt{cap:bec-bfbf8974} (at \ref{cap:bec-bfbf8974}) & \\
    \hline
    State Tracking & \texttt{cap:bec-e87b589f} (at \ref{cap:bec-e87b589f}) & \\
    \hline
    Exception Management & \texttt{cap:bec-7fd6c955} (at \ref{cap:bec-7fd6c955}) & \\
    \hline
    Outcome Reporting & \texttt{cap:bec-d56786a8} (at \ref{cap:bec-d56786a8}) & \\
    \hline
  \end{tabularx}
\end{table}

Intent. Mobile Worker Dispatch \& Safety is a candidate enabling capability within Mobile Staff Management. It frames the reusable system-supported behaviours needed to support this healthcare context without prescribing a product, interface, technology or deployment boundary.

Boundary. This capability supports the business activity and does not inherit the authority or accountability of the Business Capability it enables. Its child Features are review candidates until ancestry and semantics are adjudicated.

\subsubsection{Feature - Mobile Worker Task Dispatch}\label{cap:bec-d243c132}
\vspace{-1em}
\begin{table}[htbp]
  \centering
  \footnotesize
  \begin{tabularx}{1\textwidth}{|l|l|X|}
    \hline
    \textbf{Metadata Name} & \textbf{Metadata Value} & \textbf{Description} \\
    \hline
    \textbf{Reference Id} & \texttt{BEC-D243C132} & \\
    \hline
    \textbf{Status} & Existing Feature; candidate CT3 placement & \\
    \hline
    \textbf{Ancestry} &  & \\
    \hline
    \textbf{Source} & FEAT-HSO-16 & \\
    \hline
    \textbf{Enabled By} &  & \\
    \hline
    \textbf{Supports} &  & \\
    \hline
    \textbf{Derived From} &  & \\
    \hline
    \textbf{Reviewed Status} & Confirm existing Feature meaning and proposed CT3 ownership. & \\
    \hline
  \end{tabularx}
\end{table}

Intent. Mobile Worker Task Dispatch provides a distinct enabling behaviour in the context of Mobile Worker Dispatch \& Safety. This is an existing named Feature; its placement under the candidate CT3 remains subject to review.

\noindent\textbf{Feature boundary.} This behaviour is scoped to Mobile Worker Dispatch \& Safety; an identically named Feature elsewhere may carry different business meaning.\par
\featuretable{
Function & \texttt{FN-B0995B22} & Track Operational Staff Presence & Catalogue \\
Service & \texttt{SV-8B1ED419} & Staff Presence Telemetry Query & Linked to preceding Function in source \\
}
\subsubsection{Feature - Mobile Visit Status Tracking}\label{cap:bec-caff9e49}
\vspace{-1em}
\begin{table}[htbp]
  \centering
  \footnotesize
  \begin{tabularx}{1\textwidth}{|l|l|X|}
    \hline
    \textbf{Metadata Name} & \textbf{Metadata Value} & \textbf{Description} \\
    \hline
    \textbf{Reference Id} & \texttt{BEC-CAFF9E49} & \\
    \hline
    \textbf{Status} & Existing Feature; candidate CT3 placement & \\
    \hline
    \textbf{Ancestry} &  & \\
    \hline
    \textbf{Source} & FEAT-HSO-17 & \\
    \hline
    \textbf{Enabled By} &  & \\
    \hline
    \textbf{Supports} &  & \\
    \hline
    \textbf{Derived From} &  & \\
    \hline
    \textbf{Reviewed Status} & Confirm existing Feature meaning and proposed CT3 ownership. & \\
    \hline
  \end{tabularx}
\end{table}

Intent. Mobile Visit Status Tracking provides a distinct enabling behaviour in the context of Mobile Worker Dispatch \& Safety. This is an existing named Feature; its placement under the candidate CT3 remains subject to review.

\noindent\textbf{Feature boundary.} This behaviour is scoped to Mobile Worker Dispatch \& Safety; an identically named Feature elsewhere may carry different business meaning.\par
\featuretable{
Function & \texttt{FN-790627CC} & Receive Mobile Visit Progression & Catalogue \\
}
\subsubsection{Feature - Context Establishment}\label{cap:bec-02254531}
\vspace{-1em}
\begin{table}[htbp]
  \centering
  \footnotesize
  \begin{tabularx}{1\textwidth}{|l|l|X|}
    \hline
    \textbf{Metadata Name} & \textbf{Metadata Value} & \textbf{Description} \\
    \hline
    \textbf{Reference Id} & \texttt{BEC-02254531} & \\
    \hline
    \textbf{Status} & Generated candidate & \\
    \hline
    \textbf{Ancestry} &  & \\
    \hline
    \textbf{Source} & Semantic-pattern candidate & \\
    \hline
    \textbf{Enabled By} &  & \\
    \hline
    \textbf{Supports} &  & \\
    \hline
    \textbf{Derived From} &  & \\
    \hline
    \textbf{Reviewed Status} & Generated candidate: keep / modify / relocate / reject; verify whether independently testable. & \\
    \hline
  \end{tabularx}
\end{table}

Intent. Context Establishment is proposed as a discrete, testable enabling behaviour for Mobile Worker Dispatch \& Safety. It is included for architectural review, not as an established implementation obligation.

\noindent\textbf{Feature boundary.} This behaviour is scoped to Mobile Worker Dispatch \& Safety; an identically named Feature elsewhere may carry different business meaning.\par
\featuretable{
--- & --- & \emph{Not yet captured} & \emph{Complete during review} \\
}
\subsubsection{Feature - Eligibility Validation}\label{cap:bec-04274fdd}
\vspace{-1em}
\begin{table}[htbp]
  \centering
  \footnotesize
  \begin{tabularx}{1\textwidth}{|l|l|X|}
    \hline
    \textbf{Metadata Name} & \textbf{Metadata Value} & \textbf{Description} \\
    \hline
    \textbf{Reference Id} & \texttt{BEC-04274FDD} & \\
    \hline
    \textbf{Status} & Generated candidate & \\
    \hline
    \textbf{Ancestry} &  & \\
    \hline
    \textbf{Source} & Semantic-pattern candidate & \\
    \hline
    \textbf{Enabled By} &  & \\
    \hline
    \textbf{Supports} &  & \\
    \hline
    \textbf{Derived From} &  & \\
    \hline
    \textbf{Reviewed Status} & Generated candidate: keep / modify / relocate / reject; verify whether independently testable. & \\
    \hline
  \end{tabularx}
\end{table}

Intent. Eligibility Validation is proposed as a discrete, testable enabling behaviour for Mobile Worker Dispatch \& Safety. It is included for architectural review, not as an established implementation obligation.

\noindent\textbf{Feature boundary.} This behaviour is scoped to Mobile Worker Dispatch \& Safety; an identically named Feature elsewhere may carry different business meaning.\par
\featuretable{
--- & --- & \emph{Not yet captured} & \emph{Complete during review} \\
}
\subsubsection{Feature - Assignment \& Routing}\label{cap:bec-bfbf8974}
\vspace{-1em}
\begin{table}[htbp]
  \centering
  \footnotesize
  \begin{tabularx}{1\textwidth}{|l|l|X|}
    \hline
    \textbf{Metadata Name} & \textbf{Metadata Value} & \textbf{Description} \\
    \hline
    \textbf{Reference Id} & \texttt{BEC-BFBF8974} & \\
    \hline
    \textbf{Status} & Generated candidate & \\
    \hline
    \textbf{Ancestry} &  & \\
    \hline
    \textbf{Source} & Semantic-pattern candidate & \\
    \hline
    \textbf{Enabled By} &  & \\
    \hline
    \textbf{Supports} &  & \\
    \hline
    \textbf{Derived From} &  & \\
    \hline
    \textbf{Reviewed Status} & Generated candidate: keep / modify / relocate / reject; verify whether independently testable. & \\
    \hline
  \end{tabularx}
\end{table}

Intent. Assignment \& Routing is proposed as a discrete, testable enabling behaviour for Mobile Worker Dispatch \& Safety. It is included for architectural review, not as an established implementation obligation.

\noindent\textbf{Feature boundary.} This behaviour is scoped to Mobile Worker Dispatch \& Safety; an identically named Feature elsewhere may carry different business meaning.\par
\featuretable{
--- & --- & \emph{Not yet captured} & \emph{Complete during review} \\
}
\subsubsection{Feature - State Tracking}\label{cap:bec-e87b589f}
\vspace{-1em}
\begin{table}[htbp]
  \centering
  \footnotesize
  \begin{tabularx}{1\textwidth}{|l|l|X|}
    \hline
    \textbf{Metadata Name} & \textbf{Metadata Value} & \textbf{Description} \\
    \hline
    \textbf{Reference Id} & \texttt{BEC-E87B589F} & \\
    \hline
    \textbf{Status} & Generated candidate & \\
    \hline
    \textbf{Ancestry} &  & \\
    \hline
    \textbf{Source} & Semantic-pattern candidate & \\
    \hline
    \textbf{Enabled By} &  & \\
    \hline
    \textbf{Supports} &  & \\
    \hline
    \textbf{Derived From} &  & \\
    \hline
    \textbf{Reviewed Status} & Generated candidate: keep / modify / relocate / reject; verify whether independently testable. & \\
    \hline
  \end{tabularx}
\end{table}

Intent. State Tracking is proposed as a discrete, testable enabling behaviour for Mobile Worker Dispatch \& Safety. It is included for architectural review, not as an established implementation obligation.

\noindent\textbf{Feature boundary.} This behaviour is scoped to Mobile Worker Dispatch \& Safety; an identically named Feature elsewhere may carry different business meaning.\par
\featuretable{
--- & --- & \emph{Not yet captured} & \emph{Complete during review} \\
}
\subsubsection{Feature - Exception Management}\label{cap:bec-7fd6c955}
\vspace{-1em}
\begin{table}[htbp]
  \centering
  \footnotesize
  \begin{tabularx}{1\textwidth}{|l|l|X|}
    \hline
    \textbf{Metadata Name} & \textbf{Metadata Value} & \textbf{Description} \\
    \hline
    \textbf{Reference Id} & \texttt{BEC-7FD6C955} & \\
    \hline
    \textbf{Status} & Generated candidate & \\
    \hline
    \textbf{Ancestry} &  & \\
    \hline
    \textbf{Source} & Semantic-pattern candidate & \\
    \hline
    \textbf{Enabled By} &  & \\
    \hline
    \textbf{Supports} &  & \\
    \hline
    \textbf{Derived From} &  & \\
    \hline
    \textbf{Reviewed Status} & Generated candidate: keep / modify / relocate / reject; verify whether independently testable. & \\
    \hline
  \end{tabularx}
\end{table}

Intent. Exception Management is proposed as a discrete, testable enabling behaviour for Mobile Worker Dispatch \& Safety. It is included for architectural review, not as an established implementation obligation.

\noindent\textbf{Feature boundary.} This behaviour is scoped to Mobile Worker Dispatch \& Safety; an identically named Feature elsewhere may carry different business meaning.\par
\featuretable{
--- & --- & \emph{Not yet captured} & \emph{Complete during review} \\
}
\subsubsection{Feature - Outcome Reporting}\label{cap:bec-d56786a8}
\vspace{-1em}
\begin{table}[htbp]
  \centering
  \footnotesize
  \begin{tabularx}{1\textwidth}{|l|l|X|}
    \hline
    \textbf{Metadata Name} & \textbf{Metadata Value} & \textbf{Description} \\
    \hline
    \textbf{Reference Id} & \texttt{BEC-D56786A8} & \\
    \hline
    \textbf{Status} & Generated candidate & \\
    \hline
    \textbf{Ancestry} &  & \\
    \hline
    \textbf{Source} & Semantic-pattern candidate & \\
    \hline
    \textbf{Enabled By} &  & \\
    \hline
    \textbf{Supports} &  & \\
    \hline
    \textbf{Derived From} &  & \\
    \hline
    \textbf{Reviewed Status} & Generated candidate: keep / modify / relocate / reject; verify whether independently testable. & \\
    \hline
  \end{tabularx}
\end{table}

Intent. Outcome Reporting is proposed as a discrete, testable enabling behaviour for Mobile Worker Dispatch \& Safety. It is included for architectural review, not as an established implementation obligation.

\noindent\textbf{Feature boundary.} This behaviour is scoped to Mobile Worker Dispatch \& Safety; an identically named Feature elsewhere may carry different business meaning.\par
\featuretable{
--- & --- & \emph{Not yet captured} & \emph{Complete during review} \\
}

\section{CT2 - On-Call Management}\label{cap:bec-aa196190}
\vspace{-0.5em}
\begin{table}[htbp]
  \centering
  \footnotesize
  \begin{tabularx}{1\textwidth}{|l|l|X|}
    \hline
    \multicolumn{3}{c}{\textbf{Capability Metadata}} \\
    \hline
    \textbf{Name} & \textbf{Value} & \textbf{Comment and/or Description} \\
    \hline
    \textbf{Reference Id} & \texttt{BEC-AA196190} & \\
    \hline
    \textbf{Status} & Candidate ancestry & \\
    \hline
    \textbf{Ancestry} &  & \\
    \hline
    \textbf{Source} &  & \\
    \hline
    \textbf{Enabled By} &  & \\
    \hline
    \textbf{Supports} &  & \\
    \hline
    \textbf{Derived From} &  & \\
    \hline
    \textbf{Reviewed Status} &  & \\
    \hline
  \end{tabularx}
\end{table}

\begin{table}[htbp]
  \centering
  \footnotesize
  \begin{tabularx}{1\textwidth}{|l|l|X|}
    \hline
    \multicolumn{3}{c}{\textbf{Contained Capabilities}} \\
    \hline
    \textbf{Capability Name} & \textbf{Capability Id} & \textbf{Comment} \\
    \hline
    On-Call Directory Resolution & \texttt{cap:bec-4be71cb3} (at \ref{cap:bec-4be71cb3}) & \\
    \hline
  \end{tabularx}
\end{table}

Intent. On-Call Management is a candidate enabling capability within Health Service Operations Enablement. It frames the reusable system-supported behaviours needed to support this healthcare context without prescribing a product, interface, technology or deployment boundary.

Boundary. This capability supports the business activity and does not inherit the authority or accountability of the Business Capability it enables. Its child Features are review candidates until ancestry and semantics are adjudicated.

\subsection{CT3 - On-Call Directory Resolution}\label{cap:bec-4be71cb3}
\vspace{-0.5em}
\begin{table}[htbp]
  \centering
  \footnotesize
  \begin{tabularx}{1\textwidth}{|l|l|X|}
    \hline
    \multicolumn{3}{c}{\textbf{Capability Metadata}} \\
    \hline
    \textbf{Name} & \textbf{Value} & \textbf{Comment and/or Description} \\
    \hline
    \textbf{Reference Id} & \texttt{BEC-4BE71CB3} & \\
    \hline
    \textbf{Status} & Candidate ancestry & \\
    \hline
    \textbf{Ancestry} &  & \\
    \hline
    \textbf{Source} &  & \\
    \hline
    \textbf{Enabled By} &  & \\
    \hline
    \textbf{Supports} &  & \\
    \hline
    \textbf{Derived From} &  & \\
    \hline
    \textbf{Reviewed Status} &  & \\
    \hline
  \end{tabularx}
\end{table}

\begin{table}[htbp]
  \centering
  \footnotesize
  \begin{tabularx}{1\textwidth}{|l|l|X|}
    \hline
    \multicolumn{3}{c}{\textbf{Contained Features}} \\
    \hline
    \textbf{Feature Name} & \textbf{Feature Id} & \textbf{Comment} \\
    \hline
    Active On-Call Provider Resolution & \texttt{cap:bec-722aba6e} (at \ref{cap:bec-722aba6e}) & \\
    \hline
    Context Establishment & \texttt{cap:bec-4cfa3e90} (at \ref{cap:bec-4cfa3e90}) & \\
    \hline
    State Monitoring & \texttt{cap:bec-703e878f} (at \ref{cap:bec-703e878f}) & \\
    \hline
    Exception Management & \texttt{cap:bec-e452abbd} (at \ref{cap:bec-e452abbd}) & \\
    \hline
    Information Retrieval & \texttt{cap:bec-783c1ed2} (at \ref{cap:bec-783c1ed2}) & \\
    \hline
  \end{tabularx}
\end{table}

Intent. On-Call Directory Resolution is a candidate enabling capability within On-Call Management. It frames the reusable system-supported behaviours needed to support this healthcare context without prescribing a product, interface, technology or deployment boundary.

Boundary. This capability supports the business activity and does not inherit the authority or accountability of the Business Capability it enables. Its child Features are review candidates until ancestry and semantics are adjudicated.

\subsubsection{Feature - Active On-Call Provider Resolution}\label{cap:bec-722aba6e}
\vspace{-1em}
\begin{table}[htbp]
  \centering
  \footnotesize
  \begin{tabularx}{1\textwidth}{|l|l|X|}
    \hline
    \textbf{Metadata Name} & \textbf{Metadata Value} & \textbf{Description} \\
    \hline
    \textbf{Reference Id} & \texttt{BEC-722ABA6E} & \\
    \hline
    \textbf{Status} & Existing Feature; candidate CT3 placement & \\
    \hline
    \textbf{Ancestry} &  & \\
    \hline
    \textbf{Source} & FEAT-HSO-18 & \\
    \hline
    \textbf{Enabled By} &  & \\
    \hline
    \textbf{Supports} &  & \\
    \hline
    \textbf{Derived From} &  & \\
    \hline
    \textbf{Reviewed Status} & Confirm existing Feature meaning and proposed CT3 ownership. & \\
    \hline
  \end{tabularx}
\end{table}

Intent. Active On-Call Provider Resolution provides a distinct enabling behaviour in the context of On-Call Directory Resolution. This is an existing named Feature; its placement under the candidate CT3 remains subject to review.

\noindent\textbf{Feature boundary.} This behaviour is scoped to On-Call Directory Resolution; an identically named Feature elsewhere may carry different business meaning.\par
\featuretable{
Function & \texttt{FN-38647096} & Resolve Active On-Call Coverage & Catalogue \\
Service & \texttt{SV-93D42621} & On-Call Specialist Lookup & Linked to preceding Function in source \\
}
\subsubsection{Feature - Context Establishment}\label{cap:bec-4cfa3e90}
\vspace{-1em}
\begin{table}[htbp]
  \centering
  \footnotesize
  \begin{tabularx}{1\textwidth}{|l|l|X|}
    \hline
    \textbf{Metadata Name} & \textbf{Metadata Value} & \textbf{Description} \\
    \hline
    \textbf{Reference Id} & \texttt{BEC-4CFA3E90} & \\
    \hline
    \textbf{Status} & Generated candidate & \\
    \hline
    \textbf{Ancestry} &  & \\
    \hline
    \textbf{Source} & Semantic-pattern candidate & \\
    \hline
    \textbf{Enabled By} &  & \\
    \hline
    \textbf{Supports} &  & \\
    \hline
    \textbf{Derived From} &  & \\
    \hline
    \textbf{Reviewed Status} & Generated candidate: keep / modify / relocate / reject; verify whether independently testable. & \\
    \hline
  \end{tabularx}
\end{table}

Intent. Context Establishment is proposed as a discrete, testable enabling behaviour for On-Call Directory Resolution. It is included for architectural review, not as an established implementation obligation.

\noindent\textbf{Feature boundary.} This behaviour is scoped to On-Call Directory Resolution; an identically named Feature elsewhere may carry different business meaning.\par
\featuretable{
--- & --- & \emph{Not yet captured} & \emph{Complete during review} \\
}
\subsubsection{Feature - State Monitoring}\label{cap:bec-703e878f}
\vspace{-1em}
\begin{table}[htbp]
  \centering
  \footnotesize
  \begin{tabularx}{1\textwidth}{|l|l|X|}
    \hline
    \textbf{Metadata Name} & \textbf{Metadata Value} & \textbf{Description} \\
    \hline
    \textbf{Reference Id} & \texttt{BEC-703E878F} & \\
    \hline
    \textbf{Status} & Generated candidate & \\
    \hline
    \textbf{Ancestry} &  & \\
    \hline
    \textbf{Source} & Semantic-pattern candidate & \\
    \hline
    \textbf{Enabled By} &  & \\
    \hline
    \textbf{Supports} &  & \\
    \hline
    \textbf{Derived From} &  & \\
    \hline
    \textbf{Reviewed Status} & Generated candidate: keep / modify / relocate / reject; verify whether independently testable. & \\
    \hline
  \end{tabularx}
\end{table}

Intent. State Monitoring is proposed as a discrete, testable enabling behaviour for On-Call Directory Resolution. It is included for architectural review, not as an established implementation obligation.

\noindent\textbf{Feature boundary.} This behaviour is scoped to On-Call Directory Resolution; an identically named Feature elsewhere may carry different business meaning.\par
\featuretable{
--- & --- & \emph{Not yet captured} & \emph{Complete during review} \\
}
\subsubsection{Feature - Exception Management}\label{cap:bec-e452abbd}
\vspace{-1em}
\begin{table}[htbp]
  \centering
  \footnotesize
  \begin{tabularx}{1\textwidth}{|l|l|X|}
    \hline
    \textbf{Metadata Name} & \textbf{Metadata Value} & \textbf{Description} \\
    \hline
    \textbf{Reference Id} & \texttt{BEC-E452ABBD} & \\
    \hline
    \textbf{Status} & Generated candidate & \\
    \hline
    \textbf{Ancestry} &  & \\
    \hline
    \textbf{Source} & Semantic-pattern candidate & \\
    \hline
    \textbf{Enabled By} &  & \\
    \hline
    \textbf{Supports} &  & \\
    \hline
    \textbf{Derived From} &  & \\
    \hline
    \textbf{Reviewed Status} & Generated candidate: keep / modify / relocate / reject; verify whether independently testable. & \\
    \hline
  \end{tabularx}
\end{table}

Intent. Exception Management is proposed as a discrete, testable enabling behaviour for On-Call Directory Resolution. It is included for architectural review, not as an established implementation obligation.

\noindent\textbf{Feature boundary.} This behaviour is scoped to On-Call Directory Resolution; an identically named Feature elsewhere may carry different business meaning.\par
\featuretable{
--- & --- & \emph{Not yet captured} & \emph{Complete during review} \\
}
\subsubsection{Feature - Information Retrieval}\label{cap:bec-783c1ed2}
\vspace{-1em}
\begin{table}[htbp]
  \centering
  \footnotesize
  \begin{tabularx}{1\textwidth}{|l|l|X|}
    \hline
    \textbf{Metadata Name} & \textbf{Metadata Value} & \textbf{Description} \\
    \hline
    \textbf{Reference Id} & \texttt{BEC-783C1ED2} & \\
    \hline
    \textbf{Status} & Generated candidate & \\
    \hline
    \textbf{Ancestry} &  & \\
    \hline
    \textbf{Source} & Semantic-pattern candidate & \\
    \hline
    \textbf{Enabled By} &  & \\
    \hline
    \textbf{Supports} &  & \\
    \hline
    \textbf{Derived From} &  & \\
    \hline
    \textbf{Reviewed Status} & Generated candidate: keep / modify / relocate / reject; verify whether independently testable. & \\
    \hline
  \end{tabularx}
\end{table}

Intent. Information Retrieval is proposed as a discrete, testable enabling behaviour for On-Call Directory Resolution. It is included for architectural review, not as an established implementation obligation.

\noindent\textbf{Feature boundary.} This behaviour is scoped to On-Call Directory Resolution; an identically named Feature elsewhere may carry different business meaning.\par
\featuretable{
--- & --- & \emph{Not yet captured} & \emph{Complete during review} \\
}

\section{CT2 - Work Allocation \& Dispatch}\label{cap:bec-4a72d708}
\vspace{-0.5em}
\begin{table}[htbp]
  \centering
  \footnotesize
  \begin{tabularx}{1\textwidth}{|l|l|X|}
    \hline
    \multicolumn{3}{c}{\textbf{Capability Metadata}} \\
    \hline
    \textbf{Name} & \textbf{Value} & \textbf{Comment and/or Description} \\
    \hline
    \textbf{Reference Id} & \texttt{BEC-4A72D708} & \\
    \hline
    \textbf{Status} & Candidate ancestry & \\
    \hline
    \textbf{Ancestry} &  & \\
    \hline
    \textbf{Source} &  & \\
    \hline
    \textbf{Enabled By} &  & \\
    \hline
    \textbf{Supports} &  & \\
    \hline
    \textbf{Derived From} &  & \\
    \hline
    \textbf{Reviewed Status} &  & \\
    \hline
  \end{tabularx}
\end{table}

\begin{table}[htbp]
  \centering
  \footnotesize
  \begin{tabularx}{1\textwidth}{|l|l|X|}
    \hline
    \multicolumn{3}{c}{\textbf{Contained Capabilities}} \\
    \hline
    \textbf{Capability Name} & \textbf{Capability Id} & \textbf{Comment} \\
    \hline
    Operational Work Assignment \& Supervision & \texttt{cap:bec-09c6941c} (at \ref{cap:bec-09c6941c}) & \\
    \hline
  \end{tabularx}
\end{table}

Intent. Work Allocation \& Dispatch is a candidate enabling capability within Health Service Operations Enablement. It frames the reusable system-supported behaviours needed to support this healthcare context without prescribing a product, interface, technology or deployment boundary.

Boundary. This capability supports the business activity and does not inherit the authority or accountability of the Business Capability it enables. Its child Features are review candidates until ancestry and semantics are adjudicated.

\subsection{CT3 - Operational Work Assignment \& Supervision}\label{cap:bec-09c6941c}
\vspace{-0.5em}
\begin{table}[htbp]
  \centering
  \footnotesize
  \begin{tabularx}{1\textwidth}{|l|l|X|}
    \hline
    \multicolumn{3}{c}{\textbf{Capability Metadata}} \\
    \hline
    \textbf{Name} & \textbf{Value} & \textbf{Comment and/or Description} \\
    \hline
    \textbf{Reference Id} & \texttt{BEC-09C6941C} & \\
    \hline
    \textbf{Status} & Candidate ancestry & \\
    \hline
    \textbf{Ancestry} &  & \\
    \hline
    \textbf{Source} &  & \\
    \hline
    \textbf{Enabled By} &  & \\
    \hline
    \textbf{Supports} &  & \\
    \hline
    \textbf{Derived From} &  & \\
    \hline
    \textbf{Reviewed Status} &  & \\
    \hline
  \end{tabularx}
\end{table}

\begin{table}[htbp]
  \centering
  \footnotesize
  \begin{tabularx}{1\textwidth}{|l|l|X|}
    \hline
    \multicolumn{3}{c}{\textbf{Contained Features}} \\
    \hline
    \textbf{Feature Name} & \textbf{Feature Id} & \textbf{Comment} \\
    \hline
    Work Item Instantiation & \texttt{cap:bec-02438774} (at \ref{cap:bec-02438774}) & \\
    \hline
    Worker Matching \& Allocation & \texttt{cap:bec-2a278a2c} (at \ref{cap:bec-2a278a2c}) & \\
    \hline
    Work Dispatch Delivery & \texttt{cap:bec-11bc936a} (at \ref{cap:bec-11bc936a}) & \\
    \hline
    Work Progress Oversight \& Escalation & \texttt{cap:bec-4d7128d9} (at \ref{cap:bec-4d7128d9}) & \\
    \hline
    Context Establishment & \texttt{cap:bec-59cf5f93} (at \ref{cap:bec-59cf5f93}) & \\
    \hline
    Eligibility Validation & \texttt{cap:bec-7b45c0b7} (at \ref{cap:bec-7b45c0b7}) & \\
    \hline
    Assignment \& Routing & \texttt{cap:bec-ad553da5} (at \ref{cap:bec-ad553da5}) & \\
    \hline
    State Tracking & \texttt{cap:bec-d5aedf90} (at \ref{cap:bec-d5aedf90}) & \\
    \hline
    Exception Management & \texttt{cap:bec-c536bd61} (at \ref{cap:bec-c536bd61}) & \\
    \hline
    Outcome Reporting & \texttt{cap:bec-a3268add} (at \ref{cap:bec-a3268add}) & \\
    \hline
  \end{tabularx}
\end{table}

Intent. Operational Work Assignment \& Supervision is a candidate enabling capability within Work Allocation \& Dispatch. It frames the reusable system-supported behaviours needed to support this healthcare context without prescribing a product, interface, technology or deployment boundary.

Boundary. This capability supports the business activity and does not inherit the authority or accountability of the Business Capability it enables. Its child Features are review candidates until ancestry and semantics are adjudicated.

\subsubsection{Feature - Work Item Instantiation}\label{cap:bec-02438774}
\vspace{-1em}
\begin{table}[htbp]
  \centering
  \footnotesize
  \begin{tabularx}{1\textwidth}{|l|l|X|}
    \hline
    \textbf{Metadata Name} & \textbf{Metadata Value} & \textbf{Description} \\
    \hline
    \textbf{Reference Id} & \texttt{BEC-02438774} & \\
    \hline
    \textbf{Status} & Existing Feature; candidate CT3 placement & \\
    \hline
    \textbf{Ancestry} &  & \\
    \hline
    \textbf{Source} & FEAT-HSO-19 & \\
    \hline
    \textbf{Enabled By} &  & \\
    \hline
    \textbf{Supports} &  & \\
    \hline
    \textbf{Derived From} &  & \\
    \hline
    \textbf{Reviewed Status} & Confirm existing Feature meaning and proposed CT3 ownership. & \\
    \hline
  \end{tabularx}
\end{table}

Intent. Work Item Instantiation provides a distinct enabling behaviour in the context of Operational Work Assignment \& Supervision. This is an existing named Feature; its placement under the candidate CT3 remains subject to review.

\noindent\textbf{Feature boundary.} This behaviour is scoped to Operational Work Assignment \& Supervision; an identically named Feature elsewhere may carry different business meaning.\par
\featuretable{
Function & \texttt{FN-8AA7411F} & Dispatch Operational Work Order & Catalogue \\
Service & \texttt{SV-552F0531} & Work Dispatch Service & Linked to preceding Function in source \\
}
\subsubsection{Feature - Worker Matching \& Allocation}\label{cap:bec-2a278a2c}
\vspace{-1em}
\begin{table}[htbp]
  \centering
  \footnotesize
  \begin{tabularx}{1\textwidth}{|l|l|X|}
    \hline
    \textbf{Metadata Name} & \textbf{Metadata Value} & \textbf{Description} \\
    \hline
    \textbf{Reference Id} & \texttt{BEC-2A278A2C} & \\
    \hline
    \textbf{Status} & Existing Feature; candidate CT3 placement & \\
    \hline
    \textbf{Ancestry} &  & \\
    \hline
    \textbf{Source} & FEAT-HSO-20 & \\
    \hline
    \textbf{Enabled By} &  & \\
    \hline
    \textbf{Supports} &  & \\
    \hline
    \textbf{Derived From} &  & \\
    \hline
    \textbf{Reviewed Status} & Confirm existing Feature meaning and proposed CT3 ownership. & \\
    \hline
  \end{tabularx}
\end{table}

Intent. Worker Matching \& Allocation provides a distinct enabling behaviour in the context of Operational Work Assignment \& Supervision. This is an existing named Feature; its placement under the candidate CT3 remains subject to review.

\noindent\textbf{Feature boundary.} This behaviour is scoped to Operational Work Assignment \& Supervision; an identically named Feature elsewhere may carry different business meaning.\par
\featuretable{
Function & \texttt{FN-2C1F6A64} & Dispatch Operational Work Order & Catalogue \\
}
\subsubsection{Feature - Work Dispatch Delivery}\label{cap:bec-11bc936a}
\vspace{-1em}
\begin{table}[htbp]
  \centering
  \footnotesize
  \begin{tabularx}{1\textwidth}{|l|l|X|}
    \hline
    \textbf{Metadata Name} & \textbf{Metadata Value} & \textbf{Description} \\
    \hline
    \textbf{Reference Id} & \texttt{BEC-11BC936A} & \\
    \hline
    \textbf{Status} & Existing Feature; candidate CT3 placement & \\
    \hline
    \textbf{Ancestry} &  & \\
    \hline
    \textbf{Source} & FEAT-HSO-21 & \\
    \hline
    \textbf{Enabled By} &  & \\
    \hline
    \textbf{Supports} &  & \\
    \hline
    \textbf{Derived From} &  & \\
    \hline
    \textbf{Reviewed Status} & Confirm existing Feature meaning and proposed CT3 ownership. & \\
    \hline
  \end{tabularx}
\end{table}

Intent. Work Dispatch Delivery provides a distinct enabling behaviour in the context of Operational Work Assignment \& Supervision. This is an existing named Feature; its placement under the candidate CT3 remains subject to review.

\noindent\textbf{Feature boundary.} This behaviour is scoped to Operational Work Assignment \& Supervision; an identically named Feature elsewhere may carry different business meaning.\par
\featuretable{
Function & \texttt{FN-1D34913B} & Dispatch Operational Work Order & Catalogue \\
}
\subsubsection{Feature - Work Progress Oversight \& Escalation}\label{cap:bec-4d7128d9}
\vspace{-1em}
\begin{table}[htbp]
  \centering
  \footnotesize
  \begin{tabularx}{1\textwidth}{|l|l|X|}
    \hline
    \textbf{Metadata Name} & \textbf{Metadata Value} & \textbf{Description} \\
    \hline
    \textbf{Reference Id} & \texttt{BEC-4D7128D9} & \\
    \hline
    \textbf{Status} & Existing Feature; candidate CT3 placement & \\
    \hline
    \textbf{Ancestry} &  & \\
    \hline
    \textbf{Source} & FEAT-HSO-22 & \\
    \hline
    \textbf{Enabled By} &  & \\
    \hline
    \textbf{Supports} &  & \\
    \hline
    \textbf{Derived From} &  & \\
    \hline
    \textbf{Reviewed Status} & Confirm existing Feature meaning and proposed CT3 ownership. & \\
    \hline
  \end{tabularx}
\end{table}

Intent. Work Progress Oversight \& Escalation provides a distinct enabling behaviour in the context of Operational Work Assignment \& Supervision. This is an existing named Feature; its placement under the candidate CT3 remains subject to review.

\noindent\textbf{Feature boundary.} This behaviour is scoped to Operational Work Assignment \& Supervision; an identically named Feature elsewhere may carry different business meaning.\par
\featuretable{
Function & \texttt{FN-3922D906} & Manage Operational Work Progression & Catalogue \\
Service & \texttt{SV-858C646B} & Operational Work Status Query & Linked to preceding Function in source \\
Process & \texttt{PR-739477FF} & Operational Work Progression Process & Linked to preceding Function in source \\
}
\subsubsection{Feature - Context Establishment}\label{cap:bec-59cf5f93}
\vspace{-1em}
\begin{table}[htbp]
  \centering
  \footnotesize
  \begin{tabularx}{1\textwidth}{|l|l|X|}
    \hline
    \textbf{Metadata Name} & \textbf{Metadata Value} & \textbf{Description} \\
    \hline
    \textbf{Reference Id} & \texttt{BEC-59CF5F93} & \\
    \hline
    \textbf{Status} & Generated candidate & \\
    \hline
    \textbf{Ancestry} &  & \\
    \hline
    \textbf{Source} & Semantic-pattern candidate & \\
    \hline
    \textbf{Enabled By} &  & \\
    \hline
    \textbf{Supports} &  & \\
    \hline
    \textbf{Derived From} &  & \\
    \hline
    \textbf{Reviewed Status} & Generated candidate: keep / modify / relocate / reject; verify whether independently testable. & \\
    \hline
  \end{tabularx}
\end{table}

Intent. Context Establishment is proposed as a discrete, testable enabling behaviour for Operational Work Assignment \& Supervision. It is included for architectural review, not as an established implementation obligation.

\noindent\textbf{Feature boundary.} This behaviour is scoped to Operational Work Assignment \& Supervision; an identically named Feature elsewhere may carry different business meaning.\par
\featuretable{
--- & --- & \emph{Not yet captured} & \emph{Complete during review} \\
}
\subsubsection{Feature - Eligibility Validation}\label{cap:bec-7b45c0b7}
\vspace{-1em}
\begin{table}[htbp]
  \centering
  \footnotesize
  \begin{tabularx}{1\textwidth}{|l|l|X|}
    \hline
    \textbf{Metadata Name} & \textbf{Metadata Value} & \textbf{Description} \\
    \hline
    \textbf{Reference Id} & \texttt{BEC-7B45C0B7} & \\
    \hline
    \textbf{Status} & Generated candidate & \\
    \hline
    \textbf{Ancestry} &  & \\
    \hline
    \textbf{Source} & Semantic-pattern candidate & \\
    \hline
    \textbf{Enabled By} &  & \\
    \hline
    \textbf{Supports} &  & \\
    \hline
    \textbf{Derived From} &  & \\
    \hline
    \textbf{Reviewed Status} & Generated candidate: keep / modify / relocate / reject; verify whether independently testable. & \\
    \hline
  \end{tabularx}
\end{table}

Intent. Eligibility Validation is proposed as a discrete, testable enabling behaviour for Operational Work Assignment \& Supervision. It is included for architectural review, not as an established implementation obligation.

\noindent\textbf{Feature boundary.} This behaviour is scoped to Operational Work Assignment \& Supervision; an identically named Feature elsewhere may carry different business meaning.\par
\featuretable{
--- & --- & \emph{Not yet captured} & \emph{Complete during review} \\
}
\subsubsection{Feature - Assignment \& Routing}\label{cap:bec-ad553da5}
\vspace{-1em}
\begin{table}[htbp]
  \centering
  \footnotesize
  \begin{tabularx}{1\textwidth}{|l|l|X|}
    \hline
    \textbf{Metadata Name} & \textbf{Metadata Value} & \textbf{Description} \\
    \hline
    \textbf{Reference Id} & \texttt{BEC-AD553DA5} & \\
    \hline
    \textbf{Status} & Generated candidate & \\
    \hline
    \textbf{Ancestry} &  & \\
    \hline
    \textbf{Source} & Semantic-pattern candidate & \\
    \hline
    \textbf{Enabled By} &  & \\
    \hline
    \textbf{Supports} &  & \\
    \hline
    \textbf{Derived From} &  & \\
    \hline
    \textbf{Reviewed Status} & Generated candidate: keep / modify / relocate / reject; verify whether independently testable. & \\
    \hline
  \end{tabularx}
\end{table}

Intent. Assignment \& Routing is proposed as a discrete, testable enabling behaviour for Operational Work Assignment \& Supervision. It is included for architectural review, not as an established implementation obligation.

\noindent\textbf{Feature boundary.} This behaviour is scoped to Operational Work Assignment \& Supervision; an identically named Feature elsewhere may carry different business meaning.\par
\featuretable{
--- & --- & \emph{Not yet captured} & \emph{Complete during review} \\
}
\subsubsection{Feature - State Tracking}\label{cap:bec-d5aedf90}
\vspace{-1em}
\begin{table}[htbp]
  \centering
  \footnotesize
  \begin{tabularx}{1\textwidth}{|l|l|X|}
    \hline
    \textbf{Metadata Name} & \textbf{Metadata Value} & \textbf{Description} \\
    \hline
    \textbf{Reference Id} & \texttt{BEC-D5AEDF90} & \\
    \hline
    \textbf{Status} & Generated candidate & \\
    \hline
    \textbf{Ancestry} &  & \\
    \hline
    \textbf{Source} & Semantic-pattern candidate & \\
    \hline
    \textbf{Enabled By} &  & \\
    \hline
    \textbf{Supports} &  & \\
    \hline
    \textbf{Derived From} &  & \\
    \hline
    \textbf{Reviewed Status} & Generated candidate: keep / modify / relocate / reject; verify whether independently testable. & \\
    \hline
  \end{tabularx}
\end{table}

Intent. State Tracking is proposed as a discrete, testable enabling behaviour for Operational Work Assignment \& Supervision. It is included for architectural review, not as an established implementation obligation.

\noindent\textbf{Feature boundary.} This behaviour is scoped to Operational Work Assignment \& Supervision; an identically named Feature elsewhere may carry different business meaning.\par
\featuretable{
--- & --- & \emph{Not yet captured} & \emph{Complete during review} \\
}
\subsubsection{Feature - Exception Management}\label{cap:bec-c536bd61}
\vspace{-1em}
\begin{table}[htbp]
  \centering
  \footnotesize
  \begin{tabularx}{1\textwidth}{|l|l|X|}
    \hline
    \textbf{Metadata Name} & \textbf{Metadata Value} & \textbf{Description} \\
    \hline
    \textbf{Reference Id} & \texttt{BEC-C536BD61} & \\
    \hline
    \textbf{Status} & Generated candidate & \\
    \hline
    \textbf{Ancestry} &  & \\
    \hline
    \textbf{Source} & Semantic-pattern candidate & \\
    \hline
    \textbf{Enabled By} &  & \\
    \hline
    \textbf{Supports} &  & \\
    \hline
    \textbf{Derived From} &  & \\
    \hline
    \textbf{Reviewed Status} & Generated candidate: keep / modify / relocate / reject; verify whether independently testable. & \\
    \hline
  \end{tabularx}
\end{table}

Intent. Exception Management is proposed as a discrete, testable enabling behaviour for Operational Work Assignment \& Supervision. It is included for architectural review, not as an established implementation obligation.

\noindent\textbf{Feature boundary.} This behaviour is scoped to Operational Work Assignment \& Supervision; an identically named Feature elsewhere may carry different business meaning.\par
\featuretable{
--- & --- & \emph{Not yet captured} & \emph{Complete during review} \\
}
\subsubsection{Feature - Outcome Reporting}\label{cap:bec-a3268add}
\vspace{-1em}
\begin{table}[htbp]
  \centering
  \footnotesize
  \begin{tabularx}{1\textwidth}{|l|l|X|}
    \hline
    \textbf{Metadata Name} & \textbf{Metadata Value} & \textbf{Description} \\
    \hline
    \textbf{Reference Id} & \texttt{BEC-A3268ADD} & \\
    \hline
    \textbf{Status} & Generated candidate & \\
    \hline
    \textbf{Ancestry} &  & \\
    \hline
    \textbf{Source} & Semantic-pattern candidate & \\
    \hline
    \textbf{Enabled By} &  & \\
    \hline
    \textbf{Supports} &  & \\
    \hline
    \textbf{Derived From} &  & \\
    \hline
    \textbf{Reviewed Status} & Generated candidate: keep / modify / relocate / reject; verify whether independently testable. & \\
    \hline
  \end{tabularx}
\end{table}

Intent. Outcome Reporting is proposed as a discrete, testable enabling behaviour for Operational Work Assignment \& Supervision. It is included for architectural review, not as an established implementation obligation.

\noindent\textbf{Feature boundary.} This behaviour is scoped to Operational Work Assignment \& Supervision; an identically named Feature elsewhere may carry different business meaning.\par
\featuretable{
--- & --- & \emph{Not yet captured} & \emph{Complete during review} \\
}

\section{CT2 - Clinical Training \& Education}\label{cap:bec-06cf5b6d}
\vspace{-0.5em}
\begin{table}[htbp]
  \centering
  \footnotesize
  \begin{tabularx}{1\textwidth}{|l|l|X|}
    \hline
    \multicolumn{3}{c}{\textbf{Capability Metadata}} \\
    \hline
    \textbf{Name} & \textbf{Value} & \textbf{Comment and/or Description} \\
    \hline
    \textbf{Reference Id} & \texttt{BEC-06CF5B6D} & \\
    \hline
    \textbf{Status} & Candidate ancestry & \\
    \hline
    \textbf{Ancestry} &  & \\
    \hline
    \textbf{Source} &  & \\
    \hline
    \textbf{Enabled By} &  & \\
    \hline
    \textbf{Supports} &  & \\
    \hline
    \textbf{Derived From} &  & \\
    \hline
    \textbf{Reviewed Status} &  & \\
    \hline
  \end{tabularx}
\end{table}

\begin{table}[htbp]
  \centering
  \footnotesize
  \begin{tabularx}{1\textwidth}{|l|l|X|}
    \hline
    \multicolumn{3}{c}{\textbf{Contained Capabilities}} \\
    \hline
    \textbf{Capability Name} & \textbf{Capability Id} & \textbf{Comment} \\
    \hline
    \multicolumn{3}{|l|}{\emph{No contained capabilities are currently identified.}} \\
    \hline
  \end{tabularx}
\end{table}

Intent. Clinical Training \& Education is a candidate enabling capability within Health Service Operations Enablement. It frames the reusable system-supported behaviours needed to support this healthcare context without prescribing a product, interface, technology or deployment boundary.

Boundary. This capability supports the business activity and does not inherit the authority or accountability of the Business Capability it enables. Its child Features are review candidates until ancestry and semantics are adjudicated.

\section{CT2 - Clinical Competency Management}\label{cap:bec-de805931}
\vspace{-0.5em}
\begin{table}[htbp]
  \centering
  \footnotesize
  \begin{tabularx}{1\textwidth}{|l|l|X|}
    \hline
    \multicolumn{3}{c}{\textbf{Capability Metadata}} \\
    \hline
    \textbf{Name} & \textbf{Value} & \textbf{Comment and/or Description} \\
    \hline
    \textbf{Reference Id} & \texttt{BEC-DE805931} & \\
    \hline
    \textbf{Status} & Candidate ancestry & \\
    \hline
    \textbf{Ancestry} &  & \\
    \hline
    \textbf{Source} &  & \\
    \hline
    \textbf{Enabled By} &  & \\
    \hline
    \textbf{Supports} &  & \\
    \hline
    \textbf{Derived From} &  & \\
    \hline
    \textbf{Reviewed Status} &  & \\
    \hline
  \end{tabularx}
\end{table}

\begin{table}[htbp]
  \centering
  \footnotesize
  \begin{tabularx}{1\textwidth}{|l|l|X|}
    \hline
    \multicolumn{3}{c}{\textbf{Contained Capabilities}} \\
    \hline
    \textbf{Capability Name} & \textbf{Capability Id} & \textbf{Comment} \\
    \hline
    \multicolumn{3}{|l|}{\emph{No contained capabilities are currently identified.}} \\
    \hline
  \end{tabularx}
\end{table}

Intent. Clinical Competency Management is a candidate enabling capability within Health Service Operations Enablement. It frames the reusable system-supported behaviours needed to support this healthcare context without prescribing a product, interface, technology or deployment boundary.

Boundary. This capability supports the business activity and does not inherit the authority or accountability of the Business Capability it enables. Its child Features are review candidates until ancestry and semantics are adjudicated.

\section{CT2 - Clinical Credential Management}\label{cap:bec-9c5388a1}
\vspace{-0.5em}
\begin{table}[htbp]
  \centering
  \footnotesize
  \begin{tabularx}{1\textwidth}{|l|l|X|}
    \hline
    \multicolumn{3}{c}{\textbf{Capability Metadata}} \\
    \hline
    \textbf{Name} & \textbf{Value} & \textbf{Comment and/or Description} \\
    \hline
    \textbf{Reference Id} & \texttt{BEC-9C5388A1} & \\
    \hline
    \textbf{Status} & Candidate ancestry & \\
    \hline
    \textbf{Ancestry} &  & \\
    \hline
    \textbf{Source} &  & \\
    \hline
    \textbf{Enabled By} &  & \\
    \hline
    \textbf{Supports} &  & \\
    \hline
    \textbf{Derived From} &  & \\
    \hline
    \textbf{Reviewed Status} &  & \\
    \hline
  \end{tabularx}
\end{table}

\begin{table}[htbp]
  \centering
  \footnotesize
  \begin{tabularx}{1\textwidth}{|l|l|X|}
    \hline
    \multicolumn{3}{c}{\textbf{Contained Capabilities}} \\
    \hline
    \textbf{Capability Name} & \textbf{Capability Id} & \textbf{Comment} \\
    \hline
    Credential Verification Checkpoint & \texttt{cap:bec-e3152a69} (at \ref{cap:bec-e3152a69}) & \\
    \hline
  \end{tabularx}
\end{table}

Intent. Clinical Credential Management is a candidate enabling capability within Health Service Operations Enablement. It frames the reusable system-supported behaviours needed to support this healthcare context without prescribing a product, interface, technology or deployment boundary.

Boundary. This capability supports the business activity and does not inherit the authority or accountability of the Business Capability it enables. Its child Features are review candidates until ancestry and semantics are adjudicated.

\subsection{CT3 - Credential Verification Checkpoint}\label{cap:bec-e3152a69}
\vspace{-0.5em}
\begin{table}[htbp]
  \centering
  \footnotesize
  \begin{tabularx}{1\textwidth}{|l|l|X|}
    \hline
    \multicolumn{3}{c}{\textbf{Capability Metadata}} \\
    \hline
    \textbf{Name} & \textbf{Value} & \textbf{Comment and/or Description} \\
    \hline
    \textbf{Reference Id} & \texttt{BEC-E3152A69} & \\
    \hline
    \textbf{Status} & Candidate ancestry & \\
    \hline
    \textbf{Ancestry} &  & \\
    \hline
    \textbf{Source} &  & \\
    \hline
    \textbf{Enabled By} &  & \\
    \hline
    \textbf{Supports} &  & \\
    \hline
    \textbf{Derived From} &  & \\
    \hline
    \textbf{Reviewed Status} &  & \\
    \hline
  \end{tabularx}
\end{table}

\begin{table}[htbp]
  \centering
  \footnotesize
  \begin{tabularx}{1\textwidth}{|l|l|X|}
    \hline
    \multicolumn{3}{c}{\textbf{Contained Features}} \\
    \hline
    \textbf{Feature Name} & \textbf{Feature Id} & \textbf{Comment} \\
    \hline
    Operational Credential Check & \texttt{cap:bec-5cd4669f} (at \ref{cap:bec-5cd4669f}) & \\
    \hline
    Context Establishment & \texttt{cap:bec-d3d46bf9} (at \ref{cap:bec-d3d46bf9}) & \\
    \hline
    State Monitoring & \texttt{cap:bec-e9331a31} (at \ref{cap:bec-e9331a31}) & \\
    \hline
    Exception Management & \texttt{cap:bec-9736c9db} (at \ref{cap:bec-9736c9db}) & \\
    \hline
    Information Retrieval & \texttt{cap:bec-3de5ac43} (at \ref{cap:bec-3de5ac43}) & \\
    \hline
  \end{tabularx}
\end{table}

Intent. Credential Verification Checkpoint is a candidate enabling capability within Clinical Credential Management. It frames the reusable system-supported behaviours needed to support this healthcare context without prescribing a product, interface, technology or deployment boundary.

Boundary. This capability supports the business activity and does not inherit the authority or accountability of the Business Capability it enables. Its child Features are review candidates until ancestry and semantics are adjudicated.

\subsubsection{Feature - Operational Credential Check}\label{cap:bec-5cd4669f}
\vspace{-1em}
\begin{table}[htbp]
  \centering
  \footnotesize
  \begin{tabularx}{1\textwidth}{|l|l|X|}
    \hline
    \textbf{Metadata Name} & \textbf{Metadata Value} & \textbf{Description} \\
    \hline
    \textbf{Reference Id} & \texttt{BEC-5CD4669F} & \\
    \hline
    \textbf{Status} & Existing Feature; candidate CT3 placement & \\
    \hline
    \textbf{Ancestry} &  & \\
    \hline
    \textbf{Source} & FEAT-HSO-23 & \\
    \hline
    \textbf{Enabled By} &  & \\
    \hline
    \textbf{Supports} &  & \\
    \hline
    \textbf{Derived From} &  & \\
    \hline
    \textbf{Reviewed Status} & Confirm existing Feature meaning and proposed CT3 ownership. & \\
    \hline
  \end{tabularx}
\end{table}

Intent. Operational Credential Check provides a distinct enabling behaviour in the context of Credential Verification Checkpoint. This is an existing named Feature; its placement under the candidate CT3 remains subject to review.

\noindent\textbf{Feature boundary.} This behaviour is scoped to Credential Verification Checkpoint; an identically named Feature elsewhere may carry different business meaning.\par
\featuretable{
Function & \texttt{FN-8CD710BA} & Verify Clinical Qualification for Work Dispatch & Catalogue \\
}
\subsubsection{Feature - Context Establishment}\label{cap:bec-d3d46bf9}
\vspace{-1em}
\begin{table}[htbp]
  \centering
  \footnotesize
  \begin{tabularx}{1\textwidth}{|l|l|X|}
    \hline
    \textbf{Metadata Name} & \textbf{Metadata Value} & \textbf{Description} \\
    \hline
    \textbf{Reference Id} & \texttt{BEC-D3D46BF9} & \\
    \hline
    \textbf{Status} & Generated candidate & \\
    \hline
    \textbf{Ancestry} &  & \\
    \hline
    \textbf{Source} & Semantic-pattern candidate & \\
    \hline
    \textbf{Enabled By} &  & \\
    \hline
    \textbf{Supports} &  & \\
    \hline
    \textbf{Derived From} &  & \\
    \hline
    \textbf{Reviewed Status} & Generated candidate: keep / modify / relocate / reject; verify whether independently testable. & \\
    \hline
  \end{tabularx}
\end{table}

Intent. Context Establishment is proposed as a discrete, testable enabling behaviour for Credential Verification Checkpoint. It is included for architectural review, not as an established implementation obligation.

\noindent\textbf{Feature boundary.} This behaviour is scoped to Credential Verification Checkpoint; an identically named Feature elsewhere may carry different business meaning.\par
\featuretable{
--- & --- & \emph{Not yet captured} & \emph{Complete during review} \\
}
\subsubsection{Feature - State Monitoring}\label{cap:bec-e9331a31}
\vspace{-1em}
\begin{table}[htbp]
  \centering
  \footnotesize
  \begin{tabularx}{1\textwidth}{|l|l|X|}
    \hline
    \textbf{Metadata Name} & \textbf{Metadata Value} & \textbf{Description} \\
    \hline
    \textbf{Reference Id} & \texttt{BEC-E9331A31} & \\
    \hline
    \textbf{Status} & Generated candidate & \\
    \hline
    \textbf{Ancestry} &  & \\
    \hline
    \textbf{Source} & Semantic-pattern candidate & \\
    \hline
    \textbf{Enabled By} &  & \\
    \hline
    \textbf{Supports} &  & \\
    \hline
    \textbf{Derived From} &  & \\
    \hline
    \textbf{Reviewed Status} & Generated candidate: keep / modify / relocate / reject; verify whether independently testable. & \\
    \hline
  \end{tabularx}
\end{table}

Intent. State Monitoring is proposed as a discrete, testable enabling behaviour for Credential Verification Checkpoint. It is included for architectural review, not as an established implementation obligation.

\noindent\textbf{Feature boundary.} This behaviour is scoped to Credential Verification Checkpoint; an identically named Feature elsewhere may carry different business meaning.\par
\featuretable{
--- & --- & \emph{Not yet captured} & \emph{Complete during review} \\
}
\subsubsection{Feature - Exception Management}\label{cap:bec-9736c9db}
\vspace{-1em}
\begin{table}[htbp]
  \centering
  \footnotesize
  \begin{tabularx}{1\textwidth}{|l|l|X|}
    \hline
    \textbf{Metadata Name} & \textbf{Metadata Value} & \textbf{Description} \\
    \hline
    \textbf{Reference Id} & \texttt{BEC-9736C9DB} & \\
    \hline
    \textbf{Status} & Generated candidate & \\
    \hline
    \textbf{Ancestry} &  & \\
    \hline
    \textbf{Source} & Semantic-pattern candidate & \\
    \hline
    \textbf{Enabled By} &  & \\
    \hline
    \textbf{Supports} &  & \\
    \hline
    \textbf{Derived From} &  & \\
    \hline
    \textbf{Reviewed Status} & Generated candidate: keep / modify / relocate / reject; verify whether independently testable. & \\
    \hline
  \end{tabularx}
\end{table}

Intent. Exception Management is proposed as a discrete, testable enabling behaviour for Credential Verification Checkpoint. It is included for architectural review, not as an established implementation obligation.

\noindent\textbf{Feature boundary.} This behaviour is scoped to Credential Verification Checkpoint; an identically named Feature elsewhere may carry different business meaning.\par
\featuretable{
--- & --- & \emph{Not yet captured} & \emph{Complete during review} \\
}
\subsubsection{Feature - Information Retrieval}\label{cap:bec-3de5ac43}
\vspace{-1em}
\begin{table}[htbp]
  \centering
  \footnotesize
  \begin{tabularx}{1\textwidth}{|l|l|X|}
    \hline
    \textbf{Metadata Name} & \textbf{Metadata Value} & \textbf{Description} \\
    \hline
    \textbf{Reference Id} & \texttt{BEC-3DE5AC43} & \\
    \hline
    \textbf{Status} & Generated candidate & \\
    \hline
    \textbf{Ancestry} &  & \\
    \hline
    \textbf{Source} & Semantic-pattern candidate & \\
    \hline
    \textbf{Enabled By} &  & \\
    \hline
    \textbf{Supports} &  & \\
    \hline
    \textbf{Derived From} &  & \\
    \hline
    \textbf{Reviewed Status} & Generated candidate: keep / modify / relocate / reject; verify whether independently testable. & \\
    \hline
  \end{tabularx}
\end{table}

Intent. Information Retrieval is proposed as a discrete, testable enabling behaviour for Credential Verification Checkpoint. It is included for architectural review, not as an established implementation obligation.

\noindent\textbf{Feature boundary.} This behaviour is scoped to Credential Verification Checkpoint; an identically named Feature elsewhere may carry different business meaning.\par
\featuretable{
--- & --- & \emph{Not yet captured} & \emph{Complete during review} \\
}

\section{CT2 - Clinical Supervision}\label{cap:bec-d401771a}
\vspace{-0.5em}
\begin{table}[htbp]
  \centering
  \footnotesize
  \begin{tabularx}{1\textwidth}{|l|l|X|}
    \hline
    \multicolumn{3}{c}{\textbf{Capability Metadata}} \\
    \hline
    \textbf{Name} & \textbf{Value} & \textbf{Comment and/or Description} \\
    \hline
    \textbf{Reference Id} & \texttt{BEC-D401771A} & \\
    \hline
    \textbf{Status} & Candidate ancestry & \\
    \hline
    \textbf{Ancestry} &  & \\
    \hline
    \textbf{Source} &  & \\
    \hline
    \textbf{Enabled By} &  & \\
    \hline
    \textbf{Supports} &  & \\
    \hline
    \textbf{Derived From} &  & \\
    \hline
    \textbf{Reviewed Status} &  & \\
    \hline
  \end{tabularx}
\end{table}

\begin{table}[htbp]
  \centering
  \footnotesize
  \begin{tabularx}{1\textwidth}{|l|l|X|}
    \hline
    \multicolumn{3}{c}{\textbf{Contained Capabilities}} \\
    \hline
    \textbf{Capability Name} & \textbf{Capability Id} & \textbf{Comment} \\
    \hline
    \multicolumn{3}{|l|}{\emph{No contained capabilities are currently identified.}} \\
    \hline
  \end{tabularx}
\end{table}

Intent. Clinical Supervision is a candidate enabling capability within Health Service Operations Enablement. It frames the reusable system-supported behaviours needed to support this healthcare context without prescribing a product, interface, technology or deployment boundary.

Boundary. This capability supports the business activity and does not inherit the authority or accountability of the Business Capability it enables. Its child Features are review candidates until ancestry and semantics are adjudicated.

\section{CT2 - Patient Transport}\label{cap:bec-55fbcc7c}
\vspace{-0.5em}
\begin{table}[htbp]
  \centering
  \footnotesize
  \begin{tabularx}{1\textwidth}{|l|l|X|}
    \hline
    \multicolumn{3}{c}{\textbf{Capability Metadata}} \\
    \hline
    \textbf{Name} & \textbf{Value} & \textbf{Comment and/or Description} \\
    \hline
    \textbf{Reference Id} & \texttt{BEC-55FBCC7C} & \\
    \hline
    \textbf{Status} & Candidate ancestry & \\
    \hline
    \textbf{Ancestry} &  & \\
    \hline
    \textbf{Source} &  & \\
    \hline
    \textbf{Enabled By} &  & \\
    \hline
    \textbf{Supports} &  & \\
    \hline
    \textbf{Derived From} &  & \\
    \hline
    \textbf{Reviewed Status} &  & \\
    \hline
  \end{tabularx}
\end{table}

\begin{table}[htbp]
  \centering
  \footnotesize
  \begin{tabularx}{1\textwidth}{|l|l|X|}
    \hline
    \multicolumn{3}{c}{\textbf{Contained Capabilities}} \\
    \hline
    \textbf{Capability Name} & \textbf{Capability Id} & \textbf{Comment} \\
    \hline
    Patient Transfer Coordination & \texttt{cap:bec-854ff619} (at \ref{cap:bec-854ff619}) & \\
    \hline
  \end{tabularx}
\end{table}

Intent. Patient Transport is a candidate enabling capability within Health Service Operations Enablement. It frames the reusable system-supported behaviours needed to support this healthcare context without prescribing a product, interface, technology or deployment boundary.

Boundary. This capability supports the business activity and does not inherit the authority or accountability of the Business Capability it enables. Its child Features are review candidates until ancestry and semantics are adjudicated.

\subsection{CT3 - Patient Transfer Coordination}\label{cap:bec-854ff619}
\vspace{-0.5em}
\begin{table}[htbp]
  \centering
  \footnotesize
  \begin{tabularx}{1\textwidth}{|l|l|X|}
    \hline
    \multicolumn{3}{c}{\textbf{Capability Metadata}} \\
    \hline
    \textbf{Name} & \textbf{Value} & \textbf{Comment and/or Description} \\
    \hline
    \textbf{Reference Id} & \texttt{BEC-854FF619} & \\
    \hline
    \textbf{Status} & Candidate ancestry & \\
    \hline
    \textbf{Ancestry} &  & \\
    \hline
    \textbf{Source} &  & \\
    \hline
    \textbf{Enabled By} &  & \\
    \hline
    \textbf{Supports} &  & \\
    \hline
    \textbf{Derived From} &  & \\
    \hline
    \textbf{Reviewed Status} &  & \\
    \hline
  \end{tabularx}
\end{table}

\begin{table}[htbp]
  \centering
  \footnotesize
  \begin{tabularx}{1\textwidth}{|l|l|X|}
    \hline
    \multicolumn{3}{c}{\textbf{Contained Features}} \\
    \hline
    \textbf{Feature Name} & \textbf{Feature Id} & \textbf{Comment} \\
    \hline
    Transport Request Ingestion & \texttt{cap:bec-ec853a0f} (at \ref{cap:bec-ec853a0f}) & \\
    \hline
    Transport Dispatch \& Progress Tracking & \texttt{cap:bec-58dccbdb} (at \ref{cap:bec-58dccbdb}) & \\
    \hline
    Subject Identity Correlation & \texttt{cap:bec-05915061} (at \ref{cap:bec-05915061}) & \\
    \hline
    Subject Relationship Management & \texttt{cap:bec-f93ae731} (at \ref{cap:bec-f93ae731}) & \\
    \hline
    Subject Context Management & \texttt{cap:bec-cbf19a4e} (at \ref{cap:bec-cbf19a4e}) & \\
    \hline
    Subject Information Discovery & \texttt{cap:bec-a0e56e8c} (at \ref{cap:bec-a0e56e8c}) & \\
    \hline
  \end{tabularx}
\end{table}

Intent. Patient Transfer Coordination is a candidate enabling capability within Patient Transport. It frames the reusable system-supported behaviours needed to support this healthcare context without prescribing a product, interface, technology or deployment boundary.

Boundary. This capability supports the business activity and does not inherit the authority or accountability of the Business Capability it enables. Its child Features are review candidates until ancestry and semantics are adjudicated.

FHIR R5 conceptual analogues considered (not implementation mappings): Patient, Person.

\subsubsection{Feature - Transport Request Ingestion}\label{cap:bec-ec853a0f}
\vspace{-1em}
\begin{table}[htbp]
  \centering
  \footnotesize
  \begin{tabularx}{1\textwidth}{|l|l|X|}
    \hline
    \textbf{Metadata Name} & \textbf{Metadata Value} & \textbf{Description} \\
    \hline
    \textbf{Reference Id} & \texttt{BEC-EC853A0F} & \\
    \hline
    \textbf{Status} & Existing Feature; candidate CT3 placement & \\
    \hline
    \textbf{Ancestry} &  & \\
    \hline
    \textbf{Source} & FEAT-HSO-24 & \\
    \hline
    \textbf{Enabled By} &  & \\
    \hline
    \textbf{Supports} &  & \\
    \hline
    \textbf{Derived From} &  & \\
    \hline
    \textbf{Reviewed Status} & Confirm existing Feature meaning and proposed CT3 ownership. & \\
    \hline
  \end{tabularx}
\end{table}

Intent. Transport Request Ingestion provides a distinct enabling behaviour in the context of Patient Transfer Coordination. This is an existing named Feature; its placement under the candidate CT3 remains subject to review.

\noindent\textbf{Feature boundary.} This behaviour is scoped to Patient Transfer Coordination; an identically named Feature elsewhere may carry different business meaning.\par
\featuretable{
Function & \texttt{FN-C12DB1C6} & Manage Patient Transport Progression & Catalogue \\
Process & \texttt{PR-27224E4D} & Patient Transport Process & Linked to preceding Function in source \\
}
\subsubsection{Feature - Transport Dispatch \& Progress Tracking}\label{cap:bec-58dccbdb}
\vspace{-1em}
\begin{table}[htbp]
  \centering
  \footnotesize
  \begin{tabularx}{1\textwidth}{|l|l|X|}
    \hline
    \textbf{Metadata Name} & \textbf{Metadata Value} & \textbf{Description} \\
    \hline
    \textbf{Reference Id} & \texttt{BEC-58DCCBDB} & \\
    \hline
    \textbf{Status} & Existing Feature; candidate CT3 placement & \\
    \hline
    \textbf{Ancestry} &  & \\
    \hline
    \textbf{Source} & FEAT-HSO-25 & \\
    \hline
    \textbf{Enabled By} &  & \\
    \hline
    \textbf{Supports} &  & \\
    \hline
    \textbf{Derived From} &  & \\
    \hline
    \textbf{Reviewed Status} & Confirm existing Feature meaning and proposed CT3 ownership. & \\
    \hline
  \end{tabularx}
\end{table}

Intent. Transport Dispatch \& Progress Tracking provides a distinct enabling behaviour in the context of Patient Transfer Coordination. This is an existing named Feature; its placement under the candidate CT3 remains subject to review.

\noindent\textbf{Feature boundary.} This behaviour is scoped to Patient Transfer Coordination; an identically named Feature elsewhere may carry different business meaning.\par
\featuretable{
Function & \texttt{FN-CBDF61C2} & Manage Patient Transport Progression & Catalogue \\
Service & \texttt{SV-7E4D29AB} & Patient Transport Tracking Service & Linked to preceding Function in source \\
Process & \texttt{PR-66DCC7F2} & Patient Transport Process & Linked to preceding Function in source \\
}
\subsubsection{Feature - Subject Identity Correlation}\label{cap:bec-05915061}
\vspace{-1em}
\begin{table}[htbp]
  \centering
  \footnotesize
  \begin{tabularx}{1\textwidth}{|l|l|X|}
    \hline
    \textbf{Metadata Name} & \textbf{Metadata Value} & \textbf{Description} \\
    \hline
    \textbf{Reference Id} & \texttt{BEC-05915061} & \\
    \hline
    \textbf{Status} & Generated candidate & \\
    \hline
    \textbf{Ancestry} &  & \\
    \hline
    \textbf{Source} & FHIR-informed candidate & \\
    \hline
    \textbf{Enabled By} &  & \\
    \hline
    \textbf{Supports} &  & \\
    \hline
    \textbf{Derived From} &  & \\
    \hline
    \textbf{Reviewed Status} & Generated candidate: keep / modify / relocate / reject; verify whether independently testable. & \\
    \hline
  \end{tabularx}
\end{table}

Intent. Subject Identity Correlation is proposed to support the underlying healthcare concept represented by Patient Transfer Coordination. The proposal concerns behaviour and semantic responsibility, not a FHIR field, operation, resource instance or prescribed data model.

Conceptual evidence: HL7 FHIR R5 Patient, Person.

\noindent\textbf{Feature boundary.} This behaviour is scoped to Patient Transfer Coordination; an identically named Feature elsewhere may carry different business meaning.\par
\featuretable{
--- & --- & \emph{Not yet captured} & \emph{Complete during review} \\
}
\subsubsection{Feature - Subject Relationship Management}\label{cap:bec-f93ae731}
\vspace{-1em}
\begin{table}[htbp]
  \centering
  \footnotesize
  \begin{tabularx}{1\textwidth}{|l|l|X|}
    \hline
    \textbf{Metadata Name} & \textbf{Metadata Value} & \textbf{Description} \\
    \hline
    \textbf{Reference Id} & \texttt{BEC-F93AE731} & \\
    \hline
    \textbf{Status} & Generated candidate & \\
    \hline
    \textbf{Ancestry} &  & \\
    \hline
    \textbf{Source} & FHIR-informed candidate & \\
    \hline
    \textbf{Enabled By} &  & \\
    \hline
    \textbf{Supports} &  & \\
    \hline
    \textbf{Derived From} &  & \\
    \hline
    \textbf{Reviewed Status} & Generated candidate: keep / modify / relocate / reject; verify whether independently testable. & \\
    \hline
  \end{tabularx}
\end{table}

Intent. Subject Relationship Management is proposed to support the underlying healthcare concept represented by Patient Transfer Coordination. The proposal concerns behaviour and semantic responsibility, not a FHIR field, operation, resource instance or prescribed data model.

Conceptual evidence: HL7 FHIR R5 Patient, Person.

\noindent\textbf{Feature boundary.} This behaviour is scoped to Patient Transfer Coordination; an identically named Feature elsewhere may carry different business meaning.\par
\featuretable{
--- & --- & \emph{Not yet captured} & \emph{Complete during review} \\
}
\subsubsection{Feature - Subject Context Management}\label{cap:bec-cbf19a4e}
\vspace{-1em}
\begin{table}[htbp]
  \centering
  \footnotesize
  \begin{tabularx}{1\textwidth}{|l|l|X|}
    \hline
    \textbf{Metadata Name} & \textbf{Metadata Value} & \textbf{Description} \\
    \hline
    \textbf{Reference Id} & \texttt{BEC-CBF19A4E} & \\
    \hline
    \textbf{Status} & Generated candidate & \\
    \hline
    \textbf{Ancestry} &  & \\
    \hline
    \textbf{Source} & FHIR-informed candidate & \\
    \hline
    \textbf{Enabled By} &  & \\
    \hline
    \textbf{Supports} &  & \\
    \hline
    \textbf{Derived From} &  & \\
    \hline
    \textbf{Reviewed Status} & Generated candidate: keep / modify / relocate / reject; verify whether independently testable. & \\
    \hline
  \end{tabularx}
\end{table}

Intent. Subject Context Management is proposed to support the underlying healthcare concept represented by Patient Transfer Coordination. The proposal concerns behaviour and semantic responsibility, not a FHIR field, operation, resource instance or prescribed data model.

Conceptual evidence: HL7 FHIR R5 Patient, Person.

\noindent\textbf{Feature boundary.} This behaviour is scoped to Patient Transfer Coordination; an identically named Feature elsewhere may carry different business meaning.\par
\featuretable{
--- & --- & \emph{Not yet captured} & \emph{Complete during review} \\
}
\subsubsection{Feature - Subject Information Discovery}\label{cap:bec-a0e56e8c}
\vspace{-1em}
\begin{table}[htbp]
  \centering
  \footnotesize
  \begin{tabularx}{1\textwidth}{|l|l|X|}
    \hline
    \textbf{Metadata Name} & \textbf{Metadata Value} & \textbf{Description} \\
    \hline
    \textbf{Reference Id} & \texttt{BEC-A0E56E8C} & \\
    \hline
    \textbf{Status} & Generated candidate & \\
    \hline
    \textbf{Ancestry} &  & \\
    \hline
    \textbf{Source} & FHIR-informed candidate & \\
    \hline
    \textbf{Enabled By} &  & \\
    \hline
    \textbf{Supports} &  & \\
    \hline
    \textbf{Derived From} &  & \\
    \hline
    \textbf{Reviewed Status} & Generated candidate: keep / modify / relocate / reject; verify whether independently testable. & \\
    \hline
  \end{tabularx}
\end{table}

Intent. Subject Information Discovery is proposed to support the underlying healthcare concept represented by Patient Transfer Coordination. The proposal concerns behaviour and semantic responsibility, not a FHIR field, operation, resource instance or prescribed data model.

Conceptual evidence: HL7 FHIR R5 Patient, Person.

\noindent\textbf{Feature boundary.} This behaviour is scoped to Patient Transfer Coordination; an identically named Feature elsewhere may carry different business meaning.\par
\featuretable{
--- & --- & \emph{Not yet captured} & \emph{Complete during review} \\
}

\section{CT2 - Clinical Logistics Coordination}\label{cap:bec-bb578270}
\vspace{-0.5em}
\begin{table}[htbp]
  \centering
  \footnotesize
  \begin{tabularx}{1\textwidth}{|l|l|X|}
    \hline
    \multicolumn{3}{c}{\textbf{Capability Metadata}} \\
    \hline
    \textbf{Name} & \textbf{Value} & \textbf{Comment and/or Description} \\
    \hline
    \textbf{Reference Id} & \texttt{BEC-BB578270} & \\
    \hline
    \textbf{Status} & Candidate ancestry & \\
    \hline
    \textbf{Ancestry} &  & \\
    \hline
    \textbf{Source} &  & \\
    \hline
    \textbf{Enabled By} &  & \\
    \hline
    \textbf{Supports} &  & \\
    \hline
    \textbf{Derived From} &  & \\
    \hline
    \textbf{Reviewed Status} &  & \\
    \hline
  \end{tabularx}
\end{table}

\begin{table}[htbp]
  \centering
  \footnotesize
  \begin{tabularx}{1\textwidth}{|l|l|X|}
    \hline
    \multicolumn{3}{c}{\textbf{Contained Capabilities}} \\
    \hline
    \textbf{Capability Name} & \textbf{Capability Id} & \textbf{Comment} \\
    \hline
    Specimen \& Supply Courier Tracking & \texttt{cap:bec-999fb8ab} (at \ref{cap:bec-999fb8ab}) & \\
    \hline
  \end{tabularx}
\end{table}

Intent. Clinical Logistics Coordination is a candidate enabling capability within Health Service Operations Enablement. It frames the reusable system-supported behaviours needed to support this healthcare context without prescribing a product, interface, technology or deployment boundary.

Boundary. This capability supports the business activity and does not inherit the authority or accountability of the Business Capability it enables. Its child Features are review candidates until ancestry and semantics are adjudicated.

\subsection{CT3 - Specimen \& Supply Courier Tracking}\label{cap:bec-999fb8ab}
\vspace{-0.5em}
\begin{table}[htbp]
  \centering
  \footnotesize
  \begin{tabularx}{1\textwidth}{|l|l|X|}
    \hline
    \multicolumn{3}{c}{\textbf{Capability Metadata}} \\
    \hline
    \textbf{Name} & \textbf{Value} & \textbf{Comment and/or Description} \\
    \hline
    \textbf{Reference Id} & \texttt{BEC-999FB8AB} & \\
    \hline
    \textbf{Status} & Candidate ancestry & \\
    \hline
    \textbf{Ancestry} &  & \\
    \hline
    \textbf{Source} &  & \\
    \hline
    \textbf{Enabled By} &  & \\
    \hline
    \textbf{Supports} &  & \\
    \hline
    \textbf{Derived From} &  & \\
    \hline
    \textbf{Reviewed Status} &  & \\
    \hline
  \end{tabularx}
\end{table}

\begin{table}[htbp]
  \centering
  \footnotesize
  \begin{tabularx}{1\textwidth}{|l|l|X|}
    \hline
    \multicolumn{3}{c}{\textbf{Contained Features}} \\
    \hline
    \textbf{Feature Name} & \textbf{Feature Id} & \textbf{Comment} \\
    \hline
    Pathology Specimen Transport Tracking & \texttt{cap:bec-0bdcfa64} (at \ref{cap:bec-0bdcfa64}) & \\
    \hline
    Specimen Identity Management & \texttt{cap:bec-900a97f0} (at \ref{cap:bec-900a97f0}) & \\
    \hline
    Specimen Custody Tracking & \texttt{cap:bec-9950ba59} (at \ref{cap:bec-9950ba59}) & \\
    \hline
    Specimen Relationship Management & \texttt{cap:bec-e08c24f3} (at \ref{cap:bec-e08c24f3}) & \\
    \hline
  \end{tabularx}
\end{table}

Intent. Specimen \& Supply Courier Tracking is a candidate enabling capability within Clinical Logistics Coordination. It frames the reusable system-supported behaviours needed to support this healthcare context without prescribing a product, interface, technology or deployment boundary.

Boundary. This capability supports the business activity and does not inherit the authority or accountability of the Business Capability it enables. Its child Features are review candidates until ancestry and semantics are adjudicated.

FHIR R5 conceptual analogues considered (not implementation mappings): Specimen.

\subsubsection{Feature - Pathology Specimen Transport Tracking}\label{cap:bec-0bdcfa64}
\vspace{-1em}
\begin{table}[htbp]
  \centering
  \footnotesize
  \begin{tabularx}{1\textwidth}{|l|l|X|}
    \hline
    \textbf{Metadata Name} & \textbf{Metadata Value} & \textbf{Description} \\
    \hline
    \textbf{Reference Id} & \texttt{BEC-0BDCFA64} & \\
    \hline
    \textbf{Status} & Existing Feature; candidate CT3 placement & \\
    \hline
    \textbf{Ancestry} &  & \\
    \hline
    \textbf{Source} & FEAT-HSO-26 & \\
    \hline
    \textbf{Enabled By} &  & \\
    \hline
    \textbf{Supports} &  & \\
    \hline
    \textbf{Derived From} &  & \\
    \hline
    \textbf{Reviewed Status} & Confirm existing Feature meaning and proposed CT3 ownership. & \\
    \hline
  \end{tabularx}
\end{table}

Intent. Pathology Specimen Transport Tracking provides a distinct enabling behaviour in the context of Specimen \& Supply Courier Tracking. This is an existing named Feature; its placement under the candidate CT3 remains subject to review.

\noindent\textbf{Feature boundary.} This behaviour is scoped to Specimen \& Supply Courier Tracking; an identically named Feature elsewhere may carry different business meaning.\par
\featuretable{
Function & \texttt{FN-930305AA} & Manage Specimen Transport Progression & Catalogue \\
Service & \texttt{SV-B45BDA39} & Specimen Transit Tracking Service & Linked to preceding Function in source \\
Process & \texttt{PR-F80AFB04} & Specimen Transport Process & Linked to preceding Function in source \\
}
\subsubsection{Feature - Specimen Identity Management}\label{cap:bec-900a97f0}
\vspace{-1em}
\begin{table}[htbp]
  \centering
  \footnotesize
  \begin{tabularx}{1\textwidth}{|l|l|X|}
    \hline
    \textbf{Metadata Name} & \textbf{Metadata Value} & \textbf{Description} \\
    \hline
    \textbf{Reference Id} & \texttt{BEC-900A97F0} & \\
    \hline
    \textbf{Status} & Generated candidate & \\
    \hline
    \textbf{Ancestry} &  & \\
    \hline
    \textbf{Source} & FHIR-informed candidate & \\
    \hline
    \textbf{Enabled By} &  & \\
    \hline
    \textbf{Supports} &  & \\
    \hline
    \textbf{Derived From} &  & \\
    \hline
    \textbf{Reviewed Status} & Generated candidate: keep / modify / relocate / reject; verify whether independently testable. & \\
    \hline
  \end{tabularx}
\end{table}

Intent. Specimen Identity Management is proposed to support the underlying healthcare concept represented by Specimen \& Supply Courier Tracking. The proposal concerns behaviour and semantic responsibility, not a FHIR field, operation, resource instance or prescribed data model.

Conceptual evidence: HL7 FHIR R5 Specimen.

\noindent\textbf{Feature boundary.} This behaviour is scoped to Specimen \& Supply Courier Tracking; an identically named Feature elsewhere may carry different business meaning.\par
\featuretable{
--- & --- & \emph{Not yet captured} & \emph{Complete during review} \\
}
\subsubsection{Feature - Specimen Custody Tracking}\label{cap:bec-9950ba59}
\vspace{-1em}
\begin{table}[htbp]
  \centering
  \footnotesize
  \begin{tabularx}{1\textwidth}{|l|l|X|}
    \hline
    \textbf{Metadata Name} & \textbf{Metadata Value} & \textbf{Description} \\
    \hline
    \textbf{Reference Id} & \texttt{BEC-9950BA59} & \\
    \hline
    \textbf{Status} & Generated candidate & \\
    \hline
    \textbf{Ancestry} &  & \\
    \hline
    \textbf{Source} & FHIR-informed candidate & \\
    \hline
    \textbf{Enabled By} &  & \\
    \hline
    \textbf{Supports} &  & \\
    \hline
    \textbf{Derived From} &  & \\
    \hline
    \textbf{Reviewed Status} & Generated candidate: keep / modify / relocate / reject; verify whether independently testable. & \\
    \hline
  \end{tabularx}
\end{table}

Intent. Specimen Custody Tracking is proposed to support the underlying healthcare concept represented by Specimen \& Supply Courier Tracking. The proposal concerns behaviour and semantic responsibility, not a FHIR field, operation, resource instance or prescribed data model.

Conceptual evidence: HL7 FHIR R5 Specimen.

\noindent\textbf{Feature boundary.} This behaviour is scoped to Specimen \& Supply Courier Tracking; an identically named Feature elsewhere may carry different business meaning.\par
\featuretable{
--- & --- & \emph{Not yet captured} & \emph{Complete during review} \\
}
\subsubsection{Feature - Specimen Relationship Management}\label{cap:bec-e08c24f3}
\vspace{-1em}
\begin{table}[htbp]
  \centering
  \footnotesize
  \begin{tabularx}{1\textwidth}{|l|l|X|}
    \hline
    \textbf{Metadata Name} & \textbf{Metadata Value} & \textbf{Description} \\
    \hline
    \textbf{Reference Id} & \texttt{BEC-E08C24F3} & \\
    \hline
    \textbf{Status} & Generated candidate & \\
    \hline
    \textbf{Ancestry} &  & \\
    \hline
    \textbf{Source} & FHIR-informed candidate & \\
    \hline
    \textbf{Enabled By} &  & \\
    \hline
    \textbf{Supports} &  & \\
    \hline
    \textbf{Derived From} &  & \\
    \hline
    \textbf{Reviewed Status} & Generated candidate: keep / modify / relocate / reject; verify whether independently testable. & \\
    \hline
  \end{tabularx}
\end{table}

Intent. Specimen Relationship Management is proposed to support the underlying healthcare concept represented by Specimen \& Supply Courier Tracking. The proposal concerns behaviour and semantic responsibility, not a FHIR field, operation, resource instance or prescribed data model.

Conceptual evidence: HL7 FHIR R5 Specimen.

\noindent\textbf{Feature boundary.} This behaviour is scoped to Specimen \& Supply Courier Tracking; an identically named Feature elsewhere may carry different business meaning.\par
\featuretable{
--- & --- & \emph{Not yet captured} & \emph{Complete during review} \\
}

\section{CT2 - Discharge Management}\label{cap:bec-3c941bb0}
\vspace{-0.5em}
\begin{table}[htbp]
  \centering
  \footnotesize
  \begin{tabularx}{1\textwidth}{|l|l|X|}
    \hline
    \multicolumn{3}{c}{\textbf{Capability Metadata}} \\
    \hline
    \textbf{Name} & \textbf{Value} & \textbf{Comment and/or Description} \\
    \hline
    \textbf{Reference Id} & \texttt{BEC-3C941BB0} & \\
    \hline
    \textbf{Status} & Candidate ancestry & \\
    \hline
    \textbf{Ancestry} &  & \\
    \hline
    \textbf{Source} &  & \\
    \hline
    \textbf{Enabled By} &  & \\
    \hline
    \textbf{Supports} &  & \\
    \hline
    \textbf{Derived From} &  & \\
    \hline
    \textbf{Reviewed Status} &  & \\
    \hline
  \end{tabularx}
\end{table}

\begin{table}[htbp]
  \centering
  \footnotesize
  \begin{tabularx}{1\textwidth}{|l|l|X|}
    \hline
    \multicolumn{3}{c}{\textbf{Contained Capabilities}} \\
    \hline
    \textbf{Capability Name} & \textbf{Capability Id} & \textbf{Comment} \\
    \hline
    Discharge Progression \& Coordination & \texttt{cap:bec-3414f1eb} (at \ref{cap:bec-3414f1eb}) & \\
    \hline
  \end{tabularx}
\end{table}

Intent. Discharge Management is a candidate enabling capability within Health Service Operations Enablement. It frames the reusable system-supported behaviours needed to support this healthcare context without prescribing a product, interface, technology or deployment boundary.

Boundary. This capability supports the business activity and does not inherit the authority or accountability of the Business Capability it enables. Its child Features are review candidates until ancestry and semantics are adjudicated.

\subsection{CT3 - Discharge Progression \& Coordination}\label{cap:bec-3414f1eb}
\vspace{-0.5em}
\begin{table}[htbp]
  \centering
  \footnotesize
  \begin{tabularx}{1\textwidth}{|l|l|X|}
    \hline
    \multicolumn{3}{c}{\textbf{Capability Metadata}} \\
    \hline
    \textbf{Name} & \textbf{Value} & \textbf{Comment and/or Description} \\
    \hline
    \textbf{Reference Id} & \texttt{BEC-3414F1EB} & \\
    \hline
    \textbf{Status} & Candidate ancestry & \\
    \hline
    \textbf{Ancestry} &  & \\
    \hline
    \textbf{Source} &  & \\
    \hline
    \textbf{Enabled By} &  & \\
    \hline
    \textbf{Supports} &  & \\
    \hline
    \textbf{Derived From} &  & \\
    \hline
    \textbf{Reviewed Status} &  & \\
    \hline
  \end{tabularx}
\end{table}

\begin{table}[htbp]
  \centering
  \footnotesize
  \begin{tabularx}{1\textwidth}{|l|l|X|}
    \hline
    \multicolumn{3}{c}{\textbf{Contained Features}} \\
    \hline
    \textbf{Feature Name} & \textbf{Feature Id} & \textbf{Comment} \\
    \hline
    Discharge Milestone Tracking & \texttt{cap:bec-b778d440} (at \ref{cap:bec-b778d440}) & \\
    \hline
    Discharge Readiness \& Coordination Oversight & \texttt{cap:bec-f3326f02} (at \ref{cap:bec-f3326f02}) & \\
    \hline
    Discharge Notification Dispatch & \texttt{cap:bec-8b11d01c} (at \ref{cap:bec-8b11d01c}) & \\
    \hline
    Context Establishment & \texttt{cap:bec-377b8814} (at \ref{cap:bec-377b8814}) & \\
    \hline
    Eligibility Validation & \texttt{cap:bec-54362543} (at \ref{cap:bec-54362543}) & \\
    \hline
    Assignment \& Routing & \texttt{cap:bec-92074ce1} (at \ref{cap:bec-92074ce1}) & \\
    \hline
    State Tracking & \texttt{cap:bec-d2873e18} (at \ref{cap:bec-d2873e18}) & \\
    \hline
    Exception Management & \texttt{cap:bec-af97d3b7} (at \ref{cap:bec-af97d3b7}) & \\
    \hline
    Outcome Reporting & \texttt{cap:bec-d56cf336} (at \ref{cap:bec-d56cf336}) & \\
    \hline
  \end{tabularx}
\end{table}

Intent. Discharge Progression \& Coordination is a candidate enabling capability within Discharge Management. It frames the reusable system-supported behaviours needed to support this healthcare context without prescribing a product, interface, technology or deployment boundary.

Boundary. This capability supports the business activity and does not inherit the authority or accountability of the Business Capability it enables. Its child Features are review candidates until ancestry and semantics are adjudicated.

\subsubsection{Feature - Discharge Milestone Tracking}\label{cap:bec-b778d440}
\vspace{-1em}
\begin{table}[htbp]
  \centering
  \footnotesize
  \begin{tabularx}{1\textwidth}{|l|l|X|}
    \hline
    \textbf{Metadata Name} & \textbf{Metadata Value} & \textbf{Description} \\
    \hline
    \textbf{Reference Id} & \texttt{BEC-B778D440} & \\
    \hline
    \textbf{Status} & Existing Feature; candidate CT3 placement & \\
    \hline
    \textbf{Ancestry} &  & \\
    \hline
    \textbf{Source} & FEAT-HSO-27 & \\
    \hline
    \textbf{Enabled By} &  & \\
    \hline
    \textbf{Supports} &  & \\
    \hline
    \textbf{Derived From} &  & \\
    \hline
    \textbf{Reviewed Status} & Confirm existing Feature meaning and proposed CT3 ownership. & \\
    \hline
  \end{tabularx}
\end{table}

Intent. Discharge Milestone Tracking provides a distinct enabling behaviour in the context of Discharge Progression \& Coordination. This is an existing named Feature; its placement under the candidate CT3 remains subject to review.

\noindent\textbf{Feature boundary.} This behaviour is scoped to Discharge Progression \& Coordination; an identically named Feature elsewhere may carry different business meaning.\par
\featuretable{
Function & \texttt{FN-8867E074} & Manage Discharge Coordination Progression & Catalogue \\
Process & \texttt{PR-8068764D} & Discharge Progression Process & Linked to preceding Function in source \\
}
\subsubsection{Feature - Discharge Readiness \& Coordination Oversight}\label{cap:bec-f3326f02}
\vspace{-1em}
\begin{table}[htbp]
  \centering
  \footnotesize
  \begin{tabularx}{1\textwidth}{|l|l|X|}
    \hline
    \textbf{Metadata Name} & \textbf{Metadata Value} & \textbf{Description} \\
    \hline
    \textbf{Reference Id} & \texttt{BEC-F3326F02} & \\
    \hline
    \textbf{Status} & Existing Feature; candidate CT3 placement & \\
    \hline
    \textbf{Ancestry} &  & \\
    \hline
    \textbf{Source} & FEAT-HSO-28 & \\
    \hline
    \textbf{Enabled By} &  & \\
    \hline
    \textbf{Supports} &  & \\
    \hline
    \textbf{Derived From} &  & \\
    \hline
    \textbf{Reviewed Status} & Confirm existing Feature meaning and proposed CT3 ownership. & \\
    \hline
  \end{tabularx}
\end{table}

Intent. Discharge Readiness \& Coordination Oversight provides a distinct enabling behaviour in the context of Discharge Progression \& Coordination. This is an existing named Feature; its placement under the candidate CT3 remains subject to review.

\noindent\textbf{Feature boundary.} This behaviour is scoped to Discharge Progression \& Coordination; an identically named Feature elsewhere may carry different business meaning.\par
\featuretable{
Function & \texttt{FN-98CEF800} & Manage Discharge Coordination Progression & Catalogue \\
Service & \texttt{SV-D2F959E8} & Discharge Readiness Telemetry & Linked to preceding Function in source \\
Process & \texttt{PR-810FAC05} & Discharge Progression Process & Linked to preceding Function in source \\
}
\subsubsection{Feature - Discharge Notification Dispatch}\label{cap:bec-8b11d01c}
\vspace{-1em}
\begin{table}[htbp]
  \centering
  \footnotesize
  \begin{tabularx}{1\textwidth}{|l|l|X|}
    \hline
    \textbf{Metadata Name} & \textbf{Metadata Value} & \textbf{Description} \\
    \hline
    \textbf{Reference Id} & \texttt{BEC-8B11D01C} & \\
    \hline
    \textbf{Status} & Existing Feature; candidate CT3 placement & \\
    \hline
    \textbf{Ancestry} &  & \\
    \hline
    \textbf{Source} & FEAT-HSO-29 & \\
    \hline
    \textbf{Enabled By} &  & \\
    \hline
    \textbf{Supports} &  & \\
    \hline
    \textbf{Derived From} &  & \\
    \hline
    \textbf{Reviewed Status} & Confirm existing Feature meaning and proposed CT3 ownership. & \\
    \hline
  \end{tabularx}
\end{table}

Intent. Discharge Notification Dispatch provides a distinct enabling behaviour in the context of Discharge Progression \& Coordination. This is an existing named Feature; its placement under the candidate CT3 remains subject to review.

\noindent\textbf{Feature boundary.} This behaviour is scoped to Discharge Progression \& Coordination; an identically named Feature elsewhere may carry different business meaning.\par
\featuretable{
Function & \texttt{FN-DE97896F} & Manage Discharge Coordination Progression & Catalogue \\
Process & \texttt{PR-D4342AD3} & Discharge Progression Process & Linked to preceding Function in source \\
}
\subsubsection{Feature - Context Establishment}\label{cap:bec-377b8814}
\vspace{-1em}
\begin{table}[htbp]
  \centering
  \footnotesize
  \begin{tabularx}{1\textwidth}{|l|l|X|}
    \hline
    \textbf{Metadata Name} & \textbf{Metadata Value} & \textbf{Description} \\
    \hline
    \textbf{Reference Id} & \texttt{BEC-377B8814} & \\
    \hline
    \textbf{Status} & Generated candidate & \\
    \hline
    \textbf{Ancestry} &  & \\
    \hline
    \textbf{Source} & Semantic-pattern candidate & \\
    \hline
    \textbf{Enabled By} &  & \\
    \hline
    \textbf{Supports} &  & \\
    \hline
    \textbf{Derived From} &  & \\
    \hline
    \textbf{Reviewed Status} & Generated candidate: keep / modify / relocate / reject; verify whether independently testable. & \\
    \hline
  \end{tabularx}
\end{table}

Intent. Context Establishment is proposed as a discrete, testable enabling behaviour for Discharge Progression \& Coordination. It is included for architectural review, not as an established implementation obligation.

\noindent\textbf{Feature boundary.} This behaviour is scoped to Discharge Progression \& Coordination; an identically named Feature elsewhere may carry different business meaning.\par
\featuretable{
--- & --- & \emph{Not yet captured} & \emph{Complete during review} \\
}
\subsubsection{Feature - Eligibility Validation}\label{cap:bec-54362543}
\vspace{-1em}
\begin{table}[htbp]
  \centering
  \footnotesize
  \begin{tabularx}{1\textwidth}{|l|l|X|}
    \hline
    \textbf{Metadata Name} & \textbf{Metadata Value} & \textbf{Description} \\
    \hline
    \textbf{Reference Id} & \texttt{BEC-54362543} & \\
    \hline
    \textbf{Status} & Generated candidate & \\
    \hline
    \textbf{Ancestry} &  & \\
    \hline
    \textbf{Source} & Semantic-pattern candidate & \\
    \hline
    \textbf{Enabled By} &  & \\
    \hline
    \textbf{Supports} &  & \\
    \hline
    \textbf{Derived From} &  & \\
    \hline
    \textbf{Reviewed Status} & Generated candidate: keep / modify / relocate / reject; verify whether independently testable. & \\
    \hline
  \end{tabularx}
\end{table}

Intent. Eligibility Validation is proposed as a discrete, testable enabling behaviour for Discharge Progression \& Coordination. It is included for architectural review, not as an established implementation obligation.

\noindent\textbf{Feature boundary.} This behaviour is scoped to Discharge Progression \& Coordination; an identically named Feature elsewhere may carry different business meaning.\par
\featuretable{
--- & --- & \emph{Not yet captured} & \emph{Complete during review} \\
}
\subsubsection{Feature - Assignment \& Routing}\label{cap:bec-92074ce1}
\vspace{-1em}
\begin{table}[htbp]
  \centering
  \footnotesize
  \begin{tabularx}{1\textwidth}{|l|l|X|}
    \hline
    \textbf{Metadata Name} & \textbf{Metadata Value} & \textbf{Description} \\
    \hline
    \textbf{Reference Id} & \texttt{BEC-92074CE1} & \\
    \hline
    \textbf{Status} & Generated candidate & \\
    \hline
    \textbf{Ancestry} &  & \\
    \hline
    \textbf{Source} & Semantic-pattern candidate & \\
    \hline
    \textbf{Enabled By} &  & \\
    \hline
    \textbf{Supports} &  & \\
    \hline
    \textbf{Derived From} &  & \\
    \hline
    \textbf{Reviewed Status} & Generated candidate: keep / modify / relocate / reject; verify whether independently testable. & \\
    \hline
  \end{tabularx}
\end{table}

Intent. Assignment \& Routing is proposed as a discrete, testable enabling behaviour for Discharge Progression \& Coordination. It is included for architectural review, not as an established implementation obligation.

\noindent\textbf{Feature boundary.} This behaviour is scoped to Discharge Progression \& Coordination; an identically named Feature elsewhere may carry different business meaning.\par
\featuretable{
--- & --- & \emph{Not yet captured} & \emph{Complete during review} \\
}
\subsubsection{Feature - State Tracking}\label{cap:bec-d2873e18}
\vspace{-1em}
\begin{table}[htbp]
  \centering
  \footnotesize
  \begin{tabularx}{1\textwidth}{|l|l|X|}
    \hline
    \textbf{Metadata Name} & \textbf{Metadata Value} & \textbf{Description} \\
    \hline
    \textbf{Reference Id} & \texttt{BEC-D2873E18} & \\
    \hline
    \textbf{Status} & Generated candidate & \\
    \hline
    \textbf{Ancestry} &  & \\
    \hline
    \textbf{Source} & Semantic-pattern candidate & \\
    \hline
    \textbf{Enabled By} &  & \\
    \hline
    \textbf{Supports} &  & \\
    \hline
    \textbf{Derived From} &  & \\
    \hline
    \textbf{Reviewed Status} & Generated candidate: keep / modify / relocate / reject; verify whether independently testable. & \\
    \hline
  \end{tabularx}
\end{table}

Intent. State Tracking is proposed as a discrete, testable enabling behaviour for Discharge Progression \& Coordination. It is included for architectural review, not as an established implementation obligation.

\noindent\textbf{Feature boundary.} This behaviour is scoped to Discharge Progression \& Coordination; an identically named Feature elsewhere may carry different business meaning.\par
\featuretable{
--- & --- & \emph{Not yet captured} & \emph{Complete during review} \\
}
\subsubsection{Feature - Exception Management}\label{cap:bec-af97d3b7}
\vspace{-1em}
\begin{table}[htbp]
  \centering
  \footnotesize
  \begin{tabularx}{1\textwidth}{|l|l|X|}
    \hline
    \textbf{Metadata Name} & \textbf{Metadata Value} & \textbf{Description} \\
    \hline
    \textbf{Reference Id} & \texttt{BEC-AF97D3B7} & \\
    \hline
    \textbf{Status} & Generated candidate & \\
    \hline
    \textbf{Ancestry} &  & \\
    \hline
    \textbf{Source} & Semantic-pattern candidate & \\
    \hline
    \textbf{Enabled By} &  & \\
    \hline
    \textbf{Supports} &  & \\
    \hline
    \textbf{Derived From} &  & \\
    \hline
    \textbf{Reviewed Status} & Generated candidate: keep / modify / relocate / reject; verify whether independently testable. & \\
    \hline
  \end{tabularx}
\end{table}

Intent. Exception Management is proposed as a discrete, testable enabling behaviour for Discharge Progression \& Coordination. It is included for architectural review, not as an established implementation obligation.

\noindent\textbf{Feature boundary.} This behaviour is scoped to Discharge Progression \& Coordination; an identically named Feature elsewhere may carry different business meaning.\par
\featuretable{
--- & --- & \emph{Not yet captured} & \emph{Complete during review} \\
}
\subsubsection{Feature - Outcome Reporting}\label{cap:bec-d56cf336}
\vspace{-1em}
\begin{table}[htbp]
  \centering
  \footnotesize
  \begin{tabularx}{1\textwidth}{|l|l|X|}
    \hline
    \textbf{Metadata Name} & \textbf{Metadata Value} & \textbf{Description} \\
    \hline
    \textbf{Reference Id} & \texttt{BEC-D56CF336} & \\
    \hline
    \textbf{Status} & Generated candidate & \\
    \hline
    \textbf{Ancestry} &  & \\
    \hline
    \textbf{Source} & Semantic-pattern candidate & \\
    \hline
    \textbf{Enabled By} &  & \\
    \hline
    \textbf{Supports} &  & \\
    \hline
    \textbf{Derived From} &  & \\
    \hline
    \textbf{Reviewed Status} & Generated candidate: keep / modify / relocate / reject; verify whether independently testable. & \\
    \hline
  \end{tabularx}
\end{table}

Intent. Outcome Reporting is proposed as a discrete, testable enabling behaviour for Discharge Progression \& Coordination. It is included for architectural review, not as an established implementation obligation.

\noindent\textbf{Feature boundary.} This behaviour is scoped to Discharge Progression \& Coordination; an identically named Feature elsewhere may carry different business meaning.\par
\featuretable{
--- & --- & \emph{Not yet captured} & \emph{Complete during review} \\
}

\chapter{CT1 - Shared Information and Coordination Enablement}\label{cap:bec-03705603}
\noindent\textbf{Reference ID:} \texttt{BEC-03705603}\quad \textbf{Status:} Candidate ancestry\par
Intent. Shared Information and Coordination Enablement is a candidate enabling capability within Business Enabling Capability layer. It frames the reusable system-supported behaviours needed to support this healthcare context without prescribing a product, interface, technology or deployment boundary.

Boundary. This capability supports the business activity and does not inherit the authority or accountability of the Business Capability it enables. Its child Features are review candidates until ancestry and semantics are adjudicated.

\section{CT2 - Patient Clinical Record}\label{cap:bec-97605881}
\vspace{-0.5em}
\begin{table}[htbp]
  \centering
  \footnotesize
  \begin{tabularx}{1\textwidth}{|l|l|X|}
    \hline
    \multicolumn{3}{c}{\textbf{Capability Metadata}} \\
    \hline
    \textbf{Name} & \textbf{Value} & \textbf{Comment and/or Description} \\
    \hline
    \textbf{Reference Id} & \texttt{BEC-97605881} & \\
    \hline
    \textbf{Status} & Candidate ancestry & \\
    \hline
    \textbf{Ancestry} &  & \\
    \hline
    \textbf{Source} &  & \\
    \hline
    \textbf{Enabled By} &  & \\
    \hline
    \textbf{Supports} &  & \\
    \hline
    \textbf{Derived From} &  & \\
    \hline
    \textbf{Reviewed Status} &  & \\
    \hline
  \end{tabularx}
\end{table}

\begin{table}[htbp]
  \centering
  \footnotesize
  \begin{tabularx}{1\textwidth}{|l|l|X|}
    \hline
    \multicolumn{3}{c}{\textbf{Contained Capabilities}} \\
    \hline
    \textbf{Capability Name} & \textbf{Capability Id} & \textbf{Comment} \\
    \hline
    Longitudinal Clinical Record Governance & \texttt{cap:bec-15ab949f} (at \ref{cap:bec-15ab949f}) & \\
    \hline
  \end{tabularx}
\end{table}

Intent. Patient Clinical Record is a candidate enabling capability within Shared Information and Coordination Enablement. It frames the reusable system-supported behaviours needed to support this healthcare context without prescribing a product, interface, technology or deployment boundary.

Boundary. This capability supports the business activity and does not inherit the authority or accountability of the Business Capability it enables. Its child Features are review candidates until ancestry and semantics are adjudicated.

\subsection{CT3 - Longitudinal Clinical Record Governance}\label{cap:bec-15ab949f}
\vspace{-0.5em}
\begin{table}[htbp]
  \centering
  \footnotesize
  \begin{tabularx}{1\textwidth}{|l|l|X|}
    \hline
    \multicolumn{3}{c}{\textbf{Capability Metadata}} \\
    \hline
    \textbf{Name} & \textbf{Value} & \textbf{Comment and/or Description} \\
    \hline
    \textbf{Reference Id} & \texttt{BEC-15AB949F} & \\
    \hline
    \textbf{Status} & Candidate ancestry & \\
    \hline
    \textbf{Ancestry} &  & \\
    \hline
    \textbf{Source} &  & \\
    \hline
    \textbf{Enabled By} &  & \\
    \hline
    \textbf{Supports} &  & \\
    \hline
    \textbf{Derived From} &  & \\
    \hline
    \textbf{Reviewed Status} &  & \\
    \hline
  \end{tabularx}
\end{table}

\begin{table}[htbp]
  \centering
  \footnotesize
  \begin{tabularx}{1\textwidth}{|l|l|X|}
    \hline
    \multicolumn{3}{c}{\textbf{Contained Features}} \\
    \hline
    \textbf{Feature Name} & \textbf{Feature Id} & \textbf{Comment} \\
    \hline
    Longitudinal Health Record (LHR) Assembly & \texttt{cap:bec-b5ebe9a9} (at \ref{cap:bec-b5ebe9a9}) & \\
    \hline
    Active Clinical Record Access & \texttt{cap:bec-81c53e17} (at \ref{cap:bec-81c53e17}) & \\
    \hline
    Durable Clinical Record Preservation & \texttt{cap:bec-a048190c} (at \ref{cap:bec-a048190c}) & \\
    \hline
    Record Context Association & \texttt{cap:bec-9f1c5969} (at \ref{cap:bec-9f1c5969}) & \\
    \hline
    Record Provenance Management & \texttt{cap:bec-3149dd0e} (at \ref{cap:bec-3149dd0e}) & \\
    \hline
    Record Lifecycle Management & \texttt{cap:bec-67c43d3d} (at \ref{cap:bec-67c43d3d}) & \\
    \hline
    Record Discovery & \texttt{cap:bec-f9431119} (at \ref{cap:bec-f9431119}) & \\
    \hline
  \end{tabularx}
\end{table}

Intent. Longitudinal Clinical Record Governance is a candidate enabling capability within Patient Clinical Record. It frames the reusable system-supported behaviours needed to support this healthcare context without prescribing a product, interface, technology or deployment boundary.

Boundary. This capability supports the business activity and does not inherit the authority or accountability of the Business Capability it enables. Its child Features are review candidates until ancestry and semantics are adjudicated.

FHIR R5 conceptual analogues considered (not implementation mappings): Composition, DocumentReference, Provenance.

\subsubsection{Feature - Longitudinal Health Record (LHR) Assembly}\label{cap:bec-b5ebe9a9}
\vspace{-1em}
\begin{table}[htbp]
  \centering
  \footnotesize
  \begin{tabularx}{1\textwidth}{|l|l|X|}
    \hline
    \textbf{Metadata Name} & \textbf{Metadata Value} & \textbf{Description} \\
    \hline
    \textbf{Reference Id} & \texttt{BEC-B5EBE9A9} & \\
    \hline
    \textbf{Status} & Existing Feature; candidate CT3 placement & \\
    \hline
    \textbf{Ancestry} &  & \\
    \hline
    \textbf{Source} & FEAT-ISE-01 & \\
    \hline
    \textbf{Enabled By} &  & \\
    \hline
    \textbf{Supports} &  & \\
    \hline
    \textbf{Derived From} &  & \\
    \hline
    \textbf{Reviewed Status} & Confirm existing Feature meaning and proposed CT3 ownership. & \\
    \hline
  \end{tabularx}
\end{table}

Intent. Longitudinal Health Record (LHR) Assembly provides a distinct enabling behaviour in the context of Longitudinal Clinical Record Governance. This is an existing named Feature; its placement under the candidate CT3 remains subject to review.

\noindent\textbf{Feature boundary.} This behaviour is scoped to Longitudinal Clinical Record Governance; an identically named Feature elsewhere may carry different business meaning.\par
\featuretable{
Function & \texttt{FN-3482AB8B} & Assemble Longitudinal Clinical Record & Catalogue \\
Service & \texttt{SV-C9763495} & Longitudinal Clinical Record Query & Linked to preceding Function in source \\
}
\subsubsection{Feature - Active Clinical Record Access}\label{cap:bec-81c53e17}
\vspace{-1em}
\begin{table}[htbp]
  \centering
  \footnotesize
  \begin{tabularx}{1\textwidth}{|l|l|X|}
    \hline
    \textbf{Metadata Name} & \textbf{Metadata Value} & \textbf{Description} \\
    \hline
    \textbf{Reference Id} & \texttt{BEC-81C53E17} & \\
    \hline
    \textbf{Status} & Existing Feature; candidate CT3 placement & \\
    \hline
    \textbf{Ancestry} &  & \\
    \hline
    \textbf{Source} & FEAT-ISE-02 & \\
    \hline
    \textbf{Enabled By} &  & \\
    \hline
    \textbf{Supports} &  & \\
    \hline
    \textbf{Derived From} &  & \\
    \hline
    \textbf{Reviewed Status} & Confirm existing Feature meaning and proposed CT3 ownership. & \\
    \hline
  \end{tabularx}
\end{table}

Intent. Active Clinical Record Access provides a distinct enabling behaviour in the context of Longitudinal Clinical Record Governance. This is an existing named Feature; its placement under the candidate CT3 remains subject to review.

\noindent\textbf{Feature boundary.} This behaviour is scoped to Longitudinal Clinical Record Governance; an identically named Feature elsewhere may carry different business meaning.\par
\featuretable{
Function & \texttt{FN-1D9A1660} & Maintain Active Clinical Record & Catalogue \\
Service & \texttt{SV-BB337594} & Active Problem \& Allergy Summary Query & Linked to preceding Function in source \\
}
\subsubsection{Feature - Durable Clinical Record Preservation}\label{cap:bec-a048190c}
\vspace{-1em}
\begin{table}[htbp]
  \centering
  \footnotesize
  \begin{tabularx}{1\textwidth}{|l|l|X|}
    \hline
    \textbf{Metadata Name} & \textbf{Metadata Value} & \textbf{Description} \\
    \hline
    \textbf{Reference Id} & \texttt{BEC-A048190C} & \\
    \hline
    \textbf{Status} & Existing Feature; candidate CT3 placement & \\
    \hline
    \textbf{Ancestry} &  & \\
    \hline
    \textbf{Source} & FEAT-ISE-03 & \\
    \hline
    \textbf{Enabled By} &  & \\
    \hline
    \textbf{Supports} &  & \\
    \hline
    \textbf{Derived From} &  & \\
    \hline
    \textbf{Reviewed Status} & Confirm existing Feature meaning and proposed CT3 ownership. & \\
    \hline
  \end{tabularx}
\end{table}

Intent. Durable Clinical Record Preservation provides a distinct enabling behaviour in the context of Longitudinal Clinical Record Governance. This is an existing named Feature; its placement under the candidate CT3 remains subject to review.

\noindent\textbf{Feature boundary.} This behaviour is scoped to Longitudinal Clinical Record Governance; an identically named Feature elsewhere may carry different business meaning.\par
\featuretable{
Function & \texttt{FN-CC4B9D8F} & Preserve Durable Clinical Record & Catalogue \\
Service & \texttt{SV-F882ECD1} & Historical Clinical Record Access & Linked to preceding Function in source \\
}
\subsubsection{Feature - Record Context Association}\label{cap:bec-9f1c5969}
\vspace{-1em}
\begin{table}[htbp]
  \centering
  \footnotesize
  \begin{tabularx}{1\textwidth}{|l|l|X|}
    \hline
    \textbf{Metadata Name} & \textbf{Metadata Value} & \textbf{Description} \\
    \hline
    \textbf{Reference Id} & \texttt{BEC-9F1C5969} & \\
    \hline
    \textbf{Status} & Generated candidate & \\
    \hline
    \textbf{Ancestry} &  & \\
    \hline
    \textbf{Source} & FHIR-informed candidate & \\
    \hline
    \textbf{Enabled By} &  & \\
    \hline
    \textbf{Supports} &  & \\
    \hline
    \textbf{Derived From} &  & \\
    \hline
    \textbf{Reviewed Status} & Generated candidate: keep / modify / relocate / reject; verify whether independently testable. & \\
    \hline
  \end{tabularx}
\end{table}

Intent. Record Context Association is proposed to support the underlying healthcare concept represented by Longitudinal Clinical Record Governance. The proposal concerns behaviour and semantic responsibility, not a FHIR field, operation, resource instance or prescribed data model.

Conceptual evidence: HL7 FHIR R5 Composition, DocumentReference, Provenance.

\noindent\textbf{Feature boundary.} This behaviour is scoped to Longitudinal Clinical Record Governance; an identically named Feature elsewhere may carry different business meaning.\par
\featuretable{
--- & --- & \emph{Not yet captured} & \emph{Complete during review} \\
}
\subsubsection{Feature - Record Provenance Management}\label{cap:bec-3149dd0e}
\vspace{-1em}
\begin{table}[htbp]
  \centering
  \footnotesize
  \begin{tabularx}{1\textwidth}{|l|l|X|}
    \hline
    \textbf{Metadata Name} & \textbf{Metadata Value} & \textbf{Description} \\
    \hline
    \textbf{Reference Id} & \texttt{BEC-3149DD0E} & \\
    \hline
    \textbf{Status} & Generated candidate & \\
    \hline
    \textbf{Ancestry} &  & \\
    \hline
    \textbf{Source} & FHIR-informed candidate & \\
    \hline
    \textbf{Enabled By} &  & \\
    \hline
    \textbf{Supports} &  & \\
    \hline
    \textbf{Derived From} &  & \\
    \hline
    \textbf{Reviewed Status} & Generated candidate: keep / modify / relocate / reject; verify whether independently testable. & \\
    \hline
  \end{tabularx}
\end{table}

Intent. Record Provenance Management is proposed to support the underlying healthcare concept represented by Longitudinal Clinical Record Governance. The proposal concerns behaviour and semantic responsibility, not a FHIR field, operation, resource instance or prescribed data model.

Conceptual evidence: HL7 FHIR R5 Composition, DocumentReference, Provenance.

\noindent\textbf{Feature boundary.} This behaviour is scoped to Longitudinal Clinical Record Governance; an identically named Feature elsewhere may carry different business meaning.\par
\featuretable{
--- & --- & \emph{Not yet captured} & \emph{Complete during review} \\
}
\subsubsection{Feature - Record Lifecycle Management}\label{cap:bec-67c43d3d}
\vspace{-1em}
\begin{table}[htbp]
  \centering
  \footnotesize
  \begin{tabularx}{1\textwidth}{|l|l|X|}
    \hline
    \textbf{Metadata Name} & \textbf{Metadata Value} & \textbf{Description} \\
    \hline
    \textbf{Reference Id} & \texttt{BEC-67C43D3D} & \\
    \hline
    \textbf{Status} & Generated candidate & \\
    \hline
    \textbf{Ancestry} &  & \\
    \hline
    \textbf{Source} & FHIR-informed candidate & \\
    \hline
    \textbf{Enabled By} &  & \\
    \hline
    \textbf{Supports} &  & \\
    \hline
    \textbf{Derived From} &  & \\
    \hline
    \textbf{Reviewed Status} & Generated candidate: keep / modify / relocate / reject; verify whether independently testable. & \\
    \hline
  \end{tabularx}
\end{table}

Intent. Record Lifecycle Management is proposed to support the underlying healthcare concept represented by Longitudinal Clinical Record Governance. The proposal concerns behaviour and semantic responsibility, not a FHIR field, operation, resource instance or prescribed data model.

Conceptual evidence: HL7 FHIR R5 Composition, DocumentReference, Provenance.

\noindent\textbf{Feature boundary.} This behaviour is scoped to Longitudinal Clinical Record Governance; an identically named Feature elsewhere may carry different business meaning.\par
\featuretable{
--- & --- & \emph{Not yet captured} & \emph{Complete during review} \\
}
\subsubsection{Feature - Record Discovery}\label{cap:bec-f9431119}
\vspace{-1em}
\begin{table}[htbp]
  \centering
  \footnotesize
  \begin{tabularx}{1\textwidth}{|l|l|X|}
    \hline
    \textbf{Metadata Name} & \textbf{Metadata Value} & \textbf{Description} \\
    \hline
    \textbf{Reference Id} & \texttt{BEC-F9431119} & \\
    \hline
    \textbf{Status} & Generated candidate & \\
    \hline
    \textbf{Ancestry} &  & \\
    \hline
    \textbf{Source} & FHIR-informed candidate & \\
    \hline
    \textbf{Enabled By} &  & \\
    \hline
    \textbf{Supports} &  & \\
    \hline
    \textbf{Derived From} &  & \\
    \hline
    \textbf{Reviewed Status} & Generated candidate: keep / modify / relocate / reject; verify whether independently testable. & \\
    \hline
  \end{tabularx}
\end{table}

Intent. Record Discovery is proposed to support the underlying healthcare concept represented by Longitudinal Clinical Record Governance. The proposal concerns behaviour and semantic responsibility, not a FHIR field, operation, resource instance or prescribed data model.

Conceptual evidence: HL7 FHIR R5 Composition, DocumentReference, Provenance.

\noindent\textbf{Feature boundary.} This behaviour is scoped to Longitudinal Clinical Record Governance; an identically named Feature elsewhere may carry different business meaning.\par
\featuretable{
--- & --- & \emph{Not yet captured} & \emph{Complete during review} \\
}

\section{CT2 - HIE Integration Services}\label{cap:bec-14bbaafd}
\vspace{-0.5em}
\begin{table}[htbp]
  \centering
  \footnotesize
  \begin{tabularx}{1\textwidth}{|l|l|X|}
    \hline
    \multicolumn{3}{c}{\textbf{Capability Metadata}} \\
    \hline
    \textbf{Name} & \textbf{Value} & \textbf{Comment and/or Description} \\
    \hline
    \textbf{Reference Id} & \texttt{BEC-14BBAAFD} & \\
    \hline
    \textbf{Status} & Candidate ancestry & \\
    \hline
    \textbf{Ancestry} &  & \\
    \hline
    \textbf{Source} &  & \\
    \hline
    \textbf{Enabled By} &  & \\
    \hline
    \textbf{Supports} &  & \\
    \hline
    \textbf{Derived From} &  & \\
    \hline
    \textbf{Reviewed Status} &  & \\
    \hline
  \end{tabularx}
\end{table}

\begin{table}[htbp]
  \centering
  \footnotesize
  \begin{tabularx}{1\textwidth}{|l|l|X|}
    \hline
    \multicolumn{3}{c}{\textbf{Contained Capabilities}} \\
    \hline
    \textbf{Capability Name} & \textbf{Capability Id} & \textbf{Comment} \\
    \hline
    Multi-Model Information Exchange & \texttt{cap:bec-0fade98f} (at \ref{cap:bec-0fade98f}) & \\
    \hline
  \end{tabularx}
\end{table}

Intent. HIE Integration Services is a candidate enabling capability within Shared Information and Coordination Enablement. It frames the reusable system-supported behaviours needed to support this healthcare context without prescribing a product, interface, technology or deployment boundary.

Boundary. This capability supports the business activity and does not inherit the authority or accountability of the Business Capability it enables. Its child Features are review candidates until ancestry and semantics are adjudicated.

\subsection{CT3 - Multi-Model Information Exchange}\label{cap:bec-0fade98f}
\vspace{-0.5em}
\begin{table}[htbp]
  \centering
  \footnotesize
  \begin{tabularx}{1\textwidth}{|l|l|X|}
    \hline
    \multicolumn{3}{c}{\textbf{Capability Metadata}} \\
    \hline
    \textbf{Name} & \textbf{Value} & \textbf{Comment and/or Description} \\
    \hline
    \textbf{Reference Id} & \texttt{BEC-0FADE98F} & \\
    \hline
    \textbf{Status} & Candidate ancestry & \\
    \hline
    \textbf{Ancestry} &  & \\
    \hline
    \textbf{Source} &  & \\
    \hline
    \textbf{Enabled By} &  & \\
    \hline
    \textbf{Supports} &  & \\
    \hline
    \textbf{Derived From} &  & \\
    \hline
    \textbf{Reviewed Status} &  & \\
    \hline
  \end{tabularx}
\end{table}

\begin{table}[htbp]
  \centering
  \footnotesize
  \begin{tabularx}{1\textwidth}{|l|l|X|}
    \hline
    \multicolumn{3}{c}{\textbf{Contained Features}} \\
    \hline
    \textbf{Feature Name} & \textbf{Feature Id} & \textbf{Comment} \\
    \hline
    Direct Submission Ingestion \& Delivery & \texttt{cap:bec-a4404077} (at \ref{cap:bec-a4404077}) & \\
    \hline
    Clinical Record Retrieval Mediation & \texttt{cap:bec-3110386a} (at \ref{cap:bec-3110386a}) & \\
    \hline
    Addressed Clinical Distribution & \texttt{cap:bec-e1974228} (at \ref{cap:bec-e1974228}) & \\
    \hline
    Event-Driven Syndication & \texttt{cap:bec-05fad383} (at \ref{cap:bec-05fad383}) & \\
    \hline
    Information Validation & \texttt{cap:bec-e86070e3} (at \ref{cap:bec-e86070e3}) & \\
    \hline
    Information Correlation & \texttt{cap:bec-9924a266} (at \ref{cap:bec-9924a266}) & \\
    \hline
    Information Retrieval & \texttt{cap:bec-c048b88a} (at \ref{cap:bec-c048b88a}) & \\
    \hline
    Information Provenance & \texttt{cap:bec-f497fe5f} (at \ref{cap:bec-f497fe5f}) & \\
    \hline
    Information Lifecycle Management & \texttt{cap:bec-90eed6cf} (at \ref{cap:bec-90eed6cf}) & \\
    \hline
  \end{tabularx}
\end{table}

Intent. Multi-Model Information Exchange is a candidate enabling capability within HIE Integration Services. It frames the reusable system-supported behaviours needed to support this healthcare context without prescribing a product, interface, technology or deployment boundary.

Boundary. This capability supports the business activity and does not inherit the authority or accountability of the Business Capability it enables. Its child Features are review candidates until ancestry and semantics are adjudicated.

\subsubsection{Feature - Direct Submission Ingestion \& Delivery}\label{cap:bec-a4404077}
\vspace{-1em}
\begin{table}[htbp]
  \centering
  \footnotesize
  \begin{tabularx}{1\textwidth}{|l|l|X|}
    \hline
    \textbf{Metadata Name} & \textbf{Metadata Value} & \textbf{Description} \\
    \hline
    \textbf{Reference Id} & \texttt{BEC-A4404077} & \\
    \hline
    \textbf{Status} & Existing Feature; candidate CT3 placement & \\
    \hline
    \textbf{Ancestry} &  & \\
    \hline
    \textbf{Source} & FEAT-ISE-04 & \\
    \hline
    \textbf{Enabled By} &  & \\
    \hline
    \textbf{Supports} &  & \\
    \hline
    \textbf{Derived From} &  & \\
    \hline
    \textbf{Reviewed Status} & Confirm existing Feature meaning and proposed CT3 ownership. & \\
    \hline
  \end{tabularx}
\end{table}

Intent. Direct Submission Ingestion \& Delivery provides a distinct enabling behaviour in the context of Multi-Model Information Exchange. This is an existing named Feature; its placement under the candidate CT3 remains subject to review.

\noindent\textbf{Feature boundary.} This behaviour is scoped to Multi-Model Information Exchange; an identically named Feature elsewhere may carry different business meaning.\par
\featuretable{
Function & \texttt{FN-9DBF8DF4} & Receive Information Submission & Catalogue \\
Service & \texttt{SV-5D49B817} & HIE Submission Gateway Service & Linked to preceding Function in source \\
}
\subsubsection{Feature - Clinical Record Retrieval Mediation}\label{cap:bec-3110386a}
\vspace{-1em}
\begin{table}[htbp]
  \centering
  \footnotesize
  \begin{tabularx}{1\textwidth}{|l|l|X|}
    \hline
    \textbf{Metadata Name} & \textbf{Metadata Value} & \textbf{Description} \\
    \hline
    \textbf{Reference Id} & \texttt{BEC-3110386A} & \\
    \hline
    \textbf{Status} & Existing Feature; candidate CT3 placement & \\
    \hline
    \textbf{Ancestry} &  & \\
    \hline
    \textbf{Source} & FEAT-ISE-05 & \\
    \hline
    \textbf{Enabled By} &  & \\
    \hline
    \textbf{Supports} &  & \\
    \hline
    \textbf{Derived From} &  & \\
    \hline
    \textbf{Reviewed Status} & Confirm existing Feature meaning and proposed CT3 ownership. & \\
    \hline
  \end{tabularx}
\end{table}

Intent. Clinical Record Retrieval Mediation provides a distinct enabling behaviour in the context of Multi-Model Information Exchange. This is an existing named Feature; its placement under the candidate CT3 remains subject to review.

\noindent\textbf{Feature boundary.} This behaviour is scoped to Multi-Model Information Exchange; an identically named Feature elsewhere may carry different business meaning.\par
\featuretable{
Function & \texttt{FN-03C45006} & Mediate Information Retrieval & Catalogue \\
Service & \texttt{SV-D8396148} & Federated Health Information Query & Linked to preceding Function in source \\
}
\subsubsection{Feature - Addressed Clinical Distribution}\label{cap:bec-e1974228}
\vspace{-1em}
\begin{table}[htbp]
  \centering
  \footnotesize
  \begin{tabularx}{1\textwidth}{|l|l|X|}
    \hline
    \textbf{Metadata Name} & \textbf{Metadata Value} & \textbf{Description} \\
    \hline
    \textbf{Reference Id} & \texttt{BEC-E1974228} & \\
    \hline
    \textbf{Status} & Existing Feature; candidate CT3 placement & \\
    \hline
    \textbf{Ancestry} &  & \\
    \hline
    \textbf{Source} & FEAT-ISE-06 & \\
    \hline
    \textbf{Enabled By} &  & \\
    \hline
    \textbf{Supports} &  & \\
    \hline
    \textbf{Derived From} &  & \\
    \hline
    \textbf{Reviewed Status} & Confirm existing Feature meaning and proposed CT3 ownership. & \\
    \hline
  \end{tabularx}
\end{table}

Intent. Addressed Clinical Distribution provides a distinct enabling behaviour in the context of Multi-Model Information Exchange. This is an existing named Feature; its placement under the candidate CT3 remains subject to review.

\noindent\textbf{Feature boundary.} This behaviour is scoped to Multi-Model Information Exchange; an identically named Feature elsewhere may carry different business meaning.\par
\featuretable{
Function & \texttt{FN-AF7EAF51} & Distribute Addressed Information & Catalogue \\
Service & \texttt{SV-B3CCC084} & Addressed Document Delivery Service & Linked to preceding Function in source \\
}
\subsubsection{Feature - Event-Driven Syndication}\label{cap:bec-05fad383}
\vspace{-1em}
\begin{table}[htbp]
  \centering
  \footnotesize
  \begin{tabularx}{1\textwidth}{|l|l|X|}
    \hline
    \textbf{Metadata Name} & \textbf{Metadata Value} & \textbf{Description} \\
    \hline
    \textbf{Reference Id} & \texttt{BEC-05FAD383} & \\
    \hline
    \textbf{Status} & Existing Feature; candidate CT3 placement & \\
    \hline
    \textbf{Ancestry} &  & \\
    \hline
    \textbf{Source} & FEAT-ISE-07 & \\
    \hline
    \textbf{Enabled By} &  & \\
    \hline
    \textbf{Supports} &  & \\
    \hline
    \textbf{Derived From} &  & \\
    \hline
    \textbf{Reviewed Status} & Confirm existing Feature meaning and proposed CT3 ownership. & \\
    \hline
  \end{tabularx}
\end{table}

Intent. Event-Driven Syndication provides a distinct enabling behaviour in the context of Multi-Model Information Exchange. This is an existing named Feature; its placement under the candidate CT3 remains subject to review.

\noindent\textbf{Feature boundary.} This behaviour is scoped to Multi-Model Information Exchange; an identically named Feature elsewhere may carry different business meaning.\par
\featuretable{
Function & \texttt{FN-27FFC9FD} & Syndicate Information Change & Catalogue \\
Service & \texttt{SV-BFE23E19} & Clinical Event Syndication Stream & Linked to preceding Function in source \\
}
\subsubsection{Feature - Information Validation}\label{cap:bec-e86070e3}
\vspace{-1em}
\begin{table}[htbp]
  \centering
  \footnotesize
  \begin{tabularx}{1\textwidth}{|l|l|X|}
    \hline
    \textbf{Metadata Name} & \textbf{Metadata Value} & \textbf{Description} \\
    \hline
    \textbf{Reference Id} & \texttt{BEC-E86070E3} & \\
    \hline
    \textbf{Status} & Generated candidate & \\
    \hline
    \textbf{Ancestry} &  & \\
    \hline
    \textbf{Source} & Semantic-pattern candidate & \\
    \hline
    \textbf{Enabled By} &  & \\
    \hline
    \textbf{Supports} &  & \\
    \hline
    \textbf{Derived From} &  & \\
    \hline
    \textbf{Reviewed Status} & Generated candidate: keep / modify / relocate / reject; verify whether independently testable. & \\
    \hline
  \end{tabularx}
\end{table}

Intent. Information Validation is proposed as a discrete, testable enabling behaviour for Multi-Model Information Exchange. It is included for architectural review, not as an established implementation obligation.

\noindent\textbf{Feature boundary.} This behaviour is scoped to Multi-Model Information Exchange; an identically named Feature elsewhere may carry different business meaning.\par
\featuretable{
--- & --- & \emph{Not yet captured} & \emph{Complete during review} \\
}
\subsubsection{Feature - Information Correlation}\label{cap:bec-9924a266}
\vspace{-1em}
\begin{table}[htbp]
  \centering
  \footnotesize
  \begin{tabularx}{1\textwidth}{|l|l|X|}
    \hline
    \textbf{Metadata Name} & \textbf{Metadata Value} & \textbf{Description} \\
    \hline
    \textbf{Reference Id} & \texttt{BEC-9924A266} & \\
    \hline
    \textbf{Status} & Generated candidate & \\
    \hline
    \textbf{Ancestry} &  & \\
    \hline
    \textbf{Source} & Semantic-pattern candidate & \\
    \hline
    \textbf{Enabled By} &  & \\
    \hline
    \textbf{Supports} &  & \\
    \hline
    \textbf{Derived From} &  & \\
    \hline
    \textbf{Reviewed Status} & Generated candidate: keep / modify / relocate / reject; verify whether independently testable. & \\
    \hline
  \end{tabularx}
\end{table}

Intent. Information Correlation is proposed as a discrete, testable enabling behaviour for Multi-Model Information Exchange. It is included for architectural review, not as an established implementation obligation.

\noindent\textbf{Feature boundary.} This behaviour is scoped to Multi-Model Information Exchange; an identically named Feature elsewhere may carry different business meaning.\par
\featuretable{
--- & --- & \emph{Not yet captured} & \emph{Complete during review} \\
}
\subsubsection{Feature - Information Retrieval}\label{cap:bec-c048b88a}
\vspace{-1em}
\begin{table}[htbp]
  \centering
  \footnotesize
  \begin{tabularx}{1\textwidth}{|l|l|X|}
    \hline
    \textbf{Metadata Name} & \textbf{Metadata Value} & \textbf{Description} \\
    \hline
    \textbf{Reference Id} & \texttt{BEC-C048B88A} & \\
    \hline
    \textbf{Status} & Generated candidate & \\
    \hline
    \textbf{Ancestry} &  & \\
    \hline
    \textbf{Source} & Semantic-pattern candidate & \\
    \hline
    \textbf{Enabled By} &  & \\
    \hline
    \textbf{Supports} &  & \\
    \hline
    \textbf{Derived From} &  & \\
    \hline
    \textbf{Reviewed Status} & Generated candidate: keep / modify / relocate / reject; verify whether independently testable. & \\
    \hline
  \end{tabularx}
\end{table}

Intent. Information Retrieval is proposed as a discrete, testable enabling behaviour for Multi-Model Information Exchange. It is included for architectural review, not as an established implementation obligation.

\noindent\textbf{Feature boundary.} This behaviour is scoped to Multi-Model Information Exchange; an identically named Feature elsewhere may carry different business meaning.\par
\featuretable{
--- & --- & \emph{Not yet captured} & \emph{Complete during review} \\
}
\subsubsection{Feature - Information Provenance}\label{cap:bec-f497fe5f}
\vspace{-1em}
\begin{table}[htbp]
  \centering
  \footnotesize
  \begin{tabularx}{1\textwidth}{|l|l|X|}
    \hline
    \textbf{Metadata Name} & \textbf{Metadata Value} & \textbf{Description} \\
    \hline
    \textbf{Reference Id} & \texttt{BEC-F497FE5F} & \\
    \hline
    \textbf{Status} & Generated candidate & \\
    \hline
    \textbf{Ancestry} &  & \\
    \hline
    \textbf{Source} & Semantic-pattern candidate & \\
    \hline
    \textbf{Enabled By} &  & \\
    \hline
    \textbf{Supports} &  & \\
    \hline
    \textbf{Derived From} &  & \\
    \hline
    \textbf{Reviewed Status} & Generated candidate: keep / modify / relocate / reject; verify whether independently testable. & \\
    \hline
  \end{tabularx}
\end{table}

Intent. Information Provenance is proposed as a discrete, testable enabling behaviour for Multi-Model Information Exchange. It is included for architectural review, not as an established implementation obligation.

\noindent\textbf{Feature boundary.} This behaviour is scoped to Multi-Model Information Exchange; an identically named Feature elsewhere may carry different business meaning.\par
\featuretable{
--- & --- & \emph{Not yet captured} & \emph{Complete during review} \\
}
\subsubsection{Feature - Information Lifecycle Management}\label{cap:bec-90eed6cf}
\vspace{-1em}
\begin{table}[htbp]
  \centering
  \footnotesize
  \begin{tabularx}{1\textwidth}{|l|l|X|}
    \hline
    \textbf{Metadata Name} & \textbf{Metadata Value} & \textbf{Description} \\
    \hline
    \textbf{Reference Id} & \texttt{BEC-90EED6CF} & \\
    \hline
    \textbf{Status} & Generated candidate & \\
    \hline
    \textbf{Ancestry} &  & \\
    \hline
    \textbf{Source} & Semantic-pattern candidate & \\
    \hline
    \textbf{Enabled By} &  & \\
    \hline
    \textbf{Supports} &  & \\
    \hline
    \textbf{Derived From} &  & \\
    \hline
    \textbf{Reviewed Status} & Generated candidate: keep / modify / relocate / reject; verify whether independently testable. & \\
    \hline
  \end{tabularx}
\end{table}

Intent. Information Lifecycle Management is proposed as a discrete, testable enabling behaviour for Multi-Model Information Exchange. It is included for architectural review, not as an established implementation obligation.

\noindent\textbf{Feature boundary.} This behaviour is scoped to Multi-Model Information Exchange; an identically named Feature elsewhere may carry different business meaning.\par
\featuretable{
--- & --- & \emph{Not yet captured} & \emph{Complete during review} \\
}

\section{CT2 - Health Information Access}\label{cap:bec-0ee16b9d}
\vspace{-0.5em}
\begin{table}[htbp]
  \centering
  \footnotesize
  \begin{tabularx}{1\textwidth}{|l|l|X|}
    \hline
    \multicolumn{3}{c}{\textbf{Capability Metadata}} \\
    \hline
    \textbf{Name} & \textbf{Value} & \textbf{Comment and/or Description} \\
    \hline
    \textbf{Reference Id} & \texttt{BEC-0EE16B9D} & \\
    \hline
    \textbf{Status} & Candidate ancestry & \\
    \hline
    \textbf{Ancestry} &  & \\
    \hline
    \textbf{Source} &  & \\
    \hline
    \textbf{Enabled By} &  & \\
    \hline
    \textbf{Supports} &  & \\
    \hline
    \textbf{Derived From} &  & \\
    \hline
    \textbf{Reviewed Status} &  & \\
    \hline
  \end{tabularx}
\end{table}

\begin{table}[htbp]
  \centering
  \footnotesize
  \begin{tabularx}{1\textwidth}{|l|l|X|}
    \hline
    \multicolumn{3}{c}{\textbf{Contained Capabilities}} \\
    \hline
    \textbf{Capability Name} & \textbf{Capability Id} & \textbf{Comment} \\
    \hline
    Governed Query \& Search & \texttt{cap:bec-73452901} (at \ref{cap:bec-73452901}) & \\
    \hline
  \end{tabularx}
\end{table}

Intent. Health Information Access is a candidate enabling capability within Shared Information and Coordination Enablement. It frames the reusable system-supported behaviours needed to support this healthcare context without prescribing a product, interface, technology or deployment boundary.

Boundary. This capability supports the business activity and does not inherit the authority or accountability of the Business Capability it enables. Its child Features are review candidates until ancestry and semantics are adjudicated.

\subsection{CT3 - Governed Query \& Search}\label{cap:bec-73452901}
\vspace{-0.5em}
\begin{table}[htbp]
  \centering
  \footnotesize
  \begin{tabularx}{1\textwidth}{|l|l|X|}
    \hline
    \multicolumn{3}{c}{\textbf{Capability Metadata}} \\
    \hline
    \textbf{Name} & \textbf{Value} & \textbf{Comment and/or Description} \\
    \hline
    \textbf{Reference Id} & \texttt{BEC-73452901} & \\
    \hline
    \textbf{Status} & Candidate ancestry & \\
    \hline
    \textbf{Ancestry} &  & \\
    \hline
    \textbf{Source} &  & \\
    \hline
    \textbf{Enabled By} &  & \\
    \hline
    \textbf{Supports} &  & \\
    \hline
    \textbf{Derived From} &  & \\
    \hline
    \textbf{Reviewed Status} &  & \\
    \hline
  \end{tabularx}
\end{table}

\begin{table}[htbp]
  \centering
  \footnotesize
  \begin{tabularx}{1\textwidth}{|l|l|X|}
    \hline
    \multicolumn{3}{c}{\textbf{Contained Features}} \\
    \hline
    \textbf{Feature Name} & \textbf{Feature Id} & \textbf{Comment} \\
    \hline
    Federated Clinical Query & \texttt{cap:bec-ab8edd15} (at \ref{cap:bec-ab8edd15}) & \\
    \hline
    Access-Controlled Information Filtering & \texttt{cap:bec-4601c816} (at \ref{cap:bec-4601c816}) & \\
    \hline
    Context Establishment & \texttt{cap:bec-72a9140c} (at \ref{cap:bec-72a9140c}) & \\
    \hline
    State Monitoring & \texttt{cap:bec-474c9605} (at \ref{cap:bec-474c9605}) & \\
    \hline
    Exception Management & \texttt{cap:bec-5604785c} (at \ref{cap:bec-5604785c}) & \\
    \hline
    Information Retrieval & \texttt{cap:bec-c138a5d7} (at \ref{cap:bec-c138a5d7}) & \\
    \hline
  \end{tabularx}
\end{table}

Intent. Governed Query \& Search is a candidate enabling capability within Health Information Access. It frames the reusable system-supported behaviours needed to support this healthcare context without prescribing a product, interface, technology or deployment boundary.

Boundary. This capability supports the business activity and does not inherit the authority or accountability of the Business Capability it enables. Its child Features are review candidates until ancestry and semantics are adjudicated.

\subsubsection{Feature - Federated Clinical Query}\label{cap:bec-ab8edd15}
\vspace{-1em}
\begin{table}[htbp]
  \centering
  \footnotesize
  \begin{tabularx}{1\textwidth}{|l|l|X|}
    \hline
    \textbf{Metadata Name} & \textbf{Metadata Value} & \textbf{Description} \\
    \hline
    \textbf{Reference Id} & \texttt{BEC-AB8EDD15} & \\
    \hline
    \textbf{Status} & Existing Feature; candidate CT3 placement & \\
    \hline
    \textbf{Ancestry} &  & \\
    \hline
    \textbf{Source} & FEAT-ISE-08 & \\
    \hline
    \textbf{Enabled By} &  & \\
    \hline
    \textbf{Supports} &  & \\
    \hline
    \textbf{Derived From} &  & \\
    \hline
    \textbf{Reviewed Status} & Confirm existing Feature meaning and proposed CT3 ownership. & \\
    \hline
  \end{tabularx}
\end{table}

Intent. Federated Clinical Query provides a distinct enabling behaviour in the context of Governed Query \& Search. This is an existing named Feature; its placement under the candidate CT3 remains subject to review.

\noindent\textbf{Feature boundary.} This behaviour is scoped to Governed Query \& Search; an identically named Feature elsewhere may carry different business meaning.\par
\featuretable{
Function & \texttt{FN-F042FC3C} & Search Health Information & Catalogue \\
Service & \texttt{SV-5A341035} & Clinical Information Search Service & Linked to preceding Function in source \\
}
\subsubsection{Feature - Access-Controlled Information Filtering}\label{cap:bec-4601c816}
\vspace{-1em}
\begin{table}[htbp]
  \centering
  \footnotesize
  \begin{tabularx}{1\textwidth}{|l|l|X|}
    \hline
    \textbf{Metadata Name} & \textbf{Metadata Value} & \textbf{Description} \\
    \hline
    \textbf{Reference Id} & \texttt{BEC-4601C816} & \\
    \hline
    \textbf{Status} & Existing Feature; candidate CT3 placement & \\
    \hline
    \textbf{Ancestry} &  & \\
    \hline
    \textbf{Source} & FEAT-ISE-09 & \\
    \hline
    \textbf{Enabled By} &  & \\
    \hline
    \textbf{Supports} &  & \\
    \hline
    \textbf{Derived From} &  & \\
    \hline
    \textbf{Reviewed Status} & Confirm existing Feature meaning and proposed CT3 ownership. & \\
    \hline
  \end{tabularx}
\end{table}

Intent. Access-Controlled Information Filtering provides a distinct enabling behaviour in the context of Governed Query \& Search. This is an existing named Feature; its placement under the candidate CT3 remains subject to review.

\noindent\textbf{Feature boundary.} This behaviour is scoped to Governed Query \& Search; an identically named Feature elsewhere may carry different business meaning.\par
\featuretable{
Function & \texttt{FN-56260EE4} & Retrieve Health Information & Catalogue \\
Service & \texttt{SV-429EF479} & Clinical Resource Retrieval Service & Linked to preceding Function in source \\
}
\subsubsection{Feature - Context Establishment}\label{cap:bec-72a9140c}
\vspace{-1em}
\begin{table}[htbp]
  \centering
  \footnotesize
  \begin{tabularx}{1\textwidth}{|l|l|X|}
    \hline
    \textbf{Metadata Name} & \textbf{Metadata Value} & \textbf{Description} \\
    \hline
    \textbf{Reference Id} & \texttt{BEC-72A9140C} & \\
    \hline
    \textbf{Status} & Generated candidate & \\
    \hline
    \textbf{Ancestry} &  & \\
    \hline
    \textbf{Source} & Semantic-pattern candidate & \\
    \hline
    \textbf{Enabled By} &  & \\
    \hline
    \textbf{Supports} &  & \\
    \hline
    \textbf{Derived From} &  & \\
    \hline
    \textbf{Reviewed Status} & Generated candidate: keep / modify / relocate / reject; verify whether independently testable. & \\
    \hline
  \end{tabularx}
\end{table}

Intent. Context Establishment is proposed as a discrete, testable enabling behaviour for Governed Query \& Search. It is included for architectural review, not as an established implementation obligation.

\noindent\textbf{Feature boundary.} This behaviour is scoped to Governed Query \& Search; an identically named Feature elsewhere may carry different business meaning.\par
\featuretable{
--- & --- & \emph{Not yet captured} & \emph{Complete during review} \\
}
\subsubsection{Feature - State Monitoring}\label{cap:bec-474c9605}
\vspace{-1em}
\begin{table}[htbp]
  \centering
  \footnotesize
  \begin{tabularx}{1\textwidth}{|l|l|X|}
    \hline
    \textbf{Metadata Name} & \textbf{Metadata Value} & \textbf{Description} \\
    \hline
    \textbf{Reference Id} & \texttt{BEC-474C9605} & \\
    \hline
    \textbf{Status} & Generated candidate & \\
    \hline
    \textbf{Ancestry} &  & \\
    \hline
    \textbf{Source} & Semantic-pattern candidate & \\
    \hline
    \textbf{Enabled By} &  & \\
    \hline
    \textbf{Supports} &  & \\
    \hline
    \textbf{Derived From} &  & \\
    \hline
    \textbf{Reviewed Status} & Generated candidate: keep / modify / relocate / reject; verify whether independently testable. & \\
    \hline
  \end{tabularx}
\end{table}

Intent. State Monitoring is proposed as a discrete, testable enabling behaviour for Governed Query \& Search. It is included for architectural review, not as an established implementation obligation.

\noindent\textbf{Feature boundary.} This behaviour is scoped to Governed Query \& Search; an identically named Feature elsewhere may carry different business meaning.\par
\featuretable{
--- & --- & \emph{Not yet captured} & \emph{Complete during review} \\
}
\subsubsection{Feature - Exception Management}\label{cap:bec-5604785c}
\vspace{-1em}
\begin{table}[htbp]
  \centering
  \footnotesize
  \begin{tabularx}{1\textwidth}{|l|l|X|}
    \hline
    \textbf{Metadata Name} & \textbf{Metadata Value} & \textbf{Description} \\
    \hline
    \textbf{Reference Id} & \texttt{BEC-5604785C} & \\
    \hline
    \textbf{Status} & Generated candidate & \\
    \hline
    \textbf{Ancestry} &  & \\
    \hline
    \textbf{Source} & Semantic-pattern candidate & \\
    \hline
    \textbf{Enabled By} &  & \\
    \hline
    \textbf{Supports} &  & \\
    \hline
    \textbf{Derived From} &  & \\
    \hline
    \textbf{Reviewed Status} & Generated candidate: keep / modify / relocate / reject; verify whether independently testable. & \\
    \hline
  \end{tabularx}
\end{table}

Intent. Exception Management is proposed as a discrete, testable enabling behaviour for Governed Query \& Search. It is included for architectural review, not as an established implementation obligation.

\noindent\textbf{Feature boundary.} This behaviour is scoped to Governed Query \& Search; an identically named Feature elsewhere may carry different business meaning.\par
\featuretable{
--- & --- & \emph{Not yet captured} & \emph{Complete during review} \\
}
\subsubsection{Feature - Information Retrieval}\label{cap:bec-c138a5d7}
\vspace{-1em}
\begin{table}[htbp]
  \centering
  \footnotesize
  \begin{tabularx}{1\textwidth}{|l|l|X|}
    \hline
    \textbf{Metadata Name} & \textbf{Metadata Value} & \textbf{Description} \\
    \hline
    \textbf{Reference Id} & \texttt{BEC-C138A5D7} & \\
    \hline
    \textbf{Status} & Generated candidate & \\
    \hline
    \textbf{Ancestry} &  & \\
    \hline
    \textbf{Source} & Semantic-pattern candidate & \\
    \hline
    \textbf{Enabled By} &  & \\
    \hline
    \textbf{Supports} &  & \\
    \hline
    \textbf{Derived From} &  & \\
    \hline
    \textbf{Reviewed Status} & Generated candidate: keep / modify / relocate / reject; verify whether independently testable. & \\
    \hline
  \end{tabularx}
\end{table}

Intent. Information Retrieval is proposed as a discrete, testable enabling behaviour for Governed Query \& Search. It is included for architectural review, not as an established implementation obligation.

\noindent\textbf{Feature boundary.} This behaviour is scoped to Governed Query \& Search; an identically named Feature elsewhere may carry different business meaning.\par
\featuretable{
--- & --- & \emph{Not yet captured} & \emph{Complete during review} \\
}

\section{CT2 - Health Information Communication}\label{cap:bec-4553b47b}
\vspace{-0.5em}
\begin{table}[htbp]
  \centering
  \footnotesize
  \begin{tabularx}{1\textwidth}{|l|l|X|}
    \hline
    \multicolumn{3}{c}{\textbf{Capability Metadata}} \\
    \hline
    \textbf{Name} & \textbf{Value} & \textbf{Comment and/or Description} \\
    \hline
    \textbf{Reference Id} & \texttt{BEC-4553B47B} & \\
    \hline
    \textbf{Status} & Candidate ancestry & \\
    \hline
    \textbf{Ancestry} &  & \\
    \hline
    \textbf{Source} &  & \\
    \hline
    \textbf{Enabled By} &  & \\
    \hline
    \textbf{Supports} &  & \\
    \hline
    \textbf{Derived From} &  & \\
    \hline
    \textbf{Reviewed Status} &  & \\
    \hline
  \end{tabularx}
\end{table}

\begin{table}[htbp]
  \centering
  \footnotesize
  \begin{tabularx}{1\textwidth}{|l|l|X|}
    \hline
    \multicolumn{3}{c}{\textbf{Contained Capabilities}} \\
    \hline
    \textbf{Capability Name} & \textbf{Capability Id} & \textbf{Comment} \\
    \hline
    Boundary Protocol Mediation & \texttt{cap:bec-763876a9} (at \ref{cap:bec-763876a9}) & \\
    \hline
  \end{tabularx}
\end{table}

Intent. Health Information Communication is a candidate enabling capability within Shared Information and Coordination Enablement. It frames the reusable system-supported behaviours needed to support this healthcare context without prescribing a product, interface, technology or deployment boundary.

Boundary. This capability supports the business activity and does not inherit the authority or accountability of the Business Capability it enables. Its child Features are review candidates until ancestry and semantics are adjudicated.

\subsection{CT3 - Boundary Protocol Mediation}\label{cap:bec-763876a9}
\vspace{-0.5em}
\begin{table}[htbp]
  \centering
  \footnotesize
  \begin{tabularx}{1\textwidth}{|l|l|X|}
    \hline
    \multicolumn{3}{c}{\textbf{Capability Metadata}} \\
    \hline
    \textbf{Name} & \textbf{Value} & \textbf{Comment and/or Description} \\
    \hline
    \textbf{Reference Id} & \texttt{BEC-763876A9} & \\
    \hline
    \textbf{Status} & Candidate ancestry & \\
    \hline
    \textbf{Ancestry} &  & \\
    \hline
    \textbf{Source} &  & \\
    \hline
    \textbf{Enabled By} &  & \\
    \hline
    \textbf{Supports} &  & \\
    \hline
    \textbf{Derived From} &  & \\
    \hline
    \textbf{Reviewed Status} &  & \\
    \hline
  \end{tabularx}
\end{table}

\begin{table}[htbp]
  \centering
  \footnotesize
  \begin{tabularx}{1\textwidth}{|l|l|X|}
    \hline
    \multicolumn{3}{c}{\textbf{Contained Features}} \\
    \hline
    \textbf{Feature Name} & \textbf{Feature Id} & \textbf{Comment} \\
    \hline
    Standards-Based Boundary Exchange & \texttt{cap:bec-0dfe1c63} (at \ref{cap:bec-0dfe1c63}) & \\
    \hline
    Context Establishment & \texttt{cap:bec-8edc75dd} (at \ref{cap:bec-8edc75dd}) & \\
    \hline
    State Monitoring & \texttt{cap:bec-c8a17f6a} (at \ref{cap:bec-c8a17f6a}) & \\
    \hline
    Exception Management & \texttt{cap:bec-6f511b71} (at \ref{cap:bec-6f511b71}) & \\
    \hline
    Information Retrieval & \texttt{cap:bec-2ebb9f24} (at \ref{cap:bec-2ebb9f24}) & \\
    \hline
  \end{tabularx}
\end{table}

Intent. Boundary Protocol Mediation is a candidate enabling capability within Health Information Communication. It frames the reusable system-supported behaviours needed to support this healthcare context without prescribing a product, interface, technology or deployment boundary.

Boundary. This capability supports the business activity and does not inherit the authority or accountability of the Business Capability it enables. Its child Features are review candidates until ancestry and semantics are adjudicated.

\subsubsection{Feature - Standards-Based Boundary Exchange}\label{cap:bec-0dfe1c63}
\vspace{-1em}
\begin{table}[htbp]
  \centering
  \footnotesize
  \begin{tabularx}{1\textwidth}{|l|l|X|}
    \hline
    \textbf{Metadata Name} & \textbf{Metadata Value} & \textbf{Description} \\
    \hline
    \textbf{Reference Id} & \texttt{BEC-0DFE1C63} & \\
    \hline
    \textbf{Status} & Existing Feature; candidate CT3 placement & \\
    \hline
    \textbf{Ancestry} &  & \\
    \hline
    \textbf{Source} & FEAT-ISE-10 & \\
    \hline
    \textbf{Enabled By} &  & \\
    \hline
    \textbf{Supports} &  & \\
    \hline
    \textbf{Derived From} &  & \\
    \hline
    \textbf{Reviewed Status} & Confirm existing Feature meaning and proposed CT3 ownership. & \\
    \hline
  \end{tabularx}
\end{table}

Intent. Standards-Based Boundary Exchange provides a distinct enabling behaviour in the context of Boundary Protocol Mediation. This is an existing named Feature; its placement under the candidate CT3 remains subject to review.

\noindent\textbf{Feature boundary.} This behaviour is scoped to Boundary Protocol Mediation; an identically named Feature elsewhere may carry different business meaning.\par
\featuretable{
Function & \texttt{FN-BADDB736} & Mediate Standards-Based Health Information Communication & Catalogue \\
Service & \texttt{SV-11D6A2FC} & Standards Communication Gateway Service & Linked to preceding Function in source \\
}
\subsubsection{Feature - Context Establishment}\label{cap:bec-8edc75dd}
\vspace{-1em}
\begin{table}[htbp]
  \centering
  \footnotesize
  \begin{tabularx}{1\textwidth}{|l|l|X|}
    \hline
    \textbf{Metadata Name} & \textbf{Metadata Value} & \textbf{Description} \\
    \hline
    \textbf{Reference Id} & \texttt{BEC-8EDC75DD} & \\
    \hline
    \textbf{Status} & Generated candidate & \\
    \hline
    \textbf{Ancestry} &  & \\
    \hline
    \textbf{Source} & Semantic-pattern candidate & \\
    \hline
    \textbf{Enabled By} &  & \\
    \hline
    \textbf{Supports} &  & \\
    \hline
    \textbf{Derived From} &  & \\
    \hline
    \textbf{Reviewed Status} & Generated candidate: keep / modify / relocate / reject; verify whether independently testable. & \\
    \hline
  \end{tabularx}
\end{table}

Intent. Context Establishment is proposed as a discrete, testable enabling behaviour for Boundary Protocol Mediation. It is included for architectural review, not as an established implementation obligation.

\noindent\textbf{Feature boundary.} This behaviour is scoped to Boundary Protocol Mediation; an identically named Feature elsewhere may carry different business meaning.\par
\featuretable{
--- & --- & \emph{Not yet captured} & \emph{Complete during review} \\
}
\subsubsection{Feature - State Monitoring}\label{cap:bec-c8a17f6a}
\vspace{-1em}
\begin{table}[htbp]
  \centering
  \footnotesize
  \begin{tabularx}{1\textwidth}{|l|l|X|}
    \hline
    \textbf{Metadata Name} & \textbf{Metadata Value} & \textbf{Description} \\
    \hline
    \textbf{Reference Id} & \texttt{BEC-C8A17F6A} & \\
    \hline
    \textbf{Status} & Generated candidate & \\
    \hline
    \textbf{Ancestry} &  & \\
    \hline
    \textbf{Source} & Semantic-pattern candidate & \\
    \hline
    \textbf{Enabled By} &  & \\
    \hline
    \textbf{Supports} &  & \\
    \hline
    \textbf{Derived From} &  & \\
    \hline
    \textbf{Reviewed Status} & Generated candidate: keep / modify / relocate / reject; verify whether independently testable. & \\
    \hline
  \end{tabularx}
\end{table}

Intent. State Monitoring is proposed as a discrete, testable enabling behaviour for Boundary Protocol Mediation. It is included for architectural review, not as an established implementation obligation.

\noindent\textbf{Feature boundary.} This behaviour is scoped to Boundary Protocol Mediation; an identically named Feature elsewhere may carry different business meaning.\par
\featuretable{
--- & --- & \emph{Not yet captured} & \emph{Complete during review} \\
}
\subsubsection{Feature - Exception Management}\label{cap:bec-6f511b71}
\vspace{-1em}
\begin{table}[htbp]
  \centering
  \footnotesize
  \begin{tabularx}{1\textwidth}{|l|l|X|}
    \hline
    \textbf{Metadata Name} & \textbf{Metadata Value} & \textbf{Description} \\
    \hline
    \textbf{Reference Id} & \texttt{BEC-6F511B71} & \\
    \hline
    \textbf{Status} & Generated candidate & \\
    \hline
    \textbf{Ancestry} &  & \\
    \hline
    \textbf{Source} & Semantic-pattern candidate & \\
    \hline
    \textbf{Enabled By} &  & \\
    \hline
    \textbf{Supports} &  & \\
    \hline
    \textbf{Derived From} &  & \\
    \hline
    \textbf{Reviewed Status} & Generated candidate: keep / modify / relocate / reject; verify whether independently testable. & \\
    \hline
  \end{tabularx}
\end{table}

Intent. Exception Management is proposed as a discrete, testable enabling behaviour for Boundary Protocol Mediation. It is included for architectural review, not as an established implementation obligation.

\noindent\textbf{Feature boundary.} This behaviour is scoped to Boundary Protocol Mediation; an identically named Feature elsewhere may carry different business meaning.\par
\featuretable{
--- & --- & \emph{Not yet captured} & \emph{Complete during review} \\
}
\subsubsection{Feature - Information Retrieval}\label{cap:bec-2ebb9f24}
\vspace{-1em}
\begin{table}[htbp]
  \centering
  \footnotesize
  \begin{tabularx}{1\textwidth}{|l|l|X|}
    \hline
    \textbf{Metadata Name} & \textbf{Metadata Value} & \textbf{Description} \\
    \hline
    \textbf{Reference Id} & \texttt{BEC-2EBB9F24} & \\
    \hline
    \textbf{Status} & Generated candidate & \\
    \hline
    \textbf{Ancestry} &  & \\
    \hline
    \textbf{Source} & Semantic-pattern candidate & \\
    \hline
    \textbf{Enabled By} &  & \\
    \hline
    \textbf{Supports} &  & \\
    \hline
    \textbf{Derived From} &  & \\
    \hline
    \textbf{Reviewed Status} & Generated candidate: keep / modify / relocate / reject; verify whether independently testable. & \\
    \hline
  \end{tabularx}
\end{table}

Intent. Information Retrieval is proposed as a discrete, testable enabling behaviour for Boundary Protocol Mediation. It is included for architectural review, not as an established implementation obligation.

\noindent\textbf{Feature boundary.} This behaviour is scoped to Boundary Protocol Mediation; an identically named Feature elsewhere may carry different business meaning.\par
\featuretable{
--- & --- & \emph{Not yet captured} & \emph{Complete during review} \\
}

\section{CT2 - Health Information Control}\label{cap:bec-203e2d62}
\vspace{-0.5em}
\begin{table}[htbp]
  \centering
  \footnotesize
  \begin{tabularx}{1\textwidth}{|l|l|X|}
    \hline
    \multicolumn{3}{c}{\textbf{Capability Metadata}} \\
    \hline
    \textbf{Name} & \textbf{Value} & \textbf{Comment and/or Description} \\
    \hline
    \textbf{Reference Id} & \texttt{BEC-203E2D62} & \\
    \hline
    \textbf{Status} & Candidate ancestry & \\
    \hline
    \textbf{Ancestry} &  & \\
    \hline
    \textbf{Source} &  & \\
    \hline
    \textbf{Enabled By} &  & \\
    \hline
    \textbf{Supports} &  & \\
    \hline
    \textbf{Derived From} &  & \\
    \hline
    \textbf{Reviewed Status} &  & \\
    \hline
  \end{tabularx}
\end{table}

\begin{table}[htbp]
  \centering
  \footnotesize
  \begin{tabularx}{1\textwidth}{|l|l|X|}
    \hline
    \multicolumn{3}{c}{\textbf{Contained Capabilities}} \\
    \hline
    \textbf{Capability Name} & \textbf{Capability Id} & \textbf{Comment} \\
    \hline
    Policy \& Security Governance & \texttt{cap:bec-18459ccd} (at \ref{cap:bec-18459ccd}) & \\
    \hline
  \end{tabularx}
\end{table}

Intent. Health Information Control is a candidate enabling capability within Shared Information and Coordination Enablement. It frames the reusable system-supported behaviours needed to support this healthcare context without prescribing a product, interface, technology or deployment boundary.

Boundary. This capability supports the business activity and does not inherit the authority or accountability of the Business Capability it enables. Its child Features are review candidates until ancestry and semantics are adjudicated.

\subsection{CT3 - Policy \& Security Governance}\label{cap:bec-18459ccd}
\vspace{-0.5em}
\begin{table}[htbp]
  \centering
  \footnotesize
  \begin{tabularx}{1\textwidth}{|l|l|X|}
    \hline
    \multicolumn{3}{c}{\textbf{Capability Metadata}} \\
    \hline
    \textbf{Name} & \textbf{Value} & \textbf{Comment and/or Description} \\
    \hline
    \textbf{Reference Id} & \texttt{BEC-18459CCD} & \\
    \hline
    \textbf{Status} & Candidate ancestry & \\
    \hline
    \textbf{Ancestry} &  & \\
    \hline
    \textbf{Source} &  & \\
    \hline
    \textbf{Enabled By} &  & \\
    \hline
    \textbf{Supports} &  & \\
    \hline
    \textbf{Derived From} &  & \\
    \hline
    \textbf{Reviewed Status} &  & \\
    \hline
  \end{tabularx}
\end{table}

\begin{table}[htbp]
  \centering
  \footnotesize
  \begin{tabularx}{1\textwidth}{|l|l|X|}
    \hline
    \multicolumn{3}{c}{\textbf{Contained Features}} \\
    \hline
    \textbf{Feature Name} & \textbf{Feature Id} & \textbf{Comment} \\
    \hline
    Health Information Access \& Consent Control & \texttt{cap:bec-e140c81d} (at \ref{cap:bec-e140c81d}) & \\
    \hline
    Security Context Propagation & \texttt{cap:bec-ca2ef6f6} (at \ref{cap:bec-ca2ef6f6}) & \\
    \hline
    Non-PHI Compliance Auditing & \texttt{cap:bec-0bba26d3} (at \ref{cap:bec-0bba26d3}) & \\
    \hline
    Criteria Management & \texttt{cap:bec-9eb61778} (at \ref{cap:bec-9eb61778}) & \\
    \hline
    Evidence Association & \texttt{cap:bec-e3405f10} (at \ref{cap:bec-e3405f10}) & \\
    \hline
    Policy Evaluation & \texttt{cap:bec-5bab5ea0} (at \ref{cap:bec-5bab5ea0}) & \\
    \hline
    Decision Traceability & \texttt{cap:bec-1d2ca7d5} (at \ref{cap:bec-1d2ca7d5}) & \\
    \hline
    Exception Reporting & \texttt{cap:bec-cb8f5161} (at \ref{cap:bec-cb8f5161}) & \\
    \hline
  \end{tabularx}
\end{table}

Intent. Policy \& Security Governance is a candidate enabling capability within Health Information Control. It frames the reusable system-supported behaviours needed to support this healthcare context without prescribing a product, interface, technology or deployment boundary.

Boundary. This capability supports the business activity and does not inherit the authority or accountability of the Business Capability it enables. Its child Features are review candidates until ancestry and semantics are adjudicated.

\subsubsection{Feature - Health Information Access \& Consent Control}\label{cap:bec-e140c81d}
\vspace{-1em}
\begin{table}[htbp]
  \centering
  \footnotesize
  \begin{tabularx}{1\textwidth}{|l|l|X|}
    \hline
    \textbf{Metadata Name} & \textbf{Metadata Value} & \textbf{Description} \\
    \hline
    \textbf{Reference Id} & \texttt{BEC-E140C81D} & \\
    \hline
    \textbf{Status} & Existing Feature; candidate CT3 placement & \\
    \hline
    \textbf{Ancestry} &  & \\
    \hline
    \textbf{Source} & FEAT-ISE-11 & \\
    \hline
    \textbf{Enabled By} &  & \\
    \hline
    \textbf{Supports} &  & \\
    \hline
    \textbf{Derived From} &  & \\
    \hline
    \textbf{Reviewed Status} & Confirm existing Feature meaning and proposed CT3 ownership. & \\
    \hline
  \end{tabularx}
\end{table}

Intent. Health Information Access \& Consent Control provides a distinct enabling behaviour in the context of Policy \& Security Governance. This is an existing named Feature; its placement under the candidate CT3 remains subject to review.

\noindent\textbf{Feature boundary.} This behaviour is scoped to Policy \& Security Governance; an identically named Feature elsewhere may carry different business meaning.\par
\featuretable{
Function & \texttt{FN-BE5547FE} & Evaluate Information Access Authority & Catalogue \\
Service & \texttt{SV-6A1A73B6} & Policy Evaluation \& Authorisation Service & Linked to preceding Function in source \\
}
\subsubsection{Feature - Security Context Propagation}\label{cap:bec-ca2ef6f6}
\vspace{-1em}
\begin{table}[htbp]
  \centering
  \footnotesize
  \begin{tabularx}{1\textwidth}{|l|l|X|}
    \hline
    \textbf{Metadata Name} & \textbf{Metadata Value} & \textbf{Description} \\
    \hline
    \textbf{Reference Id} & \texttt{BEC-CA2EF6F6} & \\
    \hline
    \textbf{Status} & Existing Feature; candidate CT3 placement & \\
    \hline
    \textbf{Ancestry} &  & \\
    \hline
    \textbf{Source} & FEAT-ISE-12 & \\
    \hline
    \textbf{Enabled By} &  & \\
    \hline
    \textbf{Supports} &  & \\
    \hline
    \textbf{Derived From} &  & \\
    \hline
    \textbf{Reviewed Status} & Confirm existing Feature meaning and proposed CT3 ownership. & \\
    \hline
  \end{tabularx}
\end{table}

Intent. Security Context Propagation provides a distinct enabling behaviour in the context of Policy \& Security Governance. This is an existing named Feature; its placement under the candidate CT3 remains subject to review.

\noindent\textbf{Feature boundary.} This behaviour is scoped to Policy \& Security Governance; an identically named Feature elsewhere may carry different business meaning.\par
\featuretable{
Function & \texttt{FN-19C02743} & Evaluate Consent Constraint & Catalogue \\
Service & \texttt{SV-96F48EC6} & Consent Enforcement Service & Linked to preceding Function in source \\
}
\subsubsection{Feature - Non-PHI Compliance Auditing}\label{cap:bec-0bba26d3}
\vspace{-1em}
\begin{table}[htbp]
  \centering
  \footnotesize
  \begin{tabularx}{1\textwidth}{|l|l|X|}
    \hline
    \textbf{Metadata Name} & \textbf{Metadata Value} & \textbf{Description} \\
    \hline
    \textbf{Reference Id} & \texttt{BEC-0BBA26D3} & \\
    \hline
    \textbf{Status} & Existing Feature; candidate CT3 placement & \\
    \hline
    \textbf{Ancestry} &  & \\
    \hline
    \textbf{Source} & FEAT-ISE-13 & \\
    \hline
    \textbf{Enabled By} &  & \\
    \hline
    \textbf{Supports} &  & \\
    \hline
    \textbf{Derived From} &  & \\
    \hline
    \textbf{Reviewed Status} & Confirm existing Feature meaning and proposed CT3 ownership. & \\
    \hline
  \end{tabularx}
\end{table}

Intent. Non-PHI Compliance Auditing provides a distinct enabling behaviour in the context of Policy \& Security Governance. This is an existing named Feature; its placement under the candidate CT3 remains subject to review.

\noindent\textbf{Feature boundary.} This behaviour is scoped to Policy \& Security Governance; an identically named Feature elsewhere may carry different business meaning.\par
\featuretable{
Function & \texttt{FN-922D9EF6} & Record Security-Significant Activity & Catalogue \\
Service & \texttt{SV-CB281E3A} & Security Audit Ingress & Linked to preceding Function in source \\
}
\subsubsection{Feature - Criteria Management}\label{cap:bec-9eb61778}
\vspace{-1em}
\begin{table}[htbp]
  \centering
  \footnotesize
  \begin{tabularx}{1\textwidth}{|l|l|X|}
    \hline
    \textbf{Metadata Name} & \textbf{Metadata Value} & \textbf{Description} \\
    \hline
    \textbf{Reference Id} & \texttt{BEC-9EB61778} & \\
    \hline
    \textbf{Status} & Generated candidate & \\
    \hline
    \textbf{Ancestry} &  & \\
    \hline
    \textbf{Source} & Semantic-pattern candidate & \\
    \hline
    \textbf{Enabled By} &  & \\
    \hline
    \textbf{Supports} &  & \\
    \hline
    \textbf{Derived From} &  & \\
    \hline
    \textbf{Reviewed Status} & Generated candidate: keep / modify / relocate / reject; verify whether independently testable. & \\
    \hline
  \end{tabularx}
\end{table}

Intent. Criteria Management is proposed as a discrete, testable enabling behaviour for Policy \& Security Governance. It is included for architectural review, not as an established implementation obligation.

\noindent\textbf{Feature boundary.} This behaviour is scoped to Policy \& Security Governance; an identically named Feature elsewhere may carry different business meaning.\par
\featuretable{
--- & --- & \emph{Not yet captured} & \emph{Complete during review} \\
}
\subsubsection{Feature - Evidence Association}\label{cap:bec-e3405f10}
\vspace{-1em}
\begin{table}[htbp]
  \centering
  \footnotesize
  \begin{tabularx}{1\textwidth}{|l|l|X|}
    \hline
    \textbf{Metadata Name} & \textbf{Metadata Value} & \textbf{Description} \\
    \hline
    \textbf{Reference Id} & \texttt{BEC-E3405F10} & \\
    \hline
    \textbf{Status} & Generated candidate & \\
    \hline
    \textbf{Ancestry} &  & \\
    \hline
    \textbf{Source} & Semantic-pattern candidate & \\
    \hline
    \textbf{Enabled By} &  & \\
    \hline
    \textbf{Supports} &  & \\
    \hline
    \textbf{Derived From} &  & \\
    \hline
    \textbf{Reviewed Status} & Generated candidate: keep / modify / relocate / reject; verify whether independently testable. & \\
    \hline
  \end{tabularx}
\end{table}

Intent. Evidence Association is proposed as a discrete, testable enabling behaviour for Policy \& Security Governance. It is included for architectural review, not as an established implementation obligation.

\noindent\textbf{Feature boundary.} This behaviour is scoped to Policy \& Security Governance; an identically named Feature elsewhere may carry different business meaning.\par
\featuretable{
--- & --- & \emph{Not yet captured} & \emph{Complete during review} \\
}
\subsubsection{Feature - Policy Evaluation}\label{cap:bec-5bab5ea0}
\vspace{-1em}
\begin{table}[htbp]
  \centering
  \footnotesize
  \begin{tabularx}{1\textwidth}{|l|l|X|}
    \hline
    \textbf{Metadata Name} & \textbf{Metadata Value} & \textbf{Description} \\
    \hline
    \textbf{Reference Id} & \texttt{BEC-5BAB5EA0} & \\
    \hline
    \textbf{Status} & Generated candidate & \\
    \hline
    \textbf{Ancestry} &  & \\
    \hline
    \textbf{Source} & Semantic-pattern candidate & \\
    \hline
    \textbf{Enabled By} &  & \\
    \hline
    \textbf{Supports} &  & \\
    \hline
    \textbf{Derived From} &  & \\
    \hline
    \textbf{Reviewed Status} & Generated candidate: keep / modify / relocate / reject; verify whether independently testable. & \\
    \hline
  \end{tabularx}
\end{table}

Intent. Policy Evaluation is proposed as a discrete, testable enabling behaviour for Policy \& Security Governance. It is included for architectural review, not as an established implementation obligation.

\noindent\textbf{Feature boundary.} This behaviour is scoped to Policy \& Security Governance; an identically named Feature elsewhere may carry different business meaning.\par
\featuretable{
--- & --- & \emph{Not yet captured} & \emph{Complete during review} \\
}
\subsubsection{Feature - Decision Traceability}\label{cap:bec-1d2ca7d5}
\vspace{-1em}
\begin{table}[htbp]
  \centering
  \footnotesize
  \begin{tabularx}{1\textwidth}{|l|l|X|}
    \hline
    \textbf{Metadata Name} & \textbf{Metadata Value} & \textbf{Description} \\
    \hline
    \textbf{Reference Id} & \texttt{BEC-1D2CA7D5} & \\
    \hline
    \textbf{Status} & Generated candidate & \\
    \hline
    \textbf{Ancestry} &  & \\
    \hline
    \textbf{Source} & Semantic-pattern candidate & \\
    \hline
    \textbf{Enabled By} &  & \\
    \hline
    \textbf{Supports} &  & \\
    \hline
    \textbf{Derived From} &  & \\
    \hline
    \textbf{Reviewed Status} & Generated candidate: keep / modify / relocate / reject; verify whether independently testable. & \\
    \hline
  \end{tabularx}
\end{table}

Intent. Decision Traceability is proposed as a discrete, testable enabling behaviour for Policy \& Security Governance. It is included for architectural review, not as an established implementation obligation.

\noindent\textbf{Feature boundary.} This behaviour is scoped to Policy \& Security Governance; an identically named Feature elsewhere may carry different business meaning.\par
\featuretable{
--- & --- & \emph{Not yet captured} & \emph{Complete during review} \\
}
\subsubsection{Feature - Exception Reporting}\label{cap:bec-cb8f5161}
\vspace{-1em}
\begin{table}[htbp]
  \centering
  \footnotesize
  \begin{tabularx}{1\textwidth}{|l|l|X|}
    \hline
    \textbf{Metadata Name} & \textbf{Metadata Value} & \textbf{Description} \\
    \hline
    \textbf{Reference Id} & \texttt{BEC-CB8F5161} & \\
    \hline
    \textbf{Status} & Generated candidate & \\
    \hline
    \textbf{Ancestry} &  & \\
    \hline
    \textbf{Source} & Semantic-pattern candidate & \\
    \hline
    \textbf{Enabled By} &  & \\
    \hline
    \textbf{Supports} &  & \\
    \hline
    \textbf{Derived From} &  & \\
    \hline
    \textbf{Reviewed Status} & Generated candidate: keep / modify / relocate / reject; verify whether independently testable. & \\
    \hline
  \end{tabularx}
\end{table}

Intent. Exception Reporting is proposed as a discrete, testable enabling behaviour for Policy \& Security Governance. It is included for architectural review, not as an established implementation obligation.

\noindent\textbf{Feature boundary.} This behaviour is scoped to Policy \& Security Governance; an identically named Feature elsewhere may carry different business meaning.\par
\featuretable{
--- & --- & \emph{Not yet captured} & \emph{Complete during review} \\
}

\section{CT2 - Clinical Knowledge Services}\label{cap:bec-b622bacf}
\vspace{-0.5em}
\begin{table}[htbp]
  \centering
  \footnotesize
  \begin{tabularx}{1\textwidth}{|l|l|X|}
    \hline
    \multicolumn{3}{c}{\textbf{Capability Metadata}} \\
    \hline
    \textbf{Name} & \textbf{Value} & \textbf{Comment and/or Description} \\
    \hline
    \textbf{Reference Id} & \texttt{BEC-B622BACF} & \\
    \hline
    \textbf{Status} & Candidate ancestry & \\
    \hline
    \textbf{Ancestry} &  & \\
    \hline
    \textbf{Source} &  & \\
    \hline
    \textbf{Enabled By} &  & \\
    \hline
    \textbf{Supports} &  & \\
    \hline
    \textbf{Derived From} &  & \\
    \hline
    \textbf{Reviewed Status} &  & \\
    \hline
  \end{tabularx}
\end{table}

\begin{table}[htbp]
  \centering
  \footnotesize
  \begin{tabularx}{1\textwidth}{|l|l|X|}
    \hline
    \multicolumn{3}{c}{\textbf{Contained Capabilities}} \\
    \hline
    \textbf{Capability Name} & \textbf{Capability Id} & \textbf{Comment} \\
    \hline
    Terminology \& Reference Mapping & \texttt{cap:bec-43b6ce64} (at \ref{cap:bec-43b6ce64}) & \\
    \hline
  \end{tabularx}
\end{table}

Intent. Clinical Knowledge Services is a candidate enabling capability within Shared Information and Coordination Enablement. It frames the reusable system-supported behaviours needed to support this healthcare context without prescribing a product, interface, technology or deployment boundary.

Boundary. This capability supports the business activity and does not inherit the authority or accountability of the Business Capability it enables. Its child Features are review candidates until ancestry and semantics are adjudicated.

\subsection{CT3 - Terminology \& Reference Mapping}\label{cap:bec-43b6ce64}
\vspace{-0.5em}
\begin{table}[htbp]
  \centering
  \footnotesize
  \begin{tabularx}{1\textwidth}{|l|l|X|}
    \hline
    \multicolumn{3}{c}{\textbf{Capability Metadata}} \\
    \hline
    \textbf{Name} & \textbf{Value} & \textbf{Comment and/or Description} \\
    \hline
    \textbf{Reference Id} & \texttt{BEC-43B6CE64} & \\
    \hline
    \textbf{Status} & Candidate ancestry & \\
    \hline
    \textbf{Ancestry} &  & \\
    \hline
    \textbf{Source} &  & \\
    \hline
    \textbf{Enabled By} &  & \\
    \hline
    \textbf{Supports} &  & \\
    \hline
    \textbf{Derived From} &  & \\
    \hline
    \textbf{Reviewed Status} &  & \\
    \hline
  \end{tabularx}
\end{table}

\begin{table}[htbp]
  \centering
  \footnotesize
  \begin{tabularx}{1\textwidth}{|l|l|X|}
    \hline
    \multicolumn{3}{c}{\textbf{Contained Features}} \\
    \hline
    \textbf{Feature Name} & \textbf{Feature Id} & \textbf{Comment} \\
    \hline
    Clinical Concept Resolution & \texttt{cap:bec-b0351502} (at \ref{cap:bec-b0351502}) & \\
    \hline
    Bidirectional Terminology Mapping & \texttt{cap:bec-dc869393} (at \ref{cap:bec-dc869393}) & \\
    \hline
    Information Validation & \texttt{cap:bec-77a76c7c} (at \ref{cap:bec-77a76c7c}) & \\
    \hline
    Information Correlation & \texttt{cap:bec-cbec5f4e} (at \ref{cap:bec-cbec5f4e}) & \\
    \hline
    Information Retrieval & \texttt{cap:bec-e5775112} (at \ref{cap:bec-e5775112}) & \\
    \hline
    Information Provenance & \texttt{cap:bec-a5d9596d} (at \ref{cap:bec-a5d9596d}) & \\
    \hline
    Information Lifecycle Management & \texttt{cap:bec-760ae591} (at \ref{cap:bec-760ae591}) & \\
    \hline
  \end{tabularx}
\end{table}

Intent. Terminology \& Reference Mapping is a candidate enabling capability within Clinical Knowledge Services. It frames the reusable system-supported behaviours needed to support this healthcare context without prescribing a product, interface, technology or deployment boundary.

Boundary. This capability supports the business activity and does not inherit the authority or accountability of the Business Capability it enables. Its child Features are review candidates until ancestry and semantics are adjudicated.

\subsubsection{Feature - Clinical Concept Resolution}\label{cap:bec-b0351502}
\vspace{-1em}
\begin{table}[htbp]
  \centering
  \footnotesize
  \begin{tabularx}{1\textwidth}{|l|l|X|}
    \hline
    \textbf{Metadata Name} & \textbf{Metadata Value} & \textbf{Description} \\
    \hline
    \textbf{Reference Id} & \texttt{BEC-B0351502} & \\
    \hline
    \textbf{Status} & Existing Feature; candidate CT3 placement & \\
    \hline
    \textbf{Ancestry} &  & \\
    \hline
    \textbf{Source} & FEAT-ISE-14 & \\
    \hline
    \textbf{Enabled By} &  & \\
    \hline
    \textbf{Supports} &  & \\
    \hline
    \textbf{Derived From} &  & \\
    \hline
    \textbf{Reviewed Status} & Confirm existing Feature meaning and proposed CT3 ownership. & \\
    \hline
  \end{tabularx}
\end{table}

Intent. Clinical Concept Resolution provides a distinct enabling behaviour in the context of Terminology \& Reference Mapping. This is an existing named Feature; its placement under the candidate CT3 remains subject to review.

\noindent\textbf{Feature boundary.} This behaviour is scoped to Terminology \& Reference Mapping; an identically named Feature elsewhere may carry different business meaning.\par
\featuretable{
Function & \texttt{FN-034E310D} & Resolve Clinical Concept & Catalogue \\
Service & \texttt{SV-4537A519} & Terminology Concept Resolution Service & Linked to preceding Function in source \\
}
\subsubsection{Feature - Bidirectional Terminology Mapping}\label{cap:bec-dc869393}
\vspace{-1em}
\begin{table}[htbp]
  \centering
  \footnotesize
  \begin{tabularx}{1\textwidth}{|l|l|X|}
    \hline
    \textbf{Metadata Name} & \textbf{Metadata Value} & \textbf{Description} \\
    \hline
    \textbf{Reference Id} & \texttt{BEC-DC869393} & \\
    \hline
    \textbf{Status} & Existing Feature; candidate CT3 placement & \\
    \hline
    \textbf{Ancestry} &  & \\
    \hline
    \textbf{Source} & FEAT-ISE-15 & \\
    \hline
    \textbf{Enabled By} &  & \\
    \hline
    \textbf{Supports} &  & \\
    \hline
    \textbf{Derived From} &  & \\
    \hline
    \textbf{Reviewed Status} & Confirm existing Feature meaning and proposed CT3 ownership. & \\
    \hline
  \end{tabularx}
\end{table}

Intent. Bidirectional Terminology Mapping provides a distinct enabling behaviour in the context of Terminology \& Reference Mapping. This is an existing named Feature; its placement under the candidate CT3 remains subject to review.

\noindent\textbf{Feature boundary.} This behaviour is scoped to Terminology \& Reference Mapping; an identically named Feature elsewhere may carry different business meaning.\par
\featuretable{
Function & \texttt{FN-1E9D52B8} & Map Clinical Concept & Catalogue \\
Service & \texttt{SV-46D4C06A} & Terminology Mapping Service & Linked to preceding Function in source \\
}
\subsubsection{Feature - Information Validation}\label{cap:bec-77a76c7c}
\vspace{-1em}
\begin{table}[htbp]
  \centering
  \footnotesize
  \begin{tabularx}{1\textwidth}{|l|l|X|}
    \hline
    \textbf{Metadata Name} & \textbf{Metadata Value} & \textbf{Description} \\
    \hline
    \textbf{Reference Id} & \texttt{BEC-77A76C7C} & \\
    \hline
    \textbf{Status} & Generated candidate & \\
    \hline
    \textbf{Ancestry} &  & \\
    \hline
    \textbf{Source} & Semantic-pattern candidate & \\
    \hline
    \textbf{Enabled By} &  & \\
    \hline
    \textbf{Supports} &  & \\
    \hline
    \textbf{Derived From} &  & \\
    \hline
    \textbf{Reviewed Status} & Generated candidate: keep / modify / relocate / reject; verify whether independently testable. & \\
    \hline
  \end{tabularx}
\end{table}

Intent. Information Validation is proposed as a discrete, testable enabling behaviour for Terminology \& Reference Mapping. It is included for architectural review, not as an established implementation obligation.

\noindent\textbf{Feature boundary.} This behaviour is scoped to Terminology \& Reference Mapping; an identically named Feature elsewhere may carry different business meaning.\par
\featuretable{
--- & --- & \emph{Not yet captured} & \emph{Complete during review} \\
}
\subsubsection{Feature - Information Correlation}\label{cap:bec-cbec5f4e}
\vspace{-1em}
\begin{table}[htbp]
  \centering
  \footnotesize
  \begin{tabularx}{1\textwidth}{|l|l|X|}
    \hline
    \textbf{Metadata Name} & \textbf{Metadata Value} & \textbf{Description} \\
    \hline
    \textbf{Reference Id} & \texttt{BEC-CBEC5F4E} & \\
    \hline
    \textbf{Status} & Generated candidate & \\
    \hline
    \textbf{Ancestry} &  & \\
    \hline
    \textbf{Source} & Semantic-pattern candidate & \\
    \hline
    \textbf{Enabled By} &  & \\
    \hline
    \textbf{Supports} &  & \\
    \hline
    \textbf{Derived From} &  & \\
    \hline
    \textbf{Reviewed Status} & Generated candidate: keep / modify / relocate / reject; verify whether independently testable. & \\
    \hline
  \end{tabularx}
\end{table}

Intent. Information Correlation is proposed as a discrete, testable enabling behaviour for Terminology \& Reference Mapping. It is included for architectural review, not as an established implementation obligation.

\noindent\textbf{Feature boundary.} This behaviour is scoped to Terminology \& Reference Mapping; an identically named Feature elsewhere may carry different business meaning.\par
\featuretable{
--- & --- & \emph{Not yet captured} & \emph{Complete during review} \\
}
\subsubsection{Feature - Information Retrieval}\label{cap:bec-e5775112}
\vspace{-1em}
\begin{table}[htbp]
  \centering
  \footnotesize
  \begin{tabularx}{1\textwidth}{|l|l|X|}
    \hline
    \textbf{Metadata Name} & \textbf{Metadata Value} & \textbf{Description} \\
    \hline
    \textbf{Reference Id} & \texttt{BEC-E5775112} & \\
    \hline
    \textbf{Status} & Generated candidate & \\
    \hline
    \textbf{Ancestry} &  & \\
    \hline
    \textbf{Source} & Semantic-pattern candidate & \\
    \hline
    \textbf{Enabled By} &  & \\
    \hline
    \textbf{Supports} &  & \\
    \hline
    \textbf{Derived From} &  & \\
    \hline
    \textbf{Reviewed Status} & Generated candidate: keep / modify / relocate / reject; verify whether independently testable. & \\
    \hline
  \end{tabularx}
\end{table}

Intent. Information Retrieval is proposed as a discrete, testable enabling behaviour for Terminology \& Reference Mapping. It is included for architectural review, not as an established implementation obligation.

\noindent\textbf{Feature boundary.} This behaviour is scoped to Terminology \& Reference Mapping; an identically named Feature elsewhere may carry different business meaning.\par
\featuretable{
--- & --- & \emph{Not yet captured} & \emph{Complete during review} \\
}
\subsubsection{Feature - Information Provenance}\label{cap:bec-a5d9596d}
\vspace{-1em}
\begin{table}[htbp]
  \centering
  \footnotesize
  \begin{tabularx}{1\textwidth}{|l|l|X|}
    \hline
    \textbf{Metadata Name} & \textbf{Metadata Value} & \textbf{Description} \\
    \hline
    \textbf{Reference Id} & \texttt{BEC-A5D9596D} & \\
    \hline
    \textbf{Status} & Generated candidate & \\
    \hline
    \textbf{Ancestry} &  & \\
    \hline
    \textbf{Source} & Semantic-pattern candidate & \\
    \hline
    \textbf{Enabled By} &  & \\
    \hline
    \textbf{Supports} &  & \\
    \hline
    \textbf{Derived From} &  & \\
    \hline
    \textbf{Reviewed Status} & Generated candidate: keep / modify / relocate / reject; verify whether independently testable. & \\
    \hline
  \end{tabularx}
\end{table}

Intent. Information Provenance is proposed as a discrete, testable enabling behaviour for Terminology \& Reference Mapping. It is included for architectural review, not as an established implementation obligation.

\noindent\textbf{Feature boundary.} This behaviour is scoped to Terminology \& Reference Mapping; an identically named Feature elsewhere may carry different business meaning.\par
\featuretable{
--- & --- & \emph{Not yet captured} & \emph{Complete during review} \\
}
\subsubsection{Feature - Information Lifecycle Management}\label{cap:bec-760ae591}
\vspace{-1em}
\begin{table}[htbp]
  \centering
  \footnotesize
  \begin{tabularx}{1\textwidth}{|l|l|X|}
    \hline
    \textbf{Metadata Name} & \textbf{Metadata Value} & \textbf{Description} \\
    \hline
    \textbf{Reference Id} & \texttt{BEC-760AE591} & \\
    \hline
    \textbf{Status} & Generated candidate & \\
    \hline
    \textbf{Ancestry} &  & \\
    \hline
    \textbf{Source} & Semantic-pattern candidate & \\
    \hline
    \textbf{Enabled By} &  & \\
    \hline
    \textbf{Supports} &  & \\
    \hline
    \textbf{Derived From} &  & \\
    \hline
    \textbf{Reviewed Status} & Generated candidate: keep / modify / relocate / reject; verify whether independently testable. & \\
    \hline
  \end{tabularx}
\end{table}

Intent. Information Lifecycle Management is proposed as a discrete, testable enabling behaviour for Terminology \& Reference Mapping. It is included for architectural review, not as an established implementation obligation.

\noindent\textbf{Feature boundary.} This behaviour is scoped to Terminology \& Reference Mapping; an identically named Feature elsewhere may carry different business meaning.\par
\featuretable{
--- & --- & \emph{Not yet captured} & \emph{Complete during review} \\
}

\section{CT2 - Clinical Collaboration}\label{cap:bec-4cde0bab}
\vspace{-0.5em}
\begin{table}[htbp]
  \centering
  \footnotesize
  \begin{tabularx}{1\textwidth}{|l|l|X|}
    \hline
    \multicolumn{3}{c}{\textbf{Capability Metadata}} \\
    \hline
    \textbf{Name} & \textbf{Value} & \textbf{Comment and/or Description} \\
    \hline
    \textbf{Reference Id} & \texttt{BEC-4CDE0BAB} & \\
    \hline
    \textbf{Status} & Candidate ancestry & \\
    \hline
    \textbf{Ancestry} &  & \\
    \hline
    \textbf{Source} &  & \\
    \hline
    \textbf{Enabled By} &  & \\
    \hline
    \textbf{Supports} &  & \\
    \hline
    \textbf{Derived From} &  & \\
    \hline
    \textbf{Reviewed Status} &  & \\
    \hline
  \end{tabularx}
\end{table}

\begin{table}[htbp]
  \centering
  \footnotesize
  \begin{tabularx}{1\textwidth}{|l|l|X|}
    \hline
    \multicolumn{3}{c}{\textbf{Contained Capabilities}} \\
    \hline
    \textbf{Capability Name} & \textbf{Capability Id} & \textbf{Comment} \\
    \hline
    LHR-Aware Multidisciplinary Collaboration & \texttt{cap:bec-0c37baff} (at \ref{cap:bec-0c37baff}) & \\
    \hline
  \end{tabularx}
\end{table}

Intent. Clinical Collaboration is a candidate enabling capability within Shared Information and Coordination Enablement. It frames the reusable system-supported behaviours needed to support this healthcare context without prescribing a product, interface, technology or deployment boundary.

Boundary. This capability supports the business activity and does not inherit the authority or accountability of the Business Capability it enables. Its child Features are review candidates until ancestry and semantics are adjudicated.

\subsection{CT3 - LHR-Aware Multidisciplinary Collaboration}\label{cap:bec-0c37baff}
\vspace{-0.5em}
\begin{table}[htbp]
  \centering
  \footnotesize
  \begin{tabularx}{1\textwidth}{|l|l|X|}
    \hline
    \multicolumn{3}{c}{\textbf{Capability Metadata}} \\
    \hline
    \textbf{Name} & \textbf{Value} & \textbf{Comment and/or Description} \\
    \hline
    \textbf{Reference Id} & \texttt{BEC-0C37BAFF} & \\
    \hline
    \textbf{Status} & Candidate ancestry & \\
    \hline
    \textbf{Ancestry} &  & \\
    \hline
    \textbf{Source} &  & \\
    \hline
    \textbf{Enabled By} &  & \\
    \hline
    \textbf{Supports} &  & \\
    \hline
    \textbf{Derived From} &  & \\
    \hline
    \textbf{Reviewed Status} &  & \\
    \hline
  \end{tabularx}
\end{table}

\begin{table}[htbp]
  \centering
  \footnotesize
  \begin{tabularx}{1\textwidth}{|l|l|X|}
    \hline
    \multicolumn{3}{c}{\textbf{Contained Features}} \\
    \hline
    \textbf{Feature Name} & \textbf{Feature Id} & \textbf{Comment} \\
    \hline
    Care-Team Clinical Collaboration Coordination & \texttt{cap:bec-84d2c1df} (at \ref{cap:bec-84d2c1df}) & \\
    \hline
    In-Conversation LHR Query Resolution & \texttt{cap:bec-0b2cf342} (at \ref{cap:bec-0b2cf342}) & \\
    \hline
    Zero-PHI Collaboration Metadata Governance & \texttt{cap:bec-24657506} (at \ref{cap:bec-24657506}) & \\
    \hline
    Context Establishment & \texttt{cap:bec-3f3e28af} (at \ref{cap:bec-3f3e28af}) & \\
    \hline
    State Monitoring & \texttt{cap:bec-911087b0} (at \ref{cap:bec-911087b0}) & \\
    \hline
    Exception Management & \texttt{cap:bec-2df10aa6} (at \ref{cap:bec-2df10aa6}) & \\
    \hline
    Information Retrieval & \texttt{cap:bec-0cd2c755} (at \ref{cap:bec-0cd2c755}) & \\
    \hline
  \end{tabularx}
\end{table}

Intent. LHR-Aware Multidisciplinary Collaboration is a candidate enabling capability within Clinical Collaboration. It frames the reusable system-supported behaviours needed to support this healthcare context without prescribing a product, interface, technology or deployment boundary.

Boundary. This capability supports the business activity and does not inherit the authority or accountability of the Business Capability it enables. Its child Features are review candidates until ancestry and semantics are adjudicated.

\subsubsection{Feature - Care-Team Clinical Collaboration Coordination}\label{cap:bec-84d2c1df}
\vspace{-1em}
\begin{table}[htbp]
  \centering
  \footnotesize
  \begin{tabularx}{1\textwidth}{|l|l|X|}
    \hline
    \textbf{Metadata Name} & \textbf{Metadata Value} & \textbf{Description} \\
    \hline
    \textbf{Reference Id} & \texttt{BEC-84D2C1DF} & \\
    \hline
    \textbf{Status} & Existing Feature; candidate CT3 placement & \\
    \hline
    \textbf{Ancestry} &  & \\
    \hline
    \textbf{Source} & FEAT-ISE-16 & \\
    \hline
    \textbf{Enabled By} &  & \\
    \hline
    \textbf{Supports} &  & \\
    \hline
    \textbf{Derived From} &  & \\
    \hline
    \textbf{Reviewed Status} & Confirm existing Feature meaning and proposed CT3 ownership. & \\
    \hline
  \end{tabularx}
\end{table}

Intent. Care-Team Clinical Collaboration Coordination provides a distinct enabling behaviour in the context of LHR-Aware Multidisciplinary Collaboration. This is an existing named Feature; its placement under the candidate CT3 remains subject to review.

\noindent\textbf{Feature boundary.} This behaviour is scoped to LHR-Aware Multidisciplinary Collaboration; an identically named Feature elsewhere may carry different business meaning.\par
\featuretable{
Function & \texttt{FN-BA5F032E} & Coordinate Care-Team Collaboration & Catalogue \\
Service & \texttt{SV-AFFDEE38} & Clinical Collaboration Space Service & Linked to preceding Function in source \\
}
\subsubsection{Feature - In-Conversation LHR Query Resolution}\label{cap:bec-0b2cf342}
\vspace{-1em}
\begin{table}[htbp]
  \centering
  \footnotesize
  \begin{tabularx}{1\textwidth}{|l|l|X|}
    \hline
    \textbf{Metadata Name} & \textbf{Metadata Value} & \textbf{Description} \\
    \hline
    \textbf{Reference Id} & \texttt{BEC-0B2CF342} & \\
    \hline
    \textbf{Status} & Existing Feature; candidate CT3 placement & \\
    \hline
    \textbf{Ancestry} &  & \\
    \hline
    \textbf{Source} & FEAT-ISE-17 & \\
    \hline
    \textbf{Enabled By} &  & \\
    \hline
    \textbf{Supports} &  & \\
    \hline
    \textbf{Derived From} &  & \\
    \hline
    \textbf{Reviewed Status} & Confirm existing Feature meaning and proposed CT3 ownership. & \\
    \hline
  \end{tabularx}
\end{table}

Intent. In-Conversation LHR Query Resolution provides a distinct enabling behaviour in the context of LHR-Aware Multidisciplinary Collaboration. This is an existing named Feature; its placement under the candidate CT3 remains subject to review.

\noindent\textbf{Feature boundary.} This behaviour is scoped to LHR-Aware Multidisciplinary Collaboration; an identically named Feature elsewhere may carry different business meaning.\par
\featuretable{
Function & \texttt{FN-D3362AD2} & Resolve LHR Query within Collaboration & Catalogue \\
Service & \texttt{SV-E38C8A5E} & Collaboration Clinical Summary Resolution & Linked to preceding Function in source \\
}
\subsubsection{Feature - Zero-PHI Collaboration Metadata Governance}\label{cap:bec-24657506}
\vspace{-1em}
\begin{table}[htbp]
  \centering
  \footnotesize
  \begin{tabularx}{1\textwidth}{|l|l|X|}
    \hline
    \textbf{Metadata Name} & \textbf{Metadata Value} & \textbf{Description} \\
    \hline
    \textbf{Reference Id} & \texttt{BEC-24657506} & \\
    \hline
    \textbf{Status} & Existing Feature; candidate CT3 placement & \\
    \hline
    \textbf{Ancestry} &  & \\
    \hline
    \textbf{Source} & FEAT-ISE-18 & \\
    \hline
    \textbf{Enabled By} &  & \\
    \hline
    \textbf{Supports} &  & \\
    \hline
    \textbf{Derived From} &  & \\
    \hline
    \textbf{Reviewed Status} & Confirm existing Feature meaning and proposed CT3 ownership. & \\
    \hline
  \end{tabularx}
\end{table}

Intent. Zero-PHI Collaboration Metadata Governance provides a distinct enabling behaviour in the context of LHR-Aware Multidisciplinary Collaboration. This is an existing named Feature; its placement under the candidate CT3 remains subject to review.

\noindent\textbf{Feature boundary.} This behaviour is scoped to LHR-Aware Multidisciplinary Collaboration; an identically named Feature elsewhere may carry different business meaning.\par
\featuretable{
Function & \texttt{FN-9F72E531} & Govern Collaboration Metadata & Catalogue \\
Service & \texttt{SV-66DF870B} & Collaboration Space Governance Service & Linked to preceding Function in source \\
}
\subsubsection{Feature - Context Establishment}\label{cap:bec-3f3e28af}
\vspace{-1em}
\begin{table}[htbp]
  \centering
  \footnotesize
  \begin{tabularx}{1\textwidth}{|l|l|X|}
    \hline
    \textbf{Metadata Name} & \textbf{Metadata Value} & \textbf{Description} \\
    \hline
    \textbf{Reference Id} & \texttt{BEC-3F3E28AF} & \\
    \hline
    \textbf{Status} & Generated candidate & \\
    \hline
    \textbf{Ancestry} &  & \\
    \hline
    \textbf{Source} & Semantic-pattern candidate & \\
    \hline
    \textbf{Enabled By} &  & \\
    \hline
    \textbf{Supports} &  & \\
    \hline
    \textbf{Derived From} &  & \\
    \hline
    \textbf{Reviewed Status} & Generated candidate: keep / modify / relocate / reject; verify whether independently testable. & \\
    \hline
  \end{tabularx}
\end{table}

Intent. Context Establishment is proposed as a discrete, testable enabling behaviour for LHR-Aware Multidisciplinary Collaboration. It is included for architectural review, not as an established implementation obligation.

\noindent\textbf{Feature boundary.} This behaviour is scoped to LHR-Aware Multidisciplinary Collaboration; an identically named Feature elsewhere may carry different business meaning.\par
\featuretable{
--- & --- & \emph{Not yet captured} & \emph{Complete during review} \\
}
\subsubsection{Feature - State Monitoring}\label{cap:bec-911087b0}
\vspace{-1em}
\begin{table}[htbp]
  \centering
  \footnotesize
  \begin{tabularx}{1\textwidth}{|l|l|X|}
    \hline
    \textbf{Metadata Name} & \textbf{Metadata Value} & \textbf{Description} \\
    \hline
    \textbf{Reference Id} & \texttt{BEC-911087B0} & \\
    \hline
    \textbf{Status} & Generated candidate & \\
    \hline
    \textbf{Ancestry} &  & \\
    \hline
    \textbf{Source} & Semantic-pattern candidate & \\
    \hline
    \textbf{Enabled By} &  & \\
    \hline
    \textbf{Supports} &  & \\
    \hline
    \textbf{Derived From} &  & \\
    \hline
    \textbf{Reviewed Status} & Generated candidate: keep / modify / relocate / reject; verify whether independently testable. & \\
    \hline
  \end{tabularx}
\end{table}

Intent. State Monitoring is proposed as a discrete, testable enabling behaviour for LHR-Aware Multidisciplinary Collaboration. It is included for architectural review, not as an established implementation obligation.

\noindent\textbf{Feature boundary.} This behaviour is scoped to LHR-Aware Multidisciplinary Collaboration; an identically named Feature elsewhere may carry different business meaning.\par
\featuretable{
--- & --- & \emph{Not yet captured} & \emph{Complete during review} \\
}
\subsubsection{Feature - Exception Management}\label{cap:bec-2df10aa6}
\vspace{-1em}
\begin{table}[htbp]
  \centering
  \footnotesize
  \begin{tabularx}{1\textwidth}{|l|l|X|}
    \hline
    \textbf{Metadata Name} & \textbf{Metadata Value} & \textbf{Description} \\
    \hline
    \textbf{Reference Id} & \texttt{BEC-2DF10AA6} & \\
    \hline
    \textbf{Status} & Generated candidate & \\
    \hline
    \textbf{Ancestry} &  & \\
    \hline
    \textbf{Source} & Semantic-pattern candidate & \\
    \hline
    \textbf{Enabled By} &  & \\
    \hline
    \textbf{Supports} &  & \\
    \hline
    \textbf{Derived From} &  & \\
    \hline
    \textbf{Reviewed Status} & Generated candidate: keep / modify / relocate / reject; verify whether independently testable. & \\
    \hline
  \end{tabularx}
\end{table}

Intent. Exception Management is proposed as a discrete, testable enabling behaviour for LHR-Aware Multidisciplinary Collaboration. It is included for architectural review, not as an established implementation obligation.

\noindent\textbf{Feature boundary.} This behaviour is scoped to LHR-Aware Multidisciplinary Collaboration; an identically named Feature elsewhere may carry different business meaning.\par
\featuretable{
--- & --- & \emph{Not yet captured} & \emph{Complete during review} \\
}
\subsubsection{Feature - Information Retrieval}\label{cap:bec-0cd2c755}
\vspace{-1em}
\begin{table}[htbp]
  \centering
  \footnotesize
  \begin{tabularx}{1\textwidth}{|l|l|X|}
    \hline
    \textbf{Metadata Name} & \textbf{Metadata Value} & \textbf{Description} \\
    \hline
    \textbf{Reference Id} & \texttt{BEC-0CD2C755} & \\
    \hline
    \textbf{Status} & Generated candidate & \\
    \hline
    \textbf{Ancestry} &  & \\
    \hline
    \textbf{Source} & Semantic-pattern candidate & \\
    \hline
    \textbf{Enabled By} &  & \\
    \hline
    \textbf{Supports} &  & \\
    \hline
    \textbf{Derived From} &  & \\
    \hline
    \textbf{Reviewed Status} & Generated candidate: keep / modify / relocate / reject; verify whether independently testable. & \\
    \hline
  \end{tabularx}
\end{table}

Intent. Information Retrieval is proposed as a discrete, testable enabling behaviour for LHR-Aware Multidisciplinary Collaboration. It is included for architectural review, not as an established implementation obligation.

\noindent\textbf{Feature boundary.} This behaviour is scoped to LHR-Aware Multidisciplinary Collaboration; an identically named Feature elsewhere may carry different business meaning.\par
\featuretable{
--- & --- & \emph{Not yet captured} & \emph{Complete during review} \\
}

\section{CT2 - General Collaboration}\label{cap:bec-e1807e93}
\vspace{-0.5em}
\begin{table}[htbp]
  \centering
  \footnotesize
  \begin{tabularx}{1\textwidth}{|l|l|X|}
    \hline
    \multicolumn{3}{c}{\textbf{Capability Metadata}} \\
    \hline
    \textbf{Name} & \textbf{Value} & \textbf{Comment and/or Description} \\
    \hline
    \textbf{Reference Id} & \texttt{BEC-E1807E93} & \\
    \hline
    \textbf{Status} & Candidate ancestry & \\
    \hline
    \textbf{Ancestry} &  & \\
    \hline
    \textbf{Source} &  & \\
    \hline
    \textbf{Enabled By} &  & \\
    \hline
    \textbf{Supports} &  & \\
    \hline
    \textbf{Derived From} &  & \\
    \hline
    \textbf{Reviewed Status} &  & \\
    \hline
  \end{tabularx}
\end{table}

\begin{table}[htbp]
  \centering
  \footnotesize
  \begin{tabularx}{1\textwidth}{|l|l|X|}
    \hline
    \multicolumn{3}{c}{\textbf{Contained Capabilities}} \\
    \hline
    \textbf{Capability Name} & \textbf{Capability Id} & \textbf{Comment} \\
    \hline
    Operational Messaging & \texttt{cap:bec-2d46a392} (at \ref{cap:bec-2d46a392}) & \\
    \hline
  \end{tabularx}
\end{table}

Intent. General Collaboration is a candidate enabling capability within Shared Information and Coordination Enablement. It frames the reusable system-supported behaviours needed to support this healthcare context without prescribing a product, interface, technology or deployment boundary.

Boundary. This capability supports the business activity and does not inherit the authority or accountability of the Business Capability it enables. Its child Features are review candidates until ancestry and semantics are adjudicated.

\subsection{CT3 - Operational Messaging}\label{cap:bec-2d46a392}
\vspace{-0.5em}
\begin{table}[htbp]
  \centering
  \footnotesize
  \begin{tabularx}{1\textwidth}{|l|l|X|}
    \hline
    \multicolumn{3}{c}{\textbf{Capability Metadata}} \\
    \hline
    \textbf{Name} & \textbf{Value} & \textbf{Comment and/or Description} \\
    \hline
    \textbf{Reference Id} & \texttt{BEC-2D46A392} & \\
    \hline
    \textbf{Status} & Candidate ancestry & \\
    \hline
    \textbf{Ancestry} &  & \\
    \hline
    \textbf{Source} &  & \\
    \hline
    \textbf{Enabled By} &  & \\
    \hline
    \textbf{Supports} &  & \\
    \hline
    \textbf{Derived From} &  & \\
    \hline
    \textbf{Reviewed Status} &  & \\
    \hline
  \end{tabularx}
\end{table}

\begin{table}[htbp]
  \centering
  \footnotesize
  \begin{tabularx}{1\textwidth}{|l|l|X|}
    \hline
    \multicolumn{3}{c}{\textbf{Contained Features}} \\
    \hline
    \textbf{Feature Name} & \textbf{Feature Id} & \textbf{Comment} \\
    \hline
    Operational Team Chat Enablement & \texttt{cap:bec-1631787f} (at \ref{cap:bec-1631787f}) & \\
    \hline
    Context Establishment & \texttt{cap:bec-074f43d0} (at \ref{cap:bec-074f43d0}) & \\
    \hline
    State Monitoring & \texttt{cap:bec-d75b5d46} (at \ref{cap:bec-d75b5d46}) & \\
    \hline
    Exception Management & \texttt{cap:bec-5368e298} (at \ref{cap:bec-5368e298}) & \\
    \hline
    Information Retrieval & \texttt{cap:bec-7d688c5a} (at \ref{cap:bec-7d688c5a}) & \\
    \hline
  \end{tabularx}
\end{table}

Intent. Operational Messaging is a candidate enabling capability within General Collaboration. It frames the reusable system-supported behaviours needed to support this healthcare context without prescribing a product, interface, technology or deployment boundary.

Boundary. This capability supports the business activity and does not inherit the authority or accountability of the Business Capability it enables. Its child Features are review candidates until ancestry and semantics are adjudicated.

\subsubsection{Feature - Operational Team Chat Enablement}\label{cap:bec-1631787f}
\vspace{-1em}
\begin{table}[htbp]
  \centering
  \footnotesize
  \begin{tabularx}{1\textwidth}{|l|l|X|}
    \hline
    \textbf{Metadata Name} & \textbf{Metadata Value} & \textbf{Description} \\
    \hline
    \textbf{Reference Id} & \texttt{BEC-1631787F} & \\
    \hline
    \textbf{Status} & Existing Feature; candidate CT3 placement & \\
    \hline
    \textbf{Ancestry} &  & \\
    \hline
    \textbf{Source} & FEAT-ISE-19 & \\
    \hline
    \textbf{Enabled By} &  & \\
    \hline
    \textbf{Supports} &  & \\
    \hline
    \textbf{Derived From} &  & \\
    \hline
    \textbf{Reviewed Status} & Confirm existing Feature meaning and proposed CT3 ownership. & \\
    \hline
  \end{tabularx}
\end{table}

Intent. Operational Team Chat Enablement provides a distinct enabling behaviour in the context of Operational Messaging. This is an existing named Feature; its placement under the candidate CT3 remains subject to review.

\noindent\textbf{Feature boundary.} This behaviour is scoped to Operational Messaging; an identically named Feature elsewhere may carry different business meaning.\par
\featuretable{
Function & \texttt{FN-43B65CC5} & Coordinate Operational Collaboration & Catalogue \\
Service & \texttt{SV-69BC1ACC} & Operational Collaboration Service & Linked to preceding Function in source \\
}
\subsubsection{Feature - Context Establishment}\label{cap:bec-074f43d0}
\vspace{-1em}
\begin{table}[htbp]
  \centering
  \footnotesize
  \begin{tabularx}{1\textwidth}{|l|l|X|}
    \hline
    \textbf{Metadata Name} & \textbf{Metadata Value} & \textbf{Description} \\
    \hline
    \textbf{Reference Id} & \texttt{BEC-074F43D0} & \\
    \hline
    \textbf{Status} & Generated candidate & \\
    \hline
    \textbf{Ancestry} &  & \\
    \hline
    \textbf{Source} & Semantic-pattern candidate & \\
    \hline
    \textbf{Enabled By} &  & \\
    \hline
    \textbf{Supports} &  & \\
    \hline
    \textbf{Derived From} &  & \\
    \hline
    \textbf{Reviewed Status} & Generated candidate: keep / modify / relocate / reject; verify whether independently testable. & \\
    \hline
  \end{tabularx}
\end{table}

Intent. Context Establishment is proposed as a discrete, testable enabling behaviour for Operational Messaging. It is included for architectural review, not as an established implementation obligation.

\noindent\textbf{Feature boundary.} This behaviour is scoped to Operational Messaging; an identically named Feature elsewhere may carry different business meaning.\par
\featuretable{
--- & --- & \emph{Not yet captured} & \emph{Complete during review} \\
}
\subsubsection{Feature - State Monitoring}\label{cap:bec-d75b5d46}
\vspace{-1em}
\begin{table}[htbp]
  \centering
  \footnotesize
  \begin{tabularx}{1\textwidth}{|l|l|X|}
    \hline
    \textbf{Metadata Name} & \textbf{Metadata Value} & \textbf{Description} \\
    \hline
    \textbf{Reference Id} & \texttt{BEC-D75B5D46} & \\
    \hline
    \textbf{Status} & Generated candidate & \\
    \hline
    \textbf{Ancestry} &  & \\
    \hline
    \textbf{Source} & Semantic-pattern candidate & \\
    \hline
    \textbf{Enabled By} &  & \\
    \hline
    \textbf{Supports} &  & \\
    \hline
    \textbf{Derived From} &  & \\
    \hline
    \textbf{Reviewed Status} & Generated candidate: keep / modify / relocate / reject; verify whether independently testable. & \\
    \hline
  \end{tabularx}
\end{table}

Intent. State Monitoring is proposed as a discrete, testable enabling behaviour for Operational Messaging. It is included for architectural review, not as an established implementation obligation.

\noindent\textbf{Feature boundary.} This behaviour is scoped to Operational Messaging; an identically named Feature elsewhere may carry different business meaning.\par
\featuretable{
--- & --- & \emph{Not yet captured} & \emph{Complete during review} \\
}
\subsubsection{Feature - Exception Management}\label{cap:bec-5368e298}
\vspace{-1em}
\begin{table}[htbp]
  \centering
  \footnotesize
  \begin{tabularx}{1\textwidth}{|l|l|X|}
    \hline
    \textbf{Metadata Name} & \textbf{Metadata Value} & \textbf{Description} \\
    \hline
    \textbf{Reference Id} & \texttt{BEC-5368E298} & \\
    \hline
    \textbf{Status} & Generated candidate & \\
    \hline
    \textbf{Ancestry} &  & \\
    \hline
    \textbf{Source} & Semantic-pattern candidate & \\
    \hline
    \textbf{Enabled By} &  & \\
    \hline
    \textbf{Supports} &  & \\
    \hline
    \textbf{Derived From} &  & \\
    \hline
    \textbf{Reviewed Status} & Generated candidate: keep / modify / relocate / reject; verify whether independently testable. & \\
    \hline
  \end{tabularx}
\end{table}

Intent. Exception Management is proposed as a discrete, testable enabling behaviour for Operational Messaging. It is included for architectural review, not as an established implementation obligation.

\noindent\textbf{Feature boundary.} This behaviour is scoped to Operational Messaging; an identically named Feature elsewhere may carry different business meaning.\par
\featuretable{
--- & --- & \emph{Not yet captured} & \emph{Complete during review} \\
}
\subsubsection{Feature - Information Retrieval}\label{cap:bec-7d688c5a}
\vspace{-1em}
\begin{table}[htbp]
  \centering
  \footnotesize
  \begin{tabularx}{1\textwidth}{|l|l|X|}
    \hline
    \textbf{Metadata Name} & \textbf{Metadata Value} & \textbf{Description} \\
    \hline
    \textbf{Reference Id} & \texttt{BEC-7D688C5A} & \\
    \hline
    \textbf{Status} & Generated candidate & \\
    \hline
    \textbf{Ancestry} &  & \\
    \hline
    \textbf{Source} & Semantic-pattern candidate & \\
    \hline
    \textbf{Enabled By} &  & \\
    \hline
    \textbf{Supports} &  & \\
    \hline
    \textbf{Derived From} &  & \\
    \hline
    \textbf{Reviewed Status} & Generated candidate: keep / modify / relocate / reject; verify whether independently testable. & \\
    \hline
  \end{tabularx}
\end{table}

Intent. Information Retrieval is proposed as a discrete, testable enabling behaviour for Operational Messaging. It is included for architectural review, not as an established implementation obligation.

\noindent\textbf{Feature boundary.} This behaviour is scoped to Operational Messaging; an identically named Feature elsewhere may carry different business meaning.\par
\featuretable{
--- & --- & \emph{Not yet captured} & \emph{Complete during review} \\
}

\section{CT2 - Workflow \& Activity Coordination}\label{cap:bec-0c4a6874}
\vspace{-0.5em}
\begin{table}[htbp]
  \centering
  \footnotesize
  \begin{tabularx}{1\textwidth}{|l|l|X|}
    \hline
    \multicolumn{3}{c}{\textbf{Capability Metadata}} \\
    \hline
    \textbf{Name} & \textbf{Value} & \textbf{Comment and/or Description} \\
    \hline
    \textbf{Reference Id} & \texttt{BEC-0C4A6874} & \\
    \hline
    \textbf{Status} & Candidate ancestry & \\
    \hline
    \textbf{Ancestry} &  & \\
    \hline
    \textbf{Source} &  & \\
    \hline
    \textbf{Enabled By} &  & \\
    \hline
    \textbf{Supports} &  & \\
    \hline
    \textbf{Derived From} &  & \\
    \hline
    \textbf{Reviewed Status} &  & \\
    \hline
  \end{tabularx}
\end{table}

\begin{table}[htbp]
  \centering
  \footnotesize
  \begin{tabularx}{1\textwidth}{|l|l|X|}
    \hline
    \multicolumn{3}{c}{\textbf{Contained Capabilities}} \\
    \hline
    \textbf{Capability Name} & \textbf{Capability Id} & \textbf{Comment} \\
    \hline
    Multi-Work-Unit Orchestration & \texttt{cap:bec-a3e2c598} (at \ref{cap:bec-a3e2c598}) & \\
    \hline
  \end{tabularx}
\end{table}

Intent. Workflow \& Activity Coordination is a candidate enabling capability within Shared Information and Coordination Enablement. It frames the reusable system-supported behaviours needed to support this healthcare context without prescribing a product, interface, technology or deployment boundary.

Boundary. This capability supports the business activity and does not inherit the authority or accountability of the Business Capability it enables. Its child Features are review candidates until ancestry and semantics are adjudicated.

\subsection{CT3 - Multi-Work-Unit Orchestration}\label{cap:bec-a3e2c598}
\vspace{-0.5em}
\begin{table}[htbp]
  \centering
  \footnotesize
  \begin{tabularx}{1\textwidth}{|l|l|X|}
    \hline
    \multicolumn{3}{c}{\textbf{Capability Metadata}} \\
    \hline
    \textbf{Name} & \textbf{Value} & \textbf{Comment and/or Description} \\
    \hline
    \textbf{Reference Id} & \texttt{BEC-A3E2C598} & \\
    \hline
    \textbf{Status} & Candidate ancestry & \\
    \hline
    \textbf{Ancestry} &  & \\
    \hline
    \textbf{Source} &  & \\
    \hline
    \textbf{Enabled By} &  & \\
    \hline
    \textbf{Supports} &  & \\
    \hline
    \textbf{Derived From} &  & \\
    \hline
    \textbf{Reviewed Status} &  & \\
    \hline
  \end{tabularx}
\end{table}

\begin{table}[htbp]
  \centering
  \footnotesize
  \begin{tabularx}{1\textwidth}{|l|l|X|}
    \hline
    \multicolumn{3}{c}{\textbf{Contained Features}} \\
    \hline
    \textbf{Feature Name} & \textbf{Feature Id} & \textbf{Comment} \\
    \hline
    Work Order Progression & \texttt{cap:bec-36c99ccb} (at \ref{cap:bec-36c99ccb}) & \\
    \hline
    To Do Decision \& Approval Coordination & \texttt{cap:bec-4306ddd8} (at \ref{cap:bec-4306ddd8}) & \\
    \hline
    Synthetic Task Orchestration & \texttt{cap:bec-782cd2a4} (at \ref{cap:bec-782cd2a4}) & \\
    \hline
    Workflow Timeout \& Escalation Oversight & \texttt{cap:bec-76b22b19} (at \ref{cap:bec-76b22b19}) & \\
    \hline
    Context Establishment & \texttt{cap:bec-a1649a9c} (at \ref{cap:bec-a1649a9c}) & \\
    \hline
    Eligibility Validation & \texttt{cap:bec-2078c8d9} (at \ref{cap:bec-2078c8d9}) & \\
    \hline
    Assignment \& Routing & \texttt{cap:bec-9ffa6f7f} (at \ref{cap:bec-9ffa6f7f}) & \\
    \hline
    State Tracking & \texttt{cap:bec-d25d58d7} (at \ref{cap:bec-d25d58d7}) & \\
    \hline
    Exception Management & \texttt{cap:bec-642bd633} (at \ref{cap:bec-642bd633}) & \\
    \hline
    Outcome Reporting & \texttt{cap:bec-6b53749d} (at \ref{cap:bec-6b53749d}) & \\
    \hline
  \end{tabularx}
\end{table}

Intent. Multi-Work-Unit Orchestration is a candidate enabling capability within Workflow \& Activity Coordination. It frames the reusable system-supported behaviours needed to support this healthcare context without prescribing a product, interface, technology or deployment boundary.

Boundary. This capability supports the business activity and does not inherit the authority or accountability of the Business Capability it enables. Its child Features are review candidates until ancestry and semantics are adjudicated.

\subsubsection{Feature - Work Order Progression}\label{cap:bec-36c99ccb}
\vspace{-1em}
\begin{table}[htbp]
  \centering
  \footnotesize
  \begin{tabularx}{1\textwidth}{|l|l|X|}
    \hline
    \textbf{Metadata Name} & \textbf{Metadata Value} & \textbf{Description} \\
    \hline
    \textbf{Reference Id} & \texttt{BEC-36C99CCB} & \\
    \hline
    \textbf{Status} & Existing Feature; candidate CT3 placement & \\
    \hline
    \textbf{Ancestry} &  & \\
    \hline
    \textbf{Source} & FEAT-ISE-20 & \\
    \hline
    \textbf{Enabled By} &  & \\
    \hline
    \textbf{Supports} &  & \\
    \hline
    \textbf{Derived From} &  & \\
    \hline
    \textbf{Reviewed Status} & Confirm existing Feature meaning and proposed CT3 ownership. & \\
    \hline
  \end{tabularx}
\end{table}

Intent. Work Order Progression provides a distinct enabling behaviour in the context of Multi-Work-Unit Orchestration. This is an existing named Feature; its placement under the candidate CT3 remains subject to review.

\noindent\textbf{Feature boundary.} This behaviour is scoped to Multi-Work-Unit Orchestration; an identically named Feature elsewhere may carry different business meaning.\par
\featuretable{
Function & \texttt{FN-17E1ABC5} & Coordinate Work Order & Catalogue \\
Service & \texttt{SV-666419B8} & Work Order Coordination Service & Linked to preceding Function in source \\
Process & \texttt{PR-E16BB9F3} & Work Order Progression Process & Linked to preceding Function in source \\
}
\subsubsection{Feature - To Do Decision \& Approval Coordination}\label{cap:bec-4306ddd8}
\vspace{-1em}
\begin{table}[htbp]
  \centering
  \footnotesize
  \begin{tabularx}{1\textwidth}{|l|l|X|}
    \hline
    \textbf{Metadata Name} & \textbf{Metadata Value} & \textbf{Description} \\
    \hline
    \textbf{Reference Id} & \texttt{BEC-4306DDD8} & \\
    \hline
    \textbf{Status} & Existing Feature; candidate CT3 placement & \\
    \hline
    \textbf{Ancestry} &  & \\
    \hline
    \textbf{Source} & FEAT-ISE-21 & \\
    \hline
    \textbf{Enabled By} &  & \\
    \hline
    \textbf{Supports} &  & \\
    \hline
    \textbf{Derived From} &  & \\
    \hline
    \textbf{Reviewed Status} & Confirm existing Feature meaning and proposed CT3 ownership. & \\
    \hline
  \end{tabularx}
\end{table}

Intent. To Do Decision \& Approval Coordination provides a distinct enabling behaviour in the context of Multi-Work-Unit Orchestration. This is an existing named Feature; its placement under the candidate CT3 remains subject to review.

\noindent\textbf{Feature boundary.} This behaviour is scoped to Multi-Work-Unit Orchestration; an identically named Feature elsewhere may carry different business meaning.\par
\featuretable{
Function & \texttt{FN-95B5EB70} & Coordinate To Do & Catalogue \\
Service & \texttt{SV-23C15E70} & To Do Management Service & Linked to preceding Function in source \\
Process & \texttt{PR-C99244D8} & To Do Progression Process & Linked to preceding Function in source \\
}
\subsubsection{Feature - Synthetic Task Orchestration}\label{cap:bec-782cd2a4}
\vspace{-1em}
\begin{table}[htbp]
  \centering
  \footnotesize
  \begin{tabularx}{1\textwidth}{|l|l|X|}
    \hline
    \textbf{Metadata Name} & \textbf{Metadata Value} & \textbf{Description} \\
    \hline
    \textbf{Reference Id} & \texttt{BEC-782CD2A4} & \\
    \hline
    \textbf{Status} & Existing Feature; candidate CT3 placement & \\
    \hline
    \textbf{Ancestry} &  & \\
    \hline
    \textbf{Source} & FEAT-ISE-22 & \\
    \hline
    \textbf{Enabled By} &  & \\
    \hline
    \textbf{Supports} &  & \\
    \hline
    \textbf{Derived From} &  & \\
    \hline
    \textbf{Reviewed Status} & Confirm existing Feature meaning and proposed CT3 ownership. & \\
    \hline
  \end{tabularx}
\end{table}

Intent. Synthetic Task Orchestration provides a distinct enabling behaviour in the context of Multi-Work-Unit Orchestration. This is an existing named Feature; its placement under the candidate CT3 remains subject to review.

\noindent\textbf{Feature boundary.} This behaviour is scoped to Multi-Work-Unit Orchestration; an identically named Feature elsewhere may carry different business meaning.\par
\featuretable{
Function & \texttt{FN-F16A1D5D} & Coordinate Synthetic Task & Catalogue \\
Service & \texttt{SV-A03CAAF8} & Synthetic Task Execution Service & Linked to preceding Function in source \\
Process & \texttt{PR-3B4DB284} & Synthetic Task Progression Process & Linked to preceding Function in source \\
}
\subsubsection{Feature - Workflow Timeout \& Escalation Oversight}\label{cap:bec-76b22b19}
\vspace{-1em}
\begin{table}[htbp]
  \centering
  \footnotesize
  \begin{tabularx}{1\textwidth}{|l|l|X|}
    \hline
    \textbf{Metadata Name} & \textbf{Metadata Value} & \textbf{Description} \\
    \hline
    \textbf{Reference Id} & \texttt{BEC-76B22B19} & \\
    \hline
    \textbf{Status} & Existing Feature; candidate CT3 placement & \\
    \hline
    \textbf{Ancestry} &  & \\
    \hline
    \textbf{Source} & FEAT-ISE-23 & \\
    \hline
    \textbf{Enabled By} &  & \\
    \hline
    \textbf{Supports} &  & \\
    \hline
    \textbf{Derived From} &  & \\
    \hline
    \textbf{Reviewed Status} & Confirm existing Feature meaning and proposed CT3 ownership. & \\
    \hline
  \end{tabularx}
\end{table}

Intent. Workflow Timeout \& Escalation Oversight provides a distinct enabling behaviour in the context of Multi-Work-Unit Orchestration. This is an existing named Feature; its placement under the candidate CT3 remains subject to review.

\noindent\textbf{Feature boundary.} This behaviour is scoped to Multi-Work-Unit Orchestration; an identically named Feature elsewhere may carry different business meaning.\par
\featuretable{
Function & \texttt{FN-EB78AFD1} & Supervise Activity Timeout / Escalation & Catalogue \\
Service & \texttt{SV-A8BB5F66} & Activity Escalation Telemetry & Linked to preceding Function in source \\
}
\subsubsection{Feature - Context Establishment}\label{cap:bec-a1649a9c}
\vspace{-1em}
\begin{table}[htbp]
  \centering
  \footnotesize
  \begin{tabularx}{1\textwidth}{|l|l|X|}
    \hline
    \textbf{Metadata Name} & \textbf{Metadata Value} & \textbf{Description} \\
    \hline
    \textbf{Reference Id} & \texttt{BEC-A1649A9C} & \\
    \hline
    \textbf{Status} & Generated candidate & \\
    \hline
    \textbf{Ancestry} &  & \\
    \hline
    \textbf{Source} & Semantic-pattern candidate & \\
    \hline
    \textbf{Enabled By} &  & \\
    \hline
    \textbf{Supports} &  & \\
    \hline
    \textbf{Derived From} &  & \\
    \hline
    \textbf{Reviewed Status} & Generated candidate: keep / modify / relocate / reject; verify whether independently testable. & \\
    \hline
  \end{tabularx}
\end{table}

Intent. Context Establishment is proposed as a discrete, testable enabling behaviour for Multi-Work-Unit Orchestration. It is included for architectural review, not as an established implementation obligation.

\noindent\textbf{Feature boundary.} This behaviour is scoped to Multi-Work-Unit Orchestration; an identically named Feature elsewhere may carry different business meaning.\par
\featuretable{
--- & --- & \emph{Not yet captured} & \emph{Complete during review} \\
}
\subsubsection{Feature - Eligibility Validation}\label{cap:bec-2078c8d9}
\vspace{-1em}
\begin{table}[htbp]
  \centering
  \footnotesize
  \begin{tabularx}{1\textwidth}{|l|l|X|}
    \hline
    \textbf{Metadata Name} & \textbf{Metadata Value} & \textbf{Description} \\
    \hline
    \textbf{Reference Id} & \texttt{BEC-2078C8D9} & \\
    \hline
    \textbf{Status} & Generated candidate & \\
    \hline
    \textbf{Ancestry} &  & \\
    \hline
    \textbf{Source} & Semantic-pattern candidate & \\
    \hline
    \textbf{Enabled By} &  & \\
    \hline
    \textbf{Supports} &  & \\
    \hline
    \textbf{Derived From} &  & \\
    \hline
    \textbf{Reviewed Status} & Generated candidate: keep / modify / relocate / reject; verify whether independently testable. & \\
    \hline
  \end{tabularx}
\end{table}

Intent. Eligibility Validation is proposed as a discrete, testable enabling behaviour for Multi-Work-Unit Orchestration. It is included for architectural review, not as an established implementation obligation.

\noindent\textbf{Feature boundary.} This behaviour is scoped to Multi-Work-Unit Orchestration; an identically named Feature elsewhere may carry different business meaning.\par
\featuretable{
--- & --- & \emph{Not yet captured} & \emph{Complete during review} \\
}
\subsubsection{Feature - Assignment \& Routing}\label{cap:bec-9ffa6f7f}
\vspace{-1em}
\begin{table}[htbp]
  \centering
  \footnotesize
  \begin{tabularx}{1\textwidth}{|l|l|X|}
    \hline
    \textbf{Metadata Name} & \textbf{Metadata Value} & \textbf{Description} \\
    \hline
    \textbf{Reference Id} & \texttt{BEC-9FFA6F7F} & \\
    \hline
    \textbf{Status} & Generated candidate & \\
    \hline
    \textbf{Ancestry} &  & \\
    \hline
    \textbf{Source} & Semantic-pattern candidate & \\
    \hline
    \textbf{Enabled By} &  & \\
    \hline
    \textbf{Supports} &  & \\
    \hline
    \textbf{Derived From} &  & \\
    \hline
    \textbf{Reviewed Status} & Generated candidate: keep / modify / relocate / reject; verify whether independently testable. & \\
    \hline
  \end{tabularx}
\end{table}

Intent. Assignment \& Routing is proposed as a discrete, testable enabling behaviour for Multi-Work-Unit Orchestration. It is included for architectural review, not as an established implementation obligation.

\noindent\textbf{Feature boundary.} This behaviour is scoped to Multi-Work-Unit Orchestration; an identically named Feature elsewhere may carry different business meaning.\par
\featuretable{
--- & --- & \emph{Not yet captured} & \emph{Complete during review} \\
}
\subsubsection{Feature - State Tracking}\label{cap:bec-d25d58d7}
\vspace{-1em}
\begin{table}[htbp]
  \centering
  \footnotesize
  \begin{tabularx}{1\textwidth}{|l|l|X|}
    \hline
    \textbf{Metadata Name} & \textbf{Metadata Value} & \textbf{Description} \\
    \hline
    \textbf{Reference Id} & \texttt{BEC-D25D58D7} & \\
    \hline
    \textbf{Status} & Generated candidate & \\
    \hline
    \textbf{Ancestry} &  & \\
    \hline
    \textbf{Source} & Semantic-pattern candidate & \\
    \hline
    \textbf{Enabled By} &  & \\
    \hline
    \textbf{Supports} &  & \\
    \hline
    \textbf{Derived From} &  & \\
    \hline
    \textbf{Reviewed Status} & Generated candidate: keep / modify / relocate / reject; verify whether independently testable. & \\
    \hline
  \end{tabularx}
\end{table}

Intent. State Tracking is proposed as a discrete, testable enabling behaviour for Multi-Work-Unit Orchestration. It is included for architectural review, not as an established implementation obligation.

\noindent\textbf{Feature boundary.} This behaviour is scoped to Multi-Work-Unit Orchestration; an identically named Feature elsewhere may carry different business meaning.\par
\featuretable{
--- & --- & \emph{Not yet captured} & \emph{Complete during review} \\
}
\subsubsection{Feature - Exception Management}\label{cap:bec-642bd633}
\vspace{-1em}
\begin{table}[htbp]
  \centering
  \footnotesize
  \begin{tabularx}{1\textwidth}{|l|l|X|}
    \hline
    \textbf{Metadata Name} & \textbf{Metadata Value} & \textbf{Description} \\
    \hline
    \textbf{Reference Id} & \texttt{BEC-642BD633} & \\
    \hline
    \textbf{Status} & Generated candidate & \\
    \hline
    \textbf{Ancestry} &  & \\
    \hline
    \textbf{Source} & Semantic-pattern candidate & \\
    \hline
    \textbf{Enabled By} &  & \\
    \hline
    \textbf{Supports} &  & \\
    \hline
    \textbf{Derived From} &  & \\
    \hline
    \textbf{Reviewed Status} & Generated candidate: keep / modify / relocate / reject; verify whether independently testable. & \\
    \hline
  \end{tabularx}
\end{table}

Intent. Exception Management is proposed as a discrete, testable enabling behaviour for Multi-Work-Unit Orchestration. It is included for architectural review, not as an established implementation obligation.

\noindent\textbf{Feature boundary.} This behaviour is scoped to Multi-Work-Unit Orchestration; an identically named Feature elsewhere may carry different business meaning.\par
\featuretable{
--- & --- & \emph{Not yet captured} & \emph{Complete during review} \\
}
\subsubsection{Feature - Outcome Reporting}\label{cap:bec-6b53749d}
\vspace{-1em}
\begin{table}[htbp]
  \centering
  \footnotesize
  \begin{tabularx}{1\textwidth}{|l|l|X|}
    \hline
    \textbf{Metadata Name} & \textbf{Metadata Value} & \textbf{Description} \\
    \hline
    \textbf{Reference Id} & \texttt{BEC-6B53749D} & \\
    \hline
    \textbf{Status} & Generated candidate & \\
    \hline
    \textbf{Ancestry} &  & \\
    \hline
    \textbf{Source} & Semantic-pattern candidate & \\
    \hline
    \textbf{Enabled By} &  & \\
    \hline
    \textbf{Supports} &  & \\
    \hline
    \textbf{Derived From} &  & \\
    \hline
    \textbf{Reviewed Status} & Generated candidate: keep / modify / relocate / reject; verify whether independently testable. & \\
    \hline
  \end{tabularx}
\end{table}

Intent. Outcome Reporting is proposed as a discrete, testable enabling behaviour for Multi-Work-Unit Orchestration. It is included for architectural review, not as an established implementation obligation.

\noindent\textbf{Feature boundary.} This behaviour is scoped to Multi-Work-Unit Orchestration; an identically named Feature elsewhere may carry different business meaning.\par
\featuretable{
--- & --- & \emph{Not yet captured} & \emph{Complete during review} \\
}

\section{CT2 - Calendar Management}\label{cap:bec-78f5db56}
\vspace{-0.5em}
\begin{table}[htbp]
  \centering
  \footnotesize
  \begin{tabularx}{1\textwidth}{|l|l|X|}
    \hline
    \multicolumn{3}{c}{\textbf{Capability Metadata}} \\
    \hline
    \textbf{Name} & \textbf{Value} & \textbf{Comment and/or Description} \\
    \hline
    \textbf{Reference Id} & \texttt{BEC-78F5DB56} & \\
    \hline
    \textbf{Status} & Candidate ancestry & \\
    \hline
    \textbf{Ancestry} &  & \\
    \hline
    \textbf{Source} &  & \\
    \hline
    \textbf{Enabled By} &  & \\
    \hline
    \textbf{Supports} &  & \\
    \hline
    \textbf{Derived From} &  & \\
    \hline
    \textbf{Reviewed Status} &  & \\
    \hline
  \end{tabularx}
\end{table}

\begin{table}[htbp]
  \centering
  \footnotesize
  \begin{tabularx}{1\textwidth}{|l|l|X|}
    \hline
    \multicolumn{3}{c}{\textbf{Contained Capabilities}} \\
    \hline
    \textbf{Capability Name} & \textbf{Capability Id} & \textbf{Comment} \\
    \hline
    Operational \& Clinical Calendar Sync & \texttt{cap:bec-8cd7ea17} (at \ref{cap:bec-8cd7ea17}) & \\
    \hline
  \end{tabularx}
\end{table}

Intent. Calendar Management is a candidate enabling capability within Shared Information and Coordination Enablement. It frames the reusable system-supported behaviours needed to support this healthcare context without prescribing a product, interface, technology or deployment boundary.

Boundary. This capability supports the business activity and does not inherit the authority or accountability of the Business Capability it enables. Its child Features are review candidates until ancestry and semantics are adjudicated.

\subsection{CT3 - Operational \& Clinical Calendar Sync}\label{cap:bec-8cd7ea17}
\vspace{-0.5em}
\begin{table}[htbp]
  \centering
  \footnotesize
  \begin{tabularx}{1\textwidth}{|l|l|X|}
    \hline
    \multicolumn{3}{c}{\textbf{Capability Metadata}} \\
    \hline
    \textbf{Name} & \textbf{Value} & \textbf{Comment and/or Description} \\
    \hline
    \textbf{Reference Id} & \texttt{BEC-8CD7EA17} & \\
    \hline
    \textbf{Status} & Candidate ancestry & \\
    \hline
    \textbf{Ancestry} &  & \\
    \hline
    \textbf{Source} &  & \\
    \hline
    \textbf{Enabled By} &  & \\
    \hline
    \textbf{Supports} &  & \\
    \hline
    \textbf{Derived From} &  & \\
    \hline
    \textbf{Reviewed Status} &  & \\
    \hline
  \end{tabularx}
\end{table}

\begin{table}[htbp]
  \centering
  \footnotesize
  \begin{tabularx}{1\textwidth}{|l|l|X|}
    \hline
    \multicolumn{3}{c}{\textbf{Contained Features}} \\
    \hline
    \textbf{Feature Name} & \textbf{Feature Id} & \textbf{Comment} \\
    \hline
    Schedule Projection Integration & \texttt{cap:bec-79108df3} (at \ref{cap:bec-79108df3}) & \\
    \hline
    Temporal Event Correlation & \texttt{cap:bec-45628c87} (at \ref{cap:bec-45628c87}) & \\
    \hline
    Context Establishment & \texttt{cap:bec-9f5e17d1} (at \ref{cap:bec-9f5e17d1}) & \\
    \hline
    State Monitoring & \texttt{cap:bec-922d6fff} (at \ref{cap:bec-922d6fff}) & \\
    \hline
    Exception Management & \texttt{cap:bec-76d61f9a} (at \ref{cap:bec-76d61f9a}) & \\
    \hline
    Information Retrieval & \texttt{cap:bec-1022e17a} (at \ref{cap:bec-1022e17a}) & \\
    \hline
  \end{tabularx}
\end{table}

Intent. Operational \& Clinical Calendar Sync is a candidate enabling capability within Calendar Management. It frames the reusable system-supported behaviours needed to support this healthcare context without prescribing a product, interface, technology or deployment boundary.

Boundary. This capability supports the business activity and does not inherit the authority or accountability of the Business Capability it enables. Its child Features are review candidates until ancestry and semantics are adjudicated.

\subsubsection{Feature - Schedule Projection Integration}\label{cap:bec-79108df3}
\vspace{-1em}
\begin{table}[htbp]
  \centering
  \footnotesize
  \begin{tabularx}{1\textwidth}{|l|l|X|}
    \hline
    \textbf{Metadata Name} & \textbf{Metadata Value} & \textbf{Description} \\
    \hline
    \textbf{Reference Id} & \texttt{BEC-79108DF3} & \\
    \hline
    \textbf{Status} & Existing Feature; candidate CT3 placement & \\
    \hline
    \textbf{Ancestry} &  & \\
    \hline
    \textbf{Source} & FEAT-ISE-24 & \\
    \hline
    \textbf{Enabled By} &  & \\
    \hline
    \textbf{Supports} &  & \\
    \hline
    \textbf{Derived From} &  & \\
    \hline
    \textbf{Reviewed Status} & Confirm existing Feature meaning and proposed CT3 ownership. & \\
    \hline
  \end{tabularx}
\end{table}

Intent. Schedule Projection Integration provides a distinct enabling behaviour in the context of Operational \& Clinical Calendar Sync. This is an existing named Feature; its placement under the candidate CT3 remains subject to review.

\noindent\textbf{Feature boundary.} This behaviour is scoped to Operational \& Clinical Calendar Sync; an identically named Feature elsewhere may carry different business meaning.\par
\featuretable{
Function & \texttt{FN-DEC990C9} & Project Operational Schedule & Catalogue \\
Service & \texttt{SV-90F9835A} & Unified Calendar Projection Query & Linked to preceding Function in source \\
}
\subsubsection{Feature - Temporal Event Correlation}\label{cap:bec-45628c87}
\vspace{-1em}
\begin{table}[htbp]
  \centering
  \footnotesize
  \begin{tabularx}{1\textwidth}{|l|l|X|}
    \hline
    \textbf{Metadata Name} & \textbf{Metadata Value} & \textbf{Description} \\
    \hline
    \textbf{Reference Id} & \texttt{BEC-45628C87} & \\
    \hline
    \textbf{Status} & Existing Feature; candidate CT3 placement & \\
    \hline
    \textbf{Ancestry} &  & \\
    \hline
    \textbf{Source} & FEAT-ISE-25 & \\
    \hline
    \textbf{Enabled By} &  & \\
    \hline
    \textbf{Supports} &  & \\
    \hline
    \textbf{Derived From} &  & \\
    \hline
    \textbf{Reviewed Status} & Confirm existing Feature meaning and proposed CT3 ownership. & \\
    \hline
  \end{tabularx}
\end{table}

Intent. Temporal Event Correlation provides a distinct enabling behaviour in the context of Operational \& Clinical Calendar Sync. This is an existing named Feature; its placement under the candidate CT3 remains subject to review.

\noindent\textbf{Feature boundary.} This behaviour is scoped to Operational \& Clinical Calendar Sync; an identically named Feature elsewhere may carry different business meaning.\par
\featuretable{
Function & \texttt{FN-7421682C} & Correlate Activity with Temporal Context & Catalogue \\
Service & \texttt{SV-4B5A1F5C} & Temporal Context Resolution & Linked to preceding Function in source \\
}
\subsubsection{Feature - Context Establishment}\label{cap:bec-9f5e17d1}
\vspace{-1em}
\begin{table}[htbp]
  \centering
  \footnotesize
  \begin{tabularx}{1\textwidth}{|l|l|X|}
    \hline
    \textbf{Metadata Name} & \textbf{Metadata Value} & \textbf{Description} \\
    \hline
    \textbf{Reference Id} & \texttt{BEC-9F5E17D1} & \\
    \hline
    \textbf{Status} & Generated candidate & \\
    \hline
    \textbf{Ancestry} &  & \\
    \hline
    \textbf{Source} & Semantic-pattern candidate & \\
    \hline
    \textbf{Enabled By} &  & \\
    \hline
    \textbf{Supports} &  & \\
    \hline
    \textbf{Derived From} &  & \\
    \hline
    \textbf{Reviewed Status} & Generated candidate: keep / modify / relocate / reject; verify whether independently testable. & \\
    \hline
  \end{tabularx}
\end{table}

Intent. Context Establishment is proposed as a discrete, testable enabling behaviour for Operational \& Clinical Calendar Sync. It is included for architectural review, not as an established implementation obligation.

\noindent\textbf{Feature boundary.} This behaviour is scoped to Operational \& Clinical Calendar Sync; an identically named Feature elsewhere may carry different business meaning.\par
\featuretable{
--- & --- & \emph{Not yet captured} & \emph{Complete during review} \\
}
\subsubsection{Feature - State Monitoring}\label{cap:bec-922d6fff}
\vspace{-1em}
\begin{table}[htbp]
  \centering
  \footnotesize
  \begin{tabularx}{1\textwidth}{|l|l|X|}
    \hline
    \textbf{Metadata Name} & \textbf{Metadata Value} & \textbf{Description} \\
    \hline
    \textbf{Reference Id} & \texttt{BEC-922D6FFF} & \\
    \hline
    \textbf{Status} & Generated candidate & \\
    \hline
    \textbf{Ancestry} &  & \\
    \hline
    \textbf{Source} & Semantic-pattern candidate & \\
    \hline
    \textbf{Enabled By} &  & \\
    \hline
    \textbf{Supports} &  & \\
    \hline
    \textbf{Derived From} &  & \\
    \hline
    \textbf{Reviewed Status} & Generated candidate: keep / modify / relocate / reject; verify whether independently testable. & \\
    \hline
  \end{tabularx}
\end{table}

Intent. State Monitoring is proposed as a discrete, testable enabling behaviour for Operational \& Clinical Calendar Sync. It is included for architectural review, not as an established implementation obligation.

\noindent\textbf{Feature boundary.} This behaviour is scoped to Operational \& Clinical Calendar Sync; an identically named Feature elsewhere may carry different business meaning.\par
\featuretable{
--- & --- & \emph{Not yet captured} & \emph{Complete during review} \\
}
\subsubsection{Feature - Exception Management}\label{cap:bec-76d61f9a}
\vspace{-1em}
\begin{table}[htbp]
  \centering
  \footnotesize
  \begin{tabularx}{1\textwidth}{|l|l|X|}
    \hline
    \textbf{Metadata Name} & \textbf{Metadata Value} & \textbf{Description} \\
    \hline
    \textbf{Reference Id} & \texttt{BEC-76D61F9A} & \\
    \hline
    \textbf{Status} & Generated candidate & \\
    \hline
    \textbf{Ancestry} &  & \\
    \hline
    \textbf{Source} & Semantic-pattern candidate & \\
    \hline
    \textbf{Enabled By} &  & \\
    \hline
    \textbf{Supports} &  & \\
    \hline
    \textbf{Derived From} &  & \\
    \hline
    \textbf{Reviewed Status} & Generated candidate: keep / modify / relocate / reject; verify whether independently testable. & \\
    \hline
  \end{tabularx}
\end{table}

Intent. Exception Management is proposed as a discrete, testable enabling behaviour for Operational \& Clinical Calendar Sync. It is included for architectural review, not as an established implementation obligation.

\noindent\textbf{Feature boundary.} This behaviour is scoped to Operational \& Clinical Calendar Sync; an identically named Feature elsewhere may carry different business meaning.\par
\featuretable{
--- & --- & \emph{Not yet captured} & \emph{Complete during review} \\
}
\subsubsection{Feature - Information Retrieval}\label{cap:bec-1022e17a}
\vspace{-1em}
\begin{table}[htbp]
  \centering
  \footnotesize
  \begin{tabularx}{1\textwidth}{|l|l|X|}
    \hline
    \textbf{Metadata Name} & \textbf{Metadata Value} & \textbf{Description} \\
    \hline
    \textbf{Reference Id} & \texttt{BEC-1022E17A} & \\
    \hline
    \textbf{Status} & Generated candidate & \\
    \hline
    \textbf{Ancestry} &  & \\
    \hline
    \textbf{Source} & Semantic-pattern candidate & \\
    \hline
    \textbf{Enabled By} &  & \\
    \hline
    \textbf{Supports} &  & \\
    \hline
    \textbf{Derived From} &  & \\
    \hline
    \textbf{Reviewed Status} & Generated candidate: keep / modify / relocate / reject; verify whether independently testable. & \\
    \hline
  \end{tabularx}
\end{table}

Intent. Information Retrieval is proposed as a discrete, testable enabling behaviour for Operational \& Clinical Calendar Sync. It is included for architectural review, not as an established implementation obligation.

\noindent\textbf{Feature boundary.} This behaviour is scoped to Operational \& Clinical Calendar Sync; an identically named Feature elsewhere may carry different business meaning.\par
\featuretable{
--- & --- & \emph{Not yet captured} & \emph{Complete during review} \\
}

\section{CT2 - Presentation Services}\label{cap:bec-6aa02a96}
\vspace{-0.5em}
\begin{table}[htbp]
  \centering
  \footnotesize
  \begin{tabularx}{1\textwidth}{|l|l|X|}
    \hline
    \multicolumn{3}{c}{\textbf{Capability Metadata}} \\
    \hline
    \textbf{Name} & \textbf{Value} & \textbf{Comment and/or Description} \\
    \hline
    \textbf{Reference Id} & \texttt{BEC-6AA02A96} & \\
    \hline
    \textbf{Status} & Candidate ancestry & \\
    \hline
    \textbf{Ancestry} &  & \\
    \hline
    \textbf{Source} &  & \\
    \hline
    \textbf{Enabled By} &  & \\
    \hline
    \textbf{Supports} &  & \\
    \hline
    \textbf{Derived From} &  & \\
    \hline
    \textbf{Reviewed Status} &  & \\
    \hline
  \end{tabularx}
\end{table}

\begin{table}[htbp]
  \centering
  \footnotesize
  \begin{tabularx}{1\textwidth}{|l|l|X|}
    \hline
    \multicolumn{3}{c}{\textbf{Contained Capabilities}} \\
    \hline
    \textbf{Capability Name} & \textbf{Capability Id} & \textbf{Comment} \\
    \hline
    Decoupled Human Interaction \& Dashboards & \texttt{cap:bec-2dcf9553} (at \ref{cap:bec-2dcf9553}) & \\
    \hline
  \end{tabularx}
\end{table}

Intent. Presentation Services is a candidate enabling capability within Shared Information and Coordination Enablement. It frames the reusable system-supported behaviours needed to support this healthcare context without prescribing a product, interface, technology or deployment boundary.

Boundary. This capability supports the business activity and does not inherit the authority or accountability of the Business Capability it enables. Its child Features are review candidates until ancestry and semantics are adjudicated.

\subsection{CT3 - Decoupled Human Interaction \& Dashboards}\label{cap:bec-2dcf9553}
\vspace{-0.5em}
\begin{table}[htbp]
  \centering
  \footnotesize
  \begin{tabularx}{1\textwidth}{|l|l|X|}
    \hline
    \multicolumn{3}{c}{\textbf{Capability Metadata}} \\
    \hline
    \textbf{Name} & \textbf{Value} & \textbf{Comment and/or Description} \\
    \hline
    \textbf{Reference Id} & \texttt{BEC-2DCF9553} & \\
    \hline
    \textbf{Status} & Candidate ancestry & \\
    \hline
    \textbf{Ancestry} &  & \\
    \hline
    \textbf{Source} &  & \\
    \hline
    \textbf{Enabled By} &  & \\
    \hline
    \textbf{Supports} &  & \\
    \hline
    \textbf{Derived From} &  & \\
    \hline
    \textbf{Reviewed Status} &  & \\
    \hline
  \end{tabularx}
\end{table}

\begin{table}[htbp]
  \centering
  \footnotesize
  \begin{tabularx}{1\textwidth}{|l|l|X|}
    \hline
    \multicolumn{3}{c}{\textbf{Contained Features}} \\
    \hline
    \textbf{Feature Name} & \textbf{Feature Id} & \textbf{Comment} \\
    \hline
    Longitudinal Clinical \& Operational Presentation & \texttt{cap:bec-06ab86fa} (at \ref{cap:bec-06ab86fa}) & \\
    \hline
    Context Establishment & \texttt{cap:bec-53cda4e5} (at \ref{cap:bec-53cda4e5}) & \\
    \hline
    State Monitoring & \texttt{cap:bec-7cb4fd80} (at \ref{cap:bec-7cb4fd80}) & \\
    \hline
    Exception Management & \texttt{cap:bec-d8af8ac4} (at \ref{cap:bec-d8af8ac4}) & \\
    \hline
    Information Retrieval & \texttt{cap:bec-da832dce} (at \ref{cap:bec-da832dce}) & \\
    \hline
  \end{tabularx}
\end{table}

Intent. Decoupled Human Interaction \& Dashboards is a candidate enabling capability within Presentation Services. It frames the reusable system-supported behaviours needed to support this healthcare context without prescribing a product, interface, technology or deployment boundary.

Boundary. This capability supports the business activity and does not inherit the authority or accountability of the Business Capability it enables. Its child Features are review candidates until ancestry and semantics are adjudicated.

\subsubsection{Feature - Longitudinal Clinical \& Operational Presentation}\label{cap:bec-06ab86fa}
\vspace{-1em}
\begin{table}[htbp]
  \centering
  \footnotesize
  \begin{tabularx}{1\textwidth}{|l|l|X|}
    \hline
    \textbf{Metadata Name} & \textbf{Metadata Value} & \textbf{Description} \\
    \hline
    \textbf{Reference Id} & \texttt{BEC-06AB86FA} & \\
    \hline
    \textbf{Status} & Existing Feature; candidate CT3 placement & \\
    \hline
    \textbf{Ancestry} &  & \\
    \hline
    \textbf{Source} & FEAT-ISE-26 & \\
    \hline
    \textbf{Enabled By} &  & \\
    \hline
    \textbf{Supports} &  & \\
    \hline
    \textbf{Derived From} &  & \\
    \hline
    \textbf{Reviewed Status} & Confirm existing Feature meaning and proposed CT3 ownership. & \\
    \hline
  \end{tabularx}
\end{table}

Intent. Longitudinal Clinical \& Operational Presentation provides a distinct enabling behaviour in the context of Decoupled Human Interaction \& Dashboards. This is an existing named Feature; its placement under the candidate CT3 remains subject to review.

\noindent\textbf{Feature boundary.} This behaviour is scoped to Decoupled Human Interaction \& Dashboards; an identically named Feature elsewhere may carry different business meaning.\par
\featuretable{
Function & \texttt{FN-01E7046C} & Present Contextual Health and Operational Information & Catalogue \\
Service & \texttt{SV-1E730C1C} & Contextual Presentation Service & Linked to preceding Function in source \\
}
\subsubsection{Feature - Context Establishment}\label{cap:bec-53cda4e5}
\vspace{-1em}
\begin{table}[htbp]
  \centering
  \footnotesize
  \begin{tabularx}{1\textwidth}{|l|l|X|}
    \hline
    \textbf{Metadata Name} & \textbf{Metadata Value} & \textbf{Description} \\
    \hline
    \textbf{Reference Id} & \texttt{BEC-53CDA4E5} & \\
    \hline
    \textbf{Status} & Generated candidate & \\
    \hline
    \textbf{Ancestry} &  & \\
    \hline
    \textbf{Source} & Semantic-pattern candidate & \\
    \hline
    \textbf{Enabled By} &  & \\
    \hline
    \textbf{Supports} &  & \\
    \hline
    \textbf{Derived From} &  & \\
    \hline
    \textbf{Reviewed Status} & Generated candidate: keep / modify / relocate / reject; verify whether independently testable. & \\
    \hline
  \end{tabularx}
\end{table}

Intent. Context Establishment is proposed as a discrete, testable enabling behaviour for Decoupled Human Interaction \& Dashboards. It is included for architectural review, not as an established implementation obligation.

\noindent\textbf{Feature boundary.} This behaviour is scoped to Decoupled Human Interaction \& Dashboards; an identically named Feature elsewhere may carry different business meaning.\par
\featuretable{
--- & --- & \emph{Not yet captured} & \emph{Complete during review} \\
}
\subsubsection{Feature - State Monitoring}\label{cap:bec-7cb4fd80}
\vspace{-1em}
\begin{table}[htbp]
  \centering
  \footnotesize
  \begin{tabularx}{1\textwidth}{|l|l|X|}
    \hline
    \textbf{Metadata Name} & \textbf{Metadata Value} & \textbf{Description} \\
    \hline
    \textbf{Reference Id} & \texttt{BEC-7CB4FD80} & \\
    \hline
    \textbf{Status} & Generated candidate & \\
    \hline
    \textbf{Ancestry} &  & \\
    \hline
    \textbf{Source} & Semantic-pattern candidate & \\
    \hline
    \textbf{Enabled By} &  & \\
    \hline
    \textbf{Supports} &  & \\
    \hline
    \textbf{Derived From} &  & \\
    \hline
    \textbf{Reviewed Status} & Generated candidate: keep / modify / relocate / reject; verify whether independently testable. & \\
    \hline
  \end{tabularx}
\end{table}

Intent. State Monitoring is proposed as a discrete, testable enabling behaviour for Decoupled Human Interaction \& Dashboards. It is included for architectural review, not as an established implementation obligation.

\noindent\textbf{Feature boundary.} This behaviour is scoped to Decoupled Human Interaction \& Dashboards; an identically named Feature elsewhere may carry different business meaning.\par
\featuretable{
--- & --- & \emph{Not yet captured} & \emph{Complete during review} \\
}
\subsubsection{Feature - Exception Management}\label{cap:bec-d8af8ac4}
\vspace{-1em}
\begin{table}[htbp]
  \centering
  \footnotesize
  \begin{tabularx}{1\textwidth}{|l|l|X|}
    \hline
    \textbf{Metadata Name} & \textbf{Metadata Value} & \textbf{Description} \\
    \hline
    \textbf{Reference Id} & \texttt{BEC-D8AF8AC4} & \\
    \hline
    \textbf{Status} & Generated candidate & \\
    \hline
    \textbf{Ancestry} &  & \\
    \hline
    \textbf{Source} & Semantic-pattern candidate & \\
    \hline
    \textbf{Enabled By} &  & \\
    \hline
    \textbf{Supports} &  & \\
    \hline
    \textbf{Derived From} &  & \\
    \hline
    \textbf{Reviewed Status} & Generated candidate: keep / modify / relocate / reject; verify whether independently testable. & \\
    \hline
  \end{tabularx}
\end{table}

Intent. Exception Management is proposed as a discrete, testable enabling behaviour for Decoupled Human Interaction \& Dashboards. It is included for architectural review, not as an established implementation obligation.

\noindent\textbf{Feature boundary.} This behaviour is scoped to Decoupled Human Interaction \& Dashboards; an identically named Feature elsewhere may carry different business meaning.\par
\featuretable{
--- & --- & \emph{Not yet captured} & \emph{Complete during review} \\
}
\subsubsection{Feature - Information Retrieval}\label{cap:bec-da832dce}
\vspace{-1em}
\begin{table}[htbp]
  \centering
  \footnotesize
  \begin{tabularx}{1\textwidth}{|l|l|X|}
    \hline
    \textbf{Metadata Name} & \textbf{Metadata Value} & \textbf{Description} \\
    \hline
    \textbf{Reference Id} & \texttt{BEC-DA832DCE} & \\
    \hline
    \textbf{Status} & Generated candidate & \\
    \hline
    \textbf{Ancestry} &  & \\
    \hline
    \textbf{Source} & Semantic-pattern candidate & \\
    \hline
    \textbf{Enabled By} &  & \\
    \hline
    \textbf{Supports} &  & \\
    \hline
    \textbf{Derived From} &  & \\
    \hline
    \textbf{Reviewed Status} & Generated candidate: keep / modify / relocate / reject; verify whether independently testable. & \\
    \hline
  \end{tabularx}
\end{table}

Intent. Information Retrieval is proposed as a discrete, testable enabling behaviour for Decoupled Human Interaction \& Dashboards. It is included for architectural review, not as an established implementation obligation.

\noindent\textbf{Feature boundary.} This behaviour is scoped to Decoupled Human Interaction \& Dashboards; an identically named Feature elsewhere may carry different business meaning.\par
\featuretable{
--- & --- & \emph{Not yet captured} & \emph{Complete during review} \\
}

\section{CT2 - Information Design Governance}\label{cap:bec-233904da}
\vspace{-0.5em}
\begin{table}[htbp]
  \centering
  \footnotesize
  \begin{tabularx}{1\textwidth}{|l|l|X|}
    \hline
    \multicolumn{3}{c}{\textbf{Capability Metadata}} \\
    \hline
    \textbf{Name} & \textbf{Value} & \textbf{Comment and/or Description} \\
    \hline
    \textbf{Reference Id} & \texttt{BEC-233904DA} & \\
    \hline
    \textbf{Status} & Candidate ancestry & \\
    \hline
    \textbf{Ancestry} &  & \\
    \hline
    \textbf{Source} &  & \\
    \hline
    \textbf{Enabled By} &  & \\
    \hline
    \textbf{Supports} &  & \\
    \hline
    \textbf{Derived From} &  & \\
    \hline
    \textbf{Reviewed Status} &  & \\
    \hline
  \end{tabularx}
\end{table}

\begin{table}[htbp]
  \centering
  \footnotesize
  \begin{tabularx}{1\textwidth}{|l|l|X|}
    \hline
    \multicolumn{3}{c}{\textbf{Contained Capabilities}} \\
    \hline
    \textbf{Capability Name} & \textbf{Capability Id} & \textbf{Comment} \\
    \hline
    Canonical Semantics \& Schema Governance & \texttt{cap:bec-453a5b71} (at \ref{cap:bec-453a5b71}) & \\
    \hline
  \end{tabularx}
\end{table}

Intent. Information Design Governance is a candidate enabling capability within Shared Information and Coordination Enablement. It frames the reusable system-supported behaviours needed to support this healthcare context without prescribing a product, interface, technology or deployment boundary.

Boundary. This capability supports the business activity and does not inherit the authority or accountability of the Business Capability it enables. Its child Features are review candidates until ancestry and semantics are adjudicated.

\subsection{CT3 - Canonical Semantics \& Schema Governance}\label{cap:bec-453a5b71}
\vspace{-0.5em}
\begin{table}[htbp]
  \centering
  \footnotesize
  \begin{tabularx}{1\textwidth}{|l|l|X|}
    \hline
    \multicolumn{3}{c}{\textbf{Capability Metadata}} \\
    \hline
    \textbf{Name} & \textbf{Value} & \textbf{Comment and/or Description} \\
    \hline
    \textbf{Reference Id} & \texttt{BEC-453A5B71} & \\
    \hline
    \textbf{Status} & Candidate ancestry & \\
    \hline
    \textbf{Ancestry} &  & \\
    \hline
    \textbf{Source} &  & \\
    \hline
    \textbf{Enabled By} &  & \\
    \hline
    \textbf{Supports} &  & \\
    \hline
    \textbf{Derived From} &  & \\
    \hline
    \textbf{Reviewed Status} &  & \\
    \hline
  \end{tabularx}
\end{table}

\begin{table}[htbp]
  \centering
  \footnotesize
  \begin{tabularx}{1\textwidth}{|l|l|X|}
    \hline
    \multicolumn{3}{c}{\textbf{Contained Features}} \\
    \hline
    \textbf{Feature Name} & \textbf{Feature Id} & \textbf{Comment} \\
    \hline
    Information Standards \& Semantic Governance & \texttt{cap:bec-d57b8126} (at \ref{cap:bec-d57b8126}) & \\
    \hline
    Semantic Conformance Verification & \texttt{cap:bec-d881e00e} (at \ref{cap:bec-d881e00e}) & \\
    \hline
    Information Validation & \texttt{cap:bec-74f7b6a8} (at \ref{cap:bec-74f7b6a8}) & \\
    \hline
    Information Correlation & \texttt{cap:bec-3b311e84} (at \ref{cap:bec-3b311e84}) & \\
    \hline
    Information Retrieval & \texttt{cap:bec-e7f9d13f} (at \ref{cap:bec-e7f9d13f}) & \\
    \hline
    Information Provenance & \texttt{cap:bec-e383206b} (at \ref{cap:bec-e383206b}) & \\
    \hline
    Information Lifecycle Management & \texttt{cap:bec-915c9e8f} (at \ref{cap:bec-915c9e8f}) & \\
    \hline
  \end{tabularx}
\end{table}

Intent. Canonical Semantics \& Schema Governance is a candidate enabling capability within Information Design Governance. It frames the reusable system-supported behaviours needed to support this healthcare context without prescribing a product, interface, technology or deployment boundary.

Boundary. This capability supports the business activity and does not inherit the authority or accountability of the Business Capability it enables. Its child Features are review candidates until ancestry and semantics are adjudicated.

\subsubsection{Feature - Information Standards \& Semantic Governance}\label{cap:bec-d57b8126}
\vspace{-1em}
\begin{table}[htbp]
  \centering
  \footnotesize
  \begin{tabularx}{1\textwidth}{|l|l|X|}
    \hline
    \textbf{Metadata Name} & \textbf{Metadata Value} & \textbf{Description} \\
    \hline
    \textbf{Reference Id} & \texttt{BEC-D57B8126} & \\
    \hline
    \textbf{Status} & Existing Feature; candidate CT3 placement & \\
    \hline
    \textbf{Ancestry} &  & \\
    \hline
    \textbf{Source} & FEAT-ISE-27 & \\
    \hline
    \textbf{Enabled By} &  & \\
    \hline
    \textbf{Supports} &  & \\
    \hline
    \textbf{Derived From} &  & \\
    \hline
    \textbf{Reviewed Status} & Confirm existing Feature meaning and proposed CT3 ownership. & \\
    \hline
  \end{tabularx}
\end{table}

Intent. Information Standards \& Semantic Governance provides a distinct enabling behaviour in the context of Canonical Semantics \& Schema Governance. This is an existing named Feature; its placement under the candidate CT3 remains subject to review.

\noindent\textbf{Feature boundary.} This behaviour is scoped to Canonical Semantics \& Schema Governance; an identically named Feature elsewhere may carry different business meaning.\par
\featuretable{
Function & \texttt{FN-B8CAF869} & Govern Information Semantics & Catalogue \\
Service & \texttt{SV-AF43A3B2} & Semantic Model Definition Service & Linked to preceding Function in source \\
}
\subsubsection{Feature - Semantic Conformance Verification}\label{cap:bec-d881e00e}
\vspace{-1em}
\begin{table}[htbp]
  \centering
  \footnotesize
  \begin{tabularx}{1\textwidth}{|l|l|X|}
    \hline
    \textbf{Metadata Name} & \textbf{Metadata Value} & \textbf{Description} \\
    \hline
    \textbf{Reference Id} & \texttt{BEC-D881E00E} & \\
    \hline
    \textbf{Status} & Existing Feature; candidate CT3 placement & \\
    \hline
    \textbf{Ancestry} &  & \\
    \hline
    \textbf{Source} & FEAT-ISE-28 & \\
    \hline
    \textbf{Enabled By} &  & \\
    \hline
    \textbf{Supports} &  & \\
    \hline
    \textbf{Derived From} &  & \\
    \hline
    \textbf{Reviewed Status} & Confirm existing Feature meaning and proposed CT3 ownership. & \\
    \hline
  \end{tabularx}
\end{table}

Intent. Semantic Conformance Verification provides a distinct enabling behaviour in the context of Canonical Semantics \& Schema Governance. This is an existing named Feature; its placement under the candidate CT3 remains subject to review.

\noindent\textbf{Feature boundary.} This behaviour is scoped to Canonical Semantics \& Schema Governance; an identically named Feature elsewhere may carry different business meaning.\par
\featuretable{
Function & \texttt{FN-0C04B6EB} & Verify Semantic Conformance & Catalogue \\
Service & \texttt{SV-F964D1FC} & Semantic Conformance Verification Service & Linked to preceding Function in source \\
}
\subsubsection{Feature - Information Validation}\label{cap:bec-74f7b6a8}
\vspace{-1em}
\begin{table}[htbp]
  \centering
  \footnotesize
  \begin{tabularx}{1\textwidth}{|l|l|X|}
    \hline
    \textbf{Metadata Name} & \textbf{Metadata Value} & \textbf{Description} \\
    \hline
    \textbf{Reference Id} & \texttt{BEC-74F7B6A8} & \\
    \hline
    \textbf{Status} & Generated candidate & \\
    \hline
    \textbf{Ancestry} &  & \\
    \hline
    \textbf{Source} & Semantic-pattern candidate & \\
    \hline
    \textbf{Enabled By} &  & \\
    \hline
    \textbf{Supports} &  & \\
    \hline
    \textbf{Derived From} &  & \\
    \hline
    \textbf{Reviewed Status} & Generated candidate: keep / modify / relocate / reject; verify whether independently testable. & \\
    \hline
  \end{tabularx}
\end{table}

Intent. Information Validation is proposed as a discrete, testable enabling behaviour for Canonical Semantics \& Schema Governance. It is included for architectural review, not as an established implementation obligation.

\noindent\textbf{Feature boundary.} This behaviour is scoped to Canonical Semantics \& Schema Governance; an identically named Feature elsewhere may carry different business meaning.\par
\featuretable{
--- & --- & \emph{Not yet captured} & \emph{Complete during review} \\
}
\subsubsection{Feature - Information Correlation}\label{cap:bec-3b311e84}
\vspace{-1em}
\begin{table}[htbp]
  \centering
  \footnotesize
  \begin{tabularx}{1\textwidth}{|l|l|X|}
    \hline
    \textbf{Metadata Name} & \textbf{Metadata Value} & \textbf{Description} \\
    \hline
    \textbf{Reference Id} & \texttt{BEC-3B311E84} & \\
    \hline
    \textbf{Status} & Generated candidate & \\
    \hline
    \textbf{Ancestry} &  & \\
    \hline
    \textbf{Source} & Semantic-pattern candidate & \\
    \hline
    \textbf{Enabled By} &  & \\
    \hline
    \textbf{Supports} &  & \\
    \hline
    \textbf{Derived From} &  & \\
    \hline
    \textbf{Reviewed Status} & Generated candidate: keep / modify / relocate / reject; verify whether independently testable. & \\
    \hline
  \end{tabularx}
\end{table}

Intent. Information Correlation is proposed as a discrete, testable enabling behaviour for Canonical Semantics \& Schema Governance. It is included for architectural review, not as an established implementation obligation.

\noindent\textbf{Feature boundary.} This behaviour is scoped to Canonical Semantics \& Schema Governance; an identically named Feature elsewhere may carry different business meaning.\par
\featuretable{
--- & --- & \emph{Not yet captured} & \emph{Complete during review} \\
}
\subsubsection{Feature - Information Retrieval}\label{cap:bec-e7f9d13f}
\vspace{-1em}
\begin{table}[htbp]
  \centering
  \footnotesize
  \begin{tabularx}{1\textwidth}{|l|l|X|}
    \hline
    \textbf{Metadata Name} & \textbf{Metadata Value} & \textbf{Description} \\
    \hline
    \textbf{Reference Id} & \texttt{BEC-E7F9D13F} & \\
    \hline
    \textbf{Status} & Generated candidate & \\
    \hline
    \textbf{Ancestry} &  & \\
    \hline
    \textbf{Source} & Semantic-pattern candidate & \\
    \hline
    \textbf{Enabled By} &  & \\
    \hline
    \textbf{Supports} &  & \\
    \hline
    \textbf{Derived From} &  & \\
    \hline
    \textbf{Reviewed Status} & Generated candidate: keep / modify / relocate / reject; verify whether independently testable. & \\
    \hline
  \end{tabularx}
\end{table}

Intent. Information Retrieval is proposed as a discrete, testable enabling behaviour for Canonical Semantics \& Schema Governance. It is included for architectural review, not as an established implementation obligation.

\noindent\textbf{Feature boundary.} This behaviour is scoped to Canonical Semantics \& Schema Governance; an identically named Feature elsewhere may carry different business meaning.\par
\featuretable{
--- & --- & \emph{Not yet captured} & \emph{Complete during review} \\
}
\subsubsection{Feature - Information Provenance}\label{cap:bec-e383206b}
\vspace{-1em}
\begin{table}[htbp]
  \centering
  \footnotesize
  \begin{tabularx}{1\textwidth}{|l|l|X|}
    \hline
    \textbf{Metadata Name} & \textbf{Metadata Value} & \textbf{Description} \\
    \hline
    \textbf{Reference Id} & \texttt{BEC-E383206B} & \\
    \hline
    \textbf{Status} & Generated candidate & \\
    \hline
    \textbf{Ancestry} &  & \\
    \hline
    \textbf{Source} & Semantic-pattern candidate & \\
    \hline
    \textbf{Enabled By} &  & \\
    \hline
    \textbf{Supports} &  & \\
    \hline
    \textbf{Derived From} &  & \\
    \hline
    \textbf{Reviewed Status} & Generated candidate: keep / modify / relocate / reject; verify whether independently testable. & \\
    \hline
  \end{tabularx}
\end{table}

Intent. Information Provenance is proposed as a discrete, testable enabling behaviour for Canonical Semantics \& Schema Governance. It is included for architectural review, not as an established implementation obligation.

\noindent\textbf{Feature boundary.} This behaviour is scoped to Canonical Semantics \& Schema Governance; an identically named Feature elsewhere may carry different business meaning.\par
\featuretable{
--- & --- & \emph{Not yet captured} & \emph{Complete during review} \\
}
\subsubsection{Feature - Information Lifecycle Management}\label{cap:bec-915c9e8f}
\vspace{-1em}
\begin{table}[htbp]
  \centering
  \footnotesize
  \begin{tabularx}{1\textwidth}{|l|l|X|}
    \hline
    \textbf{Metadata Name} & \textbf{Metadata Value} & \textbf{Description} \\
    \hline
    \textbf{Reference Id} & \texttt{BEC-915C9E8F} & \\
    \hline
    \textbf{Status} & Generated candidate & \\
    \hline
    \textbf{Ancestry} &  & \\
    \hline
    \textbf{Source} & Semantic-pattern candidate & \\
    \hline
    \textbf{Enabled By} &  & \\
    \hline
    \textbf{Supports} &  & \\
    \hline
    \textbf{Derived From} &  & \\
    \hline
    \textbf{Reviewed Status} & Generated candidate: keep / modify / relocate / reject; verify whether independently testable. & \\
    \hline
  \end{tabularx}
\end{table}

Intent. Information Lifecycle Management is proposed as a discrete, testable enabling behaviour for Canonical Semantics \& Schema Governance. It is included for architectural review, not as an established implementation obligation.

\noindent\textbf{Feature boundary.} This behaviour is scoped to Canonical Semantics \& Schema Governance; an identically named Feature elsewhere may carry different business meaning.\par
\featuretable{
--- & --- & \emph{Not yet captured} & \emph{Complete during review} \\
}

\part{Enterprise Architecture Capabilities and Features}
\chapter*{Enterprise Derivation Method}\addcontentsline{toc}{chapter}{Enterprise Derivation Method}
The derived Enterprise hierarchy groups semantic recurrence across Business Enabling CT3 contexts. Feature evidence remains in the BEC part. Candidate enterprise CT1/CT2 labels are provisional; candidate CT3s are families of related reusable behaviour. The existing EC-01 to EC-14 material is added as a separate coexistence branch, even where the names or responsibilities overlap.

\chapter{CT1 - Entity and Identity}\label{cap:ec-bbad8e2c}
\noindent\textbf{Reference ID:} \texttt{EC-BBAD8E2C}\quad \textbf{Status:} Derived candidate\par
Intent. Entity and Identity is a candidate horizontally reusable architectural capability associated with BEC recurrence. It is framed around recurring enabling behaviour, not a mandatory centralised service or a particular technology.

Boundary. A common architectural capability does not remove the originating Features from their healthcare contexts. Similar wording does not establish equivalence or identical implementation.

\section{CT2 - Identity and Correlation}\label{cap:ec-c0a97a83}
\vspace{-0.5em}
\begin{table}[htbp]
  \centering
  \footnotesize
  \begin{tabularx}{1\textwidth}{|l|l|X|}
    \hline
    \multicolumn{3}{c}{\textbf{Capability Metadata}} \\
    \hline
    \textbf{Name} & \textbf{Value} & \textbf{Comment and/or Description} \\
    \hline
    \textbf{Reference Id} & \texttt{EC-C0A97A83} & \\
    \hline
    \textbf{Status} & Derived candidate & \\
    \hline
    \textbf{Ancestry} &  & \\
    \hline
    \textbf{Source} &  & \\
    \hline
    \textbf{Enabled By} &  & \\
    \hline
    \textbf{Supports} &  & \\
    \hline
    \textbf{Derived From} &  & \\
    \hline
    \textbf{Reviewed Status} &  & \\
    \hline
  \end{tabularx}
\end{table}

\begin{table}[htbp]
  \centering
  \footnotesize
  \begin{tabularx}{1\textwidth}{|l|l|X|}
    \hline
    \multicolumn{3}{c}{\textbf{Contained Capabilities}} \\
    \hline
    \textbf{Capability Name} & \textbf{Capability Id} & \textbf{Comment} \\
    \hline
    Identity and Correlation Enablement & \texttt{cap:ec-461b8f05} (at \ref{cap:ec-461b8f05}) & \\
    \hline
  \end{tabularx}
\end{table}

Intent. Identity and Correlation is a candidate horizontally reusable architectural capability associated with Entity and Identity. It is framed around recurring enabling behaviour, not a mandatory centralised service or a particular technology.

Boundary. A common architectural capability does not remove the originating Features from their healthcare contexts. Similar wording does not establish equivalence or identical implementation.

\subsection{CT3 - Identity and Correlation Enablement}\label{cap:ec-461b8f05}
\vspace{-0.5em}
\begin{table}[htbp]
  \centering
  \footnotesize
  \begin{tabularx}{1\textwidth}{|l|l|X|}
    \hline
    \multicolumn{3}{c}{\textbf{Capability Metadata}} \\
    \hline
    \textbf{Name} & \textbf{Value} & \textbf{Comment and/or Description} \\
    \hline
    \textbf{Reference Id} & \texttt{EC-461B8F05} & \\
    \hline
    \textbf{Status} & Derived candidate & \\
    \hline
    \textbf{Ancestry} &  & \\
    \hline
    \textbf{Source} &  & \\
    \hline
    \textbf{Enabled By} &  & \\
    \hline
    \textbf{Supports} &  & \\
    \hline
    \textbf{Derived From} &  & \\
    \hline
    \textbf{Reviewed Status} &  & \\
    \hline
  \end{tabularx}
\end{table}

\begin{table}[htbp]
  \centering
  \footnotesize
  \begin{tabularx}{1\textwidth}{|l|l|X|}
    \hline
    \multicolumn{3}{c}{\textbf{Contained Features}} \\
    \hline
    \textbf{Feature Name} & \textbf{Feature Id} & \textbf{Comment} \\
    \hline
    Identity Correlation & \texttt{cap:ec-1a42afef} (at \ref{cap:ec-1a42afef}) & \\
    \hline
    Identity Management & \texttt{cap:ec-00688388} (at \ref{cap:ec-00688388}) & \\
    \hline
    Device Identity Management & \texttt{cap:ec-13bdd6ae} (at \ref{cap:ec-13bdd6ae}) & \\
    \hline
    Subject Identity Correlation & \texttt{cap:ec-157baa94} (at \ref{cap:ec-157baa94}) & \\
    \hline
    Information Correlation & \texttt{cap:ec-47108aac} (at \ref{cap:ec-47108aac}) & \\
    \hline
    Diagnostic Request Correlation & \texttt{cap:ec-ea771b95} (at \ref{cap:ec-ea771b95}) & \\
    \hline
    Location Identity Management & \texttt{cap:ec-039334fa} (at \ref{cap:ec-039334fa}) & \\
    \hline
    Medication Identity Resolution & \texttt{cap:ec-5fc97327} (at \ref{cap:ec-5fc97327}) & \\
    \hline
    Medication Information Correlation & \texttt{cap:ec-d444a2c1} (at \ref{cap:ec-d444a2c1}) & \\
    \hline
    Request Status Correlation & \texttt{cap:ec-d418a306} (at \ref{cap:ec-d418a306}) & \\
    \hline
  \end{tabularx}
\end{table}

Intent. Identity and Correlation Enablement is a candidate horizontally reusable architectural capability associated with Identity and Correlation. It is framed around recurring enabling behaviour, not a mandatory centralised service or a particular technology.

Boundary. A common architectural capability does not remove the originating Features from their healthcare contexts. Similar wording does not establish equivalence or identical implementation.

\subsubsection{Feature - Identity Correlation}\label{cap:ec-1a42afef}
\vspace{-1em}
\begin{table}[htbp]
  \centering
  \footnotesize
  \begin{tabularx}{1\textwidth}{|l|l|X|}
    \hline
    \textbf{Metadata Name} & \textbf{Metadata Value} & \textbf{Description} \\
    \hline
    \textbf{Reference Id} & \texttt{EC-1A42AFEF} & \\
    \hline
    \textbf{Status} & Derived candidate & \\
    \hline
    \textbf{Ancestry} &  & \\
    \hline
    \textbf{Source} &  & \\
    \hline
    \textbf{Enabled By} &  & \\
    \hline
    \textbf{Supports} &  & \\
    \hline
    \textbf{Derived From} &  & \\
    \hline
    \textbf{Reviewed Status} & Inspect recurrence for true semantic equivalence; retain, split, rename or reject. & \\
    \hline
  \end{tabularx}
\end{table}

This candidate enterprise Feature abstracts a behaviour recurring in 7 distinct BEC CT3 contexts. The source-specific BEC Features remain in place; this entry does not assert a common software service or an identical clinical meaning.

Illustrative originating CT3 contexts: Healthcare Organisation, Organisation Contact, Organisational Structure, Person Identity, Practitioner, Practitioner Role, Product Formulary \& Catalogue.

\featuretable{--- & --- & \emph{Not yet captured} & \emph{Enterprise decomposition subject to review} \\}
\subsubsection{Feature - Identity Management}\label{cap:ec-00688388}
\vspace{-1em}
\begin{table}[htbp]
  \centering
  \footnotesize
  \begin{tabularx}{1\textwidth}{|l|l|X|}
    \hline
    \textbf{Metadata Name} & \textbf{Metadata Value} & \textbf{Description} \\
    \hline
    \textbf{Reference Id} & \texttt{EC-00688388} & \\
    \hline
    \textbf{Status} & Derived candidate & \\
    \hline
    \textbf{Ancestry} &  & \\
    \hline
    \textbf{Source} &  & \\
    \hline
    \textbf{Enabled By} &  & \\
    \hline
    \textbf{Supports} &  & \\
    \hline
    \textbf{Derived From} &  & \\
    \hline
    \textbf{Reviewed Status} & Inspect recurrence for true semantic equivalence; retain, split, rename or reject. & \\
    \hline
  \end{tabularx}
\end{table}

This candidate enterprise Feature abstracts a behaviour recurring in 6 distinct BEC CT3 contexts. The source-specific BEC Features remain in place; this entry does not assert a common software service or an identical clinical meaning.

Illustrative originating CT3 contexts: Healthcare Organisation, Organisation Contact, Organisational Structure, Person Identity, Practitioner, Practitioner Role.

\featuretable{--- & --- & \emph{Not yet captured} & \emph{Enterprise decomposition subject to review} \\}
\subsubsection{Feature - Device Identity Management}\label{cap:ec-13bdd6ae}
\vspace{-1em}
\begin{table}[htbp]
  \centering
  \footnotesize
  \begin{tabularx}{1\textwidth}{|l|l|X|}
    \hline
    \textbf{Metadata Name} & \textbf{Metadata Value} & \textbf{Description} \\
    \hline
    \textbf{Reference Id} & \texttt{EC-13BDD6AE} & \\
    \hline
    \textbf{Status} & Derived candidate & \\
    \hline
    \textbf{Ancestry} &  & \\
    \hline
    \textbf{Source} &  & \\
    \hline
    \textbf{Enabled By} &  & \\
    \hline
    \textbf{Supports} &  & \\
    \hline
    \textbf{Derived From} &  & \\
    \hline
    \textbf{Reviewed Status} & Inspect recurrence for true semantic equivalence; retain, split, rename or reject. & \\
    \hline
  \end{tabularx}
\end{table}

This candidate enterprise Feature abstracts a behaviour recurring in 4 distinct BEC CT3 contexts. The source-specific BEC Features remain in place; this entry does not assert a common software service or an identical clinical meaning.

Illustrative originating CT3 contexts: Device Association, Device Connectivity, Device Definition, Device Instance.

\featuretable{--- & --- & \emph{Not yet captured} & \emph{Enterprise decomposition subject to review} \\}
\subsubsection{Feature - Subject Identity Correlation}\label{cap:ec-157baa94}
\vspace{-1em}
\begin{table}[htbp]
  \centering
  \footnotesize
  \begin{tabularx}{1\textwidth}{|l|l|X|}
    \hline
    \textbf{Metadata Name} & \textbf{Metadata Value} & \textbf{Description} \\
    \hline
    \textbf{Reference Id} & \texttt{EC-157BAA94} & \\
    \hline
    \textbf{Status} & Derived candidate & \\
    \hline
    \textbf{Ancestry} &  & \\
    \hline
    \textbf{Source} &  & \\
    \hline
    \textbf{Enabled By} &  & \\
    \hline
    \textbf{Supports} &  & \\
    \hline
    \textbf{Derived From} &  & \\
    \hline
    \textbf{Reviewed Status} & Inspect recurrence for true semantic equivalence; retain, split, rename or reject. & \\
    \hline
  \end{tabularx}
\end{table}

This candidate enterprise Feature abstracts a behaviour recurring in 4 distinct BEC CT3 contexts. The source-specific BEC Features remain in place; this entry does not assert a common software service or an identical clinical meaning.

Illustrative originating CT3 contexts: Multidisciplinary Inpatient Care Coordination, Outpatient Flow Enablement, Patient Transfer Coordination, Remote Patient Monitoring.

\featuretable{--- & --- & \emph{Not yet captured} & \emph{Enterprise decomposition subject to review} \\}
\subsubsection{Feature - Information Correlation}\label{cap:ec-47108aac}
\vspace{-1em}
\begin{table}[htbp]
  \centering
  \footnotesize
  \begin{tabularx}{1\textwidth}{|l|l|X|}
    \hline
    \textbf{Metadata Name} & \textbf{Metadata Value} & \textbf{Description} \\
    \hline
    \textbf{Reference Id} & \texttt{EC-47108AAC} & \\
    \hline
    \textbf{Status} & Derived candidate & \\
    \hline
    \textbf{Ancestry} &  & \\
    \hline
    \textbf{Source} &  & \\
    \hline
    \textbf{Enabled By} &  & \\
    \hline
    \textbf{Supports} &  & \\
    \hline
    \textbf{Derived From} &  & \\
    \hline
    \textbf{Reviewed Status} & Inspect recurrence for true semantic equivalence; retain, split, rename or reject. & \\
    \hline
  \end{tabularx}
\end{table}

This candidate enterprise Feature abstracts a behaviour recurring in 3 distinct BEC CT3 contexts. The source-specific BEC Features remain in place; this entry does not assert a common software service or an identical clinical meaning.

Illustrative originating CT3 contexts: Canonical Semantics \& Schema Governance, Multi-Model Information Exchange, Terminology \& Reference Mapping.

\featuretable{--- & --- & \emph{Not yet captured} & \emph{Enterprise decomposition subject to review} \\}
\subsubsection{Feature - Diagnostic Request Correlation}\label{cap:ec-ea771b95}
\vspace{-1em}
\begin{table}[htbp]
  \centering
  \footnotesize
  \begin{tabularx}{1\textwidth}{|l|l|X|}
    \hline
    \textbf{Metadata Name} & \textbf{Metadata Value} & \textbf{Description} \\
    \hline
    \textbf{Reference Id} & \texttt{EC-EA771B95} & \\
    \hline
    \textbf{Status} & Derived candidate & \\
    \hline
    \textbf{Ancestry} &  & \\
    \hline
    \textbf{Source} &  & \\
    \hline
    \textbf{Enabled By} &  & \\
    \hline
    \textbf{Supports} &  & \\
    \hline
    \textbf{Derived From} &  & \\
    \hline
    \textbf{Reviewed Status} & Inspect recurrence for true semantic equivalence; retain, split, rename or reject. & \\
    \hline
  \end{tabularx}
\end{table}

This candidate enterprise Feature abstracts a behaviour recurring in 2 distinct BEC CT3 contexts. The source-specific BEC Features remain in place; this entry does not assert a common software service or an identical clinical meaning.

Illustrative originating CT3 contexts: Diagnostic Service Integration, Diagnostic Workflow Management.

\featuretable{--- & --- & \emph{Not yet captured} & \emph{Enterprise decomposition subject to review} \\}
\subsubsection{Feature - Location Identity Management}\label{cap:ec-039334fa}
\vspace{-1em}
\begin{table}[htbp]
  \centering
  \footnotesize
  \begin{tabularx}{1\textwidth}{|l|l|X|}
    \hline
    \textbf{Metadata Name} & \textbf{Metadata Value} & \textbf{Description} \\
    \hline
    \textbf{Reference Id} & \texttt{EC-039334FA} & \\
    \hline
    \textbf{Status} & Derived candidate & \\
    \hline
    \textbf{Ancestry} &  & \\
    \hline
    \textbf{Source} &  & \\
    \hline
    \textbf{Enabled By} &  & \\
    \hline
    \textbf{Supports} &  & \\
    \hline
    \textbf{Derived From} &  & \\
    \hline
    \textbf{Reviewed Status} & Inspect recurrence for true semantic equivalence; retain, split, rename or reject. & \\
    \hline
  \end{tabularx}
\end{table}

This candidate enterprise Feature abstracts a behaviour recurring in 2 distinct BEC CT3 contexts. The source-specific BEC Features remain in place; this entry does not assert a common software service or an identical clinical meaning.

Illustrative originating CT3 contexts: Healthcare Location, Location Structure.

\featuretable{--- & --- & \emph{Not yet captured} & \emph{Enterprise decomposition subject to review} \\}
\subsubsection{Feature - Medication Identity Resolution}\label{cap:ec-5fc97327}
\vspace{-1em}
\begin{table}[htbp]
  \centering
  \footnotesize
  \begin{tabularx}{1\textwidth}{|l|l|X|}
    \hline
    \textbf{Metadata Name} & \textbf{Metadata Value} & \textbf{Description} \\
    \hline
    \textbf{Reference Id} & \texttt{EC-5FC97327} & \\
    \hline
    \textbf{Status} & Derived candidate & \\
    \hline
    \textbf{Ancestry} &  & \\
    \hline
    \textbf{Source} &  & \\
    \hline
    \textbf{Enabled By} &  & \\
    \hline
    \textbf{Supports} &  & \\
    \hline
    \textbf{Derived From} &  & \\
    \hline
    \textbf{Reviewed Status} & Inspect recurrence for true semantic equivalence; retain, split, rename or reject. & \\
    \hline
  \end{tabularx}
\end{table}

This candidate enterprise Feature abstracts a behaviour recurring in 2 distinct BEC CT3 contexts. The source-specific BEC Features remain in place; this entry does not assert a common software service or an identical clinical meaning.

Illustrative originating CT3 contexts: Medication Order \& Dispense Management, Medication Therapy Provision and Monitoring.

\featuretable{--- & --- & \emph{Not yet captured} & \emph{Enterprise decomposition subject to review} \\}
\subsubsection{Feature - Medication Information Correlation}\label{cap:ec-d444a2c1}
\vspace{-1em}
\begin{table}[htbp]
  \centering
  \footnotesize
  \begin{tabularx}{1\textwidth}{|l|l|X|}
    \hline
    \textbf{Metadata Name} & \textbf{Metadata Value} & \textbf{Description} \\
    \hline
    \textbf{Reference Id} & \texttt{EC-D444A2C1} & \\
    \hline
    \textbf{Status} & Derived candidate & \\
    \hline
    \textbf{Ancestry} &  & \\
    \hline
    \textbf{Source} &  & \\
    \hline
    \textbf{Enabled By} &  & \\
    \hline
    \textbf{Supports} &  & \\
    \hline
    \textbf{Derived From} &  & \\
    \hline
    \textbf{Reviewed Status} & Inspect recurrence for true semantic equivalence; retain, split, rename or reject. & \\
    \hline
  \end{tabularx}
\end{table}

This candidate enterprise Feature abstracts a behaviour recurring in 2 distinct BEC CT3 contexts. The source-specific BEC Features remain in place; this entry does not assert a common software service or an identical clinical meaning.

Illustrative originating CT3 contexts: Medication Order \& Dispense Management, Medication Therapy Provision and Monitoring.

\featuretable{--- & --- & \emph{Not yet captured} & \emph{Enterprise decomposition subject to review} \\}
\subsubsection{Feature - Request Status Correlation}\label{cap:ec-d418a306}
\vspace{-1em}
\begin{table}[htbp]
  \centering
  \footnotesize
  \begin{tabularx}{1\textwidth}{|l|l|X|}
    \hline
    \textbf{Metadata Name} & \textbf{Metadata Value} & \textbf{Description} \\
    \hline
    \textbf{Reference Id} & \texttt{EC-D418A306} & \\
    \hline
    \textbf{Status} & Derived candidate & \\
    \hline
    \textbf{Ancestry} &  & \\
    \hline
    \textbf{Source} &  & \\
    \hline
    \textbf{Enabled By} &  & \\
    \hline
    \textbf{Supports} &  & \\
    \hline
    \textbf{Derived From} &  & \\
    \hline
    \textbf{Reviewed Status} & Inspect recurrence for true semantic equivalence; retain, split, rename or reject. & \\
    \hline
  \end{tabularx}
\end{table}

This candidate enterprise Feature abstracts a behaviour recurring in 2 distinct BEC CT3 contexts. The source-specific BEC Features remain in place; this entry does not assert a common software service or an identical clinical meaning.

Illustrative originating CT3 contexts: Closed-Loop Order Progression, Medication Order \& Dispense Management.

\featuretable{--- & --- & \emph{Not yet captured} & \emph{Enterprise decomposition subject to review} \\}

\section{CT2 - Relationship and Structure}\label{cap:ec-bac316e4}
\vspace{-0.5em}
\begin{table}[htbp]
  \centering
  \footnotesize
  \begin{tabularx}{1\textwidth}{|l|l|X|}
    \hline
    \multicolumn{3}{c}{\textbf{Capability Metadata}} \\
    \hline
    \textbf{Name} & \textbf{Value} & \textbf{Comment and/or Description} \\
    \hline
    \textbf{Reference Id} & \texttt{EC-BAC316E4} & \\
    \hline
    \textbf{Status} & Derived candidate & \\
    \hline
    \textbf{Ancestry} &  & \\
    \hline
    \textbf{Source} &  & \\
    \hline
    \textbf{Enabled By} &  & \\
    \hline
    \textbf{Supports} &  & \\
    \hline
    \textbf{Derived From} &  & \\
    \hline
    \textbf{Reviewed Status} &  & \\
    \hline
  \end{tabularx}
\end{table}

\begin{table}[htbp]
  \centering
  \footnotesize
  \begin{tabularx}{1\textwidth}{|l|l|X|}
    \hline
    \multicolumn{3}{c}{\textbf{Contained Capabilities}} \\
    \hline
    \textbf{Capability Name} & \textbf{Capability Id} & \textbf{Comment} \\
    \hline
    Relationship and Structure Enablement & \texttt{cap:ec-f50f53e6} (at \ref{cap:ec-f50f53e6}) & \\
    \hline
  \end{tabularx}
\end{table}

Intent. Relationship and Structure is a candidate horizontally reusable architectural capability associated with Entity and Identity. It is framed around recurring enabling behaviour, not a mandatory centralised service or a particular technology.

Boundary. A common architectural capability does not remove the originating Features from their healthcare contexts. Similar wording does not establish equivalence or identical implementation.

\subsection{CT3 - Relationship and Structure Enablement}\label{cap:ec-f50f53e6}
\vspace{-0.5em}
\begin{table}[htbp]
  \centering
  \footnotesize
  \begin{tabularx}{1\textwidth}{|l|l|X|}
    \hline
    \multicolumn{3}{c}{\textbf{Capability Metadata}} \\
    \hline
    \textbf{Name} & \textbf{Value} & \textbf{Comment and/or Description} \\
    \hline
    \textbf{Reference Id} & \texttt{EC-F50F53E6} & \\
    \hline
    \textbf{Status} & Derived candidate & \\
    \hline
    \textbf{Ancestry} &  & \\
    \hline
    \textbf{Source} &  & \\
    \hline
    \textbf{Enabled By} &  & \\
    \hline
    \textbf{Supports} &  & \\
    \hline
    \textbf{Derived From} &  & \\
    \hline
    \textbf{Reviewed Status} &  & \\
    \hline
  \end{tabularx}
\end{table}

\begin{table}[htbp]
  \centering
  \footnotesize
  \begin{tabularx}{1\textwidth}{|l|l|X|}
    \hline
    \multicolumn{3}{c}{\textbf{Contained Features}} \\
    \hline
    \textbf{Feature Name} & \textbf{Feature Id} & \textbf{Comment} \\
    \hline
    Device Relationship Management & \texttt{cap:ec-81917406} (at \ref{cap:ec-81917406}) & \\
    \hline
    Subject Relationship Management & \texttt{cap:ec-20119d58} (at \ref{cap:ec-20119d58}) & \\
    \hline
    Organisational Relationship Management & \texttt{cap:ec-0b8c4edc} (at \ref{cap:ec-0b8c4edc}) & \\
    \hline
    Organisational Structure Management & \texttt{cap:ec-49166b2a} (at \ref{cap:ec-49166b2a}) & \\
    \hline
    Diagnostic Result Association & \texttt{cap:ec-7a2df974} (at \ref{cap:ec-7a2df974}) & \\
    \hline
    Evidence Association & \texttt{cap:ec-279fe89f} (at \ref{cap:ec-279fe89f}) & \\
    \hline
    Location Hierarchy Management & \texttt{cap:ec-a179d91b} (at \ref{cap:ec-a179d91b}) & \\
    \hline
    Professional Relationship Management & \texttt{cap:ec-366afc9b} (at \ref{cap:ec-366afc9b}) & \\
    \hline
    Qualification Association & \texttt{cap:ec-e2f0cb40} (at \ref{cap:ec-e2f0cb40}) & \\
    \hline
    Relationship Management & \texttt{cap:ec-f160e33e} (at \ref{cap:ec-f160e33e}) & \\
    \hline
  \end{tabularx}
\end{table}

Intent. Relationship and Structure Enablement is a candidate horizontally reusable architectural capability associated with Relationship and Structure. It is framed around recurring enabling behaviour, not a mandatory centralised service or a particular technology.

Boundary. A common architectural capability does not remove the originating Features from their healthcare contexts. Similar wording does not establish equivalence or identical implementation.

\subsubsection{Feature - Device Relationship Management}\label{cap:ec-81917406}
\vspace{-1em}
\begin{table}[htbp]
  \centering
  \footnotesize
  \begin{tabularx}{1\textwidth}{|l|l|X|}
    \hline
    \textbf{Metadata Name} & \textbf{Metadata Value} & \textbf{Description} \\
    \hline
    \textbf{Reference Id} & \texttt{EC-81917406} & \\
    \hline
    \textbf{Status} & Derived candidate & \\
    \hline
    \textbf{Ancestry} &  & \\
    \hline
    \textbf{Source} &  & \\
    \hline
    \textbf{Enabled By} &  & \\
    \hline
    \textbf{Supports} &  & \\
    \hline
    \textbf{Derived From} &  & \\
    \hline
    \textbf{Reviewed Status} & Inspect recurrence for true semantic equivalence; retain, split, rename or reject. & \\
    \hline
  \end{tabularx}
\end{table}

This candidate enterprise Feature abstracts a behaviour recurring in 4 distinct BEC CT3 contexts. The source-specific BEC Features remain in place; this entry does not assert a common software service or an identical clinical meaning.

Illustrative originating CT3 contexts: Device Association, Device Connectivity, Device Definition, Device Instance.

\featuretable{--- & --- & \emph{Not yet captured} & \emph{Enterprise decomposition subject to review} \\}
\subsubsection{Feature - Subject Relationship Management}\label{cap:ec-20119d58}
\vspace{-1em}
\begin{table}[htbp]
  \centering
  \footnotesize
  \begin{tabularx}{1\textwidth}{|l|l|X|}
    \hline
    \textbf{Metadata Name} & \textbf{Metadata Value} & \textbf{Description} \\
    \hline
    \textbf{Reference Id} & \texttt{EC-20119D58} & \\
    \hline
    \textbf{Status} & Derived candidate & \\
    \hline
    \textbf{Ancestry} &  & \\
    \hline
    \textbf{Source} &  & \\
    \hline
    \textbf{Enabled By} &  & \\
    \hline
    \textbf{Supports} &  & \\
    \hline
    \textbf{Derived From} &  & \\
    \hline
    \textbf{Reviewed Status} & Inspect recurrence for true semantic equivalence; retain, split, rename or reject. & \\
    \hline
  \end{tabularx}
\end{table}

This candidate enterprise Feature abstracts a behaviour recurring in 4 distinct BEC CT3 contexts. The source-specific BEC Features remain in place; this entry does not assert a common software service or an identical clinical meaning.

Illustrative originating CT3 contexts: Multidisciplinary Inpatient Care Coordination, Outpatient Flow Enablement, Patient Transfer Coordination, Remote Patient Monitoring.

\featuretable{--- & --- & \emph{Not yet captured} & \emph{Enterprise decomposition subject to review} \\}
\subsubsection{Feature - Organisational Relationship Management}\label{cap:ec-0b8c4edc}
\vspace{-1em}
\begin{table}[htbp]
  \centering
  \footnotesize
  \begin{tabularx}{1\textwidth}{|l|l|X|}
    \hline
    \textbf{Metadata Name} & \textbf{Metadata Value} & \textbf{Description} \\
    \hline
    \textbf{Reference Id} & \texttt{EC-0B8C4EDC} & \\
    \hline
    \textbf{Status} & Derived candidate & \\
    \hline
    \textbf{Ancestry} &  & \\
    \hline
    \textbf{Source} &  & \\
    \hline
    \textbf{Enabled By} &  & \\
    \hline
    \textbf{Supports} &  & \\
    \hline
    \textbf{Derived From} &  & \\
    \hline
    \textbf{Reviewed Status} & Inspect recurrence for true semantic equivalence; retain, split, rename or reject. & \\
    \hline
  \end{tabularx}
\end{table}

This candidate enterprise Feature abstracts a behaviour recurring in 3 distinct BEC CT3 contexts. The source-specific BEC Features remain in place; this entry does not assert a common software service or an identical clinical meaning.

Illustrative originating CT3 contexts: Healthcare Organisation, Organisation Contact, Organisational Structure.

\featuretable{--- & --- & \emph{Not yet captured} & \emph{Enterprise decomposition subject to review} \\}
\subsubsection{Feature - Organisational Structure Management}\label{cap:ec-49166b2a}
\vspace{-1em}
\begin{table}[htbp]
  \centering
  \footnotesize
  \begin{tabularx}{1\textwidth}{|l|l|X|}
    \hline
    \textbf{Metadata Name} & \textbf{Metadata Value} & \textbf{Description} \\
    \hline
    \textbf{Reference Id} & \texttt{EC-49166B2A} & \\
    \hline
    \textbf{Status} & Derived candidate & \\
    \hline
    \textbf{Ancestry} &  & \\
    \hline
    \textbf{Source} &  & \\
    \hline
    \textbf{Enabled By} &  & \\
    \hline
    \textbf{Supports} &  & \\
    \hline
    \textbf{Derived From} &  & \\
    \hline
    \textbf{Reviewed Status} & Inspect recurrence for true semantic equivalence; retain, split, rename or reject. & \\
    \hline
  \end{tabularx}
\end{table}

This candidate enterprise Feature abstracts a behaviour recurring in 3 distinct BEC CT3 contexts. The source-specific BEC Features remain in place; this entry does not assert a common software service or an identical clinical meaning.

Illustrative originating CT3 contexts: Healthcare Organisation, Organisation Contact, Organisational Structure.

\featuretable{--- & --- & \emph{Not yet captured} & \emph{Enterprise decomposition subject to review} \\}
\subsubsection{Feature - Diagnostic Result Association}\label{cap:ec-7a2df974}
\vspace{-1em}
\begin{table}[htbp]
  \centering
  \footnotesize
  \begin{tabularx}{1\textwidth}{|l|l|X|}
    \hline
    \textbf{Metadata Name} & \textbf{Metadata Value} & \textbf{Description} \\
    \hline
    \textbf{Reference Id} & \texttt{EC-7A2DF974} & \\
    \hline
    \textbf{Status} & Derived candidate & \\
    \hline
    \textbf{Ancestry} &  & \\
    \hline
    \textbf{Source} &  & \\
    \hline
    \textbf{Enabled By} &  & \\
    \hline
    \textbf{Supports} &  & \\
    \hline
    \textbf{Derived From} &  & \\
    \hline
    \textbf{Reviewed Status} & Inspect recurrence for true semantic equivalence; retain, split, rename or reject. & \\
    \hline
  \end{tabularx}
\end{table}

This candidate enterprise Feature abstracts a behaviour recurring in 2 distinct BEC CT3 contexts. The source-specific BEC Features remain in place; this entry does not assert a common software service or an identical clinical meaning.

Illustrative originating CT3 contexts: Diagnostic Service Integration, Diagnostic Workflow Management.

\featuretable{--- & --- & \emph{Not yet captured} & \emph{Enterprise decomposition subject to review} \\}
\subsubsection{Feature - Evidence Association}\label{cap:ec-279fe89f}
\vspace{-1em}
\begin{table}[htbp]
  \centering
  \footnotesize
  \begin{tabularx}{1\textwidth}{|l|l|X|}
    \hline
    \textbf{Metadata Name} & \textbf{Metadata Value} & \textbf{Description} \\
    \hline
    \textbf{Reference Id} & \texttt{EC-279FE89F} & \\
    \hline
    \textbf{Status} & Derived candidate & \\
    \hline
    \textbf{Ancestry} &  & \\
    \hline
    \textbf{Source} &  & \\
    \hline
    \textbf{Enabled By} &  & \\
    \hline
    \textbf{Supports} &  & \\
    \hline
    \textbf{Derived From} &  & \\
    \hline
    \textbf{Reviewed Status} & Inspect recurrence for true semantic equivalence; retain, split, rename or reject. & \\
    \hline
  \end{tabularx}
\end{table}

This candidate enterprise Feature abstracts a behaviour recurring in 2 distinct BEC CT3 contexts. The source-specific BEC Features remain in place; this entry does not assert a common software service or an identical clinical meaning.

Illustrative originating CT3 contexts: Client Privacy, Policy \& Security Governance.

\featuretable{--- & --- & \emph{Not yet captured} & \emph{Enterprise decomposition subject to review} \\}
\subsubsection{Feature - Location Hierarchy Management}\label{cap:ec-a179d91b}
\vspace{-1em}
\begin{table}[htbp]
  \centering
  \footnotesize
  \begin{tabularx}{1\textwidth}{|l|l|X|}
    \hline
    \textbf{Metadata Name} & \textbf{Metadata Value} & \textbf{Description} \\
    \hline
    \textbf{Reference Id} & \texttt{EC-A179D91B} & \\
    \hline
    \textbf{Status} & Derived candidate & \\
    \hline
    \textbf{Ancestry} &  & \\
    \hline
    \textbf{Source} &  & \\
    \hline
    \textbf{Enabled By} &  & \\
    \hline
    \textbf{Supports} &  & \\
    \hline
    \textbf{Derived From} &  & \\
    \hline
    \textbf{Reviewed Status} & Inspect recurrence for true semantic equivalence; retain, split, rename or reject. & \\
    \hline
  \end{tabularx}
\end{table}

This candidate enterprise Feature abstracts a behaviour recurring in 2 distinct BEC CT3 contexts. The source-specific BEC Features remain in place; this entry does not assert a common software service or an identical clinical meaning.

Illustrative originating CT3 contexts: Healthcare Location, Location Structure.

\featuretable{--- & --- & \emph{Not yet captured} & \emph{Enterprise decomposition subject to review} \\}
\subsubsection{Feature - Professional Relationship Management}\label{cap:ec-366afc9b}
\vspace{-1em}
\begin{table}[htbp]
  \centering
  \footnotesize
  \begin{tabularx}{1\textwidth}{|l|l|X|}
    \hline
    \textbf{Metadata Name} & \textbf{Metadata Value} & \textbf{Description} \\
    \hline
    \textbf{Reference Id} & \texttt{EC-366AFC9B} & \\
    \hline
    \textbf{Status} & Derived candidate & \\
    \hline
    \textbf{Ancestry} &  & \\
    \hline
    \textbf{Source} &  & \\
    \hline
    \textbf{Enabled By} &  & \\
    \hline
    \textbf{Supports} &  & \\
    \hline
    \textbf{Derived From} &  & \\
    \hline
    \textbf{Reviewed Status} & Inspect recurrence for true semantic equivalence; retain, split, rename or reject. & \\
    \hline
  \end{tabularx}
\end{table}

This candidate enterprise Feature abstracts a behaviour recurring in 2 distinct BEC CT3 contexts. The source-specific BEC Features remain in place; this entry does not assert a common software service or an identical clinical meaning.

Illustrative originating CT3 contexts: Practitioner, Practitioner Role.

\featuretable{--- & --- & \emph{Not yet captured} & \emph{Enterprise decomposition subject to review} \\}
\subsubsection{Feature - Qualification Association}\label{cap:ec-e2f0cb40}
\vspace{-1em}
\begin{table}[htbp]
  \centering
  \footnotesize
  \begin{tabularx}{1\textwidth}{|l|l|X|}
    \hline
    \textbf{Metadata Name} & \textbf{Metadata Value} & \textbf{Description} \\
    \hline
    \textbf{Reference Id} & \texttt{EC-E2F0CB40} & \\
    \hline
    \textbf{Status} & Derived candidate & \\
    \hline
    \textbf{Ancestry} &  & \\
    \hline
    \textbf{Source} &  & \\
    \hline
    \textbf{Enabled By} &  & \\
    \hline
    \textbf{Supports} &  & \\
    \hline
    \textbf{Derived From} &  & \\
    \hline
    \textbf{Reviewed Status} & Inspect recurrence for true semantic equivalence; retain, split, rename or reject. & \\
    \hline
  \end{tabularx}
\end{table}

This candidate enterprise Feature abstracts a behaviour recurring in 2 distinct BEC CT3 contexts. The source-specific BEC Features remain in place; this entry does not assert a common software service or an identical clinical meaning.

Illustrative originating CT3 contexts: Practitioner, Practitioner Role.

\featuretable{--- & --- & \emph{Not yet captured} & \emph{Enterprise decomposition subject to review} \\}
\subsubsection{Feature - Relationship Management}\label{cap:ec-f160e33e}
\vspace{-1em}
\begin{table}[htbp]
  \centering
  \footnotesize
  \begin{tabularx}{1\textwidth}{|l|l|X|}
    \hline
    \textbf{Metadata Name} & \textbf{Metadata Value} & \textbf{Description} \\
    \hline
    \textbf{Reference Id} & \texttt{EC-F160E33E} & \\
    \hline
    \textbf{Status} & Derived candidate & \\
    \hline
    \textbf{Ancestry} &  & \\
    \hline
    \textbf{Source} &  & \\
    \hline
    \textbf{Enabled By} &  & \\
    \hline
    \textbf{Supports} &  & \\
    \hline
    \textbf{Derived From} &  & \\
    \hline
    \textbf{Reviewed Status} & Inspect recurrence for true semantic equivalence; retain, split, rename or reject. & \\
    \hline
  \end{tabularx}
\end{table}

This candidate enterprise Feature abstracts a behaviour recurring in 2 distinct BEC CT3 contexts. The source-specific BEC Features remain in place; this entry does not assert a common software service or an identical clinical meaning.

Illustrative originating CT3 contexts: Person Identity, Product Formulary \& Catalogue.

\featuretable{--- & --- & \emph{Not yet captured} & \emph{Enterprise decomposition subject to review} \\}

\chapter{CT1 - Information and Semantics}\label{cap:ec-47c59623}
\noindent\textbf{Reference ID:} \texttt{EC-47C59623}\quad \textbf{Status:} Derived candidate\par
Intent. Information and Semantics is a candidate horizontally reusable architectural capability associated with BEC recurrence. It is framed around recurring enabling behaviour, not a mandatory centralised service or a particular technology.

Boundary. A common architectural capability does not remove the originating Features from their healthcare contexts. Similar wording does not establish equivalence or identical implementation.

\section{CT2 - Information Lifecycle}\label{cap:ec-260111dc}
\vspace{-0.5em}
\begin{table}[htbp]
  \centering
  \footnotesize
  \begin{tabularx}{1\textwidth}{|l|l|X|}
    \hline
    \multicolumn{3}{c}{\textbf{Capability Metadata}} \\
    \hline
    \textbf{Name} & \textbf{Value} & \textbf{Comment and/or Description} \\
    \hline
    \textbf{Reference Id} & \texttt{EC-260111DC} & \\
    \hline
    \textbf{Status} & Derived candidate & \\
    \hline
    \textbf{Ancestry} &  & \\
    \hline
    \textbf{Source} &  & \\
    \hline
    \textbf{Enabled By} &  & \\
    \hline
    \textbf{Supports} &  & \\
    \hline
    \textbf{Derived From} &  & \\
    \hline
    \textbf{Reviewed Status} &  & \\
    \hline
  \end{tabularx}
\end{table}

\begin{table}[htbp]
  \centering
  \footnotesize
  \begin{tabularx}{1\textwidth}{|l|l|X|}
    \hline
    \multicolumn{3}{c}{\textbf{Contained Capabilities}} \\
    \hline
    \textbf{Capability Name} & \textbf{Capability Id} & \textbf{Comment} \\
    \hline
    Information Lifecycle Enablement & \texttt{cap:ec-ae5e9793} (at \ref{cap:ec-ae5e9793}) & \\
    \hline
  \end{tabularx}
\end{table}

Intent. Information Lifecycle is a candidate horizontally reusable architectural capability associated with Information and Semantics. It is framed around recurring enabling behaviour, not a mandatory centralised service or a particular technology.

Boundary. A common architectural capability does not remove the originating Features from their healthcare contexts. Similar wording does not establish equivalence or identical implementation.

\subsection{CT3 - Information Lifecycle Enablement}\label{cap:ec-ae5e9793}
\vspace{-0.5em}
\begin{table}[htbp]
  \centering
  \footnotesize
  \begin{tabularx}{1\textwidth}{|l|l|X|}
    \hline
    \multicolumn{3}{c}{\textbf{Capability Metadata}} \\
    \hline
    \textbf{Name} & \textbf{Value} & \textbf{Comment and/or Description} \\
    \hline
    \textbf{Reference Id} & \texttt{EC-AE5E9793} & \\
    \hline
    \textbf{Status} & Derived candidate & \\
    \hline
    \textbf{Ancestry} &  & \\
    \hline
    \textbf{Source} &  & \\
    \hline
    \textbf{Enabled By} &  & \\
    \hline
    \textbf{Supports} &  & \\
    \hline
    \textbf{Derived From} &  & \\
    \hline
    \textbf{Reviewed Status} &  & \\
    \hline
  \end{tabularx}
\end{table}

\begin{table}[htbp]
  \centering
  \footnotesize
  \begin{tabularx}{1\textwidth}{|l|l|X|}
    \hline
    \multicolumn{3}{c}{\textbf{Contained Features}} \\
    \hline
    \textbf{Feature Name} & \textbf{Feature Id} & \textbf{Comment} \\
    \hline
    Information Retrieval & \texttt{cap:ec-54c1bcab} (at \ref{cap:ec-54c1bcab}) & \\
    \hline
    Subject Information Discovery & \texttt{cap:ec-9788f116} (at \ref{cap:ec-9788f116}) & \\
    \hline
    Information Correlation & \texttt{cap:ec-bfa64aaf} (at \ref{cap:ec-bfa64aaf}) & \\
    \hline
    Information Lifecycle Management & \texttt{cap:ec-a3a099ac} (at \ref{cap:ec-a3a099ac}) & \\
    \hline
    Information Provenance & \texttt{cap:ec-10e4ea11} (at \ref{cap:ec-10e4ea11}) & \\
    \hline
    Information Validation & \texttt{cap:ec-23381315} (at \ref{cap:ec-23381315}) & \\
    \hline
    Medication Information Correlation & \texttt{cap:ec-53dff234} (at \ref{cap:ec-53dff234}) & \\
    \hline
    Request Lifecycle Management & \texttt{cap:ec-720e493e} (at \ref{cap:ec-720e493e}) & \\
    \hline
  \end{tabularx}
\end{table}

Intent. Information Lifecycle Enablement is a candidate horizontally reusable architectural capability associated with Information Lifecycle. It is framed around recurring enabling behaviour, not a mandatory centralised service or a particular technology.

Boundary. A common architectural capability does not remove the originating Features from their healthcare contexts. Similar wording does not establish equivalence or identical implementation.

\subsubsection{Feature - Information Retrieval}\label{cap:ec-54c1bcab}
\vspace{-1em}
\begin{table}[htbp]
  \centering
  \footnotesize
  \begin{tabularx}{1\textwidth}{|l|l|X|}
    \hline
    \textbf{Metadata Name} & \textbf{Metadata Value} & \textbf{Description} \\
    \hline
    \textbf{Reference Id} & \texttt{EC-54C1BCAB} & \\
    \hline
    \textbf{Status} & Derived candidate & \\
    \hline
    \textbf{Ancestry} &  & \\
    \hline
    \textbf{Source} &  & \\
    \hline
    \textbf{Enabled By} &  & \\
    \hline
    \textbf{Supports} &  & \\
    \hline
    \textbf{Derived From} &  & \\
    \hline
    \textbf{Reviewed Status} & Inspect recurrence for true semantic equivalence; retain, split, rename or reject. & \\
    \hline
  \end{tabularx}
\end{table}

This candidate enterprise Feature abstracts a behaviour recurring in 27 distinct BEC CT3 contexts. The source-specific BEC Features remain in place; this entry does not assert a common software service or an identical clinical meaning.

Illustrative originating CT3 contexts: Bed State \& Turnover Management, Boundary Protocol Mediation, Canonical Semantics \& Schema Governance, Care-Place Definition, Client Intake \& Demographic Update, Client Relationship, Clinical Documentation Lifecycle, Credential Verification Checkpoint, Decoupled Human Interaction \& Dashboards, Emergency Care Progression, Facility Demand \& Capacity Telemetry, Governed Query \& Search, and others.

\featuretable{--- & --- & \emph{Not yet captured} & \emph{Enterprise decomposition subject to review} \\}
\subsubsection{Feature - Subject Information Discovery}\label{cap:ec-9788f116}
\vspace{-1em}
\begin{table}[htbp]
  \centering
  \footnotesize
  \begin{tabularx}{1\textwidth}{|l|l|X|}
    \hline
    \textbf{Metadata Name} & \textbf{Metadata Value} & \textbf{Description} \\
    \hline
    \textbf{Reference Id} & \texttt{EC-9788F116} & \\
    \hline
    \textbf{Status} & Derived candidate & \\
    \hline
    \textbf{Ancestry} &  & \\
    \hline
    \textbf{Source} &  & \\
    \hline
    \textbf{Enabled By} &  & \\
    \hline
    \textbf{Supports} &  & \\
    \hline
    \textbf{Derived From} &  & \\
    \hline
    \textbf{Reviewed Status} & Inspect recurrence for true semantic equivalence; retain, split, rename or reject. & \\
    \hline
  \end{tabularx}
\end{table}

This candidate enterprise Feature abstracts a behaviour recurring in 4 distinct BEC CT3 contexts. The source-specific BEC Features remain in place; this entry does not assert a common software service or an identical clinical meaning.

Illustrative originating CT3 contexts: Multidisciplinary Inpatient Care Coordination, Outpatient Flow Enablement, Patient Transfer Coordination, Remote Patient Monitoring.

\featuretable{--- & --- & \emph{Not yet captured} & \emph{Enterprise decomposition subject to review} \\}
\subsubsection{Feature - Information Correlation}\label{cap:ec-bfa64aaf}
\vspace{-1em}
\begin{table}[htbp]
  \centering
  \footnotesize
  \begin{tabularx}{1\textwidth}{|l|l|X|}
    \hline
    \textbf{Metadata Name} & \textbf{Metadata Value} & \textbf{Description} \\
    \hline
    \textbf{Reference Id} & \texttt{EC-BFA64AAF} & \\
    \hline
    \textbf{Status} & Derived candidate & \\
    \hline
    \textbf{Ancestry} &  & \\
    \hline
    \textbf{Source} &  & \\
    \hline
    \textbf{Enabled By} &  & \\
    \hline
    \textbf{Supports} &  & \\
    \hline
    \textbf{Derived From} &  & \\
    \hline
    \textbf{Reviewed Status} & Inspect recurrence for true semantic equivalence; retain, split, rename or reject. & \\
    \hline
  \end{tabularx}
\end{table}

This candidate enterprise Feature abstracts a behaviour recurring in 3 distinct BEC CT3 contexts. The source-specific BEC Features remain in place; this entry does not assert a common software service or an identical clinical meaning.

Illustrative originating CT3 contexts: Canonical Semantics \& Schema Governance, Multi-Model Information Exchange, Terminology \& Reference Mapping.

\featuretable{--- & --- & \emph{Not yet captured} & \emph{Enterprise decomposition subject to review} \\}
\subsubsection{Feature - Information Lifecycle Management}\label{cap:ec-a3a099ac}
\vspace{-1em}
\begin{table}[htbp]
  \centering
  \footnotesize
  \begin{tabularx}{1\textwidth}{|l|l|X|}
    \hline
    \textbf{Metadata Name} & \textbf{Metadata Value} & \textbf{Description} \\
    \hline
    \textbf{Reference Id} & \texttt{EC-A3A099AC} & \\
    \hline
    \textbf{Status} & Derived candidate & \\
    \hline
    \textbf{Ancestry} &  & \\
    \hline
    \textbf{Source} &  & \\
    \hline
    \textbf{Enabled By} &  & \\
    \hline
    \textbf{Supports} &  & \\
    \hline
    \textbf{Derived From} &  & \\
    \hline
    \textbf{Reviewed Status} & Inspect recurrence for true semantic equivalence; retain, split, rename or reject. & \\
    \hline
  \end{tabularx}
\end{table}

This candidate enterprise Feature abstracts a behaviour recurring in 3 distinct BEC CT3 contexts. The source-specific BEC Features remain in place; this entry does not assert a common software service or an identical clinical meaning.

Illustrative originating CT3 contexts: Canonical Semantics \& Schema Governance, Multi-Model Information Exchange, Terminology \& Reference Mapping.

\featuretable{--- & --- & \emph{Not yet captured} & \emph{Enterprise decomposition subject to review} \\}
\subsubsection{Feature - Information Provenance}\label{cap:ec-10e4ea11}
\vspace{-1em}
\begin{table}[htbp]
  \centering
  \footnotesize
  \begin{tabularx}{1\textwidth}{|l|l|X|}
    \hline
    \textbf{Metadata Name} & \textbf{Metadata Value} & \textbf{Description} \\
    \hline
    \textbf{Reference Id} & \texttt{EC-10E4EA11} & \\
    \hline
    \textbf{Status} & Derived candidate & \\
    \hline
    \textbf{Ancestry} &  & \\
    \hline
    \textbf{Source} &  & \\
    \hline
    \textbf{Enabled By} &  & \\
    \hline
    \textbf{Supports} &  & \\
    \hline
    \textbf{Derived From} &  & \\
    \hline
    \textbf{Reviewed Status} & Inspect recurrence for true semantic equivalence; retain, split, rename or reject. & \\
    \hline
  \end{tabularx}
\end{table}

This candidate enterprise Feature abstracts a behaviour recurring in 3 distinct BEC CT3 contexts. The source-specific BEC Features remain in place; this entry does not assert a common software service or an identical clinical meaning.

Illustrative originating CT3 contexts: Canonical Semantics \& Schema Governance, Multi-Model Information Exchange, Terminology \& Reference Mapping.

\featuretable{--- & --- & \emph{Not yet captured} & \emph{Enterprise decomposition subject to review} \\}
\subsubsection{Feature - Information Validation}\label{cap:ec-23381315}
\vspace{-1em}
\begin{table}[htbp]
  \centering
  \footnotesize
  \begin{tabularx}{1\textwidth}{|l|l|X|}
    \hline
    \textbf{Metadata Name} & \textbf{Metadata Value} & \textbf{Description} \\
    \hline
    \textbf{Reference Id} & \texttt{EC-23381315} & \\
    \hline
    \textbf{Status} & Derived candidate & \\
    \hline
    \textbf{Ancestry} &  & \\
    \hline
    \textbf{Source} &  & \\
    \hline
    \textbf{Enabled By} &  & \\
    \hline
    \textbf{Supports} &  & \\
    \hline
    \textbf{Derived From} &  & \\
    \hline
    \textbf{Reviewed Status} & Inspect recurrence for true semantic equivalence; retain, split, rename or reject. & \\
    \hline
  \end{tabularx}
\end{table}

This candidate enterprise Feature abstracts a behaviour recurring in 3 distinct BEC CT3 contexts. The source-specific BEC Features remain in place; this entry does not assert a common software service or an identical clinical meaning.

Illustrative originating CT3 contexts: Canonical Semantics \& Schema Governance, Multi-Model Information Exchange, Terminology \& Reference Mapping.

\featuretable{--- & --- & \emph{Not yet captured} & \emph{Enterprise decomposition subject to review} \\}
\subsubsection{Feature - Medication Information Correlation}\label{cap:ec-53dff234}
\vspace{-1em}
\begin{table}[htbp]
  \centering
  \footnotesize
  \begin{tabularx}{1\textwidth}{|l|l|X|}
    \hline
    \textbf{Metadata Name} & \textbf{Metadata Value} & \textbf{Description} \\
    \hline
    \textbf{Reference Id} & \texttt{EC-53DFF234} & \\
    \hline
    \textbf{Status} & Derived candidate & \\
    \hline
    \textbf{Ancestry} &  & \\
    \hline
    \textbf{Source} &  & \\
    \hline
    \textbf{Enabled By} &  & \\
    \hline
    \textbf{Supports} &  & \\
    \hline
    \textbf{Derived From} &  & \\
    \hline
    \textbf{Reviewed Status} & Inspect recurrence for true semantic equivalence; retain, split, rename or reject. & \\
    \hline
  \end{tabularx}
\end{table}

This candidate enterprise Feature abstracts a behaviour recurring in 2 distinct BEC CT3 contexts. The source-specific BEC Features remain in place; this entry does not assert a common software service or an identical clinical meaning.

Illustrative originating CT3 contexts: Medication Order \& Dispense Management, Medication Therapy Provision and Monitoring.

\featuretable{--- & --- & \emph{Not yet captured} & \emph{Enterprise decomposition subject to review} \\}
\subsubsection{Feature - Request Lifecycle Management}\label{cap:ec-720e493e}
\vspace{-1em}
\begin{table}[htbp]
  \centering
  \footnotesize
  \begin{tabularx}{1\textwidth}{|l|l|X|}
    \hline
    \textbf{Metadata Name} & \textbf{Metadata Value} & \textbf{Description} \\
    \hline
    \textbf{Reference Id} & \texttt{EC-720E493E} & \\
    \hline
    \textbf{Status} & Derived candidate & \\
    \hline
    \textbf{Ancestry} &  & \\
    \hline
    \textbf{Source} &  & \\
    \hline
    \textbf{Enabled By} &  & \\
    \hline
    \textbf{Supports} &  & \\
    \hline
    \textbf{Derived From} &  & \\
    \hline
    \textbf{Reviewed Status} & Inspect recurrence for true semantic equivalence; retain, split, rename or reject. & \\
    \hline
  \end{tabularx}
\end{table}

This candidate enterprise Feature abstracts a behaviour recurring in 2 distinct BEC CT3 contexts. The source-specific BEC Features remain in place; this entry does not assert a common software service or an identical clinical meaning.

Illustrative originating CT3 contexts: Closed-Loop Order Progression, Medication Order \& Dispense Management.

\featuretable{--- & --- & \emph{Not yet captured} & \emph{Enterprise decomposition subject to review} \\}

\section{CT2 - Discovery and Retrieval}\label{cap:ec-2a3844c6}
\vspace{-0.5em}
\begin{table}[htbp]
  \centering
  \footnotesize
  \begin{tabularx}{1\textwidth}{|l|l|X|}
    \hline
    \multicolumn{3}{c}{\textbf{Capability Metadata}} \\
    \hline
    \textbf{Name} & \textbf{Value} & \textbf{Comment and/or Description} \\
    \hline
    \textbf{Reference Id} & \texttt{EC-2A3844C6} & \\
    \hline
    \textbf{Status} & Derived candidate & \\
    \hline
    \textbf{Ancestry} &  & \\
    \hline
    \textbf{Source} &  & \\
    \hline
    \textbf{Enabled By} &  & \\
    \hline
    \textbf{Supports} &  & \\
    \hline
    \textbf{Derived From} &  & \\
    \hline
    \textbf{Reviewed Status} &  & \\
    \hline
  \end{tabularx}
\end{table}

\begin{table}[htbp]
  \centering
  \footnotesize
  \begin{tabularx}{1\textwidth}{|l|l|X|}
    \hline
    \multicolumn{3}{c}{\textbf{Contained Capabilities}} \\
    \hline
    \textbf{Capability Name} & \textbf{Capability Id} & \textbf{Comment} \\
    \hline
    Discovery and Retrieval Enablement & \texttt{cap:ec-299d0172} (at \ref{cap:ec-299d0172}) & \\
    \hline
  \end{tabularx}
\end{table}

Intent. Discovery and Retrieval is a candidate horizontally reusable architectural capability associated with Information and Semantics. It is framed around recurring enabling behaviour, not a mandatory centralised service or a particular technology.

Boundary. A common architectural capability does not remove the originating Features from their healthcare contexts. Similar wording does not establish equivalence or identical implementation.

\subsection{CT3 - Discovery and Retrieval Enablement}\label{cap:ec-299d0172}
\vspace{-0.5em}
\begin{table}[htbp]
  \centering
  \footnotesize
  \begin{tabularx}{1\textwidth}{|l|l|X|}
    \hline
    \multicolumn{3}{c}{\textbf{Capability Metadata}} \\
    \hline
    \textbf{Name} & \textbf{Value} & \textbf{Comment and/or Description} \\
    \hline
    \textbf{Reference Id} & \texttt{EC-299D0172} & \\
    \hline
    \textbf{Status} & Derived candidate & \\
    \hline
    \textbf{Ancestry} &  & \\
    \hline
    \textbf{Source} &  & \\
    \hline
    \textbf{Enabled By} &  & \\
    \hline
    \textbf{Supports} &  & \\
    \hline
    \textbf{Derived From} &  & \\
    \hline
    \textbf{Reviewed Status} &  & \\
    \hline
  \end{tabularx}
\end{table}

\begin{table}[htbp]
  \centering
  \footnotesize
  \begin{tabularx}{1\textwidth}{|l|l|X|}
    \hline
    \multicolumn{3}{c}{\textbf{Contained Features}} \\
    \hline
    \textbf{Feature Name} & \textbf{Feature Id} & \textbf{Comment} \\
    \hline
    Information Retrieval & \texttt{cap:ec-f2753584} (at \ref{cap:ec-f2753584}) & \\
    \hline
    Service Discovery & \texttt{cap:ec-82149096} (at \ref{cap:ec-82149096}) & \\
    \hline
    Subject Information Discovery & \texttt{cap:ec-13d8a141} (at \ref{cap:ec-13d8a141}) & \\
    \hline
    Organisation Discovery & \texttt{cap:ec-50a2c171} (at \ref{cap:ec-50a2c171}) & \\
    \hline
    Location Discovery & \texttt{cap:ec-009a4a5a} (at \ref{cap:ec-009a4a5a}) & \\
    \hline
    Medication Identity Resolution & \texttt{cap:ec-09a3798f} (at \ref{cap:ec-09a3798f}) & \\
    \hline
    Practitioner Discovery & \texttt{cap:ec-26353d17} (at \ref{cap:ec-26353d17}) & \\
    \hline
  \end{tabularx}
\end{table}

Intent. Discovery and Retrieval Enablement is a candidate horizontally reusable architectural capability associated with Discovery and Retrieval. It is framed around recurring enabling behaviour, not a mandatory centralised service or a particular technology.

Boundary. A common architectural capability does not remove the originating Features from their healthcare contexts. Similar wording does not establish equivalence or identical implementation.

\subsubsection{Feature - Information Retrieval}\label{cap:ec-f2753584}
\vspace{-1em}
\begin{table}[htbp]
  \centering
  \footnotesize
  \begin{tabularx}{1\textwidth}{|l|l|X|}
    \hline
    \textbf{Metadata Name} & \textbf{Metadata Value} & \textbf{Description} \\
    \hline
    \textbf{Reference Id} & \texttt{EC-F2753584} & \\
    \hline
    \textbf{Status} & Derived candidate & \\
    \hline
    \textbf{Ancestry} &  & \\
    \hline
    \textbf{Source} &  & \\
    \hline
    \textbf{Enabled By} &  & \\
    \hline
    \textbf{Supports} &  & \\
    \hline
    \textbf{Derived From} &  & \\
    \hline
    \textbf{Reviewed Status} & Inspect recurrence for true semantic equivalence; retain, split, rename or reject. & \\
    \hline
  \end{tabularx}
\end{table}

This candidate enterprise Feature abstracts a behaviour recurring in 27 distinct BEC CT3 contexts. The source-specific BEC Features remain in place; this entry does not assert a common software service or an identical clinical meaning.

Illustrative originating CT3 contexts: Bed State \& Turnover Management, Boundary Protocol Mediation, Canonical Semantics \& Schema Governance, Care-Place Definition, Client Intake \& Demographic Update, Client Relationship, Clinical Documentation Lifecycle, Credential Verification Checkpoint, Decoupled Human Interaction \& Dashboards, Emergency Care Progression, Facility Demand \& Capacity Telemetry, Governed Query \& Search, and others.

\featuretable{--- & --- & \emph{Not yet captured} & \emph{Enterprise decomposition subject to review} \\}
\subsubsection{Feature - Service Discovery}\label{cap:ec-82149096}
\vspace{-1em}
\begin{table}[htbp]
  \centering
  \footnotesize
  \begin{tabularx}{1\textwidth}{|l|l|X|}
    \hline
    \textbf{Metadata Name} & \textbf{Metadata Value} & \textbf{Description} \\
    \hline
    \textbf{Reference Id} & \texttt{EC-82149096} & \\
    \hline
    \textbf{Status} & Derived candidate & \\
    \hline
    \textbf{Ancestry} &  & \\
    \hline
    \textbf{Source} &  & \\
    \hline
    \textbf{Enabled By} &  & \\
    \hline
    \textbf{Supports} &  & \\
    \hline
    \textbf{Derived From} &  & \\
    \hline
    \textbf{Reviewed Status} & Inspect recurrence for true semantic equivalence; retain, split, rename or reject. & \\
    \hline
  \end{tabularx}
\end{table}

This candidate enterprise Feature abstracts a behaviour recurring in 5 distinct BEC CT3 contexts. The source-specific BEC Features remain in place; this entry does not assert a common software service or an identical clinical meaning.

Illustrative originating CT3 contexts: Diagnostic Service Integration, Healthcare Service Definition, Service Availability Definition, Service Eligibility Definition, Service Provision Relationship.

\featuretable{--- & --- & \emph{Not yet captured} & \emph{Enterprise decomposition subject to review} \\}
\subsubsection{Feature - Subject Information Discovery}\label{cap:ec-13d8a141}
\vspace{-1em}
\begin{table}[htbp]
  \centering
  \footnotesize
  \begin{tabularx}{1\textwidth}{|l|l|X|}
    \hline
    \textbf{Metadata Name} & \textbf{Metadata Value} & \textbf{Description} \\
    \hline
    \textbf{Reference Id} & \texttt{EC-13D8A141} & \\
    \hline
    \textbf{Status} & Derived candidate & \\
    \hline
    \textbf{Ancestry} &  & \\
    \hline
    \textbf{Source} &  & \\
    \hline
    \textbf{Enabled By} &  & \\
    \hline
    \textbf{Supports} &  & \\
    \hline
    \textbf{Derived From} &  & \\
    \hline
    \textbf{Reviewed Status} & Inspect recurrence for true semantic equivalence; retain, split, rename or reject. & \\
    \hline
  \end{tabularx}
\end{table}

This candidate enterprise Feature abstracts a behaviour recurring in 4 distinct BEC CT3 contexts. The source-specific BEC Features remain in place; this entry does not assert a common software service or an identical clinical meaning.

Illustrative originating CT3 contexts: Multidisciplinary Inpatient Care Coordination, Outpatient Flow Enablement, Patient Transfer Coordination, Remote Patient Monitoring.

\featuretable{--- & --- & \emph{Not yet captured} & \emph{Enterprise decomposition subject to review} \\}
\subsubsection{Feature - Organisation Discovery}\label{cap:ec-50a2c171}
\vspace{-1em}
\begin{table}[htbp]
  \centering
  \footnotesize
  \begin{tabularx}{1\textwidth}{|l|l|X|}
    \hline
    \textbf{Metadata Name} & \textbf{Metadata Value} & \textbf{Description} \\
    \hline
    \textbf{Reference Id} & \texttt{EC-50A2C171} & \\
    \hline
    \textbf{Status} & Derived candidate & \\
    \hline
    \textbf{Ancestry} &  & \\
    \hline
    \textbf{Source} &  & \\
    \hline
    \textbf{Enabled By} &  & \\
    \hline
    \textbf{Supports} &  & \\
    \hline
    \textbf{Derived From} &  & \\
    \hline
    \textbf{Reviewed Status} & Inspect recurrence for true semantic equivalence; retain, split, rename or reject. & \\
    \hline
  \end{tabularx}
\end{table}

This candidate enterprise Feature abstracts a behaviour recurring in 3 distinct BEC CT3 contexts. The source-specific BEC Features remain in place; this entry does not assert a common software service or an identical clinical meaning.

Illustrative originating CT3 contexts: Healthcare Organisation, Organisation Contact, Organisational Structure.

\featuretable{--- & --- & \emph{Not yet captured} & \emph{Enterprise decomposition subject to review} \\}
\subsubsection{Feature - Location Discovery}\label{cap:ec-009a4a5a}
\vspace{-1em}
\begin{table}[htbp]
  \centering
  \footnotesize
  \begin{tabularx}{1\textwidth}{|l|l|X|}
    \hline
    \textbf{Metadata Name} & \textbf{Metadata Value} & \textbf{Description} \\
    \hline
    \textbf{Reference Id} & \texttt{EC-009A4A5A} & \\
    \hline
    \textbf{Status} & Derived candidate & \\
    \hline
    \textbf{Ancestry} &  & \\
    \hline
    \textbf{Source} &  & \\
    \hline
    \textbf{Enabled By} &  & \\
    \hline
    \textbf{Supports} &  & \\
    \hline
    \textbf{Derived From} &  & \\
    \hline
    \textbf{Reviewed Status} & Inspect recurrence for true semantic equivalence; retain, split, rename or reject. & \\
    \hline
  \end{tabularx}
\end{table}

This candidate enterprise Feature abstracts a behaviour recurring in 2 distinct BEC CT3 contexts. The source-specific BEC Features remain in place; this entry does not assert a common software service or an identical clinical meaning.

Illustrative originating CT3 contexts: Healthcare Location, Location Structure.

\featuretable{--- & --- & \emph{Not yet captured} & \emph{Enterprise decomposition subject to review} \\}
\subsubsection{Feature - Medication Identity Resolution}\label{cap:ec-09a3798f}
\vspace{-1em}
\begin{table}[htbp]
  \centering
  \footnotesize
  \begin{tabularx}{1\textwidth}{|l|l|X|}
    \hline
    \textbf{Metadata Name} & \textbf{Metadata Value} & \textbf{Description} \\
    \hline
    \textbf{Reference Id} & \texttt{EC-09A3798F} & \\
    \hline
    \textbf{Status} & Derived candidate & \\
    \hline
    \textbf{Ancestry} &  & \\
    \hline
    \textbf{Source} &  & \\
    \hline
    \textbf{Enabled By} &  & \\
    \hline
    \textbf{Supports} &  & \\
    \hline
    \textbf{Derived From} &  & \\
    \hline
    \textbf{Reviewed Status} & Inspect recurrence for true semantic equivalence; retain, split, rename or reject. & \\
    \hline
  \end{tabularx}
\end{table}

This candidate enterprise Feature abstracts a behaviour recurring in 2 distinct BEC CT3 contexts. The source-specific BEC Features remain in place; this entry does not assert a common software service or an identical clinical meaning.

Illustrative originating CT3 contexts: Medication Order \& Dispense Management, Medication Therapy Provision and Monitoring.

\featuretable{--- & --- & \emph{Not yet captured} & \emph{Enterprise decomposition subject to review} \\}
\subsubsection{Feature - Practitioner Discovery}\label{cap:ec-26353d17}
\vspace{-1em}
\begin{table}[htbp]
  \centering
  \footnotesize
  \begin{tabularx}{1\textwidth}{|l|l|X|}
    \hline
    \textbf{Metadata Name} & \textbf{Metadata Value} & \textbf{Description} \\
    \hline
    \textbf{Reference Id} & \texttt{EC-26353D17} & \\
    \hline
    \textbf{Status} & Derived candidate & \\
    \hline
    \textbf{Ancestry} &  & \\
    \hline
    \textbf{Source} &  & \\
    \hline
    \textbf{Enabled By} &  & \\
    \hline
    \textbf{Supports} &  & \\
    \hline
    \textbf{Derived From} &  & \\
    \hline
    \textbf{Reviewed Status} & Inspect recurrence for true semantic equivalence; retain, split, rename or reject. & \\
    \hline
  \end{tabularx}
\end{table}

This candidate enterprise Feature abstracts a behaviour recurring in 2 distinct BEC CT3 contexts. The source-specific BEC Features remain in place; this entry does not assert a common software service or an identical clinical meaning.

Illustrative originating CT3 contexts: Practitioner, Practitioner Role.

\featuretable{--- & --- & \emph{Not yet captured} & \emph{Enterprise decomposition subject to review} \\}

\section{CT2 - Semantic Conformance}\label{cap:ec-89e071d0}
\vspace{-0.5em}
\begin{table}[htbp]
  \centering
  \footnotesize
  \begin{tabularx}{1\textwidth}{|l|l|X|}
    \hline
    \multicolumn{3}{c}{\textbf{Capability Metadata}} \\
    \hline
    \textbf{Name} & \textbf{Value} & \textbf{Comment and/or Description} \\
    \hline
    \textbf{Reference Id} & \texttt{EC-89E071D0} & \\
    \hline
    \textbf{Status} & Derived candidate & \\
    \hline
    \textbf{Ancestry} &  & \\
    \hline
    \textbf{Source} &  & \\
    \hline
    \textbf{Enabled By} &  & \\
    \hline
    \textbf{Supports} &  & \\
    \hline
    \textbf{Derived From} &  & \\
    \hline
    \textbf{Reviewed Status} &  & \\
    \hline
  \end{tabularx}
\end{table}

\begin{table}[htbp]
  \centering
  \footnotesize
  \begin{tabularx}{1\textwidth}{|l|l|X|}
    \hline
    \multicolumn{3}{c}{\textbf{Contained Capabilities}} \\
    \hline
    \textbf{Capability Name} & \textbf{Capability Id} & \textbf{Comment} \\
    \hline
    Semantic Conformance Enablement & \texttt{cap:ec-47bece6f} (at \ref{cap:ec-47bece6f}) & \\
    \hline
  \end{tabularx}
\end{table}

Intent. Semantic Conformance is a candidate horizontally reusable architectural capability associated with Information and Semantics. It is framed around recurring enabling behaviour, not a mandatory centralised service or a particular technology.

Boundary. A common architectural capability does not remove the originating Features from their healthcare contexts. Similar wording does not establish equivalence or identical implementation.

\subsection{CT3 - Semantic Conformance Enablement}\label{cap:ec-47bece6f}
\vspace{-0.5em}
\begin{table}[htbp]
  \centering
  \footnotesize
  \begin{tabularx}{1\textwidth}{|l|l|X|}
    \hline
    \multicolumn{3}{c}{\textbf{Capability Metadata}} \\
    \hline
    \textbf{Name} & \textbf{Value} & \textbf{Comment and/or Description} \\
    \hline
    \textbf{Reference Id} & \texttt{EC-47BECE6F} & \\
    \hline
    \textbf{Status} & Derived candidate & \\
    \hline
    \textbf{Ancestry} &  & \\
    \hline
    \textbf{Source} &  & \\
    \hline
    \textbf{Enabled By} &  & \\
    \hline
    \textbf{Supports} &  & \\
    \hline
    \textbf{Derived From} &  & \\
    \hline
    \textbf{Reviewed Status} &  & \\
    \hline
  \end{tabularx}
\end{table}

\begin{table}[htbp]
  \centering
  \footnotesize
  \begin{tabularx}{1\textwidth}{|l|l|X|}
    \hline
    \multicolumn{3}{c}{\textbf{Contained Features}} \\
    \hline
    \textbf{Feature Name} & \textbf{Feature Id} & \textbf{Comment} \\
    \hline
    Eligibility Validation & \texttt{cap:ec-63e0e334} (at \ref{cap:ec-63e0e334}) & \\
    \hline
    Information Validation & \texttt{cap:ec-a86d200d} (at \ref{cap:ec-a86d200d}) & \\
    \hline
    Medication Safety Context Validation & \texttt{cap:ec-93e36e5b} (at \ref{cap:ec-93e36e5b}) & \\
    \hline
    Qualification Association & \texttt{cap:ec-db69daf0} (at \ref{cap:ec-db69daf0}) & \\
    \hline
  \end{tabularx}
\end{table}

Intent. Semantic Conformance Enablement is a candidate horizontally reusable architectural capability associated with Semantic Conformance. It is framed around recurring enabling behaviour, not a mandatory centralised service or a particular technology.

Boundary. A common architectural capability does not remove the originating Features from their healthcare contexts. Similar wording does not establish equivalence or identical implementation.

\subsubsection{Feature - Eligibility Validation}\label{cap:ec-63e0e334}
\vspace{-1em}
\begin{table}[htbp]
  \centering
  \footnotesize
  \begin{tabularx}{1\textwidth}{|l|l|X|}
    \hline
    \textbf{Metadata Name} & \textbf{Metadata Value} & \textbf{Description} \\
    \hline
    \textbf{Reference Id} & \texttt{EC-63E0E334} & \\
    \hline
    \textbf{Status} & Derived candidate & \\
    \hline
    \textbf{Ancestry} &  & \\
    \hline
    \textbf{Source} &  & \\
    \hline
    \textbf{Enabled By} &  & \\
    \hline
    \textbf{Supports} &  & \\
    \hline
    \textbf{Derived From} &  & \\
    \hline
    \textbf{Reviewed Status} & Inspect recurrence for true semantic equivalence; retain, split, rename or reject. & \\
    \hline
  \end{tabularx}
\end{table}

This candidate enterprise Feature abstracts a behaviour recurring in 11 distinct BEC CT3 contexts. The source-specific BEC Features remain in place; this entry does not assert a common software service or an identical clinical meaning.

Illustrative originating CT3 contexts: Acute Care Monitoring \& Event Coordination, Community Care Coordination, Discharge Progression \& Coordination, ED Operational Coordination, Mobile Worker Dispatch \& Safety, Multi-Work-Unit Orchestration, Operational Work Assignment \& Supervision, Primary Care Coordination, Session \& Queue Coordination, Specialized Equipment Coordination, Ward Logistics Coordination.

\featuretable{--- & --- & \emph{Not yet captured} & \emph{Enterprise decomposition subject to review} \\}
\subsubsection{Feature - Information Validation}\label{cap:ec-a86d200d}
\vspace{-1em}
\begin{table}[htbp]
  \centering
  \footnotesize
  \begin{tabularx}{1\textwidth}{|l|l|X|}
    \hline
    \textbf{Metadata Name} & \textbf{Metadata Value} & \textbf{Description} \\
    \hline
    \textbf{Reference Id} & \texttt{EC-A86D200D} & \\
    \hline
    \textbf{Status} & Derived candidate & \\
    \hline
    \textbf{Ancestry} &  & \\
    \hline
    \textbf{Source} &  & \\
    \hline
    \textbf{Enabled By} &  & \\
    \hline
    \textbf{Supports} &  & \\
    \hline
    \textbf{Derived From} &  & \\
    \hline
    \textbf{Reviewed Status} & Inspect recurrence for true semantic equivalence; retain, split, rename or reject. & \\
    \hline
  \end{tabularx}
\end{table}

This candidate enterprise Feature abstracts a behaviour recurring in 3 distinct BEC CT3 contexts. The source-specific BEC Features remain in place; this entry does not assert a common software service or an identical clinical meaning.

Illustrative originating CT3 contexts: Canonical Semantics \& Schema Governance, Multi-Model Information Exchange, Terminology \& Reference Mapping.

\featuretable{--- & --- & \emph{Not yet captured} & \emph{Enterprise decomposition subject to review} \\}
\subsubsection{Feature - Medication Safety Context Validation}\label{cap:ec-93e36e5b}
\vspace{-1em}
\begin{table}[htbp]
  \centering
  \footnotesize
  \begin{tabularx}{1\textwidth}{|l|l|X|}
    \hline
    \textbf{Metadata Name} & \textbf{Metadata Value} & \textbf{Description} \\
    \hline
    \textbf{Reference Id} & \texttt{EC-93E36E5B} & \\
    \hline
    \textbf{Status} & Derived candidate & \\
    \hline
    \textbf{Ancestry} &  & \\
    \hline
    \textbf{Source} &  & \\
    \hline
    \textbf{Enabled By} &  & \\
    \hline
    \textbf{Supports} &  & \\
    \hline
    \textbf{Derived From} &  & \\
    \hline
    \textbf{Reviewed Status} & Inspect recurrence for true semantic equivalence; retain, split, rename or reject. & \\
    \hline
  \end{tabularx}
\end{table}

This candidate enterprise Feature abstracts a behaviour recurring in 2 distinct BEC CT3 contexts. The source-specific BEC Features remain in place; this entry does not assert a common software service or an identical clinical meaning.

Illustrative originating CT3 contexts: Medication Order \& Dispense Management, Medication Therapy Provision and Monitoring.

\featuretable{--- & --- & \emph{Not yet captured} & \emph{Enterprise decomposition subject to review} \\}
\subsubsection{Feature - Qualification Association}\label{cap:ec-db69daf0}
\vspace{-1em}
\begin{table}[htbp]
  \centering
  \footnotesize
  \begin{tabularx}{1\textwidth}{|l|l|X|}
    \hline
    \textbf{Metadata Name} & \textbf{Metadata Value} & \textbf{Description} \\
    \hline
    \textbf{Reference Id} & \texttt{EC-DB69DAF0} & \\
    \hline
    \textbf{Status} & Derived candidate & \\
    \hline
    \textbf{Ancestry} &  & \\
    \hline
    \textbf{Source} &  & \\
    \hline
    \textbf{Enabled By} &  & \\
    \hline
    \textbf{Supports} &  & \\
    \hline
    \textbf{Derived From} &  & \\
    \hline
    \textbf{Reviewed Status} & Inspect recurrence for true semantic equivalence; retain, split, rename or reject. & \\
    \hline
  \end{tabularx}
\end{table}

This candidate enterprise Feature abstracts a behaviour recurring in 2 distinct BEC CT3 contexts. The source-specific BEC Features remain in place; this entry does not assert a common software service or an identical clinical meaning.

Illustrative originating CT3 contexts: Practitioner, Practitioner Role.

\featuretable{--- & --- & \emph{Not yet captured} & \emph{Enterprise decomposition subject to review} \\}

\chapter{CT1 - Context, Activity and Exchange}\label{cap:ec-b80f9118}
\noindent\textbf{Reference ID:} \texttt{EC-B80F9118}\quad \textbf{Status:} Derived candidate\par
Intent. Context, Activity and Exchange is a candidate horizontally reusable architectural capability associated with BEC recurrence. It is framed around recurring enabling behaviour, not a mandatory centralised service or a particular technology.

Boundary. A common architectural capability does not remove the originating Features from their healthcare contexts. Similar wording does not establish equivalence or identical implementation.

\section{CT2 - Context and State}\label{cap:ec-adb90088}
\vspace{-0.5em}
\begin{table}[htbp]
  \centering
  \footnotesize
  \begin{tabularx}{1\textwidth}{|l|l|X|}
    \hline
    \multicolumn{3}{c}{\textbf{Capability Metadata}} \\
    \hline
    \textbf{Name} & \textbf{Value} & \textbf{Comment and/or Description} \\
    \hline
    \textbf{Reference Id} & \texttt{EC-ADB90088} & \\
    \hline
    \textbf{Status} & Derived candidate & \\
    \hline
    \textbf{Ancestry} &  & \\
    \hline
    \textbf{Source} &  & \\
    \hline
    \textbf{Enabled By} &  & \\
    \hline
    \textbf{Supports} &  & \\
    \hline
    \textbf{Derived From} &  & \\
    \hline
    \textbf{Reviewed Status} &  & \\
    \hline
  \end{tabularx}
\end{table}

\begin{table}[htbp]
  \centering
  \footnotesize
  \begin{tabularx}{1\textwidth}{|l|l|X|}
    \hline
    \multicolumn{3}{c}{\textbf{Contained Capabilities}} \\
    \hline
    \textbf{Capability Name} & \textbf{Capability Id} & \textbf{Comment} \\
    \hline
    Context and State Enablement & \texttt{cap:ec-6f46c125} (at \ref{cap:ec-6f46c125}) & \\
    \hline
  \end{tabularx}
\end{table}

Intent. Context and State is a candidate horizontally reusable architectural capability associated with Context, Activity and Exchange. It is framed around recurring enabling behaviour, not a mandatory centralised service or a particular technology.

Boundary. A common architectural capability does not remove the originating Features from their healthcare contexts. Similar wording does not establish equivalence or identical implementation.

\subsection{CT3 - Context and State Enablement}\label{cap:ec-6f46c125}
\vspace{-0.5em}
\begin{table}[htbp]
  \centering
  \footnotesize
  \begin{tabularx}{1\textwidth}{|l|l|X|}
    \hline
    \multicolumn{3}{c}{\textbf{Capability Metadata}} \\
    \hline
    \textbf{Name} & \textbf{Value} & \textbf{Comment and/or Description} \\
    \hline
    \textbf{Reference Id} & \texttt{EC-6F46C125} & \\
    \hline
    \textbf{Status} & Derived candidate & \\
    \hline
    \textbf{Ancestry} &  & \\
    \hline
    \textbf{Source} &  & \\
    \hline
    \textbf{Enabled By} &  & \\
    \hline
    \textbf{Supports} &  & \\
    \hline
    \textbf{Derived From} &  & \\
    \hline
    \textbf{Reviewed Status} &  & \\
    \hline
  \end{tabularx}
\end{table}

\begin{table}[htbp]
  \centering
  \footnotesize
  \begin{tabularx}{1\textwidth}{|l|l|X|}
    \hline
    \multicolumn{3}{c}{\textbf{Contained Features}} \\
    \hline
    \textbf{Feature Name} & \textbf{Feature Id} & \textbf{Comment} \\
    \hline
    Context Establishment & \texttt{cap:ec-c3c11291} (at \ref{cap:ec-c3c11291}) & \\
    \hline
    State Monitoring & \texttt{cap:ec-cc8047e5} (at \ref{cap:ec-cc8047e5}) & \\
    \hline
    State Tracking & \texttt{cap:ec-d122bc94} (at \ref{cap:ec-d122bc94}) & \\
    \hline
    Service Availability Management & \texttt{cap:ec-ff26979c} (at \ref{cap:ec-ff26979c}) & \\
    \hline
    Device State Management & \texttt{cap:ec-6dce8859} (at \ref{cap:ec-6dce8859}) & \\
    \hline
    Subject Context Management & \texttt{cap:ec-1900212b} (at \ref{cap:ec-1900212b}) & \\
    \hline
    No Established Feature Decomposition (Adjacent / Reference Context) & \texttt{cap:ec-53072edd} (at \ref{cap:ec-53072edd}) & \\
    \hline
    Diagnostic Status Tracking & \texttt{cap:ec-7d34d652} (at \ref{cap:ec-7d34d652}) & \\
    \hline
    Location Availability Management & \texttt{cap:ec-582fec53} (at \ref{cap:ec-582fec53}) & \\
    \hline
    Medication Safety Context Validation & \texttt{cap:ec-2739eba1} (at \ref{cap:ec-2739eba1}) & \\
    \hline
    Request Status Correlation & \texttt{cap:ec-db0e7812} (at \ref{cap:ec-db0e7812}) & \\
    \hline
  \end{tabularx}
\end{table}

Intent. Context and State Enablement is a candidate horizontally reusable architectural capability associated with Context and State. It is framed around recurring enabling behaviour, not a mandatory centralised service or a particular technology.

Boundary. A common architectural capability does not remove the originating Features from their healthcare contexts. Similar wording does not establish equivalence or identical implementation.

\subsubsection{Feature - Context Establishment}\label{cap:ec-c3c11291}
\vspace{-1em}
\begin{table}[htbp]
  \centering
  \footnotesize
  \begin{tabularx}{1\textwidth}{|l|l|X|}
    \hline
    \textbf{Metadata Name} & \textbf{Metadata Value} & \textbf{Description} \\
    \hline
    \textbf{Reference Id} & \texttt{EC-C3C11291} & \\
    \hline
    \textbf{Status} & Derived candidate & \\
    \hline
    \textbf{Ancestry} &  & \\
    \hline
    \textbf{Source} &  & \\
    \hline
    \textbf{Enabled By} &  & \\
    \hline
    \textbf{Supports} &  & \\
    \hline
    \textbf{Derived From} &  & \\
    \hline
    \textbf{Reviewed Status} & Inspect recurrence for true semantic equivalence; retain, split, rename or reject. & \\
    \hline
  \end{tabularx}
\end{table}

This candidate enterprise Feature abstracts a behaviour recurring in 35 distinct BEC CT3 contexts. The source-specific BEC Features remain in place; this entry does not assert a common software service or an identical clinical meaning.

Illustrative originating CT3 contexts: Acute Care Monitoring \& Event Coordination, Bed State \& Turnover Management, Boundary Protocol Mediation, Care-Place Definition, Client Intake \& Demographic Update, Client Relationship, Clinical Documentation Lifecycle, Community Care Coordination, Credential Verification Checkpoint, Decoupled Human Interaction \& Dashboards, Discharge Progression \& Coordination, ED Operational Coordination, and others.

\featuretable{--- & --- & \emph{Not yet captured} & \emph{Enterprise decomposition subject to review} \\}
\subsubsection{Feature - State Monitoring}\label{cap:ec-cc8047e5}
\vspace{-1em}
\begin{table}[htbp]
  \centering
  \footnotesize
  \begin{tabularx}{1\textwidth}{|l|l|X|}
    \hline
    \textbf{Metadata Name} & \textbf{Metadata Value} & \textbf{Description} \\
    \hline
    \textbf{Reference Id} & \texttt{EC-CC8047E5} & \\
    \hline
    \textbf{Status} & Derived candidate & \\
    \hline
    \textbf{Ancestry} &  & \\
    \hline
    \textbf{Source} &  & \\
    \hline
    \textbf{Enabled By} &  & \\
    \hline
    \textbf{Supports} &  & \\
    \hline
    \textbf{Derived From} &  & \\
    \hline
    \textbf{Reviewed Status} & Inspect recurrence for true semantic equivalence; retain, split, rename or reject. & \\
    \hline
  \end{tabularx}
\end{table}

This candidate enterprise Feature abstracts a behaviour recurring in 24 distinct BEC CT3 contexts. The source-specific BEC Features remain in place; this entry does not assert a common software service or an identical clinical meaning.

Illustrative originating CT3 contexts: Bed State \& Turnover Management, Boundary Protocol Mediation, Care-Place Definition, Client Intake \& Demographic Update, Client Relationship, Clinical Documentation Lifecycle, Credential Verification Checkpoint, Decoupled Human Interaction \& Dashboards, Emergency Care Progression, Facility Demand \& Capacity Telemetry, Governed Query \& Search, Health Product Definition, and others.

\featuretable{--- & --- & \emph{Not yet captured} & \emph{Enterprise decomposition subject to review} \\}
\subsubsection{Feature - State Tracking}\label{cap:ec-d122bc94}
\vspace{-1em}
\begin{table}[htbp]
  \centering
  \footnotesize
  \begin{tabularx}{1\textwidth}{|l|l|X|}
    \hline
    \textbf{Metadata Name} & \textbf{Metadata Value} & \textbf{Description} \\
    \hline
    \textbf{Reference Id} & \texttt{EC-D122BC94} & \\
    \hline
    \textbf{Status} & Derived candidate & \\
    \hline
    \textbf{Ancestry} &  & \\
    \hline
    \textbf{Source} &  & \\
    \hline
    \textbf{Enabled By} &  & \\
    \hline
    \textbf{Supports} &  & \\
    \hline
    \textbf{Derived From} &  & \\
    \hline
    \textbf{Reviewed Status} & Inspect recurrence for true semantic equivalence; retain, split, rename or reject. & \\
    \hline
  \end{tabularx}
\end{table}

This candidate enterprise Feature abstracts a behaviour recurring in 11 distinct BEC CT3 contexts. The source-specific BEC Features remain in place; this entry does not assert a common software service or an identical clinical meaning.

Illustrative originating CT3 contexts: Acute Care Monitoring \& Event Coordination, Community Care Coordination, Discharge Progression \& Coordination, ED Operational Coordination, Mobile Worker Dispatch \& Safety, Multi-Work-Unit Orchestration, Operational Work Assignment \& Supervision, Primary Care Coordination, Session \& Queue Coordination, Specialized Equipment Coordination, Ward Logistics Coordination.

\featuretable{--- & --- & \emph{Not yet captured} & \emph{Enterprise decomposition subject to review} \\}
\subsubsection{Feature - Service Availability Management}\label{cap:ec-ff26979c}
\vspace{-1em}
\begin{table}[htbp]
  \centering
  \footnotesize
  \begin{tabularx}{1\textwidth}{|l|l|X|}
    \hline
    \textbf{Metadata Name} & \textbf{Metadata Value} & \textbf{Description} \\
    \hline
    \textbf{Reference Id} & \texttt{EC-FF26979C} & \\
    \hline
    \textbf{Status} & Derived candidate & \\
    \hline
    \textbf{Ancestry} &  & \\
    \hline
    \textbf{Source} &  & \\
    \hline
    \textbf{Enabled By} &  & \\
    \hline
    \textbf{Supports} &  & \\
    \hline
    \textbf{Derived From} &  & \\
    \hline
    \textbf{Reviewed Status} & Inspect recurrence for true semantic equivalence; retain, split, rename or reject. & \\
    \hline
  \end{tabularx}
\end{table}

This candidate enterprise Feature abstracts a behaviour recurring in 5 distinct BEC CT3 contexts. The source-specific BEC Features remain in place; this entry does not assert a common software service or an identical clinical meaning.

Illustrative originating CT3 contexts: Diagnostic Service Integration, Healthcare Service Definition, Service Availability Definition, Service Eligibility Definition, Service Provision Relationship.

\featuretable{--- & --- & \emph{Not yet captured} & \emph{Enterprise decomposition subject to review} \\}
\subsubsection{Feature - Device State Management}\label{cap:ec-6dce8859}
\vspace{-1em}
\begin{table}[htbp]
  \centering
  \footnotesize
  \begin{tabularx}{1\textwidth}{|l|l|X|}
    \hline
    \textbf{Metadata Name} & \textbf{Metadata Value} & \textbf{Description} \\
    \hline
    \textbf{Reference Id} & \texttt{EC-6DCE8859} & \\
    \hline
    \textbf{Status} & Derived candidate & \\
    \hline
    \textbf{Ancestry} &  & \\
    \hline
    \textbf{Source} &  & \\
    \hline
    \textbf{Enabled By} &  & \\
    \hline
    \textbf{Supports} &  & \\
    \hline
    \textbf{Derived From} &  & \\
    \hline
    \textbf{Reviewed Status} & Inspect recurrence for true semantic equivalence; retain, split, rename or reject. & \\
    \hline
  \end{tabularx}
\end{table}

This candidate enterprise Feature abstracts a behaviour recurring in 4 distinct BEC CT3 contexts. The source-specific BEC Features remain in place; this entry does not assert a common software service or an identical clinical meaning.

Illustrative originating CT3 contexts: Device Association, Device Connectivity, Device Definition, Device Instance.

\featuretable{--- & --- & \emph{Not yet captured} & \emph{Enterprise decomposition subject to review} \\}
\subsubsection{Feature - Subject Context Management}\label{cap:ec-1900212b}
\vspace{-1em}
\begin{table}[htbp]
  \centering
  \footnotesize
  \begin{tabularx}{1\textwidth}{|l|l|X|}
    \hline
    \textbf{Metadata Name} & \textbf{Metadata Value} & \textbf{Description} \\
    \hline
    \textbf{Reference Id} & \texttt{EC-1900212B} & \\
    \hline
    \textbf{Status} & Derived candidate & \\
    \hline
    \textbf{Ancestry} &  & \\
    \hline
    \textbf{Source} &  & \\
    \hline
    \textbf{Enabled By} &  & \\
    \hline
    \textbf{Supports} &  & \\
    \hline
    \textbf{Derived From} &  & \\
    \hline
    \textbf{Reviewed Status} & Inspect recurrence for true semantic equivalence; retain, split, rename or reject. & \\
    \hline
  \end{tabularx}
\end{table}

This candidate enterprise Feature abstracts a behaviour recurring in 4 distinct BEC CT3 contexts. The source-specific BEC Features remain in place; this entry does not assert a common software service or an identical clinical meaning.

Illustrative originating CT3 contexts: Multidisciplinary Inpatient Care Coordination, Outpatient Flow Enablement, Patient Transfer Coordination, Remote Patient Monitoring.

\featuretable{--- & --- & \emph{Not yet captured} & \emph{Enterprise decomposition subject to review} \\}
\subsubsection{Feature - No Established Feature Decomposition (Adjacent / Reference Context)}\label{cap:ec-53072edd}
\vspace{-1em}
\begin{table}[htbp]
  \centering
  \footnotesize
  \begin{tabularx}{1\textwidth}{|l|l|X|}
    \hline
    \textbf{Metadata Name} & \textbf{Metadata Value} & \textbf{Description} \\
    \hline
    \textbf{Reference Id} & \texttt{EC-53072EDD} & \\
    \hline
    \textbf{Status} & Derived candidate & \\
    \hline
    \textbf{Ancestry} &  & \\
    \hline
    \textbf{Source} &  & \\
    \hline
    \textbf{Enabled By} &  & \\
    \hline
    \textbf{Supports} &  & \\
    \hline
    \textbf{Derived From} &  & \\
    \hline
    \textbf{Reviewed Status} & Inspect recurrence for true semantic equivalence; retain, split, rename or reject. & \\
    \hline
  \end{tabularx}
\end{table}

This candidate enterprise Feature abstracts a behaviour recurring in 3 distinct BEC CT3 contexts. The source-specific BEC Features remain in place; this entry does not assert a common software service or an identical clinical meaning.

Illustrative originating CT3 contexts: Health Product Definition, Product Batch \& Serial Tracking, Product Formulary \& Catalogue.

\featuretable{--- & --- & \emph{Not yet captured} & \emph{Enterprise decomposition subject to review} \\}
\subsubsection{Feature - Diagnostic Status Tracking}\label{cap:ec-7d34d652}
\vspace{-1em}
\begin{table}[htbp]
  \centering
  \footnotesize
  \begin{tabularx}{1\textwidth}{|l|l|X|}
    \hline
    \textbf{Metadata Name} & \textbf{Metadata Value} & \textbf{Description} \\
    \hline
    \textbf{Reference Id} & \texttt{EC-7D34D652} & \\
    \hline
    \textbf{Status} & Derived candidate & \\
    \hline
    \textbf{Ancestry} &  & \\
    \hline
    \textbf{Source} &  & \\
    \hline
    \textbf{Enabled By} &  & \\
    \hline
    \textbf{Supports} &  & \\
    \hline
    \textbf{Derived From} &  & \\
    \hline
    \textbf{Reviewed Status} & Inspect recurrence for true semantic equivalence; retain, split, rename or reject. & \\
    \hline
  \end{tabularx}
\end{table}

This candidate enterprise Feature abstracts a behaviour recurring in 2 distinct BEC CT3 contexts. The source-specific BEC Features remain in place; this entry does not assert a common software service or an identical clinical meaning.

Illustrative originating CT3 contexts: Diagnostic Service Integration, Diagnostic Workflow Management.

\featuretable{--- & --- & \emph{Not yet captured} & \emph{Enterprise decomposition subject to review} \\}
\subsubsection{Feature - Location Availability Management}\label{cap:ec-582fec53}
\vspace{-1em}
\begin{table}[htbp]
  \centering
  \footnotesize
  \begin{tabularx}{1\textwidth}{|l|l|X|}
    \hline
    \textbf{Metadata Name} & \textbf{Metadata Value} & \textbf{Description} \\
    \hline
    \textbf{Reference Id} & \texttt{EC-582FEC53} & \\
    \hline
    \textbf{Status} & Derived candidate & \\
    \hline
    \textbf{Ancestry} &  & \\
    \hline
    \textbf{Source} &  & \\
    \hline
    \textbf{Enabled By} &  & \\
    \hline
    \textbf{Supports} &  & \\
    \hline
    \textbf{Derived From} &  & \\
    \hline
    \textbf{Reviewed Status} & Inspect recurrence for true semantic equivalence; retain, split, rename or reject. & \\
    \hline
  \end{tabularx}
\end{table}

This candidate enterprise Feature abstracts a behaviour recurring in 2 distinct BEC CT3 contexts. The source-specific BEC Features remain in place; this entry does not assert a common software service or an identical clinical meaning.

Illustrative originating CT3 contexts: Healthcare Location, Location Structure.

\featuretable{--- & --- & \emph{Not yet captured} & \emph{Enterprise decomposition subject to review} \\}
\subsubsection{Feature - Medication Safety Context Validation}\label{cap:ec-2739eba1}
\vspace{-1em}
\begin{table}[htbp]
  \centering
  \footnotesize
  \begin{tabularx}{1\textwidth}{|l|l|X|}
    \hline
    \textbf{Metadata Name} & \textbf{Metadata Value} & \textbf{Description} \\
    \hline
    \textbf{Reference Id} & \texttt{EC-2739EBA1} & \\
    \hline
    \textbf{Status} & Derived candidate & \\
    \hline
    \textbf{Ancestry} &  & \\
    \hline
    \textbf{Source} &  & \\
    \hline
    \textbf{Enabled By} &  & \\
    \hline
    \textbf{Supports} &  & \\
    \hline
    \textbf{Derived From} &  & \\
    \hline
    \textbf{Reviewed Status} & Inspect recurrence for true semantic equivalence; retain, split, rename or reject. & \\
    \hline
  \end{tabularx}
\end{table}

This candidate enterprise Feature abstracts a behaviour recurring in 2 distinct BEC CT3 contexts. The source-specific BEC Features remain in place; this entry does not assert a common software service or an identical clinical meaning.

Illustrative originating CT3 contexts: Medication Order \& Dispense Management, Medication Therapy Provision and Monitoring.

\featuretable{--- & --- & \emph{Not yet captured} & \emph{Enterprise decomposition subject to review} \\}
\subsubsection{Feature - Request Status Correlation}\label{cap:ec-db0e7812}
\vspace{-1em}
\begin{table}[htbp]
  \centering
  \footnotesize
  \begin{tabularx}{1\textwidth}{|l|l|X|}
    \hline
    \textbf{Metadata Name} & \textbf{Metadata Value} & \textbf{Description} \\
    \hline
    \textbf{Reference Id} & \texttt{EC-DB0E7812} & \\
    \hline
    \textbf{Status} & Derived candidate & \\
    \hline
    \textbf{Ancestry} &  & \\
    \hline
    \textbf{Source} &  & \\
    \hline
    \textbf{Enabled By} &  & \\
    \hline
    \textbf{Supports} &  & \\
    \hline
    \textbf{Derived From} &  & \\
    \hline
    \textbf{Reviewed Status} & Inspect recurrence for true semantic equivalence; retain, split, rename or reject. & \\
    \hline
  \end{tabularx}
\end{table}

This candidate enterprise Feature abstracts a behaviour recurring in 2 distinct BEC CT3 contexts. The source-specific BEC Features remain in place; this entry does not assert a common software service or an identical clinical meaning.

Illustrative originating CT3 contexts: Closed-Loop Order Progression, Medication Order \& Dispense Management.

\featuretable{--- & --- & \emph{Not yet captured} & \emph{Enterprise decomposition subject to review} \\}

\section{CT2 - Activity Coordination}\label{cap:ec-8c0e52a8}
\vspace{-0.5em}
\begin{table}[htbp]
  \centering
  \footnotesize
  \begin{tabularx}{1\textwidth}{|l|l|X|}
    \hline
    \multicolumn{3}{c}{\textbf{Capability Metadata}} \\
    \hline
    \textbf{Name} & \textbf{Value} & \textbf{Comment and/or Description} \\
    \hline
    \textbf{Reference Id} & \texttt{EC-8C0E52A8} & \\
    \hline
    \textbf{Status} & Derived candidate & \\
    \hline
    \textbf{Ancestry} &  & \\
    \hline
    \textbf{Source} &  & \\
    \hline
    \textbf{Enabled By} &  & \\
    \hline
    \textbf{Supports} &  & \\
    \hline
    \textbf{Derived From} &  & \\
    \hline
    \textbf{Reviewed Status} &  & \\
    \hline
  \end{tabularx}
\end{table}

\begin{table}[htbp]
  \centering
  \footnotesize
  \begin{tabularx}{1\textwidth}{|l|l|X|}
    \hline
    \multicolumn{3}{c}{\textbf{Contained Capabilities}} \\
    \hline
    \textbf{Capability Name} & \textbf{Capability Id} & \textbf{Comment} \\
    \hline
    Activity Coordination Enablement & \texttt{cap:ec-d05ce29d} (at \ref{cap:ec-d05ce29d}) & \\
    \hline
  \end{tabularx}
\end{table}

Intent. Activity Coordination is a candidate horizontally reusable architectural capability associated with Context, Activity and Exchange. It is framed around recurring enabling behaviour, not a mandatory centralised service or a particular technology.

Boundary. A common architectural capability does not remove the originating Features from their healthcare contexts. Similar wording does not establish equivalence or identical implementation.

\subsection{CT3 - Activity Coordination Enablement}\label{cap:ec-d05ce29d}
\vspace{-0.5em}
\begin{table}[htbp]
  \centering
  \footnotesize
  \begin{tabularx}{1\textwidth}{|l|l|X|}
    \hline
    \multicolumn{3}{c}{\textbf{Capability Metadata}} \\
    \hline
    \textbf{Name} & \textbf{Value} & \textbf{Comment and/or Description} \\
    \hline
    \textbf{Reference Id} & \texttt{EC-D05CE29D} & \\
    \hline
    \textbf{Status} & Derived candidate & \\
    \hline
    \textbf{Ancestry} &  & \\
    \hline
    \textbf{Source} &  & \\
    \hline
    \textbf{Enabled By} &  & \\
    \hline
    \textbf{Supports} &  & \\
    \hline
    \textbf{Derived From} &  & \\
    \hline
    \textbf{Reviewed Status} &  & \\
    \hline
  \end{tabularx}
\end{table}

\begin{table}[htbp]
  \centering
  \footnotesize
  \begin{tabularx}{1\textwidth}{|l|l|X|}
    \hline
    \multicolumn{3}{c}{\textbf{Contained Features}} \\
    \hline
    \textbf{Feature Name} & \textbf{Feature Id} & \textbf{Comment} \\
    \hline
    Assignment \& Routing & \texttt{cap:ec-5ee5f09e} (at \ref{cap:ec-5ee5f09e}) & \\
    \hline
    State Tracking & \texttt{cap:ec-e762705e} (at \ref{cap:ec-e762705e}) & \\
    \hline
    Device Provenance Tracking & \texttt{cap:ec-7f4beb56} (at \ref{cap:ec-7f4beb56}) & \\
    \hline
    Diagnostic Status Tracking & \texttt{cap:ec-95e97505} (at \ref{cap:ec-95e97505}) & \\
    \hline
    Medication Activity Tracking & \texttt{cap:ec-67a2e372} (at \ref{cap:ec-67a2e372}) & \\
    \hline
    Request Fulfilment Tracking & \texttt{cap:ec-a191a2bc} (at \ref{cap:ec-a191a2bc}) & \\
    \hline
  \end{tabularx}
\end{table}

Intent. Activity Coordination Enablement is a candidate horizontally reusable architectural capability associated with Activity Coordination. It is framed around recurring enabling behaviour, not a mandatory centralised service or a particular technology.

Boundary. A common architectural capability does not remove the originating Features from their healthcare contexts. Similar wording does not establish equivalence or identical implementation.

\subsubsection{Feature - Assignment \& Routing}\label{cap:ec-5ee5f09e}
\vspace{-1em}
\begin{table}[htbp]
  \centering
  \footnotesize
  \begin{tabularx}{1\textwidth}{|l|l|X|}
    \hline
    \textbf{Metadata Name} & \textbf{Metadata Value} & \textbf{Description} \\
    \hline
    \textbf{Reference Id} & \texttt{EC-5EE5F09E} & \\
    \hline
    \textbf{Status} & Derived candidate & \\
    \hline
    \textbf{Ancestry} &  & \\
    \hline
    \textbf{Source} &  & \\
    \hline
    \textbf{Enabled By} &  & \\
    \hline
    \textbf{Supports} &  & \\
    \hline
    \textbf{Derived From} &  & \\
    \hline
    \textbf{Reviewed Status} & Inspect recurrence for true semantic equivalence; retain, split, rename or reject. & \\
    \hline
  \end{tabularx}
\end{table}

This candidate enterprise Feature abstracts a behaviour recurring in 11 distinct BEC CT3 contexts. The source-specific BEC Features remain in place; this entry does not assert a common software service or an identical clinical meaning.

Illustrative originating CT3 contexts: Acute Care Monitoring \& Event Coordination, Community Care Coordination, Discharge Progression \& Coordination, ED Operational Coordination, Mobile Worker Dispatch \& Safety, Multi-Work-Unit Orchestration, Operational Work Assignment \& Supervision, Primary Care Coordination, Session \& Queue Coordination, Specialized Equipment Coordination, Ward Logistics Coordination.

\featuretable{--- & --- & \emph{Not yet captured} & \emph{Enterprise decomposition subject to review} \\}
\subsubsection{Feature - State Tracking}\label{cap:ec-e762705e}
\vspace{-1em}
\begin{table}[htbp]
  \centering
  \footnotesize
  \begin{tabularx}{1\textwidth}{|l|l|X|}
    \hline
    \textbf{Metadata Name} & \textbf{Metadata Value} & \textbf{Description} \\
    \hline
    \textbf{Reference Id} & \texttt{EC-E762705E} & \\
    \hline
    \textbf{Status} & Derived candidate & \\
    \hline
    \textbf{Ancestry} &  & \\
    \hline
    \textbf{Source} &  & \\
    \hline
    \textbf{Enabled By} &  & \\
    \hline
    \textbf{Supports} &  & \\
    \hline
    \textbf{Derived From} &  & \\
    \hline
    \textbf{Reviewed Status} & Inspect recurrence for true semantic equivalence; retain, split, rename or reject. & \\
    \hline
  \end{tabularx}
\end{table}

This candidate enterprise Feature abstracts a behaviour recurring in 11 distinct BEC CT3 contexts. The source-specific BEC Features remain in place; this entry does not assert a common software service or an identical clinical meaning.

Illustrative originating CT3 contexts: Acute Care Monitoring \& Event Coordination, Community Care Coordination, Discharge Progression \& Coordination, ED Operational Coordination, Mobile Worker Dispatch \& Safety, Multi-Work-Unit Orchestration, Operational Work Assignment \& Supervision, Primary Care Coordination, Session \& Queue Coordination, Specialized Equipment Coordination, Ward Logistics Coordination.

\featuretable{--- & --- & \emph{Not yet captured} & \emph{Enterprise decomposition subject to review} \\}
\subsubsection{Feature - Device Provenance Tracking}\label{cap:ec-7f4beb56}
\vspace{-1em}
\begin{table}[htbp]
  \centering
  \footnotesize
  \begin{tabularx}{1\textwidth}{|l|l|X|}
    \hline
    \textbf{Metadata Name} & \textbf{Metadata Value} & \textbf{Description} \\
    \hline
    \textbf{Reference Id} & \texttt{EC-7F4BEB56} & \\
    \hline
    \textbf{Status} & Derived candidate & \\
    \hline
    \textbf{Ancestry} &  & \\
    \hline
    \textbf{Source} &  & \\
    \hline
    \textbf{Enabled By} &  & \\
    \hline
    \textbf{Supports} &  & \\
    \hline
    \textbf{Derived From} &  & \\
    \hline
    \textbf{Reviewed Status} & Inspect recurrence for true semantic equivalence; retain, split, rename or reject. & \\
    \hline
  \end{tabularx}
\end{table}

This candidate enterprise Feature abstracts a behaviour recurring in 4 distinct BEC CT3 contexts. The source-specific BEC Features remain in place; this entry does not assert a common software service or an identical clinical meaning.

Illustrative originating CT3 contexts: Device Association, Device Connectivity, Device Definition, Device Instance.

\featuretable{--- & --- & \emph{Not yet captured} & \emph{Enterprise decomposition subject to review} \\}
\subsubsection{Feature - Diagnostic Status Tracking}\label{cap:ec-95e97505}
\vspace{-1em}
\begin{table}[htbp]
  \centering
  \footnotesize
  \begin{tabularx}{1\textwidth}{|l|l|X|}
    \hline
    \textbf{Metadata Name} & \textbf{Metadata Value} & \textbf{Description} \\
    \hline
    \textbf{Reference Id} & \texttt{EC-95E97505} & \\
    \hline
    \textbf{Status} & Derived candidate & \\
    \hline
    \textbf{Ancestry} &  & \\
    \hline
    \textbf{Source} &  & \\
    \hline
    \textbf{Enabled By} &  & \\
    \hline
    \textbf{Supports} &  & \\
    \hline
    \textbf{Derived From} &  & \\
    \hline
    \textbf{Reviewed Status} & Inspect recurrence for true semantic equivalence; retain, split, rename or reject. & \\
    \hline
  \end{tabularx}
\end{table}

This candidate enterprise Feature abstracts a behaviour recurring in 2 distinct BEC CT3 contexts. The source-specific BEC Features remain in place; this entry does not assert a common software service or an identical clinical meaning.

Illustrative originating CT3 contexts: Diagnostic Service Integration, Diagnostic Workflow Management.

\featuretable{--- & --- & \emph{Not yet captured} & \emph{Enterprise decomposition subject to review} \\}
\subsubsection{Feature - Medication Activity Tracking}\label{cap:ec-67a2e372}
\vspace{-1em}
\begin{table}[htbp]
  \centering
  \footnotesize
  \begin{tabularx}{1\textwidth}{|l|l|X|}
    \hline
    \textbf{Metadata Name} & \textbf{Metadata Value} & \textbf{Description} \\
    \hline
    \textbf{Reference Id} & \texttt{EC-67A2E372} & \\
    \hline
    \textbf{Status} & Derived candidate & \\
    \hline
    \textbf{Ancestry} &  & \\
    \hline
    \textbf{Source} &  & \\
    \hline
    \textbf{Enabled By} &  & \\
    \hline
    \textbf{Supports} &  & \\
    \hline
    \textbf{Derived From} &  & \\
    \hline
    \textbf{Reviewed Status} & Inspect recurrence for true semantic equivalence; retain, split, rename or reject. & \\
    \hline
  \end{tabularx}
\end{table}

This candidate enterprise Feature abstracts a behaviour recurring in 2 distinct BEC CT3 contexts. The source-specific BEC Features remain in place; this entry does not assert a common software service or an identical clinical meaning.

Illustrative originating CT3 contexts: Medication Order \& Dispense Management, Medication Therapy Provision and Monitoring.

\featuretable{--- & --- & \emph{Not yet captured} & \emph{Enterprise decomposition subject to review} \\}
\subsubsection{Feature - Request Fulfilment Tracking}\label{cap:ec-a191a2bc}
\vspace{-1em}
\begin{table}[htbp]
  \centering
  \footnotesize
  \begin{tabularx}{1\textwidth}{|l|l|X|}
    \hline
    \textbf{Metadata Name} & \textbf{Metadata Value} & \textbf{Description} \\
    \hline
    \textbf{Reference Id} & \texttt{EC-A191A2BC} & \\
    \hline
    \textbf{Status} & Derived candidate & \\
    \hline
    \textbf{Ancestry} &  & \\
    \hline
    \textbf{Source} &  & \\
    \hline
    \textbf{Enabled By} &  & \\
    \hline
    \textbf{Supports} &  & \\
    \hline
    \textbf{Derived From} &  & \\
    \hline
    \textbf{Reviewed Status} & Inspect recurrence for true semantic equivalence; retain, split, rename or reject. & \\
    \hline
  \end{tabularx}
\end{table}

This candidate enterprise Feature abstracts a behaviour recurring in 2 distinct BEC CT3 contexts. The source-specific BEC Features remain in place; this entry does not assert a common software service or an identical clinical meaning.

Illustrative originating CT3 contexts: Closed-Loop Order Progression, Medication Order \& Dispense Management.

\featuretable{--- & --- & \emph{Not yet captured} & \emph{Enterprise decomposition subject to review} \\}

\chapter{CT1 - Governance, Trust and Assurance}\label{cap:ec-8b8bc538}
\noindent\textbf{Reference ID:} \texttt{EC-8B8BC538}\quad \textbf{Status:} Derived candidate\par
Intent. Governance, Trust and Assurance is a candidate horizontally reusable architectural capability associated with BEC recurrence. It is framed around recurring enabling behaviour, not a mandatory centralised service or a particular technology.

Boundary. A common architectural capability does not remove the originating Features from their healthcare contexts. Similar wording does not establish equivalence or identical implementation.

\section{CT2 - Governed Control}\label{cap:ec-d9af0b6f}
\vspace{-0.5em}
\begin{table}[htbp]
  \centering
  \footnotesize
  \begin{tabularx}{1\textwidth}{|l|l|X|}
    \hline
    \multicolumn{3}{c}{\textbf{Capability Metadata}} \\
    \hline
    \textbf{Name} & \textbf{Value} & \textbf{Comment and/or Description} \\
    \hline
    \textbf{Reference Id} & \texttt{EC-D9AF0B6F} & \\
    \hline
    \textbf{Status} & Derived candidate & \\
    \hline
    \textbf{Ancestry} &  & \\
    \hline
    \textbf{Source} &  & \\
    \hline
    \textbf{Enabled By} &  & \\
    \hline
    \textbf{Supports} &  & \\
    \hline
    \textbf{Derived From} &  & \\
    \hline
    \textbf{Reviewed Status} &  & \\
    \hline
  \end{tabularx}
\end{table}

\begin{table}[htbp]
  \centering
  \footnotesize
  \begin{tabularx}{1\textwidth}{|l|l|X|}
    \hline
    \multicolumn{3}{c}{\textbf{Contained Capabilities}} \\
    \hline
    \textbf{Capability Name} & \textbf{Capability Id} & \textbf{Comment} \\
    \hline
    Governed Control Enablement & \texttt{cap:ec-e37e56ac} (at \ref{cap:ec-e37e56ac}) & \\
    \hline
  \end{tabularx}
\end{table}

Intent. Governed Control is a candidate horizontally reusable architectural capability associated with Governance, Trust and Assurance. It is framed around recurring enabling behaviour, not a mandatory centralised service or a particular technology.

Boundary. A common architectural capability does not remove the originating Features from their healthcare contexts. Similar wording does not establish equivalence or identical implementation.

\subsection{CT3 - Governed Control Enablement}\label{cap:ec-e37e56ac}
\vspace{-0.5em}
\begin{table}[htbp]
  \centering
  \footnotesize
  \begin{tabularx}{1\textwidth}{|l|l|X|}
    \hline
    \multicolumn{3}{c}{\textbf{Capability Metadata}} \\
    \hline
    \textbf{Name} & \textbf{Value} & \textbf{Comment and/or Description} \\
    \hline
    \textbf{Reference Id} & \texttt{EC-E37E56AC} & \\
    \hline
    \textbf{Status} & Derived candidate & \\
    \hline
    \textbf{Ancestry} &  & \\
    \hline
    \textbf{Source} &  & \\
    \hline
    \textbf{Enabled By} &  & \\
    \hline
    \textbf{Supports} &  & \\
    \hline
    \textbf{Derived From} &  & \\
    \hline
    \textbf{Reviewed Status} &  & \\
    \hline
  \end{tabularx}
\end{table}

\begin{table}[htbp]
  \centering
  \footnotesize
  \begin{tabularx}{1\textwidth}{|l|l|X|}
    \hline
    \multicolumn{3}{c}{\textbf{Contained Features}} \\
    \hline
    \textbf{Feature Name} & \textbf{Feature Id} & \textbf{Comment} \\
    \hline
    Eligibility Validation & \texttt{cap:ec-5948bfa2} (at \ref{cap:ec-5948bfa2}) & \\
    \hline
    Service Eligibility Management & \texttt{cap:ec-d98b7aae} (at \ref{cap:ec-d98b7aae}) & \\
    \hline
    Policy Evaluation & \texttt{cap:ec-840c09a8} (at \ref{cap:ec-840c09a8}) & \\
    \hline
  \end{tabularx}
\end{table}

Intent. Governed Control Enablement is a candidate horizontally reusable architectural capability associated with Governed Control. It is framed around recurring enabling behaviour, not a mandatory centralised service or a particular technology.

Boundary. A common architectural capability does not remove the originating Features from their healthcare contexts. Similar wording does not establish equivalence or identical implementation.

\subsubsection{Feature - Eligibility Validation}\label{cap:ec-5948bfa2}
\vspace{-1em}
\begin{table}[htbp]
  \centering
  \footnotesize
  \begin{tabularx}{1\textwidth}{|l|l|X|}
    \hline
    \textbf{Metadata Name} & \textbf{Metadata Value} & \textbf{Description} \\
    \hline
    \textbf{Reference Id} & \texttt{EC-5948BFA2} & \\
    \hline
    \textbf{Status} & Derived candidate & \\
    \hline
    \textbf{Ancestry} &  & \\
    \hline
    \textbf{Source} &  & \\
    \hline
    \textbf{Enabled By} &  & \\
    \hline
    \textbf{Supports} &  & \\
    \hline
    \textbf{Derived From} &  & \\
    \hline
    \textbf{Reviewed Status} & Inspect recurrence for true semantic equivalence; retain, split, rename or reject. & \\
    \hline
  \end{tabularx}
\end{table}

This candidate enterprise Feature abstracts a behaviour recurring in 11 distinct BEC CT3 contexts. The source-specific BEC Features remain in place; this entry does not assert a common software service or an identical clinical meaning.

Illustrative originating CT3 contexts: Acute Care Monitoring \& Event Coordination, Community Care Coordination, Discharge Progression \& Coordination, ED Operational Coordination, Mobile Worker Dispatch \& Safety, Multi-Work-Unit Orchestration, Operational Work Assignment \& Supervision, Primary Care Coordination, Session \& Queue Coordination, Specialized Equipment Coordination, Ward Logistics Coordination.

\featuretable{--- & --- & \emph{Not yet captured} & \emph{Enterprise decomposition subject to review} \\}
\subsubsection{Feature - Service Eligibility Management}\label{cap:ec-d98b7aae}
\vspace{-1em}
\begin{table}[htbp]
  \centering
  \footnotesize
  \begin{tabularx}{1\textwidth}{|l|l|X|}
    \hline
    \textbf{Metadata Name} & \textbf{Metadata Value} & \textbf{Description} \\
    \hline
    \textbf{Reference Id} & \texttt{EC-D98B7AAE} & \\
    \hline
    \textbf{Status} & Derived candidate & \\
    \hline
    \textbf{Ancestry} &  & \\
    \hline
    \textbf{Source} &  & \\
    \hline
    \textbf{Enabled By} &  & \\
    \hline
    \textbf{Supports} &  & \\
    \hline
    \textbf{Derived From} &  & \\
    \hline
    \textbf{Reviewed Status} & Inspect recurrence for true semantic equivalence; retain, split, rename or reject. & \\
    \hline
  \end{tabularx}
\end{table}

This candidate enterprise Feature abstracts a behaviour recurring in 5 distinct BEC CT3 contexts. The source-specific BEC Features remain in place; this entry does not assert a common software service or an identical clinical meaning.

Illustrative originating CT3 contexts: Diagnostic Service Integration, Healthcare Service Definition, Service Availability Definition, Service Eligibility Definition, Service Provision Relationship.

\featuretable{--- & --- & \emph{Not yet captured} & \emph{Enterprise decomposition subject to review} \\}
\subsubsection{Feature - Policy Evaluation}\label{cap:ec-840c09a8}
\vspace{-1em}
\begin{table}[htbp]
  \centering
  \footnotesize
  \begin{tabularx}{1\textwidth}{|l|l|X|}
    \hline
    \textbf{Metadata Name} & \textbf{Metadata Value} & \textbf{Description} \\
    \hline
    \textbf{Reference Id} & \texttt{EC-840C09A8} & \\
    \hline
    \textbf{Status} & Derived candidate & \\
    \hline
    \textbf{Ancestry} &  & \\
    \hline
    \textbf{Source} &  & \\
    \hline
    \textbf{Enabled By} &  & \\
    \hline
    \textbf{Supports} &  & \\
    \hline
    \textbf{Derived From} &  & \\
    \hline
    \textbf{Reviewed Status} & Inspect recurrence for true semantic equivalence; retain, split, rename or reject. & \\
    \hline
  \end{tabularx}
\end{table}

This candidate enterprise Feature abstracts a behaviour recurring in 2 distinct BEC CT3 contexts. The source-specific BEC Features remain in place; this entry does not assert a common software service or an identical clinical meaning.

Illustrative originating CT3 contexts: Client Privacy, Policy \& Security Governance.

\featuretable{--- & --- & \emph{Not yet captured} & \emph{Enterprise decomposition subject to review} \\}

\section{CT2 - Provenance and Accountability}\label{cap:ec-df0619b8}
\vspace{-0.5em}
\begin{table}[htbp]
  \centering
  \footnotesize
  \begin{tabularx}{1\textwidth}{|l|l|X|}
    \hline
    \multicolumn{3}{c}{\textbf{Capability Metadata}} \\
    \hline
    \textbf{Name} & \textbf{Value} & \textbf{Comment and/or Description} \\
    \hline
    \textbf{Reference Id} & \texttt{EC-DF0619B8} & \\
    \hline
    \textbf{Status} & Derived candidate & \\
    \hline
    \textbf{Ancestry} &  & \\
    \hline
    \textbf{Source} &  & \\
    \hline
    \textbf{Enabled By} &  & \\
    \hline
    \textbf{Supports} &  & \\
    \hline
    \textbf{Derived From} &  & \\
    \hline
    \textbf{Reviewed Status} &  & \\
    \hline
  \end{tabularx}
\end{table}

\begin{table}[htbp]
  \centering
  \footnotesize
  \begin{tabularx}{1\textwidth}{|l|l|X|}
    \hline
    \multicolumn{3}{c}{\textbf{Contained Capabilities}} \\
    \hline
    \textbf{Capability Name} & \textbf{Capability Id} & \textbf{Comment} \\
    \hline
    Provenance and Accountability Enablement & \texttt{cap:ec-3987b790} (at \ref{cap:ec-3987b790}) & \\
    \hline
  \end{tabularx}
\end{table}

Intent. Provenance and Accountability is a candidate horizontally reusable architectural capability associated with Governance, Trust and Assurance. It is framed around recurring enabling behaviour, not a mandatory centralised service or a particular technology.

Boundary. A common architectural capability does not remove the originating Features from their healthcare contexts. Similar wording does not establish equivalence or identical implementation.

\subsection{CT3 - Provenance and Accountability Enablement}\label{cap:ec-3987b790}
\vspace{-0.5em}
\begin{table}[htbp]
  \centering
  \footnotesize
  \begin{tabularx}{1\textwidth}{|l|l|X|}
    \hline
    \multicolumn{3}{c}{\textbf{Capability Metadata}} \\
    \hline
    \textbf{Name} & \textbf{Value} & \textbf{Comment and/or Description} \\
    \hline
    \textbf{Reference Id} & \texttt{EC-3987B790} & \\
    \hline
    \textbf{Status} & Derived candidate & \\
    \hline
    \textbf{Ancestry} &  & \\
    \hline
    \textbf{Source} &  & \\
    \hline
    \textbf{Enabled By} &  & \\
    \hline
    \textbf{Supports} &  & \\
    \hline
    \textbf{Derived From} &  & \\
    \hline
    \textbf{Reviewed Status} &  & \\
    \hline
  \end{tabularx}
\end{table}

\begin{table}[htbp]
  \centering
  \footnotesize
  \begin{tabularx}{1\textwidth}{|l|l|X|}
    \hline
    \multicolumn{3}{c}{\textbf{Contained Features}} \\
    \hline
    \textbf{Feature Name} & \textbf{Feature Id} & \textbf{Comment} \\
    \hline
    Device Provenance Tracking & \texttt{cap:ec-e41c93c7} (at \ref{cap:ec-e41c93c7}) & \\
    \hline
    Information Provenance & \texttt{cap:ec-313e85ab} (at \ref{cap:ec-313e85ab}) & \\
    \hline
    Decision Traceability & \texttt{cap:ec-829f0bf7} (at \ref{cap:ec-829f0bf7}) & \\
    \hline
    Diagnostic Provenance Capture & \texttt{cap:ec-e9aa136d} (at \ref{cap:ec-e9aa136d}) & \\
    \hline
    Evidence Association & \texttt{cap:ec-7233f33e} (at \ref{cap:ec-7233f33e}) & \\
    \hline
  \end{tabularx}
\end{table}

Intent. Provenance and Accountability Enablement is a candidate horizontally reusable architectural capability associated with Provenance and Accountability. It is framed around recurring enabling behaviour, not a mandatory centralised service or a particular technology.

Boundary. A common architectural capability does not remove the originating Features from their healthcare contexts. Similar wording does not establish equivalence or identical implementation.

\subsubsection{Feature - Device Provenance Tracking}\label{cap:ec-e41c93c7}
\vspace{-1em}
\begin{table}[htbp]
  \centering
  \footnotesize
  \begin{tabularx}{1\textwidth}{|l|l|X|}
    \hline
    \textbf{Metadata Name} & \textbf{Metadata Value} & \textbf{Description} \\
    \hline
    \textbf{Reference Id} & \texttt{EC-E41C93C7} & \\
    \hline
    \textbf{Status} & Derived candidate & \\
    \hline
    \textbf{Ancestry} &  & \\
    \hline
    \textbf{Source} &  & \\
    \hline
    \textbf{Enabled By} &  & \\
    \hline
    \textbf{Supports} &  & \\
    \hline
    \textbf{Derived From} &  & \\
    \hline
    \textbf{Reviewed Status} & Inspect recurrence for true semantic equivalence; retain, split, rename or reject. & \\
    \hline
  \end{tabularx}
\end{table}

This candidate enterprise Feature abstracts a behaviour recurring in 4 distinct BEC CT3 contexts. The source-specific BEC Features remain in place; this entry does not assert a common software service or an identical clinical meaning.

Illustrative originating CT3 contexts: Device Association, Device Connectivity, Device Definition, Device Instance.

\featuretable{--- & --- & \emph{Not yet captured} & \emph{Enterprise decomposition subject to review} \\}
\subsubsection{Feature - Information Provenance}\label{cap:ec-313e85ab}
\vspace{-1em}
\begin{table}[htbp]
  \centering
  \footnotesize
  \begin{tabularx}{1\textwidth}{|l|l|X|}
    \hline
    \textbf{Metadata Name} & \textbf{Metadata Value} & \textbf{Description} \\
    \hline
    \textbf{Reference Id} & \texttt{EC-313E85AB} & \\
    \hline
    \textbf{Status} & Derived candidate & \\
    \hline
    \textbf{Ancestry} &  & \\
    \hline
    \textbf{Source} &  & \\
    \hline
    \textbf{Enabled By} &  & \\
    \hline
    \textbf{Supports} &  & \\
    \hline
    \textbf{Derived From} &  & \\
    \hline
    \textbf{Reviewed Status} & Inspect recurrence for true semantic equivalence; retain, split, rename or reject. & \\
    \hline
  \end{tabularx}
\end{table}

This candidate enterprise Feature abstracts a behaviour recurring in 3 distinct BEC CT3 contexts. The source-specific BEC Features remain in place; this entry does not assert a common software service or an identical clinical meaning.

Illustrative originating CT3 contexts: Canonical Semantics \& Schema Governance, Multi-Model Information Exchange, Terminology \& Reference Mapping.

\featuretable{--- & --- & \emph{Not yet captured} & \emph{Enterprise decomposition subject to review} \\}
\subsubsection{Feature - Decision Traceability}\label{cap:ec-829f0bf7}
\vspace{-1em}
\begin{table}[htbp]
  \centering
  \footnotesize
  \begin{tabularx}{1\textwidth}{|l|l|X|}
    \hline
    \textbf{Metadata Name} & \textbf{Metadata Value} & \textbf{Description} \\
    \hline
    \textbf{Reference Id} & \texttt{EC-829F0BF7} & \\
    \hline
    \textbf{Status} & Derived candidate & \\
    \hline
    \textbf{Ancestry} &  & \\
    \hline
    \textbf{Source} &  & \\
    \hline
    \textbf{Enabled By} &  & \\
    \hline
    \textbf{Supports} &  & \\
    \hline
    \textbf{Derived From} &  & \\
    \hline
    \textbf{Reviewed Status} & Inspect recurrence for true semantic equivalence; retain, split, rename or reject. & \\
    \hline
  \end{tabularx}
\end{table}

This candidate enterprise Feature abstracts a behaviour recurring in 2 distinct BEC CT3 contexts. The source-specific BEC Features remain in place; this entry does not assert a common software service or an identical clinical meaning.

Illustrative originating CT3 contexts: Client Privacy, Policy \& Security Governance.

\featuretable{--- & --- & \emph{Not yet captured} & \emph{Enterprise decomposition subject to review} \\}
\subsubsection{Feature - Diagnostic Provenance Capture}\label{cap:ec-e9aa136d}
\vspace{-1em}
\begin{table}[htbp]
  \centering
  \footnotesize
  \begin{tabularx}{1\textwidth}{|l|l|X|}
    \hline
    \textbf{Metadata Name} & \textbf{Metadata Value} & \textbf{Description} \\
    \hline
    \textbf{Reference Id} & \texttt{EC-E9AA136D} & \\
    \hline
    \textbf{Status} & Derived candidate & \\
    \hline
    \textbf{Ancestry} &  & \\
    \hline
    \textbf{Source} &  & \\
    \hline
    \textbf{Enabled By} &  & \\
    \hline
    \textbf{Supports} &  & \\
    \hline
    \textbf{Derived From} &  & \\
    \hline
    \textbf{Reviewed Status} & Inspect recurrence for true semantic equivalence; retain, split, rename or reject. & \\
    \hline
  \end{tabularx}
\end{table}

This candidate enterprise Feature abstracts a behaviour recurring in 2 distinct BEC CT3 contexts. The source-specific BEC Features remain in place; this entry does not assert a common software service or an identical clinical meaning.

Illustrative originating CT3 contexts: Diagnostic Service Integration, Diagnostic Workflow Management.

\featuretable{--- & --- & \emph{Not yet captured} & \emph{Enterprise decomposition subject to review} \\}
\subsubsection{Feature - Evidence Association}\label{cap:ec-7233f33e}
\vspace{-1em}
\begin{table}[htbp]
  \centering
  \footnotesize
  \begin{tabularx}{1\textwidth}{|l|l|X|}
    \hline
    \textbf{Metadata Name} & \textbf{Metadata Value} & \textbf{Description} \\
    \hline
    \textbf{Reference Id} & \texttt{EC-7233F33E} & \\
    \hline
    \textbf{Status} & Derived candidate & \\
    \hline
    \textbf{Ancestry} &  & \\
    \hline
    \textbf{Source} &  & \\
    \hline
    \textbf{Enabled By} &  & \\
    \hline
    \textbf{Supports} &  & \\
    \hline
    \textbf{Derived From} &  & \\
    \hline
    \textbf{Reviewed Status} & Inspect recurrence for true semantic equivalence; retain, split, rename or reject. & \\
    \hline
  \end{tabularx}
\end{table}

This candidate enterprise Feature abstracts a behaviour recurring in 2 distinct BEC CT3 contexts. The source-specific BEC Features remain in place; this entry does not assert a common software service or an identical clinical meaning.

Illustrative originating CT3 contexts: Client Privacy, Policy \& Security Governance.

\featuretable{--- & --- & \emph{Not yet captured} & \emph{Enterprise decomposition subject to review} \\}

\section{CT2 - Assurance and Resilience}\label{cap:ec-74bb4ad4}
\vspace{-0.5em}
\begin{table}[htbp]
  \centering
  \footnotesize
  \begin{tabularx}{1\textwidth}{|l|l|X|}
    \hline
    \multicolumn{3}{c}{\textbf{Capability Metadata}} \\
    \hline
    \textbf{Name} & \textbf{Value} & \textbf{Comment and/or Description} \\
    \hline
    \textbf{Reference Id} & \texttt{EC-74BB4AD4} & \\
    \hline
    \textbf{Status} & Derived candidate & \\
    \hline
    \textbf{Ancestry} &  & \\
    \hline
    \textbf{Source} &  & \\
    \hline
    \textbf{Enabled By} &  & \\
    \hline
    \textbf{Supports} &  & \\
    \hline
    \textbf{Derived From} &  & \\
    \hline
    \textbf{Reviewed Status} &  & \\
    \hline
  \end{tabularx}
\end{table}

\begin{table}[htbp]
  \centering
  \footnotesize
  \begin{tabularx}{1\textwidth}{|l|l|X|}
    \hline
    \multicolumn{3}{c}{\textbf{Contained Capabilities}} \\
    \hline
    \textbf{Capability Name} & \textbf{Capability Id} & \textbf{Comment} \\
    \hline
    Assurance and Resilience Enablement & \texttt{cap:ec-b0d5a003} (at \ref{cap:ec-b0d5a003}) & \\
    \hline
  \end{tabularx}
\end{table}

Intent. Assurance and Resilience is a candidate horizontally reusable architectural capability associated with Governance, Trust and Assurance. It is framed around recurring enabling behaviour, not a mandatory centralised service or a particular technology.

Boundary. A common architectural capability does not remove the originating Features from their healthcare contexts. Similar wording does not establish equivalence or identical implementation.

\subsection{CT3 - Assurance and Resilience Enablement}\label{cap:ec-b0d5a003}
\vspace{-0.5em}
\begin{table}[htbp]
  \centering
  \footnotesize
  \begin{tabularx}{1\textwidth}{|l|l|X|}
    \hline
    \multicolumn{3}{c}{\textbf{Capability Metadata}} \\
    \hline
    \textbf{Name} & \textbf{Value} & \textbf{Comment and/or Description} \\
    \hline
    \textbf{Reference Id} & \texttt{EC-B0D5A003} & \\
    \hline
    \textbf{Status} & Derived candidate & \\
    \hline
    \textbf{Ancestry} &  & \\
    \hline
    \textbf{Source} &  & \\
    \hline
    \textbf{Enabled By} &  & \\
    \hline
    \textbf{Supports} &  & \\
    \hline
    \textbf{Derived From} &  & \\
    \hline
    \textbf{Reviewed Status} &  & \\
    \hline
  \end{tabularx}
\end{table}

\begin{table}[htbp]
  \centering
  \footnotesize
  \begin{tabularx}{1\textwidth}{|l|l|X|}
    \hline
    \multicolumn{3}{c}{\textbf{Contained Features}} \\
    \hline
    \textbf{Feature Name} & \textbf{Feature Id} & \textbf{Comment} \\
    \hline
    Exception Management & \texttt{cap:ec-50be2841} (at \ref{cap:ec-50be2841}) & \\
    \hline
    State Monitoring & \texttt{cap:ec-0f5671f1} (at \ref{cap:ec-0f5671f1}) & \\
    \hline
    Outcome Reporting & \texttt{cap:ec-09035a01} (at \ref{cap:ec-09035a01}) & \\
    \hline
    Exception Reporting & \texttt{cap:ec-20985ee7} (at \ref{cap:ec-20985ee7}) & \\
    \hline
    Medication Safety Context Validation & \texttt{cap:ec-34e6439f} (at \ref{cap:ec-34e6439f}) & \\
    \hline
    Request Exception Management & \texttt{cap:ec-85d9e54b} (at \ref{cap:ec-85d9e54b}) & \\
    \hline
  \end{tabularx}
\end{table}

Intent. Assurance and Resilience Enablement is a candidate horizontally reusable architectural capability associated with Assurance and Resilience. It is framed around recurring enabling behaviour, not a mandatory centralised service or a particular technology.

Boundary. A common architectural capability does not remove the originating Features from their healthcare contexts. Similar wording does not establish equivalence or identical implementation.

\subsubsection{Feature - Exception Management}\label{cap:ec-50be2841}
\vspace{-1em}
\begin{table}[htbp]
  \centering
  \footnotesize
  \begin{tabularx}{1\textwidth}{|l|l|X|}
    \hline
    \textbf{Metadata Name} & \textbf{Metadata Value} & \textbf{Description} \\
    \hline
    \textbf{Reference Id} & \texttt{EC-50BE2841} & \\
    \hline
    \textbf{Status} & Derived candidate & \\
    \hline
    \textbf{Ancestry} &  & \\
    \hline
    \textbf{Source} &  & \\
    \hline
    \textbf{Enabled By} &  & \\
    \hline
    \textbf{Supports} &  & \\
    \hline
    \textbf{Derived From} &  & \\
    \hline
    \textbf{Reviewed Status} & Inspect recurrence for true semantic equivalence; retain, split, rename or reject. & \\
    \hline
  \end{tabularx}
\end{table}

This candidate enterprise Feature abstracts a behaviour recurring in 35 distinct BEC CT3 contexts. The source-specific BEC Features remain in place; this entry does not assert a common software service or an identical clinical meaning.

Illustrative originating CT3 contexts: Acute Care Monitoring \& Event Coordination, Bed State \& Turnover Management, Boundary Protocol Mediation, Care-Place Definition, Client Intake \& Demographic Update, Client Relationship, Clinical Documentation Lifecycle, Community Care Coordination, Credential Verification Checkpoint, Decoupled Human Interaction \& Dashboards, Discharge Progression \& Coordination, ED Operational Coordination, and others.

\featuretable{--- & --- & \emph{Not yet captured} & \emph{Enterprise decomposition subject to review} \\}
\subsubsection{Feature - State Monitoring}\label{cap:ec-0f5671f1}
\vspace{-1em}
\begin{table}[htbp]
  \centering
  \footnotesize
  \begin{tabularx}{1\textwidth}{|l|l|X|}
    \hline
    \textbf{Metadata Name} & \textbf{Metadata Value} & \textbf{Description} \\
    \hline
    \textbf{Reference Id} & \texttt{EC-0F5671F1} & \\
    \hline
    \textbf{Status} & Derived candidate & \\
    \hline
    \textbf{Ancestry} &  & \\
    \hline
    \textbf{Source} &  & \\
    \hline
    \textbf{Enabled By} &  & \\
    \hline
    \textbf{Supports} &  & \\
    \hline
    \textbf{Derived From} &  & \\
    \hline
    \textbf{Reviewed Status} & Inspect recurrence for true semantic equivalence; retain, split, rename or reject. & \\
    \hline
  \end{tabularx}
\end{table}

This candidate enterprise Feature abstracts a behaviour recurring in 24 distinct BEC CT3 contexts. The source-specific BEC Features remain in place; this entry does not assert a common software service or an identical clinical meaning.

Illustrative originating CT3 contexts: Bed State \& Turnover Management, Boundary Protocol Mediation, Care-Place Definition, Client Intake \& Demographic Update, Client Relationship, Clinical Documentation Lifecycle, Credential Verification Checkpoint, Decoupled Human Interaction \& Dashboards, Emergency Care Progression, Facility Demand \& Capacity Telemetry, Governed Query \& Search, Health Product Definition, and others.

\featuretable{--- & --- & \emph{Not yet captured} & \emph{Enterprise decomposition subject to review} \\}
\subsubsection{Feature - Outcome Reporting}\label{cap:ec-09035a01}
\vspace{-1em}
\begin{table}[htbp]
  \centering
  \footnotesize
  \begin{tabularx}{1\textwidth}{|l|l|X|}
    \hline
    \textbf{Metadata Name} & \textbf{Metadata Value} & \textbf{Description} \\
    \hline
    \textbf{Reference Id} & \texttt{EC-09035A01} & \\
    \hline
    \textbf{Status} & Derived candidate & \\
    \hline
    \textbf{Ancestry} &  & \\
    \hline
    \textbf{Source} &  & \\
    \hline
    \textbf{Enabled By} &  & \\
    \hline
    \textbf{Supports} &  & \\
    \hline
    \textbf{Derived From} &  & \\
    \hline
    \textbf{Reviewed Status} & Inspect recurrence for true semantic equivalence; retain, split, rename or reject. & \\
    \hline
  \end{tabularx}
\end{table}

This candidate enterprise Feature abstracts a behaviour recurring in 11 distinct BEC CT3 contexts. The source-specific BEC Features remain in place; this entry does not assert a common software service or an identical clinical meaning.

Illustrative originating CT3 contexts: Acute Care Monitoring \& Event Coordination, Community Care Coordination, Discharge Progression \& Coordination, ED Operational Coordination, Mobile Worker Dispatch \& Safety, Multi-Work-Unit Orchestration, Operational Work Assignment \& Supervision, Primary Care Coordination, Session \& Queue Coordination, Specialized Equipment Coordination, Ward Logistics Coordination.

\featuretable{--- & --- & \emph{Not yet captured} & \emph{Enterprise decomposition subject to review} \\}
\subsubsection{Feature - Exception Reporting}\label{cap:ec-20985ee7}
\vspace{-1em}
\begin{table}[htbp]
  \centering
  \footnotesize
  \begin{tabularx}{1\textwidth}{|l|l|X|}
    \hline
    \textbf{Metadata Name} & \textbf{Metadata Value} & \textbf{Description} \\
    \hline
    \textbf{Reference Id} & \texttt{EC-20985EE7} & \\
    \hline
    \textbf{Status} & Derived candidate & \\
    \hline
    \textbf{Ancestry} &  & \\
    \hline
    \textbf{Source} &  & \\
    \hline
    \textbf{Enabled By} &  & \\
    \hline
    \textbf{Supports} &  & \\
    \hline
    \textbf{Derived From} &  & \\
    \hline
    \textbf{Reviewed Status} & Inspect recurrence for true semantic equivalence; retain, split, rename or reject. & \\
    \hline
  \end{tabularx}
\end{table}

This candidate enterprise Feature abstracts a behaviour recurring in 2 distinct BEC CT3 contexts. The source-specific BEC Features remain in place; this entry does not assert a common software service or an identical clinical meaning.

Illustrative originating CT3 contexts: Client Privacy, Policy \& Security Governance.

\featuretable{--- & --- & \emph{Not yet captured} & \emph{Enterprise decomposition subject to review} \\}
\subsubsection{Feature - Medication Safety Context Validation}\label{cap:ec-34e6439f}
\vspace{-1em}
\begin{table}[htbp]
  \centering
  \footnotesize
  \begin{tabularx}{1\textwidth}{|l|l|X|}
    \hline
    \textbf{Metadata Name} & \textbf{Metadata Value} & \textbf{Description} \\
    \hline
    \textbf{Reference Id} & \texttt{EC-34E6439F} & \\
    \hline
    \textbf{Status} & Derived candidate & \\
    \hline
    \textbf{Ancestry} &  & \\
    \hline
    \textbf{Source} &  & \\
    \hline
    \textbf{Enabled By} &  & \\
    \hline
    \textbf{Supports} &  & \\
    \hline
    \textbf{Derived From} &  & \\
    \hline
    \textbf{Reviewed Status} & Inspect recurrence for true semantic equivalence; retain, split, rename or reject. & \\
    \hline
  \end{tabularx}
\end{table}

This candidate enterprise Feature abstracts a behaviour recurring in 2 distinct BEC CT3 contexts. The source-specific BEC Features remain in place; this entry does not assert a common software service or an identical clinical meaning.

Illustrative originating CT3 contexts: Medication Order \& Dispense Management, Medication Therapy Provision and Monitoring.

\featuretable{--- & --- & \emph{Not yet captured} & \emph{Enterprise decomposition subject to review} \\}
\subsubsection{Feature - Request Exception Management}\label{cap:ec-85d9e54b}
\vspace{-1em}
\begin{table}[htbp]
  \centering
  \footnotesize
  \begin{tabularx}{1\textwidth}{|l|l|X|}
    \hline
    \textbf{Metadata Name} & \textbf{Metadata Value} & \textbf{Description} \\
    \hline
    \textbf{Reference Id} & \texttt{EC-85D9E54B} & \\
    \hline
    \textbf{Status} & Derived candidate & \\
    \hline
    \textbf{Ancestry} &  & \\
    \hline
    \textbf{Source} &  & \\
    \hline
    \textbf{Enabled By} &  & \\
    \hline
    \textbf{Supports} &  & \\
    \hline
    \textbf{Derived From} &  & \\
    \hline
    \textbf{Reviewed Status} & Inspect recurrence for true semantic equivalence; retain, split, rename or reject. & \\
    \hline
  \end{tabularx}
\end{table}

This candidate enterprise Feature abstracts a behaviour recurring in 2 distinct BEC CT3 contexts. The source-specific BEC Features remain in place; this entry does not assert a common software service or an identical clinical meaning.

Illustrative originating CT3 contexts: Closed-Loop Order Progression, Medication Order \& Dispense Management.

\featuretable{--- & --- & \emph{Not yet captured} & \emph{Enterprise decomposition subject to review} \\}

\chapter{CT1 - Existing Enterprise Capabilities -- Retained for Reconciliation}\label{cap:ec-1f33c0e8}
\noindent\textbf{Reference ID:} \texttt{EC-1F33C0E8}\quad \textbf{Status:} Existing catalogue; coexistence retained\par
The following entries reproduce the names and architectural scope statements represented in the current Enterprise Capability catalogue. They are intentionally not merged with the derived hierarchy, even where there is an obvious overlap.

\section{CT2 - EC-01 Managed Entity \& Relationship}\label{cap:ec-3bff62ca}
\vspace{-0.5em}
\begin{table}[htbp]
  \centering
  \footnotesize
  \begin{tabularx}{1\textwidth}{|l|l|X|}
    \hline
    \multicolumn{3}{c}{\textbf{Capability Metadata}} \\
    \hline
    \textbf{Name} & \textbf{Value} & \textbf{Comment and/or Description} \\
    \hline
    \textbf{Reference Id} & \texttt{EC-3BFF62CA} & \\
    \hline
    \textbf{Status} & Existing canonical EC & \\
    \hline
    \textbf{Ancestry} &  & \\
    \hline
    \textbf{Source} & Enterprise Capability catalogue & \\
    \hline
    \textbf{Enabled By} &  & \\
    \hline
    \textbf{Supports} &  & \\
    \hline
    \textbf{Derived From} &  & \\
    \hline
    \textbf{Reviewed Status} & Compare against derived Enterprise CT3 and Feature families; do not suppress conflicts. & \\
    \hline
  \end{tabularx}
\end{table}

\begin{table}[htbp]
  \centering
  \footnotesize
  \begin{tabularx}{1\textwidth}{|l|l|X|}
    \hline
    \multicolumn{3}{c}{\textbf{Contained Capabilities}} \\
    \hline
    \textbf{Capability Name} & \textbf{Capability Id} & \textbf{Comment} \\
    \hline
    \multicolumn{3}{|l|}{\emph{No contained capabilities are currently identified.}} \\
    \hline
  \end{tabularx}
\end{table}

Architectural scope: Identify, classify, relate, and maintain the structural integrity of managed healthcare and administrative entities.

Reconciliation boundary. Retain this existing enterprise capability alongside the newly derived candidates. Any overlap, divergence or potential consolidation is a deliberate review item, not resolved here.

\section{CT2 - EC-02 Context Management}\label{cap:ec-5de7a127}
\vspace{-0.5em}
\begin{table}[htbp]
  \centering
  \footnotesize
  \begin{tabularx}{1\textwidth}{|l|l|X|}
    \hline
    \multicolumn{3}{c}{\textbf{Capability Metadata}} \\
    \hline
    \textbf{Name} & \textbf{Value} & \textbf{Comment and/or Description} \\
    \hline
    \textbf{Reference Id} & \texttt{EC-5DE7A127} & \\
    \hline
    \textbf{Status} & Existing canonical EC & \\
    \hline
    \textbf{Ancestry} &  & \\
    \hline
    \textbf{Source} & Enterprise Capability catalogue & \\
    \hline
    \textbf{Enabled By} &  & \\
    \hline
    \textbf{Supports} &  & \\
    \hline
    \textbf{Derived From} &  & \\
    \hline
    \textbf{Reviewed Status} & Compare against derived Enterprise CT3 and Feature families; do not suppress conflicts. & \\
    \hline
  \end{tabularx}
\end{table}

\begin{table}[htbp]
  \centering
  \footnotesize
  \begin{tabularx}{1\textwidth}{|l|l|X|}
    \hline
    \multicolumn{3}{c}{\textbf{Contained Capabilities}} \\
    \hline
    \textbf{Capability Name} & \textbf{Capability Id} & \textbf{Comment} \\
    \hline
    \multicolumn{3}{|l|}{\emph{No contained capabilities are currently identified.}} \\
    \hline
  \end{tabularx}
\end{table}

Architectural scope: Establish, validate, propagate, and preserve execution, clinical, security, and administrative context.

Reconciliation boundary. Retain this existing enterprise capability alongside the newly derived candidates. Any overlap, divergence or potential consolidation is a deliberate review item, not resolved here.

\section{CT2 - EC-03 Managed State \& Lifecycle}\label{cap:ec-d9b04b9b}
\vspace{-0.5em}
\begin{table}[htbp]
  \centering
  \footnotesize
  \begin{tabularx}{1\textwidth}{|l|l|X|}
    \hline
    \multicolumn{3}{c}{\textbf{Capability Metadata}} \\
    \hline
    \textbf{Name} & \textbf{Value} & \textbf{Comment and/or Description} \\
    \hline
    \textbf{Reference Id} & \texttt{EC-D9B04B9B} & \\
    \hline
    \textbf{Status} & Existing canonical EC & \\
    \hline
    \textbf{Ancestry} &  & \\
    \hline
    \textbf{Source} & Enterprise Capability catalogue & \\
    \hline
    \textbf{Enabled By} &  & \\
    \hline
    \textbf{Supports} &  & \\
    \hline
    \textbf{Derived From} &  & \\
    \hline
    \textbf{Reviewed Status} & Compare against derived Enterprise CT3 and Feature families; do not suppress conflicts. & \\
    \hline
  \end{tabularx}
\end{table}

\begin{table}[htbp]
  \centering
  \footnotesize
  \begin{tabularx}{1\textwidth}{|l|l|X|}
    \hline
    \multicolumn{3}{c}{\textbf{Contained Capabilities}} \\
    \hline
    \textbf{Capability Name} & \textbf{Capability Id} & \textbf{Comment} \\
    \hline
    \multicolumn{3}{|l|}{\emph{No contained capabilities are currently identified.}} \\
    \hline
  \end{tabularx}
\end{table}

Architectural scope: Govern current and historical state, temporal validity, state transitions, and lifecycle progression.

Reconciliation boundary. Retain this existing enterprise capability alongside the newly derived candidates. Any overlap, divergence or potential consolidation is a deliberate review item, not resolved here.

\section{CT2 - EC-04 Information Management}\label{cap:ec-bdda8fd7}
\vspace{-0.5em}
\begin{table}[htbp]
  \centering
  \footnotesize
  \begin{tabularx}{1\textwidth}{|l|l|X|}
    \hline
    \multicolumn{3}{c}{\textbf{Capability Metadata}} \\
    \hline
    \textbf{Name} & \textbf{Value} & \textbf{Comment and/or Description} \\
    \hline
    \textbf{Reference Id} & \texttt{EC-BDDA8FD7} & \\
    \hline
    \textbf{Status} & Existing canonical EC & \\
    \hline
    \textbf{Ancestry} &  & \\
    \hline
    \textbf{Source} & Enterprise Capability catalogue & \\
    \hline
    \textbf{Enabled By} &  & \\
    \hline
    \textbf{Supports} &  & \\
    \hline
    \textbf{Derived From} &  & \\
    \hline
    \textbf{Reviewed Status} & Compare against derived Enterprise CT3 and Feature families; do not suppress conflicts. & \\
    \hline
  \end{tabularx}
\end{table}

\begin{table}[htbp]
  \centering
  \footnotesize
  \begin{tabularx}{1\textwidth}{|l|l|X|}
    \hline
    \multicolumn{3}{c}{\textbf{Contained Capabilities}} \\
    \hline
    \textbf{Capability Name} & \textbf{Capability Id} & \textbf{Comment} \\
    \hline
    \multicolumn{3}{|l|}{\emph{No contained capabilities are currently identified.}} \\
    \hline
  \end{tabularx}
\end{table}

Architectural scope: Validate, represent, persist, retrieve, version, and manage governed clinical, administrative, and operational information.

Reconciliation boundary. Retain this existing enterprise capability alongside the newly derived candidates. Any overlap, divergence or potential consolidation is a deliberate review item, not resolved here.

\section{CT2 - EC-05 Search \& Discovery}\label{cap:ec-c8c13214}
\vspace{-0.5em}
\begin{table}[htbp]
  \centering
  \footnotesize
  \begin{tabularx}{1\textwidth}{|l|l|X|}
    \hline
    \multicolumn{3}{c}{\textbf{Capability Metadata}} \\
    \hline
    \textbf{Name} & \textbf{Value} & \textbf{Comment and/or Description} \\
    \hline
    \textbf{Reference Id} & \texttt{EC-C8C13214} & \\
    \hline
    \textbf{Status} & Existing canonical EC & \\
    \hline
    \textbf{Ancestry} &  & \\
    \hline
    \textbf{Source} & Enterprise Capability catalogue & \\
    \hline
    \textbf{Enabled By} &  & \\
    \hline
    \textbf{Supports} &  & \\
    \hline
    \textbf{Derived From} &  & \\
    \hline
    \textbf{Reviewed Status} & Compare against derived Enterprise CT3 and Feature families; do not suppress conflicts. & \\
    \hline
  \end{tabularx}
\end{table}

\begin{table}[htbp]
  \centering
  \footnotesize
  \begin{tabularx}{1\textwidth}{|l|l|X|}
    \hline
    \multicolumn{3}{c}{\textbf{Contained Capabilities}} \\
    \hline
    \textbf{Capability Name} & \textbf{Capability Id} & \textbf{Comment} \\
    \hline
    \multicolumn{3}{|l|}{\emph{No contained capabilities are currently identified.}} \\
    \hline
  \end{tabularx}
\end{table}

Architectural scope: Index, search, resolve, filter, and discover information assets, entities, endpoints, and healthcare services.

Reconciliation boundary. Retain this existing enterprise capability alongside the newly derived candidates. Any overlap, divergence or potential consolidation is a deliberate review item, not resolved here.

\section{CT2 - EC-06 Policy \& Control}\label{cap:ec-22efebd8}
\vspace{-0.5em}
\begin{table}[htbp]
  \centering
  \footnotesize
  \begin{tabularx}{1\textwidth}{|l|l|X|}
    \hline
    \multicolumn{3}{c}{\textbf{Capability Metadata}} \\
    \hline
    \textbf{Name} & \textbf{Value} & \textbf{Comment and/or Description} \\
    \hline
    \textbf{Reference Id} & \texttt{EC-22EFEBD8} & \\
    \hline
    \textbf{Status} & Existing canonical EC & \\
    \hline
    \textbf{Ancestry} &  & \\
    \hline
    \textbf{Source} & Enterprise Capability catalogue & \\
    \hline
    \textbf{Enabled By} &  & \\
    \hline
    \textbf{Supports} &  & \\
    \hline
    \textbf{Derived From} &  & \\
    \hline
    \textbf{Reviewed Status} & Compare against derived Enterprise CT3 and Feature families; do not suppress conflicts. & \\
    \hline
  \end{tabularx}
\end{table}

\begin{table}[htbp]
  \centering
  \footnotesize
  \begin{tabularx}{1\textwidth}{|l|l|X|}
    \hline
    \multicolumn{3}{c}{\textbf{Contained Capabilities}} \\
    \hline
    \textbf{Capability Name} & \textbf{Capability Id} & \textbf{Comment} \\
    \hline
    \multicolumn{3}{|l|}{\emph{No contained capabilities are currently identified.}} \\
    \hline
  \end{tabularx}
\end{table}

Architectural scope: Evaluate and enforce security, consent, privacy, regulatory, authority, access-control, and information-handling policies.

Reconciliation boundary. Retain this existing enterprise capability alongside the newly derived candidates. Any overlap, divergence or potential consolidation is a deliberate review item, not resolved here.

\section{CT2 - EC-07 Provenance \& Traceability}\label{cap:ec-23c71f1a}
\vspace{-0.5em}
\begin{table}[htbp]
  \centering
  \footnotesize
  \begin{tabularx}{1\textwidth}{|l|l|X|}
    \hline
    \multicolumn{3}{c}{\textbf{Capability Metadata}} \\
    \hline
    \textbf{Name} & \textbf{Value} & \textbf{Comment and/or Description} \\
    \hline
    \textbf{Reference Id} & \texttt{EC-23C71F1A} & \\
    \hline
    \textbf{Status} & Existing canonical EC & \\
    \hline
    \textbf{Ancestry} &  & \\
    \hline
    \textbf{Source} & Enterprise Capability catalogue & \\
    \hline
    \textbf{Enabled By} &  & \\
    \hline
    \textbf{Supports} &  & \\
    \hline
    \textbf{Derived From} &  & \\
    \hline
    \textbf{Reviewed Status} & Compare against derived Enterprise CT3 and Feature families; do not suppress conflicts. & \\
    \hline
  \end{tabularx}
\end{table}

\begin{table}[htbp]
  \centering
  \footnotesize
  \begin{tabularx}{1\textwidth}{|l|l|X|}
    \hline
    \multicolumn{3}{c}{\textbf{Contained Capabilities}} \\
    \hline
    \textbf{Capability Name} & \textbf{Capability Id} & \textbf{Comment} \\
    \hline
    \multicolumn{3}{|l|}{\emph{No contained capabilities are currently identified.}} \\
    \hline
  \end{tabularx}
\end{table}

Architectural scope: Preserve required origin attribution, meaningful transformation history, significant checkpoints, audit evidence, and end-to-end traceability.

Reconciliation boundary. Retain this existing enterprise capability alongside the newly derived candidates. Any overlap, divergence or potential consolidation is a deliberate review item, not resolved here.

\section{CT2 - EC-08 Interoperability \& Exchange}\label{cap:ec-ce98d4fc}
\vspace{-0.5em}
\begin{table}[htbp]
  \centering
  \footnotesize
  \begin{tabularx}{1\textwidth}{|l|l|X|}
    \hline
    \multicolumn{3}{c}{\textbf{Capability Metadata}} \\
    \hline
    \textbf{Name} & \textbf{Value} & \textbf{Comment and/or Description} \\
    \hline
    \textbf{Reference Id} & \texttt{EC-CE98D4FC} & \\
    \hline
    \textbf{Status} & Existing canonical EC & \\
    \hline
    \textbf{Ancestry} &  & \\
    \hline
    \textbf{Source} & Enterprise Capability catalogue & \\
    \hline
    \textbf{Enabled By} &  & \\
    \hline
    \textbf{Supports} &  & \\
    \hline
    \textbf{Derived From} &  & \\
    \hline
    \textbf{Reviewed Status} & Compare against derived Enterprise CT3 and Feature families; do not suppress conflicts. & \\
    \hline
  \end{tabularx}
\end{table}

\begin{table}[htbp]
  \centering
  \footnotesize
  \begin{tabularx}{1\textwidth}{|l|l|X|}
    \hline
    \multicolumn{3}{c}{\textbf{Contained Capabilities}} \\
    \hline
    \textbf{Capability Name} & \textbf{Capability Id} & \textbf{Comment} \\
    \hline
    \multicolumn{3}{|l|}{\emph{No contained capabilities are currently identified.}} \\
    \hline
  \end{tabularx}
\end{table}

Architectural scope: Resolve destinations, transform schemas, route, transport, deliver, acknowledge, and coordinate standards-based information exchange.

Reconciliation boundary. Retain this existing enterprise capability alongside the newly derived candidates. Any overlap, divergence or potential consolidation is a deliberate review item, not resolved here.

\section{CT2 - EC-09 Event \& Subscription}\label{cap:ec-b5bce57c}
\vspace{-0.5em}
\begin{table}[htbp]
  \centering
  \footnotesize
  \begin{tabularx}{1\textwidth}{|l|l|X|}
    \hline
    \multicolumn{3}{c}{\textbf{Capability Metadata}} \\
    \hline
    \textbf{Name} & \textbf{Value} & \textbf{Comment and/or Description} \\
    \hline
    \textbf{Reference Id} & \texttt{EC-B5BCE57C} & \\
    \hline
    \textbf{Status} & Existing canonical EC & \\
    \hline
    \textbf{Ancestry} &  & \\
    \hline
    \textbf{Source} & Enterprise Capability catalogue & \\
    \hline
    \textbf{Enabled By} &  & \\
    \hline
    \textbf{Supports} &  & \\
    \hline
    \textbf{Derived From} &  & \\
    \hline
    \textbf{Reviewed Status} & Compare against derived Enterprise CT3 and Feature families; do not suppress conflicts. & \\
    \hline
  \end{tabularx}
\end{table}

\begin{table}[htbp]
  \centering
  \footnotesize
  \begin{tabularx}{1\textwidth}{|l|l|X|}
    \hline
    \multicolumn{3}{c}{\textbf{Contained Capabilities}} \\
    \hline
    \textbf{Capability Name} & \textbf{Capability Id} & \textbf{Comment} \\
    \hline
    \multicolumn{3}{|l|}{\emph{No contained capabilities are currently identified.}} \\
    \hline
  \end{tabularx}
\end{table}

Architectural scope: Publish, ingest, correlate, filter, and syndicate events and subscription-driven changes.

Reconciliation boundary. Retain this existing enterprise capability alongside the newly derived candidates. Any overlap, divergence or potential consolidation is a deliberate review item, not resolved here.

\section{CT2 - EC-10 Activity \& Execution}\label{cap:ec-29fc1d73}
\vspace{-0.5em}
\begin{table}[htbp]
  \centering
  \footnotesize
  \begin{tabularx}{1\textwidth}{|l|l|X|}
    \hline
    \multicolumn{3}{c}{\textbf{Capability Metadata}} \\
    \hline
    \textbf{Name} & \textbf{Value} & \textbf{Comment and/or Description} \\
    \hline
    \textbf{Reference Id} & \texttt{EC-29FC1D73} & \\
    \hline
    \textbf{Status} & Existing canonical EC & \\
    \hline
    \textbf{Ancestry} &  & \\
    \hline
    \textbf{Source} & Enterprise Capability catalogue & \\
    \hline
    \textbf{Enabled By} &  & \\
    \hline
    \textbf{Supports} &  & \\
    \hline
    \textbf{Derived From} &  & \\
    \hline
    \textbf{Reviewed Status} & Compare against derived Enterprise CT3 and Feature families; do not suppress conflicts. & \\
    \hline
  \end{tabularx}
\end{table}

\begin{table}[htbp]
  \centering
  \footnotesize
  \begin{tabularx}{1\textwidth}{|l|l|X|}
    \hline
    \multicolumn{3}{c}{\textbf{Contained Capabilities}} \\
    \hline
    \textbf{Capability Name} & \textbf{Capability Id} & \textbf{Comment} \\
    \hline
    \multicolumn{3}{|l|}{\emph{No contained capabilities are currently identified.}} \\
    \hline
  \end{tabularx}
\end{table}

Architectural scope: Define, instantiate, match, assign, dispatch, supervise, and progress coordinated operational and clinical activity units.

Reconciliation boundary. Retain this existing enterprise capability alongside the newly derived candidates. Any overlap, divergence or potential consolidation is a deliberate review item, not resolved here.

\section{CT2 - EC-11 Interaction \& Experience}\label{cap:ec-850ddbaa}
\vspace{-0.5em}
\begin{table}[htbp]
  \centering
  \footnotesize
  \begin{tabularx}{1\textwidth}{|l|l|X|}
    \hline
    \multicolumn{3}{c}{\textbf{Capability Metadata}} \\
    \hline
    \textbf{Name} & \textbf{Value} & \textbf{Comment and/or Description} \\
    \hline
    \textbf{Reference Id} & \texttt{EC-850DDBAA} & \\
    \hline
    \textbf{Status} & Existing canonical EC & \\
    \hline
    \textbf{Ancestry} &  & \\
    \hline
    \textbf{Source} & Enterprise Capability catalogue & \\
    \hline
    \textbf{Enabled By} &  & \\
    \hline
    \textbf{Supports} &  & \\
    \hline
    \textbf{Derived From} &  & \\
    \hline
    \textbf{Reviewed Status} & Compare against derived Enterprise CT3 and Feature families; do not suppress conflicts. & \\
    \hline
  \end{tabularx}
\end{table}

\begin{table}[htbp]
  \centering
  \footnotesize
  \begin{tabularx}{1\textwidth}{|l|l|X|}
    \hline
    \multicolumn{3}{c}{\textbf{Contained Capabilities}} \\
    \hline
    \textbf{Capability Name} & \textbf{Capability Id} & \textbf{Comment} \\
    \hline
    \multicolumn{3}{|l|}{\emph{No contained capabilities are currently identified.}} \\
    \hline
  \end{tabularx}
\end{table}

Architectural scope: Provide contextual interaction between human actors, Harmonia, and enterprise presentation services, including visualisation, alerting, and collaboration integration.

Reconciliation boundary. Retain this existing enterprise capability alongside the newly derived candidates. Any overlap, divergence or potential consolidation is a deliberate review item, not resolved here.

\section{CT2 - EC-12 Operational Assurance}\label{cap:ec-e822f3ae}
\vspace{-0.5em}
\begin{table}[htbp]
  \centering
  \footnotesize
  \begin{tabularx}{1\textwidth}{|l|l|X|}
    \hline
    \multicolumn{3}{c}{\textbf{Capability Metadata}} \\
    \hline
    \textbf{Name} & \textbf{Value} & \textbf{Comment and/or Description} \\
    \hline
    \textbf{Reference Id} & \texttt{EC-E822F3AE} & \\
    \hline
    \textbf{Status} & Existing canonical EC & \\
    \hline
    \textbf{Ancestry} &  & \\
    \hline
    \textbf{Source} & Enterprise Capability catalogue & \\
    \hline
    \textbf{Enabled By} &  & \\
    \hline
    \textbf{Supports} &  & \\
    \hline
    \textbf{Derived From} &  & \\
    \hline
    \textbf{Reviewed Status} & Compare against derived Enterprise CT3 and Feature families; do not suppress conflicts. & \\
    \hline
  \end{tabularx}
\end{table}

\begin{table}[htbp]
  \centering
  \footnotesize
  \begin{tabularx}{1\textwidth}{|l|l|X|}
    \hline
    \multicolumn{3}{c}{\textbf{Contained Capabilities}} \\
    \hline
    \textbf{Capability Name} & \textbf{Capability Id} & \textbf{Comment} \\
    \hline
    \multicolumn{3}{|l|}{\emph{No contained capabilities are currently identified.}} \\
    \hline
  \end{tabularx}
\end{table}

Architectural scope: Ensure concurrency integrity, resilience, duplicate detection, exception handling, automated recovery, and operational telemetry.

Reconciliation boundary. Retain this existing enterprise capability alongside the newly derived candidates. Any overlap, divergence or potential consolidation is a deliberate review item, not resolved here.

\section{CT2 - EC-13 Semantic Governance \& Conformance}\label{cap:ec-e7d5c47d}
\vspace{-0.5em}
\begin{table}[htbp]
  \centering
  \footnotesize
  \begin{tabularx}{1\textwidth}{|l|l|X|}
    \hline
    \multicolumn{3}{c}{\textbf{Capability Metadata}} \\
    \hline
    \textbf{Name} & \textbf{Value} & \textbf{Comment and/or Description} \\
    \hline
    \textbf{Reference Id} & \texttt{EC-E7D5C47D} & \\
    \hline
    \textbf{Status} & Existing canonical EC & \\
    \hline
    \textbf{Ancestry} &  & \\
    \hline
    \textbf{Source} & Enterprise Capability catalogue & \\
    \hline
    \textbf{Enabled By} &  & \\
    \hline
    \textbf{Supports} &  & \\
    \hline
    \textbf{Derived From} &  & \\
    \hline
    \textbf{Reviewed Status} & Compare against derived Enterprise CT3 and Feature families; do not suppress conflicts. & \\
    \hline
  \end{tabularx}
\end{table}

\begin{table}[htbp]
  \centering
  \footnotesize
  \begin{tabularx}{1\textwidth}{|l|l|X|}
    \hline
    \multicolumn{3}{c}{\textbf{Contained Capabilities}} \\
    \hline
    \textbf{Capability Name} & \textbf{Capability Id} & \textbf{Comment} \\
    \hline
    \multicolumn{3}{|l|}{\emph{No contained capabilities are currently identified.}} \\
    \hline
  \end{tabularx}
\end{table}

Architectural scope: Govern canonical models, schemas, terminologies, value sets, mappings, constraints, and semantic conformance.

Reconciliation boundary. Retain this existing enterprise capability alongside the newly derived candidates. Any overlap, divergence or potential consolidation is a deliberate review item, not resolved here.

\section{CT2 - EC-14 Service Guardian}\label{cap:ec-18c84d64}
\vspace{-0.5em}
\begin{table}[htbp]
  \centering
  \footnotesize
  \begin{tabularx}{1\textwidth}{|l|l|X|}
    \hline
    \multicolumn{3}{c}{\textbf{Capability Metadata}} \\
    \hline
    \textbf{Name} & \textbf{Value} & \textbf{Comment and/or Description} \\
    \hline
    \textbf{Reference Id} & \texttt{EC-18C84D64} & \\
    \hline
    \textbf{Status} & Existing canonical EC & \\
    \hline
    \textbf{Ancestry} &  & \\
    \hline
    \textbf{Source} & Enterprise Capability catalogue & \\
    \hline
    \textbf{Enabled By} &  & \\
    \hline
    \textbf{Supports} &  & \\
    \hline
    \textbf{Derived From} &  & \\
    \hline
    \textbf{Reviewed Status} & Compare against derived Enterprise CT3 and Feature families; do not suppress conflicts. & \\
    \hline
  \end{tabularx}
\end{table}

\begin{table}[htbp]
  \centering
  \footnotesize
  \begin{tabularx}{1\textwidth}{|l|l|X|}
    \hline
    \multicolumn{3}{c}{\textbf{Contained Capabilities}} \\
    \hline
    \textbf{Capability Name} & \textbf{Capability Id} & \textbf{Comment} \\
    \hline
    \multicolumn{3}{|l|}{\emph{No contained capabilities are currently identified.}} \\
    \hline
  \end{tabularx}
\end{table}

Architectural scope: Independently establish whether governed service behaviour, state, or outcomes remain within defined assurance expectations using applicable criteria and trustworthy evidence.

Reconciliation boundary. Retain this existing enterprise capability alongside the newly derived candidates. Any overlap, divergence or potential consolidation is a deliberate review item, not resolved here.

\appendix
\chapter{CT1 - Identifier Register}
\small\begin{longtable}{@{}p{.19\linewidth}p{.12\linewidth}p{.39\linewidth}p{.24\linewidth}@{}}\toprule Reference ID & Layer & Element & Status \\ \midrule\endhead
\texttt{BC-E85E18F4} & BC & Individual Care Delivery & Established working hierarchy \\
\texttt{BC-AA0D9B50} & BC & Individual Health Assessment & Established working hierarchy \\
\texttt{BC-6956B3FA} & BC & Health History Assessment & Established working hierarchy \\
\texttt{BC-69CE36B4} & BC & Clinical Examination & Established working hierarchy \\
\texttt{BC-3F88F485} & BC & Health Status Measurement & Established working hierarchy \\
\texttt{BC-201A989E} & BC & Clinical Determination & Established working hierarchy \\
\texttt{BC-20CA91E6} & BC & Clinical Finding Interpretation & Established working hierarchy \\
\texttt{BC-A1AC4A2A} & BC & Diagnosis \& Problem Determination & Established working hierarchy \\
\texttt{BC-75E01B0D} & BC & Health Risk \& Prognostic Determination & Established working hierarchy \\
\texttt{BC-F657B57A} & BC & Individual Care Planning & Established working hierarchy \\
\texttt{BC-E8BBA49D} & BC & Care Goal Definition & Established working hierarchy \\
\texttt{BC-3626AD34} & BC & Care Response Planning & Established working hierarchy \\
\texttt{BC-3CBC77C1} & BC & Clinical Intervention & Established working hierarchy \\
\texttt{BC-DD466A23} & BC & Individual Health Monitoring & Established working hierarchy \\
\texttt{BC-7E2744D3} & BC & Health State Monitoring & Established working hierarchy \\
\texttt{BC-02C5DBC7} & BC & Care Response Monitoring & Established working hierarchy \\
\texttt{BC-391EC3B9} & BC & Individual Care Review & Established working hierarchy \\
\texttt{BC-A52E0A33} & BC & Care Appropriateness Review & Established working hierarchy \\
\texttt{BC-31953E77} & BC & Care Effectiveness Review & Established working hierarchy \\
\texttt{BC-2C769A48} & BC & Care Access \& Coordination & Established working hierarchy \\
\texttt{BC-FA41A886} & BC & Care Access Management & Established working hierarchy \\
\texttt{BC-267F1D1B} & BC & Care Intake Management & Established working hierarchy \\
\texttt{BC-EA6B7EB9} & BC & Care Referral Management & Established working hierarchy \\
\texttt{BC-4F461126} & BC & Care Eligibility Determination & Established working hierarchy \\
\texttt{BC-FE20C7C9} & BC & Care Scheduling \& Capacity Management & Established working hierarchy \\
\texttt{BC-9000E197} & BC & Care Appointment Management & Established working hierarchy \\
\texttt{BC-A1F26127} & BC & Care-Place Allocation & Established working hierarchy \\
\texttt{BC-3284993B} & BC & Care-Place Readiness Management & Established working hierarchy \\
\texttt{BC-B53CD721} & BC & Care Resource Deployment & Established working hierarchy \\
\texttt{BC-E24E4ABE} & BC & Operational Care Flow Coordination & Established working hierarchy \\
\texttt{BC-C2D4C33E} & BC & Operational Care Capacity Management & Established working hierarchy \\
\texttt{BC-59CFFB65} & BC & Care Coordination \& Transition & Established working hierarchy \\
\texttt{BC-9C9BD5F2} & BC & Care-Team Coordination & Established working hierarchy \\
\texttt{BC-AA12612C} & BC & Care Responsibility Transition & Established working hierarchy \\
\texttt{BC-D9B48397} & BC & Care Discharge Coordination & Established working hierarchy \\
\texttt{BC-C313E0F4} & BC & Care Support Logistics Management & Established working hierarchy \\
\texttt{BC-CE1271E6} & BC & Patient Movement Coordination & Established working hierarchy \\
\texttt{BC-8924ED80} & BC & Clinical Material Logistics Coordination & Established working hierarchy \\
\texttt{BC-E0FD8D9D} & BC & Care Support Work Coordination & Established working hierarchy \\
\texttt{BC-3E1A6DCA} & BC & Health Rights, Advocacy \& Participation & Established working hierarchy \\
\texttt{BC-C705C8AD} & BC & Diagnostic, Therapeutic \& Clinical Support Services & Established working hierarchy \\
\texttt{BC-E0BE8A8A} & BC & Medication \& Pharmacy Services & Established working hierarchy \\
\texttt{BC-4DA49AEE} & BC & Prescription Medication Dispensing & Established working hierarchy \\
\texttt{BC-C67D1E73} & BC & Non-Prescription Medication Provision & Established working hierarchy \\
\texttt{BC-4E4274D4} & BC & Medication Advice & Established working hierarchy \\
\texttt{BC-F8E65CD9} & BC & Medication Dose Administration Aid Provision & Established working hierarchy \\
\texttt{BC-7B292540} & BC & Diagnostic Support Provision & Established working hierarchy \\
\texttt{BC-D7778EF0} & BC & Diagnostic Service Appointment Management & Established working hierarchy \\
\texttt{BC-20131F3B} & BC & Diagnostic Service Delivery & Established working hierarchy \\
\texttt{BC-D29EA624} & BC & Diagnostic Specimen Management & Established working hierarchy \\
\texttt{BC-47BE7761} & BC & Diagnostic Data Management & Established working hierarchy \\
\texttt{BC-30AB08BB} & BC & Diagnostic Interpretation & Established working hierarchy \\
\texttt{BC-E3C0477C} & BC & Diagnostic Report Management & Established working hierarchy \\
\texttt{BC-E001550A} & BC & Diagnostic Report Distribution & Established working hierarchy \\
\texttt{BC-C82D3CC6} & BC & Health Products \& Clinical Technology & Established working hierarchy \\
\texttt{BC-D7FFDEDC} & BC & Health Product Management & Established working hierarchy \\
\texttt{BC-377FEB37} & BC & Clinical Technology Management & Established working hierarchy \\
\texttt{BC-097846F2} & BC & Clinical Device Inventory Management & Established working hierarchy \\
\texttt{BC-E511B2DB} & BC & Clinical Device Feature \& Release Management & Established working hierarchy \\
\texttt{BC-E2E1C257} & BC & Clinical Device Loan \& Deployment Management & Established working hierarchy \\
\texttt{BC-BBE81736} & BC & Clinical Device Utilisation Support & Established working hierarchy \\
\texttt{BC-33CA15F2} & BC & Clinical Device Integration Provision & Established working hierarchy \\
\texttt{BC-4F2A8587} & BC & Clinical Device Data Integration & Established working hierarchy \\
\texttt{BC-6CB8AC63} & BC & Clinical Device Maintenance Management & Established working hierarchy \\
\texttt{BC-1BFF3F5B} & BC & Clinical Device Lifecycle Management & Established working hierarchy \\
\texttt{BC-E95C383F} & BC & Clinical Quality, Safety \& Improvement & Established working hierarchy \\
\texttt{BC-5B2A0E55} & BC & Clinical Safety Management & Established working hierarchy \\
\texttt{BC-7E3FBC64} & BC & Clinical Risk Identification & Established working hierarchy \\
\texttt{BC-5121C6F1} & BC & Clinical Risk Mitigation Management & Established working hierarchy \\
\texttt{BC-72DCF354} & BC & Clinical Risk Response Management & Established working hierarchy \\
\texttt{BC-077229CB} & BC & Clinical Quality Management & Established working hierarchy \\
\texttt{BC-807007A1} & BC & Clinical Quality Criteria Management & Established working hierarchy \\
\texttt{BC-FC14D2E2} & BC & Clinical Quality Evidence Management & Established working hierarchy \\
\texttt{BC-6884DABF} & BC & Clinical Quality Review \& Assessment Management & Established working hierarchy \\
\texttt{BC-6EEB86B4} & BC & Clinical Quality Reporting & Established working hierarchy \\
\texttt{BC-31C8DC2C} & BC & Community Health \& Wellbeing & Established working hierarchy \\
\texttt{BC-80DA636C} & BC & Preventative Health & Established working hierarchy \\
\texttt{BC-1FAD6263} & BC & Preventative Health Problem Identification & Established working hierarchy \\
\texttt{BC-B3BDD9A1} & BC & Preventative Health Programme Management & Established working hierarchy \\
\texttt{BC-57B54EBE} & BC & Preventative Health Education Services & Established working hierarchy \\
\texttt{BC-5202AEF0} & BC & Preventative Health Programme Monitoring & Established working hierarchy \\
\texttt{BC-E55D2726} & BC & Preventative Health Reporting & Established working hierarchy \\
\texttt{BC-1800BCD9} & BC & Preventative Health Cohort Management & Established working hierarchy \\
\texttt{BC-28CCEE27} & BC & Population Health \& Health-System Planning & Established working hierarchy \\
\texttt{BC-A6F4CA49} & BC & Population Health Surveillance & Established working hierarchy \\
\texttt{BC-F1D5D26C} & BC & Population Health Longitudinal Assessment & Established working hierarchy \\
\texttt{BC-AFC31F01} & BC & Population Health Surveillance Management & Established working hierarchy \\
\texttt{BC-BCA30D47} & BC & Population Health Event Identification & Established working hierarchy \\
\texttt{BC-5EAA6955} & BC & Notifiable Health Event Management & Established working hierarchy \\
\texttt{BC-D8C15299} & BC & Population Health Reporting & Established working hierarchy \\
\texttt{BC-40B6A2C7} & BC & Health-System Planning & Established working hierarchy \\
\texttt{BC-770F60DC} & BC & Health Service Demand Assessment & Established working hierarchy \\
\texttt{BC-3742724C} & BC & Health Service Capacity Planning & Established working hierarchy \\
\texttt{BC-33A5DE96} & BC & Health Service Distribution Planning & Established working hierarchy \\
\texttt{BC-D37DBD07} & BC & Health Information \& Knowledge Management & Established working hierarchy \\
\texttt{BC-E21B0F21} & BC & Health Information Management & Established working hierarchy \\
\texttt{BC-05A55E6B} & BC & Health Information Data Management & Established working hierarchy \\
\texttt{BC-822903B0} & BC & Health Information Structure Management & Established working hierarchy \\
\texttt{BC-65531F2E} & BC & Health Information Lifecycle Management & Established working hierarchy \\
\texttt{BC-2C668F82} & BC & Health Information Catalogue and Index Management & Established working hierarchy \\
\texttt{BC-AD69E2A8} & BC & Health Information Compartmentalisation & Established working hierarchy \\
\texttt{BC-39F5CC28} & BC & Health Information Data Correlation & Established working hierarchy \\
\texttt{BC-C6E88499} & BC & Health Information Data Retrieval & Established working hierarchy \\
\texttt{BC-55C02BD4} & BC & Clinical Knowledge Management & Established working hierarchy \\
\texttt{BC-16799A80} & BC & Clinical Knowledge Context Identification & Established working hierarchy \\
\texttt{BC-D45CD0AA} & BC & Clinical Knowledge Context Extraction & Established working hierarchy \\
\texttt{BC-385FCDAE} & BC & Clinical Knowledge Information Formulation & Established working hierarchy \\
\texttt{BC-AD456750} & BC & Clinical Knowledge Concept Association & Established working hierarchy \\
\texttt{BC-88D82704} & BC & Clinical Knowledge Review \& Approval Management & Established working hierarchy \\
\texttt{BC-E57A9C37} & BC & Clinical Knowledge Publication \& Sharing & Established working hierarchy \\
\texttt{BC-014363F4} & BC & Standards, Semantics \& Reference Governance & Established working hierarchy \\
\texttt{BC-C01C01B9} & BC & Health Information Standards Governance & Established working hierarchy \\
\texttt{BC-C9D7E7C2} & BC & Information Model Governance & Established working hierarchy \\
\texttt{BC-F55F16A6} & BC & Information  Model Discovery & Established working hierarchy \\
\texttt{BC-B27B41ED} & BC & Information Model Adoption & Established working hierarchy \\
\texttt{BC-C5E3E979} & BC & Information Model Education & Established working hierarchy \\
\texttt{BC-FF2BF6BE} & BC & Information Model Compliance & Established working hierarchy \\
\texttt{BC-4486449F} & BC & Standards Discovery & Established working hierarchy \\
\texttt{BC-98A00C71} & BC & Standards Impact Identification & Established working hierarchy \\
\texttt{BC-CB858814} & BC & Standards Adoption & Established working hierarchy \\
\texttt{BC-CB56B8BC} & BC & Standards Education & Established working hierarchy \\
\texttt{BC-BFC63CC1} & BC & Standards Bodies Engagement & Established working hierarchy \\
\texttt{BC-E515BE41} & BC & Standards Compliance & Established working hierarchy \\
\texttt{BC-F9F83F32} & BC & Clinical Semantics Governance & Established working hierarchy \\
\texttt{BC-9D2AC8A5} & BC & Terminology Governance & Established working hierarchy \\
\texttt{BC-9DBB18D4} & BC & Classification Governance & Established working hierarchy \\
\texttt{BC-05135118} & BC & Semantic Mapping Governance & Established working hierarchy \\
\texttt{BC-18835F13} & BC & Terminology Utilisation Oversight & Established working hierarchy \\
\texttt{BC-D4198837} & BC & Terminology Localisation & Established working hierarchy \\
\texttt{BC-A6FDE957} & BC & Classification Utilisation Oversight & Established working hierarchy \\
\texttt{BC-D0174B04} & BC & Connected Health Services & Established working hierarchy \\
\texttt{BC-1A2C0473} & BC & Health Information Exchange & Established working hierarchy \\
\texttt{BC-25F9B730} & BC & Standards-Based Interoperability Exploration & Established working hierarchy \\
\texttt{BC-5CA9556E} & BC & Standards-Based Interoperability Support & Established working hierarchy \\
\texttt{BC-9D0C5F85} & BC & Standards-Based Interoperability Publication & Established working hierarchy \\
\texttt{BC-EA4804CA} & BC & Standards-Based Interoperability Client-Support & Established working hierarchy \\
\texttt{BC-EAE75CE5} & BC & Health Information Sharing Framework Oversight & Established working hierarchy \\
\texttt{BC-F6DF2559} & BC & Health Information Sharing Client-Support & Established working hierarchy \\
\texttt{BC-93BAE720} & BC & Health Information Sharing Release Management & Established working hierarchy \\
\texttt{BC-5FD7002B} & BC & Digital Communication Channel & Established working hierarchy \\
\texttt{BC-1FA8828A} & BC & Clinical Message Conveyance & Established working hierarchy \\
\texttt{BC-BAEB9E26} & BC & Health Event Notification & Established working hierarchy \\
\texttt{BC-B4211FE9} & BC & Health Service Collaboration & Established working hierarchy \\
\texttt{BC-1079B553} & BC & Domain Health Information Services — ADJUDICATED & Established working hierarchy \\
\texttt{BC-34251AB8} & BC & Provider Directory Services & Established working hierarchy \\
\texttt{BC-667FBAA5} & BC & Enterprise Master Patient Index Services & Established working hierarchy \\
\texttt{BC-E3F3D922} & BC & Immunisation Register Services & Established working hierarchy \\
\texttt{BC-EDB9A9D1} & BC & Longitudinal Health Record Services & Established working hierarchy \\
\texttt{BC-A60C3D92} & BC & Collaboration Services & Established working hierarchy \\
\texttt{BC-3525B6FE} & BC & Security, Privacy \& Digital Trust & Established working hierarchy \\
\texttt{BC-0D570829} & BC & Digital Identity \& Access Governance & Established working hierarchy \\
\texttt{BC-51F3C14E} & BC & Identity Assurance & Established working hierarchy \\
\texttt{BC-E2872CBE} & BC & Identity Modelling \& Interoperability & Established working hierarchy \\
\texttt{BC-D6F6BD86} & BC & Identity Qualification Control & Established working hierarchy \\
\texttt{BC-2C37A4F9} & BC & Identity Identification \& Authentication & Established working hierarchy \\
\texttt{BC-92E77AA5} & BC & Access Control Context Management & Established working hierarchy \\
\texttt{BC-0E134F87} & BC & Access Control Assurance \& Reporting & Established working hierarchy \\
\texttt{BC-3A9DB312} & BC & Access Control Collaboration & Established working hierarchy \\
\texttt{BC-A78294CD} & BC & Access Control Enforcement & Established working hierarchy \\
\texttt{BC-448EE064} & BC & Health Information Protection & Established working hierarchy \\
\texttt{BC-F8DE8E42} & BC & Health Privacy Safeguard Management & Established working hierarchy \\
\texttt{BC-CBD883DA} & BC & Health Information Safeguarding & Established working hierarchy \\
\texttt{BC-88AD2C7F} & BC & Health Information Access Monitoring & Established working hierarchy \\
\texttt{BC-69A112EB} & BC & Health Information Compartment Assurance & Established working hierarchy \\
\texttt{BC-C163719D} & BC & Health Information Classification Assurance & Established working hierarchy \\
\texttt{BC-DB00A738} & BC & Digital Accountability \& Trust & Established working hierarchy \\
\texttt{BC-B2FBA48B} & BC & Governed Activity Attribution & Established working hierarchy \\
\texttt{BC-C26E8578} & BC & Security \& Information-Use Accountability & Established working hierarchy \\
\texttt{BC-C5ED32F7} & BC & Health Research \& Innovation & Established working hierarchy \\
\texttt{BC-5FD1C13A} & BC & Research Information Support & Established working hierarchy \\
\texttt{BC-3C3F6346} & BC & Governed Research Information Provision & Established working hierarchy \\
\texttt{BC-5B3EE7BD} & BC & Health Innovation & Established working hierarchy \\
\texttt{BC-E804A440} & BC & Health Innovation Evaluation & Established working hierarchy \\
\texttt{BC-A7B13259} & BC & Health Service Direction \& Stewardship & Established working hierarchy \\
\texttt{BC-25BC222B} & BC & Health Service Modelling & Established working hierarchy \\
\texttt{BC-DF112315} & BC & Health Service Need Identification & Established working hierarchy \\
\texttt{BC-A7A90EE0} & BC & Health Service Capability Identification & Established working hierarchy \\
\texttt{BC-EB94288C} & BC & Health Service Implementation Support & Established working hierarchy \\
\texttt{BC-B1242610} & BC & Health Service Continuous Improvement & Established working hierarchy \\
\texttt{BC-8229ABFD} & BC & Health Service Alignment & Established working hierarchy \\
\texttt{BC-60910383} & BC & Health Service Co-Dependency Management & Established working hierarchy \\
\texttt{BC-AC53A525} & BC & Health Service Resource Identification & Established working hierarchy \\
\texttt{BC-93FDB4F7} & BC & Health Service Financial Management & Established working hierarchy \\
\texttt{BC-EF0176F3} & BC & Health Service Implementation Planning & Established working hierarchy \\
\texttt{BC-F02F310B} & BC & Workforce \& Organisational Capability & Established working hierarchy \\
\texttt{BC-B7DE9919} & BC & Health Workforce Management & Established working hierarchy \\
\texttt{BC-875377CD} & BC & Practitioner Qualification Management & Established working hierarchy \\
\texttt{BC-93AF8CF1} & BC & Practitioner Engagement \& Privilege Management & Established working hierarchy \\
\texttt{BC-DB75D205} & BC & Workforce Service Coverage Management & Established working hierarchy \\
\texttt{BC-DC260B97} & BC & Healthcare Organisation \& Service Structure Management & Established working hierarchy \\
\texttt{BC-A170CFA2} & BC & Healthcare Organisation Identity \& Structure Management & Established working hierarchy \\
\texttt{BC-65079689} & BC & Healthcare Location Definition Management & Established working hierarchy \\
\texttt{BC-40A6EEE3} & BC & Healthcare Service Offering Management & Established working hierarchy \\
\texttt{BC-2A4887A2} & BC & Corporate Resources \& Enterprise Services & Established working hierarchy \\
\texttt{BC-3618D6D1} & BC & Financial \& Commercial Services & Established working hierarchy \\
\texttt{BC-924A98C9} & BC & Clinical Activity Charging Information Provision & Established working hierarchy \\
\texttt{BC-F411AF90} & BC & Corporate Support Services & Established working hierarchy \\
\texttt{BC-FA7FBBAD} & BC & Back-Office Services & Established working hierarchy \\
\texttt{BC-C8C4C54F} & BC & Facilities Management & Established working hierarchy \\
\texttt{BC-EDD1575C} & BC & IT Infrastructure Management & Established working hierarchy \\
\texttt{BC-08AE4743} & BC & Communications \& Publicity & Established working hierarchy \\
\texttt{BC-4406E5B2} & BC & People Management & Established working hierarchy \\
\texttt{BC-CF3A8B92} & BC & Operational Oversight & Established working hierarchy \\
\texttt{BC-1EF62F79} & BC & Health Service Governance — Harmonia-Relevant & Established working hierarchy \\
\texttt{BC-5998319B} & BC & Governance Framework Management & Established working hierarchy \\
\texttt{BC-C01FF756} & BC & Policy / Obligation Governance & Established working hierarchy \\
\texttt{BC-B3AE7E9A} & BC & Authority \& Constraint Governance & Established working hierarchy \\
\texttt{BC-15EBFFF4} & BC & Compliance Governance & Established working hierarchy \\
\texttt{BC-9BEBD16E} & BC & Compliance Assessment Definition & Established working hierarchy \\
\texttt{BC-351A6321} & BC & Governance Reporting \& Checkpoint Definition & Established working hierarchy \\
\texttt{BC-62520466} & BC & Health Service Assurance & Established working hierarchy \\
\texttt{BC-A0266334} & BC & Health Service Assurance Specification & Established working hierarchy \\
\texttt{BC-BA00F585} & BC & Assurance Approach Definition & Established working hierarchy \\
\texttt{BC-28F181C3} & BC & Assurance Evaluation Criteria Governance & Established working hierarchy \\
\texttt{BC-11ACE749} & BC & Assurance Evaluation & Established working hierarchy \\
\texttt{BC-EF506E30} & BC & Assurance Evidence Assessment & Established working hierarchy \\
\texttt{BC-BC9E1E00} & BC & Assurance Conclusion Determination & Established working hierarchy \\
\texttt{BEC-F49A0D09} & BEC & Entity Information Enablement & Candidate ancestry \\
\texttt{BEC-75CEC238} & BEC & Client Administration & Candidate ancestry \\
\texttt{BEC-31E69256} & BEC & Person Identity & Candidate ancestry \\
\texttt{BEC-1ACD822D} & BEC & Identifier Resolution & Existing Feature; candidate CT3 placement \\
\texttt{BEC-3E1D6BA5} & BEC & Cross-Authority Identifier Correlation & Existing Feature; candidate CT3 placement \\
\texttt{BEC-905C884B} & BEC & Identity Alias Association & Existing Feature; candidate CT3 placement \\
\texttt{BEC-F44A13C3} & BEC & Governed Identity Correction & Existing Feature; candidate CT3 placement \\
\texttt{BEC-60F846CB} & BEC & Identity Management & Generated candidate \\
\texttt{BEC-D5D59967} & BEC & Identity Correlation & Generated candidate \\
\texttt{BEC-44123E71} & BEC & Identity Correction & Generated candidate \\
\texttt{BEC-C7981054} & BEC & Relationship Management & Generated candidate \\
\texttt{BEC-07EB6044} & BEC & Identity Discovery & Generated candidate \\
\texttt{BEC-EDCE7C88} & BEC & Healthcare Subject & Candidate ancestry \\
\texttt{BEC-19B283F0} & BEC & Subject Context Binding & Existing Feature; candidate CT3 placement \\
\texttt{BEC-6B8C0FDF} & BEC & Demographic Context Governance & Existing Feature; candidate CT3 placement \\
\texttt{BEC-98ED0FC6} & BEC & Context Establishment & Generated candidate \\
\texttt{BEC-B3C6EFA0} & BEC & State Monitoring & Generated candidate \\
\texttt{BEC-7C07A1C7} & BEC & Exception Management & Generated candidate \\
\texttt{BEC-54EAE8CA} & BEC & Information Retrieval & Generated candidate \\
\texttt{BEC-2494865F} & BEC & Client Relationship & Candidate ancestry \\
\texttt{BEC-BD118344} & BEC & Next-of-Kin \& Guardian Association & Existing Feature; candidate CT3 placement \\
\texttt{BEC-6811FCDB} & BEC & Relationship Validity Tracking & Existing Feature; candidate CT3 placement \\
\texttt{BEC-CDB17807} & BEC & Context Establishment & Generated candidate \\
\texttt{BEC-D61B3C3B} & BEC & State Monitoring & Generated candidate \\
\texttt{BEC-4AACC1B3} & BEC & Exception Management & Generated candidate \\
\texttt{BEC-C9ABFA6F} & BEC & Information Retrieval & Generated candidate \\
\texttt{BEC-ADC0CF33} & BEC & Client Privacy & Candidate ancestry \\
\texttt{BEC-F41B7E86} & BEC & Consent Directive Evaluation & Existing Feature; candidate CT3 placement \\
\texttt{BEC-4D74B3D2} & BEC & Confidentiality Flag Enforcement & Existing Feature; candidate CT3 placement \\
\texttt{BEC-AA3166D7} & BEC & Criteria Management & Generated candidate \\
\texttt{BEC-5A06BA3C} & BEC & Evidence Association & Generated candidate \\
\texttt{BEC-EF1AF8F9} & BEC & Policy Evaluation & Generated candidate \\
\texttt{BEC-7D40E369} & BEC & Decision Traceability & Generated candidate \\
\texttt{BEC-FAF82B77} & BEC & Exception Reporting & Generated candidate \\
\texttt{BEC-A4E759E2} & BEC & Provider Administration & Candidate ancestry \\
\texttt{BEC-5BB55046} & BEC & Practitioner & Candidate ancestry \\
\texttt{BEC-95B0A32B} & BEC & National Practitioner Verification & Existing Feature; candidate CT3 placement \\
\texttt{BEC-7672710F} & BEC & Practitioner Profile Governance & Existing Feature; candidate CT3 placement \\
\texttt{BEC-A4CAE7D6} & BEC & Identity Management & Generated candidate \\
\texttt{BEC-C7BBD105} & BEC & Identity Correlation & Generated candidate \\
\texttt{BEC-F9FA105B} & BEC & Qualification Association & Generated candidate \\
\texttt{BEC-028ABD17} & BEC & Professional Relationship Management & Generated candidate \\
\texttt{BEC-E9C766CB} & BEC & Practitioner Discovery & Generated candidate \\
\texttt{BEC-2DEC073A} & BEC & Practitioner Role & Candidate ancestry \\
\texttt{BEC-116F8108} & BEC & Role \& Specialty Modeling & Existing Feature; candidate CT3 placement \\
\texttt{BEC-CBDE3B8C} & BEC & Electronic Endpoint Binding & Existing Feature; candidate CT3 placement \\
\texttt{BEC-7F9BCB4F} & BEC & Identity Management & Generated candidate \\
\texttt{BEC-314BE2D1} & BEC & Identity Correlation & Generated candidate \\
\texttt{BEC-A4596B27} & BEC & Qualification Association & Generated candidate \\
\texttt{BEC-549D792B} & BEC & Professional Relationship Management & Generated candidate \\
\texttt{BEC-E24857D9} & BEC & Practitioner Discovery & Generated candidate \\
\texttt{BEC-C410CC51} & BEC & Provider Relationship & Candidate ancestry \\
\texttt{BEC-F5495DD1} & BEC & Practitioner-Organisation Affiliation & Existing Feature; candidate CT3 placement \\
\texttt{BEC-0027885B} & BEC & Context Establishment & Generated candidate \\
\texttt{BEC-2CF899D2} & BEC & State Monitoring & Generated candidate \\
\texttt{BEC-9806E183} & BEC & Exception Management & Generated candidate \\
\texttt{BEC-737E5A3F} & BEC & Information Retrieval & Generated candidate \\
\texttt{BEC-F58DBA08} & BEC & Provider Engagement & Candidate ancestry \\
\texttt{BEC-E0307698} & BEC & Clinical Privileging Record & Existing Feature; candidate CT3 placement \\
\texttt{BEC-0C362F3D} & BEC & Context Establishment & Generated candidate \\
\texttt{BEC-141E9549} & BEC & State Monitoring & Generated candidate \\
\texttt{BEC-788A21BF} & BEC & Exception Management & Generated candidate \\
\texttt{BEC-167AF187} & BEC & Information Retrieval & Generated candidate \\
\texttt{BEC-1FF88A71} & BEC & Organisation Administration & Candidate ancestry \\
\texttt{BEC-ACE86B9B} & BEC & Healthcare Organisation & Candidate ancestry \\
\texttt{BEC-28102E35} & BEC & National Organisation Verification & Existing Feature; candidate CT3 placement \\
\texttt{BEC-DF977EAB} & BEC & Organisation Profile Management & Existing Feature; candidate CT3 placement \\
\texttt{BEC-157386A8} & BEC & Identity Management & Generated candidate \\
\texttt{BEC-9262FB93} & BEC & Identity Correlation & Generated candidate \\
\texttt{BEC-AD09D42A} & BEC & Organisational Relationship Management & Generated candidate \\
\texttt{BEC-6B51A123} & BEC & Organisational Structure Management & Generated candidate \\
\texttt{BEC-D6FB2614} & BEC & Organisation Discovery & Generated candidate \\
\texttt{BEC-FA161A55} & BEC & Organisational Structure & Candidate ancestry \\
\texttt{BEC-F5068C04} & BEC & Hierarchical Structure Navigation & Existing Feature; candidate CT3 placement \\
\texttt{BEC-48E4C563} & BEC & Identity Management & Generated candidate \\
\texttt{BEC-A50F5059} & BEC & Identity Correlation & Generated candidate \\
\texttt{BEC-676858F6} & BEC & Organisational Relationship Management & Generated candidate \\
\texttt{BEC-ACAB1168} & BEC & Organisational Structure Management & Generated candidate \\
\texttt{BEC-721B38D1} & BEC & Organisation Discovery & Generated candidate \\
\texttt{BEC-5C12BF05} & BEC & Organisation Contact & Candidate ancestry \\
\texttt{BEC-D6FA9FDF} & BEC & Directory Contact Management & Existing Feature; candidate CT3 placement \\
\texttt{BEC-3FD95735} & BEC & Identity Management & Generated candidate \\
\texttt{BEC-737C2F3E} & BEC & Identity Correlation & Generated candidate \\
\texttt{BEC-EFB7F137} & BEC & Organisational Relationship Management & Generated candidate \\
\texttt{BEC-272EFC75} & BEC & Organisational Structure Management & Generated candidate \\
\texttt{BEC-73EF81CD} & BEC & Organisation Discovery & Generated candidate \\
\texttt{BEC-D4B0D057} & BEC & Location Administration & Candidate ancestry \\
\texttt{BEC-7E83F611} & BEC & Healthcare Location & Candidate ancestry \\
\texttt{BEC-885CE86B} & BEC & Location Geospatial \& Address Resolution & Existing Feature; candidate CT3 placement \\
\texttt{BEC-15887E90} & BEC & Location Identity Management & Generated candidate \\
\texttt{BEC-C52417EE} & BEC & Location Hierarchy Management & Generated candidate \\
\texttt{BEC-0538C23A} & BEC & Location Availability Management & Generated candidate \\
\texttt{BEC-BE1B64CD} & BEC & Location Discovery & Generated candidate \\
\texttt{BEC-A66483F0} & BEC & Location Structure & Candidate ancestry \\
\texttt{BEC-C82C8D0A} & BEC & Physical Hierarchy Modeling & Existing Feature; candidate CT3 placement \\
\texttt{BEC-FC2DD94F} & BEC & Location Identity Management & Generated candidate \\
\texttt{BEC-6DCC28F3} & BEC & Location Hierarchy Management & Generated candidate \\
\texttt{BEC-310FF956} & BEC & Location Availability Management & Generated candidate \\
\texttt{BEC-4BCD6B9C} & BEC & Location Discovery & Generated candidate \\
\texttt{BEC-0C8D2717} & BEC & Care-Place Definition & Candidate ancestry \\
\texttt{BEC-6322E089} & BEC & Atomic Care-Place Modeling & Existing Feature; candidate CT3 placement \\
\texttt{BEC-705CD00C} & BEC & Context Establishment & Generated candidate \\
\texttt{BEC-67DFF33A} & BEC & State Monitoring & Generated candidate \\
\texttt{BEC-E79AFE62} & BEC & Exception Management & Generated candidate \\
\texttt{BEC-4B7196D5} & BEC & Information Retrieval & Generated candidate \\
\texttt{BEC-B3668209} & BEC & Health Service Administration & Candidate ancestry \\
\texttt{BEC-1864D570} & BEC & Healthcare Service Definition & Candidate ancestry \\
\texttt{BEC-5AF74389} & BEC & Service Catalogue Publishing & Existing Feature; candidate CT3 placement \\
\texttt{BEC-8CDB717F} & BEC & Service Definition Management & Generated candidate \\
\texttt{BEC-766061BE} & BEC & Service Eligibility Management & Generated candidate \\
\texttt{BEC-22CF71DF} & BEC & Service Availability Management & Generated candidate \\
\texttt{BEC-DEE34CB8} & BEC & Service Discovery & Generated candidate \\
\texttt{BEC-CC9FD951} & BEC & Service Provision Relationship & Candidate ancestry \\
\texttt{BEC-C68165B1} & BEC & Service-Location-Provider Binding & Existing Feature; candidate CT3 placement \\
\texttt{BEC-46B095FB} & BEC & Service Definition Management & Generated candidate \\
\texttt{BEC-99951D40} & BEC & Service Eligibility Management & Generated candidate \\
\texttt{BEC-F964694F} & BEC & Service Availability Management & Generated candidate \\
\texttt{BEC-7A743DDD} & BEC & Service Discovery & Generated candidate \\
\texttt{BEC-CF667CC1} & BEC & Service Availability Definition & Candidate ancestry \\
\texttt{BEC-A4ABC149} & BEC & Operating Schedule Specification & Existing Feature; candidate CT3 placement \\
\texttt{BEC-61F28669} & BEC & Service Definition Management & Generated candidate \\
\texttt{BEC-0AA34D18} & BEC & Service Eligibility Management & Generated candidate \\
\texttt{BEC-235C211D} & BEC & Service Availability Management & Generated candidate \\
\texttt{BEC-B8A2142A} & BEC & Service Discovery & Generated candidate \\
\texttt{BEC-8858C964} & BEC & Service Eligibility Definition & Candidate ancestry \\
\texttt{BEC-748B065D} & BEC & Eligibility \& Catchment Specification & Existing Feature; candidate CT3 placement \\
\texttt{BEC-271154E8} & BEC & Service Definition Management & Generated candidate \\
\texttt{BEC-1F711951} & BEC & Service Eligibility Management & Generated candidate \\
\texttt{BEC-FA68F289} & BEC & Service Availability Management & Generated candidate \\
\texttt{BEC-9B1E4AC5} & BEC & Service Discovery & Generated candidate \\
\texttt{BEC-60F2ADD5} & BEC & Health Product Administration & Candidate ancestry \\
\texttt{BEC-1BE2ADD6} & BEC & Health Product Definition & Candidate ancestry \\
\texttt{BEC-3700FBFB} & BEC & No established Feature decomposition (Adjacent / Reference context) & Generated candidate \\
\texttt{BEC-0EECBF8C} & BEC & Context Establishment & Generated candidate \\
\texttt{BEC-8F3D543E} & BEC & State Monitoring & Generated candidate \\
\texttt{BEC-6F386B2A} & BEC & Exception Management & Generated candidate \\
\texttt{BEC-F792B00B} & BEC & Information Retrieval & Generated candidate \\
\texttt{BEC-F3F076F6} & BEC & Product Formulary \& Catalogue & Candidate ancestry \\
\texttt{BEC-4937CA4F} & BEC & No established Feature decomposition (Adjacent / Reference context) & Generated candidate \\
\texttt{BEC-D5CD867C} & BEC & Identity Resolution & Generated candidate \\
\texttt{BEC-A315E2F3} & BEC & Identity Correlation & Generated candidate \\
\texttt{BEC-3DBACC5A} & BEC & Relationship Management & Generated candidate \\
\texttt{BEC-5C469882} & BEC & Lifecycle Management & Generated candidate \\
\texttt{BEC-9B063E58} & BEC & Discovery & Generated candidate \\
\texttt{BEC-0ACE1412} & BEC & Product Batch \& Serial Tracking & Candidate ancestry \\
\texttt{BEC-4614D02A} & BEC & No established Feature decomposition (Adjacent / Reference context) & Generated candidate \\
\texttt{BEC-A2361427} & BEC & Context Establishment & Generated candidate \\
\texttt{BEC-C2E441ED} & BEC & State Monitoring & Generated candidate \\
\texttt{BEC-DB97A930} & BEC & Exception Management & Generated candidate \\
\texttt{BEC-4C113522} & BEC & Information Retrieval & Generated candidate \\
\texttt{BEC-39C0D38D} & BEC & Clinical Device Administration & Candidate ancestry \\
\texttt{BEC-EF278D25} & BEC & Device Definition & Candidate ancestry \\
\texttt{BEC-784E5062} & BEC & Device Type Specification & Existing Feature; candidate CT3 placement \\
\texttt{BEC-3117C7A3} & BEC & Device Identity Management & Generated candidate \\
\texttt{BEC-9F1FBF5C} & BEC & Device State Management & Generated candidate \\
\texttt{BEC-018C4C3F} & BEC & Device Relationship Management & Generated candidate \\
\texttt{BEC-624E75EE} & BEC & Device Provenance Tracking & Generated candidate \\
\texttt{BEC-B3DEDCD2} & BEC & Device Instance & Candidate ancestry \\
\texttt{BEC-8FB4909A} & BEC & Device Serial Tracking & Existing Feature; candidate CT3 placement \\
\texttt{BEC-A111F20F} & BEC & Device Identity Management & Generated candidate \\
\texttt{BEC-8C9D3C83} & BEC & Device State Management & Generated candidate \\
\texttt{BEC-A8E9A77A} & BEC & Device Relationship Management & Generated candidate \\
\texttt{BEC-EB9AE886} & BEC & Device Provenance Tracking & Generated candidate \\
\texttt{BEC-868954A0} & BEC & Device Association & Candidate ancestry \\
\texttt{BEC-22B624C5} & BEC & Device-Patient/Location Assignment & Existing Feature; candidate CT3 placement \\
\texttt{BEC-DBE8DD8A} & BEC & Device Identity Management & Generated candidate \\
\texttt{BEC-D81970AA} & BEC & Device State Management & Generated candidate \\
\texttt{BEC-49A18E8C} & BEC & Device Relationship Management & Generated candidate \\
\texttt{BEC-FC971E86} & BEC & Device Provenance Tracking & Generated candidate \\
\texttt{BEC-6548F58A} & BEC & Device Connectivity & Candidate ancestry \\
\texttt{BEC-45355492} & BEC & Device Endpoint Registration & Existing Feature; candidate CT3 placement \\
\texttt{BEC-D0A68C93} & BEC & Device Identity Management & Generated candidate \\
\texttt{BEC-B24562CE} & BEC & Device State Management & Generated candidate \\
\texttt{BEC-14A1B631} & BEC & Device Relationship Management & Generated candidate \\
\texttt{BEC-41EDE6C0} & BEC & Device Provenance Tracking & Generated candidate \\
\texttt{BEC-FF45E08F} & BEC & Care Administration Enablement & Candidate ancestry \\
\texttt{BEC-DD7184DF} & BEC & Client-Centred Administration & Candidate ancestry \\
\texttt{BEC-25DB15BA} & BEC & Client Intake \& Demographic Update & Candidate ancestry \\
\texttt{BEC-4AB59D5A} & BEC & Demographic Update Ingestion & Existing Feature; candidate CT3 placement \\
\texttt{BEC-BDCCD5ED} & BEC & Demographic Change Distribution & Existing Feature; candidate CT3 placement \\
\texttt{BEC-163E98AD} & BEC & Context Establishment & Generated candidate \\
\texttt{BEC-80DFF5D9} & BEC & State Monitoring & Generated candidate \\
\texttt{BEC-F29D899B} & BEC & Exception Management & Generated candidate \\
\texttt{BEC-D55A3518} & BEC & Information Retrieval & Generated candidate \\
\texttt{BEC-C16B2202} & BEC & Referral Administration & Candidate ancestry \\
\texttt{BEC-CD3F820D} & BEC & Referral Intake \& Progression & Candidate ancestry \\
\texttt{BEC-8E9AD826} & BEC & Referral Document Ingestion & Existing Feature; candidate CT3 placement \\
\texttt{BEC-2C8F0BD3} & BEC & Referral Supporting Information & Existing Feature; candidate CT3 placement \\
\texttt{BEC-6E903B7F} & BEC & Referral Status \& Outcome Tracking & Existing Feature; candidate CT3 placement \\
\texttt{BEC-7C40F1C9} & BEC & Referral Lifecycle Management & Generated candidate \\
\texttt{BEC-A551E3F4} & BEC & Referral Routing & Generated candidate \\
\texttt{BEC-C4E89E34} & BEC & Referral Status Tracking & Generated candidate \\
\texttt{BEC-57720BC2} & BEC & Referral Outcome Association & Generated candidate \\
\texttt{BEC-2ABE98E1} & BEC & Scheduling Administration & Candidate ancestry \\
\texttt{BEC-8FEF95D4} & BEC & Appointment \& Booking Administration & Candidate ancestry \\
\texttt{BEC-8C2F56E0} & BEC & Appointment Notification Ingestion & Existing Feature; candidate CT3 placement \\
\texttt{BEC-4143272F} & BEC & Appointment Status Synchronization & Existing Feature; candidate CT3 placement \\
\texttt{BEC-262653A0} & BEC & Appointment Lifecycle Management & Generated candidate \\
\texttt{BEC-D9D79FA0} & BEC & Participant Coordination & Generated candidate \\
\texttt{BEC-425F9192} & BEC & Availability Validation & Generated candidate \\
\texttt{BEC-4D29BEFF} & BEC & Appointment Exception Management & Generated candidate \\
\texttt{BEC-63E23D60} & BEC & Episode \& Encounter Administration & Candidate ancestry \\
\texttt{BEC-1C528FC1} & BEC & Encounter Lifecycle Management & Candidate ancestry \\
\texttt{BEC-0DC29356} & BEC & Encounter Context Creation & Existing Feature; candidate CT3 placement \\
\texttt{BEC-5D8CA48C} & BEC & Encounter State Progression & Existing Feature; candidate CT3 placement \\
\texttt{BEC-847FCFDB} & BEC & Encounter Bed/Location Association & Existing Feature; candidate CT3 placement \\
\texttt{BEC-3EC99787} & BEC & Encounter Lifecycle Management & Generated candidate \\
\texttt{BEC-46431EC5} & BEC & Care Context Association & Generated candidate \\
\texttt{BEC-12EB85C6} & BEC & Participant Association & Generated candidate \\
\texttt{BEC-D42C14E5} & BEC & Encounter History Retrieval & Generated candidate \\
\texttt{BEC-B6BE6E1D} & BEC & Order Administration & Candidate ancestry \\
\texttt{BEC-5E9ECEA1} & BEC & Closed-Loop Order Progression & Candidate ancestry \\
\texttt{BEC-D0729BFB} & BEC & Order Request Ingestion & Existing Feature; candidate CT3 placement \\
\texttt{BEC-20BEDA79} & BEC & Order Destination Resolution \& Routing & Existing Feature; candidate CT3 placement \\
\texttt{BEC-E57851C6} & BEC & Order Closed-Loop Progression Tracking & Existing Feature; candidate CT3 placement \\
\texttt{BEC-84EFF3A3} & BEC & Order Cancellation \& Modification Coordination & Existing Feature; candidate CT3 placement \\
\texttt{BEC-142E39C3} & BEC & Request Lifecycle Management & Generated candidate \\
\texttt{BEC-977D6A7D} & BEC & Request Fulfilment Tracking & Generated candidate \\
\texttt{BEC-81D316C2} & BEC & Request Status Correlation & Generated candidate \\
\texttt{BEC-AA804BEF} & BEC & Request Exception Management & Generated candidate \\
\texttt{BEC-B67A0E82} & BEC & Diagnostic Administration & Candidate ancestry \\
\texttt{BEC-C0DD4810} & BEC & Diagnostic Workflow Management & Candidate ancestry \\
\texttt{BEC-53E5C256} & BEC & Diagnostic Request Correlation & Existing Feature; candidate CT3 placement \\
\texttt{BEC-4B72F961} & BEC & Diagnostic Report Ingestion \& Binding & Existing Feature; candidate CT3 placement \\
\texttt{BEC-322A3126} & BEC & Diagnostic Report Distribution & Existing Feature; candidate CT3 placement \\
\texttt{BEC-CAA7BA59} & BEC & Diagnostic Result Association & Generated candidate \\
\texttt{BEC-67FF7EC6} & BEC & Diagnostic Status Tracking & Generated candidate \\
\texttt{BEC-47552FD9} & BEC & Diagnostic Provenance Capture & Generated candidate \\
\texttt{BEC-B4278279} & BEC & Medication Administration & Candidate ancestry \\
\texttt{BEC-1EFAA849} & BEC & Medication Order \& Dispense Management & Candidate ancestry \\
\texttt{BEC-B26D4F17} & BEC & Medication Order Ingestion & Existing Feature; candidate CT3 placement \\
\texttt{BEC-0DE47133} & BEC & Dispense Event Tracking & Existing Feature; candidate CT3 placement \\
\texttt{BEC-D93D1C27} & BEC & Medication Administration Record (MAR) Ingestion & Existing Feature; candidate CT3 placement \\
\texttt{BEC-AB5055A8} & BEC & Request Lifecycle Management & Generated candidate \\
\texttt{BEC-E0ABA2EF} & BEC & Request Fulfilment Tracking & Generated candidate \\
\texttt{BEC-D9110CE8} & BEC & Request Status Correlation & Generated candidate \\
\texttt{BEC-B0E3D03A} & BEC & Request Exception Management & Generated candidate \\
\texttt{BEC-074E40CB} & BEC & Medication Identity Resolution & Generated candidate \\
\texttt{BEC-AFA2A006} & BEC & Medication Safety Context Validation & Generated candidate \\
\texttt{BEC-9E3BA088} & BEC & Medication Activity Tracking & Generated candidate \\
\texttt{BEC-7D9A630B} & BEC & Medication Information Correlation & Generated candidate \\
\texttt{BEC-3233022C} & BEC & Procedure Administration & Candidate ancestry \\
\texttt{BEC-489D3945} & BEC & Procedure \& Surgical Administration & Candidate ancestry \\
\texttt{BEC-F696528F} & BEC & Procedure Booking Ingestion & Existing Feature; candidate CT3 placement \\
\texttt{BEC-75A8D392} & BEC & Procedural Documentation Ingestion & Existing Feature; candidate CT3 placement \\
\texttt{BEC-5DBEFF67} & BEC & Context Establishment & Generated candidate \\
\texttt{BEC-6100D197} & BEC & State Monitoring & Generated candidate \\
\texttt{BEC-82E3935C} & BEC & Exception Management & Generated candidate \\
\texttt{BEC-C925BA63} & BEC & Information Retrieval & Generated candidate \\
\texttt{BEC-BF6253DC} & BEC & Clinical Record Administration & Candidate ancestry \\
\texttt{BEC-0CE70156} & BEC & Clinical Documentation Lifecycle & Candidate ancestry \\
\texttt{BEC-096B53FA} & BEC & Clinical Document Ingestion & Existing Feature; candidate CT3 placement \\
\texttt{BEC-A9A0C02B} & BEC & Document Versioning \& Supersession & Existing Feature; candidate CT3 placement \\
\texttt{BEC-2963C113} & BEC & Document Metadata Indexing & Existing Feature; candidate CT3 placement \\
\texttt{BEC-836351BA} & BEC & Context Establishment & Generated candidate \\
\texttt{BEC-CF085C02} & BEC & State Monitoring & Generated candidate \\
\texttt{BEC-A27F8453} & BEC & Exception Management & Generated candidate \\
\texttt{BEC-FDC92A6D} & BEC & Information Retrieval & Generated candidate \\
\texttt{BEC-389062F6} & BEC & Care Plan Administration & Candidate ancestry \\
\texttt{BEC-E263E2D1} & BEC & Care Plan Coordination & Candidate ancestry \\
\texttt{BEC-E3837AB5} & BEC & Care Plan Ingestion \& Distribution & Existing Feature; candidate CT3 placement \\
\texttt{BEC-D4E51095} & BEC & Care Goal Association & Generated candidate \\
\texttt{BEC-6CCF8E04} & BEC & Care Plan Lifecycle Tracking & Generated candidate \\
\texttt{BEC-065E070B} & BEC & Care Plan Context Retrieval & Generated candidate \\
\texttt{BEC-0B71E5CB} & BEC & Clinical Communication Administration & Candidate ancestry \\
\texttt{BEC-6D84C2C1} & BEC & Secure Clinical Communication & Candidate ancestry \\
\texttt{BEC-C55CEDBD} & BEC & Clinical Communication Distribution & Existing Feature; candidate CT3 placement \\
\texttt{BEC-357AFA04} & BEC & Delivery Acknowledgement Tracking & Existing Feature; candidate CT3 placement \\
\texttt{BEC-4966229D} & BEC & Communication Audit Logging & Existing Feature; candidate CT3 placement \\
\texttt{BEC-BEE7FF8B} & BEC & Communication Context Management & Generated candidate \\
\texttt{BEC-B1E25071} & BEC & Recipient Resolution & Generated candidate \\
\texttt{BEC-206D85C4} & BEC & Communication Delivery Tracking & Generated candidate \\
\texttt{BEC-09D4A645} & BEC & Funding \& Coverage Administration & Candidate ancestry \\
\texttt{BEC-47D10EB4} & BEC & Clinical Charging \& Claims Administration & Candidate ancestry \\
\texttt{BEC-3A4524B6} & BEC & Clinical Service Context Enablement & Candidate ancestry \\
\texttt{BEC-AEAB9E52} & BEC & Clinical Assessment & Candidate ancestry \\
\texttt{BEC-14FD8EDD} & BEC & Primary Care & Candidate ancestry \\
\texttt{BEC-2D7A6DDE} & BEC & Primary Care Coordination & Candidate ancestry \\
\texttt{BEC-AB31E6BC} & BEC & General Practice Shared Record Ingestion & Existing Feature; candidate CT3 placement \\
\texttt{BEC-E7F50E85} & BEC & Primary Care Notification Dispatch & Existing Feature; candidate CT3 placement \\
\texttt{BEC-DB56036E} & BEC & Context Establishment & Generated candidate \\
\texttt{BEC-8A07FDB4} & BEC & Eligibility Validation & Generated candidate \\
\texttt{BEC-CA466B56} & BEC & Assignment \& Routing & Generated candidate \\
\texttt{BEC-A945B4F5} & BEC & State Tracking & Generated candidate \\
\texttt{BEC-2737B30C} & BEC & Exception Management & Generated candidate \\
\texttt{BEC-F6FC078E} & BEC & Outcome Reporting & Generated candidate \\
\texttt{BEC-4E6DB5F9} & BEC & Acute Care & Candidate ancestry \\
\texttt{BEC-0FAC6B5A} & BEC & Acute Care Monitoring \& Event Coordination & Candidate ancestry \\
\texttt{BEC-5DD077B5} & BEC & Acute Clinical State Event Capture & Existing Feature; candidate CT3 placement \\
\texttt{BEC-870D4A16} & BEC & Acute Care Clinical Context Provision & Existing Feature; candidate CT3 placement \\
\texttt{BEC-3FB23504} & BEC & Context Establishment & Generated candidate \\
\texttt{BEC-CA660CB0} & BEC & Eligibility Validation & Generated candidate \\
\texttt{BEC-CAAB941D} & BEC & Assignment \& Routing & Generated candidate \\
\texttt{BEC-DFC60092} & BEC & State Tracking & Generated candidate \\
\texttt{BEC-1298D75B} & BEC & Exception Management & Generated candidate \\
\texttt{BEC-8ED2F7E4} & BEC & Outcome Reporting & Generated candidate \\
\texttt{BEC-729DFC66} & BEC & Emergency Care & Candidate ancestry \\
\texttt{BEC-A9EBF8FE} & BEC & Emergency Care Progression & Candidate ancestry \\
\texttt{BEC-3E1EE6B4} & BEC & Emergency Department Arrival \& Triage Ingestion & Existing Feature; candidate CT3 placement \\
\texttt{BEC-78F5EADB} & BEC & Emergency Care Cross-Facility Correlation & Existing Feature; candidate CT3 placement \\
\texttt{BEC-ED681415} & BEC & Context Establishment & Generated candidate \\
\texttt{BEC-ADC2AC77} & BEC & State Monitoring & Generated candidate \\
\texttt{BEC-0EC08E00} & BEC & Exception Management & Generated candidate \\
\texttt{BEC-1723857C} & BEC & Information Retrieval & Generated candidate \\
\texttt{BEC-EBFE3FB1} & BEC & Inpatient Care & Candidate ancestry \\
\texttt{BEC-17776772} & BEC & Multidisciplinary Inpatient Care Coordination & Candidate ancestry \\
\texttt{BEC-6F99412C} & BEC & Inpatient Rounding Context Provision & Existing Feature; candidate CT3 placement \\
\texttt{BEC-93C32399} & BEC & Inpatient Clinical Progression Tracking & Existing Feature; candidate CT3 placement \\
\texttt{BEC-826FE486} & BEC & Subject Identity Correlation & Generated candidate \\
\texttt{BEC-941ACC0E} & BEC & Subject Relationship Management & Generated candidate \\
\texttt{BEC-152BD757} & BEC & Subject Context Management & Generated candidate \\
\texttt{BEC-C7242AB0} & BEC & Subject Information Discovery & Generated candidate \\
\texttt{BEC-69DD19FD} & BEC & Ambulatory Care & Candidate ancestry \\
\texttt{BEC-6C600AC1} & BEC & Surgical \& Procedural Care & Candidate ancestry \\
\texttt{BEC-27E670BB} & BEC & Diagnostic Services & Candidate ancestry \\
\texttt{BEC-BB867C0A} & BEC & Diagnostic Service Integration & Candidate ancestry \\
\texttt{BEC-E7A7A26F} & BEC & Laboratory \& Imaging Result Distribution & Existing Feature; candidate CT3 placement \\
\texttt{BEC-821F8A39} & BEC & Diagnostic History Consolidation & Existing Feature; candidate CT3 placement \\
\texttt{BEC-48D54B88} & BEC & Service Definition Management & Generated candidate \\
\texttt{BEC-E8E1F920} & BEC & Service Eligibility Management & Generated candidate \\
\texttt{BEC-870CAEA8} & BEC & Service Availability Management & Generated candidate \\
\texttt{BEC-14C91D01} & BEC & Service Discovery & Generated candidate \\
\texttt{BEC-BDA9D14E} & BEC & Diagnostic Request Correlation & Generated candidate \\
\texttt{BEC-A299B6D0} & BEC & Diagnostic Result Association & Generated candidate \\
\texttt{BEC-D4C02A64} & BEC & Diagnostic Status Tracking & Generated candidate \\
\texttt{BEC-FCD0D5E4} & BEC & Diagnostic Provenance Capture & Generated candidate \\
\texttt{BEC-BFDB47CA} & BEC & Therapeutic Services & Candidate ancestry \\
\texttt{BEC-B48BC336} & BEC & Medication Therapy & Candidate ancestry \\
\texttt{BEC-0789035C} & BEC & Medication Therapy Provision and Monitoring & Candidate ancestry \\
\texttt{BEC-D1549684} & BEC & Medication History Reconciliation Ingestion & Existing Feature; candidate CT3 placement \\
\texttt{BEC-95E0F2E9} & BEC & Medication Administration Tracking & Existing Feature; candidate CT3 placement \\
\texttt{BEC-910194D7} & BEC & Medication Identity Resolution & Generated candidate \\
\texttt{BEC-324D8BAD} & BEC & Medication Safety Context Validation & Generated candidate \\
\texttt{BEC-0F34B38A} & BEC & Medication Activity Tracking & Generated candidate \\
\texttt{BEC-B3D9C0F4} & BEC & Medication Information Correlation & Generated candidate \\
\texttt{BEC-13B74601} & BEC & Allied Health & Candidate ancestry \\
\texttt{BEC-8C53E055} & BEC & Preventive Care & Candidate ancestry \\
\texttt{BEC-6ECB9608} & BEC & Screening and Immunisation Enablement & Candidate ancestry \\
\texttt{BEC-21483404} & BEC & Immunisation Event Submission & Existing Feature; candidate CT3 placement \\
\texttt{BEC-B609C8F5} & BEC & Screening Recall Notification Distribution & Existing Feature; candidate CT3 placement \\
\texttt{BEC-261578C2} & BEC & Context Establishment & Generated candidate \\
\texttt{BEC-C577B7C8} & BEC & State Monitoring & Generated candidate \\
\texttt{BEC-F9613D20} & BEC & Exception Management & Generated candidate \\
\texttt{BEC-336AAF9D} & BEC & Information Retrieval & Generated candidate \\
\texttt{BEC-A8289AED} & BEC & Rehabilitation & Candidate ancestry \\
\texttt{BEC-CDB89B75} & BEC & Community Care & Candidate ancestry \\
\texttt{BEC-50748227} & BEC & Community Care Coordination & Candidate ancestry \\
\texttt{BEC-70EEF999} & BEC & Community Encounter Ingestion & Existing Feature; candidate CT3 placement \\
\texttt{BEC-5CB15F94} & BEC & Community Care Plan Synchronization & Existing Feature; candidate CT3 placement \\
\texttt{BEC-D96C0908} & BEC & Context Establishment & Generated candidate \\
\texttt{BEC-A1012BD3} & BEC & Eligibility Validation & Generated candidate \\
\texttt{BEC-5D53CB17} & BEC & Assignment \& Routing & Generated candidate \\
\texttt{BEC-52C8768A} & BEC & State Tracking & Generated candidate \\
\texttt{BEC-A79A3ED2} & BEC & Exception Management & Generated candidate \\
\texttt{BEC-1725AA08} & BEC & Outcome Reporting & Generated candidate \\
\texttt{BEC-9DDD1BD3} & BEC & Outreach Care & Candidate ancestry \\
\texttt{BEC-B8EB379A} & BEC & Outreach Care Continuity & Candidate ancestry \\
\texttt{BEC-6A90478E} & BEC & Outreach Visit Context Provision & Existing Feature; candidate CT3 placement \\
\texttt{BEC-8BCC3746} & BEC & Outreach Encounter Outcome Ingestion & Existing Feature; candidate CT3 placement \\
\texttt{BEC-89F15AEF} & BEC & Context Establishment & Generated candidate \\
\texttt{BEC-59BE4E9F} & BEC & State Monitoring & Generated candidate \\
\texttt{BEC-8F38AE8D} & BEC & Exception Management & Generated candidate \\
\texttt{BEC-607A7A67} & BEC & Information Retrieval & Generated candidate \\
\texttt{BEC-F9E9DD8E} & BEC & Remote \& Virtual Care & Candidate ancestry \\
\texttt{BEC-68C27168} & BEC & Remote Patient Monitoring & Candidate ancestry \\
\texttt{BEC-FAC2E484} & BEC & Remote Telemetry Ingestion & Existing Feature; candidate CT3 placement \\
\texttt{BEC-D4A0665C} & BEC & Virtual Care Session Context Binding & Existing Feature; candidate CT3 placement \\
\texttt{BEC-08D640DD} & BEC & Subject Identity Correlation & Generated candidate \\
\texttt{BEC-2A678BE0} & BEC & Subject Relationship Management & Generated candidate \\
\texttt{BEC-0B235F68} & BEC & Subject Context Management & Generated candidate \\
\texttt{BEC-6B015BE0} & BEC & Subject Information Discovery & Generated candidate \\
\texttt{BEC-4FA35177} & BEC & Health Service Operations Enablement & Candidate ancestry \\
\texttt{BEC-ED6D98E8} & BEC & Clinic \& Practice Operations & Candidate ancestry \\
\texttt{BEC-FD9982FA} & BEC & Session \& Queue Coordination & Candidate ancestry \\
\texttt{BEC-C271D00D} & BEC & Clinic Session Tracking & Existing Feature; candidate CT3 placement \\
\texttt{BEC-1D3DC25A} & BEC & Operational Queue Progression & Existing Feature; candidate CT3 placement \\
\texttt{BEC-F2DDA45C} & BEC & Context Establishment & Generated candidate \\
\texttt{BEC-2E5F2477} & BEC & Eligibility Validation & Generated candidate \\
\texttt{BEC-91E57A4C} & BEC & Assignment \& Routing & Generated candidate \\
\texttt{BEC-28A55B14} & BEC & State Tracking & Generated candidate \\
\texttt{BEC-334D43B0} & BEC & Exception Management & Generated candidate \\
\texttt{BEC-4734D5E1} & BEC & Outcome Reporting & Generated candidate \\
\texttt{BEC-25E1A896} & BEC & Ward Operations & Candidate ancestry \\
\texttt{BEC-CE9BB804} & BEC & Ward Logistics Coordination & Candidate ancestry \\
\texttt{BEC-9408FFFA} & BEC & Ward Patient Flow Coordination & Existing Feature; candidate CT3 placement \\
\texttt{BEC-06C10C04} & BEC & Clinical Handover Context Assembly & Existing Feature; candidate CT3 placement \\
\texttt{BEC-0029BF2B} & BEC & Context Establishment & Generated candidate \\
\texttt{BEC-616FFDC2} & BEC & Eligibility Validation & Generated candidate \\
\texttt{BEC-0DA17C6A} & BEC & Assignment \& Routing & Generated candidate \\
\texttt{BEC-E6E33D5A} & BEC & State Tracking & Generated candidate \\
\texttt{BEC-E8384DB4} & BEC & Exception Management & Generated candidate \\
\texttt{BEC-4864C81B} & BEC & Outcome Reporting & Generated candidate \\
\texttt{BEC-389897E6} & BEC & Theatre Operations & Candidate ancestry \\
\texttt{BEC-D248D868} & BEC & Perioperative Logistics & Candidate ancestry \\
\texttt{BEC-98713ABB} & BEC & Operating Theatre Case Progression & Existing Feature; candidate CT3 placement \\
\texttt{BEC-B1E9630D} & BEC & Perioperative Resource Notification & Existing Feature; candidate CT3 placement \\
\texttt{BEC-63DA723B} & BEC & Context Establishment & Generated candidate \\
\texttt{BEC-9F573A4C} & BEC & State Monitoring & Generated candidate \\
\texttt{BEC-1233B732} & BEC & Exception Management & Generated candidate \\
\texttt{BEC-792B9999} & BEC & Information Retrieval & Generated candidate \\
\texttt{BEC-78F79EDB} & BEC & Emergency Department Operations & Candidate ancestry \\
\texttt{BEC-469419AB} & BEC & ED Operational Coordination & Candidate ancestry \\
\texttt{BEC-2C9F0B1A} & BEC & ED Bed \& Bay Tracking & Existing Feature; candidate CT3 placement \\
\texttt{BEC-A7E6EDD2} & BEC & ED Length of Stay Tracking & Existing Feature; candidate CT3 placement \\
\texttt{BEC-4CFE53A4} & BEC & Context Establishment & Generated candidate \\
\texttt{BEC-67497215} & BEC & Eligibility Validation & Generated candidate \\
\texttt{BEC-A2B49B7C} & BEC & Assignment \& Routing & Generated candidate \\
\texttt{BEC-704E04EA} & BEC & State Tracking & Generated candidate \\
\texttt{BEC-BCE98F58} & BEC & Exception Management & Generated candidate \\
\texttt{BEC-5E9192CC} & BEC & Outcome Reporting & Generated candidate \\
\texttt{BEC-3EA18C92} & BEC & Outpatient Operations & Candidate ancestry \\
\texttt{BEC-07398864} & BEC & Outpatient Flow Enablement & Candidate ancestry \\
\texttt{BEC-4382A684} & BEC & Outpatient Arrival Registration & Existing Feature; candidate CT3 placement \\
\texttt{BEC-53EB5F3C} & BEC & Post-Clinic Order Coordination & Existing Feature; candidate CT3 placement \\
\texttt{BEC-0D49726E} & BEC & Subject Identity Correlation & Generated candidate \\
\texttt{BEC-FCB21DFC} & BEC & Subject Relationship Management & Generated candidate \\
\texttt{BEC-CD0D5A3A} & BEC & Subject Context Management & Generated candidate \\
\texttt{BEC-17DAEA1C} & BEC & Subject Information Discovery & Generated candidate \\
\texttt{BEC-DBA54D09} & BEC & Bed \& Care-Place Management & Candidate ancestry \\
\texttt{BEC-9F88F277} & BEC & Bed State \& Turnover Management & Candidate ancestry \\
\texttt{BEC-1A021F75} & BEC & Bed Availability Tracking & Existing Feature; candidate CT3 placement \\
\texttt{BEC-ABDDBA4A} & BEC & Bed State Progression & Existing Feature; candidate CT3 placement \\
\texttt{BEC-CB31C6E4} & BEC & Bed Readiness Notification & Existing Feature; candidate CT3 placement \\
\texttt{BEC-21F00612} & BEC & Context Establishment & Generated candidate \\
\texttt{BEC-25AAE44B} & BEC & State Monitoring & Generated candidate \\
\texttt{BEC-49BB68FA} & BEC & Exception Management & Generated candidate \\
\texttt{BEC-BE0DEFDF} & BEC & Information Retrieval & Generated candidate \\
\texttt{BEC-07A2FD29} & BEC & Clinical Resource Management & Candidate ancestry \\
\texttt{BEC-556B4A3F} & BEC & Specialized Equipment Coordination & Candidate ancestry \\
\texttt{BEC-5B685445} & BEC & Mobile Clinical Equipment Tracking & Existing Feature; candidate CT3 placement \\
\texttt{BEC-E7E5A452} & BEC & Context Establishment & Generated candidate \\
\texttt{BEC-AA245FF0} & BEC & Eligibility Validation & Generated candidate \\
\texttt{BEC-A870EB50} & BEC & Assignment \& Routing & Generated candidate \\
\texttt{BEC-F1B486E3} & BEC & State Tracking & Generated candidate \\
\texttt{BEC-99E2FA68} & BEC & Exception Management & Generated candidate \\
\texttt{BEC-4D2A47AC} & BEC & Outcome Reporting & Generated candidate \\
\texttt{BEC-DB7874AF} & BEC & Service Capacity Management & Candidate ancestry \\
\texttt{BEC-D90FAB4C} & BEC & Facility Demand \& Capacity Telemetry & Candidate ancestry \\
\texttt{BEC-A5CDE7AA} & BEC & Operational Capacity Metric Aggregation & Existing Feature; candidate CT3 placement \\
\texttt{BEC-3559F81B} & BEC & Context Establishment & Generated candidate \\
\texttt{BEC-F78FC8D3} & BEC & State Monitoring & Generated candidate \\
\texttt{BEC-0203E495} & BEC & Exception Management & Generated candidate \\
\texttt{BEC-A2BE39BF} & BEC & Information Retrieval & Generated candidate \\
\texttt{BEC-20C1E1B7} & BEC & Clinical Workforce Management & Candidate ancestry \\
\texttt{BEC-F8C1B2DF} & BEC & Clinical Rostering & Candidate ancestry \\
\texttt{BEC-71740710} & BEC & Mobile Staff Management & Candidate ancestry \\
\texttt{BEC-EC43C4CB} & BEC & Mobile Worker Dispatch \& Safety & Candidate ancestry \\
\texttt{BEC-D243C132} & BEC & Mobile Worker Task Dispatch & Existing Feature; candidate CT3 placement \\
\texttt{BEC-CAFF9E49} & BEC & Mobile Visit Status Tracking & Existing Feature; candidate CT3 placement \\
\texttt{BEC-02254531} & BEC & Context Establishment & Generated candidate \\
\texttt{BEC-04274FDD} & BEC & Eligibility Validation & Generated candidate \\
\texttt{BEC-BFBF8974} & BEC & Assignment \& Routing & Generated candidate \\
\texttt{BEC-E87B589F} & BEC & State Tracking & Generated candidate \\
\texttt{BEC-7FD6C955} & BEC & Exception Management & Generated candidate \\
\texttt{BEC-D56786A8} & BEC & Outcome Reporting & Generated candidate \\
\texttt{BEC-AA196190} & BEC & On-Call Management & Candidate ancestry \\
\texttt{BEC-4BE71CB3} & BEC & On-Call Directory Resolution & Candidate ancestry \\
\texttt{BEC-722ABA6E} & BEC & Active On-Call Provider Resolution & Existing Feature; candidate CT3 placement \\
\texttt{BEC-4CFA3E90} & BEC & Context Establishment & Generated candidate \\
\texttt{BEC-703E878F} & BEC & State Monitoring & Generated candidate \\
\texttt{BEC-E452ABBD} & BEC & Exception Management & Generated candidate \\
\texttt{BEC-783C1ED2} & BEC & Information Retrieval & Generated candidate \\
\texttt{BEC-4A72D708} & BEC & Work Allocation \& Dispatch & Candidate ancestry \\
\texttt{BEC-09C6941C} & BEC & Operational Work Assignment \& Supervision & Candidate ancestry \\
\texttt{BEC-02438774} & BEC & Work Item Instantiation & Existing Feature; candidate CT3 placement \\
\texttt{BEC-2A278A2C} & BEC & Worker Matching \& Allocation & Existing Feature; candidate CT3 placement \\
\texttt{BEC-11BC936A} & BEC & Work Dispatch Delivery & Existing Feature; candidate CT3 placement \\
\texttt{BEC-4D7128D9} & BEC & Work Progress Oversight \& Escalation & Existing Feature; candidate CT3 placement \\
\texttt{BEC-59CF5F93} & BEC & Context Establishment & Generated candidate \\
\texttt{BEC-7B45C0B7} & BEC & Eligibility Validation & Generated candidate \\
\texttt{BEC-AD553DA5} & BEC & Assignment \& Routing & Generated candidate \\
\texttt{BEC-D5AEDF90} & BEC & State Tracking & Generated candidate \\
\texttt{BEC-C536BD61} & BEC & Exception Management & Generated candidate \\
\texttt{BEC-A3268ADD} & BEC & Outcome Reporting & Generated candidate \\
\texttt{BEC-06CF5B6D} & BEC & Clinical Training \& Education & Candidate ancestry \\
\texttt{BEC-DE805931} & BEC & Clinical Competency Management & Candidate ancestry \\
\texttt{BEC-9C5388A1} & BEC & Clinical Credential Management & Candidate ancestry \\
\texttt{BEC-E3152A69} & BEC & Credential Verification Checkpoint & Candidate ancestry \\
\texttt{BEC-5CD4669F} & BEC & Operational Credential Check & Existing Feature; candidate CT3 placement \\
\texttt{BEC-D3D46BF9} & BEC & Context Establishment & Generated candidate \\
\texttt{BEC-E9331A31} & BEC & State Monitoring & Generated candidate \\
\texttt{BEC-9736C9DB} & BEC & Exception Management & Generated candidate \\
\texttt{BEC-3DE5AC43} & BEC & Information Retrieval & Generated candidate \\
\texttt{BEC-D401771A} & BEC & Clinical Supervision & Candidate ancestry \\
\texttt{BEC-55FBCC7C} & BEC & Patient Transport & Candidate ancestry \\
\texttt{BEC-854FF619} & BEC & Patient Transfer Coordination & Candidate ancestry \\
\texttt{BEC-EC853A0F} & BEC & Transport Request Ingestion & Existing Feature; candidate CT3 placement \\
\texttt{BEC-58DCCBDB} & BEC & Transport Dispatch \& Progress Tracking & Existing Feature; candidate CT3 placement \\
\texttt{BEC-05915061} & BEC & Subject Identity Correlation & Generated candidate \\
\texttt{BEC-F93AE731} & BEC & Subject Relationship Management & Generated candidate \\
\texttt{BEC-CBF19A4E} & BEC & Subject Context Management & Generated candidate \\
\texttt{BEC-A0E56E8C} & BEC & Subject Information Discovery & Generated candidate \\
\texttt{BEC-BB578270} & BEC & Clinical Logistics Coordination & Candidate ancestry \\
\texttt{BEC-999FB8AB} & BEC & Specimen \& Supply Courier Tracking & Candidate ancestry \\
\texttt{BEC-0BDCFA64} & BEC & Pathology Specimen Transport Tracking & Existing Feature; candidate CT3 placement \\
\texttt{BEC-900A97F0} & BEC & Specimen Identity Management & Generated candidate \\
\texttt{BEC-9950BA59} & BEC & Specimen Custody Tracking & Generated candidate \\
\texttt{BEC-E08C24F3} & BEC & Specimen Relationship Management & Generated candidate \\
\texttt{BEC-3C941BB0} & BEC & Discharge Management & Candidate ancestry \\
\texttt{BEC-3414F1EB} & BEC & Discharge Progression \& Coordination & Candidate ancestry \\
\texttt{BEC-B778D440} & BEC & Discharge Milestone Tracking & Existing Feature; candidate CT3 placement \\
\texttt{BEC-F3326F02} & BEC & Discharge Readiness \& Coordination Oversight & Existing Feature; candidate CT3 placement \\
\texttt{BEC-8B11D01C} & BEC & Discharge Notification Dispatch & Existing Feature; candidate CT3 placement \\
\texttt{BEC-377B8814} & BEC & Context Establishment & Generated candidate \\
\texttt{BEC-54362543} & BEC & Eligibility Validation & Generated candidate \\
\texttt{BEC-92074CE1} & BEC & Assignment \& Routing & Generated candidate \\
\texttt{BEC-D2873E18} & BEC & State Tracking & Generated candidate \\
\texttt{BEC-AF97D3B7} & BEC & Exception Management & Generated candidate \\
\texttt{BEC-D56CF336} & BEC & Outcome Reporting & Generated candidate \\
\texttt{BEC-03705603} & BEC & Shared Information and Coordination Enablement & Candidate ancestry \\
\texttt{BEC-97605881} & BEC & Patient Clinical Record & Candidate ancestry \\
\texttt{BEC-15AB949F} & BEC & Longitudinal Clinical Record Governance & Candidate ancestry \\
\texttt{BEC-B5EBE9A9} & BEC & Longitudinal Health Record (LHR) Assembly & Existing Feature; candidate CT3 placement \\
\texttt{BEC-81C53E17} & BEC & Active Clinical Record Access & Existing Feature; candidate CT3 placement \\
\texttt{BEC-A048190C} & BEC & Durable Clinical Record Preservation & Existing Feature; candidate CT3 placement \\
\texttt{BEC-9F1C5969} & BEC & Record Context Association & Generated candidate \\
\texttt{BEC-3149DD0E} & BEC & Record Provenance Management & Generated candidate \\
\texttt{BEC-67C43D3D} & BEC & Record Lifecycle Management & Generated candidate \\
\texttt{BEC-F9431119} & BEC & Record Discovery & Generated candidate \\
\texttt{BEC-14BBAAFD} & BEC & HIE Integration Services & Candidate ancestry \\
\texttt{BEC-0FADE98F} & BEC & Multi-Model Information Exchange & Candidate ancestry \\
\texttt{BEC-A4404077} & BEC & Direct Submission Ingestion \& Delivery & Existing Feature; candidate CT3 placement \\
\texttt{BEC-3110386A} & BEC & Clinical Record Retrieval Mediation & Existing Feature; candidate CT3 placement \\
\texttt{BEC-E1974228} & BEC & Addressed Clinical Distribution & Existing Feature; candidate CT3 placement \\
\texttt{BEC-05FAD383} & BEC & Event-Driven Syndication & Existing Feature; candidate CT3 placement \\
\texttt{BEC-E86070E3} & BEC & Information Validation & Generated candidate \\
\texttt{BEC-9924A266} & BEC & Information Correlation & Generated candidate \\
\texttt{BEC-C048B88A} & BEC & Information Retrieval & Generated candidate \\
\texttt{BEC-F497FE5F} & BEC & Information Provenance & Generated candidate \\
\texttt{BEC-90EED6CF} & BEC & Information Lifecycle Management & Generated candidate \\
\texttt{BEC-0EE16B9D} & BEC & Health Information Access & Candidate ancestry \\
\texttt{BEC-73452901} & BEC & Governed Query \& Search & Candidate ancestry \\
\texttt{BEC-AB8EDD15} & BEC & Federated Clinical Query & Existing Feature; candidate CT3 placement \\
\texttt{BEC-4601C816} & BEC & Access-Controlled Information Filtering & Existing Feature; candidate CT3 placement \\
\texttt{BEC-72A9140C} & BEC & Context Establishment & Generated candidate \\
\texttt{BEC-474C9605} & BEC & State Monitoring & Generated candidate \\
\texttt{BEC-5604785C} & BEC & Exception Management & Generated candidate \\
\texttt{BEC-C138A5D7} & BEC & Information Retrieval & Generated candidate \\
\texttt{BEC-4553B47B} & BEC & Health Information Communication & Candidate ancestry \\
\texttt{BEC-763876A9} & BEC & Boundary Protocol Mediation & Candidate ancestry \\
\texttt{BEC-0DFE1C63} & BEC & Standards-Based Boundary Exchange & Existing Feature; candidate CT3 placement \\
\texttt{BEC-8EDC75DD} & BEC & Context Establishment & Generated candidate \\
\texttt{BEC-C8A17F6A} & BEC & State Monitoring & Generated candidate \\
\texttt{BEC-6F511B71} & BEC & Exception Management & Generated candidate \\
\texttt{BEC-2EBB9F24} & BEC & Information Retrieval & Generated candidate \\
\texttt{BEC-203E2D62} & BEC & Health Information Control & Candidate ancestry \\
\texttt{BEC-18459CCD} & BEC & Policy \& Security Governance & Candidate ancestry \\
\texttt{BEC-E140C81D} & BEC & Health Information Access \& Consent Control & Existing Feature; candidate CT3 placement \\
\texttt{BEC-CA2EF6F6} & BEC & Security Context Propagation & Existing Feature; candidate CT3 placement \\
\texttt{BEC-0BBA26D3} & BEC & Non-PHI Compliance Auditing & Existing Feature; candidate CT3 placement \\
\texttt{BEC-9EB61778} & BEC & Criteria Management & Generated candidate \\
\texttt{BEC-E3405F10} & BEC & Evidence Association & Generated candidate \\
\texttt{BEC-5BAB5EA0} & BEC & Policy Evaluation & Generated candidate \\
\texttt{BEC-1D2CA7D5} & BEC & Decision Traceability & Generated candidate \\
\texttt{BEC-CB8F5161} & BEC & Exception Reporting & Generated candidate \\
\texttt{BEC-B622BACF} & BEC & Clinical Knowledge Services & Candidate ancestry \\
\texttt{BEC-43B6CE64} & BEC & Terminology \& Reference Mapping & Candidate ancestry \\
\texttt{BEC-B0351502} & BEC & Clinical Concept Resolution & Existing Feature; candidate CT3 placement \\
\texttt{BEC-DC869393} & BEC & Bidirectional Terminology Mapping & Existing Feature; candidate CT3 placement \\
\texttt{BEC-77A76C7C} & BEC & Information Validation & Generated candidate \\
\texttt{BEC-CBEC5F4E} & BEC & Information Correlation & Generated candidate \\
\texttt{BEC-E5775112} & BEC & Information Retrieval & Generated candidate \\
\texttt{BEC-A5D9596D} & BEC & Information Provenance & Generated candidate \\
\texttt{BEC-760AE591} & BEC & Information Lifecycle Management & Generated candidate \\
\texttt{BEC-4CDE0BAB} & BEC & Clinical Collaboration & Candidate ancestry \\
\texttt{BEC-0C37BAFF} & BEC & LHR-Aware Multidisciplinary Collaboration & Candidate ancestry \\
\texttt{BEC-84D2C1DF} & BEC & Care-Team Clinical Collaboration Coordination & Existing Feature; candidate CT3 placement \\
\texttt{BEC-0B2CF342} & BEC & In-Conversation LHR Query Resolution & Existing Feature; candidate CT3 placement \\
\texttt{BEC-24657506} & BEC & Zero-PHI Collaboration Metadata Governance & Existing Feature; candidate CT3 placement \\
\texttt{BEC-3F3E28AF} & BEC & Context Establishment & Generated candidate \\
\texttt{BEC-911087B0} & BEC & State Monitoring & Generated candidate \\
\texttt{BEC-2DF10AA6} & BEC & Exception Management & Generated candidate \\
\texttt{BEC-0CD2C755} & BEC & Information Retrieval & Generated candidate \\
\texttt{BEC-E1807E93} & BEC & General Collaboration & Candidate ancestry \\
\texttt{BEC-2D46A392} & BEC & Operational Messaging & Candidate ancestry \\
\texttt{BEC-1631787F} & BEC & Operational Team Chat Enablement & Existing Feature; candidate CT3 placement \\
\texttt{BEC-074F43D0} & BEC & Context Establishment & Generated candidate \\
\texttt{BEC-D75B5D46} & BEC & State Monitoring & Generated candidate \\
\texttt{BEC-5368E298} & BEC & Exception Management & Generated candidate \\
\texttt{BEC-7D688C5A} & BEC & Information Retrieval & Generated candidate \\
\texttt{BEC-0C4A6874} & BEC & Workflow \& Activity Coordination & Candidate ancestry \\
\texttt{BEC-A3E2C598} & BEC & Multi-Work-Unit Orchestration & Candidate ancestry \\
\texttt{BEC-36C99CCB} & BEC & Work Order Progression & Existing Feature; candidate CT3 placement \\
\texttt{BEC-4306DDD8} & BEC & To Do Decision \& Approval Coordination & Existing Feature; candidate CT3 placement \\
\texttt{BEC-782CD2A4} & BEC & Synthetic Task Orchestration & Existing Feature; candidate CT3 placement \\
\texttt{BEC-76B22B19} & BEC & Workflow Timeout \& Escalation Oversight & Existing Feature; candidate CT3 placement \\
\texttt{BEC-A1649A9C} & BEC & Context Establishment & Generated candidate \\
\texttt{BEC-2078C8D9} & BEC & Eligibility Validation & Generated candidate \\
\texttt{BEC-9FFA6F7F} & BEC & Assignment \& Routing & Generated candidate \\
\texttt{BEC-D25D58D7} & BEC & State Tracking & Generated candidate \\
\texttt{BEC-642BD633} & BEC & Exception Management & Generated candidate \\
\texttt{BEC-6B53749D} & BEC & Outcome Reporting & Generated candidate \\
\texttt{BEC-78F5DB56} & BEC & Calendar Management & Candidate ancestry \\
\texttt{BEC-8CD7EA17} & BEC & Operational \& Clinical Calendar Sync & Candidate ancestry \\
\texttt{BEC-79108DF3} & BEC & Schedule Projection Integration & Existing Feature; candidate CT3 placement \\
\texttt{BEC-45628C87} & BEC & Temporal Event Correlation & Existing Feature; candidate CT3 placement \\
\texttt{BEC-9F5E17D1} & BEC & Context Establishment & Generated candidate \\
\texttt{BEC-922D6FFF} & BEC & State Monitoring & Generated candidate \\
\texttt{BEC-76D61F9A} & BEC & Exception Management & Generated candidate \\
\texttt{BEC-1022E17A} & BEC & Information Retrieval & Generated candidate \\
\texttt{BEC-6AA02A96} & BEC & Presentation Services & Candidate ancestry \\
\texttt{BEC-2DCF9553} & BEC & Decoupled Human Interaction \& Dashboards & Candidate ancestry \\
\texttt{BEC-06AB86FA} & BEC & Longitudinal Clinical \& Operational Presentation & Existing Feature; candidate CT3 placement \\
\texttt{BEC-53CDA4E5} & BEC & Context Establishment & Generated candidate \\
\texttt{BEC-7CB4FD80} & BEC & State Monitoring & Generated candidate \\
\texttt{BEC-D8AF8AC4} & BEC & Exception Management & Generated candidate \\
\texttt{BEC-DA832DCE} & BEC & Information Retrieval & Generated candidate \\
\texttt{BEC-233904DA} & BEC & Information Design Governance & Candidate ancestry \\
\texttt{BEC-453A5B71} & BEC & Canonical Semantics \& Schema Governance & Candidate ancestry \\
\texttt{BEC-D57B8126} & BEC & Information Standards \& Semantic Governance & Existing Feature; candidate CT3 placement \\
\texttt{BEC-D881E00E} & BEC & Semantic Conformance Verification & Existing Feature; candidate CT3 placement \\
\texttt{BEC-74F7B6A8} & BEC & Information Validation & Generated candidate \\
\texttt{BEC-3B311E84} & BEC & Information Correlation & Generated candidate \\
\texttt{BEC-E7F9D13F} & BEC & Information Retrieval & Generated candidate \\
\texttt{BEC-E383206B} & BEC & Information Provenance & Generated candidate \\
\texttt{BEC-915C9E8F} & BEC & Information Lifecycle Management & Generated candidate \\
\texttt{EC-BBAD8E2C} & EC & Entity and Identity & Derived candidate \\
\texttt{EC-C0A97A83} & EC & Identity and Correlation & Derived candidate \\
\texttt{EC-461B8F05} & EC & Identity and Correlation Enablement & Derived candidate \\
\texttt{EC-1A42AFEF} & EC & Identity Correlation & Derived candidate \\
\texttt{EC-00688388} & EC & Identity Management & Derived candidate \\
\texttt{EC-13BDD6AE} & EC & Device Identity Management & Derived candidate \\
\texttt{EC-157BAA94} & EC & Subject Identity Correlation & Derived candidate \\
\texttt{EC-47108AAC} & EC & Information Correlation & Derived candidate \\
\texttt{EC-EA771B95} & EC & Diagnostic Request Correlation & Derived candidate \\
\texttt{EC-039334FA} & EC & Location Identity Management & Derived candidate \\
\texttt{EC-5FC97327} & EC & Medication Identity Resolution & Derived candidate \\
\texttt{EC-D444A2C1} & EC & Medication Information Correlation & Derived candidate \\
\texttt{EC-D418A306} & EC & Request Status Correlation & Derived candidate \\
\texttt{EC-BAC316E4} & EC & Relationship and Structure & Derived candidate \\
\texttt{EC-F50F53E6} & EC & Relationship and Structure Enablement & Derived candidate \\
\texttt{EC-81917406} & EC & Device Relationship Management & Derived candidate \\
\texttt{EC-20119D58} & EC & Subject Relationship Management & Derived candidate \\
\texttt{EC-0B8C4EDC} & EC & Organisational Relationship Management & Derived candidate \\
\texttt{EC-49166B2A} & EC & Organisational Structure Management & Derived candidate \\
\texttt{EC-7A2DF974} & EC & Diagnostic Result Association & Derived candidate \\
\texttt{EC-279FE89F} & EC & Evidence Association & Derived candidate \\
\texttt{EC-A179D91B} & EC & Location Hierarchy Management & Derived candidate \\
\texttt{EC-366AFC9B} & EC & Professional Relationship Management & Derived candidate \\
\texttt{EC-E2F0CB40} & EC & Qualification Association & Derived candidate \\
\texttt{EC-F160E33E} & EC & Relationship Management & Derived candidate \\
\texttt{EC-47C59623} & EC & Information and Semantics & Derived candidate \\
\texttt{EC-260111DC} & EC & Information Lifecycle & Derived candidate \\
\texttt{EC-AE5E9793} & EC & Information Lifecycle Enablement & Derived candidate \\
\texttt{EC-54C1BCAB} & EC & Information Retrieval & Derived candidate \\
\texttt{EC-9788F116} & EC & Subject Information Discovery & Derived candidate \\
\texttt{EC-BFA64AAF} & EC & Information Correlation & Derived candidate \\
\texttt{EC-A3A099AC} & EC & Information Lifecycle Management & Derived candidate \\
\texttt{EC-10E4EA11} & EC & Information Provenance & Derived candidate \\
\texttt{EC-23381315} & EC & Information Validation & Derived candidate \\
\texttt{EC-53DFF234} & EC & Medication Information Correlation & Derived candidate \\
\texttt{EC-720E493E} & EC & Request Lifecycle Management & Derived candidate \\
\texttt{EC-2A3844C6} & EC & Discovery and Retrieval & Derived candidate \\
\texttt{EC-299D0172} & EC & Discovery and Retrieval Enablement & Derived candidate \\
\texttt{EC-F2753584} & EC & Information Retrieval & Derived candidate \\
\texttt{EC-82149096} & EC & Service Discovery & Derived candidate \\
\texttt{EC-13D8A141} & EC & Subject Information Discovery & Derived candidate \\
\texttt{EC-50A2C171} & EC & Organisation Discovery & Derived candidate \\
\texttt{EC-009A4A5A} & EC & Location Discovery & Derived candidate \\
\texttt{EC-09A3798F} & EC & Medication Identity Resolution & Derived candidate \\
\texttt{EC-26353D17} & EC & Practitioner Discovery & Derived candidate \\
\texttt{EC-89E071D0} & EC & Semantic Conformance & Derived candidate \\
\texttt{EC-47BECE6F} & EC & Semantic Conformance Enablement & Derived candidate \\
\texttt{EC-63E0E334} & EC & Eligibility Validation & Derived candidate \\
\texttt{EC-A86D200D} & EC & Information Validation & Derived candidate \\
\texttt{EC-93E36E5B} & EC & Medication Safety Context Validation & Derived candidate \\
\texttt{EC-DB69DAF0} & EC & Qualification Association & Derived candidate \\
\texttt{EC-B80F9118} & EC & Context, Activity and Exchange & Derived candidate \\
\texttt{EC-ADB90088} & EC & Context and State & Derived candidate \\
\texttt{EC-6F46C125} & EC & Context and State Enablement & Derived candidate \\
\texttt{EC-C3C11291} & EC & Context Establishment & Derived candidate \\
\texttt{EC-CC8047E5} & EC & State Monitoring & Derived candidate \\
\texttt{EC-D122BC94} & EC & State Tracking & Derived candidate \\
\texttt{EC-FF26979C} & EC & Service Availability Management & Derived candidate \\
\texttt{EC-6DCE8859} & EC & Device State Management & Derived candidate \\
\texttt{EC-1900212B} & EC & Subject Context Management & Derived candidate \\
\texttt{EC-53072EDD} & EC & No Established Feature Decomposition (Adjacent / Reference Context) & Derived candidate \\
\texttt{EC-7D34D652} & EC & Diagnostic Status Tracking & Derived candidate \\
\texttt{EC-582FEC53} & EC & Location Availability Management & Derived candidate \\
\texttt{EC-2739EBA1} & EC & Medication Safety Context Validation & Derived candidate \\
\texttt{EC-DB0E7812} & EC & Request Status Correlation & Derived candidate \\
\texttt{EC-8C0E52A8} & EC & Activity Coordination & Derived candidate \\
\texttt{EC-D05CE29D} & EC & Activity Coordination Enablement & Derived candidate \\
\texttt{EC-5EE5F09E} & EC & Assignment \& Routing & Derived candidate \\
\texttt{EC-E762705E} & EC & State Tracking & Derived candidate \\
\texttt{EC-7F4BEB56} & EC & Device Provenance Tracking & Derived candidate \\
\texttt{EC-95E97505} & EC & Diagnostic Status Tracking & Derived candidate \\
\texttt{EC-67A2E372} & EC & Medication Activity Tracking & Derived candidate \\
\texttt{EC-A191A2BC} & EC & Request Fulfilment Tracking & Derived candidate \\
\texttt{EC-8B8BC538} & EC & Governance, Trust and Assurance & Derived candidate \\
\texttt{EC-D9AF0B6F} & EC & Governed Control & Derived candidate \\
\texttt{EC-E37E56AC} & EC & Governed Control Enablement & Derived candidate \\
\texttt{EC-5948BFA2} & EC & Eligibility Validation & Derived candidate \\
\texttt{EC-D98B7AAE} & EC & Service Eligibility Management & Derived candidate \\
\texttt{EC-840C09A8} & EC & Policy Evaluation & Derived candidate \\
\texttt{EC-DF0619B8} & EC & Provenance and Accountability & Derived candidate \\
\texttt{EC-3987B790} & EC & Provenance and Accountability Enablement & Derived candidate \\
\texttt{EC-E41C93C7} & EC & Device Provenance Tracking & Derived candidate \\
\texttt{EC-313E85AB} & EC & Information Provenance & Derived candidate \\
\texttt{EC-829F0BF7} & EC & Decision Traceability & Derived candidate \\
\texttt{EC-E9AA136D} & EC & Diagnostic Provenance Capture & Derived candidate \\
\texttt{EC-7233F33E} & EC & Evidence Association & Derived candidate \\
\texttt{EC-74BB4AD4} & EC & Assurance and Resilience & Derived candidate \\
\texttt{EC-B0D5A003} & EC & Assurance and Resilience Enablement & Derived candidate \\
\texttt{EC-50BE2841} & EC & Exception Management & Derived candidate \\
\texttt{EC-0F5671F1} & EC & State Monitoring & Derived candidate \\
\texttt{EC-09035A01} & EC & Outcome Reporting & Derived candidate \\
\texttt{EC-20985EE7} & EC & Exception Reporting & Derived candidate \\
\texttt{EC-34E6439F} & EC & Medication Safety Context Validation & Derived candidate \\
\texttt{EC-85D9E54B} & EC & Request Exception Management & Derived candidate \\
\texttt{EC-1F33C0E8} & EC & Existing Enterprise Capabilities -- Retained for Reconciliation & Existing catalogue; coexistence retained \\
\texttt{EC-3BFF62CA} & EC & EC-01 Managed Entity \& Relationship & Existing canonical EC \\
\texttt{EC-5DE7A127} & EC & EC-02 Context Management & Existing canonical EC \\
\texttt{EC-D9B04B9B} & EC & EC-03 Managed State \& Lifecycle & Existing canonical EC \\
\texttt{EC-BDDA8FD7} & EC & EC-04 Information Management & Existing canonical EC \\
\texttt{EC-C8C13214} & EC & EC-05 Search \& Discovery & Existing canonical EC \\
\texttt{EC-22EFEBD8} & EC & EC-06 Policy \& Control & Existing canonical EC \\
\texttt{EC-23C71F1A} & EC & EC-07 Provenance \& Traceability & Existing canonical EC \\
\texttt{EC-CE98D4FC} & EC & EC-08 Interoperability \& Exchange & Existing canonical EC \\
\texttt{EC-B5BCE57C} & EC & EC-09 Event \& Subscription & Existing canonical EC \\
\texttt{EC-29FC1D73} & EC & EC-10 Activity \& Execution & Existing canonical EC \\
\texttt{EC-850DDBAA} & EC & EC-11 Interaction \& Experience & Existing canonical EC \\
\texttt{EC-E822F3AE} & EC & EC-12 Operational Assurance & Existing canonical EC \\
\texttt{EC-E7D5C47D} & EC & EC-13 Semantic Governance \& Conformance & Existing canonical EC \\
\texttt{EC-18C84D64} & EC & EC-14 Service Guardian & Existing canonical EC \\
\bottomrule\end{longtable}\normalsize
\chapter{CT1 - Feature-Name Recurrence Index}
Names are intentionally repeatable across distinct CT3s. This index is an inspection aid, not evidence of shared identity or mandatory reuse.

\paragraph{Assignment \& Routing} 12 occurrences.\par
\begin{itemize}
\item \texttt{BEC-CA466B56} --- Primary Care Coordination (Semantic-pattern candidate)
\item \texttt{BEC-CAAB941D} --- Acute Care Monitoring \& Event Coordination (Semantic-pattern candidate)
\item \texttt{BEC-5D53CB17} --- Community Care Coordination (Semantic-pattern candidate)
\item \texttt{BEC-91E57A4C} --- Session \& Queue Coordination (Semantic-pattern candidate)
\item \texttt{BEC-0DA17C6A} --- Ward Logistics Coordination (Semantic-pattern candidate)
\item \texttt{BEC-A2B49B7C} --- ED Operational Coordination (Semantic-pattern candidate)
\item \texttt{BEC-A870EB50} --- Specialized Equipment Coordination (Semantic-pattern candidate)
\item \texttt{BEC-BFBF8974} --- Mobile Worker Dispatch \& Safety (Semantic-pattern candidate)
\item \texttt{BEC-AD553DA5} --- Operational Work Assignment \& Supervision (Semantic-pattern candidate)
\item \texttt{BEC-92074CE1} --- Discharge Progression \& Coordination (Semantic-pattern candidate)
\item \texttt{BEC-9FFA6F7F} --- Multi-Work-Unit Orchestration (Semantic-pattern candidate)
\item \texttt{EC-5EE5F09E} --- Activity Coordination Enablement (Enterprise derived)
\end{itemize}
\paragraph{Context Establishment} 36 occurrences.\par
\begin{itemize}
\item \texttt{BEC-98ED0FC6} --- Healthcare Subject (Semantic-pattern candidate)
\item \texttt{BEC-CDB17807} --- Client Relationship (Semantic-pattern candidate)
\item \texttt{BEC-0027885B} --- Provider Relationship (Semantic-pattern candidate)
\item \texttt{BEC-0C362F3D} --- Provider Engagement (Semantic-pattern candidate)
\item \texttt{BEC-705CD00C} --- Care-Place Definition (Semantic-pattern candidate)
\item \texttt{BEC-0EECBF8C} --- Health Product Definition (Semantic-pattern candidate)
\item \texttt{BEC-A2361427} --- Product Batch \& Serial Tracking (Semantic-pattern candidate)
\item \texttt{BEC-163E98AD} --- Client Intake \& Demographic Update (Semantic-pattern candidate)
\item \texttt{BEC-5DBEFF67} --- Procedure \& Surgical Administration (Semantic-pattern candidate)
\item \texttt{BEC-836351BA} --- Clinical Documentation Lifecycle (Semantic-pattern candidate)
\item \texttt{BEC-DB56036E} --- Primary Care Coordination (Semantic-pattern candidate)
\item \texttt{BEC-3FB23504} --- Acute Care Monitoring \& Event Coordination (Semantic-pattern candidate)
\item \texttt{BEC-ED681415} --- Emergency Care Progression (Semantic-pattern candidate)
\item \texttt{BEC-261578C2} --- Screening and Immunisation Enablement (Semantic-pattern candidate)
\item \texttt{BEC-D96C0908} --- Community Care Coordination (Semantic-pattern candidate)
\item \texttt{BEC-89F15AEF} --- Outreach Care Continuity (Semantic-pattern candidate)
\item \texttt{BEC-F2DDA45C} --- Session \& Queue Coordination (Semantic-pattern candidate)
\item \texttt{BEC-0029BF2B} --- Ward Logistics Coordination (Semantic-pattern candidate)
\item \texttt{BEC-63DA723B} --- Perioperative Logistics (Semantic-pattern candidate)
\item \texttt{BEC-4CFE53A4} --- ED Operational Coordination (Semantic-pattern candidate)
\item \texttt{BEC-21F00612} --- Bed State \& Turnover Management (Semantic-pattern candidate)
\item \texttt{BEC-E7E5A452} --- Specialized Equipment Coordination (Semantic-pattern candidate)
\item \texttt{BEC-3559F81B} --- Facility Demand \& Capacity Telemetry (Semantic-pattern candidate)
\item \texttt{BEC-02254531} --- Mobile Worker Dispatch \& Safety (Semantic-pattern candidate)
\item \texttt{BEC-4CFA3E90} --- On-Call Directory Resolution (Semantic-pattern candidate)
\item \texttt{BEC-59CF5F93} --- Operational Work Assignment \& Supervision (Semantic-pattern candidate)
\item \texttt{BEC-D3D46BF9} --- Credential Verification Checkpoint (Semantic-pattern candidate)
\item \texttt{BEC-377B8814} --- Discharge Progression \& Coordination (Semantic-pattern candidate)
\item \texttt{BEC-72A9140C} --- Governed Query \& Search (Semantic-pattern candidate)
\item \texttt{BEC-8EDC75DD} --- Boundary Protocol Mediation (Semantic-pattern candidate)
\item \texttt{BEC-3F3E28AF} --- LHR-Aware Multidisciplinary Collaboration (Semantic-pattern candidate)
\item \texttt{BEC-074F43D0} --- Operational Messaging (Semantic-pattern candidate)
\item \texttt{BEC-A1649A9C} --- Multi-Work-Unit Orchestration (Semantic-pattern candidate)
\item \texttt{BEC-9F5E17D1} --- Operational \& Clinical Calendar Sync (Semantic-pattern candidate)
\item \texttt{BEC-53CDA4E5} --- Decoupled Human Interaction \& Dashboards (Semantic-pattern candidate)
\item \texttt{EC-C3C11291} --- Context and State Enablement (Enterprise derived)
\end{itemize}
\paragraph{Criteria Management} 2 occurrences.\par
\begin{itemize}
\item \texttt{BEC-AA3166D7} --- Client Privacy (Semantic-pattern candidate)
\item \texttt{BEC-9EB61778} --- Policy \& Security Governance (Semantic-pattern candidate)
\end{itemize}
\paragraph{Decision Traceability} 3 occurrences.\par
\begin{itemize}
\item \texttt{BEC-7D40E369} --- Client Privacy (Semantic-pattern candidate)
\item \texttt{BEC-1D2CA7D5} --- Policy \& Security Governance (Semantic-pattern candidate)
\item \texttt{EC-829F0BF7} --- Provenance and Accountability Enablement (Enterprise derived)
\end{itemize}
\paragraph{Device Identity Management} 5 occurrences.\par
\begin{itemize}
\item \texttt{BEC-3117C7A3} --- Device Definition (FHIR-informed candidate)
\item \texttt{BEC-A111F20F} --- Device Instance (FHIR-informed candidate)
\item \texttt{BEC-DBE8DD8A} --- Device Association (FHIR-informed candidate)
\item \texttt{BEC-D0A68C93} --- Device Connectivity (FHIR-informed candidate)
\item \texttt{EC-13BDD6AE} --- Identity and Correlation Enablement (Enterprise derived)
\end{itemize}
\paragraph{Device Provenance Tracking} 6 occurrences.\par
\begin{itemize}
\item \texttt{BEC-624E75EE} --- Device Definition (FHIR-informed candidate)
\item \texttt{BEC-EB9AE886} --- Device Instance (FHIR-informed candidate)
\item \texttt{BEC-FC971E86} --- Device Association (FHIR-informed candidate)
\item \texttt{BEC-41EDE6C0} --- Device Connectivity (FHIR-informed candidate)
\item \texttt{EC-7F4BEB56} --- Activity Coordination Enablement (Enterprise derived)
\item \texttt{EC-E41C93C7} --- Provenance and Accountability Enablement (Enterprise derived)
\end{itemize}
\paragraph{Device Relationship Management} 5 occurrences.\par
\begin{itemize}
\item \texttt{BEC-018C4C3F} --- Device Definition (FHIR-informed candidate)
\item \texttt{BEC-A8E9A77A} --- Device Instance (FHIR-informed candidate)
\item \texttt{BEC-49A18E8C} --- Device Association (FHIR-informed candidate)
\item \texttt{BEC-14A1B631} --- Device Connectivity (FHIR-informed candidate)
\item \texttt{EC-81917406} --- Relationship and Structure Enablement (Enterprise derived)
\end{itemize}
\paragraph{Device State Management} 5 occurrences.\par
\begin{itemize}
\item \texttt{BEC-9F1FBF5C} --- Device Definition (FHIR-informed candidate)
\item \texttt{BEC-8C9D3C83} --- Device Instance (FHIR-informed candidate)
\item \texttt{BEC-D81970AA} --- Device Association (FHIR-informed candidate)
\item \texttt{BEC-B24562CE} --- Device Connectivity (FHIR-informed candidate)
\item \texttt{EC-6DCE8859} --- Context and State Enablement (Enterprise derived)
\end{itemize}
\paragraph{Diagnostic Provenance Capture} 3 occurrences.\par
\begin{itemize}
\item \texttt{BEC-47552FD9} --- Diagnostic Workflow Management (FHIR-informed candidate)
\item \texttt{BEC-FCD0D5E4} --- Diagnostic Service Integration (FHIR-informed candidate)
\item \texttt{EC-E9AA136D} --- Provenance and Accountability Enablement (Enterprise derived)
\end{itemize}
\paragraph{Diagnostic Request Correlation} 3 occurrences.\par
\begin{itemize}
\item \texttt{BEC-53E5C256} --- Diagnostic Workflow Management (Existing catalogue)
\item \texttt{BEC-BDA9D14E} --- Diagnostic Service Integration (FHIR-informed candidate)
\item \texttt{EC-EA771B95} --- Identity and Correlation Enablement (Enterprise derived)
\end{itemize}
\paragraph{Diagnostic Result Association} 3 occurrences.\par
\begin{itemize}
\item \texttt{BEC-CAA7BA59} --- Diagnostic Workflow Management (FHIR-informed candidate)
\item \texttt{BEC-A299B6D0} --- Diagnostic Service Integration (FHIR-informed candidate)
\item \texttt{EC-7A2DF974} --- Relationship and Structure Enablement (Enterprise derived)
\end{itemize}
\paragraph{Diagnostic Status Tracking} 4 occurrences.\par
\begin{itemize}
\item \texttt{BEC-67FF7EC6} --- Diagnostic Workflow Management (FHIR-informed candidate)
\item \texttt{BEC-D4C02A64} --- Diagnostic Service Integration (FHIR-informed candidate)
\item \texttt{EC-7D34D652} --- Context and State Enablement (Enterprise derived)
\item \texttt{EC-95E97505} --- Activity Coordination Enablement (Enterprise derived)
\end{itemize}
\paragraph{Eligibility Validation} 13 occurrences.\par
\begin{itemize}
\item \texttt{BEC-8A07FDB4} --- Primary Care Coordination (Semantic-pattern candidate)
\item \texttt{BEC-CA660CB0} --- Acute Care Monitoring \& Event Coordination (Semantic-pattern candidate)
\item \texttt{BEC-A1012BD3} --- Community Care Coordination (Semantic-pattern candidate)
\item \texttt{BEC-2E5F2477} --- Session \& Queue Coordination (Semantic-pattern candidate)
\item \texttt{BEC-616FFDC2} --- Ward Logistics Coordination (Semantic-pattern candidate)
\item \texttt{BEC-67497215} --- ED Operational Coordination (Semantic-pattern candidate)
\item \texttt{BEC-AA245FF0} --- Specialized Equipment Coordination (Semantic-pattern candidate)
\item \texttt{BEC-04274FDD} --- Mobile Worker Dispatch \& Safety (Semantic-pattern candidate)
\item \texttt{BEC-7B45C0B7} --- Operational Work Assignment \& Supervision (Semantic-pattern candidate)
\item \texttt{BEC-54362543} --- Discharge Progression \& Coordination (Semantic-pattern candidate)
\item \texttt{BEC-2078C8D9} --- Multi-Work-Unit Orchestration (Semantic-pattern candidate)
\item \texttt{EC-63E0E334} --- Semantic Conformance Enablement (Enterprise derived)
\item \texttt{EC-5948BFA2} --- Governed Control Enablement (Enterprise derived)
\end{itemize}
\paragraph{Evidence Association} 4 occurrences.\par
\begin{itemize}
\item \texttt{BEC-5A06BA3C} --- Client Privacy (Semantic-pattern candidate)
\item \texttt{BEC-E3405F10} --- Policy \& Security Governance (Semantic-pattern candidate)
\item \texttt{EC-279FE89F} --- Relationship and Structure Enablement (Enterprise derived)
\item \texttt{EC-7233F33E} --- Provenance and Accountability Enablement (Enterprise derived)
\end{itemize}
\paragraph{Exception Management} 36 occurrences.\par
\begin{itemize}
\item \texttt{BEC-7C07A1C7} --- Healthcare Subject (Semantic-pattern candidate)
\item \texttt{BEC-4AACC1B3} --- Client Relationship (Semantic-pattern candidate)
\item \texttt{BEC-9806E183} --- Provider Relationship (Semantic-pattern candidate)
\item \texttt{BEC-788A21BF} --- Provider Engagement (Semantic-pattern candidate)
\item \texttt{BEC-E79AFE62} --- Care-Place Definition (Semantic-pattern candidate)
\item \texttt{BEC-6F386B2A} --- Health Product Definition (Semantic-pattern candidate)
\item \texttt{BEC-DB97A930} --- Product Batch \& Serial Tracking (Semantic-pattern candidate)
\item \texttt{BEC-F29D899B} --- Client Intake \& Demographic Update (Semantic-pattern candidate)
\item \texttt{BEC-82E3935C} --- Procedure \& Surgical Administration (Semantic-pattern candidate)
\item \texttt{BEC-A27F8453} --- Clinical Documentation Lifecycle (Semantic-pattern candidate)
\item \texttt{BEC-2737B30C} --- Primary Care Coordination (Semantic-pattern candidate)
\item \texttt{BEC-1298D75B} --- Acute Care Monitoring \& Event Coordination (Semantic-pattern candidate)
\item \texttt{BEC-0EC08E00} --- Emergency Care Progression (Semantic-pattern candidate)
\item \texttt{BEC-F9613D20} --- Screening and Immunisation Enablement (Semantic-pattern candidate)
\item \texttt{BEC-A79A3ED2} --- Community Care Coordination (Semantic-pattern candidate)
\item \texttt{BEC-8F38AE8D} --- Outreach Care Continuity (Semantic-pattern candidate)
\item \texttt{BEC-334D43B0} --- Session \& Queue Coordination (Semantic-pattern candidate)
\item \texttt{BEC-E8384DB4} --- Ward Logistics Coordination (Semantic-pattern candidate)
\item \texttt{BEC-1233B732} --- Perioperative Logistics (Semantic-pattern candidate)
\item \texttt{BEC-BCE98F58} --- ED Operational Coordination (Semantic-pattern candidate)
\item \texttt{BEC-49BB68FA} --- Bed State \& Turnover Management (Semantic-pattern candidate)
\item \texttt{BEC-99E2FA68} --- Specialized Equipment Coordination (Semantic-pattern candidate)
\item \texttt{BEC-0203E495} --- Facility Demand \& Capacity Telemetry (Semantic-pattern candidate)
\item \texttt{BEC-7FD6C955} --- Mobile Worker Dispatch \& Safety (Semantic-pattern candidate)
\item \texttt{BEC-E452ABBD} --- On-Call Directory Resolution (Semantic-pattern candidate)
\item \texttt{BEC-C536BD61} --- Operational Work Assignment \& Supervision (Semantic-pattern candidate)
\item \texttt{BEC-9736C9DB} --- Credential Verification Checkpoint (Semantic-pattern candidate)
\item \texttt{BEC-AF97D3B7} --- Discharge Progression \& Coordination (Semantic-pattern candidate)
\item \texttt{BEC-5604785C} --- Governed Query \& Search (Semantic-pattern candidate)
\item \texttt{BEC-6F511B71} --- Boundary Protocol Mediation (Semantic-pattern candidate)
\item \texttt{BEC-2DF10AA6} --- LHR-Aware Multidisciplinary Collaboration (Semantic-pattern candidate)
\item \texttt{BEC-5368E298} --- Operational Messaging (Semantic-pattern candidate)
\item \texttt{BEC-642BD633} --- Multi-Work-Unit Orchestration (Semantic-pattern candidate)
\item \texttt{BEC-76D61F9A} --- Operational \& Clinical Calendar Sync (Semantic-pattern candidate)
\item \texttt{BEC-D8AF8AC4} --- Decoupled Human Interaction \& Dashboards (Semantic-pattern candidate)
\item \texttt{EC-50BE2841} --- Assurance and Resilience Enablement (Enterprise derived)
\end{itemize}
\paragraph{Exception Reporting} 3 occurrences.\par
\begin{itemize}
\item \texttt{BEC-FAF82B77} --- Client Privacy (Semantic-pattern candidate)
\item \texttt{BEC-CB8F5161} --- Policy \& Security Governance (Semantic-pattern candidate)
\item \texttt{EC-20985EE7} --- Assurance and Resilience Enablement (Enterprise derived)
\end{itemize}
\paragraph{Identity Correlation} 8 occurrences.\par
\begin{itemize}
\item \texttt{BEC-D5D59967} --- Person Identity (FHIR-informed candidate)
\item \texttt{BEC-C7BBD105} --- Practitioner (FHIR-informed candidate)
\item \texttt{BEC-314BE2D1} --- Practitioner Role (FHIR-informed candidate)
\item \texttt{BEC-9262FB93} --- Healthcare Organisation (FHIR-informed candidate)
\item \texttt{BEC-A50F5059} --- Organisational Structure (FHIR-informed candidate)
\item \texttt{BEC-737C2F3E} --- Organisation Contact (FHIR-informed candidate)
\item \texttt{BEC-A315E2F3} --- Product Formulary \& Catalogue (Semantic-pattern candidate)
\item \texttt{EC-1A42AFEF} --- Identity and Correlation Enablement (Enterprise derived)
\end{itemize}
\paragraph{Identity Management} 7 occurrences.\par
\begin{itemize}
\item \texttt{BEC-60F846CB} --- Person Identity (FHIR-informed candidate)
\item \texttt{BEC-A4CAE7D6} --- Practitioner (FHIR-informed candidate)
\item \texttt{BEC-7F9BCB4F} --- Practitioner Role (FHIR-informed candidate)
\item \texttt{BEC-157386A8} --- Healthcare Organisation (FHIR-informed candidate)
\item \texttt{BEC-48E4C563} --- Organisational Structure (FHIR-informed candidate)
\item \texttt{BEC-3FD95735} --- Organisation Contact (FHIR-informed candidate)
\item \texttt{EC-00688388} --- Identity and Correlation Enablement (Enterprise derived)
\end{itemize}
\paragraph{Information Correlation} 5 occurrences.\par
\begin{itemize}
\item \texttt{BEC-9924A266} --- Multi-Model Information Exchange (Semantic-pattern candidate)
\item \texttt{BEC-CBEC5F4E} --- Terminology \& Reference Mapping (Semantic-pattern candidate)
\item \texttt{BEC-3B311E84} --- Canonical Semantics \& Schema Governance (Semantic-pattern candidate)
\item \texttt{EC-47108AAC} --- Identity and Correlation Enablement (Enterprise derived)
\item \texttt{EC-BFA64AAF} --- Information Lifecycle Enablement (Enterprise derived)
\end{itemize}
\paragraph{Information Lifecycle Management} 4 occurrences.\par
\begin{itemize}
\item \texttt{BEC-90EED6CF} --- Multi-Model Information Exchange (Semantic-pattern candidate)
\item \texttt{BEC-760AE591} --- Terminology \& Reference Mapping (Semantic-pattern candidate)
\item \texttt{BEC-915C9E8F} --- Canonical Semantics \& Schema Governance (Semantic-pattern candidate)
\item \texttt{EC-A3A099AC} --- Information Lifecycle Enablement (Enterprise derived)
\end{itemize}
\paragraph{Information Provenance} 5 occurrences.\par
\begin{itemize}
\item \texttt{BEC-F497FE5F} --- Multi-Model Information Exchange (Semantic-pattern candidate)
\item \texttt{BEC-A5D9596D} --- Terminology \& Reference Mapping (Semantic-pattern candidate)
\item \texttt{BEC-E383206B} --- Canonical Semantics \& Schema Governance (Semantic-pattern candidate)
\item \texttt{EC-10E4EA11} --- Information Lifecycle Enablement (Enterprise derived)
\item \texttt{EC-313E85AB} --- Provenance and Accountability Enablement (Enterprise derived)
\end{itemize}
\paragraph{Information Retrieval} 29 occurrences.\par
\begin{itemize}
\item \texttt{BEC-54EAE8CA} --- Healthcare Subject (Semantic-pattern candidate)
\item \texttt{BEC-C9ABFA6F} --- Client Relationship (Semantic-pattern candidate)
\item \texttt{BEC-737E5A3F} --- Provider Relationship (Semantic-pattern candidate)
\item \texttt{BEC-167AF187} --- Provider Engagement (Semantic-pattern candidate)
\item \texttt{BEC-4B7196D5} --- Care-Place Definition (Semantic-pattern candidate)
\item \texttt{BEC-F792B00B} --- Health Product Definition (Semantic-pattern candidate)
\item \texttt{BEC-4C113522} --- Product Batch \& Serial Tracking (Semantic-pattern candidate)
\item \texttt{BEC-D55A3518} --- Client Intake \& Demographic Update (Semantic-pattern candidate)
\item \texttt{BEC-C925BA63} --- Procedure \& Surgical Administration (Semantic-pattern candidate)
\item \texttt{BEC-FDC92A6D} --- Clinical Documentation Lifecycle (Semantic-pattern candidate)
\item \texttt{BEC-1723857C} --- Emergency Care Progression (Semantic-pattern candidate)
\item \texttt{BEC-336AAF9D} --- Screening and Immunisation Enablement (Semantic-pattern candidate)
\item \texttt{BEC-607A7A67} --- Outreach Care Continuity (Semantic-pattern candidate)
\item \texttt{BEC-792B9999} --- Perioperative Logistics (Semantic-pattern candidate)
\item \texttt{BEC-BE0DEFDF} --- Bed State \& Turnover Management (Semantic-pattern candidate)
\item \texttt{BEC-A2BE39BF} --- Facility Demand \& Capacity Telemetry (Semantic-pattern candidate)
\item \texttt{BEC-783C1ED2} --- On-Call Directory Resolution (Semantic-pattern candidate)
\item \texttt{BEC-3DE5AC43} --- Credential Verification Checkpoint (Semantic-pattern candidate)
\item \texttt{BEC-C048B88A} --- Multi-Model Information Exchange (Semantic-pattern candidate)
\item \texttt{BEC-C138A5D7} --- Governed Query \& Search (Semantic-pattern candidate)
\item \texttt{BEC-2EBB9F24} --- Boundary Protocol Mediation (Semantic-pattern candidate)
\item \texttt{BEC-E5775112} --- Terminology \& Reference Mapping (Semantic-pattern candidate)
\item \texttt{BEC-0CD2C755} --- LHR-Aware Multidisciplinary Collaboration (Semantic-pattern candidate)
\item \texttt{BEC-7D688C5A} --- Operational Messaging (Semantic-pattern candidate)
\item \texttt{BEC-1022E17A} --- Operational \& Clinical Calendar Sync (Semantic-pattern candidate)
\item \texttt{BEC-DA832DCE} --- Decoupled Human Interaction \& Dashboards (Semantic-pattern candidate)
\item \texttt{BEC-E7F9D13F} --- Canonical Semantics \& Schema Governance (Semantic-pattern candidate)
\item \texttt{EC-54C1BCAB} --- Information Lifecycle Enablement (Enterprise derived)
\item \texttt{EC-F2753584} --- Discovery and Retrieval Enablement (Enterprise derived)
\end{itemize}
\paragraph{Information Validation} 5 occurrences.\par
\begin{itemize}
\item \texttt{BEC-E86070E3} --- Multi-Model Information Exchange (Semantic-pattern candidate)
\item \texttt{BEC-77A76C7C} --- Terminology \& Reference Mapping (Semantic-pattern candidate)
\item \texttt{BEC-74F7B6A8} --- Canonical Semantics \& Schema Governance (Semantic-pattern candidate)
\item \texttt{EC-23381315} --- Information Lifecycle Enablement (Enterprise derived)
\item \texttt{EC-A86D200D} --- Semantic Conformance Enablement (Enterprise derived)
\end{itemize}
\paragraph{Location Availability Management} 3 occurrences.\par
\begin{itemize}
\item \texttt{BEC-0538C23A} --- Healthcare Location (FHIR-informed candidate)
\item \texttt{BEC-310FF956} --- Location Structure (FHIR-informed candidate)
\item \texttt{EC-582FEC53} --- Context and State Enablement (Enterprise derived)
\end{itemize}
\paragraph{Location Discovery} 3 occurrences.\par
\begin{itemize}
\item \texttt{BEC-BE1B64CD} --- Healthcare Location (FHIR-informed candidate)
\item \texttt{BEC-4BCD6B9C} --- Location Structure (FHIR-informed candidate)
\item \texttt{EC-009A4A5A} --- Discovery and Retrieval Enablement (Enterprise derived)
\end{itemize}
\paragraph{Location Hierarchy Management} 3 occurrences.\par
\begin{itemize}
\item \texttt{BEC-C52417EE} --- Healthcare Location (FHIR-informed candidate)
\item \texttt{BEC-6DCC28F3} --- Location Structure (FHIR-informed candidate)
\item \texttt{EC-A179D91B} --- Relationship and Structure Enablement (Enterprise derived)
\end{itemize}
\paragraph{Location Identity Management} 3 occurrences.\par
\begin{itemize}
\item \texttt{BEC-15887E90} --- Healthcare Location (FHIR-informed candidate)
\item \texttt{BEC-FC2DD94F} --- Location Structure (FHIR-informed candidate)
\item \texttt{EC-039334FA} --- Identity and Correlation Enablement (Enterprise derived)
\end{itemize}
\paragraph{Medication Activity Tracking} 3 occurrences.\par
\begin{itemize}
\item \texttt{BEC-9E3BA088} --- Medication Order \& Dispense Management (FHIR-informed candidate)
\item \texttt{BEC-0F34B38A} --- Medication Therapy Provision and Monitoring (FHIR-informed candidate)
\item \texttt{EC-67A2E372} --- Activity Coordination Enablement (Enterprise derived)
\end{itemize}
\paragraph{Medication Identity Resolution} 4 occurrences.\par
\begin{itemize}
\item \texttt{BEC-074E40CB} --- Medication Order \& Dispense Management (FHIR-informed candidate)
\item \texttt{BEC-910194D7} --- Medication Therapy Provision and Monitoring (FHIR-informed candidate)
\item \texttt{EC-5FC97327} --- Identity and Correlation Enablement (Enterprise derived)
\item \texttt{EC-09A3798F} --- Discovery and Retrieval Enablement (Enterprise derived)
\end{itemize}
\paragraph{Medication Information Correlation} 4 occurrences.\par
\begin{itemize}
\item \texttt{BEC-7D9A630B} --- Medication Order \& Dispense Management (FHIR-informed candidate)
\item \texttt{BEC-B3D9C0F4} --- Medication Therapy Provision and Monitoring (FHIR-informed candidate)
\item \texttt{EC-D444A2C1} --- Identity and Correlation Enablement (Enterprise derived)
\item \texttt{EC-53DFF234} --- Information Lifecycle Enablement (Enterprise derived)
\end{itemize}
\paragraph{Medication Safety Context Validation} 5 occurrences.\par
\begin{itemize}
\item \texttt{BEC-AFA2A006} --- Medication Order \& Dispense Management (FHIR-informed candidate)
\item \texttt{BEC-324D8BAD} --- Medication Therapy Provision and Monitoring (FHIR-informed candidate)
\item \texttt{EC-93E36E5B} --- Semantic Conformance Enablement (Enterprise derived)
\item \texttt{EC-2739EBA1} --- Context and State Enablement (Enterprise derived)
\item \texttt{EC-34E6439F} --- Assurance and Resilience Enablement (Enterprise derived)
\end{itemize}
\paragraph{No established Feature decomposition (Adjacent / Reference context)} 3 occurrences.\par
\begin{itemize}
\item \texttt{BEC-3700FBFB} --- Health Product Definition (Existing catalogue)
\item \texttt{BEC-4937CA4F} --- Product Formulary \& Catalogue (Existing catalogue)
\item \texttt{BEC-4614D02A} --- Product Batch \& Serial Tracking (Existing catalogue)
\end{itemize}
\paragraph{Organisation Discovery} 4 occurrences.\par
\begin{itemize}
\item \texttt{BEC-D6FB2614} --- Healthcare Organisation (FHIR-informed candidate)
\item \texttt{BEC-721B38D1} --- Organisational Structure (FHIR-informed candidate)
\item \texttt{BEC-73EF81CD} --- Organisation Contact (FHIR-informed candidate)
\item \texttt{EC-50A2C171} --- Discovery and Retrieval Enablement (Enterprise derived)
\end{itemize}
\paragraph{Organisational Relationship Management} 4 occurrences.\par
\begin{itemize}
\item \texttt{BEC-AD09D42A} --- Healthcare Organisation (FHIR-informed candidate)
\item \texttt{BEC-676858F6} --- Organisational Structure (FHIR-informed candidate)
\item \texttt{BEC-EFB7F137} --- Organisation Contact (FHIR-informed candidate)
\item \texttt{EC-0B8C4EDC} --- Relationship and Structure Enablement (Enterprise derived)
\end{itemize}
\paragraph{Organisational Structure Management} 4 occurrences.\par
\begin{itemize}
\item \texttt{BEC-6B51A123} --- Healthcare Organisation (FHIR-informed candidate)
\item \texttt{BEC-ACAB1168} --- Organisational Structure (FHIR-informed candidate)
\item \texttt{BEC-272EFC75} --- Organisation Contact (FHIR-informed candidate)
\item \texttt{EC-49166B2A} --- Relationship and Structure Enablement (Enterprise derived)
\end{itemize}
\paragraph{Outcome Reporting} 12 occurrences.\par
\begin{itemize}
\item \texttt{BEC-F6FC078E} --- Primary Care Coordination (Semantic-pattern candidate)
\item \texttt{BEC-8ED2F7E4} --- Acute Care Monitoring \& Event Coordination (Semantic-pattern candidate)
\item \texttt{BEC-1725AA08} --- Community Care Coordination (Semantic-pattern candidate)
\item \texttt{BEC-4734D5E1} --- Session \& Queue Coordination (Semantic-pattern candidate)
\item \texttt{BEC-4864C81B} --- Ward Logistics Coordination (Semantic-pattern candidate)
\item \texttt{BEC-5E9192CC} --- ED Operational Coordination (Semantic-pattern candidate)
\item \texttt{BEC-4D2A47AC} --- Specialized Equipment Coordination (Semantic-pattern candidate)
\item \texttt{BEC-D56786A8} --- Mobile Worker Dispatch \& Safety (Semantic-pattern candidate)
\item \texttt{BEC-A3268ADD} --- Operational Work Assignment \& Supervision (Semantic-pattern candidate)
\item \texttt{BEC-D56CF336} --- Discharge Progression \& Coordination (Semantic-pattern candidate)
\item \texttt{BEC-6B53749D} --- Multi-Work-Unit Orchestration (Semantic-pattern candidate)
\item \texttt{EC-09035A01} --- Assurance and Resilience Enablement (Enterprise derived)
\end{itemize}
\paragraph{Policy Evaluation} 3 occurrences.\par
\begin{itemize}
\item \texttt{BEC-EF1AF8F9} --- Client Privacy (Semantic-pattern candidate)
\item \texttt{BEC-5BAB5EA0} --- Policy \& Security Governance (Semantic-pattern candidate)
\item \texttt{EC-840C09A8} --- Governed Control Enablement (Enterprise derived)
\end{itemize}
\paragraph{Practitioner Discovery} 3 occurrences.\par
\begin{itemize}
\item \texttt{BEC-E9C766CB} --- Practitioner (FHIR-informed candidate)
\item \texttt{BEC-E24857D9} --- Practitioner Role (FHIR-informed candidate)
\item \texttt{EC-26353D17} --- Discovery and Retrieval Enablement (Enterprise derived)
\end{itemize}
\paragraph{Professional Relationship Management} 3 occurrences.\par
\begin{itemize}
\item \texttt{BEC-028ABD17} --- Practitioner (FHIR-informed candidate)
\item \texttt{BEC-549D792B} --- Practitioner Role (FHIR-informed candidate)
\item \texttt{EC-366AFC9B} --- Relationship and Structure Enablement (Enterprise derived)
\end{itemize}
\paragraph{Qualification Association} 4 occurrences.\par
\begin{itemize}
\item \texttt{BEC-F9FA105B} --- Practitioner (FHIR-informed candidate)
\item \texttt{BEC-A4596B27} --- Practitioner Role (FHIR-informed candidate)
\item \texttt{EC-E2F0CB40} --- Relationship and Structure Enablement (Enterprise derived)
\item \texttt{EC-DB69DAF0} --- Semantic Conformance Enablement (Enterprise derived)
\end{itemize}
\paragraph{Relationship Management} 3 occurrences.\par
\begin{itemize}
\item \texttt{BEC-C7981054} --- Person Identity (FHIR-informed candidate)
\item \texttt{BEC-3DBACC5A} --- Product Formulary \& Catalogue (Semantic-pattern candidate)
\item \texttt{EC-F160E33E} --- Relationship and Structure Enablement (Enterprise derived)
\end{itemize}
\paragraph{Request Exception Management} 3 occurrences.\par
\begin{itemize}
\item \texttt{BEC-AA804BEF} --- Closed-Loop Order Progression (FHIR-informed candidate)
\item \texttt{BEC-B0E3D03A} --- Medication Order \& Dispense Management (FHIR-informed candidate)
\item \texttt{EC-85D9E54B} --- Assurance and Resilience Enablement (Enterprise derived)
\end{itemize}
\paragraph{Request Fulfilment Tracking} 3 occurrences.\par
\begin{itemize}
\item \texttt{BEC-977D6A7D} --- Closed-Loop Order Progression (FHIR-informed candidate)
\item \texttt{BEC-E0ABA2EF} --- Medication Order \& Dispense Management (FHIR-informed candidate)
\item \texttt{EC-A191A2BC} --- Activity Coordination Enablement (Enterprise derived)
\end{itemize}
\paragraph{Request Lifecycle Management} 3 occurrences.\par
\begin{itemize}
\item \texttt{BEC-142E39C3} --- Closed-Loop Order Progression (FHIR-informed candidate)
\item \texttt{BEC-AB5055A8} --- Medication Order \& Dispense Management (FHIR-informed candidate)
\item \texttt{EC-720E493E} --- Information Lifecycle Enablement (Enterprise derived)
\end{itemize}
\paragraph{Request Status Correlation} 4 occurrences.\par
\begin{itemize}
\item \texttt{BEC-81D316C2} --- Closed-Loop Order Progression (FHIR-informed candidate)
\item \texttt{BEC-D9110CE8} --- Medication Order \& Dispense Management (FHIR-informed candidate)
\item \texttt{EC-D418A306} --- Identity and Correlation Enablement (Enterprise derived)
\item \texttt{EC-DB0E7812} --- Context and State Enablement (Enterprise derived)
\end{itemize}
\paragraph{Service Availability Management} 6 occurrences.\par
\begin{itemize}
\item \texttt{BEC-22CF71DF} --- Healthcare Service Definition (FHIR-informed candidate)
\item \texttt{BEC-F964694F} --- Service Provision Relationship (FHIR-informed candidate)
\item \texttt{BEC-235C211D} --- Service Availability Definition (FHIR-informed candidate)
\item \texttt{BEC-FA68F289} --- Service Eligibility Definition (FHIR-informed candidate)
\item \texttt{BEC-870CAEA8} --- Diagnostic Service Integration (FHIR-informed candidate)
\item \texttt{EC-FF26979C} --- Context and State Enablement (Enterprise derived)
\end{itemize}
\paragraph{Service Definition Management} 5 occurrences.\par
\begin{itemize}
\item \texttt{BEC-8CDB717F} --- Healthcare Service Definition (FHIR-informed candidate)
\item \texttt{BEC-46B095FB} --- Service Provision Relationship (FHIR-informed candidate)
\item \texttt{BEC-61F28669} --- Service Availability Definition (FHIR-informed candidate)
\item \texttt{BEC-271154E8} --- Service Eligibility Definition (FHIR-informed candidate)
\item \texttt{BEC-48D54B88} --- Diagnostic Service Integration (FHIR-informed candidate)
\end{itemize}
\paragraph{Service Discovery} 6 occurrences.\par
\begin{itemize}
\item \texttt{BEC-DEE34CB8} --- Healthcare Service Definition (FHIR-informed candidate)
\item \texttt{BEC-7A743DDD} --- Service Provision Relationship (FHIR-informed candidate)
\item \texttt{BEC-B8A2142A} --- Service Availability Definition (FHIR-informed candidate)
\item \texttt{BEC-9B1E4AC5} --- Service Eligibility Definition (FHIR-informed candidate)
\item \texttt{BEC-14C91D01} --- Diagnostic Service Integration (FHIR-informed candidate)
\item \texttt{EC-82149096} --- Discovery and Retrieval Enablement (Enterprise derived)
\end{itemize}
\paragraph{Service Eligibility Management} 6 occurrences.\par
\begin{itemize}
\item \texttt{BEC-766061BE} --- Healthcare Service Definition (FHIR-informed candidate)
\item \texttt{BEC-99951D40} --- Service Provision Relationship (FHIR-informed candidate)
\item \texttt{BEC-0AA34D18} --- Service Availability Definition (FHIR-informed candidate)
\item \texttt{BEC-1F711951} --- Service Eligibility Definition (FHIR-informed candidate)
\item \texttt{BEC-E8E1F920} --- Diagnostic Service Integration (FHIR-informed candidate)
\item \texttt{EC-D98B7AAE} --- Governed Control Enablement (Enterprise derived)
\end{itemize}
\paragraph{State Monitoring} 26 occurrences.\par
\begin{itemize}
\item \texttt{BEC-B3C6EFA0} --- Healthcare Subject (Semantic-pattern candidate)
\item \texttt{BEC-D61B3C3B} --- Client Relationship (Semantic-pattern candidate)
\item \texttt{BEC-2CF899D2} --- Provider Relationship (Semantic-pattern candidate)
\item \texttt{BEC-141E9549} --- Provider Engagement (Semantic-pattern candidate)
\item \texttt{BEC-67DFF33A} --- Care-Place Definition (Semantic-pattern candidate)
\item \texttt{BEC-8F3D543E} --- Health Product Definition (Semantic-pattern candidate)
\item \texttt{BEC-C2E441ED} --- Product Batch \& Serial Tracking (Semantic-pattern candidate)
\item \texttt{BEC-80DFF5D9} --- Client Intake \& Demographic Update (Semantic-pattern candidate)
\item \texttt{BEC-6100D197} --- Procedure \& Surgical Administration (Semantic-pattern candidate)
\item \texttt{BEC-CF085C02} --- Clinical Documentation Lifecycle (Semantic-pattern candidate)
\item \texttt{BEC-ADC2AC77} --- Emergency Care Progression (Semantic-pattern candidate)
\item \texttt{BEC-C577B7C8} --- Screening and Immunisation Enablement (Semantic-pattern candidate)
\item \texttt{BEC-59BE4E9F} --- Outreach Care Continuity (Semantic-pattern candidate)
\item \texttt{BEC-9F573A4C} --- Perioperative Logistics (Semantic-pattern candidate)
\item \texttt{BEC-25AAE44B} --- Bed State \& Turnover Management (Semantic-pattern candidate)
\item \texttt{BEC-F78FC8D3} --- Facility Demand \& Capacity Telemetry (Semantic-pattern candidate)
\item \texttt{BEC-703E878F} --- On-Call Directory Resolution (Semantic-pattern candidate)
\item \texttt{BEC-E9331A31} --- Credential Verification Checkpoint (Semantic-pattern candidate)
\item \texttt{BEC-474C9605} --- Governed Query \& Search (Semantic-pattern candidate)
\item \texttt{BEC-C8A17F6A} --- Boundary Protocol Mediation (Semantic-pattern candidate)
\item \texttt{BEC-911087B0} --- LHR-Aware Multidisciplinary Collaboration (Semantic-pattern candidate)
\item \texttt{BEC-D75B5D46} --- Operational Messaging (Semantic-pattern candidate)
\item \texttt{BEC-922D6FFF} --- Operational \& Clinical Calendar Sync (Semantic-pattern candidate)
\item \texttt{BEC-7CB4FD80} --- Decoupled Human Interaction \& Dashboards (Semantic-pattern candidate)
\item \texttt{EC-CC8047E5} --- Context and State Enablement (Enterprise derived)
\item \texttt{EC-0F5671F1} --- Assurance and Resilience Enablement (Enterprise derived)
\end{itemize}
\paragraph{State Tracking} 13 occurrences.\par
\begin{itemize}
\item \texttt{BEC-A945B4F5} --- Primary Care Coordination (Semantic-pattern candidate)
\item \texttt{BEC-DFC60092} --- Acute Care Monitoring \& Event Coordination (Semantic-pattern candidate)
\item \texttt{BEC-52C8768A} --- Community Care Coordination (Semantic-pattern candidate)
\item \texttt{BEC-28A55B14} --- Session \& Queue Coordination (Semantic-pattern candidate)
\item \texttt{BEC-E6E33D5A} --- Ward Logistics Coordination (Semantic-pattern candidate)
\item \texttt{BEC-704E04EA} --- ED Operational Coordination (Semantic-pattern candidate)
\item \texttt{BEC-F1B486E3} --- Specialized Equipment Coordination (Semantic-pattern candidate)
\item \texttt{BEC-E87B589F} --- Mobile Worker Dispatch \& Safety (Semantic-pattern candidate)
\item \texttt{BEC-D5AEDF90} --- Operational Work Assignment \& Supervision (Semantic-pattern candidate)
\item \texttt{BEC-D2873E18} --- Discharge Progression \& Coordination (Semantic-pattern candidate)
\item \texttt{BEC-D25D58D7} --- Multi-Work-Unit Orchestration (Semantic-pattern candidate)
\item \texttt{EC-D122BC94} --- Context and State Enablement (Enterprise derived)
\item \texttt{EC-E762705E} --- Activity Coordination Enablement (Enterprise derived)
\end{itemize}
\paragraph{Subject Context Management} 5 occurrences.\par
\begin{itemize}
\item \texttt{BEC-152BD757} --- Multidisciplinary Inpatient Care Coordination (FHIR-informed candidate)
\item \texttt{BEC-0B235F68} --- Remote Patient Monitoring (FHIR-informed candidate)
\item \texttt{BEC-CD0D5A3A} --- Outpatient Flow Enablement (FHIR-informed candidate)
\item \texttt{BEC-CBF19A4E} --- Patient Transfer Coordination (FHIR-informed candidate)
\item \texttt{EC-1900212B} --- Context and State Enablement (Enterprise derived)
\end{itemize}
\paragraph{Subject Identity Correlation} 5 occurrences.\par
\begin{itemize}
\item \texttt{BEC-826FE486} --- Multidisciplinary Inpatient Care Coordination (FHIR-informed candidate)
\item \texttt{BEC-08D640DD} --- Remote Patient Monitoring (FHIR-informed candidate)
\item \texttt{BEC-0D49726E} --- Outpatient Flow Enablement (FHIR-informed candidate)
\item \texttt{BEC-05915061} --- Patient Transfer Coordination (FHIR-informed candidate)
\item \texttt{EC-157BAA94} --- Identity and Correlation Enablement (Enterprise derived)
\end{itemize}
\paragraph{Subject Information Discovery} 6 occurrences.\par
\begin{itemize}
\item \texttt{BEC-C7242AB0} --- Multidisciplinary Inpatient Care Coordination (FHIR-informed candidate)
\item \texttt{BEC-6B015BE0} --- Remote Patient Monitoring (FHIR-informed candidate)
\item \texttt{BEC-17DAEA1C} --- Outpatient Flow Enablement (FHIR-informed candidate)
\item \texttt{BEC-A0E56E8C} --- Patient Transfer Coordination (FHIR-informed candidate)
\item \texttt{EC-9788F116} --- Information Lifecycle Enablement (Enterprise derived)
\item \texttt{EC-13D8A141} --- Discovery and Retrieval Enablement (Enterprise derived)
\end{itemize}
\paragraph{Subject Relationship Management} 5 occurrences.\par
\begin{itemize}
\item \texttt{BEC-941ACC0E} --- Multidisciplinary Inpatient Care Coordination (FHIR-informed candidate)
\item \texttt{BEC-2A678BE0} --- Remote Patient Monitoring (FHIR-informed candidate)
\item \texttt{BEC-FCB21DFC} --- Outpatient Flow Enablement (FHIR-informed candidate)
\item \texttt{BEC-F93AE731} --- Patient Transfer Coordination (FHIR-informed candidate)
\item \texttt{EC-20119D58} --- Relationship and Structure Enablement (Enterprise derived)
\end{itemize}
\chapter{CT1 - Outstanding Architectural Review Topics}
\paragraph{Review item.} Complete CT1/CT2/CT3 ancestry and Canonical IDs are not yet reconciled into RADS.\par
\paragraph{Review item.} 137 Feature local IDs are established; complete structural Canonical IDs remain unallocated.\par
\paragraph{Review item.} Some Feature associations are explicitly unresolved (for example Ward Patient Flow Coordination and ED Length of Stay Tracking).\par
\paragraph{Review item.} 69 enabling catalogue descriptions remain unresolved as standalone capability identities.\par
\paragraph{Review item.} Clinical handover/collaboration semantic boundary conflict remains preserved for review.\par
\paragraph{Review item.} Health Service Operational Readiness has understood semantics but unresolved canonical placement.\par
\paragraph{Review item.} Medication Compounding and Preparation remains an open capability question.\par
\paragraph{Review item.} Clinical Intervention CT3 taxonomy remains unresolved; do not infer peer types from examples.\par
\paragraph{Review item.} Domain05 B2.2 remains parked pending capability/RADS reconciliation.\par
\chapter{CT1 - Source and Evidence Register}
Primary working source: Harmonia Capability Architecture v0.4 mind-map, generated 14 October 2026. It carries established Business Capability names, candidate BEC hierarchy and 137 established Feature labels, plus Function, Service and Process associations.

Governed repository: fhirfactory/harmonia, Harmonia-GovernedAccess branch. Key sources include docs/markdown/02-strategy/capabilities/business-capabilities.md, business-enabling-capabilities.md, enterprise-capabilities.md; docs/markdown/03-business-architecture/reference-model/business-architecture.json; docs/markdown/03-business-architecture/metamodel/business-architecture-metamodel.md; docs/markdown/scratchpad/architecture-design-notes.md; and docs/markdown/governance/architectural-axioms.md.

FHIR conceptual source: HL7 FHIR Release 5 resource scope and usage descriptions, particularly Practitioner, PractitionerRole, Organization, Location, HealthcareService, Patient, Appointment, Encounter, ServiceRequest, Task, CarePlan, MedicationRequest, DiagnosticReport, Device and related resources. Candidate derivation is semantic, not data-element mapping. Individual FHIR resource pages should be verified in the offline adjudication pass.

Explicit exclusion: legacy-documentation and superseded legacy narratives were not used.

\chapter{CT1 - Generation Summary}
BEC CT3 candidates inspected: 73. Existing BEC Feature occurrences retained: 137. Additional candidate BEC Features proposed: 333. Distinct Feature display names: 222. Reference ID entries: 942.

\backmatter\printindex\end{document}
